% Auto-generated from data/segments.json + layout.json (+ transforms) — do not edit by hand.
% volume: 37Abhi09
% mode: printing
% segments: 3937
% transform_rules: 8
% toc_marks: 614

% Volume layout defaults (layout.json ``layout``).
\csromanlayoutapply{1}{1.5}{21.6pt}{8pt}{8pt}{65pt}{2.5em}%

\pitaka{อภิธมฺม\-ปิฏก}
\csromanheader{2}{\textbf{ธมฺมานุโลม}}
\csromanmatikaanchor{mtk.0}
\csromantocmarknopage{chapter}{ปฏฺฐานปาฬิ ทุติยภาค}{mtk.0}
\bookmark[dest={mtk.0},level=0]{ปฏฺฐานปาฬิ ทุติยภาค}
\gambhira{ติกปฏฺฐาน\-ปาฬิ}
\namakkaram{นโม ตสฺส ภควโต อรหโต สมฺมา\-สมฺพุทฺธสฺส.}
\csromanmarkh{6. วิตกฺกตฺติก}\csromanmarkhii{1. ปฏิจฺจวาร}\csromanheadernomark{2}{6. \textbf{วิตกฺกตฺติก 1. ปฏิจฺจวาร}}
\csromanmatikaanchor{mtk.1}
\csromanmatikaanchor{mtk.8}
\csromanmatikaanchor{mtk.137}
\csromantocmarknopage{chapter}{2. ปจฺจยปจฺจนีย 2. สงฺขฺยาวาร}{mtk.1}
\bookmark[dest={mtk.1},level=0]{2. ปจฺจยปจฺจนีย 2. สงฺขฺยาวาร}
\csromanmarkh{1. ปจฺจยานุโลม}\csromanmarkhii{1. วิภงฺควาร}\csromanheadernomark{6}{1. \textbf{ปจฺจยานุโลม 1. วิภงฺควาร}}
\csromanheader{2}{\textbf{เหตุ}}
\proseitem{1}{\csromanfolio{1}สวิตกฺก\-สวิจารํ ธมฺมํ ปฏิจฺจ สวิตกฺก\-สวิจาโร ธมฺโม อุปฺปชฺชติ เหตุปจฺจยา. สวิตกฺก\-สวิจารํ เอกํ ขนฺธํ ปฏิจฺจ ตโย ขนฺธา, ตโย ขนฺเธ ปฏิจฺจ เอโก ขนฺโธ, เทฺว ขนฺเธ ปฏิจฺจ เทฺว ขนฺธา.  ปฏิสนฺธิ\-กฺขเณ สวิตกฺก\-สวิจารํ เอกํ ขนฺธํ ปฏิจฺจ ตโย ขนฺธา ฯเปฯ เทฺว ขนฺเธ ปฏิจฺจ เทฺว ขนฺธา. (1)}

\prose{สวิตกฺก\-สวิจารํ ธมฺมํ ปฏิจฺจ อวิตกฺก\-วิจารมตฺโต ธมฺโม อุปฺปชฺชติ เหตุปจฺจยา. สวิตกฺก\-สวิจาเร ขนฺเธ ปฏิจฺจ วิตกฺโก.  ปฏิสนฺธิ\-กฺขเณ สวิตกฺก\-สวิจาเร ขนฺเธ ปฏิจฺจ วิตกฺโก. (2)}

\prose{สวิตกฺก\-สวิจารํ ธมฺมํ ปฏิจฺจ อวิตกฺก\-อวิจาโร ธมฺโม อุปฺปชฺชติ เหตุปจฺจยา. สวิตกฺก\-สวิจาเร ขนฺเธ ปฏิจฺจ จิตฺต\-สมุฏฺฐานํ รูปํ. ปฏิสนฺธิ\-กฺขเณ สวิตกฺก\-สวิจาเร ขนฺเธ ปฏิจฺจ กฏตฺตารูปํ. (3)}

\prose{\csromanfolio{2}สวิตกฺก\-สวิจารํ ธมฺมํ ปฏิจฺจ สวิตกฺก\-สวิจาโร จ อวิตกฺก\-อวิจาโร จ ธมฺมา อุปฺปชฺชนฺติ เหตุปจฺจยา. สวิตกฺก\-สวิจารํ เอกํ ขนฺธํ ปฏิจฺจ ตโย ขนฺธา จิตฺต\-สมุฏฺฐานญฺจ รูปํ ฯเปฯ เทฺว ขนฺเธ ปฏิจฺจ เทฺว ขนฺธา จิตฺต\-สมุฏฺฐานญฺจ รูปํ.  ปฏิสนฺธิ\-กฺขเณ สวิตกฺก\-สวิจารํ เอกํ ขนฺธํ ปฏิจฺจ ตโย ขนฺธา กฏตฺตา จ รูปํ ฯเปฯ เทฺว ขนฺเธ ปฏิจฺจ เทฺว ขนฺธา กฏตฺตา จ รูปํ. (4)}

\prose{สวิตกฺก\-สวิจารํ ธมฺมํ ปฏิจฺจ อวิตกฺก\-วิจารมตฺโต จ อวิตกฺก\-อวิจาโร จ ธมฺมา อุปฺปชฺชนฺติ เหตุปจฺจยา. สวิตกฺก\-สวิจาเร ขนฺเธ ปฏิจฺจ วิตกฺโก จิตฺต\-สมุฏฺฐานญฺจ รูปํ.  ปฏิสนฺธิ\-กฺขเณ สวิตกฺก\-สวิจาเร ขนฺเธ ปฏิจฺจ วิตกฺโก กฏตฺตา จ รูปํ. (5)}

\prose{สวิตกฺก\-สวิจารํ ธมฺมํ ปฏิจฺจ สวิตกฺก\-สวิจาโร จ อวิตกฺก\-วิจารมตฺโต จ ธมฺมา อุปฺปชฺชนฺติ เหตุปจฺจยา. สวิตกฺก\-สวิจารํ เอกํ ขนฺธํ ปฏิจฺจ ตโย ขนฺธา วิตกฺโก จ ฯเปฯ เทฺว ขนฺเธ ปฏิจฺจ เทฺว ขนฺธา วิตกฺโก จ.  ปฏิสนฺธิ\-กฺขเณ สวิตกฺก\-สวิจารํ เอกํ ขนฺธํ ปฏิจฺจ ตโย ขนฺธา วิตกฺโก จ ฯเปฯ เทฺว ขนฺเธ ปฏิจฺจ เทฺว ขนฺธา วิตกฺโก จ. (6)}

\prose{สวิตกฺก\-สวิจารํ ธมฺมํ ปฏิจฺจ สวิตกฺก\-สวิจาโร จ อวิตกฺก\-วิจารมตฺโต จ อวิตกฺก\-อวิจาโร จ ธมฺมา อุปฺปชฺชนฺติ เหตุปจฺจยา. สวิตกฺก\-สวิจารํ เอกํ ขนฺธํ ปฏิจฺจ ตโย ขนฺธา วิตกฺโก จ จิตฺต\-สมุฏฺฐานญฺจ รูปํ ฯเปฯ เทฺว ขนฺเธ ปฏิจฺจ เทฺว ขนฺธา วิตกฺโก จ จิตฺต\-สมุฏฺฐานญฺจ รูปํ.  ปฏิสนฺธิ\-กฺขเณ สวิตกฺก\-สวิจารํ เอกํ ขนฺธํ ปฏิจฺจ ตโย ขนฺธา วิตกฺโก จ กฏตฺตา จ รูปํ ฯเปฯ เทฺว ขนฺเธ ปฏิจฺจ เทฺว ขนฺธา วิตกฺโก จ กฏตฺตา จ รูปํ. (7)}

\proseitem{2}{อวิตกฺก\-วิจารมตฺตํ ธมฺมํ ปฏิจฺจ อวิตกฺก\-วิจารมตฺโต ธมฺโม อุปฺปชฺชติ เหตุปจฺจยา. อวิตกฺก\-วิจารมตฺตํ เอกํ ขนฺธํ ปฏิจฺจ ตโย ขนฺธา ฯเปฯ เทฺว ขนฺเธ ปฏิจฺจ เทฺว ขนฺธา. ปฏิสนฺธิ\-กฺขเณ อวิตกฺก\-วิจารมตฺตํ เอกํ ขนฺธํ ปฏิจฺจ ตโย ขนฺธา ฯเปฯ เทฺว ขนฺเธ ปฏิจฺจ เทฺว ขนฺธา. (1)}

\prose{อวิตกฺก\-วิจารมตฺตํ ธมฺมํ ปฏิจฺจ สวิตกฺก\-สวิจาโร ธมฺโม อุปฺปชฺชติ เหตุปจฺจยา. วิตกฺกํ ปฏิจฺจ สวิตกฺก\-สวิจารา ขนฺธา. ปฏิสนฺธิ\-กฺขเณ วิตกฺกํ ปฏิจฺจ สวิตกฺก\-สวิจารา ขนฺธา. (2)}

\prose{อวิตกฺก\-วิจารมตฺตํ ธมฺมํ ปฏิจฺจ อวิตกฺก\-อวิจาโร ธมฺโม อุปฺปชฺชติ เหตุปจฺจยา. อวิตกฺก\-วิจารมตฺเต ขนฺเธ ปฏิจฺจ วิจาโร จิตฺต\-สมุฏฺฐานญฺจ \csromanfolio{3}รูปํ. วิตกฺกํ ปฏิจฺจ จิตฺต\-สมุฏฺฐานํ รูปํ.  ปฏิสนฺธิ\-กฺขเณ อวิตกฺก\-วิจารมตฺเต ขนฺเธ ปฏิจฺจ วิจาโร กฏตฺตา จ รูปํ.  ปฏิสนฺธิ\-กฺขเณ วิตกฺกํ ปฏิจฺจ กฏตฺตารูปํ. (3)}

\prose{อวิตกฺก\-วิจารมตฺตํ ธมฺมํ ปฏิจฺจ สวิตกฺก\-สวิจาโร จ อวิตกฺก\-อวิจาโร จ ธมฺมา อุปฺปชฺชนฺติ เหตุปจฺจยา. วิตกฺกํ ปฏิจฺจ สวิตกฺก\-สวิจารา ขนฺธา จิตฺต\-สมุฏฺฐานญฺจ รูปํ.  ปฏิสนฺธิ\-กฺขเณ วิตกฺกํ ปฏิจฺจ สวิตกฺก\-สวิจารา ขนฺธา กฏตฺตา จ รูปํ. (4)}

\prose{อวิตกฺก\-วิจารมตฺตํ ธมฺมํ ปฏิจฺจ อวิตกฺก\-วิจารมตฺโต จ อวิตกฺก\-อวิจาโร จ ธมฺมา อุปฺปชฺชนฺติ เหตุปจฺจยา. อวิตกฺก\-วิจารมตฺตํ เอกํ ขนฺธํ ปฏิจฺจ ตโย ขนฺธา วิจาโร จ จิตฺต\-สมุฏฺฐานญฺจ รูปํ ฯเปฯ เทฺว ขนฺเธ ปฏิจฺจ เทฺว ขนฺธา วิจาโร จ จิตฺต\-สมุฏฺฐานญฺจ รูปํ.  ปฏิสนฺธิ\-กฺขเณ อวิตกฺก\-วิจารมตฺตํ เอกํ ขนฺธํ ปฏิจฺจ ตโย ขนฺธา วิจาโร จ กฏตฺตา จ รูปํ ฯเปฯ เทฺว ขนฺเธ ปฏิจฺจ เทฺว ขนฺธา วิจาโร จ กฏตฺตา จ รูปํ. (5)}

\prose{อวิตกฺก\-อวิจารํ ธมฺมํ ปฏิจฺจ อวิตกฺก\-อวิจาโร ธมฺโม อุปฺปชฺชติ เหตุปจฺจยา. อวิตกฺก\-อวิจารํ เอกํ ขนฺธํ ปฏิจฺจ ตโย ขนฺธา จิตฺต\-สมุฏฺฐานญฺจ รูปํ ฯเปฯ เทฺว ขนฺเธ ปฏิจฺจ เทฺว ขนฺธา จิตฺต\-สมุฏฺฐานญฺจ รูปํ. วิจารํ ปฏิจฺจ จิตฺต\-สมุฏฺฐานํ รูปํ.  ปฏิสนฺธิ\-กฺขเณ อวิตกฺก\-อวิจารํ เอกํ ขนฺธํ ปฏิจฺจ ตโย ขนฺธา กฏตฺตา จ รูปํ ฯเปฯ เทฺว ขนฺเธ ปฏิจฺจ เทฺว ขนฺธา กฏตฺตา จ รูปํ.  ปฏิสนฺธิ\-กฺขเณ วิจารํ ปฏิจฺจ กฏตฺตารูปํ. ขนฺเธ ปฏิจฺจ วตฺถุ, วตฺถุํ ปฏิจฺจ ขนฺธา. วิจารํ ปฏิจฺจ วตฺถุ, วตฺถุํ ปฏิจฺจ วิจาโร. เอกํ มหาภูตํ ปฏิจฺจ ตโย มหาภูตา ฯเปฯ มหาภูเต ปฏิจฺจ จิตฺต\-สมุฏฺฐานํ รูปํ กฏตฺตารูปํ อุปาทารูปํ. (1)}

\prose{อวิตกฺก\-อวิจารํ ธมฺมํ ปฏิจฺจ สวิตกฺก\-สวิจาโร ธมฺโม อุปฺปชฺชติ เหตุปจฺจยา.  ปฏิสนฺธิ\-กฺขเณ วตฺถุํ ปฏิจฺจ สวิตกฺก\-สวิจารา ขนฺธา. (2)}

\prose{อวิตกฺก\-อวิจารํ ธมฺมํ ปฏิจฺจ อวิตกฺก\-วิจารมตฺโต ธมฺโม อุปฺปชฺชติ เหตุปจฺจยา. วิจารํ ปฏิจฺจ อวิตกฺก\-วิจารมตฺตา ขนฺธา. ปฏิสนฺธิ\-กฺขเณ วิจารํ ปฏิจฺจ อวิตกฺก\-วิจารมตฺตา ขนฺธา. ปฏิสนฺธิ\-กฺขเณ วตฺถุํ ปฏิจฺจ อวิตกฺก\-วิจารมตฺตา ขนฺธา. ปฏิสนฺธิ\-กฺขเณ วตฺถุํ ปฏิจฺจ วิตกฺโก. (3)}

\proseitem{3}{\csromanfolio{4}อวิตกฺก\-อวิจารํ ธมฺมํ ปฏิจฺจ สวิตกฺก\-สวิจาโร จ อวิตกฺก\-อวิจาโร จ ธมฺมา อุปฺปชฺชนฺติ เหตุปจฺจยา.  ปฏิสนฺธิ\-กฺขเณ วตฺถุํ ปฏิจฺจ สวิตกฺก\-สวิจารา ขนฺธา. มหาภูเต ปฏิจฺจ กฏตฺตารูปํ. (4)}

\prose{อวิตกฺก\-อวิจารํ ธมฺมํ ปฏิจฺจ อวิตกฺก\-วิจารมตฺโต จ อวิตกฺก\-อวิจาโร จ ธมฺมา อุปฺปชฺชนฺติ เหตุปจฺจยา. วิจารํ ปฏิจฺจ อวิตกฺก\-วิจารมตฺตา ขนฺธา จิตฺต\-สมุฏฺฐานญฺจ รูปํ. ปฏิสนฺธิ\-กฺขเณ วิจารํ ปฏิจฺจ อวิตกฺก\-วิจารมตฺตา ขนฺธา กฏตฺตา จ รูปํ. ปฏิสนฺธิ\-กฺขเณ วตฺถุํ ปฏิจฺจ อวิตกฺก\-วิจารมตฺตา ขนฺธา. มหาภูเต ปฏิจฺจ กฏตฺตารูปํ. ปฏิสนฺธิ\-กฺขเณ วตฺถุํ ปฏิจฺจ วิตกฺโก. มหาภูเต ปฏิจฺจ กฏตฺตารูปํ. ปฏิสนฺธิ\-กฺขเณ วตฺถุํ ปฏิจฺจ อวิตกฺก\-วิจารมตฺตา ขนฺธา จ วิจาโร จ. (5)}

\prose{อวิตกฺก\-อวิจารํ ธมฺมํ ปฏิจฺจ สวิตกฺก\-สวิจาโร จ อวิตกฺก\-วิจารมตฺโต จ ธมฺมา อุปฺปชฺชนฺติ เหตุปจฺจยา. ปฏิสนฺธิ\-กฺขเณ วตฺถุํ ปฏิจฺจ สวิตกฺก\-สวิจารา ขนฺธา จ วิตกฺโก จ. (6)}

\prose{อวิตกฺก\-อวิจารํ ธมฺมํ ปฏิจฺจ สวิตกฺก\-สวิจาโร จ อวิตกฺก\-วิจารมตฺโต จ อวิตกฺก\-อวิจาโร จ ธมฺมา อุปฺปชฺชนฺติ เหตุปจฺจยา. ปฏิสนฺธิ\-กฺขเณ วตฺถุํ ปฏิจฺจ สวิตกฺก\-สวิจารา ขนฺธา จ วิตกฺโก จ. มหาภูเต ปฏิจฺจ กฏตฺตารูปํ. (7)}

\proseitem{4}{สวิตกฺกสวิจารญฺจ อวิตกฺกอวิจารญฺจ ธมฺมํ ปฏิจฺจ สวิตกฺก\-สวิจาโร ธมฺโม อุปฺปชฺชติ เหตุปจฺจยา. ปฏิสนฺธิ\-กฺขเณ สวิตกฺก\-สวิจารํ เอกํ ขนฺธญฺจ วตฺถุญฺจ ปฏิจฺจ ตโย ขนฺธา ฯเปฯ เทฺว ขนฺเธ จ วตฺถุญฺจ ปฏิจฺจ เทฺว ขนฺธา. (1)}

\prose{สวิตกฺกสวิจารญฺจ อวิตกฺกอวิจารญฺจ ธมฺมํ ปฏิจฺจ อวิตกฺก\-วิจารมตฺโต ธมฺโม อุปฺปชฺชติ เหตุปจฺจยา. ปฏิสนฺธิ\-กฺขเณ สวิตกฺก\-สวิจาเร ขนฺเธ จ วตฺถุญฺจ ปฏิจฺจ วิตกฺโก. (2)}

\prose{สวิตกฺกสวิจารญฺจ อวิตกฺกอวิจารญฺจ ธมฺมํ ปฏิจฺจ อวิตกฺก\-อวิจาโร ธมฺโม อุปฺปชฺชติ เหตุปจฺจยา. สวิตกฺก\-สวิจาเร ขนฺเธ จ มหาภูเต จ ปฏิจฺจ จิตฺต\-สมุฏฺฐานํ รูปํ. ปฏิสนฺธิ\-กฺขเณ สวิตกฺก\-สวิจาเร ขนฺเธ จ มหาภูเต จ ปฏิจฺจ กฏตฺตารูปํ. (3)}

\prose{สวิตกฺกสวิจารญฺจ อวิตกฺกอวิจารญฺจ ธมฺมํ ปฏิจฺจ สวิตกฺก\-สวิจาโร จ อวิตกฺก\-อวิจาโร จ ธมฺมา อุปฺปชฺชนฺติ เหตุปจฺจยา. ปฏิสนฺธิ\-กฺขเณ \csromanfolio{5}สวิตกฺก\-สวิจารํ เอกํ ขนฺธญฺจ วตฺถุญฺจ ปฏิจฺจ ตโย ขนฺธา ฯเปฯ เทฺว ขนฺเธ จ วตฺถุญฺจ ปฏิจฺจ เทฺว ขนฺธา. สวิตกฺก\-สวิจาเร ขนฺเธ จ มหาภูเต จ ปฏิจฺจ กฏตฺตารูปํ. (4)}

\prose{สวิตกฺกสวิจารญฺจ อวิตกฺกอวิจารญฺจ ธมฺมํ ปฏิจฺจ อวิตกฺก\-วิจารมตฺโต จ อวิตกฺก\-อวิจาโร จ ธมฺมา อุปฺปชฺชนฺติ เหตุปจฺจยา. ปฏิสนฺธิ\-กฺขเณ สวิตกฺก\-สวิจาเร ขนฺเธ จ วตฺถุญฺจ ปฏิจฺจ วิตกฺโก. สวิตกฺก\-สวิจาเร ขนฺเธ จ มหาภูเต จ ปฏิจฺจ กฏตฺตารูปํ. (5)}

\prose{สวิตกฺกสวิจารญฺจ อวิตกฺกอวิจารญฺจ ธมฺมํ ปฏิจฺจ สวิตกฺก\-สวิจาโร จ อวิตกฺก\-วิจารมตฺโต จ ธมฺมา อุปฺปชฺชนฺติ เหตุปจฺจยา. ปฏิสนฺธิ\-กฺขเณ สวิตกฺก\-สวิจารํ เอกํ ขนฺธญฺจ วตฺถุญฺจ ปฏิจฺจ ตโย ขนฺธา วิตกฺโก จ ฯเปฯ เทฺว ขนฺเธ จ วตฺถุญฺจ ปฏิจฺจ เทฺว ขนฺธา วิตกฺโก จ. (6)}

\prose{สวิตกฺกสวิจารญฺจ อวิตกฺกอวิจารญฺจ ธมฺมํ ปฏิจฺจ สวิตกฺก\-สวิจาโร จ อวิตกฺก\-วิจารมตฺโต จ อวิตกฺก\-อวิจาโร จ ธมฺมา อุปฺปชฺชนฺติ เหตุปจฺจยา. ปฏิสนฺธิ\-กฺขเณ สวิตกฺก\-สวิจารํ เอกํ ขนฺธญฺจ วตฺถุญฺจ ปฏิจฺจ ตโย ขนฺธา วิตกฺโก จ ฯเปฯ เทฺว ขนฺเธ จ วตฺถุญฺจ ปฏิจฺจ เทฺว ขนฺธา วิตกฺโก จ. สวิตกฺก\-สวิจาเร ขนฺเธ จ มหาภูเต จ ปฏิจฺจ กฏตฺตารูปํ. (7)}

\proseitem{5}{อวิตกฺกวิจารมตฺตญฺจ อวิตกฺกอวิจารญฺจ ธมฺมํ ปฏิจฺจ สวิตกฺก\-สวิจาโร ธมฺโม อุปฺปชฺชติ เหตุปจฺจยา. ปฏิสนฺธิ\-กฺขเณ วิตกฺกญฺจ วตฺถุญฺจ ปฏิจฺจ สวิตกฺก\-สวิจารา ขนฺธา. (1)}

\prose{อวิตกฺกวิจารมตฺตญฺจ อวิตกฺกอวิจารญฺจ ธมฺมํ ปฏิจฺจ อวิตกฺก\-วิจารมตฺโต ธมฺโม อุปฺปชฺชติ เหตุปจฺจยา. อวิตกฺก\-วิจารมตฺตํ เอกํ ขนฺธญฺจ วิจารญฺจ ปฏิจฺจ ตโย ขนฺธา ฯเปฯ เทฺว ขนฺเธ จ วิจารญฺจ ปฏิจฺจ เทฺว ขนฺธา. ปฏิสนฺธิ\-กฺขเณ อวิตกฺก\-วิจารมตฺตํ เอกํ ขนฺธญฺจ วิจารญฺจ ปฏิจฺจ ตโย ขนฺธา ฯเปฯ เทฺว ขนฺเธ จ วิจารญฺจ ปฏิจฺจ เทฺว ขนฺธา. ปฏิสนฺธิ\-กฺขเณ อวิตกฺก\-วิจารมตฺตํ เอกํ ขนฺธญฺจ วตฺถุญฺจ ปฏิจฺจ ตโย ขนฺธา ฯเปฯ เทฺว ขนฺเธ จ วตฺถุญฺจ ปฏิจฺจ เทฺว ขนฺธา. (2)}

\prose{อวิตกฺกวิจารมตฺตญฺจ อวิตกฺกอวิจารญฺจ ธมฺมํ ปฏิจฺจ อวิตกฺก\-อวิจาโร ธมฺโม อุปฺปชฺชติ เหตุปจฺจยา. อวิตกฺก\-วิจารมตฺเต ขนฺเธ จ วิจารญฺจ ปฏิจฺจ จิตฺต\-สมุฏฺฐานํ รูปํ. อวิตกฺก\-วิจารมตฺเต ขนฺเธ จ มหาภูเต จ ปฏิจฺจ จิตฺต\-สมุฏฺฐานํ รูปํ. วิตกฺกญฺจ มหาภูเต จ ปฏิจฺจ จิตฺต\-สมุฏฺฐานํ รูปํ. \csromanfolio{6}ปฏิสนฺธิ\-กฺขเณ อวิตกฺก\-วิจารมตฺเต ขนฺเธ จ วิจารญฺจ ปฏิจฺจ กฏตฺตารูปํ. ปฏิสนฺธิ\-กฺขเณ อวิตกฺก\-วิจารมตฺเต ขนฺเธ จ มหาภูเต จ ปฏิจฺจ กฏตฺตารูปํ. ปฏิสนฺธิ\-กฺขเณ วิตกฺกญฺจ มหาภูเต จ ปฏิจฺจ กฏตฺตารูปํ. ปฏิสนฺธิ\-กฺขเณ อวิตกฺก\-วิจารมตฺเต ขนฺเธ จ วตฺถุญฺจ ปฏิจฺจ วิจาโร. (3)}

\prose{อวิตกฺกวิจารมตฺตญฺจ อวิตกฺกอวิจารญฺจ ธมฺมํ ปฏิจฺจ สวิตกฺก\-สวิจาโร จ อวิตกฺก\-อวิจาโร จ ธมฺมา อุปฺปชฺชนฺติ เหตุปจฺจยา. ปฏิสนฺธิ\-กฺขเณ วิตกฺกญฺจ วตฺถุญฺจ ปฏิจฺจ สวิตกฺก\-สวิจารา ขนฺธา. วิตกฺกญฺจ มหาภูเต จ ปฏิจฺจ กฏตฺตารูปํ. (4)}

\prose{อวิตกฺกวิจารมตฺตญฺจ อวิตกฺกอวิจารญฺจ ธมฺมํ ปฏิจฺจ อวิตกฺก\-วิจารมตฺโต จ อวิตกฺก\-อวิจาโร จ ธมฺมา อุปฺปชฺชนฺติ เหตุปจฺจยา. อวิตกฺก\-วิจารมตฺตํ เอกํ ขนฺธญฺจ วิจารญฺจ ปฏิจฺจ ตโย ขนฺธา จิตฺต\-สมุฏฺฐานญฺจ รูปํ ฯเปฯ เทฺว ขนฺเธ จ วิจารญฺจ ปฏิจฺจ เทฺว ขนฺธา จิตฺต\-สมุฏฺฐานญฺจ รูปํ. ปฏิสนฺธิ\-กฺขเณ อวิตกฺก\-วิจารมตฺตํ เอกํ ขนฺธญฺจ วิจารญฺจ ปฏิจฺจ ตโย ขนฺธา กฏตฺตา จ รูปํ ฯเปฯ เทฺว ขนฺเธ จ วิจารญฺจ ปฏิจฺจ เทฺว ขนฺธา กฏตฺตา จ รูปํ. ปฏิสนฺธิ\-กฺขเณ อวิตกฺก\-วิจารมตฺตํ เอกํ ขนฺธญฺจ วตฺถุญฺจ ปฏิจฺจ ตโย ขนฺธา วิจาโร จ ฯเปฯ เทฺว ขนฺเธ จ วตฺถุญฺจ ปฏิจฺจ เทฺว ขนฺธา วิจาโร จ. (5)}

\proseitem{6}{สวิตกฺกสวิจารญฺจ อวิตกฺกวิจารมตฺตญฺจ ธมฺมํ ปฏิจฺจ สวิตกฺก\-สวิจาโร ธมฺโม อุปฺปชฺชติ เหตุปจฺจยา. สวิตกฺก\-สวิจารํ เอกํ ขนฺธญฺจ วิตกฺกญฺจ ปฏิจฺจ ตโย ขนฺธา ฯเปฯ เทฺว ขนฺเธ จ วิตกฺกญฺจ ปฏิจฺจ เทฺว ขนฺธา. ปฏิสนฺธิ\-กฺขเณ สวิตกฺก\-สวิจารํ เอกํ ขนฺธญฺจ วิตกฺกญฺจ ปฏิจฺจ ตโย ขนฺธา ฯเปฯ เทฺว ขนฺเธ จ วิตกฺกญฺจ ปฏิจฺจ เทฺว ขนฺธา. (1)}

\prose{สวิตกฺกสวิจารญฺจ อวิตกฺกวิจารมตฺตญฺจ ธมฺมํ ปฏิจฺจ อวิตกฺก\-อวิจาโร ธมฺโม อุปฺปชฺชติ เหตุปจฺจยา. สวิตกฺก\-สวิจาเร ขนฺเธ จ วิตกฺกญฺจ ปฏิจฺจ จิตฺต\-สมุฏฺฐานํ รูปํ. ปฏิสนฺธิ\-กฺขเณ สวิตกฺก\-สวิจาเร ขนฺเธ จ วิตกฺกญฺจ ปฏิจฺจ กฏตฺตารูปํ. (2)}

\prose{สวิตกฺกสวิจารญฺจ อวิตกฺกวิจารมตฺตญฺจ ธมฺมํ ปฏิจฺจ สวิตกฺก\-สวิจาโร จ อวิตกฺก\-อวิจาโร จ ธมฺมา อุปฺปชฺชนฺติ เหตุปจฺจยา. สวิตกฺก\-สวิจารํ เอกํ ขนฺธญฺจ วิตกฺกญฺจ ปฏิจฺจ ตโย ขนฺธา จิตฺต\-สมุฏฺฐานญฺจ รูปํ ฯเปฯ เทฺว ขนฺเธ \csromanfolio{7}จ วิตกฺกญฺจ ปฏิจฺจ เทฺว ขนฺธา จิตฺต\-สมุฏฺฐานญฺจ รูปํ. ปฏิสนฺธิ\-กฺขเณ สวิตกฺก\-สวิจารํ เอกํ ขนฺธญฺจ วิตกฺกญฺจ ปฏิจฺจ ตโย ขนฺธา กฏตฺตา จ รูปํ ฯเปฯ เทฺว ขนฺเธ จ วิตกฺกญฺจ ปฏิจฺจ เทฺว ขนฺธา กฏตฺตา จ รูปํ. (3)}

\prose{สวิตกฺกสวิจารญฺจ อวิตกฺกวิจารมตฺตญฺจ อวิตกฺกอวิจารญฺจ ธมฺมํ ปฏิจฺจ สวิตกฺก\-สวิจาโร ธมฺโม อุปฺปชฺชติ เหตุปจฺจยา. ปฏิสนฺธิ\-กฺขเณ สวิตกฺก\-สวิจารํ เอกํ ขนฺธญฺจ วิตกฺกญฺจ วตฺถุญฺจ ปฏิจฺจ ตโย ขนฺธา ฯเปฯ เทฺว ขนฺเธ จ วิตกฺกญฺจ วตฺถุญฺจ ปฏิจฺจ เทฺว ขนฺธา. (1)}

\prose{สวิตกฺกสวิจารญฺจ อวิตกฺกวิจารมตฺตญฺจ อวิตกฺกอวิจารญฺจ ธมฺมํ ปฏิจฺจ อวิตกฺก\-อวิจาโร ธมฺโม อุปฺปชฺชติ เหตุปจฺจยา. สวิตกฺก\-สวิจาเร ขนฺเธ จ วิตกฺกญฺจ มหาภูเต จ ปฏิจฺจ จิตฺต\-สมุฏฺฐานํ รูปํ. ปฏิสนฺธิ\-กฺขเณ สวิตกฺก\-สวิจาเร ขนฺเธ จ วิตกฺกญฺจ มหาภูเต จ ปฏิจฺจ กฏตฺตารูปํ. (2)}

\prose{สวิตกฺกสวิจารญฺจ อวิตกฺกวิจารมตฺตญฺจ อวิตกฺกอวิจารญฺจ ธมฺมํ ปฏิจฺจ สวิตกฺก\-สวิจาโร จ อวิตกฺก\-อวิจาโร จ ธมฺมา อุปฺปชฺชนฺติ เหตุปจฺจยา. ปฏิสนฺธิ\-กฺขเณ สวิตกฺก\-สวิจารํ เอกํ ขนฺธญฺจ วิตกฺกญฺจ วตฺถุญฺจ ปฏิจฺจ ตโย ขนฺธา, ตโย ขนฺเธ จ วิตกฺกญฺจ วตฺถุญฺจ ปฏิจฺจ เอโก ขนฺโธ, เทฺว ขนฺเธ จ วิตกฺกญฺจ วตฺถุญฺจ ปฏิจฺจ เทฺว ขนฺธา. สวิตกฺก\-สวิจาเร ขนฺเธ จ วิตกฺกญฺจ มหาภูเต จ ปฏิจฺจ กฏตฺตารูปํ. (3)}

\csromanheader{2}{\textbf{อารมฺมณ}}
\proseitem{6}{สวิตกฺก\-สวิจารํ ธมฺมํ ปฏิจฺจ สวิตกฺก\-สวิจาโร ธมฺโม อุปฺปชฺชติ อารมฺมณ\-ปจฺจยา. สวิตกฺก\-สวิจารํ เอกํ ขนฺธํ ปฏิจฺจ ตโย ขนฺธา ฯเปฯ เทฺว ขนฺเธ ปฏิจฺจ เทฺว ขนฺธา. ปฏิสนฺธิ\-กฺขเณ ฯเปฯ. (1)}

\prose{สวิตกฺก\-สวิจารํ ธมฺมํ ปฏิจฺจ อวิตกฺก\-วิจารมตฺโต ธมฺโม อุปฺปชฺชติ อารมฺมณ\-ปจฺจยา. สวิตกฺก\-สวิจาเร ขนฺเธ ปฏิจฺจ วิตกฺโก. ปฏิสนฺธิ\-กฺขเณ ฯเปฯ. (2)}

\prose{สวิตกฺก\-สวิจารํ ธมฺมํ ปฏิจฺจ สวิตกฺก\-สวิจาโร จ อวิตกฺก\-วิจารมตฺโต จ ธมฺมา อุปฺปชฺชนฺติ อารมฺมณ\-ปจฺจยา. สวิตกฺก\-สวิจารํ เอกํ ขนฺธํ ปฏิจฺจ ตโย ขนฺธา วิตกฺโก จ ฯเปฯ เทฺว ขนฺเธ ปฏิจฺจ เทฺว ขนฺธา วิตกฺโก จ. ปฏิสนฺธิ\-กฺขเณ ฯเปฯ. (3)}

\proseitem{9}{\csromanfolio{8}อวิตกฺก\-วิจารมตฺตํ ธมฺมํ ปฏิจฺจ อวิตกฺก\-วิจารมตฺโต ธมฺโม อุปฺปชฺชติ อารมฺมณ\-ปจฺจยา. อวิตกฺก\-วิจารมตฺตํ เอกํ ขนฺธํ ปฏิจฺจ ตโย ขนฺธา ฯเปฯ เทฺว ขนฺเธ ปฏิจฺจ เทฺว ขนฺธา. ปฏิสนฺธิ\-กฺขเณ ฯเปฯ. (1)}

\prose{อวิตกฺก\-วิจารมตฺตํ ธมฺมํ ปฏิจฺจ สวิตกฺก\-สวิจาโร ธมฺโม อุปฺปชฺชติ อารมฺมณ\-ปจฺจยา. วิตกฺกํ ปฏิจฺจ สวิตกฺก\-สวิจารา ขนฺธา. ปฏิสนฺธิ\-กฺขเณ ฯเปฯ. (2)}

\prose{อวิตกฺก\-วิจารมตฺตํ ธมฺมํ ปฏิจฺจ อวิตกฺก\-อวิจาโร ธมฺโม อุปฺปชฺชติ อารมฺมณ\-ปจฺจยา. อวิตกฺก\-วิจารมตฺเต ขนฺเธ ปฏิจฺจ วิจาโร. ปฏิสนฺธิ\-กฺขเณ ฯเปฯ. (3)}

\prose{อวิตกฺก\-วิจารมตฺตํ ธมฺมํ ปฏิจฺจ อวิตกฺก\-วิจารมตฺโต จ อวิตกฺก\-อวิจาโร จ ธมฺมา อุปฺปชฺชนฺติ อารมฺมณ\-ปจฺจยา. อวิตกฺก\-วิจารมตฺตํ เอกํ ขนฺธํ ปฏิจฺจ ตโย ขนฺธา วิจาโร จ ฯเปฯ เทฺว ขนฺเธ ปฏิจฺจ เทฺว ขนฺธา วิจาโร จ. ปฏิสนฺธิ\-กฺขเณ ฯเปฯ. (4)}

\proseitem{10}{อวิตกฺก\-อวิจารํ ธมฺมํ ปฏิจฺจ อวิตกฺก\-อวิจาโร ธมฺโม อุปฺปชฺชติ อารมฺมณ\-ปจฺจยา. อวิตกฺก\-อวิจารํ เอกํ ขนฺธํ ปฏิจฺจ ตโย ขนฺธา ฯเปฯ เทฺว ขนฺเธ ปฏิจฺจ เทฺว ขนฺธา. ปฏิสนฺธิ\-กฺขเณ อวิตกฺก\-อวิจารํ เอกํ ขนฺธํ ปฏิจฺจ ตโย ขนฺธา ฯเปฯ เทฺว ขนฺเธ ปฏิจฺจ เทฺว ขนฺธา. ปฏิสนฺธิ\-กฺขเณ วตฺถุํ ปฏิจฺจ ขนฺธา, วตฺถุํ ปฏิจฺจ วิจาโร. (1)}

\prose{อวิตกฺก\-อวิจารํ ธมฺมํ ปฏิจฺจ สวิตกฺก\-สวิจาโร ธมฺโม อุปฺปชฺชติ อารมฺมณ\-ปจฺจยา. ปฏิสนฺธิ\-กฺขเณ วตฺถุํ ปฏิจฺจ สวิตกฺก\-สวิจารา ขนฺธา. (2)}

\prose{อวิตกฺก\-อวิจารํ ธมฺมํ ปฏิจฺจ อวิตกฺก\-วิจารมตฺโต ธมฺโม อุปฺปชฺชติ อารมฺมณ\-ปจฺจยา. วิจารํ ปฏิจฺจ อวิตกฺก\-วิจารมตฺตา ขนฺธา. ปฏิสนฺธิ\-กฺขเณ วิจารํ ปฏิจฺจ อวิตกฺก\-วิจารมตฺตา ขนฺธา. ปฏิสนฺธิ\-กฺขเณ วตฺถุํ ปฏิจฺจ อวิตกฺก\-วิจารมตฺตา ขนฺธา. ปฏิสนฺธิ\-กฺขเณ วตฺถุํ ปฏิจฺจ วิตกฺโก. (3)}

\prose{อวิตกฺก\-อวิจารํ ธมฺมํ ปฏิจฺจ อวิตกฺก\-วิจารมตฺโต จ อวิตกฺก\-อวิจาโร จ ธมฺมา อุปฺปชฺชนฺติ อารมฺมณ\-ปจฺจยา. ปฏิสนฺธิ\-กฺขเณ วตฺถุํ ปฏิจฺจ อวิตกฺก\-วิจารมตฺตา ขนฺธา วิจาโร จ. (4)}

\prose{\csromanfolio{9}อวิตกฺก\-อวิจารํ ธมฺมํ ปฏิจฺจ สวิตกฺก\-สวิจาโร จ อวิตกฺก\-วิจารมตฺโต จ ธมฺมา อุปฺปชฺชนฺติ อารมฺมณ\-ปจฺจยา. ปฏิสนฺธิ\-กฺขเณ วตฺถุํ ปฏิจฺจ สวิตกฺก\-สวิจารา ขนฺธา จ วิตกฺโก จ. (5)}

\proseitem{11}{สวิตกฺกสวิจารญฺจ อวิตกฺกอวิจารญฺจ ธมฺมํ ปฏิจฺจ สวิตกฺก\-สวิจาโร ธมฺโม อุปฺปชฺชติ อารมฺมณ\-ปจฺจยา. ปฏิสนฺธิ\-กฺขเณ สวิตกฺก\-สวิจารํ เอกํ ขนฺธญฺจ วตฺถุญฺจ ปฏิจฺจ ตโย ขนฺธา ฯเปฯ เทฺว ขนฺเธ จ วตฺถุญฺจ ปฏิจฺจ เทฺว ขนฺธา. (1)}

\prose{สวิตกฺกสวิจารญฺจ อวิตกฺกอวิจารญฺจ ธมฺมํ ปฏิจฺจ อวิตกฺก\-วิจารมตฺโต ธมฺโม อุปฺปชฺชติ อารมฺมณ\-ปจฺจยา. ปฏิสนฺธิ\-กฺขเณ สวิตกฺก\-สวิจาเร ขนฺเธ จ วตฺถุญฺจ ปฏิจฺจ วิตกฺโก. (2)}

\prose{สวิตกฺกสวิจารญฺจ อวิตกฺกอวิจารญฺจ ธมฺมํ ปฏิจฺจ สวิตกฺก\-สวิจาโร จ อวิตกฺก\-วิจารมตฺโต จ ธมฺมา อุปฺปชฺชนฺติ อารมฺมณ\-ปจฺจยา. ปฏิสนฺธิ\-กฺขเณ สวิตกฺก\-สวิจารํ เอกํ ขนฺธญฺจ วตฺถุญฺจ ปฏิจฺจ ตโย ขนฺธา วิตกฺโก จ ฯเปฯ เทฺว ขนฺเธ จ วตฺถุญฺจ ปฏิจฺจ เทฺว ขนฺธา วิตกฺโก จ. (3)}

\proseitem{12}{อวิตกฺกวิจารมตฺตญฺจ อวิตกฺกอวิจารญฺจ ธมฺมํ ปฏิจฺจ สวิตกฺก\-สวิจาโร ธมฺโม อุปฺปชฺชติ อารมฺมณ\-ปจฺจยา. ปฏิสนฺธิ\-กฺขเณ วิตกฺกญฺจ วตฺถุญฺจ ปฏิจฺจ สวิตกฺก\-สวิจารา ขนฺธา. (1)}

\prose{อวิตกฺกวิจารมตฺตญฺจ อวิตกฺกอวิจารญฺจ ธมฺมํ ปฏิจฺจ อวิตกฺก\-วิจารมตฺโต ธมฺโม อุปฺปชฺชติ อารมฺมณ\-ปจฺจยา. อวิตกฺก\-วิจารมตฺตํ เอกํ ขนฺธญฺจ วิจารญฺจ ปฏิจฺจ ตโย ขนฺธา ฯเปฯ เทฺว ขนฺเธ จ วิจารญฺจ ปฏิจฺจ เทฺว ขนฺธา. ปฏิสนฺธิ\-กฺขเณ อวิตกฺก\-วิจารมตฺตํ เอกํ ขนฺธญฺจ วิจารญฺจ ปฏิจฺจ ตโย ขนฺธา ฯเปฯ เทฺว ขนฺเธ จ วิจารญฺจ ปฏิจฺจ เทฺว ขนฺธา. ปฏิสนฺธิ\-กฺขเณ อวิตกฺก\-วิจารมตฺตํ เอกํ ขนฺธญฺจ วตฺถุญฺจ ปฏิจฺจ ตโย ขนฺธา ฯเปฯ เทฺว ขนฺเธ จ วตฺถุญฺจ ปฏิจฺจ เทฺว ขนฺธา. (2)}

\prose{อวิตกฺกวิจารมตฺตญฺจ อวิตกฺกอวิจารญฺจ ธมฺมํ ปฏิจฺจ อวิตกฺก\-อวิจาโร ธมฺโม อุปฺปชฺชติ อารมฺมณ\-ปจฺจยา. ปฏิสนฺธิ\-กฺขเณ อวิตกฺก\-วิจารมตฺเต ขนฺเธ จ วตฺถุญฺจ ปฏิจฺจ วิจาโร. (3)}

\prose{อวิตกฺกวิจารมตฺตญฺจ อวิตกฺกอวิจารญฺจ ธมฺมํ ปฏิจฺจ อวิตกฺก\-วิจารมตฺโต จ อวิตกฺก\-อวิจาโร จ ธมฺมา อุปฺปชฺชนฺติ อารมฺมณ\-ปจฺจยา. ปฏิสนฺธิ\-กฺขเณ \csromanfolio{10}อวิตกฺก\-วิจารมตฺตํ เอกํ ขนฺธญฺจ วตฺถุญฺจ ปฏิจฺจ ตโย ขนฺธา วิจาโร จ ฯเปฯ เทฺว ขนฺเธ จ วตฺถุญฺจ ปฏิจฺจ เทฺว ขนฺธา วิจาโร จ. (4)}

\prose{สวิตกฺกสวิจารญฺจ อวิตกฺกวิจารมตฺตญฺจ ธมฺมํ ปฏิจฺจ สวิตกฺก\-สวิจาโร ธมฺโม อุปฺปชฺชติ อารมฺมณ\-ปจฺจยา. สวิตกฺก\-สวิจารํ เอกํ ขนฺธญฺจ วิตกฺกญฺจ ปฏิจฺจ ตโย ขนฺธา ฯเปฯ เทฺว ขนฺเธ จ วิตกฺกญฺจ ปฏิจฺจ เทฺว ขนฺธา. ปฏิสนฺธิ\-กฺขเณ ฯเปฯ. (1)}

\prose{สวิตกฺกสวิจารญฺจ อวิตกฺกวิจารมตฺตญฺจ อวิตกฺกอวิจารญฺจ ธมฺมํ ปฏิจฺจ สวิตกฺก\-สวิจาโร ธมฺโม อุปฺปชฺชติ อารมฺมณ\-ปจฺจยา. ปฏิสนฺธิ\-กฺขเณ สวิตกฺก\-สวิจารํ เอกํ ขนฺธญฺจ วิตกฺกญฺจ วตฺถุญฺจ ปฏิจฺจ ตโย ขนฺธา ฯเปฯ เทฺว ขนฺเธ จ วิตกฺกญฺจ วตฺถุญฺจ ปฏิจฺจ เทฺว ขนฺธา. (1)}

\prose{(เทฺว ปจฺจยา สชฺฌาย\-มคฺเคน วิภตฺตา, เอวํ อวเสสา วีสติปจฺจยา วิภชิตพฺพา.)}

\csromanheader{2}{\textbf{วิปฺปยุตฺต}}
\proseitem{14}{สวิตกฺก\-สวิจารํ ธมฺมํ ปฏิจฺจ สวิตกฺก\-สวิจาโร ธมฺโม อุปฺปชฺชติ วิปฺปยุตฺต\-ปจฺจยา. สวิตกฺก\-สวิจารํ เอกํ ขนฺธํ ปฏิจฺจ ตโย ขนฺธา ฯเปฯ. วตฺถุํ วิปฺปยุตฺต\-ปจฺจยา. ปฏิสนฺธิ\-กฺขเณ สวิตกฺก\-สวิจารํ เอกํ ขนฺธํ ปฏิจฺจ ตโย ขนฺธา ฯเปฯ. วตฺถุํ วิปฺปยุตฺต\-ปจฺจยา. (1)}

\prose{สวิตกฺก\-สวิจารํ ธมฺมํ ปฏิจฺจ อวิตกฺก\-วิจารมตฺโต ธมฺโม อุปชฺชติ วิปฺปยุตฺต\-ปจฺจยา. สวิตกฺก\-สวิจาเร ขนฺเธ ปฏิจฺจ วิตกฺโก. วตฺถุํ วิปฺปยุตฺต\-ปจฺจยา. ปฏิสนฺธิ\-กฺขเณ ฯเปฯ. (2)}

\prose{สวิตกฺก\-สวิจารํ ธมฺมํ ปฏิจฺจ อวิตกฺก\-อวิจาโร ธมฺโม อุปฺปชฺชติ วิปฺปยุตฺต\-ปจฺจยา. สวิตกฺก\-สวิจาเร ขนฺเธ ปฏิจฺจ จิตฺต\-สมุฏฺฐานํ รูปํ. ขนฺเธ วิปฺปยุตฺต\-ปจฺจยา. ปฏิสนฺธิ\-กฺขเณ ฯเปฯ. (3)}

\prose{สวิตกฺก\-สวิจารํ ธมฺมํ ปฏิจฺจ สวิตกฺก\-สวิจาโร จ อวิตกฺก\-อวิจาโร จ ธมฺมา อุปฺปชฺชนฺติ วิปฺปยุตฺต\-ปจฺจยา. สวิตกฺก\-สวิจารํ เอกํ ขนฺธํ ปฏิจฺจ ตโย ขนฺธา จิตฺต\-สมุฏฺฐานญฺจ รูปํ ฯเปฯ. ขนฺธา วตฺถุํ วิปฺปยุตฺต\-ปจฺจยา. จิตฺต\-สมุฏฺฐานํ รูปํ ขนฺเธ วิปฺปยุตฺต\-ปจฺจยา. ปฏิสนฺธิ\-กฺขเณ ฯเปฯ. (4)}

\prose{\csromanfolio{11}สวิตกฺก\-สวิจารํ ธมฺมํ ปฏิจฺจ อวิตกฺก\-วิจารมตฺโต จ อวิตกฺก\-อวิจาโร จ ธมฺมา อุปฺปชฺชนฺติ วิปฺปยุตฺต\-ปจฺจยา. สวิตกฺก\-สวิจาเร ขนฺเธ ปฏิจฺจ วิตกฺโก จ จิตฺต\-สมุฏฺฐานญฺจ รูปํ. วิตกฺโก วตฺถุํ วิปฺปยุตฺต\-ปจฺจยา. จิตฺต\-สมุฏฺฐานํ รูปํ ขนฺเธ วิปฺปยุตฺต\-ปจฺจยา. ปฏิสนฺธิ\-กฺขเณ ฯเปฯ. (5)}

\prose{สวิตกฺก\-สวิจารํ ธมฺมํ ปฏิจฺจ สวิตกฺก\-สวิจาโร จ อวิตกฺก\-วิจารมตฺโต จ ธมฺมา อุปฺปชฺชนฺติ วิปฺปยุตฺต\-ปจฺจยา. สวิตกฺก\-สวิจารํ เอกํ ขนฺธํ ปฏิจฺจ ตโย ขนฺธา วิตกฺโก จ ฯเปฯ เทฺว ขนฺเธ ปฏิจฺจ เทฺว ขนฺธา วิตกฺโก จ. วตฺถุํ วิปฺปยุตฺต\-ปจฺจยา. ปฏิสนฺธิ\-กฺขเณ ฯเปฯ. (6)}

\prose{สวิตกฺก\-สวิจารํ ธมฺมํ ปฏิจฺจ สวิตกฺก\-สวิจาโร จ อวิตกฺก\-วิจารมตฺโต จ อวิตกฺก\-อวิจาโร จ ธมฺมา อุปฺปชฺชนฺติ วิปฺปยุตฺต\-ปจฺจยา. สวิตกฺก\-สวิจารํ เอกํ ขนฺธํ ปฏิจฺจ ตโย ขนฺธา วิตกฺโก จ จิตฺต\-สมุฏฺฐานญฺจ รูปํ ฯเปฯ เทฺว ขนฺเธ ปฏิจฺจ เทฺว ขนฺธา วิตกฺโก จ จิตฺต\-สมุฏฺฐานญฺจ รูปํ. ขนฺธา จ วิตกฺโก จ วตฺถุํ วิปฺปยุตฺต\-ปจฺจยา. จิตฺต\-สมุฏฺฐานํ รูปํ ขนฺเธ วิปฺปยุตฺต\-ปจฺจยา. ปฏิสนฺธิ\-กฺขเณ ฯเปฯ. (7)}

\proseitem{15}{อวิตกฺก\-วิจารมตฺตํ ธมฺมํ ปฏิจฺจ อวิตกฺก\-วิจารมตฺโต ธมฺโม อุปฺปชฺชติ ฯเปฯ. อวิตกฺก\-วิจารมตฺตํ เอกํ ขนฺธํ ปฏิจฺจ ตโย ขนฺธา ฯเปฯ เทฺว ขนฺเธ ฯเปฯ. วตฺถุํ วิปฺปยุตฺต\-ปจฺจยา. ปฏิสนฺธิ\-กฺขเณ ฯเปฯ. (1)}

\prose{อวิตกฺก\-วิจารมตฺตํ ธมฺมํ ปฏิจฺจ สวิตกฺก\-สวิจาโร ธมฺโม อุปฺปชฺชติ วิปฺปยุตฺต\-ปจฺจยา. วิตกฺกํ ปฏิจฺจ สวิตกฺก\-สวิจารา ขนฺธา. วตฺถุํ วิปฺปยุตฺต\-ปจฺจยา. ปฏิสนฺธิ\-กฺขเณ ฯเปฯ. (2)}

\prose{อวิตกฺก\-วิจารมตฺตํ ธมฺมํ ปฏิจฺจ อวิตกฺก\-อวิจาโร ธมฺโม อุปฺปชฺชติ วิปฺปยุตฺต\-ปจฺจยา. อวิตกฺก\-วิจารมตฺเต ขนฺเธ ปฏิจฺจ วิจาโร จ จิตฺต\-สมุฏฺฐานญฺจ รูปํ, วิจาโร วตฺถุํ วิปฺปยุตฺต\-ปจฺจยา. จิตฺต\-สมุฏฺฐานํ รูปํ ขนฺเธ วิปฺปยุตฺต\-ปจฺจยา. วิตกฺกํ ปฏิจฺจ จิตฺต\-สมุฏฺฐานํ รูปํ, วิตกฺกํ วิปฺปยุตฺต\-ปจฺจยา. ปฏิสนฺธิ\-กฺขเณ ฯเปฯ. (3)}

\prose{อวิตกฺก\-วิจารมตฺตํ ธมฺมํ ปฏิจฺจ สวิตกฺก\-สวิจาโร จ อวิตกฺก\-อวิจาโร จ ธมฺมา อุปฺปชฺชนฺติ วิปฺปยุตฺต\-ปจฺจยา. วิตกฺกํ ปฏิจฺจ สวิตกฺก\-สวิจารา ขนฺธา จิตฺต\-สมุฏฺฐานญฺจ รูปํ, ขนฺธา วตฺถุํ วิปฺปยุตฺต\-ปจฺจยา. จิตฺต\-สมุฏฺฐานํ รูปํ วิตกฺกํ วิปฺปยุตฺต\-ปจฺจยา. ปฏิสนฺธิ\-กฺขเณ ฯเปฯ. (4)}

\prose{\csromanfolio{12}อวิตกฺก\-วิจารมตฺตํ ธมฺมํ ปฏิจฺจ อวิตกฺก\-วิจารมตฺโต จ อวิตกฺก\-อวิจาโร จ ธมฺมา อุปฺปชฺชนฺติ วิปฺปยุตฺต\-ปจฺจยา. อวิตกฺก\-วิจารมตฺตํ เอกํ ขนฺธํ ปฏิจฺจ ตโย ขนฺธา วิจาโร จ จิตฺต\-สมุฏฺฐานญฺจ รูปํ ฯเปฯ เทฺว ขนฺเธ ปฏิจฺจ เทฺว ขนฺธา วิจาโร จ จิตฺต\-สมุฏฺฐานญฺจ รูปํ. ขนฺธา จ วิจาโร จ วตฺถุํ วิปฺปยุตฺต\-ปจฺจยา. จิตฺต\-สมุฏฺฐานํ รูปํ ขนฺเธ วิปฺปยุตฺต\-ปจฺจยา. ปฏิสนฺธิ\-กฺขเณ ฯเปฯ. (5)}

\proseitem{16}{อวิตกฺก\-อวิจารํ ธมฺมํ ปฏิจฺจ อวิตกฺก\-อวิจาโร ธมฺโม อุปฺปชฺชติ วิปฺปยุตฺต\-ปจฺจยา. อวิตกฺก\-อวิจารํ เอกํ ขนฺธํ ปฏิจฺจ ตโย ขนฺธา จิตฺต\-สมุฏฺฐานญฺจ รูปํ ฯเปฯ เทฺว ขนฺเธ ฯเปฯ. ขนฺธา วตฺถุํ วิปฺปยุตฺต\-ปจฺจยา, จิตฺต\-สมุฏฺฐานํ รูปํ ขนฺเธ วิปฺปยุตฺต\-ปจฺจยา. วิจารํ ปฏิจฺจ จิตฺต\-สมุฏฺฐานํ รูปํ, วิจารํ วิปฺปยุตฺต\-ปจฺจยา. ปฏิสนฺธิ\-กฺขเณ ฯเปฯ. ปฏิสนฺธิ\-กฺขเณ วิจารํ ปฏิจฺจ กฏตฺตารูปํ, วิจารํ วิปฺปยุตฺต\-ปจฺจยา. ขนฺเธ ปฏิจฺจ วตฺถุ, วตฺถุํ ปฏิจฺจ ขนฺธา, ขนฺธา วตฺถุํ วิปฺปยุตฺต\-ปจฺจยา. วตฺถุ ขนฺเธ วิปฺปยุตฺต\-ปจฺจยา. วิจารํ ปฏิจฺจ วตฺถุ, วตฺถุํ ปฏิจฺจ วิจาโร, วิจาโร วตฺถุํ วิปฺปยุตฺต\-ปจฺจยา. วตฺถุ วิจารํ วิปฺปยุตฺต\-ปจฺจยา. เอกํ มหาภูตํ ปฏิจฺจ ตโย มหาภูตา ฯเปฯ มหาภูเต ปฏิจฺจ จิตฺต\-สมุฏฺฐานํ รูปํ กฏตฺตารูปํ อุปาทารูปํ. ขนฺเธ วิปฺปยุตฺต\-ปจฺจยา. (1)}

\prose{อวิตกฺก\-อวิจารํ ธมฺมํ ปฏิจฺจ สวิตกฺก\-สวิจาโร ธมฺโม อุปฺปชฺชติ วิปฺปยุตฺต\-ปจฺจยา. ปฏิสนฺธิ\-กฺขเณ วตฺถุํ ปฏิจฺจ สวิตกฺก\-สวิจารา ขนฺธา, วตฺถุํ วิปฺปยุตฺต\-ปจฺจยา. (2)}

\prose{อวิตกฺก\-อวิจารํ ธมฺมํ ปฏิจฺจ อวิตกฺก\-วิจารมตฺโต ธมฺโม อุปฺปชฺชติ วิปฺปยุตฺต\-ปจฺจยา. วิจารํ ปฏิจฺจ อวิตกฺก\-วิจารมตฺตา ขนฺธา, วตฺถุํ วิปฺปยุตฺต\-ปจฺจยา. ปฏิสนฺธิ\-กฺขเณ วิจารํ ปฏิจฺจ อวิตกฺก\-วิจารมตฺตา ขนฺธา, วตฺถุํ วิปฺปยุตฺต\-ปจฺจยา. ปฏิสนฺธิ\-กฺขเณ วตฺถุํ ปฏิจฺจ อวิตกฺก\-วิจารมตฺตา ขนฺธา, วตฺถุํ วิปฺปยุตฺต\-ปจฺจยา. ปฏิสนฺธิ\-กฺขเณ วตฺถุํ ปฏิจฺจ อวิตกฺก\-วิจารมตฺตา ขนฺธา, วตฺถุํ วิปฺปยุตฺต\-ปจฺจยา. ปฏิสนฺธิ\-กฺขเณ วตฺถุํ ปฏิจฺจ วิตกฺโก, วตฺถุํ วิปฺปยุตฺต\-ปจฺจยา. (3)}

\prose{อวิตกฺก\-อวิจารํ ธมฺมํ ปฏิจฺจ สวิตกฺก\-สวิจาโร จ อวิตกฺก\-อวิจาโร จ ธมฺมา อุปฺปชฺชนฺติ วิปฺปยุตฺต\-ปจฺจยา. ปฏิสนฺธิ\-กฺขเณ วตฺถุํ ปฏิจฺจ สวิตกฺก\-สวิจารา ขนฺธา, มหาภูเต ปฏิจฺจ กฏตฺตารูปํ, ขนฺธา วตฺถุํ วิปฺปยุตฺต\-ปจฺจยา. กฏตฺตารูปํ ขนฺเธ วิปฺปยุตฺต\-ปจฺจยา. (4)}

\prose{\csromanfolio{13}อวิตกฺก\-อวิจารํ ธมฺมํ ปฏิจฺจ อวิตกฺก\-วิจารมตฺโต จ อวิตกฺก\-อวิจาโร จ ธมฺมา อุปฺปชฺชนฺติ วิปฺปยุตฺต\-ปจฺจยา. วิจารํ ปฏิจฺจ อวิตกฺก\-วิจารมตฺตา ขนฺธา จ จิตฺต\-สมุฏฺฐานญฺจ รูปํ ฯเปฯ ขนฺธา วตฺถุํ วิปฺปยุตฺต\-ปจฺจยา. จิตฺต\-สมุฏฺฐานํ รูปํ วิจารํ วิปฺปยุตฺต\-ปจฺจยา. ปฏิสนฺธิ\-กฺขเณ วิจารํ ปฏิจฺจ อวิตกฺก\-วิจารมตฺตา ขนฺธา จ กฏตฺตา จ รูปํ. ขนฺธา วตฺถุํ วิปฺปยุตฺต\-ปจฺจยา. กฏตฺตารูปํ วิจารํ วิปฺปยุตฺต\-ปจฺจยา. ปฏิสนฺธิ\-กฺขเณ วตฺถุํ ปฏิจฺจ อวิตกฺก\-วิจารมตฺตา ขนฺธา, มหาภูเต ปฏิจฺจ กฏตฺตารูปํ, ขนฺธา วตฺถุํ วิปฺปยุตฺต\-ปจฺจยา. กฏตฺตารูปํ ขนฺเธ วิปฺปยุตฺต\-ปจฺจยา. ปฏิสนฺธิ\-กฺขเณ วตฺถุํ ปฏิจฺจ วิตกฺโก, มหาภูเต ปฏิจฺจ กฏตฺตารูปํ, วิตกฺโก วตฺถุํ วิปฺปยุตฺต\-ปจฺจยา. กฏตฺตารูปํ ขนฺเธ วิปฺปยุตฺต\-ปจฺจยา. ปฏิสนฺธิ\-กฺขเณ วตฺถุํ ปฏิจฺจ อวิตกฺก\-วิจารมตฺตา ขนฺธา จ วิจาโร จ, วตฺถุํ วิปฺปยุตฺต\-ปจฺจยา. (5)}

\prose{อวิตกฺก\-อวิจารํ ธมฺมํ ปฏิจฺจ สวิตกฺก\-สวิจาโร จ อวิตกฺก\-วิจารมตฺโต จ ธมฺมา อุปฺปชฺชนฺติ วิปฺปยุตฺต\-ปจฺจยา. ปฏิสนฺธิ\-กฺขเณ วตฺถุํ ปฏิจฺจ สวิตกฺก\-สวิจารา ขนฺธา จ วิตกฺโก จ, วตฺถุํ วิปฺปยุตฺต\-ปจฺจยา. (6)}

\prose{อวิตกฺก\-อวิจารํ ธมฺมํ ปฏิจฺจ สวิตกฺก\-สวิจาโร จ อวิตกฺก\-วิจารมตฺโต จ อวิตกฺก\-อวิจาโร จ ธมฺมา อุปฺปชฺชนฺติ วิปฺปยุตฺต\-ปจฺจยา. ปฏิสนฺธิ\-กฺขเณ วตฺถุํ ปฏิจฺจ สวิตกฺก\-สวิจารา ขนฺธา จ วิตกฺโก จ, มหาภูเต ปฏิจฺจ กฏตฺตารูปํ, ขนฺธา จ วิตกฺโก จ, วตฺถุํ วิปฺปยุตฺต\-ปจฺจยา. กฏตฺตารูปํ ขนฺเธ วิปฺปยุตฺต\-ปจฺจยา. (7)}

\proseitem{17}{สวิตกฺกสวิจารญฺจ อวิตกฺกอวิจารญฺจ ธมฺมํ ปฏิจฺจ สวิตกฺก\-สวิจาโร ธมฺโม อุปฺปชฺชติ วิปฺปยุตฺต\-ปจฺจยา. ปฏิสนฺธิ\-กฺขเณ สวิตกฺก\-สวิจารํ เอกํ ขนฺธญฺจ วตฺถุญฺจ ปฏิจฺจ ตโย ขนฺธา, วตฺถุํ วิปฺปยุตฺต\-ปจฺจยา. (1)}

\prose{สวิตกฺกสวิจารญฺจ อวิตกฺกอวิจารญฺจ ธมฺมํ ปฏิจฺจ อวิตกฺก\-วิจารมตฺโต ธมฺโม ฯเปฯ. ปฏิสนฺธิ\-กฺขเณ สวิตกฺก\-สวิจาเร ขนฺเธ จ วตฺถุญฺจ ปฏิจฺจ วิตกฺโก, วตฺถุํ วิปฺปยุตฺต\-ปจฺจยา. (2)}

\prose{สวิตกฺกสวิจารญฺจ อวิตกฺกอวิจารญฺจ ธมฺมํ ปฏิจฺจ อวิตกฺก\-อวิจาโร ธมฺโม ฯเปฯ. สวิตกฺก\-สวิจาเร ขนฺเธ จ มหาภูเต จ ปฏิจฺจ จิตฺต\-สมุฏฺฐานํ รูปํ, ขนฺเธ วิปฺปยุตฺต\-ปจฺจยา. ปฏิสนฺธิ\-กฺขเณ ฯเปฯ. (3)}

\prose{\csromanfolio{14}สวิตกฺกสวิจารญฺจ อวิตกฺกอวิจารญฺจ ธมฺมํ ปฏิจฺจ สวิตกฺก\-สวิจาโร จ อวิตกฺก\-อวิจาโร จ ธมฺมา ฯเปฯ. ปฏิสนฺธิ\-กฺขเณ สวิตกฺก\-สวิจารํ เอกํ ขนฺธญฺจ วตฺถุญฺจ ปฏิจฺจ ตโย ขนฺธา ฯเปฯ เทฺว ขนฺเธ จ วตฺถุญฺจ ปฏิจฺจ เทฺว ขนฺธา, สวิตกฺก\-สวิจาเร ขนฺเธ จ มหาภูเต จ ปฏิจฺจ กฏตฺตารูปํ, ขนฺธา วตฺถุํ วิปฺปยุตฺต\-ปจฺจยา. กฏตฺตารูปํ ขนฺเธ วิปฺปยุตฺต\-ปจฺจยา. (4)}

\prose{สวิตกฺกสวิจารญฺจ อวิตกฺกอวิจารญฺจ ธมฺมํ ปฏิจฺจ อวิตกฺก\-วิจารมตฺโต จ อวิตกฺก\-อวิจาโร จ ธมฺมา ฯเปฯ. ปฏิสนฺธิ\-กฺขเณ สวิตกฺก\-สวิจาเร ขนฺเธ จ วตฺถุญฺจ ปฏิจฺจ วิตกฺโก. สวิตกฺก\-สวิจาเร ขนฺเธ จ มหาภูเต จ ปฏิจฺจ กฏตฺตารูปํ. วิตกฺโก วตฺถุํ วิปฺปยุตฺต\-ปจฺจยา. กฏตฺตารูปํ ขนฺเธ วิปฺปยุตฺต\-ปจฺจยา. (5)}

\prose{สวิตกฺกสวิจารญฺจ อวิตกฺกอวิจารญฺจ ธมฺมํ ปฏิจฺจ สวิตกฺก\-สวิจาโร จ อวิตกฺก\-วิจารมตฺโต จ ธมฺมา ฯเปฯ. ปฏิสนฺธิ\-กฺขเณ สวิตกฺก\-สวิจารํ เอกํ ขนฺธญฺจ วตฺถุญฺจ ปฏิจฺจ ตโย ขนฺธา วิตกฺโก จ ฯเปฯ เทฺว ขนฺเธ จ วตฺถุญฺจ ปฏิจฺจ เทฺว ขนฺธา วิตกฺโก จ, วตฺถุํ วิปฺปยุตฺต\-ปจฺจยา. (6)}

\prose{สวิตกฺกสวิจารญฺจ อวิตกฺกอวิจารญฺจ ธมฺมํ ปฏิจฺจ สวิตกฺก\-สวิจาโร จ อวิตกฺก\-วิจารมตฺโต จ อวิตกฺก\-อวิจาโร จ ธมฺมา ฯเปฯ. ปฏิสนฺธิ\-กฺขเณ สวิตกฺก\-สวิจารํ เอกํ ขนฺธญฺจ วตฺถุญฺจ ปฏิจฺจ ตโย ขนฺธา วิตกฺโก จ ฯเปฯ เทฺว ขนฺเธ จ วตฺถุญฺจ ปฏิจฺจ เทฺว ขนฺธา วิตกฺโก จ. สวิตกฺก\-สวิจาเร ขนฺเธ จ มหาภูเต จ ปฏิจฺจ กฏตฺตารูปํ. ขนฺธา จ วิตกฺโก จ, วตฺถุํ วิปฺปยุตฺต\-ปจฺจยา. กฏตฺตารูปํ ขนฺเธ วิปฺปยุตฺต\-ปจฺจยา. (7)}

\proseitem{18}{อวิตกฺกวิจารมตฺตญฺจ อวิตกฺกอวิจารญฺจ ธมฺมํ ปฏิจฺจ สวิตกฺก\-สวิจาโร ธมฺโม ฯเปฯ. ปฏิสนฺธิ\-กฺขเณ วิตกฺกญฺจ วตฺถุญฺจ ปฏิจฺจ สวิตกฺก\-สวิจารา ขนฺธา, วตฺถุํ วิปฺปยุตฺต\-ปจฺจยา. (1)}

\prose{อวิตกฺกวิจารมตฺตญฺจ อวิตกฺกอวิจารญฺจ ธมฺมํ ปฏิจฺจ อวิตกฺก\-วิจารมตฺโต ธมฺโม ฯเปฯ. อวิตกฺก\-วิจารมตฺตํ เอกํ ขนฺธญฺจ วิจารญฺจ ปฏิจฺจ ตโย ขนฺธา ฯเปฯ เทฺว ขนฺเธ จ ฯเปฯ วตฺถุํ วิปฺปยุตฺต\-ปจฺจยา. ปฏิสนฺธิ\-กฺขเณ อวิตกฺก\-วิจารมตฺตํ เอกํ ขนฺธญฺจ วิจารญฺจ ปฏิจฺจ ตโย ขนฺธา ฯเปฯ เทฺว ขนฺเธ จ ฯเปฯ วตฺถุํ วิปฺปยุตฺต\-ปจฺจยา. ปฏิสนฺธิ\-กฺขเณ อวิตกฺก\-วิจารมตฺตํ เอกํ ขนฺธญฺจ วตฺถุญฺจ ปฏิจฺจ ตโย ขนฺธา ฯเปฯ เทฺว ขนฺเธ จ ฯเปฯ วตฺถุํ วิปฺปยุตฺต\-ปจฺจยา. (2)}

\prose{\csromanfolio{15}อวิตกฺกวิจารมตฺตญฺจ อวิตกฺกอวิจารญฺจ ธมฺมํ ปฏิจฺจ อวิตกฺก\-อวิจาโร ธมฺโม ฯเปฯ. อวิตกฺก\-วิจารมตฺเต ขนฺเธ จ วิจารญฺจ ปฏิจฺจ จิตฺต\-สมุฏฺฐานํ รูปํ. ขนฺเธ จ วิจารญฺจ วิปฺปยุตฺต\-ปจฺจยา. อวิตกฺก\-วิจารมตฺเต ขนฺเธ จ มหาภูเต จ ปฏิจฺจ จิตฺต\-สมุฏฺฐานํ รูปํ. ขนฺเธ วิปฺปยุตฺต\-ปจฺจยา. วิตกฺกญฺจ มหาภูเต จ ปฏิจฺจ จิตฺต\-สมุฏฺฐานํ รูปํ. วิตกฺกํ วิปฺปยุตฺต\-ปจฺจยา. ปฏิสนฺธิ\-กฺขเณ อวิตกฺก\-วิจารมตฺเต ขนฺเธ จ วิจารญฺจ ปฏิจฺจ กฏตฺตารูปํ. ขนฺเธ จ วิจารญฺจ วิปฺปยุตฺต\-ปจฺจยา. ปฏิสนฺธิ\-กฺขเณ อวิตกฺก\-วิจารมตฺเต ขนฺเธ จ มหาภูเต จ ปฏิจฺจ กฏตฺตารูปํ. ขนฺเธ วิปฺปยุตฺต\-ปจฺจยา. ปฏิสนฺธิ\-กฺขเณ วิตกฺกญฺจ มหาภูเต จ ปฏิจฺจ กฏตฺตารูปํ. วิตกฺกํ วิปฺปยุตฺต\-ปจฺจยา. ปฏิสนฺธิ\-กฺขเณ อวิตกฺก\-วิจารมตฺเต ขนฺเธ จ วตฺถุญฺจ ปฏิจฺจ วิจาโร. วตฺถุํ วิปฺปยุตฺต\-ปจฺจยา. (3)}

\prose{อวิตกฺกวิจารมตฺตญฺจ อวิตกฺกอวิจารญฺจ ธมฺมํ ปฏิจฺจ สวิตกฺก\-สวิจาโร จ อวิตกฺก\-อวิจาโร จ ธมฺมา อุปฺปชฺชนฺติ วิปฺปยุตฺต\-ปจฺจยา. ปฏิสนฺธิ\-กฺขเณ วิตกฺกญฺจ วตฺถุญฺจ ปฏิจฺจ สวิตกฺก\-สวิจารา ขนฺธา. วิตกฺกญฺจ มหาภูเต จ ปฏิจฺจ กฏตฺตารูปํ. ขนฺธา วตฺถุํ วิปฺปยุตฺต\-ปจฺจยา. กฏตฺตารูปํ วิตกฺกํ วิปฺปยุตฺต\-ปจฺจยา. (4)}

\prose{อวิตกฺกวิจารมตฺตญฺจ อวิตกฺกอวิจารญฺจ ธมฺมํ ปฏิจฺจ อวิตกฺก\-วิจารมตฺโต จ อวิตกฺก\-อวิจาโร จ ธมฺมา ฯเปฯ. อวิตกฺก\-วิจารมตฺตํ เอกํ ขนฺธญฺจ วิจารญฺจ ปฏิจฺจ ตโย ขนฺธา จิตฺต\-สมุฏฺฐานญฺจ รูปํ ฯเปฯ ขนฺธา วตฺถุํ วิปฺปยุตฺต\-ปจฺจยา. จิตฺต\-สมุฏฺฐานํ รูปํ ขนฺเธ จ วิจารญฺจ วิปฺปยุตฺต\-ปจฺจยา. ปฏิสนฺธิ\-กฺขเณ อวิตกฺก\-วิจารมตฺตํ เอกํ ขนฺธญฺจ วิจารญฺจ ปฏิจฺจ ตโย ขนฺธา กฏตฺตา จ รูปํ ฯเปฯ ขนฺธา วตฺถุํ วิปฺปยุตฺต\-ปจฺจยา. กฏตฺตารูปํ ขนฺเธ จ วิจารญฺจ วิปฺปยุตฺต\-ปจฺจยา. ปฏิสนฺธิ\-กฺขเณ อวิตกฺก\-วิจารมตฺตํ เอกํ ขนฺธญฺจ วตฺถุญฺจ ปฏิจฺจ ตโย ขนฺธา ฯเปฯ เทฺว ขนฺเธ จ ฯเปฯ. อวิตกฺก\-วิจารมตฺเต ขนฺเธ จ มหาภูเต จ ปฏิจฺจ กฏตฺตารูปํ. ขนฺธา วตฺถุํ วิปฺปยุตฺต\-ปจฺจยา. กฏตฺตารูปํ ขนฺเธ วิปฺปยุตฺต\-ปจฺจยา. ปฏิสนฺธิ\-กฺขเณ อวิตกฺก\-วิจารมตฺตํ เอกํ ขนฺธญฺจ วตฺถุญฺจ ปฏิจฺจ ตโย ขนฺธา วิจาโร จ ฯเปฯ เทฺว ขนฺเธ จ ฯเปฯ วตฺถุํ วิปฺปยุตฺต\-ปจฺจยา. (5)}

\proseitem{19}{สวิตกฺกสวิจารญฺจ อวิตกฺกวิจารมตฺตญฺจ ธมฺมํ ปฏิจฺจ สวิตกฺก\-สวิจาโร ธมฺโม อุปฺปชฺชติ ฯเปฯ. สวิตกฺก\-สวิจารํ เอกํ ขนฺธญฺจ วิตกฺกญฺจ ปฏิจฺจ \csromanfolio{16}ตโย ขนฺธา ฯเปฯ เทฺว ขนฺเธ จ ฯเปฯ วตฺถุํ วิปฺปยุตฺต\-ปจฺจยา. ปฏิสนฺธิ\-กฺขเณ ฯเปฯ วตฺถุํ วิปฺปยุตฺต\-ปจฺจยา. (1)}

\prose{สวิตกฺกสวิจารญฺจ อวิตกฺกวิจารมตฺตญฺจ ธมฺมํ ปฏิจฺจ อวิตกฺก\-อวิจาโร ธมฺโม อุปฺปชฺชติ ฯเปฯ. สวิตกฺก\-สวิจาเร ขนฺเธ จ วิตกฺกญฺจ ปฏิจฺจ จิตฺต\-สมุฏฺฐานํ รูปํ. ขนฺเธ จ วิตกฺกญฺจ วิปฺปยุตฺต\-ปจฺจยา. ปฏิสนฺธิ\-กฺขเณ ฯเปฯ ขนฺเธ จ วิตกฺกญฺจ วิปฺปยุตฺต\-ปจฺจยา. (2)}

\prose{สวิตกฺกสวิจารญฺจ อวิตกฺกวิจารมตฺตญฺจ ธมฺมํ ปฏิจฺจ สวิตกฺก\-สวิจาโร จ อวิตกฺก\-อวิจาโร จ ธมฺมา อุปฺปชฺชนฺติ ฯเปฯ. สวิตกฺก\-สวิจารํ เอกํ ขนฺธญฺจ วิตกฺกญฺจ ปฏิจฺจ ตโย ขนฺธา จิตฺต\-สมุฏฺฐานญฺจ รูปํ ฯเปฯ เทฺว ขนฺเธ ฯเปฯ ขนฺธา วตฺถุํ วิปฺปยุตฺต\-ปจฺจยา. จิตฺต\-สมุฏฺฐานํ รูปํ ขนฺเธ จ วิตกฺกญฺจ วิปฺปยุตฺต\-ปจฺจยา. ปฏิสนฺธิ\-กฺขเณ ฯเปฯ ขนฺธา วตฺถุํ วิปฺปยุตฺต\-ปจฺจยา. กฏตฺตารูปํ ขนฺเธ จ วิตกฺกญฺจ วิปฺปยุตฺต\-ปจฺจยา. (3)}

\prose{สวิตกฺกสวิจารญฺจ อวิตกฺกวิจารมตฺตญฺจ อวิตกฺกอวิจารญฺจ ธมฺมํ ปฏิจฺจ สวิตกฺก\-สวิจาโร ธมฺโม อุปฺปชฺชติ ฯเปฯ. ปฏิสนฺธิ\-กฺขเณ สวิตกฺก\-สวิจารํ เอกํ ขนฺธญฺจ วิตกฺกญฺจ วตฺถุญฺจ ปฏิจฺจ ตโย ขนฺธา ฯเปฯ เทฺว ขนฺเธ ฯเปฯ ขนฺธา วตฺถุํ วิปฺปยุตฺต\-ปจฺจยา. (1)}

\prose{สวิตกฺกสวิจารญฺจ อวิตกฺกวิจารมตฺตญฺจ อวิตกฺกอวิจารญฺจ ธมฺมํ ปฏิจฺจ อวิตกฺก\-อวิจาโร ธมฺโม อุปฺปชฺชติ ฯเปฯ. สวิตกฺก\-สวิจาเร ขนฺเธ จ วิตกฺกญฺจ มหาภูเต จ ปฏิจฺจ จิตฺต\-สมุฏฺฐานํ รูปํ. ขนฺเธ จ วิตกฺกญฺจ วิปฺปยุตฺต\-ปจฺจยา. ปฏิสนฺธิ\-กฺขเณ ฯเปฯ ขนฺเธ จ วิตกฺกญฺจ วิปฺปยุตฺต\-ปจฺจยา. (2)}

\prose{สวิตกฺกสวิจารญฺจ อวิตกฺกวิจารมตฺตญฺจ อวิตกฺกอวิจารญฺจ ธมฺมํ ปฏิจฺจ สวิตกฺก\-สวิจาโร จ อวิตกฺก\-อวิจาโร จ ธมฺมา อุปฺปชฺชนฺติ วิปฺปยุตฺต\-ปจฺจยา. ปฏิสนฺธิ\-กฺขเณ สวิตกฺก\-สวิจารํ เอกํ ขนฺธญฺจ วิตกฺกญฺจ วตฺถุญฺจ ปฏิจฺจ ตโย ขนฺธา ฯเปฯ เทฺว ขนฺเธ จ ฯเปฯ. สวิตกฺก\-สวิจาเร ขนฺเธ จ วิตกฺกญฺจ มหาภูเต จ ปฏิจฺจ กฏตฺตารูปํ. ขนฺธา วตฺถุํ วิปฺปยุตฺต\-ปจฺจยา. กฏตฺตารูปํ ขนฺเธ จ วิตกฺกญฺจ วิปฺปยุตฺต\-ปจฺจยา. (3)}

\prose{อตฺถิอวิคต}

\proseitem{19}{สวิตกฺก\-สวิจารํ ธมฺมํ ปฏิจฺจ สวิตกฺก\-สวิจาโร ธมฺโม อุปฺปชฺชติ อตฺถิปจฺจยา ฯเปฯ นตฺถิปจฺจยา, วิคตปจฺจยา, อวิคตปจฺจยา. (สํขิตฺตํ.)}

\csromanmarkh{1. ปจฺจยานุโลม}\csromanmarkhii{2. สงฺขฺยาวาร}\csromanheadernomark{2}{1. \textbf{ปจฺจยานุโลม 2. สงฺขฺยาวาร}}
\proseitem{22}{\csromanfolio{17}เหตุยา สตฺตติํส, อารมฺมเณ เอกวีส, อธิปติยา เตวีส, อนนฺตเร เอกวีส, สมนนฺตเร เอกวีส, สหชาเต สตฺตติํส, อญฺญมญฺเญ อฏฺฐวีส, นิสฺสเย สตฺตติํส, อุปนิสฺสเย เอกวีส, ปุเรชาเต เอกาทส, อาเสวเน เอกาทส, กมฺเม สตฺตติํส, วิปาเก สตฺตติํส, อาหาเร อินฺทฺริเย ฌาเน มคฺเค สตฺตติํส, สมฺปยุตฺเต เอกวีส, วิปฺปยุตฺเต สตฺตติํส, อตฺถิยา สตฺตติํส, นตฺถิยา เอกวีส, วิคเต เอกวีส, อวิคเต สตฺตติํส.}

\prosewithcloserruled{เหตุปจฺจยา อารมฺมเณ เอกวีส ฯเปฯ อวิคเต สตฺตติํส. (สํขิตฺตํ.) (ยถา กุสลตฺติเก คณนา, เอวํ คเณตพฺพํ.)}{อนุโลมํ.}
\csromanmatikaanchor{mtk.2}
\csromanmatikaanchor{mtk.14}
\csromantocmarkref{section}{สุทฺธ}{mtk.2}
\bookmark[dest={mtk.2},level=1]{สุทฺธ}
\csromanmarkh{2. ปจฺจย\-ปจฺจนีย}\csromanmarkhii{1. วิภงฺควาร}\csromanheadernomark{4}{2. ปจฺจย\-ปจฺจนีย 1. วิภงฺควาร}
\csromanheader{2}{\textbf{นเหตุ}}
\proseitem{23}{สวิตกฺก\-สวิจารํ ธมฺมํ ปฏิจฺจ สวิตกฺก\-สวิจาโร ธมฺโม อุปฺปชฺชติ นเหตุปจฺจยา. อเหตุกํ สวิตกฺก\-สวิจารํ เอกํ ขนฺธํ ปฏิจฺจ ตโย ขนฺธา ฯเปฯ เทฺว ขนฺเธ ปฏิจฺจ เทฺว ขนฺธา. อเหตุก\-ปฏิสนฺธิ\-กฺขเณ ฯเปฯ. วิจิกิจฺฉา\-สหคเต อุทฺธจฺจ\-สหคเต ขนฺเธ ปฏิจฺจ วิจิกิจฺฉา\-สหคโต อุทฺธจฺจ\-สหคโต โมโห. (1)}

\prose{สวิตกฺก\-สวิจารํ ธมฺมํ ปฏิจฺจ อวิตกฺก\-วิจารมตฺโต ธมฺโม อุปฺปชฺชติ นเหตุปจฺจยา. อเหตุเก สวิตกฺก\-สวิจาเร ขนฺเธ ปฏิจฺจ วิตกฺโก. อเหตุก\-ปฏิสนฺธิ\-กฺขเณ ฯเปฯ. (2)}

\prose{สวิตกฺก\-สวิจารํ ธมฺมํ ปฏิจฺจ อวิตกฺก\-อวิจาโร ธมฺโม อุปฺปชฺชติ นเหตุปจฺจยา. อเหตุเก สวิตกฺก\-สวิจาเร ขนฺเธ ปฏิจฺจ จิตฺต\-สมุฏฺฐานํ รูปํ. อเหตุก\-ปฏิสนฺธิ\-กฺขเณ ฯเปฯ. (3)}

\prose{\csromanfolio{18}สวิตกฺก\-สวิจารํ ธมฺมํ ปฏิจฺจ สวิตกฺก\-สวิจาโร จ อวิตกฺก\-อวิจาโร จ ธมฺมา อุปฺปชฺชนฺติ นเหตุปจฺจยา. อเหตุกํ สวิตกฺก\-สวิจารํ เอกํ ขนฺธํ ปฏิจฺจ ตโย ขนฺธา จิตฺต\-สมุฏฺฐานญฺจ รูปํ ฯเปฯ เทฺว ขนฺเธ ปฏิจฺจ เทฺว ขนฺธา จิตฺต\-สมุฏฺฐานญฺจ รูปํ. อเหตุก\-ปฏิสนฺธิ\-กฺขเณ ฯเปฯ. (4)}

\prose{สวิตกฺก\-สวิจารํ ธมฺมํ ปฏิจฺจ อวิตกฺก\-วิจารมตฺโต จ อวิตกฺก\-อวิจาโร จ ธมฺมา อุปฺปชฺชนฺติ นเหตุปจฺจยา. อเหตุเก สวิตกฺก\-สวิจาเร ขนฺเธ ปฏิจฺจ วิตกฺโก จ จิตฺต\-สมุฏฺฐานญฺจ รูปํ. อเหตุก\-ปฏิสนฺธิ\-กฺขเณ ฯเปฯ. (5)}

\prose{สวิตกฺก\-สวิจารํ ธมฺมํ ปฏิจฺจ สวิตกฺก\-สวิจาโร จ อวิตกฺก\-วิจารมตฺโต จ ธมฺมา อุปฺปชฺชนฺติ นเหตุปจฺจยา. อเหตุกํ สวิตกฺก\-สวิจารํ เอกํ ขนฺธํ ปฏิจฺจ ตโย ขนฺธา วิตกฺโก จ ฯเปฯ เทฺว ขนฺเธ ปฏิจฺจ เทฺว ขนฺธา วิตกฺโก จ. อเหตุก\-ปฏิสนฺธิ\-กฺขเณ ฯเปฯ. (6)}

\prose{สวิตกฺก\-สวิจารํ ธมฺมํ ปฏิจฺจ สวิตกฺก\-สวิจาโร จ อวิตกฺก\-วิจารมตฺโต จ อวิตกฺก\-อวิจาโร จ ธมฺมา อุปฺปชฺชนฺติ นเหตุปจฺจยา. อเหตุกํ สวิตกฺก\-สวิจารํ เอกํ ขนฺธํ ปฏิจฺจ ตโย ขนฺธา วิตกฺโก จ จิตฺต\-สมุฏฺฐานญฺจ รูปํ ฯเปฯ. อเหตุก\-ปฏิสนฺธิ\-กฺขเณ ฯเปฯ. (7)}

\proseitem{24}{อวิตกฺก\-วิจารมตฺตํ ธมฺมํ ปฏิจฺจ สวิตกฺก\-สวิจาโร ธมฺโม อุปฺปชฺชติ นเหตุปจฺจยา. อเหตุกํ วิตกฺกํ ปฏิจฺจ สวิตกฺก\-สวิจารา ขนฺธา. อเหตุก\-ปฏิสนฺธิ\-กฺขเณ วิตกฺกํ ปฏิจฺจ สวิตกฺก\-สวิจารา ขนฺธา. วิจิกิจฺฉา\-สหคตํ อุทฺธจฺจ\-สหคตํ วิตกฺกํ ปฏิจฺจ วิจิกิจฺฉา\-สหคโต อุทฺธจฺจ\-สหคโต โมโห. (1)}

\prose{อวิตกฺก\-วิจารมตฺตํ ธมฺมํ ปฏิจฺจ อวิตกฺก\-อวิจาโร ธมฺโม อุปฺปชฺชติ นเหตุปจฺจยา. อเหตุกํ วิตกฺกํ ปฏิจฺจ จิตฺต\-สมุฏฺฐานํ รูปํ. อเหตุก\-ปฏิสนฺธิ\-กฺขเณ วิตกฺกํ ปฏิจฺจ กฏตฺตารูปํ. (2)}

\prose{อวิตกฺก\-วิจารมตฺตํ ธมฺมํ ปฏิจฺจ สวิตกฺก\-สวิจาโร จ อวิตกฺก\-อวิจาโร จ ธมฺมา อุปฺปชฺชนฺติ นเหตุปจฺจยา. อเหตุกํ วิตกฺกํ ปฏิจฺจ สวิตกฺก\-สวิจารา ขนฺธา จิตฺต\-สมุฏฺฐานญฺจ รูปํ. อเหตุก\-ปฏิสนฺธิ\-กฺขเณ วิตกฺกํ ปฏิจฺจ สวิตกฺก\-สวิจารา ขนฺธา กฏตฺตา จ รูปํ. (3)}

\prose{อวิตกฺก\-อวิจารํ ธมฺมํ ปฏิจฺจ อวิตกฺก\-อวิจาโร ธมฺโม อุปฺปชฺชติ นเหตุปจฺจยา. อเหตุกํ อวิตกฺก\-อวิจารํ เอกํ ขนฺธํ ปฏิจฺจ ตโย ขนฺธา, \csromanfolio{19}ตโย ขนฺเธ ปฏิจฺจ เอโก ขนฺโธ, เทฺว ขนฺเธ ปฏิจฺจ เทฺว ขนฺธา. เอกํ มหาภูตํ ปฏิจฺจ ตโย มหาภูตา ฯเปฯ มหาภูเต ปฏิจฺจ จิตฺต\-สมุฏฺฐานํ รูปํ กฏตฺตารูปํ อุปาทารูปํ. พาหิรํ. อาหารสมุฏฺฐานํ. อุตุสมุฏฺฐานํ. อสญฺญสตฺตานํ เอกํ มหาภูตํ ปฏิจฺจ ฯเปฯ มหาภูเต ปฏิจฺจ กฏตฺตารูปํ อุปาทารูปํ. (1)}

\prose{อวิตกฺก\-อวิจารํ ธมฺมํ ปฏิจฺจ สวิตกฺก\-สวิจาโร ธมฺโม อุปฺปชฺชติ นเหตุปจฺจยา. อเหตุก\-ปฏิสนฺธิ\-กฺขเณ วตฺถุํ ปฏิจฺจ สวิตกฺก\-สวิจารา ขนฺธา. (2)}

\prose{อวิตกฺก\-อวิจารํ ธมฺมํ ปฏิจฺจ อวิตกฺก\-วิจารมตฺโต ธมฺโม อุปฺปชฺชติ นเหตุปจฺจยา. อเหตุก\-ปฏิสนฺธิ\-กฺขเณ วตฺถุํ ปฏิจฺจ วิตกฺโก. (3)}

\prose{อวิตกฺก\-อวิจารํ ธมฺมํ ปฏิจฺจ สวิตกฺก\-สวิจาโร จ อวิตกฺก\-อวิจาโร จ ธมฺมา อุปฺปชฺชนฺติ นเหตุปจฺจยา. อเหตุก\-ปฏิสนฺธิ\-กฺขเณ วตฺถุํ ปฏิจฺจ สวิตกฺก\-สวิจารา ขนฺธา. มหาภูเต ปฏิจฺจ กฏตฺตารูปํ. (4)}

\prose{อวิตกฺก\-อวิจารํ ธมฺมํ ปฏิจฺจ อวิตกฺก\-วิจารมตฺโต จ อวิตกฺก\-อวิจาโร จ ธมฺมา อุปฺปชฺชนฺติ นเหตุปจฺจยา. อเหตุก\-ปฏิสนฺธิ\-กฺขเณ วตฺถุํ ปฏิจฺจ วิตกฺโก. มหาภูเต ปฏิจฺจ กฏตฺตารูปํ. (5)}

\prose{อวิตกฺก\-อวิจารํ ธมฺมํ ปฏิจฺจ สวิตกฺก\-สวิจาโร จ อวิตกฺก\-วิจารมตฺโต จ ธมฺมา อุปฺปชฺชนฺติ นเหตุปจฺจยา. อเหตุก\-ปฏิสนฺธิ\-กฺขเณ วตฺถุํ ปฏิจฺจ สวิตกฺก\-สวิจารา ขนฺธา จ วิตกฺโก จ. (6)}

\prose{อวิตกฺก\-อวิจารํ ธมฺมํ ปฏิจฺจ สวิตกฺก\-สวิจาโร จ อวิตกฺก\-วิจารมตฺโต จ อวิตกฺก\-อวิจาโร จ ธมฺมา อุปฺปชฺชนฺติ นเหตุปจฺจยา. อเหตุก\-ปฏิสนฺธิ\-กฺขเณ วตฺถุํ ปฏิจฺจ สวิตกฺก\-สวิจารา ขนฺธา จ วิตกฺโก จ. มหาภูเต ปฏิจฺจ กฏตฺตารูปํ. (7)}

\proseitem{25}{สวิตกฺกสวิจารญฺจ อวิตกฺกอวิจารญฺจ ธมฺมํ ปฏิจฺจ สวิตกฺก\-สวิจาโร ธมฺโม อุปฺปชฺชติ นเหตุปจฺจยา. อเหตุก\-ปฏิสนฺธิ\-กฺขเณ สวิตกฺก\-สวิจารํ เอกํ ขนฺธญฺจ วตฺถุญฺจ ปฏิจฺจ ตโย ขนฺธา ฯเปฯ เทฺว ขนฺเธ จ วตฺถุญฺจ ปฏิจฺจ เทฺว ขนฺธา. (1)}

\prose{\csromanfolio{20}สวิตกฺกสวิจารญฺจ อวิตกฺกอวิจารญฺจ ธมฺมํ ปฏิจฺจ อวิตกฺก\-วิจารมตฺโต ธมฺโม อุปฺปชฺชติ นเหตุปจฺจยา. อเหตุก\-ปฏิสนฺธิ\-กฺขเณ สวิตกฺก\-สวิจาเร ขนฺเธ จ วตฺถุญฺจ ปฏิจฺจ วิตกฺโก. (2)}

\prose{สวิตกฺกสวิจารญฺจ อวิตกฺกอวิจารญฺจ ธมฺมํ ปฏิจฺจ อวิตกฺก\-อวิจาโร ธมฺโม อุปฺปชฺชติ นเหตุปจฺจยา. อเหตุเก สวิตกฺก\-สวิจาเร ขนฺเธ จ มหาภูเต จ ปฏิจฺจ จิตฺต\-สมุฏฺฐานํ รูปํ. อเหตุก\-ปฏิสนฺธิ\-กฺขเณ สวิตกฺก\-สวิจาเร ขนฺเธ จ มหาภูเต จ ปฏิจฺจ กฏตฺตารูปํ. (3)}

\prose{สวิตกฺกสวิจารญฺจ อวิตกฺกอวิจารญฺจ ธมฺมํ ปฏิจฺจ สวิตกฺก\-สวิจาโร จ อวิตกฺก\-อวิจาโร จ ธมฺมา อุปฺปชฺชนฺติ นเหตุปจฺจยา. อเหตุก\-ปฏิสนฺธิ\-กฺขเณ สวิตกฺก\-สวิจารํ เอกํ ขนฺธญฺจ วตฺถุญฺจ ปฏิจฺจ ตโย ขนฺธา ฯเปฯ เทฺว ขนฺเธ จ วตฺถุญฺจ ปฏิจฺจ เทฺว ขนฺธา. สวิตกฺก\-สวิจาเร ขนฺเธ จ มหาภูเต จ ปฏิจฺจ กฏตฺตารูปํ. (4)}

\prose{สวิตกฺกสวิจารญฺจ อวิตกฺกอวิจารญฺจ ธมฺมํ ปฏิจฺจ อวิตกฺก\-วิจารมตฺโต จ อวิตกฺก\-อวิจาโร จ ธมฺมา อุปฺปชฺชนฺติ นเหตุปจฺจยา. อเหตุก\-ปฏิสนฺธิ\-กฺขเณ สวิตกฺก\-สวิจาเร ขนฺเธ จ วตฺถุญฺจ ปฏิจฺจ วิตกฺโก. สวิตกฺก\-สวิจาเร ขนฺเธ จ มหาภูเต จ ปฏิจฺจ กฏตฺตารูปํ. (5)}

\prose{สวิตกฺกสวิจารญฺจ อวิตกฺกอวิจารญฺจ ธมฺมํ ปฏิจฺจ สวิตกฺก\-สวิจาโร จ อวิตกฺก\-วิจารมตฺโต จ ธมฺมา อุปฺปชฺชนฺติ นเหตุปจฺจยา. อเหตุก\-ปฏิสนฺธิ\-กฺขเณ สวิตกฺก\-สวิจารํ เอกํ ขนฺธญฺจ วตฺถุญฺจ ปฏิจฺจ ตโย ขนฺธา วิตกฺโก จ ฯเปฯ เทฺว ขนฺเธ จ วตฺถุญฺจ ปฏิจฺจ เทฺว ขนฺธา วิตกฺโก จ. (6)}

\prose{สวิตกฺกสวิจารญฺจ อวิตกฺกอวิจารญฺจ ธมฺมํ ปฏิจฺจ สวิตกฺก\-สวิจาโร จ อวิตกฺก\-วิจารมตฺโต จ อวิตกฺก\-อวิจาโร จ ธมฺมา อุปฺปชฺชนฺติ นเหตุปจฺจยา. อเหตุก\-ปฏิสนฺธิ\-กฺขเณ สวิตกฺก\-สวิจารํ เอกํ ขนฺธญฺจ วตฺถุญฺจ ปฏิจฺจ ตโย ขนฺธา วิตกฺโก จ ฯเปฯ เทฺว ขนฺเธ จ วตฺถุญฺจ ปฏิจฺจ เทฺว ขนฺธา วิตกฺโก จ. สวิตกฺก\-สวิจาเร ขนฺเธ จ มหาภูเต จ ปฏิจฺจ กฏตฺตารูปํ. (7)}

\proseitem{26}{อวิตกฺกวิจารมตฺตญฺจ อวิตกฺกอวิจารญฺจ ธมฺมํ ปฏิจฺจ สวิตกฺก\-สวิจาโร ธมฺโม อุปฺปชฺชติ นเหตุปจฺจยา. อเหตุก\-ปฏิสนฺธิ\-กฺขเณ วิตกฺกญฺจ วตฺถุญฺจ ปฏิจฺจ สวิตกฺก\-สวิจารา ขนฺธา. (1)}

\prose{\csromanfolio{21}อวิตกฺกวิจารมตฺตญฺจ อวิตกฺกอวิจารญฺจ ธมฺมํ ปฏิจฺจ อวิตกฺก\-อวิจาโร ธมฺโม อุปฺปชฺชติ นเหตุปจฺจยา. อเหตุกํ วิตกฺกญฺจ มหาภูเต จ ปฏิจฺจ จิตฺต\-สมุฏฺฐานํ รูปํ. อเหตุก\-ปฏิสนฺธิ\-กฺขเณ วิตกฺกญฺจ มหาภูเต จ ปฏิจฺจ กฏตฺตารูปํ. (2)}

\prose{อวิตกฺกวิจารมตฺตญฺจ อวิตกฺกอวิจารญฺจ ธมฺมํ ปฏิจฺจ สวิตกฺก\-สวิจาโร จ อวิตกฺก\-อวิจาโร จ ธมฺมา อุปฺปชฺชนฺติ นเหตุปจฺจยา. อเหตุก\-ปฏิสนฺธิ\-กฺขเณ วิตกฺกญฺจ วตฺถุญฺจ ปฏิจฺจ สวิตกฺก\-สวิจารา ขนฺธา. วิตกฺกญฺจ มหาภูเต จ ปฏิจฺจ กฏตฺตารูปํ. (3)}

\proseitem{27}{สวิตกฺกสวิจารญฺจ อวิตกฺกวิจารมตฺตญฺจ ธมฺมํ ปฏิจฺจ สวิตกฺก\-สวิจาโร ธมฺโม อุปฺปชฺชติ นเหตุปจฺจยา. อเหตุกํ สวิตกฺก\-สวิจารํ เอกํ ขนฺธญฺจ วิตกฺกญฺจ ปฏิจฺจ ตโย ขนฺธา ฯเปฯ เทฺว ขนฺเธ จ วิตกฺกญฺจ ปฏิจฺจ เทฺว ขนฺธา. อเหตุก\-ปฏิสนฺธิ\-กฺขเณ สวิตกฺก\-สวิจารํ เอกํ ขนฺธญฺจ วิตกฺกญฺจ ปฏิจฺจ ตโย ขนฺธา ฯเปฯ เทฺว ขนฺเธ จ วิตกฺกญฺจ ปฏิจฺจ เทฺว ขนฺธา. วิจิกิจฺฉา\-สหคเต อุทฺธจฺจ\-สหคเต ขนฺเธ จ วิตกฺกญฺจ ปฏิจฺจ วิจิกิจฺฉา\-สหคโต อุทฺธจฺจ\-สหคโต โมโห. (1)}

\prose{สวิตกฺกสวิจารญฺจ อวิตกฺกวิจารมตฺตญฺจ ธมฺมํ ปฏิจฺจ อวิตกฺก\-อวิจาโร ธมฺโม อุปฺปชฺชติ นเหตุปจฺจยา. อเหตุเก สวิตกฺก\-สวิจาเร ขนฺเธ จ วิตกฺกญฺจ ปฏิจฺจ จิตฺต\-สมุฏฺฐานํ รูปํ. อเหตุก\-ปฏิสนฺธิ\-กฺขเณ ฯเปฯ. (2)}

\prose{สวิตกฺกสวิจารญฺจ อวิตกฺกวิจารมตฺตญฺจ ธมฺมํ ปฏิจฺจ สวิตกฺก\-สวิจาโร จ อวิตกฺก\-อวิจาโร จ ธมฺมา อุปฺปชฺชนฺติ นเหตุปจฺจยา. อเหตุกํ สวิตกฺก\-สวิจารํ เอกํ ขนฺธญฺจ วิตกฺกญฺจ ปฏิจฺจ ตโย ขนฺธา จิตฺต\-สมุฏฺฐานญฺจ รูปํ ฯเปฯ เทฺว ขนฺเธ จ วิตกฺกญฺจ ปฏิจฺจ เทฺว ขนฺธา จิตฺต\-สมุฏฺฐานญฺจ รูปํ. อเหตุก\-ปฏิสนฺธิ\-กฺขเณ ฯเปฯ. (3)}

\proseitem{28}{สวิตกฺกสวิจารญฺจ อวิตกฺกวิจารมตฺตญฺจ อวิตกฺกอวิจารญฺจ ธมฺมํ ปฏิจฺจ สวิตกฺก\-สวิจาโร ธมฺโม อุปฺปชฺชติ นเหตุปจฺจยา. อเหตุก\-ปฏิสนฺธิ\-กฺขเณ สวิตกฺก\-สวิจารํ เอกํ ขนฺธญฺจ วิตกฺกญฺจ วตฺถุญฺจ ปฏิจฺจ ตโย ขนฺธา ฯเปฯ เทฺว ขนฺเธ จ วิตกฺกญฺจ วตฺถุญฺจ ปฏิจฺจ เทฺว ขนฺธา. (1)}

\prose{สวิตกฺกสวิจารญฺจ อวิตกฺกวิจารมตฺตญฺจ อวิตกฺกอวิจารญฺจ ธมฺมํ ปฏิจฺจ อวิตกฺก\-อวิจาโร ธมฺโม อุปฺปชฺชติ นเหตุปจฺจยา. อเหตุเก \csromanfolio{22}สวิตกฺก\-สวิจาเร ขนฺเธ จ วิตกฺกญฺจ มหาภูเต จ ปฏิจฺจ จิตฺต\-สมุฏฺฐานํ รูปํ. อเหตุก\-ปฏิสนฺธิ\-กฺขเณ ฯเปฯ. (2)}

\prose{สวิตกฺกสวิจารญฺจ อวิตกฺกวิจารมตฺตญฺจ อวิตกฺกอวิจารญฺจ ธมฺมํ ปฏิจฺจ สวิตกฺก\-สวิจาโร จ อวิตกฺก\-อวิจาโร จ ธมฺมา อุปฺปชฺชนฺติ นเหตุปจฺจยา. อเหตุก\-ปฏิสนฺธิ\-กฺขเณ สวิตกฺก\-สวิจารํ เอกํ ขนฺธญฺจ วิตกฺกญฺจ วตฺถุญฺจ ปฏิจฺจ ตโย ขนฺธา ฯเปฯ เทฺว ขนฺเธ จ วิตกฺกญฺจ วตฺถุญฺจ ปฏิจฺจ เทฺว ขนฺธา. สวิตกฺก\-สวิจาเร ขนฺเธ จ วิตกฺกญฺจ มหาภูเต จ ปฏิจฺจ กฏตฺตารูปํ. (3)}

\prose{นอารมฺมณ}

\proseitem{29}{สวิตกฺก\-สวิจารํ ธมฺมํ ปฏิจฺจ อวิตกฺก\-อวิจาโร ธมฺโม อุปฺปชฺชติ นอารมฺมณ\-ปจฺจยา. สวิตกฺก\-สวิจาเร ขนฺเธ ปฏิจฺจ จิตฺต\-สมุฏฺฐานํ รูปํ. ปฏิสนฺธิ\-กฺขเณ สวิตกฺก\-สวิจาเร ขนฺเธ ปฏิจฺจ กฏตฺตารูปํ. (1)}

\prose{อวิตกฺก\-วิจารมตฺตํ ธมฺมํ ปฏิจฺจ อวิตกฺก\-อวิจาโร ธมฺโม อุปฺปชฺชติ นอารมฺมณ\-ปจฺจยา. อวิตกฺก\-วิจารมตฺเต ขนฺเธ ปฏิจฺจ จิตฺต\-สมุฏฺฐานํ รูปํ. วิตกฺกํ ปฏิจฺจ จิตฺต\-สมุฏฺฐานํ รูปํ. ปฏิสนฺธิ\-กฺขเณ อวิตกฺก\-วิจารมตฺเต ขนฺเธ ปฏิจฺจ กฏตฺตารูปํ. ปฏิสนฺธิ\-กฺขเณ วิตกฺกํ ปฏิจฺจ กฏตฺตารูปํ. (1)}

\prose{อวิตกฺก\-อวิจารํ ธมฺมํ ปฏิจฺจ อวิตกฺก\-อวิจาโร ธมฺโม อุปฺปชฺชติ นอารมฺมณ\-ปจฺจยา. อวิตกฺก\-อวิจาเร ขนฺเธ ปฏิจฺจ จิตฺต\-สมุฏฺฐานํ รูปํ. วิจารํ ปฏิจฺจ จิตฺต\-สมุฏฺฐานํ รูปํ. ปฏิสนฺธิ\-กฺขเณ อวิตกฺก\-อวิจาเร ขนฺเธ ปฏิจฺจ กฏตฺตารูปํ. วิจารํ ปฏิจฺจ กฏตฺตารูปํ. ขนฺเธ ปฏิจฺจ วตฺถุ. วิจารํ ปฏิจฺจ วตฺถุ. เอกํ มหาภูตํ ปฏิจฺจ ฯเปฯ. พาหิรํ. อาหารสมุฏฺฐานํ. อุตุสมุฏฺฐานํ. อสญฺญสตฺตานํ เอกํ มหาภูตํ ปฏิจฺจ ฯเปฯ. (1)}

\proseitem{30}{สวิตกฺกสวิจารญฺจ อวิตกฺกอวิจารญฺจ ธมฺมํ ปฏิจฺจ อวิตกฺก\-อวิจาโร ธมฺโม อุปฺปชฺชติ นอารมฺมณ\-ปจฺจยา. สวิตกฺก\-สวิจาเร ขนฺเธ จ มหาภูเต จ ปฏิจฺจ จิตฺต\-สมุฏฺฐานํ รูปํ. ปฏิสนฺธิ\-กฺขเณ สวิตกฺก\-สวิจาเร ขนฺเธ จ มหาภูเต จ ปฏิจฺจ กฏตฺตารูปํ. (1)}

\prose{อวิตกฺกวิจารมตฺตญฺจ อวิตกฺกอวิจารญฺจ ธมฺมํ ปฏิจฺจ อวิตกฺก\-อวิจาโร ธมฺโม อุปฺปชฺชติ นอารมฺมณ\-ปจฺจยา. อวิตกฺก\-วิจารมตฺเต ขนฺเธ จ วิจารญฺจ \csromanfolio{23}ปฏิจฺจ จิตฺต\-สมุฏฺฐานํ รูปํ. อวิตกฺก\-วิจารมตฺเต ขนฺเธ จ มหาภูเต จ ปฏิจฺจ จิตฺต\-สมุฏฺฐานํ รูปํ. วิตกฺกญฺจ มหาภูเต จ ปฏิจฺจ จิตฺต\-สมุฏฺฐานํ รูปํ. ปฏิสนฺธิ\-กฺขเณ อวิตกฺก\-วิจารมตฺเต ขนฺเธ จ วิจารญฺจ ฯเปฯ กฏตฺตารูปํ. (1)}

\prose{สวิตกฺกสวิจารญฺจ อวิตกฺกวิจารมตฺตญฺจ ธมฺมํ ปฏิจฺจ อวิตกฺก\-อวิจาโร ธมฺโม อุปฺปชฺชติ นอารมฺมณ\-ปจฺจยา. สวิตกฺก\-สวิจาเร ขนฺเธ จ วิตกฺกญฺจ ปฏิจฺจ จิตฺต\-สมุฏฺฐานํ รูปํ. ปฏิสนฺธิ\-กฺขเณ ฯเปฯ. (1)}

\prose{สวิตกฺกสวิจารญฺจ อวิตกฺกวิจารมตฺตญฺจ อวิตกฺกอวิจารญฺจ ธมฺมํ ปฏิจฺจ อวิตกฺก\-อวิจาโร ธมฺโม อุปฺปชฺชติ นอารมฺมณ\-ปจฺจยา. สวิตกฺก\-สวิจาเร ขนฺเธ จ วิตกฺกญฺจ มหาภูเต จ ปฏิจฺจ จิตฺต\-สมุฏฺฐานํ รูปํ. ปฏิสนฺธิ\-กฺขเณ ฯเปฯ กฏตฺตารูปํ. (1)}

\prose{นอธิปติ}

\proseitem{31}{สวิตกฺก\-สวิจารํ ธมฺมํ ปฏิจฺจ สวิตกฺก\-สวิจาโร ธมฺโม อุปฺปชฺชติ นอธิปติ\-ปจฺจยา ฯเปฯ. สตฺต.}

\prose{อวิตกฺก\-วิจารมตฺตํ ธมฺมํ ปฏิจฺจ อวิตกฺก\-วิจารมตฺโต ธมฺโม อุปฺปชฺชติ นอธิปติ\-ปจฺจยา. อวิตกฺก\-วิจารมตฺเต ขนฺเธ ปฏิจฺจ อวิตกฺก\-วิจารมตฺตา อธิปติ. วิปากํ อวิตกฺก\-วิจารมตฺตํ เอกํ ขนฺธํ ปฏิจฺจ ตโย ขนฺธา ฯเปฯ. ปฏิสนฺธิ\-กฺขเณ ฯเปฯ. (1)}

\prose{อวิตกฺก\-วิจารมตฺตํ ธมฺมํ ปฏิจฺจ สวิตกฺก\-สวิจาโร ธมฺโม อุปฺปชฺชติ นอธิปติ\-ปจฺจยา. วิตกฺกํ ปฏิจฺจ สวิตกฺก\-สวิจารา ขนฺธา. ปฏิสนฺธิ\-กฺขเณ ฯเปฯ. (2)}

\prose{อวิตกฺก\-วิจารมตฺตํ ธมฺมํ ปฏิจฺจ อวิตกฺก\-อวิจาโร ธมฺโม อุปฺปชฺชติ นอธิปติ\-ปจฺจยา. วิปาเก อวิตกฺก\-วิจารมตฺเต ขนฺเธ ปฏิจฺจ วิจาโร จ จิตฺต\-สมุฏฺฐานญฺจ รูปํ. ปฏิสนฺธิ\-กฺขเณ ฯเปฯ. (3)}

\prose{อวิตกฺก\-วิจารมตฺตํ ธมฺมํ ปฏิจฺจ สวิตกฺก\-สวิจาโร จ อวิตกฺก\-อวิจาโร จ ธมฺมา อุปฺปชฺชนฺติ นอธิปติ\-ปจฺจยา. วิตกฺกํ ปฏิจฺจ สวิตกฺก\-สวิจารา ขนฺธา จิตฺต\-สมุฏฺฐานญฺจ รูปํ. ปฏิสนฺธิ\-กฺขเณ ฯเปฯ. (4)}

\prose{อวิตกฺก\-วิจารมตฺตํ ธมฺมํ ปฏิจฺจ อวิตกฺก\-วิจารมตฺโต จ อวิตกฺก\-อวิจาโร จ ธมฺมา อุปฺปชฺชนฺติ นอธิปติ\-ปจฺจยา. วิปากํ อวิตกฺก\-วิจารมตฺตํ เอกํ ขนฺธํ ปฏิจฺจ ตโย ขนฺธา วิจาโร จ จิตฺต\-สมุฏฺฐานญฺจ รูปํ ฯเปฯ. ปฏิสนฺธิ\-กฺขเณ ฯเปฯ. (5)}

\proseitem{32}{\csromanfolio{24}อวิตกฺก\-อวิจารํ ธมฺมํ ปฏิจฺจ อวิตกฺก\-อวิจาโร ธมฺโม อุปฺปชฺชติ นอธิปติ\-ปจฺจยา. อวิตกฺก\-อวิจาเร ขนฺเธ ปฏิจฺจ อวิตกฺก\-อวิจารา อธิปติ. วิปากํ อวิตกฺก\-อวิจารํ เอกํ ขนฺธํ ปฏิจฺจ ตโย ขนฺธา ฯเปฯ. ปฏิสนฺธิ\-กฺขเณ ฯเปฯ. (1)}

\prose{อวิตกฺก\-อวิจารํ ธมฺมํ ปฏิจฺจ สวิตกฺก\-สวิจาโร ธมฺโม อุปฺปชฺชติ นอธิปติ\-ปจฺจยา. ปฏิสนฺธิ\-กฺขเณ วตฺถุํ ปฏิจฺจ สวิตกฺก\-สวิจารา ขนฺธา. (2)}

\prose{อวิตกฺก\-อวิจารํ ธมฺมํ ปฏิจฺจ อวิตกฺก\-วิจารมตฺโต ธมฺโม อุปฺปชฺชติ นอธิปติ\-ปจฺจยา. วิจารํ ปฏิจฺจ อวิตกฺก\-วิจารมตฺตา อธิปติ. วิปากํ วิจารํ ปฏิจฺจ อวิตกฺก\-วิจารมตฺตา ขนฺธา. ปฏิสนฺธิ\-กฺขเณ ฯเปฯ. (3)}

\prose{อวิตกฺก\-อวิจารํ ธมฺมํ ปฏิจฺจ. สตฺต.}

\proseitem{33}{สวิตกฺกสวิจารญฺจ อวิตกฺกอวิจารญฺจ ธมฺมํ ปฏิจฺจ. สตฺต. (สํขิตฺตํ.)}

\prose{อวิตกฺกวิจารมตฺตญฺจ อวิตกฺกอวิจารญฺจ ธมฺมํ ปฏิจฺจ สวิตกฺก\-สวิจาโร ธมฺโม อุปฺปชฺชติ ฯเปฯ. อวิตกฺก\-วิจารมตฺโต ธมฺโม อุปฺปชฺชติ นอธิปติ\-ปจฺจยา. อวิตกฺก\-วิจารมตฺเต ขนฺเธ จ วิจารญฺจ ปฏิจฺจ อวิตกฺก\-วิจารมตฺตา อธิปติ. วิปากํ อวิตกฺก\-วิจารมตฺตํ เอกํ ขนฺธญฺจ วิจารญฺจ ปฏิจฺจ. (สํขิตฺตํ.)}

\prose{นอนนฺตราทิ}

\proseitem{34}{สวิตกฺก\-สวิจารํ ธมฺมํ ปฏิจฺจ อวิตกฺก\-อวิจาโร ธมฺโม อุปฺปชฺชติ นอนนฺตร\-ปจฺจยา. นสมนนฺตร\-ปจฺจยา. นอญฺญมญฺญ\-ปจฺจยา. นอุปนิสฺสยปจฺจยา. (นอารมฺมณ\-สทิสํ.)}

\csromanheader{2}{\textbf{นปุเรชาต}}
\proseitem{35}{สวิตกฺก\-สวิจารํ ธมฺมํ ปฏิจฺจ สวิตกฺก\-สวิจาโร ธมฺโม อุปฺปชฺชติ นปุเรชาต\-ปจฺจยา. สตฺต.}

\prose{อวิตกฺก\-วิจารมตฺตํ ธมฺมํ ปฏิจฺจ อวิตกฺก\-วิจารมตฺโต ธมฺโม อุปชฺชติ นปุเรชาต\-ปจฺจยา. อรูเป อวิตกฺก\-วิจารมตฺตํ เอกํ ขนฺธํ ปฏิจฺจ ฯเปฯ. ปฏิสนฺธิ\-กฺขเณ ฯเปฯ. (1)}

\prose{\csromanfolio{25}อวิตกฺก\-วิจารมตฺตํ ธมฺมํ ปฏิจฺจ สวิตกฺก\-สวิจาโร ธมฺโม อุปฺปชฺชติ นปุเรชาต\-ปจฺจยา. อรูเป วิตกฺกํ ปฏิจฺจ สวิตกฺก\-สวิจารา ขนฺธา. ปฏิสนฺธิ\-กฺขเณ ฯเปฯ. (2)}

\prose{อวิตกฺก\-วิจารมตฺตํ ธมฺมํ ปฏิจฺจ อวิตกฺก\-อวิจาโร ธมฺโม อุปฺปชฺชติ นปุเรชาต\-ปจฺจยา. อรูเป อวิตกฺก\-วิจารมตฺเต ขนฺเธ ปฏิจฺจ วิจาโร. อวิตกฺก\-วิจารมตฺเต ขนฺเธ ปฏิจฺจ จิตฺต\-สมุฏฺฐานํ รูปํ. วิตกฺกํ ปฏิจฺจ จิตฺต\-สมุฏฺฐานํ รูปํ. ปฏิสนฺธิ\-กฺขเณ ฯเปฯ. (3)}

\prose{อวิตกฺก\-วิจารมตฺตํ ธมฺมํ ปฏิจฺจ สวิตกฺก\-สวิจาโร จ อวิตกฺก\-อวิจาโร จ ธมฺมา อุปฺปชฺชนฺติ นปุเรชาต\-ปจฺจยา. ปฏิสนฺธิ\-กฺขเณ วิตกฺกํ ปฏิจฺจ สวิตกฺก\-สวิจารา ขนฺธา กฏตฺตา จ รูปํ. (4)}

\prose{อวิตกฺก\-วิจารมตฺตํ ธมฺมํ ปฏิจฺจ อวิตกฺก\-วิจารมตฺโต จ อวิตกฺก\-อวิจาโร จ ธมฺมา อุปฺปชฺชนฺติ นปุเรชาต\-ปจฺจยา ฯเปฯ. อรูเป อวิตกฺก\-วิจารมตฺตํ เอกํ ขนฺธํ ปฏิจฺจ ตโย ขนฺธา วิจาโร จ ฯเปฯ. ปฏิสนฺธิ\-กฺขเณ ฯเปฯ. (5)}

\proseitem{36}{อวิตกฺก\-อวิจารํ ธมฺมํ ปฏิจฺจ อวิตกฺก\-อวิจาโร ธมฺโม อุปฺปชฺชติ นปุเรชาต\-ปจฺจยา. อรูเป อวิตกฺก\-อวิจารํ เอกํ ขนฺธํ ปฏิจฺจ ตโย ขนฺธา ฯเปฯ. อวิตกฺก\-อวิจาเร ขนฺเธ ปฏิจฺจ จิตฺต\-สมุฏฺฐานํ รูปํ. วิจารํ ปฏิจฺจ จิตฺต\-สมุฏฺฐานํ รูปํ. ปฏิสนฺธิ\-กฺขเณ ฯเปฯ. (1)}

\prose{อวิตกฺก\-อวิจารํ ธมฺมํ ปฏิจฺจ สวิตกฺก\-สวิจาโร ธมฺโม อุปฺปชฺชติ นปุเรชาต\-ปจฺจยา. ปฏิสนฺธิ\-กฺขเณ วตฺถุํ ปฏิจฺจ สวิตกฺก\-สวิจารา ขนฺธา. (2)}

\prose{อวิตกฺก\-อวิจารํ ธมฺมํ ปฏิจฺจ อวิตกฺก\-วิจารมตฺโต ธมฺโม อุปฺปชฺชติ นปุเรชาต\-ปจฺจยา. อรูเป วิจารํ ปฏิจฺจ อวิตกฺก\-วิจารมตฺตา ขนฺธา. ปฏิสนฺธิ\-กฺขเณ. (สํขิตฺตํ.) (7)}

\prose{สวิตกฺกสวิจารญฺจ อวิตกฺกอวิจารญฺจ ธมฺมํ ปฏิจฺจ. สตฺต.}

\proseitem{37}{อวิตกฺกวิจารมตฺตญฺจ อวิตกฺกอวิจารญฺจ ธมฺมํ ปฏิจฺจ อวิตกฺก\-วิจารมตฺโต ธมฺโม อุปฺปชฺชติ นปุเรชาต\-ปจฺจยา. อรูเป อวิตกฺก\-วิจารมตฺตํ เอกํ ขนฺธญฺจ วิจารญฺจ ปฏิจฺจ ตโย ขนฺธา ฯเปฯ. ปฏิสนฺธิ\-กฺขเณ. (สํขิตฺตํ.) (7)}

\prose{\csromanfolio{26}อวิตกฺกวิจารมตฺตญฺจ อวิตกฺกอวิจารญฺจ ธมฺมํ ปฏิจฺจ อวิตกฺก\-วิจารมตฺโต จ อวิตกฺก\-อวิจาโร จ ธมฺมา อุปฺปชฺชนฺติ ฯเปฯ. ปฏิสนฺธิ\-กฺขเณ. (สํขิตฺตํ.)}

\prose{(นปุเรชาต\-มูลเก ยถา สุทฺธิกํ อรูปํ, ตถา อรูปา กาตพฺพา.)}

\csromanheader{2}{\textbf{นปจฺฉาชาตาทิ}}
\proseitem{38}{สวิตกฺก\-สวิจารํ ธมฺมํ ปฏิจฺจ สวิตกฺก\-สวิจาโร ธมฺโม อุปฺปชฺชติ นปจฺฉาชาต\-ปจฺจยา ฯเปฯ นอาเสวน\-ปจฺจยา. สตฺต.}

\prose{อวิตกฺก\-วิจารมตฺตํ ธมฺมํ ปฏิจฺจ อวิตกฺก\-วิจารมตฺโต ธมฺโม อุปฺปชฺชติ นอาเสวน\-ปจฺจยา. วิปากํ อวิตกฺก\-วิจารมตฺตํ เอกํ ขนฺธํ ปฏิจฺจ. (สํขิตฺตํ.) (5)}

\prose{อวิตกฺก\-อวิจารํ ธมฺมํ ปฏิจฺจ อวิตกฺก\-อวิจาโร ธมฺโม อุปฺปชฺชติ นอาเสวน\-ปจฺจยา. วิปากํ อวิตกฺก\-อวิจารํ เอกํ ขนฺธํ ปฏิจฺจ. (สํขิตฺตํ.)}

\prose{(นอาเสวน\-มูลเก อวิตกฺก\-วิจารมตฺตํ วิปาเกน สห คจฺฉนฺเตน นปุเรชาต\-สทิสํ กาตพฺพํ, อวิตกฺกวิจารมตฺตญฺจ อวิตกฺกวิจารมตฺต\-คจฺฉนฺเตน วิปาโก ทสฺเสตพฺโพ.)}

\csromanheader{2}{\textbf{นกมฺม}}
\proseitem{39}{สวิตกฺก\-สวิจารํ ธมฺมํ ปฏิจฺจ สวิตกฺก\-สวิจาโร ธมฺโม อุปฺปชฺชติ นกมฺมปจฺจยา. สวิตกฺก\-สวิจาเร ขนฺเธ ปฏิจฺจ สวิตกฺก\-สวิจารา เจตนา. (1)}

\prose{อวิตกฺก\-วิจารมตฺตํ ธมฺมํ ปฏิจฺจ อวิตกฺก\-วิจารมตฺโต ธมฺโม อุปฺปชฺชติ นกมฺมปจฺจยา. อวิตกฺก\-วิจารมตฺเต ขนฺเธ ปฏิจฺจ อวิตกฺก\-วิจารมตฺตา เจตนา. (1)}

\prose{อวิตกฺก\-วิจารมตฺตํ ธมฺมํ ปฏิจฺจ สวิตกฺก\-สวิจาโร ธมฺโม อุปฺปชฺชติ นกมฺมปจฺจยา. วิตกฺกํ ปฏิจฺจ สวิตกฺก\-สวิจารา เจตนา. (2)}

\prose{อวิตกฺก\-อวิจารํ ธมฺมํ ปฏิจฺจ อวิตกฺก\-อวิจาโร ธมฺโม อุปฺปชฺชติ นกมฺมปจฺจยา. อวิตกฺก\-อวิจาเร ขนฺเธ ปฏิจฺจ อวิตกฺก\-อวิจารา เจตนา. พาหิรํ. อาหารสมุฏฺฐานํ. อุตุสมุฏฺฐานํ เอกํ มหาภูตํ ปฏิจฺจ ฯเปฯ. (1)}

\prose{อวิตกฺก\-อวิจารํ ธมฺมํ ปฏิจฺจ อวิตกฺก\-วิจารมตฺโต ธมฺโม อุปฺปชฺชติ นกมฺมปจฺจยา. วิจารํ ปฏิจฺจ อวิตกฺก\-วิจารมตฺตา เจตนา. (2)}

\prose{\csromanfolio{27}อวิตกฺกวิจารมตฺตญฺจ อวิตกฺกอวิจารญฺจ ธมฺมํ ปฏิจฺจ อวิตกฺก\-วิจารมตฺโต ธมฺโม อุปฺปชฺชติ นกมฺมปจฺจยา. อวิตกฺก\-วิจารมตฺเต ขนฺเธ จ วิจารญฺจ ปฏิจฺจ อวิตกฺก\-วิจารมตฺตา เจตนา. (1)}

\prose{สวิตกฺกสวิจารญฺจ อวิตกฺกวิจารมตฺตญฺจ ธมฺมํ ปฏิจฺจ สวิตกฺก\-สวิจาโร ธมฺโม อุปฺปชฺชติ นกมฺมปจฺจยา. สวิตกฺก\-สวิจาเร ขนฺเธ จ วิตกฺกญฺจ ปฏิจฺจ สวิตกฺก\-สวิจารา เจตนา. (1)}

\csromanheader{2}{\textbf{นวิปากาทิ}}
\proseitem{40}{สวิตกฺก\-สวิจารํ ธมฺมํ ปฏิจฺจ สวิตกฺก\-สวิจาโร ธมฺโม อุปฺปชฺชติ นวิปาก\-ปจฺจยา ฯเปฯ นอาหารปจฺจยา. พาหิรํ. อุตุสมุฏฺฐานํ ฯเปฯ นอินฺทฺริยปจฺจยา. พาหิรํ. อาหารสมุฏฺฐานํ. อุตุสมุฏฺฐานํ ฯเปฯ. อสญฺญสตฺตานํ มหาภูเต ปฏิจฺจ รูปชีวิตินฺทฺริยํ ฯเปฯ นฌานปจฺจยา. ปญฺจ\-วิญฺญาณสหคตํ เอกํ ขนฺธํ ฯเปฯ. พาหิรํ. อาหารสมุฏฺฐานํ. อุตุสมุฏฺฐานํ. อสญฺญสตฺตานํ ฯเปฯ นมคฺคปจฺจยา. นสมฺปยุตฺต\-ปจฺจยา.}

\csromanheader{2}{\textbf{นวิปฺปยุตฺต}}
\proseitem{41}{นวิปฺปยุตฺต\-ปจฺจยา. อรูเป สวิตกฺก\-สวิจารํ เอกํ ขนฺธํ ปฏิจฺจ ตโย ขนฺธา ฯเปฯ เทฺว ขนฺเธ ปฏิจฺจ เทฺว ขนฺธา. (1)}

\prose{สวิตกฺก\-สวิจารํ ธมฺมํ ปฏิจฺจ อวิตกฺก\-วิจารมตฺโต ธมฺโม อุปฺปชฺชติ นวิปฺปยุตฺต\-ปจฺจยา. อรูเป สวิตกฺก\-สวิจาเร ขนฺเธ ปฏิจฺจ วิตกฺโก. (2)}

\prose{สวิตกฺก\-สวิจารํ ธมฺมํ ปฏิจฺจ สวิตกฺก\-สวิจาโร จ อวิตกฺก\-วิจารมตฺโต จ ธมฺมา อุปฺปชฺชนฺติ นวิปฺปยุตฺต\-ปจฺจยา. อรูเป สวิตกฺก\-สวิจารํ เอกํ ขนฺธํ ปฏิจฺจ ตโย ขนฺธา วิตกฺโก จ ฯเปฯ เทฺว ขนฺเธ ปฏิจฺจ เทฺว ขนฺธา วิตกฺโก จ. (3)}

\proseitem{42}{อวิตกฺก\-วิจารมตฺตํ ธมฺมํ ปฏิจฺจ อวิตกฺก\-วิจารมตฺโต ธมฺโม อุปฺปชฺชติ นวิปฺปยุตฺต\-ปจฺจยา. อรูเป อวิตกฺก\-วิจารมตฺตํ เอกํ ขนฺธํ ปฏิจฺจ ตโย ขนฺธา ฯเปฯ เทฺว ขนฺเธ ปฏิจฺจ เทฺว ขนฺธา. (1)}

\prose{อวิตกฺก\-วิจารมตฺตํ ธมฺมํ ปฏิจฺจ สวิตกฺก\-สวิจาโร ธมฺโม อุปฺปชฺชติ นวิปฺปยุตฺต\-ปจฺจยา. อรูเป วิตกฺกํ ปฏิจฺจ สวิตกฺก\-สวิจารา ขนฺธา. (2)}

\prose{\csromanfolio{28}อวิตกฺก\-วิจารมตฺตํ ธมฺมํ ปฏิจฺจ อวิตกฺก\-อวิจาโร ธมฺโม อุปฺปชฺชติ นวิปฺปยุตฺต\-ปจฺจยา. อรูเป อวิตกฺก\-วิจารมตฺเต ขนฺเธ ปฏิจฺจ วิจาโร. (3)}

\prose{อวิตกฺก\-วิจารมตฺตํ ธมฺมํ ปฏิจฺจ อวิตกฺก\-วิจารมตฺโต จ อวิตกฺก\-อวิจาโร จ ธมฺมา อุปฺปชฺชนฺติ นวิปฺปยุตฺต\-ปจฺจยา. อรูเป อวิตกฺก\-วิจารมตฺตํ เอกํ ขนฺธํ ปฏิจฺจ ตโย ขนฺธา วิจาโร จ ฯเปฯ เทฺว ขนฺเธ ปฏิจฺจ เทฺว ขนฺธา วิจาโร จ. (4)}

\prose{อวิตกฺก\-อวิจารํ ธมฺมํ ปฏิจฺจ อวิตกฺก\-อวิจาโร ธมฺโม อุปฺปชฺชติ นวิปฺปยุตฺต\-ปจฺจยา. อรูเป อวิตกฺก\-อวิจารํ เอกํ ขนฺธํ ปฏิจฺจ ตโย ขนฺธา ฯเปฯ เทฺว ขนฺเธ ปฏิจฺจ เทฺว ขนฺธา. พาหิรํ. อาหารสมุฏฺฐานํ. อุตุสมุฏฺฐานํ. อสญฺญสตฺตานํ เอกํ มหาภูตํ ปฏิจฺจ ฯเปฯ. (1)}

\prose{อวิตกฺก\-อวิจารํ ธมฺมํ ปฏิจฺจ อวิตกฺก\-วิจารมตฺโต ธมฺโม อุปฺปชฺชติ นวิปฺปยุตฺต\-ปจฺจยา. อรูเป วิจารํ ปฏิจฺจ อวิตกฺก\-วิจารมตฺตา ขนฺธา. (2)}

\prose{อวิตกฺกวิจารมตฺตญฺจ อวิตกฺกอวิจารญฺจ ธมฺมํ ปฏิจฺจ อวิตกฺก\-วิจารมตฺโต ธมฺโม อุปฺปชฺชติ นวิปฺปยุตฺต\-ปจฺจยา. อรูเป อวิตกฺก\-วิจารมตฺตํ เอกํ ขนฺธญฺจ วิจารญฺจ ปฏิจฺจ ตโย ขนฺธา ฯเปฯ เทฺว ขนฺเธ จ วิจารญฺจ ปฏิจฺจ เทฺว ขนฺธา. (1)}

\prose{สวิตกฺกสวิจารญฺจ อวิตกฺกวิจารมตฺตญฺจ ธมฺมํ ปฏิจฺจ สวิตกฺก\-สวิจาโร ธมฺโม อุปฺปชฺชติ นวิปฺปยุตฺต\-ปจฺจยา. อรูเป สวิตกฺก\-สวิจารํ เอกํ ขนฺธญฺจ วิตกฺกญฺจ ปฏิจฺจ ตโย ขนฺธา ฯเปฯ เทฺว ขนฺเธ จ วิตกฺกญฺจ ปฏิจฺจ เทฺว ขนฺธา. (1)}

\csromanheader{2}{\textbf{โนนตฺถิ โนวิคต}}
\proseitem{43}{สวิตกฺก\-สวิจารํ ธมฺมํ ปฏิจฺจ อวิตกฺก\-อวิจาโร ธมฺโม อุปฺปชฺชติ โนนตฺถิปจฺจยา. โนวิคต\-ปจฺจยา. (สํขิตฺตํ.)}

\csromanmarkh{2. ปจฺจย\-ปจฺจนีย}\csromanmarkhii{2. สงฺขฺยาวาร}\csromanheadernomark{2}{2. ปจฺจย\-ปจฺจนีย 2. สงฺขฺยาวาร}
\proseitemwithcloserruled{44}{นเหตุยา เตตฺติํส, นอารมฺมเณ สตฺต, นอธิปติยา สตฺตติํส, นอนนฺตเร สตฺต, นสมนนฺตเร สตฺต, นอญฺญมญฺเญ สตฺต, นอุปนิสฺสเย สตฺต, นปุเรชาเต สตฺตติํส, นปจฺฉาชาเต สตฺตติํส, นอาเสวเน สตฺตติํส, นกมฺเม สตฺต, นวิปาเก เตวีส, นอาหาเร \csromanfolio{29}เอกํ, นอินฺทฺริเย เอกํ, นฌาเน เอกํ, นมคฺเค เตตฺติํส, นสมฺปยุตฺเต สตฺต, นวิปฺปยุตฺเต เอกาทส, โนนตฺถิยา สตฺต, โนวิคเต สตฺต. (สํขิตฺตํ.) (ยถา กุสลตฺติเก ปจฺจนีย\-คณนา, เอวํ คเณตพฺพํ.)}{ปจฺจนียํ.}
\csromancenter{\textbf{เหตุ}ปจฺจยา นอารมฺมเณ สตฺต ฯเปฯ โนวิคเต สตฺต.}
\csromancenter{(ยถา กุสลตฺติเก อนุโลม\-ปจฺจนีย\-คณนา, เอวํ คเณตพฺพํ.)}
\proseitem{46}{นเหตุปจฺจยา อารมฺมเณ จุทฺทส, อนนฺตเร สมนนฺตเร จุทฺทส, สหชาเต เตตฺติํส, อญฺญมญฺเญ พาวีส, นิสฺสเย เตตฺติํส, อุปนิสฺสเย จุทฺทส, ปุเรชาเต ฉ, อาเสวเน ปญฺจ, กมฺเม เตตฺติํส ฯเปฯ ฌาเน เตตฺติํส, มคฺเค ตีณิ, สมฺปยุตฺเต จุทฺทส, วิปฺปยุตฺเต เตตฺติํส ฯเปฯ อวิคเต เตตฺติํส. (สํขิตฺตํ.) (ยถา กุสลตฺติเก ปจฺจนียา\-นุโลม\-คณนา, เอวํ คเณตพฺพํ.) ปฏิจฺจวาโร. (สหชาตวาโรปิ ปฏิจฺจวารสทิโส กาตพฺโพ.)}
\csromansectionrule

\vspace{\tipitakaparskip}
\csromancenter{(สหชาตวาโรปิ ปฏิจฺจ\-วาร\-สทิโส กาตพฺโพ.)}
\csromanmarkh{6. วิตกฺกตฺติก}\csromanmarkhii{3. ปจฺจยวาร}\csromanheadernomark{2}{6. วิตกฺกตฺติก 3. ปจฺจยวาร}
\csromanheader{2}{1. \textbf{ปจฺจยานุโลมฺ}อ}
\vspace{\tipitakaparskip}
\csromancenter{\textbf{เหตุ}}
\proseitem{47}{สวิตกฺก\-สวิจารํ ธมฺมํ ปจฺจยา สวิตกฺก\-สวิจาโร ธมฺโม อุปฺปชฺชติ เหตุปจฺจยา. สวิตกฺก\-สวิจารํ เอกํ ขนฺธํ ปจฺจยา ตโย ขนฺธา ฯเปฯ เทฺว ขนฺธา. สตฺต.}

\prose{\csromanfolio{30}อวิตกฺก\-วิจารมตฺตํ ธมฺมํ ปจฺจยา. ปญฺจ. (ปฏิจฺจ\-วาร\-สทิสา.)}

\prose{อวิตกฺก\-อวิจารํ ธมฺมํ ปจฺจยา อวิตกฺก\-อวิจาโร ธมฺโม อุปฺปชฺชติ เหตุปจฺจยา. อวิตกฺก\-อวิจารํ เอกํ ขนฺธํ ปจฺจยา ตโย ขนฺธา จิตฺต\-สมุฏฺฐานญฺจ รูปํ ฯเปฯ เทฺว ขนฺเธ ปจฺจยา ฯเปฯ วิจารํ ปจฺจยา จิตฺต\-สมุฏฺฐานํ รูปํ. ปฏิสนฺธิ\-กฺขเณ ฯเปฯ วตฺถุํ ปจฺจยา อวิตกฺก\-อวิจารา ขนฺธา. วตฺถุํ ปจฺจยา วิจาโร. (1)}

\prose{อวิตกฺก\-อวิจารํ ธมฺมํ ปจฺจยา สวิตกฺก\-สวิจาโร ธมฺโม อุปฺปชฺชติ ฯเปฯ วตฺถุํ ปจฺจยา สวิตกฺก\-สวิจารา ขนฺธา. ปฏิสนฺธิ\-กฺขเณ ฯเปฯ. (2)}

\prose{อวิตกฺก\-อวิจารํ ธมฺมํ ปจฺจยา อวิตกฺก\-วิจารมตฺโต ธมฺโม ฯเปฯ วิจารํ ปจฺจยา อวิตกฺก\-วิจารมตฺตา ขนฺธา. วตฺถุํ ปจฺจยา อวิตกฺก\-วิจารมตฺตา ขนฺธา. วตฺถุํ ปจฺจยา วิตกฺโก. ปฏิสนฺธิ\-กฺขเณ ฯเปฯ. (3)}

\prose{อวิตกฺก\-อวิจารํ ธมฺมํ ปจฺจยา สวิตกฺก\-สวิจาโร จ อวิตกฺก\-อวิจาโร จ ธมฺมา อุปฺปชฺชนฺติ ฯเปฯ วตฺถุํ ปจฺจยา สวิตกฺก\-สวิจารา ขนฺธา. มหาภูเต ปจฺจยา จิตฺต\-สมุฏฺฐานํ รูปํ. ปฏิสนฺธิ\-กฺขเณ ฯเปฯ. (4)}

\prose{อวิตกฺก\-อวิจารํ ธมฺมํ ปจฺจยา อวิตกฺก\-วิจารมตฺโต จ อวิตกฺก\-อวิจาโร จ ธมฺมา ฯเปฯ วิจารํ ปจฺจยา อวิตกฺก\-วิจารมตฺตา ขนฺธา จิตฺต\-สมุฏฺฐานญฺจ รูปํ. วตฺถุํ ปจฺจยา อวิตกฺก\-วิจารมตฺตา ขนฺธา. มหาภูเต ปจฺจยา จิตฺต\-สมุฏฺฐานํ รูปํ. วตฺถุํ ปจฺจยา วิตกฺโก. มหาภูเต ปจฺจยา จิตฺต\-สมุฏฺฐานํ รูปํ. วตฺถุํ ปจฺจยา อวิตกฺก\-วิจารมตฺตา ขนฺธา จ วิจาโร จ. ปฏิสนฺธิ\-กฺขเณ ฯเปฯ. (5)}

\prose{อวิตกฺก\-อวิจารํ ธมฺมํ ปจฺจยา สวิตกฺก\-สวิจาโร จ อวิตกฺก\-วิจารมตฺโต จ ธมฺมา ฯเปฯ. วตฺถุํ ปจฺจยา สวิตกฺก\-สวิจารา ขนฺธา จ วิตกฺโก จ. ปฏิสนฺธิ\-กฺขเณ ฯเปฯ. (6)}

\prose{อวิตกฺก\-อวิจารํ ธมฺมํ ปจฺจยา สวิตกฺก\-สวิจาโร จ อวิตกฺก\-วิจารมตฺโต จ อวิตกฺก\-อวิจาโร จ ธมฺมา ฯเปฯ. วตฺถุํ ปจฺจยา สวิตกฺก\-สวิจารา ขนฺธา จ วิตกฺโก จ. มหาภูเต ปจฺจยา จิตฺต\-สมุฏฺฐานํ รูปํ. ปฏิสนฺธิ\-กฺขเณ ฯเปฯ. (7)}

\proseitem{48}{สวิตกฺกสวิจารญฺจ อวิตกฺกอวิจารญฺจ ธมฺมํ ปจฺจยา สวิตกฺก\-สวิจาโร ธมฺโม ฯเปฯ. สวิตกฺก\-สวิจารํ เอกํ ขนฺธญฺจ วตฺถุญฺจ ปจฺจยา ตโย \csromanfolio{31}ขนฺธา ฯเปฯ เทฺว ขนฺเธ ฯเปฯ. ปฏิสนฺธิ\-กฺขเณ ฯเปฯ. สวิตกฺกสวิจารญฺจ อวิตกฺกอวิจารญฺจ ธมฺมํ ปจฺจยา อวิตกฺก\-วิจารมตฺโต ธมฺโม ฯเปฯ.}

\vspace{\tipitakaparskip}
\csromancenter{(ปฐมฺ\-เอาทาหรเณ ปวตฺเต ปฏิสนฺธิ\-กฺขเณ สตฺต ปญฺหา กาตพฺพา.)}
\prose{อวิตกฺกวิจารมตฺตญฺจ อวิตกฺกอวิจารญฺจ ธมฺมํ ปจฺจยา สวิตกฺก\-สวิจาโร ธมฺโม ฯเปฯ. วิตกฺกญฺจ วตฺถุญฺจ ปจฺจยา สวิตกฺก\-สวิจารา ขนฺธา. ปฏิสนฺธิ\-กฺขเณ ฯเปฯ. (1)}

\prose{อวิตกฺกวิจารมตฺตญฺจ อวิตกฺกอวิจารญฺจ ธมฺมํ ปจฺจยา อวิตกฺก\-วิจารมตฺโต ธมฺโม ฯเปฯ. อวิตกฺก\-วิจารมตฺตํ เอกํ ขนฺธญฺจ วิจารญฺจ ปจฺจยา ตโย ขนฺธา ฯเปฯ. อวิตกฺก\-วิจารมตฺตํ เอกํ ขนฺธญฺจ วตฺถุญฺจ ปจฺจยา ตโย ขนฺธา ฯเปฯ. ปฏิสนฺธิ\-กฺขเณ อวิตกฺก\-วิจารมตฺตํ เอกํ ขนฺธญฺจ วิจารญฺจ ปจฺจยา ตโย ขนฺธา ฯเปฯ. ปฏิสนฺธิ\-กฺขเณ อวิตกฺก\-วิจารมตฺตํ เอกํ ขนฺธญฺจ วตฺถุญฺจ ปจฺจยา ตโย ขนฺธา ฯเปฯ. (2)}

\prose{อวิตกฺกวิจารมตฺตญฺจ อวิตกฺกอวิจารญฺจ ธมฺมํ ปจฺจยา อวิตกฺก\-อวิจาโร ธมฺโม อุปฺปชฺชติ ฯเปฯ. อวิตกฺก\-วิจารมตฺเต ขนฺเธ จ วิจารญฺจ ปจฺจยา จิตฺต\-สมุฏฺฐานํ รูปํ. อวิตกฺก\-วิจารมตฺเต ขนฺเธ จ มหาภูเต จ ปจฺจยา จิตฺต\-สมุฏฺฐานํ รูปํ. วิตกฺกญฺจ มหาภูเต จ ปจฺจยา จิตฺต\-สมุฏฺฐานํ รูปํ. อวิตกฺก\-วิจารมตฺเต ขนฺเธ จ วตฺถุญฺจ ปจฺจยา วิจาโร. (เอวํ ปฏิสนฺธิ\-กฺขเณ จตฺตาโร.) (3)}

\prose{อวิตกฺกวิจารมตฺตญฺจ อวิตกฺกอวิจารญฺจ ธมฺมํ ปจฺจยา สวิตกฺก\-สวิจาโร จ อวิตกฺก\-อวิจาโร จ ธมฺมา อุปฺปชฺชนฺติ ฯเปฯ. วิตกฺกญฺจ วตฺถุญฺจ ปจฺจยา สวิตกฺก\-สวิจารา ขนฺธา. วิตกฺกญฺจ มหาภูเต จ ปจฺจยา จิตฺต\-สมุฏฺฐานํ รูปํ. ปฏิสนฺธิ\-กฺขเณ ฯเปฯ. (4)}

\prose{อวิตกฺกวิจารมตฺตญฺจ อวิตกฺกอวิจารญฺจ ธมฺมํ ปจฺจยา อวิตกฺก\-วิจารมตฺโต จ อวิตกฺก\-อวิจาโร จ ธมฺมา อุปฺปชฺชนฺติ ฯเปฯ. อวิตกฺก\-วิจารมตฺตํ เอกํ ขนฺธญฺจ วิจารญฺจ ปจฺจยา ตโย ขนฺธา จิตฺต\-สมุฏฺฐานญฺจ รูปํ ฯเปฯ. อวิตกฺก\-วิจารมตฺตํ เอกํ ขนฺธญฺจ วตฺถุญฺจ ปจฺจยา ตโย ขนฺธา ฯเปฯ. อวิตกฺก\-วิจารมตฺเต ขนฺเธ จ มหาภูเต จ ปจฺจยา จิตฺต\-สมุฏฺฐานํ รูปํ. อวิตกฺก\-วิจารมตฺตํ เอกํ ขนฺธญฺจ วตฺถุญฺจ ปจฺจยา ตโย ขนฺธา วิจาโร จ ฯเปฯ. ปฏิสนฺธิ\-กฺขเณ ตโย ขนฺธา ฯเปฯ. (5) (อวเสเสสุ ทฺวีสุ ฆฏเนสุ ปวตฺติ\-ปฏิสนฺธิ วิตฺถาเรตพฺพา.)}

\vspace{\tipitakaparskip}
\csromancenter{เหตุปจฺจยา.}
\prose{\csromanfolio{32}(เหตุปจฺจยํ อนุมชฺชนฺเตน ปจฺจยวาโร วิตฺถาเรตพฺโพ, ยถา ปฏิจฺจคณนา, เอวํ คเณตพฺพา. อธิปติยา สตฺตติํส, ปุเรชาเต จ อาเสวเน จ เอกวีส, อยํ เอตฺถ วิเสโส.)}
\csromansectionrule

\csromanheader{2}{2. ปจฺจย\-ปจฺจนีย}
\proseitem{49}{ปจฺจนีเย\csromandash{}นเหตุยา เตตฺติํส ปญฺหา, สตฺตสุ ฐาเนสุ สตฺต โมหา อุทฺธริตพฺพา มูลปเทสุ เยว.  นอารมฺมเณ สตฺต จิตฺต\-สมุฏฺฐานา อุทฺธริตพฺพา.}

\prose{นอธิปติ}

\proseitem{50}{สวิตกฺกสวิจาร\-มูลกา สตฺต ปญฺหา นอธิปติยา กาตพฺพา.}

\prose{อวิตกฺก\-วิจารมตฺตํ ธมฺมํ ปจฺจยา อวิตกฺก\-วิจารมตฺโต ธมฺโม อุปฺปชฺชติ นอธิปติ\-ปจฺจยา. อวิตกฺก\-วิจารมตฺเต ขนฺเธ ปจฺจยา อวิตกฺก\-วิจารมตฺตา อธิปติ. วิปากํ อวิตกฺก\-วิจารมตฺตํ เอกํ ขนฺธํ ปจฺจยา ตโย ขนฺธา ฯเปฯ. ปฏิสนฺธิ\-กฺขเณ ฯเปฯ. (1)}

\prose{อวิตกฺก\-วิจารมตฺตํ ธมฺมํ ปจฺจยา. (ยถา ปฏิจฺจนเย, ตถา ปญฺจ ปญฺหา กาตพฺพา.)}

\vspace{\tipitakaparskip}
\csromancenter{(ยถา ปฏิจฺจนเย, ตถา ปญฺจ ปญฺหา กาตพฺพา.)}
\proseitem{51}{อวิตกฺก\-อวิจารํ ธมฺมํ ปจฺจยา อวิตกฺก\-อวิจาโร ธมฺโม อุปฺปชฺชติ ฯเปฯ. อวิตกฺก\-อวิจาเร ขนฺเธ ปจฺจยา อวิตกฺก\-อวิจารา อธิปติ. วิปากํ อวิตกฺก\-อวิจารํ เอกํ ขนฺธํ ปจฺจยา ตโย ขนฺธา จิตฺต\-สมุฏฺฐานญฺจ รูปํ ฯเปฯ. วิปากํ วิจารํ ปจฺจยา จิตฺต\-สมุฏฺฐานํ รูปํ. ปฏิสนฺธิ\-กฺขเณ ฯเปฯ. วตฺถุํ ปจฺจยา อวิตกฺก\-อวิจารา อธิปติ. วตฺถุํ ปจฺจยา วิปากา อวิตกฺก\-อวิจารา ขนฺธา จ วิจาโร จ ฯเปฯ. (1)}

\prose{อวิตกฺก\-อวิจารํ ธมฺมํ ปจฺจยา สวิตกฺก\-สวิจาโร ธมฺโม ฯเปฯ. วตฺถุํ ปจฺจยา สวิตกฺก\-สวิจารา ขนฺธา. ปฏิสนฺธิ\-กฺขเณ ฯเปฯ. (2)}

\prose{อวิตกฺก\-อวิจารํ ธมฺมํ ปจฺจยา อวิตกฺก\-วิจารมตฺโต ธมฺโม ฯเปฯ. วิจารํ ปจฺจยา อวิตกฺก\-วิจารมตฺตา อธิปติ. วตฺถุํ ปจฺจยา อวิตกฺก\-วิจารมตฺตา อธิปติ. วิปากํ วิจารํ ปจฺจยา อวิตกฺก\-วิจารมตฺตา ขนฺธา. วตฺถุํ ปจฺจยา วิปากา อวิตกฺก\-วิจารมตฺตา ขนฺธา. ปฏิสนฺธิ\-กฺขเณ ฯเปฯ. (3)}

\prose{\csromanfolio{33}อวิตกฺก\-อวิจารํ ธมฺมํ ปจฺจยา สวิตกฺก\-สวิจาโร จ อวิตกฺก\-อวิจาโร จ ธมฺมา ฯเปฯ. วตฺถุํ ปจฺจยา สวิตกฺก\-สวิจารา ขนฺธา. มหาภูเต ปจฺจยา จิตฺต\-สมุฏฺฐานํ รูปํ. ปฏิสนฺธิ\-กฺขเณ ฯเปฯ. (4)}

\prose{อวิตกฺก\-อวิจารํ ธมฺมํ ปจฺจยา อวิตกฺก\-วิจารมตฺโต จ อวิตกฺก\-อวิจาโร จ ธมฺมา ฯเปฯ. วิปากํ วิจารํ ปจฺจยา อวิตกฺก\-วิจารมตฺตา ขนฺธา จิตฺต\-สมุฏฺฐานญฺจ รูปํ. วตฺถุํ ปจฺจยา วิปากา อวิตกฺก\-วิจารมตฺตา ขนฺธา. มหาภูเต ปจฺจยา จิตฺต\-สมุฏฺฐานํ รูปํ. วตฺถุํ ปจฺจยา วิตกฺโก. มหาภูเต ปจฺจยา จิตฺต\-สมุฏฺฐานํ รูปํ. วตฺถุํ ปจฺจยา วิปากา อวิตกฺก\-วิจารมตฺตา ขนฺธา จ วิจาโร จ. ปฏิสนฺธิ\-กฺขเณ ฯเปฯ. (5)}

\prose{อวิตกฺก\-อวิจารํ ธมฺมํ ปจฺจยา สวิตกฺก\-สวิจาโร จ อวิตกฺก\-วิจารมตฺโต จ ธมฺมา ฯเปฯ. วตฺถุํ ปจฺจยา สวิตกฺก\-สวิจารา ขนฺธา จ วิตกฺโก จ. ปฏิสนฺธิ\-กฺขเณ ฯเปฯ. (6) (ปฐม\-ฆฏนายํ สมฺปุณฺณา สตฺต ปญฺหา กาตพฺพา.)}

\proseitem{52}{อวิตกฺกวิจารมตฺตญฺจ อวิตกฺกอวิจารญฺจ ธมฺมํ ปจฺจยา สวิตกฺก\-สวิจาโร ธมฺโม อุปฺปชฺชติ ฯเปฯ. วิตกฺกญฺจ วตฺถุญฺจ ปจฺจยา สวิตกฺก\-สวิจารา ขนฺธา. ปฏิสนฺธิ\-กฺขเณ ฯเปฯ.}

\prose{อวิตกฺกวิจารมตฺตญฺจ อวิตกฺกอวิจารญฺจ ธมฺมํ ปจฺจยา อวิตกฺก\-วิจารมตฺโต ธมฺโม อุปฺปชฺชติ ฯเปฯ. อวิตกฺก\-วิจารมตฺเต ขนฺเธ จ วิจารญฺจ ปจฺจยา อวิตกฺก\-วิจารมตฺตา อธิปติ. อวิตกฺก\-วิจารมตฺเต ขนฺเธ จ วตฺถุญฺจ ปจฺจยา อวิตกฺก\-วิจารมตฺตา อธิปติ. วิปากํ อวิตกฺก\-วิจารมตฺตํ เอกํ ขนฺธญฺจ วิจารญฺจ ปจฺจยา ฯเปฯ. วิปากํ อวิตกฺก\-วิจารมตฺตํ เอกํ ขนฺธญฺจ วตฺถุญฺจ ปจฺจยา ตโย ขนฺธา ฯเปฯ.}

\prose{(ปฏิสนฺธิ\-กฺขเณ ปญฺจ ปญฺหา กาตพฺพา, ยตฺถ อวิตกฺก\-วิจารมตฺตํ อาคจฺฉติ, ตตฺถ วิปากํ กาตพฺพํ. นอธิป\-ติมูลเก สตฺตติํส ปญฺหา กาตพฺพา.)}

\prose{นอนนฺตราทิ}

\proseitem{53}{นอนนฺตรมฺปิ\csromansharedfootnote{222fca748e6b5255}{นอนนฺตเรปิ (ก)} นสมนนฺตรมฺปิ นอญฺญมญฺญมฺปิ นอุปนิสฺสยมฺปิ สตฺต ปญฺหา รูปํเยว. นปุเรชาเต สตฺตติํส, ปฏิจฺจ\-วาร\-ปจฺจนีย\-สทิสํ. นปจฺฉาชาเต \csromanfolio{34}สตฺตติํส, นอาเสวเนปิ สทิสํ. ยตฺถ อวิตกฺก\-วิจารมตฺโตปิ อาคจฺฉติ, ตตฺถ วิปากา กาตพฺพา.}

\csromanheader{2}{\textbf{นกมฺม}}
\proseitem{54}{สวิตกฺก\-สวิจารํ ธมฺมํ ปจฺจยา สวิตกฺก\-สวิจาโร ธมฺโม อุปฺปชฺชติ น กมฺมปจฺจยา. สวิตกฺก\-สวิจาเร ขนฺเธ ปจฺจยา สวิตกฺก\-สวิจารา เจตนา. (1)}

\prose{อวิตกฺก\-วิจารมตฺตํ ธมฺมํ ปจฺจยา อวิตกฺก\-วิจารมตฺโต ธมฺโม ฯเปฯ. อวิตกฺก\-วิจารมตฺเต ขนฺเธ ปจฺจยา อวิตกฺก\-วิจารมตฺตา เจตนา. สวิตกฺก\-สวิจาโร ธมฺโม ฯเปฯ. วิตกฺกํ ปจฺจยา สวิตกฺก\-สวิจารา เจตนา. (2)}

\prose{อวิตกฺก\-อวิจารํ ธมฺมํ ปจฺจยา อวิตกฺก\-อวิจาโร ธมฺโม ฯเปฯ. อวิตกฺก\-อวิจารา เจตนา ฯเปฯ. (ปริปุณฺณํ กาตพฺพํ.) สวิตกฺก\-สวิจาโร ฯเปฯ. วตฺถุํ ปจฺจยา สวิตกฺก\-สวิจารา เจตนา. อวิตกฺก\-วิจารมตฺโต ฯเปฯ. วิจารํ ปจฺจยา อวิตกฺก\-วิจารมตฺตา เจตนา. วตฺถุํ ปจฺจยา อวิตกฺก\-วิจารมตฺตา เจตนา. (3)}

\proseitem{55}{สวิตกฺกสวิจารญฺจ อวิตกฺกอวิจารญฺจ ธมฺมํ ปจฺจยา สวิตกฺก\-สวิจาโร ธมฺโม ฯเปฯ. สวิตกฺก\-สวิจาเร ขนฺเธ จ วตฺถุญฺจ ปจฺจยา สวิตกฺก\-สวิจารา เจตนา. (1)}

\prose{อวิตกฺกวิจารมตฺตญฺจ อวิตกฺกอวิจารญฺจ ธมฺมํ ปจฺจยา สวิตกฺก\-สวิจาโร ธมฺโม ฯเปฯ. วิตกฺกญฺจ วตฺถุญฺจ ปจฺจยา สวิตกฺก\-สวิจารา เจตนา. (1)}

\prose{อวิตกฺกวิจารมตฺตญฺจ อวิตกฺกอวิจารญฺจ ธมฺมํ ปจฺจยา อวิตกฺก\-วิจารมตฺโต ธมฺโม ฯเปฯ. อวิตกฺก\-วิจารมตฺเต ขนฺเธ จ วิจารญฺจ ปจฺจยา อวิตกฺก\-วิจารมตฺตา เจตนา. อวิตกฺก\-วิจารมตฺเต ขนฺเธ จ วตฺถุญฺจ ปจฺจยา อวิตกฺก\-วิจารมตฺตา เจตนา. (2)}

\prose{สวิตกฺกสวิจารญฺจ อวิตกฺกวิจารมตฺตญฺจ ธมฺมํ ปจฺจยา สวิตกฺก\-สวิจาโร ธมฺโม ฯเปฯ. สวิตกฺก\-สวิจาเร ขนฺเธ จ วิตกฺกญฺจ ปจฺจยา สวิตกฺก\-สวิจารา เจตนา. (1)}

\prose{\csromanfolio{35}สวิตกฺกสวิจารญฺจ อวิตกฺกวิจารมตฺตญฺจ อวิตกฺกอวิจารญฺจ ธมฺมํ ปจฺจยา สวิตกฺก\-สวิจาโร ธมฺโม อุปฺปชฺชติ นกมฺมปจฺจยา. สวิตกฺก\-สวิจาเร ขนฺเธ จ วิตกฺกญฺจ วตฺถุญฺจ ปจฺจยา สวิตกฺก\-สวิจารา เจตนา. (1)}

\prose{(นวิปาเก สตฺตติํส ปญฺหา กาตพฺพา. นอาหาร นอินฺทฺริย นฌาน นมคฺค นสมฺปยุตฺต นวิปฺปยุตฺต โนนตฺถิ โนวิคต\-ปจฺจยา วิตฺถาเรตพฺพา.)}

\csromanmarkh{2. ปจฺจย\-ปจฺจนีย}\csromanmarkhii{2. สงฺขฺยาวาร}\csromanheadernomark{2}{2. ปจฺจย\-ปจฺจนีย 2. สงฺขฺยาวาร}
\proseitem{56}{นเหตุยา เตตฺติํส, นอารมฺมเณ สตฺต, นอธิปติยา สตฺตติํส, นอนนฺตเร นสมนนฺตเร นอญฺญมญฺเญ นอุปนิสฺสเย สตฺต, นปุเรชาเต นปจฺฉาชาเต นอาเสวเน สตฺตติํส, นกมฺเม เอกาทส, นวิปาเก สตฺตติํส, นอาหาเร เอกํ, นอินฺทฺริเย เอกํ, นฌาเน เอกํ, นมคฺเค เตตฺติํส, นสมฺปยุตฺเต สตฺต, นวิปฺปยุตฺเต เอกาทส, โนนตฺถิยา สตฺต, โนวิคเต สตฺต.}
\csromansectionrule

\vspace{\tipitakaparskip}
\csromancenter{เหตุปจฺจยา นอารมฺมเณ สตฺต ฯเปฯ โนวิคเต สตฺต.}
\nitthitamruled{อนุโลม\-ปจฺจนียํ.}
\proseitem{58}{นเหตุปจฺจยา อารมฺมเณ อนนฺตเร สมนนฺตเร จุทฺทส, สหชาเต เตตฺติํส, อญฺญมญฺเญ พาวีส, นิสฺสเย เตตฺติํส, อุปนิสฺสเย ปุเรชาเต จุทฺทส, อาเสวเน เตรส, กมฺเม วิปาเก อาหาเร อินฺทฺริเย ฌาเน เตตฺติํส, มคฺเค ปญฺจ, สมฺปยุตฺเต จุทฺทส, วิปฺปยุตฺเต อตฺถิยา เตตฺติํส ฯเปฯ อวิคเต เตตฺติํส.}

\vspace{\tipitakaparskip}
\csromancenter{ปจฺจนียา\-นุโลมํ.}
\csromancenter{ปจฺจยวาโร.}
\csromancenter{(นิสฺสยมฺปิ นินฺนานํ.)}
\csromansectionrule
\csromanmarkh{6. วิตกฺกตฺติก}\csromanmarkhii{5. สํสฏฺฐวาร}\csromanheadernomark{2}{6. \textbf{วิตกฺกตฺติก 5. สํสฏฺฐวาร}}
\csromanheader{2}{1. \textbf{ปจฺจยานุโลม}}
\csromanheader{2}{\textbf{เหตุ}}
\proseitem{59}{\csromanfolio{36}สวิตกฺก\-สวิจารํ ธมฺมํ สํสฏฺโฐ สวิตกฺก\-สวิจาโร ธมฺโม อุปฺปชฺชติ เหตุปจฺจยา. สวิตกฺก\-สวิจารํ เอกํ ขนฺธํ สํสฏฺฐา ตโย ขนฺธา ฯเปฯ เทฺว ขนฺเธ ฯเปฯ. ปฏิสนฺธิ\-กฺขเณ ฯเปฯ. (1)}

\prose{สวิตกฺก\-สวิจารํ ธมฺมํ สํสฏฺโฐ อวิตกฺก\-วิจารมตฺโต ธมฺโม ฯเปฯ. สวิตกฺก\-สวิจาเร ขนฺเธ สํสฏฺโฐ วิตกฺโก. ปฏิสนฺธิ\-กฺขเณ ฯเปฯ. (2)}

\prose{สวิตกฺก\-สวิจารํ ธมฺมํ สํสฏฺโฐ สวิตกฺก\-สวิจาโร จ อวิตกฺก\-วิจารมตฺโต จ ธมฺมา ฯเปฯ. สวิตกฺก\-สวิจารํ เอกํ ขนฺธํ สํสฏฺฐา ตโย ขนฺธา วิตกฺโก จ ฯเปฯ เทฺว ขนฺเธ ฯเปฯ. ปฏิสนฺธิ\-กฺขเณ ฯเปฯ. (3)}

\proseitem{60}{อวิตกฺก\-วิจารมตฺตํ ธมฺมํ สํสฏฺโฐ อวิตกฺก\-วิจารมตฺโต ธมฺโม อุปฺปชฺชติ เหตุปจฺจยา. อวิตกฺก\-วิจารมตฺตํ เอกํ ขนฺธํ สํสฏฺฐา ตโย ขนฺธา ฯเปฯ เทฺว ขนฺเธ ฯเปฯ. ปฏิสนฺธิ\-กฺขเณ ฯเปฯ. (1)}

\prose{อวิตกฺก\-วิจารมตฺตํ ธมฺมํ สํสฏฺโฐ สวิตกฺก\-สวิจาโร ธมฺโม ฯเปฯ. วิตกฺกํ สํสฏฺฐา สวิตกฺก\-สวิจารา ขนฺธา. ปฏิสนฺธิ\-กฺขเณ ฯเปฯ. (2)}

\prose{อวิตกฺก\-วิจารมตฺตํ ธมฺมํ สํสฏฺโฐ อวิตกฺก\-อวิจาโร ธมฺโม ฯเปฯ. อวิตกฺก\-วิจารมตฺเต ขนฺเธ สํสฏฺโฐ วิจาโร. ปฏิสนฺธิ\-กฺขเณ อวิตกฺก\-วิจารมตฺเต ขนฺเธ สํสฏฺโฐ วิจาโร. (3)}

\prose{อวิตกฺก\-วิจารมตฺตํ ธมฺมํ สํสฏฺโฐ อวิตกฺก\-วิจารมตฺโต จ อวิตกฺก\-อวิจาโร จ ธมฺมา ฯเปฯ. อวิตกฺก\-วิจารมตฺตํ เอกํ ขนฺธํ สํสฏฺฐา ตโย ขนฺธา วิจาโร จ ฯเปฯ เทฺว ขนฺเธ ฯเปฯ. ปฏิสนฺธิ\-กฺขเณ ฯเปฯ. (4)}

\proseitem{61}{อวิตกฺก\-อวิจารํ ธมฺมํ สํสฏฺโฐ อวิตกฺก\-อวิจาโร ธมฺโม อุปฺปชฺชติ เหตุปจฺจยา. อวิตกฺก\-อวิจารํ เอกํ ขนฺธํ สํสฏฺฐา ตโย ขนฺธา ฯเปฯ เทฺว ขนฺเธ สํสฏฺฐา ฯเปฯ. ปฏิสนฺธิ\-กฺขเณ ฯเปฯ. (1)}

\prose{\csromanfolio{37}อวิตกฺก\-อวิจารํ ธมฺมํ สํสฏฺโฐ อวิตกฺก\-วิจารมตฺโต ธมฺโม ฯเปฯ. วิจารํ สํสฏฺฐา อวิตกฺก\-วิจารมตฺตา ขนฺธา. ปฏิสนฺธิ\-กฺขเณ วิจารํ สํสฏฺฐา ฯเปฯ. (2)}

\prose{อวิตกฺกวิจารมตฺตญฺจ อวิตกฺกอวิจารญฺจ ธมฺมํ สํสฏฺโฐ อวิตกฺก\-วิจารมตฺโต ธมฺโม ฯเปฯ. อวิตกฺก\-วิจารมตฺตํ เอกํ ขนฺธญฺจ วิจารญฺจ สํสฏฺฐา ตโย ขนฺธา ฯเปฯ เทฺว ขนฺเธ ฯเปฯ. ปฏิสนฺธิ\-กฺขเณ ฯเปฯ. (1)}

\prose{สวิตกฺกสวิจารญฺจ อวิตกฺกวิจารมตฺตญฺจ ธมฺมํ สํสฏฺโฐ สวิตกฺก\-สวิจาโร ธมฺโม อุปฺปชฺชติ เหตุปจฺจยา. สวิตกฺก\-สวิจารํ เอกํ ขนฺธญฺจ วิตกฺกญฺจ สํสฏฺฐา ตโย ขนฺธา ฯเปฯ เทฺว ขนฺเธ จ วิตกฺกญฺจ ฯเปฯ. ปฏิสนฺธิ\-กฺขเณ ฯเปฯ. (1) (เหตุปจฺจยํ อนุมชฺชนฺเตน สพฺเพ ปจฺจยา วิตฺถาเรตพฺพา.)}

\proseitemwithcloserruled{62}{เหตุยา เอกาทส, อารมฺมเณ อธิปติยา อนนฺตเร สมนนฺตเร สหชาเต อญฺญมญฺเญ นิสฺสเย อุปนิสฺสเย ปุเรชาเต อาเสวเน กมฺเม วิปาเก อาหาเร อินฺทฺริเย ฌาเน มคฺเค สมฺปยุตฺเต วิปฺปยุตฺเต อตฺถิยา นตฺถิยา วิคเต อวิคเต สพฺพตฺถ เอกาทส.}{อนุโลมํ.}
\csromanheader{2}{2. ปจฺจย\-ปจฺจนีย}
\vspace{\tipitakaparskip}
\csromancenter{(ปจฺจนียํ กาตพฺพํ อสมฺโมหนฺเตน.)}
\proseitemwithcloserruled{63}{นเหตุยา ฉ, นอธิปติยา เอกาทส, นปุเรชาเต เอกาทส, นปจฺฉาชาเต เอกาทส, นอาเสวเน เอกาทส, นกมฺเม สตฺต, นวิปาเก เอกาทส, นฌาเน เอกํ, นมคฺเค ฉ, นวิปฺปยุตฺเต เอกาทส.}{ปจฺจนียํ.}
\csromancenter{\textbf{เหตุ}ปจฺจยา นอธิปติยา เอกาทส ฯเปฯ นวิปฺปยุตฺเต เอกาทส. (สํขิตฺตํ.)}
\nitthitamruled{อนุโลม\-ปจฺจนียํ.}
\proseitem{65}{\csromanfolio{38}นเหตุปจฺจยา อารมฺมเณ ฉ ฯเปฯ ปุเรชาเต ฉ, อาเสวเน ปญฺจ, กมฺเม ฉ ฯเปฯ ฌาเน ฉ, มคฺเค ตีณิ, สมฺปยุตฺเต ฉ ฯเปฯ อวิคเต ฉ.}

\vspace{\tipitakaparskip}
\csromancenter{(สมฺปยุตฺตวาโรปิ วิตฺถาเรตพฺโพ.)}
\csromansectionrule
\csromanmarkh{6. วิตกฺกตฺติก}\csromanmarkhii{7. ปญฺหาวาร}\csromanheadernomark{2}{6. \textbf{วิตกฺกตฺติก 7. ปญฺหาวาร}}
\csromanmarkh{1. ปจฺจยานุโลม}\csromanmarkhii{1. วิภงฺควาร}\csromanheadernomark{2}{1. \textbf{ปจฺจยานุโลม 1. วิภงฺควาร}}
\vspace{\tipitakaparskip}
\csromancenter{\textbf{เหตุ}}
\proseitem{66}{สวิตกฺก\-สวิจาโร ธมฺโม สวิตกฺก\-สวิจารสฺส ธมฺมสฺส เหตุปจฺจเยน ปจฺจโย. สวิตกฺก\-สวิจารา เหตู สมฺปยุตฺตกานํ ขนฺธานํ เหตุปจฺจเยน ปจฺจโย. ปฏิสนฺธิ\-กฺขเณ ฯเปฯ. (1)}

\prose{สวิตกฺก\-สวิจาโร ธมฺโม อวิตกฺก\-วิจารมตฺตสฺส ธมฺมสฺส เหตุปจฺจเยน ปจฺจโย. สวิตกฺก\-สวิจารา เหตู วิตกฺกสฺส เหตุปจฺจเยน ปจฺจโย. ปฏิสนฺธิ\-กฺขเณ ฯเปฯ. (2)}

\prose{สวิตกฺก\-สวิจาโร ธมฺโม อวิตกฺก\-อวิจารสฺส ธมฺมสฺส เหตุปจฺจเยน ปจฺจโย. สวิตกฺก\-สวิจารา เหตู จิตฺต\-สมุฏฺฐานานํ รูปานํ เหตุปจฺจเยน ปจฺจโย. ปฏิสนฺธิ\-กฺขเณ สวิตกฺก\-สวิจารา เหตู กฏตฺตารูปานํ เหตุปจฺจเยน ปจฺจโย. (3)}

\prose{\csromanfolio{39}สวิตกฺก\-สวิจาโร ธมฺโม สวิตกฺก\-สวิจารสฺส จ อวิตกฺก\-อวิจารสฺส จ ธมฺมสฺส เหตุปจฺจเยน ปจฺจโย. สวิตกฺก\-สวิจารา เหตู สมฺปยุตฺตกานํ ขนฺธานํ จิตฺตสมุฏฺฐานานญฺจ รูปานํ เหตุปจฺจเยน ปจฺจโย. ปฏิสนฺธิ\-กฺขเณ สวิตกฺก\-สวิจารา เหตู สมฺปยุตฺตา\-กานํ ขนฺธานํ กฏตฺตา จ รูปานํ เหตุปจฺจเยน ปจฺจโย. (4)}

\prose{สวิตกฺก\-สวิจาโร ธมฺโม อวิตกฺก\-วิจารมตฺตสฺส จ อวิตกฺก\-อวิจารสฺส จ ธมฺมสฺส เหตุปจฺจเยน ปจฺจโย. สวิตกฺก\-สวิจารา เหตู วิตกฺกสฺส จิตฺตสมุฏฺฐานานญฺจ รูปานํ เหตุปจฺจเยน ปจฺจโย. ปฏิสนฺธิ\-กฺขเณ สวิตกฺก\-สวิจารา เหตู วิตกฺกสฺส กฏตฺตา จ รูปานํ เหตุปจฺจเยน ปจฺจโย. (5)}

\prose{สวิตกฺก\-สวิจาโร ธมฺโม สวิตกฺก\-สวิจารสฺส จ อวิตกฺก\-วิจารมตฺตสฺส จ ธมฺมสฺส เหตุปจฺจเยน ปจฺจโย. สวิตกฺก\-สวิจารา เหตู สมฺปยุตฺตกานํ ขนฺธานํ วิตกฺกสฺส จ เหตุปจฺจเยน ปจฺจโย. ปฏิสนฺธิ\-กฺขเณ สวิตกฺก\-สวิจารา เหตู สมฺปยุตฺตกานํ ขนฺธานํ วิตกฺกสฺส จ เหตุปจฺจเยน ปจฺจโย. (6)}

\prose{สวิตกฺก\-สวิจาโร ธมฺโม สวิตกฺก\-สวิจารสฺส จ อวิตกฺก\-วิจารมตฺตสฺส จ อวิตกฺก\-อวิจารสฺส จ ธมฺมสฺส เหตุปจฺจเยน ปจฺจโย. สวิตกฺก\-สวิจารา เหตู สมฺปยุตฺตกานํ ขนฺธานํ วิตกฺกสฺส จ จิตฺตสมุฏฺฐานานญฺจ รูปานํ เหตุปจฺจเยน ปจฺจโย. ปฏิสนฺธิ\-กฺขเณ สวิตกฺก\-สวิจารา เหตู สมฺปยุตฺตกานํ ขนฺธานํ วิตกฺกสฺส จ กฏตฺตา จ รูปานํ เหตุปจฺจเยน ปจฺจโย. (7)}

\proseitem{67}{อวิตกฺก\-วิจารมตฺโต ธมฺโม อวิตกฺก\-วิจารมตฺตสฺส ธมฺมสฺส เหตุปจฺจเยน ปจฺจโย. อวิตกฺก\-วิจารมตฺตา เหตู สมฺปยุตฺตกานํ ขนฺธานํ เหตุปจฺจเยน ปจฺจโย. ปฏิสนฺธิ\-กฺขเณ อวิตกฺก\-วิจารมตฺตา เหตู สมฺปยุตฺตกานํ ขนฺธานํ เหตุปจฺจเยน ปจฺจโย. (1)}

\prose{อวิตกฺก\-วิจารมตฺโต ธมฺโม อวิตกฺก\-อวิจารสฺส ธมฺมสฺส เหตุปจฺจเยน ปจฺจโย. อวิตกฺก\-วิจารมตฺตา เหตู วิจารสฺส จิตฺตสมุฏฺฐานานญฺจ รูปานํ เหตุปจฺจเยน ปจฺจโย. ปฏิสนฺธิ\-กฺขเณ อวิตกฺก\-วิจารมตฺตา เหตู วิจารสฺส จ กฏตฺตา จ รูปานํ เหตุปจฺจเยน ปจฺจโย. (2)}

\prose{\csromanfolio{40}อวิตกฺก\-วิจารมตฺโต ธมฺโม อวิตกฺก\-วิจารมตฺตสฺส จ อวิตกฺก\-อวิจารสฺส จ ธมฺมสฺส เหตุปจฺจเยน ปจฺจโย. อวิตกฺก\-วิจารมตฺตา เหตู สมฺปยุตฺตกานํ ขนฺธานํ วิจารสฺส จ จิตฺตสมุฏฺฐานานญฺจ รูปานํ เหตุปจฺจเยน ปจฺจโย. ปฏิสนฺธิ\-กฺขเณ อวิตกฺก\-วิจารมตฺตา เหตู สมฺปยุตฺตกานํ ขนฺธานํ วิจารสฺส จ กฏตฺตา จ รูปานํ เหตุปจฺจเยน ปจฺจโย. (3)}

\proseitem{68}{อวิตกฺก\-อวิจาโร ธมฺโม อวิตกฺก\-อวิจารสฺส ธมฺมสฺส เหตุปจฺจเยน ปจฺจโย. อวิตกฺก\-อวิจารา เหตู สมฺปยุตฺตกานํ ขนฺธานํ จิตฺตสมุฏฺฐานานญฺจ รูปานํ เหตุปจฺจเยน ปจฺจโย. ปฏิสนฺธิ\-กฺขเณ อวิตกฺก\-อวิจารา เหตู สมฺปยุตฺตกานํ ขนฺธานํ กฏตฺตา จ รูปานํ เหตุปจฺจเยน ปจฺจโย. (1)}

\csromanheader{2}{\textbf{อารมฺมณ}}
\proseitem{69}{สวิตกฺก\-สวิจาโร ธมฺโม สวิตกฺก\-สวิจารสฺส ธมฺมสฺส อารมฺมณ\-ปจฺจเยน ปจฺจโย. ทานํ ทตฺวา, สีลํ สมาทิยิตฺวา, อุโปสถกมฺมํ กตฺวา ตํ ปจฺจเวกฺขติ, ปุพฺเพ สุจิณฺณานิ ปจฺจเวกฺขติ, สวิตกฺก\-สวิจารา ฌานา วุฏฺฐหิตฺวา, มคฺคา วุฏฺฐหิตฺวา, ผลา วุฏฺฐหิตฺวา ผลํ ปจฺจเวกฺขติ. อริยา ปหีเน กิเลเส ปจฺจเวกฺขนฺติ, วิกฺขมฺภิเต กิเลเส ปจฺจเวกฺขนฺติ, ปุพฺเพ สมุทาจิณฺเณ กิเลเส ชานนฺติ. สวิตกฺก\-สวิจาเร ขนฺเธ อนิจฺจโต ทุกฺขโต อนตฺตโต วิปสฺสนฺติ, อสฺสาเทนฺติ อภินนฺทนฺติ, ตํ อารพฺภ ราโค อุปฺปชฺชติ ฯเปฯ โทมนสฺสํ อุปฺปชฺชติ. สวิตกฺก\-สวิจาเร ขนฺเธ อารพฺภ สวิตกฺก\-สวิจารา ขนฺธา อุปฺปชฺชนฺติ. (1)}

\prose{สวิตกฺก\-สวิจาโร ธมฺโม อวิตกฺก\-วิจารมตฺตสฺส ธมฺมสฺส อารมฺมณ\-ปจฺจเยน ปจฺจโย. ทานํ ทตฺวา, สีลํ สมาทิยิตฺวา, อุโปสถกมฺมํ กตฺวา ตํ ปจฺจเวกฺขติ, ตํ อารพฺภ วิตกฺโก อุปฺปชฺชติ. ปุพฺเพ สุจิณฺณานิ ปจฺจเวกฺขติ, สวิตกฺก\-สวิจารา ฌานา วุฏฺฐหิตฺวา, มคฺคา วุฏฺฐหิตฺวา, ผลา วุฏฺฐหิตฺวา ผลํ ปจฺจเวกฺขติ, ตํ อารพฺภ วิตกฺโก อุปฺปชฺชติ. อริยา ปหีเน กิเลเส ปจฺจเวกฺขนฺติ, วิกฺขมฺภิเต กิเลเส ปจฺจเวกฺขนฺติ, ปุพฺเพ สมุทาจิณฺเณ กิเลเส ชานนฺติ. สวิตกฺก\-สวิจาเร ขนฺเธ อนิจฺจโต \csromanfolio{41}ทุกฺขโต อนตฺตโต วิปสฺสนฺติ, อสฺสาเทนฺติ อภินนฺทนฺติ, ตํ อารพฺภ วิตกฺโก อุปฺปชฺชติ. สวิตกฺก\-สวิจาเร ขนฺเธ อารพฺภ วิตกฺโก อุปฺปชฺชติ. (2)}

\prose{สวิตกฺก\-สวิจาโร ธมฺโม อวิตกฺก\-อวิจารสฺส ธมฺมสฺส อารมฺมณ\-ปจฺจเยน ปจฺจโย. เจโตปริย\-ญาเณน สวิตกฺกสวิจาร\-จิตฺต\-สมงฺคิสฺส จิตฺตํ ชานาติ. สวิตกฺก\-สวิจารา ขนฺธา เจโตปริย\-ญาณสฺส, ปุพฺเพ\-นิวาสา\-นุสฺสติ\-ญาณสฺส, ยถากมฺมูปคญาณสฺส, อนาคตํสญาณสฺส อารมฺมณ\-ปจฺจเยน ปจฺจโย. สวิตกฺก\-สวิจาเร ขนฺเธ อารพฺภ อวิตกฺก\-อวิจารา ขนฺธา อุปฺปชฺชนฺติ. (3)}

\prose{สวิตกฺก\-สวิจาโร ธมฺโม สวิตกฺก\-สวิจารสฺส จ อวิตกฺก\-วิจารมตฺตสฺส จ ธมฺมสฺส อารมฺมณ\-ปจฺจเยน ปจฺจโย. ทานํ ทตฺวา, สีลํ สมาทิยิตฺวา, อุโปสถกมฺมํ กตฺวา ตํ ปจฺจเวกฺขติ, ตํ อารพฺภ สวิตกฺก\-สวิจารา ขนฺธา จ วิตกฺโก จ อุปฺปชฺชนฺติ. ปุพฺเพ สุจิณฺณานิ ปจฺจเวกฺขติ, สวิตกฺก\-สวิจารา ฌานา วุฏฺฐหิตฺวา, มคฺคา วุฏฺฐหิตฺวา, ผลา วุฏฺฐหิตฺวา ผลํ ปจฺจเวกฺขติ, ตํ อารพฺภ สวิตกฺก\-สวิจารา ขนฺธา จ วิตกฺโก จ อุปฺปชฺชนฺติ. อริยา ปหีเน กิเลเส ปจฺจเวกฺขนฺติ, วิกฺขมฺภิเต กิเลเส ปจฺจเวกฺขนฺติ, ปุพฺเพ สมุทาจิณฺเณ กิเลเส ชานนฺติ. สวิตกฺก\-สวิจาเร ขนฺเธ อนิจฺจโต ทุกฺขโต อนตฺตโต วิปสฺสนฺติ, อสฺสาเทนฺติ อภินนฺทนฺติ, ตํ อารพฺภ สวิตกฺก\-สวิจารา ขนฺธา จ วิตกฺโก จ อุปฺปชฺชนฺติ. สวิตกฺก\-สวิจาเร ขนฺเธ อารพฺภ สวิตกฺก\-สวิจารา ขนฺธา จ วิตกฺโก จ อุปฺปชฺชนฺติ. (4)}

\proseitem{70}{อวิตกฺก\-วิจารมตฺโต ธมฺโม อวิตกฺก\-วิจารมตฺตสฺส ธมฺมสฺส อารมฺมณ\-ปจฺจเยน ปจฺจโย. อวิตกฺก\-วิจารมตฺตา ฌานา วุฏฺฐหิตฺวา, มคฺคา วุฏฺฐหิตฺวา, ผลา วุฏฺฐหิตฺวา ผลํ ปจฺจเวกฺขติ, ตํ อารพฺภ วิตกฺโก อุปฺปชฺชติ. อวิตกฺก\-วิจารมตฺเต ขนฺเธ จ วิตกฺกญฺจ อนิจฺจโต ทุกฺขโต อนตฺตโต วิปสฺสติ, อสฺสาเทติ อภินนฺทติ, ตํ อารพฺภ วิตกฺโก อุปฺปชฺชติ. อวิตกฺก\-วิจารมตฺเต ขนฺเธ จ วิตกฺกญฺจ อารพฺภ วิตกฺโก อุปฺปชฺชติ. (1)}

\prose{อวิตกฺก\-วิจารมตฺโต ธมฺโม สวิตกฺก\-สวิจารสฺส ธมฺมสฺส อารมฺมณ\-ปจฺจเยน ปจฺจโย. อวิตกฺก\-วิจารมตฺตา ฌานา วุฏฺฐหิตฺวา, มคฺคา วุฏฺฐหิตฺวา, ผลา วุฏฺฐหิตฺวา ผลํ ปจฺจเวกฺขนฺติ, ตํ อารพฺภ สวิตกฺก\-สวิจารา ขนฺธา อุปฺปชฺชนฺติ. อวิตกฺก\-วิจารมตฺเต ขนฺเธ จ วิตกฺกญฺจ อนิจฺจโต ทุกฺขโต \csromanfolio{42}อนตฺตโต วิปสฺสติ, อสฺสาเทติ อภินนฺทติ, ตํ อารพฺภ ราโค อุปฺปชฺชติ ฯเปฯ โทมนสฺสํ อุปฺปชฺชติ. อวิตกฺก\-วิจารมตฺเต ขนฺเธ จ วิตกฺกญฺจ อารพฺภ สวิตกฺก\-สวิจารา ขนฺธา อุปฺปชฺชนฺติ. (2)}

\prose{อวิตกฺก\-วิจารมตฺโต ธมฺโม อวิตกฺก\-อวิจารสฺส ธมฺมสฺส อารมฺมณ\-ปจฺจเยน ปจฺจโย. เจโตปริย\-ญาเณน อวิตกฺกวิจารมตฺต\-จิตฺต\-สมงฺคิสฺส จิตฺตํ ชานาติ. อวิตกฺก\-วิจารมตฺตา ขนฺธา เจโตปริย\-ญาณสฺส, ปุพฺเพ\-นิวาสา\-นุสฺสติ\-ญาณสฺส, ยถากมฺมูปคญาณสฺส, อนาคตํสญาณสฺส อารมฺมณ\-ปจฺจเยน ปจฺจโย. อวิตกฺก\-วิจารมตฺเต ขนฺเธ จ วิตกฺกญฺจ อารพฺภ อวิตกฺก\-อวิจารา ขนฺธา อุปฺปชฺชนฺติ. (3)}

\prose{อวิตกฺก\-วิจารมตฺโต ธมฺโม สวิตกฺก\-สวิจารสฺส จ อวิตกฺก\-วิจารมตฺตสฺส จ ธมฺมสฺส อารมฺมณ\-ปจฺจเยน ปจฺจโย. อวิตกฺก\-วิจารมตฺตา ฌานา วุฏฺฐหิตฺวา, มคฺคา วุฏฺฐหิตฺวา, ผลา วุฏฺฐหิตฺวา ผลํ ปจฺจเวกฺขติ, ตํ อารพฺภ สวิตกฺก\-สวิจารา ขนฺธา จ วิตกฺโก จ อุปฺปชฺชนฺติ. อวิตกฺก\-วิจารมตฺเต ขนฺเธ จ วิตกฺกญฺจ อนิจฺจโต ทุกฺขโต อนตฺตโต วิปสฺสติ, อสฺสาเทติ อภินนฺทติ, ตํ อารพฺภ สวิตกฺก\-สวิจารา ขนฺธา จ วิตกฺโก จ อุปฺปชฺชนฺติ. อวิตกฺก\-วิจารมตฺเต ขนฺเธ จ วิตกฺกญฺจ อารพฺภ สวิตกฺก\-สวิจารา ขนฺธา จ วิตกฺโก จ อุปฺปชฺชนฺติ. (4)}

\proseitem{71}{อวิตกฺก\-อวิจาโร ธมฺโม อวิตกฺก\-อวิจารสฺส ธมฺมสฺส อารมฺมณ\-ปจฺจเยน ปจฺจโย. นิพฺพานํ อวิตกฺก\-อวิจารสฺส มคฺคสฺส ผลสฺส วิจารสฺส จ อารมฺมณ\-ปจฺจเยน ปจฺจโย. ทิพฺเพน จกฺขุนา รูปํ ปสฺสติ. ทิพฺพาย โสตธาตุยา สทฺทํ สุณาติ. เจโตปริย\-ญาเณน อวิตกฺกอวิจาร\-จิตฺต\-สมงฺคิสฺส จิตฺตํ ชานาติ.  อากาสา\-นญฺจา\-ยตนํ วิญฺญาณญฺจา\-ยตนสฺส ฯเปฯ. อากิญฺจญฺญา\-ยตนํ เนว\-สญฺญา\-นา\-สญฺญา\-ยตนสฺส อารมฺมณ\-ปจฺจเยน ปจฺจโย. รูปายตนํ จกฺขุ\-วิญฺญาณสฺส ฯเปฯ. โผฏฺฐพฺพา\-ยตนํ กายวิญฺญาณสฺส อารมฺมณ\-ปจฺจเยน ปจฺจโย. อวิตกฺก\-อวิจารา ขนฺธา อิทฺธิวิธ\-ญาณสฺส, เจโตปริย\-ญาณสฺส, ปุพฺเพ\-นิวาสา\-นุสฺสติ\-ญาณสฺส, ยถากมฺมูปคญาณสฺส, อนาคตํสญาณสฺส อารมฺมณ\-ปจฺจเยน ปจฺจโย. อวิตกฺก\-อวิจาเร ขนฺเธ จ วิจารญฺจ อารพฺภ อวิตกฺก\-อวิจารา ขนฺธา อุปฺปชฺชนฺติ. (1)}

\prose{\csromanfolio{43}อวิตกฺก\-อวิจาโร ธมฺโม สวิตกฺก\-สวิจารสฺส ธมฺมสฺส อารมฺมณ\-ปจฺจเยน ปจฺจโย. อริยา อวิตกฺก\-อวิจารา ฌานา วุฏฺฐหิตฺวา, มคฺคา วุฏฺฐหิตฺวา, ผลา วุฏฺฐหิตฺวา ผลํ ปจฺจเวกฺขนฺติ, ตํ อารพฺภ สวิตกฺก\-สวิจารา ขนฺธา อุปฺปชฺชนฺติ. อริยา นิพฺพานํ ปจฺจเวกฺขนฺติ, นิพฺพานํ โคตฺรภุสฺส โวทานสฺส สวิตกฺก\-สวิจารสฺส มคฺคสฺส ผลสฺส อาวชฺชนาย อารมฺมณ\-ปจฺจเยน ปจฺจโย. จกฺขุํ อนิจฺจโต ทุกฺขโต อนตฺตโต วิปสฺสติ, อสฺสาเทติ อภินนฺทติ, ตํ อารพฺภ ราโค อุปฺปชฺชติ ฯเปฯ โทมนสฺสํ อุปฺปชฺชติ. โสตํ. ฆานํ. ชิวฺหํ, กายํ. รูเป. สทฺเท. คนฺเธ. รเส. โผฏฺฐพฺเพ. วตฺถุํ. อวิตกฺก\-อวิจาเร ขนฺเธ จ วิจารญฺจ อนิจฺจโต ทุกฺขโต อนตฺตโต วิปสฺสติ, อสฺสาเทติ อภินนฺทติ, ตํ อารพฺภ ราโค อุปฺปชฺชติ ฯเปฯ โทมนสฺสํ อุปฺปชฺชติ. อวิตกฺก\-อวิจาเร ขนฺเธ จ วิจารญฺจ อารพฺภ สวิตกฺก\-สวิจารา ขนฺธา อุปฺปชฺชนฺติ. (2)}

\prose{อวิตกฺก\-อวิจาโร ธมฺโม อวิตกฺก\-วิจารมตฺตสฺส ธมฺมสฺส อารมฺมณ\-ปจฺจเยน ปจฺจโย. อริยา อวิตกฺก\-อวิจารา ฌานา วุฏฺฐหิตฺวา, มคฺคา วุฏฺฐหิตฺวา, ผลา วุฏฺฐหิตฺวา ผลํ ปจฺจเวกฺขนฺติ, ตํ อารพฺภ วิตกฺโก อุปฺปชฺชติ. อริยา นิพฺพานํ ปจฺจเวกฺขนฺติ, นิพฺพานํ อวิตกฺก\-วิจารมตฺตสฺส มคฺคสฺส ผลสฺส วิตกฺกสฺส จ อารมฺมณ\-ปจฺจเยน ปจฺจโย. จกฺขุํ อนิจฺจโต ทุกฺขโต อนตฺตโต ฯเปฯ. วตฺถุํ. อวิตกฺก\-อวิจาเร ขนฺเธ จ วิจารญฺจ อนิจฺจโต ทุกฺขโต อนตฺตโต วิปสฺสติ, อสฺสาเทติ อภินนฺทติ, ตํ อารพฺภ วิตกฺโก อุปฺปชฺชติ. อวิตกฺก\-อวิจาเร ขนฺเธ จ วิจารญฺจ อารพฺภ วิตกฺโก อุปฺปชฺชติ. (3)}

\prose{อวิตกฺก\-อวิจาโร ธมฺโม อวิตกฺก\-วิจารมตฺตสฺส จ อวิตกฺก\-อวิจารสฺส จ ธมฺมสฺส อารมฺมณ\-ปจฺจเยน ปจฺจโย. นิพฺพานํ อวิตกฺก\-วิจารมตฺตสฺส มคฺคสฺส ผลสฺส วิจารสฺส จ อารมฺมณ\-ปจฺจเยน ปจฺจโย. (4)}

\prose{อวิตกฺก\-อวิจาโร ธมฺโม สวิตกฺก\-สวิจารสฺส จ อวิตกฺก\-วิจารมตฺตสฺส จ ธมฺมสฺส อารมฺมณ\-ปจฺจเยน ปจฺจโย. อริยา อวิตกฺก\-อวิจารา ฌานา วุฏฺฐหิตฺวา, มคฺคา วุฏฺฐหิตฺวา, ผลา วุฏฺฐหิตฺวา ผลํ ปจฺจเวกฺขนฺติ, ตํ อารพฺภ สวิตกฺก\-สวิจารา ขนฺธา จ วิตกฺโก จ อุปฺปชฺชนฺติ. อริยา นิพฺพานํ ปจฺจเวกฺขนฺติ, นิพฺพานํ โคตฺรภุสฺส วิตกฺกสฺส จ, โวทานสฺส วิตกฺกสฺส จ, สวิตกฺก\-สวิจารสฺส มคฺคสฺส วิตกฺกสฺส จ, สวิตกฺก\-สวิจารสฺส ผลสฺส วิตกฺกสฺส จ, อาวชฺชนาย วิตกฺกสฺส จ \csromanfolio{44}อารมฺมณ\-ปจฺจเยน ปจฺจโย. จกฺขุํ อนิจฺจโต ทุกฺขโต อนตฺตโต วิปสฺสติ, อสฺสาเทติ อภินนฺทติ, ตํ อารพฺภ สวิตกฺก\-สวิจารา ขนฺธา จ วิตกฺโก จ อุปฺปชฺชนฺติ. โสตํ ฯเปฯ. โผฏฺฐพฺพํ. วตฺถุํ. อวิตกฺก\-อวิจาเร ขนฺเธ จ วิจารญฺจ อนิจฺจโต ทุกฺขโต อนตฺตโต วิปสฺสติ, อสฺสาเทติ อภินนฺทติ, ตํ อารพฺภ สวิตกฺก\-สวิจารา ขนฺธา จ วิตกฺโก จ อุปฺปชฺชนฺติ. (5)}

\proseitem{72}{อวิตกฺก\-วิจารมตฺโต จ อวิตกฺก\-อวิจาโร จ ธมฺมา สวิตกฺก\-สวิจารสฺส ธมฺมสฺส อารมฺมณ\-ปจฺจเยน ปจฺจโย. อวิตกฺก\-วิจารมตฺเต ขนฺเธ จ วิจารญฺจ อารพฺภ สวิตกฺก\-สวิจารา ขนฺธา อุปฺปชฺชนฺติ. (1)}

\prose{อวิตกฺก\-วิจารมตฺโต จ อวิตกฺก\-อวิจาโร จ ธมฺมา อวิตกฺก\-วิจารมตฺตสฺส ธมฺมสฺส อารมฺมณ\-ปจฺจเยน ปจฺจโย. อวิตกฺก\-วิจารมตฺเต ขนฺเธ จ วิจารญฺจ อารพฺภ วิตกฺโก อุปฺปชฺชติ. (2)}

\prose{อวิตกฺก\-วิจารมตฺโต จ อวิตกฺก\-อวิจาโร จ ธมฺมา อวิตกฺก\-อวิจารสฺส ธมฺมสฺส อารมฺมณ\-ปจฺจเยน ปจฺจโย. อวิตกฺก\-วิจารมตฺตา ขนฺธา จ วิจาโร จ เจโตปริย\-ญาณสฺส, ปุพฺเพ\-นิวาสา\-นุสฺสติ\-ญาณสฺส, ยถากมฺมูปคญาณสฺส, อนาคตํสญาณสฺส อารมฺมณ\-ปจฺจเยน ปจฺจโย. อวิตกฺก\-วิจารมตฺเต ขนฺเธ จ วิจารญฺจ อารพฺภ อวิตกฺก\-อวิจารา ขนฺธา อุปฺปชฺชนฺติ. (3)}

\prose{อวิตกฺก\-วิจารมตฺโต จ อวิตกฺก\-อวิจาโร จ ธมฺมา สวิตกฺก\-สวิจารสฺส จ อวิตกฺก\-วิจารมตฺตสฺส จ ธมฺมสฺส อารมฺมณ\-ปจฺจเยน ปจฺจโย. อวิตกฺก\-วิจารมตฺเต ขนฺเธ จ วิจารญฺจ อารพฺภ สวิตกฺก\-สวิจารา ขนฺธา จ วิตกฺโก จ อุปฺปชฺชนฺติ. (4)}

\proseitem{73}{สวิตกฺก\-สวิจาโร จ อวิตกฺก\-วิจารมตฺโต จ ธมฺมา สวิตกฺก\-สวิจารสฺส ธมฺมสฺส อารมฺมณ\-ปจฺจเยน ปจฺจโย. สวิตกฺก\-สวิจาเร ขนฺเธ จ วิตกฺกญฺจ อารพฺภ สวิตกฺก\-สวิจารา ขนฺธา อุปฺปชฺชนฺติ. (1)}

\prose{สวิตกฺก\-สวิจาโร จ อวิตกฺก\-วิจารมตฺโต จ ธมฺมา อวิตกฺก\-วิจารมตฺตสฺส อารมฺมณ\-ปจฺจเยน ปจฺจโย. สวิตกฺก\-สวิจาเร ขนฺเธ จ วิตกฺกญฺจ อารพฺภ วิตกฺโก อุปฺปชฺชติ. (2)}

\prose{\csromanfolio{45}สวิตกฺก\-สวิจาโร จ อวิตกฺก\-วิจารมตฺโต จ ธมฺมา อวิตกฺก\-อวิจารสฺส ธมฺมสฺส อารมฺมณ\-ปจฺจเยน ปจฺจโย. สวิตกฺก\-สวิจารา ขนฺธา จ วิตกฺโก จ เจโตปริย\-ญาณสฺส,ปุพฺเพ\-นิวาสา\-นุสฺสติ\-ญาณสฺส, ยถากมฺมูปคญาณสฺส, อนาคตํสญาณสฺส อารมฺมณ\-ปจฺจเยน ปจฺจโย. สวิตกฺก\-สวิจาเร ขนฺเธ จ วิตกฺกญฺจ อารพฺภ อวิตกฺก\-อวิจารา ขนฺธา อุปฺปชฺชนฺติ. (3)}

\prose{สวิตกฺก\-สวิจาโร จ อวิตกฺก\-วิจารมตฺโต จ ธมฺมา สวิตกฺก\-สวิจารสฺส จ อวิตกฺก\-วิจารมตฺตสฺส จ ธมฺมสฺส อารมฺมณ\-ปจฺจเยน ปจฺจโย. สวิตกฺก\-สวิจาเร ขนฺเธ จ วิตกฺกญฺจ อารพฺภ สวิตกฺก\-สวิจารา ขนฺธา จ วิตกฺโก จ อุปฺปชฺชนฺติ. (4)}

\proseitem{74}{สวิตกฺก\-สวิจาโร ธมฺโม สวิตกฺก\-สวิจารสฺส ธมฺมสฺส อธิปติ\-ปจฺจเยน ปจฺจโย. อารมฺมณา\-ธิปติ, สหชาตา\-ธิปติ.  อารมฺมณา\-ธิปติ\csromandash{}ทานํ ทตฺวา, สีลํ สมาทิยิตฺวา, อุโปสถกมฺมํ กตฺวา ตํ ครุํ กตฺวา ปจฺจเวกฺขติ, ปุพฺเพ สุจิณฺณานิ ครุํ กตฺวา ปจฺจเวกฺขติ. สวิตกฺก\-สวิจารา ฌานา วุฏฺฐหิตฺวา, มคฺคา วุฏฺฐหิตฺวา, ผลา วุฏฺฐหิตฺวา ผลํ ครุํ กตฺวา ปจฺจเวกฺขติ. สวิตกฺก\-สวิจาเร ขนฺเธ ครุํ กตฺวา อสฺสาเทติ อภินนฺทติ, ตํ ครุํ กตฺวา ราโค อุปฺปชฺชติ, ทิฏฺฐิ อุปฺปชฺชติ.  สหชาตา\-ธิปติ\csromandash{}สวิตกฺก\-สวิจารา อธิปติ สมฺปยุตฺตกานํ ขนฺธานํ อธิปติ\-ปจฺจเยน ปจฺจโย. (1)}

\prose{สวิตกฺก\-สวิจาโร ธมฺโม อวิตกฺก\-วิจารมตฺตสฺส ธมฺมสฺส อธิปติ\-ปจฺจเยน ปจฺจโย. อารมฺมณา\-ธิปติ, สหชาตา\-ธิปติ.  อารมฺมณา\-ธิปติ\csromandash{}ทานํ ทตฺวา, สีลํ สมาทิยิตฺวา, อุโปสถกมฺมํ กตฺวา ตํ ครุํ กตฺวา ปจฺจเวกฺขติ, ตํ ครุํ กตฺวา วิตกฺโก อุปฺปชฺชติ. ปุพฺเพ สุจิณฺณานิ ครุํ กตฺวา ปจฺจเวกฺขติ. สวิตกฺก\-สวิจารา ฌานา วุฏฺฐหิตฺวา, มคฺคา วุฏฺฐหิตฺวา, ผลา วุฏฺฐหิตฺวา ผลํ ครุํ กตฺวา ปจฺจเวกฺขนฺติ, ตํ ครุํ กตฺวา วิตกฺโก อุปฺปชฺชติ. สวิตกฺก\-สวิจาเร ขนฺเธ ครุํ กตฺวา อสฺสาเทติ อภินนฺทติ, ตํ ครุํ กตฺวา วิตกฺโก อุปฺปชฺชติ.  สหชาตา\-ธิปติ\csromandash{}สวิตกฺก\-สวิจารา อธิปติ วิตกฺกสฺส อธิปติ\-ปจฺจเยน ปจฺจโย. (2)}

\prose{\csromanfolio{46}สวิตกฺก\-สวิจาโร ธมฺโม อวิตกฺก\-อวิจารสฺส ธมฺมสฺส อธิปติ\-ปจฺจเยน ปจฺจโย. สหชาตา\-ธิปติ\csromandash{}สวิตกฺก\-สวิจารา อธิปติ จิตฺต\-สมุฏฺฐานานํ รูปานํ อธิปติ\-ปจฺจเยน ปจฺจโย. (3)}

\prose{สวิตกฺก\-สวิจาโร ธมฺโม สวิตกฺก\-สวิจารสฺส จ อวิตกฺก\-อวิจารสฺส จ ธมฺมสฺส อธิปติ\-ปจฺจเยน ปจฺจโย. สหชาตา\-ธิปติ\csromandash{} สวิตกฺก\-สวิจารา อธิปติ สมฺปยุตฺตกานํ ขนฺธานํ จิตฺตสมุฏฺฐานานญฺจ รูปานํ อธิปติ\-ปจฺจเยน ปจฺจโย. (4)}

\prose{สวิตกฺก\-สวิจาโร ธมฺโม อวิตกฺก\-วิจารมตฺตสฺส จ อวิตกฺก\-อวิจารสฺส จ ธมฺมสฺส อธิปติ\-ปจฺจเยน ปจฺจโย. สหชาตา\-ธิปติ\csromandash{}สวิตกฺก\-สวิจารา อธิปติ วิตกฺกสฺส จ จิตฺตสมุฏฺฐานานญฺจ รูปานํ อธิปติ\-ปจฺจเยน ปจฺจโย. (5)}

\prose{สวิตกฺก\-สวิจาโร ธมฺโม สวิตกฺก\-สวิจารสฺส จ อวิตกฺก\-วิจารมตฺตสฺส จ ธมฺมสฺส อธิปติ\-ปจฺจเยน ปจฺจโย. อารมฺมณา\-ธิปติ, สหชาตา\-ธิปติ.  อารมฺมณา\-ธิปติ\csromandash{}ทานํ ทตฺวา, สีลํ สมาทิยิตฺวา, อุโปสถกมฺมํ กตฺวา ตํ ครุํ กตฺวา ปจฺจเวกฺขติ, ตํ ครุํ กตฺวา สวิตกฺก\-สวิจารา ขนฺธา จ วิตกฺโก จ อุปฺปชฺชนฺติ, ปุพฺเพ สุจิณฺณานิ ครุํ กตฺวา ปจฺจเวกฺขติ, สวิตกฺก\-สวิจารา ฌานา วุฏฺฐหิตฺวา, มคฺคา วุฏฺฐหิตฺวา, ผลา วุฏฺฐหิตฺวา ผลํ ครุํ กตฺวา ปจฺจเวกฺขติ, ตํ ครุํ กตฺวา สวิตกฺก\-สวิจารา ขนฺธา จ วิตกฺโก จ อุปฺปชฺชนฺติ. สวิตกฺก\-สวิจาเร ขนฺเธ ครุํ กตฺวา อสฺสาเทติ อภินนฺทติ, ตํ ครุํ กตฺวา สวิตกฺก\-สวิจารา ขนฺธา จ วิตกฺโก จ อุปฺปชฺชนฺติ.  สหชาตา\-ธิปติ\csromandash{}สวิตกฺก\-สวิจารา อธิปติ สมฺปยุตฺตกานํ ขนฺธานํ วิตกฺกสฺส จ อธิปติ\-ปจฺจเยน ปจฺจโย. (6)}

\prose{สวิตกฺก\-สวิจาโร ธมฺโม สวิตกฺก\-สวิจารสฺส จ อวิตกฺก\-วิจารมตฺตสฺส จ อวิตกฺก\-อวิจารสฺส จ ธมฺมสฺส อธิปติ\-ปจฺจเยน ปจฺจโย. สหชาตา\-ธิปติ\csromandash{}สวิตกฺก\-สวิจารา อธิปติ สมฺปยุตฺตกานํ ขนฺธานํ วิตกฺกสฺส จ จิตฺตสมุฏฺฐานานญฺจ รูปานํ อธิปติ\-ปจฺจเยน ปจฺจโย. (7)}

\proseitem{75}{อวิตกฺก\-วิจารมตฺโต ธมฺโม อวิตกฺก\-วิจารมตฺตสฺส ธมฺมสฺส อธิปติ\-ปจฺจเยน ปจฺจโย. อารมฺมณา\-ธิปติ, สหชาตา\-ธิปติ.  อารมฺมณา\-ธิปติ\csromandash{}อวิตกฺก\-วิจารมตฺตา ฌานา วุฏฺฐหิตฺวา, มคฺคา วุฏฺฐหิตฺวา, \csromanfolio{47}ผลา วุฏฺฐหิตฺวา ผลํ ครุํ กตฺวา ปจฺจเวกฺขติ, ตํ ครุํ กตฺวา วิตกฺโก อุปฺปชฺชติ. อวิตกฺก\-วิจารมตฺเต ขนฺเธ จ วิตกฺกญฺจ ครุํ กตฺวา อสฺสาเทติ อภินนฺทติ, ตํ ครุํ กตฺวา วิตกฺโก อุปฺปชฺชติ.  สหชาตา\-ธิปติ\csromandash{} อวิตกฺก\-วิจารมตฺตา อธิปติ สมฺปยุตฺตกานํ ขนฺธานํ อธิปติ\-ปจฺจเยน ปจฺจโย. (1)}

\prose{อวิตกฺก\-วิจารมตฺโต ธมฺโม สวิตกฺก\-สวิจารสฺส ธมฺมสฺส อธิปติ\-ปจฺจเยน ปจฺจโย. อารมฺมณา\-ธิปติ\csromandash{}อวิตกฺก\-วิจารมตฺตา ฌานา วุฏฺฐหิตฺวา, มคฺคา วุฏฺฐหิตฺวา, ผลา วุฏฺฐหิตฺวา ผลํ ครุํ กตฺวา ปจฺจเวกฺขติ, ตํ ครุํ กตฺวา สวิตกฺก\-สวิจารา ขนฺธา อุปฺปชฺชนฺติ. อวิตกฺก\-วิจารมตฺเต ขนฺเธ จ วิตกฺกญฺจ ครุํ กตฺวา อสฺสาเทติ อภินนฺทติ, ตํ ครุํ กตฺวา ราโค อุปฺปชฺชติ, ทิฏฺฐิ อุปฺปชฺชติ. (2)}

\prose{อวิตกฺก\-วิจารมตฺโต ธมฺโม อวิตกฺก\-อวิจารสฺส ธมฺมสฺส อธิปติ\-ปจฺจเยน ปจฺจโย. สหชาตา\-ธิปติ\csromandash{}อวิตกฺก\-วิจารมตฺตา อธิปติ วิจารสฺส จิตฺตสมุฏฺฐานานญฺจ รูปานํ อธิปติ\-ปจฺจเยน ปจฺจโย. (3)}

\prose{อวิตกฺก\-วิจารมตฺโต ธมฺโม อวิตกฺก\-วิจารมตฺตสฺส จ อวิตกฺก\-อวิจารสฺส จ ธมฺมสฺส อธิปติ\-ปจฺจเยน ปจฺจโย. สหชาตา\-ธิปติ\csromandash{}อวิตกฺก\-วิจารมตฺตา อธิปติ สมฺปยุตฺตกานํ ขนฺธานํ วิจารสฺส จ จิตฺตสมุฏฺฐานานญฺจ รูปานํ อธิปติ\-ปจฺจเยน ปจฺจโย. (4)}

\prose{อวิตกฺก\-วิจารมตฺโต ธมฺโม สวิตกฺก\-สวิจารสฺส จ อวิตกฺก\-วิจารมตฺตสฺส จ ธมฺมสฺส อธิปติ\-ปจฺจเยน ปจฺจโย. อารมฺมณา\-ธิปติ\csromandash{}อวิตกฺก\-วิจารมตฺตา ฌานา วุฏฺฐหิตฺวา, มคฺคา วุฏฺฐหิตฺวา, ผลา วุฏฺฐหิตฺวา ผลํ ครุํ กตฺวา ปจฺจเวกฺขติ, ตํ ครุํ กตฺวา สวิตกฺก\-สวิจารา ขนฺธา จ วิตกฺโก จ อุปฺปชฺชนฺติ. อวิตกฺก\-วิจารมตฺเต ขนฺเธ จ วิตกฺกญฺจ ครุํ กตฺวา อสฺสาเทติ อภินนฺทติ, ตํ ครุํ กตฺวา สวิตกฺก\-สวิจารา ขนฺธา จ วิตกฺโก จ อุปฺปชฺชนฺติ. (5)}

\proseitem{76}{อวิตกฺก\-อวิจาโร ธมฺโม อวิตกฺก\-อวิจารสฺส ธมฺมสฺส อธิปติ\-ปจฺจเยน ปจฺจโย. อารมฺมณา\-ธิปติ, สหชาตา\-ธิปติ.  อารมฺมณา\-ธิปติ\csromandash{}นิพฺพานํ อวิตกฺก\-อวิจารสฺส มคฺคสฺส ผลสฺส วิจารสฺส จ อธิปติ\-ปจฺจเยน ปจฺจโย. สหชาตา\-ธิปติ\csromandash{}อวิตกฺก\-อวิจารา \csromanfolio{48}อธิปติ สมฺปยุตฺตกานํ ขนฺธานํ จิตฺตสมุฏฺฐานานญฺจ รูปานํ อธิปติ\-ปจฺจเยน ปจฺจโย. (1)}

\prose{อวิตกฺก\-อวิจาโร ธมฺโม สวิตกฺก\-สวิจารสฺส ธมฺมสฺส อธิปติ\-ปจฺจเยน ปจฺจโย. อารมฺมณา\-ธิปติ\csromandash{}อริยา อวิตกฺก\-อวิจารา ฌานา วุฏฺฐหิตฺวา, มคฺคา วุฏฺฐหิตฺวา, ผลา วุฏฺฐหิตฺวา ผลํ ครุํ กตฺวา ปจฺจเวกฺขนฺติ, ตํ ครุํ กตฺวา สวิตกฺก\-สวิจารา ขนฺธา อุปฺปชฺชนฺติ. อริยา นิพฺพานํ ครุํ กตฺวา ปจฺจเวกฺขนฺติ, นิพฺพานํ โคตฺรภุสฺส โวทานสฺส สวิตกฺก\-สวิจารสฺส มคฺคสฺส ผลสฺส อธิปติ\-ปจฺจเยน ปจฺจโย. จกฺขุํ ครุํ กตฺวา อสฺสาเทติ อภินนฺทติ, ตํ ครุํ กตฺวา ราโค อุปฺปชฺชติ, ทิฏฺฐิ อุปฺปชฺชติ. โสตํ. ฆานํ. ชิวฺหํ. กายํ. รูเป. สทฺเท. คนฺเธ. รเส. โผฏฺฐพฺเพ. วตฺถุํ. อวิตกฺก\-อวิจาเร ขนฺเธ จ วิจารญฺจ ครุํ กตฺวา อสฺสาเทติ อภินนฺทติ, ตํ ครุํ กตฺวา ราโค อุปฺปชฺชติ, ทิฏฺฐิ อุปฺปชฺชติ. (2)}

\prose{อวิตกฺก\-อวิจาโร ธมฺโม อวิตกฺก\-วิจารมตฺตสฺส ธมฺมสฺส อธิปติ\-ปจฺจเยน ปจฺจโย. อารมฺมณา\-ธิปติ\csromandash{}อริยา อวิตกฺก\-อวิจารา ฌานา วุฏฺฐหิตฺวา, มคฺคา วุฏฺฐหิตฺวา, ผลา วุฏฺฐหิตฺวา ผลํ ครุํ กตฺวา ปจฺจเวกฺขนฺติ, ตํ ครุํ กตฺวา วิตกฺโก อุปฺปชฺชติ. อริยา นิพฺพานํ ครุํ กตฺวา ปจฺจเวกฺขนฺติ, นิพฺพานํ อวิตกฺก\-วิจารมตฺตสฺส มคฺคสฺส ผลสฺส วิตกฺกสฺส จ อธิปติ\-ปจฺจเยน ปจฺจโย. จกฺขุํ ฯเปฯ. วตฺถุํ. อวิตกฺก\-อวิจาเร ขนฺเธ จ วิจารญฺจ ครุํ กตฺวา อสฺสาเทติ อภินนฺทติ, ตํ ครุํ กตฺวา วิตกฺโก อุปฺปชฺชติ. (3)}

\prose{อวิตกฺก\-อวิจาโร ธมฺโม อวิตกฺก\-วิจารมตฺตสฺส จ อวิตกฺก\-อวิจารสฺส จ ธมฺมสฺส อธิปติ\-ปจฺจเยน ปจฺจโย. อารมฺมณา\-ธิปติ\csromandash{}นิพฺพานํ อวิตกฺก\-วิจารมตฺตสฺส มคฺคสฺส ผลสฺส วิจารสฺส จ อธิปติ\-ปจฺจเยน ปจฺจโย. (4)}

\prose{อวิตกฺก\-อวิจาโร ธมฺโม สวิตกฺก\-สวิจารสฺส จ อวิตกฺก\-วิจารมตฺตสฺส จ ธมฺมสฺส อธิปติ\-ปจฺจเยน ปจฺจโย. อารมฺมณา\-ธิปติ\csromandash{}อริยา อวิตกฺก\-อวิจารา ฌานา วุฏฺฐหิตฺวา, มคฺคา วุฏฺฐหิตฺวา, ผลา วุฏฺฐหิตฺวา ผลํ ครุํ กตฺวา ปจฺจเวกฺขนฺติ, ตํ ครุํ กตฺวา สวิตกฺก\-สวิจารา ขนฺธา จ วิตกฺโก จ อุปฺปชฺชนฺติ. อริยา นิพฺพานํ ครุํ กตฺวา ปจฺจเวกฺขนฺติ, นิพฺพานํ โคตฺรภุสฺส วิตกฺกสฺส จ, โวทานสฺส วิตกฺกสฺส จ, สวิตกฺก\-สวิจารสฺส \csromanfolio{49}มคฺคสฺส วิตกฺกสฺส จ, สวิตกฺก\-สวิจารสฺส ผลสฺส วิตกฺกสฺส จ อธิปติ\-ปจฺจเยน ปจฺจโย. จกฺขุํ ครุํ กตฺวา ฯเปฯ. วตฺถุํ. อวิตกฺก\-อวิจาเร ขนฺเธ จ วิจารญฺจ ครุํ กตฺวา อสฺสาเทติ อภินนฺทติ, ตํ ครุํ กตฺวา สวิตกฺก\-สวิจารา ขนฺธา จ วิตกฺโก จ อุปฺปชฺชนฺติ. (5)}

\proseitem{77}{อวิตกฺก\-วิจารมตฺโต จ อวิตกฺก\-อวิจาโร จ ธมฺมา สวิตกฺก\-สวิจารสฺส ธมฺมสฺส อธิปติ\-ปจฺจเยน ปจฺจโย. อารมฺมณา\-ธิปติ\csromandash{}อวิตกฺก\-วิจารมตฺเต ขนฺเธ จ วิจารญฺจ ครุํ กตฺวา สวิตกฺก\-สวิจารา ขนฺธา อุปฺปชฺชนฺติ. (1)}

\prose{อวิตกฺก\-วิจารมตฺโต จ อวิตกฺก\-อวิจาโร จ ธมฺมา อวิตกฺก\-วิจารมตฺตสฺส ธมฺมสฺส อธิปติ\-ปจฺจเยน ปจฺจโย. อารมฺมณา\-ธิปติ\csromandash{}อวิตกฺก\-วิจารมตฺเต ขนฺเธ จ วิจารญฺจ ครุํ กตฺวา วิตกฺโก อุปฺปชฺชติ. (2)}

\prose{อวิตกฺก\-วิจารมตฺโต จ อวิตกฺก\-อวิจาโร จ ธมฺมา สวิตกฺก\-สวิจารสฺส จ อวิตกฺก\-วิจารมตฺตสฺส จ ธมฺมสฺส อธิปติ\-ปจฺจเยน ปจฺจโย. อารมฺมณา\-ธิปติ\csromandash{}อวิตกฺก\-วิจารมตฺเต ขนฺเธ จ วิจารญฺจ ครุํ กตฺวา สวิตกฺก\-สวิจารา ขนฺธา จ วิตกฺโก จ อุปฺปชฺชนฺติ. (3)}

\proseitem{78}{สวิตกฺก\-สวิจาโร จ อวิตกฺก\-วิจารมตฺโต จ ธมฺมา สวิตกฺก\-สวิจารสฺส ธมฺมสฺส อธิปติ\-ปจฺจเยน ปจฺจโย. อารมฺมณา\-ธิปติ\csromandash{}สวิตกฺก\-สวิจาเร ขนฺเธ จ วิตกฺกญฺจ ครุํ กตฺวา สวิตกฺก\-สวิจารา ขนฺธา อุปฺปชฺชนฺติ. (1)}

\prose{สวิตกฺก\-สวิจาโร จ อวิตกฺก\-วิจารมตฺโต จ ธมฺมา อวิตกฺก\-วิจารมตฺตสฺส ธมฺมสฺส อธิปติ\-ปจฺจเยน ปจฺจโย. อารมฺมณา\-ธิปติ\csromandash{}สวิตกฺก\-สวิจาเร ขนฺเธ จ วิตกฺกญฺจ ครุํ กตฺวา วิตกฺโก อุปฺปชฺชติ. (2)}

\prose{สวิตกฺก\-สวิจาโร จ อวิตกฺก\-วิจารมตฺโต จ ธมฺมา สวิตกฺก\-สวิจารสฺส จ อวิตกฺก\-วิจารมตฺตสฺส จ ธมฺมสฺส อธิปติ\-ปจฺจเยน ปจฺจโย. อารมฺมณา\-ธิปติ\csromandash{}สวิตกฺก\-สวิจาเร ขนฺเธ จ วิตกฺกญฺจ ครุํ กตฺวา สวิตกฺก\-สวิจารา ขนฺธา จ วิตกฺโก จ อุปฺปชฺชนฺติ. (3)}

\csromanheader{2}{\textbf{อนนฺตร}}
\proseitem{79}{\csromanfolio{50}สวิตกฺก\-สวิจาโร ธมฺโม สวิตกฺก\-สวิจารสฺส ธมฺมสฺส อนนฺตร\-ปจฺจเยน ปจฺจโย. ปุริมา ปุริมา สวิตกฺก\-สวิจารา ขนฺธา ปจฺฉิมานํ ปจฺฉิมานํ สวิตกฺก\-สวิจารานํ ขนฺธานํ อนนฺตร\-ปจฺจเยน ปจฺจโย. อนุโลมํ โคตฺรภุสฺส. อนุโลมํ โวทานสฺส. โคตฺรภุ สวิตกฺก\-สวิจารสฺส มคฺคสฺส. โวทานํ สวิตกฺก\-สวิจารสฺส มคฺคสฺส. สวิตกฺก\-สวิจาโร มคฺโค สวิตกฺก\-สวิจารสฺส ผลสฺส. สวิตกฺก\-สวิจารํ ผลํ สวิตกฺก\-สวิจารสฺส ผลสฺส. อนุโลมํ สวิตกฺก\-สวิจาราย ผลสมาปตฺติยา อนนฺตร\-ปจฺจเยน ปจฺจโย. (1)}

\prose{สวิตกฺก\-สวิจาโร ธมฺโม อวิตกฺก\-วิจารมตฺตสฺส ธมฺมสฺส อนนฺตร\-ปจฺจเยน ปจฺจโย. ปุริมา ปุริมา สวิตกฺก\-สวิจารา ขนฺธา ปจฺฉิมสฺส ปจฺฉิมสฺส วิตกฺกสฺส อนนฺตร\-ปจฺจเยน ปจฺจโย. สวิตกฺก\-สวิจารํ จุติจิตฺตํ อวิตกฺก\-วิจารมตฺตสฺส อุปปตฺติ\-จิตฺตสฺส ฯเปฯ. สวิตกฺก\-สวิจารา ขนฺธา อวิตกฺก\-วิจารมตฺตสฺส วุฏฺฐานสฺส วิตกฺกสฺส จ ฯเปฯ. อวิตกฺก\-วิจารมตฺตสฺส ฌานสฺส ปริกมฺมํ อวิตกฺก\-วิจารมตฺตสฺส ฌานสฺส. โคตฺรภุ อวิตกฺก\-วิจารมตฺตสฺส มคฺคสฺส. โวทานํ อวิตกฺก\-วิจารมตฺตสฺส มคฺคสฺส. อนุโลมํ อวิตกฺก\-วิจารมตฺตาย ผลสมาปตฺติยา วิตกฺกสฺส จ อนนฺตร\-ปจฺจเยน ปจฺจโย. (2)}

\prose{สวิตกฺก\-สวิจาโร ธมฺโม อวิตกฺก\-อวิจารสฺส ธมฺมสฺส อนนฺตร\-ปจฺจเยน ปจฺจโย. สวิตกฺก\-สวิจารํ จุจิจิตฺตํ อวิตกฺก\-อวิจารสฺส อุปปตฺติ\-จิตฺตสฺส วิจารสฺส จ อนนฺตร\-ปจฺจเยน ปจฺจโย. อาวชฺชนา ปญฺจนฺนํ วิญฺญาณานํ อนนฺตร\-ปจฺจเยน ปจฺจโย. สวิตกฺก\-สวิจารา ขนฺธา อวิตกฺก\-อวิจารสฺส วุฏฺฐานสฺส วิจารสฺส จ ฯเปฯ. ทุติยสฺส ฌานสฺส ปริกมฺมํ ทุติเย ฌาเน วิจารสฺส อนนฺตร\-ปจฺจเยน ปจฺจโย. ตติยสฺส ฌานสฺส ปริกมฺมํ ฯเปฯ. จตุตฺถสฺส ฌานสฺส ปริกมฺมํ ฯเปฯ. อากาสา\-นญฺจา\-ยตนสฺส ปริกมฺมํ ฯเปฯ. วิญฺญาณญฺจา\-ยตนสฺส ปริกมฺมํ ฯเปฯ. อากิญฺจญฺญา\-ยตนสฺส ปริกมฺมํ ฯเปฯ. เนว\-สญฺญา\-นา\-สญฺญา\-ยตนสฺส ปริกมฺมํ ฯเปฯ. ทิพฺพสฺส จกฺขุสฺส ปริกมฺมํ ฯเปฯ. ทิพฺพย โสตธาตุยา ปริกมฺมํ ฯเปฯ. อิทฺธิวิธ\-ญาณสฺส ปริกมฺมํ ฯเปฯ. เจโตปริย\-ญาณสฺส ปริกมฺมํ ฯเปฯ. ปุพฺเพ\-นิวาสา\-นุสฺสติ\-ญาณสฺส ปริกมฺมํ ฯเปฯ. ยถากมฺมูปคญาณสฺส ปริกมฺมํ ฯเปฯ. อนาคตํสญาณสฺส ปริกมฺมํ ฯเปฯ. โคตฺรภุ อวิตกฺก\-อวิจารสฺส มคฺคสฺส วิจารสฺส จ. โวทานํ อวิตกฺก\-อวิจารสฺส \csromanfolio{51}มคฺคสฺส วิจารสฺส จ. อนุโลมํ อวิตกฺก\-อวิจาราย ผลสมาปตฺติยา วิจารสฺส จ อนนฺตร\-ปจฺจเยน ปจฺจโย. (3)}

\prose{สวิตกฺก\-สวิจาโร ธมฺโม อวิตกฺก\-วิจารมตฺตสฺส จ อวิตกฺก\-อวิจารสฺส จ ธมฺมสฺส อนนฺตร\-ปจฺจเยน ปจฺจโย. สวิตกฺก\-สวิจารํ จุติจิตฺตํ อวิตกฺก\-วิจารมตฺตสฺส อุปปตฺติ\-จิตฺตสฺส วิจารสฺส จ อนนฺตร\-ปจฺจเยน ปจฺจโย. สวิตกฺก\-สวิจารา ขนฺธา อวิตกฺก\-วิจารมตฺตสฺส วุฏฺฐานสฺส วิจารสฺส จ อนนฺตร\-ปจฺจเยน ปจฺจโย. อวิตกฺก\-วิจารมตฺตสฺส ฌานสฺส ปริกมฺมํ อวิตกฺก\-วิจารมตฺตสฺส ฌานสฺส วิจารสฺส จ อนนฺตร\-ปจฺจเยน ปจฺจโย. โคตฺรภุ อวิตกฺก\-วิจารมตฺตสฺส มคฺคสฺส วิจารสฺส จ. โวทานํ อวิตกฺก\-วิจารมตฺตสฺส มคฺคสฺส วิจารสฺส จ. อนุโลมํ อวิตกฺก\-วิจารมตฺตาย ผลสมาปตฺติยา วิจารสฺส จ อนนฺตร\-ปจฺจเยน ปจฺจโย. (4)}

\prose{สวิตกฺก\-สวิจาโร ธมฺโม สวิตกฺก\-สวิจารสฺส จ อวิตกฺก\-วิจารมตฺตสฺส จ ธมฺมสฺส อนนฺตร\-ปจฺจเยน ปจฺจโย. ปุริมา ปุริมา สวิตกฺก\-สวิจารา ขนฺธา ปจฺฉิมานํ ปจฺฉิมานํ สวิตกฺก\-สวิจารานํ ขนฺธานํ วิตกฺกสฺส จ อนนฺตร\-ปจฺจเยน ปจฺจโย. อนุโลมํ โคตฺรภุสฺส วิตกฺกสฺส จ. อนุโลมํ โวทานสฺส วิตกฺกสฺส จ. โคตฺรภุ สวิตกฺก\-สวิจารสฺส มคฺคสฺส วิตกฺกสฺส จ. โวทานํ สวิตกฺก\-สวิจารสฺส มคฺคสฺส วิตกฺกสฺส จ. สวิตกฺก\-สวิจาโร มคฺโค สวิตกฺก\-สวิจารสฺส ผลสฺส วิตกฺกสฺส จ. สวิตกฺก\-สวิจารํ ผลํ สวิตกฺก\-สวิจารสฺส ผลสฺส วิตกฺกสฺส จ. อนุโลมํ สวิตกฺก\-สวิจาราย ผลสมาปตฺติยา วิตกฺกสฺส จ อนนฺตร\-ปจฺจเยน ปจฺจโย. (5)}

\proseitem{80}{อวิตกฺก\-วิจารมตฺโต ธมฺโม อวิตกฺก\-วิจารมตฺตสฺส ธมฺมสฺส อนนฺตร\-ปจฺจเยน ปจฺจโย. ปุริโม ปุริโม วิตกฺโก ปจฺฉิมสฺส ปจฺฉิมสฺส วิตกฺกสฺส อนนฺตร\-ปจฺจเยน ปจฺจโย. ปุริมา ปุริมา อวิตกฺก\-วิจารมตฺตา ขนฺธา ปจฺฉิมานํ ปจฺฉิมานํ อวิตกฺก\-วิจารมตฺตานํ ขนฺธานํ อนนฺตร\-ปจฺจเยน ปจฺจโย. อวิตกฺก\-วิจารมตฺโต มคฺโค อวิตกฺก\-วิจารมตฺตสฺส ผลสฺส. อวิตกฺก\-วิจารมตฺตํ ผลํ อวิตกฺก\-วิจารมตฺตสฺส ผลสฺส อนนฺตร\-ปจฺจเยน ปจฺจโย. (1)}

\prose{\csromanfolio{52}อวิตกฺก\-วิจารมตฺโต ธมฺโม สวิตกฺก\-สวิจารสฺส ธมฺมสฺส อนนฺตร\-ปจฺจเยน ปจฺจโย. ปุริโม ปุริโม วิตกฺโก ปจฺฉิมานํ ปจฺฉิมานํ สวิตกฺก\-สวิจารานํ ขนฺธานํ อนนฺตร\-ปจฺจเยน ปจฺจโย. อวิตกฺก\-วิจารมตฺตํ จุติจิตฺตํ สวิตกฺก\-สวิจารสฺส อุปปตฺติ\-จิตฺตสฺส อนนฺตร\-ปจฺจเยน ปจฺจโย. อวิตกฺก\-วิจารมตฺตํ ภวงฺคํ อาวชฺชนาย อนนฺตร\-ปจฺจเยน ปจฺจโย. อวิตกฺก\-วิจารมตฺตา ขนฺธา สวิตกฺก\-สวิจารสฺส วุฏฺฐานสฺส อนนฺตร\-ปจฺจเยน ปจฺจโย. (2)}

\prose{อวิตกฺก\-วิจารมตฺโต ธมฺโม อวิตกฺก\-อวิจารสฺส ธมฺมสฺส อนนฺตร\-ปจฺจเยน ปจฺจโย. ปุริมา ปุริมา อวิตกฺก\-วิจารมตฺตา ขนฺธา ปจฺฉิมสฺส ปจฺฉิมสฺส วิจารสฺส อนนฺตร\-ปจฺจเยน ปจฺจโย. อวิตกฺก\-วิจารมตฺตํ จุติจิตฺตํ วิตกฺโก จ อวิตกฺก\-อวิจารสฺส อุปปตฺติ\-จิตฺตสฺส วิจารสฺส จ อนนฺตร\-ปจฺจเยน ปจฺจโย. อวิตกฺก\-วิจารมตฺตา ขนฺธา วิตกฺโก จ อวิตกฺก\-อวิจารสฺส วุฏฺฐานสฺส วิจารสฺส จ อนนฺตร\-ปจฺจเยน ปจฺจโย. (3)}

\prose{อวิตกฺก\-วิจารมตฺโต ธมฺโม อวิตกฺก\-วิจารมตฺตสฺส จ อวิตกฺก\-อวิจารสฺส จ ธมฺมสฺส อนนฺตร\-ปจฺจเยน ปจฺจโย. ปุริมา ปุริมา อวิตกฺก\-วิจารมตฺตา ขนฺธา ปจฺฉิมานํ ปจฺฉิมานํ อวิตกฺก\-วิจารมตฺตานํ ขนฺธานํ วิจารสฺส จ อนนฺตร\-ปจฺจเยน ปจฺจโย. อวิตกฺก\-วิจารมตฺโต มคฺโค อวิตกฺก\-วิจารมตฺตสฺส ผลสฺส วิจารสฺส จ. อวิตกฺก\-วิจารมตฺตํ ผลํ อวิตกฺก\-วิจารมตฺตสฺส ผลสฺส วิจารสฺส จ อนนฺตร\-ปจฺจเยน ปจฺจโย. (4)}

\prose{อวิตกฺก\-วิจารมตฺโต ธมฺโม สวิตกฺก\-สวิจารสฺส จ อวิตกฺก\-วิจารมตฺตสฺส จ ธมฺมสฺส อนนฺตร\-ปจฺจเยน ปจฺจโย. ปุริโม ปุริโม วิตกฺโก ปจฺฉิมานํ ปจฺฉิมานํ สวิตกฺก\-สวิจารานํ ขนฺธานํ วิตกฺกสฺส จ อนนฺตร\-ปจฺจเยน ปจฺจโย. อวิตกฺก\-วิจารมตฺตํ จุติจิตฺตํ สวิตกฺก\-สวิจารสฺส อุปปตฺติ\-จิตฺตสฺส วิตกฺกสฺส จ อนนฺตร\-ปจฺจเยน ปจฺจโย. อวิตกฺก\-วิจารมตฺตํ ภวงฺคํ อาวชฺชนาย วิตกฺกสฺส จ อนนฺตร\-ปจฺจเยน ปจฺจโย. อวิตกฺก\-วิจารมตฺตา ขนฺธา สวิตกฺก\-สวิจารสฺส วุฏฺฐานสฺส วิตกฺกสฺส จ อนนฺตร\-ปจฺจเยน ปจฺจโย. (5)}

\proseitem{81}{อวิตกฺก\-อวิจาโร ธมฺโม อวิตกฺก\-อวิจารสฺส ธมฺมสฺส อนนฺตร\-ปจฺจเยน ปจฺจโย. ปุริโม ปุริโม วิจาโร ปจฺฉิมสฺส ปจฺฉิมสฺส วิจารสฺส อนนฺตร\-ปจฺจเยน ปจฺจโย. ปุริมา ปุริมา อวิตกฺก\-อวิจารา ขนฺธา ปจฺฉิมานํ \csromanfolio{53}ปจฺฉิมานํ อวิตกฺก\-อวิจารานํ ขนฺธานํ อนนฺตร\-ปจฺจเยน ปจฺจโย. อวิตกฺก\-อวิจาโร มคฺโค อวิตกฺก\-อวิจารสฺส ผลสฺส. อวิตกฺก\-อวิจารํ ผลํ อวิตกฺก\-อวิจารสฺส ผลสฺส อนนฺตร\-ปจฺจเยน ปจฺจโย. นิโรธา วุฏฺฐหนฺตสฺส เนว\-สญฺญา\-นา\-สญฺญา\-ยตนํ อวิตกฺก\-อวิจาราย ผลสมาปตฺติยา วิจารสฺส จ อนนฺตร\-ปจฺจเยน ปจฺจโย. (1)}

\prose{อวิตกฺก\-อวิจาโร ธมฺโม สวิตกฺก\-สวิจารสฺส ธมฺมสฺส อนนฺตร\-ปจฺจเยน ปจฺจโย. อวิตกฺก\-อวิจารํ จุติจิตฺตํ วิจาโร จ สวิตกฺก\-สวิจารสฺส อุปปตฺติ\-จิตฺตสฺส อนนฺตร\-ปจฺจเยน ปจฺจโย. อวิตกฺก\-อวิจารํ ภวงฺคํ วิจาโร จ อาวชฺชนาย อนนฺตร\-ปจฺจเยน ปจฺจโย. อวิตกฺก\-อวิจารา ขนฺธา วิจาโร จ สวิตกฺก\-สวิจารสฺส วุฏฺฐานสฺส อนนฺตร\-ปจฺจเยน ปจฺจโย. นิโรธา วุฏฺฐหนฺตสฺส เนว\-สญฺญา\-นา\-สญฺญา\-ยตนํ สวิตกฺก\-สวิจาราย ผลสมาปตฺติยา อนนฺตร\-ปจฺจเยน ปจฺจโย. (2)}

\prose{อวิตกฺก\-อวิจาโร ธมฺโม อวิตกฺก\-วิจารมตฺตสฺส ธมฺมสฺส อนนฺตร\-ปจฺจเยน ปจฺจโย. ปุริโม ปุริโม วิจาโร ปจฺฉิมานํ ปจฺฉิมานํ อวิตกฺก\-วิจารมตฺตานํ ขนฺธานํ อนนฺตร\-ปจฺจเยน ปจฺจโย. อวิตกฺก\-อวิจารํ จุติจิตฺตํ วิจาโร จ อวิตกฺก\-วิจารมตฺตสฺส อุปปตฺติ\-จิตฺตสฺส วิตกฺกสฺส จ อนนฺตร\-ปจฺจเยน ปจฺจโย. อวิตกฺก\-อวิจารา ขนฺธา วิจาโร จ อวิตกฺก\-วิจารมตฺตสฺส วุฏฺฐานสฺส วิตกฺกสฺส จ อนนฺตร\-ปจฺจเยน ปจฺจโย. นิโรธา วุฏฺฐหนฺตสฺส เนว\-สญฺญา\-นา\-สญฺญา\-ยตนํ อวิตกฺก\-วิจารมตฺตาย ผลสมาปตฺติยา วิตกฺกสฺส จ อนนฺตร\-ปจฺจเยน ปจฺจโย. (3)}

\prose{อวิตกฺก\-อวิจาโร ธมฺโม อวิตกฺก\-วิจารมตฺตสฺส จ อวิตกฺก\-อวิจารสฺส จ ธมฺมสฺส อนนฺตร\-ปจฺจเยน ปจฺจโย. ปุริโม ปุริโม วิจาโร ปจฺฉิมานํ ปจฺฉิมานํ อวิตกฺก\-วิจารมตฺตานํ ขนฺธานํ วิจารสฺส จ อนนฺตร\-ปจฺจเยน ปจฺจโย. อวิตกฺก\-อวิจารํ จุติจิตฺตํ อวิตกฺก\-วิจารมตฺตสฺส อุปปตฺติ\-จิตฺตสฺส วิจารสฺส จ อนนฺตร\-ปจฺจเยน ปจฺจโย. อวิตกฺก\-อวิจารา ขนฺธา อวิตกฺก\-วิจารมตฺตสฺส วุฏฺฐานสฺส วิจารสฺส จ อนนฺตร\-ปจฺจเยน ปจฺจโย. นิโรธา วุฏฺฐหนฺตสฺส เนว\-สญฺญา\-นา\-สญฺญา\-ยตนํ อวิตกฺก\-วิจารมตฺตาย ผลสมาปตฺติยา วิจารสฺส จ อนนฺตร\-ปจฺจเยน ปจฺจโย. (4)}

\prose{อวิตกฺก\-อวิจาโร ธมฺโม สวิตกฺก\-สวิจารสฺส จ อวิตกฺก\-วิจารมตฺตสฺส จ ธมฺมสฺส อนนฺตร\-ปจฺจเยน ปจฺจโย. อวิตกฺก\-อวิจารํ จุติจิตฺตํ \csromanfolio{54}วิจาโร จ สวิตกฺก\-สวิจารสฺส อุปปตฺติ\-จิตฺตสฺส วิตกฺกสฺส จ อนนฺตร\-ปจฺจเยน ปจฺจโย. อวิตกฺก\-อวิจารํ ภวงฺคญฺจ วิจาโร จ อาวชฺชนาย วิตกฺกสฺส จ อนนฺตร\-ปจฺจเยน ปจฺจโย. อวิตกฺก\-อวิจารา ขนฺธา วิจาโร จ สวิตกฺก\-สวิจารสฺส วุฏฺฐานสฺส วิตกฺกสฺส จ อนนฺตร\-ปจฺจเยน ปจฺจโย. นิโรธา วุฏฺฐหนฺตสฺส เนว\-สญฺญา\-นา\-สญฺญา\-ยตนํ สวิตกฺก\-สวิจาราย ผลสมาปตฺติยา วิตกฺกสฺส จ อนนฺตร\-ปจฺจเยน ปจฺจโย. (5)}

\proseitem{82}{อวิตกฺก\-วิจารมตฺโต จ อวิตกฺก\-อวิจาโร จ ธมฺมา สวิตกฺก\-สวิจารสฺส ธมฺมสฺส อนนฺตร\-ปจฺจเยน ปจฺจโย. อวิตกฺก\-วิจารมตฺตํ จุติจิตฺตญฺจ วิจาโร จ สวิตกฺก\-สวิจารสฺส อุปปตฺติ\-จิตฺตสฺส อนนฺตร\-ปจฺจเยน ปจฺจโย. อวิตกฺก\-วิจารมตฺตํ ภวงฺคญฺจ วิจาโร จ อาวชฺชนาย อนนฺตร\-ปจฺจเยน ปจฺจโย. อวิตกฺก\-วิจารมตฺตา ขนฺธา จ วิจาโร จ สวิตกฺก\-สวิจารสฺส วุฏฺฐานสฺส อนนฺตร\-ปจฺจเยน ปจฺจโย. (1)}

\prose{อวิตกฺก\-วิจารมตฺโต จ อวิตกฺก\-อวิจาโร จ ธมฺมา อวิตกฺก\-วิจารมตฺตสฺส ธมฺมสฺส อนนฺตร\-ปจฺจเยน ปจฺจโย. ปุริมา ปุริมา อวิตกฺก\-วิจารมตฺตา ขนฺธา จ วิจาโร จ ปจฺฉิมานํ ปจฺฉิมานํ อวิตกฺก\-วิจารมตฺตานํ ขนฺธานํ อนนฺตร\-ปจฺจเยน ปจฺจโย. อวิตกฺก\-วิจารมตฺโต มคฺโค จ วิจาโร จ อวิตกฺก\-วิจารมตฺตสฺส ผลสฺส อนนฺตร\-ปจฺจเยน ปจฺจโย. อวิตกฺก\-วิจารมตฺตํ ผลญฺจ วิจาโร จ อวิตกฺก\-วิจารมตฺตสฺส ผลสฺส อนนฺตร\-ปจฺจเยน ปจฺจโย. (2)}

\prose{อวิตกฺก\-วิจารมตฺโต จ อวิตกฺก\-อวิจาโร จ ธมฺมา อวิตกฺก\-อวิจารสฺส ธมฺมสฺส อนนฺตร\-ปจฺจเยน ปจฺจโย. ปุริมา ปุริมา อวิตกฺก\-วิจารมตฺตา ขนฺธา จ วิจาโร จ ปจฺฉิมสฺส ปจฺฉิมสฺส วิจารสฺส อนนฺตร\-ปจฺจเยน ปจฺจโย. อวิตกฺก\-วิจารมตฺตํ จุติจิตฺตญฺจ วิจาโร จ อวิตกฺก\-อวิจารสฺส อุปปตฺติ\-จิตฺตสฺส อนนฺตร\-ปจฺจเยน ปจฺจโย. อวิตกฺก\-วิจารมตฺตา ขนฺธา จ วิจาโร จ อวิตกฺก\-อวิจารสฺส วุฏฺฐนสฺส อนนฺตร\-ปจฺจเยน ปจฺจโย. (3)}

\prose{อวิตกฺก\-วิจารมตฺโต จ อวิตกฺก\-อวิจาโร จ ธมฺมา อวิตกฺก\-วิจารมตฺตสฺส จ อวิตกฺก\-อวิจารสฺส จ ธมฺมสฺส อนนฺตร\-ปจฺจเยน ปจฺจโย. ปุริมา ปุริมา อวิตกฺก\-วิจารมตฺตา ขนฺธา จ วิจาโร จ ปจฺฉิมานํ ปจฺฉิมานํ อวิตกฺก\-วิจารมตฺตานํ ขนฺธานํ วิจารสฺส จ อนนฺตร\-ปจฺจเยน ปจฺจโย. อวิตกฺก\-วิจารมตฺโต มคฺโค จ วิจาโร จ อวิตกฺก\-วิจารมตฺตสฺส ผลสฺส \csromanfolio{55}วิจารสฺส จ อนนฺตร\-ปจฺจเยน ปจฺจโย. อวิตกฺก\-วิจารมตฺตํ ผลญฺจ วิจาโร จ อวิตกฺก\-วิจารมตฺตสฺส ผลสฺส จ วิจารสฺส จ อนนฺตร\-ปจฺจเยน ปจฺจโย. (4)}

\prose{อวิตกฺก\-วิจารมตฺโต จ อวิตกฺก\-อวิจาโร จ ธมฺมา สวิตกฺก\-สวิจารสฺส จ อวิตกฺก\-วิจารมตฺตสฺส จ ธมฺมสฺส อนนฺตร\-ปจฺจเยน ปจฺจโย. อวิตกฺก\-วิจารมตฺตํ จุติจิตฺตญฺจ วิจาโร จ สวิตกฺก\-สวิจารสฺส อุปปตฺติ\-จิตฺตสฺส วิตกฺกสฺส จ อนนฺตร\-ปจฺจเยน ปจฺจโย. อวิตกฺก\-วิจารมตฺตํ ภวงฺคญฺจ วิจาโร จ อาวชฺชนาย วิตกฺกสฺส จ อนนฺตร\-ปจฺจเยน ปจฺจโย. อวิตกฺก\-วิจารมตฺตา ขนฺธา จ วิจาโร จ สวิตกฺก\-สวิจารสฺส วุฏฺฐานสฺส วิตกฺกสฺส จ อนนฺตร\-ปจฺจเยน ปจฺจโย. (5)}

\proseitem{83}{สวิตกฺก\-สวิจาโร จ อวิตกฺก\-วิจารมตฺโต จ ธมฺมา สวิตกฺก\-สวิจารสฺส ธมฺมสฺส อนนฺตร\-ปจฺจเยน ปจฺจโย. ปุริมา ปุริมา สวิตกฺก\-สวิจารา ขนฺธา จ วิตกฺโก จ ปจฺฉิมานํ ปจฺฉิมานํ สวิตกฺก\-สวิจารานํ ขนฺธานํ อนนฺตร\-ปจฺจเยน ปจฺจโย. อนุโลมญฺจ วิตกฺโก จ โคตฺรภุสฺส. อนุโลมญฺจ วิตกฺโกจ โวทานสฺส. โคตฺรภุ จ วิตกฺโก จ สวิตกฺก\-สวิจารสฺส มคฺคสฺส. โวทานญฺจ วิตกฺโก จ สวิตกฺก\-สวิจารสฺส มคฺคสฺส. สวิตกฺก\-สวิจาโร มคฺโค จ วิตกฺโก จ สวิตกฺก\-สวิจารสฺส ผลสฺส. สวิตกฺก\-สวิจารํ ผลญฺจ วิตกฺโก จ สวิตกฺก\-สวิจารสฺส ผลสฺส. อนุโลมญฺจ วิตกฺโก จ สวิตกฺก\-สวิจาราย ผลสมาปตฺติยา อนนฺตร\-ปจฺจเยน ปจฺจโย. (1)}

\prose{สวิตกฺก\-สวิจาโร จ อวิตกฺก\-วิจารมตฺโต จ ธมฺมา อวิตกฺก\-วิจารมตฺตสฺส ธมฺมสฺส อนนฺตร\-ปจฺจเยน ปจฺจโย. ปุริมา ปุริมา สวิตกฺก\-สวิจารา ขนฺธา จ วิตกฺโก จ ปจฺฉิมสฺส ปจฺฉิมสฺส วิตกฺกสฺส อนนฺตร\-ปจฺจเยน ปจฺจโย. สวิตกฺก\-สวิจารํ จุติจิตฺตญฺจ วิตกฺโก จ อวิตกฺก\-วิจารมตฺตสฺส อุปปตฺติ\-จิตฺตสฺส อนนฺตร\-ปจฺจเยน ปจฺจโย. สวิตกฺก\-สวิจารา ขนฺธา จ วิตกฺโก จ อวิตกฺก\-วิจารมตฺตสฺส วุฏฺฐานสฺส อนนฺตร\-ปจฺจเยน ปจฺจโย. อวิตกฺก\-วิจารมตฺตสฺส ฌานสฺส ปริกมฺมญฺจ วิตกฺโก จ อวิตกฺก\-วิจารมตฺตสฺส ฌานสฺส อนนฺตร\-ปจฺจเยน ปจฺจโย. โคตฺรภุ จ วิตกฺโก จ อวิตกฺก\-วิจารมตฺตสฺส มคฺคสฺส. โวทานญฺจ วิตกฺโก จ อวิตกฺก\-วิจารมตฺตสฺส \csromanfolio{56}มคฺคสฺส. อนุโลมญฺจ วิตกฺโก จ อวิตกฺก\-วิจารมตฺตาย ผลสมาปตฺติยา อนนฺตร\-ปจฺจเยน ปจฺจโย. (2)}

\prose{สวิตกฺก\-สวิจาโร จ อวิตกฺก\-วิจารมตฺโต จ ธมฺมา อวิตกฺก\-อวิจารสฺส ธมฺมสฺส อนนฺตร\-ปจฺจเยน ปจฺจโย. สวิตกฺก\-สวิจารํ จุติจิตฺตญฺจ วิตกฺโก จ อวิตกฺก\-อวิจารสฺส อุปปตฺติ\-จิตฺตสฺส วิจารสฺส จ อนนฺตร\-ปจฺจเยน ปจฺจโย. อาวชฺชนา จ วิตกฺโก จ ปญฺจนฺนํ วิญฺญาณานํ อนนฺตร\-ปจฺจเยน ปจฺจโย. สวิตกฺก\-สวิจารา ขนฺธา จ วิตกฺโก จ อวิตกฺก\-อวิจารสฺส วุฏฺฐานสฺส จ วิจารสฺส จ อนนฺตร\-ปจฺจเยน ปจฺจโย. ทุติยสฺส ฌานสฺส ปริกมฺมญฺจ วิตกฺโก จ ทุติเย ฌาเน วิจารสฺส อนนฺตร\-ปจฺจเยน ปจฺจโย. ตติยสฺส ฌานสฺส ปริกมฺมญฺจ วิตกฺโก จ ฯเปฯ. จตุตฺถสฺส ฌานสฺส ปริกมฺมญฺจ วิตกฺโก จ ฯเปฯ. อากาสา\-นญฺจา\-ยตนสฺส ปริกมฺมญฺจ วิตกฺโก จ ฯเปฯ. วิญฺญาณญฺจา\-ยตนสฺส ปริกมฺมญฺจ วิตกฺโก จ ฯเปฯ. อากิญฺจญฺญา\-ยตนสฺส ปริกมฺมญฺจ วิตกฺโก จ ฯเปฯ. เนว\-สญฺญา\-นา\-สญฺญา\-ยตนสฺส ปริกมฺมญฺจ วิตกฺโก จ ฯเปฯ. ทิพฺพสฺส จกฺขุสฺส ปริกมฺมญฺจ วิตกฺโก จ ฯเปฯ. ทิพฺพาย โสตธาตุยา ปริกมฺมญฺจ วิตกฺโก จ ฯเปฯ. อิทฺธิวิธ\-ญาณสฺส ปริกมฺมญฺจ วิตกฺโก จ ฯเปฯ. เจโตปริย\-ญาณสฺส ปริกมฺมญฺจ วิตกฺโก จ ฯเปฯ. ปุพฺเพ\-นิวาสา\-นุสฺสติ\-ญาณสฺส ปริกมฺมญฺจ วิตกฺโก จ ฯเปฯ. ยถากมฺมูปคญาณสฺส ปริกมฺมญฺจ วิตกฺโก จ ฯเปฯ. อนาคตํสญาณสฺส ปริกมฺมญฺจ วิตกฺโก จ ฯเปฯ. โคตฺรภุ จ วิตกฺโก จ อวิตกฺก\-อวิจารสฺส มคฺคสฺส วิจารสฺส จ. โวทานญฺจ วิตกฺโก จ อวิตกฺก\-อวิจารสฺส มคฺคสฺส วิจารสฺส จ. อนุโลมญฺจ วิตกฺโก จ อวิตกฺก\-อวิจาราย ผลสมาปตฺติยา วิจารสฺส จ อนนฺตร\-ปจฺจเยน ปจฺจโย. (3)}

\prose{สวิตกฺก\-สวิจาโร จ อวิตกฺก\-วิจารมตฺโต จ ธมฺมา อวิตกฺก\-วิจารมตฺตสฺส จ อวิตกฺก\-อวิจารสฺส จ ธมฺมสฺส อนนฺตร\-ปจฺจเยน ปจฺจโย. สวิตกฺก\-สวิจารํ จุติจิตฺตญฺจ วิตกฺโก จ อวิตกฺก\-วิจารมตฺตสฺส อุปปตฺติ\-จิตฺตสฺส จ วิจารสฺส จ อนนฺตร\-ปจฺจเยน ปจฺจโย. สวิตกฺก\-สวิจารา ขนฺธา จ วิตกฺโก จ อวิตกฺก\-วิจารมตฺตสฺส วุฏฺฐานสฺส จ วิจารสฺส จ อนนฺตร\-ปจฺจเยน ปจฺจโย. อวิตกฺก\-วิจารมตฺตสฺส ฌานสฺส ปริกมฺมญฺจ วิตกฺโก จ อวิตกฺก\-วิจารมตฺตสฺส ฌานสฺส จ วิจารสฺส จ อนนฺตร\-ปจฺจเยน ปจฺจโย. โคตฺรภุ จ วิตกฺโก จ อวิตกฺก\-วิจารมตฺตสฺส มคฺคสฺส จ วิจารสฺส จ อนนฺตร\-ปจฺจเยน ปจฺจโย. โวทานญฺจ วิตกฺโก จ \csromanfolio{57}อวิตกฺก\-วิจารมตฺตสฺส มคฺคสฺส จ วิจารสฺส จ. อนุโลมญฺจ วิตกฺโก จ อวิตกฺก\-วิจารมตฺตาย ผลสมาปตฺติยา จ วิจารสฺส จ อนนฺตร\-ปจฺจเยน ปจฺจโย. (4)}

\prose{สวิตกฺก\-สวิจาโร จ อวิตกฺก\-วิจารมตฺโต จ ธมฺมา สวิตกฺก\-สวิจารสฺส จ อวิตกฺก\-วิจารมตฺตสฺส จ ธมฺมสฺส อนนฺตร\-ปจฺจเยน ปจฺจโย. ปุริมา ปุริมา สวิตกฺก\-สวิจารา ขนฺธา จ วิตกฺโก จ ปจฺฉิมานํ ปจฺฉิมานํ สวิตกฺก\-สวิจารานํ ขนฺธานํ วิตกฺกสฺส จ อนนฺตร\-ปจฺจเยน ปจฺจโย. อนุโลมญฺจ วิตกฺโก จ โคตฺรภุสฺส จ วิตกฺกสฺส จ. อนุโลมญฺจ วิตกฺโก จ โวทานสฺส จ วิตกฺกสฺส จ. โคตฺรภุ จ วิตกฺโก จ สวิตกฺก\-สวิจารสฺส จ มคฺคสฺส จ วิตกฺกสฺส จ. โวทานญฺจ วิตกฺโก จ สวิตกฺก\-สวิจารสฺส มคฺคสฺส จ วิตกฺกสฺส จ อนนฺตร\-ปจฺจเยน ปจฺจโย. สวิตกฺก\-สวิจาโร มคฺโค จ วิตกฺโก จ สวิตกฺก\-สวิจารสฺส ผลสฺส จ วิตกฺกสฺส จ. สวิตกฺก\-สวิจารํ ผลญฺจ วิตกฺโก จ สวิตกฺก\-สวิจารสฺส ผลสฺส จ วิตกฺกสฺส จ. อนุโลมญฺจ วิตกฺโก จ สวิตกฺก\-สวิจาราย ผลสมาปตฺติยา จ วิตกฺกสฺส จ อนนฺตร\-ปจฺจเยน ปจฺจโย. (5)}

\csromanheader{2}{\textbf{สมนนฺตร}}
\proseitem{84}{สวิตกฺก\-สวิจาโร ธมฺโม สวิตกฺก\-สวิจารสฺส ธมฺมสฺส สมนนฺตร\-ปจฺจเยน ปจฺจโย. (อนนฺตรปจฺจโยปิ สมนนฺตรปจฺจโยปิ สทิโส.)}

\vspace{\tipitakaparskip}
\csromancenter{(อนนฺตร\-ปจฺจโยปิ สมนนฺตรปจฺจโยปิ สทิโส.)}
\csromanheader{2}{\textbf{สหชาต}}
\proseitem{85}{สวิตกฺก\-สวิจาโร ธมฺโม สวิตกฺก\-สวิจารสฺส ธมฺมสฺส สหชาต\-ปจฺจเยน ปจฺจโย. สวิตกฺก\-สวิจาโร เอโก ขนฺโธ ติณฺณนฺนํ ขนฺธานํ สหชาต\-ปจฺจเยน ปจฺจโย. ตโย ขนฺธา เอกสฺส ขนฺธสฺส สหชาต\-ปจฺจเยน ปจฺจโย. เทฺว ขนฺธา ทฺวินฺนํ ขนฺธานํ สหชาต\-ปจฺจเยน ปจฺจโย. ปฏิสนฺธิ\-กฺขเณ สวิตกฺก\-สวิจาโร เอโก ขนฺโธ ติณฺณนฺนํ ขนฺธานํ ฯเปฯ เทฺว ขนฺธา ทฺวินฺนํ ขนฺธานํ ฯเปฯ. (1)}

\prose{สวิตกฺก\-สวิจาโร ธมฺโม อวิตกฺก\-วิจารมตฺตสฺส ธมฺมสฺส สหชาต\-ปจฺจเยน ปจฺจโย. สวิตกฺก\-สวิจารา ขนฺธา วิตกฺกสฺส สหชาต\-ปจฺจเยน ปจฺจโย. ปฏิสนฺธิ\-กฺขเณ ฯเปฯ. (2)}

\prose{\csromanfolio{58}สวิตกฺก\-สวิจาโร ธมฺโม อวิตกฺก\-อวิจารสฺส ธมฺมสฺส สหชาต\-ปจฺจเยน ปจฺจโย. สวิตกฺก\-สวิจารา ขนฺธา จิตฺต\-สมุฏฺฐานานํ รูปานํ สหชาต\-ปจฺจเยน ปจฺจโย. ปฏิสนฺธิ\-กฺขเณ ฯเปฯ กฏตฺตารูปานํ ฯเปฯ. (3)}

\prose{สวิตกฺก\-สวิจาโร ธมฺโม สวิตกฺก\-สวิจารสฺส จ อวิตกฺก\-อวิจารสฺส จ ธมฺมสฺส สหชาต\-ปจฺจเยน ปจฺจโย. สวิตกฺก\-สวิจาโร เอโก ขนฺโธ ติณฺณนฺนํ ขนฺธานํ จิตฺตสมุฏฺฐานานญฺจ รูปานํ สหชาต\-ปจฺจเยน ปจฺจโย ฯเปฯ เทฺว ขนฺธา ทฺวินฺนํ ขนฺธานํ จิตฺตสมุฏฺฐานานญฺจ รูปานํ สหชาต\-ปจฺจเยน ปจฺจโย. ปฏิสนฺธิ\-กฺขเณ ฯเปฯ. (4)}

\prose{สวิตกฺก\-สวิจาโร ธมฺโม อวิตกฺก\-วิจารมตฺตสฺส จ อวิตกฺก\-อวิจารสฺส จ ธมฺมสฺส สหชาต\-ปจฺจเยน ปจฺจโย. สวิตกฺก\-สวิจารา ขนฺธา วิตกฺกสฺส จ จิตฺตสมุฏฺฐานานญฺจ รูปานํ สหชาต\-ปจฺจเยน ปจฺจโย. ปฏิสนฺธิ\-กฺขเณ ฯเปฯ. (5)}

\prose{สวิตกฺก\-สวิจาโร ธมฺโม สวิตกฺก\-สวิจารสฺส จ อวิตกฺก\-วิจารมตฺตสฺส จ ธมฺมสฺส สหชาต\-ปจฺจเยน ปจฺจโย. สวิตกฺก\-สวิจาโร เอโก ขนฺโธ ติณฺณนฺนํ ขนฺธานํ วิตกฺกสฺส จ สหชาต\-ปจฺจเยน ปจฺจโย ฯเปฯ เทฺว ขนฺธา ทฺวินฺนํ ขนฺธานํ วิตกฺกสฺส จ สหชาต\-ปจฺจเยน ปจฺจโย. ปฏิสนฺธิ\-กฺขเณ ฯเปฯ. (6)}

\prose{สวิตกฺก\-สวิจาโร ธมฺโม สวิตกฺก\-สวิจารสฺส จ อวิตกฺก\-วิจารมตฺตสฺส จ อวิตกฺก\-อวิจารสฺส จ ธมฺมสฺส สหชาต\-ปจฺจเยน ปจฺจโย. สวิตกฺก\-สวิจาโร เอโก ขนฺโธ ติณฺณนฺนํ ขนฺธานํ วิตกฺกสฺส จ จิตฺตสมุฏฺฐานานญฺจ รูปานํ สหชาต\-ปจฺจเยน ปจฺจโย ฯเปฯ เทฺว ขนฺธา ทฺวินฺนํ ขนฺธานํ วิตกฺกสฺส จ จิตฺตสมุฏฺฐานานญฺจ รูปานํ ฯเปฯ. ปฏิสนฺธิ\-กฺขเณ ฯเปฯ. (7)}

\proseitem{86}{อวิตกฺก\-วิจารมตฺโต ธมฺโม อวิตกฺก\-วิจารมตฺตสฺส ธมฺมสฺส สหชาต\-ปจฺจเยน ปจฺจโย. อวิตกฺก\-วิจารมตฺโต เอโก ขนฺโธ ติณฺณนฺนํ ขนฺธานํ สหชาต\-ปจฺจเยน ปจฺจโย ฯเปฯ เทฺว ขนฺธา ทฺวินฺนํ ขนฺธานํ ฯเปฯ. ปฏิสนฺธิ\-กฺขเณ ฯเปฯ. (1)}

\prose{อวิตกฺก\-วิจารมตฺโต ธมฺโม สวิตกฺก\-สวิจารสฺส ธมฺมสฺส สหชาต\-ปจฺจเยน ปจฺจโย. วิตกฺโก สวิตกฺก\-สวิจารานํ ขนฺธานํ สหชาต\-ปจฺจเยน ปจฺจโย. ปฏิสนฺธิ\-กฺขเณ ฯเปฯ. (2)}

\prose{\csromanfolio{59}อวิตกฺก\-วิจารมตฺโต ธมฺโม อวิตกฺก\-อวิจารสฺส ธมฺมสฺส สหชาต\-ปจฺจเยน ปจฺจโย. อวิตกฺก\-วิจารมตฺตา ขนฺธา วิจารสฺส จิตฺตสมุฏฺฐานานญฺจ รูปานํ สหชาต\-ปจฺจเยน ปจฺจโย. วิตกฺโก จิตฺต\-สมุฏฺฐานานํ รูปานํ สหชาต\-ปจฺจเยน ปจฺจโย. ปฏิสนฺธิ\-กฺขเณ ฯเปฯ. (3)}

\prose{อวิตกฺก\-วิจารมตฺโต ธมฺโม สวิตกฺก\-สวิจารสฺส จ อวิตกฺก\-อวิจารสฺส จ ธมฺมสฺส สหชาต\-ปจฺจเยน ปจฺจโย. วิตกฺโก สวิตกฺก\-สวิจารานํ ขนฺธานํ จิตฺตสมุฏฺฐานานญฺจ รูปานํ สหชาต\-ปจฺจเยน ปจฺจโย. ปฏิสนฺธิ\-กฺขเณ ฯเปฯ. (4)}

\prose{อวิตกฺก\-วิจารมตฺโต ธมฺโม อวิตกฺก\-วิจารมตฺตสฺส จ อวิตกฺก\-อวิจารสฺส จ ธมฺมสฺส สหชาต\-ปจฺจเยน ปจฺจโย. อวิตกฺก\-วิจารมตฺโต เอโก ขนฺโธ ติณฺณนฺนํ ขนฺธานํ วิจารสฺส จ จิตฺตสมุฏฺฐานานญฺจ รูปานํ สหชาต\-ปจฺจเยน ปจฺจโย ฯเปฯ เทฺว ขนฺธา ทฺวินฺนํ ขนฺธานํ วิจารสฺส จ จิตฺตสมุฏฺฐานานญฺจ รูปานํ ฯเปฯ. ปฏิสนฺธิ\-กฺขเณ ฯเปฯ. (5)}

\proseitem{87}{อวิตกฺก\-อวิจาโร ธมฺโม อวิตกฺก\-อวิจารสฺส ธมฺมสฺส สหชาต\-ปจฺจเยน ปจฺจโย. อวิตกฺก\-อวิจาโร เอโก ขนฺโธ ติณฺณนฺนํ ขนฺธานํ จิตฺตสมุฏฺฐานานญฺจ รูปานํ สหชาต\-ปจฺจเยน ปจฺจโย ฯเปฯ เทฺว ขนฺธา ทฺวินฺนํ ขนฺธานํ จิตฺตสมุฏฺฐานานญฺจ รูปานํ สหชาต\-ปจฺจเยน ปจฺจโย. วิจาโร จิตฺต\-สมุฏฺฐานานํ รูปานํ ฯเปฯ. ปฏิสนฺธิ\-กฺขเณ อวิตกฺก\-อวิจาโร เอโก ขนฺโธ ติณฺณนฺนํ ขนฺธานํ กฏตฺตา จ รูปานํ ฯเปฯ เทฺว ขนฺธา ทฺวินฺนํ ขนฺธานํ กฏตฺตา จ รูปานํ ฯเปฯ. วิจาโร กฏตฺตารูปานํ ฯเปฯ. ขนฺธา วตฺถุสฺส ฯเปฯ. วตฺถุ ขนฺธานํ ฯเปฯ. วิจาโร วตฺถุสฺส ฯเปฯ. วตฺถุ วิจารสฺส ฯเปฯ. เอกํ มหาภูตํ ติณฺณนฺนํ มหาภูตานํ ฯเปฯ มหาภูตา จิตฺต\-สมุฏฺฐานานํ รูปานํ กฏตฺตารูปานํ อุปาทารูปานํ ฯเปฯ. พาหิรํ. อาหารสมุฏฺฐานํ. อุตุสมุฏฺฐานํ. อสญฺญสตฺตานํ เอกํ มหาภูตํ ฯเปฯ มหาภูตา กฏตฺตารูปานํ อุปาทารูปานํ สหชาต\-ปจฺจเยน ปจฺจโย. (1)}

\prose{อวิตกฺก\-อวิจาโร ธมฺโม สวิตกฺก\-สวิจารสฺส ธมฺมสฺส สหชาต\-ปจฺจเยน ปจฺจโย. ปฏิสนฺธิ\-กฺขเณ วตฺถุ สวิตกฺก\-สวิจารานํ ขนฺธานํ สหชาต\-ปจฺจเยน ปจฺจโย. (2)}

\prose{\csromanfolio{60}อวิตกฺก\-อวิจาโร ธมฺโม อวิตกฺก\-วิจารมตฺตสฺส ธมฺมสฺส สหชาต\-ปจฺจเยน ปจฺจโย. วิจาโร อวิตกฺก\-วิจารมตฺตานํ ขนฺธานํ สหชาต\-ปจฺจเยน ปจฺจโย. ปฏิสนฺธิ\-กฺขเณ วิจาโร อวิตกฺก\-วิจารมตฺตานํ ขนฺธานํ สหชาต\-ปจฺจเยน ปจฺจโย. ปฏิสนฺธิ\-กฺขเณ วตฺถุ อวิตกฺก\-วิจารมตฺตานํ ขนฺธานํ สหชาต\-ปจฺจเยน ปจฺจโย. ปฏิสนฺธิ\-กฺขเณ วตฺถุ วิตกฺกสฺส สหชาต\-ปจฺจเยน ปจฺจโย. (3)}

\prose{อวิตกฺก\-อวิจาโร ธมฺโม อวิตกฺก\-วิจารมตฺตสฺส จ อวิตกฺก\-อวิจารสฺส จ ธมฺมสฺส สหชาต\-ปจฺจเยน ปจฺจโย. วิจาโร อวิตกฺก\-วิจารมตฺตานํ ขนฺธานํ จิตฺตสมุฏฺฐานานญฺจ รูปานํ สหชาต\-ปจฺจเยน ปจฺจโย. ปฏิสนฺธิ\-กฺขเณ วิจาโร อวิตกฺก\-วิจารมตฺตานํ ขนฺธานํ กฏตฺตา จ รูปานํ ฯเปฯ. ปฏิสนฺธิ\-กฺขเณ วตฺถุ อวิตกฺก\-วิจารมตฺตานํ ขนฺธานํ วิจารสฺส จ สหชาต\-ปจฺจเยน ปจฺจโย. (4)}

\prose{อวิตกฺก\-อวิจาโร ธมฺโม สวิตกฺก\-สวิจารสฺส จ อวิตกฺก\-วิจารมตฺตสฺส จ ธมฺมสฺส สหชาต\-ปจฺจเยน ปจฺจโย. ปฏิสนฺธิ\-กฺขเณ วตฺถุ สวิตกฺก\-สวิจารานํ ขนฺธานํ วิตกฺกสฺส จ สหชาต\-ปจฺจเยน ปจฺจโย. (5)}

\proseitem{88}{สวิตกฺก\-สวิจาโร จ อวิตกฺก\-อวิจาโร จ ธมฺมา สวิตกฺก\-สวิจารสฺส ธมฺมสฺส สหชาต\-ปจฺจเยน ปจฺจโย. ปฏิสนฺธิ\-กฺขเณ สวิตกฺก\-สวิจาโร เอโก ขนฺโธ จ วตฺถุ จ ติณฺณนฺนํ ขนฺธานํ สหชาต\-ปจฺจเยน ปจฺจโย ฯเปฯ เทฺว ขนฺธา จ วตฺถุ จ ทฺวินฺนํ ขนฺธานํ สหชาต\-ปจฺจเยน ปจฺจโย. (1)}

\prose{สวิตกฺก\-สวิจาโร จ อวิตกฺก\-อวิจาโร จ ธมฺมา อวิตกฺก\-วิจารมตฺตสฺส ธมฺมสฺส สหชาต\-ปจฺจเยน ปจฺจโย. ปฏิสนฺธิ\-กฺขเณ สวิตกฺก\-สวิจารา ขนฺธา จ วตฺถุ จ วิตกฺกสฺส สหชาต\-ปจฺจเยน ปจฺจโย. (2)}

\prose{สวิตกฺก\-สวิจาโร จ อวิตกฺก\-อวิจาโร จ ธมฺมา อวิตกฺก\-อวิจารสฺส ธมฺมสฺส สหชาต\-ปจฺจเยน ปจฺจโย. สวิตกฺก\-สวิจารา ขนฺธา จ มหาภูตา จ จิตฺต\-สมุฏฺฐานานํ รูปานํ สหชาต\-ปจฺจเยน ปจฺจโย. ปฏิสนฺธิ\-กฺขเณ สวิตกฺก\-สวิจารา ขนฺธา จ มหาภูตา จ กฏตฺตารูปานํ สหชาต\-ปจฺจเยน ปจฺจโย. (3)}

\prose{\csromanfolio{61}สวิตกฺก\-สวิจาโร จ อวิตกฺก\-อวิจาโร จ ธมฺมา สวิตกฺก\-สวิจารสฺส จ อวิตกฺก\-วิจารมตฺตสฺส จ ธมฺมสฺส สหชาต\-ปจฺจเยน ปจฺจโย. ปฏิสนฺธิ\-กฺขเณ สวิตกฺก\-สวิจาโร เอโก ขนฺโธ จ วตฺถุ จ ติณฺณนฺนํ ขนฺธานํ วิตกฺกสฺส จ ฯเปฯ เทฺว ขนฺธา จ วตฺถุ จ ทฺวินฺนํ ขนฺธานํ วิตกฺกสฺส จ สหชาต\-ปจฺจเยน ปจฺจโย. (4)}

\proseitem{89}{อวิตกฺก\-วิจารมตฺโต จ อวิตกฺก\-อวิจาโร จ ธมฺมา สวิตกฺก\-สวิจารสฺส ธมฺมสฺส สหชาต\-ปจฺจเยน ปจฺจโย. ปฏิสนฺธิ\-กฺขเณ วิตกฺโก จ วตฺถุ จ สวิตกฺก\-สวิจารานํ ขนฺธานํ สหชาต\-ปจฺจเยน ปจฺจโย. (1)}

\prose{อวิตกฺก\-วิจารมตฺโต จ อวิตกฺก\-อวิจาโร จ ธมฺมา อวิตกฺก\-วิจารมตฺตสฺส ธมฺมสฺส สหชาต\-ปจฺจเยน ปจฺจโย. อวิตกฺก\-วิจารมตฺโต เอโก ขนฺโธ จ วิจาโร จ ติณฺณนฺนํ ขนฺธานํ ฯเปฯ เทฺว ขนฺธา จ วิจาโร จ ทฺวินฺนํ ขนฺธานํ สหชาต\-ปจฺจเยน ปจฺจโย. ปฏิสนฺธิ\-กฺขเณ อวิตกฺก\-วิจารมตฺโต เอโก ขนฺโธ จ วิจาโร จ ติณฺณนฺนํ ขนฺธานํ ฯเปฯ เทฺว ขนฺธา จ วิจาโร จ ทฺวินฺนํ ขนฺธานํ ฯเปฯ. ปฏิสนฺธิ\-กฺขเณ อวิตกฺก\-วิจารมตฺโต เอโก ขนฺโธ จ วตฺถุ จ ติณฺณนฺนํ ขนฺธานํ สหชาต\-ปจฺจเยน ปจฺจโย. (2)}

\prose{อวิตกฺก\-วิจารมตฺโต จ อวิตกฺก\-อวิจาโร จ ธมฺมา อวิตกฺก\-อวิจารสฺส ธมฺมสฺส สหชาต\-ปจฺจเยน ปจฺจโย. อวิตกฺก\-วิจารมตฺตา ขนฺธา จ วิจาโร จ จิตฺต\-สมุฏฺฐานานํ รูปานํ สหชาต\-ปจฺจเยน ปจฺจโย. อวิตกฺก\-วิจารมตฺตา ขนฺธา จ มหาภูตา จ จิตฺต\-สมุฏฺฐานานํ รูปานํ ฯเปฯ. วิตกฺโก จ มหาภูตา จ จิตฺต\-สมุฏฺฐานานํ รูปานํ ฯเปฯ. ปฏิสนฺธิ\-กฺขเณ อวิตกฺก\-วิจารมตฺตา ขนฺธา จ วิจาโร จ กฏตฺตารูปานํ ฯเปฯ. ปฏิสนฺธิ\-กฺขเณ อวิตกฺก\-วิจารมตฺตา ขนฺธา จ มหาภูตา จ กฏตฺตารูปานํ ฯเปฯ. ปฏิสนฺธิ\-กฺขเณ วิตกฺโก จ มหาภูตา จ กฏตฺตารูปานํ ฯเปฯ. ปฏิสนฺธิ\-กฺขเณ อวิตกฺก\-วิจารมตฺตา ขนฺธา จ วตฺถุ จ วิจารสฺส สหชาต\-ปจฺจเยน ปจฺจโย. (3)}

\prose{อวิตกฺก\-วิจารมตฺโต จ อวิตกฺก\-อวิจาโร จ ธมฺมา อวิตกฺก\-วิจารมตฺตสฺส จ อวิตกฺก\-อวิจารสฺส จ ธมฺมสฺส สหชาต\-ปจฺจเยน ปจฺจโย. อวิตกฺก\-วิจารมตฺโต เอโก ขนฺโธ จ วิจาโร จ ติณฺณนฺนํ ขนฺธานํ จิตฺตสมุฏฺฐานานญฺจ รูปานํ สหชาต\-ปจฺจเยน ปจฺจโย. ตโย ขนฺธา จ วิจาโร จ เอกสฺส ขนฺธสฺส จิตฺตสมุฏฺฐานานญฺจ รูปานํ สหชาต\-ปจฺจเยน ปจฺจโย. \csromanfolio{62}เทฺว ขนฺธา จ วิจาโร จ ทฺวินฺนํ ขนฺธานํ จิตฺตสมุฏฺฐานานญฺจ รูปานํ สหชาต\-ปจฺจเยน ปจฺจโย. ปฏิสนฺธิ\-กฺขเณ อวิตกฺก\-วิจารมตฺโต เอโก ขนฺโธ จ วิจาโร จ ติณฺณนฺนํ ขนฺธานํ กฏตฺตา จ รูปานํ สหชาต\-ปจฺจเยน ปจฺจโย ฯเปฯ เทฺว ขนฺธา ฯเปฯ. ปฏิสนฺธิ\-กฺขเณ อวิตกฺก\-วิจารมตฺโต เอโก ขนฺโธ จ วตฺถุ จ ติณฺณนฺนํ ขนฺธานํ วิจารสฺส จ สหชาต\-ปจฺจเยน ปจฺจโย ฯเปฯ เทฺว ขนฺธา จ วตฺถุ จ ทฺวินฺนํ ขนฺธานํ วิจารสฺส จ สหชาต\-ปจฺจเยน ปจฺจโย. (4)}

\proseitem{90}{สวิตกฺก\-สวิจาโร จ อวิตกฺก\-วิจารมตฺโต จ ธมฺมา สวิตกฺก\-สวิจารสฺส ธมฺมสฺส สหชาต\-ปจฺจเยน ปจฺจโย. สวิตกฺก\-สวิจาโร เอโก ขนฺโธ จ วิตกฺโก จ ติณฺณนฺนํ ขนฺธานํ สหชาต\-ปจฺจเยน ปจฺจโย ฯเปฯ เทฺว ขนฺธา จ วิตกฺโก จ ทฺวินฺนํ ขนฺธานํ สหชาต\-ปจฺจเยน ปจฺจโย. ปฏิสนฺธิ\-กฺขเณ ฯเปฯ. (1)}

\prose{สวิตกฺก\-สวิจาโร จ อวิตกฺก\-วิจารมตฺโต จ ธมฺมา อวิตกฺก\-อวิจารสฺส ธมฺมสฺส สหชาต\-ปจฺจเยน ปจฺจโย. สวิตกฺก\-สวิจารา ขนฺธา จ วิตกฺโก จ จิตฺต\-สมุฏฺฐานานํ รูปานํ สหชาต\-ปจฺจเยน ปจฺจโย. ปฏิสนฺธิ\-กฺขเณ ฯเปฯ. (2)}

\prose{สวิตกฺก\-สวิจาโร จ อวิตกฺก\-วิจารมตฺโต จ ธมฺมา สวิตกฺก\-สวิจารสฺส จ อวิตกฺก\-อวิจารสฺส จ ธมฺมสฺส สหชาต\-ปจฺจเยน ปจฺจโย. สวิตกฺก\-สวิจาโร เอโก ขนฺโธ จ วิตกฺโก จ ติณฺณนฺนํ ขนฺธานํ จิตฺตสมุฏฺฐานานญฺจ รูปานํ สหชาต\-ปจฺจเยน ปจฺจโย ฯเปฯ เทฺว ขนฺธา จ วิตกฺโก จ ทฺวินฺนํ ขนฺธานํ จิตฺตสมุฏฺฐานานญฺจ รูปานํ สหชาต\-ปจฺจเยน ปจฺจโย. ปฏิสนฺธิ\-กฺขเณ ฯเปฯ. (3)}

\proseitem{91}{สวิตกฺก\-สวิจาโร จ อวิตกฺก\-วิจารมตฺโต จ อวิตกฺก\-อวิจาโร จ ธมฺมา สวิตกฺก\-สวิจารสฺส ธมฺมสฺส สหชาต\-ปจฺจเยน ปจฺจโย. ปฏิสนฺธิ\-กฺขเณ สวิตกฺก\-สวิจาโร เอโก ขนฺโธ จ วิตกฺโก จ วตฺถุ จ ติณฺณนฺนํ ขนฺธานํ สหชาต\-ปจฺจเยน ปจฺจโย ฯเปฯ เทฺว ขนฺธา จ วิตกฺโก จ วตฺถุ จ ทฺวินฺนํ ขนฺธานํ สหชาต\-ปจฺจเยน ปจฺจโย. (1)}

\prose{สวิตกฺก\-สวิจาโร จ อวิตกฺก\-วิจารมตฺโต จ อวิตกฺก\-อวิจาโร จ ธมฺมา อวิตกฺก\-อวิจารสฺส ธมฺมสฺส สหชาต\-ปจฺจเยน ปจฺจโย. สวิตกฺก\-สวิจารา ขนฺธา จ วิตกฺโก จ มหาภูตา จ จิตฺต\-สมุฏฺฐานานํ รูปานํ \csromanfolio{63}สหชาต\-ปจฺจเยน ปจฺจโย. ปฏิสนฺธิ\-กฺขเณ สวิตกฺก\-สวิจารา ขนฺธา จ วิตกฺโก จ มหาภูตา จ กฏตฺตารูปานํ สหชาต\-ปจฺจเยน ปจฺจโย. (2)}

\csromanheader{2}{\textbf{อญฺญมญฺญ}}
\proseitem{92}{สวิตกฺก\-สวิจาโร ธมฺโม สวิตกฺก\-สวิจารสฺส ธมฺมสฺส อญฺญมญฺญ\-ปจฺจเยน ปจฺจโย. สวิตกฺก\-สวิจาโร เอโก ขนฺโธ ติณฺณนฺนํ ขนฺธานํ อญฺญมญฺญ\-ปจฺจเยน ปจฺจโย ฯเปฯ. ปฏิสนฺธิ\-กฺขเณ สวิตกฺก\-สวิจาโร เอโก ขนฺโธ ติณฺณนฺนํ ขนฺธานํ อญฺญมญฺญ\-ปจฺจเยน ปจฺจโย ฯเปฯ เทฺว ขนฺธา ทฺวินฺนํ ขนฺธานํ อญฺญมญฺญ\-ปจฺจเยน ปจฺจโย. (1)}

\prose{สวิตกฺก\-สวิจาโร ธมฺโม อวิตกฺก\-วิจารมตฺตสฺส ธมฺมสฺส อญฺญมญฺญ\-ปจฺจเยน ปจฺจโย. สวิตกฺก\-สวิจารา ขนฺธา วิตกฺกสฺส อญฺญมญฺญ\-ปจฺจเยน ปจฺจโย. ปฏิสนฺธิ\-กฺขเณ ฯเปฯ. (2)}

\prose{สวิตกฺก\-สวิจาโร ธมฺโม อวิตกฺก\-อวิจารสฺส ธมฺมสฺส อญฺญมญฺญ\-ปจฺจเยน ปจฺจโย. ปฏิสนฺธิ\-กฺขเณ สวิตกฺก\-สวิจารา ขนฺธา วตฺถุสฺส อญฺญมญฺญ\-ปจฺจเยน ปจฺจโย. (3)}

\prose{สวิตกฺก\-สวิจาโร ธมฺโม สวิตกฺก\-สวิจารสฺส จ อวิตกฺก\-อวิจารสฺส จ ธมฺมสฺส อญฺญมญฺญ\-ปจฺจเยน ปจฺจโย. ปฏิสนฺธิ\-กฺขเณ สวิตกฺก\-สวิจาโร เอโก ขนฺโธ ติณฺณนฺนํ ขนฺธานํ วตฺถุสฺส จ อญฺญมญฺญ\-ปจฺจเยน ปจฺจโย ฯเปฯ เทฺว ขนฺธา ทฺวินฺนํ ขนฺธานํ วตฺถุสฺส จ อญฺญมญฺญ\-ปจฺจเยน ปจฺจโย. (4)}

\prose{สวิตกฺก\-สวิจาโร ธมฺโม อวิตกฺก\-วิจารมตฺตสฺส จ อวิตกฺก\-อวิจารสฺส จ ธมฺมสฺส อญฺญมญฺญ\-ปจฺจเยน ปจฺจโย. ปฏิสนฺธิ\-กฺขเณ สวิตกฺก\-สวิจารา ขนฺธา วิตกฺกสฺส จ วตฺถุสฺส จ อญฺญมญฺญ\-ปจฺจเยน ปจฺจโย. (5)}

\prose{สวิตกฺก\-สวิจาโร ธมฺโม สวิตกฺก\-สวิจารสฺส จ อวิตกฺก\-วิจารมตฺตสฺส จ ธมฺมสฺส อญฺญมญฺญ\-ปจฺจเยน ปจฺจโย. สวิตกฺก\-สวิจาโร เอโก ขนฺโธ ติณฺณนฺนํ ขนฺธานํ วิตกฺกสฺส จ อญฺญมญฺญ\-ปจฺจเยน ปจฺจโย ฯเปฯ เทฺว ขนฺธา ทฺวินฺนํ ขนฺธานํ วิตกฺกสฺส จ อญฺญมญฺญ\-ปจฺจเยน ปจฺจโย. ปฏิสนฺธิ\-กฺขเณ ฯเปฯ. (6)}

\prose{สวิตกฺก\-สวิจาโร ธมฺโม สวิตกฺก\-สวิจารสฺส จ อวิตกฺก\-วิจารมตฺตสฺส จ อวิตกฺก\-อวิจารสฺส จ ธมฺมสฺส อญฺญมญฺญ\-ปจฺจเยน ปจฺจโย. ปฏิสนฺธิ\-กฺขเณ สวิตกฺก\-สวิจาโร เอโก ขนฺโธ ติณฺณนฺนํ ขนฺธานํ วิตกฺกสฺส \csromanfolio{64}จ วตฺถุสฺส จ อญฺญมญฺญ\-ปจฺจเยน ปจฺจโย ฯเปฯ เทฺว ขนฺธา ทฺวินฺนํ ขนฺธานํ วิตกฺกสฺส จ วตฺถุสฺส จ อญฺญมญฺญ\-ปจฺจเยน ปจฺจโย. (7)}

\proseitem{93}{อวิตกฺก\-วิจารมตฺโต ธมฺโม อวิตกฺก\-วิจารมตฺตสฺส ธมฺมสฺส อญฺญมญฺญ\-ปจฺจเยน ปจฺจโย. อวิตกฺก\-วิจารมตฺโต เอโก ขนฺโธ ติณฺณนฺนํ ขนฺธานํ อญฺญมญฺญ\-ปจฺจเยน ปจฺจโย ฯเปฯ เทฺว ขนฺธา ทฺวินฺนํ ขนฺธานํ อญฺญมญฺญ\-ปจฺจเยน ปจฺจโย. ปฏิสนฺธิ\-กฺขเณ ฯเปฯ. (1)}

\prose{อวิตกฺก\-วิจารมตฺโต ธมฺโม สวิตกฺก\-สวิจารสฺส ธมฺมสฺส อญฺญมญฺญ\-ปจฺจเยน ปจฺจโย. วิตกฺโก สวิตกฺก\-สวิจารานํ ขนฺธานํ อญฺญมญฺญ\-ปจฺจเยน ปจฺจโย. ปฏิสนฺธิ\-กฺขเณ ฯเปฯ. (2)}

\prose{อวิตกฺก\-วิจารมตฺโต ธมฺโม อวิตกฺก\-อวิจารสฺส ธมฺมสฺส อญฺญมญฺญ\-ปจฺจเยน ปจฺจโย. อวิตกฺก\-วิจารมตฺตา ขนฺธา วิจารสฺส อญฺญมญฺญ\-ปจฺจเยน ปจฺจโย. ปฏิสนฺธิ\-กฺขเณ อวิตกฺก\-วิจารมตฺตา ขนฺธา วิจารสฺส จ วตฺถุสฺส จ อญฺญมญฺญ\-ปจฺจเยน ปจฺจโย. ปฏิสนฺธิ\-กฺขเณ วิตกฺโก วตฺถุสฺส อญฺญมญฺญ\-ปจฺจเยน ปจฺจโย. (3)}

\prose{อวิตกฺก\-วิจารมตฺโต ธมฺโม สวิตกฺก\-สวิจารสฺส จ อวิตกฺก\-อวิจารสฺส จ ธมฺมสฺส อญฺญมญฺญ\-ปจฺจเยน ปจฺจโย. ปฏิสนฺธิ\-กฺขเณ วิตกฺโก สวิตกฺก\-สวิจารานํ ขนฺธานํ วตฺถุสฺส จ อญฺญมญฺญ\-ปจฺจเยน ปจฺจโย. (4)}

\prose{อวิตกฺก\-วิจารมตฺโต ธมฺโม อวิตกฺก\-วิจารมตฺตสฺส จ อวิตกฺก\-อวิจารสฺส จ ธมฺมสฺส อญฺญมญฺญ\-ปจฺจเยน ปจฺจโย. อวิตกฺก\-วิจารมตฺโต เอโก ขนฺโธ ติณฺณนฺนํ ขนฺธานํ วิจารสฺส จ อญฺญมญฺญ\-ปจฺจเยน ปจฺจโย ฯเปฯ เทฺว ขนฺธา ทฺวินฺนํ ขนฺธานํ วิจารสฺส จ อญฺญมญฺญ\-ปจฺจเยน ปจฺจโย. ปฏิสนฺธิ\-กฺขเณ อวิตกฺก\-วิจารมตฺโต เอโก ขนฺโธ ติณฺณนฺนํ ขนฺธานํ วิจารสฺส จ วตฺถุสฺส จ อญฺญมญฺญ\-ปจฺจเยน ปจฺจโย ฯเปฯ เทฺว ขนฺธา ทฺวินฺนํ ขนฺธานํ วิจารสฺส จ วตฺถุสฺส จ อญฺญมญฺญ\-ปจฺจเยน ปจฺจโย. (5)}

\proseitem{94}{อวิตกฺก\-อวิจาโร ธมฺโม อวิตกฺก\-อวิจารสฺส ธมฺมสฺส อญฺญมญฺญ\-ปจฺจเยน ปจฺจโย. อวิตกฺก\-อวิจาโร เอโก ขนฺโธ ติณฺณนฺนํ ขนฺธานํ อญฺญมญฺญ\-ปจฺจเยน ปจฺจโย ฯเปฯ เทฺว ขนฺธา ทฺวินฺนํ ขนฺธานํ อญฺญมญฺญ\-ปจฺจเยน ปจฺจโย. ปฏิสนฺธิ\-กฺขเณ อวิตกฺก\-อวิจาโร เอโก ขนฺโธ ติณฺณนฺนํ ขนฺธานํ \csromanfolio{65}วตฺถุสฺส จ อญฺญมญฺญ\-ปจฺจเยน ปจฺจโย ฯเปฯ เทฺว ขนฺธา ทฺวินฺนํ ขนฺธานํ วตฺถุสฺส จ อญฺญมญฺญ\-ปจฺจเยน ปจฺจโย. ขนฺธา วตฺถุสฺส ฯเปฯ. วตฺถุ ขนฺธานํ ฯเปฯ. วิจาโร วตฺถุสฺส ฯเปฯ. วตฺถุ วิจารสฺส ฯเปฯ. เอกํ มหาภูตํ ติณฺณนฺนํ มหาภูตานํ ฯเปฯ. พาหิรํ. อาหารสมุฏฺฐานํ. อุตุสมุฏฺฐานํ. อสญฺญสตฺตานํ เอกํ มหาภูตํ ฯเปฯ. (1)}

\prose{อวิตกฺก\-อวิจาโร ธมฺโม สวิตกฺก\-สวิจารสฺส ธมฺมสฺส อญฺญมญฺญ\-ปจฺจเยน ปจฺจโย. ปฏิสนฺธิ\-กฺขเณ วตฺถุ สวิตกฺก\-สวิจารานํ ขนฺธานํ อญฺญมญฺญ\-ปจฺจเยน ปจฺจโย. (2)}

\prose{อวิตกฺก\-อวิจาโร ธมฺโม อวิตกฺก\-วิจารมตฺตสฺส ธมฺมสฺส อญฺญมญฺญ\-ปจฺจเยน ปจฺจโย. วิจาโร อวิตกฺก\-วิจารมตฺตานํ ขนฺธานํ อญฺญมญฺญ\-ปจฺจเยน ปจฺจโย. ปฏิสนฺธิ\-กฺขเณ วิจาโร อวิตกฺก\-วิจารมตฺตานํ ขนฺธานํ อญฺญมญฺญ\-ปจฺจเยน ปจฺจโย. ปฏิสนฺธิ\-กฺขเณ วตฺถุ อวิตกฺก\-วิจารมตฺตานํ ขนฺธานํ อญฺญมญฺญ\-ปจฺจเยน ปจฺจโย. ปฏิสนฺธิ\-กฺขเณ วตฺถุ วิตกฺกสฺส อญฺญมญฺญ\-ปจฺจเยน ปจฺจโย. (3)}

\prose{อวิตกฺก\-อวิจาโร ธมฺโม อวิตกฺก\-วิจารมตฺตสฺส จ อวิตกฺก\-อวิจารสฺส จ ธมฺมสฺส อญฺญมญฺญ\-ปจฺจเยน ปจฺจโย. ปฏิสนฺธิ\-กฺขเณ วิจาโร อวิตกฺก\-วิจารมตฺตานํ ขนฺธานํ วตฺถุสฺส จ อญฺญมญฺญ\-ปจฺจเยน ปจฺจโย. ปฏิสนฺธิ\-กฺขเณ วตฺถุ อวิตกฺก\-วิจารมตฺตานํ ขนฺธานํ วิจารสฺส จ อญฺญมญฺญ\-ปจฺจเยน ปจฺจโย. (4)}

\prose{อวิตกฺก\-อวิจาโร ธมฺโม สวิตกฺก\-สวิจารสฺส จ อวิตกฺก\-วิจารมตฺตสฺส จ ธมฺมสฺส อญฺญมญฺญ\-ปจฺจเยน ปจฺจโย. ปฏิสนฺธิ\-กฺขเณ วตฺถุ สวิตกฺก\-สวิจารานํ ขนฺธานํ วิตกฺกสฺส จ อญฺญมญฺญ\-ปจฺจเยน ปจฺจโย. (5)}

\proseitem{95}{สวิตกฺก\-สวิจาโร จ อวิตกฺก\-อวิจาโร จ ธมฺมา สวิตกฺก\-สวิจารสฺส ธมฺมสฺส อญฺญมญฺญ\-ปจฺจเยน ปจฺจโย. ปฏิสนฺธิ\-กฺขเณ สวิตกฺก\-สวิจาโร เอโก ขนฺโธ จ วตฺถุ จ ติณฺณนฺนํ ขนฺธานํ ฯเปฯ เทฺว ขนฺธา จ วตฺถุ จ ทฺวินฺนํ ขนฺธานํ อญฺญมญฺญ\-ปจฺจเยน ปจฺจโย. (1)}

\prose{\csromanfolio{66}สวิตกฺก\-สวิจาโร จ อวิตกฺก\-อวิจาโร จ ธมฺมา อวิตกฺก\-วิจารมตฺตสฺส ธมฺมสฺส อญฺญมญฺญ\-ปจฺจเยน ปจฺจโย. ปฏิสนฺธิ\-กฺขเณ สวิตกฺก\-สวิจารา ขนฺธา จ วตฺถุ จ วิตกฺกสฺส อญฺญมญฺญ\-ปจฺจเยน ปจฺจโย. (2)}

\prose{สวิตกฺก\-สวิจาโร จ อวิตกฺก\-อวิจาโร จ ธมฺมา สวิตกฺก\-สวิจารสฺส จ อวิตกฺก\-วิจารมตฺตสฺส จ ธมฺมสฺส อญฺญมญฺญ\-ปจฺจเยน ปจฺจโย. ปฏิสนฺธิ\-กฺขเณ สวิตกฺก\-สวิจาโร เอโก ขนฺโธ จ วตฺถุ จ ติณฺณนฺนํ ขนฺธานํ วิตกฺกสฺส จ อญฺญมญฺญ\-ปจฺจเยน ปจฺจโย ฯเปฯ เทฺว ขนฺธา จ วตฺถุ จ ทฺวินฺนํ ขนฺธานํ วิตกฺกสฺส จ อญฺญมญฺญ\-ปจฺจเยน ปจฺจโย. (3)}

\proseitem{96}{อวิตกฺก\-วิจารมตฺโต จ อวิตกฺก\-อวิจาโร จ ธมฺมา สวิตกฺก\-สวิจารสฺส ธมฺมสฺส อญฺญมญฺญ\-ปจฺจเยน ปจฺจโย. ปฏิสนฺธิ\-กฺขเณ วิตกฺโก จ วตฺถุ จ สวิตกฺก\-สวิจารานํ ขนฺธานํ อญฺญมญฺญ\-ปจฺจเยน ปจฺจโย. (1)}

\prose{อวิตกฺก\-วิจารมตฺโต จ อวิตกฺก\-อวิจาโร จ ธมฺมา อวิตกฺก\-วิจารมตฺตสฺส ธมฺมสฺส อญฺญมญฺญ\-ปจฺจเยน ปจฺจโย. อวิตกฺก\-วิจารมตฺโต เอโก ขนฺโธ จ วิจาโร จ ติณฺณนฺนํ ขนฺธานํ อญฺญมญฺญ\-ปจฺจเยน ปจฺจโย ฯเปฯ เทฺว ขนฺธา จ วิจาโร จ ทฺวินฺนํ ขนฺธานํ อญฺญมญฺญ\-ปจฺจเยน ปจฺจโย. ปฏิสนฺธิ\-กฺขเณ อวิตกฺก\-วิจารมตฺโต เอโก ขนฺโธ จ วิจาโร จ วตฺถุ จ ติณฺณนฺนํ ขนฺธานํ ฯเปฯ เทฺว ขนฺธา จ วิจาโร จ วตฺถุ จ ทฺวินฺนํ ขนฺธานํ อญฺญมญฺญ\-ปจฺจเยน ปจฺจโย. (2)}

\prose{อวิตกฺก\-วิจารมตฺโต จ อวิตกฺก\-อวิจาโร จ ธมฺมา อวิตกฺก\-อวิจารสฺส ธมฺมสฺส อญฺญมญฺญ\-ปจฺจเยน ปจฺจโย. ปฏิสนฺธิ\-กฺขเณ อวิตกฺก\-วิจารมตฺตา ขนฺธา จ วิจาโร จ วตฺถุสฺส อญฺญมญฺญ\-ปจฺจเยน ปจฺจโย. ปฏิสนฺธิ\-กฺขเณ อวิตกฺก\-วิจารมตฺตา ขนฺธา จ วตฺถุ จ วิจารสฺส อญฺญมญฺญ\-ปจฺจเยน ปจฺจโย. (3)}

\prose{อวิตกฺก\-วิจารมตฺโต จ อวิตกฺก\-อวิจาโร จ ธมฺมา อวิตกฺก\-วิจารมตฺตสฺส จ อวิตกฺก\-อวิจารสฺส จ ธมฺมสฺส อญฺญมญฺญ\-ปจฺจเยน ปจฺจโย. ปฏิสนฺธิ\-กฺขเณ อวิตกฺก\-วิจารมตฺโต เอโก ขนฺโธ จ วิจาโร จ ติณฺณนฺนํ ขนฺธานํ วตฺถุสฺส จ อญฺญมญฺญ\-ปจฺจเยน ปจฺจโย ฯเปฯ เทฺว ขนฺธา จ วิจาโร จ ทฺวินฺนํ ขนฺธานํ วตฺถุสฺส จ อญฺญมญฺญ\-ปจฺจเยน ปจฺจโย. ปฏิสนฺธิ\-กฺขเณ อวิตกฺก\-วิจารมตฺโต เอโก ขนฺโธ จ วตฺถุ จ ติณฺณนฺนํ ขนฺธานํ วิจารสฺส จ \csromanfolio{67}อญฺญมญฺญ\-ปจฺจเยน ปจฺจโย ฯเปฯ เทฺว ขนฺธา จ วตฺถุ จ ทฺวินฺนํ ขนฺธานํ วิจารสฺส จ อญฺญมญฺญ\-ปจฺจเยน ปจฺจโย. (4)}

\proseitem{97}{สวิตกฺก\-สวิจาโร จ อวิตกฺก\-วิจารมตฺโต จ ธมฺมา สวิตกฺก\-สวิจารสฺส ธมฺมสฺส อญฺญมญฺญ\-ปจฺจเยน ปจฺจโย. สวิตกฺก\-สวิจาโร เอโก ขนฺโธ จ วิตกฺโก จ ติณฺณนฺนํ ขนฺธานํ อญฺญมญฺญ\-ปจฺจเยน ปจฺจโย ฯเปฯ เทฺว ขนฺธา จ วิตกฺโก จ ทฺวินฺนํ ขนฺธานํ อญฺญมญฺญ\-ปจฺจเยน ปจฺจโย. ปฏิสนฺธิ\-กฺขเณ ฯเปฯ. (1)}

\prose{สวิตกฺก\-สวิจาโร จ อวิตกฺก\-วิจารมตฺโต จ ธมฺมา อวิตกฺก\-อวิจารสฺส ธมฺมสฺส อญฺญมญฺญ\-ปจฺจเยน ปจฺจโย. ปฏิสนฺธิ\-กฺขเณ สวิตกฺก\-สวิจารา ขนฺธา จ วิตกฺโก จ วตฺถุสฺส อญฺญมญฺญ\-ปจฺจเยน ปจฺจโย. (2)}

\prose{สวิตกฺก\-สวิจาโร จ อวิตกฺก\-วิจารมตฺโต จ ธมฺมา สวิตกฺก\-สวิจารสฺส จ อวิตกฺก\-อวิจารสฺส จ ธมฺมสฺส อญฺญมญฺญ\-ปจฺจเยน ปจฺจโย. ปฏิสนฺธิ\-กฺขเณ สวิตกฺก\-สวิจาโร เอโก ขนฺโธ จ วิตกฺโก จ ติณฺณนฺนํ ขนฺธานํ วตฺถุสฺส จ อญฺญมญฺญ\-ปจฺจเยน ปจฺจโย ฯเปฯ เทฺว ขนฺธา จ วิตกฺโก จ ทฺวินฺนํ ขนฺธานํ วตฺถุสฺส จ อญฺญมญฺญ\-ปจฺจเยน ปจฺจโย. (3)}

\prose{สวิตกฺก\-สวิจาโร จ อวิตกฺก\-วิจารมตฺโต จ อวิตกฺก\-อวิจาโร จ ธมฺมา สวิตกฺก\-สวิจารสฺส ธมฺมสฺส อญฺญมญฺญ\-ปจฺจเยน ปจฺจโย. ปฏิสนฺธิ\-กฺขเณ สวิตกฺก\-สวิจาโร เอโก ขนฺโธ จ วิตกฺโก จ วตฺถุ จ ติณฺณนฺนํ ขนฺธานํ อญฺญมญฺญ\-ปจฺจเยน ปจฺจโย ฯเปฯ เทฺว ขนฺธา จ วิตกฺโก จ วตฺถุ จ ทฺวินฺนํ ขนฺธานํ อญฺญมญฺญ\-ปจฺจเยน ปจฺจโย. (1)}

\csromanheader{2}{\textbf{นิสฺสย}}
\proseitem{98}{สวิตกฺก\-สวิจาโร ธมฺโม สวิตกฺก\-สวิจารสฺส ธมฺมสฺส นิสฺสย\-ปจฺจเยน ปจฺจโย. สวิตกฺก\-สวิจาโร เอโก ขนฺโธ ติณฺณนฺนํ ขนฺธานํ. (สํขิตฺตํ.) สตฺต.}

\prose{อวิตกฺก\-วิจารมตฺโต ธมฺโม อวิตกฺก\-วิจารมตฺตสฺส ธมฺมสฺส นิสฺสย ปจฺจเยน ปจฺจโย. (สํขิตฺตํ.) ปญฺจ.}

\prose{อวิตกฺก\-อวิจาโร ธมฺโม อวิตกฺก\-อวิจารสฺส ธมฺมสฺส นิสฺสย\-ปจฺจเยน ปจฺจโย. อวิตกฺก\-อวิจาโร เอโก ขนฺโธ ติณฺณนฺนํ ขนฺธานํ \csromanfolio{68}จิตฺตสมุฏฺฐานานญฺจ รูปานํ นิสฺสย\-ปจฺจเยน ปจฺจโย ฯเปฯ เทฺว ขนฺธา ทฺวินฺนํ ขนฺธานํ จิตฺตสมุฏฺฐานานญฺจ รูปานํ นิสฺสย\-ปจฺจเยน ปจฺจโย. วิจาโร จิตฺต\-สมุฏฺฐานานํ รูปานํ นิสฺสย\-ปจฺจเยน ปจฺจโย. ปฏิสนฺธิ\-กฺขเณ อวิตกฺก\-อวิจาโร เอโก ขนฺโธ ติณฺณนฺนํ ขนฺธานํ กฏตฺตา จ รูปานํ นิสฺสย\-ปจฺจเยน ปจฺจโย ฯเปฯ เทฺว ขนฺธา ทฺวินฺนํ ขนฺธานํ ฯเปฯ. อสญฺญสตฺตานํ เอกํ มหาภูตํ ฯเปฯ. จกฺขายตนํ จกฺขุ\-วิญฺญาณสฺส ฯเปฯ. กายายตนํ กายวิญฺญาณสฺส ฯเปฯ. วตฺถุ อวิตกฺก\-อวิจารานํ ขนฺธานํ วิจารสฺส จ นิสฺสย\-ปจฺจเยน ปจฺจโย. (1)}

\prose{อวิตกฺก\-อวิจาโร ธมฺโม สวิตกฺก\-สวิจารสฺส ธมฺมสฺส นิสฺสย\-ปจฺจเยน ปจฺจโย. วตฺถุ สวิตกฺก\-สวิจารานํ ขนฺธานํ นิสฺสย\-ปจฺจเยน ปจฺจโย. ปฏิสนฺธิ\-กฺขเณ วตฺถุ ฯเปฯ. (2)}

\prose{อวิตกฺก\-อวิจาโร ธมฺโม อวิตกฺก\-วิจารมตฺตสฺส ธมฺมสฺส นิสฺสย\-ปจฺจเยน ปจฺจโย. วิจาโร อวิตกฺก\-วิจารมตฺตานํ ขนฺธานํ นิสฺสย\-ปจฺจเยน ปจฺจโย. วตฺถุ อวิตกฺก\-วิจารมตฺตานํ ขนฺธานํ วิตกฺกสฺส จ นิสฺสย\-ปจฺจเยน ปจฺจโย. ปฏิสนฺธิ\-กฺขเณ วิจาโร ฯเปฯ. (3)}

\prose{อวิตกฺก\-อวิจาโร ธมฺโม อวิตกฺก\-วิจารมตฺตสฺส จ อวิตกฺก\-อวิจารสฺส จ ธมฺมสฺส นิสฺสย\-ปจฺจเยน ปจฺจโย. วิจาโร อวิตกฺก\-วิจารมตฺตานํ ขนฺธานํ จิตฺตสมุฏฺฐานานญฺจ รูปานํ ฯเปฯ. วตฺถุ อวิตกฺก\-วิจารมตฺตานํ ขนฺธานํ วิจารสฺส จ นิสฺสย\-ปจฺจเยน ปจฺจโย. ปฏิสนฺธิ\-กฺขเณ วิจาโร ฯเปฯ. (4)}

\prose{อวิตกฺก\-อวิจาโร ธมฺโม สวิตกฺก\-สวิจารสฺส จ อวิตกฺก\-วิจารมตฺตสฺส จ ธมฺมสฺส นิสฺสย\-ปจฺจเยน ปจฺจโย. วตฺถุ สวิตกฺก\-สวิจารานํ ขนฺธานํ วิตกฺกสฺส จ นิสฺสย\-ปจฺจเยน ปจฺจโย. ปฏิสนฺธิ\-กฺขเณ วตฺถุ ฯเปฯ. (5)}

\proseitem{99}{สวิตกฺก\-สวิจาโร จ อวิตกฺก\-อวิจาโร จ ธมฺมา สวิตกฺก\-สวิจารสฺส ธมฺมสฺส นิสฺสย\-ปจฺจเยน ปจฺจโย. สวิตกฺก\-สวิจาโร เอโก ขนฺโธ จ วตฺถุ จ ติณฺณนฺนํ ขนฺธานํ นิสฺสย\-ปจฺจเยน ปจฺจโย ฯเปฯ. (ปวตฺติปิ, ปฏิสนฺธิปิ ทีเปตพฺพา.) (1)}

\prose{\csromanfolio{69}สวิตกฺก\-สวิจาโร จ อวิตกฺก\-อวิจาโร จ ธมฺมา อวิตกฺก\-วิจารมตฺตสฺส ธมฺมสฺส นิสฺสย\-ปจฺจเยน ปจฺจโย. สวิตกฺก\-สวิจารา ขนฺธา จ วตฺถุ จ วิตกฺกสฺส ฯเปฯ. ปฏิสนฺธิ\-กฺขเณ ฯเปฯ. (2)}

\prose{สวิตกฺก\-สวิจาโร จ อวิตกฺก\-อวิจาโร จ ธมฺมา อวิตกฺก\-อวิจารสฺส ธมฺมสฺส นิสฺสย\-ปจฺจเยน ปจฺจโย. สวิตกฺก\-สวิจารา ขนฺธา จ มหาภูตา จ จิตฺต\-สมุฏฺฐานานํ รูปานํ นิสฺสย\-ปจฺจเยน ปจฺจโย. ปฏิสนฺธิ\-กฺขเณ ฯเปฯ. (3)}

\prose{สวิตกฺก\-สวิจาโร จ อวิตกฺก\-อวิจาโร จ ธมฺมา สวิตกฺก\-สวิจารสฺส จ อวิตกฺก\-วิจารมตฺตสฺส จ ธมฺมสฺส นิสฺสย\-ปจฺจเยน ปจฺจโย. สวิตกฺก\-สวิจาโร เอโก ขนฺโธ จ วตฺถุ จ ติณฺณนฺนํ ขนฺธานํ วิตกฺกสฺส จ นิสฺสย\-ปจฺจเยน ปจฺจโย ฯเปฯ. ปฏิสนฺธิ\-กฺขเณ ฯเปฯ. (4)}

\proseitem{100}{อวิตกฺก\-วิจารมตฺโต จ อวิตกฺก\-อวิจาโร จ ธมฺมา สวิตกฺก\-สวิจารสฺส ธมฺมสฺส นิสฺสย\-ปจฺจเยน ปจฺจโย. วิตกฺโก จ วตฺถุ จ สวิตกฺก\-สวิจารานํ ขนฺธานํ ฯเปฯ. ปฏิสนฺธิ\-กฺขเณ ฯเปฯ. (1)}

\prose{อวิตกฺก\-วิจารมตฺโต จ อวิตกฺก\-อวิจาโร จ ธมฺมา อวิตกฺก\-วิจารมตฺตสฺส ธมฺมสฺส นิสฺสย\-ปจฺจเยน ปจฺจโย. อวิตกฺก\-วิจารมตฺโต เอโก ขนฺโธ จ วิจาโร จ ติณฺณนฺนํ ขนฺธานํ ฯเปฯ. อวิตกฺก\-วิจารมตฺโต เอโก ขนฺโธ จ วตฺถุ จ ติณฺณนฺนํ ขนฺธานํ ฯเปฯ. ปฏิสนฺธิ\-กฺขเณ ฯเปฯ. (2)}

\prose{อวิตกฺก\-วิจารมตฺโต จ อวิตกฺก\-อวิจาโร จ ธมฺมา อวิตกฺก\-อวิจารสฺส ธมฺมสฺส นิสฺสย\-ปจฺจเยน ปจฺจโย. อวิตกฺก\-วิจารมตฺตา ขนฺธา จ วิจาโร จ จิตฺต\-สมุฏฺฐานานํ รูปานํ นิสฺสย\-ปจฺจเยน ปจฺจโย. อวิตกฺก\-วิจารมตฺตา ขนฺธา จ มหาภูตา จ จิตฺต\-สมุฏฺฐานานํ รูปานํ ฯเปฯ. วิตกฺโก จ มหาภูตา จ จิตฺต\-สมุฏฺฐานานํ รูปานํ ฯเปฯ. อวิตกฺก\-วิจารมตฺตา ขนฺธา จ วตฺถุ จ วิจารสฺส นิสฺสย\-ปจฺจเยน ปจฺจโย. (ปฏิสนฺธิกานิ จตฺตาริ. สํขิตฺตํ.) (3)}

\prose{อวิตกฺก\-วิจารมตฺโต จ อวิตกฺก\-อวิจาโร จ ธมฺมา อวิตกฺก\-วิจารมตฺตสฺส จ อวิตกฺก\-อวิจารสฺส จ ธมฺมสฺส นิสฺสย\-ปจฺจเยน ปจฺจโย. อวิตกฺก\-วิจารมตฺโต เอโก ขนฺโธ จ วิจาโร จ ติณฺณนฺนํ ขนฺธานํ จิตฺตสมุฏฺฐานานญฺจ รูปานํ นิสฺสย\-ปจฺจเยน ปจฺจโย ฯเปฯ เทฺว ขนฺธา จ วิจาโร \csromanfolio{70}จ ทฺวินฺนํ ขนฺธานํ จิตฺตสมุฏฺฐานานญฺจ รูปานํ นิสฺสย\-ปจฺจเยน ปจฺจโย. อวิตกฺก\-วิจารมตฺโต เอโก ขนฺโธ จ วตฺถุ จ ติณฺณนฺนํ ขนฺธานํ วิจารสฺส จ ฯเปฯ เทฺว ขนฺธา จ วตฺถุ จ ทฺวินฺนํ ขนฺธานํ วิจารสฺส จ นิสฺสย\-ปจฺจเยน ปจฺจโย. ปฏิสนฺธิ\-กฺขเณ ฯเปฯ. (4)}

\proseitem{101}{สวิตกฺก\-สวิจาโร จ อวิตกฺก\-วิจารมตฺโต จ ธมฺมา สวิตกฺก\-สวิจารสฺส ธมฺมสฺส ฯเปฯ อวิตกฺก\-อวิจารสฺส ธมฺมสฺส ฯเปฯ สวิตกฺก\-สวิจารสฺส จ อวิตกฺก\-อวิจารสฺส จ ธมฺมสฺส ฯเปฯ. ตีณิ.}

\prose{สวิตกฺก\-สวิจาโร จ อวิตกฺก\-วิจารมตฺโต จ อวิตกฺก\-อวิจาโร จ ธมฺมา สวิตกฺก\-สวิจารสฺส ธมฺมสฺส ฯเปฯ อวิตกฺก\-อวิจารสฺส ธมฺมสฺส นิสฺสย\-ปจฺจเยน ปจฺจโย. (เทฺว วารา วิตฺถาเรตพฺพา.)}

\csromanheader{2}{\textbf{อุปนิสฺสย}}
\proseitem{102}{สวิตกฺก\-สวิจาโร ธมฺโม สวิตกฺก\-สวิจารสฺส ธมฺมสฺส อุปนิสฺสย\-ปจฺจเยน ปจฺจโย. อารมฺมณู\-ปนิสฺสโย, อนนฺตรู\-ปนิสฺสโย, ปกตู\-ปนิสฺสโย ฯเปฯ. ปกตู\-ปนิสฺสโย\csromandash{}สวิตกฺก\-สวิจารํ สทฺธํ อุปนิสฺสาย ทานํ เทติ, สีลํ สมาทิยติ, อุโปสถกมฺมํ กโรติ, สวิตกฺก\-สวิจารํ ฌานํ อุปฺปาเทติ, วิปสฺสนํ อุปฺปาเทติ, มคฺคํ ฯเปฯ สมาปตฺติํ อุปฺปาเทติ. มานํ ชปฺเปติ, ทิฏฺฐิํ คณฺหาติ.  สวิตกฺก\-สวิจารํ สีลํ. สุตํ. จาคํ. ปญฺญํ. ราคํ. โทสํ. โมหํ. มานํ. ทิฏฺฐิํ. ปตฺถนํ อุปนิสฺสาย ทานํ เทติ, สีลํ สมาทิยติ, อุโปสถกมฺมํ กโรติ. สวิตกฺก\-สวิจารํ ฌานํ อุปฺปาเทติ, วิปสฺสนํ อุปฺปาเทติ, มคฺคํ ฯเปฯ สมาปตฺติํ ฯเปฯ. ปาณํ หนติ ฯเปฯ สํฆํ ภินฺทติ.  สวิตกฺก\-สวิจารา สทฺธา. สีลํ. สุตํ. จาโค. ปญฺญา. ราโค. โทโส. โมโห. มาโน. ทิฏฺฐิ. ปตฺถนา สวิตกฺก\-สวิจาราย สทฺธาย. สีลสฺส. สุตสฺส. จาคสฺส. ปญฺญาย. ราคสฺส. โทสสฺส. โมหสฺส. มานสฺส. ทิฏฺฐิยา. ปตฺถนาย อุปนิสฺสย\-ปจฺจเยน ปจฺจโย. (1)}

\prose{สวิตกฺก\-สวิจาโร ธมฺโม อวิตกฺก\-วิจารมตฺตสฺส ธมฺมสฺส อุปนิสฺสย\-ปจฺจเยน ปจฺจโย. อารมฺมณู\-ปนิสฺสโย, อนนฺตรู\-ปนิสฺสโย, ปกตู\-ปนิสฺสโย ฯเปฯ. ปกตู\-ปนิสฺสโย\csromandash{}สวิตกฺก\-สวิจารํ สทฺธํ อุปนิสฺสาย อวิตกฺก\-วิจารมตฺตํ ฌานํ อุปฺปาเทติ, มคฺคํ ฯเปฯ สมาปตฺติํ ฯเปฯ. สวิตกฺก\-สวิจารํ \csromanfolio{71}สีลํ ฯเปฯ ปตฺถนํ อุปนิสฺสาย อวิตกฺก\-วิจารมตฺตํ ฌานํ อุปฺปาเทติ, มคฺคํ ฯเปฯ สมาปตฺติํ ฯเปฯ. สวิตกฺก\-สวิจารา สทฺธา ฯเปฯ ปตฺถนา อวิตกฺก\-วิจารมตฺตาย สทฺธาย. สีลสฺส. สุตสฺส. จาคสฺส. ปญฺญาย วิตกฺกสฺส จ อุปนิสฺสย\-ปจฺจเยน ปจฺจโย. (2)}

\prose{สวิตกฺก\-สวิจาโร ธมฺโม อวิตกฺก\-อวิจารสฺส ธมฺมสฺส อุปนิสฺสย\-ปจฺจเยน ปจฺจโย. อนนฺตรู\-ปนิสฺสโย, ปกตู\-ปนิสฺสโย ฯเปฯ. ปกตู\-ปนิสฺสโย\csromandash{}สวิตกฺก\-สวิจารํ สทฺธํ อุปนิสฺสาย อวิตกฺก\-อวิจารํ ฌานํ อุปฺปาเทติ, มคฺคํ ฯเปฯ อภิญฺญํ ฯเปฯ สมาปตฺติํ อุปฺปาเทติ. สวิตกฺก\-สวิจารํ สีลํ ฯเปฯ ปตฺถนํ อุปนิสฺสาย อวิตกฺก\-อวิจารํ ฌานํ อุปฺปาเทติ, มคฺคํ ฯเปฯ อภิญฺญํ ฯเปฯ สมาปตฺติํ อุปฺปาเทติ. สวิตกฺก\-สวิจารา สทฺธา ฯเปฯ ปตฺถนา อวิตกฺก\-อวิจาราย สทฺธาย. สีลสฺส. สุตสฺส. จาคสฺส. ปญฺญาย วิจารสฺส จ. กายิกสฺส สุขสฺส, กายิกสฺส ทุกฺขสฺส อุปนิสฺสย\-ปจฺจเยน ปจฺจโย. (3)}

\prose{สวิตกฺก\-สวิจาโร ธมฺโม อวิตกฺก\-วิจารมตฺตสฺส จ อวิตกฺก\-อวิจารสฺส จ ธมฺมสฺส อุปนิสฺสย\-ปจฺจเยน ปจฺจโย. อนนฺตรู\-ปนิสฺสโย, ปกตู\-ปนิสฺสโย ฯเปฯ. ปกตู\-ปนิสฺสโย\csromandash{} สวิตกฺก\-สวิจารา สทฺธา ฯเปฯ ปตฺถนา อวิตกฺก\-วิจารมตฺตาย สทฺธาย. สีลสฺส. สุตสฺส. จาคสฺส. ปญฺญาย วิจารสฺส จ อุปนิสฺสย\-ปจฺจเยน ปจฺจโย. (4)}

\prose{สวิตกฺก\-สวิจาโร ธมฺโม สวิตกฺก\-สวิจารสฺส จ อวิตกฺก\-วิจารมตฺตสฺส จ ธมฺมสฺส อุปนิสฺสย\-ปจฺจเยน ปจฺจโย. อารมฺมณู\-ปนิสฺสโย, อนนฺตรู\-ปนิสฺสโย, ปกตู\-ปนิสฺสโย ฯเปฯ. ปกตู\-ปนิสฺสโย\csromandash{}สวิตกฺก\-สวิจารา สทฺธา ฯเปฯ ปตฺถนา สวิตกฺก\-สวิจาราย สทฺธาย ฯเปฯ ปตฺถนาย วิตกฺกสฺส จ อุปนิสฺสย\-ปจฺจเยน ปจฺจโย. (5)}

\proseitem{103}{อวิตกฺก\-วิจารมตฺโต ธมฺโม อวิตกฺก\-วิจารมตฺตสฺส ธมฺมสฺส อุปนิสฺสย\-ปจฺจเยน ปจฺจโย. อารมฺมณู\-ปนิสฺสโย, อนนฺตรู\-ปนิสฺสโย, ปกตู\-ปนิสฺสโย ฯเปฯ. ปกตู\-ปนิสฺสโย\csromandash{}อวิตกฺก\-วิจารมตฺตํ สทฺธํ อุปนิสฺสาย อวิตกฺก\-วิจารมตฺตํ ฌานํ อุปฺปาเทติ, มคฺคํ ฯเปฯ สมาปตฺติํ อุปฺปาเทติ. อวิตกฺก\-วิจารมตฺตํ สีลํ. สุตํ. จาคํ. ปญฺญํ. วิตกฺกํ อุปนิสฺสาย อวิตกฺก\-วิจารมตฺตํ ฌานํ อุปฺปาเทติ, มคฺคํ ฯเปฯ สมาปตฺติํ อุปฺปาเทติ. อวิตกฺก\-วิจารมตฺตา สทฺธา. สีลํ. สุตํ. จาโค. ปญฺญา วิตกฺโก จ \csromanfolio{72}อวิตกฺก\-วิจารมตฺตาย สทฺธาย. สีลสฺส. สุตสฺส. จาคสฺส. ปญฺญาย วิตกฺกสฺส จ อุปนิสฺสย\-ปจฺจเยน ปจฺจโย. (1)}

\prose{อวิตกฺก\-วิจารมตฺโต ธมฺโม สวิตกฺก\-สวิจารสฺส ธมฺมสฺส อุปนิสฺสย\-ปจฺจเยน ปจฺจโย. อารมฺมณู\-ปนิสฺสโย, อนนฺตรู\-ปนิสฺสโย, ปกตู\-ปนิสฺสโย ฯเปฯ. ปกตู\-ปนิสฺสโย\csromandash{}อวิตกฺก\-วิจารมตฺตํ สทฺธํ อุปนิสฺสาย ทานํ เทติ, สีลํ สมาทิยติ, อุโปสถกมฺมํ กโรติ, สวิตกฺก\-สวิจารํ ฌานํ อุปฺปาเทติ, วิปสฺสนํ ฯเปฯ มคฺคํ ฯเปฯ สมาปตฺติํ อุปฺปาเทติ, มานํ ชปฺเปติ, ทิฏฺฐิํ คณฺหาติ.  อวิตกฺก\-วิจารมตฺตํ สีลํ. สุตํ. จาคํ. ปญฺญํ. วิตกฺกํ อุปนิสฺสาย ทานํ เทติ, สีลํ สมาทิยติ, อุโปสถกมฺมํ กโรติ, สวิตกฺก\-สวิจารํ ฌานํ อุปฺปาเทติ, วิปสฺสนํ อุปฺปาเทติ, มคฺคํ อุปฺปาเทติ, สมาปตฺติํ อุปฺปาเทติ, ปาณํ หนติ ฯเปฯ สํฆํ ภินฺทติ.  อวิตกฺก\-วิจารมตฺตา สทฺธา. สีลํ. สุตํ. จาโค. ปญฺญา วิตกฺโก จ สวิตกฺก\-สวิจาราย สทฺธาย ฯเปฯ ปตฺถนาย อุปนิสฺสย\-ปจฺจเยน ปจฺจโย. (2)}

\prose{อวิตกฺก\-วิจารมตฺโต ธมฺโม อวิตกฺก\-อวิจารสฺส ธมฺมสฺส อุปนิสฺสย\-ปจฺจเยน ปจฺจโย. อนนฺตรู\-ปนิสฺสโย, ปกตู\-ปนิสฺสโย ฯเปฯ. ปกตู\-ปนิสฺสโย\csromandash{}อวิตกฺก\-วิจารมตฺตํ สทฺธํ อุปนิสฺสาย อวิตกฺก\-อวิจารํ ฌานํ อุปฺปาเทติ, มคฺคํ ฯเปฯ อภิญฺญํ ฯเปฯ สมาปตฺติํ อุปฺปาเทติ. อวิตกฺก\-วิจารมตฺตํ สีลํ. สุตํ. จาคํ. ปญฺญํ. วิตกฺกํ อุปนิสฺสาย อวิตกฺก\-อวิจารํ ฌานํ อุปฺปาเทติ, มคฺคํ ฯเปฯ อภิญฺญํ ฯเปฯ สมาปตฺติํ อุปฺปาเทติ. อวิตกฺก\-วิจารมตฺตา สทฺธา. สีลํ. สุตํ. จาโค. ปญฺญา วิตกฺโก จ อวิตกฺก\-อวิจาราย สทฺธาย. สีลสฺส. สุตสฺส. จาคสฺส. ปญฺญาย วิจารสฺส จ. กายิกสฺส สุขสฺส, กายิกสฺส ทุกฺขสฺส อุปนิสฺสย\-ปจฺจเยน ปจฺจโย. (3)}

\prose{อวิตกฺก\-วิจารมตฺโต ธมฺโม อวิตกฺก\-วิจารมตฺตสฺส จ อวิตกฺก\-อวิจารสฺส จ ธมฺมสฺส อุปนิสฺสย\-ปจฺจเยน ปจฺจโย. อนนฺตรู\-ปนิสฺสโย, ปกตู\-ปนิสฺสโย ฯเปฯ. ปกตู\-ปนิสฺสโย\csromandash{} อวิตกฺก\-วิจารมตฺตา สทฺธา. สีลํ. สุตํ. จาโค. ปญฺญา วิตกฺโก จ อวิตกฺก\-วิจารมตฺตาย สทฺธาย. สีลสฺส. สุตสฺส. จาคสฺส. ปญฺญาย วิจารสฺส จ อุปนิสฺสย\-ปจฺจเยน ปจฺจโย. (4)}

\prose{\csromanfolio{73}อวิตกฺก\-วิจารมตฺโต ธมฺโม สวิตกฺก\-สวิจารสฺส จ อวิตกฺก\-วิจารมตฺตสฺส จ ธมฺมสฺส อุปนิสฺสย\-ปจฺจเยน ปจฺจโย. อารมฺมณู\-ปนิสฺสโย, อนนฺตรู\-ปนิสฺสโย, ปกตู\-ปนิสฺสโย ฯเปฯ. ปกตู\-ปนิสฺสโย\csromandash{}อวิตกฺก\-วิจารมตฺตา สทฺธา. สีลํ. สุตํ. จาโค. ปญฺญา วิตกฺโก จ สวิตกฺก\-สวิจาราย สทฺธาย ฯเปฯ ปตฺถนาย วิตกฺกสฺส จ อุปนิสฺสย\-ปจฺจเยน ปจฺจโย. (5)}

\proseitem{104}{อวิตกฺก\-อวิจาโร ธมฺโม อวิตกฺก\-อวิจารสฺส ธมฺมสฺส อุปนิสฺสย\-ปจฺจเยน ปจฺจโย. อารมฺมณู\-ปนิสฺสโย, อนนฺตรู\-ปนิสฺสโย, ปกตู\-ปนิสฺสโย ฯเปฯ. ปกตู\-ปนิสฺสโย\csromandash{}อวิตกฺก\-อวิจารํ สทฺธํ อุปนิสฺสาย อวิตกฺก\-อวิจารํ ฌานํ อุปฺปาเทติ, มคฺคํ ฯเปฯ อภิญฺญํ ฯเปฯ สมาปตฺติํ อุปฺปาเทติ. อวิตกฺก\-อวิจารํ สีลํ. สุตํ. จาคํ. ปญฺญํ. วิจารํ. กายิกํ สุขํ. กายิกํ ทุกฺขํ. อุตุํ. โภชนํ. เสนาสนํ อุปนิสฺสาย อวิตกฺก\-อวิจารํ ฌานํ อุปฺปาเทติ, มคฺคํ ฯเปฯ อภิญฺญํ ฯเปฯ สมาปตฺติํ อุปฺปาเทติ. อวิตกฺก\-อวิจารา สทฺธา. สีลํ. สุตํ. จาโค. ปญฺญา. วิจาโร. กายิกํ สุขํ. กายิกํ ทุกฺขํ. อุตุ. โภชนํ. เสนาสนํ อวิตกฺก\-อวิจาราย สทฺธาย. สีลสฺส. สุตสฺส. จาคสฺส. ปญฺญาย. วิจารสฺส. กายิกสฺส สุขสฺส, กายิกสฺส ทุกฺขสฺส อุปนิสฺสย\-ปจฺจเยน ปจฺจโย. (1)}

\prose{อวิตกฺก\-อวิจาโร ธมฺโม สวิตกฺก\-สวิจารสฺส ธมฺมสฺส อุปนิสฺสย\-ปจฺจเยน ปจฺจโย. อารมฺมณู\-ปนิสฺสโย, อนนฺตรู\-ปนิสฺสโย, ปกตู\-ปนิสฺสโย ฯเปฯ. ปกตู\-ปนิสฺสโย\csromandash{}อวิตกฺก\-อวิจารํ สทฺธํ อุปนิสฺสาย ทานํ เทติ, สีลํ สมาทิยติ, อุโปสถกมฺมํ กโรติ, สวิตกฺก\-สวิจารํ ฌานํ อุปฺปาเทติ, วิปสฺสนํ ฯเปฯ มคฺคํ ฯเปฯ สมาปตฺติํ อุปฺปาเทติ. มานํ ชปฺเปติ, ทิฏฺฐิํ คณฺหาติ. อวิตกฺก\-อวิจารํ สีลํ. สุตํ. จาคํ. ปญฺญํ. วิจารํ. กายิกํ สุขํ. กายิกํ ทุกฺขํ. อุตุํ. โภชนํ. เสนาสนํ อุปนิสฺสาย ทานํ เทติ, สีลํ สมาทิยติ, อุโปสถกมฺมํ กโรติ, สวิตกฺก\-สวิจารํ ฌานํ อุปฺปาเทติ, วิปสฺสนํ ฯเปฯ มคฺคํ ฯเปฯ สมาปตฺติํ อุปฺปาเทติ, ปาณํ หนติ ฯเปฯ สํฆํ ภินฺทติ. อวิตกฺก\-อวิจารา สทฺธา ฯเปฯ เสนาสนํ สวิตกฺก\-สวิจาราย สทฺธาย. สีลสฺส ฯเปฯ ปตฺถนาย อุปนิสฺสย\-ปจฺจเยน ปจฺจโย. (2)}

\prose{อวิตกฺก\-อวิจาโร ธมฺโม อวิตกฺก\-วิจารมตฺตสฺส ธมฺมสฺส อุปนิสฺสย\-ปจฺจเยน ปจฺจโย. อารมฺมณู\-ปนิสฺสโย, อนนฺตรู\-ปนิสฺสโย, \csromanfolio{74}ปกตู\-ปนิสฺสโย ฯเปฯ. ปกตู\-ปนิสฺสโย\csromandash{}อวิตกฺก\-อวิจารํ สทฺธํ อุปนิสฺสาย อวิตกฺก\-วิจารมตฺตํ ฌานํ อุปฺปาเทติ, วิปสฺสนํ ฯเปฯ มคฺคํ ฯเปฯ สมาปตฺติํ อุปฺปาเทติ. อวิตกฺก\-อวิจารํ สีลํ ฯเปฯ. เสนาสนํ อุปนิสฺสาย อวิตกฺก\-วิจารมตฺตํ ฌานํ อุปฺปาเทติ, วิปสฺสนํ ฯเปฯ มคฺคํ ฯเปฯ สมาปตฺติํ อุปฺปาเทติ. อวิตกฺก\-อวิจารา สทฺธา ฯเปฯ. เสนาสนํ อวิตกฺก\-วิจารมตฺตาย สทฺธาย. สีลสฺส. สุตสฺส. จาคสฺส. ปญฺญาย วิตกฺกสฺส จ อุปนิสฺสย\-ปจฺจเยน ปจฺจโย. (3)}

\prose{อวิตกฺก\-อวิจาโร ธมฺโม อวิตกฺก\-วิจารมตฺตสฺส จ อวิตกฺก\-อวิจารสฺส จ ธมฺมสฺส อุปนิสฺสย\-ปจฺจเยน ปจฺจโย. อารมฺมณู\-ปนิสฺสโย, อนนฺตรู\-ปนิสฺสโย, ปกตู\-ปนิสฺสโย ฯเปฯ. ปกตู\-ปนิสฺสโย\csromandash{}อวิตกฺก\-อวิจารา สทฺธา ฯเปฯ. เสนาสนํ อวิตกฺก\-วิจารมตฺตาย สทฺธาย. สีลสฺส. สุตสฺส. จาคสฺส. ปญฺญาย วิจารสฺส จ อุปนิสฺสย\-ปจฺจเยน ปจฺจโย. (4)}

\prose{อวิตกฺก\-อวิจาโร ธมฺโม สวิตกฺก\-สวิจารสฺส จ อวิตกฺก\-วิจารมตฺตสฺส จ ธมฺมสฺส อุปนิสฺสย\-ปจฺจเยน ปจฺจโย. อารมฺมณู\-ปนิสฺสโย, อนนฺตรู\-ปนิสฺสโย, ปกตู\-ปนิสฺสโย ฯเปฯ. ปกตู\-ปนิสฺสโย\csromandash{}อวิตกฺก\-อวิจารา สทฺธา ฯเปฯ. เสนาสนํ สวิตกฺก\-สวิจาราย สทฺธาย. สีลสฺส ฯเปฯ. ปตฺถนาย วิตกฺกสฺส จ อุปนิสฺสย\-ปจฺจเยน ปจฺจโย. (5)}

\proseitem{105}{อวิตกฺก\-วิจารมตฺโต จ อวิตกฺก\-อวิจาโร จ ธมฺมา สวิตกฺก\-สวิจารสฺส ธมฺมสฺส อุปนิสฺสย\-ปจฺจเยน ปจฺจโย. อารมฺมณู\-ปนิสฺสโย, อนนฺตรู\-ปนิสฺสโย, ปกตู\-ปนิสฺสโย ฯเปฯ. ปกตู\-ปนิสฺสโย\csromandash{}อวิตกฺก\-วิจารมตฺตา สทฺธา. สีลํ. สุตํ. จาโค. ปญฺญา วิจาโร จ สวิตกฺก\-สวิจาราย สทฺธาย ฯเปฯ ปตฺถนาย อุปนิสฺสย\-ปจฺจเยน ปจฺจโย. (1)}

\prose{อวิตกฺก\-วิจารมตฺโต จ อวิตกฺก\-อวิจาโร จ ธมฺมา อวิตกฺก\-วิจารมตฺตสฺส ธมฺมสฺส อุปนิสฺสย\-ปจฺจเยน ปจฺจโย. อารมฺมณู\-ปนิสฺสโย, อนนฺตรู\-ปนิสฺสโย, ปกตู\-ปนิสฺสโย ฯเปฯ. ปกตู\-ปนิสฺสโย\csromandash{}อวิตกฺก\-วิจารมตฺตา สทฺธา. สีลํ. สุตํ. จาโค. ปญฺญา วิจาโร จ อวิตกฺก\-วิจารมตฺตาย สทฺธาย. สีลสฺส. สุตสฺส. จาคสฺส. ปญฺญาย วิตกฺกสฺส จ อุปนิสฺสย\-ปจฺจเยน ปจฺจโย. (2)}

\prose{\csromanfolio{75}อวิตกฺก\-วิจารมตฺโต จ อวิตกฺก\-อวิจาโร จ ธมฺมา อวิตกฺก\-อวิจารสฺส ธมฺมสฺส อุปนิสฺสย\-ปจฺจเยน ปจฺจโย. อนนฺตรู\-ปนิสฺสโย, ปกตู\-ปนิสฺสโย ฯเปฯ. ปกตู\-ปนิสฺสโย\csromandash{} อวิตกฺก\-วิจารมตฺตา สทฺธา. สีลํ. สุตํ. จาโค. ปญฺญา วิจาโร จ อวิตกฺก\-อวิจาราย สทฺธาย. สีลสฺส. สุตสฺส. จาคสฺส. ปญฺญาย วิจารสฺส จ. กายิกสฺส สุขสฺส, กายิกสฺส ทุกฺขสฺส อุปนิสฺสย\-ปจฺจเยน ปจฺจโย. (3)}

\prose{อวิตกฺก\-วิจารมตฺโต จ อวิตกฺก\-อวิจาโร จ ธมฺมา อวิตกฺก\-วิจารมตฺตสฺส จ อวิตกฺก\-อวิจารสฺส จ ธมฺมสฺส อุปนิสฺสย\-ปจฺจเยน ปจฺจโย. อนนฺตรู\-ปนิสฺสโย, ปกตู\-ปนิสฺสโย ฯเปฯ. ปกตู\-ปนิสฺสโย\csromandash{}อวิตกฺก\-วิจารมตฺตา สทฺธา. สีลํ. สุตํ. จาโค. ปญฺญา วิจาโร จ อวิตกฺก\-วิจารมตฺตาย สทฺธาย. สีลสฺส. สุตสฺส. จาคสฺส. ปญฺญาย วิจารสฺส จ อุปนิสฺสย\-ปจฺจเยน ปจฺจโย. (4)}

\prose{อวิตกฺก\-วิจารมตฺโต จ อวิตกฺก\-อวิจาโร จ ธมฺมา สวิตกฺก\-สวิจารสฺส จ อวิตกฺก\-วิจารมตฺตสฺส จ ธมฺมสฺส อุปนิสฺสย\-ปจฺจเยน ปจฺจโย. อารมฺมณู\-ปนิสฺสโย, อนนฺตรู\-ปนิสฺสโย, ปกตู\-ปนิสฺสโย ฯเปฯ. ปกตู\-ปนิสฺสโย\csromandash{}อวิตกฺก\-วิจารมตฺตา สทฺธา. สีลํ. สุตํ. จาโค. ปญฺญา วิจาโร จ สวิตกฺก\-สวิจาราย สทฺธาย. สีลสฺส. สุตสฺส. จาคสฺส. ปญฺญาย. ราคสฺส. โทสสฺส. โมหสฺส. มานสฺส. ทิฏฺฐิยา. ปตฺถนาย วิตกฺกสฺส จ อุปนิสฺสย\-ปจฺจเยน ปจฺจโย. (5)}

\proseitem{106}{สวิตกฺก\-สวิจาโร จ อวิตกฺก\-วิจารมตฺโต จ ธมฺมา สวิตกฺก\-สวิจารสฺส ธมฺมสฺส อุปนิสฺสย\-ปจฺจเยน ปจฺจโย. อารมฺมณู\-ปนิสฺสโย, อนนฺตรู\-ปนิสฺสโย, ปกตู\-ปนิสฺสโย ฯเปฯ. ปกตู\-ปนิสฺสโย\csromandash{}สวิตกฺก\-สวิจารา สทฺธา. สีลํ. สุตํ. จาโค. ปญฺญา. ราโค. โทโส. โมโห. มาโน. ทิฏฺฐิ. ปตฺถนา วิตกฺโก จ สวิตกฺก\-สวิจาราย สทฺธาย. สีลสฺส ฯเปฯ ปตฺถนาย อุปนิสฺสย\-ปจฺจเยน ปจฺจโย. (1)}

\prose{สวิตกฺก\-สวิจาโร จ อวิตกฺก\-วิจารมตฺโต จ ธมฺมา อวิตกฺก\-วิจารมตฺตสฺส ธมฺมสฺส อุปนิสฺสย\-ปจฺจเยน ปจฺจโย. อารมฺมณู\-ปนิสฺสโย, อนนฺตรู\-ปนิสฺสโย, ปกตู\-ปนิสฺสโย ฯเปฯ. ปกตู\-ปนิสฺสโย\csromandash{} \csromanfolio{76}สวิตกฺก\-สวิจารา สทฺธา. สีลํ ฯเปฯ ปตฺถนา วิตกฺโก จ อวิตกฺก\-วิจารมตฺตาย สทฺธาย. สีลสฺส. สุตสฺส. จาคสฺส. ปญฺญาย วิตกฺกสฺส จ อุปนิสฺสย\-ปจฺจเยน ปจฺจโย. (2)}

\prose{สวิตกฺก\-สวิจาโร จ อวิตกฺก\-วิจารมตฺโต จ ธมฺมา อวิตกฺก\-อวิจารสฺส ธมฺมสฺส อุปนิสฺสย\-ปจฺจเยน ปจฺจโย. อนนฺตรู\-ปนิสฺสโย, ปกตู\-ปนิสฺสโย ฯเปฯ. ปกตู\-ปนิสฺสโย\csromandash{} สวิตกฺก\-สวิจารา สทฺธา. สีลํ. สุตํ. จาโค. ปญฺญา ฯเปฯ ปตฺถนา วิตกฺโก จ อวิตกฺก\-อวิจาราย สทฺธาย. สีลสฺส. สุตสฺส. จาคสฺส. ปญฺญาย วิจารสฺส จ. กายิกสฺส สุขสฺส, กายิกสฺส ทุกฺขสฺส อุปนิสฺสย\-ปจฺจเยน ปจฺจโย. (3)}

\prose{สวิตกฺก\-สวิจาโร จ อวิตกฺก\-วิจารมตฺโต จ ธมฺมา อวิตกฺก\-วิจารมตฺตสฺส จ อวิตกฺก\-อวิจารสฺส จ ธมฺมสฺส อุปนิสฺสย\-ปจฺจเยน ปจฺจโย. อนนฺตรู\-ปนิสฺสโย, ปกตู\-ปนิสฺสโย ฯเปฯ. ปกตู\-ปนิสฺสโย\csromandash{}สวิตกฺก\-สวิจารา สทฺธา. สีลํ. สุตํ. จาโค. ปญฺญา ฯเปฯ. ปตฺถนา วิตกฺโก จ อวิตกฺก\-วิจารมตฺตาย สทฺธาย. สีลสฺส. สุตสฺส. จาคสฺส. ปญฺญาย วิจารสฺส จ อุปนิสฺสย\-ปจฺจเยน ปจฺจโย. (4)}

\prose{สวิตกฺก\-สวิจาโร จ อวิตกฺก\-วิจารมตฺโต จ ธมฺมา สวิตกฺก\-สวิจารสฺส จ อวิตกฺก\-วิจารมตฺตสฺส จ ธมฺมสฺส อุปนิสฺสย\-ปจฺจเยน ปจฺจโย. อารามฺมณูปนิสฺสโย, อนนฺตรู\-ปนิสฺสโย, ปกตู\-ปนิสฺสโย ฯเปฯ. ปกตู\-ปนิสฺสโย\csromandash{}สวิตกฺก\-สวิจารา สทฺธา. สีลํ. สุตํ. จาโค. ปญฺญา. ราโค. โทโส. โมโห. มาโน. ทิฏฺฐิ. ปตฺถนา วิตกฺโก จ สวิตกฺก\-สวิจาราย สทฺธาย ฯเปฯ ปตฺถนาย วิตกฺกสฺส จ อุปนิสฺสย\-ปจฺจเยน ปจฺจโย. (5)}

\csromanheader{2}{\textbf{ปุเรชาต}}
\proseitem{107}{อวิตกฺก\-อวิจาโร ธมฺโม อวิตกฺก\-อวิจารสฺส ธมฺมสฺส ปุเรชาต\-ปจฺจเยน ปจฺจโย. อารมฺมณ\-ปุเรชาตํ, วตฺถุ\-ปุเรชาตํ.  อารมฺมณ\-ปุเรชาตํ\csromandash{}ทิพฺเพน จกฺขุนา รูปํ ปสฺสติ. ทิพฺพาย โสตธาตุยา สทฺทํ สุณาติ. รูปายตนํ จกฺขุ\-วิญฺญาณสฺส ฯเปฯ. โผฏฺฐพฺพา\-ยตนํ กายวิญฺญาณสฺส ปุเรชาต\-ปจฺจเยน ปจฺจโย.  วตฺถุ\-ปุเรชาตํ\csromandash{}จกฺขายตนํ จกฺขุ\-วิญฺญาณสฺส ฯเปฯ กายายตนํ กายวิญฺญาณสฺส ฯเปฯ. วตฺถุ อวิตกฺก\-อวิจารานํ ขนฺธานํ วิจารสฺส จ ปุเรชาต\-ปจฺจเยน ปจฺจโย. (1)}

\prose{\csromanfolio{77}อวิตกฺก\-อวิจาโร ธมฺโม สวิตกฺก\-สวิจารสฺส ธมฺมสฺส ปุเรชาต\-ปจฺจเยน ปจฺจโย. อารมฺมณ\-ปุเรชาตํ, วตฺถุ\-ปุเรชาตํ.  อารมฺมณ\-ปุเรชาตํ\csromandash{}จกฺขุํ อนิจฺจโต ทุกฺขโต อนตฺตโต วิปสฺสติ ฯเปฯ โผฏฺฐพฺเพ. วตฺถุํ อนิจฺจโต ทุกฺขโต อนตฺตโต วิปสฺสติ, อสฺสาเทติ อภินนฺทติ, ตํ อารพฺภ ราโค อุปฺปชฺชติ ฯเปฯ โทมนสฺสํ อุปฺปชฺชติ.  วตฺถุ\-ปุเรชาตํ\csromandash{}วตฺถุ สวิตกฺก\-สวิจารานํ ขนฺธานํ ปุเรชาต\-ปจฺจเยน ปจฺจโย. (2)}

\prose{อวิตกฺก\-อวิจาโร ธมฺโม อวิตกฺก\-วิจารมตฺตสฺส ธมฺมสฺส ปุเรชาต\-ปจฺจเยน ปจฺจโย. อารมฺมณ\-ปุเรชาตํ, วตฺถุ\-ปุเรชาตํ.  อารมฺมณ\-ปุเรชาตํ\csromandash{}จกฺขุํ อนิจฺจโต ทุกฺขโต อนตฺตโต วิปสฺสติ, อสฺสาเทติ อภินนฺทติ, ตํ อารพฺภ วิตกฺโก อุปชฺชติ ฯเปฯ วตฺถุํ อนิจฺจโต ทุกฺขโต อนตฺตโต วิปสฺสติ, อสฺสาเทติ อภินนฺทติ, ตํ อารพฺภ วิตกฺโก อุปฺปชฺชติ.  วตฺถุ\-ปุเรชาตํ\csromandash{}วตฺถุ อวิตกฺก\-วิจารมตฺตานํ ขนฺธานํ วิตกฺกสฺส จ ปุเรชาต\-ปจฺจเยน ปจฺจโย. (3)}

\prose{อวิตกฺก\-อวิจาโร ธมฺโม อวิตกฺก\-วิจารมตฺตสฺส จ อวิตกฺก\-อวิจารสฺส จ ธมฺมสฺส ปุเรชาต\-ปจฺจเยน ปจฺจโย. วตฺถุ\-ปุเรชาตํ\csromandash{}วตฺถุ อวิตกฺก\-วิจารมตฺตานํ ขนฺธานํ วิจารสฺส จ ปุเรชาต\-ปจฺจเยน ปจฺจโย. (4)}

\prose{อวิตกฺก\-อวิจาโร ธมฺโม สวิตกฺก\-สวิจารสฺส จ อวิตกฺก\-วิจารมตฺตสฺส จ ธมฺมสฺส ปุเรชาต\-ปจฺจเยน ปจฺจโย. อารมฺมณ\-ปุเรชาตํ, วตฺถุ\-ปุเรชาตํ.  อารมฺมณ\-ปุเรชาตํ\csromandash{}จกฺขุํ อนิจฺจโต ทุกฺขโต อนตฺตโต วิปสฺสติ, อสฺสาเทติ อภินนฺทติ, ตํ อารพฺภ สวิตกฺก\-สวิจารา ขนฺธา จ วิตกฺโก จ อุปฺปชฺชนฺติ. โสตํ. ฆานํ. ชิวฺหํ. กายํ. รูเป. สทฺเท. คนฺเธ. รเส. โผฏฺฐพฺเพ. วตฺถุํ อนิจฺจโต ทุกฺขโต อนตฺตโต วิปสฺสติ, อสฺสาเทติ อภินนฺทติ, ตํ อารพฺภ สวิตกฺก\-สวิจารา ขนฺธา จ วิตกฺโก จ อุปฺปชฺชนฺติ.  วตฺถุ\-ปุเรชาตํ\csromandash{} วตฺถุ สวิตกฺก\-สวิจารานํ ขนฺธานํ วิตกฺกสฺส จ ปุเรชาต\-ปจฺจเยน ปจฺจโย. (5)}

\csromanheader{2}{\textbf{ปจฺฉาชาต}}
\proseitem{108}{สวิตกฺก\-สวิจาโร ธมฺโม อวิตกฺก\-อวิจารสฺส ธมฺมสฺส ปจฺฉาชาต\-ปจฺจเยน ปจฺจโย. ปจฺฉาชาตา สวิตกฺก\-สวิจารา ขนฺธา ปุเรชาตสฺส อิมสฺส กายสฺส ปจฺฉาชาต\-ปจฺจเยน ปจฺจโย. (1)}

\prose{\csromanfolio{78}อวิตกฺก\-วิจารมตฺโต ธมฺโม อวิตกฺก\-อวิจารสฺส ธมฺมสฺส ปจฺฉาชาต\-ปจฺจเยน ปจฺจโย. ปจฺฉาชาตา อวิตกฺก\-วิจารมตฺตา ขนฺธา จ วิตกฺโก จ ปุเรชาตสฺส อิมสฺส กายสฺส ปจฺฉาชาต\-ปจฺจเยน ปจฺจโย. (1)}

\prose{อวิตกฺก\-อวิจาโร ธมฺโม อวิตกฺก\-อวิจารสฺส ธมฺมสฺส ปจฺฉาชาต\-ปจฺจเยน ปจฺจโย. ปจฺฉาชาตา อวิตกฺก\-อวิจารา ขนฺธา จ วิจาโร จ ปุเรชาตสฺส อิมสฺส กายสฺส ปจฺฉาชาต\-ปจฺจเยน ปจฺจโย. (1)}

\prose{อวิตกฺก\-วิจารมตฺโต จ อวิตกฺก\-อวิจาโร จ ธมฺมา อวิตกฺก\-อวิจารสฺส ธมฺมสฺส ปจฺฉาชาต\-ปจฺจเยน ปจฺจโย. ปจฺฉาชาตา อวิตกฺก\-วิจารมตฺตา ขนฺธา จ วิจาโร จ ปุเรชาตสฺส อิมสฺส กายสฺส ปจฺฉาชาต\-ปจฺจเยน ปจฺจโย. (1)}

\prose{สวิตกฺก\-สวิจาโร จ อวิตกฺก\-วิจารมตฺโต จ ธมฺมา อวิตกฺก\-อวิจารสฺส ธมฺมสฺส ปจฺฉาชาต\-ปจฺจเยน ปจฺจโย. ปจฺฉาชาตา สวิตกฺก\-สวิจารา ขนฺธา จ วิตกฺโก จ ปุเรชาตสฺส อิมสฺส กายสฺส ปจฺฉาชาต\-ปจฺจเยน ปจฺจโย. (1)}

\csromanheader{2}{\textbf{อาเสวน}}
\proseitem{109}{สวิตกฺก\-สวิจาโร ธมฺโม สวิตกฺก\-สวิจารสฺส ธมฺมสฺส อาเสวน\-ปจฺจเยน ปจฺจโย. ปุริมา ปุริมา สวิตกฺก\-สวิจารา ขนฺธา ปจฺฉิมานํ ปจฺฉิมานํ สวิตกฺก\-สวิจารานํ ขนฺธานํ อาเสวน\-ปจฺจเยน ปจฺจโย. อนุโลมํ โคตฺรภุสฺส. อนุโลมํ โวทานสฺส. โคตฺรภุ สวิตกฺก\-สวิจารสฺส มคฺคสฺส. โวทานํ สวิตกฺก\-สวิจารสฺส มคฺคสฺส อาเสวน\-ปจฺจเยน ปจฺจโย. (1)}

\prose{สวิตกฺก\-สวิจาโร ธมฺโม อวิตกฺก\-วิจารมตฺตสฺส ธมฺมสฺส อาเสวน\-ปจฺจเยน ปจฺจโย. ปุริมา ปุริมา สวิตกฺก\-สวิจารา ขนฺธา ปจฺฉิมสฺส ปจฺฉิมสฺส วิตกฺกสฺส อาเสวน\-ปจฺจเยน ปจฺจโย. อวิตกฺก\-วิจารมตฺตสฺส ฌานสฺส ปริกมฺมํ อวิตกฺก\-วิจารมตฺตสฺส ฌานสฺส อาเสวน\-ปจฺจเยน ปจฺจโย. โคตฺรภุ อวิตกฺก\-วิจารมตฺตสฺส มคฺคสฺส อาเสวน\-ปจฺจเยน ปจฺจโย. โวทานํ อวิตกฺก\-วิจารมตฺตสฺส มคฺคสฺส อาเสวน\-ปจฺจเยน ปจฺจโย. (2)}

\prose{สวิตกฺก\-สวิจาโร ธมฺโม อวิตกฺก\-อวิจารสฺส ธมฺมสฺส อาเสวน\-ปจฺจเยน ปจฺจโย. ทุติยสฺส ฌานสฺส ปริกมฺมํ ทุติเย ฌาเน วิจารสฺส \csromanfolio{79}อาเสวน\-ปจฺจเยน ปจฺจโย. ตติยสฺส ฌานสฺส ปริกมฺมํ ตติยสฺส ฌานสฺส ฯเปฯ. จตุตฺถสฺส ฌานสฺส ปริกมฺมํ จตุตฺถสฺส ฌานสฺส ฯเปฯ. อากาสา\-นญฺจา\-ยตนสฺส ปริกมฺมํ อากาสา\-นญฺจา\-ยตนสฺส ฯเปฯ. วิญฺญาณญฺจา\-ยตนสฺส ปริกมฺมํ วิญฺญาณญฺจา\-ยตนสฺส ฯเปฯ. อากิญฺจญฺญา\-ยตนสฺส ปริกมฺมํ อากิญฺจญฺญา\-ยตนสฺส ฯเปฯ. เนว\-สญฺญา\-นา\-สญฺญา\-ยตนสฺส ปริกมฺมํ เนว\-สญฺญา\-นา\-สญฺญา\-ยตนสฺส ฯเปฯ. ทิพฺพสฺส จกฺขุสฺส ปริกมฺมํ ทิพฺพสฺส จกฺขุสฺส ฯเปฯ. ทิพฺพาย โสตธาตุยา ปริกมฺมํ ทิพฺพาย โสตธาตุยา ฯเปฯ. อิทฺธิวิธ\-ญาณสฺส ปริกมฺมํ ฯเปฯ. เจโตปริย\-ญาณสฺส ปริกมฺมํ ฯเปฯ. ปุพฺเพ\-นิวาสา\-นุสฺสติ\-ญาณสฺส ปริกมฺมํ ฯเปฯ. ยถากมฺมูปคญาณสฺส ปริกมฺมํ ฯเปฯ. อนาคตํสญาณสฺส ปริกมฺมํ ฯเปฯ. โคตฺรภุ อวิตกฺก\-อวิจารสฺส มคฺคสฺส วิจารสฺส จ. โวทานํ อวิตกฺก\-อวิจารสฺส มคฺคสฺส วิจารสฺส จ อาเสวน\-ปจฺจเยน ปจฺจโย. (3)}

\prose{สวิตกฺก\-สวิจาโร ธมฺโม อวิตกฺก\-วิจารมตฺตสฺส จ อวิตกฺก\-อวิจารสฺส จ ธมฺมสฺส อาเสวน\-ปจฺจเยน ปจฺจโย. อวิตกฺก\-วิจารมตฺตสฺส ฌานสฺส ปริกมฺมํ อวิตกฺก\-วิจารมตฺตสฺส ฌานสฺส วิจารสฺส จ อาเสวน\-ปจฺจเยน ปจฺจโย. โคตฺรภุ อวิตกฺก\-วิจารมตฺตสฺส มคฺคสฺส วิจารสฺส จ. โวทานํ อวิตกฺก\-วิจารมตฺตสฺส มคฺคสฺส วิจารสฺส จ อาเสวน\-ปจฺจเยน ปจฺจโย. (4)}

\prose{สวิตกฺก\-สวิจาโร ธมฺโม สวิตกฺก\-สวิจารสฺส จ อวิตกฺก\-วิจารมตฺตสฺส จ ธมฺมสฺส อาเสวน\-ปจฺจเยน ปจฺจโย. ปุริมา ปุริมา สวิตกฺก\-สวิจารา ขนฺธา ปจฺฉิมานํ ปจฺฉิมานํ สวิตกฺก\-สวิจารานํ ขนฺธานํ วิตกฺกสฺส จ อาเสวน\-ปจฺจเยน ปจฺจโย. อนุโลมํ โคตฺรภุสฺส วิตกฺกสฺส จ. อนุโลมํ โวทานสฺส วิตกฺกสฺส จ. โคตฺรภุ สวิตกฺก\-สวิจารสฺส มคฺคสฺส วิตกฺกสฺส จ. โวทานํ สวิตกฺก\-สวิจารสฺส มคฺคสฺส วิตกฺกสฺส จ อาเสวน\-ปจฺจเยน ปจฺจโย. (5)}

\proseitem{110}{อวิตกฺก\-วิจารมตฺโต ธมฺโม อวิตกฺก\-วิจารมตฺตสฺส ธมฺมสฺส อาเสวน\-ปจฺจเยน ปจฺจโย. ปุริโม ปุริโม วิตกฺโก ปจฺฉิมสฺส ปจฺฉิมสฺส วิตกฺกสฺส อาเสวน\-ปจฺจเยน ปจฺจโย. ปุริมา ปุริมา อวิตกฺก\-วิจารมตฺตา ขนฺธา ปจฺฉิมานํ ปจฺฉิมานํ อวิตกฺก\-วิจารมตฺตานํ ขนฺธานํ อาเสวน\-ปจฺจเยน ปจฺจโย. (1)}

\prose{\csromanfolio{80}อวิตกฺก\-วิจารมตฺโต ธมฺโม สวิตกฺก\-สวิจารสฺส ธมฺมสฺส อาเสวน\-ปจฺจเยน ปจฺจโย. ปุริโม ปุริโม วิตกฺโก ปจฺฉิมานํ ปจฺฉิมานํ สวิตกฺก\-สวิจารานํ ขนฺธานํ อาเสวน\-ปจฺจเยน ปจฺจโย. (2)}

\prose{อวิตกฺก\-วิจารมตฺโต ธมฺโม อวิตกฺก\-อวิจารสฺส ธมฺมสฺส อาเสวน\-ปจฺจเยน ปจฺจโย. ปุริมา ปุริมา อวิตกฺก\-วิจารมตฺตา ขนฺธา ปจฺฉิมสฺส ปจฺฉิมสฺส วิจารสฺส อาเสวน\-ปจฺจเยน ปจฺจโย. (3)}

\prose{อวิตกฺก\-วิจารมตฺโต ธมฺโม อวิตกฺก\-วิจารมตฺตสฺส จ อวิตกฺก\-อวิจารสฺส จ ธมฺมสฺส อาเสวน\-ปจฺจเยน ปจฺจโย. ปุริมา ปุริมา อวิตกฺก\-วิจารมตฺตา ขนฺธา ปจฺฉิมานํ ปจฺฉิมานํ อวิตกฺก\-วิจารมตฺตานํ ขนฺธานํ วิจารสฺส จ อาเสวน\-ปจฺจเยน ปจฺจโย. (4)}

\prose{อวิตกฺก\-วิจารมตฺโต ธมฺโม สวิตกฺก\-สวิจารสฺส จ อวิตกฺก\-วิจารมตฺตสฺส จ ธมฺมสฺส อาเสวน\-ปจฺจเยน ปจฺจโย. ปุริโม ปุริโม วิตกฺโก ปจฺฉิมานํ ปจฺฉิมานํ สวิตกฺก\-สวิจารานํ ขนฺธานํ วิตกฺกสฺส จ อาเสวน\-ปจฺจเยน ปจฺจโย. (5)}

\proseitem{111}{อวิตกฺก\-อวิจาโร ธมฺโม อวิตกฺก\-อวิจารสฺส ธมฺมสฺส อาเสวน\-ปจฺจเยน ปจฺจโย. ปุริโม ปุริโม วิจาโร ปจฺฉิมสฺส ปจฺฉิมสฺส วิจารสฺส อาเสวน\-ปจฺจเยน ปจฺจโย. ปุริมา ปุริมา อวิตกฺก\-อวิจารา ขนฺธา ปจฺฉิมานํ ปจฺฉิมานํ อวิตกฺก\-อวิจารานํ ขนฺธานํ อาเสวน\-ปจฺจเยน ปจฺจโย. (1)}

\prose{อวิตกฺก\-อวิจาโร ธมฺโม อวิตกฺก\-วิจารมตฺตสฺส ธมฺมสฺส อาเสวน\-ปจฺจเยน ปจฺจโย. ปุริโม ปุริโม วิจาโร ปจฺฉิมานํ ปจฺฉิมานํ อวิตกฺก\-วิจารมตฺตานํ ขนฺธานํ อาเสวน\-ปจฺจเยน ปจฺจโย. (2)}

\prose{อวิตกฺก\-อวิจาโร ธมฺโม อวิตกฺก\-วิจารมตฺตสฺส จ อวิตกฺก\-อวิจารสฺส จ ธมฺมสฺส อาเสวน\-ปจฺจเยน ปจฺจโย. ปุริโม ปุริโม วิจาโร ปจฺฉิมานํ ปจฺฉิมานํ อวิตกฺก\-วิจารมตฺตานํ ขนฺธานํ วิจารสฺส จ อาเสวน\-ปจฺจเยน ปจฺจโย. (3)}

\proseitem{112}{\csromanfolio{81}อวิตกฺก\-วิจารมตฺโต จ อวิตกฺก\-อวิจาโร จ ธมฺมา อวิตกฺก\-วิจารมตฺตสฺส ธมฺมสฺส อาเสวน\-ปจฺจเยน ปจฺจโย. ปุริมา ปุริมา อวิตกฺก\-วิจารมตฺตา ขนฺธา จ วิจาโร จ ปจฺฉิมานํ ปจฺฉิมานํ อวิตกฺก\-วิจารมตฺตานํ ขนฺธานํ อาเสวน\-ปจฺจเยน ปจฺจโย. (1)}

\prose{อวิตกฺก\-วิจารมตฺโต จ อวิตกฺก\-อวิจาโร จ ธมฺมา อวิตกฺก\-อวิจารสฺส ธมฺมสฺส อาเสวน\-ปจฺจเยน ปจฺจโย. ปุริมา ปุริมา อวิตกฺก\-วิจารมตฺตา ขนฺธา จ วิจาโร จ ปจฺฉิมสฺส ปจฺฉิมสฺส วิจารสฺส อาเสวน\-ปจฺจเยน ปจฺจโย. (2)}

\prose{อวิตกฺก\-วิจารมตฺโต จ อวิตกฺก\-อวิจาโร จ ธมฺมา อวิตกฺก\-วิจารมตฺตสฺส จ อวิตกฺก\-อวิจารสฺส จ ธมฺมสฺส อาเสวน\-ปจฺจเยน ปจฺจโย. ปุริมา ปุริมา อวิตกฺก\-วิจารมตฺตา ขนฺธา จ วิจาโร จ ปจฺฉิมานํ ปจฺฉิมานํ อวิตกฺก\-วิจารมตฺตานํ ขนฺธานํ วิจารสฺส จ อาเสวน\-ปจฺจเยน ปจฺจโย. (3)}

\proseitem{113}{สวิตกฺก\-สวิจาโร จ อวิตกฺก\-วิจารมตฺโต จ ธมฺมา สวิตกฺก\-สวิจารสฺส ธมฺมสฺส อาเสวน\-ปจฺจเยน ปจฺจโย. ปุริมา ปุริมา สวิตกฺก\-สวิจารา ขนฺธา จ วิตกฺโก จ ปจฺฉิมานํ ปจฺฉิมานํ สวิตกฺก\-สวิจารานํ ขนฺธานํ อาเสวน\-ปจฺจเยน ปจฺจโย. อนุโลมญฺจ วิตกฺโก จ โคตฺรภุสฺส. อนุโลมญฺจ วิตกฺโก จ โวทานสฺส. โคตฺรภุ จ วิตกฺโก จ สวิตกฺก\-สวิจารสฺส มคฺคสฺส. โวทานญฺจ วิตกฺโก จ สวิตกฺก\-สวิจารสฺส มคฺคสฺส อาเสวน\-ปจฺจเยน ปจฺจโย. (1)}

\prose{สวิตกฺก\-สวิจาโร จ อวิตกฺก\-วิจารมตฺโต จ ธมฺมา อวิตกฺก\-วิจารมตฺตสฺส ธมฺมสฺส อาเสวน\-ปจฺจเยน ปจฺจโย. ปุริมา ปุริมา สวิตกฺก\-สวิจารา ขนฺธา จ วิตกฺโก จ ปจฺฉิมสฺส ปจฺฉิมสฺส วิตกฺกสฺส อาเสวน\-ปจฺจเยน ปจฺจโย. อวิตกฺก\-วิจารมตฺตสฺส ฌานสฺส ปริกมฺมญฺจ วิตกฺโก จ อวิตกฺก\-วิจารมตฺตสฺส ฌานสฺส อาเสวน\-ปจฺจเยน ปจฺจโย. โคตฺรภุ จ วิตกฺโก จ อวิตกฺก\-วิจารมตฺตสฺส มคฺคสฺส. โวทานญฺจ วิตกฺโก จ อวิตกฺก\-วิจารมตฺตสฺส มคฺคสฺส อาเสวน\-ปจฺจเยน ปจฺจโย. (2)}

\prose{สวิตกฺก\-สวิจาโร จ อวิตกฺก\-วิจารมตฺโต จ ธมฺมา อวิตกฺก\-อวิจารสฺส ธมฺมสฺส อาเสวน\-ปจฺจเยน ปจฺจโย. ทุติยสฺส ฌานสฺส ปริกมฺมญฺจ \csromanfolio{82}วิตกฺโก จ ทุติเย ฌาเน วิจารสฺส อาเสวน\-ปจฺจเยน ปจฺจโย ฯเปฯ. เนว\-สญฺญา\-นา\-สญฺญา\-ยตนสฺส ปริกมฺมญฺจ วิตกฺโก จ ฯเปฯ. ทิพฺพสฺส จกฺขุสฺส ปริกมฺมญฺจ วิตกฺโก จ ฯเปฯ. อนาคตํสญาณสฺส ปริกมฺมญฺจ วิตกฺโก จ อนาคตํสญาณสฺส อาเสวน\-ปจฺจเยน ปจฺจโย. โคตฺรภุ จ วิตกฺโก จ อวิตกฺก\-อวิจารสฺส มคฺคสฺส วิจารสฺส จ อาเสวน\-ปจฺจเยน ปจฺจโย. โวทานญฺจ วิตกฺโก จ อวิตกฺก\-อวิจารสฺส มคฺคสฺส วิจารสฺส จ อาเสวน\-ปจฺจเยน ปจฺจโย. (3)}

\prose{สวิตกฺก\-สวิจาโร จ อวิตกฺก\-วิจารมตฺโต จ ธมฺมา อวิตกฺก\-วิจารมตฺตสฺส จ อวิตกฺก\-อวิจารสฺส จ ธมฺมสฺส อาเสวน\-ปจฺจเยน ปจฺจโย. อวิตกฺก\-วิจารมตฺตสฺส ฌานสฺส ปริกมฺมญฺจ วิตกฺโก จ อวิตกฺก\-วิจารมตฺตสฺส ฌานสฺส วิจารสฺส จ. โคตฺรภุ จ วิตกฺโก จ อวิตกฺก\-วิจารมตฺตสฺส มคฺคสฺส วิจารสฺส จ. โวทานญฺจ วิตกฺโก จ อวิตกฺก\-วิจารมตฺตสฺส มคฺคสฺส วิจารสฺส จ อาเสวน\-ปจฺจเยน ปจฺจโย. (4)}

\prose{สวิตกฺก\-สวิจาโร จ อวิตกฺก\-วิจารมตฺโต จ ธมฺมา สวิตกฺก\-สวิจารสฺส จ อวิตกฺก\-วิจารมตฺตสฺส จ ธมฺมสฺส อาเสวน\-ปจฺจเยน ปจฺจโย. ปุริมา ปุริมา สวิตกฺก\-สวิจารา ขนฺธา จ วิตกฺโก จ ปจฺฉิมานํ ปจฺฉิมานํ สวิตกฺก\-สวิจารานํ ขนฺธานํ วิตกฺกสฺส จ อาเสวน\-ปจฺจเยน ปจฺจโย. อนุโลมญฺจ วิตกฺโก จ โคตฺรภุสฺส วิตกฺกสฺส จ. อนุโลมญฺจ วิตกฺโก จ โวทานสฺส วิตกฺกสฺส จ. โคตฺรภุ จ วิตกฺโก จ สวิตกฺก\-สวิจารสฺส มคฺคสฺส วิตกฺกสฺส จ. โวทานญฺจ วิตกฺโก จ สวิตกฺก\-สวิจารสฺส มคฺคสฺส วิตกฺกสฺส จ อาเสวน\-ปจฺจเยน ปจฺจโย. (5)}

\csromanheader{2}{\textbf{กมฺม}}
\proseitem{114}{สวิตกฺก\-สวิจาโร ธมฺโม สวิตกฺก\-สวิจารสฺส ธมฺมสฺส กมฺมปจฺจเยน ปจฺจโย. สหชาตา, นานากฺขณิกา.  สหชาตา สวิตกฺก\-สวิจารา เจตนา สมฺปยุตฺตกานํ ขนฺธานํ กมฺมปจฺจเยน ปจฺจโย. ปฏิสนฺธิ\-กฺขเณ สวิตกฺก\-สวิจารา เจตนา สมฺปยุตฺตกานํ ขนฺธานํ กมฺมปจฺจเยน ปจฺจโย.  นานากฺขณิกา สวิตกฺก\-สวิจารา เจตนา วิปากานํ สวิตกฺก\-สวิจารานํ ขนฺธานํ กมฺมปจฺจเยน ปจฺจโย. (1)}

\prose{\csromanfolio{83}สวิตกฺก\-สวิจาโร ธมฺโม อวิตกฺก\-วิจารมตฺตสฺส ธมฺมสฺส กมฺมปจฺจเยน ปจฺจโย. สหชาตา, นานากฺขณิกา.  สหชาตา สวิตกฺก\-สวิจารา เจตนา วิตกฺกสฺส กมฺมปจฺจเยน ปจฺจโย. ปฏิสนฺธิ\-กฺขเณ สวิตกฺก\-สวิจารา เจตนา วิตกฺกสฺส กมฺมปจฺจเยน ปจฺจโย.  นานากฺขณิกา สวิตกฺก\-สวิจารา เจตนา วิปากสฺส วิตกฺกสฺส กมฺมปจฺจเยน ปจฺจโย. (2)}

\prose{สวิตกฺก\-สวิจาโร ธมฺโม อวิตกฺก\-อวิจารสฺส ธมฺมสฺส กมฺมปจฺจเยน ปจฺจโย. สหชาตา, นานากฺขณิกา.  สหชาตา สวิตกฺก\-สวิจารา เจตนา จิตฺต\-สมุฏฺฐานานํ รูปานํ กมฺมปจฺจเยน ปจฺจโย. ปฏิสนฺธิ\-กฺขเณ สวิตกฺก\-สวิจารา เจตนา กฏตฺตารูปานํ กมฺมปจฺจเยน ปจฺจโย.  นานากฺขณิกา สวิตกฺก\-สวิจารา เจตนา วิปากานํ อวิตกฺก\-อวิจารานํ ขนฺธานํ กฏตฺตา จ รูปานํ กมฺมปจฺจเยน ปจฺจโย. (3)}

\prose{สวิตกฺก\-สวิจาโร ธมฺโม สวิตกฺก\-สวิจารสฺส จ อวิตกฺก\-อวิจารสฺส จ ธมฺมสฺส กมฺมปจฺจเยน ปจฺจโย. สหชาตา, นานากฺขณิกา.  สหชาตา สวิตกฺก\-สวิจารา เจตนา สมฺปยุตฺตกานํ ขนฺธานํ จิตฺตสมุฏฺฐานานญฺจ รูปานํ กมฺมปจฺจเยน ปจฺจโย. ปฏิสนฺธิ\-กฺขเณ สวิตกฺก\-สวิจารา เจตนา สมฺปยุตฺตกานํ ขนฺธานํ กฏตฺตา จ รูปานํ กมฺมปจฺจเยน ปจฺจโย.  นานากฺขณิกา สวิตกฺก\-สวิจารา เจตนา วิปากานํ สวิตกฺก\-สวิจารานํ ขนฺธานํ กฏตฺตา จ รูปานํ กมฺมปจฺจเยน ปจฺจโย. (4)}

\prose{สวิตกฺก\-สวิจาโร ธมฺโม อวิตกฺก\-วิจารมตฺตสฺส จ อวิตกฺก\-อวิจารสฺส จ ธมฺมสฺส กมฺมปจฺจเยน ปจฺจโย. สหชาตา, นานากฺขณิกา.  สหชาตา สวิตกฺก\-สวิจารา เจตนา วิตกฺกสฺส จ จิตฺตสมุฏฺฐานานญฺจ รูปานํ กมฺมปจฺจเยน ปจฺจโย. ปฏิสนฺธิ\-กฺขเณ สวิตกฺก\-สวิจารา เจตนา วิตกฺกสฺส จ กฏตฺตา จ รูปานํ กมฺมปจฺจเยน ปจฺจโย.  นานากฺขณิกา สวิตกฺก\-สวิจารา เจตนา วิปากสฺส วิตกฺกสฺส จ กฏตฺตา จ รูปานํ กมฺมปจฺจเยน ปจฺจโย. (5)}

\prose{สวิตกฺก\-สวิจาโร ธมฺโม สวิตกฺก\-สวิจารสฺส จ อวิตกฺก\-วิจารมตฺตสฺส จ ธมฺมสฺส กมฺมปจฺจเยน ปจฺจโย. สหชาตา, นานากฺขณิกา.  สหชาตา สวิตกฺก\-สวิจารา เจตนา สมฺปยุตฺตกานํ ขนฺธานํ วิตกฺกสฺส จ กมฺมปจฺจเยน ปจฺจโย. ปฏิสนฺธิ\-กฺขเณ สวิตกฺก\-สวิจารา เจตนา \csromanfolio{84}สมฺปยุตฺตกานํ ขนฺธานํ วิตกฺกสฺส จ กมฺมปจฺจเยน ปจฺจโย.  นานากฺขณิกา สวิตกฺก\-สวิจารา เจตนา วิปากานํ สวิตกฺก\-สวิจารานํ ขนฺธานํ วิตกฺกสฺส จ กมฺมปจฺจเยน ปจฺจโย. (6)}

\prose{สวิตกฺก\-สวิจาโร ธมฺโม สวิตกฺก\-สวิจารสฺส จ อวิตกฺก\-วิจารมตฺตสฺส จ อวิตกฺก\-อวิจารสฺส จ ธมฺมสฺส กมฺมปจฺจเยน ปจฺจโย. สหชาตา, นานากฺขณิกา.  สหชาตา สวิตกฺก\-สวิจารา เจตนา สมฺปยุตฺตกานํ ขนฺธานํ วิตกฺกสฺส จ จิตฺตสมุฏฺฐานานญฺจ รูปานํ กมฺมปจฺจเยน ปจฺจโย. ปฏิสนฺธิ\-กฺขเณ สวิตกฺก\-สวิจารา เจตนา สมฺปยุตฺตกานํ ขนฺธานํ วิตกฺกสฺส จ กฏตฺตา จ รูปานํ กมฺมปจฺจเยน ปจฺจโย.  นานากฺขณิกา สวิตกฺก\-สวิจารา เจตนา วิปากานํ สวิตกฺก\-สวิจารานํ ขนฺธานํ วิตกฺกสฺส จ กฏตฺตา จ รูปานํ กมฺมปจฺจเยน ปจฺจโย. (7)}

\proseitem{115}{อวิตกฺก\-วิจารมตฺโต ธมฺโม อวิตกฺก\-วิจารมตฺตสฺส ธมฺมสฺส กมฺมปจฺจเยน ปจฺจโย. สหชาตา, นานากฺขณิกา.  สหชาตา อวิตกฺก\-วิจารมตฺตา เจตนา สมฺปยุตฺตกานํ ขนฺธานํ กมฺมปจฺจเยน ปจฺจโย. ปฏิสนฺธิ\-กฺขเณ ฯเปฯ. นานากฺขณิกา อวิตกฺก\-วิจารมตฺตา เจตนา วิปากานํ อวิตกฺก\-วิจารมตฺตานํ ขนฺธานํ กมฺมปจฺจเยน ปจฺจโย. (1)}

\prose{อวิตกฺก\-วิจารมตฺโต ธมฺโม อวิตกฺก\-อวิจารสฺส ธมฺมสฺส กมฺมปจฺจเยน ปจฺจโย. สหชาตา, นานากฺขณิกา.  สหชาตา อวิตกฺก\-วิจารมตฺตา เจตนา วิจารสฺส จ จิตฺตสมุฏฺฐานานญฺจ รูปานํ กมฺมปจฺจเยน ปจฺจโย. ปฏิสนฺธิ\-กฺขเณ อวิตกฺก\-วิจารมตฺตา เจตนา วิปากสฺส วิจารสฺส จ กฏตฺตา จ รูปานํ กมฺมปจฺจเยน ปจฺจโย.  นานากฺขณิกา อวิตกฺก\-วิจารมตฺตา เจตนา วิปากสฺส วิจารสฺส จ กฏตฺตา จ รูปานํ กมฺมปจฺจเยน ปจฺจโย. (2)}

\prose{อวิตกฺก\-วิจารมตฺโต ธมฺโม อวิตกฺก\-วิจารมตฺตสฺส จ อวิตกฺก\-อวิจารสฺส จ ธมฺมสฺส กมฺมปจฺจเยน ปจฺจโย. สหชาตา นานากฺขณิกา.  สหชาตา อวิตกฺก\-วิจารมตฺตา เจตนา สมฺปยุตฺตกานํ ขนฺธานํ วิจารสฺส จ จิตฺตสมุฏฺฐานานญฺจ รูปานํ กมฺมปจฺจเยน ปจฺจโย. ปฏิสนฺธิ\-กฺขเณ อวิตกฺก\-วิจารมตฺตา เจตนา สมฺปยุตฺตกานํ ขนฺธานํ วิจารสฺส จ กฏตฺตา จ รูปานํ กมฺมปจฺจเยน ปจฺจโย.   \csromanfolio{85}นานากฺขณิกา อวิตกฺก\-วิจารมตฺตา เจตนา วิปากานํ อวิตกฺก\-วิจารมตฺตานํ ขนฺธานํ วิจารสฺส จ กฏตฺตา จ รูปานํ กมฺมปจฺจเยน ปจฺจโย. (3)}

\prose{อวิตกฺก\-อวิจาโร ธมฺโม อวิตกฺก\-อวิจารสฺส ธมฺมสฺส กมฺมปจฺจเยน ปจฺจโย. สหชาตา, นานากฺขณิกา.  สหชาตา อวิตกฺก\-อวิจารา เจตนา สมฺปยุตฺตกานํ ขนฺธานํ จิตฺตสมุฏฺฐานานญฺจ รูปานํ กมฺมปจฺจเยน ปจฺจโย. ปฏิสนฺธิ\-กฺขเณ อวิตกฺก\-อวิจารา เจตนา สมฺปยุตฺตกานํ ขนฺธานํ กฏตฺตา จ รูปานํ กมฺมปจฺจเยน ปจฺจโย.  นานากฺขณิกา อวิตกฺก\-อวิจารา เจตนา วิปากานํ อวิตกฺก\-อวิจารานํ ขนฺธานํ กฏตฺตา จ รูปานํ กมฺมปจฺจเยน ปจฺจโย. (1)}

\csromanheader{2}{\textbf{วิปาก}}
\proseitem{116}{สวิตกฺก\-สวิจาโร ธมฺโม สวิตกฺก\-สวิจารสฺส ธมฺมสฺส วิปาก\-ปจฺจเยน ปจฺจโย. วิปาโก สวิตกฺก\-สวิจาโร เอโก ขนฺโธ ติณฺณนฺนํ ขนฺธานํ วิปาก\-ปจฺจเยน ปจฺจโย ฯเปฯ. ปฏิสนฺธิ\-กฺขเณ สวิตกฺก\-สวิจาโร เอโก ขนฺโธ ติณฺณนฺนํ ขนฺธานํ วิปาก\-ปจฺจเยน ปจฺจโย ฯเปฯ.}

\prose{สวิตกฺก\-สวิจาโร ธมฺโม อวิตกฺก\-วิจารมตฺตสฺส ธมฺมสฺส วิปาก\-ปจฺจเยน ปจฺจโย. วิปากา สวิตกฺก\-สวิจารา ขนฺธา วิตกฺกสฺส วิปาก\-ปจฺจเยน ปจฺจโย. ปฏิสนฺธิ\-กฺขเณ ฯเปฯ.}

\vspace{\tipitakaparskip}
\csromancenter{(สวิตกฺกสวิจาร\-มูลกา สตฺตปิ ปญฺหา ปริปุณฺณา.)}
\prose{อวิตกฺก\-วิจารมตฺโต ธมฺโม อวิตกฺก\-วิจารมตฺตสฺส ธมฺมสฺส วิปาก\-ปจฺจเยน ปจฺจโย. วิปาโก อวิตกฺก\-วิจารมตฺโต เอโก ขนฺโธ ติณฺณนฺนํ ขนฺธานํ วิปาก\-ปจฺจเยน ปจฺจโย ฯเปฯ เทฺว ขนฺธา ฯเปฯ. ปฏิสนฺธิ\-กฺขเณ วิปาโก อวิตกฺก\-วิจารมตฺโต เอโก ขนฺโธ ฯเปฯ.}

\prose{(อวิตกฺกวิจารมตฺต\-มูลกา ปญฺจ ปญฺหา กาตพฺพา, วิปากนฺติ นิยาเมตพฺพา.)}

\proseitem{117}{อวิตกฺก\-อวิจาโร ธมฺโม อวิตกฺก\-อวิจารสฺส ธมฺมสฺส วิปาก\-ปจฺจเยน ปจฺจโย. วิปาโก อวิตกฺก\-อวิจาโร เอโก ขนฺโธ ติณฺณนฺนํ จิตฺตสมุฏฺฐานานญฺจ รูปานํ วิปาก\-ปจฺจเยน ปจฺจโย ฯเปฯ เทฺว \csromanfolio{86}ขนฺธา ฯเปฯ. วิปาโก วิจาโร จิตฺต\-สมุฏฺฐานานํ รูปานํ วิปาก\-ปจฺจเยน ปจฺจโย. ปฏิสนฺธิ\-กฺขเณ ฯเปฯ. ขนฺธา วตฺถุสฺส วิปาก\-ปจฺจเยน ปจฺจโย. วิจาโร วตฺถุสฺส วิปาก\-ปจฺจเยน ปจฺจโย. (1)}

\prose{อวิตกฺก\-อวิจาโร ธมฺโม อวิตกฺก\-วิจารมตฺตสฺส ธมฺมสฺส วิปาก\-ปจฺจเยน ปจฺจโย. วิปาโก วิจาโร อวิตกฺก\-วิจารมตฺตานํ ขนฺธานํ วิปาก\-ปจฺจเยน ปจฺจโย. ปฏิสนฺธิ\-กฺขเณ วิปาโก วิจาโร อวิตกฺก\-วิจารมตฺตานํ ขนฺธานํ วิปาก\-ปจฺจเยน ปจฺจโย. (2)}

\prose{อวิตกฺก\-อวิจาโร ธมฺโม อวิตกฺก\-วิจารมตฺตสฺส จ อวิตกฺก\-อวิจารสฺส จ ธมฺมสฺส วิปาก\-ปจฺจเยน ปจฺจโย. วิปาโก วิจาโร อวิตกฺก\-วิจารมตฺตานํ ขนฺธานํ จิตฺตสมุฏฺฐานานญฺจ รูปานํ วิปาก\-ปจฺจเยน ปจฺจโย. ปฏิสนฺธิ\-กฺขเณ วิปาโก วิจาโร อวิตกฺก\-วิจารมตฺตานํ ขนฺธานํ กฏตฺตา จ รูปานํ วิปาก\-ปจฺจเยน ปจฺจโย. (3)}

\proseitem{118}{อวิตกฺก\-วิจารมตฺโต จ อวิตกฺก\-อวิจาโร จ ธมฺมา อวิตกฺก\-วิจารมตฺตสฺส ธมฺมสฺส วิปาก\-ปจฺจเยน ปจฺจโย. วิปาโก อวิตกฺก\-วิจารมตฺโต เอโก ขนฺโธ จ วิจาโร จ ติณฺณนฺนํ ขนฺธานํ วิปาก\-ปจฺจเยน ปจฺจโย ฯเปฯ เทฺว ขนฺธา ฯเปฯ. ปฏิสนฺธิ\-กฺขเณ ฯเปฯ. (1)}

\prose{อวิตกฺก\-วิจารมตฺโต จ อวิตกฺก\-อวิจาโร จ ธมฺมา อวิตกฺก\-อวิจารสฺส ธมฺมสฺส วิปาก\-ปจฺจเยน ปจฺจโย. วิปากา อวิตกฺก\-วิจารมตฺตา ขนฺธา จ วิจาโร จ จิตฺต\-สมุฏฺฐานานํ รูปานํ วิปาก\-ปจฺจเยน ปจฺจโย. ปฏิสนฺธิ\-กฺขเณ อวิตกฺก\-วิจารมตฺตา ขนฺธา จ วิจาโร จ กฏตฺตารูปานํ วิปาก\-ปจฺจเยน ปจฺจโย. (2)}

\prose{อวิตกฺก\-วิจารมตฺโต จ อวิตกฺก\-อวิจาโร จ ธมฺมา อวิตกฺก\-วิจารมตฺตสฺส จ อวิตกฺก\-อวิจารสฺส จ ธมฺมสฺส วิปาก\-ปจฺจเยน ปจฺจโย. วิปาโก อวิตกฺก\-วิจารมตฺโต เอโก ขนฺโธ จ วิจาโร จ ติณฺณนฺนํ ขนฺธานํ จิตฺตสมุฏฺฐานานญฺจ รูปานํ วิปาก\-ปจฺจเยน ปจฺจโย ฯเปฯ. ปฏิสนฺธิ\-กฺขเณ วิปาโก อวิตกฺก\-วิจารมตฺโต เอโก ขนฺโธ จ วิจาโร จ ติณฺณนฺนํ ขนฺธานํ กฏตฺตา จ รูปานํ วิปาก\-ปจฺจเยน ปจฺจโย. (3)}

\prose{สวิตกฺก\-สวิจาโร จ อวิตกฺก\-วิจารมตฺโต จ ธมฺมา สวิตกฺก\-สวิจารสฺส ธมฺมสฺส วิปาก\-ปจฺจเยน ปจฺจโย. วิปาโก สวิตกฺก\-สวิจาโร \csromanfolio{87}เอโก ขนฺโธ จ วิตกฺโก จ ติณฺณนฺนํ ขนฺธานํ วิปาก\-ปจฺจเยน ปจฺจโย ฯเปฯ เทฺว ขนฺธา ฯเปฯ. ปฏิสนฺธิ\-กฺขเณ ฯเปฯ. (1)}

\prose{สวิตกฺก\-สวิจาโร จ อวิตกฺก\-วิจารมตฺโต จ ธมฺมา อวิตกฺก\-อวิจารสฺส ธมฺมสฺส วิปาก\-ปจฺจเยน ปจฺจโย. วิปากา สวิตกฺก\-สวิจารา ขนฺธา จ วิตกฺโก จ จิตฺต\-สมุฏฺฐานานํ รูปานํ วิปาก\-ปจฺจเยน ปจฺจโย. ปฏิสนฺธิ\-กฺขเณ ฯเปฯ. (2)}

\prose{สวิตกฺก\-สวิจาโร จ อวิตกฺก\-วิจารมตฺโต จ ธมฺมา สวิตกฺก\-สวิจารสฺส จ อวิตกฺก\-อวิจารสฺส จ ธมฺมสฺส วิปาก\-ปจฺจเยน ปจฺจโย. วิปาโก สวิตกฺก\-สวิจาโร เอโก ขนฺโธ จ วิตกฺโก จ ติณฺณนฺนํ ขนฺธานํ จิตฺตสมุฏฺฐานานญฺจ รูปานํ วิปาก\-ปจฺจเยน ปจฺจโย. ปฏิสนฺธิ\-กฺขเณ ฯเปฯ. (3)}

\csromanheader{2}{\textbf{อาหาร}}
\proseitem{119}{สวิตกฺก\-สวิจาโร ธมฺโม สวิตกฺก\-สวิจารสฺส ธมฺมสฺส อาหารปจฺจเยน ปจฺจโย. สวิตกฺก\-สวิจารา อาหารา สมฺปยุตฺตกานํ ขนฺธานํ อาหารปจฺจเยน ปจฺจโย. ปฏิสนฺธิ\-กฺขเณ ฯเปฯ. (1)}

\prose{สวิตกฺก\-สวิจาโร ธมฺโม อวิตกฺก\-วิจารมตฺตสฺส ธมฺมสฺส อาหารปจฺจเยน ปจฺจโย. สวิตกฺก\-สวิจารา อาหารา วิตกฺกสฺส อาหารปจฺจเยน ปจฺจโย. ปฏิสนฺธิ\-กฺขเณ ฯเปฯ.}

\prose{(สวิตกฺกสวิจาร\-มูลกา อิมินา การเณน สตฺต ปญฺหา วิภชิตพฺพา.)}

\prose{อวิตกฺก\-วิจารมตฺโต ธมฺโม อวิตกฺก\-วิจารมตฺตสฺส ธมฺมสฺส อาหารปจฺจเยน ปจฺจโย. อวิตกฺก\-วิจารมตฺตา อาหารา สมฺปยุตฺตกานํ ขนฺธานํ อาหารปจฺจเยน ปจฺจโย. ปฏิสนฺธิ\-กฺขเณ ฯเปฯ. (1)}

\prose{อวิตกฺก\-วิจารมตฺโต ธมฺโม อวิตกฺก\-อวิจารสฺส ธมฺมสฺส อาหารปจฺจเยน ปจฺจโย. อวิตกฺก\-วิจารมตฺตา อาหารา วิจารสฺส จ จิตฺตสมุฏฺฐานานญฺจ รูปานํ อาหารปจฺจเยน ปจฺจโย. ปฏิสนฺธิ\-กฺขเณ ฯเปฯ. (2)}

\prose{อวิตกฺก\-วิจารมตฺโต ธมฺโม อวิตกฺก\-วิจารมตฺตสฺส จ อวิตกฺก\-อวิจารสฺส จ ธมฺมสฺส อาหารปจฺจเยน ปจฺจโย. อวิตกฺก\-วิจารมตฺตา \csromanfolio{88}อาหารา สมฺปยุตฺตกานํ ขนฺธานํ วิจารสฺส จ จิตฺตสมุฏฺฐานานญฺจ รูปานํ อาหารปจฺจเยน ปจฺจโย. ปฏิสนฺธิ\-กฺขเณ ฯเปฯ. (3)}

\prose{อวิตกฺก\-อวิจาโร ธมฺโม อวิตกฺก\-อวิจารสฺส ธมฺมสฺส อาหารปจฺจเยน ปจฺจโย. อวิตกฺก\-อวิจารา อาหารา สมฺปยุตฺตกานํ ขนฺธานํ จิตฺตสมุฏฺฐานานญฺจ รูปานํ อาหารปจฺจเยน ปจฺจโย. ปฏิสนฺธิ\-กฺขเณ ฯเปฯ. กพฬีกาโร อาหาโร อิมสฺส กายสฺส อาหารปจฺจเยน ปจฺจโย.}

\csromanheader{2}{\textbf{อินฺทฺริย}}
\proseitem{120}{สวิตกฺก\-สวิจาโร ธมฺโม สวิตกฺก\-สวิจารสฺส ธมฺมสฺส อินฺทฺริย\-ปจฺจเยน ปจฺจโย. สวิตกฺก\-สวิจารา อินฺทฺริยา สมฺปยุตฺตกานํ ขนฺธานํ อินฺทฺริย\-ปจฺจเยน ปจฺจโย. ปฏิสนฺธิ\-กฺขเณ ฯเปฯ. (1)}

\prose{สวิตกฺก\-สวิจาโร ธมฺโม อวิตกฺก\-วิจารมตฺตสฺส ธมฺมสฺส อินฺทฺริย\-ปจฺจเยน ปจฺจโย. สวิตกฺก\-สวิจารา อินฺทฺริยา วิตกฺกสฺส อินฺทฺริย\-ปจฺจเยน ปจฺจโย. ปฏิสนฺธิ\-กฺขเณ ฯเปฯ. (7)}

\prose{(สวิตกฺกสวิจาร\-มูลกา สตฺต ปญฺหา อิมินา การเณน วิภชิตพฺพา.)}

\prose{อวิตกฺก\-วิจารมตฺโต ธมฺโม อวิตกฺก\-วิจารมตฺตสฺส ธมฺมสฺส อินฺทฺริย\-ปจฺจเยน ปจฺจโย. อวิตกฺก\-วิจารมตฺตา อินฺทฺริยา สมฺปยุตฺตกานํ ขนฺธานํ อินฺทฺริย\-ปจฺจเยน ปจฺจโย. ปฏิสนฺธิ\-กฺขเณ ฯเปฯ. (1)}

\prose{อวิตกฺก\-วิจารมตฺโต ธมฺโม อวิตกฺก\-อวิจารสฺส ธมฺมสฺส อินฺทฺริย\-ปจฺจเยน ปจฺจโย. อวิตกฺก\-วิจารมตฺตา อินฺทฺริยา วิจารสฺส จ จิตฺตสมุฏฺฐานานญฺจ รูปานํ อินฺทฺริย\-ปจฺจเยน ปจฺจโย. ปฏิสนฺธิ\-กฺขเณ ฯเปฯ. (2)}

\prose{อวิตกฺก\-วิจารมตฺโต ธมฺโม อวิตกฺก\-วิจารมตฺตสฺส จ อวิตกฺก\-อวิจารสฺส จ ธมฺมสฺส อินฺทฺริย\-ปจฺจเยน ปจฺจโย. อวิตกฺก\-วิจารมตฺตา อินฺทฺริยา สมฺปยุตฺตกานํ ขนฺธานํ วิจารสฺส จ จิตฺตสมุฏฺฐานานญฺจ รูปานํ อินฺทฺริย\-ปจฺจเยน ปจฺจโย. ปฏิสนฺธิ\-กฺขเณ ฯเปฯ. (3)}

\prose{อวิตกฺก\-อวิจาโร ธมฺโม อวิตกฺก\-อวิจารสฺส ธมฺมสฺส อินฺทฺริย\-ปจฺจเยน ปจฺจโย. อวิตกฺก\-อวิจารา อินฺทฺริยา สมฺปยุตฺตกานํ ขนฺธานํ จิตฺตสมุฏฺฐานานญฺจ รูปานํ อินฺทฺริย\-ปจฺจเยน ปจฺจโย. ปฏิสนฺธิ\-กฺขเณ ฯเปฯ. จกฺขุนฺทฺริยํ \csromanfolio{89}จกฺขุ\-วิญฺญาณสฺส ฯเปฯ. กายินฺทฺริยํ กายวิญฺญาณสฺส อินฺทฺริย\-ปจฺจเยน ปจฺจโย. รูปชีวิตินฺทฺริยํ กฏตฺตารูปานํ อินฺทฺริย\-ปจฺจเยน ปจฺจโย. (1)}

\csromanheader{2}{\textbf{ฌาน}}
\proseitem{121}{สวิตกฺก\-สวิจาโร ธมฺโม สวิตกฺก\-สวิจารสฺส ธมฺมสฺส ฌานปจฺจเยน ปจฺจโย. สวิตกฺก\-สวิจารานิ ฌานงฺคานิ สมฺปยุตฺตกานํ ขนฺธานํ ฌานปจฺจเยน ปจฺจโย. ปฏิสนฺธิ\-กฺขเณ ฯเปฯ. (7)}

\prose{(สวิตกฺกสวิจาร\-มูลกา สตฺต ปญฺหา อิมินา การเณน วิภชิตพฺพา.)}

\prose{อวิตกฺก\-วิจารมตฺโต ธมฺโม อวิตกฺก\-วิจารมตฺตสฺส ธมฺมสฺส ฌานปจฺจเยน ปจฺจโย. อวิตกฺก\-วิจารมตฺตานิ ฌานงฺคานิ สมฺปยุตฺตกานํ ขนฺธานํ ฌานปจฺจเยน ปจฺจโย. ปฏิสนฺธิ\-กฺขเณ ฯเปฯ. (5)}

\prose{(อวิตกฺกวิจารมตฺต\-มูลกา ปญฺจ ปญฺหา อิมินา การเณน วิภชิตพฺพา.)}

\prose{อวิตกฺก\-อวิจาโร ธมฺโม อวิตกฺก\-อวิจารสฺส ธมฺมสฺส ฌานปจฺจเยน ปจฺจโย. อวิตกฺก\-อวิจารานิ ฌานงฺคานิ สมฺปยุตฺตกานํ ขนฺธานํ จิตฺตสมุฏฺฐานานญฺจ รูปานํ ฌานปจฺจเยน ปจฺจโย. วิจาโร จิตฺต\-สมุฏฺฐานานํ รูปานํ ฌานปจฺจเยน ปจฺจโย. ปฏิสนฺธิ\-กฺขเณ วิจาโร กฏตฺตารูปานํ ฌานปจฺจเยน ปจฺจโย. วิจาโร วตฺถุสฺส ฌานปจฺจเยน ปจฺจโย. (1)}

\prose{อวิตกฺก\-อวิจาโร ธมฺโม อวิตกฺก\-วิจารมตฺตสฺส ธมฺมสฺส ฌานปจฺจเยน ปจฺจโย. วิจาโร อวิตกฺก\-วิจารมตฺตานํ ขนฺธานํ ฌานปจฺจเยน ปจฺจโย. ปฏิสนฺธิ\-กฺขเณ วิจาโร อวิตกฺก\-วิจารมตฺตานํ ขนฺธานํ ฌานปจฺจเยน ปจฺจโย. (2)}

\prose{อวิตกฺก\-อวิจาโร ธมฺโม อวิตกฺก\-วิจารมตฺตสฺส จ อวิตกฺก\-อวิจารสฺส จ ธมฺมสฺส ฌานปจฺจเยน ปจฺจโย. วิจาโร อวิตกฺก\-วิจารมตฺตานํ ขนฺธานํ จิตฺตสมุฏฺฐานานญฺจ รูปานํ ฌานปจฺจเยน ปจฺจโย. ปฏิสนฺธิ\-กฺขเณ วิจาโร อวิตกฺก\-วิจารมตฺตานํ ขนฺธานํ กฏตฺตา จ รูปานํ ฌานปจฺจเยน ปจฺจโย. (3)}

\proseitem{122}{\csromanfolio{90}อวิตกฺก\-วิจารมตฺโต จ อวิตกฺก\-อวิจาโร จ ธมฺมา อวิตกฺก\-วิจารมตฺตสฺส ธมฺมสฺส ฌานปจฺจเยน ปจฺจโย. อวิตกฺก\-วิจารมตฺตานิ ฌานงฺคานิ วิจาโร จ สมฺปยุตฺตกานํ ขนฺธานํ ฌานปจฺจเยน ปจฺจโย. ปฏิสนฺธิ\-กฺขเณ ฯเปฯ. (1)}

\prose{อวิตกฺก\-วิจารมตฺโต จ อวิตกฺก\-อวิจาโร จ ธมฺมา อวิตกฺก\-อวิจารสฺส ธมฺมสฺส ฌานปจฺจเยน ปจฺจโย. อวิตกฺก\-วิจารมตฺตานิ ฌานงฺคานิ วิจาโร จ จิตฺต\-สมุฏฺฐานานํ รูปานํ ฌานปจฺจเยน ปจฺจโย. ปฏิสนฺธิ\-กฺขเณ ฯเปฯ. (2)}

\prose{อวิตกฺก\-วิจารมตฺโต จ อวิตกฺก\-อวิจาโร จ ธมฺมา อวิตกฺก\-วิจารมตฺตสฺส จ อวิตกฺก\-อวิจารสฺส จ ธมฺมสฺส ฌานปจฺจเยน ปจฺจโย. อวิตกฺก\-วิจารมตฺตานิ ฌานงฺคานิ วิจาโร จ สมฺปยุตฺตกานํ ขนฺธานํ จิตฺตสมุฏฺฐานานญฺจ รูปานํ ฌานปจฺจเยน ปจฺจโย. ปฏิสนฺธิ\-กฺขเณ อวิตกฺก\-วิจารมตฺตานิ ฌานงฺคานิ วิจาโร จ สมฺปยุตฺตกานํ ขนฺธานํ กฏตฺตา จ รูปานํ ฌานปจฺจเยน ปจฺจโย. (3)}

\prose{สวิตกฺก\-สวิจาโร จ อวิตกฺก\-วิจารมตฺโต จ ธมฺมา สวิตกฺก\-สวิจารสฺส ธมฺมสฺส ฌานปจฺจเยน ปจฺจโย. สวิตกฺก\-สวิจารานิ ฌานงฺคานิ วิตกฺโก จ สมฺปยุตฺตกานํ ขนฺธานํ ฌานปจฺจเยน ปจฺจโย. ปฏิสนฺธิ\-กฺขเณ ฯเปฯ. (1)}

\prose{สวิตกฺก\-สวิจาโร จ อวิตกฺก\-วิจารมตฺโต จ ธมฺมา อวิตกฺก\-อวิจารสฺส ธมฺมสฺส ฌานปจฺจเยน ปจฺจโย. สวิตกฺก\-สวิจารานิ ฌานงฺคานิ วิตกฺโก จ จิตฺต\-สมุฏฺฐานานํ รูปานํ ฌานปจฺจเยน ปจฺจโย. ปฏิสนฺธิ\-กฺขเณ ฯเปฯ. (2)}

\prose{สวิตกฺก\-สวิจาโร จ อวิตกฺก\-วิจารมตฺโต จ ธมฺมา สวิตกฺก\-สวิจารสฺส จ อวิตกฺก\-อวิจารสฺส จ ธมฺมสฺส ฌานปจฺจเยน ปจฺจโย. สวิตกฺก\-สวิจารานิ ฌานงฺคานิ วิตกฺโก จ สมฺปยุตฺตกานํ ขนฺธานํ จิตฺตสมุฏฺฐานานญฺจ รูปานํ ฌานปจฺจเยน ปจฺจโย. ปฏิสนฺธิ\-กฺขเณ ฯเปฯ. (3)}

\csromanheader{2}{\textbf{มคฺค}}
\proseitem{123}{\csromanfolio{91}สวิตกฺก\-สวิจาโร ธมฺโม สวิตกฺก\-สวิจารสฺส ธมฺมสฺส มคฺคปจฺจเยน ปจฺจโย. สวิตกฺก\-สวิจารานิ มคฺคงฺคานิ สมฺปยุตฺตกานํ ขนฺธานํ มคฺคปจฺจเยน ปจฺจโย. ปฏิสนฺธิ\-กฺขเณ ฯเปฯ.}

\prose{(สวิตกฺกสวิจาร\-มูลกา สตฺต ปญฺหา อิมินา การเณน วิภชิตพฺพา.)}

\prose{อวิตกฺก\-วิจารมตฺโต ธมฺโม อวิตกฺก\-วิจารมตฺตสฺส ธมฺมสฺส มคฺคปจฺจเยน ปจฺจโย. อวิตกฺก\-วิจารมตฺตานิ มคฺคงฺคานิ สมฺปยุตฺตกานํ ขนฺธานํ มคฺคปจฺจเยน ปจฺจโย. ปฏิสนฺธิ\-กฺขเณ ฯเปฯ.}

\prose{(อวิตกฺกวิจารมตฺต\-มูลกา ปญฺจ ปญฺหา อิมินา การเณน วิภชิตพฺพา.)}

\prose{อวิตกฺก\-อวิจาโร ธมฺโม อวิตกฺก\-อวิจารสฺส ธมฺมสฺส มคฺคปจฺจเยน ปจฺจโย. อวิตกฺก\-อวิจารานิ มคฺคงฺคานิ สมฺปยุตฺตกานํ ขนฺธานํ จิตฺตสมุฏฺฐานานญฺจ รูปานํ มคฺคปจฺจเยน ปจฺจโย. ปฏิสนฺธิ\-กฺขเณ ฯเปฯ. (1)}

\prose{สวิตกฺก\-สวิจาโร จ อวิตกฺก\-วิจารมตฺโต จ ธมฺมา สวิตกฺก\-สวิจารสฺส ธมฺมสฺส มคฺคปจฺจเยน ปจฺจโย. สวิตกฺก\-สวิจารานิ มคฺคงฺคานิ วิตกฺโก จ สมฺปยุตฺตกานํ ขนฺธานํ มคฺคปจฺจเยน ปจฺจโย. ปฏิสนฺธิ\-กฺขเณ ฯเปฯ. (1)}

\prose{สวิตกฺก\-สวิจาโร จ อวิตกฺก\-วิจารมตฺโต จ ธมฺมา อวิตกฺก\-อวิจารสฺส ธมฺมสฺส มคฺคปจฺจเยน ปจฺจโย. สวิตกฺก\-สวิจารานิ มคฺคงฺคานิ วิตกฺโก จ จิตฺต\-สมุฏฺฐานานํ รูปานํ มคฺคปจฺจเยน ปจฺจโย. ปฏิสนฺธิ\-กฺขเณ ฯเปฯ. (2)}

\prose{สวิตกฺก\-สวิจาโร จ อวิตกฺก\-วิจารมตฺโต จ ธมฺมา สวิตกฺก\-สวิจารสฺส จ อวิตกฺก\-อวิจารสฺส จ ธมฺมสฺส มคฺคปจฺจเยน ปจฺจโย. สวิตกฺก\-สวิจารานิ มคฺคงฺคานิ วิตกฺโก จ สมฺปยุตฺตกานํ ขนฺธานํ จิตฺตสมุฏฺฐานานญฺจ รูปานํ มคฺคปจฺจเยน ปจฺจโย. ปฏิสนฺธิ\-กฺขเณ ฯเปฯ. (3)}

\csromanheader{2}{\textbf{สมฺปยุตฺต}}
\proseitem{124}{สวิตกฺก\-สวิจาโร ธมฺโม สวิตกฺก\-สวิจารสฺส ธมฺมสฺส สมฺปยุตฺต\-ปจฺจเยน ปจฺจโย. สวิตกฺก\-สวิจาโร เอโก ขนฺโธ ติณฺณนฺนํ \csromanfolio{92}ขนฺธานํ สมฺปยุตฺต\-ปจฺจเยน ปจฺจโย ฯเปฯ เทฺว ขนฺธา ทฺวินฺนํ ขนฺธานํ ฯเปฯ. ปฏิสนฺธิ\-กฺขเณ ฯเปฯ. (1)}

\prose{สวิตกฺก\-สวิจาโร ธมฺโม อวิตกฺก\-วิจารมตฺตสฺส ธมฺมสฺส สมฺปยุตฺต\-ปจฺจเยน ปจฺจโย. สวิตกฺก\-สวิจารา ขนฺธา วิตกฺกสฺส สมฺปยุตฺต\-ปจฺจเยน ปจฺจโย. ปฏิสนฺธิ\-กฺขเณ ฯเปฯ. (2)}

\prose{สวิตกฺก\-สวิจาโร ธมฺโม สวิตกฺก\-สวิจารสฺส จ อวิตกฺก\-วิจารมตฺตสฺส จ ธมฺมสฺส สมฺปยุตฺต\-ปจฺจเยน ปจฺจโย. สวิตกฺก\-สวิจาโร เอโก ขนฺโธ ติณฺณนฺนํ ขนฺธานํ วิตกฺกสฺส จ สมฺปยุตฺต\-ปจฺจเยน ปจฺจโย ฯเปฯ เทฺว ขนฺธา ทฺวินฺนํ ขนฺธานํ วิตกฺกสฺส จ ฯเปฯ. ปฏิสนฺธิ\-กฺขเณ ฯเปฯ. (3)}

\proseitem{125}{อวิตกฺก\-วิจารมตฺโต ธมฺโม อวิตกฺก\-วิจารมตฺตสฺส ธมฺมสฺส สมฺปยุตฺต\-ปจฺจเยน ปจฺจโย. อวิตกฺก\-วิจารมตฺโต เอโก ขนฺโธ ติณฺณนฺนํ ขนฺธานํ สมฺปยุตฺต\-ปจฺจเยน ปจฺจโย ฯเปฯ เทฺว ขนฺธา ทฺวินฺนํ ขนฺธานํ ฯเปฯ. ปฏิสนฺธิ\-กฺขเณ ฯเปฯ. (1)}

\prose{อวิตกฺก\-วิจารมตฺโต ธมฺโม สวิตกฺก\-สวิจารสฺส ธมฺมสฺส สมฺปยุตฺต\-ปจฺจเยน ปจฺจโย. วิตกฺโก สวิตกฺก\-สวิจารานํ ขนฺธานํ สมฺปยุตฺต\-ปจฺจเยน ปจฺจโย. ปฏิสนฺธิ\-กฺขเณ ฯเปฯ. (2)}

\prose{อวิตกฺก\-วิจารมตฺโต ธมฺโม อวิตกฺก\-อวิจารสฺส ธมฺมสฺส สมฺปยุตฺต\-ปจฺจเยน ปจฺจโย. อวิตกฺก\-วิจารมตฺตา ขนฺธา วิจารสฺส สมฺปยุตฺต\-ปจฺจเยน ปจฺจโย. ปฏิสนฺธิ\-กฺขเณ ฯเปฯ. (3)}

\prose{อวิตกฺก\-วิจารมตฺโต ธมฺโม อวิตกฺก\-วิจารมตฺตสฺส จ อวิตกฺก\-อวิจารสฺส จ ธมฺมสฺส สมฺปยุตฺต\-ปจฺจเยน ปจฺจโย. อวิตกฺก\-วิจารมตฺโต เอโก ขนฺโธ ติณฺณนฺนํ ขนฺธานํ วิจารสฺส จ สมฺปยุตฺต\-ปจฺจเยน ปจฺจโย ฯเปฯ เทฺว ขนฺธา ทฺวินฺนํ ขนฺธานํ วิจารสฺส จ ฯเปฯ. ปฏิสนฺธิ\-กฺขเณ ฯเปฯ. (4)}

\proseitem{126}{อวิตกฺก\-อวิจาโร ธมฺโม อวิตกฺก\-อวิจารสฺส ธมฺมสฺส สมฺปยุตฺต\-ปจฺจเยน ปจฺจโย. อวิตกฺก\-อวิจาโร เอโก ขนฺโธ ติณฺณนฺนํ ขนฺธานํ สมฺปยุตฺต\-ปจฺจเยน ปจฺจโย ฯเปฯ เทฺว ขนฺธา ทฺวินฺนํ ขนฺธานํ ฯเปฯ. ปฏิสนฺธิ\-กฺขเณ ฯเปฯ. (1)}

\prose{\csromanfolio{93}อวิตกฺก\-อวิจาโร ธมฺโม อวิตกฺก\-วิจารมตฺตสฺส ธมฺมสฺส สมฺปยุตฺต\-ปจฺจเยน ปจฺจโย. วิจาโร อวิตกฺก\-วิจารมตฺตานํ ขนฺธานํ สมฺปยุตฺต\-ปจฺจเยน ปจฺจโย. ปฏิสนฺธิ\-กฺขเณ ฯเปฯ. (2)}

\prose{อวิตกฺก\-วิจารมตฺโต จ อวิตกฺก\-อวิจาโร จ ธมฺมา อวิตกฺก\-วิจารมตฺตสฺส ธมฺมสฺส สมฺปยุตฺต\-ปจฺจเยน ปจฺจโย. อวิตกฺก\-วิจารมตฺโต เอโก ขนฺโธ จ วิจาโร จ ติณฺณนฺนํ ขนฺธานํ สมฺปยุตฺต\-ปจฺจเยน ปจฺจโย ฯเปฯ เทฺว ขนฺธา จ วิจาโร จ ทฺวินฺนํ ขนฺธานํ ฯเปฯ. ปฏิสนฺธิ\-กฺขเณ ฯเปฯ. (1)}

\prose{สวิตกฺก\-สวิจาโร จ อวิตกฺก\-วิจารมตฺโต จ ธมฺมา สวิตกฺก\-สวิจารสฺส ธมฺมสฺส สมฺปยุตฺต\-ปจฺจเยน ปจฺจโย. สวิตกฺก\-สวิจาโร เอโก ขนฺโธ จ วิตกฺโก จ ติณฺณนฺนํ ขนฺธานํ สมฺปยุตฺต\-ปจฺจเยน ปจฺจโย ฯเปฯ เทฺว ขนฺธา จ วิตกฺโก จ ทฺวินฺนํ ขนฺธานํ ฯเปฯ. ปฏิสนฺธิ\-กฺขเณ ฯเปฯ. (1)}

\csromanheader{2}{\textbf{วิปฺปยุตฺต}}
\proseitem{127}{สวิตกฺก\-สวิจาโร ธมฺโม อวิตกฺก\-อวิจารสฺส ธมฺมสฺส วิปฺปยุตฺต\-ปจฺจเยน ปจฺจโย. สหชาตํ, ปจฺฉาชาตํ.  สหชาตา สวิตกฺก\-สวิจารา ขนฺธา จิตฺต\-สมุฏฺฐานานํ รูปานํ วิปฺปยุตฺต\-ปจฺจเยน ปจฺจโย. ปฏิสนฺธิ\-กฺขเณ สวิตกฺก\-สวิจารา ขนฺธา กฏตฺตารูปานํ วิปฺปยุตฺต\-ปจฺจเยน ปจฺจโย.  ปจฺฉาชาตา สวิตกฺก\-สวิจารา ขนฺธา ปุเรชาตสฺส อิมสฺส กายสฺส วิปฺปยุตฺต\-ปจฺจเยน ปจฺจโย. (1)}

\prose{อวิตกฺก\-วิจารมตฺโต ธมฺโม อวิตกฺก\-อวิจารสฺส ธมฺมสฺส วิปฺปยุตฺต\-ปจฺจเยน ปจฺจโย. สหชาตํ, ปจฺฉาชาตํ.  สหชาตา อวิตกฺก\-วิจารมตฺตา ขนฺธา จิตฺต\-สมุฏฺฐานานํ รูปานํ วิปฺปยุตฺต\-ปจฺจเยน ปจฺจโย. วิตกฺโก จิตฺต\-สมุฏฺฐานานํ รูปานํ วิปฺปยุตฺต\-ปจฺจเยน ปจฺจโย. ปฏิสนฺธิ\-กฺขเณ อวิตกฺก\-วิจารมตฺตา ขนฺธา กฏตฺตารูปานํ วิปฺปยุตฺต\-ปจฺจเยน ปจฺจโย. วิตกฺโก กฏตฺตารูปานํ วิปฺปยุตฺต\-ปจฺจเยน ปจฺจโย.  ปจฺฉาชาตา อวิตกฺก\-วิจารมตฺตา ขนฺธา จ วิตกฺโก จ ปุเรชาตสฺส อิมสฺส กายสฺส วิปฺปยุตฺต\-ปจฺจเยน ปจฺจโย. (1)}

\proseitem{128}{อวิตกฺก\-อวิจาโร ธมฺโม อวิตกฺก\-อวิจารสฺส ธมฺมสฺส วิปฺปยุตฺต\-ปจฺจเยน ปจฺจโย. สหชาตํ, ปุเรชาตํ, ปจฺฉาชาตํ.   \csromanfolio{94}สหชาตา อวิตกฺก\-อวิจารา ขนฺธา จิตฺต\-สมุฏฺฐานานํ รูปานํ วิปฺปยุตฺต\-ปจฺจเยน ปจฺจโย. วิจาโร จิตฺต\-สมุฏฺฐานานํ รูปานํ วิปฺปยุตฺต\-ปจฺจเยน ปจฺจโย. ปฏิสนฺธิ\-กฺขเณ อวิตกฺก\-อวิจารา ขนฺธา กฏตฺตารูปานํ วิปฺปยุตฺต\-ปจฺจเยน ปจฺจโย. วิจาโร กฏตฺตารูปานํ วิปฺปยุตฺต\-ปจฺจเยน ปจฺจโย. ขนฺธา วตฺถุสฺส วิปฺปยุตฺต\-ปจฺจเยน ปจฺจโย. วตฺถุ ขนฺธานํ วิปฺปยุตฺต\-ปจฺจเยน ปจฺจโย. วิจาโร วตฺถุสฺส วิปฺปยุตฺต\-ปจฺจเยน ปจฺจโย. วตฺถุ วิจารสฺส วิปฺปยุตฺต\-ปจฺจเยน ปจฺจโย.  ปุเรชาตํ จกฺขายตนํ จกฺขุ\-วิญฺญาณสฺส วิปฺปยุตฺต\-ปจฺจเยน ปจฺจโย ฯเปฯ. กายายตนํ กายวิญฺญาณสฺส วิปฺปยุตฺต\-ปจฺจเยน ปจฺจโย. วตฺถุ อวิตกฺก\-อวิจารานํ ขนฺธานํ วิจารสฺส จ วิปฺปยุตฺต\-ปจฺจเยน ปจฺจโย.  ปจฺฉาชาตา อวิตกฺก\-อวิจารา ขนฺธา จ วิจาโร จ ปุเรชาตสฺส อิมสฺส กายสฺส วิปฺปยุตฺต\-ปจฺจเยน ปจฺจโย. (1)}

\prose{อวิตกฺก\-อวิจาโร ธมฺโม สวิตกฺก\-สวิจารสฺส ธมฺมสฺส วิปฺปยุตฺต\-ปจฺจเยน ปจฺจโย. สหชาตํ, ปุเรชาตํ.  สหชาตํ ปฏิสนฺธิ\-กฺขเณ วตฺถุ สวิตกฺก\-สวิจารานํ ขนฺธานํ วิปฺปยุตฺต\-ปจฺจเยน ปจฺจโย. ปุเรชาตํ วตฺถุ สวิตกฺก\-สวิจารานํ ขนฺธานํ วิปฺปยุตฺต\-ปจฺจเยน ปจฺจโย. (2)}

\prose{อวิตกฺก\-อวิจาโร ธมฺโม อวิตกฺก\-วิจารมตฺตสฺส ธมฺมสฺส วิปฺปยุตฺต\-ปจฺจเยน ปจฺจโย. สหชาตํ, ปุเรชาตํ.  สหชาตํ ปฏิสนฺธิ\-กฺขเณ วตฺถุ อวิตกฺก\-วิจารมตฺตานํ ขนฺธานํ วิตกฺกสฺส จ วิปฺปยุตฺต\-ปจฺจเยน ปจฺจโย. ปุเรชาตํ วตฺถุ อวิตกฺก\-วิจารมตฺตานํ ขนฺธานํ วิตกฺกสฺส จ วิปฺปยุตฺต\-ปจฺจเยน ปจฺจโย. (3)}

\prose{อวิตกฺก\-อวิจาโร ธมฺโม อวิตกฺก\-วิจารมตฺตสฺส จ อวิตกฺก\-อวิจารสฺส จ ธมฺมสฺส วิปฺปยุตฺต\-ปจฺจเยน ปจฺจโย. สหชาตํ, ปุเรชาตํ.  สหชาตํ ปฏิสนฺธิ\-กฺขเณ วตฺถุ อวิตกฺก\-วิจารมตฺตานํ ขนฺธานํ วิจารสฺส จ วิปฺปยุตฺต\-ปจฺจเยน ปจฺจโย. ปุเรชาตํ วตฺถุ อวิตกฺก\-วิจารมตฺตานํ ขนฺธานํ วิจารสฺส จ วิปฺปยุตฺต\-ปจฺจเยน ปจฺจโย. (4)}

\prose{อวิตกฺก\-อวิจาโร ธมฺโม สวิตกฺก\-สวิจารสฺส จ อวิตกฺก\-วิจารมตฺตสฺส จ ธมฺมสฺส วิปฺปยุตฺต\-ปจฺจเยน ปจฺจโย. สหชาตํ, ปุเรชาตํ.  สหชาตํ ปฏิสนฺธิ\-กฺขเณ วตฺถุ สวิตกฺก\-สวิจารานํ ขนฺธานํ วิตกฺกสฺส จ วิปฺปยุตฺต\-ปจฺจเยน ปจฺจโย. ปุเรชาตํ วตฺถุ สวิตกฺก\-สวิจารานํ ขนฺธานํ วิตกฺกสฺส จ วิปฺปยุตฺต\-ปจฺจเยน ปจฺจโย. (5)}

\proseitem{129}{\csromanfolio{95}อวิตกฺก\-วิจารมตฺโต จ อวิตกฺก\-อวิจาโร จ ธมฺมา อวิตกฺก\-อวิจารสฺส ธมฺมสฺส วิปฺปยุตฺต\-ปจฺจเยน ปจฺจโย. สหชาตํ, ปจฺฉาชาตํ.  สหชาตา อวิตกฺก\-วิจารมตฺตา ขนฺธา จ วิจาโร จ จิตฺต\-สมุฏฺฐานานํ รูปานํ วิปฺปยุตฺต\-ปจฺจเยน ปจฺจโย. ปฏิสนฺธิ\-กฺขเณ อวิตกฺก\-วิจารมตฺตา ขนฺธา จ วิจาโร จ กฏตฺตารูปานํ วิปฺปยุตฺต\-ปจฺจเยน ปจฺจโย.  ปจฺฉาชาตา อวิตกฺก\-วิจารมตฺตา ขนฺธา จ วิจาโร จ ปุเรชาตสฺส อิมสฺส กายสฺส วิปฺปยุตฺต\-ปจฺจเยน ปจฺจโย. (1)}

\prose{สวิตกฺก\-สวิจาโร จ อวิตกฺก\-วิจารมตฺโต จ ธมฺมา อวิตกฺก\-อวิจารสฺส ธมฺมสฺส วิปฺปยุตฺต\-ปจฺจเยน ปจฺจโย. สหชาตํ, ปจฺฉาชาตํ.  สหชาตา สวิตกฺก\-สวิจารา ขนฺธา จ วิตกฺโก จ จิตฺต\-สมุฏฺฐานานํ รูปานํ วิปฺปยุตฺต\-ปจฺจเยน ปจฺจโย. ปฏิสนฺธิ\-กฺขเณ สวิตกฺก\-สวิจารา ขนฺธา จ วิตกฺโก จ กฏตฺตารูปานํ วิปฺปยุตฺต\-ปจฺจเยน ปจฺจโย.  ปจฺฉาชาตา สวิตกฺก\-สวิจารา ขนฺธา จ วิตกฺโก จ ปุเรชาตสฺส อิมสฺส กายสฺส วิปฺปยุตฺต\-ปจฺจเยน ปจฺจโย. (1)}

\csromanheader{2}{\textbf{อตฺถิ}}
\proseitem{130}{สวิตกฺก\-สวิจาโร ธมฺโม สวิตกฺก\-สวิจารสฺส ธมฺมสฺส อตฺถิปจฺจเยน ปจฺจโย. สวิตกฺก\-สวิจาโร เอโก ขนฺโธ ติณฺณนฺนํ ขนฺธานํ อตฺถิปจฺจเยน ปจฺจโย ฯเปฯ เทฺว ขนฺธา ทฺวินฺนํ ขนฺธานํ ฯเปฯ. ปฏิสนฺธิ\-กฺขเณ ฯเปฯ. (1)}

\prose{สวิตกฺก\-สวิจาโร ธมฺโม อวิตกฺก\-วิจารมตฺตสฺส ธมฺมสฺส อตฺถิปจฺจเยน ปจฺจโย. สวิตกฺก\-สวิจารา ขนฺธา วิตกฺกสฺส อตฺถิปจฺจเยน ปจฺจโย. ปฏิสนฺธิ\-กฺขเณ ฯเปฯ. (2)}

\prose{สวิตกฺก\-สวิจาโร ธมฺโม อวิตกฺก\-อวิจารสฺส ธมฺมสฺส อตฺถิปจฺจเยน ปจฺจโย. สหชาตํ, ปจฺฉาชาตํ.  สหชาตา สวิตกฺก\-สวิจารา ขนฺธา จิตฺต\-สมุฏฺฐานานํ รูปานํ อตฺถิปจฺจเยน ปจฺจโย. ปฏิสนฺธิ\-กฺขเณ สวิตกฺก\-สวิจารา ขนฺธา กฏตฺตารูปานํ อตฺถิปจฺจเยน ปจฺจโย.  ปจฺฉาชาตา สวิตกฺก\-สวิจารา ขนฺธา ปุเรชาตสฺส อิมสฺส กายสฺส อตฺถิปจฺจเยน ปจฺจโย. (3) (สวิตกฺกสวิจาร\-มูลเก อวเสสา ปญฺหา สหชาตปจฺจย\-สทิสา.)}

\proseitem{131}{\csromanfolio{96}อวิตกฺก\-วิจารมตฺโต ธมฺโม อวิตกฺก\-วิจารมตฺตสฺส ธมฺมสฺส อตฺถิปจฺจเยน ปจฺจโย. อวิตกฺก\-วิจารมตฺโต เอโก ขนฺโธ ติณฺณนฺนํ ขนฺธานํ อตฺถิปจฺจเยน ปจฺจโย ฯเปฯ เทฺว ขนฺธา ทฺวินฺนํ ขนฺธานํ ฯเปฯ. ปฏิสนฺธิ\-กฺขเณ ฯเปฯ. (1)}

\prose{อวิตกฺก\-วิจารมตฺโต ธมฺโม สวิตกฺก\-สวิจารสฺส ธมฺมสฺส อตฺถิปจฺจเยน ปจฺจโย. วิตกฺโก สวิตกฺก\-สวิจารานํ ขนฺธานํ อตฺถิปจฺจเยน ปจฺจโย. ปฏิสนฺธิ\-กฺขเณ ฯเปฯ. (2)}

\prose{อวิตกฺก\-วิจารมตฺโต ธมฺโม อวิตกฺก\-อวิจารสฺส ธมฺมสฺส อตฺถิปจฺจเยน ปจฺจโย. สหชาตํ, ปจฺฉาชาตํ.  สหชาตา อวิตกฺก\-วิจารมตฺตา ขนฺธา วิจารสฺส จิตฺตสมุฏฺฐานานญฺจ รูปานํ อตฺถิปจฺจเยน ปจฺจโย. วิตกฺโก จิตฺต\-สมุฏฺฐานานํ รูปานํ อตฺถิปจฺจเยน ปจฺจโย. ปฏิสนฺธิ\-กฺขเณ อวิตกฺก\-วิจารมตฺตา ขนฺธา วิจารสฺส กฏตฺตา จ รูปานํ อตฺถิปจฺจเยน ปจฺจโย. วิตกฺโก กฏตฺตารูปานํ อตฺถิปจฺจเยน ปจฺจโย.  ปจฺฉาชาตา อวิตกฺก\-วิจารมตฺตา ขนฺธา จ วิตกฺโก จ ปุเรชาตสฺส อิมสฺส กายสฺส อตฺถิปจฺจเยน ปจฺจโย. (3)}

\prose{(อวิตกฺกวิจารมตฺต\-มูลกา ปญฺจ ปญฺหา, อวเสสา สหชาตปจฺจย\-สทิสา.)}

\proseitem{132}{อวิตกฺก\-อวิจาโร ธมฺโม อวิตกฺก\-อวิจารสฺส ธมฺมสฺส อตฺถิปจฺจเยน ปจฺจโย. สหชาตํ, ปุเรชาตํ, ปจฺฉาชาตํ, อาหารํ, อินฺทฺริยํ.  สหชาโต อวิตกฺก\-อวิจาโร เอโก ขนฺโธ ติณฺณนฺนํ ขนฺธานํ จิตฺตสมุฏฺฐานานญฺจ รูปานํ อตฺถิปจฺจเยน ปจฺจโย ฯเปฯ เทฺว ขนฺธา ทฺวินฺนํ ขนฺธานํ ฯเปฯ. วิจาโร จิตฺต\-สมุฏฺฐานานํ รูปานํ อตฺถิปจฺจเยน ปจฺจโย. ปฏิสนฺธิ\-กฺขเณ อวิตกฺก\-อวิจาโร เอโก ขนฺโธ ติณฺณนฺนํ ขนฺธานํ กฏตฺตา จ รูปานํ อตฺถิปจฺจเยน ปจฺจโย ฯเปฯ. ขนฺธา วตฺถุสฺส อตฺถิปจฺจเยน ปจฺจโย. วตฺถุ ขนฺธานํ อตฺถิปจฺจเยน ปจฺจโย. วิจาโร วตฺถุสฺส อตฺถิปจฺจเยน ปจฺจโย. วตฺถุ วิจารสฺส อตฺถิปจฺจเยน ปจฺจโย. เอกํ มหาภูตํ ติณฺณนฺนํ มหาภูตานํ อตฺถิปจฺจเยน ปจฺจโย ฯเปฯ มหาภูตา จิตฺต\-สมุฏฺฐานานํ รูปานํ กฏตฺตารูปานํ อุปาทารูปานํ อตฺถิปจฺจเยน ปจฺจโย. พาหิรํ. อาหารสมุฏฺฐานํ. อุตุสมุฏฺฐานํ. อสญฺญสตฺตานํ เอกํ มหาภูตํ ติณฺณนฺนํ มหาภูตานํ ฯเปฯ มหาภูตา กฏตฺตารูปานํ อุปาทารูปานํ อตฺถิปจฺจเยน ปจฺจโย.  ปุเรชาตํ ทิพฺเพน จกฺขุนา รูปํ ปสฺสติ. ทิพฺพาย \csromanfolio{97}โสตธาตุยา สทฺทํ สุณาติ. รูปายตนํ จกฺขุ\-วิญฺญาณสฺส ฯเปฯ. โผฏฺฐพฺพา\-ยตนํ กายวิญฺญาณสฺส อตฺถิปจฺจเยน ปจฺจโย. จกฺขายตนํ จกฺขุ\-วิญฺญาณสฺส ฯเปฯ. กายายตนํ กายวิญฺญาณสฺส อตฺถิปจฺจเยน ปจฺจโย. วตฺถุ อวิตกฺก\-อวิจารานํ ขนฺธานํ วิจารสฺส จ อตฺถิปจฺจเยน ปจฺจโย.  ปจฺฉาชาตา อวิตกฺก\-อวิจารา ขนฺธา จ วิจาโร จ ปุเรชาตสฺส อิมสฺส กายสฺส อตฺถิปจฺจเยน ปจฺจโย. กพฬีกาโร อาหาโร อิมสฺส กายสฺส อตฺถิปจฺจเยน ปจฺจโย. รูปชีวิตินฺทฺริยํ กฏตฺตารูปานํ อตฺถิปจฺจเยน ปจฺจโย. (1)}

\prose{อวิตกฺก\-อวิจาโร ธมฺโม สวิตกฺก\-สวิจารสฺส ธมฺมสฺส อตฺถิปจฺจเยน ปจฺจโย. สหชาตํ, ปุเรชาตํ.  สหชาตํ ปฏิสนฺธิ\-กฺขเณ วตฺถุ สวิตกฺก\-สวิจารานํ ขนฺธานํ อตฺถิปจฺจเยน ปจฺจโย.  ปุเรชาตํ จกฺขุํ อนิจฺจโต ทุกฺขโต อนตฺตโต วิปสฺสติ, อสฺสาเทติ อภินนฺทติ, ตํ อารพฺภ ราโค อุปฺปชฺชติ ฯเปฯ โทมนสฺสํ อุปฺปชฺชติ. โสตํ. ฆานํ. ชิวฺหํ. กายํ. รูเป. สทฺเท. คนฺเธ. รเส. โผฏฺฐพฺเพ. วตฺถุํ อนิจฺจโต ทุกฺขโต อนตฺตโต วิปสฺสติ ฯเปฯ โทมนสฺสํ อุปฺปชฺชติ. วตฺถุ สวิตกฺก\-สวิจารานํ ขนฺธานํ อตฺถิปจฺจเยน ปจฺจโย. (2)}

\prose{อวิตกฺก\-อวิจาโร ธมฺโม อวิตกฺก\-วิจารมตฺตสฺส ธมฺมสฺส อตฺถิปจฺจเยน ปจฺจโย. สหชาตํ, ปุเรชาตํ.  สหชาโต วิจาโร อวิตกฺก\-วิจารมตฺตานํ ขนฺธานํ อตฺถิปจฺจเยน ปจฺจโย. ปฏิสนฺธิ\-กฺขเณ วิจาโร อวิตกฺก\-วิจารมตฺตานํ ขนฺธานํ อตฺถิปจฺจเยน ปจฺจโย. ปฏิสนฺธิ\-กฺขเณ วตฺถุ อวิตกฺก\-วิจารมตฺตานํ ขนฺธานํ วิตกฺกสฺส จ อตฺถิปจฺจเยน ปจฺจโย.  ปุเรชาตํ จกฺขุํ อนิจฺจโต ทุกฺขโต อนตฺตโต วิปสฺสติ, อสฺสาเทติ อภินนฺทติ, ตํ อารพฺภ วิตกฺโก อุปฺปชฺชติ, โสตํ. ฆานํ. ชิวฺหํ. กายํ. รูเป. สทฺเท. คนฺเธ. รเส. โผฏฺฐพฺเพ. วตฺถุํ อนิจฺจโต ทุกฺขโต อนตฺตโต วิปสฺสติ, อสฺสาเทติ อภินนฺทติ, ตํ อารพฺภ วิตกฺโก อุปฺปชฺชติ. วตฺถุ อวิตกฺก\-วิจารมตฺตานํ ขนฺธานํ วิตกฺกสฺส จ อตฺถิปจฺจเยน ปจฺจโย. (3)}

\prose{อวิตกฺก\-อวิจาโร ธมฺโม อวิตกฺก\-วิจารมตฺตสฺส จ อวิตกฺก\-อวิจารสฺส จ ธมฺมสฺส อตฺถิปจฺจเยน ปจฺจโย. สหชาตํ, ปุเรชาตํ.  สหชาโต วิจาโร อวิตกฺก\-วิจารมตฺตานํ ขนฺธานํ จิตฺตสมุฏฺฐานานญฺจ รูปานํ อตฺถิปจฺจเยน ปจฺจโย. ปฏิสนฺธิ\-กฺขเณ วิจาโร อวิตกฺก\-วิจารมตฺตานํ \csromanfolio{98}ขนฺธานํ กฏตฺตา จ รูปานํ อตฺถิปจฺจเยน ปจฺจโย. ปฏิสนฺธิ\-กฺขเณ วตฺถุ อวิตกฺก\-วิจารมตฺตานํ ขนฺธานํ วิจารสฺส จ อตฺถิปจฺจเยน ปจฺจโย.  ปุเรชาตํ วตฺถุ อวิตกฺก\-วิจารมตฺตานํ ขนฺธานํ วิจารสฺส จ อตฺถิปจฺจเยน ปจฺจโย. (4)}

\prose{อวิตกฺก\-อวิจาโร ธมฺโม สวิตกฺก\-สวิจารสฺส จ อวิตกฺก\-วิจารมตฺตสฺส จ ธมฺมสฺส อตฺถิปจฺจเยน ปจฺจโย. สหชาตํ, ปุเรชาตํ.  สหชาตํ ปฏิสนฺธิ\-กฺขเณ วตฺถุ สวิตกฺก\-สวิจารานํ ขนฺธานํ วิตกฺกสฺส จ อตฺถิปจฺจเยน ปจฺจโย.  ปุเรชาตํ จกฺขุํ อนิจฺจโต ทุกฺขโต อนตฺตโต วิปสฺสติ, อสฺสาเทติ อภินนฺทติ, ตํ อารพฺภ สวิตกฺก\-สวิจารา ขนฺธา จ วิตกฺโก จ อุปฺปชฺชนฺติ. โสตํ. ฆานํ. ชิวฺหํ. กายํ ฯเปฯ. วตฺถุํ อนิจฺจโต ทุกฺขโต อนตฺตโต วิปสฺสติ, อสฺสาเทติ อภินนฺทติ, ตํ อารพฺภ สวิตกฺก\-สวิจารา ขนฺธา จ วิตกฺโก จ อุปฺปชฺชนฺติ. วตฺถุ สวิตกฺก\-สวิจารานํ ขนฺธานํ วิตกฺกสฺส จ อตฺถิปจฺจเยน ปจฺจโย. (5)}

\proseitem{133}{สวิตกฺก\-สวิจาโร จ อวิตกฺก\-อวิจาโร จ ธมฺมา สวิตกฺก\-สวิจารสฺส ธมฺมสฺส อตฺถิปจฺจเยน ปจฺจโย. สหชาตํ, ปุเรชาตํ.  สหชาโต สวิตกฺก\-สวิจาโร เอโก ขนฺโธ จ วตฺถุ จ ติณฺณนฺนํ ขนฺธานํ อตฺถิปจฺจเยน ปจฺจโย ฯเปฯ เทฺว ขนฺธา จ วตฺถุ จ ทฺวินฺนํ ขนฺธานํ อตฺถิปจฺจเยน ปจฺจโย. ปฏิสนฺธิ\-กฺขเณ สวิตกฺก\-สวิจาโร เอโก ขนฺโธ จ วตฺถุ จ ติณฺณนฺนํ ขนฺธานํ อตฺถิปจฺจเยน ปจฺจโย ฯเปฯ เทฺว ขนฺธา จ วตฺถุ จ ทฺวินฺนํ ขนฺธานํ อตฺถิปจฺจเยน ปจฺจโย. (1)}

\prose{สวิตกฺก\-สวิจาโร จ อวิตกฺก\-อวิจาโร จ ธมฺมา อวิตกฺก\-วิจารมตฺตสฺส ธมฺมสฺส อตฺถิปจฺจเยน ปจฺจโย. สหชาตํ, ปุเรชาตํ. สหชาตา สวิตกฺก\-สวิจารา ขนฺธา จ วตฺถุ จ วิตกฺกสฺส อตฺถิปจฺจเยน ปจฺจโย. ปฏิสนฺธิ\-กฺขเณ สวิตกฺก\-สวิจารา ขนฺธา จ วตฺถุ จ วิตกฺกสฺส อตฺถิปจฺจเยน ปจฺจโย. (2)}

\prose{สวิตกฺก\-สวิจาโร จ อวิตกฺก\-อวิจาโร จ ธมฺมา อวิตกฺก\-อวิจารสฺส ธมฺมสฺส อตฺถิปจฺจเยน ปจฺจโย. สหชาตํ, ปจฺฉาชาตํ, อาหารํ, อินฺทฺริยํ.  สหชาตา สวิตกฺก\-สวิจารา ขนฺธา จ มหาภูตา จ จิตฺต\-สมุฏฺฐานานํ รูปานํ อตฺถิปจฺจเยน ปจฺจโย. ปฏิสนฺธิ\-กฺขเณ สวิตกฺก\-สวิจารา ขนฺธา จ มหาภูตา จ กฏตฺตารูปานํ อตฺถิปจฺจเยน ปจฺจโย.   \csromanfolio{99}ปจฺฉาชาตา สวิตกฺก\-สวิจารา ขนฺธา จ กพฬีกาโร อาหาโร จ ปุเรชาตสฺส อิมสฺส กายสฺส อตฺถิปจฺจเยน ปจฺจโย.  ปจฺฉาชาตา สวิตกฺก\-สวิจารา ขนฺธา จ รูปชีวิตินฺทฺริยญฺจ กฏตฺตารูปานํ อตฺถิปจฺจเยน ปจฺจโย. (3)}

\prose{สวิตกฺก\-สวิจาโร จ อวิตกฺก\-อวิจาโร จ ธมฺมา สวิตกฺก\-สวิจารสฺส จ อวิตกฺก\-วิจารมตฺตสฺส จ ธมฺมสฺส อตฺถิปจฺจเยน ปจฺจโย. สหชาตํ, ปุเรชาตํ.  สหชาโต สวิตกฺก\-สวิจาโร เอโก ขนฺโธ จ วตฺถุ จ ติณฺณนฺนํ ขนฺธานํ วิตกฺกสฺส จ อตฺถิปจฺจเยน ปจฺจโย ฯเปฯ เทฺว ขนฺธา จ วตฺถุ จ ทฺวินฺนํ ขนฺธานํ วิตกฺกสฺส จ อตฺถิปจฺจเยน ปจฺจโย. ปฏิสนฺธิ\-กฺขเณ ฯเปฯ. (4)}

\proseitem{134}{อวิตกฺก\-วิจารมตฺโต จ อวิตกฺก\-อวิจาโร จ ธมฺมา สวิตกฺก\-สวิจารสฺส ธมฺมสฺส อตฺถิปจฺจเยน ปจฺจโย. สหชาตํ, ปุเรชาตํ.  สหชาโต วิตกฺโก จ วตฺถุ จ สวิตกฺก\-สวิจารานํ ขนฺธานํ อตฺถิปจฺจเยน ปจฺจโย ปฏิสนฺธิ\-กฺขเณ วิตกฺโก จ วตฺถุ จ สวิตกฺก\-สวิจารานํ ขนฺธานํ อตฺถิปจฺจเยน ปจฺจโย. (1)}

\prose{อวิตกฺก\-วิจารมตฺโต จ อวิตกฺก\-อวิจาโร จ ธมฺมา อวิตกฺก\-วิจารมตฺตสฺส ธมฺมสฺส อตฺถิปจฺจเยน ปจฺจโย. สหชาตํ, ปุเรชาตํ.  สหชาโต อวิตกฺก\-วิจารมตฺโต เอโก ขนฺโธ จ วิจาโร จ ติณฺณนฺนํ ขนฺธานํ อตฺถิปจฺจเยน ปจฺจโย ฯเปฯ เทฺว ขนฺธา จ วิจาโร จ ทฺวินฺนํ ขนฺธานํ. อวิตกฺก\-วิจารมตฺโต เอโก ขนฺโธ จ วตฺถุ จ ติณฺณนฺนํ ขนฺธานํ อตฺถิปจฺจเยน ปจฺจโย ฯเปฯ เทฺว ขนฺธา จ วตฺถุ จ ทฺวินฺนํ ขนฺธานํ อตฺถิปจฺจเยน ปจฺจโย. ปฏิสนฺธิ\-กฺขเณ อวิตกฺก\-วิจารมตฺโต เอโก ขนฺโธ จ วิจาโร จ ติณฺณนฺนํ ขนฺธานํ อตฺถิปจฺจเยน ปจฺจโย ฯเปฯ เทฺว ขนฺธา จ วิจาโร จ ทฺวินฺนํ ขนฺธานํ ฯเปฯ. ปฏิสนฺธิ\-กฺขเณ อวิตกฺก\-วิจารมตฺโต เอโก ขนฺโธ จ วตฺถุ จ ติณฺณนฺนํ ขนฺธานํ ฯเปฯ เทฺว ขนฺธา จ วตฺถุ จ ทฺวินฺนํ ขนฺธานํ อตฺถิปจฺจเยน ปจฺจโย. (2)}

\prose{อวิตกฺก\-วิจารมตฺโต จ อวิตกฺก\-อวิจาโร จ ธมฺมา อวิตกฺก\-อวิจารสฺส ธมฺมสฺส อตฺถิปจฺจเยน ปจฺจโย. สหชาตํ, ปุเรชาตํ, ปจฺฉาชาตํ, อาหารํ, อินฺทฺริยํ.  สหชาตา อวิตกฺก\-วิจารมตฺตา ขนฺธา จ วิจาโร จ จิตฺต\-สมุฏฺฐานานํ รูปานํ อตฺถิปจฺจเยน ปจฺจโย.   \csromanfolio{100}สหชาตา อวิตกฺก\-วิจารมตฺตา ขนฺธา จ มหาภูตา จ จิตฺต\-สมุฏฺฐานานํ รูปานํ อตฺถิปจฺจเยน ปจฺจโย. สหชาโต วิตกฺโก จ มหาภูตา จ จิตฺต\-สมุฏฺฐานานํ รูปานํ อตฺถิปจฺจเยน ปจฺจโย.  สหชาตา อวิตกฺก\-วิจารมตฺตา ขนฺธา จ วตฺถุ จ วิจารสฺส อตฺถิปจฺจเยน ปจฺจโย. ปฏิสนฺธิ\-กฺขเณ อวิตกฺก\-วิจารมตฺตา ขนฺธา จ วิจาโร จ กฏตฺตารูปานํ อตฺถิปจฺจเยน ปจฺจโย. ปฏิสนฺธิ\-กฺขเณ อวิตกฺก\-วิจารมตฺตา ขนฺธา จ มหาภูตา จ กฏตฺตารูปานํ อตฺถิปจฺจเยน ปจฺจโย. ปฏิสนฺธิ\-กฺขเณ วิตกฺโก จ มหาภูตา จ กฏตฺตารูปานํ อตฺถิปจฺจเยน ปจฺจโย. ปฏิสนฺธิ\-กฺขเณ อวิตกฺก\-วิจารมตฺตา ขนฺธา จ วตฺถุ จ วิจารสฺส อตฺถิปจฺจเยน ปจฺจโย.  ปจฺฉาชาตา อวิตกฺก\-วิจารมตฺตา ขนฺธา จ วิจาโร จ ปุเรชาตสฺส อิมสฺส กายสฺส อตฺถิปจฺจเยน ปจฺจโย.  ปจฺฉาชาตา อวิตกฺก\-วิจารมตฺตา ขนฺธา จ วิตกฺโก จ กพฬีกาโร อาหาโร จ ปุเรชาตสฺส อิมสฺส กายสฺส อตฺถิปจฺจเยน ปจฺจโย.  ปจฺฉาชาตา อวิตกฺก\-วิจารมตฺตา ขนฺธา จ วิตกฺโก จ รูปชีวิตินฺทฺริยญฺจ กฏตฺตารูปานํ อตฺถิปจฺจเยน ปจฺจโย. (3)}

\prose{อวิตกฺก\-วิจารมตฺโต จ อวิตกฺก\-อวิจาโร จ ธมฺมา อวิตกฺก\-วิจารมตฺตสฺส จ อวิตกฺก\-อวิจารสฺส จ ธมฺมสฺส อตฺถิปจฺจเยน ปจฺจโย. สหชาตํ, ปุเรชาตํ.  สหชาโต อวิตกฺก\-วิจารมตฺโต เอโก ขนฺโธ จ วิจาโร จ ติณฺณนฺนํ ขนฺธานํ จิตฺตสมุฏฺฐานานญฺจ รูปานํ อตฺถิปจฺจเยน ปจฺจโย ฯเปฯ เทฺว ขนฺธา จ วิจาโร จ ฯเปฯ. อวิตกฺก\-วิจารมตฺโต เอโก ขนฺโธ จ วตฺถุ จ ติณฺณนฺนํ ขนฺธานํ วิจารสฺส จ อตฺถิปจฺจเยน ปจฺจโย ฯเปฯ เทฺว ขนฺธา จ วตฺถุ จ ฯเปฯ. ปฏิสนฺธิ\-กฺขเณ ฯเปฯ. (4)}

\proseitem{135}{สวิตกฺก\-สวิจาโร จ อวิตกฺก\-วิจารมตฺโต จ ธมฺมา สวิตกฺก\-สวิจารสฺส ธมฺมสฺส อตฺถิปจฺจเยน ปจฺจโย. สวิตกฺก\-สวิจาโร เอโก ขนฺโธ จ วิตกฺโก จ ติณฺณนฺนํ ขนฺธานํ อตฺถิปจฺจเยน ปจฺจโย ฯเปฯ เทฺว ขนฺธา จ วิตกฺโก จ ฯเปฯ. ปฏิสนฺธิ\-กฺขเณ ฯเปฯ. (1)}

\prose{สวิตกฺก\-สวิจาโร จ อวิตกฺก\-วิจารมตฺโต จ ธมฺมา อวิตกฺก\-อวิจารสฺส ธมฺมสฺส อตฺถิปจฺจเยน ปจฺจโย. สหชาตํ, ปจฺฉาชาตํ.  สหชาตา สวิตกฺก\-สวิจารา ขนฺธา จ วิตกฺโก จ จิตฺต\-สมุฏฺฐานานํ รูปานํ อตฺถิปจฺจเยน ปจฺจโย. ปฏิสนฺธิ\-กฺขเณ สวิตกฺก\-สวิจารา ขนฺธา จ \csromanfolio{101}วิตกฺโก จ กฏตฺตารูปานํ อตฺถิปจฺจเยน ปจฺจโย. ปจฺฉาชาตา สวิตกฺก\-สวิจารา ขนฺธา จ วิตกฺโก จ ปุเรชาตสฺส อิมสฺส กายสฺส อตฺถิปจฺจเยน ปจฺจโย. (2)}

\prose{สวิตกฺก\-สวิจาโร จ อวิตกฺก\-วิจารมตฺโต จ ธมฺมา สวิตกฺก\-สวิจารสฺส จ อวิตกฺก\-อวิจารสฺส จ ธมฺมสฺส อตฺถิปจฺจเยน ปจฺจโย. สวิตกฺก\-สวิจาโร เอโก ขนฺโธ จ วิตกฺโก จ ติณฺณนฺนํ ขนฺธานํ จิตฺตสมุฏฺฐานานญฺจ รูปานํ อตฺถิปจฺจเยน ปจฺจโย ฯเปฯ เทฺว ขนฺธา จ วิตกฺโก จ ทฺวินฺนํ ขนฺธานํ ฯเปฯ. ปฏิสนฺธิ\-กฺขเณ ฯเปฯ. (3)}

\proseitem{136}{สวิตกฺก\-สวิจาโร จ อวิตกฺก\-วิจารมตฺโต จ อวิตกฺก\-อวิจาโร จ ธมฺมา สวิตกฺก\-สวิจารสฺส ธมฺมสฺส อตฺถิปจฺจเยน ปจฺจโย. สหชาตํ, ปุเรชาตํ.  สหชาโต สวิตกฺก\-สวิจาโร เอโก ขนฺโธ จ วิตกฺโก จ วตฺถุ จ ติณฺณนฺนํ ขนฺธานํ อตฺถิปจฺจเยน ปจฺจโย ฯเปฯ. ปฏิสนฺธิ\-กฺขเณ ฯเปฯ. (1)}

\prose{สวิตกฺก\-สวิจาโร จ อวิตกฺก\-วิจารมตฺโต จ อวิตกฺก\-อวิจาโร จ ธมฺมา อวิตกฺก\-อวิจารสฺส ธมฺมสฺส อตฺถิปจฺจเยน ปจฺจโย. สหชาตํ, ปจฺฉาชาตํ, อาหารํ, อินฺทฺริยํ.  สหชาตา สวิตกฺก\-สวิจารา ขนฺธา จ วิตกฺโก จ มหาภูตา จ จิตฺต\-สมุฏฺฐานานํ รูปานํ อตฺถิปจฺจเยน ปจฺจโย. ปฏิสนฺธิ\-กฺขเณ ฯเปฯ. ปจฺฉาชาตา สวิตกฺก\-สวิจารา ขนฺธา จ วิตกฺโก จ กพฬีกาโร อาหาโร จ ปุเรชาตสฺส อิมสฺส กายสฺส อตฺถิปจฺจเยน ปจฺจโย.  ปจฺฉาชาตา สวิตกฺก\-สวิจารา ขนฺธา จ วิตกฺโก จ รูปชีวิตินฺทฺริยญฺจ กฏตฺตารูปานํ อตฺถิปจฺจเยน ปจฺจโย. (2)}

\csromanheader{2}{\textbf{นตฺถิ วิคตาวิคต}}
\proseitem{137}{สวิตกฺก\-สวิจาโร ธมฺโม สวิตกฺก\-สวิจารสฺส ธมฺมสฺส นตฺถิ\-ปจฺจเยน ปจฺจโย. วิคต\-ปจฺจเยน ปจฺจโย.}

\prose{(นตฺถิ\-ปจฺจยญฺจ วิคต\-ปจฺจยญฺจ อนนฺตร\-สทิสํ, อวิคตํ อตฺถิสทิสํ.)}

\csromanmarkh{1. ปจฺจยานุโลม}\csromanmarkhii{2. สงฺขฺยาวาร}\csromanheadernomark{2}{1. \textbf{ปจฺจยานุโลม 2. สงฺขฺยาวาร}}
\proseitem{138}{\csromanfolio{102}เหตุยา เอกาทส, อารมฺมเณ เอกวีส, อธิปติยา เตวีส, อนนฺตเร ปญฺจวีส, สมนนฺตเร ปญฺจวีส, สหชาเต ติํส, อญฺญมญฺเญ อฏฺฐวีส, นิสฺสเย ติํส, อุปนิสฺสเย ปญฺจวีส, ปุเรชาเต ปญฺจ, ปจฺฉาชาเต ปญฺจ, อาเสวเน เอกวีส, กมฺเม เอกาทส, วิปาเก เอกวีส, อาหาเร เอกาทส, อินฺทฺริเย เอกาทส, ฌาเน เอกวีส, มคฺเค โสฬส, สมฺปยุตฺเต เอกาทส, วิปฺปยุตฺเต นว, อตฺถิยา ติํส, นตฺถิยา ปญฺจวีส, วิคเต ปญฺจวีส, อวิคเต ติํส.}

\prosewithcloserruled{(ฆฏนา กุสล\-ตฺติก\-สทิสาเยว, ปญฺหา\-วาร\-คณนํ เอวํ อสมฺโมหนฺเตน คเณตพฺพํ.)}{อนุโลมํ.}
\csromanheader{2}{2. ปจฺจนียุ\-ทฺธาร}
\proseitem{139}{สวิตกฺก\-สวิจาโร ธมฺโม สวิตกฺก\-สวิจารสฺส ธมฺมสฺส อารมฺมณ\-ปจฺจเยน ปจฺจโย. สหชาต\-ปจฺจเยน ปจฺจโย. อุปนิสฺสย\-ปจฺจเยน ปจฺจโย. กมฺมปจฺจเยน ปจฺจโย. (1)}

\prose{สวิตกฺก\-สวิจาโร ธมฺโม อวิตกฺก\-วิจารมตฺตสฺส ธมฺมสฺส อารมฺมณ\-ปจฺจเยน ปจฺจโย. สหชาต\-ปจฺจเยน ปจฺจโย. อุปนิสฺสย\-ปจฺจเยน ปจฺจโย. กมฺมปจฺจเยน ปจฺจโย. (2)}

\prose{สวิตกฺก\-สวิจาโร ธมฺโม อวิตกฺก\-อวิจารสฺส ธมฺมสฺส อารมฺมณ\-ปจฺจเยน ปจฺจโย. สหชาต\-ปจฺจเยน ปจฺจโย. อุปนิสฺสย\-ปจฺจเยน ปจฺจโย. ปจฺฉาชาต\-ปจฺจเยน ปจฺจโย. กมฺมปจฺจเยน ปจฺจโย. (3)}

\prose{สวิตกฺก\-สวิจาโร ธมฺโม สวิตกฺก\-สวิจารสฺส จ อวิตกฺก\-อวิจารสฺส จ ธมฺมสฺส สหชาต\-ปจฺจเยน ปจฺจโย. กมฺมปจฺจเยน ปจฺจโย. (4)}

\prose{\csromanfolio{103}สวิตกฺก\-สวิจาโร ธมฺโม อวิตกฺก\-วิจารมตฺตสฺส จ อวิตกฺก\-อวิจารสฺส จ ธมฺมสฺส สหชาต\-ปจฺจเยน ปจฺจโย. อุปนิสฺสย\-ปจฺจเยน ปจฺจโย. กมฺมปจฺจเยน ปจฺจโย. (5)}

\prose{สวิตกฺก\-สวิจาโร ธมฺโม สวิตกฺก\-สวิจารสฺส จ อวิตกฺก\-วิจารมตฺตสฺส จ ธมฺมสฺส อารมฺมณ\-ปจฺจเยน ปจฺจโย. สหชาต\-ปจฺจเยน ปจฺจโย. อุปนิสฺสย\-ปจฺจเยน ปจฺจโย. กมฺมปจฺจเยน ปจฺจโย. (6)}

\prose{สวิตกฺก\-สวิจาโร ธมฺโม สวิตกฺก\-สวิจารสฺส จ อวิตกฺก\-วิจารมตฺตสฺส จ อวิตกฺก\-อวิจารสฺส จ ธมฺมสฺส สหชาต\-ปจฺจเยน ปจฺจโย. กมฺมปจฺจเยน ปจฺจโย. (7)}

\proseitem{140}{อวิตกฺก\-วิจารมตฺโต ธมฺโม อวิตกฺก\-วิจารมตฺตสฺส ธมฺมสฺส อารมฺมณ\-ปจฺจเยน ปจฺจโย. สหชาต\-ปจฺจเยน ปจฺจโย. อุปนิสฺสย\-ปจฺจเยน ปจฺจโย. (1)}

\prose{อวิตกฺก\-วิจารมตฺโต ธมฺโม สวิตกฺก\-สวิจารสฺส ธมฺมสฺส อารมฺมณ\-ปจฺจเยน ปจฺจโย. สหชาต\-ปจฺจเยน ปจฺจโย. อุปนิสฺสย\-ปจฺจเยน ปจฺจโย. (2)}

\prose{อวิตกฺก\-วิจารมตฺโต ธมฺโม อวิตกฺก\-อวิจารสฺส ธมฺมสฺส อารมฺมณ\-ปจฺจเยน ปจฺจโย. สหชาต\-ปจฺจเยน ปจฺจโย. อุปนิสฺสย\-ปจฺจเยน ปจฺจโย. ปจฺฉาชาต\-ปจฺจเยน ปจฺจโย. กมฺมปจฺจเยน ปจฺจโย. (3)}

\prose{อวิตกฺก\-วิจารมตฺโต ธมฺโม สวิตกฺก\-สวิจารสฺส จ อวิตกฺก\-อวิจารสฺส จ ธมฺมสฺส สหชาต\-ปจฺจเยน ปจฺจโย. (4)}

\prose{อวิตกฺก\-วิจารมตฺโต ธมฺโม อวิตกฺก\-วิจารมตฺตสฺส จ อวิตกฺก\-อวิจารสฺส จ ธมฺมสฺส สหชาต\-ปจฺจเยน ปจฺจโย. อุปนิสฺสย\-ปจฺจเยน ปจฺจโย. กมฺมปจฺจเยน ปจฺจโย. (5)}

\prose{อวิตกฺก\-วิจารมตฺโต ธมฺโม สวิตกฺก\-สวิจารสฺส จ อวิตกฺก\-วิจารมตฺตสฺส จ ธมฺมสฺส อารมฺมณ\-ปจฺจเยน ปจฺจโย. อุปนิสฺสย\-ปจฺจเยน ปจฺจโย. (6)}

\proseitem{141}{\csromanfolio{104}อวิตกฺก\-อวิจาโร ธมฺโม อวิตกฺก\-อวิจารสฺส ธมฺมสฺส อารมฺมณ\-ปจฺจเยน ปจฺจโย. สหชาต\-ปจฺจเยน ปจฺจโย. อุปนิสฺสย\-ปจฺจเยน ปจฺจโย. ปุเรชาต\-ปจฺจเยน ปจฺจโย. ปจฺฉาชาต\-ปจฺจเยน ปจฺจโย. กมฺมปจฺจเยน ปจฺจโย. อาหารปจฺจเยน ปจฺจโย. อินฺทฺริย\-ปจฺจเยน ปจฺจโย. (1)}

\prose{อวิตกฺก\-อวิจาโร ธมฺโม สวิตกฺก\-สวิจารสฺส ธมฺมสฺส อารมฺมณ\-ปจฺจเยน ปจฺจโย. สหชาต\-ปจฺจเยน ปจฺจโย. อุปนิสฺสย\-ปจฺจเยน ปจฺจโย. ปุเรชาต\-ปจฺจเยน ปจฺจโย. (2)}

\prose{อวิตกฺก\-อวิจาโร ธมฺโม อวิตกฺก\-วิจารมตฺตสฺส ธมฺมสฺส อารมฺมณ\-ปจฺจเยน ปจฺจโย. สหชาต\-ปจฺจเยน ปจฺจโย. อุปนิสฺสย\-ปจฺจเยน ปจฺจโย. ปุเรชาต\-ปจฺจเยน ปจฺจโย. (3)}

\prose{อวิตกฺก\-อวิจาโร ธมฺโม อวิตกฺก\-วิจารมตฺตสฺส จ อวิตกฺก\-อวิจารสฺส จ ธมฺมสฺส อารมฺมณ\-ปจฺจเยน ปจฺจโย. สหชาต\-ปจฺจเยน ปจฺจโย. อุปนิสฺสย\-ปจฺจเยน ปจฺจโย. ปุเรชาต\-ปจฺจเยน ปจฺจโย. (4)}

\prose{อวิตกฺก\-อวิจาโร ธมฺโม สวิตกฺก\-สวิจารสฺส จ อวิตกฺก\-วิจารมตฺตสฺส จ ธมฺมสฺส อารมฺมณ\-ปจฺจเยน ปจฺจโย. สหชาต\-ปจฺจเยน ปจฺจโย. อุปนิสฺสย\-ปจฺจเยน ปจฺจโย. ปุเรชาต\-ปจฺจเยน ปจฺจโย. (5)}

\proseitem{142}{สวิตกฺก\-สวิจาโร จ อวิตกฺก\-อวิจาโร จ ธมฺมา สวิตกฺก\-สวิจารสฺส ธมฺมสฺส สหชาตํ. ปุเรชาตํ. (1)}

\prose{สวิตกฺก\-สวิจาโร จ อวิตกฺก\-อวิจาโร จ ธมฺมา อวิตกฺก\-วิจารมตฺตสฺส ธมฺมสฺส สหชาตํ. ปุเรชาตํ. (2)}

\prose{สวิตกฺก\-สวิจาโร จ อวิตกฺก\-อวิจาโร จ ธมฺมา อวิตกฺก\-อวิจารสฺส ธมฺมสฺส สหชาตํ. ปจฺฉาชาตํ. อาหารํ. อินฺทฺริยํ. (3)}

\prose{สวิตกฺก\-สวิจาโร จ อวิตกฺก\-อวิจาโร จ ธมฺมา สวิตกฺก\-สวิจารสฺส จ อวิตกฺก\-วิจารมตฺตสฺส จ ธมฺมสฺส สหชาตํ. ปุเรชาตํ. (4)}

\proseitem{143}{อวิตกฺก\-วิจารมตฺโต จ อวิตกฺก\-อวิจาโร จ ธมฺมา สวิตกฺก\-สวิจารสฺส ธมฺมสฺส อารมฺมณ\-ปจฺจเยน ปจฺจโย. สหชาต\-ปจฺจเยน ปจฺจโย. อุปนิสฺสย\-ปจฺจเยน ปจฺจโย. ปุเรชาต\-ปจฺจเยน ปจฺจโย. (1)}

\prose{\csromanfolio{105}อวิตกฺก\-วิจารมตฺโต จ อวิตกฺก\-อวิจาโร จ ธมฺมา อวิตกฺก\-วิจารมตฺตสฺส ธมฺมสฺส อารมฺมณ\-ปจฺจเยน ปจฺจโย. สหชาต\-ปจฺจเยน ปจฺจโย. อุปนิสฺสย\-ปจฺจเยน ปจฺจโย. ปุเรชาต\-ปจฺจเยน ปจฺจโย. (2)}

\prose{อวิตกฺก\-วิจารมตฺโต จ อวิตกฺก\-อวิจาโร จ ธมฺมา อวิตกฺก\-อวิจารสฺส ธมฺมสฺส อารมฺมณ\-ปจฺจเยน ปจฺจโย. สหชาต\-ปจฺจเยน ปจฺจโย. อุปนิสฺสย\-ปจฺจเยน ปจฺจโย. ปุเรชาต\-ปจฺจเยน ปจฺจโย. ปจฺฉาชาต\-ปจฺจเยน ปจฺจโย. อาหารปจฺจเยน ปจฺจโย. อินฺทฺริย\-ปจฺจเยน ปจฺจโย. (3)}

\prose{อวิตกฺก\-วิจารมตฺโต จ อวิตกฺก\-อวิจาโร จ ธมฺมา อวิตกฺก\-วิจารมตฺตสฺส จ อวิตกฺก\-อวิจารสฺส จ ธมฺมสฺส สหชาต\-ปจฺจเยน ปจฺจโย. อุปนิสฺสย\-ปจฺจเยน ปจฺจโย. ปุเรชาต\-ปจฺจเยน ปจฺจโย. (4)}

\prose{อวิตกฺก\-วิจารมตฺโต จ อวิตกฺก\-อวิจาโร จ ธมฺมา สวิตกฺก\-สวิจารสฺส จ อวิตกฺก\-วิจารมตฺตสฺส จ ธมฺมสฺส อารมฺมณ\-ปจฺจเยน ปจฺจโย. อุปนิสฺสย\-ปจฺจเยน ปจฺจโย. (5)}

\proseitem{144}{สวิตกฺก\-สวิจาโร จ อวิตกฺก\-วิจารมตฺโต จ ธมฺมา สวิตกฺก\-สวิจารสฺส ธมฺมสฺส อารมฺมณ\-ปจฺจเยน ปจฺจโย. สหชาต\-ปจฺจเยน ปจฺจโย. อุปนิสฺสย\-ปจฺจเยน ปจฺจโย. (1)}

\prose{สวิตกฺก\-สวิจาโร จ อวิตกฺก\-วิจารมตฺโต จ ธมฺมา อวิตกฺก\-วิจารมตฺตสฺส ธมฺมสฺส อารมฺมณ\-ปจฺจเยน ปจฺจโย. อุปนิสฺสย\-ปจฺจเยน ปจฺจโย. (2)}

\prose{สวิตกฺก\-สวิจาโร จ อวิตกฺก\-วิจารมตฺโต จ ธมฺมา อวิตกฺก\-อวิจารสฺส ธมฺมสฺส อารมฺมณ\-ปจฺจเยน ปจฺจโย. สหชาต\-ปจฺจเยน ปจฺจโย. อุปนิสฺสย\-ปจฺจเยน ปจฺจโย. ปจฺฉาชาต\-ปจฺจเยน ปจฺจโย. (3)}

\prose{สวิตกฺก\-สวิจาโร จ อวิตกฺก\-วิจารมตฺโต จ ธมฺมา สวิตกฺก\-สวิจารสฺส จ อวิตกฺก\-อวิจารสฺส จ ธมฺมสฺส สหชาต\-ปจฺจเยน ปจฺจโย. (4)}

\prose{สวิตกฺก\-สวิจาโร จ อวิตกฺก\-วิจารมตฺโต จ ธมฺมา อวิตกฺก\-วิจารมตฺตสฺส จ อวิตกฺก\-อวิจารสฺส จ ธมฺมสฺส อุปนิสฺสย\-ปจฺจเยน ปจฺจโย. (5)}

\prose{\csromanfolio{106}สวิตกฺก\-สวิจาโร จ อวิตกฺก\-วิจารมตฺโต จ ธมฺมา สวิตกฺก\-สวิจารสฺส จ อวิตกฺก\-วิจารมตฺตสฺส จ ธมฺมสฺส อารมฺมณ\-ปจฺจเยน ปจฺจโย. อุปนิสฺสย\-ปจฺจเยน ปจฺจโย. (6)}

\prose{สวิตกฺก\-สวิจาโร จ อวิตกฺก\-วิจารมตฺโต จ อวิตกฺก\-อวิจาโร จ ธมฺมา สวิตกฺก\-สวิจารสฺส ธมฺมสฺส สหชาตํ. ปุเรชาตํ. (1)}

\prose{สวิตกฺก\-สวิจาโร จ อวิตกฺก\-วิจารมตฺโต จ อวิตกฺก\-อวิจาโร จ ธมฺมา อวิตกฺก\-อวิจารสฺส ธมฺมสฺส สหชาตํ. ปจฺฉาชาตํ. อาหารํ. อินฺทฺริยํ. (2) ปจฺจนียุ\-ทฺธาโร.}

\csromanmarkh{2. ปจฺจย\-ปจฺจนีย}\csromanmarkhii{2. สงฺขฺยาวาร}\csromanheadernomark{2}{2. ปจฺจย\-ปจฺจนีย 2. สงฺขฺยาวาร}
\proseitem{145}{นเหตุยา ปญฺจติํส, นอารมฺมเณ ปญฺจติํส, นอธิปติยา ปญฺจติํส, นอนนฺตเร ปญฺจติํส, นสมนนฺตเร ปญฺจติํส, นสหชาเต เอกูนติํส, นอญฺญมญฺเญ เอกูนติํส, นนิสฺสเย เอกูนติํส, นอุปนิสฺสเย จตุตฺติํส, นปุเรชาเต ปญฺจติํส, นปจฺฉาชาเต นอาเสวเน นกมฺเม นวิปาเก นอาหาเร นอินฺทฺริเย นฌาเน นมคฺเค ปญฺจติํส, นสมฺปยุตฺเต เอกูนติํส, นวิปฺปยุตฺเต สตฺตวีส, โนอตฺถิยา สตฺตวีส, โนนตฺถิยา ปญฺจติํส, โนวิคเต ปญฺจติํส, โนอวิคเต สตฺตวีส.}

\prosewithcloserruled{(ปจฺจนียํ คเณนฺเตน อิมานิ ปทานิ อนุมชฺชนฺเตน คเณตพฺพานิ.)}{ปจฺจนียํ.}
\proseitem{146}{เหตุปจฺจยา นอารมฺมเณ เอกาทส, นอธิปติยา เอกาทส, นอนนฺตเร นสมนนฺตเร เอกาทส, นอญฺญมญฺเญ ตีณิ, นอุปนิสฺสเย เอกาทส, นปุเรชาเต นปจฺฉาชาเต นอาเสวเน นกมฺเม นวิปาเก \csromanfolio{107}นอาหาเร นอินฺทฺริเย นฌาเน นมคฺเค สพฺเพ เอกาทส, นสมฺปยุตฺเต ตีณิ, นวิปฺปยุตฺเต สตฺต, โนนตฺถิยา เอกาทส, โนวิคเต เอกาทส.}

\vspace{\tipitakaparskip}
\csromancenter{(อนุโลม\-ปจฺจนีย\-คณนา อิมินา การเณน คเณตพฺพา.)}
\nitthitamruled{อนุโลม\-ปจฺจนียํ.}
\proseitem{147}{นเหตุปจฺจยา อารมฺมเณ เอกวีส, อธิปติยา เตวีส, อนนฺตเร ปญฺจวีส, สมนนฺตเร ปญฺจวีส, สหชาเต ติํส, อญฺญมญฺเญ อฏฺฐวีส, นิสฺสเย ติํส, อุปนิสฺสเย ปญฺจวีส, ปุเรชาเต ปญฺจ, ปจฺฉาชาเต ปญฺจ, อาเสวเน เอกวีส, กมฺเม เอกาทส, วิปาเก เอกวีส, อาหาเร เอกาทส, อินฺทฺริเย เอกาทส, ฌาเน เอกวีส, มคฺเค โสฬส, สมฺปยุตฺเต เอกาทส, วิปฺปยุตฺเต นว, อตฺถิยา ติํส, นตฺถิยา ปญฺจวีส, วิคเต ปญฺจวีส, อวิคเต ติํส.}

\vspace{\tipitakaparskip}
\csromancenter{(ปจฺจนียา\-นุโลมํ อิมินา การเณน วิภชิตพฺพํ.)}
\csromancenter{ปจฺจนียา\-นุโลมํ.}
\nitthitam{วิตกฺกตฺติกํ นิฏฺฐิตํ.}
\csromanmarkh{7. ปีติตฺติก}\csromanmarkhii{1. ปฏิจฺจวาร}\csromanheadernomark{2}{7. \textbf{ปีติตฺติก 1. ปฏิจฺจวาร}}
\csromanmarkh{1. ปจฺจยานุโลม}\csromanmarkhii{1. วิภงฺควาร}\csromanheadernomark{2}{1. \textbf{ปจฺจยานุโลม 1. วิภงฺควาร}}
\csromanheader{2}{\textbf{เหตุ}}
\proseitem{1}{\csromanfolio{108}ปีติสหคตํ ธมฺมํ ปฏิจฺจ ปีติสหคโต ธมฺโม อุปฺปชฺชติ เหตุปจฺจยา. ปีติสหคตํ เอกํ ขนฺธํ ปฏิจฺจ ตโย ขนฺธา, ตโย ขนฺเธ ปฏิจฺจ เอโก ขนฺโธ, เทฺว ขนฺเธ ปฏิจฺจ เทฺว ขนฺธา. ปฏิสนฺธิ\-กฺขเณ ปีติสหคตํ เอกํ ขนฺธํ ปฏิจฺจ ตโย ขนฺธา ฯเปฯ เทฺว ขนฺเธ ปฏิจฺจ เทฺว ขนฺธา. (1)}

\prose{ปีติสหคตํ ธมฺมํ ปฏิจฺจ สุขสหคโต ธมฺโม อุปฺปชฺชติ เหตุปจฺจยา. ปีติสหคตํ เอกํ ขนฺธํ ปฏิจฺจ สุขสหคตา ตโย ขนฺธา ฯเปฯ เทฺว ขนฺเธ ปฏิจฺจ เทฺว ขนฺธา. ปฏิสนฺธิ\-กฺขเณ ปีติสหคตํ เอกํ ขนฺธํ ปฏิจฺจ สุขสหคตา ตโย ขนฺธา ฯเปฯ เทฺว ขนฺเธ ปฏิจฺจ เทฺว ขนฺธา. (2)}

\prose{ปีติสหคตํ ธมฺมํ ปฏิจฺจ ปีติสหคโต จ สุขสหคโต จ ธมฺมา อุปฺปชฺชนฺติ เหตุปจฺจยา. ปีติสหคตํ เอกํ ขนฺธํ ปฏิจฺจ ปีติสหคตา จ สุขสหคตา จ ตโย ขนฺธา ฯเปฯ เทฺว ขนฺเธ ปฏิจฺจ เทฺว ขนฺธา. ปฏิสนฺธิ\-กฺขเณ ปีติสหคตํ เอกํ ขนฺธํ ปฏิจฺจ ปีติสหคตา จ สุขสหคตา จ ตโย ขนฺธา ฯเปฯ เทฺว ขนฺเธ ปฏิจฺจ เทฺว ขนฺธา. (3)}

\proseitem{2}{สุขสหคตํ ธมฺมํ ปฏิจฺจ สุขสหคโต ธมฺโม อุปฺปชฺชติ เหตุปจฺจยา. สุขสหคตํ เอกํ ขนฺธํ ปฏิจฺจ เทฺว ขนฺธา, เทฺว ขนฺเธ ปฏิจฺจ เอโก ขนฺโธ. ปฏิสนฺธิ\-กฺขเณ สุขสหคตํ เอกํ ขนฺธํ ปฏิจฺจ เทฺว ขนฺธา, เทฺว ขนฺเธ ปฏิจฺจ เอโก ขนฺโธ. (1)}

\prose{สุขสหคตํ ธมฺมํ ปฏิจฺจ ปีติสหคโต ธมฺโม อุปฺปชฺชติ เหตุปจฺจยา. สุขสหคตํ เอกํ ขนฺธํ ปฏิจฺจ ปีติสหคตา ตโย ขนฺธา ฯเปฯ เทฺว ขนฺเธ ปฏิจฺจ เทฺว ขนฺธา. ปฏิสนฺธิ\-กฺขเณ สุขสหคตํ เอกํ ขนฺธํ ปฏิจฺจ ปีติสหคตา ตโย ขนฺธา ฯเปฯ เทฺว ขนฺเธ ปฏิจฺจ เทฺว ขนฺธา. (2)}

\prose{สุขสหคตํ ธมฺมํ ปฏิจฺจ ปีติสหคโต จ สุขสหคโต จ ธมฺมา อุปฺปชฺชนฺติ เหตุปจฺจยา. สุขสหคตํ เอกํ ขนฺธํ ปฏิจฺจ ปีติสหคตา จ สุขสหคตา จ เทฺว ขนฺธา, เทฺว ขนฺเธ ปฏิจฺจ เอโก ขนฺโธ. ปฏิสนฺธิ\-กฺขเณ \csromanfolio{109}สุขสหคตํ เอกํ ขนฺธํ ปฏิจฺจ ปีติสหคตา จ สุขสหคตา จ เทฺว ขนฺธา, เทฺว ขนฺเธ ปฏิจฺจ เอโก ขนฺโธ. (3)}

\prose{อุเปกฺขา\-สหคตํ ธมฺมํ ปฏิจฺจ อุเปกฺขา\-สหคโต ธมฺโม อุปฺปชฺชติ เหตุปจฺจยา. อุเปกฺขา\-สหคตํ เอกํ ขนฺธํ ปฏิจฺจ เทฺว ขนฺธา, เทฺว ขนฺเธ ปฏิจฺจ เอโก ขนฺโธ. ปฏิสนฺธิ\-กฺขเณ ฯเปฯ. (1)}

\prose{ปีติสหคตญฺจ สุขสหคตญฺจ ธมฺมํ ปฏิจฺจ ปีติสหคโต ธมฺโม อุปฺปชฺชติ เหตุปจฺจยา. ปีติสหคตญฺจ สุขสหคตญฺจ เอกํ ขนฺธํ ปฏิจฺจ ปีติสหคตา ตโย ขนฺธา ฯเปฯ เทฺว ขนฺเธ ปฏิจฺจ เทฺว ขนฺธา. ปฏิสนฺธิ\-กฺขเณ ปีติสหคตญฺจ สุขสหคตญฺจ เอกํ ขนฺธํ ปฏิจฺจ ปีติสหคตา ตโย ขนฺธา ฯเปฯ เทฺว ขนฺเธ ปฏิจฺจ เทฺว ขนฺธา. (1)}

\prose{ปีติสหคตญฺจ สุขสหคตญฺจ ธมฺมํ ปฏิจฺจ สุขสหคโต ธมฺโม อุปฺปชฺชติ เหตุปจฺจยา. ปีติสหคตญฺจ สุขสหคตญฺจ เอกํ ขนฺธํ ปฏิจฺจ สุขสหคตา เทฺว ขนฺธา, เทฺว ขนฺเธ ปฏิจฺจ เอโก ขนฺโธ. ปฏิสนฺธิ\-กฺขเณ ปีติสหคตญฺจ สุขสหคตญฺจ เอกํ ขนฺธํ ปฏิจฺจ สุขสหคตา เทฺว ขนฺธา, เทฺว ขนฺเธ ปฏิจฺจ เอโก ขนฺโธ. (2)}

\prose{ปีติสหคตญฺจ สุขสหคตญฺจ ธมฺมํ ปฏิจฺจ ปีติสหคโต จ สุขสหคโต จ ธมฺมา อุปฺปชฺชนฺติ เหตุปจฺจยา. ปีติสหคตญฺจ สุขสหคตญฺจ เอกํ ขนฺธํ ปฏิจฺจ ปีติสหคตา จ สุขสหคตา จ เทฺว ขนฺธา, เทฺว ขนฺเธ ปฏิจฺจ เอโก ขนฺโธ. ปฏิสนฺธิ\-กฺขเณ ปีติสหคตญฺจ สุขสหคตญฺจ เอกํ ขนฺธํ ปฏิจฺจ ปีติสหคตา จ สุขสหคตา จ เทฺว ขนฺธา, เทฺว ขนฺเธ ปฏิจฺจ เอโก ขนฺโธ. (3)}

\csromanheader{2}{\textbf{อารมฺมณาทิ}}
\proseitem{2}{ปีติสหคตํ ธมฺมํ ปฏิจฺจ ปีติสหคโต ธมฺโม อุปฺปชฺชติ อารมฺมณ\-ปจฺจยา. อธิปติ\-ปจฺจยา. (ปฏิสนฺธิ\-กฺขเณ นตฺถิ.) อนนฺตร\-ปจฺจยา. สมนนฺตร\-ปจฺจยา. สหชาต\-ปจฺจยา. อญฺญมญฺญ\-ปจฺจยา. นิสฺสยปจฺจยา. อุปนิสฺสย\-ปจฺจยา. ปุเรชาต\-ปจฺจยา. (ปุเรชาเต ปฏิสนฺธิ\-กฺขเณ นตฺถิ.) อาเสวนปจฺจยา. (อาเสวเน วิปากํ นตฺถิ.) กมฺมปจฺจยา. วิปากปจฺจยา. อาหาร ฯเปฯ. อินฺทฺริย. ฌาน. มคฺค. สมฺปยุตฺต. วิปฺปยุตฺต. อตฺถิ. นตฺถิ. วิคต. อวิคตปจฺจยา.}

\csromanmarkh{1. ปจฺจยานุโลม}\csromanmarkhii{2. สงฺขฺยาวาร}\csromanheadernomark{2}{1. \textbf{ปจฺจยานุโลม 2. สงฺขฺยาวาร}}
\proseitemwithcloserruled{5}{\csromanfolio{110}เหตุยา ทส, อารมฺมเณ ทส, อธิปติยา ทส, อนนฺตเร สมนนฺตเร สหชาเต อญฺญมญฺเญ นิสฺสเย อุปนิสฺสเย ปุเรชาเต อาเสวเน กมฺเม วิปาเก อาหาเร อินฺทฺริเย ฌาเน มคฺเค สมฺปยุตฺเต วิปฺปยุตฺเต อตฺถิยา นตฺถิยา วิคเต อวิคเต สพฺพตฺถ ทส. (เอวํ อนุโลมคณนา คเณตพฺพา.)}{อนุโลมํ.}
\csromanmarkh{2. ปจฺจย\-ปจฺจนีย}\csromanmarkhii{1. วิภงฺควาร}\csromanheadernomark{2}{2. ปจฺจย\-ปจฺจนีย 1. วิภงฺควาร}
\csromanheader{2}{\textbf{นเหตุ}}
\proseitem{6}{ปีติสหคตํ ธมฺมํ ปฏิจฺจ ปีติสหคโต ธมฺโม อุปฺปชฺชติ นเหตุปจฺจยา. อเหตุกํ ปีติสหคตํ เอกํ ขนฺธํ ปฏิจฺจ ตโย ขนฺธา ฯเปฯ เทฺว ขนฺเธ ปฏิจฺจ เทฺว ขนฺธา. (1)}

\prose{ปีติสหคตํ ธมฺมํ ปฏิจฺจ สุขสหคโต ธมฺโม อุปฺปชฺชติ นเหตุปจฺจยา. อเหตุกํ ปีติสหคตํ เอกํ ขนฺธํ ปฏิจฺจ สุขสหคตา ตโย ขนฺธา ฯเปฯ เทฺว ขนฺเธ ปฏิจฺจ เทฺว ขนฺธา. (2)}

\prose{ปีติสหคตํ ธมฺมํ ปฏิจฺจ ปีติสหคโต จ สุขสหคโต จ ธมฺมา อุปฺปชฺชนฺติ นเหตุปจฺจยา. อเหตุกํ ปีติสหคตํ เอกํ ขนฺธํ ปฏิจฺจ ปีติสหคตา จ สุขสหคตา จ ตโย ขนฺธา ฯเปฯ เทฺว ขนฺเธ ปฏิจฺจ เทฺว ขนฺธา. (3)}

\proseitem{7}{สุขสหคตํ ธมฺมํ ปฏิจฺจ สุขสหคโต ธมฺโม อุปฺปชฺชติ นเหตุปจฺจยา. อเหตุกํ สุขสหคตํ เอกํ ขนฺธํ ปฏิจฺจ เทฺว ขนฺธา, เทฺว ขนฺเธ ปฏิจฺจ เอโก ขนฺโธ. (1)}

\prose{สุขสหคตํ ธมฺมํ ปฏิจฺจ ปีติสหคโต ธมฺโม อุปฺปชฺชติ นเหตุปจฺจยา. อเหตุกํ สุขสหคตํ เอกํ ขนฺธํ ปฏิจฺจ ปีติสหคตา ตโย ขนฺธา ฯเปฯ เทฺว ขนฺเธ ปฏิจฺจ เทฺว ขนฺธา. (2)}

\prose{\csromanfolio{111}สุขสหคตํ ธมฺมํ ปฏิจฺจ ปีติสหคโต จ สุขสหคโต จ ธมฺมา อุปฺปชฺชนฺติ นเหตุปจฺจยา. อเหตุกํ สุขสหคตํ เอกํ ขนฺธํ ปฏิจฺจ ปีติสหคตา จ สุขสหคตา จ เทฺว ขนฺธา, เทฺว ขนฺเธ ปฏิจฺจ เอโก ขนฺโธ. (3)}

\proseitem{8}{อุเปกฺขา\-สหคตํ ธมฺมํ ปฏิจฺจ อุเปกฺขา\-สหคโต ธมฺโม อุปฺปชฺชติ นเหตุปจฺจยา. อเหตุกํ อุเปกฺขา\-สหคตํ เอกํ ขนฺธํ ปฏิจฺจ เทฺว ขนฺธา, เทฺว ขนฺเธ ปฏิจฺจ เอโก ขนฺโธ. อเหตุก\-ปฏิสนฺธิ\-กฺขเณ อุเปกฺขา\-สหคตํ เอกํ ขนฺธํ ปฏิจฺจ เทฺว ขนฺธา, เทฺว ขนฺเธ ปฏิจฺจ เอโก ขนฺโธ. วิจิกิจฺฉา\-สหคเต อุทฺธจฺจ\-สหคเต ขนฺเธ ปฏิจฺจ วิจิกิจฺฉา\-สหคโต อุทฺธจฺจ\-สหคโต โมโห. (1)}

\proseitem{9}{ปีติสหคตญฺจ สุขสหคตญฺจ ธมฺมํ ปฏิจฺจ ปีติสหคโต ธมฺโม อุปฺปชฺชติ นเหตุปจฺจยา. อเหตุกํ ปีติสหคตญฺจ สุขสหคตญฺจ เอกํ ขนฺธํ ปฏิจฺจ ปีติสหคตา ตโย ขนฺธา ฯเปฯ เทฺว ขนฺเธ ปฏิจฺจ เทฺว ขนฺธา. (1)}

\prose{ปีติสหคตญฺจ สุขสหคตญฺจ ธมฺมํ ปฏิจฺจ สุขสหคโต ธมฺโม อุปฺปชฺชติ นเหตุปจฺจยา. อเหตุกํ ปีติสหคตญฺจ สุขสหคตญฺจ เอกํ ขนฺธํ ปฏิจฺจ สุขสหคตา เทฺว ขนฺธา, เทฺว ขนฺเธ ปฏิจฺจ เอโก ขนฺโธ. (2)}

\prose{ปีติสหคตญฺจ สุขสหคตญฺจ ธมฺมํ ปฏิจฺจ ปีติสหคโต จ สุขสหคโต จ ธมฺมา อุปฺปชฺชนฺติ นเหตุปจฺจยา. อเหตุกํ ปีติสหคตญฺจ สุขสหคตญฺจ เอกํ ขนฺธํ ปฏิจฺจ ปีติสหคตา จ สุขสหคตา จ เทฺว ขนฺธา, เทฺว ขนฺเธ ปฏิจฺจ เอโก ขนฺโธ. (3)}

\prose{นอธิปติ\csromandash{}นอาเสวน}

\proseitem{10}{ปีติสหคตํ ธมฺมํ ปฏิจฺจ ปีติสหคโต ธมฺโม อุปฺปชฺชติ น อธิปติ\-ปจฺจยา. (นอธิปติ ปฏิสนฺธิ\-กฺขเณ ปริปุณฺณํ.) นปุเรชาต\-ปจฺจยา. (“อรูเป”ติ นิยาเมตพฺพํ “ปฏิสนฺธิ\-กฺขเณ”ติ จ.) นปจฺฉาชาต\-ปจฺจยา. นอาเสวน\-ปจฺจยา.}

\csromanheader{2}{\textbf{นกมฺม}}
\proseitem{11}{ปีติสหคตํ ธมฺมํ ปฏิจฺจ ปีติสหคโต ธมฺโม อุปฺปชฺชติ นกมฺมปจฺจยา. ปีติสหคเต ขนฺเธ ปฏิจฺจ ปีติสหคตา เจตนา.}

\prose{\csromanfolio{112}ปีติสหคตํ ธมฺมํ ปฏิจฺจ สุขสหคโต ธมฺโม อุปฺปชฺชติ นกมฺมปจฺจยา. ปีติสหคเต ขนฺเธ ปฏิจฺจ สุขสหคตา เจตนา. (อิมินา การเณน ทส ปญฺหา วิตฺถาเรตพฺพา.)}

\vspace{\tipitakaparskip}
\csromancenter{(อิมินา การเณน ทส ปญฺหา วิตฺถาเรตพฺพา.)}
\proseitem{12}{ปีติสหคตํ ธมฺมํ ปฏิจฺจ ปีติสหคโต ธมฺโม อุปฺปชฺชติ นวิปาก\-ปจฺจยา ฯเปฯ. (ปริปุณฺณํ, ปฏิสนฺธิ นตฺถิ.)}

\csromanheader{2}{\textbf{นฌานาทิ}}
\proseitem{13}{สุขสหคตํ ธมฺมํ ปฏิจฺจ สุขสหคโต ธมฺโม อุปฺปชฺชติ นฌานปจฺจยา. สุขสหคตํ กาย\-วิญฺญาณสหคตํ เอกํ ขนฺธํ ปฏิจฺจ เทฺว ขนฺธา, เทฺว ขนฺเธ ปฏิจฺจ เอโก ขนฺโธ. (1)}

\prose{อุเปกฺขา\-สหคตํ ธมฺมํ ปฏิจฺจ อุเปกฺขา\-สหคโต ธมฺโม อุปฺปชฺชติ นฌานปจฺจยา. จตุวิญฺญาณสหคตํ เอกํ ขนฺธํ ปฏิจฺจ เทฺว ขนฺธา, เทฺว ขนฺเธ ปฏิจฺจ เอโก ขนฺโธ. (1)}

\prose{(นมคฺคปจฺจยา นเหตุปจฺจย\-สทิสํ. โมโห นตฺถิ. นวิปฺปยุตฺต\-ปจฺจยา ปริปุณฺณํ. อรูปปญฺหเมว.)}

\csromanmarkh{2. ปจฺจย\-ปจฺจนีย}\csromanmarkhii{2. สงฺขฺยาวาร}\csromanheadernomark{2}{2. ปจฺจย\-ปจฺจนีย 2. สงฺขฺยาวาร}
\proseitemwithcloserruled{14}{นเหตุยา ทส, นอธิปติยา ทส, นปุเรชาเต นปจฺฉาชาเต นอาเสวเน นกมฺเม นวิปาเก ทส, นฌาเน เทฺว, นมคฺเค ทส, นวิปฺปยุตฺเต ทส. (ปจฺจนียํ ปริปุณฺณํ กาตพฺพํ.)}{ปจฺจนียํ.}
\csromanmatikaanchor{mtk.3}
\csromantocmarknopage{chapter}{3. ปจฺจยานุโลมปจฺจนีย}{mtk.3}
\bookmark[dest={mtk.3},level=0]{3. ปจฺจยานุโลมปจฺจนีย}
\chapterheadpageread{7. ปีติตฺติก 3. ปจฺจยา\-นุโลม\-ปจฺจนีย}
\vspace{\tipitakaparskip}
\csromancenter{\textbf{เหตุ}ปจฺจยา นอธิปติยา ทส, นปุเรชาเต ทส, นปจฺฉาชาเต นอาเสวเน นกมฺเม นวิปาเก นวิปฺปยุตฺเต ทส. (อนุโลมปจฺจนียํ วิตฺถาเรน คเณตพฺพํ.)}
\csromancenter{(อนุโลม\-ปจฺจนียํ วิตฺถาเรน คเณตพฺพํ.)}
\nitthitamruled{อนุโลม\-ปจฺจนียํ.}
\proseitem{16}{\csromanfolio{113}นเหตุปจฺจยา อารมฺมเณ ทส, อนนฺตเร ทส, สมนนฺตเร ทส, สหชาเต อญฺญมญฺเญ นิสฺสเย อุปนิสฺสเย ปุเรชาเต อาเสวเน กมฺเม วิปาเก อาหาเร อินฺทฺริเย ฌาเน สพฺเพ ทส, มคฺเค เอกํ, สมฺปยุตฺเต ทส, วิปฺปยุตฺเต อตฺถิยา นตฺถิยา วิคเต อวิคเต สพฺเพ ทส. ปจฺจนียา\-นุโลมํ. ปฏิจฺจวาโร.}

\prose{(สหชาตวาโรปิ ปจฺจยวาโรปิ นิสฺสยวาโรปิ สํสฏฺฐ\-วาโรปิ สมฺปยุตฺตวาโรปิ ปฏิจฺจ\-วาร\-สทิสา.)}
\csromansectionrule

\csromanmarkh{7. ปีติตฺติก}\csromanmarkhii{7. ปญฺหาวาร}\csromanheadernomark{2}{7. ปีติตฺติก 7. ปญฺหาวาร}
\csromanmarkh{1. ปจฺจยานุโลม}\csromanmarkhii{1. วิภงฺควาร}\csromanheadernomark{2}{1. \textbf{ปจฺจยานุโลม 1. วิภงฺควาร}}
\vspace{\tipitakaparskip}
\csromancenter{\textbf{เหตุ}}
\proseitem{17}{ปีติสหคโต ธมฺโม ปีติสหคตสฺส ธมฺมสฺส เหตุปจฺจเยน ปจฺจโย. ปีติสหคตา เหตู สมฺปยุตฺตกานํ ขนฺธานํ เหตุปจฺจเยน ปจฺจโย. ปฏิสนฺธิ\-กฺขเณ ปีติสหคตา เหตู สมฺปยุตฺตกานํ ขนฺธานํ เหตุปจฺจเยน ปจฺจโย. (1)}

\prose{\csromanfolio{114}ปีติสหคโต ธมฺโม สุขสหคตสฺส ธมฺมสฺส เหตุปจฺจเยน ปจฺจโย. ปีติสหคตา เหตู สมฺปยุตฺตกานํ สุขสหคตานํ ขนฺธานํ เหตุปจฺจเยน ปจฺจโย. ปฏิสนฺธิ\-กฺขเณ ฯเปฯ. (2)}

\prose{ปีติสหคโต ธมฺโม ปีติสหคตสฺส จ สุขสหคตสฺส จ ธมฺมสฺส เหตุปจฺจเยน ปจฺจโย. ปีติสหคตา เหตู สมฺปยุตฺตกานํ ปีติสหคตานญฺจ สุขสหคตานญฺจ ขนฺธานํ เหตุปจฺจเยน ปจฺจโย. ปฏิสนฺธิ\-กฺขเณ ฯเปฯ. (3)}

\proseitem{18}{สุขสหคโต ธมฺโม สุขสหคตสฺส ธมฺมสฺส ฯเปฯ ปีติสหคตสฺส ธมฺมสฺส ฯเปฯ ปีติสหคตสฺส จ สุขสหคตสฺส จ ธมฺมสฺส ฯเปฯ. (สุขมูเล ตีณิ.)}

\prose{อุเปกฺขา\-สหคโต ธมฺโม อุเปกฺขา\-สหคตสฺส ธมฺมสฺส เหตุปจฺจเยน ปจฺจโย. อุเปกฺขา\-สหคตา เหตู สมฺปยุตฺตกานํ ขนฺธานํ เหตุปจฺจเยน ปจฺจโย. ปฏิสนฺธิ\-กฺขเณ ฯเปฯ. (1)}

\prose{ปีติสหคโต จ สุขสหคโต จ ธมฺมา ปีติสหคตสฺส ธมฺมสฺส ฯเปฯ. สุขสหคตสฺส ธมฺมสฺส ฯเปฯ. ปีติสหคตสฺส จ สุขสหคตสฺส จ ธมฺมสฺส เหตุปจฺจเยน ปจฺจโย. ปีติสหคตา จ สุขสหคตา จ เหตู สมฺปยุตฺตกานํ ปีติสหคตานญฺจ สุขสหคตานญฺจ ขนฺธานํ เหตุปจฺจเยน ปจฺจโย. ปฏิสนฺธิ\-กฺขเณ ฯเปฯ. (3)}

\csromanheader{2}{\textbf{อารมฺมณ}}
\proseitem{19}{ปีติสหคโต ธมฺโม ปีติสหคตสฺส ธมฺมสฺส อารมฺมณ\-ปจฺจเยน ปจฺจโย. ปีติสหคเตน จิตฺเตน ทานํ ทตฺวา, สีลํ สมาทิยิตฺวา, อุโปสถกมฺมํ กตฺวา ตํ ปีติสหคเตน จิตฺเตน ปจฺจเวกฺขติ, ปีติสหคตา ฌานา วุฏฺฐหิตฺวา, มคฺคา วุฏฺฐหิตฺวา, ผลา วุฏฺฐหิตฺวา ตํ ปีติสหคเตน จิตฺเตน ปจฺจเวกฺขติ. อริยา ปีติสหคเตน จิตฺเตน ปีติสหคเต ปหีเน กิเลเส ปจฺจเวกฺขนฺติ, วิกฺขมฺภิเต กิเลเส ปจฺจเวกฺขนฺติ, ปุพฺเพ สมุทาจิณฺเณ กิเลเส ชานนฺติ. ปีติสหคเต ขนฺเธ ปีติสหคเตน จิตฺเตน อนิจฺจโต ทุกฺขโต อนตฺตโต วิปสฺสนฺติ, อสฺสาเทนฺติ อภินนฺทนฺติ, ตํ อารพฺภ ปีติสหคโต ราโค อุปฺปชฺชติ, ทิฏฺฐิ อุปฺปชฺชติ. ปีติสหคเต ขนฺเธ อารพฺภ ปีติสหคตา ขนฺธา อุปฺปชฺชนฺติ. (1)}

\prose{\csromanfolio{115}ปีติสหคโต ธมฺโม สุขสหตสฺส ธมฺมสฺส อารมฺมณ\-ปจฺจเยน ปจฺจโย. ปีติสหคเตน จิตฺเตน ทานํ ทตฺวา, สีลํ สมาทิยิตฺวา, อุโปสถกมฺมํ กตฺวา ตํ สุขสหคเตน จิตฺเตน ปจฺจเวกฺขติ, ปีติสหคตา ฌานา วุฏฺฐหิตฺวา, มคฺคา วุฏฺฐหิตฺวา, ผลา วุฏฺฐหิตฺวา ตํ สุขสหคเตน จิตฺเตน ปจฺจเวกฺขติ. อริยา สุขสหคเตน จิตฺเตน ปีติสหคเต ปหีเน กิเลเส ปจฺจเวกฺขนฺติ, วิกฺขมฺภิเต กิเลเส ปจฺจเวกฺขนฺติ, ปุพฺเพ สมุทาจิณฺเณ กิเลเส ชานนฺติ. ปีติสหคเต ขนฺเธ สุขสหคเตน จิตฺเตน อนิจฺจโต ทุกฺขโต อนตฺตโต วิปสฺสนฺติ, อสฺสาเทนฺติ อภินนฺทนฺติ, ตํ อารพฺภ สุขสหคโต ราโค อุปฺปชฺชติ, ทิฏฺฐิ อุปฺปชฺชติ. ปีติสหคเต ขนฺเธ อารพฺภ สุขสหคตา ขนฺธา อุปฺปชฺชนฺติ. (2)}

\prose{ปีติสหคโต ธมฺโม อุเปกฺขา\-สหคตสฺส ธมฺมสฺส อารมฺมณ\-ปจฺจเยน ปจฺจโย. ปีติสหคเตน จิตฺเตน ทานํ ทตฺวา, สีลํ สมาทิยิตฺวา, อุโปสถกมฺมํ กตฺวา ตํ อุเปกฺขา\-สหคเตน จิตฺเตน ปจฺจเวกฺขติ, ปีติสหคตา ฌานา วุฏฺฐหิตฺวา, มคฺคา วุฏฺฐหิตฺวา, ผลา วุฏฺฐหิตฺวา ตํ อุเปกฺขา\-สหคเตน จิตฺเตน ปจฺจเวกฺขติ. อริยา อุเปกฺขา\-สหคเตน จิตฺเตน ปีติสหคเต ปหีเน กิเลเส ปจฺจเวกฺขนฺติ, วิกฺขมฺภิเต กิเลเส ปจฺจเวกฺขนฺติ, ปุพฺเพ สมุทาจิณฺเณ กิเลเส ชานนฺติ. ปีติสหคเต ขนฺเธ อุเปกฺขา\-สหคเตน จิตฺเตน อนิจฺจโต ทุกฺขโต อนตฺตโต วิปสฺสนฺติ, อสฺสาเทนฺติ อภินนฺทนฺติ, ตํ อารพฺภ อุเปกฺขา\-สหคโต ราโค อุปฺปชฺชติ, ทิฏฺฐิ อุปฺปชฺชติ, วิจิกิจฺฉา อุปฺปชฺชติ, อุทฺธจฺจํ อุปฺปชฺชติ. เจโตปริย\-ญาเณน ปีติสหคต\-จิตฺต\-สมงฺคิสฺส จิตฺตํ ชานนฺติ, ปีติสหคตา ขนฺธา เจโตปริย\-ญาณสฺส, ปุพฺเพ\-นิวาสา\-นุสฺสติ\-ญาณสฺส, ยถากมฺมูปคญาณสฺส, อนาคตํสญาณสฺส, อาวชฺชนาย อารมฺมณ\-ปจฺจเยน ปจฺจโย. ปีติสหคเต ขนฺเธ อารพฺภ อุเปกฺขา\-สหคตา ขนฺธา อุปฺปชฺชนฺติ. (3)}

\prose{ปีติสหคโต ธมฺโม ปีติสหคตสฺส จ สุขสหคตสฺส จ ธมฺมสฺส อารมฺมณ\-ปจฺจเยน ปจฺจโย. ปีติสหคเตน จิตฺเตน ทานํ ทตฺวา, สีลํ สมาทิยิตฺวา, อุโปสถกมฺมํ กตฺวา ตํ ปีติสหคเตน จ สุขสหคเตน จ จิตฺเตน ปจฺจเวกฺขติ, ปีติสหคตา ฌานา วุฏฺฐหิตฺวา, มคฺคา วุฏฺฐหิตฺวา, ผลา วุฏฺฐหิตฺวา ตํ ปีติสหคเตน จ สุขสหคเตน จ จิตฺเตน ปจฺจเวกฺขติ. อริยา ปีติสหคเตน จ สุขสหคเตน จ จิตฺเตน ปีติสหคเต ปหีเน กิเลเส ปจฺจเวกฺขนฺติ, วิกฺขมฺภิเต กิเลเส}

\prose{\csromanfolio{116}ปจฺจเวกฺขนฺติ, ปุพฺเพ สมุทาจิณฺเณ กิเลเส ชานนฺติ. ปีติสหคเต ขนฺเธ ปีติสหคเตน จ สุขสหคเตน จ จิตฺเตน อนิจฺจโต ทุกฺขโต อนตฺตโต วิปสฺสนฺติ, อสฺสาเทนฺติ อภินนฺทนฺติ, ตํ อารพฺภ ปีติสหคโต จ สุขสหคโต จ ราโค อุปฺปชฺชติ, ทิฏฺฐิ อุปฺปชฺชติ. ปีติสหคเต ขนฺเธ อารพฺภ ปีติสหคตา จ สุขสหคตา จ ขนฺธา อุปฺปชฺชนฺติ. (4)}

\proseitem{20}{สุขสหคโต ธมฺโม สุขสหคตสฺส ธมฺมสฺส อารมฺมณ\-ปจฺจเยน ปจฺจโย. สุขสหคโต ธมฺโม ปีติสหคตสฺส ธมฺมสฺส ฯเปฯ อุเปกฺขา\-สหคตสฺส ธมฺมสฺส ฯเปฯ ปีติสหคตสฺส จ สุขสหคตสฺส จ ธมฺมสฺส อารมฺมณ\-ปจฺจเยน ปจฺจโย. สุขสหคเต ขนฺเธ อารพฺภ ปีติสหคตา จ สุขสหคตา จ ขนฺธา อุปฺปชฺชนฺติ. (4)}

\prose{อุเปกฺขา\-สหคโต ธมฺโม อุเปกฺขา\-สหคตสฺส ธมฺมสฺส อารมฺมณ\-ปจฺจเยน ปจฺจโย. อุเปกฺขา\-สหคเตน จิตฺเตน ทานํ ทตฺวา, สีลํ สมาทิยิตฺวา, อุโปสถกมฺมํ กตฺวา ตํ อุเปกฺขา\-สหคเตน จิตฺเตน ปจฺจเวกฺขติ, อุเปกฺขา\-สหคตา ฌานา วุฏฺฐหิตฺวา, มคฺคา วุฏฺฐหิตฺวา, ผลา วุฏฺฐหิตฺวา ตํ อุเปกฺขา\-สหคเตน จิตฺเตน ปจฺจเวกฺขติ. อริยา อุเปกฺขา\-สหคเตน จิตฺเตน อุเปกฺขา\-สหคเต ปหีเน กิเลเส ปจฺจเวกฺขนฺติ, วิกฺขมฺภิเต กิเลเส ปจฺจเวกฺขนฺติ, ปุพฺเพ สมุทาจิณฺเณ กิเลเส ชานนฺติ. อุเปกฺขา\-สหคเต ขนฺเธ อุเปกฺขา\-สหคเตน จิตฺเตน อนิจฺจโต ทุกฺขโต อนตฺตโต วิปสฺสนฺติ, อสฺสาเทนฺติ อภินนฺทนฺติ, ตํ อารพฺภ อุเปกฺขา\-สหคโต ราโค อุปฺปชฺชติ, ทิฏฺฐิ อุปฺปชฺชติ, วิจิกิจฺฉา อุปฺปชฺชติ, อุทฺธจฺจํ อุปฺปชฺชติ. เจโตปริย\-ญาเณน อุเปกฺขาสหคต\-จิตฺต\-สมงฺคิสฺส จิตฺตํ ชานาติ.  อากาสา\-นญฺจา\-ยตนํ วิญฺญาณญฺจา\-ยตนสฺส อารมฺมณ\-ปจฺจเยน ปจฺจโย ฯเปฯ. อากิญฺจญฺญา\-ยตนํ เนว\-สญฺญา\-นา\-สญฺญา\-ยตนสฺส อารมฺมณ\-ปจฺจเยน ปจฺจโย. อุเปกฺขา\-สหคตา ขนฺธา อิทฺธิวิธ\-ญาณสฺส, เจโตปริย\-ญาณสฺส, ปุพฺเพ\-นิวาสา\-นุสฺสติ\-ญาณสฺส, ยถากมฺมูปคญาณสฺส, อนาคตํสญาณสฺส, อาวชฺชนาย อารมฺมณ\-ปจฺจเยน ปจฺจโย. อุเปกฺขา\-สหคเต ขนฺเธ อารพฺภ อุเปกฺขา\-สหคตา ขนฺธา อุปฺปชฺชนฺติ. (1)}

\prose{อุเปกฺขา\-สหคโต ธมฺโม ปีติสหคตสฺส ธมฺมสฺส ฯเปฯ สุขสหคตสฺส ธมฺมสฺส ฯเปฯ ปีติสหคตสฺส จ สุขสหคตสฺส จ ธมฺมสฺส อารมฺมณ\-ปจฺจเยน ปจฺจโย. อุเปกฺขา\-สหคเตน จิตฺเตน ทานํ ทตฺวา, \csromanfolio{117}สีลํ สมาทิยิตฺวา, อุโปสถกมฺมํ กตฺวา ตํ ปีติสหคเตน จ สุขสหคเตน จ จิตฺเตน ปจฺจเวกฺขติ, อุเปกฺขา\-สหคตา ฌานา วุฏฺฐหิตฺวา, มคฺคา วุฏฺฐหิตฺวา, ผลา วุฏฺฐหิตฺวา ตํ ปีติสหคเตน จ สุขสหคเตน จ จิตฺเตน ปจฺจเวกฺขติ. อริยา ปีติสหคเตน จ สุขสหคเตน จ จิตฺเตน อุเปกฺขา\-สหคเต ปหีเน กิเลเส ปจฺจเวกฺขนฺติ, วิกฺขมฺภิเต กิเลเส ปจฺจเวกฺขนฺติ, ปุพฺเพ สมุทาจิณฺเณ กิเลเส ชานนฺติ. อุเปกฺขา\-สหคเต ขนฺเธ ปีติสหคเตน จ สุขสหคเตน จ จิตฺเตน อนิจฺจโต ทุกฺขโต อนตฺตโต วิปสฺสนฺติ, อสฺสาเทนฺติ อภินนฺทนฺติ, ตํ อารพฺภ ปีติสหคโต จ สุขสหคโต จ ราโค อุปฺปชฺชติ, ทิฏฺฐิ อุปฺปชฺชติ. อุเปกฺขา\-สหคเต ขนฺเธ อารพฺภ ปีติสหคตา จ สุขสหคตา จ ขนฺธา อุปฺปชฺชนฺติ. (3)}

\prose{ปีติสหคโต จ สุขสหคโต จ ธมฺมา ปีติสหคตสฺส ธมฺมสฺส ฯเปฯ สุขสหคตสฺส ธมฺมสฺส ฯเปฯ อุเปกฺขา\-สหคตสฺส ธมฺมสฺส อารมฺมณ\-ปจฺจเยน ปจฺจโย. ปีติสหคเตน จ สุขสหคเตน จ จิตฺเตน ทานํ ทตฺวา, สีลํ สมาทิยิตฺวา, อุโปสถกมฺมํ กตฺวา ฯเปฯ. ปีติสหคเต จ สุขสหคเต จ ขนฺเธ อุเปกฺขา\-สหคเตน จิตฺเตน อนิจฺจโต ทุกฺขโต อนตฺตโต วิปสฺสติ, อสฺสาเทติ อภินนฺทติ, ตํ อารพฺภ อุเปกฺขา\-สหคโต ราโค อุปฺปชฺชติ, ทิฏฺฐิ อุปฺปชฺชติ, วิจิกิจฺฉา อุปฺปชฺชติ, อุทฺธจฺจํ อุปฺปชฺชติ. เจโตปริย\-ญาเณน ปีติสหคต\-สุขสหคต\-จิตฺต\-สมงฺคิสฺส จิตฺตํ ชานาติ. ปีติสหคตา จ สุขสหคตา จ ขนฺธา เจโตปริย\-ญาณสฺส, ปุพฺเพ\-นิวาสา\-นุสฺสติ\-ญาณสฺส, ยถากมฺมูปคญาณสฺส, อนาคตํสญาณสฺส, อาวชฺชนาย อารมฺมณ\-ปจฺจเยน ปจฺจโย. ปีติสหคเต จ สุขสหคเต จ ขนฺเธ อารพฺภ อุเปกฺขา\-สหคตา ขนฺธา อุปฺปชฺชนฺติ. (3)}

\prose{ปีติสหคโต จ สุขสหคโต จ ธมฺมา ปีติสหคตสฺส จ สุขสหคตสฺส จ ธมฺมสฺส อารมฺมณ\-ปจฺจเยน ปจฺจโย. (สํขิตฺตํ.) (4)}

\prose{ปีติสหคโต ธมฺโม ปีติสหคตสฺส ธมฺมสฺส อธิปติ\-ปจฺจเยน ปจฺจโย. อารมฺมณา\-ธิปติ, สหชาตา\-ธิปติ.  อารมฺมณา\-ธิปติ\csromandash{} ปีติสหคเตน จิตฺเตน ทานํ ทตฺวา, สีลํ สมาทิยิตฺวา, อุโปสถกมฺมํ \csromanfolio{118}กตฺวา ปีติสหคเตน จิตฺเตน ตํ ครุํ กตฺวา ปจฺจเวกฺขติ, ปีติสหคตา ฌานา วุฏฺฐหิตฺวา, มคฺคา วุฏฺฐหิตฺวา, ผลา วุฏฺฐหิตฺวา ปีติสหคเตน จิตฺเตน ตํ ครุํ กตฺวา ปจฺจเวกฺขติ. ปีติสหคเต ขนฺเธ ปีติสหคเตน จิตฺเตน ครุํ กตฺวา อสฺสาเทติ อภินนฺทติ, ตํ ครุํ กตฺวา ปีติสหคโต ราโค อุปฺปชฺชติ, ทิฏฺฐิ อุปฺปชฺชติ. สหชาตา\-ธิปติ\csromandash{}ปีติสหคตา\-ธิปติ สมฺปยุตฺตกานํ ขนฺธานํ อธิปติ\-ปจฺจเยน ปจฺจโย. (1)}

\proseitem{22}{ปีติสหคโต ธมฺโม สุขสหคตสฺส ธมฺมสฺส อธิปติ\-ปจฺจเยน ปจฺจโย. อารมฺมณา\-ธิปติ, สหชาตา\-ธิปติ.  อารมฺมณา\-ธิปติ\csromandash{} ปีติสหคเตน จิตฺเตน ทานํ ทตฺวา ฯเปฯ. สหชาตา\-ธิปติ\csromandash{} ปีติสหคตา\-ธิปติ สมฺปยุตฺตกานํ สุขสหคตานํ ขนฺธานํ อธิปติ\-ปจฺจเยน ปจฺจโย. (2)}

\prose{ปีติสหคโต ธมฺโม อุเปกฺขา\-สหคตสฺส ธมฺมสฺส อธิปติ\-ปจฺจเยน ปจฺจโย. อารมฺมณา\-ธิปติ\csromandash{}ปีติสหคเตน จิตฺเตน ทานํ ทตฺวา ฯเปฯ อุเปกฺขา\-สหคเตน จิตฺเตน. (สํขิตฺตํ.) (3)}

\prose{ปีติสหคโต ธมฺโม ปีติสหคตสฺส จ สุขสหคตสฺส จ ธมฺมสฺส อธิปติ\-ปจฺจเยน ปจฺจโย. อารมฺมณา\-ธิปติ, สหชาตา\-ธิปติ.  อารมฺมณา\-ธิปติ\csromandash{}ปีติสหคเตน จิตฺเตน ทานํ ทตฺวา ฯเปฯ. สหชาตา\-ธิปติ\csromandash{}ปีติสหคตา\-ธิปติ สมฺปยุตฺตกานํ ปีติสหคตานญฺจ สุขสหคตานญฺจ ขนฺธานํ อธิปติ\-ปจฺจเยน ปจฺจโย. (4)}

\proseitem{23}{สุขสหคโต ธมฺโม สุขสหคตสฺส ธมฺมสฺส อธิปติ\-ปจฺจเยน ปจฺจโย. อารมฺมณา\-ธิปติ, สหชาตา\-ธิปติ.  อารมฺมณา\-ธิปติ\csromandash{}สุขสหคเตน จิตฺเตน ทานํ ทตฺวา ฯเปฯ. สหชาตา\-ธิปติ\csromandash{}สุขสหคตา\-ธิปติ สมฺปยุตฺตกานํ ขนฺธานํ อธิปติ\-ปจฺจเยน ปจฺจโย. (1)}

\prose{สุขสหคโต ธมฺโม ปีติสหคตสฺส ธมฺมสฺส อธิปติ\-ปจฺจเยน ปจฺจโย. อารมฺมณา\-ธิปติ, สหชาตา\-ธิปติ ฯเปฯ. สหชาตา\-ธิปติ\csromandash{} สุขสหคตา\-ธิปติ สมฺปยุตฺตกานํ ปีติสหคตานํ ขนฺธานํ อธิปติ\-ปจฺจเยน ปจฺจโย. (2)}

\prose{สุขสหคโต ธมฺโม อุเปกฺขา\-สหคตสฺส ธมฺมสฺส อธิปติ\-ปจฺจเยน ปจฺจโย. อารมฺมณา\-ธิปติ. (สํขิตฺตํ.) (3)}

\prose{\csromanfolio{119}สุขสหคโต ธมฺโม ปีติสหคตสฺส จ สุขสหคตสฺส จ ธมฺมสฺส อธิปติ\-ปจฺจเยน ปจฺจโย. อารมฺมณา\-ธิปติ, สหชาตา\-ธิปติ ฯเปฯ. สหชาตา\-ธิปติ\csromandash{}สุขสหคตา\-ธิปติ สมฺปยุตฺตกานํ ปีติสหคตานญฺจ สุขสหคตานญฺจ ขนฺธานํ อธิปติ\-ปจฺจเยน ปจฺจโย. (4)}

\proseitem{24}{อุเปกฺขา\-สหคโต ธมฺโม อุเปกฺขา\-สหคตสฺส ธมฺมสฺส อธิปติ\-ปจฺจเยน ปจฺจโย. อารมฺมณา\-ธิปติ, สหชาตา\-ธิปติ ฯเปฯ. สหชาตา\-ธิปติ\csromandash{}อุเปกฺขา\-สหคตาธิปติ สมฺปยุตฺตกานํ ขนฺธานํ อธิปติ\-ปจฺจเยน ปจฺจโย. (1)}

\prose{อุเปกฺขา\-สหคโต ธมฺโม ปีติสหคตสฺส ธมฺมสฺส อธิปติ\-ปจฺจเยน ปจฺจโย. อารมฺมณา\-ธิปติ. (สํขิตฺตํ.) (2)}

\prose{อุเปกฺขา\-สหคโต ธมฺโม สุขสหคตสฺส ธมฺมสฺส อธิปติ\-ปจฺจเยน ปจฺจโย. อารมฺมณา\-ธิปติ. (สํขิตฺตํ.) (3)}

\prose{อุเปกฺขา\-สหคโต ธมฺโม ปีติสหคตสฺส จ สุขสหคตสฺส จ ธมฺมสฺส อธิปติ\-ปจฺจเยน ปจฺจโย. อารมฺมณา\-ธิปติ. (สํขิตฺตํ.) (4)}

\proseitem{25}{ปีติสหคโต จ สุขสหคโต จ ธมฺมา ปีติสหคตสฺส ธมฺมสฺส อธิปติ\-ปจฺจเยน ปจฺจโย. อารมฺมณา\-ธิปติ, สหชาตา\-ธิปติ ฯเปฯ. สหชาตา\-ธิปติ\csromandash{}ปีติสหคตา จ สุขสหคตา จ อธิปติ สมฺปยุตฺตกานํ ปีติสหคตานํ ขนฺธานํ อธิปติ\-ปจฺจเยน ปจฺจโย. (1)}

\prose{ปีติสหคโต จ สุขสหคโต จ ธมฺมา สุขสหคตสฺส ธมฺมสฺส อธิปติ\-ปจฺจเยน ปจฺจโย. อารมฺมณา\-ธิปติ, สหชาตา\-ธิปติ ฯเปฯ. สหชาตา\-ธิปติ\csromandash{}ปีติสหคตา จ สุขสหคตา จ อธิปติ สมฺปยุตฺตกานํ สุขสหคตานํ ขนฺธานํ อธิปติ\-ปจฺจเยน ปจฺจโย. (2)}

\prose{ปีติสหคโต จ สุขสหคโต จ ธมฺมฺอา อุเปกฺขา\-สหคตสฺส ธมฺมสฺส อธิปติ\-ปจฺจเยน ปจฺจโย. อารมฺมณา\-ธิปติ. (สํขิตฺตํ.) (3)}

\prose{ปีติสหคโต จ สุขสหคโต จ ธมฺมา ปีติสหคตสฺส จ สุขสหคตสฺส จ ธมฺมสฺส อธิปติ\-ปจฺจเยน ปจฺจโย. อารมฺมณา\-ธิปติ, \csromanfolio{120}สหชาตา\-ธิปติ ฯเปฯ. สหชาตา\-ธิปติ\csromandash{}ปีติสหคตา จ สุขสหคตา จ อธิปติ สมฺปยุตฺตกานํ ปีติสหคตานญฺจ สุขสหคตานญฺจ ขนฺธานํ อธิปติ\-ปจฺจเยน ปจฺจโย. (4)}

\csromanheader{2}{\textbf{อนนฺตร}}
\proseitem{26}{ปีติสหคโต ธมฺโม ปีติสหคตสฺส ธมฺมสฺส อนนฺตร\-ปจฺจเยน ปจฺจโย. ปุริมา ปุริมา ปีติสหคตา ขนฺธา ปจฺฉิมานํ ปจฺฉิมานํ ปีติสหคตานํ ขนฺธานํ อนนฺตร\-ปจฺจเยน ปจฺจโย. ปีติสหคตํ อนุโลมํ โคตฺรภุสฺส อนนฺตร\-ปจฺจเยน ปจฺจโย. (อิมินา การเณน สพฺเพสํ ปทานํ ปจฺจโยติ ทีเปตพฺโพ.) อนุโลมํ โวทานสฺส. โคตฺรภุ มคฺคสฺส. โวทานํ มคฺคสฺส. มคฺโค ผลสฺส. ผลํ ผลสฺส. อนุโลมํ ปีติสหคตาย ผลสมาปตฺติยา อนนฺตร\-ปจฺจเยน ปจฺจโย. (1)}

\prose{ปีติสหคโต ธมฺโม สุขสหคตสฺส ธมฺมสฺส อนนฺตร\-ปจฺจเยน ปจฺจโย. ปุริมา ปุริมา ปีติสหคตา ขนฺธา ปจฺฉิมานํ ปจฺฉิมานํ สุขสหคตานํ ขนฺธานํ อนนฺตร\-ปจฺจเยน ปจฺจโย. ปีติสหคตํ อนุโลมํ สุขสหคตสฺส โคตฺรภุสฺส อนนฺตร\-ปจฺจเยน ปจฺจโย. ปีติสหคตํ อนุโลมํ สุขสหคตสฺส โวทานสฺส อนนฺตร\-ปจฺจเยน ปจฺจโย. ปีติสหคตํ อนุโลมํ สุขสหคตาย ผลสมาปตฺติยา อนนฺตร\-ปจฺจเยน ปจฺจโย. (2)}

\prose{ปีติสหคโต ธมฺโม อุเปกฺขา\-สหคตสฺส ธมฺมสฺส อนนฺตร\-ปจฺจเยน ปจฺจโย. ปีติสหคตํ จุติจิตฺตํ อุเปกฺขา\-สหคตสฺส อุปปตฺติ\-จิตฺตสฺส อนนฺตร\-ปจฺจเยน ปจฺจโย. ปีติสหคตํ ภวงฺคํ อาวชฺชนาย อนนฺตร\-ปจฺจเยน ปจฺจโย. ปีติสหคตา วิปาก\-มโนวิญฺญาณ\-ธาตุ กิริย\-มโนวิญฺญาณ\-ธาตุยา อนนฺตร\-ปจฺจเยน ปจฺจโย. ปีติสหคตํ ภวงฺคํ อุเปกฺขา\-สหคตสฺส ภวงฺคสฺส อนนฺตร\-ปจฺจเยน ปจฺจโย. ปีติสหคตํ กุสลากุสลํ อุเปกฺขา\-สหคตสฺส วุฏฺฐานสฺส. กิริยํ วุฏฺฐานสฺส. ผลํ วุฏฺฐานสฺส อนนฺตร\-ปจฺจเยน ปจฺจโย. (3)}

\prose{ปีติสหคโต ธมฺโม ปีติสหคตสฺส จ สุขสหคตสฺส จ ธมฺมสฺส อนนฺตร\-ปจฺจเยน ปจฺจโย. ปุริมา ปุริมา ปีติสหคตา ขนฺธา ปจฺฉิมานํ ปจฺฉิมานํ ปีติสหคตานญฺจ สุขสหคตานญฺจ ขนฺธานํ อนนฺตร\-ปจฺจเยน \csromanfolio{121}ปจฺจโย. ปีติสหคตํ อนุโลมํ ปีติสหคตสฺส จ สุขสหคตสฺส จ โคตฺรภุสฺส อนนฺตร\-ปจฺจเยน ปจฺจโย ฯเปฯ. ปีติสหคตํ อนุโลมํ ปีติสหคตาย จ สุขสหคตาย จ ผลสมาปตฺติยา อนนฺตร\-ปจฺจเยน ปจฺจโย. (4)}

\proseitem{27}{สุขสหคโต ธมฺโม สุขสหคตสฺส ธมฺมสฺส อนนฺตร\-ปจฺจเยน ปจฺจโย. ปุริมา ปุริมา สุขสหคตา ขนฺธา ปจฺฉิมานํ ปจฺฉิมานํ สุขสหคตานํ ขนฺธานํ อนนฺตร\-ปจฺจเยน ปจฺจโย. สุขสหคตํ อนุโลมํ สุขสหคตสฺส โคตฺรภุสฺส อนนฺตร\-ปจฺจเยน ปจฺจโย ฯเปฯ. สุขสหคตํ อนุโลมํ สุขสหคตาย ผลสมาปตฺติยา อนนฺตร\-ปจฺจเยน ปจฺจโย. (1)}

\prose{สุขสหคโต ธมฺโม ปีติสหคตสฺส ธมฺมสฺส อนนฺตร\-ปจฺจเยน ปจฺจโย. ปุริมา ปุริมา สุขสหคตา ขนฺธา ปจฺฉิมานํ ปจฺฉิมานํ ปีติสหคตานํ ขนฺธานํ อนนฺตร\-ปจฺจเยน ปจฺจโย ฯเปฯ. สุขสหคตํ อนุโลมํ ปีติสหคตาย ผลสมาปตฺติยา อนนฺตร\-ปจฺจเยน ปจฺจโย. (2)}

\prose{สุขสหคโต ธมฺโม อุเปกฺขา\-สหคตสฺส ธมฺมสฺส อนนฺตร\-ปจฺจเยน ปจฺจโย. สุขสหคตํ จุติจิตฺตํ อุเปกฺขา\-สหคตสฺส อุปปตฺติ\-จิตฺตสฺส อนนฺตร\-ปจฺจเยน ปจฺจโย. สุขสหคตํ ภวงฺคํ อาวชฺชนาย อนนฺตร\-ปจฺจเยน ปจฺจโย. สุขสหคตํ กายวิญฺญาณํ วิปาก\-มโนธาตุยา อนนฺตร\-ปจฺจเยน ปจฺจโย. สุขสหคตา วิปาก\-มโนวิญฺญาณ\-ธาตุ กิริย\-มโนวิญฺญาณ\-ธาตุยา อนนฺตร\-ปจฺจเยน ปจฺจโย. สุขสหคตํ ภวงฺคํ อุเปกฺขา\-สหคตสฺส ภวงฺคสฺส อนนฺตร\-ปจฺจเยน ปจฺจโย. สุขสหคตํ กุสลากุสลํ อุเปกฺขา\-สหคตสฺส วุฏฺฐานสฺส. กิริยํ วุฏฺฐานสฺส. ผลํ วุฏฺฐานสฺส อนนฺตร\-ปจฺจเยน ปจฺจโย. (3)}

\prose{สุขสหคโต ธมฺโม ปีติสหคตสฺส จ สุขสหคตสฺส จ ธมฺมสฺส อนนฺตร\-ปจฺจเยน ปจฺจโย. ปุริมา ปุริมา สุขสหคตา ขนฺธา ปจฺฉิมานํ ปจฺฉิมานํ ปีติสหคตานญฺจ สุขสหคตานญฺจ ขนฺธานํ อนนฺตร\-ปจฺจเยน ปจฺจโย ฯเปฯ. สุขสหคตํ อนุโลมํ ปีติสหคตาย จ สุขสหคตาย จ ผลสมาปตฺติยา อนนฺตร\-ปจฺจเยน ปจฺจโย. (4)}

\proseitem{28}{\csromanfolio{122}อุเปกฺขา\-สหคโต ธมฺโม อุเปกฺขา\-สหคตสฺส ธมฺมสฺส อนนฺตร\-ปจฺจเยน ปจฺจโย. ปุริมา ปุริมา อุเปกฺขา\-สหคตา ขนฺธา ปจฺฉิมานํ ปจฺฉิมานํ อุเปกฺขา\-สหคตานํ ขนฺธานํ ฯเปฯ. อาวชฺชนา ปญฺจนฺนํ วิญฺญาณานํ อนนฺตร\-ปจฺจเยน ปจฺจโย. อุเปกฺขา\-สหคตํ อนุโลมํ อุเปกฺขา\-สหคตาย ผลสมาปตฺติยา. นิโรธา วุฏฺฐหนฺตสฺส เนว\-สญฺญา\-นา\-สญฺญา\-ยตนํ อุเปกฺขา\-สหคตาย ผลสมาปตฺติยา อนนฺตร\-ปจฺจเยน ปจฺจโย. (1)}

\prose{อุเปกฺขา\-สหคโต ธมฺโม ปีติสหคตสฺส ธมฺมสฺส อนนฺตร\-ปจฺจเยน ปจฺจโย. อุเปกฺขา\-สหคตํ จุติจิตฺตํ ปีติสหคตสฺส อุปปตฺติ\-จิตฺตสฺส ฯเปฯ อาวชฺชนา ปีติสหคตานํ ขนฺธานํ ฯเปฯ. วิปาก\-มโนธาตุ ปีติสหคตาย วิปาก\-มโนวิญฺญาณ\-ธาตุยา ฯเปฯ. อุเปกฺขา\-สหคตํ ภวงฺคํ ปีติสหคตสฺส ภวงฺคสฺส ฯเปฯ. อุเปกฺขา\-สหคตํ กุสลากุสลํ ปีติสหคตสฺส วุฏฺฐานสฺส. กิริยํ วุฏฺฐานสฺส. ผลํ วุฏฺฐานสฺส. นิโรธา วุฏฺฐหนฺตสฺส เนว\-สญฺญา\-นา\-สญฺญา\-ยตนํ ปีติสหคตาย ผลสมาปตฺติยา อนนฺตร\-ปจฺจเยน ปจฺจโย. (2)}

\prose{อุเปกฺขา\-สหคโต ธมฺโม สุขสหคตสฺส ธมฺมสฺส ฯเปฯ ปีติสหคตสฺส จ สุขสหคตสฺส จ ธมฺมสฺส อนนฺตร\-ปจฺจเยน ปจฺจโย. (4) (ตานิเยว จ คมนานิ นิยาเมตพฺพานิ.)}

\proseitem{29}{ปีติสหคโต จ สุขสหคโต จ ธมฺมา ปีติสหคตสฺส ธมฺมสฺส ฯเปฯ สุขสหคตสฺส ธมฺมสฺส ฯเปฯ อุเปกฺขา\-สหคตสฺส ธมฺมสฺส อนนฺตร\-ปจฺจเยน ปจฺจโย. ปีติสหคตญฺจ สุขสหคตญฺจ จุติจิตฺตํ อุเปกฺขา\-สหคตสฺส อุปปตฺติ\-จิตฺตสฺส ฯเปฯ. ปีติสหคตญฺจ สุขสหคตญฺจ ภวงฺคํ อาวชฺชนาย ฯเปฯ. ปีติสหคตา จ สุขสหคตา จ วิปาก\-มโนวิญฺญาณ\-ธาตุ กิริย\-มโนวิญฺญาณ\-ธาตุยา ฯเปฯ. ปีติสหคตญฺจ สุขสหคตญฺจ ภวงฺคํ อุเปกฺขา\-สหคตสฺส ภวงฺคสฺส ฯเปฯ. ปีติสหคตญฺจ สุขสหคตญฺจ กุสลากุสลํ อุเปกฺขา\-สหคตสฺส วุฏฺฐานสฺส. กิริยํ วุฏฺฐานสฺส. ผลํ วุฏฺฐานสฺส อนนฺตร\-ปจฺจเยน ปจฺจโย. (3)}

\prose{ปีติสหคโต จ สุขสหคโต จ ธมฺมา ปีติสหคตสฺส จ สุขสหคตสฺส จ ธมฺมสฺส อนนฺตร\-ปจฺจเยน ปจฺจโย. ปุริมา ปุริมา \csromanfolio{123}ปีติสหคตา จ สุขสหคตา จ ขนฺธา ปจฺฉิมานํ ปจฺฉิมานํ ปีติสหคตานญฺจ สุขสหคตานญฺจ ขนฺธานํ ฯเปฯ. ปีติสหคตญฺจ สุขสหคตญฺจ อนุโลมํ ปีติสหคตาย จ สุขสหคตาย จ ผลสมาปตฺติยา อนนฺตร\-ปจฺจเยน ปจฺจโย. (4)}

\csromanheader{2}{\textbf{สมนนฺตร}}
\proseitem{30}{ปีติสหคโต ธมฺโม ปีติสหคตสฺส ธมฺมสฺส สมนนฺตร\-ปจฺจเยน ปจฺจโย. (อนนฺตรปจฺจย\-สทิสํ.)}

\csromanheader{2}{\textbf{สหชาต}}
\proseitem{31}{ปีติสหคโต ธมฺโม ปีติสหคตสฺส ธมฺมสฺส สหชาต\-ปจฺจเยน ปจฺจโย. ปีติสหคโต เอโก ขนฺโธ ติณฺณนฺนํ ขนฺธานํ สหชาต\-ปจฺจเยน ปจฺจโย ฯเปฯ เทฺว ขนฺธา ทฺวินฺนํ ขนฺธานํ สหชาต\-ปจฺจเยน ปจฺจโย ฯเปฯ. (ปฏิจฺจสทิสํ, สหชาเต ทส ปญฺหา.)}

\csromanheader{2}{\textbf{อญฺญมญฺญ นิสฺสย}}
\proseitem{32}{ปีติสหคโต ธมฺโม ปีติสหคตสฺส ธมฺมสฺส อญฺญมญฺญ\-ปจฺจเยน ปจฺจโย. นิสฺสย\-ปจฺจเยน ปจฺจโย. (ทส ปญฺหา กาตพฺพา.)}

\csromanheader{2}{\textbf{อุปนิสฺสย}}
\proseitem{33}{ปีติสหคโต ธมฺโม ปีติสหคตสฺส ธมฺมสฺส อุปนิสฺสย\-ปจฺจเยน ปจฺจโย. อารมฺมณู\-ปนิสฺสโย, อนนฺตรู\-ปนิสฺสโย, ปกตู\-ปนิสฺสโย ฯเปฯ. ปกตู\-ปนิสฺสโย\csromandash{}ปีติสหคตํ สทฺธํ อุปนิสฺสาย ปีติสหคเตน จิตฺเตน ทานํ เทติ, สีลํ สมาทิยติ, อุโปสถกมฺมํ กโรติ, ปีติสหคตํ ฌานํ อุปฺปาเทติ, วิปสฺสนํ อุปฺปาเทติ, มคฺคํ อุปฺปาเทติ, สมาปตฺติํ อุปฺปาเทติ, มานํ ชปฺเปติ, ทิฏฺฐิํ คณฺหาติ. ปีติสหคตํ สีลํ. สุตํ. จาคํ. ปญฺญํ อุปนิสฺสาย ปีติสหคเตน จิตฺเตน ทานํ เทติ, สีลํ สมาทิยติ ฯเปฯ มานํ ชปฺเปติ, ทิฏฺฐิํ คณฺหาติ. ปีติสหคตํ ราคํ. โมหํ. มานํ. ทิฏฺฐิํ. ปตฺถนํ อุปนิสฺสาย ปีติสหคเตน จิตฺเตน ทานํ เทติ, สีลํ สมาทิยติ, อุโปสถกมฺมํ กโรติ, ปีติสหคตํ ฌานํ อุปฺปาเทติ ฯเปฯ สมาปตฺติํ อุปฺปาเทติ. ปีติสหคเตน จิตฺเตน อทินฺนํ อาทิยติ, มุสา \csromanfolio{124}ภณติ, ปิสุณํ ภณติ, สมฺผํ ปลปติ, สนฺธิํ ฉินฺทติ, นิลฺโลปํ หรติ, เอกาคาริกํ กโรติ, ปริปนฺเถ ติฏฺฐติ, ปรทารํ คจฺฉติ, คามฆาตํ กโรติ, นิคมฆาตํ กโตติ.  ปีติสหคตา สทฺธา. สีลํ. สุตํ. จาโค. ปญฺญา. ราโค. โมโห. มาโน. ทิฏฺฐิ. ปตฺถนา ปีติสหคตาย สทฺธาย. สีลสฺส. สุตสฺส. จาคสฺส. ปญฺญาย. ราคสฺส. โมหสฺส. มานสฺส. ทิฏฺฐิยา. ปตฺถนาย อุปนิสฺสย\-ปจฺจเยน ปจฺจโย. (1)}

\prose{ปีติสหคโต ธมฺโม สุขสหคตสฺส ธมฺมสฺส อุปนิสฺสย\-ปจฺจเยน ปจฺจโย. อารมฺมณู\-ปนิสฺสโย, อนนฺตรู\-ปนิสฺสโย, ปกตู\-ปนิสฺสโย ฯเปฯ. ปกตู\-ปนิสฺสโย\csromandash{}ปีติสหคตํ สทฺธํ อุปนิสฺสาย สุขสหคเตน จิตฺเตน ทานํ เทติ ฯเปฯ สมาปตฺติํ อุปฺปาเทติ, มานํ ชปฺเปติ, ทิฏฺฐิํ คณฺหาติ. ปีติสหคตํ สีลํ. สุตํ. จาคํ. ปญฺญํ. ราคํ. โมหํ. มานํ. ทิฏฺฐิํ. ปตฺถนํ อุปนิสฺสาย สุขสหคเตน จิตฺเตน ทานํ เทติ ฯเปฯ สมาปตฺติํ อุปฺปาเทติ. สุขสหคเตน จิตฺเตน อทินฺนํ อาทิยติ ฯเปฯ นิคมฆาตํ กโรติ. ปีติสหคตา สทฺธา ฯเปฯ ปตฺถนา สุขสหคตาย สทฺธาย ฯเปฯ ปตฺถนาย สุขสหคตสฺส กายวิญฺญาณสฺส อุปนิสฺสย\-ปจฺจเยน ปจฺจโย. (2)}

\prose{ปีติสหคโต ธมฺโม อุเปกฺขา\-สหคตสฺส ธมฺมสฺส อุปนิสฺสย\-ปจฺจเยน ปจฺจโย. อารมฺมณู\-ปนิสฺสโย, อนนฺตรู\-ปนิสฺสโย, ปกตู\-ปนิสฺสโย ฯเปฯ. ปกตู\-ปนิสฺสโย\csromandash{}ปีติสหคตํ สทฺธํ อุปนิสฺสาย อุเปกฺขา\-สหคเตน จิตฺเตน ทานํ เทติ ฯเปฯ อภิญฺญํ อุปฺปาเทติ, สมาปตฺติํ อุปฺปาเทติ. มานํ ชปฺเปติ, ทิฏฺฐิํ คณฺหาติ. ปีติสหคตํ สีลํ ฯเปฯ. ปตฺถนํ อุปนิสฺสาย อุเปกฺขา\-สหคเตน จิตฺเตน ทานํ เทติ ฯเปฯ นิคมฆาตํ กโรติ. ปีติสหคตา สทฺธา ฯเปฯ ปตฺถนา อุเปกฺขา\-สหคตาย สทฺธาย ฯเปฯ ปตฺถนาย อุปนิสฺสย\-ปจฺจเยน ปจฺจโย. (3)}

\prose{ปีติสหคโต ธมฺโม ปีติสหคตสฺส จ สุขสหคตสฺส จ ธมฺมสฺส อุปนิสฺสย\-ปจฺจเยน ปจฺจโย. อารมฺมณู\-ปนิสฺสโย, อนนฺตรู\-ปนิสฺสโย, ปกตู\-ปนิสฺสโย ฯเปฯ. ปกตู\-ปนิสฺสโย\csromandash{}ปีติสหคตํ สทฺธํ อุปนิสฺสาย ปีติสหคเตน จ สุขสหคเตน จ จิตฺเตน ทานํ เทติ ฯเปฯ ทิฏฺฐิํ คณฺหาติ. ปีติสหคตํ สีลํ ฯเปฯ ปตฺถนํ อุปนิสฺสาย ปีติสหคเตน จ สุขสหคเตน จ จิตฺเตน ทานํ เทติ ฯเปฯ นิคมฆาตํ กโรติ. \csromanfolio{125}ปีติสหคตา สทฺธา ฯเปฯ. ปตฺถนา ปีติสหคตาย จ สุขสหคตาย จ สทฺธาย ฯเปฯ ปตฺถนาย อุปนิสฺสย\-ปจฺจเยน ปจฺจโย. (4)}

\proseitem{34}{สุขสหคโต ธมฺโม สุขสหคตสฺส ธมฺมสฺส อุปนิสฺสย\-ปจฺจเยน ปจฺจโย. อารมฺมณู\-ปนิสฺสโย, อนนฺตรู\-ปนิสฺสโย, ปกตู\-ปนิสฺสโย ฯเปฯ. ปกตู\-ปนิสฺสโย\csromandash{}สุขสหคตํ สทฺธํ อุปนิสฺสาย สุขสหคเตน จิตฺเตน ทานํ เทติ ฯเปฯ ทิฏฺฐิํ คณฺหาติ. สุขสหคตํ สีลํ ฯเปฯ ปตฺถนํ สุขสหคตํ กายวิญฺญาณํ อุปนิสฺสาย สุขสหคเตน จิตฺเตน ทานํ เทติ ฯเปฯ นิคมฆาตํ กโรติ. สุขสหคตา สทฺธา ฯเปฯ ปตฺถนา สุขสหคตํ กายวิญฺญาณํ สุขสหคตาย สทฺธาย ฯเปฯ ปตฺถนาย สุขสหคตสฺส กายวิญฺญาณสฺส อุปนิสฺสย\-ปจฺจเยน ปจฺจโย. (1)}

\prose{สุขสหคโต ธมฺโม ปีติสหคตสฺส ธมฺมสฺส อุปนิสฺสย\-ปจฺจเยน ปจฺจโย. อารมฺมณู\-ปนิสฺสโย, อนนฺตรู\-ปนิสฺสโย, ปกตู\-ปนิสฺสโย ฯเปฯ. ปกตู\-ปนิสฺสโย\csromandash{}สุขสหคตํ สทฺธํ อุปนิสฺสาย ปีติสหคเตน จิตฺเตน ทานํ เทติ ฯเปฯ ทิฏฺฐิํ คณฺหาติ. สุขสหคตํ สีลํ ฯเปฯ ปตฺถนํ สุขสหคตํ กายวิญฺญาณํ อุปนิสฺสาย ปีติสหคเตน จิตฺเตน ทานํ เทติ ฯเปฯ นิคมฆาตํ กโรติ. สุขสหคตา สทฺธา ฯเปฯ ปตฺถนา สุขสหคตํ กายวิญฺญาณํ ปีติสหคตาย สทฺธาย ฯเปฯ ปตฺถนาย อุปนิสฺสย\-ปจฺจเยน ปจฺจโย. (2)}

\prose{สุขสหคโต ธมฺโม อุเปกฺขา\-สหคตสฺส ธมฺมสฺส อุปนิสฺสย\-ปจฺจเยน ปจฺจโย. อารมฺมณู\-ปนิสฺสโย, อนนฺตรู\-ปนิสฺสโย, ปกตู\-ปนิสฺสโย ฯเปฯ. ปกตู\-ปนิสฺสโย\csromandash{}สุขสหคตํ สทฺธํ อุปนิสฺสาย อุเปกฺขา\-สหคเตน จิตฺเตน ทานํ เทติ ฯเปฯ อภิญฺญํ อุปฺปาเทติ ฯเปฯ ทิฏฺฐิํ คณฺหาติ. สุขสหคตํ สีลํ ฯเปฯ ปตฺถนํ สุขสหคตํ กายวิญฺญาณํ อุปนิสฺสาย อุเปกฺขา\-สหคเตน จิตฺเตน ทานํ เทติ ฯเปฯ นิคมฆาตํ กโรติ. สุขสหคตา สทฺธา ฯเปฯ ปตฺถนา สุขสหคตํ กายวิญฺญาณํ อุเปกฺขา\-สหคตาย สทฺธาย ฯเปฯ ปตฺถนาย อุปนิสฺสย\-ปจฺจเยน ปจฺจโย. (3)}

\prose{สุขสหคโต ธมฺโม ปีติสหคตสฺส จ สุขสหคตสฺสจ ธมฺมสฺส อุปนิสฺสย\-ปจฺจเยน ปจฺจโย. อารมฺมณู\-ปนิสฺสโย, อนนฺตรู\-ปนิสฺสโย, \csromanfolio{126}ปกตู\-ปนิสฺสโย ฯเปฯ. ปกตู\-ปนิสฺสโย\csromandash{}สุขสหคตํ สทฺธํ อุปนิสฺสาย ปีติสหคเตน จ สุขสหคเตน จ จิตฺเตน ทานํ เทติ ฯเปฯ ทิฏฺฐิํ คณฺหาติ. สุขสหคตํ สีลํ ฯเปฯ ปตฺถนํ สุขสหคตํ กายวิญฺญาณํ อุปนิสฺสาย ปีติสหคเตน จ สุขสหคเตน จ จิตฺเตน ทานํ เทติ ฯเปฯ นิคมฆาตํ กโรติ. สุขสหคตา สทฺธา ฯเปฯ ปตฺถนา สุขสหคตํ กายวิญฺญาณํ ปีติสหคตาย จ สุขสหคตาย จ สทฺธาย ฯเปฯ ปตฺถนาย อุปนิสฺสย\-ปจฺจเยน ปจฺจโย. (4)}

\proseitem{35}{อุเปกฺขา\-สหคโต ธมฺโม อุเปกฺขา\-สหคตสฺส ธมฺมสฺส อุปนิสฺสย\-ปจฺจเยน ปจฺจโย. อารมฺมณู\-ปนิสฺสโย, อนนฺตรู\-ปนิสฺสโย, ปกตู\-ปนิสฺสโย ฯเปฯ. ปกตู\-ปนิสฺสโย\csromandash{}อุเปกฺขา\-สหคตํ สทฺธํ อุปนิสฺสาย อุเปกฺขา\-สหคเตน จิตฺเตน ทานํ เทติ ฯเปฯ อภิญฺญํ อุปฺปาเทติ ฯเปฯ ทิฏฺฐิํ คณฺหาติ. อุเปกฺขา\-สหคตํ สีลํ ฯเปฯ ปตฺถนํ อุปนิสฺสาย อุเปกฺขา\-สหคเตน จิตฺเตน ทานํ เทติ ฯเปฯ นิคมฆาตํ กโรติ. อุเปกฺขา\-สหคตา สทฺธา ฯเปฯ ปตฺถนา อุเปกฺขา\-สหคตาย สทฺธาย ฯเปฯ ปตฺถนาย อุปนิสฺสย\-ปจฺจเยน ปจฺจโย. (1)}

\prose{อุเปกฺขา\-สหคโต ธมฺโม ปีติสหคตสฺส ธมฺมสฺส อุปนิสฺสย\-ปจฺจเยน ปจฺจโย. อารมฺมณู\-ปนิสฺสโย, อนนฺตรู\-ปนิสฺสโย, ปกตู\-ปนิสฺสโย ฯเปฯ. ปกตู\-ปนิสฺสโย\csromandash{}อุเปกฺขา\-สหคตํ สทฺธํ อุปนิสฺสาย ปีติสหคเตน จิตฺเตน ทานํ เทติ ฯเปฯ ทิฏฺฐิํ คณฺหาติ. อุเปกฺขา\-สหคตํ สีลํ ฯเปฯ ปตฺถนํ อุปนิสฺสาย ปีติสหคเตน จิตฺเตน ทานํ เทติ ฯเปฯ นิคมฆาตํ กโรติ. อุเปกฺขา\-สหคตา สทฺธา ฯเปฯ ปตฺถนา ปีติสหคตาย สทฺธาย ฯเปฯ ปตฺถนาย อุปนิสฺสย\-ปจฺจเยน ปจฺจโย. (2)}

\prose{อุเปกฺขา\-สหคโต ธมฺโม สุขสหคตสฺส ธมฺมสฺส อุปนิสฺสย\-ปจฺจเยน ปจฺจโย. อารมฺมณู\-ปนิสฺสโย, อนนฺตรู\-ปนิสฺสโย, ปกตู\-ปนิสฺสโย ฯเปฯ. ปกตู\-ปนิสฺสโย\csromandash{}อุเปกฺขา\-สหคตํ สทฺธํ อุปนิสฺสาย สุขสหคเตน จิตฺเตน ทานํ เทติ ฯเปฯ ทิฏฺฐิํ คณฺหาติ. อุเปกฺขา\-สหคตํ สีลํ ฯเปฯ ปตฺถนํ อุปนิสฺสาย สุขสหคเตน จิตฺเตน ทานํ เทติ ฯเปฯ นิคมฆาตํ กโรติ. อุเปกฺขา\-สหคตา สทฺธา ฯเปฯ ปตฺถนา สุขสหคตาย สทฺธาย ฯเปฯ ปตฺถนาย สุขสหคตสฺส กายวิญฺญาณสฺส อุปนิสฺสย\-ปจฺจเยน ปจฺจโย. (3)}

\prose{\csromanfolio{127}อุเปกฺขา\-สหคโต ธมฺโม ปีติสหคตสฺส จ สุขสหคตสฺส จ ธมฺมสฺส อุปนิสฺสย\-ปจฺจเยน ปจฺจโย. อารมฺมณู\-ปนิสฺสโย, อนนฺตรู\-ปนิสฺสโย, ปกตู\-ปนิสฺสโย ฯเปฯ. ปกตู\-ปนิสฺสโย\csromandash{} อุเปกฺขา\-สหคตํ สทฺธํ อุปนิสฺสาย ปีติสหคเตน จ สุขสหคเตน จ จิตฺเตน ทานํ เทติ ฯเปฯ ทิฏฺฐิํ คณฺหาติ. อุเปกฺขา\-สหคตํ สีลํ ฯเปฯ ปตฺถนํ อุปนิสฺสาย ปีติสหคเตน จ สุขสหคเตน จ จิตฺเตน ทานํ เทติ ฯเปฯ นิคมฆาตํ กโรติ. อุเปกฺขา\-สหคตา สทฺธา ฯเปฯ ปตฺถนา ปีติสหคตาย จ สุขสหคตาย จ สทฺธาย ฯเปฯ ปตฺถนาย อุปนิสฺสย\-ปจฺจเยน ปจฺจโย. (4)}

\proseitem{36}{ปีติสหคโต จ สุขสหคโต จ ธมฺมา ปีติสหคตสฺส ธมฺมสฺส อุปนิสฺสย\-ปจฺจเยน ปจฺจโย. อารมฺมณู\-ปนิสฺสโย, อนนฺตรู\-ปนิสฺสโย, ปกตู\-ปนิสฺสโย ฯเปฯ. ปกตู\-ปนิสฺสโย\csromandash{}ปีติสหคตญฺจ สุขสหคตญฺจ สทฺธํ อุปนิสฺสาย ปีติสหคเตน จิตฺเตน ทานํ เทติ ฯเปฯ ทิฏฺฐิํ คณฺหาติ. ปีติสหคตญฺจ สุขสหคตญฺจ สีลํ ฯเปฯ ปตฺถนํ อุปนิสฺสาย ปีติสหคเตน จิตฺเตน ทานํ เทติ ฯเปฯ นิคมฆาตํ กโรติ. ปีติสหคตา จ สุขสหคตา จ สทฺธา ฯเปฯ. ปตฺถนา ปีติสหคตาย สทฺธาย ฯเปฯ ปตฺถนาย อุปนิสฺสย\-ปจฺจเยน ปจฺจโย. (1)}

\prose{ปีติสหคโต จ สุขสหคโต จ ธมฺมา สุขสหคตสฺส ธมฺมสฺส อุปนิสฺสย\-ปจฺจเยน ปจฺจโย. อารมฺมณู\-ปนิสฺสโย, อนนฺตรู\-ปนิสฺสโย, ปกตู\-ปนิสฺสโย ฯเปฯ. ปกตู\-ปนิสฺสโย\csromandash{}ปีติสหคตญฺจ สุขสหคตญฺจ สทฺธํ อุปนิสฺสาย สุขสหคเตน จิตฺเตน ทานํ เทติ ฯเปฯ ทิฏฺฐิํ คณฺหาติ. ปีติสหคตญฺจ สุขสหคตญฺจ สีลํ ฯเปฯ ปตฺถนํ อุปนิสฺสาย สุขสหคเตน จิตฺเตน ทานํ เทติ ฯเปฯ นิคมฆาตํ กโรติ. ปีติสหคตา จ สุขสหคตา จ สทฺธา ฯเปฯ ปตฺถนา สุขสหคตาย สทฺธาย ฯเปฯ ปตฺถนาย สุขสหคตสฺส กายวิญฺญาณสฺส อุปนิสฺสย\-ปจฺจเยน ปจฺจโย. (2)}

\prose{ปีติสหคโต จ สุขสหคโต จ ธมฺมา อุเปกฺขา\-สหคตสฺส ธมฺมสฺส อุปนิสฺสย\-ปจฺจเยน ปจฺจโย. อารมฺมณู\-ปนิสฺสโย, อนนฺตรู\-ปนิสฺสโย, ปกตู\-ปนิสฺสโย ฯเปฯ. ปกตู\-ปนิสฺสโย\csromandash{}ปีติสหคตญฺจ สุขสหคตญฺจ สทฺธํ อุปนิสฺสาย อุเปกฺขา\-สหคเตน จิตฺเตน ทานํ \csromanfolio{128}เทติ ฯเปฯ อภิญฺญํ อุปฺปาเทติ ฯเปฯ ทิฏฺฐิํ คณฺหาติ. ปีติสหคตญฺจ สุขสหคตญฺจ สีลํ ฯเปฯ ปตฺถนํ อุปนิสฺสาย อุเปกฺขา\-สหคเตน จิตฺเตน ทานํ เทติ ฯเปฯ นิคมฆาตํ กโรติ. ปีติสหคตา จ สุขสหคตา จ สทฺธา ฯเปฯ ปตฺถนา อุเปกฺขา\-สหคตาย สทฺธาย ฯเปฯ ปตฺถนาย อุปนิสฺสย\-ปจฺจเยน ปจฺจโย. (3)}

\prose{ปีติสหคโต จ สุขสหคโต จ ธมฺมา ปีติสหคตสฺส จ สุขสหคตสฺส จ ธมฺมสฺส อุปนิสฺสย\-ปจฺจเยน ปจฺจโย. อารมฺมณู\-ปนิสฺสโย, อนนฺตรู\-ปนิสฺสโย, ปกตู\-ปนิสฺสโย ฯเปฯ. ปกตู\-ปนิสฺสโย\csromandash{}ปีติสหคตญฺจ สุขสหคตญฺจ สทฺธํ อุปนิสฺสาย ปีติสหคเตน จ สุขสหคเตน จ จิตฺเตน ทานํ เทติ, สีลํ สมาทิยติ, อุโปสถกมฺมํ กโรติ. ปีติสหคตญฺจ สุขสหคตญฺจ ฌานํ อุปฺปาเทติ, วิปสฺสนํ อุปฺปาเทติ, มคฺคํ อุปฺปาเทติ, สมาปตฺติํ อุปฺปาเทติ, มานํ ชปฺเปติ, ทิฏฺฐิํ คณฺหาติ. ปีติสหคตญฺจ สุขสหคตญฺจ สีลํ. สุตํ. จาคํ. ปญฺญํ. ราคํ. โมหํ. มานํ. ทิฏฺฐิํ. ปตฺถนํ อุปนิสฺสาย ปีติสหคเตน จ สุขสหคเตน จ จิตฺเตน ทานํ เทติ, สีลํ สมาทิยติ, อุโปสถกมฺมํ กโรติ. ปีติสหคตญฺจ สุขสหคตญฺจ ฌานํ อุปฺปาเทติ ฯเปฯ สมาปตฺติํ อุปฺปาเทติ. ปีติสหคเตน จ สุขสหคเตน จ จิตฺเตน อทินฺนํ อาทิยติ, มุสา ภณติ, ปิสุณํ ภณติ, สมฺผํ ปลปติ, สนฺธิํ ฉินฺทติ, นิลฺโลปํ หรติ, เอกาคาริกํ กโรติ, ปริปนฺเถ ติฏฺฐติ, ปรทานํ คจฺฉติ, คามฆาตํ กโรติ, นิคมฆาตํ กโรติ. ปีติสหคตา จ สุขสหคตา จ สทฺธา ฯเปฯ ปตฺถนา ปีติสหคตาย จ สุขสหคตาย จ สทฺธาย ฯเปฯ ปตฺถนาย อุปนิสฺสย\-ปจฺจเยน ปจฺจโย. (4)}

\csromanheader{2}{\textbf{อาเสวน}}
\proseitem{37}{ปีติสหคโต ธมฺโม ปีติสหคตสฺส ธมฺมสฺส อาเสวน\-ปจฺจเยน ปจฺจโย. ปุริมา ปุริมา ปีติสหคตา ขนฺธา ปจฺฉิมานํ ปจฺฉิมานํ ปีติสหคตานํ ขนฺธานํ อาเสวน\-ปจฺจเยน ปจฺจโย. ปีติสหคตํ อนุโลมํ ปีติสหคตสฺส โคตฺรภุสฺส. อนุโลมํ โวทานสฺส. โคตฺรภุ มคฺคสฺส. โวทานํ มคฺคสฺส อาเสวน\-ปจฺจเยน ปจฺจโย. (1)}

\prose{ปีติสหคโต ธมฺโม สุขสหคตสฺส ธมฺมสฺส อาเสวน\-ปจฺจเยน ปจฺจโย. ปุริมา ปุริมา ปีติสหคตา ขนฺธา ปจฺฉิมานํ ปจฺฉิมานํ \csromanfolio{129}สุขสหคตานํ ขนฺธานํ อาเสวน\-ปจฺจเยน ปจฺจโย. ปีติสหคตํ อนุโลมํ สุขสหคตสฺส โคตฺรภุสฺส อาเสวน\-ปจฺจเยน ปจฺจโย. ปีติสหคตํ อนุโลมํ สุขสหคตสฺส โวทานสฺส อาเสวน\-ปจฺจเยน ปจฺจโย. ปีติสหคตํ โคตฺรภุ สุขสหคตสฺส มคฺคสฺส. ปีติสหคตํ โวทานํ สุขสหคตสฺส มคฺคสฺส อาเสวน\-ปจฺจเยน ปจฺจโย. (2)}

\prose{ปีติสหคโต ธมฺโม ปีติสหคตสฺส จ สุขสหคตสฺส จ ธมฺมสฺส อาเสวน\-ปจฺจเยน ปจฺจโย. ปุริมา ปุริมา ปีติสหคตา ขนฺธา ปจฺฉิมานํ ปจฺฉิมานํ ปีติสหคตานญฺจ สุขสหคตานญฺจ ขนฺธานํ อาเสวน\-ปจฺจเยน ปจฺจโย ฯเปฯ. ปีติสหคตํ โวทานํ ปีติสหคตสฺส จ สุขสหคตสฺส จ มคฺคสฺส อาเสวน\-ปจฺจเยน ปจฺจโย. (3)}

\proseitem{38}{สุขสหคโต ธมฺโม สุขสหคตสฺส ธมฺมสฺส ฯเปฯ ปีติสหคตสฺส ธมฺมสฺส ฯเปฯ ปีติสหคตสฺส จ สุขสหคตสฺส จ ธมฺมสฺส อาเสวน\-ปจฺจเยน ปจฺจโย. (สํขิตฺตํ.) (ปีตินยํ ปสฺสิตฺวา กาตพฺพํ.)}

\prose{อุเปกฺขา\-สหคโต ธมฺโม อุเปกฺขา\-สหคตสฺส ธมฺมสฺส อาเสวน\-ปจฺจเยน ปจฺจโย. ปุริมา ปุริมา อุเปกฺขา\-สหคตา ขนฺธา ปจฺฉิมานํ ปจฺฉิมานํ อุเปกฺขา\-สหคตานํ ขนฺธานํ ฯเปฯ. อุเปกฺขา\-สหคตํ โวทานํ อุเปกฺขา\-สหคตสฺส มคฺคสฺส อาเสวน\-ปจฺจเยน ปจฺจโย. (1)}

\proseitem{39}{ปีติสหคโต จ สุขสหคโต จ ธมฺมา ปีติสหคตสฺส ธมฺมสฺส ฯเปฯ สุขสหคตสฺส ธมฺมสฺส ฯเปฯ ปีติสหคตสฺส จ สุขสหคตสฺส จ ธมฺมสฺส อาเสวน\-ปจฺจเยน ปจฺจโย. ปุริมา ปุริมา ปีติสหคตา จ สุขสหคตา จ ขนฺธา ปจฺฉิมานํ ปจฺฉิมานํ ปีติสหคตานญฺจ สุขสหคตานญฺจ ขนฺธานํ อาเสวน\-ปจฺจเยน ปจฺจโย ฯเปฯ. ปีติสหคตญฺจ สุขสหคตญฺจ โวทานํ ปีติสหคตสฺส จ สุขสหคตสฺส จ มคฺคสฺส อาเสวน\-ปจฺจเยน ปจฺจโย. (3)}

\csromanheader{2}{\textbf{กมฺม}}
\proseitem{40}{ปีติสหคโต ธมฺโม ปีติสหคตสฺส ธมฺมสฺส กมฺมปจฺจเยน ปจฺจโย. สหชาตา, นานากฺขณิกา\csromansharedfootnote{da31f9fb8bf64194}{นานาขณิกา (สฺยา, ก)}. สหชาตา ปีติสหคตา เจตนา สมฺปยุตฺตกานํ ขนฺธานํ กมฺมปจฺจเยน ปจฺจโย. ปฏิสนฺธิ\-กฺขเณ \csromanfolio{130}ปีติสหคตา เจตนา สมฺปยุตฺตกานํ ขนฺธานํ กมฺมปจฺจเยน ปจฺจโย.  นานากฺขณิกา ปีติสหคตา เจตนา วิปากานํ ปีติสหคตานํ ขนฺธานํ กมฺมปจฺจเยน ปจฺจโย. (1)}

\prose{ปีติสหคโต ธมฺโม สุขสหคตสฺส ธมฺมสฺส กมฺมปจฺจเยน ปจฺจโย. สหชาตา, นานากฺขณิกา.  สหชาตา ปีติสหคตา เจตนา สมฺปยุตฺตกานํ สุขสหคตานํ ขนฺธานํ กมฺมปจฺจเยน ปจฺจโย. ปฏิสนฺธิ\-กฺขเณ ปีติสหคตา เจตนา สมฺปยุตฺตกานํ สุขสหคตานํ ขนฺธานํ กมฺมปจฺจเยน ปจฺจโย.  นานากฺขณิกา ปีติสหคตา เจตนา วิปากานํ สุขสหคตานํ ขนฺธานํ กมฺมปจฺจเยน ปจฺจโย. (2)}

\prose{ปีติสหคโต ธมฺโม อุเปกฺขา\-สหคตสฺส ธมฺมสฺส กมฺมปจฺจเยน ปจฺจโย. นานากฺขณิกา ปีติสหคตา เจตนา วิปากานํ อุเปกฺขา\-สหคตานํ ขนฺธานํ กมฺมปจฺจเยน ปจฺจโย. (3)}

\prose{ปีติสหคโต ธมฺโม ปีติสหคตสฺส จ สุขสหคตสฺส จ ธมฺมสฺส กมฺมปจฺจเยน ปจฺจโย. สหชาตา, นานากฺขณิกา.  สหชาตา ปีติสหคตา เจตนา สมฺปยุตฺตกานํ ปีติสหคตานญฺจ สุขสหคตานญฺจ ขนฺธานํ กมฺมปจฺจเยน ปจฺจโย. ปฏิสนฺธิ\-กฺขเณ ฯเปฯ. นานากฺขณิกา ปีติสหคตา เจตนา วิปากานํ ปีติสหคตานญฺจ สุขสหคตานญฺจ ขนฺธานํ กมฺมปจฺจเยน ปจฺจโย. (4)}

\proseitem{41}{สุขสหคโต ธมฺโม สุขสหคตสฺส ธมฺมสฺส. (จตฺตาริปิ คณนานิ ปสฺสิตฺวา กาตพฺพานิ.)}

\proseitem{42}{อุเปกฺขา\-สหคโต ธมฺโม อุเปกฺขา\-สหคตสฺส ธมฺมสฺส กมฺมปจฺจเยน ปจฺจโย. สหชาตา, นานากฺขณิกา ฯเปฯ.}

\prose{อุเปกฺขา\-สหคโต ธมฺโม ปีติสหคตสฺส ธมฺมสฺส กมฺมปจฺจเยน ปจฺจโย. นานากฺขณิกา อุเปกฺขา\-สหคตา เจตนา ฯเปฯ.}

\prose{อุเปกฺขา\-สหคโต ธมฺโม สุขสหคตสฺส ธมฺมสฺส กมฺมปจฺจเยน ปจฺจโย. นานากฺขณิกา อุเปกฺขา\-สหคตา เจตนา ฯเปฯ.}

\prose{\csromanfolio{131}อุเปกฺขา\-สหคโต ธมฺโม ปีติสหคตสฺส จ สุขสหคตสฺส จ ธมฺมสฺส กมฺมปจฺจเยน ปจฺจโย. นานากฺขณิกา อุเปกฺขา\-สหคตา เจตนา ฯเปฯ. (4)}

\prose{ปีติสหคโต จ สุขสหคโต จ ธมฺมา ปีติสหคตสฺส ธมฺมสฺส. (จตฺตาริ กาตพฺพานิ, ปีติสหคตํ อนุมชฺชนฺเตน วิภชิตพฺพํ.) (4)}

\csromanheader{2}{\textbf{วิปาก}}
\proseitem{43}{ปีติสหคโต ธมฺโม ปีติสหคตสฺส ธมฺมสฺส วิปาก\-ปจฺจเยน ปจฺจโย. ปีติสหคโต วิปาโก เอโก ขนฺโธ ติณฺณนฺนํ ขนฺธานํ วิปาก\-ปจฺจเยน ปจฺจโย ฯเปฯ เทฺว ขนฺธา ทฺวินฺนํ ขนฺธานํ ฯเปฯ. ปฏิสนฺธิ\-กฺขเณ ปีติสหคโต เอโก ขนฺโธ ติณฺณนฺนํ ขนฺธานํ ฯเปฯ เทฺว ขนฺธา ทฺวินฺนํ ขนฺธานํ ฯเปฯ.}

\prose{(ยถา ปฏิจฺจวาเร เหตุปจฺจเย, เอวํ วิตฺถาเรตพฺพา ทส ปญฺหา.)}

\csromanheader{2}{\textbf{อาหาราทิ}}
\proseitem{44}{ปีติสหคโต ธมฺโม ปีติสหคตสฺส ธมฺมสฺส อาหารปจฺจเยน ปจฺจโย. อินฺทฺริย\-ปจฺจเยน ปจฺจโย. ฌานปจฺจเยน ปจฺจโย. มคฺคปจฺจเยน ปจฺจโย. สมฺปยุตฺต\-ปจฺจเยน ปจฺจโย. อตฺถิปจฺจเยน ปจฺจโย. (ทส ปญฺหา วิตฺถาเรตพฺพา.) นตฺถิ\-ปจฺจเยน ปจฺจโย. วิคต\-ปจฺจเยน ปจฺจโย. (นตฺถิปิ วิคตมฺปิ อนนฺตร\-สทิสํ.) อวิคต\-ปจฺจเยน ปจฺจโย.}

\csromanmarkh{1. ปจฺจยานุโลม}\csromanmarkhii{2. สงฺขฺยาวาร}\csromanheadernomark{2}{1. \textbf{ปจฺจยานุโลม 2. สงฺขฺยาวาร}}
\proseitem{45}{เหตุยา ทส, อารมฺมเณ โสฬส, อธิปติยา โสฬส, อนนฺตเร โสฬส, สมนนฺตเร โสฬส, สหชาเต ทส, อญฺญมญฺเญ ทส, นิสฺสเย ทส, อุปนิสฺสเย โสฬส, อาเสวเน ทส, กมฺเม โสฬส, วิปาเก ทส, อาหาเร อินฺทฺริเย ฌาเน มคฺเค สมฺปยุตฺเต อตฺถิยา ทส, นตฺถิยา โสฬส, วิคเต โสฬส, อวิคเต ทส.}

\vspace{\tipitakaparskip}
\csromancenter{(กุสลตฺติกํ อนุโลมํ อนุมชฺชนฺเตน คเณตพฺพํ.)}
\nitthitamruled{\textbf{อนุโลมํ.}}
\csromanheader{2}{2. ปจฺจนียุ\-ทฺธาร}
\proseitem{46}{\csromanfolio{132}ปีติสหคโต ธมฺโม ปีติสหคตสฺส ธมฺมสฺส อารมฺมณ\-ปจฺจเยน ปจฺจโย. สหชาต\-ปจฺจเยน ปจฺจโย. อุปนิสฺสย\-ปจฺจเยน ปจฺจโย. กมฺมปจฺจเยน ปจฺจโย. (1)}

\prose{ปีติสหคโต ธมฺโม สุขสหคตสฺส ธมฺมสฺส อารมฺมณ\-ปจฺจเยน ปจฺจโย. สหชาต\-ปจฺจเยน ปจฺจโย. อุปนิสฺสย\-ปจฺจเยน ปจฺจโย. กมฺมปจฺจเยน ปจฺจโย. (2)}

\prose{ปีติสหคโต ธมฺโม อุเปกฺขา\-สหคตสฺส ธมฺมสฺส อารมฺมณ\-ปจฺจเยน ปจฺจโย. อุปนิสฺสย\-ปจฺจเยน ปจฺจโย. กมฺมปจฺจเยน ปจฺจโย. (3)}

\prose{ปีติสหคโต ธมฺโม ปีติสหคตสฺส จ สุขสหคตสฺส จ ธมฺมสฺส อารมฺมณ\-ปจฺจเยน ปจฺจโย. สหชาต\-ปจฺจเยน ปจฺจโย. อุปนิสฺสย\-ปจฺจเยน ปจฺจโย. กมฺมปจฺจเยน ปจฺจโย. (4)}

\proseitem{47}{สุขสหคโต ธมฺโม สุขสหคตสฺส ธมฺมสฺส อารมฺมณ\-ปจฺจเยน ปจฺจโย. สหชาต\-ปจฺจเยน ปจฺจโย. อุปนิสฺสย\-ปจฺจเยน ปจฺจโย. กมฺมปจฺจเยน ปจฺจโย. (1)}

\prose{สุขสหคโต ธมฺโม ปีติสหคตสฺส ธมฺมสฺส อารมฺมณ\-ปจฺจเยน ปจฺจโย. สหชาต\-ปจฺจเยน ปจฺจโย. อุปนิสฺสย\-ปจฺจเยน ปจฺจโย. กมฺมปจฺจเยน ปจฺจโย. (2)}

\prose{สุขสหคโต ธมฺโม อุเปกฺขา\-สหคตสฺส ธมฺมสฺส อารมฺมณ\-ปจฺจเยน ปจฺจโย. อุปนิสฺสย\-ปจฺจเยน ปจฺจโย. กมฺมปจฺจเยน ปจฺจโย. (3)}

\prose{สุขสหคโต ธมฺโม ปีติสหคตสฺส จ สุขสหคตสฺส จ ธมฺมสฺส อารมฺมณ\-ปจฺจเยน ปจฺจโย. สหชาต\-ปจฺจเยน ปจฺจโย. อุปนิสฺสย\-ปจฺจเยน ปจฺจโย. กมฺมปจฺจเยน ปจฺจโย. (4)}

\proseitem{48}{อุเปกฺขา\-สหคโต ธมฺโม อุเปกฺขา\-สหคตสฺส ธมฺมสฺส อารมฺมณ\-ปจฺจเยน ปจฺจโย. สหชาต\-ปจฺจเยน ปจฺจโย. อุปนิสฺสย\-ปจฺจเยน ปจฺจโย. กมฺมปจฺจเยน ปจฺจโย. (1)}

\prose{อุเปกฺขา\-สหคโต ธมฺโม ปีติสหคตสฺส ธมฺมสฺส อารมฺมณ\-ปจฺจเยน ปจฺจโย. อุปนิสฺสย\-ปจฺจเยน ปจฺจโย. กมฺมปจฺจเยน ปจฺจโย. (2)}

\prose{\csromanfolio{133}อุเปกฺขา\-สหคโต ธมฺโม สุขสหคตสฺส ธมฺมสฺส อารมฺมณ\-ปจฺจเยน ปจฺจโย. อุปนิสฺสย\-ปจฺจเยน ปจฺจโย. กมฺมปจฺจเยน ปจฺจโย. (3)}

\prose{อุเปกฺขา\-สหคโต ธมฺโม ปีติสหคตสฺส จ สุขสหคตสฺส จ ธมฺมสฺส อารมฺมณ\-ปจฺจเยน ปจฺจโย. อุปนิสฺสย\-ปจฺจเยน ปจฺจโย. กมฺมปจฺจเยน ปจฺจโย. (4)}

\proseitem{49}{ปีติสหคโต จ สุขสหคโต จ ธมฺมา ปีติสหคตสฺส ธมฺมสฺส อารมฺมณ\-ปจฺจเยน ปจฺจโย. สหชาต\-ปจฺจเยน ปจฺจโย. อุปนิสฺสย\-ปจฺจเยน ปจฺจโย. กมฺมปจฺจเยน ปจฺจโย. (1)}

\prose{ปีติสหคโต จ สุขสหคโต จ ธมฺมา สุขสหคตสฺส ธมฺมสฺส อารมฺมณ\-ปจฺจเยน ปจฺจโย. สหชาต\-ปจฺจเยน ปจฺจโย. อุปนิสฺสย\-ปจฺจเยน ปจฺจโย. กมฺมปจฺจเยน ปจฺจโย. (2)}

\prose{ปีติสหคโต จ สุขสหคโต จ ธมฺมา อุเปกฺขา\-สหคตสฺส ธมฺมสฺส อารมฺมณ\-ปจฺจเยน ปจฺจเยน. อุปนิสฺสย\-ปจฺจเยน ปจฺจโย. กมฺมปจฺจเยน ปจฺจโย. (3)}

\prose{ปีติสหคโต จ สุขสหคโต จ ธมฺมา ปีติสหคตสฺส จ สุขสหคตสฺส จ ธมฺมสฺส อารมฺมณ\-ปจฺจเยน ปจฺจโย. สหชาต\-ปจฺจเยน ปจฺจโย. อุปนิสฺสย\-ปจฺจเยน ปจฺจโย. กมฺมปจฺจเยน ปจฺจโย. (4)}

\csromanmarkh{2. ปจฺจย\-ปจฺจนีย}\csromanmarkhii{2. สงฺขฺยาวาร}\csromanheadernomark{2}{2. ปจฺจย\-ปจฺจนีย 2. สงฺขฺยาวาร}
\proseitemwithcloserruled{50}{นเหตุยา โสฬส, นอารมฺมเณ นอธิปติยา นอนนฺตเร นสมนนฺตเร นสหชาเต นอญฺญมญฺเญ นนิสฺสเย นอุปนิสฺสเย นปุเรชาเต นปจฺฉาชาเต นอาเสวเน นกมฺเม นวิปาเก นอาหาเร นอินฺทฺริเย นฌาเน นมคฺเค นสมฺปยุตฺเต นวิปฺปยุตฺเต โนอตฺถิยา โนนตฺถิยา โนวิคเต โนอวิคเต สพฺพตฺถ โสฬส. (ปจฺจนียํ อนุมชฺชนฺเตน คเณตพฺพํ.)}{ปจฺจนียํ.}
\proseitem{51}{\csromanfolio{134}เหตุปจฺจยา นอารมฺมเณ ทส, นอธิปติยา ทส, นอนนฺตเร นสมนนฺตเร นอุปนิสฺสเย นปุเรชาเต นปจฺฉาชาเต นอาเสวเน นกมฺเม นวิปาเก นอาหาเร นอินฺทฺริเย นฌาเน นมคฺเค นวิปฺปยุตฺเต โนนตฺถิยา โนวิคเต สพฺพตฺถ ทส. (อนุโลมปจฺจนียํ อนุมชฺชนฺเตน คเณตพฺพํ.)}

\vspace{\tipitakaparskip}
\csromancenter{(อนุโลม\-ปจฺจนียํ อนุมชฺชนฺเตน คเณตพฺพํ.)}
\nitthitamruled{อนุโลม\-ปจฺจนียํ.}
\proseitem{52}{นเหตุปจฺจยา อารมฺมเณ โสฬส, อธิปติยา อนนฺตเร สมนนฺตเร โสฬส, สหชาเต ทส, อญฺญมญฺเญ ทส, นิสฺสเย ทส, อุปนิสฺสเย โสฬส, อาเสวเน ทส, กมฺเม โสฬส, วิปาเก ทส, อาหาเร ทส, อินฺทฺริเย ทส, ฌาเน ทส, มคฺเค ทส, สมฺปยุตฺเต ทส, อตฺถิยา ทส, นตฺถิยา โสฬส, วิคเต โสฬส, อวิคเต ทส. (ปจฺจนียานุโลมํ อนุมชฺชนฺเตน คเณตพฺพํ.)}

\vspace{\tipitakaparskip}
\csromancenter{(ปจฺจนียา\-นุโลมํ อนุมชฺชนฺเตน คเณตพฺพํ.)}
\csromancenter{ปจฺจนียา\-นุโลมํ.}
\nitthitam{ปีติตฺติกํ นิฏฺฐิตํ.}
\csromancenter{8. ทสฺสเนน\-ปหาตพฺพ\-ตฺติก 1. ปฏิจฺจวาร}
\csromanmarkh{1. ปจฺจยานุโลม}\csromanmarkhii{1. วิภงฺควาร}\csromanheadernomark{2}{1. \textbf{ปจฺจยานุโลม 1. วิภงฺควาร}}
\csromanheader{2}{\textbf{เหตุ}}
\proseitem{1}{\csromanfolio{135}ทสฺสเนน ปหาตพฺพํ ธมฺมํ ปฏิจฺจ ทสฺสเนน ปหาตพฺโพ ธมฺโม อุปฺปชฺชติ เหตุปจฺจยา. ทสฺสเนน ปหาตพฺพํ เอกํ ขนฺธํ ปฏิจฺจ ตโย ขนฺธา ฯเปฯ เทฺว ขนฺเธ ปฏิจฺจ เทฺว ขนฺธา. (1)}

\prose{ทสฺสเนน ปหาตพฺพํ ธมฺมํ ปฏิจฺจ เนวทสฺสเนน นภาวนาย ปหาตพฺโพ ธมฺโม อุปฺปชฺชติ เหตุปจฺจยา. ทสฺสเนน ปหาตพฺเพ ขนฺเธ ปฏิจฺจ จิตฺต\-สมุฏฺฐานํ รูปํ. (2)}

\prose{ทสฺสเนน ปหาตพฺพํ ธมฺมํ ปฏิจฺจ ทสฺสเนน ปหาตพฺโพ จ เนวทสฺสเนน นภาวนาย ปหาตพฺโพ จ ธมฺมา อุปฺปชฺชนฺติ เหตุปจฺจยา. ทสฺสเนน ปหาตพฺพํ เอกํ ขนฺธํ ปฏิจฺจ ตโย ขนฺธา จิตฺต\-สมุฏฺฐานญฺจ รูปํ ฯเปฯ เทฺว ขนฺเธ ปฏิจฺจ เทฺว ขนฺธา จิตฺต\-สมุฏฺฐานญฺจ รูปํ. (3)}

\proseitem{2}{ภาวนาย ปหาตพฺพํ ธมฺมํ ปฏิจฺจ ภาวนาย ปหาตพฺโพ ธมฺโม อุปฺปชฺชติ เหตุปจฺจยา. ภาวนาย ปหาตพฺพํ เอกํ ขนฺธํ ปฏิจฺจ ตโย ขนฺธา ฯเปฯ เทฺว ขนฺเธ ปฏิจฺจ เทฺว ขนฺธา. (1)}

\prose{ภาวนาย ปหาตพฺพํ ธมฺมํ ปฏิจฺจ เนวทสฺสเนน นภาวนาย ปหาตพฺโพ ธมฺโม อุปฺปชฺชติ เหตุปจฺจยา. ภาวนาย ปหาตพฺเพ ขนฺเธ ปฏิจฺจ จิตฺต\-สมุฏฺฐานํ รูปํ. (2)}

\prose{ภาวนาย ปหาตพฺพํ ธมฺมํ ปฏิจฺจ ภาวนาย ปหาตพฺโพ จ เนวทสฺสเนน นภาวนาย ปหาตพฺโพ จ ธมฺมา อุปฺปชฺชนฺติ เหตุปจฺจยา. ภาวนาย ปหาตพฺพํ เอกํ ขนฺธํ ปฏิจฺจ ตโย ขนฺธา จิตฺต\-สมุฏฺฐานญฺจ รูปํ ฯเปฯ เทฺว ขนฺเธ ปฏิจฺจ เทฺว ขนฺธา จิตฺต\-สมุฏฺฐานญฺจ รูปํ. (3)}

\proseitem{3}{เนวทสฺสเนน นภาวนาย ปหาตพฺพํ ธมฺมํ ปฏิจฺจ เนวทสฺสเนน นภาวนาย ปหาตพฺโพ ธมฺโม อุปฺปชฺชติ เหตุปจฺจยา. เนวทสฺสเนน นภาวนาย ปหาตพฺพํ เอกํ ขนฺธํ ปฏิจฺจ ตโย ขนฺธา จิตฺต\-สมุฏฺฐานญฺจ รูปํ ฯเปฯ. \csromanfolio{136}เทฺว ขนฺเธ ปฏิจฺจ เทฺว ขนฺธา จิตฺต\-สมุฏฺฐานญฺจ รูปํ. ปฏิสนฺธิ\-กฺขเณ เนวทสฺสเนน นภาวนาย ปหาตพฺพํ เอกํ ขนฺธํ ปฏิจฺจ ตโย ขนฺธา กฏตฺตา จ รูปํ ฯเปฯ เทฺว ขนฺเธ ปฏิจฺจ เทฺว ขนฺธา กฏตฺตา จ รูปํ. ขนฺเธ ปฏิจฺจ วตฺถุ, วตฺถุํ ปฏิจฺจ ขนฺธา. เอกํ มหาภูตํ ปฏิจฺจ ตโย มหาภูตา ฯเปฯ เทฺว มหาภูเต ปฏิจฺจ เทฺว มหาภูตา. มหาภูเต ปฏิจฺจ จิตฺต\-สมุฏฺฐานํ รูปํ, กฏตฺตารูปํ, อุปาทารูปํ. (1)}

\proseitem{4}{ทสฺสเนน ปหาตพฺพญฺจ เนวทสฺสเนน นภาวนาย ปหาตพฺพญฺจ ธมฺมํ ปฏิจฺจ เนวทสฺสเนน นภาวนาย ปหาตพฺโพ ธมฺโม อุปฺปชฺชติ เหตุปจฺจยา. ทสฺสเนน ปหาตพฺเพ ขนฺเธ จ มหาภูเต จ ปฏิจฺจ จิตฺต\-สมุฏฺฐานํ รูปํ. (1)}

\prose{ภาวนาย ปหาตพฺพญฺจ เนวทสฺสเนน นภาวนาย ปหาตพฺพญฺจ ธมฺมํ ปฏิจฺจ เนวทสฺสเนน นภาวนาย ปหาตพฺโพ ธมฺโม อุปฺปชฺชติ เหตุปจฺจยา. ภาวนาย ปหาตพฺเพ ขนฺเธ จ มหาภูเต จ ปฏิจฺจ จิตฺต\-สมุฏฺฐานํ รูปํ. (1)}

\csromanheader{2}{\textbf{อารมฺมณ}}
\proseitem{5}{ทสฺสเนน ปหาตพฺพํ ธมฺมํ ปฏิจฺจ ทสฺสเนน ปหาตพฺโพ ธมฺโม อุปฺปชฺชติ อารมฺมณ\-ปจฺจยา. ทสฺสเนน ปหาตพฺพํ เอกํ ขนฺธํ ปฏิจฺจ ตโย ขนฺธา ฯเปฯ เทฺว ขนฺเธ ปฏิจฺจ เทฺว ขนฺธา. (1)}

\prose{ภาวนาย ปหาตพฺพํ ธมฺมํ ปฏิจฺจ ภาวนาย ปหาตพฺโพ ธมฺโม อุปฺปชฺชติ อารมฺมณ\-ปจฺจยา. ภาวนาย ปหาตพฺพํ เอกํ ขนฺธํ ปฏิจฺจ ตโย ขนฺธา ฯเปฯ เทฺว ขนฺเธ ปฏิจฺจ เทฺว ขนฺธา. (1)}

\prose{เนวทสฺสเนน นภาวนาย ปหาตพฺพํ ธมฺมํ ปฏิจฺจ เนวทสฺสเนน นภาวนาย ปหาตพฺโพ ธมฺโม อุปฺปชฺชติ อารมฺมณ\-ปจฺจยา. เนวทสฺสเนน นภาวนาย ปหาตพฺพํ เอกํ ขนฺธํ ปฏิจฺจ ตโย ขนฺธา ฯเปฯ เทฺว ขนฺเธ ปฏิจฺจ เทฺว ขนฺธา. ปฏิสนฺธิ\-กฺขเณ เนวทสฺสเนน นภาวนาย ปหาตพฺพํ เอกํ ขนฺธํ ปฏิจฺจ ตโย ขนฺธา ฯเปฯ เทฺว ขนฺเธ ปฏิจฺจ เทฺว ขนฺธา. วตฺถุํ ปฏิจฺจ ขนฺธา. (1)}

\proseitem{6}{ทสฺสเนน ปหาตพฺพํ ธมฺมํ ปฏิจฺจ ทสฺสเนน ปหาตพฺโพ ธมฺโม อุปฺปชฺชติ อธิปติ\-ปจฺจยา. ตีณิ.}

\prose{\csromanfolio{137}ภาวนาย ปหาตพฺพํ ธมฺมํ ปฏิจฺจ ภาวนาย ปหาตพฺโพ ธมฺโม อุปฺปชฺชติ อธิปติ\-ปจฺจยา. ตีณิ.}

\prose{เนวทสฺสเนน นภาวนาย ปหาตพฺพํ ธมฺมํ ปฏิจฺจ เนวทสฺสเนน นภาวนาย ปหาตพฺโพ ธมฺโม อุปฺปชฺชติ อธิปติ\-ปจฺจยา. เนวทสฺสเนน นภาวนาย ปหาตพฺพํ เอกํ ขนฺธํ ปฏิจฺจ ตโย ขนฺธา จิตฺต\-สมุฏฺฐานญฺจ รูปํ ฯเปฯ เทฺว ขนฺเธ ปฏิจฺจ เทฺว ขนฺธา จิตฺต\-สมุฏฺฐานญฺจ รูปํ. เอกํ มหาภูตํ ปฏิจฺจ ตโย มหาภูตา ฯเปฯ เทฺว มหาภูเต ปฏิจฺจ เทฺว มหาภูตา. มหาภูเต ปฏิจฺจ จิตฺต\-สมุฏฺฐานํ รูปํ, อุปาทารูปํ. (1)}

\proseitem{7}{ทสฺสเนน ปหาตพฺพญฺจ เนวทสฺสเนน นภาวนาย ปหาตพฺพญฺจ ธมฺมํ ปฏิจฺจ เนวทสฺสเนน นภาวนาย ปหาตพฺโพ ธมฺโม อุปฺปชฺชติ อธิปติ\-ปจฺจยา. ทสฺสเนน ปหาตพฺเพ ขนฺเธ จ มหาภูเต จ ปฏิจฺจ จิตฺต\-สมุฏฺฐานํ รูปํ. (1)}

\prose{ภาวนาย ปหาตพฺพญฺจ เนวทสฺสเนน นภาวนาย ปหาตพฺพญฺจ ธมฺมํ ปฏิจฺจ เนวทสฺสเนน นภาวนาย ปหาตพฺโพ ธมฺโม อุปฺปชฺชติ อธิปติ\-ปจฺจยา. ภาวนาย ปหาตพฺเพ ขนฺเธ จ มหาภูเต จ ปฏิจฺจ จิตฺต\-สมุฏฺฐานํ รูปํ. (1)}

\csromanheader{2}{อนนฺตร\-สมนนฺตร}
\proseitem{8}{ทสฺสเนน ปหาตพฺพํ ธมฺมํ ปฏิจฺจ ทสฺสเนน ปหาตพฺโพ ธมฺโม อุปฺปชฺชติ อนนฺตร\-ปจฺจยา. สมนนฺตร\-ปจฺจยา. (อารมฺมณ\-สทิสํ.)}

\csromanheader{2}{\textbf{สหชาต}}
\proseitem{9}{ทสฺสเนน ปหาตพฺพํ ธมฺมํ ปฏิจฺจ ทสฺสเนน ปหาตพฺโพ ธมฺโม อุปฺปชฺชติ สหชาต\-ปจฺจยา. ตีณิ.}

\prose{ภาวนาย ปหาตพฺพํ ธมฺมํ ปฏิจฺจ ภาวนาย ปหาตพฺโพ ธมฺโม อุปฺปชฺชติ สหชาต\-ปจฺจยา. ตีณิ.}

\prose{เนวทสฺสเนน นภาวนาย ปหาตพฺพํ ธมฺมํ ปฏิจฺจ เนวทสฺสเนน นภาวนาย ปหาตพฺโพ ธมฺโม อุปฺปชฺชติ สหชาต\-ปจฺจยา ฯเปฯ. เอกํ มหาภูตํ ปฏิจฺจ ฯเปฯ. พาหิรํ. อาหารสมุฏฺฐานํ. อุตุสมุฏฺฐานํ. อสญฺญสตฺตานํ เอกํ มหาภูตํ ฯเปฯ. (1)}

\proseitem{10}{\csromanfolio{138}ทสฺสเนน ปหาตพฺพญฺจ เนวทสฺสเนน นภาวนาย ปหาตพฺพญฺจ ธมฺมํ ปฏิจฺจ เนวทสฺสเนน นภาวนาย ปหาตพฺโพ ธมฺโม อุปฺปชฺชติ สหชาต\-ปจฺจยา. ทสฺสเนน ปหาตพฺเพ ขนฺเธ จ มหาภูเต จ ปฏิจฺจ จิตฺต\-สมุฏฺฐานํ รูปํ. (1)}

\prose{ภาวนาย ปหาตพฺพญฺจ เนวทสฺสเนน นภาวนาย ปหาตพฺพญฺจ ธมฺมํ ปฏิจฺจ เนวทสฺสเนน นภาวนาย ปหาตพฺโพ ธมฺโม อุปฺปชฺชติ สหชาต\-ปจฺจยา. ภาวนาย ปหาตพฺเพ ขนฺเธ จ มหาภูเต จ ปฏิจฺจ จิตฺต\-สมุฏฺฐานํ รูปํ. (1)}

\csromanheader{2}{\textbf{อญฺญมญฺญ}}
\proseitem{11}{ทสฺสเนน ปหาตพฺพํ ธมฺมํ ปฏิจฺจ ทสฺสเนน ปหาตพฺโพ ธมฺโม อุปฺปชฺชติ อญฺญมญฺญ\-ปจฺจยา. เอกํ.}

\prose{ภาวนาย ปหาตพฺพํ ธมฺมํ ปฏิจฺจ ภาวนาย ปหาตพฺโพ ธมฺโม อุปฺปชฺชติ สหชาต\-ปจฺจยา. เอกํ.}

\prose{เนวทสฺสเนน นภาวนาย ปหาตพฺพํ ธมฺมํ ปฏิจฺจ เนวทสฺสเนน นภาวนาย ปหาตพฺโพ ธมฺโม อุปฺปชฺชติ อญฺญมญฺญ\-ปจฺจยา. เนวทสฺสเนน นภาวนาย ปหาตพฺพํ เอกํ ขนฺธํ ปฏิจฺจ ตโย ขนฺธา ฯเปฯ เทฺว ขนฺเธ ปฏิจฺจ เทฺว ขนฺธา. ปฏิสนฺธิ\-กฺขเณ เนวทสฺสเนน นภาวนาย ปหาตพฺพํ เอกํ ขนฺธํ ปฏิจฺจ ตโย ขนฺธา วตฺถุ จ ฯเปฯ เทฺว ขนฺเธ ปฏิจฺจ เทฺว ขนฺธา วตฺถุ จ. ขนฺเธ ปฏิจฺจ วตฺถุ, วตฺถุํ ปฏิจฺจ ขนฺธา. เอกํ มหาภูตํ ปฏิจฺจ ตโย มหาภูตา ฯเปฯ เทฺว มหาภูเต ปฏิจฺจ เทฺว มหาภูตา. พาหิรํ. อาหารสมุฏฺฐานํ. อุตุสมุฏฺฐานํ. อสญฺญสตฺตานํ ฯเปฯ เทฺว มหาภูเต ปฏิจฺจ เทฺว มหาภูตา.}

\csromanheader{2}{\textbf{นิสฺสยาทิ}}
\proseitem{12}{ทสฺสเนน ปหาตพฺพํ ธมฺมํ ปฏิจฺจ ทสฺสเนน ปหาตพฺโพ ธมฺโม อุปฺปชฺชติ นิสฺสยปจฺจยา. (เหตุปจฺจย\-สทิสํ.) อุปนิสฺสย\-ปจฺจยา. ตีณิ. ปุเรชาต\-ปจฺจยา. ตีณิ. (ปฏิสนฺธิ นตฺถิ.) อาเสวนปจฺจยา. (วิปาก\-ปฏิสนฺธิ นตฺถิ.) กมฺมปจฺจยา (ปริปุณฺณํ. อชฺฌตฺติกา จ อสญฺญสตฺตานญฺจ มหาภูตา.)}

\csromanheader{2}{\textbf{วิปาก}}
\proseitem{13}{เนวทสฺสเนน นภาวนาย ปหาตพฺพํ ธมฺมํ ปฏิจฺจ เนวทสฺสเนน นภาวนาย ปหาตพฺโพ ธมฺโม อุปฺปชฺชติ วิปากปจฺจยา. วิปากํ เอกํ ขนฺธํ \csromanfolio{139}ปฏิจฺจ ตโย ขนฺธา จิตฺต\-สมุฏฺฐานญฺจ รูปํ ฯเปฯ. ปฏิสนฺธิ\-กฺขเณ ฯเปฯ ขนฺเธ ปฏิจฺจ วตฺถุ, วตฺถุํ ปฏิจฺจ ขนฺธา. เอกํ มหาภูตํ ฯเปฯ มหาภูเต ปฏิจฺจ จิตฺต\-สมุฏฺฐานํ รูปํ, กฏตฺตารูปํ, อุปาทารูปํ. (1)}

\csromanheader{2}{\textbf{อาหาราทิ}}
\proseitem{14}{ทสฺสเนน ปหาตพฺพํ ธมฺมํ ปฏิจฺจ ทสฺสเนน ปหาตพฺโพ ธมฺโม อุปฺปชฺชติ อาหารปจฺจยา. (ปริปุณฺณํ. อชฺฌตฺติกา มหาภูตา จ อาหารสมุฏฺฐานญฺจ) อินฺทฺริยปจฺจยา. (กมฺม\-ปจฺจย\-สทิสํ.) ฌานปจฺจยา. มคฺคปจฺจยา. (เหตุปจฺจย\-สทิสํ.) สมฺปยุตฺต\-ปจฺจยา. (อารมฺมณปจฺจย\-สทิสํ.) วิปฺปยุตฺต\-ปจฺจยา. (กุสลตฺติเก วิปฺปยุตฺตปจฺจย\-สทิสํ.) อตฺถิปจฺจยา. (สหชาตปจฺจย\-สทิสํ.) นตฺถิปจฺจยา. วิคตปจฺจยา. อวิคตปจฺจยา.}

\csromanmarkh{1. ปจฺจยานุโลม}\csromanmarkhii{2. สงฺขฺยาวาร}\csromanheadernomark{2}{1. \textbf{ปจฺจยานุโลม 2. สงฺขฺยาวาร}}
\proseitem{13}{เหตุยา นว, อารมฺมเณ ตีณิ, อธิปติยา นว, อนนฺตเร ตีณิ, สมนนฺตเร ตีณิ, สหชาเต นว, อญฺญมญฺเญ ตีณิ, นิสฺสเย นว, อุปนิสฺสเย ตีณิ, ปุเรชาเต ตีณิ, อาเสวเน ตีณิ, กมฺเม นว, วิปาเก เอกํ, อาหาเร นว, อินฺทฺริเย นว, ฌาเน นว, มคฺเค นว, สมฺปยุตฺเต ตีณิ, วิปฺปยุตฺเต นว, อตฺถิยา นว, นตฺถิยา ตีณิ, วิคเต ตีณิ, อวิคเต นว.}

\vspace{\tipitakaparskip}
\csromancenter{(อิมานิ ปทานิ อนุมชฺชนฺเตน อนุโลมํ คเณตพฺพํ.)}
\nitthitamruled{อนุโลมํ.}
\csromanmarkh{2. ปจฺจย\-ปจฺจนีย}\csromanmarkhii{1. วิภงฺควาร}\csromanheadernomark{2}{2. ปจฺจย\-ปจฺจนีย 1. วิภงฺควาร}
\csromanheader{2}{\textbf{นเหตุ}}
\proseitem{13}{ทสฺสเนน ปหาตพฺพํ ธมฺมํ ปฏิจฺจ ทสฺสเนน ปหาตพฺโพ ธมฺโม อุปฺปชฺชติ นเหตุปจฺจยา. วิจิกิจฺฉา\-สหคเต ขนฺเธ ปฏิจฺจ วิจิกิจฺฉา\-สหคโต โมโห. (1)}

\proseitem{16}{\csromanfolio{140}ภาวนาย ปหาตพฺพํ ธมฺมํ ปฏิจฺจ ภาวนาย ปหาตพฺโพ ธมฺโม อุปฺปชฺชติ นเหตุปจฺจยา. อุทฺธจฺจ\-สหคเต ขนฺเธ ปฏิจฺจ อุทฺธจฺจ\-สหคโต โมโห. (1)}

\prose{เนวทสฺสเนน นภาวนาย ปหาตพฺพํ ธมฺมํ ปฏิจฺจ เนวทสฺสเนน นภาวนาย ปหาตพฺโพ ธมฺโม อุปฺปชฺชติ นเหตุปจฺจยา. อเหตุกํ เนวทสฺสเนน นภาวนาย ปหาตพฺพํ เอกํ ขนฺธํ ปฏิจฺจ ตโย ขนฺธา จิตฺต\-สมุฏฺฐานญฺจ รูปํ ฯเปฯ เทฺว ขนฺเธ ปฏิจฺจ ฯเปฯ. อเหตุก\-ปฏิสนฺธิ\-กฺขเณ เนวทสฺสเนน นภาวนาย ปหาตพฺพํ เอกํ ขนฺธํ ปฏิจฺจ ตโย ขนฺธา กฏตฺตา จ รูปํ ฯเปฯ เทฺว ขนฺเธ ปฏิจฺจ เทฺว ขนฺธา ฯเปฯ. ขนฺเธ ปฏิจฺจ วตฺถุ, วตฺถุํ ปฏิจฺจ ขนฺธา. เอกํ มหาภูตํ ปฏิจฺจ ตโย มหาภูตา ฯเปฯ มหาภูเต ปฏิจฺจ จิตฺต\-สมุฏฺฐานํ รูปํ, กฏตฺตารูปํ, อุปาทารูปํ. พาหิรํ. อาหารสมุฏฺฐานํ. อุตุสมุฏฺฐานํ. อสญฺญสตฺตานํ เอกํ มหาภูตํ ปฏิจฺจ ตโย มหาภูตา ฯเปฯ มหาภูเต ปฏิจฺจ กฏตฺตารูปํ อุปาทารูปํ. (1)}

\prose{นอารมฺมณ}

\proseitem{17}{ทสฺสเนน ปหาตพฺพํ ธมฺมํ ปฏิจฺจ เนวทสฺสเนน นภาวนาย ปหาตพฺโพ ธมฺโม อุปฺปชฺชติ นอารมฺมณ\-ปจฺจยา. ทสฺสเนน ปหาตพฺเพ ขนฺเธ ปฏิจฺจ จิตฺต\-สมุฏฺฐานํ รูปํ. (1)}

\prose{ภาวนาย ปหาตพฺพํ ธมฺมํ ปฏิจฺจ เนวทสฺสเนน นภาวนาย ปหาตพฺโพ ธมฺโม อุปฺปชฺชติ นอารมฺมณ\-ปจฺจยา. ภาวนาย ปหาตพฺเพ ขนฺเธ ปฏิจฺจ จิตฺต\-สมุฏฺฐานํ รูปํ. (1)}

\prose{เนวทสฺสเนน นภาวนาย ปหาตพฺพํ ธมฺมํ ปฏิจฺจ เนวทสฺสเนน นภาวนาย ปหาตพฺโพ ธมฺโม อุปฺปชฺชติ นอารมฺมณ\-ปจฺจยา. เนวทสฺสเนน นภาวนาย ปหาตพฺเพ ขนฺเธ ปฏิจฺจ จิตฺต\-สมุฏฺฐานํ รูปํ. ปฏิสนฺธิ\-กฺขเณ เนวทสฺสเนน นภาวนาย ปหาตพฺเพ ขนฺเธ ปฏิจฺจ กฏตฺตารูปํ. ขนฺเธ ปฏิจฺจ วตฺถุ. เอกํ มหาภูตํ ฯเปฯ. พาหิรํ. อาหารสมุฏฺฐานํ. อุตุสมุฏฺฐานํ อสญฺญสตฺตานํ เอกํ มหาภูตํ ฯเปฯ. (1)}

\proseitem{18}{ทสฺสเนน ปหาตพฺพญฺจ เนวทสฺสเนน นภาวนาย ปหาตพฺพญฺจ ธมฺมํ ปฏิจฺจ เนวทสฺสเนน นภาวนาย ปหาตพฺโพ ธมฺโม อุปฺปชฺชติ}

\prose{\csromanfolio{141}นอารมฺมณ\-ปจฺจยา. ทสฺสเนน ปหาตพฺเพ ขนฺเธ จ มหาภูเต จ ปฏิจฺจ จิตฺต\-สมุฏฺฐานํ รูปํ. (1)}

\prose{ภาวนาย ปหาตพฺพญฺจ เนวทสฺสเนน นภาวนาย ปหาตพฺพญฺจ ธมฺมํ ปฏิจฺจ เนวทสฺสเนน นภาวนาย ปหาตพฺโพ ธมฺโม อุปฺปชฺชติ นอารมฺมณ\-ปจฺจยา. ภาวนาย ปหาตพฺเพ ขนฺเธ จ มหาภูเต จ ปฏิจฺจ จิตฺต\-สมุฏฺฐานํ รูปํ. (1)}

\prose{นอธิปตฺยาทิ}

\proseitem{18}{ทสฺสเนน ปหาตพฺพํ ธมฺมํ ปฏิจฺจ ทสฺสเนน ปหาตพฺโพ ธมฺโม อุปฺปชฺชติ นอธิปติ\-ปจฺจยา. (ปริปุณฺณํ เหตุปจฺจย\-สทิสํ.) นอนนฺตร\-ปจฺจยา. นสมนนฺตร\-ปจฺจยา. นอญฺญมญฺญ\-ปจฺจยา. นอุปนิสฺสยปจฺจยา.}

\csromanheader{2}{\textbf{นปุเรชาต}}
\proseitem{18}{ทสฺสเนน ปหาตพฺพํ ธมฺมํ ปฏิจฺจ ทสฺสเนน ปหาตพฺโพ ธมฺโม อุปฺปชฺชติ นปุเรชาต\-ปจฺจยา. อรูเป ทสฺสเนน ปหาตพฺพํ เอกํ ขนฺธํ ปฏิจฺจ ตโย ขนฺธา ฯเปฯ. (1)}

\prose{ทสฺสเนน ปหาตพฺพํ ธมฺมํ ปฏิจฺจ เนวทสฺสเนน นภาวนาย ปหาตพฺโพ ธมฺโม อุปฺปชฺชติ นปุเรชาต\-ปจฺจยา. ทสฺสเนน ปหาตพฺเพ ขนฺเธ ปฏิจฺจ จิตฺต\-สมุฏฺฐานํ รูปํ. (2)}

\prose{ภาวนาย ปหาตพฺพํ ธมฺมํ ปฏิจฺจ ภาวนาย ปหาตพฺโพ ฯเปฯ. อรูเป ภาวนาย ปหาตพฺพํ เอกํ ขนฺธํ ปฏิจฺจ ตโย ขนฺธา ฯเปฯ เทฺว ขนฺเธ ปฏิจฺจ เทฺว ขนฺธา. (1)}

\prose{ภาวนาย ปหาตพฺพํ ธมฺมํ ปฏิจฺจ เนวทสฺสเนน นภาวนาย ปหาตพฺโพ ฯเปฯ. ภาวนาย ปหาตพฺเพ ขนฺเธ ปฏิจฺจ จิตฺต\-สมุฏฺฐานํ รูปํ. (2)}

\prose{เนวทสฺสเนน นภาวนาย ปหาตพฺพํ ธมฺมํ ปฏิจฺจ เนวทสฺสเนน นภาวนาย ปหาตพฺโพ ธมฺโม อุปฺปชฺชติ นปุเรชาต\-ปจฺจยา. อรูเป เนวทสฺสเนน นภาวนาย ปหาตพฺพํ เอกํ ขนฺธํ ปฏิจฺจ ตโย ขนฺธา ฯเปฯ เทฺว ขนฺเธ ปฏิจฺจ เทฺว ขนฺธา. เนวทสฺสเนน นภาวนาย ปหาตพฺเพ ขนฺเธ ปฏิจฺจ จิตฺต\-สมุฏฺฐานํ รูปํ. ปฏิสนฺธิ\-กฺขเณ เนวทสฺสเนน นภาวนาย ปหาตพฺพํ \csromanfolio{142}เอกํ ขนฺธํ ปฏิจฺจ ตโย ขนฺธา กฏตฺตา จ รูปํ ฯเปฯ เทฺว ขนฺเธ ปฏิจฺจ เทฺว ขนฺธา กฏตฺตา จ รูปํ. ขนฺเธ ปฏิจฺจ วตฺถุ, วตฺถุํ ปฏิจฺจ ขนฺธา. เอกํ มหาภูตํ ปฏิจฺจ ตโย มหาภูตา ฯเปฯ. พาหิรํ. อาหารสมุฏฺฐานํ. อุตุสมุฏฺฐานํ. อสญฺญสตฺตานํ เอกํ มหาภูตํ ฯเปฯ. (1)}

\proseitem{21}{ทสฺสเนน ปหาตพฺพญฺจ เนวทสฺสเนน นภาวนาย ปหาตพฺพญฺจ ธมฺมํ ปฏิจฺจ เนวทสฺสเนน นภาวนาย ปหาตพฺโพ ธมฺโม ฯเปฯ. ทสฺสเนน ปหาตพฺเพ ขนฺเธ จ มหาภูเต จ ปฏิจฺจ จิตฺต\-สมุฏฺฐานํ รูปํ. (1)}

\prose{ภาวนาย ปหาตพฺพญฺจ เนวทสฺสเนน นภาวนาย ปหาตพฺพญฺจ ธมฺมํ ปฏิจฺจ เนวทสฺสเนน นภาวนาย ปหาตพฺโพ ธมฺโม อุปฺปชฺชติ นปุเรชาต\-ปจฺจยา. ภาวนาย ปหาตพฺเพ ขนฺเธ จ มหาภูเต จ ปฏิจฺจ จิตฺต\-สมุฏฺฐานํ รูปํ. (1)}

\csromanheader{2}{\textbf{นปจฺฉาชาตาทิ}}
\proseitem{21}{ทสฺสเนน ปหาตพฺพํ ธมฺมํ ปฏิจฺจ ทสฺสเนน ปหาตพฺโพ ธมฺโม อุปฺปชฺชติ นปจฺฉาชาต\-ปจฺจยา. นอาเสวน\-ปจฺจยา.}

\csromanheader{2}{\textbf{นกมฺม}}
\proseitem{23}{ทสฺสเนน ปหาตพฺพํ ธมฺมํ ปฏิจฺจ ทสฺสเนน ปหาตพฺโพ ธมฺโม อุปฺปชฺชติ นกมฺมปจฺจยา. ทสฺสเนน ปหาตพฺเพ ขนฺเธ ปฏิจฺจ ทสฺสเนน ปหาตพฺพา เจตนา. (1)}

\prose{ภาวนาย ปหาตพฺพํ ธมฺมํ ปฏิจฺจ ภาวนาย ปหาตพฺโพ ธมฺโม อุปฺปชฺชติ นกมฺมปจฺจยา. ภาวนาย ปหาตพฺเพ ขนฺเธ ปฏิจฺจ ภาวนาย ปหาตพฺพา เจตนา. (1)}

\prose{เนวทสฺสเนน นภาวนาย ปหาตพฺพํ ธมฺมํ ปฏิจฺจ เนวทสฺสเนน นภาวนาย ปหาตพฺโพ ธมฺโม อุปฺปชฺชติ นกมฺมปจฺจยา. เนวทสฺสเนน นภาวนาย ปหาตพฺเพ ขนฺเธ ปฏิจฺจ เนวทสฺสเนน นภาวนาย ปหาตพฺพา เจตนา. พาหิรํ. อาหารสมุฏฺฐานํ. อุตุสมุฏฺฐานํ. อสญฺญสตฺตานํ เอกํ มหาภูตํ ฯเปฯ. (1)}

\csromanheader{2}{8. ทสฺสเนน\-ปหาตพฺพ\-ตฺติก นวิปาก}
\proseitem{24}{\csromanfolio{143}ทสฺสเนน ปหาตพฺพํ ธมฺมํ ปฏิจฺจ ทสฺสเนน ปหาตพฺโพ ธมฺโม อุปฺปชฺชติ นวิปาก\-ปจฺจยา. (นอธิปติปจฺจยสทิสํ, ปฏิสนฺธิ นตฺถิ.)}

\vspace{\tipitakaparskip}
\csromancenter{(นอธิปติ\-ปจฺจย\-สทิสํ, ปฏิสนฺธิ นตฺถิ.)}
\prose{นอาหาร}

\proseitem{25}{เนวทสฺสเนน นภาวนาย ปหาตพฺพํ ธมฺมํ ปฏิจฺจ เนวทสฺสเนน นภาวนาย ปหาตพฺโพ ธมฺโม อุปฺปชฺชติ นอาหารปจฺจยา. พาหิรํ. อุตุสมุฏฺฐานํ. อสญฺญสตฺตานํ เอกํ มหาภูตํ ฯเปฯ.}

\prose{นอินฺทฺริย}

\proseitem{26}{เนวทสฺสเนน นภาวนาย ปหาตพฺพํ ฯเปฯ นอินฺทฺริยปจฺจยา. พาหิรํ. อาหารสมุฏฺฐานํ. อุตุสมุฏฺฐานํ เอกํ มหาภูตํ ฯเปฯ. อสญฺญสตฺตานํ มหาภูเต ปฏิจฺจ รูปชีวิตินฺทฺริยํ.}

\csromanheader{2}{\textbf{นฌาน}}
\proseitem{27}{เนวทสฺสเนน นภาวนาย ปหาตพฺพํ ธมฺมํ ฯเปฯ นฌานปจฺจยา, ปญฺจ\-วิญฺญาณสหคตํ เอกํ ขนฺธํ ปฏิจฺจ ตโย ขนฺธา ฯเปฯ เทฺว ขนฺเธ ฯเปฯ. พาหิรํ, อาหารสมุฏฺฐานํ. อุตุสมุฏฺฐานํ. อสญฺญสตฺตานํ เอกํ มหาภูตํ ฯเปฯ.}

\csromanheader{2}{\textbf{นมคฺค}}
\proseitem{28}{เนวทสฺสเนน นภาวนาย ปหาตพฺพํ ฯเปฯ นมคฺคปจฺจยา. อเหตุกํ เนวทสฺสเนน นภาวนาย ปหาตพฺพํ ฯเปฯ. อเหตุก\-ปฏิสนฺธิ\-กฺขเณ ฯเปฯ. เอกํ มหาภูตํ ฯเปฯ. พาหิรํ. อาหารสมุฏฺฐานํ. อุตุสมุฏฺฐานํ. อสญฺญสตฺตานํ เอกํ มหาภูตํ ฯเปฯ.}

\csromanheader{2}{\textbf{นสมฺปยุตฺต}}
\proseitem{29}{ทสฺสเนน ปหาตพฺพํ ธมฺมํ ปฏิจฺจ เนวทสฺสเนน นภาวนาย ปหาตพฺโพ ธมฺโม อุปฺปชฺชติ นสมฺปยุตฺต\-ปจฺจยา. (นอารมฺมณ\-สทิสํ.)}

\csromanheader{2}{\textbf{นวิปฺปยุตฺต}}
\proseitem{30}{ทสฺสเนน ปหาตพฺพํ ฯเปฯ นวิปฺปยุตฺต\-ปจฺจยา. อรูเป ทสฺสเนน ปหาตพฺพํ เอกํ ขนฺธํ ปฏิจฺจ ฯเปฯ. (1)}

\prose{\csromanfolio{144}ภาวนาย ปหาตพฺพํ ฯเปฯ นวิปฺปยุตฺต\-ปจฺจยา. อรูเป ภาวนาย ปหาตพฺพํ เอกํ ขนฺธํ ฯเปฯ. (1)}

\prose{เนวทสฺสเนน นภาวนาย ปหาตพฺพํ ฯเปฯ นวิปฺปยุตฺต\-ปจฺจยา. อรูเป เนวทสฺสเนน นภาวนาย ปหาตพฺพํ เอกํ ขนฺธํ ปฏิจฺจ ฯเปฯ. พาหิรํ. อาหารสมุฏฺฐานํ. อุตุสมุฏฺฐานํ. อสญฺญสตฺตานํ ฯเปฯ. (1)}

\csromanheader{2}{\textbf{โนนตฺถิ โนวิคต}}
\proseitem{31}{ทสฺสเนน ปหาตพฺพํ ธมฺมํ ปฏิจฺจ เนวทสฺสเนน นภาวนาย ฯเปฯ โนนตฺถิปจฺจยา. โนวิคต\-ปจฺจยา. (นอารมฺมณ\-สทิสํ.)}

\csromanmarkh{2. ปจฺจย\-ปจฺจนีย}\csromanmarkhii{2. สงฺขฺยาวาร}\csromanheadernomark{2}{2. ปจฺจย\-ปจฺจนีย 2. สงฺขฺยาวาร}
\proseitemwithcloserruled{32}{นเหตุยา ตีณิ, นอารมฺมเณ ปญฺจ, นอธิปติยา นว, นอนนฺตเร ปญฺจ, นสมนนฺตเร ปญฺจ, นอญฺญมญฺเญ ปญฺจ, นอุปนิสฺสเย ปญฺจ, นปุเรชาเต สตฺต, นปจฺฉาชาเต นว, นอาเสวเน นว, นกมฺเม ตีณิ, นวิปาเก นว, นอาหาเร เอกํ, นอินฺทฺริเย เอกํ, นฌาเน เอกํ, นมคฺเค เอกํ, นสมฺปยุตฺเต ปญฺจ, นวิปฺปยุตฺเต ตีณิ, โนนตฺถิยา ปญฺจ, โนวิคเต ปญฺจ. (ญตฺวา คเณตพฺพํ.)}{ปจฺจนียํ.}
\proseitemwithcloserruled{33}{เหตุปจฺจยา นอารมฺมเณ ปญฺจ, นอธิปติยา นว, นอนนฺตเร ปญฺจ, นสมนนฺตเร ปญฺจ, นอญฺญมญฺเญ ปญฺจ, นอุปนิสฺสเย ปญฺจ, นปุเรชาเต สตฺต, นปจฺฉาชาเต นว, นอาเสวเน นว, นกมฺเม ตีณิ, นวิปาเก นว, นสมฺปยุตฺเต ปญฺจ, นวิปฺปยุตฺเต ตีณิ, โนนตฺถิยา ปญฺจ, โนวิคเต ปญฺจ. (เอวํ อนุมชฺชนฺเตน คเณตพฺพํ.)}{อนุโลม\-ปจฺจนียํ.}
\csromanmatikaanchor{mtk.4}
\csromanmatikaanchor{mtk.5}
\csromantocmarkref{section}{เหตุทุก}{mtk.4}
\bookmark[dest={mtk.4},level=1]{เหตุทุก}
\csromantocmarknopage{chapter}{4. ปจฺจยปจฺจนียานุโลม}{mtk.5}
\bookmark[dest={mtk.5},level=0]{4. ปจฺจยปจฺจนียานุโลม}
\chapterheadpageread{8. ทสฺสเนน\-ปหาตพฺพ\-ตฺติก 4. ปจฺจย\-ปจฺจนียา\-นุโลม}
\proseitem{34}{\csromanfolio{145}นเหตุปจฺจยา อารมฺมเณ ตีณิ, อนนฺตเร ตีณิ, สมนนฺตเร ตีณิ, สหชาเต ตีณิ, อญฺญมญฺเญ ตีณิ, นิสฺสเย ตีณิ, อุปนิสฺสเย ตีณิ, ปุเรชาเต ตีณิ, อาเสวเน ตีณิ, กมฺเม ตีณิ, วิปาเก เอกํ, อาหาเร ตีณิ, อินฺทฺริเย ตีณิ, ฌาเน ตีณิ, มคฺเค เทฺว, สมฺปยุตฺเต ตีณิ, วิปฺปยุตฺเต ตีณิ, อตฺถิยา ตีณิ, นตฺถิยา ตีณิ, วิคเต ตีณิ, อวิคเต ตีณิ. (เอวํ อนุมชฺชนฺเตน คเณตพฺพํ.) ปจฺจนียา\-นุโลมํ. ปฏิจฺจวาโร.}
\csromansectionrule

\csromanmarkh{8. ทสฺสเนน\-ปหาตพฺพ\-ตฺติก}\csromanmarkhii{2. สหชาตวาร}\csromanheadernomark{2}{8. ทสฺสเนน\-ปหาตพฺพ\-ตฺติก 2. สหชาตวาร}
\csromanheader{2}{1-4. ปจฺจยา\-นุโลมาทิ}
\proseitem{35}{ทสฺสเนน ปหาตพฺพํ ธมฺมํ สหชาโต ทสฺสเนน ปหาตพฺโพ ธมฺโม อุปฺปชฺชติ เหตุปจฺจยา. ทสฺสเนน ปหาตพฺพํ เอกํ ขนฺธํ สหชาตา ตโย ขนฺธา ฯเปฯ เทฺว ขนฺเธ สหชาตา ฯเปฯ. (สหชาตวาโร ปฏิจฺจ\-วาร\-สทิโส.)}
\csromansectionrule

\csromanmarkh{8. ทสฺสเนน\-ปหาตพฺพ\-ตฺติก}\csromanmarkhii{3. ปจฺจยวาร}\csromanheadernomark{2}{8. ทสฺสเนน\-ปหาตพฺพ\-ตฺติก 3. ปจฺจยวาร}
\csromanmarkh{1. ปจฺจยานุโลม}\csromanmarkhii{1. วิภงฺควาร}\csromanheadernomark{2}{1. \textbf{ปจฺจยานุโลม 1. วิภงฺควาร}}
\csromanheader{2}{\textbf{เหตุ}}
\proseitem{36}{ทสฺสเนน ปหาตพฺพํ ธมฺมํ ปจฺจยา ทสฺสเนน ปหาตพฺโพ ธมฺโม อุปฺปชฺชติ เหตุปจฺจยา. ตีณิ.}

\prose{ภาวนาย ปหาตพฺพํ ธมฺมํ ปจฺจยา ภาวนาย. ตีณิ.}

\proseitem{37}{\csromanfolio{146}เนวทสฺสเนน นภาวนาย ปหาตพฺพํ ธมฺมํ ปจฺจยา เนวทสฺสเนน นภาวนาย ปหาตพฺโพ ธมฺโม อุปฺปชฺชติ เหตุปจฺจยา. เนวทสฺสเนน นภาวนาย ปหาตพฺพํ เอกํ ขนฺธํ ปจฺจยา ตโย ขนฺธา จิตฺต\-สมุฏฺฐานญฺจ รูปํ ฯเปฯ เทฺว ขนฺเธ ฯเปฯ. ปฏิสนฺธิ\-กฺขเณ เนวทสฺสเนน นภาวนาย ปหาตพฺพํ เอกํ ขนฺธํ ปจฺจยา ตโย ขนฺธา กฏตฺตา จ รูปํ ฯเปฯ เทฺว ขนฺเธ ปจฺจยา เทฺว ขนฺธา ฯเปฯ. ขนฺเธ ปจฺจยา วตฺถุ, วตฺถุํ ปจฺจยา ขนฺธา. เอกํ มหาภูตํ ปจฺจยา ตโย มหาภูตา ฯเปฯ มหาภูเต ปจฺจยา จิตฺต\-สมุฏฺฐานํ รูปํ, กฏตฺตารูปํ, อุปาทารูปํ. วตฺถุํ ปจฺจยา เนวทสฺสเนน นภาวนาย ปหาตพฺพา ขนฺธา. (1)}

\prose{เนวทสฺสเนน นภาวนาย ปหาตพฺพํ ธมฺมํ ปจฺจยา ทสฺสเนน ปหาตพฺโพ ธมฺโม อุปฺปชฺชติ เหตุปจฺจยา. วตฺถุํ ปจฺจยา ทสฺสเนน ปหาตพฺพา ขนฺธา. (2)}

\prose{เนวทสฺสเนน นภาวนาย ปหาตพฺพํ ธมฺมํ ปจฺจยา ภาวนาย ปหาตพฺโพ ธมฺโม อุปฺปชฺชติ เหตุปจฺจยา. วตฺถุํ ปจฺจยา ภาวนาย ปหาตพฺพา ขนฺธา. (3)}

\prose{เนวทสฺสเนน นภาวนาย ปหาตพฺพํ ธมฺมํ ปจฺจยา ทสฺสเนน ปหาตพฺโพ จ เนวทสฺสเนน นภาวนาย ปหาตพฺโพ จ ธมฺมา อุปฺปชฺชนฺติ เหตุปจฺจยา. วตฺถุํ ปจฺจยา ทสฺสเนน ปหาตพฺพา ขนฺธา. มหาภูเต ปจฺจยา จิตฺต\-สมุฏฺฐานํ รูปํ. (4)}

\prose{เนวทสฺสเนน นภาวนาย ปหาตพฺพํ ธมฺมํ ปจฺจยา ภาวนาย ปหาตพฺโพ จ เนวทสฺสเนน นภาวนาย ปหาตพฺโพ จ ธมฺมา อุปฺปชฺชนฺติ เหตุปจฺจยา. วตฺถุํ ปจฺจยา ภาวนาย ปหาตพฺพา ขนฺธา. มหาภูเต ปจฺจยา จิตฺต\-สมุฏฺฐานํ รูปํ. (5)}

\proseitem{38}{ทสฺสเนน ปหาตพฺพญฺจ เนวทสฺสเนน นภาวนาย ปหาตพฺพญฺจ ธมฺมํ ปจฺจยา ทสฺสเนน ปหาตพฺโพ ธมฺโม อุปฺปชฺชติ เหตุปจฺจยา. ทสฺสเนน ปหาตพฺพํ เอกํ ขนฺธญฺจ วตฺถุญฺจ ปจฺจยา ตโย ขนฺธา ฯเปฯ เทฺว ขนฺเธ จ วตฺถุญฺจ ปจฺจโย เทฺว ขนฺธา. (1)}

\prose{ทสฺสเนน ปหาตพฺพญฺจ เนวทสฺสเนน นภาวนาย ปหาตพฺพญฺจ ธมฺมํ ปจฺจยา เนวทสฺสเนน นภาวนาย ปหาตพฺโพ ธมฺโม อุปฺปชฺชติ \csromanfolio{147}เหตุปจฺจยา. ทสฺสเนน ปหาตพฺเพ ขนฺเธ จ มหาภูเต จ ปจฺจยา จิตฺต\-สมุฏฺฐานํ รูปํ. (2)}

\prose{ทสฺสเนน ปหาตพฺพญฺจ เนวทสฺสเนน นภาวนาย ปหาตพฺพญฺจ ธมฺมํ ปจฺจยา ทสฺสเนน ปหาตพฺโพ จ เนวทสฺสเนน นภาวนาย ปหาตพฺโพ จ ธมฺมา อุปฺปชฺชนฺติ เหตุปจฺจยา. ทสฺสเนน ปหาตพฺพํ เอกํ ขนฺธญฺจ วตฺถุญฺจ ปจฺจยา ตโย ขนฺธา ฯเปฯ เทฺว ขนฺเธ จ วตฺถุญฺจ ปจฺจยา เทฺว ขนฺธา. ทสฺสเนน ปหาตพฺเพ ขนฺเธ จ มหาภูเต จ ปจฺจยา จิตฺต\-สมุฏฺฐานํ รูปํ. (3)}

\proseitem{39}{ภาวนาย ปหาตพฺพญฺจ เนวทสฺสเนน นภาวนาย ปหาตพฺพญฺจ ธมฺมํ ปจฺจยา ภาวนาย ปหาตพฺโพ ธมฺโม อุปฺปชฺชติ เหตุปจฺจยา. ภาวนาย ปหาตพฺพํ เอกํ ขนฺธญฺจ วตฺถุญฺจ ปจฺจยา ตโย ขนฺธา ฯเปฯ เทฺว ขนฺเธ จ วตฺถุญฺจ ฯเปฯ. (1)}

\prose{ภาวนาย ปหาตพฺพญฺจ เนวทสฺสเนน นภาวนาย ปหาตพฺพญฺจ ธมฺมํ ปจฺจยา เนวทสฺสเนน นภาวนาย ปหาตพฺโพ ธมฺโม อุปฺปชฺชติ เหตุปจฺจยา. ภาวนาย ปหาตพฺเพ ขนฺเธ จ มหาภูเต จ ปจฺจยา จิตฺต\-สมุฏฺฐานํ รูปํ. (2)}

\prose{ภาวนาย ปหาตพฺพญฺจ เนวทสฺสเนน นภาวนาย ปหาตพฺพญฺจ ธมฺมํ ปจฺจยา ภาวนาย ปหาตพฺโพ จ เนวทสฺสเนน นภาวนาย ปหาตพฺโพ จ ธมฺมา อุปฺปชฺชนฺติ เหตุปจฺจยา. ภาวนาย ปหาตพฺพํ เอกํ ขนฺธญฺจ วตฺถุญฺจ ปจฺจยา ตโย ขนฺธา ฯเปฯ เทฺว ขนฺเธ จ วตฺถุญฺจ ฯเปฯ. ภาวนาย ปหาตพฺเพ ขนฺเธ จ มหาภูเต จ ปจฺจยา จิตฺต\-สมุฏฺฐานํ รูปํ. (3)}

\csromanheader{2}{\textbf{อารมฺมณ}}
\proseitem{40}{ทสฺสเนน ปหาตพฺพํ ธมฺมํ ปจฺจยา ทสฺสเนน ปหาตพฺโพ ธมฺโม อุปฺปชฺชติ อารมฺมณ\-ปจฺจยา. ทสฺสเนน ปหาตพฺพํ เอกํ ขนฺธํ ปจฺจยา ตโย ขนฺธา ฯเปฯ เทฺว ขนฺเธ ฯเปฯ. (1)}

\prose{ภาวนาย ปหาตพฺพํ ธมฺมํ ปจฺจยา ภาวนาย ปหาตพฺโพ ธมฺโม อุปฺปชฺชติ อารมฺมณ\-ปจฺจยา. ภาวนาย ปหาตพฺพํ เอกํ ขนฺธํ ปจฺจยา ตโย ขนฺธา ฯเปฯ. (1)}

\prose{\csromanfolio{148}เนวทสฺสเนน นภาวนาย ปหาตพฺพํ ธมฺมํ ปจฺจยา เนวทสฺสเนน นภาวนาย ปหาตพฺโพ ธมฺโม อุปฺปชฺชติ อารมฺมณ\-ปจฺจยา. เนวทสฺสเนน นภาวนาย ปหาตพฺพํ เอกํ ขนฺธํ ปจฺจยา ตโย ขนฺธา ฯเปฯ เทฺว ขนฺเธ ฯเปฯ. ปฏิสนฺธิ\-กฺขเณ เนวทสฺสเนน นภาวนาย ปหาตพฺพํ เอกํ ขนฺธํ ปจฺจยา ตโย ขนฺธา ฯเปฯ เทฺว ขนฺเธ ฯเปฯ. วตฺถุํ ปจฺจยา ขนฺธา. จกฺขายตนํ ปจฺจยา จกฺขุวิญฺญาณํ ฯเปฯ. กายายตนํ ปจฺจยา กายวิญฺญาณํ. วตฺถุํ ปจฺจยา เนวทสฺสเนน นภาวนาย ปหาตพฺพา ขนฺธา. (1)}

\prose{เนวทสฺสเนน นภาวนาย ปหาตพฺพํ ธมฺมํ ปจฺจยา ทสฺสเนน ปหาตพฺโพ ธมฺโม อุปฺปชฺชติ อารมฺมณ\-ปจฺจยา. วตฺถุํ ปจฺจยา ทสฺสเนน ปหาตพฺพา ขนฺธา. (2)}

\prose{เนวทสฺสเนน นภาวนาย ปหาตพฺพํ ธมฺมํ ปจฺจยา ภาวนาย ปหาตพฺโพ ธมฺโม อุปฺปชฺชติ อารมฺมณ\-ปจฺจยา. วตฺถุํ ปจฺจยา ภาวนาย ปหาตพฺพา ขนฺธา. (3)}

\proseitem{41}{ทสฺสเนน ปหาตพฺพญฺจ เนวทสฺสเนน นภาวนาย ปหาตพฺพญฺจ ธมฺมํ ปจฺจยา ทสฺสเนน ปหาตพฺโพ ธมฺโม อุปฺปชฺชติ อารมฺมณ\-ปจฺจยา. ทสฺสเนน ปหาตพฺพํ เอกํ ขนฺธญฺจ วตฺถุญฺจ ปจฺจยา ตโย ขนฺธา ฯเปฯ เทฺว ขนฺเธ จ วตฺถุญฺจ ปจฺจยา ฯเปฯ. (1)}

\prose{ภาวนาย ปหาตพฺพญฺจ เนวทสฺสเนน นภาวนาย ปหาตพฺพญฺจ ธมฺมํ ปจฺจยา ภาวนาย ปหาตพฺโพ ธมฺโม อุปฺปชฺชติ อารมฺมณ\-ปจฺจยา. ภาวนาย ปหาตพฺพํ เอกํ ขนฺธญฺจ วตฺถุญฺจ ปจฺจยา ตโย ขนฺธา ฯเปฯ เทฺว ขนฺเธ จ วตฺถุญฺจ ปจฺจยา เทฺว ขนฺธา. (1)}

\csromanheader{2}{\textbf{อธิปตฺยาทิ}}
\proseitem{42}{ทสฺสเนน ปหาตพฺพํ ธมฺมํ ปจฺจยา ทสฺสเนน ปหาตพฺโพ ธมฺโม อุปฺปชฺชติ อธิปติ\-ปจฺจยา. (ปริปุณฺณํ, ปฏิสนฺธิ นตฺถิ.) อนนฺตร\-ปจฺจยา. สมนนฺตร\-ปจฺจยา. (อารมฺมณ\-สทิสํ.)}

\csromanheader{2}{8. ทสฺสเนน\-ปหาตพฺพ\-ตฺติก สหชาต}
\prose{\csromanfolio{149}ทสฺสเนน ปหาตพฺพํ ธมฺมํ ปจฺจยา ทสฺสเนน ปหาตพฺโพ ธมฺโม อุปฺปชฺชติ สหชาต\-ปจฺจยา. ทสฺสเนน ปหาตพฺพํ เอกํ ขนฺธํ. ตีณิ.}

\prose{ภาวนาย ปหาตพฺพํ ธมฺมํ ปจฺจยา. ตีณิ.}

\proseitem{43}{เนวทสฺสเนน นภาวนาย ปหาตพฺพํ ธมฺมํ ปจฺจยา เนวทสฺสเนน นภาวนาย ปหาตพฺโพ ธมฺโม อุปฺปชฺชติ สหชาต\-ปจฺจยา. เนวทสฺสเนน นภาวนาย ปหาตพฺพํ เอกํ ขนฺธํ ปจฺจยา ตโย ขนฺธา จิตฺต\-สมุฏฺฐานญฺจ รูปํ ฯเปฯ เทฺว ขนฺเธ ปจฺจยา เทฺว ขนฺธา ฯเปฯ. ปฏิสนฺธิ\-กฺขเณ ฯเปฯ. ขนฺเธ ปจฺจยา วตฺถุ, วตฺถุํ ปจฺจยา ขนฺธา. เอกํ มหาภูตํ ปจฺจยา ตโย มหาภูตา ฯเปฯ มหาภูเต ปจฺจยา ฯเปฯ. พาหิรํ. อาหารสมุฏฺฐานํ. อุตุสมุฏฺฐานํ. อสญฺญสตฺตานํ เอกํ ฯเปฯ. จกฺขายตนํ ปจฺจยา จกฺขุวิญฺญาณํ ฯเปฯ. กายายตนํ ปจฺจยา กายวิญฺญาณํ. วตฺถุํ ปจฺจยา เนวทสฺสเนน นภาวนาย ปหาตพฺพา ขนฺธา.}

\prose{เนวทสฺสเนน นภาวนาย ปหาตพฺพํ ธมฺมํ ปจฺจยา ทสฺสเนน ปหาตพฺโพ ธมฺโม อุปฺปชฺชติ สหชาต\-ปจฺจยา. (อวเสสา เหตุปจฺจย\-สทิสา.)}

\csromanheader{2}{\textbf{อญฺญมญฺญาทิ}}
\proseitem{44}{ทสฺสเนน ปหาตพฺพํ ธมฺมํ ปจฺจยา ทสฺสเนน ปหาตพฺโพ ธมฺโม อุปฺปชฺชติ อญฺญมญฺญ\-ปจฺจยา. นิสฺสยปจฺจยา. อุปนิสฺสย\-ปจฺจยา. ปุเรชาต\-ปจฺจยา. (ปฏิสนฺธิ นตฺถิ.) อาเสวนปจฺจยา. (ปฏิสนฺธิ นตฺถิ, วิปากญฺจ.) กมฺมปจฺจยา. วิปากปจฺจยา. อาหารปจฺจยา. อินฺทฺริยปจฺจยา. ฌานปจฺจยา. มคฺคปจฺจยา. สมฺปยุตฺต\-ปจฺจยา. วิปฺปยุตฺต\-ปจฺจยา. อตฺถิปจฺจยา. นตฺถิปจฺจยา. วิคตปจฺจยา. อวิคตปจฺจยา.}

\csromanmarkh{1. ปจฺจยานุโลม}\csromanmarkhii{2. สงฺขฺยาวาร}\csromanheadernomark{2}{1. \textbf{ปจฺจยานุโลม 2. สงฺขฺยาวาร}}
\proseitemwithcloserruled{45}{เหตุยา สตฺตรส, อารมฺมเณ สตฺต, อธิปติยา สตฺตรส, อนนฺตเร สตฺต, สมนนฺตเร สตฺต, สหชาเต สตฺตรส, อญฺญมญฺเญ สตฺต, นิสฺสเย สตฺตรส, อุปนิสฺสเย สตฺต, ปุเรชาเต สตฺต, \csromanfolio{150}อาเสวเน สตฺต, กมฺเม สตฺตรส, วิปาเก เอกํ, อาหาเร สตฺตรส, อินฺทฺริเย สตฺตรส, ฌาเน สตฺตรส, มคฺเค สตฺตรส, สมฺปยุตฺเต สตฺต, วิปฺปยุตฺเต สตฺตรส, อตฺถิยา สตฺตรส, นตฺถิยา สตฺต, วิคเต สตฺต, อวิคเต สตฺตรส. (เอวํ คเณตพฺพํ.)}{อนุโลมํ.}
\csromanmarkh{2. ปจฺจย\-ปจฺจนีย}\csromanmarkhii{1. วิภงฺควาร}\csromanheadernomark{2}{2. ปจฺจย\-ปจฺจนีย 1. วิภงฺควาร}
\csromanheader{2}{\textbf{นเหตุ}}
\proseitem{46}{ทสฺสเนน ปหาตพฺพํ ธมฺมํ ปจฺจยา ทสฺสเนน ปหาตพฺโพ ธมฺโม อุปฺปชฺชติ นเหตุปจฺจยา. วิจิกิจฺฉา\-สหคเต ขนฺเธ ปจฺจยา วิจิกิจฺฉา\-สหคโต โมโห. (1)}

\prose{ภาวนาย ปหาตพฺพํ ธมฺมํ ปจฺจยา ภาวนาย ปหาตพฺโพ ธมฺโม อุปฺปชฺชติ นเหตุปจฺจยา. อุทฺธจฺจ\-สหคเต ขนฺเธ ปจฺจยา อุทฺธจฺจ\-สหคโต โมโห. (1)}

\prose{เนวทสฺสเนน นภาวนาย ปหาตพฺพํ ธมฺมํ ปจฺจยา เนวทสฺสเนน นภาวนาย ปหาตพฺโพ ธมฺโม อุปฺปชฺชติ นเหตุปจฺจยา. อเหตุกํ เนวทสฺสเนน นภาวนาย ปหาตพฺพํ เอกํ ขนฺธํ ปจฺจยา ตโย ขนฺธา จิตฺต\-สมุฏฺฐานญฺจ รูปํ ฯเปฯ เทฺว ขนฺเธ ฯเปฯ. อเหตุก\-ปฏิสนฺธิ\-กฺขเณ ฯเปฯ. ขนฺเธ ปจฺจยา วตฺถุ, วตฺถุํ ปจฺจยา ขนฺธา. เอกํ มหาภูตํ ฯเปฯ. พาหิรํ. อาหารสมุฏฺฐานํ. อุตุสมุฏฺฐานํ. อสญฺญสตฺตานํ ฯเปฯ. จกฺขายตนํ ปจฺจยา จกฺขุวิญฺญาณํ ฯเปฯ. กายายตนํ ปจฺจยา กายวิญฺญาณํ. วตฺถุํ ปจฺจยา อเหตุกา เนวทสฺสเนน นภาวนาย ปหาตพฺพา ขนฺธา. (1)}

\prose{เนวทสฺสเนน นภาวนาย ปหาตพฺพํ ธมฺมํ ปจฺจยา ทสฺสเนน ปหาตพฺโพ ธมฺโม อุปฺปชฺชติ นเหตุปจฺจยา. วตฺถุํ ปจฺจยา วิจิกิจฺฉา\-สหคโต โมโห. (2)}

\prose{เนวทสฺสเนน นภาวนาย ปหาตพฺพํ ธมฺมํ ปจฺจยา ภาวนาย ปหาตพฺโพ ธมฺโม อุปฺปชฺชติ นเหตุปจฺจยา. วตฺถุํ ปจฺจยา อุทฺธจฺจ\-สหคโต โมโห. (3)}

\proseitem{47}{\csromanfolio{151}ทสฺสเนน ปหาตพฺพญฺจ เนวทสฺสเนน นภาวนาย ปหาตพฺพญฺจ ธมฺมํ ปจฺจยา ทสฺสเนน ปหาตพฺโพ ธมฺโม อุปฺปชฺชติ นเหตุปจฺจยา. วิจิกิจฺฉา\-สหคเต ขนฺเธ จ วตฺถุญฺจ ปจฺจยา วิจิกิจฺฉา\-สหคโต โมโห. (1)}

\prose{ภาวนาย ปหาตพฺพญฺจ เนวทสฺสเนน นภาวนาย ปหาตพฺพญฺจ ธมฺมํ ปจฺจยา ภาวนาย ปหาตพฺโพ ธมฺโม อุปฺปชฺชติ นเหตุปจฺจยา. อุทฺธจฺจ\-สหคเต ขนฺเธ จ วตฺถุญฺจ ปจฺจยา อุทฺธจฺจ\-สหคโต โมโห. (1)}

\prose{นอารมฺมณ}

\proseitem{48}{ทสฺสเนน ปหาตพฺพํ ธมฺมํ ปจฺจยา เนวทสฺสเนน นภาวนาย ปหาตพฺโพ ธมฺโม อุปฺปชฺชติ นอารมฺมณ\-ปจฺจยา. ทสฺสเนน ปหาตพฺเพ ขนฺเธ ปจฺจยา จิตฺต\-สมุฏฺฐานํ รูปํ. (1)}

\prose{ภาวนาย ปหาตพฺพํ ธมฺมํ ปจฺจยา เนวทสฺสเนน นภาวนาย ปหาตพฺโพ ธมฺโม อุปฺปชฺชติ นอารมฺมณ\-ปจฺจยา. ภาวนาย ปหาตพฺเพ ขนฺเธ ปจฺจยา จิตฺต\-สมุฏฺฐานํ รูปํ. (1)}

\prose{เนวทสฺสเนน นภาวนาย ปหาตพฺพํ ธมฺมํ ปจฺจยา เนวทสฺสเนน นภาวนาย ปหาตพฺโพ ธมฺโม อุปฺปชฺชติ นอารมฺมณ\-ปจฺจยา. เนวทสฺสเนน นภาวนาย ปหาตพฺเพ ขนฺเธ ปจฺจยา จิตฺต\-สมุฏฺฐานํ รูปํ. ปฏิสนฺธิ\-กฺขเณ เนวทสฺสเนน นภาวนาย ปหาตพฺเพ ขนฺเธ ปจฺจยา กฏตฺตารูปํ. ขนฺเธ ปจฺจยา วตฺถุ. เอกํ มหาภูตํ ปจฺจยา ฯเปฯ. พาหิรํ. อาหารสมุฏฺฐานํ. อุตุสมุฏฺฐานํ. อสญฺญสตฺตานํ ฯเปฯ. (1)}

\proseitem{49}{ทสฺสเนน ปหาตพฺพญฺจ เนวทสฺสเนน นภาวนาย ปหาตพฺพญฺจ ธมฺมํ ปจฺจยา ทสฺสเนน ปหาตพฺโพ ธมฺโม อุปฺปชฺชติ นอารมฺมณ\-ปจฺจยา. ทสฺสเนน ปหาตพฺเพ ขนฺเธ จ มหาภูเต จ ปจฺจยา จิตฺต\-สมุฏฺฐานํ รูปํ. (1)}

\prose{ภาวนาย ปหาตพฺพญฺจ เนวทสฺสเนน นภาวนาย ปหาตพฺพญฺจ ธมฺมํ ปจฺจยา เนวทสฺสเนน นภาวนาย ปหาตพฺโพ ธมฺโม อุปฺปชฺชติ นอารมฺมณ\-ปจฺจยา. ภาวนาย ปหาตพฺเพ ขนฺเธ จ มหาภูเต จ ปจฺจยา จิตฺต\-สมุฏฺฐานํ รูปํ. (1)}

\prose{\csromanfolio{152}นอธิปตฺยาทิ}

\proseitem{50}{ทสฺสเนน ปหาตพฺพํ ธมฺมํ ปจฺจยา ทสฺสเนน ปหาตพฺโพ ธมฺโม อุปฺปชฺชติ นอธิปติ\-ปจฺจยา. (สหชาต\-สทิสํ.) นอนนฺตร\-ปจฺจยา. นสมนนฺตร\-ปจฺจยา. นอญฺญมญฺญ\-ปจฺจยา. นอุปนิสฺสยปจฺจยา.}

\csromanheader{2}{\textbf{นปุเรชาต}}
\proseitem{51}{ทสฺสเนน ปหาตพฺพํ ธมฺมํ ปจฺจยา ทสฺสเนน ปหาตพฺโพ ธมฺโม อุปฺปชฺชติ นปุเรชาต\-ปจฺจยา. อรูเป ทสฺสเนน ปหาตพฺพํ เอกํ ขนฺธํ ปจฺจยา ฯเปฯ. (1)}

\prose{ทสฺสเนน ปหาตพฺพํ ธมฺมํ ปจฺจยา เนวทสฺสเนน นภาวนาย ปหาตพฺโพ ธมฺโม อุปฺปชฺชติ นปุเรชาต\-ปจฺจยา. ทสฺสเนน ปหาตพฺเพ ขนฺเธ ปจฺจยา จิตฺต\-สมุฏฺฐานํ รูปํ. (2)}

\prose{ภาวนาย ปหาตพฺพํ ธมฺมํ ปจฺจยา ภาวนาย ปหาตพฺโพ ธมฺโม อุปฺปชฺชติ นปุเรชาต\-ปจฺจยา. อรูเป ภาวนาย ปหาตพฺพํ เอกํ ขนฺธํ ปจฺจยา ฯเปฯ. (1)}

\prose{ภาวนาย ปหาตพฺพํ ธมฺมํ ปจฺจยา เนวทสฺสเนน นภาวนาย ปหาตพฺโพ ธมฺโม อุปฺปชฺชติ นปุเรชาต\-ปจฺจยา. ภาวนาย ปหาตพฺเพ ขนฺเธ ปจฺจยา จิตฺต\-สมุฏฺฐานํ รูปํ. (2)}

\proseitem{52}{เนวทสฺสเนน นภาวนาย ปหาตพฺพํ ธมฺมํ ปจฺจยา เนวทสฺสเนน นภาวนาย ปหาตพฺโพ ธมฺโม อุปฺปชฺชติ นปุเรชาต\-ปจฺจยา. อรูเป เนวทสฺสเนน นภาวนาย ปหาตพฺพํ เอกํ ขนฺธํ ปจฺจยา ตโย ขนฺธา ฯเปฯ. เนวทสฺสเนน นภาวนาย ปหาตพฺเพ ขนฺเธ ปจฺจยา จิตฺต\-สมุฏฺฐานํ รูปํ. ปฏิสนฺธิ\-กฺขเณ ฯเปฯ. ขนฺเธ ปจฺจยา กฏตฺตารูปํ. ขนฺเธ ปจฺจยา วตฺถุ, วตฺถุํ ปจฺจยา ขนฺธา. เอกํ มหาภูตํ ฯเปฯ. พาหิรํ. อาหารสมุฏฺฐานํ. อุตุสมุฏฺฐานํ. อสญฺญสตฺตานํ ฯเปฯ. (1)}

\prose{ทสฺสเนน ปหาตพฺพญฺจ เนวทสฺสเนน นภาวนาย ปหาตพฺพญฺจ ธมฺมํ ปจฺจยา เนวทสฺสเนน นภาวนาย ปหาตพฺโพ ธมฺโม อุปฺปชฺชติ นปุเรชาต\-ปจฺจยา. ทสฺสเนน ปหาตพฺเพ ขนฺเธ จ มหาภูเต จ ปจฺจยา จิตฺต\-สมุฏฺฐานํ รูปํ. (1)}

\prose{\csromanfolio{153}ภาวนาย ปหาตพฺพญฺจ เนวทสฺสเนน นภาวนาย ปหาตพฺพญฺจ ธมฺมํ ปจฺจยา เนวทสฺสเนน นภาวนาย ปหาตพฺโพ ธมฺโม อุปฺปชฺชติ นปุเรชาต\-ปจฺจยา. ภาวนาย ปหาตพฺเพ ขนฺเธ จ มหาภูเต จ ปจฺจยา จิตฺต\-สมุฏฺฐานํ รูปํ. (1)}

\csromanheader{2}{\textbf{นปจฺฉาชาตาทิ}}
\proseitem{53}{ทสฺสเนน ปหาตพฺพํ ธมฺมํ ปจฺจยา ทสฺสเนน ปหาตพฺโพ ธมฺโม อุปฺปชฺชติ นปจฺฉาชาต\-ปจฺจยา. นอาเสวน\-ปจฺจยา.}

\csromanheader{2}{\textbf{นกมฺม}}
\proseitem{54}{ทสฺสเนน ปหาตพฺพํ ธมฺมํ ปจฺจยา ทสฺสเนน ปหาตพฺโพ ธมฺโม อุปฺปชฺชติ นกมฺมปจฺจยา. ทสฺสเนน ปหาตพฺเพ ขนฺเธ ปจฺจยา ทสฺสเนน ปหาตพฺพา เจตนา. (1)}

\prose{ภาวนาย ปหาตพฺพํ ธมฺมํ ฯเปฯ นกมฺมปจฺจยา. ภาวนาย ปหาตพฺเพ ขนฺเธ ปจฺจยา ภาวนาย ปหาตพฺพา เจตนา. (1)}

\prose{เนวทสฺสเนน นภาวนาย ปหาตพฺพํ ธมฺมํ ฯเปฯ นกมฺมปจฺจยา. เนวทสฺสเนน นภาวนาย ปหาตพฺเพ ขนฺเธ ปจฺจยา เนวทสฺสเนน นภาวนาย ปหาตพฺพา เจตนา. พาหิรํ. อาหารสมุฏฺฐานํ. อุตุสมุฏฺฐานํ ฯเปฯ. (1)}

\prose{เนวทสฺสเนน นภาวนาย ปหาตพฺพํ ธมฺมํ ปจฺจยา ทสฺสเนน ปหาตพฺโพ ธมฺโม อุปฺปชฺชติ นกมฺมปจฺจยา. วตฺถุํ ปจฺจยา ทสฺสเนน ปหาตพฺพา เจตนา. (2)}

\prose{เนวทสฺสเนน นภาวนาย ปหาตพฺพํ ธมฺมํ ปจฺจยา ภาวนาย ปหาตพฺโพ ธมฺโม อุปฺปชฺชติ นกมฺมปจฺจยา. วตฺถุํ ปจฺจยา ภาวนาย ปหาตพฺพา เจตนา. (3)}

\proseitem{55}{ทสฺสเนน ปหาตพฺพญฺจ เนวทสฺสเนน นภาวนาย ปหาตพฺพญฺจ ธมฺมํ ปจฺจยา ทสฺสเนน ปหาตพฺโพ ธมฺโม อุปฺปชฺชติ นกมฺมปจฺจยา. ทสฺสเนน ปหาตพฺเพ ขนฺเธ จ วตฺถุญฺจ ปจฺจยา ทสฺสเนน ปหาตพฺพา เจตนา. (1)}

\prose{ภาวนาย ปหาตพฺพญฺจ เนวทสฺสเนน นภาวนาย ปหาตพฺพญฺจ ธมฺมํ ปจฺจยา ภาวนาย ปหาตพฺโพ ธมฺโม อุปฺปชฺชติ นกมฺมปจฺจยา. \csromanfolio{154}ภาวนาย ปหาตพฺเพ ขนฺเธ จ วตฺถุญฺจ ปจฺจยา ภาวนาย ปหาตพฺพา เจตนา. (1)}

\csromanheader{2}{\textbf{นวิปากาทิ}}
\proseitem{56}{ทสฺสเนน ปหาตพฺพํ ธมฺมํ ปจฺจยา ทสฺสเนน ปหาตพฺโพ ธมฺโม อุปฺปชฺชติ นวิปาก\-ปจฺจยา. (ปริปุณฺณํ, ปฏิสนฺธิ นตฺถิ.) นอาหารปจฺจยา. พาหิรํ. อุตุสมุฏฺฐานํ. อสญฺญสตฺตานํ ฯเปฯ นอินฺทฺริยปจฺจยา. พาหิรํ. อาหารสมุฏฺฐานํ. อุตุสมุฏฺฐานํ. อสญฺญสตฺตานํ ฯเปฯ มหาภูเต ปจฺจยา รูปชีวิตินฺทฺริยํ. นฌานปจฺจยา. ปญฺจวิญฺญาณํ ฯเปฯ. พาหิรํ ฯเปฯ. อสญฺญสตฺตานํ ฯเปฯ นมคฺคปจฺจยา. อเหตุกํ เนวทสฺสเนน นภาวนาย ปหาตพฺพํ เอกํ ขนฺธํ ปจฺจยา ตโย ขนฺธา จิตฺต\-สมุฏฺฐานญฺจ รูปํ ฯเปฯ. อเหตุก\-ปฏิสนฺธิ\-กฺขเณ ฯเปฯ. เอกํ มหาภูตํ ฯเปฯ. อสญฺญสตฺตานํ ฯเปฯ นสมฺปยุตฺต\-ปจฺจยา. นวิปฺปยุตฺต\-ปจฺจยา. อรูเป ทสฺสเนน ปหาตพฺพํ เอกํ ขนฺธํ ปจฺจยา ฯเปฯ. อรูเป ภาวนาย ปหาตพฺพํ เอกํ ขนฺธํ ปจฺจยา ตโย ขนฺธา ฯเปฯ. เนวทสฺสเนน นภาวนาย ปหาตพฺพํ ธมฺมํ ปจฺจยา เนวทสฺสเนน ฯเปฯ. อรูเป เนวทสฺสเนน นภาวนาย ปหาตพฺพํ เอกํ ขนฺธํ ปจฺจยา ตโย ขนฺธา ฯเปฯ เทฺว ขนฺเธ ฯเปฯ. พาหิรํ. อาหารสมุฏฺฐานํ. อุตุสมุฏฺฐานํ. อสญฺญสตฺตานํ ฯเปฯ โนนตฺถิปจฺจยา. โนวิคต\-ปจฺจยา.}

\csromanmarkh{2. ปจฺจย\-ปจฺจนีย}\csromanmarkhii{2. สงฺขฺยาวาร}\csromanheadernomark{2}{2. ปจฺจย\-ปจฺจนีย 2. สงฺขฺยาวาร}
\proseitemwithcloserruled{57}{นเหตุยา สตฺต, นอารมฺมเณ ปญฺจ, นอธิปติยา สตฺตรส, นอนนฺตเร ปญฺจ, นสมนนฺตเร ปญฺจ, นอญฺญมญฺเญ ปญฺจ, นอุปนิสฺสเย ปญฺจ, นปุเรชาเต สตฺต, นปจฺฉาชาเต สตฺตรส, นอาเสวเน สตฺตรส, นกมฺเม สตฺต, นวิปาเก สตฺตรส, นอาหาเร เอกํ, นอินฺทฺริเย เอกํ, นฌาเน เอกํ, นมคฺเค เอกํ, นสมฺปยุตฺเต ปญฺจ, นวิปฺปยุตฺเต ตีณิ, โนนตฺถิยา ปญฺจ, โนวิคเต ปญฺจ.}{ปจฺจนียํ.}
\proseitemwithcloserruled{58}{เหตุปจฺจยา นอารมฺมเณ ปญฺจ, นอธิปติยา สตฺตรส, นอนนฺตเร ปญฺจ, นสมนนฺตเร ปญฺจ, นอญฺญมญฺเญ ปญฺจ, นอุปนิสฺสเย ปญฺจ, \csromanfolio{155}นปุเรชาเต สตฺต, นปจฺฉาชาเต สตฺตรส, นอาเสวเน สตฺตรส, นกมฺเม สตฺต, นวิปาเก สตฺตรส, นสมฺปยุตฺเต ปญฺจ, นวิปฺปยุตฺเต ตีณิ, โนนตฺถิยา ปญฺจ, โนวิคเต ปญฺจ.}{อนุโลม\-ปจฺจนียํ.}
\proseitem{59}{นเหตุปจฺจยา อารมฺมเณ สตฺต, อนนฺตเร สตฺต, สมนนฺตเร สตฺต, สหชาเต สตฺต, อญฺญมญฺเญ สตฺต, นิสฺสเย สตฺต, อุปนิสฺสเย สตฺต, ปุเรชาเต สตฺต, อาเสวเน สตฺต, กมฺเม สตฺต, วิปาเก เอกํ, อาหาเร สตฺต, อินฺทฺริเย สตฺต, ฌาเน สตฺต, มคฺเค ฉ, สมฺปยุตฺเต สตฺต, วิปฺปยุตฺเต สตฺต, อตฺถิยา สตฺต, นตฺถิยา สตฺต, วิคเต สตฺต, อวิคเต สตฺต. (เอวํ คเณตพฺพํ.) ปจฺจนียา\-นุโลมํ. ปจฺจยวาโร. (นิสฺสยวาโร ปจฺจยสทิโส กาตพฺโพ.)}
\csromansectionrule

\csromanmarkh{8. ทสฺสเนน\-ปหาตพฺพ\-ตฺติก}\csromanmarkhii{5. สํสฏฺฐวาร}\csromanheadernomark{2}{8. ทสฺสเนน\-ปหาตพฺพ\-ตฺติก 5. สํสฏฺฐวาร}
\csromanmarkh{1. ปจฺจยานุโลม}\csromanmarkhii{1. วิภงฺควาร}\csromanheadernomark{2}{1. \textbf{ปจฺจยานุโลม 1. วิภงฺควาร}}
\csromanheader{2}{\textbf{เหตุ}}
\proseitem{60}{ทสฺสเนน ปหาตพฺพํ ธมฺมํ สํสฏฺโฐ ทสฺสเนน ปหาตพฺโพ ธมฺโม อุปฺปชฺชติ เหตุปจฺจยา. ทสฺสเนน ปหาตพฺพํ เอกํ ขนฺธํ สํสฏฺฐา ตโย ขนฺธา ฯเปฯ เทฺว ขนฺเธ สํสฏฺฐา เทฺว ขนฺธา. (1)}

\prose{ภาวนาย ปหาตพฺพํ ธมฺมํ สํสฏฺโฐ ภาวนาย ปหาตพฺโพ ธมฺโม อุปฺปชฺชติ เหตุปจฺจยา. ภาวนาย ปหาตพฺพํ เอกํ ขนฺธํ สํสฏฺฐา ฯเปฯ. (1)}

\prose{เนวทสฺสเนน นภาวนาย ปหาตพฺพํ ธมฺมํ สํสฏฺโฐ เนวทสฺสเนน นภาวนาย ปหาตพฺโพ ธมฺโม อุปฺปชฺชติ เหตุปจฺจยา. เนวทสฺสเนน \csromanfolio{156}นภาวนาย ปหาตพฺพํ เอกํ ขนฺธํ สํสฏฺฐา ตโย ขนฺธา ฯเปฯ. ปฏิสนฺธิ\-กฺขเณ เนวทสฺสเนน นภาวนาย ปหาตพฺพํ เอกํ ขนฺธํ สํสฏฺฐา ฯเปฯ. (1)}

\csromanheader{2}{\textbf{อารมฺมณ}}
\proseitem{61}{ทสฺสเนน ปหาตพฺพํ ธมฺมํ สํสฏฺโฐ ทสฺสเนน ปหาตพฺโพ ธมฺโม อุปฺปชฺชติ อารมฺมณ\-ปจฺจยา. (สพฺพานิ ปทานิ วิตฺถาเร\-ตพฺพานิ. ตีณิ ตีณิ.)}

\csromanmarkh{1. ปจฺจยานุโลม}\csromanmarkhii{2. สงฺขฺยาวาร}\csromanheadernomark{2}{1. \textbf{ปจฺจยานุโลม 2. สงฺขฺยาวาร}}
\proseitemwithcloserruled{62}{เหตุยา ตีณิ, อารมฺมเณ ตีณิ, อธิปติยา ตีณิ, อนนฺตเร ตีณิ, สมนนฺตเร ตีณิ, สหชาเต ตีณิ, อญฺญมญฺเญ ตีณิ, นิสฺสเย ตีณิ, อุปนิสฺสเย ตีณิ, ปุเรชาเต ตีณิ, อาเสวเน ตีณิ, กมฺเม ตีณิ, วิปาเก เอกํ, อาหาเร ตีณิ, อินฺทฺริเย ตีณิ, ฌาเน ตีณิ, มคฺเค ตีณิ, สมฺปยุตฺเต ตีณิ, วิปฺปยุตฺเต ตีณิ, อตฺถิยา ตีณิ, นตฺถิยา ตีณิ, วิคเต ตีณิ, อวิคเต ตีณิ.}{อนุโลมํ.}
\csromanmarkh{2. ปจฺจย\-ปจฺจนีย}\csromanmarkhii{1. วิภงฺควาร}\csromanheadernomark{2}{2. ปจฺจย\-ปจฺจนีย 1. วิภงฺควาร}
\csromanheader{2}{\textbf{นเหตุ}}
\proseitem{63}{ทสฺสเนน ปหาตพฺพํ ธมฺมํ สํสฏฺโฐ ทสฺสเนน ปหาตพฺโพ ธมฺโม อุปฺปชฺชติ นเหตุปจฺจยา. วิจิกิจฺฉา\-สหคเต ขนฺเธ สํสฏฺโฐ วิจิกิจฺฉา\-สหคโต โมโห. (1)}

\prose{ภาวนาย ปหาตพฺพํ ธมฺมํ สํสฏฺโฐ ภาวนาย ปหาตพฺโพ ธมฺโม อุปฺปชฺชติ นเหตุปจฺจยา. อุทฺธจฺจ\-สหคเต ขนฺเธ สํสฏฺโฐ อุทฺธจฺจ\-สหคโต โมโห. (1)}

\prose{เนวทสฺสเนน นภาวนาย ปหาตพฺพํ ธมฺมํ สํสฏฺโฐ เนวทสฺสเนน นภาวนาย ปหาตพฺโพ ธมฺโม อุปฺปชฺชติ นเหตุปจฺจยา. อเหตุกํ เนวทสฺสเนน นภาวนาย ปหาตพฺพํ เอกํ ขนฺธํ สํสฏฺฐา ตโย ขนฺธา ฯเปฯ. อเหตุก\-ปฏิสนฺธิ\-กฺขเณ ฯเปฯ. (1)}

\csromanheader{2}{8. ทสฺสเนน\-ปหาตพฺพ\-ตฺติก นอธิปตฺยาทิ}
\proseitem{64}{\csromanfolio{157}ทสฺสเนน ปหาตพฺพํ ธมฺมํ สํสฏฺโฐ ทสฺสเนน ปหาตพฺโพ ธมฺโม อุปฺปชฺชติ นอธิปติ\-ปจฺจยา. นปุเรชาต\-ปจฺจยา. นปจฺฉาชาต\-ปจฺจยา. นอาเสวน\-ปจฺจยา. นกมฺมปจฺจยา. นวิปาก\-ปจฺจยา. เนวทสฺสเนน นภาวนาย ฯเปฯ นฌานปจฺจยา. ปญฺจวิญฺญาณํ ฯเปฯ นมคฺคปจฺจยา. อเหตุกํ ฯเปฯ. เนวทสฺสเนน นภาวนาย ปหาตพฺโพ ฯเปฯ นวิปฺปยุตฺต\-ปจฺจยา. ตีณิ.}

\csromanmarkh{2. ปจฺจย\-ปจฺจนีย}\csromanmarkhii{2. สงฺขฺยาวาร}\csromanheadernomark{2}{2. ปจฺจย\-ปจฺจนีย 2. สงฺขฺยาวาร}
\proseitemwithcloserruled{65}{นเหตุยา ตีณิ, นอธิปติยา ตีณิ, นปุเรชาเต ตีณิ, นปจฺฉาชาเต ตีณิ, นอาเสวเน ตีณิ, นกมฺเม ตีณิ, นวิปาเก ตีณิ, นฌาเน เอกํ, นมคฺเค เอกํ, นวิปฺปยุตฺเต ตีณิ.}{ปจฺจนียํ.}
\proseitemwithcloserruled{66}{เหตุปจฺจยา นอธิปติยา ตีณิ, นปุเรชาเต ตีณิ, นปจฺฉาชาเต ตีณิ, นอาเสวเน ตีณิ, นกมฺเม ตีณิ, นวิปาเก ตีณิ, นวิปฺปยุตฺเต ตีณิ.}{อนุโลม\-ปจฺจนียํ.}
\proseitem{67}{นเหตุปจฺจยา อารมฺมเณ ตีณิ, อนนฺตเร ตีณิ, สมนนฺตเร ตีณิ, สหชาเต ตีณิ, อญฺญมญฺเญ ตีณิ, นิสฺสเย ตีณิ, อุปนิสฺสเย ตีณิ, ปุเรชาเต ตีณิ, อาเสวเน ตีณิ, กมฺเม ตีณิ, วิปาเก เอกํ, อาหาเร ตีณิ, อินฺทฺริเย ตีณิ, ฌาเน ตีณิ, มคฺเค เทฺว, สมฺปยุตฺเต ตีณิ, วิปฺปยุตฺเต ตีณิ, อตฺถิยา ตีณิ, นตฺถิยา ตีณิ, วิคเต ตีณิ, อวิคเต ตีณิ.}

\vspace{\tipitakaparskip}
\csromancenter{ปจฺจนียา\-นุโลมํ.}
\nitthitamruled{สํสฏฺฐวาโร.}
\csromanmarkh{8. ทสฺสเนน\-ปหาตพฺพ\-ตฺติก}\csromanmarkhii{6. สมฺปยุตฺตวาร}\csromanheadernomark{2}{8. ทสฺสเนน\-ปหาตพฺพ\-ตฺติก 6. สมฺปยุตฺตวาร}
\csromanheader{2}{1. ปจฺจยา\-นุโลมาทิ}
\proseitem{68}{\csromanfolio{158}ทสฺสเนน ปหาตพฺพํ ธมฺมํ สมฺปยุตฺโต ทสฺสเนน ปหาตพฺโพ ธมฺโม อุปฺปชฺชติ เหตุปจฺจยา. (สมฺปยุตฺตวาโร สํสฏฺฐวาร\-สทิโส.)}
\csromansectionrule

\csromanmarkh{8. ทสฺสเนน\-ปหาตพฺพ\-ตฺติก}\csromanmarkhii{7. ปญฺหาวาร}\csromanheadernomark{2}{8. ทสฺสเนน\-ปหาตพฺพ\-ตฺติก 7. ปญฺหาวาร}
\csromanmarkh{1. ปจฺจยานุโลม}\csromanmarkhii{1. วิภงฺควาร}\csromanheadernomark{2}{1. \textbf{ปจฺจยานุโลม 1. วิภงฺควาร}}
\csromanheader{2}{\textbf{เหตุ}}
\proseitem{69}{ทสฺสเนน ปหาตพฺโพ ธมฺโม ทสฺสเนน ปหาตพฺพสฺส ธมฺมสฺส เหตุปจฺจเยน\-ปจฺจโย. ทสฺสเนน ปหาตพฺพา เหตู สมฺปยุตฺตกานํ ขนฺธานํ เหตุปจฺจเยน ปจฺจโย. (1)}

\prose{ทสฺสเนน ปหาตพฺโพ ธมฺโม เนวทสฺสเนน นภาวนาย ปหาตพฺพสฺส ธมฺมสฺส เหตุปจฺจเยน ปจฺจโย. ทสฺสเนน ปหาตพฺพา เหตู จิตฺต\-สมุฏฺฐานานํ รูปานํ เหตุปจฺจเยน ปจฺจโย. (2)}

\prose{ทสฺสเนน ปหาตพฺโพ ธมฺโม ทสฺสเนน ปหาตพฺพสฺส จ เนวทสฺสเนน นภาวนาย ปหาตพฺพสฺส จ ธมฺมสฺส เหตุปจฺจเยน ปจฺจโย. ทสฺสเนน ปหาตพฺพา เหตู สมฺปยุตฺตกานํ ขนฺธานํ จิตฺตสมุฏฺฐานานญฺจ รูปานํ เหตุปจฺจเยน ปจฺจโย. (3)}

\proseitem{70}{ภาวนาย ปหาตพฺโพ ธมฺโม ฯเปฯ ธมฺมสฺส เหตุปจฺจเยน ปจฺจโย. ภาวนาย ปหาตพฺพา เหตู สมฺปยุตฺตกานํ ขนฺธานํ เหตุปจฺจเยน ปจฺจโย. (1)}

\prose{ภาวนาย ปหาตพฺโพ ธมฺโม เนวทสฺสเนน นภาวนาย ปหาตพฺพสฺส ฯเปฯ. ภาวนาย ปหาตพฺพา เหตู จิตฺต\-สมุฏฺฐานานํ รูปานํ เหตุปจฺจเยน ปจฺจโย. (2)}

\prose{ภาวนาย ปหาตพฺโพ ธมฺโม ภาวนาย ปหาตพฺพสฺส จ เนวทสฺสเนน นภาวนาย ปหาตพฺพสฺส จ ธมฺมสฺส ฯเปฯ. ภาวนาย ปหาตพฺพา \csromanfolio{159}เหตู สมฺปยุตฺตกานํ ขนฺธานํ จิตฺตสมุฏฺฐานานญฺจ รูปานํ เหตุปจฺจเยน ปจฺจโย. (3)}

\prose{เนวทสฺสเนน นภาวนาย ปหาตพฺโพ ธมฺโม เนวทสฺสเนน นภาวนาย ปหาตพฺพสฺส ธมฺมสฺส ฯเปฯ. เนวทสฺสเนน นภาวนาย ปหาตพฺพา เหตู สมฺปยุตฺตกานํ ขนฺธานํ จิตฺตสมุฏฺฐานานญฺจ รูปานํ เหตุปจฺจเยน ปจฺจโย. ปฏิสนฺธิ\-กฺขเณ เนวทสฺสเนน นภาวนาย ปหาตพฺพา เหตู สมฺปยุตฺตกานํ ขนฺธานํ กฏตฺตา จ รูปานํ เหตุปจฺจเยน ปจฺจโย. (1)}

\csromanheader{2}{\textbf{อารมฺมณ}}
\proseitem{71}{ทสฺสเนน ปหาตพฺโพ ธมฺโม ทสฺสเนน ปหาตพฺพสฺส ธมฺมสฺส อารมฺมณ\-ปจฺจเยน ปจฺจโย. ทสฺสเนน ปหาตพฺพํ ราคํ อสฺสาเทติ อภินนฺทติ, ตํ อารพฺภ ทสฺสเนน ปหาตพฺโพ ราโค อุปฺปชฺชติ, ทิฏฺฐิ อุปฺปชฺชติ, วิจิกิจฺฉา อุปฺปชฺชติ. ทสฺสเนน ปหาตพฺพํ โทมนสฺสํ อุปฺปชฺชติ, ทิฏฺฐิํ อสฺสาเทติ อภินนฺทติ, ตํ อารพฺภ ทสฺสเนน ปหาตพฺโพ ราโค อุปฺปชฺชติ, ทิฏฺฐิ อุปฺปชฺชติ, วิจิกิจฺฉา อุปฺปชฺชติ. ทสฺสเนน ปหาตพฺพํ โทมนสฺสํ อุปฺปชฺชติ, วิจิกิจฺฉํ อารพฺภ วิจิกิจฺฉา อุปฺปชฺชติ, ทิฏฺฐิ อุปฺปชฺชติ, ทสฺสเนน ปหาตพฺพํ โทมนสฺสํ อุปฺปชฺชติ. ทสฺสเนน ปหาตพฺพํ โทมนสฺสํ อารพฺภ ทสฺสเนน ปหาตพฺพํ โทมนสฺสํ อุปฺปชฺชติ, ทิฏฺฐิ อุปฺปชฺชติ, วิจิกิจฺฉา อุปฺปชฺชติ. (1)}

\prose{ทสฺสเนน ปหาตพฺโพ ธมฺโม เนวทสฺสเนน นภาวนาย ปหาตพฺพสฺส ธมฺมสฺส อารมฺมณ\-ปจฺจเยน ปจฺจโย. อริยา ทสฺสเนน ปหาตพฺเพ ปหีเน กิเลเส ปจฺจเวกฺขนฺติ, ปุพฺเพ สมุทาจิณฺเณ กิเลเส ชานนฺติ. ทสฺสเนน ปหาตพฺเพ ขนฺเธ อนิจฺจโต ทุกฺขโต อนตฺตโต วิปสฺสนฺติ. เจโตปริย\-ญาเณน ทสฺสเนน ปหาตพฺพ\-จิตฺต\-สมงฺคิสฺส จิตฺตํ ชานนฺติ. ทสฺสเนน ปหาตพฺพา ขนฺธา เจโตปริย\-ญาณสฺส, ปุพฺเพ\-นิวาสา\-นุสฺสติ\-ญาณสฺส, ยถากมฺมูปคญาณสฺส, อนาคตํสญาณสฺส, อาวชฺชนาย อารมฺมณ\-ปจฺจเยน ปจฺจโย. (2)}

\proseitem{72}{ภาวนาย ปหาตพฺโพ ธมฺโม ภาวนาย ปหาตพฺพสฺส ธมฺมสฺส อารมฺมณ\-ปจฺจเยน ปจฺจโย. ภาวนาย ปหาตพฺพํ ราคํ อสฺสาเทติ อภินนฺทติ, ตํ อารพฺภ ภาวนาย ปหาตพฺโพ ราโค อุปฺปชฺชติ, อุทฺธจฺจํ \csromanfolio{160}อุปฺปชฺชติ. ภาวนาย ปหาตพฺพํ โทมนสฺสํ อุปฺปชฺชติ, อุทฺธจฺจํ อารพฺภ อุทฺธจฺจํ อุปฺปชฺชติ. ภาวนาย ปหาตพฺพํ โทมนสฺสํ อุปฺปชฺชติ, ภาวนาย ปหาตพฺพํ โทมนสฺสํ อารพฺภ ภาวนาย ปหาตพฺพํ โทมนสฺสํ อุปฺปชฺชติ, อุทฺธจฺจํ อุปฺปชฺชติ. (1)}

\prose{ภาวนาย ปหาตพฺโพ ธมฺโม ทสฺสเนน ปหาตพฺพสฺส ธมฺมสฺส อารมฺมณ\-ปจฺจเยน ปจฺจโย. ภาวนาย ปหาตพฺพํ ราคํ อสฺสาเทติ อภินนฺทติ, ตํ อารพฺภ ทสฺสเนน ปหาตพฺโพ ราโค อุปฺปชฺชติ, ทิฏฺฐิ อุปฺปชฺชติ, วิจิกิจฺฉา อุปฺปชฺชติ. ทสฺสเนน ปหาตพฺพํ โทมนสฺสํ อุปฺปชฺชติ. อุทฺธจฺจํ อารพฺภ ทิฏฺฐิ อุปฺปชฺชติ, วิจิกิจฺฉา อุปฺปชฺชติ. ทสฺสเนน ปหาตพฺพํ โทมนสฺสํ อุปฺปชฺชติ. ภาวนาย ปหาตพฺพํ โทมนสฺสํ อารพฺภ ทสฺสเนน ปหาตพฺพํ โทมนสฺสํ อุปฺปชฺชติ, ทิฏฺฐิ อุปฺปชฺชติ, วิจิกิจฺฉา อุปฺปชฺชติ. (2)}

\prose{ภาวนาย ปหาตพฺโพ ธมฺโม เนวทสฺสเนน นภาวนาย ปหาตพฺพสฺส ธมฺมสฺส อารมฺมณ\-ปจฺจเยน ปจฺจโย. อริยา ภาวนาย ปหาตพฺเพ ปหีเน กิเลเส ปจฺจเวกฺขนฺติ, วิกฺขมฺภิเต กิเลเส ปจฺจเวกฺขนฺติ, ปุพฺเพ สมุทาจิณฺเณ กิเลเส ชานนฺติ. ภาวนาย ปหาตพฺเพ ขนฺเธ อนิจฺจโต ทุกฺขโต อนตฺตโต วิปสฺสนฺติ. เจโตปริย\-ญาเณน ภาวนาย ปหาตพฺพ\-จิตฺต\-สมงฺคิสฺส จิตฺตํ ชานนฺติ. ภาวนาย ปหาตพฺพา ขนฺธา เจโตปริย\-ญาณสฺส, ปุพฺเพ\-นิวาสา\-นุสฺสติ\-ญาณสฺส, ยถากมฺมูปคญาณสฺส, อนาคตํสญาณสฺส, อาวชฺชนาย อารมฺมณ\-ปจฺจเยน ปจฺจโย. (3)}

\proseitem{73}{เนวทสฺสเนน นภาวนาย ปหาตพฺโพ ธมฺโม เนวทสฺสเนน นภาวนาย ปหาตพฺพสฺส ธมฺมสฺส อารมฺมณ\-ปจฺจเยน ปจฺจโย. ทานํ ทตฺวา, สีลํ สมาทิยิตฺวา, อุโปสถกมฺมํ กตฺวา ตํ ปจฺจเวกฺขติ, ปุพฺเพ สุจิณฺณานิ ปจฺจเวกฺขติ, ฌานา วุฏฺฐหิตฺวา ฌานํ ปจฺจเวกฺขติ. อริยา มคฺคา วุฏฺฐหิตฺวา มคฺคํ ปจฺจเวกฺขนฺติ, ผลํ ปจฺจเวกฺขนฺติ, นิพฺพานํ ปจฺจเวกฺขนฺติ. นิพฺพานํ โคตฺรภุสฺส, โวทานสฺส, มคฺคสฺส, ผลสฺส, อาวชฺชนาย อารมฺมณ\-ปจฺจเยน ปจฺจโย. จกฺขุํ อนิจฺจโต ทุกฺขโต อนตฺตโต วิปสฺสติ. โสตํ. ฆานํ. ชิวฺหํ. กายํ. รูเป. สทฺเท. คนฺเธ. รเส. โผฏฺฐพฺเพ. วตฺถุํ. เนวทสฺสเนน นภาวนาย ปหาตพฺเพ ขนฺเธ อนิจฺจโต ทุกฺขโต อนตฺตโต วิปสฺสติ. ทิพฺเพน จกฺขุนา รูปํ ปสฺสติ. ทิพฺพาย โสตธาตุยา สทฺทํ สุณาติ. เจโตปริย\-ญาเณน เนวทสฺสเนน นภาวนาย ปหาตพฺพ\-จิตฺต\-สมงฺคิสฺส จิตฺตํ \csromanfolio{161}ชานนฺติ. อากาสา\-นญฺจา\-ยตนํ วิญฺญาณญฺจา\-ยตนสฺส ฯเปฯ. อากิญฺจญฺญา\-ยตนํ เนว\-สญฺญา\-นา\-สญฺญา\-ยตนสฺส ฯเปฯ. รูปายตนํ จกฺขุ\-วิญฺญาณสฺส อารมฺมณ\-ปจฺจเยน ปจฺจโย ฯเปฯ. โผฏฺฐพฺพา\-ยตนํ กายวิญฺญาณสฺส อารมฺมณ\-ปจฺจเยน ปจฺจโย ฯเปฯ. เนวทสฺสเนน นภาวนาย ปหาตพฺพา ขนฺธา อิทฺธิวิธ\-ญาณสฺส, เจโตปริย\-ญาณสฺส, ปุพฺเพ\-นิวาสา\-นุสฺสติ\-ญาณสฺส, ยถากมฺมูปคญาณสฺส, อนาคตํสญาณสฺส, อาวชฺชนาย อารมฺมณ\-ปจฺจเยน ปจฺจโย. (1)}

\prose{เนวทสฺสเนน นภาวนาย ปหาตพฺโพ ธมฺโม ทสฺสเนน ปหาตพฺพสฺส ธมฺมสฺส อารมฺมณ\-ปจฺจเยน ปจฺจโย. ทานํ ทตฺวา, สีลํ สมาทิยิตฺวา, อุโปสถกมฺมํ กตฺวา ตํ อสฺสาเทติ อภินนฺทติ, ตํ อารพฺภ ทสฺสเนน ปหาตพฺโพ ราโค อุปฺปชฺชติ, ทิฏฺฐิ อุปฺปชฺชติ, วิจิกิจฺฉา อุปฺปชฺชติ, ทสฺสเนน ปหาตพฺพํ โทมนสฺสํ อุปฺปชฺชติ. ปุพฺเพ สุจิณฺณานิ อสฺสาเทติ อภินนฺทติ ฯเปฯ ฌานา วุฏฺฐหิตฺวา ฌานํ อสฺสาเทติ อภินนฺทติ, ตํ อารพฺภ ทสฺสเนน ปหาตพฺโพ ราโค ฯเปฯ ทิฏฺฐิ ฯเปฯ วิจิกิจฺฉา ฯเปฯ ฌาเน ปริหีเน วิปฺปฏิสาริสฺส ทสฺสเนน ปหาตพฺพํ โทมนสฺสํ อุปฺปชฺชติ. จกฺขุํ อสฺสาเทติ อภินนฺทติ ฯเปฯ. โสตํ. ฆานํ. ชิวฺหํ. กายํ. รูเป. สทฺเท. คนฺเธ. รเส. โผฏฺฐพฺเพ. วตฺถุํ. เนวทสฺสเนน นภาวนาย ปหาตพฺเพ ขนฺเธ อสฺสาเทติ อภินนฺทติ, ตํ อารพฺภ ทสฺสเนน ปหาตพฺโพ ราโค ฯเปฯ ทิฏฺฐิ ฯเปฯ วิจิกิจฺฉา ฯเปฯ ทสฺสเนน ปหาตพฺพํ โทมนสฺสํ อุปฺปชฺชติ. (2)}

\prose{เนวทสฺสเนน นภาวนาย ปหาตพฺโพ ธมฺโม ภาวนาย ปหาตพฺพสฺส ธมฺมสฺส อารมฺมณ\-ปจฺจเยน ปจฺจโย. ทานํ ทตฺวา, สีลํ สมาทิยิตฺวา, อุโปสถกมฺมํ กตฺวา ตํ อสฺสาเทติ อภินนฺทติ, ตํ อารพฺภ ภาวนาย ปหาตพฺโพ ราโค อุปฺปชฺชติ, อุทฺธจฺจํ อุปฺปชฺชติ, ภาวนาย ปหาตพฺพํ โทมนสฺสํ อุปฺปชฺชติ. ปุพฺเพ สุจิณฺณานิ ฯเปฯ ฌานา วุฏฺฐหิตฺวา ฯเปฯ จกฺขุํ ฯเปฯ วตฺถุํ. เนวทสฺสเนน นภาวนาย ปหาตพฺเพ ขนฺเธ อสฺสาเทติ ฯเปฯ ตํ อารพฺภ ภาวนาย ปหาตพฺโพ ราโค อุปฺปชฺชติ, อุทฺธจฺจํ อุปฺปชฺชติ. ภาวนาย ปหาตพฺพํ โทมนสฺสํ อุปฺปชฺชติ. (3)}

\proseitem{74}{ทสฺสเนน ปหาตพฺโพ ธมฺโม ทสฺสเนน ปหาตพฺพสฺส ธมฺมสฺส อธิปติ\-ปจฺจเยน ปจฺจโย. อารมฺมณา\-ธิปติ, สหชาตา\-ธิปติ.   \csromanfolio{162}อารมฺมณา\-ธิปติ\csromandash{}ทสฺสเนน ปหาตพฺพํ ราคํ ครุํ กตฺวา อสฺสาเทติ อภินนฺทติ, ตํ ครุํ กตฺวา ทสฺสเนน ปหาตพฺโพ ราโค อุปฺปชฺชติ, ทิฏฺฐิ อุปฺปชฺชติ. ทิฏฺฐิํ ครุํ กตฺวา อสฺสาเทติ อภินนฺทติ, ตํ ครุํ กตฺวา ทสฺสเนน ปหาตพฺโพ ราโค อุปฺปชฺชติ, ทิฏฺฐิ อุปฺปชฺชติ.  สหชาตา\-ธิปติ\csromandash{} ทสฺสเนน ปหาตพฺพา\-ธิปติ สมฺปยุตฺตกานํ ขนฺธานํ อธิปติ\-ปจฺจเยน ปจฺจโย. (1)}

\prose{ทสฺสเนน ปหาตพฺโพ ธมฺโม เนวทสฺสเนน นภาวนาย ปหาตพฺพสฺส ธมฺมสฺส อธิปติ\-ปจฺจเยน ปจฺจโย. สหชาตา\-ธิปติ\csromandash{} ทสฺสเนน ปหาตพฺพา\-ธิปติ จิตฺต\-สมุฏฺฐานานํ รูปานํ อธิปติ\-ปจฺจเยน ปจฺจโย. (2)}

\prose{ทสฺสเนน ปหาตพฺโพ ธมฺโม ทสฺสเนน ปหาตพฺพสฺส จ เนวทสฺสเนน นภาวนาย ปหาตพฺพสฺส จ ธมฺมสฺส อธิปติ\-ปจฺจเยน ปจฺจโย. สหชาตา\-ธิปติ\csromandash{}ทสฺสเนน ปหาตพฺพา\-ธิปติ สมฺปยุตฺตกานํ ขนฺธานํ จิตฺตสมุฏฺฐานานญฺจ รูปานํ อธิปติ\-ปจฺจเยน ปจฺจโย. (3)}

\proseitem{75}{ภาวนาย ปหาตพฺโพ ธมฺโม ภาวนาย ปหาตพฺพสฺส ธมฺมสฺส อธิปติ\-ปจฺจเยน ปจฺจโย. อารมฺมณา\-ธิปติ, สหชาตา\-ธิปติ.  อารมฺมณา\-ธิปติ\csromandash{}ภาวนาย ปหาตพฺพํ ราคํ ครุํ กตฺวา อสฺสาเทติ อภินนฺทติ, ตํ ครุํ กตฺวา ภาวนาย ปหาตพฺโพ ราโค อุปฺปชฺชติ.  สหชาตา\-ธิปติ\csromandash{}ภาวนาย ปหาตพฺพา\-ธิปติ สมฺปยุตฺตกานํ ขนฺธานํ อธิปติ\-ปจฺจเยน ปจฺจโย. (1)}

\prose{ภาวนาย ปหาตพฺโพ ธมฺโม ทสฺสเนน ปหาตพฺพสฺส ธมฺมสฺส อธิปติ\-ปจฺจเยน ปจฺจโย. อารมฺมณา\-ธิปติ\csromandash{}ภาวนาย ปหาตพฺพํ ราคํ ครุํ กตฺวา อสฺสาเทติ อภินนฺทติ, ตํ ครุํ กตฺวา ทสฺสเนน ปหาตพฺโพ ราโค อุปฺปชฺชติ, ทิฏฺฐิ อุปฺปชฺชติ. (2)}

\prose{ภาวนาย ปหาตพฺโพ ธมฺโม เนวทสฺสเนน นภาวนาย ปหาตพฺพสฺส ธมฺมสฺส อธิปติ\-ปจฺจเยน ปจฺจโย. สหชาตา\-ธิปติ\csromandash{} ภาวนาย ปหาตพฺพา\-ธิปติ จิตฺต\-สมุฏฺฐานานํ รูปานํ อธิปติ\-ปจฺจเยน ปจฺจโย. (3)}

\prose{ภาวนาย ปหาตพฺโพ ธมฺโม ภาวนาย ปหาตพฺพสฺส จ เนวทสฺสเนน นภาวนาย ปหาตพฺพสฺส จ ธมฺมสฺส อธิปติ\-ปจฺจเยน ปจฺจโย. สหชาตา\-ธิปติ\csromandash{}ภาวนาย ปหาตพฺพา\-ธิปติ สมฺปยุตฺตกานํ ขนฺธานํ จิตฺตสมุฏฺฐานานญฺจ รูปานํ อธิปติ\-ปจฺจเยน ปจฺจโย. (4)}

\proseitem{76}{\csromanfolio{163}เนวทสฺสเนน นภาวนาย ปหาตพฺโพ ธมฺโม เนวทสฺสเนน นภาวนาย ปหาตพฺพสฺส ธมฺมสฺส อธิปติ\-ปจฺจเยน ปจฺจโย. อารมฺมณา\-ธิปติ, สหชาตา\-ธิปติ.  อารมฺมณา\-ธิปติ\csromandash{}ทานํ ทตฺวา, สีลํ สมาทิยิตฺวา, อุโปสถกมฺมํ กตฺวา ตํ ครุํ กตฺวา ปจฺจเวกฺขติ. ปุพฺเพ สุจิณฺณานิ ฯเปฯ ฌานา วุฏฺฐหิตฺวา ฌานํ ครุํ กตฺวา ปจฺจเวกฺขติ. อริยา มคฺคา วุฏฺฐหิตฺวา มคฺคํ ครุํ กตฺวา ปจฺจเวกฺขนฺติ, ผลํ ครุํ กตฺวา ปจฺจเวกฺขนฺติ, นิพฺพานํ ครุํ กตฺวา ปจฺจเวกฺขนฺติ, นิพฺพานํ โคตฺรภุสฺส, โวทานสฺส, มคฺคสฺส, ผลสฺส อธิปติ\-ปจฺจเยน ปจฺจโย.  สหชาตา\-ธิปติ\csromandash{}เนวทสฺสเนน นภาวนาย ปหาตพฺพา\-ธิปติ สมฺปยุตฺตกานํ ขนฺธานํ จิตฺตสมุฏฺฐานานญฺจ รูปานํ อธิปติ\-ปจฺจเยน ปจฺจโย. (1)}

\prose{เนวทสฺสเนน นภาวนาย ปหาตพฺโพ ธมฺโม ทสฺสเนน ปหาตพฺพสฺส ธมฺมสฺส อธิปติ\-ปจฺจเยน ปจฺจโย. อารมฺมณา\-ธิปติ\csromandash{} ทานํ ทตฺวา, สีลํ สมาทิยิตฺวา, อุโปสถกมฺมํ กตฺวา ตํ ครุํ กตฺวา อสฺสาเทติ อภินนฺทติ, ตํ ครุํ กตฺวา ทสฺสเนน ปหาตพฺโพ ราโค อุปฺปชฺชติ, ทิฏฺฐิ อุปฺปชฺชติ. ปุพฺเพ สุจิณฺณานิ ครุํ กตฺวา ฯเปฯ. ฌานา วุฏฺฐหิตฺวา ฌานํ ครุํ กตฺวา ฯเปฯ. จกฺขุํ ฯเปฯ. วตฺถุํ เนวทสฺสเนน นภาวนาย ปหาตพฺเพ ขนฺเธ ครุํ กตฺวา อสฺสาเทติ อภินนฺทติ, ตํ ครุํ กตฺวา ทสฺสเนน ปหาตพฺโพ ราโค อุปฺปชฺชติ, ทิฏฺฐิ อุปฺปชฺชติ. (2)}

\prose{เนวทสฺสเนน นภาวนาย ปหาตพฺโพ ธมฺโม ภาวนาย ปหาตพฺพสฺส ธมฺมสฺส อธิปติ\-ปจฺจเยน ปจฺจโย. อารมฺมณา\-ธิปติ\csromandash{} ทานํ ทตฺวา, สีลํ สมาทิยิตฺวา, อุโปสถกมฺมํ กตฺวา ตํ ครุํ กตฺวา อสฺสาเทติ อภินนฺทติ, ตํ ครุํ กตฺวา ภาวนาย ปหาตพฺโพ ราโค อุปฺปชฺชติ. ปุพฺเพ ฯเปฯ. เนวทสฺสเนน นภาวนาย ปหาตพฺเพ ขนฺเธ ครุํ กตฺวา อสฺสาเทติ อภินนฺทติ, ตํ ครุํ กตฺวา ภาวนาย ปหาตพฺโพ ราโค อุปฺปชฺชติ. (3)}

\prose{\textbf{อนนฺตร}}

\proseitem{77}{ทสฺสเนน ปหาตพฺโพ ธมฺโม ทสฺสเนน ปหาตพฺพสฺส ธมฺมสฺส อนนฺตร\-ปจฺจเยน ปจฺจโย. ปุริมา ปุริมา ทสฺสเนน ปหาตพฺพา ขนฺธา ปจฺฉิมานํ ปจฺฉิมานํ ทสฺสเนน ปหาตพฺพานํ ขนฺธานํ อนนฺตร\-ปจฺจเยน ปจฺจโย.}

\prose{\csromanfolio{164}ทสฺสเนน ปหาตพฺโพ ธมฺโม เนวทสฺสเนน นภาวนาย ปหาตพฺพสฺส ธมฺมสฺส อนนฺตร\-ปจฺจเยน ปจฺจโย. ทสฺสเนน ปหาตพฺพา ขนฺธา วุฏฺฐานสฺส อนนฺตร\-ปจฺจเยน ปจฺจโย. (2)}

\prose{ภาวนาย ปหาตพฺโพ ธมฺโม ภาวนาย ปหาตพฺพสฺส ธมฺมสฺส อนนฺตร\-ปจฺจเยน ปจฺจโย. ปุริมา ปุริมา ภาวนาย ปหาตพฺพา ขนฺธา ปจฺฉิมานํ ปจฺฉิมานํ ภาวนาย ปหาตพฺพานํ ขนฺธานํ อนนฺตร\-ปจฺจเยน ปจฺจโย. (1)}

\prose{ภาวนาย ปหาตพฺโพ ธมฺโม เนวทสฺสเนน นภาวนาย ปหาตพฺพสฺส ธมฺมสฺส อนนฺตร\-ปจฺจเยน ปจฺจโย. ภาวนาย ปหาตพฺพา ขนฺธา วุฏฺฐานสฺส อนนฺตร\-ปจฺจเยน ปจฺจโย. (2)}

\prose{เนวทสฺสเนน นภาวนาย ปหาตพฺโพ ธมฺโม เนวทสฺสเนน นภาวนาย ปหาตพฺพสฺส ธมฺมสฺส อนนฺตร\-ปจฺจเยน ปจฺจโย. ปุริมา ปุริมา เนวทสฺสเนน นภาวนาย ปหาตพฺพา ขนฺธา ปจฺฉิมานํ ปจฺฉิมานํ เนวทสฺสเนน นภาวนาย ปหาตพฺพานํ ขนฺธานํ อนนฺตร\-ปจฺจเยน ปจฺจโย. อนุโลมํ โคตฺรภุสฺส. อนุโลมํ โวทานสฺส. โคตฺรภุ มคฺคสฺส. โวทานํ มคฺคสฺส. มคฺโค ผลสฺส. ผลํ ผลสฺส. อนุโลมํ ผลสมาปตฺติยา. นิโรธา วุฏฺฐหนฺตสฺส เนว\-สญฺญา\-นา\-สญฺญา\-ยตนํ ผลสมาปตฺติยา อนนฺตร\-ปจฺจเยน ปจฺจโย. (1)}

\prose{เนวทสฺสเนน นภาวนาย ปหาตพฺโพ ธมฺโม ทสฺสเนน ปหาตพฺพสฺส ธมฺมสฺส ฯเปฯ. อาวชฺชนา ทสฺสเนน ปหาตพฺพานํ ขนฺธานํ อนนฺตร\-ปจฺจเยน ปจฺจโย. (2)}

\prose{เนวทสฺสเนน นภาวนาย ปหาตพฺโพ ธมฺโม ภาวนาย ปหาตพฺพสฺส ธมฺมสฺส อนนฺตร\-ปจฺจเยน ปจฺจโย. อาวชฺชนา ภาวนาย ปหาตพฺพานํ ขนฺธานํ อนนฺตร\-ปจฺจเยน ปจฺจโย. (3)}

\csromanheader{2}{\textbf{สมนนฺตร}}
\proseitem{78}{ทสฺสเนน ปหาตพฺโพ ธมฺโม ทสฺสเนน ปหาตพฺพสฺส ธมฺมสฺส สมนนฺตร\-ปจฺจเยน ปจฺจโย. (อนนฺตร\-สทิสํ.)}

\csromanheader{2}{8. ทสฺสเนน\-ปหาตพฺพ\-ตฺติก สหชาต}
\proseitem{79}{\csromanfolio{165}ทสฺสเนน ปหาตพฺโพ ธมฺโม ทสฺสเนน ปหาตพฺพสฺส ธมฺมสฺส สหชาต\-ปจฺจเยน ปจฺจโย. ตีณิ.}

\prose{ภาวนาย ปหาตพฺโพ ธมฺโม ภาวนาย ปหาตพฺพสฺส ธมฺมสฺส. ตีณิ.}

\prose{เนวทสฺสเนน นภาวนาย ปหาตพฺโพ ธมฺโม เนวทสฺสเนน นภาวนาย ปหาตพฺพสฺส ธมฺมสฺส สหชาต\-ปจฺจเยน ปจฺจโย. เนวทสฺสเนน นภาวนาย ปหาตพฺโพ เอโก ขนฺโธ ติณฺณนฺนํ ขนฺธานํ จิตฺตสมุฏฺฐานานญฺจ รูปานํ สหชาต\-ปจฺจเยน ปจฺจโย ฯเปฯ เทฺว ขนฺธา ฯเปฯ. ปฏิสนฺธิ\-กฺขเณ ฯเปฯ. ขนฺธา วตฺถุสฺส สหชาต\-ปจฺจเยน ปจฺจโย. วตฺถุ ขนฺธานํ ฯเปฯ. เอกํ มหาภูตํ ติณฺณนฺนํ มหาภูตานํ ฯเปฯ มหาภูตา จิตฺต\-สมุฏฺฐานานํ รูปานํ ฯเปฯ. พาหิรํ. อาหารสมุฏฺฐานํ. อุตุสมุฏฺฐานํ. อสญฺญสตฺตานํ ฯเปฯ. (1)}

\prose{ทสฺสเนน ปหาตพฺโพ จ เนวทสฺสเนน นภาวนาย ปหาตพฺโพ จ ธมฺมา เนวทสฺสเนน นภาวนาย ปหาตพฺพสฺส ธมฺมสฺส สหชาต\-ปจฺจเยน ปจฺจโย. ทสฺสเนน ปหาตพฺพา ขนฺธา จ มหาภูตา จ จิตฺต\-สมุฏฺฐานานํ รูปานํ สหชาต\-ปจฺจเยน ปจฺจโย. (1)}

\prose{ภาวนาย ปหาตพฺโพ จ เนวทสฺสเนน นภาวนาย ปหาตพฺโพ จ ธมฺมา เนวทสฺสเนน นภาวนาย ปหาตพฺพสฺส ธมฺมสฺส สหชาต\-ปจฺจเยน ปจฺจโย. ภาวนาย ปหาตพฺพา ขนฺธา จ มหาภูตา จ จิตฺต\-สมุฏฺฐานานํ รูปานํ สหชาต\-ปจฺจเยน ปจฺจโย. (1)}

\csromanheader{2}{\textbf{อญฺญมญฺญ}}
\proseitem{80}{ทสฺสเนน ปหาตพฺโพ ธมฺโม ทสฺสเนน ปหาตพฺพสฺส ธมฺมสฺส อญฺญมญฺญ\-ปจฺจเยน ปจฺจโย. ทสฺสเนน ปหาตพฺโพ เอโก ขนฺโธ ติณฺณนฺนํ ขนฺธานํ ฯเปฯ. (1)}

\prose{ภาวนาย ปหาตพฺโพ ธมฺโม ภาวนาย ปหาตพฺพสฺส ธมฺมสฺส ฯเปฯ. ภาวนาย ปหาตพฺโพ เอโก ขนฺโธ ติณฺณนฺนํ ขนฺธานํ ฯเปฯ. (1)}

\prose{\csromanfolio{166}เนวทสฺสเนน นภาวนาย ปหาตพฺโพ ธมฺโม เนวทสฺสเนน นภาวนาย ปหาตพฺพสฺส ธมฺมสฺส ฯเปฯ. เนวทสฺสเนน นภาวนาย ปหาตพฺโพ เอโก ขนฺโธ ติณฺณนฺนํ ขนฺธานํ อญฺญมญฺญ\-ปจฺจเยน ปจฺจโย ฯเปฯ เทฺว ขนฺธา ฯเปฯ. ปฏิสนฺธิ\-กฺขเณ เนวทสฺสเนน นภาวนาย ปหาตพฺโพ เอโก ขนฺโธ ติณฺณนฺนํ ขนฺธานํ วตฺถุสฺส จ อญฺญมญฺญ\-ปจฺจเยน ปจฺจโย ฯเปฯ เทฺว ขนฺธา ฯเปฯ. ขนฺธา วตฺถุสฺส ฯเปฯ. วตฺถุ ขนฺธานํ ฯเปฯ. เอกํ มหาภูตํ ติณฺณนฺนํ มหาภูตานํ อญฺญมญฺญ\-ปจฺจเยน ปจฺจโย ฯเปฯ. อสญฺญสตฺตานํ ฯเปฯ. (1)}

\csromanheader{2}{\textbf{นิสฺสย}}
\proseitem{81}{ทสฺสเนน ปหาตพฺโพ ธมฺโม ทสฺสเนน ปหาตพฺพสฺส ธมฺมสฺส นิสฺสย\-ปจฺจเยน ปจฺจโย. ตีณิ.}

\prose{ภาวนาย ปหาตพฺโพ ธมฺโม ภาวนาย ปหาตพฺพสฺส ธมฺมสฺส. ตีณิ.}

\prose{เนวทสฺสเนน นภาวนาย ปหาตพฺโพ ธมฺโม เนวทสฺสเนน นภาวนาย ปหาตพฺพสฺส ธมฺมสฺส นิสฺสย\-ปจฺจเยน ปจฺจโย. เนวทสฺสเนน นภาวนาย ปหาตพฺโพ เอโก ขนฺโธ ติณฺณนฺนํ ขนฺธานํ จิตฺตสมุฏฺฐานานญฺจ รูปานํ ฯเปฯ เทฺว ขนฺธา ฯเปฯ. ปฏิสนฺธิ\-กฺขเณ ฯเปฯ. ขนฺธา วตฺถุสฺส ฯเปฯ. วตฺถุ ขนฺธานํ ฯเปฯ. เอกํ มหาภูตํ ฯเปฯ. อสญฺญสตฺตานํ ฯเปฯ. จกฺขายตนํ จกฺขุ\-วิญฺญาณสฺส ฯเปฯ. กายายตนํ กายวิญฺญาณสฺส ฯเปฯ. วตฺถุ ฯเปฯ. (1)}

\prose{เนวทสฺสเนน นภาวนาย ปหาตพฺโพ ธมฺโม ทสฺสเนน ปหาตพฺพสฺส ธมฺมสฺส นิสฺสย\-ปจฺจเยน ปจฺจโย. วตฺถุ ทสฺสเนน ปหาตพฺพานํ ขนฺธานํ นิสฺสย\-ปจฺจเยน ปจฺจโย. (2)}

\prose{เนวทสฺสเนน นภาวนาย ปหาตพฺโพ ธมฺโม ภาวนาย ปหาตพฺพสฺส ธมฺมสฺส นิสฺสย\-ปจฺจเยน ปจฺจโย. วตฺถุ ภาวนาย ปหาตพฺพานํ ขนฺธานํ นิสฺสย\-ปจฺจเยน ปจฺจโย. (3)}

\proseitem{82}{ทสฺสเนน ปหาตพฺโพ จ เนวทสฺสเนน นภาวนาย ปหาตพฺโพ จ ธมฺมา ทสฺสเนน ปหาตพฺพสฺส ธมฺมสฺส นิสฺสย\-ปจฺจเยน ปจฺจโย. ทสฺสเนน ปหาตพฺโพ เอโก ขนฺโธ จ วตฺถุ จ ติณฺณนฺนํ ขนฺธานํ นิสฺสย\-ปจฺจเยน ปจฺจโย ฯเปฯ เทฺว ขนฺธา ฯเปฯ. (1)}

\prose{ทสฺสเนน ปหาตพฺโพ จ เนวทสฺสเนน นภาวนาย ปหาตพฺโพ จ ธมฺมา เนวทสฺสเนน นภาวนาย ปหาตพฺพสฺส ธมฺมสฺส นิสฺสย\-ปจฺจเยน \csromanfolio{167}ปจฺจโย. ทสฺสเนน ปหาตพฺพา ขนฺธา จ มหาภูตา จ จิตฺต\-สมุฏฺฐานานํ รูปานํ นิสฺสย\-ปจฺจเยน ปจฺจโย. (2)}

\prose{ภาวนาย ปหาตพฺโพ จ เนวทสฺสเนน นภาวนาย ปหาตพฺโพ จ ธมฺมา ภาวนาย ปหาตพฺพสฺส ธมฺมสฺส นิสฺสย\-ปจฺจเยน ปจฺจโย. ภาวนาย ปหาตพฺโพ เอโก ขนฺโธ จ วตฺถุ จ ติณฺณนฺนํ ขนฺธานํ นิสฺสย\-ปจฺจเยน ปจฺจโย ฯเปฯ เทฺว ขนฺธา ฯเปฯ. (1)}

\prose{ภาวนาย ปหาตพฺโพ จ เนวทสฺสเนน น ภาวนาย ปหาตพฺโพ จ ธมฺมา เนวทสฺสเนน นภาวนาย ปหาตพฺพสฺส ธมฺมสฺส นิสฺสย\-ปจฺจเยน ปจฺจโย. ภาวนาย ปหาตพฺพา ขนฺธา จ มหาภูตา จ จิตฺต\-สมุฏฺฐานานํ รูปานํ นิสฺสย\-ปจฺจเยน ปจฺจโย. (2)}

\csromanheader{2}{\textbf{อุปนิสฺสย}}
\proseitem{83}{ทสฺสเนน ปหาตพฺโพ ธมฺโม ทสฺสเนน ปหาตพฺพสฺส ธมฺมสฺส อุปนิสฺสย\-ปจฺจเยน ปจฺจโย. อารมฺมณู\-ปนิสฺสโย, อนนฺตรู\-ปนิสฺสโย, ปกตู\-ปนิสฺสโย ฯเปฯ. ปกตู\-ปนิสฺสโย\csromandash{}ทสฺสเนน ปหาตพฺพํ ราคํ อุปนิสฺสาย ปาณํ หนติ, อทินฺนํ อาทิยติ ฯเปฯ สํฆํ ภินฺทติ. ทสฺสเนน ปหาตพฺพํ โทสํ. โมหํ. ทิฏฺฐิํ. ปตฺถนํ อุปนิสฺสาย ปาณํ หนติ ฯเปฯ สํฆํ ภินฺทติ. ทสฺสเนน ปหาตพฺโพ ราโค. โทโส. โมโห. ทิฏฺฐิ. ปตฺถนา ทสฺสเนน ปหาตพฺพสฺส ราคสฺส. โทสสฺส. โมหสฺส. ทิฏฺฐิยา. ปตฺถนาย อุปนิสฺสย\-ปจฺจเยน ปจฺจโย. (1)}

\prose{ทสฺสเนน ปหาตพฺโพ ธมฺโม เนวทสฺสเนน นภาวนาย ปหาตพฺพสฺส ธมฺมสฺส อุปนิสฺสย\-ปจฺจเยน ปจฺจโย. อนนฺตรู\-ปนิสฺสโย, ปกตู\-ปนิสฺสโย ฯเปฯ. ปกตู\-ปนิสฺสโย\csromandash{}ทสฺสเนน ปหาตพฺพํ ราคํ อุปนิสฺสาย ทานํ เทติ, สีลํ สมาทิยติ, อุโปสถกมฺมํ กโรติ ฯเปฯ สมาปตฺติํ อุปฺปาเทติ. ทสฺสเนน ปหาตพฺพํ โทสํ ฯเปฯ. ปตฺถนํ อุปนิสฺสาย ทานํ เทติ ฯเปฯ สมาปตฺติํ อุปฺปาเทติ. ทสฺสเนน ปหาตพฺโพ ราโค. โทโส. โมโห. ทิฏฺฐิ. ปตฺถนา สทฺธาย ฯเปฯ ปญฺญาย. กายิกสฺส สุขสฺส, กายิกสฺส ทุกฺขสฺส, ผลสมาปตฺติยา อุปนิสฺสย\-ปจฺจเยน ปจฺจโย. (2)}

\proseitem{84}{\csromanfolio{168}ภาวนาย ปหาตพฺโพ ธมฺโม ภาวนาย ปหาตพฺพสฺส ธมฺมสฺส อุปนิสฺสย\-ปจฺจเยน ปจฺจโย. อารมฺมณู\-ปนิสฺสโย, อนนฺตรู\-ปนิสฺสโย, ปกตู\-ปนิสฺสโย ฯเปฯ. ปกตู\-ปนิสฺสโย\csromandash{}ภาวนาย ปหาตพฺโพ ราโค. โทโส. โมโห. มาโน. ปตฺถนา ภาวนาย ปหาตพฺพสฺส ราคสฺส. โทสสฺส. โมหสฺส. มานสฺส. ปตฺถนาย อุปนิสฺสย\-ปจฺจเยน ปจฺจโย. (1)}

\prose{ภาวนาย ปหาตพฺโพ ธมฺโม ทสฺสเนน ปหาตพฺพสฺส ธมฺมสฺส อุปนิสฺสย\-ปจฺจเยน ปจฺจโย. อารมฺมณู\-ปนิสฺสโย, ปกตู\-ปนิสฺสโย ฯเปฯ. ปกตู\-ปนิสฺสโย\csromandash{}ภาวนาย ปหาตพฺพํ ราคํ อุปนิสฺสาย ปาณํ หนติ ฯเปฯ สํฆํ ภินฺทติ. ภาวนาย ปหาตพฺพํ โทสํ. โมหํ. มานํ. ปตฺถนํ อุปนิสฺสาย ปาณํ หนติ ฯเปฯ สํฆํ ภินฺทติ. ภาวนาย ปหาตพฺโพ ราโค. โทโส. โมโห. มาโน. ปตฺถนา ทสฺสเนน ปหาตพฺพสฺส ราคสฺส. โทสสฺส. โมหสฺส. ทิฏฺฐิยา. ปตฺถนาย อุปนิสฺสย\-ปจฺจเยน ปจฺจโย. สกภณฺเฑ ฉนฺทราโค ปรภณฺเฑ ฉนฺทราคสฺส อุปนิสฺสย\-ปจฺจเยน ปจฺจโย. สกปริคฺคเห ฉนฺทราโค ปรปริคฺคเห ฉนฺทราคสฺส อุปนิสฺสย\-ปจฺจเยน ปจฺจโย. (2)}

\prose{ภาวนาย ปหาตพฺโพ ธมฺโม เนวทสฺสเนน นภาวนาย ปหาตพฺพสฺส ธมฺมสฺส อุปนิสฺสย\-ปจฺจเยน ปจฺจโย. อนนฺตรู\-ปนิสฺสโย, ปกตู\-ปนิสฺสโย ฯเปฯ. ปกตู\-ปนิสฺสโย\csromandash{}ภาวนาย ปหาตพฺพํ ราคํ อุปนิสฺสาย ทานํ เทติ ฯเปฯ สมาปตฺติํ อุปฺปาเทติ. ภาวนาย ปหาตพฺพํ โทสํ. โมหํ. มานํ. ปตฺถนํ อุปนิสฺสาย ทานํ เทติ ฯเปฯ สมาปตฺติํ อุปฺปาเทติ. ภาวนาย ปหาตพฺโพ ราโค. โทโส. โมโห. มาโน. ปตฺถนา สทฺธาย ฯเปฯ. ผลสมาปตฺติยา อุปนิสฺสย\-ปจฺจเยน ปจฺจโย. (3)}

\proseitem{85}{เนวทสฺสเนน นภาวนาย ปหาตพฺโพ ธมฺโม เนวทสฺสเนน นภาวนาย ปหาตพฺพสฺส ธมฺมสฺส อุปนิสฺสย\-ปจฺจเยน ปจฺจโย. อารมฺมณู\-ปนิสฺสโย, อนนฺตรู\-ปนิสฺสโย, ปกตู\-ปนิสฺสโย ฯเปฯ. ปกตู\-ปนิสฺสโย\csromandash{}สทฺธํ อุปนิสฺสาย ทานํ เทติ ฯเปฯ สมาปตฺติํ อุปฺปาเทติ. สีลํ. สุตํ. จาคํ. ปญฺญํ. กายิกํ สุขํ. กายิกํ ทุกฺขํ. อุตุํ. โภชนํ. เสนาสนํ อุปนิสฺสาย ทานํ เทติ ฯเปฯ สมาปตฺติํ อุปฺปาเทติ. สทฺธา. สีลํ. สุตํ. \csromanfolio{169}จาโค. ปญฺญา. กายิกํ สุขํ. กายิกํ ทุกฺขํ. อุตุ. โภชนํ. เสนาสนํ สทฺธาย ฯเปฯ ปญฺญาย. กายิกสฺส สุขสฺส, กายิกสฺส ทุกฺขสฺส, ผลสมาปตฺติยา อุปนิสฺสย\-ปจฺจเยน ปจฺจโย. (1)}

\prose{เนวทสฺสเนน นภาวนาย ปหาตพฺโพ ธมฺโม ทสฺสเนน ปหาตพฺพสฺส ธมฺมสฺส อุปนิสฺสย\-ปจฺจเยน ปจฺจโย. อารมฺมณู\-ปนิสฺสโย, อนนฺตรู\-ปนิสฺสโย, ปกตู\-ปนิสฺสโย ฯเปฯ. ปกตู\-ปนิสฺสโย\csromandash{}สทฺธํ อุปนิสฺสาย ทิฏฺฐิํ คณฺหาติ. สีลํ ฯเปฯ. เสนาสนํ อุปนิสฺสาย ปาณํ หนติ ฯเปฯ สํฆํ ภินฺทติ. สทฺธา ฯเปฯ. เสนาสนํ ทสฺสเนน ปหาตพฺพสฺส ราคสฺส. โทสสฺส. โมหสฺส. ทิฏฺฐิยา. ปตฺถนาย อุปนิสฺสย\-ปจฺจเยน ปจฺจโย. (2)}

\prose{เนวทสฺสเนน นภาวนาย ปหาตพฺโพ ธมฺโม ภาวนาย ปหาตพฺพสฺส ธมฺมสฺส อุปนิสฺสย\-ปจฺจเยน ปจฺจโย. อารมฺมณู\-ปนิสฺสโย, อนนฺตรู\-ปนิสฺสโย, ปกตู\-ปนิสฺสโย ฯเปฯ. ปกตู\-ปนิสฺสโย\csromandash{}สทฺธํ อุปนิสฺสาย มานํ ชปฺเปติ. สีลํ ฯเปฯ. ปญฺญํ. กายิกํ สุขํ. กายิกํ ทุกฺขํ. อุตุํ. โภชนํ. เสนาสนํ อุปนิสฺสาย มานํ ชปฺเปติ. สทฺธา ฯเปฯ. ปญฺญา. กายิกํ สุขํ. กายิกํ ทุกฺขํ. อุตุ. โภชนํ. เสนาสนํ ภาวนาย ปหาตพฺพสฺส ราคสฺส. โทสสฺส. โมหสฺส. มานสฺส. ปตฺถนาย อุปนิสฺสย\-ปจฺจเยน ปจฺจโย. (3)}

\csromanheader{2}{\textbf{ปุเรชาต}}
\proseitem{86}{เนวทสฺสเนน นภาวนาย ปหาตพฺโพ ธมฺโม เนวทสฺสเนน นภาวนาย ปหาตพฺพสฺส ธมฺมสฺส ปุเรชาต\-ปจฺจเยน ปจฺจโย. อารมฺมณ\-ปุเรชาตํ, วตฺถุ\-ปุเรชาตํ.  อารมฺมณ\-ปุเรชาตํ\csromandash{}จกฺขุํ อนิจฺจโต ทุกฺขโต อนตฺตโต วิปสฺสติ. โสตํ. ฆานํ. ชิวฺหํ. กายํ. รูเป. สทฺเท. คนฺเธ. รเส. โผฏฺฐพฺเพ. วตฺถุํ ฯเปฯ. ทิพฺเพน จกฺขุนา รูปํ ปสฺสติ. ทิพฺพาย โสตธาตุยา สทฺทํ สุณาติ. รูปายตนํ จกฺขุ\-วิญฺญาณสฺส ฯเปฯ. โผฏฺฐพฺพา\-ยตนํ กายวิญฺญาณสฺส ปุเรชาต\-ปจฺจเยน ปจฺจโย.  วตฺถุ\-ปุเรชาตํ\csromandash{}จกฺขายตนํ จกฺขุ\-วิญฺญาณสฺส ฯเปฯ. กายายตนํ กายวิญฺญาณสฺส ฯเปฯ. วตฺถุ เนวทสฺสเนน นภาวนาย ปหาตพฺพานํ ขนฺธานํ ปุเรชาต\-ปจฺจเยน ปจฺจโย. (1)}

\prose{\csromanfolio{170}เนวทสฺสเนน นภาวนาย ปหาตพฺโพ ธมฺโม ทสฺสเนน ปหาตพฺพสฺส ธมฺมสฺส ปุเรชาต\-ปจฺจเยน ปจฺจโย. อารมฺมณ\-ปุเรชาตํ, วตฺถุ\-ปุเรชาตํ.  อารมฺมณ\-ปุเรชาตํ\csromandash{}จกฺขุํ อสฺสาเทติ อภินนฺทติ, ตํ อารพฺภ ทสฺสเนน ปหาตพฺโพ ราโค อุปฺปชฺชติ, ทิฏฺฐิ อุปฺปชฺชติ, วิจิกิจฺฉา อุปฺปชฺชติ. ทสฺสเนน ปหาตพฺพํ โทมนสฺสํ อุปฺปชฺชติ ฯเปฯ. วตฺถุํ อสฺสาเทติ อภินนฺทติ, ตํ อารพฺภ ทสฺสเนน ปหาตพฺโพ ราโค อุปฺปชฺชติ, ทิฏฺฐิ อุปฺปชฺชติ, วิจิกิจฺฉา อุปฺปชฺชติ. ทสฺสเนน ปหาตพฺพํ โทมนสฺสํ อุปฺปชฺชติ.  วตฺถุ\-ปุเรชาตํ\csromandash{}วตฺถุ ทสฺสเนน ปหาตพฺพานํ ขนฺธานํ ปุเรชาต\-ปจฺจเยน ปจฺจโย. (2)}

\prose{เนวทสฺสเนน นภาวนาย ปหาตพฺโพ ธมฺโม ภาวนาย ปหาตพฺพสฺส ธมฺมสฺส ปุเรชาต\-ปจฺจเยน ปจฺจโย. อารมฺมณ\-ปุเรชาตํ, วตฺถุ\-ปุเรชาตํ.  อารมฺมณ\-ปุเรชาตํ\csromandash{}จกฺขุํ อสฺสาเทติ อภินนฺทติ, ตํ อารพฺภ ภาวนาย ปหาตพฺโพ ราโค อุปฺปชฺชติ. อุทฺธจฺจํ อุปฺปชฺชติ. ภาวนาย ปหาตพฺพํ โทมนสฺสํ อุปฺปชฺชติ. โสตํ ฯเปฯ. กายํ. รูเป ฯเปฯ. โผฏฺฐพฺเพ. วตฺถุํ อสฺสาเทติ อภินนฺทติ, ตํ อารพฺภ ภาวนาย ปหาตพฺโพ ราโค อุปฺปชฺชติ, อุทฺธจฺจํ อุปฺปชฺชติ. ภาวนาย ปหาตพฺพํ โทมนสฺสํ อุปฺปชฺชติ.  วตฺถุ\-ปุเรชาตํ\csromandash{}วตฺถุ ภาวนาย ปหาตพฺพานํ ขนฺธานํ ปุเรชาต\-ปจฺจเยน ปจฺจโย. (3)}

\csromanheader{2}{\textbf{ปจฺฉาชาต}}
\proseitem{87}{ทสฺสเนน ปหาตพฺโพ ธมฺโม เนวทสฺสเนน นภาวนาย ปหาตพฺพสฺส ธมฺมสฺส ปจฺฉาชาต\-ปจฺจเยน ปจฺจโย. ปจฺฉาชาตา ทสฺสเนน ปหาตพฺพา ขนฺธา ปุเรชาตสฺส อิมสฺส กายสฺส ปจฺฉาชาต\-ปจฺจเยน ปจฺจโย. (1)}

\prose{ภาวนาย ปหาตพฺโพ ธมฺโม เนวทสฺสเนน นภาวนาย ปหาตพฺพสฺส ธมฺมสฺส ปจฺฉาชาต\-ปจฺจเยน ปจฺจโย. ปจฺฉาชาตา ภาวนาย ปหาตพฺพา ขนฺธา ปุเรชาตสฺส อิมสฺส กายสฺส ปจฺฉาชาต\-ปจฺจเยน ปจฺจโย. (1)}

\prose{เนวทสฺสเนน นภาวนาย ปหาตพฺโพ ธมฺโม เนวทสฺสเนน นภาวนาย ปหาตพฺพสฺส ธมฺมสฺส ปจฺฉาชาต\-ปจฺจเยน ปจฺจโย. \csromanfolio{171}ปจฺฉาชาตา เนวทสฺสเนน นภาวนาย ปหาตพฺพา ขนฺธา ปุเรชาตสฺส อิมสฺส กายสฺส ปจฺฉาชาต\-ปจฺจเยน ปจฺจโย. (1)}

\csromanheader{2}{\textbf{อาเสวน}}
\proseitem{88}{ทสฺสเนน ปหาตพฺโพ ธมฺโม ทสฺสเนน ปหาตพฺพสฺส ธมฺมสฺส อาเสวน\-ปจฺจเยน ปจฺจโย. ปุริมา ปุริมา ทสฺสเนน ปหาตพฺพา ขนฺธา ปจฺฉิมานํ ปจฺฉิมานํ ฯเปฯ. อาเสวน\-ปจฺจเยน ปจฺจโย. (1)}

\prose{ภาวนาย ปหาตพฺโพ ธมฺโม ภาวนาย ปหาตพฺพสฺส ธมฺมสฺส อาเสวน\-ปจฺจเยน ปจฺจโย. ปุริมา ปุริมา ภาวนาย ปหาตพฺพา ขนฺธา ปจฺฉิมานํ ปจฺฉิมานํ ฯเปฯ. อาเสวน\-ปจฺจเยน ปจฺจโย. (1)}

\prose{เนวทสฺสเนน นภาวนาย ปหาตพฺโพ ธมฺโม เนวทสฺสเนน นภาวนาย ปหาตพฺพสฺส ธมฺมสฺส อาเสวน\-ปจฺจเยน ปจฺจโย. ปุริมา ปุริมา เนวทสฺสเนน นภาวนาย ปหาตพฺพา ขนฺธา ปจฺฉิมานํ ปจฺฉิมานํ ฯเปฯ. อาเสวน\-ปจฺจเยน ปจฺจโย. อนุโลมํ โคตฺรภุสฺส. อนุโลมํ โวทานสฺส. โคตฺรภุ มคฺคสฺส. โวทานํ มคฺคสฺส อาเสวน\-ปจฺจเยน ปจฺจโย. (1)}

\csromanheader{2}{\textbf{กมฺม}}
\proseitem{89}{ทสฺสเนน ปหาตพฺโพ ธมฺโม ทสฺสเนน ปหาตพฺพสฺส ธมฺมสฺส กมฺมปจฺจเยน ปจฺจโย. ทสฺสเนน ปหาตพฺพา เจตนา สมฺปยุตฺตกานํ ขนฺธานํ กมฺมปจฺจเยน ปจฺจโย. (1)}

\prose{ทสฺสเนน ปหาตพฺโพ ธมฺโม เนวทสฺสเนน นภาวนาย ปหาตพฺพสฺส ธมฺมสฺส กมฺมปจฺจเยน ปจฺจโย. สหชาตา, นานากฺขณิกา.  สหชาตา ทสฺสเนน ปหาตพฺพา เจตนา จิตฺต\-สมุฏฺฐานานํ รูปานํ กมฺมปจฺจเยน ปจฺจโย.  นานากฺขณิกา ทสฺสเนน ปหาตพฺพา เจตนา วิปากานํ ขนฺธานํ กฏตฺตา จ รูปานํ กมฺมปจฺจเยน ปจฺจโย. (2)}

\prose{ทสฺสเนน ปหาตพฺโพ ธมฺโม ทสฺสเนน ปหาตพฺพสฺส จ เนวทสฺสเนน นภาวนาย ปหาตพฺพสฺส จ ธมฺมสฺส กมฺมปจฺจเยน ปจฺจโย. ทสฺสเนน ปหาตพฺพา เจตนา สมฺปยุตฺตา\-กานํ ขนฺธานํ จิตฺตสมุฏฺฐานานญฺจ รูปานํ กมฺมปจฺจเยน ปจฺจโย. (3)}

\prose{\csromanfolio{172}ภาวนาย ปหาตพฺโพ ธมฺโม ภาวนาย ปหาตพฺพสฺส ธมฺมสฺส กมฺมปจฺจเยน ปจฺจโย. ภาวนาย ปหาตพฺพา เจตนา สมฺปยุตฺตกานํ ขนฺธานํ กมฺมปจฺจเยน ปจฺจโย. (1)}

\prose{ภาวนาย ปหาตพฺโพ ธมฺโม เนวทสฺสเนน นภาวนาย ปหาตพฺพสฺส ธมฺมสฺส กมฺมปจฺจเยน ปจฺจโย. ภาวนาย ปหาตพฺพา เจตนา จิตฺต\-สมุฏฺฐานานํ รูปานํ กมฺมปจฺจเยน ปจฺจโย. (2)}

\prose{ภาวนาย ปหาตพฺโพ ธมฺโม ภาวนาย ปหาตพฺพสฺส จ เนวทสฺสเนน นภาวนาย ปหาตพฺพสฺส จ ธมฺมสฺส กมฺมปจฺจเยน ปจฺจโย. ภาวนาย ปหาตพฺพา เจตนา สมฺปยุตฺตกานํ ขนฺธานํ จิตฺตสมุฏฺฐานานญฺจ รูปานํ กมฺมปจฺจเยน ปจฺจโย. (3)}

\prose{เนวทสฺสเนน นภาวนาย ปหาตพฺโพ ธมฺโม เนวทสฺสเนน นภาวนาย ปหาตพฺพสฺส ธมฺมสฺส กมฺมปจฺจเยน ปจฺจโย. สหชาตา, นานากฺขณิกา.  สหชาตา เนวทสฺสเนน นภาวนาย ปหาตพฺพา เจตนา สมฺปยุตฺตกานํ ขนฺธานํ จิตฺตสมุฏฺฐานานญฺจ รูปานํ กมฺมปจฺจเยน ปจฺจโย. ปฏิสนฺธิ\-กฺขเณ เนวทสฺสเนน นภาวนาย ปหาตพฺพา เจตนา สมฺปยุตฺตกานํ ขนฺธานํ กฏตฺตา จ รูปานํ ฯเปฯ. นานากฺขณิกา เนวทสฺสเนน นภาวนาย ปหาตพฺพา เจตนา วิปากานํ ขนฺธานํ กฏตฺตา จ รูปานํ กมฺมปจฺจเยน ปจฺจโย. (1)}

\csromanheader{2}{\textbf{วิปาก}}
\proseitem{90}{เนวทสฺสเนน นภาวนาย ปหาตพฺโพ ธมฺโม เนวทสฺสเนน นภาวนาย ปหาตพฺพสฺส ธมฺมสฺส วิปาก\-ปจฺจเยน ปจฺจโย. วิปาโก เนวทสฺสเนน นภาวนาย ปหาตพฺโพ เอโก ขนฺโธ ติณฺณนฺนํ ขนฺธานํ จิตฺตสมุฏฺฐานานญฺจ รูปานํ วิปาก\-ปจฺจเยน ปจฺจโย ฯเปฯ. ปฏิสนฺธิ\-กฺขเณ ฯเปฯ. ขนฺธา วตฺถุสฺส วิปาก\-ปจฺจเยน ปจฺจโย.}

\csromanheader{2}{\textbf{อาหาราทิ}}
\proseitem{91}{ทสฺสเนน ปหาตพฺโพ ธมฺโม ทสฺสเนน ปหาตพฺพสฺส ธมฺมสฺส อาหารปจฺจเยน ปจฺจโย. (สํขิตฺตํ.) กพฬีกาโร. สตฺต ปญฺหา. \csromanfolio{173}อินฺทฺริย\-ปจฺจเยน ปจฺจโย. จกฺขุนฺทฺริยญฺจ ฯเปฯ. รูปชีวิตินฺทฺริยญฺจ ฯเปฯ. สตฺต ปญฺหา. ฌานปจฺจเยน ปจฺจโย. มคฺคปจฺจเยน ปจฺจโย. สมฺปยุตฺต\-ปจฺจเยน ปจฺจโย.}

\csromanheader{2}{\textbf{วิปฺปยุตฺต}}
\proseitem{92}{ทสฺสเนน ปหาตพฺโพ ธมฺโม ทสฺสเนน ปหาตพฺพสฺส ธมฺมสฺส วิปฺปยุตฺต\-ปจฺจเยน ปจฺจโย. สหชาตํ, ปจฺฉาชาตํ.  สหชาตา ทสฺสเนน ปหาตพฺพา ขนฺธา จิตฺต\-สมุฏฺฐานานํ รูปานํ วิปฺปยุตฺต\-ปจฺจเยน ปจฺจโย.  ปจฺฉาชาตา ทสฺสเนน ปหาตพฺพา ขนฺธา ปุเรชาตสฺส อิมสฺส กายสฺส วิปฺปยุตฺต\-ปจฺจเยน ปจฺจโย. (1)}

\prose{ภาวนาย ปหาตพฺโพ ธมฺโม เนวทสฺสเนน นภาวนาย ปหาตพฺพสฺส ธมฺมสฺส วิปฺปยุตฺต\-ปจฺจเยน ปจฺจโย. สหชาตํ, ปจฺฉาชาตํ. (อิทมฺปิ ทสฺสเนน สทิสํ.)}

\prose{เนวทสฺสเนน นภาวนาย ปหาตพฺโพ ธมฺโม เนวทสฺสเนน นภาวนาย ปหาตพฺพสฺส ธมฺมสฺส วิปฺปยุตฺต\-ปจฺจเยน ปจฺจโย. สหชาตํ, ปุเรชาตํ, ปจฺฉาชาตํ.  สหชาตา เนวทสฺสเนน นภาวนาย ปหาตพฺพา ขนฺธา จิตฺต\-สมุฏฺฐานานํ รูปานํ วิปฺปยุตฺต\-ปจฺจเยน ปจฺจโย. ปฏิสนฺธิ\-กฺขเณ เนวทสฺสเนน นภาวนาย ปหาตพฺพา ขนฺธา กฏตฺตารูปานํ วิปฺปยุตฺต\-ปจฺจเยน ปจฺจโย. ขนฺธา วตฺถุสฺส วิปฺปยุตฺต\-ปจฺจเยน ปจฺจโย. วตฺถุ ขนฺธานํ วิปฺปยุตฺต\-ปจฺจเยน ปจฺจโย.  ปุเรชาตํ จกฺขายตนํ จกฺขุ\-วิญฺญาณสฺส ฯเปฯ. กายายตนํ กายวิญฺญาณสฺส ฯเปฯ. วตฺถุ เนวทสฺสเนน นภาวนาย ปหาตพฺพานํ ขนฺธานํ วิปฺปยุตฺต\-ปจฺจเยน ปจฺจโย.  ปจฺฉาชาตา เนวทสฺสเนน นภาวนาย ปหาตพฺพา ขนฺธา ปุเรชาตสฺส อิมสฺส กายสฺส วิปฺปยุตฺต\-ปจฺจเยน ปจฺจโย. (1)}

\prose{เนวทสฺสเนน นภาวนาย ปหาตพฺโพ ธมฺโม ทสฺสเนน ปหาตพฺพสฺส ธมฺมสฺส วิปฺปยุตฺต\-ปจฺจเยน ปจฺจโย. ปุเรชาตํ วตฺถุ ทสฺสเนน ปหาตพฺพานํ ขนฺธานํ วิปฺปยุตฺต\-ปจฺจเยน ปจฺจโย. (2)}

\prose{เนวทสฺสเนน นภาวนาย ปหาตพฺโพ ธมฺโม ภาวนาย ปหาตพฺพสฺส ธมฺมสฺส วิปฺปยุตฺต\-ปจฺจเยน ปจฺจโย. ปุเรชาตํ วตฺถุ ภาวนาย ปหาตพฺพานํ ขนฺธานํ วิปฺปยุตฺต\-ปจฺจเยน ปจฺจโย. (3)}

\csromanheader{2}{\textbf{อตฺถิ}}
\proseitem{93}{\csromanfolio{174}ทสฺสเนน ปหาตพฺโพ ธมฺโม ทสฺสเนน ปหาตพฺพสฺส ธมฺมสฺส อตฺถิปจฺจเยน ปจฺจโย. ทสฺสเนน ปหาตพฺโพ เอโก ขนฺโธ ติณฺณนฺนํ ขนฺธานํ อตฺถิปจฺจเยน ปจฺจโย ฯเปฯ เทฺว ขนฺธา ฯเปฯ. (1)}

\prose{ทสฺสเนน ปหาตพฺโพ ธมฺโม เนวทสฺสเนน นภาวนาย ปหาตพฺพสฺส ธมฺมสฺส อตฺถิปจฺจเยน ปจฺจโย. สหชาตํ, ปจฺฉาชาตํ.  สหชาตา ทสฺสเนน ปหาตพฺพา ขนฺธา จิตฺต\-สมุฏฺฐานานํ รูปานํ อตฺถิปจฺจเยน ปจฺจโย.  ปจฺฉาชาตา ทสฺสเนน ปหาตพฺพา ขนฺธา ปุเรชาตสฺส อิมสฺส กายสฺส อตฺถิปจฺจเยน ปจฺจโย. (2)}

\prose{ทสฺสเนน ปหาตพฺโพ ธมฺโม ทสฺสเนน ปหาตพฺพสฺส จ เนวทสฺสเนน นภาวนาย ปหาตพฺพสฺส จ ธมฺมสฺส อตฺถิปจฺจเยน ปจฺจโย. ทสฺสเนน ปหาตพฺโพ เอโก ขนฺโธ ติณฺณนฺนํ ขนฺธานํ จิตฺตสมุฏฺฐานานญฺจ รูปานํ ฯเปฯ เทฺว ขนฺธา ฯเปฯ. (3)}

\prose{ภาวนาย ปหาตพฺโพ ธมฺโม ภาวนาย ปหาตพฺพสฺส. ตีณิ. (ทสฺสเนน สทิสํ กาตพฺพํ.)}

\proseitem{94}{เนวทสฺสเนน นภาวนาย ปหาตพฺโพ ธมฺโม เนวทสฺสเนน นภาวนาย ปหาตพฺพสฺส ธมฺมสฺส อตฺถิปจฺจเยน ปจฺจโย. สหชาตํ, ปุเรชาตํ, ปจฺฉาชาตํ, อาหารํ, อินฺทฺริยํ.  สหชาโต เนวทสฺสเนน นภาวนาย ปหาตพฺโพ เอโก ขนฺโธ ติณฺณนฺนํ ขนฺธานํ จิตฺตสมุฏฺฐานานญฺจ รูปานํ อตฺถิปจฺจเยน ปจฺจโย ฯเปฯ. ปฏิสนฺธิ\-กฺขเณ ฯเปฯ. ขนฺธา วตฺถุสฺส ฯเปฯ. วตฺถุ ขนฺธานํ อตฺถิปจฺจเยน ปจฺจโย. เอกํ มหาภูตํ ฯเปฯ. อสญฺญสตฺตานํ เอกํ มหาภูตํ ฯเปฯ. ปุเรชาตํ จกฺขุํ อนิจฺจโต ทุกฺขโต อนตฺตโต วิปสฺสติ. โสตํ ฯเปฯ. กายํ. รูเป ฯเปฯ. โผฏฺฐพฺเพ. วตฺถุํ อนิจฺจโต ทุกฺขโต อนตฺตโต วิปสฺสติ. ทิพฺเพน จกฺขุนา รูปํ ปสฺสติ. ทิพฺพาย โสตธาตุยา สทฺทํ สุณาติ. รูปายตนํ จกฺขุ\-วิญฺญาณสฺส ฯเปฯ. โผฏฺฐพฺพา\-ยตนํ กายวิญฺญาณสฺส ฯเปฯ. จกฺขายตนํ จกฺขุ\-วิญฺญาณสฺส ฯเปฯ. กายายตนํ กายวิญฺญาณสฺส ฯเปฯ. วตฺถุ เนวทสฺสเนน นภาวนาย ปหาตพฺพานํ ขนฺธานํ อตฺถิปจฺจเยน ปจฺจโย.  ปจฺฉาชาตา เนวทสฺสเนน นภาวนาย ปหาตพฺพา ขนฺธา ปุเรชาตสฺส อิมสฺส กายสฺส \csromanfolio{175}อตฺถิปจฺจเยน ปจฺจโย. กพฬีกาโร อาหาโร อิมสฺส กายสฺส ฯเปฯ. รูปชีวิตินฺทฺริยํ กฏตฺตารูปานํ อตฺถิปจฺจเยน ปจฺจโย. (1)}

\prose{เนวทสฺสเนน นภาวนาย ปหาตพฺโพ ธมฺโม ทสฺสเนน ปหาตพฺพสฺส ธมฺมสฺส อตฺถิปจฺจเยน ปจฺจโย. \textbf{ปุเรชาตํ} จกฺขุํ อสฺสาเทติ อภินนฺทติ, ตํ อารพฺภ ทสฺสเนน ปหาตพฺโพ ราโค อุปฺปชฺชติ, ทิฏฺฐิ อุปฺปชฺชติ, วิจิกิจฺฉา อุปฺปชฺชติ. ทสฺสเนน ปหาตพฺพํ โทมนสฺสํ อุปฺปชฺชติ. โสตํ ฯเปฯ. วตฺถุํ อสฺสาเทติ ฯเปฯ. วตฺถุ ทสฺสเนน ปหาตพฺพานํ ขนฺธานํ อตฺถิปจฺจเยน ปจฺจโย. (2)}

\prose{เนวทสฺสเนน นภาวนาย ปหาตพฺโพ ธมฺโม ภาวนาย ปหาตพฺพสฺส ธมฺมสฺส อตฺถิปจฺจเยน ปจฺจโย. \textbf{ปุเรชาตํ} จกฺขุํ อสฺสาเทติ อภินนฺทติ, ตํ อารพฺภ ภาวนาย ปหาตพฺโพ ราโค อุปฺปชฺชติ, อุทฺธจฺจํ อุปฺปชฺชติ. ภาวนาย ปหาตพฺพํ โทมนสฺสํ อุปฺปชฺชติ. โสตํ ฯเปฯ. วตฺถุํ อสฺสาเทติ อภินนฺทติ ฯเปฯ. วตฺถุ ภาวนาย ปหาตพฺพานํ ขนฺธานํ อตฺถิปจฺจเยน ปจฺจโย. (3)}

\proseitem{95}{ทสฺสเนน ปหาตพฺโพ จ เนวทสฺสเนน นภาวนาย ปหาตพฺโพ จ ธมฺมา ทสฺสเนน ปหาตพฺพสฺส ธมฺมสฺส อตฺถิปจฺจเยน ปจฺจโย. สหชาตํ, ปุเรชาตํ.  สหชาโต ทสฺสเนน ปหาตพฺโพ เอโก ขนฺโธ จ วตฺถุ จ ติณฺณนฺนํ ขนฺธานํ อตฺถิปจฺจเยน ปจฺจโย ฯเปฯ เทฺว ขนฺธา จ วตฺถุ จ ฯเปฯ. (1)}

\prose{ทสฺสเนน ปหาตพฺโพ จ เนวทสฺสเนน นภาวนาย ปหาตพฺโพ จ ธมฺมา เนวทสฺสเนน นภาวนาย ปหาตพฺพสฺส ธมฺมสฺส อตฺถิปจฺจเยน ปจฺจโย. สหชาตํ, ปจฺฉาชาตํ, อาหารํ, อินฺทฺริยํ.  สหชาตา ทสฺสเนน ปหาตพฺพา ขนฺธา จ มหาภูตา จ จิตฺต\-สมุฏฺฐานานํ รูปานํ อตฺถิปจฺจเยน ปจฺจโย.  ปจฺฉาชาตา ทสฺสเนน ปหาตพฺพา ขนฺธา จ กพฬีกาโร อาหาโร จ อิมสฺส กายสฺส อตฺถิปจฺจเยน ปจฺจโย.  ปจฺฉาชาตา ทสฺสเนน ปหาตพฺพา ขนฺธา จ รูปชีวิตินฺทฺริยญฺจ กฏตฺตารูปานํ อตฺถิปจฺจเยน ปจฺจโย. (2)}

\prose{ภาวนาย ปหาตพฺโพ จ เนวทสฺสเนน นภาวนาย ปหาตพฺโพ จ ธมฺมา ภาวนาย ปหาตพฺพสฺส ฯเปฯ. (เทฺว ปญฺหา กาตพฺพา.)}

\csromanheader{2}{\textbf{นตฺถิ วิคตาวิคต}}
\proseitem{96}{\csromanfolio{176}ทสฺสเนน ปหาตพฺโพ ธมฺโม ทสฺสเนน ปหาตพฺพสฺส ธมฺมสฺส นตฺถิ\-ปจฺจเยน ปจฺจโย. วิคต\-ปจฺจเยน ปจฺจโย. อวิคต\-ปจฺจเยน ปจฺจโย.}

\csromanmarkh{1. ปจฺจยานุโลม}\csromanmarkhii{2. สงฺขฺยาวาร}\csromanheadernomark{2}{1. \textbf{ปจฺจยานุโลม 2. สงฺขฺยาวาร}}
\proseitemwithcloserruled{97}{เหตุยา สตฺต, อารมฺมเณ อฏฺฐ, อธิปติยา ทส, อนนฺตเร สตฺต, สมนนฺตเร สตฺต, สหชาเต นว, อญฺญมญฺเญ ตีณิ, นิสฺสเย เตรส, อุปนิสฺสเย อฏฺฐ, ปุเรชาเต ตีณิ, ปจฺฉาชาเต ตีณิ, อาเสวเน ตีณิ, กมฺเม สตฺต, วิปาเก เอกํ, อาหาเร สตฺต, อินฺทฺริเย สตฺต, ฌาเน สตฺต, มคฺเค สตฺต, สมฺปยุตฺเต ตีณิ, วิปฺปยุตฺเต ปญฺจ, อตฺถิยา เตรส, นตฺถิยา สตฺต, วิคเต สตฺต, อวิคเต เตรส. (เอวํ คเณตพฺพํ.)}{\textbf{อนุโลมํ.}}
\csromanheader{2}{2. ปจฺจนียุ\-ทฺธาร}
\proseitem{98}{ทสฺสเนน ปหาตพฺโพ ธมฺโม ทสฺสเนน ปหาตพฺพสฺส ธมฺมสฺส อารมฺมณ\-ปจฺจเยน ปจฺจโย. สหชาต\-ปจฺจเยน ปจฺจโย. อุปนิสฺสย\-ปจฺจเยน ปจฺจโย. (1)}

\prose{ทสฺสเนน ปหาตพฺโพ ธมฺโม เนวทสฺสเนน นภาวนาย ปหาตพฺพสฺส ธมฺมสฺส อารมฺมณ\-ปจฺจเยน ปจฺจโย. สหชาต\-ปจฺจเยน ปจฺจโย. อุปนิสฺสย\-ปจฺจเยน ปจฺจโย. ปจฺฉาชาต\-ปจฺจเยน ปจฺจโย. กมฺมปจฺจเยน ปจฺจโย. (2)}

\prose{ทสฺสเนน ปหาตพฺโพ ธมฺโม ทสฺสเนน ปหาตพฺพสฺส จ เนวทสฺสเนน นภาวนาย ปหาตพฺพสฺส จ ธมฺมสฺส สหชาต\-ปจฺจเยน ปจฺจโย. (3)}

\proseitem{99}{ภาวนาย ปหาตพฺโพ ธมฺโม ภาวนาย ปหาตพฺพสฺส ธมฺมสฺส อารมฺมณ\-ปจฺจเยน ปจฺจโย. สหชาต\-ปจฺจเยน ปจฺจโย. อุปนิสฺสย\-ปจฺจเยน ปจฺจโย. (1)}

\prose{ภาวนาย ปหาตพฺโพ ธมฺโม ทสฺสเนน ปหาตพฺพสฺส ธมฺมสฺส อารมฺมณ\-ปจฺจเยน ปจฺจโย. อุปนิสฺสย\-ปจฺจเยน ปจฺจโย. (2)}

\prose{\csromanfolio{177}ภาวนาย ปหาตพฺโพ ธมฺโม เนวทสฺสเนน นภาวนาย ปหาตพฺพสฺส ธมฺมสฺส อารมฺมณ\-ปจฺจเยน ปจฺจโย. สหชาต\-ปจฺจเยน ปจฺจโย. อุปนิสฺสย\-ปจฺจเยน ปจฺจโย. ปจฺฉาชาต\-ปจฺจเยน ปจฺจโย. (3)}

\prose{ภาวนาย ปหาตพฺโพ ธมฺโม ภาวนาย ปหาตพฺพสฺส จ เนวทสฺสเนน นภาวนาย ปหาตพฺพสฺส จ ธมฺมสฺส สหชาต\-ปจฺจเยน ปจฺจโย. (4)}

\proseitem{100}{เนวทสฺสเนน นภาวนาย ปหาตพฺโพ ธมฺโม เนวทสฺสเนน นภาวนาย ปหาตพฺพสฺส ธมฺมสฺส อารมฺมณ\-ปจฺจเยน ปจฺจโย. สหชาต\-ปจฺจเยน ปจฺจโย. อุปนิสฺสย\-ปจฺจเยน ปจฺจโย. ปุเรชาต\-ปจฺจเยน ปจฺจโย. ปจฺฉาชาต\-ปจฺจเยน ปจฺจโย. กมฺมปจฺจเยน ปจฺจโย. อาหารปจฺจเยน ปจฺจโย. อินฺทฺริย\-ปจฺจเยน ปจฺจโย. (1)}

\prose{เนวทสฺสเนน นภาวนาย ปหาตพฺโพ ธมฺโม ทสฺสเนน ปหาตพฺพสฺส ธมฺมสฺส อารมฺมณ\-ปจฺจเยน ปจฺจโย. อุปนิสฺสย\-ปจฺจเยน ปจฺจโย. ปุเรชาต\-ปจฺจเยน ปจฺจโย. (2)}

\prose{เนวทสฺสเนน นภาวนาย ปหาตพฺโพ ธมฺโม ภาวนาย ปหาตพฺพสฺส ธมฺมสฺส อารมฺมณ\-ปจฺจเยน ปจฺจโย. อุปนิสฺสย\-ปจฺจเยน ปจฺจโย. ปุเรชาต\-ปจฺจเยน ปจฺจโย. (3)}

\proseitem{101}{ทสฺสเนน ปหาตพฺโพ จ เนวทสฺสเนน นภาวนาย ปหาตพฺโพ จ ธมฺมา ทสฺสเนน ปหาตพฺพสฺส ธมฺมสฺส สหชาตํ, ปุเรชาตํ. (1)}

\prose{ทสฺสเนน ปหาตพฺโพ จ เนวทสฺสเนน นภาวนาย ปหาตพฺโพ จ ธมฺมา เนวทสฺสเนน นภาวนาย ปหาตพฺพสฺส ธมฺมสฺส สหชาตํ, ปจฺฉาชาตํ, อาหารํ, อินฺทฺริยํ. (2)}

\prose{ภาวนาย ปหาตพฺโพ จ เนวทสฺสเนน นภาวนาย ปหาตพฺโพ จ ธมฺมา ภาวนาย ปหาตพฺพสฺส ธมฺมสฺส สหชาตํ, ปุเรชาตํ. (1)}

\prose{ภาวนาย ปหาตพฺโพ จ เนวทสฺสเนน นภาวนาย ปหาตพฺโพ จ ธมฺมา เนวทสฺสเนน นภาวนาย ปหาตพฺพสฺส ธมฺมสฺส สหชาตํ, ปจฺฉาชาตํ, อาหารํ, อินฺทฺริยํ. (2)}

\csromanmarkh{2. ปจฺจย\-ปจฺจนีย}\csromanmarkhii{2. สงฺขฺยาวาร}\csromanheadernomark{2}{2. ปจฺจย\-ปจฺจนีย 2. สงฺขฺยาวาร}
\proseitemwithcloserruled{102}{\csromanfolio{178}นเหตุยา จุทฺทส, นอารมฺมเณ จุทฺทส, นอธิปติยา จุทฺทส, นอนนฺตเร จุทฺทส, นสมนนฺตเร จุทฺทส, นสหชาเต ทส, นอญฺญมญฺเญ ทส, นนิสฺสเย ทส, นอุปนิสฺสเย จุทฺทส, นปุเรชาเต ทฺวาทส, นปจฺฉาชาเต จุทฺทส, นอาเสวเน จุทฺทส, นกมฺเม จุทฺทส, นวิปาเก จุทฺทส, นอาหาเร จุทฺทส, นอินฺทฺริเย จุทฺทส, นฌาเน จุทฺทส, นมคฺเค จุทฺทส, นสมฺปยุตฺเต ทส, นวิปฺปยุตฺเต อฏฺฐ, โนอตฺถิยา อฏฺฐ, โนนตฺถิยา จุทฺทส, โนวิคเต จุทฺทส, โนอวิคเต อฏฺฐ. (เอวํ คเณตพฺพํ.)}{ปจฺจนียํ.}
\proseitemwithcloserruled{103}{เหตุปจฺจยา นอารมฺมเณ สตฺต, นอธิปติยา สตฺต, นอนนฺตเร สตฺต, นสมนนฺตเร สตฺต, นอญฺญมญฺเญ ตีณิ, นอุปนิสฺสเย สตฺต, นปุเรชาเต สตฺต, นปจฺฉาชาเต สตฺต, นอาเสวเน สตฺต, นกมฺเม สตฺต, นวิปาเก สตฺต, นอาหาเร สตฺต, นอินฺทฺริเย สตฺต, นฌาเน สตฺต, นมคฺเค สตฺต, นสมฺปยุตฺเต ตีณิ, นวิปฺปยุตฺเต ตีณิ, โนนตฺถิยา สตฺต, โนวิคเต สตฺต. (เอวํ คเณตพฺพํ.)}{อนุโลม\-ปจฺจนียํ.}
\proseitem{104}{นเหตุปจฺจยา อารมฺมเณ อฏฺฐ, อธิปติยา ทส, อนนฺตเร สตฺต, สมนนฺตเร สตฺต, สหชาเต นว, อญฺญมญฺเญ ตีณิ, นิสฺสเย เตรส, อุปนิสฺสเย อฏฺฐ, ปุเรชาเต ตีณิ, ปจฺฉาชาเต ตีณิ, อาเสวเน ตีณิ, กมฺเม สตฺต, วิปาเก เอกํ, อาหาเร สตฺต, อินฺทฺริเย สตฺต, ฌาเน สตฺต, มคฺเค สตฺต, สมฺปยุตฺเต ตีณิ, วิปฺปยุตฺเต ปญฺจ, อตฺถิยา เตรส, นตฺถิยา สตฺต, วิคเต สตฺต, อวิคเต เตรส. (เอวํ คเณตพฺพํ.)}

\vspace{\tipitakaparskip}
\csromancenter{ปจฺจนียา\-นุโลมํ.}
\nitthitamruled{ทสฺสเนน ปหาตพฺพตฺติกํ นิฏฺฐิตํ.}
\csromancenter{9. ทสฺสเนน\-ปหาตพฺพ\-เหตุ\-กตฺติก 1. ปฏิจฺจวาร}
\csromanmarkh{1. ปจฺจยานุโลม}\csromanmarkhii{1. วิภงฺควาร}\csromanheadernomark{2}{1. \textbf{ปจฺจยานุโลม 1. วิภงฺควาร}}
\csromanheader{2}{\textbf{เหตุ}}
\proseitem{1}{\csromanfolio{179}ทสฺสเนน ปหาตพฺพ\-เหตุกํ ธมฺมํ ปฏิจฺจ ทสฺสเนน ปหาตพฺพ\-เหตุโก ธมฺโม อุปฺปชฺชติ เหตุปจฺจยา. ทสฺสเนน ปหาตพฺพ\-เหตุกํ เอกํ ขนฺธํ ปฏิจฺจ ตโย ขนฺธา ฯเปฯ เทฺว ขนฺธา. (1)}

\prose{ทสฺสเนน ปหาตพฺพ\-เหตุกํ ธมฺมํ ปฏิจฺจ เนวทสฺสเนน นภาวนาย ปหาตพฺพ\-เหตุโก ธมฺโม อุปฺปชฺชติ เหตุปจฺจยา. ทสฺสเนน ปหาตพฺพ\-เหตุเก ขนฺเธ ปฏิจฺจ จิตฺต\-สมุฏฺฐานํ รูปํ. (2)}

\prose{ทสฺสเนน ปหาตพฺพ\-เหตุกํ ธมฺมํ ปฏิจฺจ ทสฺสเนน ปหาตพฺพ\-เหตุโก จ เนวทสฺสเนน นภาวนาย ปหาตพฺพ\-เหตุโก จ ธมฺมา อุปฺปชฺชนฺติ เหตุปจฺจยา. ทสฺสเนน ปหาตพฺพ\-เหตุกํ เอกํ ขนฺธํ ปฏิจฺจ ตโย ขนฺธา จิตฺต\-สมุฏฺฐานญฺจ รูปํ ฯเปฯ เทฺว ขนฺธา ฯเปฯ. (3)}

\proseitem{2}{ภาวนาย ปหาตพฺพ\-เหตุกํ ธมฺมํ ปฏิจฺจ ภาวนาย ปหาตพฺพ\-เหตุโก ธมฺโม อุปฺปชฺชติ เหตุปจฺจยา. ภาวนาย ปหาตพฺพ\-เหตุกํ เอกํ ขนฺธํ ปฏิจฺจ ตโย ขนฺธา ฯเปฯ เทฺว ขนฺธา. (1)}

\prose{ภาวนาย ปหาตพฺพ\-เหตุกํ ธมฺมํ ปฏิจฺจ เนวทสฺสเนน นภาวนาย ปหาตพฺพ\-เหตุโก ธมฺโม อุปฺปชฺชติ เหตุปจฺจยา. ภาวนาย ปหาตพฺพ\-เหตุเก ขนฺเธ ปฏิจฺจ จิตฺต\-สมุฏฺฐานํ รูปํ. (2)}

\prose{ภาวนาย ปหาตพฺพ\-เหตุกํ ธมฺมํ ปฏิจฺจ ภาวนาย ปหาตพฺพ\-เหตุโก จ เนวทสฺสเนน นภาวนาย ปหาตพฺพ\-เหตุโก จ ธมฺมา อุปฺปชฺชนฺติ เหตุปจฺจยา. ภาวนาย ปหาตพฺพ\-เหตุกํ เอกํ ขนฺธํ ปฏิจฺจ ตโย ขนฺธา จิตฺต\-สมุฏฺฐานญฺจ รูปํ ฯเปฯ เทฺว ขนฺเธ ฯเปฯ. (3)}

\proseitem{3}{เนวทสฺสเนน นภาวนาย ปหาตพฺพ\-เหตุกํ ธมฺมํ ปฏิจฺจ เนวทสฺสเนน นภาวนาย ปหาตพฺพ\-เหตุโก ธมฺโม อุปฺปชฺชติ เหตุปจฺจยา. เนวทสฺสเนน นภาวนาย ปหาตพฺพ\-เหตุกํ เอกํ ขนฺธํ ปฏิจฺจ ตโย \csromanfolio{180}ขนฺธา จิตฺต\-สมุฏฺฐานญฺจ รูปํ ฯเปฯ เทฺว ขนฺเธ ปฏิจฺจ เทฺว ขนฺธา จิตฺต\-สมุฏฺฐานญฺจ รูปํ. วิจิกิจฺฉา\-สหคตํ อุทฺธจฺจ\-สหคตํ โมหํ ปฏิจฺจ จิตฺต\-สมุฏฺฐานํ รูปํ. ปฏิสนฺธิ\-กฺขเณ เนวทสฺสเนน นภาวนาย ปหาตพฺพ\-เหตุกํ เอกํ ขนฺธํ ปฏิจฺจ ตโย ขนฺธา กฏตฺตา จ รูปํ ฯเปฯ เทฺว ขนฺธา ฯเปฯ. ขนฺเธ ปฏิจฺจ วตฺถุ, วตฺถุํ ปฏิจฺจ ขนฺธา. เอกํ มหาภูตํ ปฏิจฺจ ฯเปฯ. มหาภูเต ปฏิจฺจ จิตฺต\-สมุฏฺฐานํ รูปํ, กฏตฺตารูปํ, อุปาทารูปํ. (1)}

\prose{เนวทสฺสเนน นภาวนาย ปหาตพฺพ\-เหตุกํ ธมฺมํ ปฏิจฺจ ทสฺสเนน ปหาตพฺพ\-เหตุโก ธมฺโม อุปฺปชฺชติ เหตุปจฺจยา. วิจิกิจฺฉา\-สหคตํ โมหํ ปฏิจฺจ สมฺปยุตฺตกา ขนฺธา. (2)}

\prose{เนวทสฺสเนน นภาวนาย ปหาตพฺพ\-เหตุกํ ธมฺมํ ปฏิจฺจ ภาวนาย ปหาตพฺพ\-เหตุโก ธมฺโม อุปฺปชฺชติ เหตุปจฺจยา. อุทฺธจฺจ\-สหคตํ โมหํ ปฏิจฺจ สมฺปยุตฺตกา ขนฺธา. (3)}

\prose{เนวทสฺสเนน นภาวนาย ปหาตพฺพ\-เหตุกํ ธมฺมํ ปฏิจฺจ ทสฺสเนน ปหาตพฺพ\-เหตุโก จ เนวทสฺสเนน นภาวนาย ปหาตพฺพ\-เหตุโก จ ธมฺมา อุปฺปชฺชนฺติ เหตุปจฺจยา. วิจิกิจฺฉา\-สหคตํ โมหํ ปฏิจฺจ สมฺปยุตฺตกา ขนฺธา จิตฺต\-สมุฏฺฐานญฺจ รูปํ. (4)}

\prose{เนวทสฺสเนน นภาวนาย ปหาตพฺพ\-เหตุกํ ธมฺมํ ปฏิจฺจ ภาวนาย ปหาตพฺพ\-เหตุโก จ เนวทสฺสเนน นภาวนาย ปหาตพฺพ\-เหตุโก จ ธมฺมา อุปฺปชฺชนฺติ เหตุปจฺจยา. อุทฺธจฺจ\-สหคตํ โมหํ ปฏิจฺจ สมฺปยุตฺตกา ขนฺธา จิตฺต\-สมุฏฺฐานญฺจ รูปํ. (5)}

\prose{ทสฺสเนน ปหาตพฺพ\-เหตุ\-กญฺจ เนวทสฺสเนน นภาวนาย ปหาตพฺพ\-เหตุ\-กญฺจ ธมฺมํ ปฏิจฺจ ทสฺสเนน ปหาตพฺพ\-เหตุโก ธมฺโม อุปฺปชฺชติ เหตุปจฺจยา. วิจิกิจฺฉา\-สหคตํ เอกํ ขนฺธญฺจ โมหญฺจ ปฏิจฺจ ตโย ขนฺธา ฯเปฯ เทฺว ขนฺเธ จ โมหญฺจ ปฏิจฺจ เทฺว ขนฺธา. (1)}

\prose{ทสฺสเนน ปหาตพฺพ\-เหตุ\-กญฺจ เนวทสฺสเนน นภาวนาย ปหาตพฺพ\-เหตุ\-กญฺจ ธมฺมํ ปฏิจฺจ เนวทสฺสเนน นภาวนาย ปหาตพฺพ\-เหตุโก ธมฺโม อุปฺปชฺชติ เหตุปจฺจยา. ทสฺสเนน ปหาตพฺพ\-เหตุเก ขนฺเธ จ มหาภูเต จ ปฏิจฺจ จิตฺต\-สมุฏฺฐานํ รูปํ. วิจิกิจฺฉา\-สหคเต ขนฺเธ จ โมหญฺจ ปฏิจฺจ จิตฺต\-สมุฏฺฐานํ รูปํ. (2)}

\prose{\csromanfolio{181}ทสฺสเนน ปหาตพฺพ\-เหตุ\-กญฺจ เนวทสฺสเนน นภาวนาย ปหาตพฺพ\-เหตุ\-กญฺจ ธมฺมํ ปฏิจฺจ ทสฺสเนน ปหาตพฺพ\-เหตุโก จ เนวทสฺสเนน นภาวนาย ปหาตพฺพ\-เหตุโก จ ธมฺมา อุปฺปชฺชนฺติ เหตุปจฺจยา. วิจิกิจฺฉา\-สหคตํ เอกํ ขนฺธญฺจ โมหญฺจ ปฏิจฺจ ตโย ขนฺธา จิตฺต\-สมุฏฺฐานญฺจ รูปํ ฯเปฯ เทฺว ขนฺเธ จ โมหญฺจ ปฏิจฺจ เทฺว ขนฺธา ฯเปฯ. (3)}

\proseitem{5}{ภาวนาย ปหาตพฺพ\-เหตุ\-กญฺจ เนวทสฺสเนน นภาวนาย ปหาตพฺพ\-เหตุ\-กญฺจ ธมฺมํ ปฏิจฺจ ภาวนาย ปหาตพฺพ\-เหตุโก ธมฺโม อุปฺปชฺชติ เหตุปจฺจยา. อุทฺธจฺจ\-สหคตํ เอกํ ขนฺธญฺจ โมหญฺจ ปฏิจฺจ ตโย ขนฺธา ฯเปฯ เทฺว ขนฺเธ จ โมหญฺจ ปฏิจฺจ เทฺว ขนฺธา. (1)}

\prose{ภาวนาย ปหาตพฺพ\-เหตุ\-กญฺจ เนวทสฺสเนน นภาวนาย ปหาตพฺพ\-เหตุ\-กญฺจ ธมฺมํ ปฏิจฺจ เนวทสฺสเนน นภาวนาย ปหาตพฺพ\-เหตุโก ธมฺโม อุปฺปชฺชติ เหตุปจฺจยา. ภาวนาย ปหาตพฺพ\-เหตุเก ขนฺเธ จ มหาภูเต จ ปฏิจฺจ จิตฺต\-สมุฏฺฐานํ รูปํ. อุทฺธจฺจ\-สหคเต ขนฺเธ จ โมหญฺจ ปฏิจฺจ จิตฺต\-สมุฏฺฐานํ รูปํ. (2)}

\prose{ภาวนาย ปหาตพฺพ\-เหตุ\-กญฺจ เนวทสฺสเนน นภาวนาย ปหาตพฺพ\-เหตุ\-กญฺจ ธมฺมํ ปฏิจฺจ ภาวนาย ปหาตพฺพ\-เหตุโก จ เนวทสฺสเนน นภาวนาย ปหาตพฺพ\-เหตุโก จ ธมฺมา อุปฺปชฺชนฺติ เหตุปจฺจยา. อุทฺธจฺจ\-สหคตํ เอกํ ขนฺธญฺจ โมหญฺจ ปฏิจฺจ ตโย ขนฺธา จิตฺต\-สมุฏฺฐานญฺจ รูปํ ฯเปฯ. (3)}

\csromanheader{2}{\textbf{อารมฺมณ}}
\proseitem{6}{ทสฺสเนน ปหาตพฺพ\-เหตุกํ ธมฺมํ ปฏิจฺจ ทสฺสเนน ปหาตพฺพ\-เหตุโก ธมฺโม อุปฺปชฺชติ อารมฺมณ\-ปจฺจยา. ทสฺสเนน ปหาตพฺพ\-เหตุกํ เอกํ ขนฺธํ ปฏิจฺจ ตโย ขนฺธา ฯเปฯ เทฺว ขนฺเธ ปฏิจฺจ เทฺว ขนฺธา. (1)}

\prose{ทสฺสเนน ปหาตพฺพ\-เหตุกํ ธมฺมํ ปฏิจฺจ เนวทสฺสเนน นภาวนาย ปหาตพฺพ\-เหตุโก ธมฺโม อุปฺปชฺชติ อารมฺมณ\-ปจฺจยา. วิจิกิจฺฉา\-สหคเต ขนฺเธ ปฏิจฺจ วิจิกิจฺฉา\-สหคโต โมโห. (2)}

\prose{ทสฺสเนน ปหาตพฺพ\-เหตุกํ ธมฺมํ ปฏิจฺจ ทสฺสเนน ปหาตพฺพ\-เหตุโก จ เนวทสฺสเนน นภาวนาย ปหาตพฺพ\-เหตุโก จ ธมฺมา อุปฺปชฺชนฺติ อารมฺมณ\-ปจฺจยา. วิจิกิจฺฉา\-สหคตํ เอกํ ขนฺธํ ปฏิจฺจ ตโย ขนฺธา จ โมโห จ ฯเปฯ เทฺว ขนฺเธ ปฏิจฺจ เทฺว ขนฺธา จ โมโห จ. (3)}

\prose{\csromanfolio{182}ภาวนาย ปหาตพฺพ\-เหตุกํ ธมฺมํ ปฏิจฺจ ภาวนาย ปหาตพฺพ\-เหตุโก ธมฺโม อุปฺปชฺชติ อารมฺมณ\-ปจฺจยา. ตีณิ. (ทสฺสเนน สทิสํ วิภชิตพฺพํ.)}

\proseitem{7}{เนวทสฺสเนน นภาวนาย ปหาตพฺพ\-เหตุกํ ธมฺมํ ปฏิจฺจ เนวทสฺสเนน นภาวนาย ปหาตพฺพ\-เหตุโก ธมฺโม อุปฺปชฺชติ อารมฺมณ\-ปจฺจยา. เนวทสฺสเนน นภาวนาย ปหาตพฺพ\-เหตุกํ เอกํ ขนฺธํ ปฏิจฺจ ตโย ขนฺธา ฯเปฯ เทฺว ขนฺเธ ปฏิจฺจ เทฺว ขนฺธา. ปฏิสนฺธิ\-กฺขเณ ฯเปฯ. วตฺถุํ ปฏิจฺจ ขนฺธา. (1)}

\prose{เนวทสฺสเนน นภาวนาย ปหาตพฺพ\-เหตุกํ ธมฺมํ ปฏิจฺจ ทสฺสเนน ปหาตพฺพ\-เหตุโก ธมฺโม อุปฺปชฺชติ อารมฺมณ\-ปจฺจยา. วิจิกิจฺฉา\-สหคตํ โมหํ ปฏิจฺจ สมฺปยุตฺตกา ขนฺธา. (2)}

\prose{เนวทสฺสเนน นภาวนาย ปหาตพฺพ\-เหตุกํ ธมฺมํ ปฏิจฺจ ภาวนาย ปหาตพฺพ\-เหตุโก ธมฺโม อุปฺปชฺชติ อารมฺมณ\-ปจฺจยา. อุทฺธจฺจ\-สหคตํ โมหํ ปฏิจฺจ สมฺปยุตฺตกา ขนฺธา. (3)}

\proseitem{8}{ทสฺสเนน ปหาตพฺพ\-เหตุ\-กญฺจ เนวทสฺสเนน นภาวนาย ปหาตพฺพ\-เหตุ\-กญฺจ ธมฺมํ ปฏิจฺจ ทสฺสเนน ปหาตพฺพ\-เหตุโก ธมฺโม อุปฺปชฺชติ อารมฺมณ\-ปจฺจยา. วิจิกิจฺฉา\-สหคตํ เอกํ ขนฺธญฺจ โมหญฺจ ปฏิจฺจ ตโย ขนฺธา ฯเปฯ เทฺว ขนฺเธ จ โมหญฺจ ปฏิจฺจ เทฺว ขนฺธา. (1)}

\prose{ภาวนาย ปหาตพฺพ\-เหตุ\-กญฺจ เนวทสฺสเนน นภาวนาย ปหาตพฺพ\-เหตุ\-กญฺจ ธมฺมํ ปฏิจฺจ ภาวนาย ปหาตพฺพ\-เหตุโก ธมฺโม อุปฺปชฺชติ อารมฺมณ\-ปจฺจยา. อุทฺธจฺจ\-สหคตํ เอกํ ขนฺธญฺจ โมหญฺจ ปฏิจฺจ ตโย ขนฺธา ฯเปฯ เทฺว ขนฺเธ จ โมหญฺจ ปฏิจฺจ เทฺว ขนฺธา. (1)}

\proseitem{9}{ทสฺสเนน ปหาตพฺพ\-เหตุกํ ธมฺมํ ปฏิจฺจ ทสฺสเนน ปหาตพฺพ\-เหตุโก ธมฺโม อุปฺปชฺชติ อธิปติ\-ปจฺจยา. ตีณิ. (เหตุสทิสา.)}

\prose{ภาวนาย ปหาตพฺพ\-เหตุกํ ธมฺมํ ปฏิจฺจ ฯเปฯ. ตีณิ. (เหตุสทิสา อธิปติยา โมโห นตฺถิ.)}

\prose{\csromanfolio{183}เนวทสฺสเนน นภาวนาย ปหาตพฺพ\-เหตุกํ ธมฺมํ ปฏิจฺจ เนวทสฺสเนน นภาวนาย ปหาตพฺพ\-เหตุโก ธมฺโม อุปฺปชฺชติ อธิปติ\-ปจฺจยา. เนวทสฺสเนน นภาวนาย ปหาตพฺพ\-เหตุกํ เอกํ ขนฺธํ ปฏิจฺจ ตโย ขนฺธา จิตฺต\-สมุฏฺฐานญฺจ รูปํ ฯเปฯ เทฺว ขนฺธา. เอกํ มหาภูตํ ปฏิจฺจ ตโย มหาภูตา ฯเปฯ มหาภูเต ปฏิจฺจ จิตฺต\-สมุฏฺฐานํ รูปํ, อุปาทารูปํ. (1)}

\proseitem{10}{ทสฺสเนน ปหาตพฺพ\-เหตุ\-กญฺจ เนวทสฺสเนน นภาวนาย ปหาตพฺพ\-เหตุ\-กญฺจ ธมฺมํ ปฏิจฺจ เนวทสฺสเนน นภาวนาย ปหาตพฺพ\-เหตุโก ธมฺโม อุปฺปชฺชติ อธิปติ\-ปจฺจยา. ทสฺสเนน ปหาตพฺพ\-เหตุเก ขนฺเธ จ มหาภูเต จ ปฏิจฺจ จิตฺต\-สมุฏฺฐานํ รูปํ. (1)}

\prose{ภาวนาย ปหาตพฺพ\-เหตุ\-กญฺจ เนวทสฺสเนน นภาวนาย ปหาตพฺพ\-เหตุ\-กญฺจ ธมฺมํ ปฏิจฺจ เนวทสฺสเนน นภาวนาย ปหาตพฺพ\-เหตุโก ธมฺโม อุปฺปชฺชติ อธิปติ\-ปจฺจยา. ภาวนาย ปหาตพฺพ\-เหตุเก ขนฺเธ จ มหาภูเต จ ปฏิจฺจ จิตฺต\-สมุฏฺฐานํ รูปํ. (1)}

\csromanheader{2}{\textbf{อนนฺตร สมนนฺตร}}
\proseitem{11}{ทสฺสเนน ปหาตพฺพ\-เหตุกํ ธมฺมํ ปฏิจฺจ ทสฺสเนน ปหาตพฺพ\-เหตุโก ธมฺโม อุปฺปชฺชติ อนนฺตร\-ปจฺจยา. สมนนฺตร\-ปจฺจยา. (อารมฺมณ\-สทิสํ.)}

\csromanheader{2}{\textbf{สหชาต}}
\proseitem{12}{ทสฺสเนน ปหาตพฺพ\-เหตุกํ ธมฺมํ ปฏิจฺจ ทสฺสเนน ปหาตพฺพ\-เหตุโก ธมฺโม อุปฺปชฺชติ สหชาต\-ปจฺจยา. ทสฺสเนน ปหาตพฺพ\-เหตุกํ เอกํ ขนฺธํ ปฏิจฺจ ตโย ขนฺธา ฯเปฯ เทฺว ขนฺเธ ปฏิจฺจ เทฺว ขนฺธา. (1)}

\prose{ทสฺสเนน ปหาตพฺพ\-เหตุกํ ธมฺมํ ปฏิจฺจ เนวทสฺสเนน นภาวนาย ปหาตพฺพ\-เหตุโก ธมฺโม อุปฺปชฺชติ สหชาต\-ปจฺจยา. ทสฺสเนน ปหาตพฺพ\-เหตุเก ขนฺเธ ปฏิจฺจ จิตฺต\-สมุฏฺฐานํ รูปํ. วิจิกิจฺฉา\-สหคเต ขนฺเธ ปฏิจฺจ โมโห จิตฺต\-สมุฏฺฐานญฺจ รูปํ. (2)}

\prose{ทสฺสเนน ปหาตพฺพ\-เหตุกํ ธมฺมํ ปฏิจฺจ ทสฺสเนน ปหาตพฺพ\-เหตุโก จ เนวทสฺสเนน นภาวนาย ปหาตพฺพ\-เหตุโก จ ธมฺมา อุปฺปชฺชนฺติ \csromanfolio{184}สหชาต\-ปจฺจยา. ทสฺสเนน ปหาตพฺพ\-เหตุกํ เอกํ ขนฺธํ ปฏิจฺจ ตโย ขนฺธา จิตฺต\-สมุฏฺฐานญฺจ รูปํ ฯเปฯ เทฺว ขนฺเธ ปฏิจฺจ เทฺว ขนฺธา จิตฺต\-สมุฏฺฐานญฺจ รูปํ. วิจิกิจฺฉา\-สหคตํ เอกํ ขนฺธํ ปฏิจฺจ ตโย ขนฺธา โมโห จ จิตฺต\-สมุฏฺฐานญฺจ รูปํ ฯเปฯ. (3)}

\prose{ภาวนาย ปหาตพฺพ\-เหตุกํ ธมฺมํ ปฏิจฺจ ภาวนาย ปหาตพฺพ\-เหตุโก ธมฺโม อุปฺปชฺชติ สหชาต\-ปจฺจยา. ภาวนาย ปหาตพฺพ\-เหตุกํ เอกํ ขนฺธํ ปฏิจฺจ ตโย ขนฺธา ฯเปฯ เทฺว ขนฺเธ ปฏิจฺจ เทฺว ขนฺธา. (1)}

\prose{ภาวนาย ปหาตพฺพ\-เหตุกํ ธมฺมํ ปฏิจฺจ เนวทสฺสเนน นภาวนาย ปหาตพฺพ\-เหตุโก ธมฺโม อุปฺปชฺชติ สหชาต\-ปจฺจยา. ภาวนาย ปหาตพฺพ\-เหตุเก ขนฺเธ ปฏิจฺจ จิตฺต\-สมุฏฺฐานํ รูปํ. อุทฺธจฺจ\-สหคเต ขนฺเธ ปฏิจฺจ โมโห จ จิตฺต\-สมุฏฺฐานญฺจ รูปํ. (2)}

\prose{ภาวนาย ปหาตพฺพ\-เหตุกํ ธมฺมํ ปฏิจฺจ ภาวนาย ปหาตพฺพ\-เหตุโก จ เนวทสฺสเนน นภาวนาย ปหาตพฺพ\-เหตุโก จ ธมฺมา อุปฺปชฺชนฺติ สหชาต\-ปจฺจยา. ภาวนาย ปหาตพฺพ\-เหตุกํ เอกํ ขนฺธํ ปฏิจฺจ ตโย ขนฺธา จิตฺต\-สมุฏฺฐานญฺจ รูปํ ฯเปฯ เทฺว ขนฺเธ ปฏิจฺจ เทฺว ขนฺธา จิตฺต\-สมุฏฺฐานญฺจ รูปํ. อุทฺธจฺจ\-สหคตํ เอกํ ขนฺธํ ปฏิจฺจ ตโย ขนฺธา โมโห จ จิตฺต\-สมุฏฺฐานญฺจ รูปํ ฯเปฯ. (3)}

\proseitem{14}{เนวทสฺสเนน นภาวนาย ปหาตพฺพ\-เหตุกํ ธมฺมํ ปฏิจฺจ เนวทสฺสเนน นภาวนาย ปหาตพฺพ\-เหตุโก ธมฺโม อุปฺปชฺชติ สหชาต\-ปจฺจยา. เนวทสฺสเนน นภาวนาย ปหาตพฺพ\-เหตุกํ เอกํ ขนฺธํ ปฏิจฺจ ตโย ขนฺธา จิตฺต\-สมุฏฺฐานญฺจ รูปํ ฯเปฯ เทฺว ขนฺเธ ฯเปฯ. วิจิกิจฺฉา\-สหคตํ อุทฺธจฺจ\-สหคตํ โมหํ ปฏิจฺจ จิตฺต\-สมุฏฺฐานํ รูปํ. ปฏิสนฺธิ\-กฺขเณ ฯเปฯ. ขนฺเธ ปฏิจฺจ วตฺถุ, วตฺถุํ ปฏิจฺจ ขนฺธา. เอกํ มหาภูตํ ปฏิจฺจ ตโย มหาภูตา ฯเปฯ. พาหิรํ. อาหารสมุฏฺฐานํ. อุตุสมุฏฺฐานํ. อสญฺญสตฺตานํ เอกํ ฯเปฯ. (1)}

\prose{เนวทสฺสเนน นภาวนาย ปหาตพฺพ\-เหตุกํ ธมฺมํ ปฏิจฺจ ทสฺสเนน ปหาตพฺพ\-เหตุโก ธมฺโม อุปฺปชฺชติ สหชาต\-ปจฺจยา. วิจิกิจฺฉา\-สหคตํ โมหํ ปฏิจฺจ สมฺปยุตฺตกา ขนฺธา. (2)}

\prose{\csromanfolio{185}เนวทสฺสเนน นภาวนาย ปหาตพฺพ\-เหตุกํ ธมฺมํ ปฏิจฺจ ภาวนาย ปหาตพฺพ\-เหตุโก ธมฺโม อุปฺปชฺชติ สหชาต\-ปจฺจยา. (สํขิตฺตํ.) (เหตุสทิสํ กาตพฺพํ.) (3)}

\csromanheader{2}{\textbf{อญฺญมญฺญาทิ}}
\proseitem{15}{ทสฺสเนน ปหาตพฺพ\-เหตุกํ ธมฺมํ ปฏิจฺจ ทสฺสเนน ปหาตพฺพ\-เหตุโก ธมฺโม อุปฺปชฺชติ อญฺญมญฺญ\-ปจฺจยา. นิสฺสยปจฺจยา. อุปนิสฺสย\-ปจฺจยา. ปุเรชาต\-ปจฺจยา. อาเสวนปจฺจยา. กมฺมปจฺจยา. วิปากปจฺจยา. อาหารปจฺจยา. อินฺทฺริยปจฺจยา. ฌานปจฺจยา. มคฺคปจฺจยา. สมฺปยุตฺต\-ปจฺจยา. วิปฺปยุตฺต\-ปจฺจยา. อตฺถิปจฺจยา. นตฺถิปจฺจยา. วิคตปจฺจยา. อวิคตปจฺจยา.}

\csromanmarkh{1. ปจฺจยานุโลม}\csromanmarkhii{2. สงฺขฺยาวาร}\csromanheadernomark{2}{1. \textbf{ปจฺจยานุโลม 2. สงฺขฺยาวาร}}
\proseitemwithcloserruled{16}{เหตุยา สตฺตรส, อารมฺมเณ เอกาทส, อธิปติยา นว, อนนฺตเร เอกาทส, สมนนฺตเร เอกาทส, สหชาเต สตฺตรส, อญฺญมญฺเญ เอกาทส, นิสฺสเย สตฺตรส, อุปนิสฺสเย เอกาทส, ปุเรชาเต เอกาทส, อาเสวเน เอกาทส, กมฺเม สตฺตรส, วิปาเก เอกํ, อาหาเร สตฺตรส, อินฺทฺริเย สตฺตรส, ฌาเน สตฺตรส, มคฺเค สตฺตรส, สมฺปยุตฺเต เอกาทส, วิปฺปยุตฺเต สตฺตรส, อตฺถิยา สตฺตรส, นตฺถิยา เอกาทส, วิคเต เอกาทส, อวิคเต สตฺตรส. (เอวํ คเณตพฺพํ.)}{อนุโลมํ.}
\csromanmarkh{2. ปจฺจย\-ปจฺจนีย}\csromanmarkhii{1. วิภงฺควาร}\csromanheadernomark{2}{2. ปจฺจย\-ปจฺจนีย 1. วิภงฺควาร}
\csromanheader{2}{\textbf{นเหตุ}}
\proseitem{17}{ทสฺสเนน ปหาตพฺพ\-เหตุกํ ธมฺมํ ปฏิจฺจ เนวทสฺสเนน นภาวนาย ปหาตพฺพ\-เหตุโก ธมฺโม อุปฺปชฺชติ นเหตุปจฺจยา. วิจิกิจฺฉา\-สหคเต ขนฺเธ ปฏิจฺจ วิจิกิจฺฉา\-สหคโต โมโห. (1)}

\prose{\csromanfolio{186}ภาวนาย ปหาตพฺพ\-เหตุกํ ธมฺมํ ปฏิจฺจ เนวทสฺสเนน นภาวนาย ปหาตพฺพ\-เหตุโก ธมฺโม อุปฺปชฺชติ นเหตุปจฺจยา. อุทฺธจฺจ\-สหคเต ขนฺเธ ปฏิจฺจ อุทฺธจฺจ\-สหคโต โมโห. (1)}

\prose{เนวทสฺสเนน นภาวนาย ปหาตพฺพ\-เหตุกํ ธมฺมํ ปฏิจฺจ เนวทสฺสเนน นภาวนาย ปหาตพฺพ\-เหตุโก ธมฺโม อุปฺปชฺชติ นเหตุปจฺจยา. อเหตุกํ เนวทสฺสเนน นภาวนาย ปหาตพฺพ\-เหตุกํ เอกํ ขนฺธํ ปฏิจฺจ ตโย ขนฺธา จิตฺต\-สมุฏฺฐานญฺจ รูปํ ฯเปฯ เทฺว ขนฺธา. อเหตุก\-ปฏิสนฺธิ\-กฺขเณ ฯเปฯ. ขนฺเธ ปฏิจฺจ วตฺถุ, วตฺถุํ ปฏิจฺจ ขนฺธา. เอกํ มหาภูตํ ปฏิจฺจ ฯเปฯ. พาหิรํ. อาหารสมุฏฺฐานํ. อุตุสมุฏฺฐานํ. อสญฺญสตฺตานํ ฯเปฯ. (1)}

\prose{นอารมฺมณ}

\proseitem{18}{ทสฺสเนน ปหาตพฺพ\-เหตุกํ ธมฺมํ ปฏิจฺจ เนวทสฺสเนน นภาวนาย ปหาตพฺพ\-เหตุโก ธมฺโม อุปฺปชฺชติ นอารมฺมณ\-ปจฺจยา. ทสฺสเนน ปหาตพฺพ\-เหตุเก ขนฺเธ ปฏิจฺจ จิตฺต\-สมุฏฺฐานํ รูปํ. (1)}

\prose{ภาวนาย ปหาตพฺพ\-เหตุกํ ธมฺมํ ปฏิจฺจ เนวทสฺสเนน นภาวนาย ปหาตพฺพ\-เหตุโก ธมฺโม อุปฺปชฺชติ นอารมฺมณ\-ปจฺจยา. ภาวนาย ปหาตพฺพ\-เหตุเก ขนฺเธ ปฏิจฺจ จิตฺต\-สมุฏฺฐานํ รูปํ. (1)}

\prose{เนวทสฺสเนน นภาวนาย ปหาตพฺพ\-เหตุกํ ธมฺมํ ปฏิจฺจ เนวทสฺสเนน นภาวนาย ปหาตพฺพ\-เหตุโก ธมฺโม อุปฺปชฺชติ นอารมฺมณ\-ปจฺจยา. เนวทสฺสเนน นภาวนาย ปหาตพฺพ\-เหตุเก ขนฺเธ ปฏิจฺจ จิตฺต\-สมุฏฺฐานํ รูปํ. วิจิกิจฺฉา\-สหคตํ อุทฺธจฺจ\-สหคตํ โมหํ ปฏิจฺจ จิตฺต\-สมุฏฺฐานํ รูปํ. ปฏิสนฺธิ\-กฺขเณ เนวทสฺสเนน นภาวนาย ปหาตพฺพ\-เหตุเก ขนฺเธ ปฏิจฺจ กฏตฺตารูปํ. ขนฺเธ ปฏิจฺจ วตฺถุ ฯเปฯ. เอกํ มหาภูตํ ฯเปฯ. อสญฺญสตฺตานํ ฯเปฯ. (1)}

\proseitem{19}{ทสฺสเนน ปหาตพฺพ\-เหตุ\-กญฺจ เนวทสฺสเนน นภาวนาย ปหาตพฺพ\-เหตุ\-กญฺจ ธมฺมํ ปฏิจฺจ เนวทสฺสเนน นภาวนาย ปหาตพฺพ\-เหตุโก ธมฺโม อุปฺปชฺชติ นอารมฺมณ\-ปจฺจยา. ทสฺสเนน ปหาตพฺพ\-เหตุเก ขนฺเธ จ มหาภูเต จ ปฏิจฺจ จิตฺต\-สมุฏฺฐานํ รูปํ. วิจิกิจฺฉา\-สหคเต ขนฺเธ จ โมหญฺจ ปฏิจฺจ จิตฺต\-สมุฏฺฐานํ รูปํ. (1)}

\prose{\csromanfolio{187}ภาวนาย ปหาตพฺพ\-เหตุ\-กญฺจ เนวทสฺสเนน นภาวนาย ปหาตพฺพ\-เหตุ\-กญฺจ ธมฺมํ ปฏิจฺจ เนวทสฺสเนน นภาวนาย ปหาตพฺพ\-เหตุโก ธมฺโม อุปฺปชฺชติ นอารมฺมณ\-ปจฺจยา. ภาวนาย ปหาตพฺพ\-เหตุเก ขนฺเธ จ มหาภูเต จ ปฏิจฺจ จิตฺต\-สมุฏฺฐานํ รูปํ. อุทฺธจฺจ\-สหคเต ขนฺเธ จ โมหญฺจ ปฏิจฺจ จิตฺต\-สมุฏฺฐานํ รูปํ. (1)}

\prose{นอธิปตฺยาทิ}

\proseitem{20}{ทสฺสเนน ปหาตพฺพ\-เหตุกํ ธมฺมํ ปฏิจฺจ เนวทสฺสเนน นภาวนาย ปหาตพฺพ\-เหตุโก ธมฺโม อุปฺปชฺชติ นอธิปติ\-ปจฺจยา. (สหชาต\-สทิสํ.) นอนนฺตร\-ปจฺจยา. นสมนนฺตร\-ปจฺจยา. นอญฺญมญฺญ\-ปจฺจยา. นอุปนิสฺสยปจฺจยา.}

\csromanheader{2}{\textbf{นปุเรชาต}}
\proseitem{21}{ทสฺสเนน ปหาตพฺพ\-เหตุกํ ธมฺมํ ปฏิจฺจ ทสฺสเนน ปหาตพฺพ\-เหตุโก ธมฺโม อุปฺปชฺชติ นปุเรชาต\-ปจฺจยา. อรูเป ทสฺสเนน ปหาตพฺพ\-เหตุกํ เอกํ ขนฺธํ ปฏิจฺจ ฯเปฯ. (1)}

\prose{ทสฺสเนน ปหาตพฺพ\-เหตุกํ ธมฺมํ ปฏิจฺจ เนวทสฺสเนน นภาวนาย ปหาตพฺพ\-เหตุโก ธมฺโม อุปฺปชฺชติ นปุเรชาต\-ปจฺจยา. อรูเป วิจิกิจฺฉา\-สหคเต ขนฺเธ ปฏิจฺจ วิจิกิจฺฉา\-สหคโต โมโห. ทสฺสเนน ปหาตพฺพ\-เหตุเก ขนฺเธ ปฏิจฺจ จิตฺต\-สมุฏฺฐานํ รูปํ. (2)}

\prose{ทสฺสเนน ปหาตพฺพ\-เหตุกํ ธมฺมํ ปฏิจฺจ ทสฺสเนน ปหาตพฺพ\-เหตุโก จ เนวทสฺสเนน นภาวนาย ปหาตพฺพ\-เหตุโก จ ธมฺมา อุปฺปชฺชนฺติ นปุเรชาต\-ปจฺจยา. อรูเป วิจิกิจฺฉา\-สหคตํ เอกํ ขนฺธํ ปฏิจฺจ ตโย ขนฺธา โมโห จ ฯเปฯ เทฺว ขนฺเธ ฯเปฯ. (3)}

\prose{ภาวนาย ปหาตพฺพ\-เหตุกํ ธมฺมํ. ตีณิ. (ทสฺสเนน สทิสํ.)}

\proseitem{22}{เนวทสฺสเนน นภาวนาย ปหาตพฺพ\-เหตุกํ ธมฺมํ ปฏิจฺจ เนวทสฺสเนน นภาวนาย ปหาตพฺพ\-เหตุโก ธมฺโม อุปฺปชฺชติ นปุเรชาต\-ปจฺจยา. อรูเป เนวทสฺสเนน นภาวนาย ปหาตพฺพ\-เหตุกํ เอกํ ขนฺธํ ปฏิจฺจ ตโย ขนฺธา ฯเปฯ เทฺว ขนฺธา. เนวทสฺสเนน นภาวนาย ปหาตพฺพ\-เหตุเก ขนฺเธ ปฏิจฺจ จิตฺต\-สมุฏฺฐานํ รูปํ. วิจิกิจฺฉา\-สหคตํ \csromanfolio{188}อุทฺธจฺจ\-สหคตํ โมหํ ปฏิจฺจ จิตฺต\-สมุฏฺฐานํ รูปํ. ปฏิสนฺธิ\-กฺขเณ ฯเปฯ. อสญฺญสตฺตานํ ฯเปฯ. (1)}

\prose{เนวทสฺสเนน นภาวนาย ปหาตพฺพ\-เหตุกํ ธมฺมํ ปฏิจฺจ ทสฺสเนน ปหาตพฺพ\-เหตุโก ธมฺโม อุปฺปชฺชติ นปุเรชาต\-ปจฺจยา. อรูเป วิจิกิจฺฉา\-สหคตํ โมหํ ปฏิจฺจ สมฺปยุตฺตกา ขนฺธา. (2)}

\prose{เนวทสฺสเนน นภาวนาย ปหาตพฺพ\-เหตุกํ ธมฺมํ ปฏิจฺจ ภาวนาย ปหาตพฺพ\-เหตุโก ธมฺโม อุปฺปชฺชติ นปุเรชาต\-ปจฺจยา. อรูเป อุทฺธสหคตํ โมหํ ปฏิจฺจ สมฺปยุตฺตกา ขนฺธา. (3)}

\proseitem{23}{ทสฺสเนน ปหาตพฺพ\-เหตุ\-กญฺจ เนวทสฺสเนน นภาวนาย ปหาตพฺพ\-เหตุ\-กญฺจ ธมฺมํ ปฏิจฺจ ทสฺสเนน ปหาตพฺพ\-เหตุโก ธมฺโม อุปฺปชฺชติ นปุเรชาต\-ปจฺจยา. อรูเป วิจิกิจฺฉา\-สหคตํ เอกํ ขนฺธญฺจ โมหญฺจ ปฏิจฺจ ตโย ขนฺธา ฯเปฯ. (1)}

\prose{ทสฺสเนน ปหาตพฺพ\-เหตุ\-กญฺจ เนวทสฺสเนน นภาวนาย ปหาตพฺพ\-เหตุ\-กญฺจ ธมฺมํ ปฏิจฺจ เนวทสฺสเนน นภาวนาย ปหาตพฺพ\-เหตุโก ธมฺโม อุปฺปชฺชติ นปุเรชาต\-ปจฺจยา. ทสฺสเนน ปหาตพฺพ\-เหตุเก ขนฺเธ จ มหาภูเต จ ปฏิจฺจ จิตฺต\-สมุฏฺฐานํ รูปํ. วิจิกิจฺฉา\-สหคเต ขนฺเธ จ โมหญฺจ ปฏิจฺจ จิตฺต\-สมุฏฺฐานํ รูปํ. (2)}

\prose{ภาวนาย ปหาตพฺพ\-เหตุ\-กญฺจ เนวทสฺสเนน นภาวนาย ปหาตพฺพ\-เหตุ\-กญฺจ ธมฺมํ ปฏิจฺจ ภาวนาย ปหาตพฺพ\-เหตุโก ธมฺโม อุปฺปชฺชติ นปุเรชาต\-ปจฺจยา. (อิเมปิ เทฺว กาตพฺพา.)}

\csromanheader{2}{\textbf{นปจฺฉาชาตาทิ}}
\proseitem{24}{ทสฺสเนน ปหาตพฺพ\-เหตุกํ ธมฺมํ ปฏิจฺจ ทสฺสเนน ปหาตพฺพ\-เหตุโก ธมฺโม อุปฺปชฺชติ นปจฺฉาชาต\-ปจฺจยา. นอาเสวน\-ปจฺจยา.}

\csromanheader{2}{\textbf{นกมฺม}}
\proseitem{25}{ทสฺสเนน ปหาตพฺพ\-เหตุกํ ธมฺมํ ปฏิจฺจ ทสฺสเนน ปหาตพฺพ\-เหตุโก ธมฺโม อุปฺปชฺชติ นกมฺมปจฺจยา. ทสฺสเนน ปหาตพฺพ\-เหตุเก ขนฺเธ ปฏิจฺจ ทสฺสเนน ปหาตพฺพ\-เหตุกา เจตนา. (1)}

\prose{\csromanfolio{189}ภาวนาย ปหาตพฺพ\-เหตุกํ ธมฺมํ ปฏิจฺจ ภาวนาย ปหาตพฺพ\-เหตุโก ธมฺโม อุปฺปชฺชติ นกมฺมปจฺจยา. ภาวนาย ปหาตพฺพ\-เหตุเก ขนฺเธ ปฏิจฺจ ภาวนาย ปหาตพฺพ\-เหตุกา เจตนา. (1)}

\prose{เนวทสฺสเนน นภาวนาย ปหาตพฺพ\-เหตุกํ ธมฺมํ ปฏิจฺจ เนวทสฺสเนน นภาวนาย ปหาตพฺพ\-เหตุโก ธมฺโม อุปฺปชฺชติ นกมฺมปจฺจยา. เนวทสฺสเนน นภาวนาย ปหาตพฺพ\-เหตุเก ขนฺเธ ปฏิจฺจ เนวทสฺสเนน นภาวนาย ปหาตพฺพ\-เหตุกา เจตนา. พาหิรํ. อาหารสมุฏฺฐานํ. อุตุสมุฏฺฐานํ ฯเปฯ. (1)}

\prose{เนวทสฺสเนน นภาวนาย ปหาตพฺพ\-เหตุกํ ธมฺมํ ปฏิจฺจ ทสฺสเนน ปหาตพฺพ\-เหตุโก ธมฺโม อุปฺปชฺชติ นกมฺมปจฺจยา. วิจิกิจฺฉา\-สหคตํ โมหํ ปฏิจฺจ สมฺปยุตฺตกา เจตนา. (2)}

\prose{เนวทสฺสเนน นภาวนาย ปหาตพฺพ\-เหตุกํ ธมฺมํ ปฏิจฺจ ภาวนาย ปหาตพฺพ\-เหตุโก ธมฺโม อุปฺปชฺชติ นกมฺมปจฺจยา. อุทฺธจฺจ\-สหคตํ โมหํ ปฏิจฺจ สมฺปยุตฺตกา เจตนา. (3)}

\proseitem{26}{ทสฺสเนน ปหาตพฺพ\-เหตุ\-กญฺจ เนวทสฺสเนน นภาวนาย ปหาตพฺพ\-เหตุ\-กญฺจ ธมฺมํ ปฏิจฺจ ทสฺสเนน ปหาตพฺพ\-เหตุโก ธมฺโม อุปฺปชฺชติ นกมฺมปจฺจยา. วิจิกิจฺฉา\-สหคเต ขนฺเธ จ โมหญฺจ ปฏิจฺจ สมฺปยุตฺตกา เจตนา. (1)}

\prose{ภาวนาย ปหาตพฺพ\-เหตุ\-กญฺจ เนวทสฺสเนน นภาวนาย ปหาตพฺพ\-เหตุ\-กญฺจ ธมฺมํ ปฏิจฺจ ภาวนาย ปหาตพฺพ\-เหตุโก ธมฺโม อุปฺปชฺชติ นกมฺมปจฺจยา. อุทฺธจฺจ\-สหคเต ขนฺเธ จ โมหญฺจ ปฏิจฺจ สมฺปยุตฺตกา เจตนา. (1)}

\proseitem{27}{ทสฺสเนน ปหาตพฺพ\-เหตุกํ ธมฺมํ ปฏิจฺจ ทสฺสเนน ปหาตพฺพ\-เหตุโก ธมฺโม อุปฺปชฺชติ นวิปาก\-ปจฺจยา. (ปฏิสนฺธิ นตฺถิ.)}

\prose{นอาหาราทิ}

\proseitem{28}{เนวทสฺสเนน นภาวนาย ปหาตพฺพ\-เหตุกํ ธมฺมํ ปฏิจฺจ เนวทสฺสเนน นภาวนาย ปหาตพฺพ\-เหตุโก ธมฺโม อุปฺปชฺชติ นอาหารปจฺจยา. พาหิรํ. อุตุสมุฏฺฐานํ. อสญฺญสตฺตานํ ฯเปฯ นอินฺทฺริยปจฺจยา. พาหิรํ. \csromanfolio{190}อาหารสมุฏฺฐานํ. อุตุสมุฏฺฐานํ. อสญฺญสตฺตานํ ฯเปฯ มหาภูเต ปฏิจฺจ รูปชีวิตินฺทฺริยํ. นฌานปจฺจยา. ปญฺจวิญฺญาณํ ฯเปฯ. (มหาภูตา กาตพฺพา.) นมคฺคปจฺจยา. อเหตุกํ เนวทสฺสเนน นภาวนาย ปหาตพฺพ\-เหตุกํ เอกํ ขนฺธํ ฯเปฯ. อสญฺญสตฺตานํ ฯเปฯ นสมฺปยุตฺต\-ปจฺจยา.}

\csromanheader{2}{\textbf{นวิปฺปยุตฺตาทิ}}
\proseitem{29}{ทสฺสเนน ปหาตพฺพ\-เหตุกํ ธมฺมํ ปฏิจฺจ ทสฺสเนน ปหาตพฺพ\-เหตุโก ธมฺโม อุปฺปชฺชติ นวิปฺปยุตฺต\-ปจฺจยา. อรูเป ทสฺสเนน ปหาตพฺพ\-เหตุกํ เอกํ ขนฺธํ ปฏิจฺจ ตโย ขนฺธา ฯเปฯ. (1)}

\prose{ทสฺสเนน ปหาตพฺพ\-เหตุกํ ธมฺมํ ปฏิจฺจ เนวทสฺสเนน นภาวนาย ปหาตพฺพ\-เหตุโก ธมฺโม อุปฺปชฺชติ นวิปฺปยุตฺต\-ปจฺจยา. อรูเป วิจิกิจฺฉา\-สหคเต ขนฺเธ ปฏิจฺจ วิจิกิจฺฉา\-สหคโต โมโห. (2)}

\prose{ทสฺสเนน ปหาตพฺพ\-เหตุกํ ธมฺมํ ปฏิจฺจ ทสฺสเนน ปหาตพฺพ\-เหตุโก จ เนวทสฺสเนน นภาวนาย ปหาตพฺพ\-เหตุโก จ ธมฺมา อุปฺปชฺชนฺติ นวิปฺปยุตฺต\-ปจฺจยา. อรูเป วิจิกิจฺฉา\-สหคตํ เอกํ ขนฺธํ ปฏิจฺจ ตโย ขนฺธา โมโห จ ฯเปฯ เทฺว ขนฺเธ ฯเปฯ. (3)}

\prose{ภาวนาย ปหาตพฺพ\-เหตุกํ ธมฺมํ ปฏิจฺจ ภาวนาย ปหาตพฺพ\-เหตุโก ธมฺโม อุปฺปชฺชติ นวิปฺปยุตฺต\-ปจฺจยา. อรูเป ภาวนาย. ตีณิ.}

\proseitem{30}{เนวทสฺสเนน นภาวนาย ปหาตพฺพ\-เหตุกํ ธมฺมํ ปฏิจฺจ เนวทสฺสเนน นภาวนาย ปหาตพฺพ\-เหตุโก ธมฺโม อุปฺปชฺชติ นวิปฺปยุตฺต\-ปจฺจยา. อรูเป เนวทสฺสเนน นภาวนาย ปหาตพฺพ\-เหตุกํ เอกํ ขนฺธํ ปฏิจฺจ ตโย ขนฺธา ฯเปฯ เทฺว ขนฺเธ ปฏิจฺจ เทฺว ขนฺธา. พาหิรํ. อาหารสมุฏฺฐานํ. อุตุสมุฏฺฐานํ. อสญฺญสตฺตานํ ฯเปฯ. (1)}

\prose{เนวทสฺสเนน นภาวนาย ปหาตพฺพ\-เหตุกํ ธมฺมํ ปฏิจฺจ ทสฺสเนน ปหาตพฺพ\-เหตุโก ธมฺโม อุปฺปชฺชติ นวิปฺปยุตฺต\-ปจฺจยา. อรูเป วิจิกิจฺฉา สหคตํ โมหํ ปฏิจฺจ สมฺปยุตฺตกา ขนฺธา. (2)}

\prose{เนวทสฺสเนน นภาวนาย ปหาตพฺพ\-เหตุกํ ธมฺมํ ปฏิจฺจ ภาวนาย ปหาตพฺพ\-เหตุโก ธมฺโม อุปฺปชฺชติ นวิปฺปยุตฺต\-ปจฺจยา. อรูเป อุทฺธจฺจ\-สหคตํ โมหํ ปฏิจฺจ สมฺปยุตฺตกา ขนฺธา. (3)}

\proseitem{31}{\csromanfolio{191}ทสฺสเนน ปหาตพฺพ\-เหตุ\-กญฺจ เนวทสฺสเนน นภาวนาย ปหาตพฺพ\-เหตุ\-กญฺจ ธมฺมํ ปฏิจฺจ ทสฺสเนน ปหาตพฺพ\-เหตุโก ธมฺโม อุปฺปชฺชติ นวิปฺปยุตฺต\-ปจฺจยา. อรูเป วิจิกิจฺฉา\-สหคตํ เอกํ ขนฺธญฺจ โมหญฺจ ปฏิจฺจ ตโย ขนฺธา ฯเปฯ เทฺว ขนฺธา. (1)}

\prose{ภาวนาย ปหาตพฺพ\-เหตุ\-กญฺจ เนวทสฺสเนน นภาวนาย ปหาตพฺพ\-เหตุ\-กญฺจ ธมฺมํ ปฏิจฺจ ภาวนาย ปหาตพฺพ\-เหตุโก ธมฺโม อุปฺปชฺชติ นวิปฺปยุตฺต\-ปจฺจยา. อรูเป อุทฺธจฺจ\-สหคตํ เอกํ ขนฺธญฺจ โมหญฺจ ปฏิจฺจ ตโย ขนฺธา ฯเปฯ เทฺว ขนฺธา. โนนตฺถิปจฺจยา. โนวิคต\-ปจฺจยา.}

\csromanmarkh{2. ปจฺจย\-ปจฺจนีย}\csromanmarkhii{2. สงฺขฺยาวาร}\csromanheadernomark{2}{2. ปจฺจย\-ปจฺจนีย 2. สงฺขฺยาวาร}
\proseitemwithcloserruled{32}{นเหตุยา ตีณิ, นอารมฺมเณ ปญฺจ, นอธิปติยา สตฺตรส, นอนนฺตเร ปญฺจ, นสมนนฺตเร ปญฺจ, นอญฺญมญฺเญ ปญฺจ, นอุปนิสฺสเย ปญฺจ, นปุเรชาเต เตรส, นปจฺฉาชาเต สตฺตรส, นอาเสวเน สตฺตรส, นกมฺเม สตฺต, นวิปาเก สตฺตรส, นอาหาเร เอกํ, นอินฺทฺริเย เอกํ, นฌาเน เอกํ, นมคฺเค เอกํ, นสมฺปยุตฺเต ปญฺจ, นวิปฺปยุตฺเต เอกาทส, โนนตฺถิยา ปญฺจ, โนวิคเต ปญฺจ. (เอวํ คเณตพฺพํ.)}{ปจฺจนียํ.}
\proseitemwithcloserruled{33}{เหตุปจฺจยา นอารมฺมเณ ปญฺจ, นอธิปติยา สตฺตรส, นอนนฺตเร ปญฺจ, นสมนนฺตเร ปญฺจ, นอญฺญมญฺเญ ปญฺจ, นอุปนิสฺสเย ปญฺจ, นปุเรชาเต เตรส, นปจฺฉาชาเต สตฺตรส, นอาเสวเน สตฺตรส, นกมฺเม สตฺต, นวิปาเก สตฺตรส, นสมฺปยุตฺเต ปญฺจ, นวิปฺปยุตฺเต เอกาทส, โนนตฺถิยา ปญฺจ, โนวิคเต ปญฺจ. (เอวํ คเณตพฺพํ.)}{อนุโลม\-ปจฺจนียํ.}
\proseitem{34}{\csromanfolio{192}นเหตุปจฺจยา อารมฺมเณ ตีณิ, อนนฺตเร ตีณิ, สมนนฺตเร ตีณิ, สหชาเต ตีณิ, อญฺญมญฺเญ ตีณิ, นิสฺสเย ตีณิ, อุปนิสฺสเย ตีณิ, ปุเรชาเต ตีณิ, อาเสวเน ตีณิ, กมฺเม ตีณิ, วิปาเก เอกํ, อาหาเร ตีณิ, อินฺทฺริเย ตีณิ, ฌาเน ตีณิ, มคฺเค เทฺว, สมฺปยุตฺเต ตีณิ, วิปฺปยุตฺเต ตีณิ, อตฺถิยา ตีณิ, นตฺถิยา ตีณิ, วิคเต ตีณิ, อวิคเต ตีณิ. (เอวํ คเณตพฺพํ.) ปจฺจนียา\-นุโลมํ. ปฏิจฺจวาโร. (สหชาตวาโร ปฏิจฺจ\-วาร\-สทิโส.)}
\csromansectionrule

\csromanmarkh{9. ทสฺสเนน\-ปหาตพฺพ\-เหตุ\-กตฺติก}\csromanmarkhii{3. ปจฺจยวาร}\csromanheadernomark{2}{9. ทสฺสเนน\-ปหาตพฺพ\-เหตุ\-กตฺติก 3. ปจฺจยวาร}
\csromanmarkh{1. ปจฺจยานุโลม}\csromanmarkhii{1. วิภงฺควาร}\csromanheadernomark{2}{1. \textbf{ปจฺจยานุโลม 1. วิภงฺควาร}}
\csromanheader{2}{\textbf{เหตุ}}
\proseitem{35}{ทสฺสเนน ปหาตพฺพ\-เหตุกํ ธมฺมํ ปจฺจยา ทสฺสเนน ปหาตพฺพ\-เหตุโก ธมฺโม อุปฺปชฺชติ เหตุปจฺจยา. ตีณิ. (ปฏิจฺจ\-วาร\-สทิสํ.)}

\prose{ภาวนาย ปหาตพฺพ\-เหตุกํ ธมฺมํ ปจฺจยา ตีณิ. (ปฏิจฺจ\-วาร\-สทิสํ.)}

\prose{เนวทสฺสเนน นภาวนาย ปหาตพฺพ\-เหตุกํ ธมฺมํ ปจฺจยา เนวทสฺสเนน นภาวนาย. เอกํ. (ปฏิจฺจ\-วาร\-สทิสํ.) วตฺถุํ ปจฺจยา เนวทสฺสเนน นภาวนาย ปหาตพฺพ\-เหตุกา ขนฺธา. (1)}

\prose{เนวทสฺสเนน นภาวนาย ปหาตพฺพ\-เหตุกํ ธมฺมํ ปจฺจยา ทสฺสเนน ปหาตพฺพ\-เหตุโก ธมฺโม อุปฺปชฺชติ เหตุปจฺจยา. วตฺถุํ ปจฺจยา ทสฺสเนน ปหาตพฺพ\-เหตุกา ขนฺธา. วิจิกิจฺฉา\-สหคตํ โมหํ ปจฺจยา สมฺปยุตฺตกา ขนฺธา. (2)}

\prose{เนวทสฺสเนน นภาวนาย ปหาตพฺพ\-เหตุกํ ธมฺมํ ปจฺจยา ภาวนาย ปหาตพฺพ\-เหตุโก ธมฺโม อุปฺปชฺชติ เหตุปจฺจยา. วตฺถุํ ปจฺจยา ภาวนาย \csromanfolio{193}ปหาตพฺพ\-เหตุกา ขนฺธา. อุทฺธจฺจ\-สหคตํ โมหํ ปจฺจยา สมฺปยุตฺตกา ขนฺธา. (3)}

\prose{เนวทสฺสเนน นภาวนาย ปหาตพฺพ\-เหตุกํ ธมฺมํ ปจฺจยา ทสฺสเนน ปหาตพฺพ\-เหตุโก จ เนวทสฺสเนน นภาวนาย ปหาตพฺพ\-เหตุโก จ ธมฺมา อุปฺปชฺชนฺติ เหตุปจฺจยา. วตฺถุํ ปจฺจยา ทสฺสเนน ปหาตพฺพ\-เหตุกา ขนฺธา. มหาภูเต ปจฺจยา จิตฺต\-สมุฏฺฐานํ รูปํ. วิจิกิจฺฉา\-สหคตํ โมหํ ปจฺจยา สมฺปยุตฺตกา ขนฺธา จิตฺต\-สมุฏฺฐานญฺจ รูปํ. (4)}

\prose{เนวทสฺสเนน นภาวนาย ปหาตพฺพ\-เหตุกํ ธมฺมํ ปจฺจยา ภาวนาย ปหาตพฺพ\-เหตุโก จ เนวทสฺสเนน นภาวนาย ปหาตพฺพ\-เหตุโก จ ธมฺมา อุปฺปชฺชนฺติ เหตุปจฺจยา. วตฺถุํ ปจฺจยา ภาวนาย ปหาตพฺพ\-เหตุกา ขนฺธา. มหาภูเต ปจฺจยา จิตฺต\-สมุฏฺฐานํ รูปํ. อุทฺธจฺจ\-สหคตํ โมหํ ปจฺจยา สมฺปยุตฺตกา ขนฺธา จิตฺต\-สมุฏฺฐานญฺจ รูปํ. (5)}

\proseitem{36}{ทสฺสเนน ปหาตพฺพ\-เหตุ\-กญฺจ เนวทสฺสเนน นภาวนาย ปหาตพฺพ\-เหตุ\-กญฺจ ธมฺมํ ปจฺจยา ทสฺสเนน ปหาตพฺพ\-เหตุโก ธมฺโม อุปฺปชฺชติ เหตุปจฺจยา. ทสฺสเนน ปหาตพฺพ\-เหตุกํ เอกํ ขนฺธญฺจ วตฺถุญฺจ ปจฺจยา ตโย ขนฺธา ฯเปฯ เทฺว ขนฺเธ จ ฯเปฯ. วิจิกิจฺฉา\-สหคตํ เอกํ ขนฺธญฺจ โมหญฺจ ปจฺจยา ตโย ขนฺธา ฯเปฯ เทฺว ขนฺเธ จ โมหญฺจ ปจฺจยา เทฺว ขนฺธา. (1)}

\prose{ทสฺสเนน ปหาตพฺพ\-เหตุ\-กญฺจ เนวทสฺสเนน นภาวนาย ปหาตพฺพ\-เหตุ\-กญฺจ ธมฺมํ ปจฺจยา เนวทสฺสเนน นภาวนาย ปหาตพฺพ\-เหตุโก ธมฺโม อุปฺปชฺชติ เหตุปจฺจยา. ทสฺสเนน ปหาตพฺพ\-เหตุเก ขนฺเธ จ มหาภูเต จ ปจฺจยา จิตฺต\-สมุฏฺฐานํ รูปํ. วิจิกิจฺฉา\-สหคเต ขนฺเธ จ โมหญฺจ ปจฺจยา จิตฺต\-สมุฏฺฐานํ รูปํ. (2)}

\prose{ทสฺสเนน ปหาตพฺพ\-เหตุ\-กญฺจ เนวทสฺสเนน นภาวนาย ปหาตพฺพ\-เหตุ\-กญฺจ ธมฺมํ ปจฺจยา ทสฺสเนน ปหาตพฺพ\-เหตุโก จ เนวทสฺสเนน นภาวนาย ปหาตพฺพ\-เหตุโก จ ธมฺมา อุปฺปชฺชนฺติ เหตุปจฺจยา. ทสฺสเนน ปหาตพฺพ\-เหตุกํ เอกํ ขนฺธญฺจ วตฺถุญฺจ ปจฺจยา ตโย ขนฺธา ฯเปฯ เทฺว ขนฺเธ จ วตฺถุญฺจ ปจฺจยา เทฺว ขนฺธา. ทสฺสเนน ปหาตพฺพ\-เหตุเก ขนฺเธ จ มหาภูเต จ ปจฺจยา จิตฺต\-สมุฏฺฐานํ รูปํ. วิจิกิจฺฉา\-สหคตํ เอกํ ขนฺธญฺจ โมหญฺจ ปจฺจยา ตโย ขนฺธา จิตฺต\-สมุฏฺฐานญฺจ รูปํ ฯเปฯ. (3)}

\prose{\csromanfolio{194}ภาวนาย ปหาตพฺพ\-เหตุ\-กญฺจ เนวทสฺสเนน นภาวนาย ปหาตพฺพ\-เหตุ\-กญฺจ ธมฺมํ ปจฺจยา ภาวนาย ปหาตพฺพ\-เหตุโก ธมฺโม อุปฺปชฺชติ เหตุปจฺจยา. ตีณิ.}

\csromanheader{2}{\textbf{อารมฺมณ}}
\proseitem{37}{ทสฺสเนน ปหาตพฺพ\-เหตุกํ ธมฺมํ ปจฺจยา ทสฺสเนน ปหาตพฺพ\-เหตุโก ธมฺโม อุปฺปชฺชติ อารมฺมณ\-ปจฺจยา. ตีณิ. (ปฏิจฺจวาเร อารมฺมณสทิสา.)}

\prose{ภาวนาย ปหาตพฺพ\-เหตุกํ ธมฺมํ ปจฺจยา ตีณิ. (ปฏิจฺจ\-วาร\-สทิสา.)}

\proseitem{38}{เนวทสฺสเนน นภาวนาย ปหาตพฺพ\-เหตุกํ ธมฺมํ ปจฺจยา เนวทสฺสเนน นภาวนาย ปหาตพฺพ\-เหตุโก ธมฺโม อุปฺปชฺชติ อารมฺมณ\-ปจฺจยา. เนวทสฺสเนน นภาวนาย ปหาตพฺพ\-เหตุกํ เอกํ ขนฺธํ ปจฺจยา ตโย ขนฺธา ฯเปฯ เทฺว ขนฺธา. ปฏิสนฺธิ\-กฺขเณ ฯเปฯ. วตฺถุํ ปจฺจยา ขนฺธา. จกฺขายตนํ ปจฺจยา จกฺขุวิญฺญาณํ ฯเปฯ. กายายตนํ ปจฺจยา กายวิญฺญาณํ. วตฺถุํ ปจฺจยา เนวทสฺสเนน นภาวนาย ปหาตพฺพ\-เหตุกา ขนฺธา. (1)}

\prose{เนวทสฺสเนน นภาวนาย ปหาตพฺพ\-เหตุกํ ธมฺมํ ปจฺจยา ทสฺสเนน ปหาตพฺพ\-เหตุโก ธมฺโม อุปฺปชฺชติ อารมฺมณ\-ปจฺจยา. วตฺถุํ ปจฺจยา ทสฺสเนน ปหาตพฺพ\-เหตุกา ขนฺธา. วิจิกิจฺฉา\-สหคตํ โมหํ ปจฺจยา สมฺปยุตฺตกา ขนฺธา. (2)}

\prose{เนวทสฺสเนน นภาวนาย ปหาตพฺพ\-เหตุกํ ธมฺมํ ปจฺจยา ภาวนาย ปหาตพฺพ\-เหตุโก ธมฺโม อุปฺปชฺชติ อารมฺมณ\-ปจฺจยา. วตฺถุํ ปจฺจยา ภาวนาย ปหาตพฺพ\-เหตุกา ขนฺธา. อุทฺธจฺจ\-สหคตํ โมหํ ปจฺจยา สมฺปยุตฺตกา ขนฺธา. (3)}

\prose{เนวทสฺสเนน นภาวนาย ปหาตพฺพ\-เหตุกํ ธมฺมํ ปจฺจยา ทสฺสเนน ปหาตพฺพ\-เหตุโก จ เนวทสฺสเนน นภาวนาย ปหาตพฺพ\-เหตุโก จ ธมฺมา อุปฺปชฺชนฺติ อารมฺมณ\-ปจฺจยา. วตฺถุํ ปจฺจยา วิจิกิจฺฉา\-สหคตา ขนฺธา จ โมโห จ. (4)}

\prose{\csromanfolio{195}เนวทสฺสเนน นภาวนาย ปหาตพฺพ\-เหตุกํ ธมฺมํ ปจฺจยา ภาวนาย ปหาตพฺพ\-เหตุโก จ เนวทสฺสเนน นภาวนาย ปหาตพฺพ\-เหตุโก จ ธมฺมา อุปฺปชฺชนฺติ อารมฺมณ\-ปจฺจยา. วตฺถุํ ปจฺจยา อุทฺธจฺจ\-สหคตา ขนฺธา จ โมโห จ. (5)}

\proseitem{39}{ทสฺสเนน ปหาตพฺพ\-เหตุ\-กญฺจ เนวทสฺสเนน นภาวนาย ปหาตพฺพ\-เหตุ\-กญฺจ ธมฺมํ ปจฺจยา ทสฺสเนน ปหาตพฺพ\-เหตุโก ธมฺโม อุปฺปชฺชติ อารมฺมณ\-ปจฺจยา. ทสฺสเนน ปหาตพฺพ\-เหตุกํ เอกํ ขนฺธญฺจ วตฺถุญฺจ ปจฺจยา ตโย ขนฺธา ฯเปฯ เทฺว ขนฺธา. วิจิกิจฺฉา\-สหคตํ เอกํ ขนฺธญฺจ โมหญฺจ ปจฺจยา ตโย ขนฺธา ฯเปฯ เทฺว ขนฺเธ จ โมหญฺจ ปจฺจยา เทฺว ขนฺธา. (1)}

\prose{ทสฺสเนน ปหาตพฺพ\-เหตุ\-กญฺจ เนวทสฺสเนน นภาวนาย ปหาตพฺพ\-เหตุ\-กญฺจ ธมฺมํ ปจฺจยา เนวทสฺสเนน นภาวนาย ปหาตพฺพ\-เหตุโก ธมฺโม อุปฺปชฺชติ อารมฺมณ\-ปจฺจยา. วิจิกิจฺฉา\-สหคเต ขนฺเธ จ วตฺถุญฺจ ปจฺจยา วิจิกิจฺฉา\-สหคโต โมโห. (2)}

\prose{ทสฺสเนน ปหาตพฺพ\-เหตุ\-กญฺจ เนวทสฺสเนน นภาวนาย ปหาตพฺพ\-เหตุ\-กญฺจ ธมฺมํ ปจฺจยา ทสฺสเนน ปหาตพฺพ\-เหตุโก จ เนวทสฺสเนน นภาวนาย ปหาตพฺพ\-เหตุโก จ ธมฺมา อุปฺปชฺชนฺติ อารมฺมณ\-ปจฺจยา. วิจิกิจฺฉา\-สหคตํ เอกํ ขนฺธญฺจ วตฺถุญฺจ ปจฺจยา ตโย ขนฺธา โมโห จ ฯเปฯ เทฺว ขนฺเธ จ วตฺถุญฺจ ฯเปฯ. (3)}

\proseitem{40}{ภาวนาย ปหาตพฺพ\-เหตุ\-กญฺจ เนวทสฺสเนน นภาวนาย ปหาตพฺพ\-เหตุ\-กญฺจ ธมฺมํ ปจฺจยา ภาวนาย ปหาตพฺพ\-เหตุโก ธมฺโม อุปฺปชฺชติ อารมฺมณ\-ปจฺจยา. ภาวนาย ปหาตพฺพ\-เหตุกํ เอกํ ขนฺธญฺจ วตฺถุญฺจ ปจฺจยา ตโย ขนฺธา ฯเปฯ เทฺว ขนฺธา. อุทฺธจฺจ\-สหคตํ เอกํ ขนฺธญฺจ โมหญฺจ ปจฺจยา ตโย ขนฺธา ฯเปฯ เทฺว ขนฺธา. (1)}

\prose{ภาวนาย ปหาตพฺพ\-เหตุ\-กญฺจ เนวทสฺสเนน นภาวนาย ปหาตพฺพ\-เหตุ\-กญฺจ ธมฺมํ ปจฺจยา เนวทสฺสเนน นภาวนาย ปหาตพฺพ\-เหตุโก ธมฺโม อุปฺปชฺชติ อารมฺมณ\-ปจฺจยา. อุทฺธจฺจ\-สหคเต ขนฺเธ จ วตฺถุญฺจ ปจฺจยา อุทฺธจฺจ\-สหคโต โมโห. (2)}

\prose{ภาวนาย ปหาตพฺพ\-เหตุ\-กญฺจ เนวทสฺสเนน นภาวนาย ปหาตพฺพ\-เหตุ\-กญฺจ ธมฺมํ ปจฺจยา ภาวนาย ปหาตพฺพ\-เหตุโก จ เนวทสฺสเนน \csromanfolio{196}นภาวนาย ปหาตพฺพ\-เหตุโก จ ธมฺมา อุปฺปชฺชนฺติ อารมฺมณ\-ปจฺจยา. อุทฺธจฺจ\-สหคตํ เอกํ ขนฺธญฺจ วตฺถุญฺจ ปจฺจยา ตโย ขนฺธา โมโห จ ฯเปฯ เทฺว ขนฺเธ จ ฯเปฯ. (3)}

\csromanheader{2}{\textbf{อธิปตฺยาทิ}}
\proseitem{41}{ทสฺสเนน ปหาตพฺพ\-เหตุกํ ธมฺมํ ปจฺจยา ทสฺสเนน ปหาตพฺพ\-เหตุโก ธมฺโม อุปฺปชฺชติ อธิปติ\-ปจฺจยา. ตีณิ.}

\prose{ภาวนาย ปหาตพฺพ\-เหตุกํ ธมฺมํ ปจฺจยา. ตีณิ.}

\prose{เนวทสฺสเนน นภาวนาย ปหาตพฺพ\-เหตุกํ ธมฺมํ ปจฺจยา เนวทสฺสเนน นภาวนาย ปหาตพฺพ\-เหตุโก ธมฺโม อุปฺปชฺชติ อธิปติ\-ปจฺจยา. เอกํ ฯเปฯ. วตฺถุํ ปจฺจยา เนวทสฺสเนน นภาวนาย ปหาตพฺพ\-เหตุกา ขนฺธา. (1)}

\prose{เนวทสฺสเนน นภาวนาย ปหาตพฺพ\-เหตุกํ ธมฺมํ ปจฺจยา ทสฺสเนน ปหาตพฺพ\-เหตุโก ธมฺโม อุปฺปชฺชติ อธิปติ\-ปจฺจยา. วตฺถุํ ปจฺจยา ทสฺสเนน ปหาตพฺพ\-เหตุกา ขนฺธา. (2)}

\prose{เนวทสฺสเนน นภาวนาย ปหาตพฺพ\-เหตุกํ ธมฺมํ ปจฺจยา ภาวนาย ปหาตพฺพ\-เหตุโก ธมฺโม อุปฺปชฺชติ อธิปติ\-ปจฺจยา. วตฺถุํ ปจฺจยา ภาวนาย ปหาตพฺพ\-เหตุกา ขนฺธา. (3)}

\prose{เนวทสฺสเนน นภาวนาย ปหาตพฺพ\-เหตุกํ ธมฺมํ ปจฺจยา ทสฺสเนน ปหาตพฺพ\-เหตุโก จ เนวทสฺสเนน นภาวนาย ปหาตพฺพ\-เหตุโก จ ธมฺมา อุปฺปชฺชนฺติ อธิปติ\-ปจฺจยา. วตฺถุํ ปจฺจยา ทสฺสเนน ปหาตพฺพ\-เหตุกา ขนฺธา. มหาภูเต ปจฺจยา จิตฺต\-สมุฏฺฐานํ รูปํ. (4)}

\prose{เนวทสฺสเนน นภาวนาย ปหาตพฺพ\-เหตุกํ ธมฺมํ ปจฺจยา ภาวนาย ปหาตพฺพ\-เหตุโก จ เนวทสฺสเนน นภาวนาย ปหาตพฺพ\-เหตุโก จ ธมฺมา อุปฺปชฺชนฺติ อธิปติ\-ปจฺจยา. วตฺถุํ ปจฺจยา ภาวนาย ปหาตพฺพ\-เหตุกา ขนฺธา. มหาภูเต ปจฺจยา จิตฺต\-สมุฏฺฐานํ รูปํ. (5)}

\proseitem{42}{ทสฺสเนน ปหาตพฺพ\-เหตุ\-กญฺจ เนวทสฺสเนน นภา\-วนาย\-ปหาตพฺพ\-เหตุ\-กญฺจ ธมฺมํ ปจฺจยา ทสฺสเนน ปหาตพฺพ\-เหตุโก ธมฺโม อุปฺปชฺชติ อธิปติ\-ปจฺจยา. ทสฺสเนน ปหาตพฺพ\-เหตุกํ เอกํ ขนฺธญฺจ วตฺถุญฺจ ปจฺจยา ตโย ขนฺธา ฯเปฯ เทฺว ขนฺเธ จ ฯเปฯ. (1)}

\prose{\csromanfolio{197}ทสฺสเนน ปหาตพฺพ\-เหตุ\-กญฺจ เนวทสฺสเนน นภาวนาย ปหาตพฺพ\-เหตุ\-กญฺจ ธมฺมํ ปจฺจยา เนวทสฺสเนน นภาวนาย ปหาตพฺพ\-เหตุโก ธมฺโม อุปฺปชฺชติ อธิปติ\-ปจฺจยา. ทสฺสเนน ปหาตพฺพ\-เหตุเก ขนฺเธ จ มหาภูเต จ ปจฺจยา จิตฺต\-สมุฏฺฐานํ รูปํ. (2)}

\prose{ทสฺสเนน ปหาตพฺพ\-เหตุ\-กญฺจ เนวทสฺสเนน นภาวนาย ปหาตพฺพ\-เหตุ\-กญฺจ ธมฺมํ ปจฺจยา ทสฺสเนน ปหาตพฺพ\-เหตุโก จ เนวทสฺสเนน นภาวนาย ปหาตพฺพ\-เหตุโก จ ธมฺมา อุปฺปชฺชนฺติ อธิปติ\-ปจฺจยา. ทสฺสเนน ปหาตพฺพ\-เหตุกํ เอกํ ขนฺธญฺจ วตฺถุญฺจ ปจฺจยา ตโย ขนฺธา ฯเปฯ เทฺว ขนฺเธ จ ฯเปฯ. ทสฺสเนน ปหาตพฺพ\-เหตุเก ขนฺเธ จ มหาภูเต จ ปจฺจยา จิตฺต\-สมุฏฺฐานํ รูปํ. (3)}

\prose{ภาวนาย ปหาตพฺพ\-เหตุ\-กญฺจ เนวทสฺสเนน นภาวนาย ปหาตพฺพ\-เหตุ\-กญฺจ ธมฺมํ ปจฺจยา ภาวนาย ปหาตพฺพ\-เหตุโก ธมฺโม อุปฺปชฺชติ อธิปติ\-ปจฺจยา. ภาวนาย ปหาตพฺพ\-เหตุกํ เอกํ ขนฺธญฺจ วตฺถุญฺจ ปจฺจยา ตโย ขนฺธา ฯเปฯ. ตีณิ. (ทสฺสเนน สทิสา.) อนนฺตร\-ปจฺจยา. สมนนฺตร\-ปจฺจยา.}

\csromanheader{2}{\textbf{สหชาต}}
\proseitem{43}{ทสฺสเนน ปหาตพฺพ\-เหตุกํ ธมฺมํ ปจฺจยา ทสฺสเนน ปหาตพฺพ\-เหตุโก ธมฺโม อุปฺปชฺชติ สหชาต\-ปจฺจยา. ทสฺสเนน ปหาตพฺพ\-เหตุกํ เอกํ ขนฺธํ ปจฺจยา ตโย ขนฺธา ฯเปฯ. (1)}

\prose{ทสฺสเนน ปหาตพฺพ\-เหตุกํ ธมฺมํ ปจฺจยา เนวทสฺสเนน นภาวนาย ปหาตพฺพ\-เหตุโก ธมฺโม อุปฺปชฺชติ สหชาต\-ปจฺจยา. ทสฺสเนน ปหาตพฺพ\-เหตุเก ขนฺเธ ปจฺจยา จิตฺต\-สมุฏฺฐานํ รูปํ. วิจิกิจฺฉา\-สหคเต ขนฺเธ ปจฺจยา วิจิกิจฺฉา\-สหคโต โมโห จิตฺต\-สมุฏฺฐานญฺจ รูปํ. (2)}

\prose{ทสฺสเนน ปหาตพฺพ\-เหตุกํ ธมฺมํ ปจฺจยา ทสฺสเนน ปหาตพฺพ\-เหตุโก จ เนวทสฺสเนน นภาวนาย ปหาตพฺพ\-เหตุโก จ ธมฺมา อุปฺปชฺชนฺติ สหชาต\-ปจฺจยา. ทสฺสเนน ปหาตพฺพ\-เหตุกํ เอกํ ขนฺธํ ปจฺจยา ตโย ขนฺธา จิตฺต\-สมุฏฺฐานญฺจ รูปํ ฯเปฯ. วิจิกิจฺฉา\-สหคตํ เอกํ ขนฺธํ ปจฺจยา ตโย ขนฺธา โมโห จ จิตฺต\-สมุฏฺฐานญฺจ รูปํ ฯเปฯ เทฺว ขนฺเธ ฯเปฯ. (3)}

\prose{\csromanfolio{198}ภาวนาย ปหาตพฺพ\-เหตุกํ ธมฺมํ ปจฺจยา ตีณิ. (สํขิตฺตํ.) (ทสฺสเนน สทิสา.)}

\proseitem{44}{เนวทสฺสเนน นภาวนาย ปหาตพฺพ\-เหตุกํ ธมฺมํ ปจฺจยา เนวทสฺสเนน นภาวนาย ปหาตพฺพ\-เหตุโก ธมฺโม อุปฺปชฺชติ สหชาต\-ปจฺจยา. เนวทสฺสเนน นภาวนาย ปหาตพฺพ\-เหตุกํ เอกํ ขนฺธํ ปจฺจยา ตโย ขนฺธา จิตฺต\-สมุฏฺฐานญฺจ รูปํ ฯเปฯ. วิจิกิจฺฉา\-สหคตํ อุทฺธจฺจ\-สหคตํ โมหํ ปจฺจยา จิตฺต\-สมุฏฺฐานํ รูปํ. ปฏิสนฺธิ\-กฺขเณ ฯเปฯ. ขนฺเธ ปจฺจยา วตฺถุ, วตฺถุํ ปจฺจยา ขนฺธา. เอกํ มหาภูตํ ปจฺจยา ตโย มหาภูตา ฯเปฯ. อสญฺญสตฺตานํ ฯเปฯ. จกฺขายตนํ ปจฺจยา จกฺขุวิญฺญาณํ ฯเปฯ. กายายตนํ ปจฺจยา กายวิญฺญาณํ. วตฺถุํ ปจฺจยา เนวทสฺสเนน นภาวนาย ปหาตพฺพ\-เหตุกา ขนฺธา. (1)}

\prose{เนวทสฺสเนน นภาวนาย ปหาตพฺพ\-เหตุกํ ธมฺมํ ปจฺจยา ทสฺสเนน ปหาตพฺพ\-เหตุโก ธมฺโม อุปฺปชฺชติ สหชาต\-ปจฺจยา. วตฺถุํ ปจฺจยา ทสฺสเนน ปหาตพฺพ\-เหตุกา ขนฺธา. วิจิกิจฺฉา\-สหคตํ โมหํ ปจฺจยา สมฺปยุตฺตกา ขนฺธา. (2)}

\prose{เนวทสฺสเนน นภาวนาย ปหาตพฺพ\-เหตุกํ ธมฺมํ ปจฺจยา ภาวนาย ปหาตพฺพ\-เหตุโก ธมฺโม อุปฺปชฺชติ สหชาต\-ปจฺจยา วตฺถุํ ปจฺจยา ภาวนาย ปหาตพฺพ\-เหตุกา ขนฺธา. อุทฺธจฺจ\-สหคตํ โมหํ ปจฺจยา สมฺปยุตฺตกา ขนฺธา. (3)}

\prose{เนวทสฺสเนน นภาวนาย ปหาตพฺพ\-เหตุกํ ธมฺมํ ปจฺจยา ทสฺสเนน ปหาตพฺพ\-เหตุโก จ เนวทสฺสเนน นภาวนาย ปหาตพฺพ\-เหตุโก จ ธมฺมา อุปฺปชฺชนฺติ สหชาต\-ปจฺจยา. วตฺถุํ ปจฺจยา ทสฺสเนน ปหาตพฺพ\-เหตุกา ขนฺธา. มหาภูเต ปจฺจยา จิตฺต\-สมุฏฺฐานํ รูปํ. วิจิกิจฺฉา\-สหคตํ โมหํ ปจฺจยา สมฺปยุตฺตกา ขนฺธา จิตฺต\-สมุฏฺฐานญฺจ รูปํ. วตฺถุํ ปจฺจยา วิจิกิจฺฉา\-สหคตา ขนฺธา จ โมโห จ. (4)}

\prose{เนวทสฺสเนน นภาวนาย ปหาตพฺพ\-เหตุกํ ธมฺมํ ปจฺจยา ภาวนาย ปหาตพฺพ\-เหตุโก จ เนวทสฺสเนน นภาวนาย ปหาตพฺพ\-เหตุโก จ ธมฺมา อุปฺปชฺชนฺติ สหชาต\-ปจฺจยา. วตฺถุํ ปจฺจยา ภาวนาย ปหาตพฺพ\-เหตุกา ขนฺธา. มหาภูเต ปจฺจยา จิตฺต\-สมุฏฺฐานํ รูปํ. อุทฺธจฺจ\-สหคตํ \csromanfolio{199}โมหํ ปจฺจยา สมฺปยุตฺตกา ขนฺธา จิตฺต\-สมุฏฺฐานญฺจ รูปํ. วตฺถุํ ปจฺจยา อุทฺธจฺจ\-สหคตา ขนฺธา จ โมโห จ. (5)}

\proseitem{45}{ทสฺสเนน ปหาตพฺพ\-เหตุ\-กญฺจ เนวทสฺสเนน นภาวนาย ปหาตพฺพ\-เหตุ\-กญฺจ ธมฺมํ ปจฺจยา ทสฺสเนน ปหาตพฺพ\-เหตุโก ธมฺโม อุปฺปชฺชติ สหชาต\-ปจฺจยา. ทสฺสเนน ปหาตพฺพ\-เหตุกํ เอกํ ขนฺธญฺจ วตฺถุญฺจ ปจฺจยา ตโย ขนฺธา ฯเปฯ. วิจิกิจฺฉา\-สหคตํ เอกํ ขนฺธญฺจ โมหญฺจ ปจฺจยา ตโย ขนฺธา ฯเปฯ. (1)}

\prose{ทสฺสเนน ปหาตพฺพ\-เหตุ\-กญฺจ เนวทสฺสเนน นภาวนาย ปหาตพฺพ\-เหตุ\-กญฺจ ธมฺมํ ปจฺจยา เนวทสฺสเนน นภาวนาย ปหาตพฺพ\-เหตุโก ธมฺโม อุปฺปชฺชติ สหชาต\-ปจฺจยา. ทสฺสเนน ปหาตพฺพ\-เหตุเก ขนฺเธ จ มหาภูเต จ ปจฺจยา จิตฺต\-สมุฏฺฐานํ รูปํ. วิจิกิจฺฉา\-สหคเต ขนฺเธ จ โมหญฺจ ปจฺจยา จิตฺต\-สมุฏฺฐานํ รูปํ. วิจิกิจฺฉา\-สหคเต ขนฺเธ จ วตฺถุญฺจ ปจฺจยา วิจิกิจฺฉา\-สหคโต โมโห. (2)}

\prose{ทสฺสเนน ปหาตพฺพ\-เหตุ\-กญฺจ เนวทสฺสเนน นภาวนาย ปหาตพฺพ\-เหตุ\-กญฺจ ธมฺมํ ปจฺจยา ทสฺสเนน ปหาตพฺพ\-เหตุโก จ เนวทสฺสเนน นภาวนาย ปหาตพฺพ\-เหตุโก จ ธมฺมา อุปฺปชฺชนฺติ สหชาต\-ปจฺจยา. ทสฺสเนน ปหาตพฺพ\-เหตุกํ เอกํ ขนฺธญฺจ วตฺถุญฺจ ปจฺจยา ตโย ขนฺธา ฯเปฯ เทฺว ขนฺเธ ฯเปฯ. ทสฺสเนน ปหาตพฺพ\-เหตุเก ขนฺเธ จ มหาภูเต จ ปจฺจยา จิตฺต\-สมุฏฺฐานํ รูปํ. วิจิกิจฺฉา\-สหคตํ เอกํ ขนฺธญฺจ โมหญฺจ ปจฺจยา ตโย ขนฺธา จิตฺต\-สมุฏฺฐานญฺจ รูปํ ฯเปฯ เทฺว ขนฺเธ ฯเปฯ. วิจิกิจฺฉา\-สหคตํ เอกํ ขนฺธญฺจ วตฺถุญฺจ ปจฺจยา ตโย ขนฺธา โมโห จ ฯเปฯ เทฺว ขนฺเธ จ วตฺถุญฺจ ปจฺจยา เทฺว ขนฺธา โมโห จ. (3)}

\prose{ภาวนาย ปหาตพฺพ\-เหตุ\-กญฺจ เนวทสฺสเนน นภาวนาย ปหาตพฺพ\-เหตุ\-กญฺจ ธมฺมํ ปจฺจยา ภาวนาย ปหาตพฺพ\-เหตุโก ธมฺโม อุปฺปชฺชติ สหชาต\-ปจฺจยา. ตีณิ.}

\csromanheader{2}{\textbf{อญฺญมญฺญาทิ}}
\proseitem{46}{ทสฺสเนน ปหาตพฺพ\-เหตุกํ ธมฺมํ ปจฺจยา ทสฺสเนน ปหาตพฺพ\-เหตุโก ธมฺโม อุปฺปชฺชติ อญฺญมญฺญ\-ปจฺจยา. นิสฺสยปจฺจยา. อุปนิสฺสย\-ปจฺจยา. ปุเรชาต\-ปจฺจยา. อาเสวนปจฺจยา. กมฺมปจฺจยา. วิปากปจฺจยา. อาหารปจฺจยา. อินฺทฺริยปจฺจยา. ฌานปจฺจยา. มคฺคปจฺจยา. สมฺปยุตฺต\-ปจฺจยา.}

\csromanheader{2}{\textbf{วิปฺปยุตฺต}}
\proseitem{47}{\csromanfolio{200}ทสฺสเนน ปหาตพฺพ\-เหตุกํ ธมฺมํ ปจฺจยา ทสฺสเนน ปหาตพฺพ\-เหตุโก ธมฺโม อุปฺปชฺชติ วิปฺปยุตฺต\-ปจฺจยา. ทสฺสเนน ปหาตพฺพ\-เหตุกํ เอกํ ขนฺธํ ปจฺจยา ตโย ขนฺธา ฯเปฯ เทฺว ขนฺธา. ขนฺธา วตฺถุํ วิปฺปยุตฺต\-ปจฺจยา. (1)}

\prose{ทสฺสเนน ปหาตพฺพ\-เหตุกํ ธมฺมํ ปจฺจยา เนวทสฺสเนน นภาวนาย ปหาตพฺพ\-เหตุโก ธมฺโม อุปฺปชฺชติ วิปฺปยุตฺต\-ปจฺจยา. ทสฺสเนน ปหาตพฺพ\-เหตุเก ขนฺเธ ปจฺจยา จิตฺต\-สมุฏฺฐานํ รูปํ. ขนฺเธ วิปฺปยุตฺต\-ปจฺจยา. วิจิกิจฺฉา\-สหคเต ขนฺเธ ปจฺจยา โมโห จิตฺต\-สมุฏฺฐานญฺจ รูปํ. โมโห วตฺถุํ วิปฺปยุตฺต\-ปจฺจยา. จิตฺต\-สมุฏฺฐานํ รูปํ ขนฺเธ วิปฺปยุตฺต\-ปจฺจยา. (2)}

\prose{ทสฺสเนน ปหาตพฺพ\-เหตุกํ ธมฺมํ ปจฺจยา ทสฺสเนน ปหาตพฺพ\-เหตุโก จ เนวทสฺสเนน นภาวนาย ปหาตพฺพ\-เหตุโก จ ธมฺมา อุปฺปชฺชนฺติ วิปฺปยุตฺต\-ปจฺจยา. ทสฺสเนน ปหาตพฺพ\-เหตุกํ เอกํ ขนฺธํ ปจฺจยา ตโย ขนฺธา จิตฺต\-สมุฏฺฐานญฺจ รูปํ ฯเปฯ เทฺว ขนฺเธ ฯเปฯ. ขนฺธา วตฺถุํ วิปฺปยุตฺต\-ปจฺจยา. จิตฺต\-สมุฏฺฐานํ รูปํ ขนฺเธ วิปฺปยุตฺต\-ปจฺจยา. วิจิกิจฺฉา\-สหคตํ เอกํ ขนฺธํ ปจฺจยา ตโย ขนฺธา โมโห จ จิตฺต\-สมุฏฺฐานญฺจ รูปํ ฯเปฯ เทฺว ขนฺเธ ฯเปฯ. ขนฺธา จ โมโห จ วตฺถุํ วิปฺปยุตฺต\-ปจฺจยา. จิตฺต\-สมุฏฺฐานํ รูปํ ขนฺเธ วิปฺปยุตฺต\-ปจฺจยา. (3)}

\prose{ภาวนาย ปหาตพฺพ\-เหตุกํ ธมฺมํ ปจฺจยา ภาวนาย ปหาตพฺพ\-เหตุโก ธมฺโม อุปฺปชฺชติ วิปฺปยุตฺต\-ปจฺจยา. ตีณิ. (ทสฺสเนน สทิสา.)}

\proseitem{48}{เนวทสฺสเนน นภาวนาย ปหาตพฺพ\-เหตุกํ ธมฺมํ เนวทสฺสเนน นภาวนาย ปหาตพฺพ\-เหตุโก ธมฺโม อุปฺปชฺชติ วิปฺปยุตฺต\-ปจฺจยา. เนวทสฺสเนน นภาวนาย ปหาตพฺพ\-เหตุกํ เอกํ ขนฺธํ ปจฺจยา ตโย ขนฺธา จิตฺต\-สมุฏฺฐานญฺจ รูปํ. เทฺว ขนฺเธ ฯเปฯ. ขนฺธา วตฺถุํ วิปฺปยุตฺต\-ปจฺจยา. จิตฺต\-สมุฏฺฐานํ รูปํ ขนฺเธ วิปฺปยุตฺต\-ปจฺจยา. วิจิกิจฺฉา\-สหคตํ อุทฺธจฺจ\-สหคตํ โมหํ ปจฺจยา จิตฺต\-สมุฏฺฐานํ รูปํ. โมหํ วิปฺปยุตฺต\-ปจฺจยา. ปฏิสนฺธิ\-กฺขเณ ฯเปฯ. ขนฺเธ ปจฺจยา วตฺถุ, วตฺถุํ ปจฺจยา ขนฺธา. ขนฺธา วตฺถุํ วิปฺปยุตฺต\-ปจฺจยา. วตฺถุ ขนฺเธ วิปฺปยุตฺต\-ปจฺจยา. เอกํ มหาภูตํ ปจฺจยา ตโย มหาภูตา ฯเปฯ มหาภูเต ปจฺจยา จิตฺต\-สมุฏฺฐานํ รูปํ, กฏตฺตารูปํ, \csromanfolio{201}อุปาทารูปํ. ขนฺเธ วิปฺปยุตฺต\-ปจฺจยา. จกฺขายตนํ ปจฺจยา จกฺขุวิญฺญาณํ ฯเปฯ. กายายตนํ ปจฺจยา กายวิญฺญาณํ. วตฺถุํ ปจฺจยา เนวทสฺสเนน นภาวนาย ปหาตพฺพ\-เหตุกา ขนฺธา. (1)}

\prose{เนวทสฺสเนน นภาวนาย ปหาตพฺพ\-เหตุกํ ธมฺมํ ปจฺจยา ทสฺสเนน ปหาตพฺพ\-เหตุโก ธมฺโม อุปฺปชฺชติ วิปฺปยุตฺต\-ปจฺจยา. วตฺถุํ ปจฺจยา ทสฺสเนน ปหาตพฺพ\-เหตุกา ขนฺธา. วตฺถุํ วิปฺปยุตฺต\-ปจฺจยา. วิจิกิจฺฉา\-สหคตํ โมหํ ปจฺจยา สมฺปยุตฺตกา ขนฺธา. วตฺถุํ วิปฺปยุตฺต\-ปจฺจยา. (2)}

\prose{เนวทสฺสเนน นภาวนาย ปหาตพฺพ\-เหตุกํ ธมฺมํ ปจฺจยา ภาวนาย ปหาตพฺพ\-เหตุโก ธมฺโม อุปฺปชฺชติ วิปฺปยุตฺต\-ปจฺจยา. วตฺถุํ ปจฺจยา ภาวนาย ปหาตพฺพ\-เหตุกา ขนฺธา. วตฺถุํ วิปฺปยุตฺต\-ปจฺจยา. อุทฺธจฺจ\-สหคตํ โมหํ ปจฺจยา สมฺปยุตฺตกา ขนฺธา. วตฺถุํ วิปฺปยุตฺต\-ปจฺจยา. (3)}

\prose{เนวทสฺสเนน นภาวนาย ปหาตพฺพ\-เหตุกํ ธมฺมํ ปจฺจยา ทสฺสเนน ปหาตพฺพ\-เหตุโก จ เนวทสฺสเนน นภาวนาย ปหาตพฺพ\-เหตุโก จ ธมฺมา อุปฺปชฺชนฺติ วิปฺปยุตฺต\-ปจฺจยา. วตฺถุํ ปจฺจยา ทสฺสเนน ปหาตพฺพ\-เหตุกา ขนฺธา. มหาภูเต ปจฺจยา จิตฺต\-สมุฏฺฐานํ รูปํ. ขนฺธา วตฺถุํ วิปฺปยุตฺต\-ปจฺจยา. จิตฺต\-สมุฏฺฐานํ รูปํ. ขนฺเธ วิปฺปยุตฺต\-ปจฺจยา. วิจิกิจฺฉา\-สหคตํ โมหํ ปจฺจยา สมฺปยุตฺตกา ขนฺธา จิตฺต\-สมุฏฺฐานญฺจ รูปํ. ขนฺธา วตฺถุํ วิปฺปยุตฺต\-ปจฺจยา. จิตฺต\-สมุฏฺฐานํ รูปํ โมหํ วิปฺปยุตฺต\-ปจฺจยา. วตฺถุํ ปจฺจยา วิจิกิจฺฉา\-สหคตา ขนฺธา จ โมโห จ. วตฺถุํ วิปฺปยุตฺต\-ปจฺจยา. (4)}

\prose{เนวทสฺสเนน นภาวนาย ปหาตพฺพ\-เหตุกํ ธมฺมํ ปจฺจยา ภาวนาย ปหาตพฺพ\-เหตุโก จ เนวทสฺสเนน นภาวนาย ปหาตพฺพ\-เหตุโก จ ธมฺมา อุปฺปชฺชนฺติ วิปฺปยุตฺต\-ปจฺจยา. วตฺถุํ ปจฺจยา ภาวนาย ปหาตพฺพ\-เหตุกา ขนฺธา ฯเปฯ. (ทสฺสเนน สทิสํ.) (5)}

\proseitem{49}{ทสฺสเนน ปหาตพฺพ\-เหตุ\-กญฺจ เนวทสฺสเนน นภาวนาย ปหาตพฺพ\-เหตุ\-กญฺจ ธมฺมํ ปจฺจยา ทสฺสเนน ปหาตพฺพ\-เหตุโก ธมฺโม อุปฺปชฺชติ วิปฺปยุตฺต\-ปจฺจยา. ทสฺสเนน ปหาตพฺพ\-เหตุกํ เอกํ ขนฺธญฺจ วตฺถุญฺจ ปจฺจยา ตโย ขนฺธา ฯเปฯ เทฺว ขนฺเธ ฯเปฯ. วตฺถุํ วิปฺปยุตฺต\-ปจฺจยา วิจิกิจฺฉา\-สหคตํ เอกํ ขนฺธญฺจ โมหญฺจ ปจฺจยา ตโย ขนฺธา ฯเปฯ เทฺว ขนฺเธ ฯเปฯ. วตฺถุํ วิปฺปยุตฺต\-ปจฺจยา. (1)}

\prose{\csromanfolio{202}ทสฺสเนน ปหาตพฺพ\-เหตุ\-กญฺจ เนวทสฺสเนน นภาวนาย ปหาตพฺพ\-เหตุ\-กญฺจ ธมฺมํ ปจฺจยา. เนวทสฺสเนน นภาวนาย ปหาตพฺพ\-เหตุโก ธมฺโม อุปฺปชฺชติ วิปฺปยุตฺต\-ปจฺจยา. ทสฺสเนน ปหาตพฺพ\-เหตุเก ขนฺเธ จ มหาภูเต จ ปจฺจยา จิตฺต\-สมุฏฺฐานํ รูปํ. ขนฺเธ วิปฺปยุตฺต\-ปจฺจยา. วิจิกิจฺฉา\-สหคเต ขนฺเธ จ โมหญฺจ ปจฺจยา จิตฺต\-สมุฏฺฐานํ รูปํ. ขนฺเธ จ โมหญฺจ วิปฺปยุตฺต\-ปจฺจยา. วิจิกิจฺฉา\-สหคเต ขนฺเธ จ วตฺถุญฺจ ปจฺจยา วิจิกิจฺฉา\-สหคโต โมโห. วตฺถุํ วิปฺปยุตฺต\-ปจฺจยา. (2)}

\prose{ทสฺสเนน ปหาตพฺพ\-เหตุ\-กญฺจ เนวทสฺสเนน นภาวนาย ปหาตพฺพ\-เหตุ\-กญฺจ ธมฺมํ ปจฺจยา ทสฺสเนน ปหาตพฺพ\-เหตุโก จ เนวทสฺสเนน นภาวนาย ปหาตพฺพ\-เหตุโก จ ธมฺมา อุปฺปชฺชนฺติ วิปฺปยุตฺต\-ปจฺจยา. ทสฺสเนน ปหาตพฺพ\-เหตุกํ เอกํ ขนฺธญฺจ วตฺถุญฺจ ปจฺจยา ตโย ขนฺธา ฯเปฯ เทฺว ขนฺเธ ฯเปฯ. ทสฺสเนน ปหาตพฺพ\-เหตุเก ขนฺเธ จ มหาภูเต จ ปจฺจยา จิตฺต\-สมุฏฺฐานํ รูปํ. ขนฺธา วตฺถุํ วิปฺปยุตฺต\-ปจฺจยา. จิตฺต\-สมุฏฺฐานํ รูปํ ขนฺเธ วิปฺปยุตฺต\-ปจฺจยา. วิจิกิจฺฉา\-สหคตํ เอกํ ขนฺธญฺจ โมหญฺจ ปจฺจยา ตโย ขนฺธา จิตฺต\-สมุฏฺฐานญฺจ รูปํ ฯเปฯ เทฺว ขนฺเธ จ ฯเปฯ. ขนฺธา วตฺถุํ วิปฺปยุตฺต\-ปจฺจยา. จิตฺต\-สมุฏฺฐานํ รูปํ ขนฺเธ จ โมหญฺจ วิปฺปยุตฺต\-ปจฺจยา. วิจิกิจฺฉา\-สหคตํ เอกํ ขนฺธญฺจ วตฺถุญฺจ ปจฺจยา ตโย ขนฺธา โมโห จ ฯเปฯ เทฺว ขนฺเธ จ ฯเปฯ. วตฺถุํ วิปฺปยุตฺต\-ปจฺจยา. (3)}

\prose{ภาวนาย ปหาตพฺพ\-เหตุ\-กญฺจ ฯเปฯ. ตีณิ. (ทสฺสเนน สทิสา.)}

\prose{อตฺถิอาทิ}

\proseitem{50}{ทสฺสเนน ปหาตพฺพ\-เหตุกํ ธมฺมํ ปจฺจยา ทสฺสเนน ปหาตพฺพ\-เหตุโก ธมฺโม อุปฺปชฺชติ อตฺถิปจฺจยา. นตฺถิปจฺจยา. วิคตปจฺจยา. อวิคตปจฺจยา.}

\csromanmarkh{1. ปจฺจยานุโลม}\csromanmarkhii{2. สงฺขฺยาวาร}\csromanheadernomark{2}{1. \textbf{ปจฺจยานุโลม 2. สงฺขฺยาวาร}}
\proseitemwithcloserruled{51}{เหตุยา สตฺตรส, อารมฺมเณ สตฺตรส, อธิปติยา สตฺตรส, อนนฺตเร สตฺตรส, สมนนฺตเร สตฺตรส, สหชาเต สตฺตรส, อญฺญมญฺเญ สตฺตรส, นิสฺสเย สตฺตรส, อุปนิสฺสเย สตฺตรส, \csromanfolio{203}ปุเรชาเต สตฺตรส, อาเสวเน สตฺตรส, กมฺเม สตฺตรส, วิปาเก เอกํ, อาหาเร สตฺตรส, อินฺทฺริเย สตฺตรส, ฌาเน สตฺตรส, มคฺเค สตฺตรส, สมฺปยุตฺเต สตฺตรส, วิปฺปยุตฺเต สตฺตรส, อตฺถิยา สตฺตรส, นตฺถิยา สตฺตรส, วิคเต สตฺตรส, อวิคเต สตฺตรส. (เอวํ คเณตพฺพํ.)}{อนุโลมํ.}
\csromanmarkh{2. ปจฺจย\-ปจฺจนีย}\csromanmarkhii{1. วิภงฺควาร}\csromanheadernomark{2}{2. ปจฺจย\-ปจฺจนีย 1. วิภงฺควาร}
\csromanheader{2}{\textbf{นเหตุ}}
\proseitem{52}{ทสฺสเนน ปหาตพฺพ\-เหตุกํ ธมฺมํ ปจฺจยา เนวทสฺสเนน นภาวนาย ปหาตพฺพ\-เหตุโก ธมฺโม อุปฺปชฺชติ นเหตุปจฺจยา. วิจิกิจฺฉา\-สหคเต ขนฺเธ ปจฺจยา วิจิกิจฺฉา\-สหคโต โมโห. (1)}

\prose{ภาวนาย ปหาตพฺพ\-เหตุกํ ธมฺมํ ปจฺจยา เนวทสฺสเนน นภาวนาย ปหาตพฺพ\-เหตุโก ธมฺโม อุปฺปชฺชติ นเหตุปจฺจยา. อุทฺธจฺจ\-สหคเต ขนฺเธ ปจฺจยา อุทฺธจฺจ\-สหคโต โมโห. (1)}

\prose{เนวทสฺสเนน นภาวนาย ปหาตพฺพ\-เหตุกํ ธมฺมํ ปจฺจยา เนวทสฺสเนน นภาวนาย ปหาตพฺพ\-เหตุโก ธมฺโม อุปฺปชฺชติ นเหตุปจฺจยา. อเหตุกํ เนวทสฺสเนน นภาวนาย ปหาตพฺพ\-เหตุกํ เอกํ ขนฺธํ ปจฺจยา ตโย ขนฺธา จิตฺต\-สมุฏฺฐานญฺจ รูปํ ฯเปฯ. อเหตุก\-ปฏิสนฺธิ\-กฺขเณ (ปริปุณฺณํ.) จกฺขายตนํ ปจฺจยา จกฺขุวิญฺญาณํ ฯเปฯ. กายายตนํ ปจฺจยา กายวิญฺญาณํ. วตฺถุํ ปจฺจยา อเหตุกา เนวทสฺสเนน นภาวนาย ปหาตพฺพ\-เหตุกา ขนฺธา. วตฺถุํ ปจฺจยา วิจิกิจฺฉา\-สหคโต อุทฺธจฺจ\-สหคโต โมโห. (1)}

\proseitem{53}{ทสฺสเนน ปหาตพฺพ\-เหตุ\-กญฺจ เนวทสฺสเนน นภาวนาย ปหาตพฺพ\-เหตุ\-กญฺจ ธมฺมํ ปจฺจยา เนวทสฺสเนน นภาวนาย ปหาตพฺพ\-เหตุโก ธมฺโม อุปฺปชฺชติ นเหตุปจฺจยา. วิจิกิจฺฉา\-สหคเต ขนฺเธ จ วตฺถุญฺจ ปจฺจยา วิจิกิจฺฉา\-สหคโต โมโห. (1)}

\prose{ภาวนาย ปหาตพฺพ\-เหตุ\-กญฺจ เนวทสฺสเนน นภาวนาย ปหาตพฺพ\-เหตุ\-กญฺจ ธมฺมํ ปจฺจยา เนวทสฺสเนน นภาวนาย ปหาตพฺพ\-เหตุโก \csromanfolio{204}ธมฺโม อุปฺปชฺชติ นเหตุปจฺจยา. อุทฺธจฺจ\-สหคเต ขนฺเธ จ วตฺถุญฺจ ปจฺจยา อุทฺธจฺจ\-สหคโต โมโห. (1)}

\prose{นอารมฺมณ}

\proseitem{54}{ทสฺสเนน ปหาตพฺพ\-เหตุกํ ธมฺมํ ปจฺจยา เนวทสฺสเนน นภาวนาย ปหาตพฺพ\-เหตุโก ธมฺโม อุปฺปชฺชติ นอารมฺมณ\-ปจฺจยา. ทสฺสเนน ปหาตพฺพ\-เหตุเก ขนฺเธ ปจฺจยา จิตฺต\-สมุฏฺฐานํ รูปํ. (1)}

\prose{ภาวนาย ปหาตพฺพ\-เหตุกํ ธมฺมํ ปจฺจยา เนวทสฺสเนน นภาวนาย ปหาตพฺพ\-เหตุโก ธมฺโม อุปฺปชฺชติ นอารมฺมณ\-ปจฺจยา. ภาวนาย ปหาตพฺพ\-เหตุเก ขนฺเธ ปจฺจยา จิตฺต\-สมุฏฺฐานํ รูปํ. (1)}

\prose{เนวทสฺสเนน นภาวนาย ปหาตพฺพ\-เหตุกํ ธมฺมํ ปจฺจยา เนวทสฺสเนน นภาวนาย ปหาตพฺพ\-เหตุโก ธมฺโม อุปฺปชฺชติ นอารมฺมณ\-ปจฺจยา. เนวทสฺสเนน นภาวนาย ปหาตพฺพ\-เหตุเก ขนฺเธ ปจฺจยา จิตฺต\-สมุฏฺฐานํ รูปํ. ปฏิสนฺธิ\-กฺขเณ เนวทสฺสเนน นภาวนาย ปหาตพฺพ\-เหตุเก ขนฺเธ ปจฺจยา กฏตฺตารูปํ. ขนฺเธ ปจฺจยา วตฺถุ ฯเปฯ. เอกํ มหาภูตํ ฯเปฯ. พาหิรํ. อาหารสมุฏฺฐานํ. อุตุสมุฏฺฐานํ. อสญฺญสตฺตานํ ฯเปฯ. (1)}

\proseitem{55}{ทสฺสเนน ปหาตพฺพ\-เหตุ\-กญฺจ เนวทสฺสเนน นภาวนาย ปหาตพฺพ\-เหตุ\-กญฺจ ธมฺมํ ปจฺจยา เนวทสฺสเนน นภาวนาย ปหาตพฺพ\-เหตุโก ธมฺโม อุปฺปชฺชติ นอารมฺมณ\-ปจฺจยา. ทสฺสเนน ปหาตพฺพ\-เหตุเก ขนฺเธ จ มหาภูเต จ ปจฺจยา จิตฺต\-สมุฏฺฐานํ รูปํ. วิจิกิจฺฉา\-สหคเต ขนฺเธ จ โมหญฺจ ปจฺจยา จิตฺต\-สมุฏฺฐานํ รูปํ. (1)}

\prose{ภาวนาย ปหาตพฺพ\-เหตุ\-กญฺจ เนวทสฺสเนน นภาวนาย ปหาตพฺพ\-เหตุ\-กญฺจ ธมฺมํ ปจฺจยา เนวทสฺสเนน นภาวนาย ปหาตพฺพ\-เหตุโก ธมฺโม อุปฺปชฺชติ นอารมฺมณ\-ปจฺจยา. ภาวนาย ปหาตพฺพ\-เหตุเก ขนฺเธ จ มหาภูเต จ ปจฺจยา จิตฺต\-สมุฏฺฐานํ รูปํ. อุทฺธจฺจ\-สหคเต ขนฺเธ จ โมหญฺจ ปจฺจยา จิตฺต\-สมุฏฺฐานํ รูปํ. (1)}

\prose{นอธิปตฺยาทิ}

\proseitem{56}{ทสฺสเนน ปหาตพฺพ\-เหตุกํ ธมฺมํ ปจฺจยา ทสฺสเนน ปหาตพฺพ\-เหตุโก ธมฺโม อุปฺปชฺชติ นอธิปติ\-ปจฺจยา. (สหชาต\-สทิสํ.)}

\prose{\csromanfolio{205}นอนนฺตร\-ปจฺจยา. นสมนนฺตร\-ปจฺจยา. นอญฺญมญฺญ\-ปจฺจยา. นฺเอาปนิสฺสย\-ปจฺจยา. นปุเรชาต\-ปจฺจยา. (ปฏิจฺจวาเร ปจฺจนีย\-สทิสํ, เตรส ปญฺหา. นินฺนานํ.) นปจฺฉาชาต\-ปจฺจยา. นอาเสวน\-ปจฺจยา.}

\csromanheader{2}{\textbf{นกมฺม}}
\proseitem{57}{ทสฺสเนน ปหาตพฺพ\-เหตุกํ ธมฺมํ ปจฺจยา ทสฺสเนน ปหาตพฺพ\-เหตุโก ธมฺโม อุปฺปชฺชติ นกมฺมปจฺจยา. ทสฺสเนน ปหาตพฺพ\-เหตุเก ขนฺเธ ปจฺจยา ทสฺสเนน ปหาตพฺพ\-เหตุกา เจตนา. (1)}

\prose{ภาวนาย ปหาตพฺพ\-เหตุกํ ธมฺมํ ปจฺจยา ภาวนาย ปหาตพฺพ\-เหตุโก ธมฺโม อุปฺปชฺชติ นกมฺมปจฺจยา. ภาวนาย ปหาตพฺพ\-เหตุเก ขนฺเธ ปจฺจยา ภาวนาย ปหาตพฺพ\-เหตุกา เจตนา. (1)}

\prose{เนวทสฺสเนน นภาวนาย ปหาตพฺพ\-เหตุกํ ธมฺมํ ปจฺจยา เนวทสฺสเนน นภาวนาย ปหาตพฺพ\-เหตุโก ธมฺโม อุปฺปชฺชติ น กมฺมปจฺจยา. เนวทสฺสเนน นภาวนาย ปหาตพฺพ\-เหตุเก ขนฺเธ ปจฺจยา เนวทสฺสเนน นภาวนาย ปหาตพฺพ\-เหตุกา เจตนา. พาหิรํ. อาหารสมุฏฺฐานํ. อุตุสมุฏฺฐานํ ฯเปฯ. วตฺถุํ ปจฺจยา เนวทสฺสเนน นภาวนาย ปหาตพฺพ\-เหตุกา เจตนา. (1)}

\prose{เนวทสฺสเนน นภาวนาย ปหาตพฺพ\-เหตุกํ ธมฺมํ ปจฺจยา ทสฺสเนน ปหาตพฺพ\-เหตุโก ธมฺโม อุปฺปชฺชติ นกมฺมปจฺจยา. วตฺถุํ ปจฺจยา ทสฺสเนน ปหาตพฺพ\-เหตุกา เจตนา. วิจิกิจฺฉา\-สหคตํ โมหํ ปจฺจยา สมฺปยุตฺตกา เจตนา. (2)}

\prose{เนวทสฺสเนน นภาวนาย ปหาตพฺพ\-เหตุกํ ธมฺมํ ปจฺจยา ภาวนาย ปหาตพฺพ\-เหตุโก ธมฺโม อุปฺปชฺชติ นกมฺมปจฺจยา. วตฺถุํ ปจฺจยา ภาวนาย ปหาตพฺพ\-เหตุกา เจตนา. อุทฺธจฺจ\-สหคตํ โมหํ ปจฺจยา สมฺปยุตฺตกา เจตนา. (3)}

\proseitem{58}{ทสฺสเนน ปหาตพฺพ\-เหตุ\-กญฺจ เนวทสฺสเนน นภาวนาย ปหาตพฺพ\-เหตุ\-กญฺจ ธมฺมํ ปจฺจยา ทสฺสเนน ปหาตพฺพ\-เหตุโก ธมฺโม อุปฺปชฺชติ นกมฺมปจฺจยา. ทสฺสเนน ปหาตพฺพ\-เหตุเก ขนฺเธ จ วตฺถุญฺจ ปจฺจยา ทสฺสเนน ปหาตพฺพ\-เหตุกา เจตนา. วิจิกิจฺฉา\-สหคเต ขนฺเธ จ โมหญฺจ ปจฺจยา สมฺปยุตฺตกา เจตนา. (1)}

\prose{\csromanfolio{206}ภาวนาย ปหาตพฺพ\-เหตุ\-กญฺจ เนวทสฺสเนน นภาวนาย ปหาตพฺพ\-เหตุ\-กญฺจ ธมฺมํ ปจฺจยา ภาวนาย ปหาตพฺพ\-เหตุโก ธมฺโม อุปฺปชฺชติ นกมฺมปจฺจยา. ภาวนาย ปหาตพฺพ\-เหตุเก ขนฺเธ จ วตฺถุญฺจ ปจฺจยา ภาวนาย ปหาตพฺพ\-เหตุกา เจตนา. อุทฺธจฺจ\-สหคเต ขนฺเธ จ โมหญฺจ ปจฺจยา สมฺปยุตฺตกา เจตนา. (1)}

\csromanheader{2}{\textbf{นวิปากาทิ}}
\proseitem{59}{ทสฺสเนน ปหาตพฺพ\-เหตุกํ ธมฺมํ ปจฺจยา ทสฺสเนน ปหาตพฺพ\-เหตุโก ธมฺโม อุปฺปชฺชติ นวิปาก\-ปจฺจยา. (ปริปุณฺณํ, ปฏิสนฺธิ นตฺถิ.) นอาหารปจฺจยา. พาหิรํ. อุตุสมุฏฺฐานํ. อสญฺญสตฺตานํ ฯเปฯ นอินฺทฺริยปจฺจยา. พาหิรํ. อาหารสมุฏฺฐานํ. อุตุสมุฏฺฐานํ. อสญฺญสตฺตานํ ฯเปฯ มหาภูเต ปจฺจยา รูปชีวิตินฺทฺริยํ. นฌานปจฺจยา. ปญฺจ\-วิญฺญาณสหคตํ เอกํ ขนฺธํ ฯเปฯ. พาหิรํ. อาหารสมุฏฺฐานํ. อุตุสมุฏฺฐานํ. อสญฺญสตฺตานํ ฯเปฯ นมคฺคปจฺจยา. อเหตุกํ เอกํ ฯเปฯ นสมฺปยุตฺต\-ปจฺจยา. นวิปฺปยุตฺต\-ปจฺจยา. (ปฏิจฺจ\-วาร\-ปจฺจนีเย นวิปฺปยุตฺต\-สทิสํ, นินฺนานํ. เอกาทส.) โนนตฺถิปจฺจยา. โนวิคต\-ปจฺจยา.}

\csromanmarkh{2. ปจฺจย\-ปจฺจนีย}\csromanmarkhii{2. สงฺขฺยาวาร}\csromanheadernomark{2}{2. ปจฺจย\-ปจฺจนีย 2. สงฺขฺยาวาร}
\proseitemwithcloserruled{60}{นเหตุยา ปญฺจ, นอารมฺมเณ ปญฺจ, นอธิปติยา สตฺตรส, นอนนฺตเร ปญฺจ, นสมนนฺตเร ปญฺจ, นอญฺญมญฺเญ ปญฺจ, นอุปนิสฺสเย ปญฺจ, นปุเรชาเต เตรส, นปจฺฉาชาเต สตฺตรส, นอาเสวเน สตฺตรส, นกมฺเม สตฺต, นวิปาเก สตฺตรส, นอาหาเร เอกํ, นอินฺทฺริเย เอกํ, นฌาเน เอกํ, นมคฺเค เอกํ, นสมฺปยุตฺเต ปญฺจ, นวิปฺปยุตฺเต เอกาทส, โนนตฺถิยา ปญฺจ, โนวิคเต ปญฺจ.  (เอวํ คเณตพฺพํ.)}{ปจฺจนียํ.}
\proseitem{61}{เหตุปจฺจยา นอารมฺมเณ ปญฺจ, นอธิปติยา สตฺตรส, นอนนฺตเร ปญฺจ, นสมนนฺตเร ปญฺจ, นอญฺญมญฺเญ ปญฺจ, นอุปนิสฺสเย ปญฺจ, \csromanfolio{207}นปุเรชาเต เตรส, นปจฺฉาชาเต สตฺตรส, นอาเสวเน สตฺตรส, นกมฺเม สตฺต, นวิปาเก สตฺตรส, นสมฺปยุตฺเต ปญฺจ, นวิปฺปยุตฺเต เอกาทส, โนนตฺถิยา ปญฺจ, โนวิคเต ปญฺจ.  (เอวํ คเณตพฺพํ.) อนุโลม\-ปจฺจนียํ.}

\proseitem{62}{นเหตุปจฺจยา อารมฺมเณ ปญฺจ, อนนฺตเร ปญฺจ, สมนนฺตเร ปญฺจ, สหชาเต ปญฺจ, อญฺญมญฺเญ ปญฺจ, นิสฺสเย ปญฺจ, อุปนิสฺสเย ปญฺจ, ปุเรชาเต ปญฺจ, อาเสวเน ปญฺจ, กมฺเม ปญฺจ, วิปาเก เอกํ, อาหาเร ปญฺจ, อินฺทฺริเย ปญฺจ, ฌาเน ปญฺจ, มคฺเค ปญฺจ, สมฺปยุตฺเต ปญฺจ, วิปฺปยุตฺเต ปญฺจ, อตฺถิยา ปญฺจ, นตฺถิยา ปญฺจ, วิคเต ปญฺจ, อวิคเต ปญฺจ.  (เอวํ คเณตพฺพํ.) ปจฺจนียา\-นุโลมํ. ปจฺจยวาโร. (นิสฺสยวาโร ปจฺจยวาร\-สทิโส.)}
\csromansectionrule

\csromanmarkh{9. ทสฺสเนน\-ปหาตพฺพ\-เหตุ\-กตฺติก}\csromanmarkhii{5. สํสฏฺฐวาร}\csromanheadernomark{2}{9. ทสฺสเนน\-ปหาตพฺพ\-เหตุ\-กตฺติก 5. สํสฏฺฐวาร}
\csromanmarkh{1. ปจฺจยานุโลม}\csromanmarkhii{1. วิภงฺควาร}\csromanheadernomark{2}{1. \textbf{ปจฺจยานุโลม 1. วิภงฺควาร}}
\csromanheader{2}{\textbf{เหตุ}}
\proseitem{63}{ทสฺสเนน ปหาตพฺพ\-เหตุกํ ธมฺมํ สํสฏฺโฐ ทสฺสเนน ปหาตพฺพ\-เหตุโก ธมฺโม อุปฺปชฺชติ เหตุปจฺจยา. ทสฺสเนน ปหาตพฺพ\-เหตุกํ เอกํ ขนฺธํ สํสฏฺฐา ตโย ขนฺธา ฯเปฯ เทฺว ขนฺเธ สํสฏฺฐา เทฺว ขนฺธา. (1)}

\prose{ภาวนาย ปหาตพฺพ\-เหตุกํ ธมฺมํ สํสฏฺโฐ ภาวนาย ปหาตพฺพ\-เหตุโก ธมฺโม อุปฺปชฺชติ เหตุปจฺจยา. ภาวนาย ปหาตพฺพ\-เหตุกํ เอกํ ขนฺธํ สํสฏฺฐา ตโย ขนฺธา ฯเปฯ เทฺว ขนฺเธ สํสฏฺฐา เทฺว ขนฺธา. (1)}

\prose{\csromanfolio{208}เนวทสฺสเนน นภาวนาย ปหาตพฺพ\-เหตุกํ ธมฺมํ สํสฏฺโฐ เนวทสฺสเนน นภาวนาย ปหาตพฺพ\-เหตุโก ธมฺโม อุปฺปชฺชติ เหตุปจฺจยา. เนวทสฺสเนน นภาวนาย ปหาตพฺพ\-เหตุกํ เอกํ ขนฺธํ สํสฏฺฐา ตโย ขนฺธา ฯเปฯ. ปฏิสนฺธิ\-กฺขเณ ฯเปฯ. (1)}

\prose{เนวทสฺสเนน นภาวนาย ปหาตพฺพ\-เหตุกํ ธมฺมํ สํสฏฺโฐ ทสฺสเนน ปหาตพฺพ\-เหตุโก ธมฺโม อุปฺปชฺชติ เหตุปจฺจยา. วิจิกิจฺฉา\-สหคตํ โมหํ สํสฏฺฐา สมฺปยุตฺตกา ขนฺธา. (2)}

\prose{เนวทสฺสเนน นภาวนาย ปหาตพฺพ\-เหตุกํ ธมฺมํ สํสฏฺโฐ ภาวนาย ปหาตพฺพ\-เหตุโก ธมฺโม อุปฺปชฺชติ เหตุปจฺจยา. อุทฺธจฺจ\-สหคตํ โมหํ สํสฏฺฐา สมฺปยุตฺตกา ขนฺธา. (3)}

\proseitem{64}{ทสฺสเนน ปหาตพฺพ\-เหตุ\-กญฺจ เนวทสฺสเนน นภาวนาย ปหาตพฺพ\-เหตุ\-กญฺจ ธมฺมํ สํสฏฺโฐ ทสฺสเนน ปหาตพฺพ\-เหตุโก ธมฺโม อุปฺปชฺชติ เหตุปจฺจยา. วิจิกิจฺฉา\-สหคตํ เอกํ ขนฺธญฺจ โมหญฺจ สํสฏฺฐา ตโย ขนฺธา ฯเปฯ เทฺว ขนฺธา. (1)}

\prose{ภาวนาย ปหาตพฺพ\-เหตุ\-กญฺจ เนวทสฺสเนน นภาวนาย ปหาตพฺพ\-เหตุ\-กญฺจ ธมฺมํ ภาวนาย ปหาตพฺพ\-เหตุโก ธมฺโม อุปฺปชฺชติ เหตุปจฺจยา. อุทฺธจฺจ\-สหคตํ เอกํ ขนฺธญฺจ โมหญฺจ สํสฏฺฐา ตโย ขนฺธา ฯเปฯ เทฺว ขนฺธา. (1)}

\csromanheader{2}{\textbf{อารมฺมณ}}
\proseitem{65}{ทสฺสเนน ปหาตพฺพ\-เหตุกํ ธมฺมํ สํสฏฺโฐ ทสฺสเนน ปหาตพฺพ\-เหตุโก ธมฺโม อุปฺปชฺชติ อารมฺมณ\-ปจฺจยา. ทสฺสเนน ปหาตพฺพ\-เหตุกํ เอกํ ขนฺธํ สํสฏฺฐา ตโย ขนฺธา ฯเปฯ เทฺว ขนฺธา. (1)}

\prose{ทสฺสเนน ปหาตพฺพ\-เหตุกํ ธมฺมํ สํสฏฺโฐ เนวทสฺสเนน นภาวนาย ปหาตพฺพ\-เหตุโก ธมฺโม อุปฺปชฺชติ อารมฺมณ\-ปจฺจยา. วิจิกิจฺฉา\-สหคเต ขนฺเธ สํสฏฺโฐ วิจิกิจฺฉา\-สหคโต โมโห. (2)}

\prose{ทสฺสเนน ปหาตพฺพ\-เหตุกํ ธมฺมํ สํสฏฺโฐ ทสฺสเนน ปหาตพฺพ\-เหตุโก จ เนวทสฺสเนน นภาวนย ปหาตพฺพ\-เหตุโก จ ธมฺมา อุปฺปชฺชนฺติ อารมฺมณ\-ปจฺจยา. วิจิกิจฺฉา\-สหคตํ เอกํ ขนฺธํ สํสฏฺฐา ตโย ขนฺธา โมโห จ ฯเปฯ เทฺว ขนฺเธ ฯเปฯ. (3)}

\prose{\csromanfolio{209}ภาวนาย ปหาตพฺพ\-เหตุกํ ธมฺมํ สํสฏฺโฐ. ตีณิ.}

\proseitem{66}{เนวทสฺสเนน นภาวนาย ปหาตพฺพ\-เหตุกํ ธมฺมํ สํสฏฺโฐ เนวทสฺสเนน นภาวนาย ปหาตพฺพ\-เหตุโก ธมฺโม อุปฺปชฺชติ อารมฺมณ\-ปจฺจยา. เนวทสฺสเนน นภาวนาย ปหาตพฺพ\-เหตุกํ เอกํ ขนฺธํ สํสฏฺฐา ตโย ขนฺธา ฯเปฯ เทฺว ขนฺเธ สํสฏฺฐา เทฺว ขนฺธา. ปฏิสนฺธิ\-กฺขเณ ฯเปฯ. (1)}

\prose{เนวทสฺสเนน นภาวนาย ปหาตพฺพ\-เหตุกํ ธมฺมํ สํสฏฺโฐ ทสฺสเนน ปหาตพฺพ\-เหตุโก ธมฺโม อุปฺปชฺชติ อารมฺมณ\-ปจฺจยา. วิจิกิจฺฉา\-สหคตํ โมหํ สํสฏฺฐา สมฺปยุตฺตกา ขนฺธา. (2)}

\prose{เนวทสฺสเนน นภาวนาย ปหาตพฺพ\-เหตุกํ ธมฺมํ สํสฏฺโฐ ภาวนาย ปหาตพฺพ\-เหตุโก ธมฺโม อุปฺปชฺชติ อารมฺมณ\-ปจฺจยา. อุทฺธจฺจ\-สหคตํ โมหํ สํสฏฺฐา สมฺปยุตฺตกา ขนฺธา. (3)}

\prose{ทสฺสเนน ปหาตพฺพ\-เหตุ\-กญฺจ เนวทสฺสเนน นภาวนาย ปหาตพฺพ\-เหตุ\-กญฺจ ธมฺมํ สํสฏฺโฐ ทสฺสเนน ปหาตพฺพ\-เหตุโก ธมฺโม อุปฺปชฺชติ, อารมฺมณ\-ปจฺจยา. วิจิกิจฺฉา\-สหคตํ เอกํ ขนฺธญฺจ โมหญฺจ สํสฏฺฐา ตโย ขนฺธา ฯเปฯ เทฺว ขนฺเธ จ โมหญฺจ สํสฏฺฐา เทฺว ขนฺธา. (1)}

\prose{ภาวนาย ปหาตพฺพ\-เหตุ\-กญฺจ เนวทสฺสเนน นภาวนาย ปหาตพฺพ\-เหตุ\-กญฺจ ธมฺมํ สํสฏฺโฐ ภาวนาย ปหาตพฺพ\-เหตุโก ธมฺโม อุปฺปชฺชติ อารมฺมณ\-ปจฺจยา. อุทฺธจฺจ\-สหคตํ เอกํ ขนฺธญฺจ โมหญฺจ สํสฏฺฐา ตโย ขนฺธา ฯเปฯ เทฺว ขนฺธา. (1)}

\csromanheader{2}{\textbf{อธิปตฺยาทิ}}
\proseitem{67}{ทสฺสเนน ปหาตพฺพ\-เหตุกํ ธมฺมํ สํสฏฺโฐ ทสฺสเนน ปหาตพฺพ\-เหตุโก ธมฺโม อุปฺปชฺชติ อธิปติ\-ปจฺจยา. ทสฺสเนน ปหาตพฺพ\-เหตุกํ เอกํ ขนฺธํ สํสฏฺฐา ตโย ขนฺธา ฯเปฯ เทฺว ขนฺธา. (1)}

\prose{ภาวนาย ปหาตพฺพ\-เหตุกํ ธมฺมํ สํสฏฺโฐ. เอกํ.}

\prose{เนวทสฺสเนน นภาวนาย ปหาตพฺพ\-เหตุกํ ธมฺมํ สํสฏฺโฐ เนวทสฺสเนน นภาวนาย ปหาตพฺพ\-เหตุโก ธมฺโม อุปฺปชฺชติ อธิปติ\-ปจฺจยา. เนวทสฺสเนน นภาวนาย ปหาตพฺพ\-เหตุกํ เอกํ ขนฺธํ สํสฏฺฐา ตโย ขนฺธา ฯเปฯ เทฺว ขนฺธา. อนนฺตร\-ปจฺจยา. สมนนฺตร\-ปจฺจยา.}

\csromanheader{2}{\textbf{สหชาตาทิ}}
\proseitem{68}{\csromanfolio{210}ทสฺสเนน ปหาตพฺพ\-เหตุกํ ธมฺมํ สํสฏฺโฐ ทสฺสเนน ปหาตพฺพ\-เหตุโก ธมฺโม อุปฺปชฺชติ สหชาต\-ปจฺจยา. อญฺญมญฺญ\-ปจฺจยา. นิสฺสยปจฺจยา. อุปนิสฺสย\-ปจฺจยา. ปุเรชาต\-ปจฺจยา. อาเสวนปจฺจยา. กมฺมปจฺจยา. วิปากปจฺจยา. อาหารปจฺจยา. อินฺทฺริยปจฺจยา. ฌานปจฺจยา. มคฺคปจฺจยา. สมฺปยุตฺต\-ปจฺจยา. วิปฺปยุตฺต\-ปจฺจยา. อตฺถิปจฺจยา. นตฺถิปจฺจยา. วิคตปจฺจยา. อวิคตปจฺจยา.}

\csromanmarkh{1. ปจฺจยานุโลม}\csromanmarkhii{2. สงฺขฺยาวาร}\csromanheadernomark{2}{1. \textbf{ปจฺจยานุโลม 2. สงฺขฺยาวาร}}
\proseitemwithcloserruled{69}{เหตุยา สตฺต, อารมฺมเณ เอกาทส, อธิปติยา ตีณิ, อนนฺตเร เอกาทส, สมนนฺตเร เอกาทส, สหชาเต เอกาทส, อญฺญมญฺเญ เอกาทส, นิสฺสเย เอกาทส, อุปนิสฺสเย เอกาทส, ปุเรชาเต เอกาทส, อาเสวเน เอกาทส, กมฺเม เอกาทส, วิปาเก เอกํ, อาหาเร เอกาทส, อินฺทฺริเย เอกาทส, ฌาเน เอกาทส, มคฺเค เอกาทส, สมฺปยุตฺเต เอกาทส, วิปฺปยุตฺเต เอกาทส, อตฺถิยา เอกาทส, นตฺถิยา เอกาทส, วิคเต เอกาทส, อวิคเต เอกาทส. (เอวํ คเณตพฺพํ.)}{อนุโลมํ.}
\csromanmarkh{2. ปจฺจย\-ปจฺจนีย}\csromanmarkhii{1. วิภงฺควาร}\csromanheadernomark{2}{2. ปจฺจย\-ปจฺจนีย 1. วิภงฺควาร}
\csromanheader{2}{\textbf{นเหตุ}}
\proseitem{70}{ทสฺสเนน ปหาตพฺพ\-เหตุกํ ธมฺมํ สํสฏฺโฐ เนวทสฺสเนน นภาวนาย ปหาตพฺพ\-เหตุโก ธมฺโม อุปฺปชฺชติ นเหตุปจฺจยา. วิจิกิจฺฉา\-สหคเต ขนฺเธ สํสฏฺโฐ วิจิกิจฺฉา\-สหคโต โมโห. (1)}

\prose{ภาวนาย ปหาตพฺพ\-เหตุกํ ธมฺมํ สํสฏฺโฐ เนวทสฺสเนน นภาวนาย ปหาตพฺพ\-เหตุโก ธมฺโม อุปฺปชฺชติ นเหตุปจฺจยา. อุทฺธจฺจ\-สหคเต ขนฺเธ สํสฏฺโฐ อุทฺธจฺจ\-สหคโต โมโห. (1)}

\prose{\csromanfolio{211}เนวทสฺสเนน นภาวนาย ปหาตพฺพ\-เหตุกํ ธมฺมํ สํสฏฺโฐ เนวทสฺสเนน นภาวนาย ปหาตพฺพ\-เหตุโก ธมฺโม อุปฺปชฺชติ นเหตุปจฺจยา. อเหตุกํ เนวทสฺสเนน นภาวนาย ปหาตพฺพ\-เหตุกํ เอกํ ขนฺธํ สํสฏฺฐา ตโย ขนฺธา ฯเปฯ เทฺว ขนฺเธ สํสฏฺฐา เทฺว ขนฺธา. อเหตุก\-ปฏิสนฺธิ\-กฺขเณ ฯเปฯ. (1)}

\prose{นอธิปตฺยาทิ}

\proseitem{71}{ทสฺสเนน ปหาตพฺพ\-เหตุกํ ธมฺมํ สํสฏฺโฐ ฯเปฯ นอธิปติ\-ปจฺจยา. (สหชาต\-สทิสํ.) นปุเรชาต\-ปจฺจยา. นปจฺฉาชาต\-ปจฺจยา. นอาเสวน\-ปจฺจยา. นกมฺมปจฺจยา. สตฺต. นวิปาก\-ปจฺจยา. นฌานปจฺจยา. นมคฺคปจฺจยา. นวิปฺปยุตฺต\-ปจฺจยา.}

\csromanmarkh{2. ปจฺจย\-ปจฺจนีย}\csromanmarkhii{2. สงฺขฺยาวาร}\csromanheadernomark{2}{2. ปจฺจย\-ปจฺจนีย 2. สงฺขฺยาวาร}
\proseitemwithcloserruled{72}{นเหตุยา ตีณิ, นอธิปติยา เอกาทส, นปุเรชาเต เอกาทส, นปจฺฉาชาเต เอกาทส, นอาเสวเน เอกาทส, นกมฺเม สตฺต, นวิปาเก เอกาทส, นฌาเน เอกํ, นมคฺเค เอกํ, นวิปฺปยุตฺเต เอกาทส. (เอวํ คเณตพฺพํ.)}{ปจฺจนียํ.}
\csromanheader{2}{3. \textbf{ปจฺจยานุโลม ปจฺจนีย}}
\proseitemwithcloserruled{73}{เหตุปจฺจยา นอธิปติยา สตฺต, นปุเรชาเต สตฺต, นปจฺฉาชาเต สตฺต, นอาเสวเน สตฺต, นกมฺเม สตฺต, นวิปาเก สตฺต, นวิปฺปยุตฺเต สตฺต. (เอวํ คเณตพฺพํ.)}{อนุโลม\-ปจฺจนียํ.}
\proseitem{74}{นเหตุปจฺจยา อารมฺมเณ ตีณิ, อนนฺตเร ตีณิ, สมนนฺตเร ตีณิ, สหชาเต ตีณิ, อญฺญมญฺเญ ตีณิ, นิสฺสเย ตีณิ, อุปนิสฺสเย ตีณิ, ปุเรชาเต ตีณิ, อาเสวเน ตีณิ, กมฺเม ตีณิ, วิปาเก \csromanfolio{212}เอกํ, อาหาเร ตีณิ, อินฺทฺริเย ตีณิ, ฌาเน ตีณิ, มคฺเค เทฺว, สมฺปยุตฺเต ตีณิ, วิปฺปยุตฺเต ตีณิ, อตฺถิยา ตีณิ, นตฺถิยา ตีณิ, วิคเต ตีณิ, อวิคเต ตีณิ. (เอวํ คเณตพฺพํ.) ปจฺจนียา\-นุโลมํ. สํสฏฺฐวาโร. (สมฺปยุตฺตวาโร สํสฏฺฐวาร\-สทิโส.)}

\csromanmarkh{9. ทสฺสเนน\-ปหาตพฺพ\-เหตุ\-กตฺติก}\csromanmarkhii{7. ปญฺหาวาร}\csromanheadernomark{2}{9. ทสฺสเนน\-ปหาตพฺพ\-เหตุ\-กตฺติก 7. ปญฺหาวาร}
\csromanmarkh{1. ปจฺจยานุโลม}\csromanmarkhii{1. วิภงฺควาร}\csromanheadernomark{2}{1. \textbf{ปจฺจยานุโลม 1. วิภงฺควาร}}
\csromanheader{2}{\textbf{เหตุ}}
\proseitem{75}{ทสฺสเนน ปหาตพฺพ\-เหตุโก ธมฺโม ทสฺสเนน ปหาตพฺพ\-เหตุกสฺส ธมฺมสฺส เหตุปจฺจเยน ปจฺจโย. ทสฺสเนน ปหาตพฺพ\-เหตุกา เหตู สมฺปยุตฺตกานํ ขนฺธานํ เหตุปจฺจเยน ปจฺจโย. (1)}

\prose{ทสฺสเนน ปหาตพฺพ\-เหตุโก ธมฺโม เนวทสฺสเนน นภาวนาย ปหาตพฺพ\-เหตุกสฺส ธมฺมสฺส เหตุปจฺจเยน ปจฺจโย. ทสฺสเนน ปหาตพฺพ\-เหตุกา เหตู จิตฺต\-สมุฏฺฐานานํ รูปานํ เหตุปจฺจเยน ปจฺจโย. (2)}

\prose{ทสฺสเนน ปหาตพฺพ\-เหตุโก ธมฺโม ทสฺสเนน ปหาตพฺพ\-เหตุกสฺส จ เนวทสฺสเนน นภาวนาย ปหาตพฺพ\-เหตุกสฺส จ ธมฺมสฺส เหตุปจฺจเยน ปจฺจโย. ทสฺสเนน ปหาตพฺพ\-เหตุกา เหตู สมฺปยุตฺตกานํ ขนฺธานํ จิตฺตสมุฏฺฐานานญฺจ รูปานํ เหตุปจฺจเยน ปจฺจโย. (3)}

\prose{ภาวนาย ปหาตพฺพ\-เหตุโก ธมฺโม. ตีณิ.}

\proseitem{76}{เนวทสฺสเนน นภาวนาย ปหาตพฺพ\-เหตุโก ธมฺโม เนวทสฺสเนน นภาวนาย ปหาตพฺพ\-เหตุกสฺส ธมฺมสฺส เหตุปจฺจเยน ปจฺจโย ฯเปฯ.}

\prose{เนวทสฺสเนน นภาวนาย ปหาตพฺพ\-เหตุโก ธมฺโม ทสฺสเนน ปหาตพฺพ\-เหตุกสฺส ธมฺมสฺส เหตุปจฺจเยน ปจฺจโย. วิจิกิจฺฉา\-สหคโต โมโห สมฺปยุตฺตกานํ ขนฺธานํ เหตุปจฺจเยน ปจฺจโย. (2)}

\prose{\csromanfolio{213}เนวทสฺสเนน นภาวนาย ปหาตพฺพ\-เหตุโก ธมฺโม ภาวนาย ปหาตพฺพ\-เหตุกสฺส ธมฺมสฺส เหตุปจฺจเยน ปจฺจโย. อุทฺธจฺจ\-สหคโต โมโห สมฺปยุตฺตกานํ ขนฺธานํ เหตุปจฺจเยน ปจฺจโย. (3)}

\prose{เนวทสฺสเนน นภาวนาย ปหาตพฺพ\-เหตุโก ธมฺโม ทสฺสเนน ปหาตพฺพ\-เหตุกสฺส จ เนวทสฺสเนน นภาวนาย ปหาตพฺพ\-เหตุกสฺส จ ธมฺมสฺส เหตุปจฺจเยน ปจฺจโย. วิจิกิจฺฉา\-สหคโต โมโห สมฺปยุตฺตกานํ ขนฺธานํ จิตฺตสมุฏฺฐานานญฺจ รูปานํ เหตุปจฺจเยน ปจฺจโย. (4)}

\prose{เนวทสฺสเนน นภาวนาย ปหาตพฺพ\-เหตุโก ธมฺโม ภาวนาย ปหาตพฺพ\-เหตุกสฺส จ เนวทสฺสเนน นภาวนาย ปหาตพฺพ\-เหตุกสฺส จ ธมฺมสฺส เหตุปจฺจเยน ปจฺจโย. อุทฺธจฺจ\-สหคโต โมโห สมฺปยุตฺตกานํ ขนฺธานํ จิตฺตสมุฏฺฐานานญฺจ รูปานํ เหตุปจฺจเยน ปจฺจโย. (5)}

\csromanheader{2}{\textbf{อารมฺมณ}}
\proseitem{77}{ทสฺสเนน ปหาตพฺพ\-เหตุโก ธมฺโม ทสฺสเนน ปหาตพฺพ\-เหตุกสฺส ธมฺมสฺส อารมฺมณ\-ปจฺจเยน ปจฺจโย. ทสฺสเนน ปหาตพฺพ\-เหตุกํ ราคํ อสฺสาเทติ อภินนฺทติ, ตํ อารพฺภ ทสฺสเนน ปหาตพฺพ\-เหตุโก ราโค อุปฺปชฺชติ, ทิฏฺฐิ อุปฺปชฺชติ, วิจิกิจฺฉา อุปฺปชฺชติ, ทสฺสเนน ปหาตพฺพ\-เหตุกํ โทมนสฺสํ อุปฺปชฺชติ. ทิฏฺฐิํ อสฺสาเทติ อภินนฺทติ, ตํ อารพฺภ ทสฺสเนน ปหาตพฺพ\-เหตุโก ราโค อุปฺปชฺชติ, ทิฏฺฐิ อุปฺปชฺชติ, วิจิกิจฺฉา อุปฺปชฺชติ, ทสฺสเนน ปหาตพฺพ\-เหตุกํ โทมนสฺสํ อุปฺปชฺชติ. วิจิกิจฺฉํ อารพฺภ วิจิกิจฺฉา อุปฺปชฺชติ, ทิฏฺฐิ อุปฺปชฺชติ, ทสฺสเนน ปหาตพฺพ\-เหตุกํ โทมนสฺสํ อุปฺปชฺชติ. ทสฺสเนน ปหาตพฺพ\-เหตุกํ โทมนสฺสํ อารพฺภ ทสฺสเนน ปหาตพฺพ\-เหตุกํ โทมนสฺสํ อุปฺปชฺชติ, ทิฏฺฐิ อุปฺปชฺชติ, วิจิกิจฺฉา อุปฺปชฺชติ. (1)}

\prose{ทสฺสเนน ปหาตพฺพ\-เหตุโก ธมฺโม เนวทสฺสเนน นภาวนาย ปหาตพฺพ\-เหตุกสฺส ธมฺมสฺส อารมฺมณ\-ปจฺจเยน ปจฺจโย. อริยา ทสฺสเนน ปหาตพฺพ\-เหตุเก ปหีเน กิเลเส ปจฺจเวกฺขนฺติ. ปุพฺเพ สมุทาจิณฺเณ กิเลเส ชานนฺติ. ทสฺสเนน ปหาตพฺพ\-เหตุเก ขนฺเธ อนิจฺจโต ฯเปฯ. เจโตปริย\-ญาเณน ฯเปฯ. ทสฺสเนน ปหาตพฺพ\-เหตุกา \csromanfolio{214}ขนฺธา เจโตปริย\-ญาณสฺส, ปุพฺเพ\-นิวาสา\-นุสฺสติ\-ญาณสฺส, ยถากมฺมูปคญาณสฺส, อนาคตํสญาณสฺส, อาวชฺชนาย โมหสฺส จ อารมฺมณ\-ปจฺจเยน ปจฺจโย. (2)}

\prose{ทสฺสเนน ปหาตพฺพ\-เหตุโก ธมฺโม ทสฺสเนน ปหาตพฺพ\-เหตุกสฺส จ เนวทสฺสเนน นภาวนาย ปหาตพฺพ\-เหตุกสฺส จ ธมฺมสฺส อารมฺมณ\-ปจฺจเยน ปจฺจโย. ทสฺสเนน ปหาตพฺพ\-เหตุเก ขนฺเธ อารพฺภ วิจิกิจฺฉา\-สหคตา ขนฺธา จ โมโห จ อุปฺปชฺชนฺติ. (3)}

\proseitem{78}{ภาวนาย ปหาตพฺพ\-เหตุโก ธมฺโม ภาวนาย ปหาตพฺพ\-เหตุกสฺส ธมฺมสฺส อารมฺมณ\-ปจฺจเยน ปจฺจโย. ภาวนาย ปหาตพฺพ\-เหตุกํ ราคํ อสฺสาเทติ อภินนฺทติ, ตํ อารพฺภ ภาวนาย ปหาตพฺพ\-เหตุโก ราโค อุปฺปชฺชติ, อุทฺธจฺจํ อุปฺปชฺชติ, ภาวนาย ปหาตพฺพ\-เหตุกํ โทมนสฺสํ อุปฺปชฺชติ. อุทฺธจฺจํ อารพฺภ อุทฺธจฺจํ อุปฺปชฺชติ, ภาวนาย ปหาตพฺพ\-เหตุกํ โทมนสฺสํ อุปฺปชฺชติ. ภาวนาย ปหาตพฺพ\-เหตุกํ โทมนสฺสํ อารพฺภ ภาวนาย ปหาตพฺพ\-เหตุกํ โทมนสฺสํ อุปฺปชฺชติ, อุทฺธจฺจํ อุปฺปชฺชติ. (1)}

\prose{ภาวนาย ปหาตพฺพ\-เหตุโก ธมฺโม ทสฺสเนน ปหาตพฺพ\-เหตุกสฺส ธมฺมสฺส อารมฺมณ\-ปจฺจเยน ปจฺจโย. ภาวนาย ปหาตพฺพ\-เหตุกํ ราคํ อสฺสาเทติ อภินนฺทติ, ตํ อารพฺภ ทสฺสเนน ปหาตพฺพ\-เหตุโก ราโค อุปฺปชฺชติ, ทิฏฺฐิ อุปฺปชฺชติ, วิจิกิจฺฉา อุปฺปชฺชติ. ทสฺสเนน ปหาตพฺพ\-เหตุกํ โทมนสฺสํ อุปฺปชฺชติ. อุทฺธจฺจํ อารพฺภ ทิฏฺฐิ อุปฺปชฺชติ, วิจิกิจฺฉา อุปฺปชฺชติ, ทสฺสเนน ปหาตพฺพ\-เหตุกํ โทมนสฺสํ อุปฺปชฺชติ. ภาวนาย ปหาตพฺพ\-เหตุกํ โทมนสฺสํ อารพฺภ ทสฺสเนน ปหาตพฺพ\-เหตุกํ โทมนสฺสํ อุปฺปชฺชติ, ทิฏฺฐิ อุปฺปชฺชติ, วิจิกิจฺฉา อุปฺปชฺชติ. (2)}

\prose{ภาวนาย ปหาตพฺพ\-เหตุโก ธมฺโม เนวทสฺสเนน นภาวนาย ปหาตพฺพ\-เหตุกสฺส ธมฺมสฺส อารมฺมณ\-ปจฺจเยน ปจฺจโย. อริยา ภาวนาย ปหาตพฺพ\-เหตุเก ปหีเน กิเลเส ปจฺจเวกฺขนฺติ, วิกฺขมฺภิเต กิเลเส ปจฺจเวกฺขนฺติ. ปุพฺเพ สมุทาจิณฺเณ กิเลเส ชานนฺติ. ภาวนาย ปหาตพฺพ\-เหตุเก ขนฺเธ อนิจฺจโต ฯเปฯ. เจโตปริย\-ญาเณน ฯเปฯ. \csromanfolio{215}ภาวนาย ปหาตพฺพ\-เหตุกา ขนฺธา เจโตปริย\-ญาณสฺส, ปุพฺเพ\-นิวาสา\-นุสฺสติ\-ญาณสฺส, ยถากมฺมูปคญาณสฺส, อนาคตํสญาณสฺส, อาวชฺชนาย, โมหสฺส จ อารมฺมณ\-ปจฺจเยน ปจฺจโย. (3)}

\prose{ภาวนาย ปหาตพฺพ\-เหตุโก ธมฺโม ทสฺสเนน ปหาตพฺพ\-เหตุกสฺส จ เนวทสฺสเนน นภาวนาย ปหาตพฺพ\-เหตุกสฺส จ ธมฺมสฺส อารมฺมณ\-ปจฺจเยน ปจฺจโย. ภาวนาย ปหาตพฺพ\-เหตุเก ขนฺเธ อารพฺภ วิจิกิจฺฉา\-สหคตา ขนฺธา จ โมโห จ อุปฺปชฺชนฺติ. (4)}

\prose{ภาวนาย ปหาตพฺพ\-เหตุโก ธมฺโม ภาวนาย ปหาตพฺพ\-เหตุกสฺส จ เนวทสฺสเนน นภาวนาย ปหาตพฺพ\-เหตุกสฺส จ ธมฺมสฺส อารมฺมณ\-ปจฺจเยน ปจฺจโย. ภาวนาย ปหาตพฺพ\-เหตุเก ขนฺเธ อารพฺภ อุทฺธจฺจ\-สหคตา ขนฺธา จ โมโห จ อุปฺปชฺชนฺติ. (5)}

\proseitem{79}{เนวทสฺสเนน นภาวนาย ปหาตพฺพ\-เหตุโก ธมฺโม เนวทสฺสเนน นภาวนาย ปหาตพฺพ\-เหตุกสฺส ธมฺมสฺส อารมฺมณ\-ปจฺจเยน ปจฺจโย. ทานํ ทตฺวา. (วิตฺถาเรตพฺพํ. ทสฺสน\-ตฺติก\-สทิสํ.) อาวชฺชนาย โมหสฺส จ อารมฺมณ\-ปจฺจเยน ปจฺจโย. (1)}

\prose{เนวทสฺสเนน นภาวนาย ปหาตพฺพ\-เหตุโก ธมฺโม ทสฺสเนน ปหาตพฺพ\-เหตุกสฺส ธมฺมสฺส อารมฺมณ\-ปจฺจเยน ปจฺจโย. ทานํ ทตฺวา. (ยถา ทสฺสนตฺติกํ.) (2)}

\prose{เนวทสฺสเนน นภาวนาย ปหาตพฺพ\-เหตุโก ธมฺโม ภาวนาย ปหาตพฺพ\-เหตุกสฺส ธมฺมสฺส อารมฺมณ\-ปจฺจเยน ปจฺจโย. ทานํ ทตฺวา. (ยถา ทสฺสนตฺติกํ.) (3)}

\prose{เนวทสฺสเนน นภาวนาย ปหาตพฺพ\-เหตุโก ธมฺโม ทสฺสเนน ปหาตพฺพ\-เหตุกสฺส จ เนวทสฺสเนน นภาวนาย ปหาตพฺพ\-เหตุกสฺส จ ธมฺมสฺส อารมฺมณ\-ปจฺจเยน ปจฺจโย. จกฺขุํ อารพฺภ วิจิกิจฺฉา\-สหคตา ขนฺธา จ โมโห จ อุปฺปชฺชนฺติ. โสตํ ฯเปฯ. วตฺถุํ. เนวทสฺสเนน นภาวนาย ปหาตพฺพ\-เหตุเก ขนฺเธ อารพฺภ วิจิกิจฺฉา\-สหคตา ขนฺธา จ โมโห จ อุปฺปชฺชนฺติ. (4)}

\prose{เนวทสฺสเนน นภาวนาย ปหาตพฺพ\-เหตุโก ธมฺโม ภาวนาย ปหาตพฺพ\-เหตุกสฺส จ เนวทสฺสเนน นภาวนาย ปหาตพฺพ\-เหตุกสฺส \csromanfolio{216}จ ธมฺมสฺส อารมฺมณ\-ปจฺจเยน ปจฺจโย. จกฺขุํ ฯเปฯ. วตฺถุํ. เนวทสฺสเนน นภาวนาย ปหาตพฺพ\-เหตุเก ขนฺเธ อารพฺภ อุทฺธจฺจ\-สหคตา ขนฺธา จ โมโห จ อุปฺปชฺชนฺติ. (5)}

\proseitem{80}{ทสฺสเนน ปหาตพฺพ\-เหตุโก จ เนวทสฺสเนน นภาวนาย ปหาตพฺพ\-เหตุโก จ ธมฺมา ทสฺสเนน ปหาตพฺพ\-เหตุกสฺส ธมฺมสฺส อารมฺมณ\-ปจฺจเยน ปจฺจโย. วิจิกิจฺฉา\-สหคเต ขนฺเธ จ โมหญฺจ อารพฺภ ทสฺสเนน ปหาตพฺพ\-เหตุกา ขนฺธา อุปฺปชฺชนฺติ. (1)}

\prose{ทสฺสเนน ปหาตพฺพ\-เหตุโก จ เนวทสฺสเนน นภาวนาย ปหาตพฺพ\-เหตุโก จ ธมฺมา เนวทสฺสเนน นภาวนาย ปหาตพฺพ\-เหตุกสฺส ธมฺมสฺส อารมฺมณ\-ปจฺจเยน ปจฺจโย. วิจิกิจฺฉา\-สหคเต ขนฺเธ จ โมหญฺจ อารพฺภ เนวทสฺสเนน นภาวนาย ปหาตพฺพ\-เหตุกา ขนฺธา จ โมโห จ อุปฺปชฺชนฺติ. (2)}

\prose{ทสฺสเนน ปหาตพฺพ\-เหตุโก จ เนวทสฺสเนน นภาวนาย ปหาตพฺพ\-เหตุโก จ ธมฺมา ทสฺสเนน ปหาตพฺพ\-เหตุกสฺส จ เนวทสฺสเนน นภาวนาย ปหาตพฺพ\-เหตุกสฺส จ ธมฺมสฺส อารมฺมณ\-ปจฺจเยน ปจฺจโย. วิจิกิจฺฉา\-สหคเต ขนฺเธ จ โมหญฺจ อารพฺภ วิจิกิจฺฉา\-สหคตา ขนฺธา จ โมโห จ อุปฺปชฺชนฺติ. (3)}

\proseitem{81}{ภาวนาย ปหาตพฺพ\-เหตุโก จ เนวทสฺสเนน นภาวนาย ปหาตพฺพ\-เหตุโก จ ธมฺมา ทสฺสเนน ปหาตพฺพ\-เหตุกสฺส ธมฺมสฺส อารมฺมณ\-ปจฺจเยน ปจฺจโย. อุทฺธจฺจ\-สหคเต ขนฺเธ จ โมหญฺจ อารพฺภ ทสฺสเนน ปหาตพฺพ\-เหตุกา ขนฺธา อุปฺปชฺชนฺติ. (1)}

\prose{ภาวนาย ปหาตพฺพ\-เหตุโก จ เนวทสฺสเนน นภาวนาย ปหาตพฺพ\-เหตุโก จ ธมฺมา ภาวนาย ปหาตพฺพ\-เหตุกสฺส ธมฺมสฺส อารมฺมณ\-ปจฺจเยน ปจฺจโย. อุทฺธจฺจ\-สหคเต ขนฺเธ จ โมหญฺจ อารพฺภ ภาวนาย ปหาตพฺพ\-เหตุกา ขนฺธา อุปฺปชฺชนฺติ. (2)}

\prose{ภาวนาย ปหาตพฺพ\-เหตุโก จ เนวทสฺสเนน นภาวนาย ปหาตพฺพ\-เหตุโก จ ธมฺมา เนวทสฺสเนน นภาวนาย ปหาตพฺพ\-เหตุกสฺส ธมฺมสฺส อารมฺมณ\-ปจฺจเยน ปจฺจโย. อุทฺธจฺจ\-สหคเต ขนฺเธ จ โมหญฺจ \csromanfolio{217}อารพฺภ เนวทสฺสเนน นภาวนาย ปหาตพฺพ\-เหตุกา ขนฺธา จ โมโห จ อุปฺปชฺชนฺติ. (3)}

\prose{ภาวนาย ปหาตพฺพ\-เหตุโก จ เนวทสฺสเนน นภาวนาย ปหาตพฺพ\-เหตุโก จ ธมฺมา ทสฺสเนน ปหาตพฺพ\-เหตุกสฺส จ เนวทสฺสเนน นภาวนาย ปหาตพฺพ\-เหตุกสฺส จ ธมฺมสฺส อารมฺมณ\-ปจฺจเยน ปจฺจโย. อุทฺธจฺจ\-สหคเต ขนฺเธ จ โมหญฺจ อารพฺภ วิจิกิจฺฉา\-สหคตา ขนฺธา จ โมโห จ อุปฺปชฺชนฺติ. (4)}

\prose{ภาวนาย ปหาตพฺพ\-เหตุโก จ เนวทสฺสเนน นภาวนาย ปหาตพฺพ\-เหตุโก จ ธมฺมา ภาวนาย ปหาตพฺพ\-เหตุกสฺส จ เนวทสฺสเนน นภาวนาย ปหาตพฺพ\-เหตุกสฺส จ ธมฺมสฺส อารมฺมณ\-ปจฺจเยน ปจฺจโย. อุทฺธจฺจ\-สหคเต ขนฺเธ จ โมหญฺจ อารพฺภ อุทฺธจฺจ\-สหคตา ขนฺธา จ โมโห จ อุปฺปชฺชนฺติ. (5)}

\proseitem{82}{ทสฺสเนน ปหาตพฺพ\-เหตุโก ธมฺโม ทสฺสเนน ปหาตพฺพ\-เหตุกสฺส ธมฺมสฺส อธิปติ\-ปจฺจเยน ปจฺจโย. (ทสฺสน\-ตฺติก\-สทิสํ, ทส ปญฺหา.)}

\csromanheader{2}{\textbf{อนนฺตร}}
\proseitem{83}{ทสฺสเนน ปหาตพฺพ\-เหตุโก ธมฺโม ทสฺสเนน ปหาตพฺพ\-เหตุกสฺส ธมฺมสฺส อนนฺตร\-ปจฺจเยน ปจฺจโย. ปุริมา ปุริมา ทสฺสเนน ปหาตพฺพ\-เหตุกา ขนฺธา ปจฺฉิมานํ ปจฺฉิมานํ ทสฺสเนน ปหาตพฺพ\-เหตุกานํ ขนฺธานํ อนนฺตร\-ปจฺจเยน ปจฺจโย. (1)}

\prose{ทสฺสเนน ปหาตพฺพ\-เหตุโก ธมฺโม เนวทสฺสเนน นภาวนาย ปหาตพฺพ\-เหตุกสฺส ธมฺมสฺส อนนฺตร\-ปจฺจเยน ปจฺจโย. ปุริมา ปุริมา วิจิกิจฺฉา\-สหคตา ขนฺธา ปจฺฉิมสฺส ปจฺฉิมสฺส โมหสฺส อนนฺตร\-ปจฺจเยน ปจฺจโย. ทสฺสเนน ปหาตพฺพ\-เหตุกา ขนฺธา วุฏฺฐานสฺส อนนฺตร\-ปจฺจเยน ปจฺจโย. (2)}

\prose{ทสฺสเนน ปหาตพฺพ\-เหตุโก ธมฺโม ทสฺสเนน ปหาตพฺพ\-เหตุกสฺส จ เนวทสฺสเนน นภาวนาย ปหาตพฺพ\-เหตุกสฺส จ \csromanfolio{218}ธมฺมสฺส อนนฺตร\-ปจฺจเยน ปจฺจโย. ปุริมา ปุริมา วิจิกิจฺฉา\-สหคตา ขนฺธา ปจฺฉิมานํ ปจฺฉิมานํ วิจิกิจฺฉา\-สหคตานํ ขนฺธานํ โมหสฺส จ อนนฺตร\-ปจฺจเยน ปจฺจโย. (3)}

\proseitem{84}{ภาวนาย ปหาตพฺพ\-เหตุโก ธมฺโม ภาวนาย ปหาตพฺพ\-เหตุกสฺส ธมฺมสฺส อนนฺตร\-ปจฺจเยน ปจฺจโย. ปุริมา ปุริมา ภาวนาย ปหาตพฺพ\-เหตุกา ขนฺธา ปจฺฉิมานํ ปจฺฉิมานํ ภาวนาย ปหาตพฺพ\-เหตุกานํ ขนฺธานํ อนนฺตร\-ปจฺจเยน ปจฺจโย. (1)}

\prose{ภาวนาย ปหาตพฺพ\-เหตุโก ธมฺโม เนวทสฺสเนน นภาวนาย ปหาตพฺพ\-เหตุกสฺส ธมฺมสฺส อนนฺตร\-ปจฺจเยน ปจฺจโย. ปุริมา ปุริมา อุทฺธจฺจ\-สหคตา ขนฺธา ปจฺฉิมสฺส ปจฺฉิมสฺส โมหสฺส อนนฺตร\-ปจฺจเยน ปจฺจโย. ภาวนาย ปหาตพฺพ\-เหตุกา ขนฺธา วุฏฺฐานสฺส อนนฺตร\-ปจฺจเยน ปจฺจโย. (2)}

\prose{ภาวนาย ปหาตพฺพ\-เหตุโก ธมฺโม ภาวนาย ปหาตพฺพ\-เหตุกสฺส จ เนวทสฺสเนน นภาวนาย ปหาตพฺพ\-เหตุกสฺส จ ธมฺมสฺส อนนฺตร\-ปจฺจเยน ปจฺจโย. ปุริมา ปุริมา อุทฺธจฺจ\-สหคตา ขนฺธา ปจฺฉิมานํ ปจฺฉิมานํ อุทฺธจฺจ\-สหคตานํ ขนฺธานํ โมหสฺส จ อนนฺตร\-ปจฺจเยน ปจฺจโย. (3)}

\proseitem{85}{เนวทสฺสเนน นภาวนาย ปหาตพฺพ\-เหตุโก ธมฺโม เนวทสฺสเนน นภาวนาย ปหาตพฺพ\-เหตุกสฺส ธมฺมสฺส อนนฺตร\-ปจฺจเยน ปจฺจโย. ปุริโม ปุริโม วิจิกิจฺฉา\-สหคโต อุทฺธจฺจ\-สหคโต โมโห ปจฺฉิมสฺส ปจฺฉิมสฺส วิจิกิจฺฉา\-สหคตสฺส อุทฺธจฺจ\-สหคตสฺส โมหสฺส อนนฺตร\-ปจฺจเยน ปจฺจโย. ปุริมา ปุริมา เนวทสฺสเนน นภาวนาย ปหาตพฺพ\-เหตุกา ขนฺธา ปจฺฉิมานํ ปจฺฉิมานํ เนวทสฺสเนน นภาวนาย ปหาตพฺพ\-เหตุกานํ ขนฺธานํ อนนฺตร\-ปจฺจเยน ปจฺจโย. อนุโลมํ โคตฺรภุสฺส. อนุโลมํ โวทานสฺส ฯเปฯ. นิโรธา วุฏฺฐหนฺตสฺส เนว\-สญฺญา\-นา\-สญฺญา\-ยตนํ ผลสมาปตฺติยา อนนฺตร\-ปจฺจเยน ปจฺจโย. (1)}

\prose{เนวทสฺสเนน นภาวนาย ปหาตพฺพ\-เหตุโก ธมฺโม ทสฺสเนน ปหาตพฺพ\-เหตุกสฺส ธมฺมสฺส อนนฺตร\-ปจฺจเยน ปจฺจโย. ปุริโม ปุริโม \csromanfolio{219}วิจิกิจฺฉา\-สหคโต โมโห ปจฺฉิมานํ ปจฺฉิมานํ วิจิกิจฺฉา\-สหคตานํ ขนฺธานํ อนนฺตร\-ปจฺจเยน ปจฺจโย. อาวชฺชนา ทสฺสเนน ปหาตพฺพ\-เหตุกานํ ขนฺธานํ อนนฺตร\-ปจฺจเยน ปจฺจโย. (2)}

\prose{เนวทสฺสเนน นภาวนาย ปหาตพฺพ\-เหตุโก ธมฺโม ภาวนาย ปหาตพฺพ\-เหตุกสฺส ธมฺมสฺส อนนฺตร\-ปจฺจเยน ปจฺจโย. ปุริโม ปุริโม อุทฺธจฺจ\-สหคโต โมโห ปจฺฉิมานํ ปจฺฉิมานํ อุทฺธจฺจ\-สหคตานํ ขนฺธานํ อนนฺตร\-ปจฺจเยน ปจฺจโย. อาวชฺชนา ภาวนาย ปหาตพฺพ\-เหตุกานํ ขนฺธานํ อนนฺตร\-ปจฺจเยน ปจฺจโย. (3)}

\prose{เนวทสฺสเนน นภาวนาย ปหาตพฺพ\-เหตุโก ธมฺโม ทสฺสเนน ปหาตพฺพ\-เหตุกสฺส จ เนวทสฺสเนน นภาวนาย ปหาตพฺพ\-เหตุกสฺส จ ธมฺมสฺส อนนฺตร\-ปจฺจเยน ปจฺจโย. ปุริโม ปุริโม วิจิกิจฺฉา\-สหคโต โมโห ปจฺฉิมานํ ปจฺฉิมานํ วิจิกิจฺฉา\-สหคตานํ ขนฺธานํ โมหสฺส จ อนนฺตร\-ปจฺจเยน ปจฺจโย. อาวชฺชนา วิจิกิจฺฉา\-สหคตานํ ขนฺธานํ โมหสฺส จ อนนฺตร\-ปจฺจเยน ปจฺจโย. (4)}

\prose{เนวทสฺสเนน นภาวนาย ปหาตพฺพ\-เหตุโก ธมฺโม ภาวนาย ปหาตพฺพ\-เหตุกสฺส จ เนวทสฺสเนน นภาวนาย ปหาตพฺพ\-เหตุกสฺส จ ธมฺมสฺส อนนฺตร\-ปจฺจเยน ปจฺจโย. ปุริโม ปุริโม อุทฺธจฺจ\-สหคโต โมโห ปจฺฉิมานํ ปจฺฉิมานํ อุทฺธจฺจ\-สหคตานํ ขนฺธานํ โมหสฺส จ อนนฺตร\-ปจฺจเยน ปจฺจโย. อาวชฺชนา อุทฺธจฺจ\-สหคตานํ ขนฺธานํ โมหสฺส จ อนนฺตร\-ปจฺจเยน ปจฺจโย. (5)}

\proseitem{86}{ทสฺสเนน ปหาตพฺพ\-เหตุโก จ เนวทสฺสเนน นภาวนาย ปหาตพฺพ\-เหตุโก จ ธมฺมา ทสฺสเนน ปหาตพฺพ\-เหตุกสฺส ธมฺมสฺส อนนฺตร\-ปจฺจเยน ปจฺจโย. ปุริมา ปุริมา วิจิกิจฺฉา\-สหคตา ขนฺธา จ โมโห จ ปจฺฉิมานํ ปจฺฉิมานํ วิจิกิจฺฉา\-สหคตานํ ขนฺธานํ อนนฺตร\-ปจฺจเยน ปจฺจโย. (1)}

\prose{ทสฺสเนน ปหาตพฺพ\-เหตุโก จ เนวทสฺสเนน นภาวนาย ปหาตพฺพ\-เหตุโก จ ธมฺมา เนวทสฺสเนน นภาวนาย ปหาตพฺพ\-เหตุกสฺส ธมฺมสฺส อนนฺตร\-ปจฺจเยน ปจฺจโย. ปุริมา ปุริมา วิจิกิจฺฉา\-สหคตา ขนฺธา จ โมโห จ ปจฺฉิมสฺส ปจฺฉิมสฺส โมหสฺส \csromanfolio{220}อนนฺตร\-ปจฺจเยน ปจฺจโย. วิจิกิจฺฉา\-สหคตา ขนฺธา จ โมโห จ วุฏฺฐานสฺส อนนฺตร\-ปจฺจเยน ปจฺจโย. (2)}

\prose{ทสฺสเนน ปหาตพฺพ\-เหตุโก จ เนวทสฺสเนน นภาวนาย ปหาตพฺพ\-เหตุโก จ ธมฺมา ทสฺสเนน ปหาตพฺพ\-เหตุกสฺส จ เนวทสฺสเนน นภาวนาย ปหาตพฺพ\-เหตุกสฺส จ ธมฺมสฺส อนนฺตร\-ปจฺจเยน ปจฺจโย. ปุริมา ปุริมา วิจิกิจฺฉา\-สหคตา ขนฺธา จ โมโห จ ปจฺฉิมานํ ปจฺฉิมานํ วิจิกิจฺฉา\-สหคตานํ ขนฺธานํ โมหสฺส จ อนนฺตร\-ปจฺจเยน ปจฺจโย. (3)}

\prose{ภาวนาย ปหาตพฺพ\-เหตุโก จ เนวทสฺสเนน นภาวนาย ปหาตพฺพ\-เหตุโก จ ธมฺมา ภาวนาย ปหาตพฺพ\-เหตุกสฺส ธมฺมสฺส อนนฺตร\-ปจฺจเยน ปจฺจโย. ตีณิ. (ทสฺสเนน สทิสํ คมนํ.)}

\csromanheader{2}{\textbf{สมนนฺตราทิ}}
\proseitem{87}{ทสฺสเนน ปหาตพฺพ\-เหตุโก ธมฺโม ทสฺสเนน ปหาตพฺพ\-เหตุกสฺส ธมฺมสฺส สมนนฺตร\-ปจฺจเยน ปจฺจโย. (อนนฺตร\-สทิสํ.) สหชาต\-ปจฺจเยน ปจฺจโย. (สํขิตฺตํ. ปฏิจฺจวาเร สหชาต\-สทิสํ.) อญฺญมญฺญ\-ปจฺจเยน ปจฺจโย. (สํขิตฺตํ. ปฏิจฺจวาเร อญฺญมญฺญ\-สทิสํ.) นิสฺสย\-ปจฺจเยน ปจฺจโย. (สํขิตฺตํ. ปจฺจยวาเร นิสฺสย\-วาร\-สทิสํ. วิสุํ ฆฏนา นตฺถิ.)}

\csromanheader{2}{\textbf{อุปนิสฺสย}}
\proseitem{88}{ทสฺสเนน ปหาตพฺพ\-เหตุโก ธมฺโม ทสฺสเนน ปหาตพฺพ\-เหตุกสฺส ธมฺมสฺส อุปนิสฺสย\-ปจฺจเยน ปจฺจโย. อารมฺมณู\-ปนิสฺสโย, อนนฺตรู\-ปนิสฺสโย, ปกตู\-ปนิสฺสโย ฯเปฯ. ปกตู\-ปนิสฺสโย\csromandash{}ทสฺสเนน ปหาตพฺพ\-เหตุกํ ราคํ อุปนิสฺสาย ปาณํ หนติ ฯเปฯ สํฆํ ภินฺทติ. ทสฺสเนน ปหาตพฺพ\-เหตุกํ โทสํ. โมหํ. ทิฏฺฐิํ. ปตฺถนํ อุปนิสฺสาย ปาณํ หนติ ฯเปฯ สํฆํ ภินฺทติ. ทสฺสเนน ปหาตพฺพ\-เหตุโก ราโค. โทโส. โมโห. ทิฏฺฐิ. ปตฺถนา ทสฺสเนน ปหาตพฺพ\-เหตุกสฺส ราคสฺส ฯเปฯ. ปตฺถนาย อุปนิสฺสย\-ปจฺจเยน ปจฺจโย. (1)}

\prose{\csromanfolio{221}ทสฺสเนน ปหาตพฺพ\-เหตุโก ธมฺโม เนวทสฺสเนน นภาวนาย ปหาตพฺพ\-เหตุกสฺส ธมฺมสฺส อุปนิสฺสย\-ปจฺจเยน ปจฺจโย. อนนฺตรู\-ปนิสฺสโย, ปกตู\-ปนิสฺสโย ฯเปฯ. ปกตู\-ปนิสฺสโย\csromandash{}ทสฺสเนน ปหาตพฺพ\-เหตุกํ ราคํ อุปนิสฺสาย ทานํ เทติ ฯเปฯ สมาปตฺติํ อุปฺปาเทติ. ทสฺสเนน ปหาตพฺพ\-เหตุกํ โทสํ. โมหํ. ทิฏฺฐิํ. ปตฺถนํ อุปนิสฺสาย ทานํ เทติ ฯเปฯ สมาปตฺติํ อุปฺปาเทติ. ทสฺสเนน ปหาตพฺพ\-เหตุโก ราโค ฯเปฯ. ปตฺถนา. สทฺธาย ฯเปฯ ปญฺญาย. กายิกสฺส สุขสฺส, กายิกสฺส ทุกฺขสฺส, ผลสมาปตฺติยา, โมหสฺส จ อุปนิสฺสย\-ปจฺจเยน ปจฺจโย. (2)}

\prose{ทสฺสเนน ปหาตพฺพ\-เหตุโก ธมฺโม ทสฺสเนน ปหาตพฺพ\-เหตุกสฺส จ เนวทสฺสเนน นภาวนาย ปหาตพฺพ\-เหตุกสฺส จ ธมฺมสฺส อุปนิสฺสย\-ปจฺจเยน ปจฺจโย. อนนฺตรู\-ปนิสฺสโย, ปกตู\-ปนิสฺสโย ฯเปฯ. ปกตู\-ปนิสฺสโย\csromandash{}ทสฺสเนน ปหาตพฺพ\-เหตุโก ราโค. โทโส. โมโห. ทิฏฺฐิ. ปตฺถนา วิจิกิจฺฉา\-สหคตานํ ขนฺธานํ โมหสฺส จ อุปนิสฺสย\-ปจฺจเยน ปจฺจโย. (3)}

\proseitem{89}{ภาวนาย ปหาตพฺพ\-เหตุโก ธมฺโม ภาวนาย ปหาตพฺพ\-เหตุกสฺส ธมฺมสฺส อุปนิสฺสย\-ปจฺจเยน ปจฺจโย. อารมฺมณู\-ปนิสฺสโย, อนนฺตรู\-ปนิสฺสโย, ปกตู\-ปนิสฺสโย ฯเปฯ. ปกตู\-ปนิสฺสโย\csromandash{}ภาวนาย ปหาตพฺพ\-เหตุโก ราโค. โทโส. โมโห. มาโน. ปตฺถนา ภาวนาย ปหาตพฺพ\-เหตุกสฺส ราคสฺส. โทสสฺส. โมหสฺส. มานสฺส. ปตฺถนาย อุปนิสฺสย\-ปจฺจเยน ปจฺจโย. (1)}

\prose{ภาวนาย ปหาตพฺพ\-เหตุโก ธมฺโม ทสฺสเนน ปหาตพฺพ\-เหตุกสฺส ธมฺมสฺส อุปนิสฺสย\-ปจฺจเยน ปจฺจโย. อารมฺมณู\-ปนิสฺสโย, ปกตู\-ปนิสฺสโย ฯเปฯ. ปกตู\-ปนิสฺสโย\csromandash{}ภาวนาย ปหาตพฺพ\-เหตุกํ ราคํ อุปนิสฺสาย ปาณํ หนติ ฯเปฯ สํฆํ ภินฺทติ. ภาวนาย ปหาตพฺพ\-เหตุกํ โทสํ. โมหํ. มานํ. ปตฺถนํ อุปนิสฺสาย ปาณํ หนติ ฯเปฯ สํฆํ ภินฺทติ. ภาวนาย ปหาตพฺพ\-เหตุโก ราโค ฯเปฯ ปตฺถนา ทสฺสเนน ปหาตพฺพ\-เหตุกสฺส ราคสฺส. โทสสฺส. โมหสฺส. ทิฏฺฐิยา. ปตฺถนาย อุปนิสฺสย\-ปจฺจเยน ปจฺจโย. สกภณฺเฑ ฉนฺทราโค \csromanfolio{222}ปรภณฺเฑ ฉนฺทราคสฺส อุปนิสฺสย\-ปจฺจเยน ปจฺจโย. สกปริคฺคเห ฉนฺทราโค ปรปริคฺคเห ฉนฺทราคสฺส อุปนิสฺสย\-ปจฺจเยน ปจฺจโย. (2)}

\prose{ภาวนาย ปหาตพฺพ\-เหตุโก ธมฺโม เนวทสฺสเนน นภาวนาย ปหาตพฺพ\-เหตุกสฺส ธมฺมสฺส อุปนิสฺสย\-ปจฺจเยน ปจฺจโย. อนนฺตรู\-ปนิสฺสโย, ปกตู\-ปนิสฺสโย ฯเปฯ. ปกตู\-ปนิสฺสโย\csromandash{}ภาวนาย ปหาตพฺพ\-เหตุกํ ราคํ อุปนิสฺสาย ทานํ เทติ ฯเปฯ สมาปตฺติํ อุปฺปาเทติ. ภาวนาย ปหาตพฺพ\-เหตุกํ โทสํ. โมหํ. มานํ. ปตฺถนํ อุปนิสฺสาย ทานํ เทติ ฯเปฯ สมาปตฺติํ อุปฺปาเทติ. ภาวนาย ปหาตพฺพ\-เหตุโก ราโค ฯเปฯ. ปตฺถนา สทฺธาย ฯเปฯ ปญฺญาย. กายิกสฺส สุขสฺส, กายิกสฺส ทุกฺขสฺส, ผลสมาปตฺติยา, โมหสฺส จ อุปนิสฺสย\-ปจฺจเยน ปจฺจโย. (3)}

\prose{ภาวนาย ปหาตพฺพ\-เหตุโก ธมฺโม ทสฺสเนน ปหาตพฺพ\-เหตุกสฺส จ เนวทสฺสเนน นภาวนาย ปหาตพฺพ\-เหตุกสฺส จ ธมฺมสฺส อุปนิสฺสย\-ปจฺจเยน ปจฺจโย.  ปกตู\-ปนิสฺสโย\csromandash{}ภาวนาย ปหาตพฺพ\-เหตุโก ราโค ฯเปฯ ปตฺถนา วิจิกิจฺฉา\-สหคตานํ ขนฺธานํ โมหสฺส จ อุปนิสฺสย\-ปจฺจเยน ปจฺจโย. (4)}

\prose{ภาวนาย ปหาตพฺพ\-เหตุโก ธมฺโม ภาวนาย ปหาตพฺพ\-เหตุกสฺส จ เนวทสฺสเนน นภาวนาย ปหาตพฺพ\-เหตุกสฺส จ ธมฺมสฺส อุปนิสฺสย\-ปจฺจเยน ปจฺจโย. อนนฺตรู\-ปนิสฺสโย, ปกตู\-ปนิสฺสโย ฯเปฯ. ปกตู\-ปนิสฺสโย\csromandash{}ภาวนาย ปหาตพฺพ\-เหตุโก ราโค ฯเปฯ ปตฺถนา อุทฺธจฺจ\-สหคตานํ ขนฺธานํ โมหสฺส จ อุปนิสฺสย\-ปจฺจเยน ปจฺจโย. (5)}

\proseitem{90}{เนวทสฺสเนน นภาวนาย ปหาตพฺพ\-เหตุโก ธมฺโม เนวทสฺสเนน นภาวนาย ปหาตพฺพ\-เหตุกสฺส ธมฺมสฺส อุปนิสฺสย\-ปจฺจเยน ปจฺจโย. อารมฺมณู\-ปนิสฺสโย, อนนฺตรู\-ปนิสฺสโย, ปกตู\-ปนิสฺสโย ฯเปฯ. ปกตู\-ปนิสฺสโย\csromandash{}สทฺธํ อุปนิสฺสาย ทานํ เทติ ฯเปฯ. สมาปตฺติํ อุปฺปาเทติ. สีลํ ฯเปฯ. ปญฺญํ. กายิกํ สุขํ. กายิกํ ทุกฺขํ. อุตุํ. โภชนํ. เสนาสนํ. โมหํ อุปนิสฺสาย ทานํ เทติ ฯเปฯ. สทฺธา ฯเปฯ. โมโห สทฺธาย ฯเปฯ ผลสมาปตฺติยา โมหสฺส จ อุปนิสฺสย\-ปจฺจเยน ปจฺจโย. (1)}

\prose{\csromanfolio{223}เนวทสฺสเนน นภาวนาย ปหาตพฺพ\-เหตุโก ธมฺโม ทสฺสเนน ปหาตพฺพ\-เหตุกสฺส ธมฺมสฺส อุปนิสฺสย\-ปจฺจเยน ปจฺจโย. อารมฺมณู\-ปนิสฺสโย, อนนฺตรู\-ปนิสฺสโย, ปกตู\-ปนิสฺสโย ฯเปฯ. ปกตู\-ปนิสฺสโย\csromandash{}สทฺธํ อุปนิสฺสาย ทิฏฺฐิํ คณฺหาติ. สีลํ ฯเปฯ. ปญฺญํ. กายิกํ สุขํ. กายิกํ ทุกฺขํ ฯเปฯ. เสนาสนํ. โมหํ อุปนิสฺสาย ปาณํ หนติ ฯเปฯ สํฆํ ภินฺทติ. สทฺธา ฯเปฯ. เสนาสนํ โมโห จ ทสฺสเนน ปหาตพฺพ\-เหตุกสฺส ราคสฺส ฯเปฯ ปตฺถนาย อุปนิสฺสย\-ปจฺจเยน ปจฺจโย. (2)}

\prose{เนวทสฺสเนน นภาวนาย ปหาตพฺพ\-เหตุโก ธมฺโม ภาวนาย ปหาตพฺพ\-เหตุกสฺส ธมฺมสฺส อุปนิสฺสย\-ปจฺจเยน ปจฺจโย. อารมฺมณู\-ปนิสฺสโย, อนนฺตรู\-ปนิสฺสโย, ปกตู\-ปนิสฺสโย ฯเปฯ. ปกตู\-ปนิสฺสโย\csromandash{}สทฺธํ อุปนิสฺสาย มานํ ชปฺเปติ ฯเปฯ. โมหํ อุปนิสฺสาย มานํ ชปฺเปติ. สทฺธา ฯเปฯ. เสนาสนํ โมโห จ ภาวนาย ปหาตพฺพ\-เหตุกสฺส ราคสฺส ฯเปฯ ปตฺถนาย อุปนิสฺสย\-ปจฺจเยน ปจฺจโย. (3)}

\prose{เนวทสฺสเนน นภาวนาย ปหาตพฺพ\-เหตุโก ธมฺโม ทสฺสเนน ปหาตพฺพ\-เหตุกสฺส จ เนวทสฺสเนน นภาวนาย ปหาตพฺพ\-เหตุกสฺส จ ธมฺมสฺส อุปนิสฺสย\-ปจฺจเยน ปจฺจโย. อนนฺตรู\-ปนิสฺสโย, ปกตู\-ปนิสฺสโย ฯเปฯ. ปกตู\-ปนิสฺสโย\csromandash{}สทฺธา ฯเปฯ. ปญฺญา. กายิกํ สุขํ. กายิกํ ทุกฺขํ ฯเปฯ. เสนาสนํ โมโห จ วิจิกิจฺฉา\-สหคตานํ ขนฺธานํ โมหสฺส จ อุปนิสฺสย\-ปจฺจเยน ปจฺจโย. (4)}

\prose{เนวทสฺสเนน นภาวนาย ปหาตพฺพ\-เหตุโก ธมฺโม ภาวนาย ปหาตพฺพ\-เหตุกสฺส จ เนวทสฺสเนน นภาวนาย ปหาตพฺพ\-เหตุกสฺส จ ธมฺมสฺส อุปนิสฺสย\-ปจฺจเยน ปจฺจโย. อนนฺตรู\-ปนิสฺสโย, ปกตู\-ปนิสฺสโย ฯเปฯ. ปกตู\-ปนิสฺสโย\csromandash{}สทฺธา ฯเปฯ. เสนาสนํ โมโห จ อุทฺธจฺจ\-สหคตานํ ขนฺธานํ โมหสฺส จ อุปนิสฺสย\-ปจฺจเยน ปจฺจโย. (5)}

\proseitem{91}{ทสฺสเนน ปหาตพฺพ\-เหตุโก จ เนวทสฺสเนน นภาวนาย ปหาตพฺพ\-เหตุโก จ ธมฺมา ทสฺสเนน ปหาตพฺพ\-เหตุกสฺส ธมฺมสฺส อุปนิสฺสย\-ปจฺจเยน ปจฺจโย. อนนฺตรู\-ปนิสฺสโย, ปกตู\-ปนิสฺสโย ฯเปฯ. ปกตู\-ปนิสฺสโย\csromandash{}วิจิกิจฺฉา\-สหคตา ขนฺธา จ โมโห จ ทสฺสเนน \csromanfolio{224}ปหาตพฺพ\-เหตุกสฺส ราคสฺส ฯเปฯ ปตฺถนาย อุปนิสฺสย\-ปจฺจเยน ปจฺจโย. (1)}

\prose{ทสฺสเนน ปหาตพฺพ\-เหตุโก จ เนวทสฺสเนน นภาวนาย ปหาตพฺพ\-เหตุโก จ ธมฺมา เนวทสฺสเนน นภาวนาย ปหาตพฺพ\-เหตุกสฺส ธมฺมสฺส อุปนิสฺสย\-ปจฺจเยน ปจฺจโย. อนนฺตรู\-ปนิสฺสโย, ปกตู\-ปนิสฺสโย ฯเปฯ. ปกตู\-ปนิสฺสโย\csromandash{} วิจิกิจฺฉา\-สหคตา ขนฺธา จ โมโห จ สทฺธาย ฯเปฯ ปญฺญาย. กายิกสฺส สุขสฺส, กายิกสฺส ทุกฺขสฺส, ผลสมาปตฺติยา, โมหสฺส จ อุปนิสฺสย\-ปจฺจเยน ปจฺจโย. (2)}

\prose{ทสฺสเนน ปหาตพฺพ\-เหตุโก จ เนวทสฺสเนน นภาวนาย ปหาตพฺพ\-เหตุโก จ ธมฺมา ทสฺสเนน ปหาตพฺพ\-เหตุกสฺส จ เนวทสฺสเนน นภาวนาย ปหาตพฺพ\-เหตุกสฺส จ ธมฺมสฺส อุปนิสฺสย\-ปจฺจเยน ปจฺจโย. อนนฺตรู\-ปนิสฺสโย, ปกตู\-ปนิสฺสโย ฯเปฯ. ปกตู\-ปนิสฺสโย\csromandash{}วิจิกิจฺฉา\-สหคตา ขนฺธา จ โมโห จ วิจิกิจฺฉา\-สหคตานํ ขนฺธานํ โมหสฺส จ อุปนิสฺสย\-ปจฺจเยน ปจฺจโย. (3)}

\proseitem{92}{ภาวนาย ปหาตพฺพ\-เหตุโก จ เนวทสฺสเนน นภาวนาย ปหาตพฺพ\-เหตุโก จ ธมฺมา ทสฺสเนน ปหาตพฺพ\-เหตุกสฺส ธมฺมสฺส อุปนิสฺสย\-ปจฺจเยน ปจฺจโย. ปกตู\-ปนิสฺสโย\csromandash{}อุทฺธจฺจ\-สหคตา ขนฺธา จ โมโห จ ทสฺสเนน ปหาตพฺพ\-เหตุกสฺส ราคสฺส ฯเปฯ ปตฺถนาย อุปนิสฺสย\-ปจฺจเยน ปจฺจโย. (1)}

\prose{ภาวนาย ปหาตพฺพ\-เหตุโก จ เนวทสฺสเนน นภาวนาย ปหาตพฺพ\-เหตุโก จ ธมฺมา ภาวนาย ปหาตพฺพ\-เหตุกสฺส ธมฺมสฺส อุปนิสฺสย\-ปจฺจเยน ปจฺจโย. อนนฺตรู\-ปนิสฺสโย, ปกตู\-ปนิสฺสโย ฯเปฯ. ปกตู\-ปนิสฺสโย\csromandash{}อุทฺธจฺจ\-สหคตา ขนฺธา จ โมโห จ ภาวนาย ปหาตพฺพ\-เหตุกสฺส ราคสฺส ฯเปฯ ปตฺถนาย อุปนิสฺสายปจฺจเยน ปจฺจโย. (2)}

\prose{ภาวนาย ปหาตพฺพ\-เหตุโก จ เนวทสฺสเนน นภาวนาย ปหาตพฺพ\-เหตุโก จ ธมฺมา เนวทสฺสเนน นภาวนาย ปหตพฺพเหตุกสฺส ธมฺมสฺส อุปนิสฺสย\-ปจฺจเยน ปจฺจโย. อนนฺตรู\-ปนิสฺสโย, ปกตู\-ปนิสฺสโย ฯเปฯ. ปกตู\-ปนิสฺสโย\csromandash{} อุทฺธจฺจ\-สหคตา ขนฺธา จ \csromanfolio{225}โมโห จ สทฺธาย ฯเปฯ ผลสมาปตฺติยา โมหสฺส จ อุปนิสฺสย\-ปจฺจเยน ปจฺจโย. (3)}

\prose{ภาวนาย ปหาตพฺพ\-เหตุโก จ เนวทสฺสเนน นภาวนาย ปหาตพฺพ\-เหตุโก จ ธมฺมา ทสฺสเนน ปหาตพฺพ\-เหตุกสฺส จ เนวทสฺสเนน นภาวนาย ปหาตพฺพ\-เหตุกสฺส จ ธมฺมสฺส อุปนิสฺสย\-ปจฺจเยน ปจฺจโย. ปกตู\-ปนิสฺสโย\csromandash{}อุทฺธจฺจ\-สหคตา ขนฺธา จ โมโห จ วิจิกิจฺฉา\-สหคตานํ ขนฺธานํ โมหสฺส จ อุปนิสฺสย\-ปจฺจเยน ปจฺจโย. (4)}

\prose{ภาวนาย ปหาตพฺพ\-เหตุโก จ เนวทสฺสเนน นภาวนาย ปหาตพฺพ\-เหตุโก จ ธมฺมา ภาวนาย ปหาตพฺพ\-เหตุกสฺส จ เนวทสฺสเนน นภาวนาย ปหาตพฺพ\-เหตุกสฺส จ ธมฺมสฺส อุปนิสฺสย\-ปจฺจเยน ปจฺจโย. อนนฺตรู\-ปนิสฺสโย, ปกตู\-ปนิสฺสโย ฯเปฯ. ปกตู\-ปนิสฺสโย\csromandash{}อุทฺธจฺจ\-สหคตา ขนฺธา จ โมโห จ อุทฺธจฺจ\-สหคตานํ ขนฺธานํ โมหสฺส จ อุปนิสฺสย\-ปจฺจเยน ปจฺจโย. (5)}

\csromanheader{2}{\textbf{ปุเรชาต}}
\proseitem{93}{เนวทสฺสเนน นภาวนาย ปหาตพฺพ\-เหตุโก ธมฺโม เนวทสฺสเนน นภาวนาย ปหาตพฺพ\-เหตุกสฺส ธมฺมสฺส ปุเรชาต\-ปจฺจเยน ปจฺจโย. อารมฺมณ\-ปุเรชาตํ, วตฺถุ\-ปุเรชาตํ.  อารมฺมณ\-ปุเรชาตํ\csromandash{}จกฺขุํ อนิจฺจโต ฯเปฯ วิปสฺสติ. โสตํ ฯเปฯ. วตฺถุํ อนิจฺจโต ฯเปฯ วิปสฺสติ. ทิพฺเพน จกฺขุนา รูปํ ปสฺสติ. ทิพฺพาย โสตธาตุยา สทฺทํ สุณาติ. รูปายตนํ จกฺขุ\-วิญฺญาณสฺส ฯเปฯ. โผฏฺฐพฺพา\-ยตนํ กายวิญฺญาณสฺส ปุเรชาต\-ปจฺจเยน ปจฺจโย.  วตฺถุ\-ปุเรชาตํ\csromandash{}จกฺขายตนํ จกฺขุ\-วิญฺญาณสฺส ฯเปฯ. กายายตนํ กายวิญฺญาณสฺส ฯเปฯ. วตฺถุ เนวทสฺสเนน นภาวนาย ปหาตพฺพ\-เหตุกานํ ขนฺธานํ โมหสฺส จ ปุเรชาต\-ปจฺจเยน ปจฺจโย. (1)}

\prose{เนวทสฺสเนน นภาวนาย ปหาตพฺพ\-เหตุโก ธมฺโม ทสฺสเนน ปหาตพฺพ\-เหตุกสฺส ธมฺมสฺส ปุเรชาต\-ปจฺจเยน ปจฺจโย. อารมฺมณ\-ปุเรชาตํ, วตฺถุ\-ปุเรชาตํ.  อารมฺมณ\-ปุเรชาตํ\csromandash{}จกฺขุํ ฯเปฯ. วตฺถุํ อสฺสาเทติ อภินนฺทติ, ตํ อารพฺภ ทสฺสเนน ปหาตพฺพ\-เหตุโก ราโค ฯเปฯ ทิฏฺฐิ ฯเปฯ วิจิกิจฺฉา ฯเปฯ ทสฺสเนน ปหาตพฺพ\-เหตุกํ โทมนสฺสํ \csromanfolio{226}อุปฺปชฺชติ. วตฺถุ\-ปุเรชาตํ\csromandash{}วตฺถุ ทสฺสเนน ปหาตพฺพ\-เหตุกานํ ขนฺธานํ ปุเรชาต\-ปจฺจเยน ปจฺจโย. (2)}

\prose{เนวทสฺสเนน นภาวนาย ปหาตพฺพ\-เหตุโก ธมฺโม ภาวนาย ปหาตพฺพ\-เหตุกสฺส ธมฺมสฺส ปุเรชาต\-ปจฺจเยน ปจฺจโย. อารมฺมณ\-ปุเรชาตํ, วตฺถุ\-ปุเรชาตํ. อารมฺมณ\-ปุเรชาตํ\csromandash{}จกฺขุํ ฯเปฯ. วตฺถุํ อสฺสาเทติ อภินนฺทติ, ตํ อารพฺภ ภาวนาย ปหาตพฺพ\-เหตุโก ราโค อุปฺปชฺชติ, อุทฺธจฺจํ อุปฺปชฺชติ. ภาวนาย ปหาตพฺพ\-เหตุกํ โทมนสฺสํ อุปฺปชฺชติ.  วตฺถุ\-ปุเรชาตํ\csromandash{}วตฺถุ ภาวนาย ปหาตพฺพ\-เหตุกานํ ขนฺธานํ ปุเรชาต\-ปจฺจเยน ปจฺจโย. (3)}

\prose{เนวทสฺสเนน นภาวนาย ปหาตพฺพ\-เหตุโก ธมฺโม ทสฺสเนน ปหาตพฺพ\-เหตุกสฺส จ เนวทสฺสเนน นภาวนาย ปหาตพฺพ\-เหตุกสฺส จ ธมฺมสฺส ปุเรชาต\-ปจฺจเยน ปจฺจโย. อารมฺมณ\-ปุเรชาตํ, วตฺถุ\-ปุเรชาตํ.  อารมฺมณ\-ปุเรชาตํ\csromandash{}จกฺขุํ ฯเปฯ. วตฺถุํ อารพฺภ วิจิกิจฺฉา\-สหคตา ขนฺธา จ โมโห จ อุปฺปชฺชนฺติ. วตฺถุ\-ปุเรชาตํ\csromandash{} วตฺถุ วิจิกิจฺฉา\-สหคตานํ ขนฺธานํ โมหสฺส จ ปุเรชาต\-ปจฺจเยน ปจฺจโย. (4)}

\prose{เนวทสฺสเนน นภาวนาย ปหาตพฺพ\-เหตุโก ธมฺโม ภาวนาย ปหาตพฺพ\-เหตุกสฺส จ เนวทสฺสเนน นภาวนาย ปหาตพฺพ\-เหตุกสฺส จ ธมฺมสฺส ปุเรชาต\-ปจฺจเยน ปจฺจโย. อารมฺมณ\-ปุเรชาตํ, วตฺถุ\-ปุเรชาตํ.  อารมฺมณ\-ปุเรชาตํ\csromandash{}จกฺขุํ ฯเปฯ. วตฺถุํ อารพฺภ อุทฺธจฺจ\-สหคตา ขนฺธา จ โมโห จ อุปฺปชฺชนฺติ.  วตฺถุ\-ปุเรชาตํ\csromandash{} วตฺถุ อุทฺธจฺจ\-สหคตานํ ขนฺธานํ โมหสฺส จ ปุเรชาต\-ปจฺจเยน ปจฺจโย. (5)}

\csromanheader{2}{\textbf{ปจฺฉาชาต}}
\proseitem{94}{ทสฺสเนน ปหาตพฺพ\-เหตุโก ธมฺโม เนวทสฺสเนน นภาวนาย ปหาตพฺพ\-เหตุกสฺส ธมฺมสฺส ปจฺฉาชาต\-ปจฺจเยน ปจฺจโย. ปจฺฉาชาตา ทสฺสเนน ปหาตพฺพ\-เหตุกา ขนฺธา ปุเรชาตสฺส อิมสฺส กายสฺส ปจฺฉาชาต\-ปจฺจเยน ปจฺจโย. (1)}

\prose{ภาวนาย ปหาตพฺพ\-เหตุโก ธมฺโม เนวทสฺสเนน นภาวนาย ปหาตพฺพ\-เหตุกสฺส ธมฺมสฺส ปจฺฉาชาต\-ปจฺจเยน ปจฺจโย. ปจฺฉาชาตา \csromanfolio{227}ภาวนาย ปหาตพฺพ\-เหตุกา ขนฺธา ปุเรชาตสฺส อิมสฺส กายสฺส ปจฺฉาชาต\-ปจฺจเยน ปจฺจโย. (1)}

\prose{เนวทสฺสเนน นภาวนาย ปหาตพฺพ\-เหตุโก ธมฺโม เนวทสฺสเนน นภาวนาย ปหาตพฺพ\-เหตุกสฺส ธมฺมสฺส ปจฺฉาชาตปจฺจเยเน ปจฺจโย. ปจฺฉาชาตา เนวทสฺสเนน นภาวนาย ปหาตพฺพ\-เหตุกา ขนฺธา ปุเรชาตสฺส อิมสฺส กายสฺส ปจฺฉาชาต\-ปจฺจเยน ปจฺจโย. (1)}

\prose{ทสฺสเนน ปหาตพฺพ\-เหตุโก จ เนวทสฺสเนน นภาวนาย ปหาตพฺพ\-เหตุโก จ ธมฺมา เนวทสฺสเนน นภาวนาย ปหาตพฺพ\-เหตุกสฺส ธมฺมสฺส ปจฺฉาชาต\-ปจฺจเยน ปจฺจโย. ปจฺฉาชาตา วิจิกิจฺฉา\-สหคตา ขนฺธา จ โมโห จ ปุเรชาตสฺส อิมสฺส กายสฺส ปจฺฉาชาต\-ปจฺจเยน ปจฺจโย. (1)}

\prose{ภาวนาย ปหาตพฺพ\-เหตุโก จ เนวทสฺสเนน นภาวนาย ปหาตพฺพ\-เหตุโก จ ธมฺมา เนวทสฺสเนน นภาวนาย ปหาตพฺพ\-เหตุกสฺส ธมฺมสฺส ปจฺฉาชาต\-ปจฺจเยน ปจฺจโย. ปจฺฉาชาตา อุทฺธจฺจ\-สหคตา ขนฺธา จ โมโห จ ปุเรชาตสฺส อิมสฺส กายสฺส ปจฺฉาชาต\-ปจฺจเยน ปจฺจโย. (1)}

\csromanheader{2}{\textbf{อาเสวน}}
\proseitem{95}{ทสฺสเนน ปหาตพฺพ\-เหตุโก ธมฺโม ทสฺสเนน ปหาตพฺพ\-เหตุกสฺส ธมฺมสฺส อาเสวน\-ปจฺจเยน ปจฺจโย. ปุริมา ปุริมา ทสฺสเนน ปหาตพฺพ\-เหตุกา ขนฺธา ปจฺฉิมานํ ปจฺฉิมานํ ทสฺสเนน ปหาตพฺพ\-เหตุกานํ ขนฺธานํ อาเสวน\-ปจฺจเยน ปจฺจโย. (1)}

\prose{ทสฺสเนน ปหาตพฺพ\-เหตุโก ธมฺโม เนวทสฺสเนน นภาวนาย ปหาตพฺพ\-เหตุกสฺส ธมฺมสฺส อาเสวน\-ปจฺจเยน ปจฺจโย. ปุริมา ปุริมา วิจิกิจฺฉา\-สหคตา ขนฺธา ปจฺฉิมสฺส ปจฺฉิมสฺส โมหสฺส อาเสวน\-ปจฺจเยน ปจฺจโย. (2)}

\prose{ทสฺสเนน ปหาตพฺพ\-เหตุโก ธมฺโม ทสฺสเนน ปหาตพฺพ\-เหตุกสฺส จ เนวทสฺสเนน นภาวนาย ปหาตพฺพ\-เหตุกสฺส จ ธมฺมสฺส อาเสวน\-ปจฺจเยน ปจฺจโย. ปุริมา ปุริมา วิจิกิจฺฉา\-สหคตา ขนฺธา \csromanfolio{228}ปจฺฉิมานํ ปจฺฉิมานํ วิจิกิจฺฉา\-สหคตานํ ขนฺธานํ โมหสฺส จ อาเสวน\-ปจฺจเยน ปจฺจโย.}

\prose{ภาวนาย ปหาตพฺพ\-เหตุโก ธมฺโม ภาวนาย ปหาตพฺพ\-เหตุกสฺส ธมฺมสฺส. (สํขิตฺตํ.) ตีณิ.}

\prose{เนวทสฺสเนน นภาวนาย ปหาตพฺพ\-เหตุโก ธมฺโม ฯเปฯ.}

\vspace{\tipitakaparskip}
\csromancenter{(อาเสวนมูลเก วุฏฺฐานสฺสปิ อาวชฺชนายปิ ปหาตพฺพํ, สตฺตรส ปญฺหา ปริปุณฺณา, อนนฺตรสทิสา.)}
\csromanheader{2}{\textbf{กมฺม}}
\proseitem{96}{ทสฺสเนน ปหาตพฺพ\-เหตุโก ธมฺโม ทสฺสเนน ปหาตพฺพ\-เหตุกสฺส ธมฺมสฺส กมฺมปจฺจเยน ปจฺจโย. ทสฺสเนน ปหาตพฺพ\-เหตุกา เจตนา สมฺปยุตฺตกานํ ขนฺธานํ กมฺมปจฺจเยน ปจฺจโย. (1)}

\prose{ทสฺสเนน ปหาตพฺพ\-เหตุโก ธมฺโม เนวทสฺสเนน นภาวนาย ปหาตพฺพ\-เหตุกสฺส ธมฺมสฺส กมฺมปจฺจเยน ปจฺจโย. สหชาตา, นานากฺขณิกา.  สหชาตา ทสฺสเนน ปหาตพฺพ\-เหตุกา เจตนา โมหสฺส จิตฺตสมุฏฺฐานานญฺจ รูปานํ กมฺมปจฺจเยน ปจฺจโย.  นานากฺขณิกา ทสฺสเนน ปหาตพฺพ\-เหตุกา เจตนา วิปากานํ ขนฺธานํ กฏตฺตา จ รูปานํ กมฺมปจฺจเยน ปจฺจโย. (2)}

\prose{ทสฺสเนน ปหาตพฺพ\-เหตุโก ธมฺโม ทสฺสเนน ปหาตพฺพ\-เหตุกสฺส จ เนวทสฺสเนน นภาวนาย ปหาตพฺพ\-เหตุกสฺส จ ธมฺมสฺส กมฺมปจฺจเยน ปจฺจโย. ทสฺสเนน ปหาตพฺพ\-เหตุกา เจตนา สมฺปยุตฺตกานํ ขนฺธานํ โมหสฺส จ จิตฺตสมุฏฺฐานานญฺจ รูปานํ กมฺมปจฺจเยน ปจฺจโย. (3)}

\proseitem{97}{ภาวนาย ปหาตพฺพ\-เหตุโก ธมฺโม ภาวนาย ปหาตพฺพ\-เหตุกสฺส ธมฺมสฺส กมฺมปจฺจเยน ปจฺจโย. ภาวนาย ปหาตพฺพ\-เหตุกา เจตนา สมฺปยุตฺตกานํ ขนฺธานํ กมฺมปจฺจเยน ปจฺจโย. (1)}

\prose{ภาวนาย ปหาตพฺพ\-เหตุโก ธมฺโม เนวทสฺสเนน นภาวนาย ปหาตพฺพ\-เหตุกสฺส ธมฺมสฺส กมฺมปจฺจเยน ปจฺจโย. ภาวนาย ปหาตพฺพ\-เหตุกา เจตนา โมหสฺส จิตฺตสมุฏฺฐานานญฺจ รูปานํ กมฺมปจฺจเยน ปจฺจโย. (2)}

\prose{\csromanfolio{229}ภาวนาย ปหาตพฺพ\-เหตุโก ธมฺโม ภาวนาย ปหาตพฺพ\-เหตุกสฺส จ เนวทสฺสเนน นภาวนาย ปหาตพฺพ\-เหตุกสฺส จ ธมฺมสฺส กมฺมปจฺจเยน ปจฺจโย. ภาวนาย ปหาตพฺพ\-เหตุกา เจตนา สมฺปยุตฺตกานํ ขนฺธานํ โมหสฺส จ จิตฺตสมุฏฺฐานานญฺจ รูปานํ กมฺมปจฺจเยน ปจฺจโย. (3)}

\proseitem{98}{เนวทสฺสเนน นภาวนาย ปหาตพฺพ\-เหตุโก ธมฺโม เนวทสฺสเนน นภาวนาย ปหาตพฺพ\-เหตุกสฺส ธมฺมสฺส กมฺมปจฺจเยน ปจฺจโย. สหชาตา, นานากฺขณิกา.  สหชาตา เนวทสฺสเนน นภาวนาย ปหาตพฺพ\-เหตุกา เจตนา สมฺปยุตฺตกานํ ขนฺธานํ จิตฺตสมุฏฺฐานานญฺจ รูปานํ กมฺมปจฺจเยน ปจฺจโย. ปฏิสนฺธิ\-กฺขเณ ฯเปฯ. นานากฺขณิกา เนวทสฺสเนน นภาวนาย ปหาตพฺพ\-เหตุกา เจตนา วิปากานํ ขนฺธานํ กฏตฺตา จ รูปานํ กมฺมปจฺจเยน ปจฺจโย. (1)}

\csromanheader{2}{\textbf{วิปาก}}
\proseitem{99}{เนวทสฺสเนน นภาวนาย ปหาตพฺพ\-เหตุโก ธมฺโม เนวทสฺสเนน นภาวนาย ปหาตพฺพ\-เหตุกสฺส ธมฺมสฺส วิปาก\-ปจฺจเยน ปจฺจโย. (ปวตฺติ\-ปฏิสนฺธิ) วิปากา ขนฺธา วตฺถุสฺส ฯเปฯ.}

\csromanheader{2}{\textbf{อาหาร}}
\proseitem{100}{ทสฺสเนน ปหาตพฺพ\-เหตุโก ธมฺโม ทสฺสเนน ปหาตพฺพ\-เหตุกสฺส ธมฺมสฺส อาหารปจฺจเยน ปจฺจโย. ทสฺสเนน ปหาตพฺพ\-เหตุกา อาหารา สมฺปยุตฺตกานํ ขนฺธานํ อาหารปจฺจเยน ปจฺจโย. (1)}

\prose{ทสฺสเนน ปหาตพฺพ\-เหตุโก ธมฺโม เนวทสฺสเนน นภาวนาย ปหาตพฺพ\-เหตุกสฺส ธมฺมสฺส อาหารปจฺจเยน ปจฺจโย. ทสฺสเนน ปหาตพฺพ\-เหตุกา อาหารา โมหสฺส จิตฺตสมุฏฺฐานานญฺจ รูปานํ อาหารปจฺจเยน ปจฺจโย. (2)}

\prose{ทสฺสเนน ปหาตพฺพ\-เหตุโก ธมฺโม ทสฺสเนน ปหาตพฺพ\-เหตุกสฺส จ เนวทสฺสเนน นภาวนาย ปหาตพฺพ\-เหตุกสฺส จ ธมฺมสฺส อาหารปจฺจเยน ปจฺจโย. ทสฺสเนน ปหาตพฺพ\-เหตุกา อาหารา สมฺปยุตฺตกานํ ขนฺธานํ โมหสฺส จ จิตฺตสมุฏฺฐานานญฺจ รูปานํ อาหารปจฺจเยน ปจฺจโย. (3)}

\prose{\csromanfolio{230}ภาวนาย ปหาตพฺพ\-เหตุโก ธมฺโม. ตีณิ. (ทสฺสเนน สทิสํ.)}

\prose{เนวทสฺสเนน นภาวนาย ปหาตพฺพ\-เหตุโก ธมฺโม เนวทสฺสเนน นภาวนาย ปหาตพฺพ\-เหตุกสฺส ธมฺมสฺส อาหารปจฺจเยน ปจฺจโย. เนวทสฺสเนน นภาวนาย ปหาตพฺพ\-เหตุกา อาหารา สมฺปยุตฺตกานํ ขนฺธานํ จิตฺตสมุฏฺฐานานญฺจ รูปานํ อาหารปจฺจเยน ปจฺจโย. ปฏิสนฺธิ\-กฺขเณ ฯเปฯ. กพฬีกาโร อาหาโร อิมสฺส กายสฺส อาหารปจฺจเยน ปจฺจโย.}

\csromanheader{2}{\textbf{อินฺทฺริยาทิ}}
\proseitem{101}{ทสฺสเนน ปหาตพฺพ\-เหตุโก ธมฺโม ทสฺสเนน ปหาตพฺพ\-เหตุกสฺส ธมฺมสฺส อินฺทฺริย\-ปจฺจเยน ปจฺจโย. ตีณิ. (อาหารสทิสํ. โมโห กาตพฺโพ.)}

\prose{ภาวนาย ปหาตพฺพ\-เหตุโก ธมฺโม. ตีณิ.}

\prose{เนวทสฺสเนน นภาวนาย ปหาตพฺพ\-เหตุโก ธมฺโม เนวทสฺสเนน นภาวนาย ปหาตพฺพ\-เหตุกสฺส ธมฺมสฺส อินฺทฺริย\-ปจฺจเยน ปจฺจโย. เนวทสฺสเนน นภาวนาย ปหาตพฺพ\-เหตุกา อินฺทฺริยา สมฺปยุตฺตกานํ ขนฺธานํ ฯเปฯ. จกฺขุนฺทฺริยํ จกฺขุ\-วิญฺญาณสฺส ฯเปฯ. กายินฺทฺริยํ กายวิญฺญาณสฺส ฯเปฯ. รูปชีวิตินฺทฺริยํ กฏตฺตารูปานํ อินฺทฺริย\-ปจฺจเยน ปจฺจโย. ฌานปจฺจเยน ปจฺจโย. มคฺคปจฺจเยน ปจฺจโย. (อิเม สเหตุกา กาตพฺพา.) สมฺปยุตฺต\-ปจฺจเยน ปจฺจโย. (ปฏิจฺจวาเร สมฺปยุตฺต\-วาร\-สทิสํ.)}

\csromanheader{2}{\textbf{วิปฺปยุตฺต}}
\proseitem{102}{ทสฺสเนน ปหาตพฺพ\-เหตุโก ธมฺโม เนวทสฺสเนน นภาวนาย ปหาตพฺพ\-เหตุกสฺส ธมฺมสฺส วิปฺปยุตฺต\-ปจฺจเยน ปจฺจโย. สหชาตํ, ปจฺฉาชาตํ. (ทสฺสน\-ตฺติก\-สทิสํ.)}

\prose{ภาวนาย ปหาตพฺพ\-เหตุโก ธมฺโม เนวทสฺสเนน นภาวนาย ปหาตพฺพ\-เหตุกสฺส ธมฺมสฺส วิปฺปยุตฺต\-ปจฺจเยน ปจฺจโย. สหชาตํ, ปจฺฉาชาตํ. (ทสฺสน\-ตฺติก\-สทิสํ.)}

\prose{เนวทสฺสเนน นภาวนาย ปหาตพฺพ\-เหตุโก ธมฺโม เนวทสฺสเนน นภาวนาย ปหาตพฺพ\-เหตุกสฺส ธมฺมสฺส วิปฺปยุตฺต\-ปจฺจเยน ปจฺจโย. สหชาตํ, ปุเรชาตํ, ปจฺฉาชาตํ. (ทสฺสน\-ตฺติก\-สทิสํ.) ปจฺฉาชาตา \csromanfolio{231}เนวทสฺสเนน นภาวนาย ปหาตพฺพ\-เหตุกา ขนฺธา จ โมโห จ ปุเรชาตสฺส อิมสฺส กายสฺส ฯเปฯ. (1)}

\prose{เนวทสฺสเนน นภาวนาย ปหาตพฺพ\-เหตุโก ธมฺโม ทสฺสเนน ปหาตพฺพ\-เหตุกสฺส ธมฺมสฺส วิปฺปยุตฺต\-ปจฺจเยน ปจฺจโย. ปุเรชาตํ วตฺถุ ทสฺสเนน ปหาตพฺพ\-เหตุกานํ ขนฺธานํ ฯเปฯ. (2)}

\prose{เนวทสฺสเนน นภาวนาย ปหาตพฺพ\-เหตุโก ธมฺโม ภาวนาย ปหาตพฺพ\-เหตุกสฺส ธมฺมสฺส วิปฺปยุตฺต\-ปจฺจเยน ปจฺจโย. ปุเรชาตํ วตฺถุ ภาวนาย ปหาตพฺพ\-เหตุกานํ ขนฺธานํ ฯเปฯ. (3)}

\prose{เนวทสฺสเนน นภาวนาย ปหาตพฺพ\-เหตุโก ธมฺโม ทสฺสเนน ปหาตพฺพ\-เหตุกสฺส จ เนวทสฺสเนน นภาวนาย ปหาตพฺพ\-เหตุกสฺส จ ธมฺมสฺส วิปฺปยุตฺต\-ปจฺจเยน ปจฺจโย. ปุเรชาตํ วตฺถุ วิจิกิจฺฉา\-สหคตานํ ขนฺธานํ โมหสฺส จ วิปฺปยุตฺต\-ปจฺจเยน ปจฺจโย. (4)}

\prose{เนวทสฺสเนน นภาวนาย ปหาตพฺพ\-เหตุโก ธมฺโม ภาวนาย ปหาตพฺพ\-เหตุกสฺส จ เนวทสฺสเนน นภาวนาย ปหาตพฺพ\-เหตุกสฺส จ ธมฺมสฺส วิปฺปยุตฺต\-ปจฺจเยน ปจฺจโย. ปุเรชาตํ วตฺถุ อุทฺธจฺจ\-สหคตานํ ขนฺธานํ โมหสฺส จ วิปฺปยุตฺต\-ปจฺจเยน ปจฺจโย. (5)}

\proseitem{103}{ทสฺสเนน ปหาตพฺพ\-เหตุโก จ เนวทสฺสเนน นภาวนาย ปหาตพฺพ\-เหตุโก จ ธมฺมา เนวทสฺสเนน นภาวนาย ปหาตพฺพ\-เหตุกสฺส ธมฺมสฺส วิปฺปยุตฺต\-ปจฺจเยน ปจฺจโย. สหชาตํ, ปจฺฉาชาตํ.  สหชาตา วิจิกิจฺฉา\-สหคตา ขนฺธา จ โมโห จ จิตฺต\-สมุฏฺฐานานํ รูปานํ วิปฺปยุตฺต\-ปจฺจเยน ปจฺจโย.  ปจฺฉาชาตา วิจิกิจฺฉา\-สหคตา ขนฺธา จ โมโห จ ปุเรชาตสฺส อิมสฺส กายสฺส วิปฺปยุตฺต\-ปจฺจเยน ปจฺจโย. (1)}

\prose{ภาวนาย ปหาตพฺพ\-เหตุโก จ เนวทสฺสเนน นภาวนาย ปหาตพฺพ\-เหตุโก จ ธมฺมา เนวทสฺสเนน นภาวนาย ปหาตพฺพ\-เหตุกสฺส ธมฺมสฺส วิปฺปยุตฺต\-ปจฺจเยน ปจฺจโย. สหชาตํ, ปจฺฉาชาตํ.  สหชาตา อุทฺธจฺจ\-สหคตา ขนฺธา จ โมโห จ จิตฺต\-สมุฏฺฐานานํ รูปานํ ฯเปฯ. ปจฺฉาชาตา อุทฺธจฺจ\-สหคตา ขนฺธา จ โมโห จ ปุเรชาตสฺส อิมสฺส กายสฺส วิปฺปยุตฺต\-ปจฺจเยน ปจฺจโย. (1)}

\prose{\csromanfolio{232}อตฺถิอาทิ}

\proseitem{104}{ทสฺสเนน ปหาตพฺพ\-เหตุโก ธมฺโม ทสฺสเนน ปหาตพฺพ\-เหตุกสฺส ธมฺมสฺส อตฺถิปจฺจเยน ปจฺจโย. ทสฺสเนน ปหาตพฺพ\-เหตุโก เอโก ขนฺโธ ติณฺณนฺนํ ขนฺธานํ ฯเปฯ. (1)}

\prose{ทสฺสเนน ปหาตพฺพ\-เหตุโก ธมฺโม เนวทสฺสเนน นภาวนาย ปหาตพฺพ\-เหตุกสฺส ธมฺมสฺส อตฺถิปจฺจเยน ปจฺจโย. สหชาตํ, ปจฺฉาชาตํ.  สหชาตา ทสฺสเนน ปหาตพฺพ\-เหตุกา ขนฺธา จิตฺต\-สมุฏฺฐานานํ รูปานํ อตฺถิปจฺจเยน ปจฺจโย. สหชาตา วิจิกิจฺฉา\-สหคตา ขนฺธา โมหสฺส จิตฺตสมุฏฺฐานานญฺจ รูปานํ อตฺถิปจฺจเยน ปจฺจโย.  ปจฺฉาชาตา ทสฺสเนน ปหาตพฺพ\-เหตุกา ขนฺธา ปุเรชาตสฺส อิมสฺส กายสฺส อตฺถิปจฺจเยน ปจฺจโย. (2)}

\prose{ทสฺสเนน ปหาตพฺพ\-เหตุโก ธมฺโม ทสฺสเนน ปหาตพฺพ\-เหตุกสฺส จ เนวทสฺสเนน นภาวนาย ปหาตพฺพ\-เหตุกสฺส จ ธมฺมสฺส อตฺถิปจฺจเยน ปจฺจโย. ทสฺสเนน ปหาตพฺพ\-เหตุโก เอโก ขนฺโธ ติณฺณนฺนํ ขนฺธานํ จิตฺตสมุฏฺฐานานญฺจ รูปานํ อตฺถิปจฺจเยน ปจฺจโย. วิจิกิจฺฉา\-สหคโต เอโก ขนฺโธ ติณฺณนฺนํ ขนฺธานํ โมหสฺส จ จิตฺตสมุฏฺฐานานญฺจ รูปานํ อตฺถิปจฺจเยน ปจฺจโย ฯเปฯ. (3)}

\prose{ภาวนาย ปหาตพฺพ\-เหตุโก ธมฺโม. ตีณิ.}

\proseitem{105}{เนวทสฺสเนน นภาวนาย ปหาตพฺพ\-เหตุโก ธมฺโม เนวทสฺสเนน นภาวนาย ปหาตพฺพ\-เหตุกสฺส ธมฺมสฺส อตฺถิปจฺจเยน ปจฺจโย. สหชาตํ, ปุเรชาตํ, ปจฺฉาชาตํ, อาหารํ, อินฺทฺริยํ.  สหชาโต เนวทสฺสเนน นภาวนาย ปหาตพฺพ\-เหตุโก เอโก ขนฺโธ ติณฺณนฺนํ ขนฺธานํ จิตฺตสมุฏฺฐานานญฺจ รูปานํ อตฺถิปจฺจเยน ปจฺจโย ฯเปฯ. วิจิกิจฺฉา\-สหคโต อุทฺธจฺจ\-สหคโต โมโห จิตฺต\-สมุฏฺฐานานํ รูปานํ อตฺถิปจฺจเยน ปจฺจโย. ปฏิสนฺธิ\-กฺขเณ ฯเปฯ. อสญฺญสตฺตานํ ฯเปฯ. ปุเรชาตํ จกฺขุํ ฯเปฯ. วตฺถุํ อนิจฺจโต ฯเปฯ. ทิพฺเพน จกฺขุนา รูปํ ปสฺสติ. ทิพฺพาย โสตธาตุยา สทฺทํ สุณาติ. รูปายตนํ จกฺขุ\-วิญฺญาณสฺส ฯเปฯ. โผฏฺฐพฺพา\-ยตนํ กายวิญฺญาณสฺส ฯเปฯ. จกฺขายตนํ ฯเปฯ. กายายตนํ ฯเปฯ. วตฺถุ เนวทสฺสเนน นภาวนาย ปหาตพฺพ\-เหตุกานํ ขนฺธานํ โมหสฺส จ อตฺถิปจฺจเยน ปจฺจโย.  ปจฺฉาชาตา เนวทสฺสเนน นภาวนาย \csromanfolio{233}ปหาตพฺพ\-เหตุกา ขนฺธา จ โมโห จ ปุเรชาตสฺส อิมสฺส กายสฺส อตฺถิปจฺจเยน ปจฺจโย. กพฬีกาโร อาหาโร อิมสฺส กายสฺส ฯเปฯ. รูปชีวิตินฺทฺริยํ กฏตฺตารูปานํ ฯเปฯ. (1)}

\prose{เนวทสฺสเนน นภาวนาย ปหาตพฺพ\-เหตุโก ธมฺโม ทสฺสเนน ปหาตพฺพ\-เหตุกสฺส ธมฺมสฺส อตฺถิปจฺจเยน ปจฺจโย. สหชาตํ, ปุเรชาตํ.  สหชาโต วิจิกิจฺฉา\-สหคโต โมโห สมฺปยุตฺตกานํ ขนฺธานํ อตฺถิปจฺจเยน ปจฺจโย.  ปุเรชาตํ จกฺขุํ ฯเปฯ วตฺถุํ อสฺสาเทติ อภินนฺทติ, ตํ อารพฺภ ทสฺสเนน ปหาตพฺพ\-เหตุโก ราโค อุปฺปชฺชติ, ทิฏฺฐิ อุปฺปชฺชติ. วิจิกิจฺฉา อุปฺปชฺชติ. ทสฺสเนน ปหาตพฺพ\-เหตุกํ โทมนสฺสํ อุปฺปชฺชติ. วตฺถุ ทสฺสเนน ปหาตพฺพ\-เหตุกานํ ขนฺธานํ อตฺถิปจฺจเยน ปจฺจโย. (2)}

\prose{เนวทสฺสเนน นภาวนาย ปหาตพฺพ\-เหตุโก ธมฺโม ภาวนาย ปหาตพฺพ\-เหตุกสฺส ธมฺมสฺส อตฺถิปจฺจเยน ปจฺจโย. สหชาตํ, ปุเรชาตํ.  สหชาโต อุทฺธจฺจ\-สหคโต โมโห สมฺปยุตฺตกานํ ขนฺธานํ อตฺถิปจฺจเยน ปจฺจโย.  ปุเรชาตํ จกฺขุํ ฯเปฯ วตฺถุํ อสฺสาเทติ อภินนฺทติ ฯเปฯ. วตฺถุ ภาวนาย ปหาตพฺพ\-เหตุกานํ ขนฺธานํ อตฺถิปจฺจเยน ปจฺจโย. (3)}

\prose{เนวทสฺสเนน นภาวนาย ปหาตพฺพ\-เหตุโก ธมฺโม ทสฺสเนน ปหาตพฺพ\-เหตุกสฺส จ เนวทสฺสเนน นภาวนาย ปหาตพฺพ\-เหตุกสฺส จ ธมฺมสฺส อตฺถิปจฺจเยน ปจฺจโย. สหชาตํ, ปุเรชาตํ.  สหชาโต วิจิกิจฺฉา\-สหคโต โมโห สมฺปยุตฺตกานํ ขนฺธานํ จิตฺตสมุฏฺฐานานญฺจ รูปานํ อตฺถิปจฺจเยน ปจฺจโย.  ปุเรชาตํ จกฺขุํ อารพฺภ วิจิกิจฺฉา\-สหคตา ขนฺธา จ โมโห จ อุปฺปชฺชนฺติ ฯเปฯ วตฺถุํ อารพฺภ ฯเปฯ. วตฺถุ วิจิกิจฺฉา\-สหคตานํ ขนฺธานํ โมหสฺส จ อตฺถิปจฺจเยน ปจฺจโย. (4)}

\prose{เนวทสฺสเนน นภาวนาย ปหาตพฺพ\-เหตุโก ธมฺโม ภาวนาย ปหาตพฺพ\-เหตุกสฺส จ เนวทสฺสเนน นภาวนาย ปหาตพฺพ\-เหตุกสฺส จ ธมฺมสฺส อตฺถิปจฺจเยน ปจฺจโย. สหชาตํ, ปุเรชาตํ.  สหชาโต อุทฺธจฺจ\-สหคโต โมโห สมฺปยุตฺตกานํ ขนฺธานํ จิตฺตสมุฏฺฐานานญฺจ รูปานํ อตฺถิปจฺจเยน ปจฺจโย.  ปุเรชาตํ จกฺขุํ อารพฺภ \csromanfolio{234}อุทฺธจฺจ\-สหคตา ขนฺธา จ โมโห จ อุปฺปชฺชนฺติ ฯเปฯ วตฺถุํ อารพฺภ ฯเปฯ. วตฺถุ อุทฺธจฺจ\-สหคตานํ ขนฺธานํ โมหสฺส จ อตฺถิปจฺจเยน ปจฺจโย. (5)}

\proseitem{106}{ทสฺสเนน ปหาตพฺพ\-เหตุโก จ เนวทสฺสเนน นภาวนาย ปหาตพฺพ\-เหตุโก จ ธมฺมา ทสฺสเนน ปหาตพฺพ\-เหตุกสฺส ธมฺมสฺส อตฺถิปจฺจเยน ปจฺจโย. สหชาตํ, ปุเรชาตํ.  สหชาโต ทสฺสเนน ปหาตพฺพ\-เหตุโก เอโก ขนฺโธ จ วตฺถุ จ ติณฺณนฺนํ ขนฺธานํ อตฺถิปจฺจเยน ปจฺจโย ฯเปฯ เทฺว ขนฺธา ฯเปฯ. วิจิกิจฺฉา\-สหคโต เอโก ขนฺโธ จ โมโห จ ติณฺณนฺนํ ขนฺธานํ อตฺถิปจฺจเยน ปจฺจโย ฯเปฯ เทฺว ขนฺธา ฯเปฯ. (1)}

\prose{ทสฺสเนน ปหาตพฺพ\-เหตุโก จ เนวทสฺสเนน นภาวนาย ปหาตพฺพ\-เหตุโก จ ธมฺมา เนวทสฺสเนน นภาวนาย ปหาตพฺพ\-เหตุกสฺส ธมฺมสฺส อตฺถิปจฺจเยน ปจฺจโย. สหชาตํ, ปุเรชาตํ, ปจฺฉาชาตํ, อาหารํ, อินฺทฺริยํ.  สหชาตา ทสฺสเนน ปหาตพฺพ\-เหตุกา ขนฺธา จ มหาภูตา จ จิตฺต\-สมุฏฺฐานานํ รูปานํ อตฺถิปจฺจเยน ปจฺจโย. สหชาตา วิจิกิจฺฉา\-สหคตา ขนฺธา จ โมโห จ จิตฺต\-สมุฏฺฐานานํ รูปานํ อตฺถิปจฺจเยน ปจฺจโย. สหชาตา วิจิกิจฺฉา\-สหคตา ขนฺธา จ วตฺถุ จ โมหสฺส อตฺถิปจฺจเยน ปจฺจโย.  ปจฺฉาชาตา วิจิกิจฺฉา\-สหคตา ขนฺธา จ โมโห จ ปุเรชาตสฺส อิมสฺส กายสฺส อตฺถิปจฺจเยน ปจฺจโย. ปจฺฉาชาตา ทสฺสเนน ปหาตพฺพ\-เหตุกา ขนฺธา จ กพฬีกาโร อาหาโร จ อิมสฺส กายสฺส อตฺถิปจฺจเยน ปจฺจโย. ปจฺฉาชาตา ทสฺสเนน ปหาตพฺพ\-เหตุกา ขนฺธา จ รูปชีวิตินฺทฺริยญฺจ กฏตฺตารูปานํ อตฺถิปจฺจเยน ปจฺจโย. (2)}

\prose{ทสฺสเนน ปหาตพฺพ\-เหตุโก จ เนวทสฺสเนน นภาวนาย ปหาตพฺพ\-เหตุโก จ ธมฺมา ทสฺสเนน ปหาตพฺพ\-เหตุกสฺส จ เนวทสฺสเนน นภาวนาย ปหาตพฺพ\-เหตุกสฺส จ ธมฺมสฺส อตฺถิปจฺจเยน ปจฺจโย. สหชาตํ, ปุเรชาตํ.  สหชาโต วิจิกิจฺฉา\-สหคโต เอโก ขนฺโธ จ วตฺถุ จ ติณฺณนฺนํ ขนฺธานํ โมหสฺส จ อตฺถิปจฺจเยน ปจฺจโย ฯเปฯ เทฺว ขนฺธา ฯเปฯ. วิจิกิจฺฉา\-สหคโต เอโก ขนฺโธ จ โมโห จ ติณฺณนฺนํ ขนฺธานํ จิตฺตสมุฏฺฐานานญฺจ รูปานํ อตฺถิปจฺจเยน ปจฺจโย ฯเปฯ เทฺว ขนฺธา จ โมโห จ ฯเปฯ. (3)}

\prose{\csromanfolio{235}ภาวนาย ปหาตพฺพ\-เหตุโก จ เนวทสฺสเนน นภาวนาย ปหาตพฺพ\-เหตุโก จ ธมฺมา ภาวนาย ปหาตพฺพ\-เหตุกสฺส ธมฺมสฺส อตฺถิปจฺจเยน ปจฺจโย. (สํขิตฺตํ. ติสฺโส ปญฺหา, ทสฺสเนน นเยน วิภชิตพฺพา, “อุทฺธจฺจนฺ”ติ นิยาเมตพฺพํ.) นตฺถิ\-ปจฺจเยน ปจฺจโย. วิคต\-ปจฺจเยน ปจฺจโย. อวิคต\-ปจฺจเยน ปจฺจโย.}

\csromanmarkh{1. ปจฺจยานุโลม}\csromanmarkhii{2. สงฺขฺยาวาร}\csromanheadernomark{2}{1. \textbf{ปจฺจยานุโลม 2. สงฺขฺยาวาร}}
\proseitem{107}{เหตุยา เอกาทส, อารมฺมเณ เอกวีส, อธิปติยา ทส, อนนฺตเร สตฺตรส, สมนนฺตเร สตฺตรส, สหชาเต สตฺตรส, อญฺญมญฺเญ เอกาทส, นิสฺสเย สตฺตรส, อุปนิสฺสเย เอกวีส, ปุเรชาเต ปญฺจ, ปจฺฉาชาเต ปญฺจ, อาเสวเน สตฺตรส, กมฺเม สตฺต, วิปาเก เอกํ, อาหาเร สตฺต, อินฺทฺริเย สตฺต, ฌาเน สตฺต, มคฺเค สตฺต, สมฺปยุตฺเต เอกาทส, วิปฺปยุตฺเต นว, อตฺถิยา สตฺตรส, นตฺถิยา สตฺตรส, วิคเต สตฺตรส, อวิคเต สตฺตรส. (เอวํ คเณตพฺพํ.) อนุโลมํ.}
\csromansectionrule

\csromanheader{2}{2. ปจฺจนียุ\-ทฺธาร}
\proseitem{108}{ทสฺสเนน ปหาตพฺพ\-เหตุโก ธมฺโม ทสฺสเนน ปหาตพฺพ\-เหตุกสฺส ธมฺมสฺส อารมฺมณ\-ปจฺจเยน ปจฺจโย. สหชาต\-ปจฺจเยน ปจฺจโย. อุปนิสฺสย\-ปจฺจเยน ปจฺจโย. (1)}

\prose{ทสฺสเนน ปหาตพฺพ\-เหตุโก ธมฺโม เนวทสฺสเนน นภาวนาย ปหาตพฺพ\-เหตุกสฺส ธมฺมสฺส อารมฺมณ\-ปจฺจเยน ปจฺจโย. สหชาต\-ปจฺจเยน ปจฺจโย. อุปนิสฺสย\-ปจฺจเยน ปจฺจโย. ปจฺฉาชาต\-ปจฺจเยน ปจฺจโย. กมฺมปจฺจเยน ปจฺจโย. (2)}

\prose{ทสฺสเนน ปหาตพฺพ\-เหตุโก ธมฺโม ทสฺสเนน ปหาตพฺพ\-เหตุกสฺส จ เนวทสฺสเนน นภาวนาย ปหาตพฺพ\-เหตุกสฺส จ ธมฺมสฺส อารมฺมณ\-ปจฺจเยน ปจฺจโย. สหชาต\-ปจฺจเยน ปจฺจโย. อุปนิสฺสย\-ปจฺจเยน ปจฺจโย. (3)}

\proseitem{109}{\csromanfolio{236}ภาวนาย ปหาตพฺพ\-เหตุโก ธมฺโม ภาวนาย ปหาตพฺพ\-เหตุกสฺส ธมฺมสฺส อารมฺมณ\-ปจฺจเยน ปจฺจโย. สหชาต\-ปจฺจเยน ปจฺจโย. อุปนิสฺสย\-ปจฺจเยน ปจฺจโย. (1)}

\prose{ภาวนาย ปหาตพฺพ\-เหตุโก ธมฺโม ทสฺสเนน ปหาตพฺพ\-เหตุกสฺส ธมฺมสฺส อารมฺมณ\-ปจฺจเยน ปจฺจโย. อุปนิสฺสย\-ปจฺจเยน ปจฺจโย. (2)}

\prose{ภาวนาย ปหาตพฺพ\-เหตุโก ธมฺโม เนวทสฺสเนน นภาวนาย ปหาตพฺพ\-เหตุกสฺส ธมฺมสฺส อารมฺมณ\-ปจฺจเยน ปจฺจโย. สหชาต\-ปจฺจเยน ปจฺจโย. อุปนิสฺสย\-ปจฺจเยน ปจฺจโย. ปจฺฉาชาต\-ปจฺจเยน ปจฺจโย. (3)}

\prose{ภาวนาย ปหาตพฺพ\-เหตุโก ธมฺโม ทสฺสเนน ปหาตพฺพ\-เหตุกสฺส จ เนวทสฺสเนน นภาวนาย ปหาตพฺพ\-เหตุกสฺส จ ธมฺมสฺส อารมฺมณ\-ปจฺจเยน ปจฺจโย. อุปนิสฺสย\-ปจฺจเยน ปจฺจโย. (4)}

\prose{ภาวนาย ปหาตพฺพ\-เหตุโก ธมฺโม ภาวนาย ปหาตพฺพ\-เหตุกสฺส จ เนวทสฺสเนน นภาวนาย ปหาตพฺพ\-เหตุกสฺส จ ธมฺมสฺส อารมฺมณ\-ปจฺจเยน ปจฺจโย. สหชาต\-ปจฺจเยน ปจฺจโย. อุปนิสฺสย\-ปจฺจเยน ปจฺจโย. (5)}

\proseitem{110}{เนวทสฺสเนน นภาวนาย ปหาตพฺพ\-เหตุโก ธมฺโม เนวทสฺสเนน นภาวนาย ปหาตพฺพ\-เหตุกสฺส ธมฺมสฺส อารมฺมณ\-ปจฺจเยน ปจฺจโย. สหชาต\-ปจฺจเยน ปจฺจโย. อุปนิสฺสย\-ปจฺจเยน ปจฺจโย. ปุเรชาต\-ปจฺจเยน ปจฺจโย. ปจฺฉาชาต\-ปจฺจเยน ปจฺจโย. กมฺมปจฺจเยน ปจฺจโย. อาหารปจฺจเยน ปจฺจโย. อินฺทฺริย\-ปจฺจเยน ปจฺจโย. (1)}

\prose{เนวทสฺสเนน นภาวนาย ปหาตพฺพ\-เหตุโก ธมฺโม ทสฺสเนน ปหาตพฺพ\-เหตุกสฺส ธมฺมสฺส อารมฺมณ\-ปจฺจเยน ปจฺจโย. สหชาต\-ปจฺจเยน ปจฺจโย. อุปนิสฺสย\-ปจฺจเยน ปจฺจโย. ปุเรชาต\-ปจฺจเยน ปจฺจโย. (2)}

\prose{\csromanfolio{237}เนวทสฺสเนน นภาวนาย ปหาตพฺพ\-เหตุโก ธมฺโม ภาวนาย ปหาตพฺพ\-เหตุกสฺส ธมฺมสฺส อารมฺมณ\-ปจฺจเยน ปจฺจโย. สหชาต\-ปจฺจเยน ปจฺจโย. อุปนิสฺสย\-ปจฺจเยน ปจฺจโย. ปุเรชาต\-ปจฺจเยน ปจฺจโย. (3)}

\prose{เนวทสฺสเนน นภาวนาย ปหาตพฺพ\-เหตุโก ธมฺโม ทสฺสเนน ปหาตพฺพ\-เหตุกสฺส จ เนวทสฺสเนน นภาวนาย ปหาตพฺพ\-เหตุกสฺส จ ธมฺมสฺส อารมฺมณ\-ปจฺจเยน ปจฺจโย. สหชาต\-ปจฺจเยน ปจฺจโย. อุปนิสฺสย\-ปจฺจเยน ปจฺจโย. ปุเรชาต\-ปจฺจเยน ปจฺจโย. (4)}

\prose{เนวทสฺสเนน นภาวนาย ปหาตพฺพ\-เหตุโก ธมฺโม ภาวนาย ปหาตพฺพ\-เหตุกสฺส จ เนวทสฺสเนน นภาวนาย ปหาตพฺพ\-เหตุกสฺส จ ธมฺมสฺส อารมฺมณ\-ปจฺจเยน ปจฺจโย. สหชาต\-ปจฺจเยน ปจฺจโย. อุปนิสฺสย\-ปจฺจเยน ปจฺจโย. ปุเรชาต\-ปจฺจเยน ปจฺจโย. (5)}

\proseitem{111}{ทสฺสเนน ปหาตพฺพ\-เหตุโก จ เนวทสฺสเนน นภาวนาย ปหาตพฺพ\-เหตุโก จ ธมฺมา ทสฺสเนน ปหาตพฺพ\-เหตุกสฺส ธมฺมสฺส อารมฺมณ\-ปจฺจเยน ปจฺจโย. สหชาต\-ปจฺจเยน ปจฺจโย. อุปนิสฺสย\-ปจฺจเยน ปจฺจโย. (1)}

\prose{(อิธ สหชาตํ ปุเรชาตํ มิสฺสคตํ อตฺถิ, ปาฬิยํ กาตพฺพํ. คณนาย อุปธาเรตฺวา คเณตพฺพํ.)}

\prose{ทสฺสเนน ปหาตพฺพ\-เหตุโก จ เนวทสฺสเนน นภาวนาย ปหาตพฺพ\-เหตุโก จ ธมฺมา เนวทสฺสเนน นภาวนาย ปหาตพฺพ\-เหตุกสฺส ธมฺมสฺส สหชาตํ, ปุเรชาตํ, ปจฺฉาชาตํ, อาหารํ, อินฺทฺริยํ. (2)}

\prose{(อิธาปิ อารมฺมณ\-ปจฺจยา อุปนิสฺสย\-ปจฺจยา อตฺถิ, ปาฬิยํ นตฺถิ. คเณนฺเตน อุปธาเรตฺวา คเณตพฺพํ.)}

\prose{ทสฺสเนน ปหาตพฺพ\-เหตุโก จ เนวทสฺสเนน นภาวนาย ปหาตพฺพ\-เหตุโก จ ธมฺมา ทสฺสเนน ปหาตพฺพ\-เหตุกสฺส จ เนวทสฺสเนน นภาวนาย ปหาตพฺพ\-เหตุกสฺส จ ธมฺมสฺส อารมฺมณ\-ปจฺจเยน ปจฺจโย. สหชาต\-ปจฺจเยน ปจฺจโย. อุปนิสฺสย\-ปจฺจเยน ปจฺจโย. (3)}

\prose{(อิธาปิ “สหชาตํ ปุเรชาตํ” ยํ มิสฺสกปญฺหา อตฺถิ, ปาฬิยํ กาตพฺพํ.)}

\proseitem{112}{\csromanfolio{238}ภาวนาย ปหาตพฺพ\-เหตุโก จ เนวทสฺสเนน นภาวนาย ปหาตพฺพ\-เหตุโก จ ธมฺมา ทสฺสเนน ปหาตพฺพ\-เหตุกสฺส ธมฺมสฺส อารมฺมณ\-ปจฺจเยน ปจฺจโย. อุปนิสฺสย\-ปจฺจเยน ปจฺจโย. (1)}

\prose{ภาวนาย ปหาตพฺพ\-เหตุโก จ เนวทสฺสเนน นภาวนาย ปหาตพฺพ\-เหตุโก จ ธมฺมา ภาวนาย ปหาตพฺพ\-เหตุกสฺส ธมฺมสฺส อารมฺมณ\-ปจฺจเยน ปจฺจโย. สหชาต\-ปจฺจเยน ปจฺจโย. อุปนิสฺสย\-ปจฺจเยน ปจฺจโย. (อิธาปิ “สหชาตํ ปุเรชาตํ” ยํ มิสฺสกปญฺหา อตฺถิ.) (2)}

\prose{ภาวนาย ปหาตพฺพ\-เหตุโก จ เนวทสฺสเนน นภาวนาย ปหาตพฺพ\-เหตุโก จ ธมฺมา เนวทสฺสเนน นภาวนาย ปหาตพฺพ\-เหตุกสฺส ธมฺมสฺส สหชาตํ, ปุเรชาตํ, ปจฺฉาชาตํ, อาหารํ, อินฺทฺริยํ. (อิธาปิ อารมฺมณอุปนิสฺสยา อตฺถิ.) (3)}

\prose{ภาวนาย ปหาตพฺพ\-เหตุโก จ เนวทสฺสเนน นภาวนาย ปหาตพฺพ\-เหตุโก จ ธมฺมา ทสฺสเนน ปหาตพฺพ\-เหตุกสฺส จ เนวทสฺสเนน นภาวนาย ปหาตพฺพ\-เหตุกสฺส จ ธมฺมสฺส อารมฺมณ\-ปจฺจเยน ปจฺจโย. อุปนิสฺสย\-ปจฺจเยน ปจฺจโย. (4)}

\prose{ภาวนาย ปหาตพฺพ\-เหตุโก จ เนวทสฺสเนน นภาวนาย ปหาตพฺพ\-เหตุโก จ ธมฺมา ภาวนาย ปหาตพฺพ\-เหตุกสฺส จ เนวทสฺสเนน นภาวนาย ปหาตพฺพ\-เหตุกสฺส จ ธมฺมสฺส อารมฺมณ\-ปจฺจเยน ปจฺจโย. สหชาต\-ปจฺจเยน ปจฺจโย. อุปนิสฺสย\-ปจฺจเยน ปจฺจโย. (5)}

\prose{(อิธาปิ สหชาตํ ปุเรชาตํ อตฺถิ. เย เต ปญฺหา น ลิขิตา, เต ปาฬิยํ คเณนฺตานํ พฺยญฺชเนน น สเมนฺติ. เต ปาฬิยํ น ลิขิตา, คณนา ปากฏา โหนฺติ. ยทิ สํสโย อุปฺปชฺชติ, อนุโลเม อตฺถิปจฺจเย เปกฺขิตพฺพํ.)}

\csromanmarkh{2. ปจฺจย\-ปจฺจนีย}\csromanmarkhii{2. สงฺขฺยาวาร}\csromanheadernomark{2}{2. ปจฺจย\-ปจฺจนีย 2. สงฺขฺยาวาร}
\proseitemwithcloserruled{113}{นเหตุยา เอกวีส, นอารมฺมเณ, นอธิปติยา, นอนนฺตเร, นสมนนฺตเร, นสหชาเต, นอญฺญมญฺเญ, นนิสฺสเย, นอุปนิสฺสเย, \csromanfolio{239}นปุเรชาเต, นปจฺฉาชาเต, นอาเสวเน, นกมฺเม, นวิปาเก, นอาหาเร, ไนนฺทฺริเย, นฌาเน, นมคฺเค, นสมฺปยุตฺเต, นวิปฺปยุตฺเต, โนอตฺถิยา, โนนตฺถิยา, โนวิคเต, โนอวิคเต, สพฺพตฺถ เอกวีส. (เอวํ คเณตพฺพํ.)}{ปจฺจนียํ.}
\proseitemwithcloserruled{114}{เหตุปจฺจยา นอารมฺมเณ เอกาทส, นอธิปติยา, นอนนฺตเร, นสมนนฺตเร เอกาทส, นอญฺญมญฺเญ ตีณิ, นอุปนิสฺสเย, นปุเรชาเต, นปจฺฉาชาเต, นอาเสวเน, นกมฺเม, นวิปาเก, นอาหาเร, นอินฺทฺริเย, นฌาเน, นมคฺเค เอกาทส, นสมฺปยุตฺเต ตีณิ, นวิปฺปยุตฺเต ปญฺจ, โนนตฺถิยา เอกาทส, โนวิคเต เอกาทส. (เอวํ คเณตพฺพํ.)}{อนุโลม\-ปจฺจนียํ.}
\proseitemwithcloser{115}{นเหตุปจฺจยา อารมฺมเณ เอกวีส, อธิปติยา ทส, อนนฺตเร สตฺตรส, สมนนฺตเร สตฺตรส, สหชาเต สตฺตรส, อญฺญมญฺเญ เอกาทส, นิสฺสเย สตฺตรส, อุปนิสฺสเย เอกวีส, ปุเรชาเต ปญฺจ, ปจฺฉาชาเต ปญฺจ, อาเสวเน สตฺตรส, กมฺเม สตฺต, วิปาเก เอกํ, อาหาเร สตฺต, อินฺทฺริเย สตฺต, ฌาเน สตฺต, มคฺเค สตฺต, สมฺปยุตฺเต เอกาทส, วิปฺปยุตฺเต นว, อตฺถิยา สตฺตรส, นตฺถิยา สตฺตรส, วิคเต สตฺตรส, อวิคเต สตฺตรส. (เอวํ คเณตพฺพํ.) ปจฺจนียา\-นุโลมํ. ปญฺหาวาโร.}{ทสฺสเนน\-ปหาตพฺพ\-เหตุ\-กตฺติกํ นิฏฺฐิตํ.}
\csromanmarkh{10. อาจยคามิตฺติก}\csromanmarkhii{1. ปฏิจฺจวาร}\csromanheadernomark{2}{10. \textbf{อาจยคามิตฺติก 1. ปฏิจฺจวาร}}
\csromanmarkh{1. ปจฺจยานุโลม}\csromanmarkhii{1. วิภงฺควาร}\csromanheadernomark{2}{1. \textbf{ปจฺจยานุโลม 1. วิภงฺควาร}}
\csromanheader{2}{\textbf{เหตุ}}
\proseitem{1}{\csromanfolio{240}อาจยคามิํ ธมฺมํ ปฏิจฺจ อาจยคามี ธมฺโม อุปฺปชฺชติ เหตุปจฺจยา. อาจยคามิํ เอกํ ขนฺธํ ปฏิจฺจ ตโย ขนฺธา ฯเปฯ เทฺว ขนฺเธ ปฏิจฺจ เทฺว ขนฺธา. (1)}

\prose{อาจยคามิํ ธมฺมํ ปฏิจฺจ เนวา\-จยคามิ\-นา\-ปจยคามี ธมฺโม อุปฺปชฺชติ เหตุปจฺจยา. อาจยคามี ขนฺเธ ปฏิจฺจ จิตฺต\-สมุฏฺฐานํ รูปํ. (2)}

\prose{อาจยคามิํ ธมฺมํ ปฏิจฺจ อาจยคามี จ เนวา\-จยคามิ\-นา\-ปจยคามี จ ธมฺมา อุปฺปชฺชนฺติ เหตุปจฺจยา. อาจยคามิํ เอกํ ขนฺธํ ปฏิจฺจ ตโย ขนฺธา จิตฺต\-สมุฏฺฐานญฺจ รูปํ ฯเปฯ เทฺว ขนฺเธ ปฏิจฺจ เทฺว ขนฺธา จิตฺต\-สมุฏฺฐานญฺจ รูปํ. (3)}

\proseitem{2}{อปจยคามิํ ธมฺมํ ปฏิจฺจ อปจยคามี ธมฺโม อุปฺปชฺชติ เหตุปจฺจยา. อปจยคามิํ เอกํ ขนฺธํ ปฏิจฺจ ตโย ขนฺธา ฯเปฯ เทฺว ขนฺเธ ปฏิจฺจ เทฺว ขนฺธา. (1)}

\prose{อปจยคามิํ ธมฺมํ ปฏิจฺจ เนวา\-จยคามิ\-นา\-ปจยคามี ธมฺโม อุปฺปชฺชติ เหตุปจฺจยา. อปจยคามี ขนฺเธ ปฏิจฺจ จิตฺต\-สมุฏฺฐานํ รูปํ. (2)}

\prose{อปจยคามิํ ธมฺมํ ปฏิจฺจ อปจยคามี จ เนวา\-จยคามิ\-นา\-ปจยคามี จ ธมฺมา อุปฺปชฺชนฺติ เหตุปจฺจยา. อปจยคามิํ เอกํ ขนฺธํ ปฏิจฺจ ตโย ขนฺธา จิตฺต\-สมุฏฺฐานํ รูปํ ฯเปฯ เทฺว ขนฺเธ ปฏิจฺจ เทฺว ขนฺธา จิตฺตสมุฏฺฐานานญฺจ รูปํ. (3)}

\proseitem{3}{เนวา\-จยคามิ\-นา\-ปจยคามิํ ธมฺมํ ปฏิจฺจ เนวา\-จยคามิ\-นา\-ปจยคามี ธมฺโม อุปฺปชฺชติ เหตุปจฺจยา. เนวา\-จยคามิ\-นา\-ปจยคามิํ เอกํ ขนฺธํ ปฏิจฺจ ตโย ขนฺธา จิตฺต\-สมุฏฺฐานญฺจ รูปํ ฯเปฯ เทฺว ขนฺเธ ฯเปฯ. ปฏิสนฺธิ\-กฺขเณ เนวา\-จยคามิ\-นา\-ปจยคามิํ เอกํ ขนฺธํ ปฏิจฺจ ตโย ขนฺธา กฏตฺตา จ รูปํ ฯเปฯ เทฺว ขนฺเธ ฯเปฯ. ขนฺเธ ปฏิจฺจ วตฺถุ, วตฺถุํ ปฏิจฺจ ขนฺธา. เอกํ มหาภูตํ ปฏิจฺจ ตโย มหาภูตา ฯเปฯ เทฺว มหาภูเต ปฏิจฺจ เทฺว มหาภูตา. มหาภูเต ปฏิจฺจ จิตฺต\-สมุฏฺฐานํ รูปํ, กฏตฺตารูปํ, อุปาทารูปํ. (1)}

\prose{\csromanfolio{241}อาจยคามิญฺจ เนวา\-จยคามิ\-นา\-ปจยคามิญฺจ ธมฺมํ ปฏิจฺจ เนวา\-จยคามิ\-นา\-ปจยคามี ธมฺโม อุปฺปชฺชติ เหตุปจฺจยา. อาจยคามี ขนฺเธ จ มหาภูเต จ ปฏิจฺจ จิตฺต\-สมุฏฺฐานํ รูปํ. (1)}

\prose{อปจยคามิญฺจ เนวา\-จยคามิ\-นา\-ปจยคามิญฺจ ธมฺมํ ปฏิจฺจ เนวา\-จยคามิ\-นา\-ปจยคามี ธมฺโม อุปฺปชฺชติ เหตุปจฺจยา. อปจยคามี ขนฺเธ จ มหาภูเต จ ปฏิจฺจ จิตฺต\-สมุฏฺฐานํ รูปํ. (1)}

\csromanheader{2}{\textbf{อารมฺมณ}}
\proseitem{4}{อาจยคามิํ ธมฺมํ ปฏิจฺจ อาจยคามี ธมฺโม อุปฺปชฺชติ อารมฺมณ\-ปจฺจยา. อาจยคามิํ เอกํ ขนฺธํ ปฏิจฺจ ตโย ขนฺธา ฯเปฯ เทฺว ขนฺเธ ฯเปฯ. (1)}

\prose{อปจยคามิํ ธมฺมํ ปฏิจฺจ อปจยคามี ธมฺโม อุปฺปชฺชติ อารมฺมณ\-ปจฺจยา. อปจยคามิํ เอกํ ขนฺธํ ปฏิจฺจ ตโย ขนฺธา ฯเปฯ เทฺว ขนฺเธ ฯเปฯ. (1)}

\prose{เนวา\-จยคามิ\-นา\-ปจยคามิํ ธมฺมํ ปฏิจฺจ เนวา\-จยคามิ\-นา\-ปจยคามี ธมฺโม อุปฺปชฺชติ อารมฺมณ\-ปจฺจยา. เนวา\-จยคามิ\-นา\-ปจยคามิํ เอกํ ขนฺธํ ปฏิจฺจ ตโย ขนฺธา ฯเปฯ เทฺว ขนฺเธ ปฏิจฺจ เทฺว ขนฺธา. ปฏิสนฺธิ\-กฺขเณ ฯเปฯ. วตฺถุํ ปฏิจฺจ ขนฺธา. (1)}

\proseitem{5}{อาจยคามิํ ธมฺมํ ปฏิจฺจ อาจยคามี ธมฺโม อุปฺปชฺชติ อธิปติ\-ปจฺจยา. ตีณิ.}

\prose{อปจยคามิํ ธมฺมํ ปฏิจฺจ อปจยคามี ธมฺโม อุปฺปชฺชติ อธิปติ\-ปจฺจยา. ตีณิ.}

\prose{เนวา\-จยคามิ\-นา\-ปจยคามิํ ธมฺมํ ปฏิจฺจ เนวา\-จยคามิ\-นา\-ปจยคามี ธมฺโม. เอกํ. (ปฏิสนฺธิ นตฺถิ.) เอกํ มหาภูตํ ปฏิจฺจ ตโย มหาภูตา ฯเปฯ. มหาภูเต ปฏิจฺจ จิตฺต\-สมุฏฺฐานํ รูปํ, อุปาทารูปํ. (1)}

\prose{อาจยคามิญฺจ เนวา\-จยคามิ\-นา\-ปจยคามิญฺจ ธมฺมํ ปฏิจฺจ เนวา\-จยคามิ\-นา\-ปจยคามี ธมฺโม อุปฺปชฺชติ อธิปติ\-ปจฺจยา. อาจยคามี ขนฺเธ จ มหาภูเต จ ปฏิจฺจ จิตฺต\-สมุฏฺฐาน รูปํ. (1)}

\prose{\csromanfolio{242}อปจยคามิญฺจ เนวา\-จยคามิ\-นา\-ปจยคามิญฺจ ธมฺมํ ปฏิจฺจ เนวา\-จยคามิ\-นา\-ปจยคามี ธมฺโม อุปฺปชฺชติ อธิปติ\-ปจฺจยา. อปจยคามี ขนฺเธ จ มหาภูเต จ ปฏิจฺจ จิตฺต\-สมุฏฺฐานํ รูปํ. (1)}

\csromanheader{2}{\textbf{อนนฺตราทิ}}
\proseitem{6}{อาจยคามิํ ธมฺมํ ปฏิจฺจ อาจยคามี ธมฺโม อุปฺปชฺชติ อนนฺตร\-ปจฺจยา. สมนนฺตร\-ปจฺจยา. สหชาต\-ปจฺจยา. (สพฺเพปิ มหาภูตา กาตพฺพา.) อญฺญมญฺญ\-ปจฺจยา. (จิตฺต\-สมุฏฺฐานมฺปิ กฏตฺตารูปมฺปิ อุปาทารูปมฺปิ นตฺถิ.) นิสฺสยปจฺจยา. อุปนิสฺสย\-ปจฺจยา. ปุเรชาต\-ปจฺจยา. อาเสวนปจฺจยา. กมฺมปจฺจยา. วิปากปจฺจยา. อาหารปจฺจยา. อินฺทฺริยปจฺจยา. ฌานปจฺจยา. มคฺคปจฺจยา. สมฺปยุตฺต\-ปจฺจยา. วิปฺปยุตฺต\-ปจฺจยา. อตฺถิปจฺจยา. นตฺถิปจฺจยา. วิคตปจฺจยา. อวิคตปจฺจยา.}

\csromanmarkh{1. ปจฺจยานุโลม}\csromanmarkhii{2. สงฺขฺยาวาร}\csromanheadernomark{2}{1. \textbf{ปจฺจยานุโลม 2. สงฺขฺยาวาร}}
\proseitemwithcloserruled{7}{เหตุยา นว, อารมฺมเณ ตีณิ, อธิปติยา นว, อนนฺตเร ตีณิ, สมนนฺตเร ตีณิ, สหชาเต นว, อญฺญมญฺเญ ตีณิ, นิสฺสเย นว, อุปนิสฺสเย ตีณิ, ปุเรชาเต ตีณิ, อาเสวเน ตีณิ, กมฺเม นว, วิปาเก เอกํ, อาหาเร นว, อินฺทฺริเย นว, ฌาเน นว, มคฺเค นว, สมฺปยุตฺเต ตีณิ, วิปฺปยุตฺเต นว, อตฺถิยา นว, นตฺถิยา ตีณิ, วิคเต ตีณิ, อวิคเต นว. (เอวํ คเณตพฺพํ.)}{อนุโลมํ.}
\csromanmarkh{2. ปจฺจย\-ปจฺจนีย}\csromanmarkhii{1. วิภงฺควาร}\csromanheadernomark{2}{2. ปจฺจย\-ปจฺจนีย 1. วิภงฺควาร}
\csromanheader{2}{\textbf{นเหตุ}}
\proseitem{8}{อาจยคามิํ ธมฺมํ ปฏิจฺจ อาจยคามี ธมฺโม อุปฺปชฺชติ นเหตุปจฺจยา. วิจิกิจฺฉา\-สหคเต อุทฺธจฺจ\-สหคเต ขนฺเธ ปฏิจฺจ วิจิกิจฺฉา\-สหคโต อุทฺธจฺจ\-สหคโต โมโห. (1)}

\prose{\csromanfolio{243}เนวา\-จยคามิ\-นา\-ปจยคามิํ ธมฺมํ ปฏิจฺจ เนวา\-จยคามิ\-นา\-ปจยคามี ธมฺโม อุปฺปชฺชติ นเหตุปจฺจยา. อเหตุกํ เนวา\-จยคามิ\-นา\-ปจยคามิํ เอกํ ขนฺธํ ปฏิจฺจ ตโย ขนฺธา จิตฺต\-สมุฏฺฐานญฺจ รูปํ ฯเปฯ เทฺว ขนฺเธ ฯเปฯ. อเหตุก\-ปฏิสนฺธิ\-กฺขเณ ฯเปฯ. ขนฺเธ ปฏิจฺจ วตฺถุ, วตฺถุํ ปฏิจฺจ ขนฺธา. เอกํ มหาภูตํ ฯเปฯ. พาหิรํ. อาหารสมุฏฺฐานํ. อุตุสมุฏฺฐานํ. อสญฺญสตฺตานํ เอกํ มหาภูตํ ฯเปฯ. (1)}

\prose{นอารมฺมณ}

\proseitem{9}{อาจยคามิํ ธมฺมํ ปฏิจฺจ เนวา\-จยคามิ\-นา\-ปจยคามี ธมฺโม อุปฺปชฺชติ นอารมฺมณ\-ปจฺจยา. อาจยคามี ขนฺเธ ปฏิจฺจ จิตฺต\-สมุฏฺฐานํ รูปํ. (1)}

\prose{อปจยคามิํ ธมฺมํ ปฏิจฺจ เนวา\-จยคามิ\-นา\-ปจยคามี ธมฺโม อุปฺปชฺชติ นอารมฺมณ\-ปจฺจยา. อปจยคามี ขนฺเธ ปฏิจฺจ จิตฺต\-สมุฏฺฐานํ รูปํ. (1)}

\prose{เนวา\-จยคามิ\-นา\-ปจยคามิํ ธมฺมํ ปฏิจฺจ เนวา\-จยคามิ\-นา\-ปจยคามี ธมฺโม อุปฺปชฺชติ นอารมฺมณ\-ปจฺจยา. เนวา\-จยคามิ\-นา\-ปจยคามี ขนฺเธ ปฏิจฺจ จิตฺต\-สมุฏฺฐานํ รูปํ. ปฏิสนฺธิ\-กฺขเณ ฯเปฯ. ขนฺเธ ปฏิจฺจ วตฺถุ ฯเปฯ. เอกํ มหาภูตํ ฯเปฯ. พาหิรํ. อาหารสมุฏฺฐานํ. อุตุสมุฏฺฐานํ. อสญฺญสตฺตานํ ฯเปฯ. (1)}

\prose{อาจยคามิญฺจ เนวา\-จยคามิ\-นา\-ปจยคามิญฺจ ธมฺมํ ปฏิจฺจ เนวา\-จยคามิ\-นา\-ปจยคามี ธมฺโม อุปฺปชฺชติ นอารมฺมณ\-ปจฺจยา. อาจยคามี ขนฺเธ จ มหาภูเต จ ปฏิจฺจ จิตฺต\-สมุฏฺฐานํ รูปํ. (1)}

\prose{อปจยคามิญฺจ เนวา\-จยคามิ\-นา\-ปจยคามิญฺจ ธมฺมํ ปฏิจฺจ เนวา\-จยคามิ\-นา\-ปจยคามี ธมฺโม อุปฺปชฺชติ นอารมฺมณ\-ปจฺจยา. อปจยคามี ขนฺเธ จ มหาภูเต จ ปฏิจฺจ จิตฺต\-สมุฏฺฐานํ รูปํ. (1)}

\prose{นอธิปติ}

\proseitem{10}{อาจยคามิํ ธมฺมํ ปฏิจฺจ อาจยคามี ธมฺโม อุปฺปชฺชติ นอธิปติ\-ปจฺจยา. ตีณิ.}

\prose{อปจยคามิํ ธมฺมํ ปฏิจฺจ อปจยคามี ธมฺโม อุปฺปชฺชติ นอธิปติ\-ปจฺจยา. อปจยคามี ขนฺเธ ปฏิจฺจ อปจยคามี อธิปติ. (1)}

\prose{\csromanfolio{244}เนวา\-จยคามิ\-นา\-ปจยคามิํ ธมฺมํ ปฏิจฺจ เนวา\-จยคามิ\-นา\-ปจยคามี ธมฺโม อุปฺปชฺชติ นอธิปติ\-ปจฺจยา. เนวา\-จยคามิ\-นา\-ปจยคามิํ เอกํ ขนฺธํ ปฏิจฺจ ตโย ขนฺธา จิตฺต\-สมุฏฺฐานญฺจ รูปํ ฯเปฯ เทฺว ขนฺเธ ฯเปฯ. ปฏิสนฺธิ\-กฺขเณ ฯเปฯ. ขนฺเธ ปฏิจฺจ วตฺถุ, วตฺถุํ ปฏิจฺจ ขนฺธา. เอกํ มหาภูตํ ฯเปฯ. พาหิรํ. อาหารสมุฏฺฐานํ. อุตุสมุฏฺฐานํ. อสญฺญสตฺตานํ เอกํ มหาภูตํ ฯเปฯ. (1)}

\prose{อาจยคามิญฺจ เนวา\-จยคามิ\-นา\-ปจยคามิญฺจ ธมฺมํ ปฏิจฺจ เนวา\-จยคามิ\-นา\-ปจยคามี ธมฺโม อุปฺปชฺชติ นอธิปติ\-ปจฺจยา. อาจยคามี ขนฺเธ จ มหาภูเต จ ปฏิจฺจ จิตฺต\-สมุฏฺฐานํ รูปํ. (1)}

\prose{นอนนฺตราทิ}

\proseitem{11}{อาจยคามิํ ธมฺมํ ปฏิจฺจ เนวา\-จยคามิ\-นา\-ปจยคามี ธมฺโม อุปฺปชฺชติ นอนนฺตร\-ปจฺจยา. นสมนนฺตร\-ปจฺจยา. นอญฺญมญฺญ\-ปจฺจยา. นอุปนิสฺสยปจฺจยา. นปุเรชาต\-ปจฺจยา. (กุสล\-ตฺติก\-สทิสา สตฺต ปญฺหา.) นปจฺฉาชาต\-ปจฺจยา.}

\prose{นอาเสวน}

\proseitem{12}{อาจยคามิํ ธมฺมํ ปฏิจฺจ อาจยคามี ธมฺโม อุปฺปชฺชติ นอาเสวน\-ปจฺจยา. ตีณิ.}

\prose{อปจยคามิํ ธมฺมํ ปฏิจฺจ เนวา\-จยคามิ\-นา\-ปจยคามี ธมฺโม อุปฺปชฺชติ นอาเสวน\-ปจฺจยา. อปจยคามี ขนฺเธ ปฏิจฺจ จิตฺต\-สมุฏฺฐานํ รูปํ. (1)}

\prose{เนวา\-จยคามิ\-นา\-ปจยคามิํ ธมฺมํ ปฏิจฺจ เนวา\-จยคามิ\-นา\-ปจยคามี ธมฺโม อุปฺปชฺชติ นอาเสวน\-ปจฺจยา. (เอกา ปญฺหา. สพฺเพ มหาภูตา กาตพฺพา.)}

\prose{อาจยคามิญฺจ เนวา\-จยคามิ\-นา\-ปจยคามิญฺจ ธมฺมํ ปฏิจฺจ เนวา\-จยคามิ\-นา\-ปจยคามี ธมฺโม อุปฺปชฺชติ นอาเสวน\-ปจฺจยา. อาจยคามี ขนฺเธ จ มหาภูเต จ ปฏิจฺจ จิตฺต\-สมุฏฺฐานํ รูปํ. (1)}

\prose{อปจยคามิญฺจ เนวา\-จยคามิ\-นา\-ปจยคามิญฺจ ธมฺมํ ปฏิจฺจ เนวา\-จยคามิ\-นา\-ปจยคามี ธมฺโม อุปฺปชฺชติ นอาเสวน\-ปจฺจยา. อปจยคามี ขนฺเธ จ มหาภูเต จ ปฏิจฺจ จิตฺต\-สมุฏฺฐานํ รูปํ. (1)}

\csromanheader{2}{10. อาจยคามิตฺติก \textbf{นกมฺม}}
\proseitem{13}{\csromanfolio{245}อาจยคามิํ ธมฺมํ ปฏิจฺจ อาจยคามี ธมฺโม อุปฺปชฺชติ นกมฺมปจฺจยา. อาจยคามี ขนฺเธ ปฏิจฺจ อาจยคามี เจตนา. (1)}

\prose{อปจยคามิํ ธมฺมํ ปฏิจฺจ อปจยคามี ธมฺโม อุปฺปชฺชติ นกมฺมปจฺจยา. อปจยคามี ขนฺเธ ปฏิจฺจ อปจยคามี เจตนา. (1)}

\prose{เนวา\-จยคามิ\-นา\-ปจยคามิํ ธมฺมํ ปฏิจฺจ เนวา\-จยคามิ\-นา\-ปจยคามี ธมฺโม อุปฺปชฺชติ นกมฺมปจฺจยา. เนวา\-จยคามิ\-นา\-ปจยคามี ขนฺเธ ปฏิจฺจ เนวา\-จยคามิ\-นา\-ปจยคามี เจตนา. พาหิรํ. อาหารสมุฏฺฐานํ. อุตุสมุฏฺฐานํ. เอกํ มหาภูตํ ฯเปฯ. (1)}

\csromanheader{2}{\textbf{นวิปากาทิ}}
\proseitem{14}{อาจยคามิํ ธมฺมํ ปฏิจฺจ อาจยคามี ธมฺโม อุปฺปชฺชติ นวิปาก\-ปจฺจยา. (ปริปุณฺณํ, ปฏิสนฺธิ นตฺถิ.) นอาหารปจฺจยา. นอินฺทฺริยปจฺจยา. นฌานปจฺจยา. นมคฺคปจฺจยา. นสมฺปยุตฺต\-ปจฺจยา. นวิปฺปยุตฺต\-ปจฺจยา. (ตีณิ.) โนนตฺถิปจฺจยา. โนวิคต\-ปจฺจยา.}

\csromanmarkh{2. ปจฺจย\-ปจฺจนีย}\csromanmarkhii{2. สงฺขฺยาวาร}\csromanheadernomark{2}{2. ปจฺจย\-ปจฺจนีย 2. สงฺขฺยาวาร}
\proseitemwithcloserruled{15}{นเหตุยา เทฺว, นอารมฺมเณ ปญฺจ, นอธิปติยา ฉ, นอนนฺตเร ปญฺจ, นสมนนฺตเร ปญฺจ, นอญฺญมญฺเญ ปญฺจ, นอุปนิสฺสเย ปญฺจ, นปุเรชาเต สตฺต, นปจฺฉาชาเต นว, นอาเสวเน สตฺต, นกมฺเม ตีณิ, นวิปาเก นว, นอาหาเร เอกํ, นอินฺทฺริเย เอกํ, นฌาเน เอกํ, นมคฺเค เอกํ, นสมฺปยุตฺเต ปญฺจ, นวิปฺปยุตฺเต ตีณิ, โนนตฺถิยา ปญฺจ, โนวิคเต ปญฺจ. (เอวํ คเณตพฺพํ.)}{ปจฺจนียํ.}
\proseitemwithcloserruled{16}{เหตุปจฺจยา นอารมฺมเณ ปญฺจ, นอธิปติยา ฉ, นอนนฺตเร ปญฺจ, นสมนนฺตเร ปญฺจ, นอญฺญมญฺเญ ปญฺจ, นอุปนิสฺสเย ปญฺจ, นปุเรชาเต \csromanfolio{246}สตฺต, นปจฺฉาชาเต นว, นอาเสวเน สตฺต, นกมฺเม ตีณิ, นวิปาเก นว, นสมฺปยุตฺเต ปญฺจ, นวิปฺปยุตฺเต ตีณิ, โนนตฺถิยา ปญฺจ, โนวิคเต ปญฺจ. (เอวํ คเณตพฺพํ.)}{อนุโลม\-ปจฺจนียํ.}
\prose{นเหตุปจฺจยา อารมฺมเณ เทฺว, อนนฺตเร เทฺว, สมนนฺตเร เทฺว, สหชาเต เทฺว, อญฺญมญฺเญ, นิสฺสเย, อุปนิสฺสเย, ปุเรชาเต, อาเสวเน, กมฺเม เทฺว, วิปาเก เอกํ, อาหาเร เทฺว, อินฺทฺริเย เทฺว, ฌาเน เทฺว, มคฺเค เอกํ, สมฺปยุตฺเต เทฺว, วิปฺปยุตฺเต, อตฺถิยา, นตฺถิยา, วิคเต, อวิคเต เทฺว. (เอวํ คเณตพฺพํ.) ปจฺจนียา\-นุโลมํ. ปฏิจฺจวาโร. (สหชาตวาโร ปฏิจฺจ\-วาร\-สทิโส.)}
\csromansectionrule

\csromanmarkh{10. อาจยคามิตฺติก}\csromanmarkhii{3. ปจฺจยวาร}\csromanheadernomark{2}{10. \textbf{อาจยคามิตฺติก 3. ปจฺจยวาร}}
\csromanmarkh{1. ปจฺจยานุโลม}\csromanmarkhii{1. วิภงฺควาร}\csromanheadernomark{2}{1. \textbf{ปจฺจยานุโลม 1. วิภงฺควาร}}
\csromanheader{2}{\textbf{เหตุ}}
\proseitem{16}{อาจยคามิํ ธมฺมํ ปจฺจยา อาจยคามี ธมฺโม อุปฺปชฺชติ เหตุปจฺจยา. อาจยคามิํ เอกํ ขนฺธํ ปจฺจยา ตโย ขนฺธา ฯเปฯ เทฺว ขนฺเธ ปจฺจยา เทฺว ขนฺธา. (1)}

\prose{อาจยคามิํ ธมฺมํ ปจฺจยา เนวา\-จยคามิ\-นา\-ปจยคามี ธมฺโม อุปฺปชฺชติ เหตุปจฺจยา. อาจยคามี ขนฺเธ ปจฺจยา จิตฺต\-สมุฏฺฐานํ รูปํ. (2)}

\prose{\csromanfolio{247}อาจยคามิํ ธมฺมํ ปจฺจยา อาจยคามี จ เนวา\-จยคามิ\-นา\-ปจยคามี จ ธมฺมา อุปฺปชฺชนฺติ เหตุปจฺจยา. อาจยคามิํ เอกํ ขนฺธํ ปจฺจยา ตโย ขนฺธา จิตฺต\-สมุฏฺฐานญฺจ รูปํ ฯเปฯ เทฺว ขนฺเธ ฯเปฯ. (3)}

\prose{อปจยคามิํ ธมฺมํ ปจฺจยา อปจยคามี ธมฺโม. ตีณิ.}

\proseitem{19}{เนวา\-จยคามิ\-นา\-ปจยคามิํ ธมฺมํ ปจฺจยา เนวา\-จยคามิ\-นา\-ปจยคามี ธมฺโม อุปฺปชฺชติ เหตุปจฺจยา. เนวา\-จยคามิ\-นา\-ปจยคามิํ เอกํ ขนฺธํ ปจฺจยา ตโย ขนฺธา จิตฺต\-สมุฏฺฐานญฺจ รูปํ ฯเปฯ เทฺว ขนฺเธ ฯเปฯ. ปฏิสนฺธิ\-กฺขเณ ฯเปฯ. ขนฺเธ ปจฺจยา วตฺถุ, วตฺถุํ ปจฺจยา ขนฺธา. เอกํ มหาภูตํ ปจฺจยา ฯเปฯ. วตฺถุํ ปจฺจยา เนวา\-จยคามิ\-นา\-ปจยคามี ขนฺธา. (1)}

\prose{เนวา\-จยคามิ\-นา\-ปจยคามิํ ธมฺมํ ปจฺจยา อาจยคามี ธมฺโม อุปฺปชฺชติ เหตุปจฺจยา. วตฺถุํ ปจฺจยา อาจยคามี ขนฺธา. (2)}

\prose{เนวา\-จยคามิ\-นา\-ปจยคามิํ ธมฺมํ ปจฺจยา อปจยคามี ธมฺโม อุปฺปชฺชติ เหตุปจฺจยา. วตฺถุํ ปจฺจยา อปจยคามี ขนฺธา. (3)}

\prose{เนวา\-จยคามิ\-นา\-ปจยคามิํ ธมฺมํ ปจฺจยา อาจยคามี จ เนวา\-จยคามิ\-นา\-ปจยคามี จ ธมฺมา อุปฺปชฺชนฺติ เหตุปจฺจยา. วตฺถุํ ปจฺจยา อาจยคามี ขนฺธา. มหาภูเต ปจฺจยา จิตฺต\-สมุฏฺฐานํ รูปํ. (4)}

\prose{เนวา\-จยคามิ\-นา\-ปจยคามิํ ธมฺมํ ปจฺจยา อปจยคามี จ เนวา\-จยคามิ\-นา\-ปจยคามี จ ธมฺมา อุปฺปชฺชนฺติ เหตุปจฺจยา. วตฺถุํ ปจฺจยา อปจยคามี ขนฺธา. มหาภูเต ปจฺจยา จิตฺต\-สมุฏฺฐานํ รูปํ. (5)}

\proseitem{20}{อาจยคามิญฺจ เนวา\-จยคามิ\-นา\-ปจยคามิญฺจ ธมฺมํ ปจฺจยา อาจยคามี ธมฺโม อุปฺปชฺชติ เหตุปจฺจยา. อาจยคามิํ เอกํ ขนฺธญฺจ วตฺถุญฺจ ปจฺจยา ตโย ขนฺธา ฯเปฯ เทฺว ขนฺเธ ฯเปฯ. (1)}

\prose{อาจยคามิญฺจ เนวา\-จยคามิ\-นา\-ปจยคามิญฺจ ธมฺมํ ปจฺจยา เนวา\-จยคามิ\-นา\-ปจยคามี ธมฺโม อุปฺปชฺชติ เหตุปจฺจยา. อาจยคามี ขนฺเธ จ มหาภูเต จ ปจฺจยา จิตฺต\-สมุฏฺฐานํ รูปํ. (2)}

\prose{อาจยคามิญฺจ เนวา\-จยคามิ\-นา\-ปจยคามิญฺจ ธมฺมํ ปจฺจยา อาจยคามี จ เนวา\-จยคามิ\-นา\-ปจยคามี จ ธมฺมา อุปฺปชฺชนฺติ เหตุปจฺจยา. อาจยคามิํ \csromanfolio{248}เอกํ ขนฺธญฺจ วตฺถุญฺจ ปจฺจยา ตโย ขนฺธา ฯเปฯ เทฺว ขนฺเธ ฯเปฯ. อาจยคามี ขนฺเธ จ มหาภูเต จ ปจฺจยา จิตฺต\-สมุฏฺฐานํ รูปํ. (3)}

\prose{อปจยคามิญฺจ เนวา\-จยคามิ\-นา\-ปจยคามิญฺจ ธมฺมํ ปจฺจยา อปจยคามี ธมฺโม. ตีณิ.}

\csromanheader{2}{\textbf{อารมฺมณ}}
\proseitem{20}{อาจยคามิํ ธมฺมํ ปจฺจยา อาจยคามี ธมฺโม อุปฺปชฺชติ อารมฺมณ\-ปจฺจยา. อาจยคามิํ เอกํ ขนฺธํ ปจฺจยา ฯเปฯ. (1)}

\prose{อปจยคามิํ ธมฺมํ ปจฺจยา อปจยคามี ธมฺโม. เอกํ.}

\prose{เนวา\-จยคามิ\-นา\-ปจยคามิํ ธมฺมํ ปจฺจยา เนวา\-จยคามิ\-นา\-ปจยคามี ธมฺโม อุปฺปชฺชติ อารมฺมณ\-ปจฺจยา. เนวา\-จยคามิ\-นา\-ปจยคามิํ เอกํ ขนฺธํ ปจฺจยา ตโย ขนฺธา ฯเปฯ เทฺว ขนฺเธ ฯเปฯ. ปฏิสนฺธิ\-กฺขเณ ฯเปฯ. วตฺถุํ ปจฺจยา ขนฺธา. จกฺขายตนํ ปจฺจยา จกฺขุวิญฺญาณํ ฯเปฯ. กายายตนํ ปจฺจยา กายวิญฺญาณํ. วตฺถุํ ปจฺจยา เนวา\-จยคามิ\-นา\-ปจยคามี ขนฺธา. (1)}

\prose{เนวา\-จยคามิ\-นา\-ปจยคามิํ ธมฺมํ ปจฺจยา อาจยคามี ธมฺโม อุปฺปชฺชติ อารมฺมณ\-ปจฺจยา. วตฺถุํ ปจฺจยา อาจยคามี ขนฺธา. (2)}

\prose{เนวา\-จยคามิ\-นา\-ปจยคามิํ ธมฺมํ ปจฺจยา อปจยคามี ธมฺโม อุปฺปชฺชติ อารมฺมณ\-ปจฺจยา. วตฺถุํปจฺจยา อปจยคามี ขนฺธา. (3)}

\prose{อาจยคามิญฺจ เนวา\-จยคามิ\-นา\-ปจยคามิญฺจ ธมฺมํ ปจฺจยา อาจยคามี ธมฺโม อุปฺปชฺชติ อารมฺมณ\-ปจฺจยา. อาจยคามิํ เอกํ ขนฺธญฺจ วตฺถุญฺจ ปจฺจยา ตโย ขนฺธา ฯเปฯ เทฺว ขนฺเธ ฯเปฯ. (1)}

\prose{อปจยคามิญฺจ เนวา\-จยคามิ\-นา\-ปจยคามิญฺจ ธมฺมํ ปจฺจยา อปจยคามี ธมฺโม อุปฺปชฺชติ อารมฺมณ\-ปจฺจยา. อปจยคามิํ เอกํ ขนฺธญฺจ วตฺถุญฺจ ปจฺจยา ตโย ขนฺธา ฯเปฯ เทฺว ขนฺเธ ฯเปฯ. (1)}

\prose{อาจยคามิํ ธมฺมํ ปจฺจยา อาจยคามี ธมฺโม อุปฺปชฺชติ อธิปติ\-ปจฺจยา. ตีณิ.}

\prose{อปจยคามิํ ธมฺมํ ปจฺจยา อปจยคามี ธมฺโม. ตีณิ.}

\prose{\csromanfolio{249}เนวา\-จยคามิ\-นา\-ปจยคามิํ ธมฺมํ ปจฺจยา เนวา\-จยคามิ\-นา\-ปจยคามี ธมฺโม. เอกํ ฯเปฯ. วตฺถุํ ปจฺจยา เนวา\-จยคามิ\-นา\-ปจยคามี ขนฺธา.}

\prose{เนวา\-จยคามิ\-นา\-ปจยคามิํ ธมฺมํ ปจฺจยา อาจยคามี ธมฺโม ฯเปฯ. (อิธาปิ ฆฏนา เหตุสทิสา.)}

\csromanheader{2}{\textbf{อนนฺตราทิ}}
\proseitem{23}{อาจยคามิํ ธมฺมํ ปจฺจยา อาจยคามี ธมฺโม อุปฺปชฺชติ อนนฺตร\-ปจฺจยา. สมนนฺตร\-ปจฺจยา. สหชาต\-ปจฺจยา. ตีณิ.}

\prose{อปจยคามิํ ธมฺมํ ปจฺจยา อปจยคามี ธมฺโม. ตีณิ.}

\prose{เนวา\-จยคามิ\-นา\-ปจยคามิํ ธมฺมํ ปจฺจยา เนวา\-จยคามิ\-นา\-ปจยคามี ธมฺโม อุปฺปชฺชติ สหชาต\-ปจฺจยา. เนวา\-จยคามิ\-นา\-ปจยคามิํ เอกํ ขนฺธํ ปจฺจยา ตโย ขนฺธา จิตฺต\-สมุฏฺฐานญฺจ รูปํ ฯเปฯ เทฺว ขนฺเธ ฯเปฯ. ปฏิสนฺธิ\-กฺขเณ ฯเปฯ. อสญฺญสตฺตานํ เอกํ มหาภูตํ ฯเปฯ. จกฺขายตนํ ปจฺจยา จกฺขุวิญฺญาณํ ฯเปฯ. กายายตนํ ปจฺจยา กายวิญฺญาณํ. วตฺถุํ ปจฺจยา ฯเปฯ. (1)}

\prose{เนวา\-จยคามิ\-นา\-ปจยคามิํ ธมฺมํ ปจฺจยา อาจยคามี ธมฺโม อุปฺปชฺชติ สหชาต\-ปจฺจยา. (สํขิตฺตํ. สพฺเพ ฆฏนา กาตพฺพา.)}

\csromanheader{2}{\textbf{อญฺญมญฺญาทิ}}
\proseitem{24}{อปจยคามิํ ธมฺมํ ปจฺจยา อปจยคามี ธมฺโม อุปฺปชฺชติ อญฺญมญฺญ\-ปจฺจยา. นิสฺสยปจฺจยา. อุปนิสฺสย\-ปจฺจยา. ปุเรชาต\-ปจฺจยา. อาเสวนปจฺจยา. กมฺมปจฺจยา. วิปากปจฺจยา. อาหารปจฺจยา. อินฺทฺริยปจฺจยา. ฌานปจฺจยา. มคฺคปจฺจยา. สมฺปยุตฺต\-ปจฺจยา. วิปฺปยุตฺต\-ปจฺจยา. อตฺถิปจฺจยา. นตฺถิปจฺจยา. วิคตปจฺจยา. อวิคตปจฺจยา.}

\csromanmarkh{1. ปจฺจยานุโลม}\csromanmarkhii{2. สงฺขฺยาวาร}\csromanheadernomark{2}{1. \textbf{ปจฺจยานุโลม 2. สงฺขฺยาวาร}}
\proseitemwithcloserruled{25}{เหตุยา สตฺตรส, อารมฺมเณ สตฺต, อธิปติยา สตฺตรส, อนนฺตเร สตฺต, สมนนฺตเร สตฺต, สหชาเต สตฺตรส, อญฺญมญฺเญ สตฺต, นิสฺสเย สตฺตรส, อุปนิสฺสเย สตฺต, ปุเรชาเต สตฺต, \csromanfolio{250}อาเสวเน สตฺต, กมฺเม สตฺตรส, วิปาเก เอกํ, อาหาเร สตฺตรส, อินฺทฺริเย, ฌาเน, มคฺเค สตฺตรส, สมฺปยุตฺเต สตฺต, วิปฺปยุตฺเต สตฺตรส, อตฺถิยา สตฺตรส, นตฺถิยา สตฺต, วิคเต สตฺต, อวิคเต สตฺตรส. (เอวํ คเณตพฺพํ.)}{อนุโลมํ.}
\csromanmarkh{2. ปจฺจย\-ปจฺจนีย}\csromanmarkhii{1. วิภงฺควาร}\csromanheadernomark{2}{2. ปจฺจย\-ปจฺจนีย 1. วิภงฺควาร}
\csromanheader{2}{\textbf{นเหตุ}}
\proseitem{26}{อาจยคามิํ ธมฺมํ ปจฺจยา อาจยคามี ธมฺโม อุปฺปชฺชติ นเหตุปจฺจยา, วิจิกิจฺฉา\-สหคเต อุทฺธจฺจ\-สหคเต ขนฺเธ ปจฺจยา วิจิกิจฺฉา\-สหคโต อุทฺธจฺจ\-สหคโต โมโห. (1)}

\prose{เนวา\-จยคามิ\-นา\-ปจยคามิํ ธมฺมํ ปจฺจยา เนวา\-จยคามิ\-นา\-ปจยคามี ธมฺโม อุปฺปชฺชติ นเหตุปจฺจยา. อเหตุกํ เนวา\-จยคามิ\-นา\-ปจยคามิํ เอกํ ขนฺธํ ปจฺจยา ตโย ขนฺธา จิตฺต\-สมุฏฺฐานญฺจ รูปํ ฯเปฯ เทฺว ขนฺเธ ฯเปฯ. อเหตุก\-ปฏิสนฺธิ\-กฺขเณ ฯเปฯ. อสญฺญสตฺตานํ เอกํ มหาภูตํ ฯเปฯ. จกฺขายตนํ ปจฺจยา จกฺขุวิญฺญาณํ ฯเปฯ. กายายตนํ ปจฺจยา กายวิญฺญาณํ. วตฺถุํ ปจฺจยา อเหตุกา เนวา\-จยคามิ\-นา\-ปจยคามี ขนฺธา. (1)}

\prose{เนวา\-จยคามิ\-นา\-ปจยคามิํ ธมฺมํ ปจฺจยา อาจยคามี ธมฺโม อุปฺปชฺชติ นเหตุปจฺจยา. วตฺถุํ ปจฺจยา วิจิกิจฺฉา\-สหคโต อุทฺธจฺจ\-สหคโต โมโห. (2)}

\prose{อาจยคามิญฺจ เนวา\-จยคามิ\-นา\-ปจยคามิญฺจ ธมฺมํ ปจฺจยา อาจยคามี ธมฺโม อุปฺปชฺชติ นเหตุปจฺจยา. วิจิกิจฺฉา\-สหคเต อุทฺธจฺจ\-สหคเต ขนฺเธ จ วตฺถุญฺจ ปจฺจยา วิจิกิจฺฉา\-สหคโต อุทฺธจฺจ\-สหคโต โมโห. (1)}

\prose{นอารมฺมณ}

\proseitem{27}{อาจยคามิํ ธมฺมํ ปจฺจยา เนวา\-จยคามิ\-นา\-ปจยคามี ธมฺโม อุปฺปชฺชติ นอารมฺมณ\-ปจฺจยา. (สํขิตฺตํ. ปฏิจฺจ\-วาร\-สทิสํ.)}

\csromanheader{2}{10. อาจยคามิตฺติก นอธิปติ}
\proseitem{28}{\csromanfolio{251}อาจยคามิํ ธมฺมํ ปจฺจยา อาจยคามี ธมฺโม อุปฺปชฺชติ นอธิปติ\-ปจฺจยา. ตีณิ.}

\prose{อปจยคามิํ ธมฺมํ ปจฺจยา อปจยคามี ธมฺโม อุปฺปชฺชติ นอธิปติ\-ปจฺจยา. อปจยคามี ขนฺเธ ปจฺจยา อปจยคามี อธิปติ. (1)}

\prose{เนวา\-จยคามิ\-นา\-ปจยคามิํ ธมฺมํ ปจฺจยา เนวา\-จยคามิ\-นา\-ปจยคามี ธมฺโม อุปฺปชฺชติ นอธิปติ\-ปจฺจยา ฯเปฯ. อสญฺญสตฺตานํ ฯเปฯ. จกฺขายตนํ ปจฺจยา จกฺขุวิญฺญาณํ ฯเปฯ. กายายตนํ ปจฺจยา กายวิญฺญาณํ. วตฺถุํ ปจฺจยา เนวา\-จยคามิ\-นา\-ปจยคามี ขนฺธา. (1)}

\prose{เนวา\-จยคามิ\-นา\-ปจยคามิํ ธมฺมํ ปจฺจยา อาจยคามี ธมฺโม อุปฺปชฺชติ นอธิปติ\-ปจฺจยา. วตฺถุํ ปจฺจยา อาจยคามี ขนฺธา. (2)}

\prose{เนวา\-จยคามิ\-นา\-ปจยคามิํ ธมฺมํ ปจฺจยา อปจยคามี ธมฺโม อุปฺปชฺชติ นอธิปติ\-ปจฺจยา. วตฺถุํ ปจฺจยา อปจยคามี อธิปติ. (3)}

\prose{เนวา\-จยคามิ\-นา\-ปจยคามิํ ธมฺมํ ปจฺจยา อาจยคามี จ เนวา\-จยคามิ\-นา\-ปจยคามี จ ธมฺมา อุปฺปชฺชนฺติ นอธิปติ\-ปจฺจยา. วตฺถุํ ปจฺจยา อาจยคามี ขนฺธา. มหาภูเต ปจฺจยา จิตฺต\-สมุฏฺฐานํ รูปํ. (4)}

\proseitem{29}{อาจยคามิญฺจ เนวา\-จยคามิ\-นา\-ปจยคามิญฺจ ธมฺมํ ปจฺจยา อาจยคามี ธมฺโม อุปฺปชฺชติ นอธิปติ\-ปจฺจยา. อาจยคามิํ เอกํ ขนฺธญฺจ วตฺถุญฺจ ปจฺจยา ตโย ขนฺธา ฯเปฯ เทฺว ขนฺเธ ฯเปฯ. (1)}

\prose{อาจยคามิญฺจ เนวา\-จยคามิ\-นา\-ปจยคามิญฺจ ธมฺมํ ปจฺจยา เนวา\-จยคามิ\-นา\-ปจยคามี ธมฺโม อุปฺปชฺชติ นอธิปติ\-ปจฺจยา. อาจยคามี ขนฺเธ จ มหาภูเต จ ปจฺจยา จิตฺต\-สมุฏฺฐานํ รูปํ. (2)}

\prose{อาจยคามิญฺจ เนวา\-จยคามิ\-นา\-ปจยคามิญฺจ ธมฺมํ ปจฺจยา อาจยคามี จ เนวา\-จยคามิ\-นา\-ปจยคามี จ ธมฺมา อุปฺปชฺชนฺติ นอธิปติ\-ปจฺจยา. อาจยคามิํ เอกํ ขนฺธญฺจ วตฺถุญฺจ ปจฺจยา ตโย ขนฺธา ฯเปฯ เทฺว ขนฺเธ ฯเปฯ. อาจยคามี ขนฺเธ จ มหาภูเต จ ปจฺจยา จิตฺต\-สมุฏฺฐานํ รูปํ. (3)}

\prose{\csromanfolio{252}อปจยคามิญฺจ เนวา\-จยคามิ\-นา\-ปจยคามิญฺจ ธมฺมํ ปจฺจยา อปจยคามี ธมฺโม อุปฺปชฺชติ นอธิปติ\-ปจฺจยา. อปจยคามี ขนฺเธ จ วตฺถุญฺจ ปจฺจยา อปจยคามี อธิปติ. (1)}

\prose{นอนนฺตราทิ}

\proseitem{30}{นอนนฺตร\-ปจฺจยา. นสมนนฺตร\-ปจฺจยา. นอญฺญมญฺญ\-ปจฺจยา. นอุปนิสฺสยปจฺจยา. นปุเรชาต\-ปจฺจยา. (ปฏิจฺจ\-วาร\-สทิสา. สตฺต ปญฺหา.) นปจฺฉาชาต\-ปจฺจยา. (ปริปุณฺณํ.)}

\prose{นอาเสวน}

\proseitem{31}{อาจยคามิํ ธมฺมํ ปจฺจยา อาจยคามี ธมฺโม อุปฺปชฺชติ นอาเสวน\-ปจฺจยา. ตีณิ.}

\prose{อปจยคามิํ ธมฺมํ ปจฺจยา เนวา\-จยคามิ\-นา\-ปจยคามี ธมฺโม อุปฺปชฺชติ นอาเสวน\-ปจฺจยา. อปจยคามี ขนฺเธ ปจฺจยา จิตฺต\-สมุฏฺฐานํ รูปํ. (1)}

\prose{เนวา\-จยคามิ\-นา\-ปจยคามิํ ธมฺมํ ปจฺจยา เนวา\-จยคามิ\-นา\-ปจยคามี ธมฺโม อุปฺปชฺชติ นอาเสวน\-ปจฺจยา. อสญฺญสตฺตานํ ฯเปฯ. จกฺขายตนํ ปจฺจยา จกฺขุวิญฺญาณํ ฯเปฯ. กายายตนํ ปจฺจยา กายวิญฺญาณํ. วตฺถุํ ปจฺจยา เนวา\-จยคามิ\-นา\-ปจยคามี ขนฺธา. (1)}

\prose{เนวา\-จยคามิ\-นา\-ปจยคามิํ ธมฺมํ ปจฺจยา อาจยคามี ธมฺโม อุปฺปชฺชติ นอาเสวน\-ปจฺจยา. วตฺถุํ ปจฺจยา อาจยคามี ขนฺธา. (2)}

\prose{เนวา\-จยคามิ\-นา\-ปจยคามิํ ธมฺมํ ปจฺจยา อาจยคามี จ เนวา\-จยคามิ\-นา\-ปจยคามี จ ธมฺมา อุปฺปชฺชนฺติ นอาเสวน\-ปจฺจยา. วตฺถุํ ปจฺจยา อาจยคามี ขนฺธา. มหาภูเต ปจฺจยา จิตฺต\-สมุฏฺฐานํ รูปํ. (3)}

\prose{อาจยคามิญฺจ เนวา\-จยคามิ\-นา\-ปจยคามิญฺจ ธมฺมํ ปจฺจยา อาจยคามี ธมฺโม อุปฺปชฺชติ นอาเสวน\-ปจฺจยา. อาจยคามิํ เอกํ ขนฺธญฺจ วตฺถุญฺจ ปจฺจยา ตโย ขนฺธา ฯเปฯ เทฺว ขนฺเธ ฯเปฯ. (1)}

\prose{อาจยคามิญฺจ เนวา\-จยคามิ\-นา\-ปจยคามิญฺจ ธมฺมํ ปจฺจยา เนวา\-จยคามิ\-นา\-ปจยคามี ธมฺโม อุปฺปชฺชติ นอาเสวน\-ปจฺจยา. อาจยคามี ขนฺเธ จ มหาภูเต จ ปจฺจยา จิตฺต\-สมุฏฺฐานํ รูปํ. (2)}

\prose{\csromanfolio{253}อาจยคามิญฺจ เนวา\-จยคามิ\-นา\-ปจยคามิญฺจ ธมฺมํ ปจฺจยา อาจยคามี จ เนวา\-จยคามิ\-นา\-ปจยคามี จ ธมฺมา อุปฺปชฺชนฺติ นอาเสวน\-ปจฺจยา. อาจยคามิํ เอกํ ขนฺธญฺจ วตฺถุญฺจ ปจฺจยา ตโย ขนฺธา ฯเปฯ เทฺว ขนฺเธ ฯเปฯ. อาจยคามี ขนฺเธ จ มหาภูเต จ ปจฺจยา จิตฺต\-สมุฏฺฐานํ รูปํ. (3)}

\prose{อปจยคามิญฺจ เนวา\-จยคามิ\-นา\-ปจยคามิญฺจ ธมฺมํ ปจฺจยา เนวา\-จยคามิ\-นา\-ปจยคามี ธมฺโม อุปฺปชฺชติ นอาเสวน\-ปจฺจยา. อปจยคามี ขนฺเธ จ มหาภูเต จ ปจฺจยา จิตฺต\-สมุฏฺฐานํ รูปํ. (1)}

\csromanheader{2}{\textbf{นกมฺม}}
\proseitem{32}{อาจยคามิํ ธมฺมํ ปจฺจยา อาจยคามี ธมฺโม อุปฺปชฺชติ นกมฺมปจฺจยา. อาจยคามี ขนฺเธ ปจฺจยา อาจยคามี เจตนา. (1)}

\prose{อปจยคามิํ ธมฺมํ ปจฺจยา อปจยคามี ธมฺโม อุปฺปชฺชติ นกมฺมปจฺจยา. อปจยคามี ขนฺเธ ปจฺจยา อปจยคามี เจตนา. (1)}

\prose{เนวา\-จยคามิ\-นา\-ปจยคามิํ ธมฺมํ ปจฺจยา เนวา\-จยคามิ\-นา\-ปจยคามี ธมฺโม อุปฺปชฺชติ นกมฺมปจฺจยา. เนวา\-จยคามิ\-นา\-ปจยคามี ขนฺเธ ปจฺจยา เนวา\-จยคามิ\-นา\-ปจยคามี เจตนา. พาหิรํ. อาหารสมุฏฺฐานํ. อุตุสมุฏฺฐานํ ฯเปฯ. วตฺถุํ ปจฺจยา เนวา\-จยคามิ\-นา\-ปจยคามี เจตนา. (1)}

\prose{เนวา\-จยคามิ\-นา\-ปจยคามิํ ธมฺมํ ปจฺจยา อาจยคามี ธมฺโม อุปฺปชฺชติ นกมฺมปจฺจยา. วตฺถุํ ปจฺจยา อาจยคามี เจตนา. (2)}

\prose{เนวา\-จยคามิ\-นา\-ปจยคามิํ ธมฺมํ ปจฺจยา อปจยคามี ธมฺโม อุปฺปชฺชติ นกมฺมปจฺจยา. วตฺถุํ ปจฺจยา อปจยคามี เจตนา. (3)}

\prose{อาจยคามิญฺจ เนวา\-จยคามิ\-นา\-ปจยคามิญฺจ ธมฺมํ ปจฺจยา อาจยคามี ธมฺโม อุปฺปชฺชติ นกมฺมปจฺจยา. อาจยคามี ขนฺเธ จ วตฺถุญฺจ ปจฺจยา อาจยคามี เจตนา. (1)}

\prose{อปจยคามิญฺจ เนวา\-จยคามิ\-นา\-ปจยคามิญฺจ ธมฺมํ ปจฺจยา อปจยคามี ธมฺโม อุปฺปชฺชติ นกมฺมปจฺจยา. อปจยคามี ขนฺเธ จ วตฺถุญฺจ ปจฺจยา อปจยคามี เจตนา. (1)}

\csromanheader{2}{\textbf{นวิปากาทิ}}
\proseitem{33}{\csromanfolio{254}อาจยคามิํ ธมฺมํ ปจฺจยา อาจยคามี ธมฺโม อุปฺปชฺชติ นวิปาก\-ปจฺจยา. (ปริปุณฺณํ กาตพฺพํ. ปฏิสนฺธิ\-กฺขเณ นตฺถิ.)}

\prose{เนวา\-จยคามิ\-นา\-ปจยคามิํ ธมฺมํ ปจฺจยา เนวา\-จยคามิ\-นา\-ปจยคามี ธมฺโม อุปฺปชฺชติ นอาหารปจฺจยา. พาหิรํ. อุตุสมุฏฺฐานํ. อสญฺญสตฺตานํ ฯเปฯ นอินฺทฺริยปจฺจยา. พาหิรํ. อาหารสมุฏฺฐานํ. อุตุสมุฏฺฐานํ. อสญฺญสตฺตานํ ฯเปฯ. มหาภูเต ปจฺจยา รูปชีวิตินฺทฺริยํ. นฌานปจฺจยา. ปญฺจวิญฺญาณํ ฯเปฯ. พาหิรํ. อาหารสมุฏฺฐานํ. อุตุสมุฏฺฐานํ. อสญฺญสตฺตานํ ฯเปฯ. จกฺขายตนํ ปจฺจยา จกฺขุวิญฺญาณํ ฯเปฯ. กายายตนํ ปจฺจยา กายวิญฺญาณํ ฯเปฯ นมคฺคปจฺจยา. อเหตุกา เนวา\-จยคามิ\-นา\-ปจยคามี ฯเปฯ. อสญฺญสตฺตานํ เอกํ มหาภูตํ ฯเปฯ. จกฺขายตนํ ปจฺจยา จกฺขุวิญฺญาณํ ฯเปฯ. กายายตนํ ปจฺจยา กายวิญฺญาณํ. วตฺถุํ ปจฺจยา อเหตุกา เนวา\-จยคามิ\-นา\-ปจยคามี ฯเปฯ นสมฺปยุตฺต\-ปจฺจยา. นวิปฺปยุตฺต\-ปจฺจยา. (ปฏิจฺจ\-วาร\-สทิสํ. ตีณิ.) โนนตฺถิปจฺจยา. โนวิคต\-ปจฺจยา.}

\csromanmarkh{2. ปจฺจย\-ปจฺจนีย}\csromanmarkhii{2. สงฺขฺยาวาร}\csromanheadernomark{2}{2. ปจฺจย\-ปจฺจนีย 2. สงฺขฺยาวาร}
\proseitemwithcloserruled{34}{นเหตุยา จตฺตาริ, นอารมฺมเณ ปญฺจ, นอธิปติยา ทฺวาทส, นอนนฺตเร ปญฺจ, นสมนนฺตเร, นอญฺญมญฺเญ, นอุปนิสฺสเย ปญฺจ, นปุเรชาเต สตฺต, นปจฺฉาชาเต สตฺตรส, นอาเสวเน เอกาทส, นกมฺเม สตฺต, นวิปาเก สตฺตรส, นอาหาเร, นอินฺทฺริเย, นฌาเน, นมคฺเค เอกํ, นสมฺปยุตฺเต ปญฺจ, นวิปฺปยุตฺเต ตีณิ, โนนตฺถิยา, โนวิคเต ปญฺจ. (เอวํ คเณตพฺพํ.)}{ปจฺจนียํ.}
\proseitem{35}{เหตุปจฺจยา นอารมฺมเณ ปญฺจ, นอธิปติยา ทฺวาทส, นอนนฺตเร, นสมนนฺตเร, นอญฺญมญฺเญ, นอุปนิสฺสเย ปญฺจ, นปุเรชาเต สตฺต,}

\prosewithcloserruled{\csromanfolio{255}นปจฺฉาชาเต สตฺตรส, นอาเสวเน เอกาทส, นกมฺเม สตฺต, นวิปาเก สตฺตรส, นสมฺปยุตฺเต ปญฺจ, นวิปฺปยุตฺเต ตีณิ, โนนตฺถิยา โนวิคเต ปญฺจ. (เอวํ คเณตพฺพํ.)}{อนุโลม\-ปจฺจนียํ.}
\proseitem{36}{นเหตุปจฺจยา อารมฺมเณ จตฺตาริ, อนนฺตเร, สมนนฺตเร, สหชาเต, อญฺญมญฺเญ, นิสฺสเย, อุปนิสฺสเย, ปุเรชาเต, อาเสวเน, กมฺเม จตฺตาริ, วิปาเก เอกํ, อาหาเร จตฺตาริ, อินฺทฺริเย, ฌาเน จตฺตาริ, มคฺเค ตีณิ, สมฺปยุตฺเต, วิปฺปยุตฺเต, อตฺถิยา, นตฺถิยา, วิคเต จตฺตาริ, อวิคเต จตฺตาริ. (เอวํ คเณตพฺพํ.) ปจฺจนียา\-นุโลมํ. ปจฺจยวาโร. (นิสฺสยวาโร ปจฺจยวาร\-สทิโส.)}
\csromansectionrule

\csromanmarkh{10. อาจยคามิตฺติก}\csromanmarkhii{5. สํสฏฺฐวาร}\csromanheadernomark{2}{10. อาจยคามิตฺติก 5. สํสฏฺฐวาร}
\csromanmarkh{1. ปจฺจยานุโลม}\csromanmarkhii{1. วิภงฺควาร}\csromanheadernomark{2}{1. \textbf{ปจฺจยานุโลม 1. วิภงฺควาร}}
\csromanheader{2}{\textbf{เหตุ}}
\proseitem{37}{อาจยคามิํ ธมฺมํ สํสฏฺโฐ อาจยคามี ธมฺโม อุปฺปชฺชติ เหตุปจฺจยา. อาจยคามิํ เอกํ ขนฺธํ สํสฏฺฐา ตโย ขนฺธา ฯเปฯ เทฺว ขนฺเธ สํสฏฺฐา เทฺว ขนฺธา. (1)}

\prose{อปจยคามิํ ธมฺมํ สํสฏฺโฐ อปจยคามี ธมฺโม อุปฺปชฺชติ เหตุปจฺจยา. อปจยคามิํ เอกํ ขนฺธํ สํสฏฺฐา ตโย ขนฺธา ฯเปฯ เทฺว ขนฺเธ ฯเปฯ. (1)}

\prose{\csromanfolio{256}เนวา\-จยคามิ\-นา\-ปจยคามิํ ธมฺมํ สํสฏฺโฐ เนวา\-จยคามิ\-นา\-ปจยคามี ธมฺโม อุปฺปชฺชติ เหตุปจฺจยา. เนวา\-จยคามิ\-นา\-ปจยคามิํ เอกํ ขนฺธํ สํสฏฺฐา ตโย ขนฺธา ฯเปฯ เทฺว ขนฺเธ ฯเปฯ. ปฏิสนฺธิ\-กฺขเณ ฯเปฯ. (1)}

\csromanheader{2}{\textbf{อารมฺมณาทิ}}
\proseitem{38}{อาจยคามิํ ธมฺมํ สํสฏฺโฐ อาจยคามี ธมฺโม อุปฺปชฺชติ อารมฺมณ\-ปจฺจยา. อธิปติ\-ปจฺจยา. อนนฺตร\-ปจฺจยา. สมนนฺตร\-ปจฺจยา. สหชาต\-ปจฺจยา. อญฺญมญฺญ\-ปจฺจยา. นิสฺสยปจฺจยา. อุปนิสฺสย\-ปจฺจยา. ปุเรชาต\-ปจฺจยา. อาเสวนปจฺจยา. กมฺมปจฺจยา. วิปากปจฺจยา. อาหารปจฺจยา. อินฺทฺริยปจฺจยา. ฌานปจฺจยา. มคฺคปจฺจยา. สมฺปยุตฺต\-ปจฺจยา. วิปฺปยุตฺต\-ปจฺจยา. อตฺถิปจฺจยา. นตฺถิปจฺจยา. วิคตปจฺจยา. อวิคตปจฺจยา.}

\csromanmarkh{1. ปจฺจยานุโลม}\csromanmarkhii{2. สงฺขฺยาวาร}\csromanheadernomark{2}{1. \textbf{ปจฺจยานุโลม 2. สงฺขฺยาวาร}}
\proseitemwithcloserruled{39}{เหตุยา ตีณิ, อารมฺมเณ, อธิปติยา, อนนฺตเร, สมนนฺตเร, สหชาเต, อญฺญมญฺเญ, นิสฺสเย, อุปนิสฺสเย, ปุเรชาเต, อาเสวเน, กมฺเม สพฺพตฺถ ตีณิ, วิปาเก เอกํ, อาหาเร ตีณิ, อินฺทฺริเย, ฌาเน, มคฺเค, สมฺปยุตฺเต, วิปฺปยุตฺเต, อตฺถิยา, นตฺถิยา, วิคเต, อวิคเต ตีณิ. (เอวํ คเณตพฺพํ.)}{อนุโลมํ.}
\csromanmarkh{2. ปจฺจย\-ปจฺจนีย}\csromanmarkhii{1. วิภงฺควาร}\csromanheadernomark{2}{2. ปจฺจย\-ปจฺจนีย 1. วิภงฺควาร}
\csromanheader{2}{\textbf{นเหตุ}}
\proseitem{40}{อาจยคามิํ ธมฺมํ สํสฏฺโฐ อาจยคามี ธมฺโม อุปฺปชฺชติ นเหตุปจฺจยา. วิจิกิจฺฉา\-สหคเต อุทฺธจฺจ\-สหคเต ขนฺเธ สํสฏฺโฐ วิจิกิจฺฉา\-สหคโต อุทฺธจฺจ\-สหคโต โมโห. (1)}

\prose{เนวา\-จยคามิ\-นา\-ปจยคามิํ ธมฺมํ สํสฏฺโฐ เนวา\-จยคามิ\-นา\-ปจยคามี ธมฺโม อุปฺปชฺชติ นเหตุปจฺจยา. อเหตุกํ เนวา\-จยคามิ\-นา\-ปจยคามิํ เอกํ ขนฺธํ สํสฏฺฐา ตโย ขนฺธา ฯเปฯ เทฺว ขนฺเธ ฯเปฯ. อเหตุก\-ปฏิสนฺธิ\-กฺขเณ ฯเปฯ. (1)}

\csromanheader{2}{10. อาจยคามิตฺติก นอธิปตฺยาทิ}
\proseitem{41}{\csromanfolio{257}อาจยคามิํ ธมฺมํ สํสฏฺโฐ อาจยคามี ธมฺโม อุปฺปชฺชติ นอธิปติ\-ปจฺจยา. นปุเรชาต\-ปจฺจยา. นปจฺฉาชาต\-ปจฺจยา. นอาเสวน\-ปจฺจยา. อาจยคามิํ เอกํ ขนฺธํ สํสฏฺฐา ตโย ขนฺธา ฯเปฯ เทฺว ขนฺเธ ฯเปฯ.}

\prose{เนวา\-จยคามิ\-นา\-ปจยคามิํ ธมฺมํ สํสฏฺโฐ เนวา\-จยคามิ\-นา\-ปจยคามี ธมฺโม อุปฺปชฺชติ นอาเสวน\-ปจฺจยา. เนวา\-จยคามิ\-นา\-ปจยคามิํ เอกํ ขนฺธํ สํสฏฺฐา ฯเปฯ. ปฏิสนฺธิ\-กฺขเณ ฯเปฯ. นกมฺมปจฺจยา. นวิปาก\-ปจฺจยา. นฌานปจฺจยา. นมคฺคปจฺจยา. นวิปฺปยุตฺต\-ปจฺจยา.}

\csromanmarkh{2. ปจฺจย\-ปจฺจนีย}\csromanmarkhii{2. สงฺขฺยาวาร}\csromanheadernomark{2}{2. ปจฺจย\-ปจฺจนีย 2. สงฺขฺยาวาร}
\proseitemwithcloserruled{42}{นเหตุยา เทฺว, นอธิปติยา ตีณิ, นปุเรชาเต ตีณิ, นปจฺฉาชาเต ตีณิ, นอาเสวเน เทฺว, นกมฺเม ตีณิ, นวิปาเก ตีณิ, นฌาเน เอกํ, นมคฺเค เอกํ, นวิปฺปยุตฺเต ตีณิ.  (เอวํ คเณตพฺพํ.)}{ปจฺจนียํ.}
\proseitemwithcloserruled{43}{เหตุปจฺจยา นอธิปติยา ตีณิ, นปุเรชาเต ตีณิ, นปจฺฉาชาเต ตีณิ, นอาเสวเน เทฺว, นกมฺเม ตีณิ, นวิปาเก ตีณิ, นวิปฺปยุตฺเต ตีณิ.  (เอวํ คเณตพฺพํ.)}{อนุโลม\-ปจฺจนียํ.}
\proseitem{44}{นเหตุปจฺจยา อารมฺมเณ เทฺว, อนนฺตเร, สมนนฺตเร, สหชาเต, อญฺญมญฺเญ, นิสฺสเย, อุปนิสฺสเย, ปุเรชาเต, อาเสวเน, \csromanfolio{258}กมฺเม สพฺพตฺถ เทฺว, วิปาเก เอกํ, อาหาเร เทฺว, อินฺทฺริเย เทฺว, ฌาเน เทฺว, มคฺเค เอกํ, สมฺปยุตฺเต เทฺว, วิปฺปยุตฺเต, อตฺถิยา, นตฺถิยา, วิคเต, อวิคเต เทฺว.  (เอวํ คเณตพฺพํ.) ปจฺจนียา\-นุโลมํ. สํสฏฺฐวาโร. (สมฺปยุตฺตวาโร สํสฏฺฐวาร\-สทิโส.)}

\csromanmarkh{10. อาจยคามิตฺติก}\csromanmarkhii{7. ปญฺหาวาร}\csromanheadernomark{2}{10. \textbf{อาจยคามิตฺติก 7. ปญฺหาวาร}}
\csromanmarkh{1. ปจฺจยานุโลม}\csromanmarkhii{1. วิภงฺควาร}\csromanheadernomark{2}{1. \textbf{ปจฺจยานุโลม 1. วิภงฺควาร}}
\csromanheader{2}{\textbf{เหตุ}}
\proseitem{45}{อาจยคามี\csromansharedfootnote{dc6672d3ba6513ec}{อาจยคามิ (สี, สฺยา) เอวมุปริปิ.} ธมฺโม อาจยคามิสฺส ธมฺมสฺส เหตุปจฺจเยน ปจฺจโย. อาจยคามี เหตู สมฺปยุตฺตกานํ ขนฺธานํ เหตุปจฺจเยน ปจฺจโย. (1)}

\prose{อาจยคามี ธมฺโม เนวา\-จยคามิ\-นา\-ปจยคามิสฺส ธมฺมสฺส เหตุปจฺจเยน ปจฺจโย. อาจยคามี เหตู จิตฺต\-สมุฏฺฐานานํ รูปานํ เหตุปจฺจเยน ปจฺจโย. (2)}

\prose{อาจยคามี ธมฺโม อาจยคามิสฺส จ เนวา\-จยคามิ\-นา\-ปจยคามิสฺส จ ธมฺมสฺส เหตุปจฺจเยน ปจฺจโย. อาจยคามี เหตู สมฺปยุตฺถกานํ ขนฺธานํ จิตฺตสมุฏฺฐานานญฺจ รูปานํ เหตุปจฺจเยน ปจฺจโย. (3)}

\prose{อปจยคามี ธมฺโม อปจยคามิสฺส ธมฺมสฺส เหตุปจฺจเยน ปจฺจโย. ตีณิ.}

\prose{เนวา\-จยคามิ\-นา\-ปจยคามี ธมฺโม เนวา\-จยคามิ\-นา\-ปจยคามิสฺส ธมฺมสฺส เหตุปจฺจเยน ปจฺจโย. เนวา\-จยคามิ\-นา\-ปจยคามี เหตู สมฺปยุตฺตกานํ ขนฺธานํ จิตฺตสมุฏฺฐานานญฺจ รูปานํ เหตุปจฺจเยน ปจฺจโย. ปฏิสนฺธิ\-กฺขเณ เนวา\-จยคามิ\-นา\-ปจยคามี เหตู สมฺปยุตฺตกานํ ขนฺธานํ กฏตฺตา จ รูปานํ เหตุปจฺจเยน ปจฺจโย. (1)}

\csromanheader{2}{10. อาจยคามิตฺติก \textbf{อารมฺมณ}}
\proseitem{46}{\csromanfolio{259}อาจยคามี ธมฺโม อาจยคามิสฺส ธมฺมสฺส อารมฺมณ\-ปจฺจเยน ปจฺจโย. ทานํ ทตฺวา, สีลํ สมาทิยิตฺวา, อุโปสถกมฺมํ กตฺวา ตํ ปจฺจเวกฺขติ, ปุพฺเพ สุจิณฺณานิ ปจฺจเวกฺขติ, ฌานา วุฏฺฐหิตฺวา ฌานํ ปจฺจเวกฺขติ. เสกฺขา ปหีเน กิเลเส ปจฺจเวกฺขนฺติ, วิกฺขมฺภิเต กิเลเส ปจฺจเวกฺขนฺติ, ปุพฺเพ สมุทาจิณฺเณ กิเลเส ชานนฺติ. เสกฺขา วา ปุถุชฺชนา วา อาจยคามี ขนฺเธ อนิจฺจโต ทุกฺขโต อนตฺตโต วิปสฺสนฺติ, อสฺสาเทนฺติ อภินนฺทนฺติ, ตํ อารพฺภ ราโค อุปฺปชฺชติ. ทิฏฺฐิ ฯเปฯ วิจิกิจฺฉา ฯเปฯ อุทฺธจฺจํ ฯเปฯ โทมนสฺสํ อุปฺปชฺชติ. เจโตปริย\-ญาเณน อาจยคามิ\-จิตฺต\-สมงฺคิสฺส จิตฺตํ ชานนฺติ. อากาสา\-นญฺจา\-ยตนกุสลํ วิญฺญาณญฺจา\-ยตนกุสลสฺส อารมฺมณ\-ปจฺจเยน ปจฺจโย. อากิญฺจญฺญายตน\-กุสลํ เนว\-สญฺญา\-นา\-สญฺญา\-ยตนกุสลสฺส ฯเปฯ. อาจยคามี ขนฺธา อิทฺธิวิธ\-ญาณสฺส, เจโตปริย\-ญาณสฺส, ปุพฺเพ\-นิวาสา\-นุสฺสติ\-ญาณสฺส, ยถากมฺมูปคญาณสฺส, อนาคตํสญาณสฺส อารมฺมณ\-ปจฺจเยน ปจฺจโย. (1)}

\prose{อาจยคามี ธมฺโม เนวา\-จยคามิ\-นา\-ปจยคามิสฺส ธมฺมสฺส อารมฺมณ\-ปจฺจเยน ปจฺจโย. อรหา ปหีเน กิเลเส ปจฺจเวกฺขติ, ปุพฺเพ สมุทาจิณฺเณ กิเลเส ชานาติ. อาจยคามี ขนฺเธ อนิจฺจโต ทุกฺขโต อนตฺตโต วิปสฺสติ. เจโตปริย\-ญาเณน อาจยคามิ\-จิตฺต\-สมงฺคิสฺส จิตฺตํ ชานาติ. เสกฺขา วา ปุถุชฺชนา วา อาจยคามี ขนฺเธ อนิจฺจโต ทุกฺขโต อนตฺตโต วิปสฺสนฺติ. กุสเล นิรุทฺเธ วิปาโก ตทารมฺมณตา อุปฺปชฺชติ. อาจยคามี ขนฺเธ อสฺสาเทติ อภินนฺทติ, ตํ อารพฺภ ราโค อุปฺปชฺชติ ฯเปฯ โทมนสฺสํ อุปฺปชฺชติ. อกุสเล นิรุทฺเธ วิปาโก ตทารมฺมณตา อุปฺปชฺชติ. อากาสา\-นญฺจา\-ยตนกุสลํ วิญฺญาณญฺจา\-ยตนวิปากสฺส จ กิริยสฺส จ อารมฺมณ\-ปจฺจเยน ปจฺจโย. อากิญฺจญฺญายตน\-กุสลํ เนว\-สญฺญา\-นา\-สญฺญา\-ยตนวิปากสฺส จ กิริยสฺส จ อารมฺมณ\-ปจฺจเยน ปจฺจโย. อาจยคามี ขนฺธา เจโตปริย\-ญาณสฺส, ปุพฺเพ\-นิวาสา\-นุสฺสติ\-ญาณสฺส, ยถากมฺมูปคญาณสฺส, อนาคตํสญาณสฺส, อาวชฺชนาย อารมฺมณ\-ปจฺจเยน ปจฺจโย. (2)}

\proseitem{47}{อปจยคามี ธมฺโม อาจยคามิสฺส ธมฺมสฺส อารมฺมณ\-ปจฺจเยน ปจฺจโย. เสกฺขา มคฺคา วุฏฺฐหิตฺวา มคฺคํ ปจฺจเวกฺขนฺติ. เจโตปริย\-ญาเณน \csromanfolio{260}อปจยคามิ\-จิตฺต\-สมงฺคิสฺส จิตฺตํ ชานนฺติ. อปจยคามี ขนฺธา เจโตปริย\-ญาณสฺส, ปุพฺเพ\-นิวาสา\-นุสฺสติ\-ญาณสฺส, อนาคตํสญาณสฺส อารมฺมณ\-ปจฺจเยน ปจฺจโย. (1)}

\prose{อปจยคามี ธมฺโม เนวา\-จยคามิ\-นา\-ปจยคามิสฺส ธมฺมสฺส อารมฺมณ\-ปจฺจเยน ปจฺจโย. อรหา มคฺคา วุฏฺฐหิตฺวา มคฺคํ ปจฺจเวกฺขติ. เจโตปริย\-ญาเณน อปจยคามิ\-จิตฺต\-สมงฺคิสฺส จิตฺตํ ชานาติ. อปจยคามี ขนฺธา เจโตปริย\-ญาณสฺส, ปุพฺเพ\-นิวาสา\-นุสฺสติ\-ญาณสฺส, อนาคตํสญาณสฺส, อาวชฺชนาย อารมฺมณ\-ปจฺจเยน ปจฺจโย. (2)}

\proseitem{48}{เนวา\-จยคามิ\-นา\-ปจยคามี ธมฺโม เนวา\-จยคามิ\-นา\-ปจยคามิสฺส ธมฺมสฺส อารมฺมณ\-ปจฺจเยน ปจฺจโย. อรหา ผลํ ปจฺจเวกฺขติ, นิพฺพานํ ปจฺจเวกฺขติ. นิพฺพานํ ผลสฺส อาวชฺชนาย อารมฺมณ\-ปจฺจเยน ปจฺจโย. อรหา จกฺขุํ อนิจฺจโต ทุกฺขโต อนตฺตโต วิปสฺสติ. โสตํ ฯเปฯ. วตฺถุํ. เนวา\-จยคามิ\-นา\-ปจยคามี ขนฺเธ อนิจฺจโต ทุกฺขโต อนตฺตโต วิปสฺสติ. ทิพฺเพน จกฺขุนา รูปํ ปสฺสติ. ทิพฺพาย โสตธาตุยา สทฺทํ สุณาติ. เจโตปริย\-ญาเณน เนวาจยคามินาปจยคามิจิตฺตสมงฺคิสฺส จิตฺตํ ชานาติ. อากาสา\-นญฺจา\-ยตนกิริยํ วิญฺญาณญฺจา\-ยตนกิริยสฺส อารมฺมณ\-ปจฺจเยน ปจฺจโย. อากิญฺจญฺญายตน\-กิริยํ เนว\-สญฺญา\-นา\-สญฺญา\-ยตนกิริยสฺส ฯเปฯ. รูปายตนํ จกฺขุ\-วิญฺญาณสฺส ฯเปฯ. โผฏฺฐพฺพา\-ยตนํ กายวิญฺญาณสฺส ฯเปฯ. เนวา\-จยคามิ\-นา\-ปจยคามี ขนฺธา อิทฺธิวิธ\-ญาณสฺส เจโตปริย\-ญาณสฺส, ปุพฺเพ\-นิวาสา\-นุสฺสติ\-ญาณสฺส, อนาคตํสญาณสฺส, อาวชฺชนาย อารมฺมณ\-ปจฺจเยน ปจฺจโย. (1)}

\prose{เนวา\-จยคามิ\-นา\-ปจยคามี ธมฺโม อาจยคามิสฺส ธมฺมสฺส อารมฺมณ\-ปจฺจเยน ปจฺจโย. เสกฺขา ผลํ ปจฺจเวกฺขนฺติ, นิพฺพานํ ปจฺจเวกฺขนฺติ. นิพฺพานํ โคตฺรภุสฺส, โวทานสฺส อารมฺมณ\-ปจฺจเยน ปจฺจโย. เสกฺขา วา ปุถุชฺชนา วา จกฺขุํ อนิจฺจโต ทุกฺขโต อนตฺตโต วิปสฺสนฺติ, อสฺสาเทนฺติ อภินนฺทนฺติ, ตํ อารพฺภ ราโค อุปฺปชฺชติ ฯเปฯ โทมนสฺสํ อุปฺปชฺชติ. โสตํ ฯเปฯ. วตฺถุํ. เนวา\-จยคามิ\-นา\-ปจยคามี ขนฺเธ อนิจฺจโต ทุกฺขโต อนตฺตโต วิปสฺสนฺติ, อสฺสาเทนฺติ อภินนฺทนฺติ, ตํ อารพฺภ ราโค อุปฺปชฺชติ. ทิฏฺฐิ อุปฺปชฺชติ. วิจิกิจฺฉา ฯเปฯ อุทฺธจฺจํ ฯเปฯ โทมนสฺสํ ฯเปฯ. ทิพฺเพน จกฺขุนา รูปํ ปสฺสนฺติ. ทิพฺพาย โสตธาตุยา สทฺทํ สุณนฺติ. \csromanfolio{261}เจโตปริย\-ญาเณน เนวาจยคามินาปจยคามิจิตฺตสมงฺคิสฺส จิตฺตํ ชานนฺติ. เนวา\-จยคามิ\-นา\-ปจยคามี ขนฺธา อิทฺธิวิธ\-ญาณสฺส, เจโตปริย\-ญาณสฺส, ปุพฺเพ\-นิวาสา\-นุสฺสติ\-ญาณสฺส, อนาคตํสญาณสฺส อารมฺมณ\-ปจฺจเยน ปจฺจโย. (2)}

\prose{เนวา\-จยคามิ\-นา\-ปจยคามี ธมฺโม อปจยคามิสฺส ธมฺมสฺส อารมฺมณ\-ปจฺจเยน ปจฺจโย. นิพฺพานํ มคฺคสฺส อารมฺมณ\-ปจฺจเยน ปจฺจโย. (3)}

\proseitem{49}{อาจยคามี ธมฺโม อาจยคามิสฺส ธมฺมสฺส อธิปติ\-ปจฺจเยน ปจฺจโย. อารมฺมณา\-ธิปติ, สหชาตา\-ธิปติ.  อารมฺมณา\-ธิปติ\csromandash{}ทานํ ทตฺวา, สีลํ สมาทิยิตฺวา, อุโปสถกมฺมํ กตฺวา ตํ ครุํ กตฺวา ปจฺจเวกฺขติ, ปุพฺเพ สุจิณฺณานิ ครุํ กตฺวา ปจฺจเวกฺขติ, ฌานา วุฏฺฐหิตฺวา ฌานํ ครุํ กตฺวา ปจฺจเวกฺขติ. อาจยคามี ขนฺเธ ครุํ กตฺวา อสฺสาเทติ อภินนฺทติ, ตํ ครุํ กตฺวา ราโค อุปฺปชฺชติ, ทิฏฺฐิ อุปฺปชฺชติ. สหชาตา\-ธิปติ\csromandash{}อาจยคามี อธิปติ สมฺปยุตฺตกานํ ขนฺธานํ อธิปติ\-ปจฺจเยน ปจฺจโย. (1)}

\prose{อาจยคามี ธมฺโม เนวา\-จยคามิ\-นา\-ปจยคามิสฺส ธมฺมสฺส อธิปติ\-ปจฺจเยน ปจฺจโย. สหชาตา\-ธิปติ\csromandash{}อาจยคามี อธิปติ จิตฺต\-สมุฏฺฐานานํ รูปานํ อธิปติ\-ปจฺจเยน ปจฺจโย. (2)}

\prose{อาจยคามี ธมฺโม อาจยคามิสฺส จ เนวา\-จยคามิ\-นา\-ปจยคามิสฺส จ ธมฺมสฺส อธิปติ\-ปจฺจเยน ปจฺจโย. สหชาตา\-ธิปติ\csromandash{}อาจยคามี อธิปติ สมฺปยุตฺตกานํ ขนฺธานํ จิตฺตสมุฏฺฐานานญฺจ รูปานํ อธิปติ\-ปจฺจเยน ปจฺจโย. (3)}

\proseitem{50}{อปจยคามี ธมฺโม อปจยคามิสฺส ธมฺมสฺส อธิปติ\-ปจฺจเยน ปจฺจโย. สหชาตา\-ธิปติ\csromandash{}อปจยคามี อธิปติ สมฺปยุตฺตกานํ ขนฺธานํ อธิปติ\-ปจฺจเยน ปจฺจโย. (1)}

\prose{อปจยคามี ธมฺโม อาจยคามิสฺส ธมฺมสฺส อธิปติ\-ปจฺจเยน ปจฺจโย. อารมฺมณา\-ธิปติ\csromandash{}เสกฺขา มคฺคา วุฏฺฐหิตฺวา มคฺคํ ครุํ กตฺวา ปจฺจเวกฺขนฺติ. (2)}

\prose{\csromanfolio{262}อปจยคามี ธมฺโม เนวา\-จยคามิ\-นา\-ปจยคามิสฺส ธมฺมสฺส อธิปติ\-ปจฺจเยน ปจฺจโย. อารมฺมณา\-ธิปติ. สหชาตา\-ธิปติ. อารมฺมณา\-ธิปติ\csromandash{}อรหา มคฺคา วุฏฺฐหิตฺวา มคฺคํ ครุํ กตฺวา ปจฺจเวกฺขติ. สหชาตา\-ธิปติ\csromandash{}อปจยคามี อธิปติ จิตฺต\-สมุฏฺฐานานํ รูปานํ อธิปติ\-ปจฺจเยน ปจฺจโย. (3)}

\prose{อปจยคามี ธมฺโม อปจยคามิสฺส จ เนวา\-จยคามิ\-นา\-ปจยคามิสฺส จ ธมฺมสฺส อธิปติ\-ปจฺจเยน ปจฺจโย. สหชาตา\-ธิปติ\csromandash{}อปจยคามี อธิปติ สมฺปยุตฺตกานํ ขนฺธานํ จิตฺตสมุฏฺฐานานญฺจ รูปานํ อธิปติ\-ปจฺจเยน ปจฺจโย. (4)}

\proseitem{51}{เนวา\-จยคามิ\-นา\-ปจยคามี ธมฺโม เนวา\-จยคามิ\-นา\-ปจยคามิสฺส ธมฺมสฺส อธิปติ\-ปจฺจเยน ปจฺจโย. อารมฺมณา\-ธิปติ, สหชาตา\-ธิปติ.  อารมฺมณา\-ธิปติ\csromandash{}อรหา ผลํ ครุํ กตฺวา ปจฺจเวกฺขติ, นิพฺพานํ ครุํ กตฺวา ปจฺจเวกฺขติ. นิพฺพานํ ผลสฺส อธิปติ\-ปจฺจเยน ปจฺจโย. สหชาตา\-ธิปติ\csromandash{} เนวา\-จยคามิ\-นา\-ปจยคามี อธิปติ สมฺปยุตฺตกานํ ขนฺธานํ จิตฺตสมุฏฺฐานานญฺจ รูปานํ อธิปติ\-ปจฺจเยน ปจฺจโย. (1)}

\prose{เนวา\-จยคามิ\-นา\-ปจยคามี ธมฺโม อาจยคามิสฺส ธมฺมสฺส อธิปติ\-ปจฺจเยน ปจฺจโย. อารมฺมณา\-ธิปติ\csromandash{}เสกฺขา ผลํ ครุํ กตฺวา ปจฺจเวกฺขนฺติ, นิพฺพานํ ครุํ กตฺวา ปจฺจเวกฺขนฺติ. นิพฺพานํ โคตฺรภุสฺส, โวทานสฺส อธิปติ\-ปจฺจเยน ปจฺจโย. จกฺขุํ ครุํ กตฺวา อสฺสาเทติ ฯเปฯ. วตฺถุํ. เนวา\-จยคามิ\-นา\-ปจยคามี ขนฺเธ ครุํ กตฺวา อสฺสาเทติ อภินนฺทติ, ตํ ครุํ กตฺวา ราโค อุปฺปชฺชติ. ทิฏฺฐิ อุปฺปชฺชติ. (2)}

\prose{เนวา\-จยคามิ\-นา\-ปจยคามี ธมฺโม อปจยคามิสฺส ธมฺมสฺส อธิปติ\-ปจฺจเยน ปจฺจโย. อารมฺมณา\-ธิปติ\csromandash{}นิพฺพานํ มคฺคสฺส อธิปติ\-ปจฺจเยน ปจฺจโย. (3)}

\csromanheader{2}{\textbf{อนนฺตร}}
\proseitem{52}{อาจยคามี ธมฺโม อาจยคามิสฺส ธมฺมสฺส อนนฺตร\-ปจฺจเยน ปจฺจโย. ปุริมา ปุริมา อาจยคามี ขนฺธา ปจฺฉิมานํ ปจฺฉิมานํ อาจยคามีนํ ขนฺธานํ อนนฺตร\-ปจฺจเยน ปจฺจโย. อนุโลมํ โคตฺรภุสฺส. อนุโลมํ โวทานสฺส อนนฺตร\-ปจฺจเยน ปจฺจโย. (1)}

\prose{\csromanfolio{263}อาจยคามี ธมฺโม อปจยคามิสฺส ธมฺมสฺส อนนฺตร\-ปจฺจเยน ปจฺจโย. โคตฺรภุ มคฺคสฺส. โวทานํ มคฺคสฺส อนนฺตร\-ปจฺจเยน ปจฺจโย. (2)}

\prose{อาจยคามี ธมฺโม เนวา\-จยคามิ\-นา\-ปจยคามิสฺส ธมฺมสฺส อนนฺตร\-ปจฺจเยน ปจฺจโย. อาจยคามี ขนฺธา วุฏฺฐานสฺส อนนฺตร\-ปจฺจเยน ปจฺจโย. เสกฺขานํ อนุโลมํ ผลสมาปตฺติยา. นิโรธา วุฏฺฐหนฺตสฺส เนว\-สญฺญา\-นา\-สญฺญา\-ยตนกุสลํ ผลสมาปตฺติยา อนนฺตร\-ปจฺจเยน ปจฺจโย. (3)}

\prose{อปจยคามี ธมฺโม เนวา\-จยคามิ\-นา\-ปจยคามิสฺส ธมฺมสฺส อนนฺตร\-ปจฺจเยน ปจฺจโย. มคฺโค ผลสฺส อนนฺตร\-ปจฺจเยน ปจฺจโย. (1)}

\prose{เนวา\-จยคามิ\-นา\-ปจยคามี ธมฺโม เนวา\-จยคามิ\-นา\-ปจยคามิสฺส ธมฺมสฺส อนนฺตร\-ปจฺจเยน ปจฺจโย. ปุริมา ปุริมา เนวา\-จยคามิ\-นา\-ปจยคามี ขนฺธา ปจฺฉิมานํ ปจฺฉิมานํ เนวา\-จยคามิ\-นา\-ปจยคามีนํ ขนฺธานํ อนนฺตร\-ปจฺจเยน ปจฺจโย. ภวงฺคํ อาวชฺชนาย. กิริยํ วุฏฺฐานสฺส. อรหโต อนุโลมํ ผลสมาปตฺติยา. นิโรธา วุฏฺฐหนฺตสฺส เนว\-สญฺญา\-นา\-สญฺญา\-ยตนกิริยํ ผลสมาปตฺติยา อนนฺตร\-ปจฺจเยน ปจฺจโย. (1)}

\prose{เนวา\-จยคามิ\-นา\-ปจยคามี ธมฺโม อาจยคามิสฺส ธมฺมสฺส อนนฺตร\-ปจฺจเยน ปจฺจโย. อาวชฺชนา อาจยคามีนํ ขนฺธานํ อนนฺตร\-ปจฺจเยน ปจฺจโย. (2)}

\csromanheader{2}{\textbf{สมนนฺตราทิ}}
\proseitem{53}{อาจยคามี ธมฺโม อาจยคามิสฺส ธมฺมสฺส สมนนฺตร\-ปจฺจเยน ปจฺจโย ฯเปฯ. (อนนฺตร\-สทิสํ.) (สหชาต\-ปจฺจเย ปฏิจฺจวาเร สหชาต\-วาร\-สทิสา นว ปญฺหา. อญฺญมญฺญ\-ปจฺจเย ปฏิจฺจวาเร อญฺญมญฺญ\-สทิสํ. ตีณิ. นิสฺสยปจฺจเย ปจฺจยวาเร นิสฺสย\-วาร\-สทิสํ. จตฺตาริปิ หิ วิสุํ ฆฏนา นตฺถิ. เตรส ปญฺหา.)}

\csromanheader{2}{\textbf{อุปนิสฺสย}}
\proseitem{54}{อาจยคามี ธมฺโม อาจยคามิสฺส ธมฺมสฺส อุปนิสฺสย\-ปจฺจเยน ปจฺจโย. อารมฺมณู\-ปนิสฺสโย, อนนฺตรู\-ปนิสฺสโย, ปกตู\-ปนิสฺสโย ฯเปฯ. ปกตู\-ปนิสฺสโย\csromandash{}อาจยคามิํ สทฺธํ อุปนิสฺสาย ทานํ เทติ. สีลํ ฯเปฯ อุโปสถกมฺมํ ฯเปฯ ฌานํ ฯเปฯ วิปสฺสนํ ฯเปฯ อภิญฺญํ ฯเปฯ \csromanfolio{264}สมาปตฺติํ อุปฺปาเทติ, มานํ ชปฺเปติ, ทิฏฺฐิํ คณฺหาติ, อาจยคามิํ สีลํ. สุตํ. จาคํ. ปญฺญํ. ราคํ. โทสํ. โมหํ. มานํ. ทิฏฺฐิํ. ปตฺถนํ อุปนิสฺสาย ทานํ เทติ. สีลํ ฯเปฯ อุโปสถกมฺมํ ฯเปฯ ฌานํ ฯเปฯ วิปสฺสนํ ฯเปฯ อภิญฺญํ ฯเปฯ สมาปตฺติํ ฯเปฯ ปาณํ หนติ ฯเปฯ สํฆํ ภินฺทติ. อาจยคามี สทฺธา ฯเปฯ. ปญฺญา. ราโค ฯเปฯ. ปตฺถนา อาจยคามิยา สทฺธาย ฯเปฯ ปญฺญาย. ราคสฺส ฯเปฯ ปตฺถนาย อุปนิสฺสย\-ปจฺจเยน ปจฺจโย. ปฐมสฺส ฌานสฺส ปริกมฺมํ ปฐมสฺส ฌานสฺส อุปนิสฺสย\-ปจฺจเยน ปจฺจโย ฯเปฯ. เนว\-สญฺญา\-นา\-สญฺญา\-ยตนสฺส ปริกมฺมํ เนว\-สญฺญา\-นา\-สญฺญา\-ยตนสฺส อุปนิสฺสย\-ปจฺจเยน ปจฺจโย. ปฐมํ ฌานํ ทุติยสฺส ฌานสฺส ฯเปฯ. อากิญฺจญฺญา\-ยตนํ เนว\-สญฺญา\-นา\-สญฺญา\-ยตนสฺส อุปนิสฺสย\-ปจฺจเยน ปจฺจโย. (1)}

\prose{อาจยคามี ธมฺโม อปจยคามิสฺส ธมฺมสฺส อุปนิสฺสย\-ปจฺจเยน ปจฺจโย. อนนฺตรู\-ปนิสฺสโย, ปกตู\-ปนิสฺสโย ฯเปฯ. ปกตู\-ปนิสฺสโย\csromandash{} ปฐมสฺส มคฺคสฺส ปริกมฺมํ ปฐมสฺส มคฺคสฺส ฯเปฯ. จตุตฺถสฺส มคฺคสฺส ปริกมฺมํ จตุตฺถสฺส มคฺคสฺส อุปนิสฺสย\-ปจฺจเยน ปจฺจโย. (2)}

\prose{อาจยคามี ธมฺโม เนวา\-จยคามิ\-นา\-ปจยคามิสฺส ธมฺมสฺส อุปนิสฺสย\-ปจฺจเยน ปจฺจโย. อนนฺตรู\-ปนิสฺสโย, ปกตู\-ปนิสฺสโย ฯเปฯ. ปกตู\-ปนิสฺสโย\csromandash{}อาจยคามิํ สทฺธํ อุปนิสฺสาย อตฺตานํ อาตาเปติ ปริตาเปติ, ปริยิฏฺฐิมูลกํ ทุกฺขํ ปจฺจนุโภติ. อาจยคามิํ สีลํ ฯเปฯ ปญฺญํ. ราคํ ฯเปฯ. ปตฺถนํ อุปนิสฺสาย อตฺตานํ อาตาเปติ ปริตาเปติ, ปริยิฏฺฐิมูลกํ ทุกฺขํ ปจฺจนุโภติ. อาจยคามี สทฺธา ฯเปฯ ปญฺญา. ราโค ฯเปฯ ปตฺถนา, กายิกสฺส สุขสฺส, กายิกสฺส ทุกฺขสฺส, ผลสมาปตฺติยา อุปนิสฺสย\-ปจฺจเยน ปจฺจโย. กุสลากุสลํ กมฺมํ วิปากสฺส อุปนิสฺสย\-ปจฺจเยน ปจฺจโย. (3)}

\proseitem{55}{อปจยคามี ธมฺโม อปจยคามิสฺส ธมฺมสฺส อุปนิสฺสย ปจฺจเยน ปจฺจโย. ปกตู\-ปนิสฺสโย\csromandash{}ปฐโม มคฺโค ทุติยสฺส มคฺคสฺส ฯเปฯ. ตติโย มคฺโค จตุตฺถสฺส มคฺคสฺส อุปนิสฺสย\-ปจฺจเยน ปจฺจโย. (1)}

\prose{อปจยคามี ธมฺโม อาจยคามิสฺส ธมฺมสฺส อุปนิสฺสย\-ปจฺจเยน ปจฺจโย. อารมฺมณู\-ปนิสฺสโย, ปกตู\-ปนิสฺสโย ฯเปฯ. ปกตู\-ปนิสฺสโย\csromandash{} เสกฺขา มคฺคํ อุปนิสฺสาย อนุปฺปนฺนํ กุสล\-สมาปตฺติํ อุปฺปาเทนฺติ, อุปฺปนฺนํ \csromanfolio{265}สมาปชฺชนฺติ, สงฺขาเร อนิจฺจโต ทุกฺขโต อนตฺตโต วิปสฺสนฺติ, มคฺโค เสกฺขานํ อตฺถ\-ปฏิสมฺภิทาย ฯเปฯ. ปฏิภาน\-ปฏิสมฺภิทาย. ฐานาฐาน\-โกสลฺลสฺส อุปนิสฺสย\-ปจฺจเยน ปจฺจโย. (2)}

\prose{อปจยคามี ธมฺโม เนวา\-จยคามิ\-นา\-ปจยคามิสฺส ธมฺมสฺส อุปนิสฺสย\-ปจฺจเยน ปจฺจโย. อารมฺมณู\-ปนิสฺสโย, อนนฺตรู\-ปนิสฺสโย, ปกตู\-ปนิสฺสโย ฯเปฯ. ปกตู\-ปนิสฺสโย\csromandash{}อรหา มคฺคํ อุปนิสฺสาย อนุปฺปนฺนํ กิริย\-สมาปตฺติํ อุปฺปาเทติ, อุปฺปนฺนํ สมาปชฺชติ ฯเปฯ. ฐานาฐาน\-โกสลฺลสฺส อุปนิสฺสย\-ปจฺจเยน ปจฺจโย. มคฺโค ผลสมาปตฺติยา อุปนิสฺสย\-ปจฺจเยน ปจฺจโย. (3)}

\proseitem{56}{เนวา\-จยคามิ\-นา\-ปจยคามี ธมฺโม เนวา\-จยคามิ\-นา\-ปจยคามิสฺส ธมฺมสฺส อุปนิสฺสย\-ปจฺจเยน ปจฺจโย. อารมฺมณู\-ปนิสฺสโย, อนนฺตรู\-ปนิสฺสโย, ปกตู\-ปนิสฺสโย ฯเปฯ. ปกตู\-ปนิสฺสโย\csromandash{}กายิกํ สุขํ อุปนิสฺสาย อตฺตานํ อาตาเปติ ปริตาเปติ, ปริยิฏฺฐิมูลกํ ทุกฺขํ ปจฺจนุโภติ. กายิกํ ทุกฺขํ. อุตุํ. โภชนํ. เสนาสนํ อุปนิสฺสาย อตฺตานํ อาตาเปติ ปริตาเปติ ฯเปฯ. กายิกํ สุขํ. กายิกํ ทุกฺขํ. อุตุ. โภชนํ. เสนาสนํ กายิกสฺส สุขสฺส, กายิกสฺส ทุกฺขสฺส, ผลสมาปตฺติยา อุปนิสฺสย\-ปจฺจเยน ปจฺจโย. อรหา กายิกํ สุขํ อุปนิสฺสาย อนุปฺปนฺนํ กิริย\-สมาปตฺติํ อุปฺปาเทติ ฯเปฯ วิปสฺสติ. กายิกํ ทุกฺขํ. อุตุํ. โภชนํ. เสนาสนํ อุปนิสฺสาย ฯเปฯ วิปสฺสติ. (1)}

\prose{เนวา\-จยคามิ\-นา\-ปจยคามี ธมฺโม อาจยคามิสฺส ธมฺมสฺส อุปนิสฺสย\-ปจฺจเยน ปจฺจโย. อารมฺมณู\-ปนิสฺสโย, อนนฺตรู\-ปนิสฺสโย, ปกตู\-ปนิสฺสโย ฯเปฯ. ปกตู\-ปนิสฺสโย\csromandash{}กายิกํ สุขํ อุปนิสฺสาย ทานํ เทติ ฯเปฯ สมาปตฺติํ อุปฺปาเทติ, ปาณํ หนติ ฯเปฯ สํฆํ ภินฺทติ. กายิกํ ทุกฺขํ. อุตุํ. โภชนํ. เสนาสนํ อุปนิสฺสาย ทานํ เทติ ฯเปฯ สํฆํ ภินฺทติ, กายิกํ สุขํ ฯเปฯ. เสนาสนํ อาจยคามิยา สทฺธาย ฯเปฯ ปญฺญาย. ราคสฺส ฯเปฯ ปตฺถนาย อุปนิสฺสย\-ปจฺจเยน ปจฺจโย. (2)}

\prose{เนวา\-จยคามิ\-นา\-ปจยคามี ธมฺโม อปจยคามิสฺส ธมฺมสฺส อุปนิสฺสย\-ปจฺจเยน ปจฺจโย. อารมฺมณู\-ปนิสฺสโย, ปกตู\-ปนิสฺสโย ฯเปฯ. ปกตู\-ปนิสฺสโย\csromandash{}กายิกํ สุขํ อุปนิสฺสาย มคฺคํ อุปฺปาเทติ. กายิกํ \csromanfolio{266}ทุกฺขํ ฯเปฯ. เสนาสนํ อุปนิสฺสาย มคฺคํ อุปฺปาเทติ, กายิกํ สุขํ, กายิกํ ทุกฺขํ ฯเปฯ. เสนาสนํ มคฺคสฺส อุปนิสฺสย\-ปจฺจเยน ปจฺจโย. (3)}

\csromanheader{2}{\textbf{ปุเรชาต}}
\proseitem{57}{เนวา\-จยคามิ\-นา\-ปจยคามี ธมฺโม เนวา\-จยคามิ\-นา\-ปจยคามิสฺส ธมฺมสฺส ปุเรชาต\-ปจฺจเยน ปจฺจโย. อารมฺมณ\-ปุเรชาตํ, วตฺถุ\-ปุเรชาตํ.  อารมฺมณ\-ปุเรชาตํ\csromandash{}อรหา จกฺขุํ ฯเปฯ วตฺถุํ อนิจฺจโต ทุกฺขโต อนตฺตโต วิปสฺสติ. ทิพฺเพน จกฺขุนา รูปํ ปสฺสติ. ทิพฺพาย โสตธาตุยา สทฺทํ สุณาติ. รูปายตนํ จกฺขุ\-วิญฺญาณสฺส ฯเปฯ. โผฏฺฐพฺพา\-ยตนํ กายวิญฺญาณสฺส ปุเรชาต\-ปจฺจเยน ปจฺจโย. วตฺถุ\-ปุเรชาตํ\csromandash{}จกฺขายตนํ จกฺขุ\-วิญฺญาณสฺส ฯเปฯ. กายายตนํ กายวิญฺญาณสฺส ฯเปฯ. วตฺถุ เนวา\-จยคามิ\-นา\-ปจยคามีนํ ขนฺธานํ ปุเรชาต\-ปจฺจเยน ปจฺจโย. (1)}

\prose{เนวา\-จยคามิ\-นา\-ปจยคามี ธมฺโม อาจยคามิสฺส ธมฺมสฺส ปุเรชาต\-ปจฺจเยน ปจฺจโย. อารมฺมณ\-ปุเรชาตํ, วตฺถุ\-ปุเรชาตํ.  อารมฺมณ\-ปุเรชาตํ\csromandash{}เสกฺขา วา ปุถุชฺชนา วา จกฺขุํ ฯเปฯ. วตฺถุํ อนิจฺจโต ทุกฺขโต อนตฺตโต วิปสฺสนฺติ, อสฺสาเทนฺติ อภินนฺทนฺติ, ตํ อารพฺภ ราโค อุปฺปชฺชติ ฯเปฯ โทมนสฺสํ อุปฺปชฺชติ. ทิพฺเพน จกฺขุนา รูปํ ปสฺสติ. ทิพฺพาย โสตธาตุยา สทฺทํ สุณาติ. วตฺถุ\-ปุเรชาตํ\csromandash{}วตฺถุ อาจยคามีนํ ขนฺธานํ ปุเรชาต\-ปจฺจเยน ปจฺจโย. (2)}

\prose{เนวา\-จยคามิ\-นา\-ปจยคามี ธมฺโม อปจยคามิสฺส ธมฺมสฺส ปุเรชาต\-ปจฺจเยน ปจฺจโย. วตฺถุ\-ปุเรชาตํ\csromandash{}วตฺถุ อปจยคามีนํ ขนฺธานํ ปุเรชาต\-ปจฺจเยน ปจฺจโย. (3)}

\csromanheader{2}{\textbf{ปจฺฉาชาต}}
\proseitem{58}{อาจยคามี ธมฺโม เนวา\-จยคามิ\-นา\-ปจยคามิสฺส ธมฺมสฺส ปจฺฉาชาต\-ปจฺจเยน ปจฺจโย. ปจฺฉาชาตา อาจยคามี ขนฺธา ปุเรชาตสฺส อิมสฺส กายสฺส ปจฺฉาชาต\-ปจฺจเยน ปจฺจโย. (1)}

\prose{อปจยคามี ธมฺโม เนวา\-จยคามิ\-นา\-ปจยคามิสฺส ธมฺมสฺส ปจฺฉาชาต\-ปจฺจเยน ปจฺจโย. ปจฺฉาชาตา อปจยคามี ขนฺธา ปุเรชาตสฺส อิมสฺส กายสฺส ปจฺฉาชาต\-ปจฺจเยน ปจฺจโย. (1)}

\prose{\csromanfolio{267}เนวา\-จยคามิ\-นา\-ปจยคามี ธมฺโม เนวา\-จยคามิ\-นา\-ปจยคามิสฺส ธมฺมสฺส ปจฺฉาชาต\-ปจฺจเยน ปจฺจโย. ปจฺฉาชาตา เนวา\-จยคามิ\-นา\-ปจยคามี ขนฺธา ปุเรชาตสฺส อิมสฺส กายสฺส ปจฺฉาชาต\-ปจฺจเยน ปจฺจโย.}

\csromanheader{2}{\textbf{อาเสวน}}
\proseitem{59}{อาจยคามี ธมฺโม อาจยคามิสฺส ธมฺมสฺส อาเสวน\-ปจฺจเยน ปจฺจโย. ปุริมา ปุริมา อาจยคามี ขนฺธา ปจฺฉิมานํ ปจฺฉิมานํ อาจยคามีนํ ขนฺธานํ อาเสวน\-ปจฺจเยน ปจฺจโย. อนุโลมํ โคตฺรภุสฺส. อนุโลมํ โวทานสฺส อาเสวน\-ปจฺจเยน ปจฺจโย. (1)}

\prose{อาจยคามี ธมฺโม อปจยคามิสฺส ธมฺมสฺส อาเสวน\-ปจฺจเยน ปจฺจโย. โคตฺรภุ มคฺคสฺส. โวทานํ มคฺคสฺส อาเสวน\-ปจฺจเยน ปจฺจโย. (2)}

\prose{เนวา\-จยคามิ\-นา\-ปจยคามี ธมฺโม เนวา\-จยคามิ\-นา\-ปจยคามิสฺส ธมฺมสฺส อาเสวน\-ปจฺจเยน ปจฺจโย. ปุริมา ปุริมา เนวา\-จยคามิ\-นา\-ปจยคามี ขนฺธา ปจฺฉิมานํ ปจฺฉิมานํ เนวา\-จยคามิ\-นา\-ปจยคามีนํ ขนฺธานํ อาเสวน\-ปจฺจเยน ปจฺจโย. (1)}

\csromanheader{2}{\textbf{กมฺม}}
\proseitem{60}{อาจยคามี ธมฺโม อาจยคามิสฺส ธมฺมสฺส กมฺมปจฺจเยน ปจฺจโย. อาจยคามี เจตนา สมฺปยุตฺตกานํ ขนฺธานํ กมฺมปจฺจเยน ปจฺจโย. (1)}

\prose{อาจยคามี ธมฺโม เนวา\-จยคามิ\-นา\-ปจยคามิสฺส ธมฺมสฺส กมฺมปจฺจเยน ปจฺจโย. สหชาตา, นานากฺขณิกา.  สหชาตา อาจยคามี เจตนา จิตฺต\-สมุฏฺฐานานํ รูปานํ กมฺมปจฺจเยน ปจฺจโย. นานากฺขณิกา อาจยคามี เจตนา วิปากานํ ขนฺธานํ กฏตฺตา จ รูปานํ กมฺมปจฺจเยน ปจฺจโย. (2)}

\prose{อาจยคามี ธมฺโม อาจยคามิสฺส จ เนวา\-จยคามิ\-นา\-ปจยคามิสฺส จ ธมฺมสฺส กมฺมปจฺจเยน ปจฺจโย. อาจยคามี เจตนา สมฺปยุตฺตกานํ ขนฺธานํ จิตฺตสมุฏฺฐานานญฺจ รูปานํ กมฺมปจฺจเยน ปจฺจโย. (3)}

\proseitem{61}{อปจยคามี ธมฺโม อปจยคามิสฺส ธมฺมสฺส กมฺมปจฺจเยน ปจฺจโย. อปจยคามี เจตนา สมฺปยุตฺตกานํ ขนฺธานํ กมฺมปจฺจเยน ปจฺจโย. (1)}

\prose{\csromanfolio{268}อปจยคามี ธมฺโม เนวา\-จยคามิ\-นา\-ปจยคามิสฺส ธมฺมสฺส กมฺมปจฺจเยน ปจฺจโย. สหชาตา, นานากฺขณิกา.  สหชาตา อปจยคามี เจตนา จิตฺต\-สมุฏฺฐานานํ รูปานํ กมฺมปจฺจเยน ปจฺจโย. นานากฺขณิกา อปจยคามี เจตนา วิปากานํ ขนฺธานํ กมฺมปจฺจเยน ปจฺจโย. (2)}

\prose{อปจยคามี ธมฺโม อปจยคามิสฺส จ เนวา\-จยคามิ\-นา\-ปจยคามิสฺส จ ธมฺมสฺส กมฺมปจฺจเยน ปจฺจโย. อปจยคามี เจตนา สมฺปยุตฺตกานํ ขนฺธานํ จิตฺตสมุฏฺฐานานญฺจ รูปานํ กมฺมปจฺจเยน ปจฺจโย. (3)}

\prose{เนวา\-จยคามิ\-นา\-ปจยคามี ธมฺโม เนวา\-จยคามิ\-นา\-ปจยคามิสฺส ธมฺมสฺส กมฺมปจฺจเยน ปจฺจโย. เนวา\-จยคามิ\-นา\-ปจยคามี เจตนา สมฺปยุตฺตกานํ ขนฺธานํ จิตฺตสมุฏฺฐานานญฺจ รูปานํ กมฺมปจฺจเยน ปจฺจโย. ปฏิสนฺธิ\-กฺขเณ เนวา\-จยคามิ\-นา\-ปจยคามี เจตนา สมฺปยุตฺตกานํ ขนฺธานํ กฏตฺตา จ รูปานํ กมฺมปจฺจเยน ปจฺจโย. (1)}

\csromanheader{2}{\textbf{วิปาก}}
\proseitem{62}{เนวา\-จยคามิ\-นา\-ปจยคามี ธมฺโม เนวา\-จยคามิ\-นา\-ปจยคามิสฺส ธมฺมสฺส วิปาก\-ปจฺจเยน ปจฺจโย. วิปาโก เนวา\-จยคามิ\-นา\-ปจยคามี เอโก ขนฺโธ ติณฺณนฺนํ ขนฺธานํ ฯเปฯ. ปฏิสนฺธิ\-กฺขเณ ฯเปฯ. ขนฺธา วตฺถุสฺส วิปาก\-ปจฺจเยน ปจฺจโย. (1)}

\csromanheader{2}{\textbf{อาหาราทิ}}
\proseitem{63}{อาจยคามี ธมฺโม อาจยคามิสฺส ธมฺมสฺส อาหารปจฺจเยน ปจฺจโย. อินฺทฺริย\-ปจฺจเยน ปจฺจโย. ฌานปจฺจเยน ปจฺจโย. มคฺคปจฺจเยน ปจฺจโย. สมฺปยุตฺต\-ปจฺจเยน ปจฺจโย.}

\csromanheader{2}{\textbf{วิปฺปยุตฺต}}
\proseitem{64}{อาจยคามี ธมฺโม เนวา\-จยคามิ\-นา\-ปจยคามิสฺส ธมฺมสฺส วิปฺปยุตฺต\-ปจฺจเยน ปจฺจโย. สหชาตํ, ปจฺฉาชาตํ.  สหชาตา อาจยคามี ขนฺธา จิตฺต\-สมุฏฺฐานานํ รูปานํ วิปฺปยุตฺต\-ปจฺจเยน ปจฺจโย. ปจฺฉาชาตา อาจยคามี ขนฺธา ปุเรชาตสฺส อิมสฺส กายสฺส วิปฺปยุตฺต\-ปจฺจเยน ปจฺจโย. (1)}

\prose{\csromanfolio{269}อปจยคามี ธมฺโม เนวา\-จยคามิ\-นา\-ปจยคามิสฺส ธมฺมสฺส วิปฺปยุตฺต\-ปจฺจเยน ปจฺจโย. สหชาตํ, ปจฺฉาชาตํ. สหชาตา อปจยคามี ขนฺธา จิตฺต\-สมุฏฺฐานานํ รูปานํ วิปฺปยุตฺต\-ปจฺจเยน ปจฺจโย. ปจฺฉาชาตา อปจยคามี ขนฺธา ปุเรชาตสฺส อิมสฺส กายสฺส วิปฺปยุตฺต\-ปจฺจเยน ปจฺจโย. (1)}

\prose{เนวา\-จยคามิ\-นา\-ปจยคามี ธมฺโม เนวา\-จยคามิ\-นา\-ปจยคามิสฺส ธมฺมสฺส วิปฺปยุตฺต\-ปจฺจเยน ปจฺจโย. สหชาตํ, ปุเรชาตํ, ปจฺฉาชาตํ.  สหชาตา เนวา\-จยคามิ\-นา\-ปจยคามี ขนฺธา จิตฺต\-สมุฏฺฐานานํ รูปานํ วิปฺปยุตฺต\-ปจฺจเยน ปจฺจโย. ปฏิสนฺธิ\-กฺขเณ เนวา\-จยคามิ\-นา\-ปจยคามี ขนฺธา กฏตฺตารูปานํ วิปฺปยุตฺต\-ปจฺจเยน ปจฺจโย. ขนฺธา วตฺถุสฺส วิปฺปยุตฺต\-ปจฺจเยน ปจฺจโย. วตฺถุ ขนฺธานํ วิปฺปยุตฺต\-ปจฺจเยน ปจฺจโย. ปุเรชาตํ จกฺขายตนํ จกฺขุ\-วิญฺญาณสฺส ฯเปฯ. กายายตนํ กายวิญฺญาณสฺส ฯเปฯ. วตฺถุ เนวา\-จยคามิ\-นา\-ปจยคามีนํ ขนฺธานํ วิปฺปยุตฺต\-ปจฺจเยน ปจฺจโย. ปจฺฉาชาตา เนวา\-จยคามิ\-นา\-ปจยคามี ขนฺธา ปุเรชาตสฺส อิมสฺส กายสฺส ฯเปฯ. (1)}

\prose{เนวา\-จยคามิ\-นา\-ปจยคามี ธมฺโม อาจยคามิสฺส ธมฺมสฺส วิปฺปยุตฺต\-ปจฺจเยน ปจฺจโย. ปุเรชาตํ วตฺถุ อาจยคามีนํ ขนฺธานํ วิปฺปยุตฺต\-ปจฺจเยน ปจฺจโย. (2)}

\prose{เนวา\-จยคามิ\-นา\-ปจยคามี ธมฺโม อปจยคามิสฺส ธมฺมสฺส วิปฺปยุตฺต\-ปจฺจเยน ปจฺจโย. ปุเรชาตํ วตฺถุ อปจยคามีนํ ขนฺธานํ วิปฺปยุตฺต\-ปจฺจเยน ปจฺจโย. (3)}

\prose{อตฺถิอาทิ}

\proseitem{65}{อาจยคามี ธมฺโม อาจยคามิสฺส ธมฺมสฺส อตฺถิปจฺจเยน ปจฺจโย. อาจยคามี เอโก ขนฺโธ ติณฺณนฺนํ ขนฺธานํ ฯเปฯ. (1)}

\prose{อาจยคามี ธมฺโม เนวา\-จยคามิ\-นา\-ปจยคามิสฺส ธมฺมสฺส อตฺถิปจฺจเยน ปจฺจโย. สหชาตํ, ปจฺฉาชาตํ.  สหชาตา อาจยคามี ขนฺธา จิตฺต\-สมุฏฺฐานานํ รูปานํ อตฺถิปจฺจเยน ปจฺจโย. ปจฺฉาชาตา อาจยคามี ขนฺธา ปุเรชาตสฺส อิมสฺส กายสฺส อตฺถิปจฺจเยน ปจฺจโย. (2)}

\prose{\csromanfolio{270}อาจยคามี ธมฺโม อาจยคามิสฺส จ เนวา\-จยคามิ\-นา\-ปจยคามิสฺส จ ธมฺมสฺส อตฺถิปจฺจเยน ปจฺจโย. อาจยคามี เอโก ขนฺโธ ติณฺณนฺนํ ขนฺธานํ จิตฺตสมุฏฺฐานานญฺจ รูปานํ อตฺถิปจฺจเยน ปจฺจโย ฯเปฯ เทฺว ขนฺธา ฯเปฯ. (3)}

\prose{อปจยคามี ธมฺโม. ตีณิ. (อาจยคามิ\-นเยน กาตพฺพํ.)}

\proseitem{66}{เนวา\-จยคามิ\-นา\-ปจยคามี ธมฺโม เนวา\-จยคามิ\-นา\-ปจยคามิสฺส ธมฺมสฺส อตฺถิปจฺจเยน ปจฺจโย. สหชาตํ, ปุเรชาตํ, ปจฺฉาชาตํ, อาหารํ, อินฺทฺริยํ.  สหชาโต เนวา\-จยคามิ\-นา\-ปจยคามี เอโก ขนฺโธ ติณฺณนฺนํ ขนฺธานํ จิตฺตสมุฏฺฐานานญฺจ รูปานํ อตฺถิปจฺจเยน ปจฺจโย ฯเปฯ เทฺว ขนฺธา ฯเปฯ. ปฏิสนฺธิ\-กฺขเณ ฯเปฯ. ขนฺธา วตฺถุสฺส อตฺถิปจฺจเยน ปจฺจโย. วตฺถุ ขนฺธานํ อตฺถิปจฺจเยน ปจฺจโย. เอกํ มหาภูตํ ฯเปฯ. พาหิรํ. อาหารสมุฏฺฐานํ. อุตุสมุฏฺฐานํ. อสญฺญสตฺตานํ ฯเปฯ. ปุเรชาตํ อรหา จกฺขุํ ฯเปฯ วตฺถุํ อนิจฺจโต ทุกฺขโต อนตฺตโต วิปสฺสติ. ทิพฺเพน จกฺขุนา รูปํ ปสฺสติ. ทิพฺพาย โสตธาตุยา สทฺทํ สุณาติ. รูปายตนํ จกฺขุ\-วิญฺญาณสฺส ฯเปฯ. โผฏฺฐพฺพา\-ยตนํ กายวิญฺญาณสฺส ฯเปฯ. จกฺขายตนํ จกฺขุ\-วิญฺญาณสฺส ฯเปฯ. กายายตนํ กายวิญฺญาณสฺส ฯเปฯ. วตฺถุ เนวา\-จยคามิ\-นา\-ปจยคามีนํ ขนฺธานํ อตฺถิปจฺจเยน ปจฺจโย. ปจฺฉาชาตา เนวา\-จยคามิ\-นา\-ปจยคามี ขนฺธา ปุเรชาตสฺส อิมสฺส กายสฺส อตฺถิปจฺจเยน ปจฺจโย. กพฬีกาโร อาหาโร อิมสฺส กายสฺส ฯเปฯ. รูปชีวิตินฺทฺริยํ กฏตฺตารูปานํ ฯเปฯ. (1)}

\prose{เนวา\-จยคามิ\-นา\-ปจยคามี ธมฺโม อาจยคามิสฺส ธมฺมสฺส อตฺถิปจฺจเยน ปจฺจโย. ปุเรชาตํ เสกฺขา วา ปุถุชฺชนา วา จกฺขุํ อนิจฺจโต ทุกฺขโต อนตฺตโต วิปสฺสนฺติ, อสฺสาเทนฺติ อภินนฺทนฺติ, ตํ อารพฺภ ราโค ฯเปฯ โทมนสฺสํ อุปฺปชฺชติ, โสตํ ฯเปฯ วตฺถุํ อนิจฺจโต ฯเปฯ วิปสฺสนฺติ, อสฺสาเทนฺติ อภินนฺทนฺติ, ตํ อารพฺภ ราโค ฯเปฯ โทมนสฺสํ อุปฺปชฺชติ. ทิพฺเพน จกฺขุนา รูปํ ปสฺสติ. ทิพฺพาย โสตธาตุยา สทฺทํ สุณาติ. วตฺถุ อาจยคามีนํ ขนฺธานํ อตฺถิปจฺจเยน ปจฺจโย. (2)}

\prose{เนวา\-จยคามิ\-นา\-ปจยคามี ธมฺโม อปจยคามิสฺส ธมฺมสฺส อตฺถิปจฺจเยน ปจฺจโย. ปุเรชาตํ วตฺถุ อปจยคามีนํ ขนฺธานํ อตฺถิปจฺจเยน ปจฺจโย. (3)}

\proseitem{67}{\csromanfolio{271}อาจยคามี จ เนวา\-จยคามิ\-นา\-ปจยคามี จ ธมฺมา อาจยคามิสฺส ธมฺมสฺส อตฺถิปจฺจเยน ปจฺจโย. สหชาตํ, ปุเรชาตํ. สหชาโต อาจยคามี เอโก ขนฺโธ จ วตฺถุ จ ติณฺณนฺนํ ขนฺธานํ อตฺถิปจฺจเยน ปจฺจโย ฯเปฯ เทฺว ขนฺธา ฯเปฯ. (1)}

\prose{อาจยคามี จ เนวา\-จยคามิ\-นา\-ปจยคามี จ ธมฺมา เนวา\-จยคามิ\-นา\-ปจยคามิสฺส ธมฺมสฺส อตฺถิปจฺจเยน ปจฺจโย. สหชาตํ, ปจฺฉาชาตํ, อาหารํ, อินฺทฺริยํ.  สหชาตา อาจยคามี ขนฺธา จ มหาภูตา จ จิตฺต\-สมุฏฺฐานานํ รูปานํ อตฺถิปจฺจเยน ปจฺจโย. ปจฺฉาชาตา อาจยคามี ขนฺธา จ กพฬีกาโร อาหาโร จ อิมสฺส กายสฺส อตฺถิปจฺจเยน ปจฺจโย. ปจฺฉาชาตา อาจยคามี ขนฺธา จ รูปชีวิตินฺทฺริยญฺจ กฏตฺตารูปานํ อตฺถิปจฺจเยน ปจฺจโย. (2)}

\prose{อปจยคามี จ เนวา\-จยคามิ\-นา\-ปจยคามี จ ธมฺมา อปจยคามิสฺส ธมฺมสฺส อตฺถิปจฺจเยน ปจฺจโย. (เทฺว กาตพฺพา ทสฺสิตนเยน.) นตฺถิ\-ปจฺจเยน ปจฺจโย, วิคต\-ปจฺจเยน ปจฺจโย, อวิคต\-ปจฺจเยน ปจฺจโย.}

\csromanmarkh{1. ปจฺจยานุโลม}\csromanmarkhii{2. สงฺขฺยาวาร}\csromanheadernomark{2}{1. \textbf{ปจฺจยานุโลม 2. สงฺขฺยาวาร}}
\proseitemwithcloserruled{68}{เหตุยา สตฺต, อารมฺมเณ สตฺต, อธิปติยา ทส, อนนฺตเร ฉ, สมนนฺตเร ฉ, สหชาเต นว, อญฺญมญฺเญ ตีณิ, นิสฺสเย เตรส, อุปนิสฺสเย นว, ปุเรชาเต ตีณิ, ปจฺฉาชาเต ตีณิ, อาเสวเน ตีณิ, กมฺเม สตฺต, วิปาเก เอกํ, อาหาเร สตฺต, อินฺทฺริเย สตฺต, ฌาเน สตฺต, มคฺเค สตฺต, สมฺปยุตฺเต ตีณิ, วิปฺปยุตฺเต ปญฺจ, อตฺถิยา เตรส, นตฺถิยา ฉ, วิคเต ฉ, อวิคเต เตรส.}{อนุโลมํ.}
\csromanheader{2}{2. ปจฺจนียุ\-ทฺธาร}
\proseitem{69}{อาจยคามี ธมฺโม อาจยคามิสฺส ธมฺมสฺส อารมฺมณ\-ปจฺจเยน ปจฺจโย. สหชาต\-ปจฺจเยน ปจฺจโย. อุปนิสฺสย\-ปจฺจเยน ปจฺจโย. (1)}

\prose{\csromanfolio{272}อาจยคามี ธมฺโม อปจยคามิสฺส ธมฺมสฺส อุปนิสฺสย\-ปจฺจเยน ปจฺจโย. (2)}

\prose{อาจยคามี ธมฺโม เนวา\-จยคามิ\-นา\-ปจยคามิสฺส ธมฺมสฺส อารมฺมณ\-ปจฺจเยน ปจฺจโย. สหชาต\-ปจฺจเยน ปจฺจโย. อุปนิสฺสย\-ปจฺจเยน ปจฺจโย. ปจฺฉาชาต\-ปจฺจเยน ปจฺจโย. กมฺมปจฺจเยน ปจฺจโย. (3)}

\prose{อาจยคามี ธมฺโม อาจยคามิสฺส จ เนวา\-จยคามิ\-นา\-ปจยคามิสฺส จ ธมฺมสฺส สหชาต\-ปจฺจเยน ปจฺจโย. (4)}

\proseitem{70}{อปจยคามี ธมฺโม อปจยคามิสฺส ธมฺมสฺส สหชาต\-ปจฺจเยน ปจฺจโย. อุปนิสฺสย\-ปจฺจเยน ปจฺจโย. (1)}

\prose{อปจยคามี ธมฺโม อาจยคามิสฺส ธมฺมสฺส อารมฺมณ\-ปจฺจเยน ปจฺจโย. อุปนิสฺสย\-ปจฺจเยน ปจฺจโย. (2)}

\prose{อปจยคามี ธมฺโม เนวา\-จยคามิ\-นา\-ปจยคามิสฺส ธมฺมสฺส อารมฺมณ\-ปจฺจเยน ปจฺจโย. สหชาต\-ปจฺจเยน ปจฺจโย. อุปนิสฺสย\-ปจฺจเยน ปจฺจโย. ปจฺฉาชาต\-ปจฺจเยน ปจฺจโย. (3)}

\prose{อปจยคามี ธมฺโม อปจยคามิสฺส จ เนวา\-จยคามิ\-นา\-ปจยคามิสฺส จ ธมฺมสฺส สหชาต\-ปจฺจเยน ปจฺจโย. (4)}

\proseitem{71}{เนวา\-จยคามิ\-นา\-ปจยคามี ธมฺโม เนวา\-จยคามิ\-นา\-ปจยคามิสฺส ธมฺมสฺส อารมฺมณ\-ปจฺจเยน ปจฺจโย. สหชาต\-ปจฺจเยน ปจฺจโย. อุปนิสฺสย\-ปจฺจเยน ปจฺจโย. ปุเรชาต\-ปจฺจเยน ปจฺจโย. ปจฺฉาชาต\-ปจฺจเยน ปจฺจโย. อาหารปจฺจเยน ปจฺจโย. อินฺทฺริย\-ปจฺจเยน ปจฺจโย. (1)}

\prose{เนวา\-จยคามิ\-นา\-ปจยคามี ธมฺโม อาจยคามิสฺส ธมฺมสฺส อารมฺมณ\-ปจฺจเยน ปจฺจโย. อุปนิสฺสย\-ปจฺจเยน ปจฺจโย. ปุเรชาต\-ปจฺจเยน ปจฺจโย. (2)}

\prose{เนวา\-จยคามิ\-นา\-ปจยคามี ธมฺโม อปจยคามิสฺส ธมฺมสฺส อุปนิสฺสย\-ปจฺจเยน ปจฺจโย. ปุเรชาต\-ปจฺจเยน ปจฺจโย. (3)}

\proseitem{72}{\csromanfolio{273}อาจยคามี จ เนวา\-จยคามิ\-นา\-ปจยคามี จ ธมฺมา อาจยคามิสฺส ธมฺมสฺส สหชาตํ, ปุเรชาตํ. (1)}

\prose{อาจยคามี จ เนวา\-จยคามิ\-นา\-ปจยคามี จ ธมฺมา เนวา\-จยคามิ\-นา\-ปจยคามิสฺส ธมฺมสฺส สหชาตํ, ปจฺฉาชาตํ, อาหารํ, อินฺทฺริยํ. (2)}

\prose{อปจยคามี จ เนวา\-จยคามิ\-นา\-ปจยคามี จ ธมฺมา อปจยคามิสฺส ธมฺมสฺส สหชาตํ, ปุเรชาตํ. (1)}

\prose{อปจยคามี จ เนวา\-จยคามิ\-นา\-ปจยคามี จ ธมฺมา เนวา\-จยคามิ\-นา\-ปจยคามิสฺส ธมฺมสฺส สหชาตํ, ปจฺฉาชาตํ, อาหารํ, อินฺทฺริยํ. (2)}

\csromanmarkh{2. ปจฺจย\-ปจฺจนีย}\csromanmarkhii{2. สงฺขฺยาวาร}\csromanheadernomark{2}{2. ปจฺจย\-ปจฺจนีย 2. สงฺขฺยาวาร}
\proseitemwithcloserruled{73}{นเหตุยา ปนฺนรส, นอารมฺมเณ, นอธิปติยา, นอนนฺตเร, นสมนนฺตเร ปนฺนรส, นสหชาเต เอกาทส, นอญฺญมญฺเญ เอกาทส, นนิสฺสเย เอกาทส, นอุปนิสฺสเย จุทฺทส, นปุเรชาเต เตรส, นปจฺฉาชาเต ปนฺนรส, นอาเสวเน, นกมฺเม, นวิปาเก, นอาหาเร, นอินฺทฺริเย, นฌาเน, นมคฺเค ปนฺนรส, นสมฺปยุตฺเต เอกาทส, นวิปฺปยุตฺเต นว, โนอตฺถิยา นว, โนนตฺถิยา ปนฺนรส, โนวิคเต ปนฺนรส, โนอวิคเต นว. (เอวํ คเณตพฺพํ.)}{ปจฺจนียํ.}
\proseitem{74}{เหตุปจฺจยา นอารมฺมเณ สตฺต, นอธิปติยา, นอนนฺตเร, นสมนนฺตเร สตฺต, นอญฺญมญฺเญ ตีณิ, นอุปนิสฺสเย สตฺต, นปุเรชาเต, นปจฺฉาชาเต, นอาเสวเน, นกมฺเม, นวิปาเก, นอาหาเร, นอินฺทฺริเย, นฌาเน, นมคฺเค สตฺต, นสมฺปยุตฺเต ตีณิ, นวิปฺปยุตฺเต ตีณิ, โนนตฺถิยา สตฺต, โนวิคเต สตฺต. (เอวํ คเณตพฺพํ.)}

\vspace{\tipitakaparskip}
\csromancenter{อนุโลม\-ปจฺจนียํ.}
\proseitem{75}{\csromanfolio{274}นเหตุปจฺจยา อารมฺมเณ สตฺต, อธิปติยา ทส, อนนฺตเร ฉ, สมนนฺตเร ฉ, สหชาเต นว, อญฺญมญฺเญ ตีณิ, นิสฺสเย เตรส, อุปนิสฺสเย นว, ปุเรชาเต ตีณิ, ปจฺฉาชาเต ตีณิ, อาเสวเน ตีณิ, กมฺเม สตฺต, วิปาเก เอกํ, อาหาเร สตฺต, อินฺทฺริเย, ฌาเน, มคฺเค สตฺต, สมฺปยุตฺเต ตีณิ, วิปฺปยุตฺเต ปญฺจ, อตฺถิยา เตรส, นตฺถิยา ฉ, วิคเต ฉ, อวิคเต เตรส. (เอวํ คเณตพฺพํ.)}

\vspace{\tipitakaparskip}
\csromancenter{ปจฺจนียา\-นุโลมํ.}
\nitthitam{อาจยคามิ\-ตฺติกํ นิฏฺฐิตํ.}
\csromanmarkh{11. เสกฺขตฺติก}\csromanmarkhii{11. เสกฺขตฺติก 1. ปฏิจฺจวาร}\csromanheadernomark{2}{11. เสกฺขตฺติก \textbf{11. เสกฺขตฺติก 1. ปฏิจฺจวาร}}
\csromanmarkh{1. ปจฺจยานุโลม}\csromanmarkhii{1. วิภงฺควาร}\csromanheadernomark{2}{1. \textbf{ปจฺจยานุโลม 1. วิภงฺควาร}}
\csromanheader{2}{\textbf{เหตุ}}
\proseitem{1}{\csromanfolio{275}เสกฺขํ ธมฺมํ ปฏิจฺจ เสกฺโข ธมฺโม อุปฺปชฺชติ เหตุปจฺจยา. เสกฺขํ เอกํ ขนฺธํ ปฏิจฺจ ตโย ขนฺธา ฯเปฯ เทฺว ขนฺเธ ฯเปฯ. (1)}

\prose{เสกฺขํ ธมฺมํ ปฏิจฺจ เนว\-เสกฺข\-นา\-เสกฺโข\csromansharedfootnote{5bbbe1cdee105b6a}{เนวเสขานาเสโข (สี), เนวเสกฺขานาเสกฺโข (สฺยา, ก)} ธมฺโม อุปฺปชฺชติ เหตุปจฺจยา. เสกฺเข ขนฺเธ ปฏิจฺจ จิตฺต\-สมุฏฺฐานํ รูปํ. (2)}

\prose{เสกฺขํ ธมฺมํ ปฏิจฺจ เสกฺโข จ เนว\-เสกฺข\-นา\-เสกฺโข จ ธมฺมา อุปฺปชฺชนฺติ เหตุปจฺจยา. เสกฺขํ เอกํ ขนฺธํ ปฏิจฺจ ตโย ขนฺธา จิตฺต\-สมุฏฺฐานญฺจ รูปํ ฯเปฯ เทฺว ขนฺเธ ฯเปฯ. (3)}

\proseitem{2}{อเสกฺขํ ธมฺมํ ปฏิจฺจ อเสกฺโข ธมฺโม อุปฺปชฺชติ เหตุปจฺจยา. อเสกฺขํ เอกํ ขนฺธํ ปฏิจฺจ ตโย ขนฺธา ฯเปฯ. (1)}

\prose{อเสกฺขํ ธมฺมํ ปฏิจฺจ เนว\-เสกฺข\-นา\-เสกฺโข ธมฺโม อุปฺปชฺชติ เหตุปจฺจยา. อเสกฺเข ขนฺเธ ปฏิจฺจ จิตฺต\-สมุฏฺฐานํ รูปํ. (2)}

\prose{อเสกฺขํ ธมฺมํ ปฏิจฺจ อเสกฺโข จ เนว\-เสกฺข\-นา\-เสกฺโข จ ธมฺมา ฯเปฯ. อเสกฺขํ เอกํ ขนฺธํ ปฏิจฺจ ตโย ขนฺธา จิตฺต\-สมุฏฺฐานญฺจ รูปํ ฯเปฯ เทฺว ขนฺเธ ฯเปฯ. (3)}

\proseitem{3}{เนว\-เสกฺข\-นาเสกฺขํ ธมฺมํ ปฏิจฺจ เนว\-เสกฺข\-นา\-เสกฺโข ธมฺโม อุปฺปชฺชติ เหตุปจฺจยา. เนว\-เสกฺข\-นาเสกฺขํ เอกํ ขนฺธํ ปฏิจฺจ ตโย ขนฺธา จิตฺต\-สมุฏฺฐานญฺจ รูปํ ฯเปฯ เทฺว ขนฺเธ ฯเปฯ. ปฏิสนฺธิ\-กฺขเณ ฯเปฯ. ขนฺเธ ปฏิจฺจ วตฺถุ, วตฺถุํ ปฏิจฺจ ขนฺธา. เอกํ มหาภูตํ ปฏิจฺจ ตโย มหาภูตา ฯเปฯ เทฺว มหาภูเต ปฏิจฺจ เทฺว มหาภูตา, มหาภูเต ปฏิจฺจ จิตฺต\-สมุฏฺฐานํ รูปํ, กฏตฺตารูปํ, อุปาทารูปํ. (1)}

\prose{\csromanfolio{276}เสกฺขญฺจ เนว\-เสกฺข\-นาเสกฺขญฺจ ธมฺมํ ปฏิจฺจ เนว\-เสกฺข\-นา\-เสกฺโข ธมฺโม อุปฺปชฺชติ เหตุปจฺจยา. เสกฺเข ขนฺเธ จ มหาภูเต จ ปฏิจฺจ จิตฺต\-สมุฏฺฐานํ รูปํ. (1)}

\prose{อเสกฺขญฺจ เนว\-เสกฺข\-นาเสกฺขญฺจ ธมฺมํ ปฏิจฺจ เนว\-เสกฺข\-นา\-เสกฺโข ธมฺโม อุปฺปชฺชติ เหตุปจฺจยา. อเสกฺเข ขนฺเธ จ มหาภูเต จ ปฏิจฺจ จิตฺต\-สมุฏฺฐานํ รูปํ. (1)}

\csromanheader{2}{\textbf{อารมฺมณาทิ}}
\proseitem{4}{เสกฺขํ ธมฺมํ ปฏิจฺจ เสกฺโข ธมฺโม อุปฺปชฺชติ อารมฺมณ\-ปจฺจยา. อธิปติ\-ปจฺจยา. (ปฏิสนฺธิ นตฺถิ.) อนนฺตร\-ปจฺจยา. สมนนฺตร\-ปจฺจยา. สหชาต\-ปจฺจยา. (สพฺเพ มหาภูตา กาตพฺพา.) อญฺญมญฺญ\-ปจฺจยา. นิสฺสยปจฺจยา. อุปนิสฺสย\-ปจฺจยา. ปุเรชาต\-ปจฺจยา. อาเสวนปจฺจยา. เสกฺขํ เอกํ ขนฺธํ ปฏิจฺจ ตโย ขนฺธา ฯเปฯ เทฺว ขนฺเธ ฯเปฯ.}

\prose{เนว\-เสกฺข\-นาเสกฺขํ ธมฺมํ ปฏิจฺจ เนว\-เสกฺข\-นา\-เสกฺโข ธมฺโม อุปฺปชฺชติ อาเสวนปจฺจยา. เนว\-เสกฺข\-นาเสกฺขํ เอกํ ขนฺธํ ปฏิจฺจ ตโย ขนฺธา ฯเปฯ เทฺว ขนฺเธ ฯเปฯ. กมฺมปจฺจยา. วิปากปจฺจยา. วิปากํ เสกฺขํ เอกํ ขนฺธํ ปฏิจฺจ ตโย ขนฺธา ฯเปฯ เทฺว ขนฺเธ ฯเปฯ. (ตีณิ. ปริปุณฺณํ.)}

\prose{อเสกฺขํ ธมฺมํ ปฏิจฺจ อเสกฺโข ธมฺโม อุปฺปชฺชติ วิปากปจฺจยา. อเสกฺขํ เอกํ ขนฺธํ ปฏิจฺจ. ตีณิ.}

\prose{เนว\-เสกฺข\-นาเสกฺขํ ธมฺมํ ปฏิจฺจ เนว\-เสกฺข\-นา\-เสกฺโข ธมฺโม อุปฺปชฺชติ วิปากปจฺจยา. วิปากํ เนว\-เสกฺข\-นาเสกฺขํ เอกํ ขนฺธํ ปฏิจฺจ ตโย ขนฺธา จิตฺต\-สมุฏฺฐานญฺจ รูปํ ฯเปฯ เทฺว ขนฺเธ ฯเปฯ. ปฏิสนฺธิ\-กฺขเณ ฯเปฯ. ขนฺเธ ปฏิจฺจ วตฺถุ, วตฺถุํ ปฏิจฺจ ขนฺธา. เอกํ มหาภูตํ ฯเปฯ. (1)}

\prose{เสกฺขญฺจ เนว\-เสกฺข\-นาเสกฺขญฺจ ธมฺมํ ปฏิจฺจ เนว\-เสกฺข\-นา\-เสกฺโข ธมฺโม อุปฺปชฺชติ วิปากปจฺจยา. วิปาเก เสกฺเข ขนฺเธ จ มหาภูเต จ ปฏิจฺจ จิตฺต\-สมุฏฺฐานํ รูปํ. (1)}

\prose{อเสกฺขญฺจ เนว\-เสกฺข\-นาเสกฺขญฺจ ธมฺมํ ปฏิจฺจ เนว\-เสกฺข\-นา\-เสกฺโข ธมฺโม อุปฺปชฺชติ วิปากปจฺจยา. อเสกฺเข ขนฺเธ จ มหาภูเต จ ปฏิจฺจ จิตฺต\-สมุฏฺฐานํ รูปํ. (1)}

\csromanheader{2}{11. เสกฺขตฺติก \textbf{อาหาราทิ}}
\proseitem{5}{\csromanfolio{277}เสกฺขํ ธมฺมํ ปฏิจฺจ เสกฺโข ธมฺโม อุปฺปชฺชติ อาหารปจฺจยา. อินฺทฺริยปจฺจยา. ฌานปจฺจยา. มคฺคปจฺจยา. สมฺปยุตฺต\-ปจฺจยา. วิปฺปยุตฺต\-ปจฺจยา. อตฺถิปจฺจยา. นตฺถิปจฺจยา. วิคตปจฺจยา. อวิคตปจฺจยา.}

\csromanmarkh{1. ปจฺจยานุโลม}\csromanmarkhii{2. สงฺขฺยาวาร}\csromanheadernomark{2}{1. \textbf{ปจฺจยานุโลม 2. สงฺขฺยาวาร}}
\proseitemwithcloserruled{6}{เหตุยา นว, อารมฺมเณ ตีณิ, อธิปติยา นว, อนนฺตเร ตีณิ, สมนนฺตเร ตีณิ, สหชาเต นว, อญฺญมญฺเญ ตีณิ, นิสฺสเย นว, อุปนิสฺสเย ตีณิ, ปุเรชาเต ตีณิ, อาเสวเน เทฺว, กมฺเม นว, วิปาเก, อาหาเร, อินฺทฺริเย, ฌาเน, มคฺเค นว, สมฺปยุตฺเต ตีณิ, วิปฺปยุตฺเต นว, อตฺถิยา นว, นตฺถิยา ตีณิ, วิคเต ตีณิ, อวิคเต นว. (เอวํ คเณตพฺพํ.)}{อนุโลมํ.}
\csromanmarkh{2. ปจฺจย\-ปจฺจนีย}\csromanmarkhii{1. วิภงฺควาร}\csromanheadernomark{2}{2. ปจฺจย\-ปจฺจนีย 1. วิภงฺควาร}
\csromanheader{2}{\textbf{นเหตุ}}
\proseitem{7}{เนว\-เสกฺข\-นาเสกฺขํ ธมฺมํ ปฏิจฺจ เนว\-เสกฺข\-นา\-เสกฺโข ธมฺโม อุปฺปชฺชติ นเหตุปจฺจยา. อเหตุกํ เนว\-เสกฺข\-นาเสกฺขํ เอกํ ขนฺธํ ปฏิจฺจ ตโย ขนฺธา จิตฺต\-สมุฏฺฐานญฺจ รูปํ ฯเปฯ เทฺว ขนฺเธ ฯเปฯ. อเหตุก\-ปฏิสนฺธิ\-กฺขเณ ฯเปฯ. ขนฺเธ ปฏิจฺจ วตฺถุ, วตฺถุํ ปฏิจฺจ ขนฺธา, เอกํ มหาภูตํ ฯเปฯ. พาหิรํ. อาหารสมุฏฺฐานํ. อุตุสมุฏฺฐานํ. อสญฺญสตฺตานํ เอกํ มหาภูตํ ฯเปฯ. วิจิกิจฺฉา\-สหคเต อุทฺธจฺจ\-สหคเต ขนฺเธ ปฏิจฺจ วิจิกิจฺฉา\-สหคโต อุทฺธจฺจ\-สหคโต โมโห. (1)}

\prose{นอารมฺมณ}

\proseitem{8}{เสกฺขํ ธมฺมํ ปฏิจฺจ เนว\-เสกฺข\-นา\-เสกฺโข ธมฺโม อุปฺปชฺชติ นอารมฺมณ\-ปจฺจยา. เสกฺเข ขนฺเธ ปฏิจฺจ จิตฺต\-สมุฏฺฐานํ รูปํ. (1)}

\prose{\csromanfolio{278}อเสกฺขํ ธมฺมํ ปฏิจฺจ เนว\-เสกฺข\-นา\-เสกฺโข ธมฺโม อุปฺปชฺชติ นอารมฺมณ\-ปจฺจยา. อเสกฺเข ขนฺเธ ปฏิจฺจ จิตฺต\-สมุฏฺฐานํ รูปํ. (1)}

\prose{เนว\-เสกฺข\-นาเสกฺขํ ธมฺมํ ปฏิจฺจ เนว\-เสกฺข\-นา\-เสกฺโข ธมฺโม อุปฺปชฺชติ นอารมฺมณ\-ปจฺจยา. เนว\-เสกฺข\-นาเสกฺเข ขนฺเธ ปฏิจฺจ จิตฺต\-สมุฏฺฐานํ รูปํ ฯเปฯ. ปฏิสนฺธิ\-กฺขเณ ฯเปฯ. ขนฺเธ ปฏิจฺจ วตฺถุ ฯเปฯ. เอกํ มหาภูตํ ฯเปฯ. อสญฺญสตฺตานํ ฯเปฯ. (1)}

\prose{เสกฺขญฺจ เนว\-เสกฺข\-นาเสกฺขญฺจ ธมฺมํ ปฏิจฺจ เนว\-เสกฺข\-นา\-เสกฺโข ธมฺโม อุปฺปชฺชติ นอารมฺมณ\-ปจฺจยา. เสกฺเข ขนฺเธ จ มหาภูเต จ ปฏิจฺจ จิตฺต\-สมุฏฺฐานํ รูปํ. (1)}

\prose{อเสกฺขญฺจ เนว\-เสกฺข\-นาเสกฺขญฺจ ธมฺมํ ปฏิจฺจ เนว\-เสกฺข\-นา\-เสกฺโข ธมฺโม อุปฺปชฺชติ นอารมฺมณ\-ปจฺจยา. อเสกฺเข ขนฺเธ จ มหาภูเต จ ปฏิจฺจ จิตฺต\-สมุฏฺฐานํ รูปํ. (1)}

\prose{นอธิปตฺยาทิ}

\proseitem{9}{เสกฺขํ ธมฺมํ ปฏิจฺจ เสกฺโข ธมฺโม อุปฺปชฺชติ นอธิปติ\-ปจฺจยา. เสกฺเข ขนฺเธ ปฏิจฺจ เสกฺโข อธิปติ. (1)}

\prose{อเสกฺขํ ธมฺมํ ปฏิจฺจ อเสกฺโข ธมฺโม อุปฺปชฺชติ นอธิปติ\-ปจฺจยา. อเสกฺเข ขนฺเธ ปฏิจฺจ อเสกฺโข อธิปติ. (1)}

\prose{เนว\-เสกฺข\-นาเสกฺขํ ธมฺมํ ปฏิจฺจ เนว\-เสกฺข\-นา\-เสกฺโข ธมฺโม อุปฺปชฺชติ นอธิปติ\-ปจฺจยา. (ปริปุณฺณํ. ปฏิสนฺธิปิ มหาภูตาปิ สพฺเพ.) นอนนฺตร\-ปจฺจยา. นสมนนฺตร\-ปจฺจยา. นอญฺญมญฺญ\-ปจฺจยา. นอุปนิสฺสยปจฺจยา. นปุเรชาต\-ปจฺจยา. สตฺต. (กุสล\-ตฺติก\-สทิสา.) นปจฺฉาชาต\-ปจฺจยา ฯเปฯ นอาเสวน\-ปจฺจยา. วิปากํ เสกฺขํ เอกํ ขนฺธํ ปฏิจฺจ ตโย ขนฺธา ฯเปฯ เทฺว ขนฺเธ ฯเปฯ. (1)}

\proseitem{10}{เสกฺขํ ธมฺมํ ปฏิจฺจ เนว\-เสกฺข\-นา\-เสกฺโข ธมฺโม อุปฺปชฺชติ นอาเสวน\-ปจฺจยา. เสกฺเข ขนฺเธ ปฏิจฺจ จิตฺต\-สมุฏฺฐานํ รูปํ. (2)}

\prose{เสกฺขํ ธมฺมํ ปฏิจฺจ เสกฺโข จ เนว\-เสกฺข\-นา\-เสกฺโข จ ธมฺมา อุปฺปชฺชนฺติ นอาเสวน\-ปจฺจยา. วิปากํ เสกฺขํ เอกํ ขนฺธํ ปฏิจฺจ ตโย ขนฺธา จิตฺต\-สมุฏฺฐานญฺจ รูปํ ฯเปฯ เทฺว ขนฺเธ ฯเปฯ. (3)}

\prose{\csromanfolio{279}อเสกฺขํ ธมฺมํ ปฏิจฺจ อเสกฺโข ธมฺโม. ตีณิ.}

\prose{เนว\-เสกฺข\-นาเสกฺขํ ธมฺมํ ปฏิจฺจ เนว\-เสกฺข\-นา\-เสกฺโข ธมฺโม อุปฺปชฺชติ นอาเสวน\-ปจฺจยา. เนว\-เสกฺข\-นาเสกฺขํ (เอกํ, ปริปุณฺณํ, เสกฺขญฺจ, เนว\-เสกฺข\-นาเสกฺขญฺจ, ฆฏนา ปริปุณฺณา, เทฺวปิ กาตพฺพา. นว.) นกมฺมปจฺจยา. เสกฺเข ขนฺเธ ปฏิจฺจ เสกฺขา เจตนา.}

\prose{เนว\-เสกฺข\-นาเสกฺขํ ธมฺมํ ปฏิจฺจ เนว\-เสกฺข\-นา\-เสกฺโข ธมฺโม อุปฺปชฺชติ นกมฺมปจฺจยา. เนว\-เสกฺข\-นาเสกฺเข ขนฺเธ ปฏิจฺจ เนว\-เสกฺข\-นาเสกฺขา เจตนา. พาหิรํ. อาหารสมุฏฺฐานํ. อุตุสมุฏฺฐานํ. เอกํ มหาภูตํ ฯเปฯ.}

\proseitem{11}{เสกฺขํ ธมฺมํ ปฏิจฺจ เสกฺโข ธมฺโม อุปฺปชฺชติ นวิปาก\-ปจฺจยา. เสกฺขํ เอกํ ขนฺธํ ปฏิจฺจ ตโย ขนฺธา ฯเปฯ. (1)}

\prose{เสกฺขํ ธมฺมํ ปฏิจฺจ เนว\-เสกฺข\-นา\-เสกฺโข ธมฺโม อุปฺปชฺชติ นวิปาก\-ปจฺจยา. เสกฺเข ขนฺเธ ปฏิจฺจ จิตฺต\-สมุฏฺฐานํ รูปํ. (2)}

\prose{เสกฺขํ ธมฺมํ ปฏิจฺจ เสกฺโข จ เนว\-เสกฺข\-นา\-เสกฺโข จ ธมฺมา อุปฺปชฺชนฺติ นวิปาก\-ปจฺจยา. เสกฺขํ เอกํ ขนฺธํ ปฏิจฺจ ตโย ขนฺธา จิตฺต\-สมุฏฺฐานญฺจ รูปํ ฯเปฯ. (3)}

\proseitem{12}{เนว\-เสกฺข\-นาเสกฺขํ ธมฺมํ ปฏิจฺจ เนว\-เสกฺข\-นา\-เสกฺโข ธมฺโม อุปฺปชฺชติ นวิปาก\-ปจฺจยา. (ปริปุณฺณํ. ปฏิสนฺธิ นตฺถิ.)}

\prose{เสกฺขญฺจ เนว\-เสกฺข\-นาเสกฺขญฺจ ธมฺมํ ปฏิจฺจ เนว\-เสกฺข\-นา\-เสกฺโข ธมฺโม อุปฺปชฺชติ นวิปาก\-ปจฺจยา. เสกฺเข ขนฺเธ จ มหาภูเต จ ปฏิจฺจ จิตฺต\-สมุฏฺฐานํ รูปํ. (1)}

\prose{นอาหาราทิ}

\proseitem{13}{เนว\-เสกฺข\-นาเสกฺขํ ธมฺมํ ปฏิจฺจ เนว\-เสกฺข\-นา\-เสกฺโข ธมฺโม อุปฺปชฺชติ นอาหารปจฺจยา. นอินฺทฺริยปจฺจยา. นฌานปจฺจยา นมคฺคปจฺจยา.}

\csromanheader{2}{\textbf{นสมฺปยุตฺตาทิ}}
\proseitem{14}{เสกฺขํ ธมฺมํ ปฏิจฺจ เนว\-เสกฺข\-นา\-เสกฺโข ธมฺโม อุปฺปชฺชติ นสมฺปยุตฺต\-ปจฺจยา ฯเปฯ. (นอารมฺมณปจฺจย\-สทิสํ)}

\csromanheader{2}{\textbf{นวิปฺปยุตฺตาทิ}}
\prose{\csromanfolio{280}เสกฺขํ ธมฺมํ ปฏิจฺจ เสกฺโข ธมฺโม อุปฺปชฺชติ นวิปฺปยุตฺต\-ปจฺจยา. อรูเป เสกฺขํ เอกํ ขนฺธํ ปฏิจฺจ ตโย ขนฺธา ฯเปฯ. (1)}

\prose{อเสกฺขํ ธมฺมํ ปฏิจฺจ อเสกฺโข ธมฺโม อุปฺปชฺชติ นวิปฺปยุตฺต\-ปจฺจยา. อรูเป อเสกฺขํ เอกํ ขนฺธํ ปฏิจฺจ ตโย ขนฺธา ฯเปฯ. (1)}

\prose{เนว\-เสกฺข\-นาเสกฺขํ ธมฺมํ ปฏิจฺจ เนว\-เสกฺข\-นา\-เสกฺโข ธมฺโม อุปฺปชฺชติ นวิปฺปยุตฺต\-ปจฺจยา. อรูเป เนว\-เสกฺข\-นาเสกฺขํ เอกํ ขนฺธํ ปฏิจฺจ ตโย ขนฺธา ฯเปฯ. พาหิรํ. อาหารสมุฏฺฐานํ. อุตุสมุฏฺฐานํ. อสญฺญสตฺตานํ ฯเปฯ โนนตฺถิปจฺจยา. โนวิคต\-ปจฺจยา.}

\csromanmarkh{2. ปจฺจย\-ปจฺจนีย}\csromanmarkhii{2. สงฺขฺยาวาร}\csromanheadernomark{2}{2. ปจฺจย\-ปจฺจนีย 2. สงฺขฺยาวาร}
\proseitemwithcloserruled{15}{นเหตุยา เอกํ, นอารมฺมเณ ปญฺจ, นอธิปติยา ตีณิ, นอนนฺตเร, นสมนนฺตเร, นอญฺญมญฺเญ, นอุปนิสฺสเย ปญฺจ, นปุเรชาเต สตฺต, นปจฺฉาชาเต นว, นอาเสวเน นว, นกมฺเม เทฺว, นวิปาเก ปญฺจ, นอาหาเร เอกํ, นอินฺทฺริเย เอกํ, นฌาเน เอกํ, นมคฺเค เอกํ, นสมฺปยุตฺเต ปญฺจ, นวิปฺปยุตฺเต ตีณิ, โนนตฺถิยา ปญฺจ, โนวิคเต ปญฺจ. (เอวํ คเณตพฺพํ.)}{\textbf{ปจฺจนียํ.}}
\proseitem{16}{เหตุปจฺจยา นอารมฺมเณ ปญฺจ, นอธิปติยา ตีณิ, นอนนฺตเร, นสมนนฺตเร, นอญฺญมญฺเญ, นอุปนิสฺสเย ปญฺจ, นปุเรชาเต สตฺต, นปจฺฉาชาเต, นอาเสวเน นว, นกมฺเม เทฺว, นวิปาเก ปญฺจ, นสมฺปยุตฺเต ปญฺจ, นวิปฺปยุตฺเต ตีณิ, โนนตฺถิยา ปญฺจ, โนวิคเต ปญฺจ. (เอวํ คเณตพฺพํ.)}

\vspace{\tipitakaparskip}
\csromancenter{อนุโลม\-ปจฺจนียํ.}
\csromanmarkh{11. เสกฺขตฺติก}\csromanmarkhii{4. ปจฺจย\-ปจฺจนียา\-นุโลม}\csromanheadernomark{2}{11. เสกฺขตฺติก 4. ปจฺจย\-ปจฺจนียา\-นุโลม}
\proseitem{17}{\csromanfolio{281}นเหตุปจฺจยา อารมฺมเณ เอกํ, อนนฺตเร, สมนนฺตเร, สหชาเต, อญฺญมญฺเญ, นิสฺสเย, อุปนิสฺสเย, ปุเรชาเต, อาเสวเน, กมฺเม, วิปาเก, อาหาเร, อินฺทฺริเย, ฌาเน, มคฺเค, สมฺปยุตฺเต, วิปฺปยุตฺเต, อตฺถิยา, นตฺถิยา, วิคเต, อวิคเต เอกํ. (เอวํ คเณตพฺพํ.) ปจฺจนียา\-นุโลมํ. ปฏิจฺจวาโร. (สหชาตวาโร ปฏิจฺจ\-วาร\-สทิโส.)}
\csromansectionrule

\csromanmatikaanchor{mtk.6}
\csromanmatikaanchor{mtk.7}
\csromantocmarkref{section}{นเหตุทุก}{mtk.6}
\bookmark[dest={mtk.6},level=1]{นเหตุทุก}
\csromantocmarknopage{section}{11. เสกฺขตฺติก 3. ปจฺจยวาร}{mtk.7}
\bookmark[dest={mtk.7},level=1]{11. เสกฺขตฺติก 3. ปจฺจยวาร}
\csromantocmarknopage{subsection}{1. ปจฺจยานุโลม 1. วิภงฺควาร}{mtk.8}
\bookmark[dest={mtk.8},level=2]{1. ปจฺจยานุโลม 1. วิภงฺควาร}
\csromanmarkh{11. เสกฺขตฺติก}\csromanmarkhii{3. ปจฺจยวาร}\csromanheadernomark{1}{11. เสกฺขตฺติก 3. ปจฺจยวาร}
\csromanmarkh{1. ปจฺจยานุโลม}\csromanmarkhii{1. วิภงฺควาร}\csromanheadernomark{2}{1. \textbf{ปจฺจยานุโลม 1. วิภงฺควาร}}
\csromanmatikaanchor{mtk.9}
\csromantocmarkref{subsubsection}{เหตุ}{mtk.9}
\bookmark[dest={mtk.9},level=3]{เหตุ}
\csromanheader{3}{\textbf{เหตุ}}
\proseitem{18}{เสกฺขํ ธมฺมํ ปจฺจยา เสกฺโข ธมฺโม อุปฺปชฺชติ เหตุปจฺจยา. ตีณิ. (ปฏิจฺจ\-วาร\-สทิสํ.)}

\prose{อเสกฺขํ ธมฺมํ ปจฺจยา อเสกฺโข ธมฺโม อุปฺปชฺชติ เหตุปจฺจยา. ตีณิ. (ปฏิจฺจ\-วาร\-สทิสํ.)}

\proseitem{19}{เนว\-เสกฺข\-นาเสกฺขํ ธมฺมํ ปจฺจยา เนว\-เสกฺข\-นา\-เสกฺโข ธมฺโม อุปฺปชฺชติ เหตุปจฺจยา. (ปริปุณฺณํ.) มหาภูเต ปจฺจยา จิตฺต\-สมุฏฺฐานํ รูปํ, กฏตฺตารูปํ, อุปาทารูปํ. วตฺถุํ ปจฺจยา เนว\-เสกฺข\-นาเสกฺขา ขนฺธา. (1)}

\prose{เนว\-เสกฺข\-นาเสกฺขํ ธมฺมํ ปจฺจยา เสกฺโข ธมฺโม อุปฺปชฺชติ เหตุปจฺจยา. วตฺถุํ ปจฺจยา เสกฺขา ขนฺธา. (2)}

\prose{เนว\-เสกฺข\-นาเสกฺขํ ธมฺมํ ปจฺจยา อเสกฺโข ธมฺโม อุปฺปชฺชติ เหตุปจฺจยา. วตฺถุํ ปจฺจยา อเสกฺขา ขนฺธา. (3)}

\prose{\csromanfolio{282}เนว\-เสกฺข\-นาเสกฺขํ ธมฺมํ ปจฺจยา เสกฺโข จ เนว\-เสกฺข\-นา\-เสกฺโข จ ธมฺมา อุปฺปชฺชนฺติ เหตุปจฺจยา. วตฺถุํ ปจฺจยา เสกฺขา ขนฺธา. มหาภูเต ปจฺจยา จิตฺต\-สมุฏฺฐานํ รูปํ. (4)}

\prose{เนว\-เสกฺข\-นาเสกฺขํ ธมฺมํ ปจฺจยา อเสกฺโข จ เนว\-เสกฺข\-นา\-เสกฺโข จ ธมฺมา อุปฺปชฺชนฺติ เหตุปจฺจยา. วตฺถุํ ปจฺจยา อเสกฺขา ขนฺธา. มหาภูเต ปจฺจยา จิตฺต\-สมุฏฺฐานํ รูปํ. (5)}

\proseitem{20}{เสกฺขญฺจ เนว\-เสกฺข\-นาเสกฺขญฺจ ธมฺมํ ปจฺจยา เสกฺโข ธมฺโม อุปฺปชฺชติ เหตุปจฺจยา. เสกฺขํ เอกํ ขนฺธญฺจ วตฺถุญฺจ ปจฺจยา ตโย ขนฺธา ฯเปฯ เทฺว ขนฺธา. (1)}

\prose{เสกฺขญฺจ เนว\-เสกฺข\-นาเสกฺขญฺจ ธมฺมํ ปจฺจยา เนว\-เสกฺข\-นา\-เสกฺโข ธมฺโม อุปฺปชฺชติ เหตุปจฺจยา. เสกฺเข ขนฺเธ จ มหาภูเต จ ปจฺจยา จิตฺต\-สมุฏฺฐานํ รูปํ. (2)}

\prose{เสกฺขญฺจ เนว\-เสกฺข\-นาเสกฺขญฺจ ธมฺมํ ปจฺจยา เสกฺโข จ เนว\-เสกฺข\-นา\-เสกฺโข จ ธมฺมา อุปฺปชฺชนฺติ เหตุปจฺจยา. เสกฺขํ เอกํ ขนฺธญฺจ วตฺถุญฺจ ปจฺจยา ตโย ขนฺธา ฯเปฯ เทฺว ขนฺธา. เสกฺเข ขนฺเธ จ มหาภูเต จ ปจฺจยา จิตฺต\-สมุฏฺฐานํ รูปํ. (3)}

\prose{อเสกฺขญฺจ เนว\-เสกฺข\-นาเสกฺขญฺจ ธมฺมํ ปจฺจยา. ตีณิ. (เสกฺขสทิสา.)}

\csromanmatikaanchor{mtk.10}
\csromantocmarkref{subsubsection}{อารมฺมณ}{mtk.10}
\bookmark[dest={mtk.10},level=3]{อารมฺมณ}
\csromanheader{3}{\textbf{อารมฺมณ}}
\proseitem{21}{เสกฺขํ ธมฺมํ ปจฺจยา เสกฺโข ธมฺโม อุปฺปชฺชติ อารมฺมณ\-ปจฺจยา. เอกํ.}

\prose{อเสกฺขํ ธมฺมํ ปจฺจยา อเสกฺโข ธมฺโม. เอกํ.}

\prose{เนว\-เสกฺข\-นาเสกฺขํ ธมฺมํ ปจฺจยา. เนว\-เสกฺข\-นา\-เสกฺโข ธมฺโม อุปฺปชฺชติ อารมฺมณ\-ปจฺจยา. เอกํ. วตฺถุํ ปจฺจยา เนว\-เสกฺข\-นาเสกฺขา ขนฺธา. จกฺขายตนํ ปจฺจยา จกฺขุวิญฺญาณํ ฯเปฯ. กายายตนํ ปจฺจยา กายวิญฺญาณํ. วตฺถุํ ปจฺจยา เนว\-เสกฺข\-นาเสกฺขา ขนฺธา. (1)}

\prose{เนว\-เสกฺข\-นาเสกฺขํ ธมฺมํ ปจฺจยา เสกฺโข ธมฺโม อุปฺปชฺชติ อารมฺมณ\-ปจฺจยา. วตฺถุํ ปจฺจยา เสกฺขา ขนฺธา. (2)}

\prose{\csromanfolio{283}เนว\-เสกฺข\-นาเสกฺขํ ธมฺมํ ปจฺจยา อเสกฺโข ธมฺโม อุปฺปชฺชติ อารมฺมณ\-ปจฺจยา. วตฺถุํ ปจฺจยา อเสกฺขา ขนฺธา. (3)}

\prose{เสกฺขญฺจ เนว\-เสกฺข\-นาเสกฺขญฺจ ธมฺมํ ปจฺจยา เสกฺโข ธมฺโม อุปฺปชฺชติ อารมฺมณ\-ปจฺจยา. เสกฺขํ เอกํ ขนฺธญฺจ วตฺถุญฺจ ปจฺจยา ตโย ขนฺธา ฯเปฯ เทฺว ขนฺธา. (1)}

\prose{อเสกฺขญฺจ เนว\-เสกฺข\-นาเสกฺขญฺจ ธมฺมํ ปจฺจยา อเสกฺโข ธมฺโม อุปฺปชฺชติ อารมฺมณ\-ปจฺจยา. อเสกฺขํ เอกํ ขนฺธญฺจ วตฺถุญฺจ ปจฺจยา ตโย ขนฺธา ฯเปฯ เทฺว ขนฺธา. (1)}

\csromanmatikaanchor{mtk.11}
\csromantocmarkref{subsubsection}{อธิปตฺยาทิ}{mtk.11}
\bookmark[dest={mtk.11},level=3]{อธิปตฺยาทิ}
\csromanheader{3}{\textbf{อธิปตฺยาทิ}}
\proseitem{22}{เสกฺขํ ธมฺมํ ปจฺจยา เสกฺโข ธมฺโม อุปฺปชฺชติ อธิปติ\-ปจฺจยา. อนนฺตร\-ปจฺจยา. สมนนฺตร\-ปจฺจยา. สหชาต\-ปจฺจยา. อญฺญมญฺญ\-ปจฺจยา. นิสฺสยปจฺจยา. อุปนิสฺสย\-ปจฺจยา. ปุเรชาต\-ปจฺจยา. อาเสวนปจฺจยา. เสกฺขํ เอกํ ขนฺธํ ปจฺจยา ตโย ขนฺธา ฯเปฯ.}

\prose{เนว\-เสกฺข\-นาเสกฺขํ ธมฺมํ ปจฺจยา เนว\-เสกฺข\-นา\-เสกฺโข ธมฺโม อุปฺปชฺชติ อาเสวนปจฺจยา. เนว\-เสกฺข\-นาเสกฺขํ เอกํ ขนฺธํ ปจฺจยา ฯเปฯ. วตฺถุํ ปจฺจยา เนว\-เสกฺข\-นาเสกฺขา ขนฺธา. (1)}

\prose{เนว\-เสกฺข\-นาเสกฺขํ ธมฺมํ ปจฺจยา เสกฺโข ธมฺโม อุปฺปชฺชติ อาเสวนปจฺจยา. วตฺถุํ ปจฺจยา เสกฺขา ขนฺธา. (2)}

\prose{เสกฺขญฺจ เนว\-เสกฺข\-นาเสกฺขญฺจ ธมฺมํ ปจฺจยา เสกฺโข ธมฺโม อุปฺปชฺชติ อาเสวนปจฺจยา. เสกฺขํ เอกํ ขนฺธญฺจ วตฺถุญฺจ ปจฺจยา ตโย ขนฺธา ฯเปฯ เทฺว ขนฺธา. (1)}

\csromanheader{2}{\textbf{กมฺมาทิ}}
\proseitem{23}{เสกฺขํ ธมฺมํ ปจฺจยา เสกฺโข ธมฺโม อุปฺปชฺชติ กมฺมปจฺจยา. วิปากปจฺจยา. วิปากํ เสกฺขํ เอกํ ขนฺธํ ฯเปฯ. อาหารปจฺจยา. อินฺทฺริยปจฺจยา. ฌานปจฺจยา. มคฺคปจฺจยา. สมฺปยุตฺต\-ปจฺจยา. วิปฺปยุตฺต\-ปจฺจยา. อตฺถิปจฺจยา. นตฺถิปจฺจยา. วิคตปจฺจยา. อวิคตปจฺจยา.}

\csromanmarkh{1. ปจฺจยานุโลม}\csromanmarkhii{2. สงฺขฺยาวาร}\csromanheadernomark{2}{1. \textbf{ปจฺจยานุโลม 2. สงฺขฺยาวาร}}
\csromanmatikaanchor{mtk.12}
\csromanmatikaanchor{mtk.13}
\csromanmatikaanchor{mtk.16}
\csromantocmarknopage{paragraph}{1. ปจฺจยานุโลม 2. สงฺขฺยาวาร}{mtk.12}
\bookmark[dest={mtk.12},level=4]{1. ปจฺจยานุโลม 2. สงฺขฺยาวาร}
\csromantocmarkref{subparagraph}{สุทฺธ}{mtk.13}
\bookmark[dest={mtk.13},level=5]{สุทฺธ}
\csromantocmarknopage{paragraph}{2. ปจฺจยปจฺจนีย 1. วิภงฺควาร}{mtk.14}
\bookmark[dest={mtk.14},level=4]{2. ปจฺจยปจฺจนีย 1. วิภงฺควาร}
\proseitemwithcloserruled{24}{\csromanfolio{284}เหตุยา สตฺตรส, อารมฺมเณ สตฺต, อธิปติยา สตฺตรส, อนนฺตเร สตฺต, สมนนฺตเร สตฺต, สหชาเต สตฺตรส, อญฺญมญฺเญ สตฺต, นิสฺสเย สตฺตรส, อุปนิสฺสเย สตฺต, ปุเรชาเต สตฺต, อาเสวเน จตฺตาริ, กมฺเม สตฺตรส, วิปาเก สตฺตรส, อาหาเร สตฺตรส, อินฺทฺริเย, ฌาเน, มคฺเค สตฺตรส, สมฺปยุตฺเต สตฺต, วิปฺปยุตฺเต สตฺตรส, อตฺถิยา สตฺตรส, นตฺถิยา สตฺต, วิคเต สตฺต, อวิคเต สตฺตรส. (เอวํ คเณตพฺพํ.)}{อนุโลมํ.}
\csromanmarkh{2. ปจฺจย\-ปจฺจนีย}\csromanmarkhii{1. วิภงฺควาร}\csromanheadernomark{2}{2. ปจฺจย\-ปจฺจนีย 1. วิภงฺควาร}
\csromanmatikaanchor{mtk.15}
\csromantocmarkref{subparagraph}{นเหตุ}{mtk.15}
\bookmark[dest={mtk.15},level=5]{นเหตุ}
\csromantocmarkref{subparagraph}{นอารมฺมณาทิ}{mtk.16}
\bookmark[dest={mtk.16},level=5]{นอารมฺมณาทิ}
\csromanheader{5}{\textbf{นเหตุ}}
\proseitem{25}{เนว\-เสกฺข\-นาเสกฺขํ ธมฺมํ ปจฺจยา เนว\-เสกฺข\-นา\-เสกฺโข ธมฺโม อุปฺปชฺชติ นเหตุปจฺจยา. อเหตุกํ เนว\-เสกฺข\-นาเสกฺขํ เอกํ ขนฺธํ ปจฺจยา ตโย ขนฺธา จิตฺต\-สมุฏฺฐานญฺจ รูปํ ฯเปฯ เทฺว ขนฺธา. อเหตุก\-ปฏิสนฺธิ\-กฺขเณ ฯเปฯ. ขนฺเธ ปจฺจยา วตฺถุ, วตฺถุํ ปจฺจยา ขนฺธา. เอกํ มหาภูตํ ปจฺจยา ฯเปฯ. พาหิรํ. อาหารสมุฏฺฐานํ. อุตุสมุฏฺฐานํ. อสญฺญสตฺตานํ ฯเปฯ. จกฺขายตนํ ปจฺจยา จกฺขุวิญฺญาณํ ฯเปฯ. กายายตนํ ปจฺจยา กายวิญฺญาณํ. วตฺถุํ ปจฺจยา อเหตุกา เนว\-เสกฺข\-นาเสกฺขา ขนฺธา. วิจิกิจฺฉา\-สหคเต อุทฺธจฺจ\-สหคเต ขนฺเธ จ วตฺถุญฺจ ปจฺจยา วิจิกิจฺฉา\-สหคโต อุทฺธจฺจ\-สหคโต โมโห. (1)}

\prose{นอารมฺมณาทิ}

\proseitem{26}{เสกฺขํ ธมฺมํ ปจฺจยา เนว\-เสกฺข\-นา\-เสกฺโข ธมฺโม อุปฺปชฺชติ นอารมฺมณ\-ปจฺจยา ฯเปฯ. (1)}

\prose{เสกฺขํ ธมฺมํ ปจฺจยา เสกฺโข ธมฺโม อุปฺปชฺชติ นอธิปติ\-ปจฺจยา. เสกฺเข ขนฺเธ ปจฺจยา เสกฺโข อธิปติ. (1)}

\prose{\csromanfolio{285}อเสกฺขํ ธมฺมํ ปจฺจยา อเสกฺโข ธมฺโม อุปฺปชฺชติ นอธิปติ\-ปจฺจยา. อเสกฺเข ขนฺเธ ปจฺจยา อเสกฺโข อธิปติ. (1)}

\prose{เนว\-เสกฺข\-นาเสกฺขํ ธมฺมํ ปจฺจยา เนว\-เสกฺข\-นา\-เสกฺโข ธมฺโม อุปฺปชฺชติ นอธิปติ\-ปจฺจยา. (ปริปุณฺณํ.) อสญฺญสตฺตานํ ฯเปฯ. จกฺขายตนํ ฯเปฯ. วตฺถุํ ปจฺจยา เนว\-เสกฺข\-นาเสกฺขา ขนฺธา. (1)}

\prose{เนว\-เสกฺข\-นาเสกฺขํ ธมฺมํ ปจฺจยา เสกฺโข ธมฺโม อุปฺปชฺชติ นอธิปติ\-ปจฺจยา. วตฺถุํ ปจฺจยา เสกฺโข อธิปติ. (2)}

\prose{เนว\-เสกฺข\-นาเสกฺขํ ธมฺมํ ปจฺจยา อเสกฺโข ธมฺโม อุปฺปชฺชติ นอธิปติ\-ปจฺจยา. วตฺถุํ ปจฺจยา อเสกฺโข อธิปติ. (3)}

\prose{เสกฺขญฺจ เนว\-เสกฺข\-นาเสกฺขญฺจ ธมฺมํ ปจฺจยา เสกฺโข ธมฺโม อุปฺปชฺชติ นอธิปติ\-ปจฺจยา. เสกฺเข ขนฺเธ จ วตฺถุญฺจ ปจฺจยา เสกฺโข อธิปติ. (1)}

\prose{อเสกฺขญฺจ เนว\-เสกฺข\-นาเสกฺขญฺจ ธมฺมํ ปจฺจยา อเสกฺโข ธมฺโม อุปฺปชฺชติ นอธิปติ\-ปจฺจยา. อเสกฺเข ขนฺเธ จ วตฺถุญฺจ ปจฺจยา อเสกฺโข อธิปติ. (1)}

\prose{นอนนฺตราทิ}

\proseitem{27}{เสกฺขํ ธมฺมํ ปจฺจยา เนว\-เสกฺข\-นา\-เสกฺโข ธมฺโม อุปฺปชฺชติ นอนนฺตร\-ปจฺจยา. นสมนนฺตร\-ปจฺจยา. นอญฺญมญฺญ\-ปจฺจยา. นอุปนิสฺสยปจฺจยา. นปุเรชาต\-ปจฺจยา. นปจฺฉาชาต\-ปจฺจยา. (สตฺต.) นอาเสวน\-ปจฺจยา. นกมฺมปจฺจยา. เสกฺเข ขนฺเธ ปจฺจยา เสกฺขา เจตนา.}

\prose{เนว\-เสกฺข\-นาเสกฺขํ ธมฺมํ ปจฺจยา เนว\-เสกฺข\-นา\-เสกฺโข ธมฺโม อุปฺปชฺชติ นกมฺมปจฺจยา. เนว\-เสกฺข\-นาเสกฺเข ขนฺเธ ปจฺจยา เนว\-เสกฺข\-นาเสกฺขา เจตนา. พาหิรํ. อาหารสมุฏฺฐานํ. อุตุสมุฏฺฐานํ ฯเปฯ. วตฺถุํ ปจฺจยา เนว\-เสกฺข\-นาเสกฺขา เจตนา. (1)}

\prose{เนว\-เสกฺข\-นาเสกฺขํ ธมฺมํ ปจฺจยา เสกฺโข ธมฺโม อุปฺปชฺชติ นกมฺมปจฺจยา. วตฺถุํ ปจฺจยา เสกฺขา เจตนา. (2)}

\prose{เสกฺขญฺจ เนว\-เสกฺข\-นาเสกฺขญฺจ ธมฺมํ ปจฺจยา เสกฺโข ธมฺโม อุปฺปชฺชติ นกมฺมปจฺจยา. เสกฺเข ขนฺเธ จ วตฺถุญฺจ ปจฺจยา เสกฺขา เจตนา. (1)}

\csromanmatikaanchor{mtk.17}
\csromanmatikaanchor{mtk.18}
\csromanmatikaanchor{mtk.19}
\csromanmatikaanchor{mtk.20}
\csromantocmarknopage{subsection}{2. ปจฺจยปจฺจนีย 2. สงฺขฺยาวาร}{mtk.17}
\bookmark[dest={mtk.17},level=2]{2. ปจฺจยปจฺจนีย 2. สงฺขฺยาวาร}
\csromantocmarkref{subsubsection}{สุทฺธ}{mtk.18}
\bookmark[dest={mtk.18},level=3]{สุทฺธ}
\csromantocmarknopage{subsection}{3. ปจฺจยานุโลมปจฺจนีย}{mtk.19}
\bookmark[dest={mtk.19},level=2]{3. ปจฺจยานุโลมปจฺจนีย}
\csromantocmarkref{subsubsection}{เหตุทุก}{mtk.20}
\bookmark[dest={mtk.20},level=3]{เหตุทุก}
\proseitem{28}{\csromanfolio{286}เสกฺขํ ธมฺมํ ปจฺจยา เสกฺโข ธมฺโม อุปฺปชฺชติ นวิปาก\-ปจฺจยา. (เสกฺขมูลเก ตีณิ.)}

\prose{เนว\-เสกฺข\-นาเสกฺขํ ธมฺมํ ปจฺจยา เนว\-เสกฺข\-นา\-เสกฺโข ธมฺโม อุปฺปชฺชติ นวิปาก\-ปจฺจยา. (เนวเสกฺขนาเสกฺข\-มูลเก ตีณิ.)}

\prose{เสกฺขญฺจ เนว\-เสกฺข\-นาเสกฺขญฺจ ธมฺมํ ปจฺจยา เสกฺโข ธมฺโม อุปฺปชฺชติ นวิปาก\-ปจฺจยา. (เสกฺข\-ฆฏเนสุ ตีณิ.)}

\prose{นอาหาราทิ}

\proseitem{29}{เนว\-เสกฺข\-นาเสกฺขํ ธมฺมํ ปจฺจยา เนว\-เสกฺข\-นา\-เสกฺโข ธมฺโม อุปฺปชฺชติ นอาหารปจฺจยา. นอินฺทฺริยปจฺจยา. นฌานปจฺจยา. นมคฺคปจฺจยา. นสมฺปยุตฺต\-ปจฺจยา. นวิปฺปยุตฺต\-ปจฺจยา. โนนตฺถิปจฺจยา. โนวิคต\-ปจฺจยา ฯเปฯ.}

\csromanmarkh{2. ปจฺจย\-ปจฺจนีย}\csromanmarkhii{2. สงฺขฺยาวาร}\csromanheadernomark{2}{2. ปจฺจย\-ปจฺจนีย 2. สงฺขฺยาวาร}
\proseitemwithcloserruled{30}{นเหตุยา เอกํ, นอารมฺมเณ ปญฺจ, นอธิปติยา สตฺต, นอนนฺตเร ปญฺจ, นสมนนฺตเร ปญฺจ, นอญฺญมญฺเญ ปญฺจ, นอุปนิสฺสเย ปญฺจ, นปุเรชาเต สตฺต, นปจฺฉาชาเต สตฺตรส, นอาเสวเน สตฺตรส, นกมฺเม จตฺตาริ, นวิปาเก นว, นอาหาเร เอกํ, นอินฺทฺริเย เอกํ, นฌาเน เอกํ, นมคฺเค เอกํ, นสมฺปยุตฺเต ปญฺจ, นวิปฺปยุตฺเต ตีณิ, โนนตฺถิยา ปญฺจ, โนวิคเต ปญฺจ. (เอวํ คเณตพฺพํ.)}{ปจฺจนียํ.}
\proseitem{31}{เหตุปจฺจยา นอารมฺมเณ ปญฺจ, นอธิปติยา สตฺต, นอนนฺตเร, นสมนนฺตเร, นอญฺญมญฺเญ, นอุปนิสฺสเย ปญฺจ, นปุเรชาเต สตฺต,}

\csromanmatikaanchor{mtk.21}
\csromantocmarkref{subsection}{4. ปจฺจยปจฺจนียานุโลม ...}{mtk.21}
\bookmark[dest={mtk.21},level=2]{4. ปจฺจยปจฺจนียานุโลม ...}
\prosewithcloserruled{\csromanfolio{287}นปจฺฉาชาเต สตฺตรส, นอาเสวเน สตฺตรส, นกมฺเม จตฺตาริ, นวิปาเก นว, นสมฺปยุตฺเต ปญฺจ, นวิปฺปยุตฺเต ตีณิ, โนนตฺถิยา ปญฺจ, โนวิคเต ปญฺจ. (เอวํ คเณตพฺพํ.)}{อนุโลม\-ปจฺจนียํ.}
\proseitem{32}{นเหตุปจฺจยา อารมฺมเณ เอกํ, อนนฺตเร, สมนนฺตเร, สหชาเต, อญฺญมญฺเญ เอกํ ฯเปฯ อวิคเต เอกํ. (เอวํ คเณตพฺพํ.) ปจฺจนียา\-นุโลมํ. ปจฺจยวาโร. (นิสฺสยวาโร ปจฺจยวาร\-สทิโส.)}
\csromansectionrule

\csromanmatikaanchor{mtk.22}
\csromantocmarknopage{paragraph}{11. เสกฺขตฺติก 5. สํสฏฺฐวาร}{mtk.22}
\bookmark[dest={mtk.22},level=4]{11. เสกฺขตฺติก 5. สํสฏฺฐวาร}
\csromanmarkh{11. เสกฺขตฺติก}\csromanmarkhii{5. สํสฏฺฐวาร}\csromanheadernomark{4}{11. เสกฺขตฺติก 5. สํสฏฺฐวาร}
\csromanmarkh{1. ปจฺจยานุโลม}\csromanmarkhii{1. วิภงฺควาร}\csromanheadernomark{2}{1. \textbf{ปจฺจยานุโลม 1. วิภงฺควาร}}
\csromanmatikaanchor{mtk.23}
\csromanmatikaanchor{mtk.24}
\csromantocmarknopage{subparagraph}{1. ปจฺจยานุโลม 1. วิภงฺควาร}{mtk.23}
\bookmark[dest={mtk.23},level=5]{1. ปจฺจยานุโลม 1. วิภงฺควาร}
\csromantocmarkref{subparagraph}{เหตุ}{mtk.24}
\bookmark[dest={mtk.24},level=5]{เหตุ}
\csromanheader{6}{\textbf{เหตุ}}
\proseitem{33}{เสกฺขํ ธมฺมํ สํสฏฺโฐ เสกฺโข ธมฺโม อุปฺปชฺชติ เหตุปจฺจยา. เสกฺขํ เอกํ ขนฺธํ สํสฏฺฐา ตโย ขนฺธา ฯเปฯ เทฺว ขนฺเธ ฯเปฯ. (1)}

\prose{อเสกฺขํ ธมฺมํ สํสฏฺโฐ อเสกฺโข ธมฺโม อุปฺปชฺชติ เหตุปจฺจยา. อเสกฺขํ เอกํ ขนฺธํ สํสฏฺฐา ตโย ขนฺธา ฯเปฯ เทฺว ขนฺเธ ฯเปฯ. (1)}

\prose{เนว\-เสกฺข\-นาเสกฺขํ ธมฺมํ สํสฏฺโฐ เนว\-เสกฺข\-นา\-เสกฺโข ธมฺโม อุปฺปชฺชติ เหตุปจฺจยา. เนว\-เสกฺข\-นาเสกฺขํ เอกํ ขนฺธํ สํสฏฺฐา ตโย ขนฺธา ฯเปฯ เทฺว ขนฺเธ ฯเปฯ. ปฏิสนฺธิ\-กฺขเณ ฯเปฯ. (1)}

\csromanheader{2}{\textbf{อารมฺมณาทิ}}
\proseitem{34}{\csromanfolio{288}เสกฺขํ ธมฺมํ สํสฏฺโฐ เสกฺโข ธมฺโม อุปฺปชฺชติ อารมฺมณ\-ปจฺจยา. อธิปติ\-ปจฺจยา ฯเปฯ. ปุเรชาต\-ปจฺจยา. อาเสวนปจฺจยา. (เทฺว กาตพฺพา) ฯเปฯ. อวิคตปจฺจยา.}

\csromanmatikaanchor{mtk.25}
\csromantocmarkref{subparagraph}{1. ปจฺจยานุโลม 2. สงฺขฺยาวาร}{mtk.25}
\bookmark[dest={mtk.25},level=5]{1. ปจฺจยานุโลม 2. สงฺขฺยาวาร}
\csromanmarkh{1. ปจฺจยานุโลม}\csromanmarkhii{2. สงฺขฺยาวาร}\csromanheadernomark{6}{1. \textbf{ปจฺจยานุโลม 2. สงฺขฺยาวาร}}
\proseitemwithcloserruled{35}{เหตุยา ตีณิ, อารมฺมเณ ตีณิ, อธิปติยา ตีณิ, อนนฺตเร, สมนนฺตเร, สหชาเต, อญฺญมญฺเญ, นิสฺสเย, อุปนิสฺสเย, ปุเรชาเต ตีณิ, อาเสวเน เทฺว, กมฺเม ตีณิ ฯเปฯ อวิคเต ตีณิ. (เอวํ คเณตพฺพํ.)}{อนุโลมํ.}
\csromanmatikaanchor{mtk.26}
\csromantocmarkref{subparagraph}{2. ปจฺจยปจฺจนีย 1. วิภงฺควาร}{mtk.26}
\bookmark[dest={mtk.26},level=5]{2. ปจฺจยปจฺจนีย 1. วิภงฺควาร}
\csromanmarkh{2. ปจฺจย\-ปจฺจนีย}\csromanmarkhii{1. วิภงฺควาร}\csromanheadernomark{6}{2. ปจฺจย\-ปจฺจนีย 1. วิภงฺควาร}
\csromanheader{2}{\textbf{นเหตุ}}
\proseitem{36}{เนว\-เสกฺข\-นาเสกฺขํ ธมฺมํ สํสฏฺโฐ เนว\-เสกฺข\-นา\-เสกฺโข ธมฺโม อุปฺปชฺชติ นเหตุปจฺจยา. อเหตุกํ เนว\-เสกฺข\-นาเสกฺขํ เอกํ ขนฺธํ สํสฏฺฐา ตโย ขนฺธา ฯเปฯ เทฺว ขนฺเธ ฯเปฯ. อเหตุก\-ปฏิสนฺธิ\-กฺขเณ ฯเปฯ. วิจิกิจฺฉา\-สหคเต อุทฺธจฺจ\-สหคเต ขนฺเธ สํสฏฺโฐ วิจิกิจฺฉา\-สหคโต อุทฺธจฺจ\-สหคโต โมโห. (1)}

\prose{นอธิปติ}

\proseitem{37}{เสกฺขํ ธมฺมํ สํสฏฺโฐ เสกฺโข ธมฺโม อุปฺปชฺชติ นอธิปติ\-ปจฺจยา. เสกฺเข ขนฺเธ สํสฏฺโฐ เสกฺโข อธิปติ. (1)}

\prose{อเสกฺขํ ธมฺมํ สํสฏฺโฐ อเสกฺโข ธมฺโม อุปฺปชฺชติ นอธิปติ\-ปจฺจยา อเสกฺเข ขนฺเธ สํสฏฺโฐ อเสกฺโข อธิปติ. (1)}

\prose{เนว\-เสกฺข\-นาเสกฺขํ ธมฺมํ สํสฏฺโฐ เนว\-เสกฺข\-นา\-เสกฺโข ธมฺโม อุปฺปชฺชติ นอธิปติ\-ปจฺจยา. (ปริปุณฺณํ. เอกํ.)}

\csromanheader{2}{11. เสกฺขตฺติก \textbf{นปุเรชาตาทิ}}
\csromanmatikaanchor{mtk.27}
\csromanmatikaanchor{mtk.28}
\csromanmatikaanchor{mtk.29}
\csromantocmarkref{subparagraph}{2. ปจฺจยปจฺจนีย 2. สงฺขฺยาวาร}{mtk.27}
\bookmark[dest={mtk.27},level=5]{2. ปจฺจยปจฺจนีย 2. สงฺขฺยาวาร}
\csromantocmarkref{subparagraph}{3. ปจฺจยานุโลมปจฺจนีย}{mtk.28}
\bookmark[dest={mtk.28},level=5]{3. ปจฺจยานุโลมปจฺจนีย}
\csromantocmarkref{subparagraph}{4. ปจฺจยปจฺจนียานุโลม}{mtk.29}
\bookmark[dest={mtk.29},level=5]{4. ปจฺจยปจฺจนียานุโลม}
\proseitem{38}{\csromanfolio{289}เสกฺขํ ธมฺมํ สํสฏฺโฐ เสกฺโข ธมฺโม อุปฺปชฺชติ นปุเรชาต\-ปจฺจยา. นปจฺฉาชาต\-ปจฺจยา. นอาเสวน\-ปจฺจยา. นกมฺมปจฺจยา. (เทฺว กาตพฺพา.) นวิปาก\-ปจฺจยา. (เทฺว กาตพฺพา.) นฌานปจฺจยา. นมคฺคปจฺจยา, นวิปฺปยุตฺต\-ปจฺจยา ฯเปฯ.}

\csromanmarkh{2. ปจฺจย\-ปจฺจนีย}\csromanmarkhii{2. สงฺขฺยาวาร}\csromanheadernomark{5}{2. ปจฺจย\-ปจฺจนีย 2. สงฺขฺยาวาร}
\proseitemwithcloserruled{39}{นเหตุยา เอกํ, นอธิปติยา ตีณิ, นปุเรชาเต ตีณิ, นปจฺฉาชาเต ตีณิ, นอาเสวเน ตีณิ, นกมฺเม เทฺว, นวิปาเก เทฺว, นฌาเน เอกํ, นมคฺเค เอกํ, นวิปฺปยุตฺเต ตีณิ.  (เอวํ คเณตพฺพํ.)}{ปจฺจนียํ.}
\proseitemwithcloserruled{40}{เหตุปจฺจยา นอธิปติยา ตีณิ, นปุเรชาเต ตีณิ, นปจฺฉาชาเต, นอาเสวเน ตีณิ, นกมฺเม เทฺว, นวิปาเก เทฺว, นวิปฺปยุตฺเต ตีณิ.  (เอวํ คเณตพฺพํ.)}{อนุโลม\-ปจฺจนียํ.}
\proseitem{41}{นเหตุปจฺจยา อารมฺมเณ เอกํ, อนนฺตเร, สมนนฺตเร, สหชาเต, อญฺญมญฺเญ, นิสฺสเย, อุปนิสฺสเย, ปุเรชาเต, อาเสวเน, กมฺเม, วิปาเก, อาหาเร, อินฺทฺริเย, ฌาเน, มคฺเค, สมฺปยุตฺเต, วิปฺปยุตฺเต, อตฺถิยา, นตฺถิยา, วิคเต, อวิคเต เอกํ.  (เอวํ คเณตพฺพํ.) ปจฺจนียา\-นุโลมํ. สํสฏฺฐวาโร. (สมฺปยุตฺตวาโร สํสฏฺฐวาร\-สทิโส.)}

\csromanmatikaanchor{mtk.30}
\csromantocmarknopage{subparagraph}{11. เสกฺขตฺติก 7. ปญฺหาวาร}{mtk.30}
\bookmark[dest={mtk.30},level=5]{11. เสกฺขตฺติก 7. ปญฺหาวาร}
\csromanmarkh{11. เสกฺขตฺติก}\csromanmarkhii{7. ปญฺหาวาร}\csromanheadernomark{6}{11. \textbf{เสกฺขตฺติก 7. ปญฺหาวาร}}
\csromanmarkh{1. ปจฺจยานุโลม}\csromanmarkhii{1. วิภงฺควาร}\csromanheadernomark{2}{1. \textbf{ปจฺจยานุโลม 1. วิภงฺควาร}}
\csromanmatikaanchor{mtk.31}
\csromanmatikaanchor{mtk.32}
\csromantocmarknopage{subparagraph}{1. ปจฺจยานุโลม 1. วิภงฺควาร}{mtk.31}
\bookmark[dest={mtk.31},level=5]{1. ปจฺจยานุโลม 1. วิภงฺควาร}
\csromantocmarkref{subparagraph}{เหตุ}{mtk.32}
\bookmark[dest={mtk.32},level=5]{เหตุ}
\csromanheader{6}{\textbf{เหตุ}}
\csromanmatikaanchor{mtk.33}
\csromantocmarkref{subparagraph}{อารมฺมณาทิ}{mtk.33}
\bookmark[dest={mtk.33},level=5]{อารมฺมณาทิ}
\proseitem{42}{\csromanfolio{290}เสกฺโข ธมฺโม เสกฺขสฺส ธมฺมสฺส เหตุปจฺจเยน ปจฺจโย. เสกฺขา เหตู สมฺปยุตฺตกานํ ขนฺธานํ เหตุปจฺจเยน ปจฺจโย. (1)}

\prose{เสกฺโข ธมฺโม เนว\-เสกฺข\-นาเสกฺขสฺส ธมฺมสฺส เหตุปจฺจเยน ปจฺจโย. เสกฺขา เหตู จิตฺต\-สมุฏฺฐานานํ รูปานํ เหตุปจฺจเยน ปจฺจโย. (2)}

\prose{เสกฺโข ธมฺโม เสกฺขสฺส จ เนว\-เสกฺข\-นาเสกฺขสฺส จ ธมฺมสฺส เหตุปจฺจเยน ปจฺจโย. เสกฺขา เหตู สมฺปยุตฺตกานํ ขนฺธานํ จิตฺตสมุฏฺฐานานญฺจ รูปานํ เหตุปจฺจเยน ปจฺจโย. (3)}

\prose{อเสกฺโข ธมฺโม อเสกฺขสฺส ธมฺมสฺส. (ตีณิ.)}

\prose{เนว\-เสกฺข\-นา\-เสกฺโข ธมฺโม เนว\-เสกฺข\-นาเสกฺขสฺส ธมฺมสฺส เหตุปจฺจเยน ปจฺจโย. เนว\-เสกฺข\-นาเสกฺขา เหตู สมฺปยุตฺตกานํ ขนฺธานํ จิตฺตสมุฏฺฐานานญฺจ รูปานํ เหตุปจฺจเยน ปจฺจโย. ปฏิสนฺธิ\-กฺขเณ ฯเปฯ. (1)}

\csromanheader{2}{\textbf{อารมฺมณ}}
\proseitem{43}{เสกฺโข ธมฺโม เนว\-เสกฺข\-นาเสกฺขสฺส ธมฺมสฺส อารมฺมณ\-ปจฺจเยน ปจฺจโย. อริยา มคฺคา วุฏฺฐหิตฺวา มคฺคํ ปจฺจเวกฺขนฺติ, เสกฺขํ ผลํ ปจฺจเวกฺขนฺติ. เจโตปริย\-ญาเณน เสกฺข\-จิตฺต\-สมงฺคิสฺส จิตฺตํ ชานนฺติ. เสกฺขา ขนฺธา เจโตปริย\-ญาณสฺส, ปุพฺเพ\-นิวาสา\-นุสฺสติ\-ญาณสฺส, อนาคตํสญาณสฺส, อาวชฺชนาย อารมฺมณ\-ปจฺจเยน ปจฺจโย. (1)}

\prose{อเสกฺโข ธมฺโม เนว\-เสกฺข\-นาเสกฺขสฺส ธมฺมสฺส อารมฺมณ\-ปจฺจเยน ปจฺจโย. อรหา อเสกฺขํ ผลํ ปจฺจเวกฺขติ, เจโตปริย\-ญาเณน อเสกฺข\-จิตฺต\-สมงฺคิสฺส จิตฺตํ ชานาติ. อเสกฺขา ขนฺธา เจโตปริย\-ญาณสฺส, ปุพฺเพ\-นิวาสา\-นุสฺสติ\-ญาณสฺส, อนาคตํสญาณสฺส, อาวชฺชนาย อารมฺมณ\-ปจฺจเยน ปจฺจโย. (1)}

\proseitem{44}{\csromanfolio{291}เนว\-เสกฺข\-นา\-เสกฺโข ธมฺโม เนว\-เสกฺข\-นาเสกฺขสฺส ธมฺมสฺส อารมฺมณ\-ปจฺจเยน ปจฺจโย. ทานํ ทตฺวา, สีลํ สมาทิยิตฺวา, อุโปสถกมฺมํ กตฺวา ตํ ปจฺจเวกฺขติ, ปุพฺเพ สุจิณฺณานิ ฯเปฯ. ฌานา วุฏฺฐหิตฺวา ฌานํ ปจฺจเวกฺขติ, อริยา นิพฺพานํ ปจฺจเวกฺขนฺติ. นิพฺพานํ โคตฺรภุสฺส, โวทานสฺส, อาวชฺชนาย อารมฺมณ\-ปจฺจเยน ปจฺจโย. อริยา ปหีเน กิเลเส ปจฺจเวกฺขนฺติ, วิกฺขมฺภิเต กิเลเส ปจฺจเวกฺขนฺติ, ปุพฺเพ สมุทาจิณฺเณ กิเลเส ชานนฺติ. จกฺขุํ อนิจฺจโต ทุกฺขโต อนตฺตโต วิปสฺสนฺติ, อสฺสาเทนฺติ อภินนฺทนฺติ, ตํ อารพฺภ ราโค อุปฺปชฺชติ ฯเปฯ โทมนสฺสํ อุปฺปชฺชติ. โสตํ ฯเปฯ. วตฺถุํ เนว\-เสกฺข\-นาเสกฺเข ขนฺเธ อนิจฺจโต ทุกฺขโต อนตฺตโต วิปสฺสนฺติ, อสฺสาเทนฺติ ฯเปฯ โทมนสฺสํ อุปฺปชฺชติ. ทิพฺเพน จกฺขุนา รูปํ ปสฺสติ. ทิพฺพาย โสตธาตุยา สทฺทํ สุณาติ. เจโตปริย\-ญาเณน เนวเสกฺขนาเสกฺข\-จิตฺต\-สมงฺคิสฺส จิตฺตํ ชานาติ. อากาสา\-นญฺจา\-ยตนํ วิญฺญาณญฺจา\-ยตนสฺส ฯเปฯ. อากิญฺจญฺญา\-ยตนํ เนว\-สญฺญา\-นา\-สญฺญา\-ยตนสฺส อารมฺมณ\-ปจฺจเยน ปจฺจโย. รูปายตนํ จกฺขุ\-วิญฺญาณสฺส ฯเปฯ. โผฏฺฐพฺพา\-ยตนํ กายวิญฺญาณสฺส ฯเปฯ. เนว\-เสกฺข\-นาเสกฺขา ขนฺธา อิทฺธิวิธ\-ญาณสฺส, เจโตปริย\-ญาณสฺส, ปุพฺเพ\-นิวาสา\-นุสฺสติ\-ญาณสฺส, ยถากมฺมูปคญาณสฺส, อนาคตํสญาณสฺส, อาวชฺชนาย อารมฺมณ\-ปจฺจเยน ปจฺจโย. (1)}

\prose{เนว\-เสกฺข\-นา\-เสกฺโข ธมฺโม เสกฺขสฺส ธมฺมสฺส อารมฺมณ\-ปจฺจเยน ปจฺจโย. นิพฺพานํ มคฺคสฺส, เสกฺขสฺส ผลสฺส อารมฺมณ\-ปจฺจเยน ปจฺจโย. (2)}

\prose{เนว\-เสกฺข\-นา\-เสกฺโข ธมฺโม อเสกฺขสฺส ธมฺมสฺส อารมฺมณ\-ปจฺจเยน ปจฺจโย. นิพฺพานํ อเสกฺขสฺส ผลสฺส อารมฺมณ\-ปจฺจเยน ปจฺจโย. (3)}

\proseitem{45}{เสกฺโข ธมฺโม เสกฺขสฺส ธมฺมสฺส อธิปติ\-ปจฺจเยน ปจฺจโย. สหชาตา\-ธิปติ\csromandash{}เสกฺขาธิปติ สมฺปยุตฺตกานํ ขนฺธานํ อธิปติ\-ปจฺจเยน ปจฺจโย. (1)}

\prose{เสกฺโข ธมฺโม เนว\-เสกฺข\-นาเสกฺขสฺส ธมฺมสฺส อธิปติ\-ปจฺจเยน ปจฺจโย. อารมฺมณา\-ธิปติ, สหชาตา\-ธิปติ.  อารมฺมณา\-ธิปติ\csromandash{}อริยา มคฺคา วุฏฺฐหิตฺวา มคฺคํ ครุํ กตฺวา ปจฺจเวกฺขนฺติ, เสกฺขํ ผลํ ครุํ กตฺวา \csromanfolio{292}ปจฺจเวกฺขนฺติ. สหชาตา\-ธิปติ\csromandash{}เสกฺขาธิปติ จิตฺต\-สมุฏฺฐานานํ รูปานํ อธิปติ\-ปจฺจเยน ปจฺจโย. (2)}

\prose{เสกฺโข ธมฺโม เสกฺขสฺส จ เนว\-เสกฺข\-นาเสกฺขสฺส จ ธมฺมสฺส อธิปติ\-ปจฺจเยน ปจฺจโย. สหชาตา\-ธิปติ\csromandash{}เสกฺขาธิปติ สมฺปยุตฺตกานํ ขนฺธานํ จิตฺตสมุฏฺฐานานญฺจ รูปานํ อธิปติ\-ปจฺจเยน ปจฺจโย. (3)}

\proseitem{46}{อเสกฺโข ธมฺโม อเสกฺขสฺส ธมฺมสฺส อธิปติ\-ปจฺจเยน ปจฺจโย. สหชาตา\-ธิปติ\csromandash{}อเสกฺขาธิปติ สมฺปยุตฺตกานํ ขนฺธานํ อธิปติ\-ปจฺจเยน ปจฺจโย. (1)}

\prose{อเสกฺโข ธมฺโม เนว\-เสกฺข\-นาเสกฺขสฺส ธมฺมสฺส อธิปติ\-ปจฺจเยน ปจฺจโย. อารมฺมณา\-ธิปติ, สหชาตา\-ธิปติ.  อารมฺมณา\-ธิปติ\csromandash{}อรหา อเสกฺขํ ผลํ ครุํ กตฺวา ปจฺจเวกฺขติ. สหชาตา\-ธิปติ\csromandash{}อเสกฺขาธิปติ จิตฺต\-สมุฏฺฐานานํ รูปานํ อธิปติ\-ปจฺจเยน ปจฺจโย. (2)}

\prose{อเสกฺโข ธมฺโม อเสกฺขสฺส จ เนว\-เสกฺข\-นาเสกฺขสฺส จ ธมฺมสฺส อธิปติ\-ปจฺจเยน ปจฺจโย. สหชาตา\-ธิปติ\csromandash{}อเสกฺขาธิปติ สมฺปยุตฺตกานํ ขนฺธานํ จิตฺตสมุฏฺฐานานญฺจ รูปานํ อธิปติ\-ปจฺจเยน ปจฺจโย. (3)}

\proseitem{47}{เนว\-เสกฺข\-นา\-เสกฺโข ธมฺโม เนว\-เสกฺข\-นาเสกฺขสฺส ธมฺมสฺส อธิปติ\-ปจฺจเยน ปจฺจโย. อารมฺมณา\-ธิปติ, สหชาตา\-ธิปติ.  อารมฺมณา\-ธิปติ\csromandash{}ทานํ ทตฺวา, สีลํ สมาทิยิตฺวา, อุโปสถกมฺมํ กตฺวา ตํ ครุํ กตฺวา ปจฺจเวกฺขติ, ปุพฺเพ สุจิณฺณานิ ครุํ กตฺวา ปจฺจเวกฺขติ, ฌานา วุฏฺฐหิตฺวา ฌานํ ครุํ กตฺวา ปจฺจเวกฺขติ. อริยา นิพฺพานํ ครุํ กตฺวา ปจฺจเวกฺขนฺติ, นิพฺพานํ โคตฺรภุสฺส, โวทานสฺส อธิปติ\-ปจฺจเยน ปจฺจโย. จกฺขุํ ครุํ กตฺวา อสฺสาเทติ อภินนฺทติ, ตํ ครุํ กตฺวา ราโค อุปฺปชฺชติ, ทิฏฺฐิ อุปฺปชฺชติ. โสตํ ฯเปฯ. วตฺถุํ เนว\-เสกฺข\-นาเสกฺเข ขนฺเธ ครุํ กตฺวา อสฺสาเทติ อภินนฺทติ, ตํ ครุํ กตฺวา ราโค อุปฺปชฺชติ, ทิฏฺฐิ อุปฺปชฺชติ. สหชาตา\-ธิปติ\csromandash{} เนว\-เสกฺข\-นาเสกฺขาธิปติ สมฺปยุตฺตกานํ ขนฺธานํ จิตฺตสมุฏฺฐานานญฺจ รูปานํ อธิปติ\-ปจฺจเยน ปจฺจโย. (1)}

\prose{เนว\-เสกฺข\-นา\-เสกฺโข ธมฺโม เสกฺขสฺส ธมฺมสฺส อธิปติ\-ปจฺจเยน ปจฺจโย. อารมฺมณา\-ธิปติ\csromandash{}นิพฺพานํ มคฺคสฺส, เสกฺขสฺส ผลสฺส อธิปติ\-ปจฺจเยน ปจฺจโย. (2)}

\prose{\csromanfolio{293}เนว\-เสกฺข\-นา\-เสกฺโข ธมฺโม อเสกฺขสฺส ธมฺมสฺส อธิปติ\-ปจฺจเยน ปจฺจโย. อารมฺมณา\-ธิปติ\csromandash{}นิพฺพานํ อเสกฺขสฺส ผลสฺส อธิปติ\-ปจฺจเยน ปจฺจโย. (3)}

\csromanheader{2}{\textbf{อนนฺตร}}
\proseitem{48}{เสกฺโข ธมฺโม เสกฺขสฺส ธมฺมสฺส อนนฺตร\-ปจฺจเยน ปจฺจโย. ปุริมา ปุริมา เสกฺขา ขนฺธา ปจฺฉิมานํ ปจฺฉิมานํ เสกฺขานํ ขนฺธานํ อนนฺตร\-ปจฺจเยน ปจฺจโย. มคฺโค เสกฺขสฺส ผลสฺส. เสกฺขํ ผลํ เสกฺขสฺส ผลสฺส อนนฺตร\-ปจฺจเยน ปจฺจโย. (1)}

\prose{เสกฺโข ธมฺโม อเสกฺขสฺส ธมฺมสฺส อนนฺตร\-ปจฺจเยน ปจฺจโย. มคฺโค อเสกฺขสฺส ผลสฺส อนนฺตร\-ปจฺจเยน ปจฺจโย. (2)}

\prose{เสกฺโข ธมฺโม เนว\-เสกฺข\-นาเสกฺขสฺส ธมฺมสฺส อนนฺตร\-ปจฺจเยน ปจฺจโย. เสกฺขํ ผลํ วุฏฺฐานสฺส อนนฺตร\-ปจฺจเยน ปจฺจโย. (3)}

\proseitem{49}{อเสกฺโข ธมฺโม อเสกฺขสฺส ธมฺมสฺส อนนฺตร\-ปจฺจเยน ปจฺจโย. ปุริมา ปุริมา อเสกฺขา ขนฺธา ปจฺฉิมานํ ปจฺฉิมานํ อเสกฺขานํ ขนฺธานํ อนนฺตร\-ปจฺจเยน ปจฺจโย. อเสกฺขํ ผลํ อเสกฺขสฺส ผลสฺส อนนฺตร\-ปจฺจเยน ปจฺจโย. (1)}

\prose{อเสกฺโข ธมฺโม เนว\-เสกฺข\-นาเสกฺขสฺส ธมฺมสฺส อนนฺตร\-ปจฺจเยน ปจฺจโย. อเสกฺขํ ผลํ วุฏฺฐานสฺส อนนฺตร\-ปจฺจเยน ปจฺจโย. (2)}

\proseitem{50}{เนว\-เสกฺข\-นา\-เสกฺโข ธมฺโม เนว\-เสกฺข\-นาเสกฺขสฺส ธมฺมสฺส อนนฺตร\-ปจฺจเยน ปจฺจโย. ปุริมา ปุริมา เนว\-เสกฺข\-นาเสกฺขา ขนฺธา ปจฺฉิมานํ ปจฺฉิมานํ เนว\-เสกฺข\-นาเสกฺขานํ ขนฺธานํ ฯเปฯ. อนุโลมํ โคตฺรภุสฺส. อนุโลมํ โวทานสฺส. อาวชฺชนา เนว\-เสกฺข\-นาเสกฺขานํ ขนฺธานํ อนนฺตร\-ปจฺจเยน ปจฺจโย. (1)}

\prose{เนว\-เสกฺข\-นา\-เสกฺโข ธมฺโม เสกฺขสฺส ธมฺมสฺส อนนฺตร\-ปจฺจเยน ปจฺจโย. โคตฺรภุ มคฺคสฺส. โวทานํ มคฺคสฺส. อนุโลมํ เสกฺขาย ผลสมาปตฺติยา. นิโรธา วุฏฺฐหนฺตสฺส เนว\-สญฺญา\-นา\-สญฺญา\-ยตนกุสลํ เสกฺขาย ผลสมาปตฺติยา อนนฺตร\-ปจฺจเยน ปจฺจโย. (2)}

\prose{\csromanfolio{294}เนว\-เสกฺข\-นา\-เสกฺโข ธมฺโม อเสกฺขสฺส ธมฺมสฺส อนนฺตร\-ปจฺจเยน ปจฺจโย. อนุโลมํ อเสกฺขาย ผลสมาปตฺติยา. นิโรธา วุฏฺฐหนฺตสฺส เนว\-สญฺญา\-นา\-สญฺญา\-ยตนกิริยํ อเสกฺขาย ผลสมาปตฺติยา อนนฺตร\-ปจฺจเยน ปจฺจโย. (3)}

\csromanheader{2}{\textbf{สมนนฺตร}}
\proseitem{51}{เสกฺโข ธมฺโม เสกฺขสฺส ธมฺมสฺส สมนนฺตร\-ปจฺจเยน ปจฺจโย ฯเปฯ. (อนนฺตร\-สทิสํ, อฏฺฐ ปญฺหา.)}

\csromanheader{2}{\textbf{สหชาตาทิ}}
\proseitem{52}{เสกฺโข ธมฺโม เสกฺขสฺส ธมฺมสฺส สหชาต\-ปจฺจเยน ปจฺจโย ฯเปฯ. (ปฏิจฺจวาเร สหชาต\-สทิสํ, นว ปญฺหา.) อญฺญมญฺญ\-ปจฺจเยน ปจฺจโย. (ปฏิจฺจวาเร อญฺญมญฺญ\-สทิสํ. ตีณิ. นิสฺสยปจฺจเย กุสลตฺติเก นิสฺสยปจฺจย\-สทิสํ, เตรส ปญฺหา.)}

\csromanheader{2}{\textbf{อุปนิสฺสย}}
\proseitem{53}{เสกฺโข ธมฺโม เสกฺขสฺส ธมฺมสฺส อุปนิสฺสย\-ปจฺจเยน ปจฺจโย. อนนฺตรู\-ปนิสฺสโย, ปกตู\-ปนิสฺสโย ฯเปฯ. ปกตู\-ปนิสฺสโย\csromandash{} ปฐโม มคฺโค ทุติยสฺส มคฺคสฺส อุปนิสฺสย\-ปจฺจเยน ปจฺจโย. ทุติโย มคฺโค ตติยสฺส มคฺคสฺส อุปนิสฺสย\-ปจฺจเยน ปจฺจโย. ตติโย มคฺโค จตุตฺถสฺส มคฺคสฺส อุปนิสฺสย\-ปจฺจเยน ปจฺจโย. มคฺโค เสกฺขาย ผลสมาปตฺติยา อุปนิสฺสย\-ปจฺจเยน ปจฺจโย. (1)}

\prose{เสกฺโข ธมฺโม อเสกฺขสฺส ธมฺมสฺส อุปนิสฺสย\-ปจฺจเยน ปจฺจโย. อนนฺตรู\-ปนิสฺสโย, ปกตู\-ปนิสฺสโย ฯเปฯ. ปกตู\-ปนิสฺสโย\csromandash{}มคฺโค อเสกฺขาย ผลสมาปตฺติยา อุปนิสฺสย\-ปจฺจเยน ปจฺจโย. (2)}

\prose{เสกฺโข ธมฺโม เนว\-เสกฺข\-นาเสกฺขสฺส ธมฺมสฺส อุปนิสฺสย\-ปจฺจเยน ปจฺจโย. อารมฺมณู\-ปนิสฺสโย, อนนฺตรู\-ปนิสฺสโย, ปกตู\-ปนิสฺสโย ฯเปฯ. ปกตู\-ปนิสฺสโย\csromandash{}อริยา มคฺคํ อุปนิสฺสาย อนุปฺปนฺนํ กุสล\-สมาปตฺติํ อุปฺปาเทนฺติ, อุปฺปนฺนํ สมาปชฺชนฺติ, สงฺขาเร อนิจฺจโต ฯเปฯ วิปสฺสนฺติ. มคฺโค อริยานํ อตฺถ\-ปฺปฏิสมฺภิทาย ฯเปฯ. ปฏิภาน\-ปฺปฏิสมฺภิทาย. \csromanfolio{295}ฐานาฐาน\-โกสลฺลสฺส อุปนิสฺสย\-ปจฺจเยน ปจฺจโย. เสกฺขา ผลสมาปตฺติ กายิกสฺส สุขสฺส อุปนิสฺสย\-ปจฺจเยน ปจฺจโย. (3)}

\proseitem{54}{อเสกฺโข ธมฺโม อเสกฺขสฺส ธมฺมสฺส อุปนิสฺสย\-ปจฺจเยน ปจฺจโย. อนนฺตรู\-ปนิสฺสโย\csromandash{}ปุริมา ปุริมา อเสกฺขา ขนฺธา ปจฺฉิมานํ ปจฺฉิมานํ อเสกฺขานํ ขนฺธานํ ฯเปฯ. อเสกฺขํ ผลํ อเสกฺขสฺส ผลสฺส อุปนิสฺสย\-ปจฺจเยน ปจฺจโย. (1)}

\prose{อเสกฺโข ธมฺโม เนว\-เสกฺข\-นาเสกฺขสฺส ธมฺมสฺส อุปนิสฺสย\-ปจฺจเยน ปจฺจโย. อารมฺมณู\-ปนิสฺสโย, อนนฺตรู\-ปนิสฺสโย, ปกตู\-ปนิสฺสโย ฯเปฯ. ปกตู\-ปนิสฺสโย\csromandash{}อเสกฺขา ผลสมาปตฺติ กายิกสฺส สุขสฺส อุปนิสฺสย\-ปจฺจเยน ปจฺจโย. (2)}

\proseitem{55}{เนว\-เสกฺข\-นา\-เสกฺโข ธมฺโม เนว\-เสกฺข\-นาเสกฺขสฺส ธมฺมสฺส อุปนิสฺสย\-ปจฺจเยน ปจฺจโย. อารมฺมณู\-ปนิสฺสโย, อนนฺตรู\-ปนิสฺสโย, ปกตู\-ปนิสฺสโย ฯเปฯ. ปกตู\-ปนิสฺสโย\csromandash{}เนว\-เสกฺข\-นาเสกฺขํ สทฺธํ อุปนิสฺสาย ทานํ เทติ, สีลํ ฯเปฯ อุโปสถกมฺมํ ฯเปฯ ฌานํ ฯเปฯ วิปสฺสนํ ฯเปฯ อภิญฺญํ ฯเปฯ สมาปตฺติํ อุปฺปาเทติ, มานํ ชปฺเปติ, ทิฏฺฐิํ คณฺหาติ. เนว\-เสกฺข\-นาเสกฺขํ สีลํ ฯเปฯ ปญฺญํ. ราคํ ฯเปฯ ปตฺถนํ. กายิกํ สุขํ ฯเปฯ. อุตุํ. โภชนํ. เสนาสนํ อุปนิสฺสาย ทานํ เทติ, สีลํ ฯเปฯ สมาปตฺติํ อุปฺปาเทติ. ปาณํ หนติ ฯเปฯ สํฆํ ภินฺทติ. เนว\-เสกฺข\-นาเสกฺขา สทฺธา ฯเปฯ ปญฺญา. ราโค ฯเปฯ ปตฺถนา. กายิกํ สุขํ ฯเปฯ. เสนาสนํ เนว\-เสกฺข\-นาเสกฺขาย สทฺธาย ฯเปฯ ปญฺญาย. ราคสฺส ฯเปฯ ปตฺถนาย. กายิกสฺส สุขสฺส, กายิกสฺส ทุกฺขสฺส อุปนิสฺสย\-ปจฺจเยน ปจฺจโย. ปฐมสฺส ฌานสฺส ปริกมฺมํ ปฐมสฺส ฌานสฺส อุปนิสฺสย\-ปจฺจเยน ปจฺจโย ฯเปฯ. เนว\-สญฺญา\-นา\-สญฺญา\-ยตนสฺส ปริกมฺมํ เนว\-สญฺญา\-นา\-สญฺญา\-ยตนสฺส ฯเปฯ. ปฐมํ ฌานํ ทุติยสฺส ฌานสฺส อุปนิสฺสย\-ปจฺจเยน ปจฺจโย ฯเปฯ. อากิญฺจญฺญา\-ยตนํ เนว\-สญฺญา\-นา\-สญฺญา\-ยตนสฺส อุปนิสฺสย\-ปจฺจเยน ปจฺจโย. (1)}

\prose{เนว\-เสกฺข\-นา\-เสกฺโข ธมฺโม เสกฺขสฺส ธมฺมสฺส อุปนิสฺสย\-ปจฺจเยน ปจฺจโย. อารมฺมณู\-ปนิสฺสโย, อนนฺตรู\-ปนิสฺสโย, ปกตู\-ปนิสฺสโย ฯเปฯ. ปกตู\-ปนิสฺสโย\csromandash{}ปฐมสฺส มคฺคสฺส ปริกมฺมํ ปฐมสฺส มคฺคสฺส อุปนิสฺสย\-ปจฺจเยน ปจฺจโย ฯเปฯ. จตุตฺถสฺส มคฺคสฺส ปริกมฺมํ จตุตฺถสฺส มคฺคสฺส อุปนิสฺสย\-ปจฺจเยน ปจฺจโย. (2)}

\prose{\csromanfolio{296}เนว\-เสกฺข\-นา\-เสกฺโข ธมฺโม อเสกฺขสฺส ธมฺมสฺส อุปนิสฺสย\-ปจฺจเยน ปจฺจโย. อารมฺมณู\-ปนิสฺสโย, อนนฺตรู\-ปนิสฺสโย, ปกตู\-ปนิสฺสโย ฯเปฯ. ปกตู\-ปนิสฺสโย\csromandash{}กายิกํ สุขํ. กายิกํ ทุกฺขํ. อุตุํ. โภชนํ. เสนาสนํ อเสกฺขาย ผลสมาปตฺติยา อุปนิสฺสย\-ปจฺจเยน ปจฺจโย. (3)}

\csromanheader{2}{\textbf{ปุเรชาต}}
\proseitem{56}{เนว\-เสกฺข\-นา\-เสกฺโข ธมฺโม เนว\-เสกฺข\-นาเสกฺขสฺส ธมฺมสฺส ปุเรชาต\-ปจฺจเยน ปจฺจโย. อารมฺมณ\-ปุเรชาตํ, วตฺถุ\-ปุเรชาตํ.  อารมฺมณ\-ปุเรชาตํ\csromandash{}จกฺขุํ อนิจฺจโต ฯเปฯ วิปสฺสติ, อสฺสาเทติ อภินนฺทติ, ตํ อารพฺภ ราโค อุปฺปชฺชติ ฯเปฯ โทมนสฺสํ อุปฺปชฺชติ. โสตํ ฯเปฯ. วตฺถุํ อนิจฺจโต ทุกฺขโต อนตฺตโต วิปสฺสติ ฯเปฯ โทมนสฺสํ อุปฺปชฺชติ. ทิพฺเพน จกฺขุนา รูปํ ปสฺสติ. ทิพฺพาย โสตธาตุยา สทฺทํ สุณาติ. รูปายตนํ จกฺขุ\-วิญฺญาณสฺส ฯเปฯ. โผฏฺฐพฺพา\-ยตนํ กายวิญฺญาณสฺส ปุเรชาต\-ปจฺจเยน ปจฺจโย. วตฺถุ\-ปุเรชาตํ\csromandash{} จกฺขายตนํ จกฺขุ\-วิญฺญาณสฺส ฯเปฯ. กายายตนํ กายวิญฺญาณสฺส. วตฺถุ เนว\-เสกฺข\-นาเสกฺขานํ ขนฺธานํ ปุเรชาต\-ปจฺจเยน ปจฺจโย. (1)}

\prose{เนว\-เสกฺข\-นา\-เสกฺโข ธมฺโม เสกฺขสฺส ธมฺมสฺส ปุเรชาต\-ปจฺจเยน ปจฺจโย. วตฺถุ\-ปุเรชาตํ\csromandash{}วตฺถุ เสกฺขานํ ขนฺธานํ ปุเรชาต\-ปจฺจเยน ปจฺจโย. (2)}

\prose{เนว\-เสกฺข\-นา\-เสกฺโข ธมฺโม อเสกฺขสฺส ธมฺมสฺส ปุเรชาต\-ปจฺจเยน ปจฺจโย. วตฺถุ\-ปุเรชาตํ\csromandash{}วตฺถุ อเสกฺขานํ ขนฺธานํ ปุเรชาต\-ปจฺจเยน ปจฺจโย. (3)}

\csromanheader{2}{\textbf{ปจฺฉาชาต}}
\proseitem{57}{เสกฺโข ธมฺโม เนว\-เสกฺข\-นาเสกฺขสฺส ธมฺมสฺส ปจฺฉาชาต\-ปจฺจเยน ปจฺจโย. ปจฺฉาชาตา เสกฺขา ขนฺธา ปุเรชาตสฺส อิมสฺส กายสฺส ปจฺฉาชาต\-ปจฺจเยน ปจฺจโย. (1)}

\prose{อเสกฺโข ธมฺโม เนว\-เสกฺข\-นาเสกฺขสฺส ธมฺมสฺส ปจฺฉาชาต\-ปจฺจเยน ปจฺจโย. ปจฺฉาชาตา อเสกฺขา ขนฺธา ปุเรชาตสฺส อิมสฺส กายสฺส ปจฺฉาชาต\-ปจฺจเยน ปจฺจโย. (1)}

\prose{\csromanfolio{297}เนว\-เสกฺข\-นา\-เสกฺโข ธมฺโม เนว\-เสกฺข\-นาเสกฺขสฺส ธมฺมสฺส ปจฺฉาชาต\-ปจฺจเยน ปจฺจโย. ปจฺฉาชาตา เนว\-เสกฺข\-นาเสกฺขา ขนฺธา ปุเรชาตสฺส อิมสฺส กายสฺส ปจฺฉาชาต\-ปจฺจเยน ปจฺจโย. (1)}

\csromanheader{2}{\textbf{อาเสวน}}
\proseitem{58}{เนว\-เสกฺข\-นา\-เสกฺโข ธมฺโม เนว\-เสกฺข\-นาเสกฺขสฺส ธมฺมสฺส อาเสวน\-ปจฺจเยน ปจฺจโย. ปุริมา ปุริมา เนว\-เสกฺข\-นาเสกฺขา ขนฺธา ปจฺฉิมานํ ปจฺฉิมานํ เนว\-เสกฺข\-นาเสกฺขานํ ขนฺธานํ อาเสวน\-ปจฺจเยน ปจฺจโย. อนุโลมํ โคตฺรภุสฺส. อนุโลมํ โวทานสฺส อาเสวน\-ปจฺจเยน ปจฺจโย. (1)}

\prose{เนว\-เสกฺข\-นา\-เสกฺโข ธมฺโม เสกฺขสฺส ธมฺมสฺส อาเสวน\-ปจฺจเยน ปจฺจโย. โคตฺรภุ มคฺคสฺส. โวทานํ มคฺคสฺส อาเสวน\-ปจฺจเยน ปจฺจโย. (2)}

\csromanheader{2}{\textbf{กมฺม}}
\proseitem{59}{เสกฺโข ธมฺโม เสกฺขสฺส ธมฺมสฺส กมฺมปจฺจเยน ปจฺจโย. \textbf{สหชาตา}, นานากฺขณิกา. สหชาตา เสกฺขา เจตนา สมฺปยุตฺตกานํ ขนฺธานํ กมฺมปจฺจเยน ปจฺจโย. \textbf{นานากฺขณิกา} เสกฺขา เจตนา วิปากานํ เสกฺขานํ ขนฺธานํ กมฺมปจฺจเยน ปจฺจโย. (1)}

\prose{เสกฺโข ธมฺโม อเสกฺขสฺส ธมฺมสฺส กมฺมปจฺจเยน ปจฺจโย. \textbf{นานากฺขณิกา} เสกฺขา เจตนา อเสกฺขานํ ขนฺธานํ กมฺมปจฺจเยน ปจฺจโย. (2)}

\prose{เสกฺโข ธมฺโม เนว\-เสกฺข\-นาเสกฺขสฺส ธมฺมสฺส กมฺมปจฺจเยน ปจฺจโย. สหชาตา เสกฺขา เจตนา จิตฺต\-สมุฏฺฐานานํ รูปานํ กมฺมปจฺจเยน ปจฺจโย. (3)}

\prose{เสกฺโข ธมฺโม เสกฺขสฺส จ เนว\-เสกฺข\-นาเสกฺขสฺส จ ธมฺมสฺส กมฺมปจฺจเยน ปจฺจโย. เสกฺขา เจตนา สมฺปยุตฺตกานํ ขนฺธานํ จิตฺตสมุฏฺฐานานญฺจ รูปานํ กมฺมปจฺจเยน ปจฺจโย. (4)}

\proseitem{60}{อเสกฺโข ธมฺโม อเสกฺขสฺส ธมฺมสฺส กมฺมปจฺจเยน ปจฺจโย. อเสกฺขา เจตนา สมฺปยุตฺตกานํ ขนฺธานํ กมฺมปจฺจเยน ปจฺจโย. (1)}

\prose{\csromanfolio{298}อเสกฺโข ธมฺโม เนว\-เสกฺข\-นาเสกฺขสฺส ธมฺมสฺส กมฺมปจฺจเยน ปจฺจโย. อเสกฺขา เจตนา จิตฺต\-สมุฏฺฐานานํ รูปานํ กมฺมปจฺจเยน ปจฺจโย. (2)}

\prose{อเสกฺโข ธมฺโม อเสกฺขสฺส จ เนว\-เสกฺข\-นาเสกฺขสฺส จ ธมฺมสฺส กมฺมปจฺจเยน ปจฺจโย. อเสกฺขา เจตนา สมฺปยุตฺตกานํ ขนฺธานํ จิตฺตสมุฏฺฐานานญฺจ รูปานํ กมฺมปจฺจเยน ปจฺจโย. (3)}

\proseitem{61}{เนว\-เสกฺข\-นา\-เสกฺโข ธมฺโม เนว\-เสกฺข\-นาเสกฺขสฺส ธมฺมสฺส กมฺมปจฺจเยน ปจฺจโย. สหชาตา. นานากฺขณิกา.  สหชาตา เนว\-เสกฺข\-นาเสกฺขา เจตนา สมฺปยุตฺตกานํ ขนฺธานํ จิตฺตสมุฏฺฐานานญฺจ รูปานํ กมฺมปจฺจเยน ปจฺจโย. ปฏิสนฺธิ\-กฺขเณ ฯเปฯ. นานากฺขณิกา เนว\-เสกฺข\-นาเสกฺขา เจตนา วิปากานํ เนว\-เสกฺข\-นาเสกฺขานํ ขนฺธานํ กฏตฺตา จ รูปานํ กมฺมปจฺจเยน ปจฺจโย. (1)}

\csromanheader{2}{\textbf{วิปาก}}
\proseitem{62}{เสกฺโข ธมฺโม เสกฺขสฺส ธมฺมสฺส วิปาก\-ปจฺจเยน ปจฺจโย. วิปาโก เสกฺโข เอโก ขนฺโธ ติณฺณนฺนํ ขนฺธานํ ฯเปฯ. (เสกฺขมูลเก ตีณิ.)}

\prose{อเสกฺโข ธมฺโม อเสกฺขสฺส ธมฺมสฺส วิปาก\-ปจฺจเยน ปจฺจโย. อเสกฺโข เอโก ขนฺโธ ติณฺณนฺนํ ขนฺธานํ ฯเปฯ. (อเสกฺขมูลเก ตีณิ.)}

\prose{เนว\-เสกฺข\-นา\-เสกฺโข ธมฺโม เนว\-เสกฺข\-นาเสกฺขสฺส ธมฺมสฺส วิปาก\-ปจฺจเยน ปจฺจโย. วิปาโก เนว\-เสกฺข\-นา\-เสกฺโข เอโก ขนฺโธ ติณฺณนฺนํ ขนฺธานํ ฯเปฯ. ปฏิสนฺธิ\-กฺขเณ ฯเปฯ. ขนฺธา วตฺถุสฺส วิปาก\-ปจฺจเยน ปจฺจโย. (1)}

\csromanheader{2}{\textbf{อาหาราทิ}}
\proseitem{63}{เสกฺโข ธมฺโม เสกฺขสฺส ธมฺมสฺส อาหารปจฺจเยน ปจฺจโย. อินฺทฺริย\-ปจฺจเยน ปจฺจโย. ฌานปจฺจเยน ปจฺจโย. มคฺคปจฺจเยน ปจฺจโย. สมฺปยุตฺต\-ปจฺจเยน ปจฺจโย ฯเปฯ.}

\csromanheader{2}{11. เสกฺขตฺติก \textbf{วิปฺปยุตฺต}}
\proseitem{64}{\csromanfolio{299}เสกฺโข ธมฺโม เนว\-เสกฺข\-นาเสกฺขสฺส ธมฺมสฺส วิปฺปยุตฺต\-ปจฺจเยน ปจฺจโย. สหชาตํ, ปจฺฉาชาตํ.  สหชาตา เสกฺขา ขนฺธา จิตฺต\-สมุฏฺฐานานํ รูปานํ วิปฺปยุตฺต\-ปจฺจเยน ปจฺจโย. ปจฺฉาชาตา เสกฺขา ขนฺธา ปุเรชาตสฺส อิมสฺส กายสฺส วิปฺปยุตฺต\-ปจฺจเยน ปจฺจโย. (1)}

\prose{อเสกฺโข ธมฺโม เนว\-เสกฺข\-นาเสกฺขสฺส ธมฺมสฺส วิปฺปยุตฺต\-ปจฺจเยน ปจฺจโย. สหชาตํ, ปจฺฉาชาตํ. (เสกฺขสทิสํ.)}

\prose{เนว\-เสกฺข\-นา\-เสกฺโข ธมฺโม เนว\-เสกฺข\-นาเสกฺขสฺส ธมฺมสฺส วิปฺปยุตฺต\-ปจฺจเยน ปจฺจโย. สหชาตํ, ปุเรชาตํ, ปจฺฉาชาตํ.  สหชาตา เนว\-เสกฺข\-นาเสกฺขา ขนฺธา จิตฺต\-สมุฏฺฐานานํ รูปานํ วิปฺปยุตฺต\-ปจฺจเยน ปจฺจโย. ปฏิสนฺธิ\-กฺขเณ เนว\-เสกฺข\-นาเสกฺขา ขนฺธา กฏตฺตารูปานํ วิปฺปยุตฺต\-ปจฺจเยน ปจฺจโย. ขนฺธา วตฺถุสฺส วิปฺปยุตฺต\-ปจฺจเยน ปจฺจโย. วตฺถุ ขนฺธานํ วิปฺปยุตฺต\-ปจฺจเยน ปจฺจโย. ปุเรชาตํ จกฺขายตนํ จกฺขุ\-วิญฺญาณสฺส ฯเปฯ. กายายตนํ กายวิญฺญาณสฺส. วตฺถุ เนว\-เสกฺข\-นาเสกฺขานํ ขนฺธานํ วิปฺปยุตฺต\-ปจฺจเยน ปจฺจโย. ปจฺฉาชาตา เนว\-เสกฺข\-นาเสกฺขา ขนฺธา ปุเรชาตสฺส อิมสฺส กายสฺส วิปฺปยุตฺต\-ปจฺจเยน ปจฺจโย. (1)}

\prose{เนว\-เสกฺข\-นา\-เสกฺโข ธมฺโม เสกฺขสฺส ธมฺมสฺส วิปฺปยุตฺต\-ปจฺจเยน ปจฺจโย. ปุเรชาตํ วตฺถุ เสกฺขานํ ขนฺธานํ วิปฺปยุตฺต\-ปจฺจเยน ปจฺจโย. (2)}

\prose{เนว\-เสกฺข\-นา\-เสกฺโข ธมฺโม อเสกฺขสฺส ธมฺมสฺส วิปฺปยุตฺต\-ปจฺจเยน ปจฺจโย. ปุเรชาตํ วตฺถุ อเสกฺขานํ ขนฺธานํ วิปฺปยุตฺต\-ปจฺจเยน ปจฺจโย. (3)}

\csromanheader{2}{\textbf{อตฺถิ}}
\proseitem{65}{เสกฺโข ธมฺโม เสกฺขสฺส ธมฺมสฺส อตฺถิปจฺจเยน ปจฺจโย. เสกฺโข เอโก ขนฺโธ ติณฺณนฺนํ ขนฺธานํ ฯเปฯ. (1)}

\prose{เสกฺโข ธมฺโม เนว\-เสกฺข\-นาเสกฺขสฺส ธมฺมสฺส อตฺถิปจฺจเยน ปจฺจโย. สหชาตํ, ปจฺฉาชาตํ.  สหชาตา เสกฺขา ขนฺธา จิตฺต\-สมุฏฺฐานานํ รูปานํ อตฺถิปจฺจเยน ปจฺจโย. ปจฺฉาชาตา เสกฺขา ขนฺธา ปุเรชาตสฺส อิมสฺส กายสฺส อตฺถิปจฺจเยน ปจฺจโย. (2)}

\prose{\csromanfolio{300}เสกฺโข ธมฺโม เสกฺขสฺส จ เนว\-เสกฺข\-นาเสกฺขสฺส จ ธมฺมสฺส อตฺถิปจฺจเยน ปจฺจโย. เสกฺโข เอโก ขนฺโธ ติณฺณนฺนํ ขนฺธานํ จิตฺตสมุฏฺฐานานญฺจ รูปานํ อตฺถิปจฺจเยน ปจฺจโย ฯเปฯ เทฺว ขนฺธา ฯเปฯ. (3)}

\prose{อเสกฺโข ธมฺโม อเสกฺขสฺส ธมฺมสฺส อตฺถิปจฺจเยน ปจฺจโย ฯเปฯ. ตีณิ. (เสกฺขสทิสํ.)}

\proseitem{66}{เนว\-เสกฺข\-นา\-เสกฺโข ธมฺโม เนว\-เสกฺข\-นาเสกฺขสฺส ธมฺมสฺส อตฺถิปจฺจเยน ปจฺจโย. สหชาตํ, ปุเรชาตํ, ปจฺฉาชาตํ, อาหารํ, อินฺทฺริยํ.  สหชาโต เนว\-เสกฺข\-นา\-เสกฺโข เอโก ขนฺโธ ติณฺณนฺนํ ขนฺธานํ จิตฺตสมุฏฺฐานานญฺจ รูปานํ อตฺถิปจฺจเยน ปจฺจโย ฯเปฯ เทฺว ขนฺธา ฯเปฯ. ปฏิสนฺธิ\-กฺขเณ ฯเปฯ. ขนฺธา วตฺถุสฺส อตฺถิปจฺจเยน ปจฺจโย. วตฺถุ ขนฺธานํ อตฺถิปจฺจเยน ปจฺจโย. เอกํ มหาภูตํ ฯเปฯ. พาหิรํ ฯเปฯ. อสญฺญสตฺตานํ ฯเปฯ. ปุเรชาตํ จกฺขุํ อนิจฺจโต ทุกฺขโต อนตฺตโต วิปสฺสติ, อสฺสาเทติ อภินนฺทติ, ตํ อารพฺภ ราโค อุปฺปชฺชติ ฯเปฯ โทมนสฺสํ อุปฺปชฺชติ. โสตํ ฯเปฯ. วตฺถุํ อนิจฺจโต ฯเปฯ วิปสฺสติ. ทิพฺเพน จกฺขุนา รูปํ ปสฺสติ. ทิพฺพาย โสตธาตุยา สทฺทํ สุณาติ. รูปายตนํ จกฺขุ\-วิญฺญาณสฺส ฯเปฯ. โผฏฺฐพฺพา\-ยตนํ กายวิญฺญาณสฺส ฯเปฯ. จกฺขายตนํ จกฺขุ\-วิญฺญาณสฺส ฯเปฯ. กายายตนํ กายวิญฺญาณสฺส. วตฺถุ เนว\-เสกฺข\-นาเสกฺขานํ ขนฺธานํ อตฺถิปจฺจเยน ปจฺจโย. ปจฺฉาชาตา เนว\-เสกฺข\-นาเสกฺขา ขนฺธา ปุเรชาตสฺส อิมสฺส กายสฺส อตฺถิปจฺจเยน ปจฺจโย. กพฬีกาโร อาหาโร อิมสฺส กายสฺส อตฺถิปจฺจเยน ปจฺจโย. รูปชีวิตินฺทฺริยํ กฏตฺตารูปานํ อตฺถิปจฺจเยน ปจฺจโย. (1)}

\prose{เนว\-เสกฺข\-นา\-เสกฺโข ธมฺโม เสกฺขสฺส ธมฺมสฺส อตฺถิปจฺจเยน ปจฺจโย. ปุเรชาตํ วตฺถุ เสกฺขานํ ขนฺธานํ อตฺถิปจฺจเยน ปจฺจโย. (2)}

\prose{เนว\-เสกฺข\-นา\-เสกฺโข ธมฺโม อเสกฺขสฺส ธมฺมสฺส อตฺถิปจฺจเยน ปจฺจโย. ปุเรชาตํ วตฺถุ อเสกฺขานํ ขนฺธานํ ขนฺธานํ อตฺถิปจฺจเยน ปจฺจโย. (3)}

\proseitem{67}{เสกฺโข จ เนว\-เสกฺข\-นา\-เสกฺโข จ ธมฺมา เสกฺขสฺส ธมฺมสฺส อตฺถิปจฺจเยน ปจฺจโย. สหชาตํ, ปุเรชาตํ.  สหชาโต\csromandash{} \csromanfolio{301}เสกฺโข เอโก ขนฺโธ จ วตฺถุ จ ติณฺณนฺนํ ขนฺธานํ อตฺถิปจฺจเยน ปจฺจโย ฯเปฯ เทฺว ขนฺธา ฯเปฯ. (1)}

\prose{เสกฺโข จ เนว\-เสกฺข\-นา\-เสกฺโข จ ธมฺมา เนว\-เสกฺข\-นาเสกฺขสฺส ธมฺมสฺส อตฺถิปจฺจเยน ปจฺจโย. สหชาตํ, ปจฺฉาชาตํ, อาหารํ, อินฺทฺริยํ.  สหชาตา เสกฺขา ขนฺธา จ มหาภูตา จ จิตฺต\-สมุฏฺฐานานํ รูปานํ อตฺถิปจฺจเยน ปจฺจโย. ปจฺฉาชาตา เสกฺขา ขนฺธา จ กพฬีกาโร อาหาโร จ อิมสฺส กายสฺส อตฺถิปจฺจเยน ปจฺจโย. ปจฺฉาชาตา เสกฺขา ขนฺธา จ รูปชีวิตินฺทฺริยญฺจ กฏตฺตารูปานํ อตฺถิปจฺจเยน ปจฺจโย. (2)}

\prose{อเสกฺโข จ เนว\-เสกฺข\-นา\-เสกฺโข จ ธมฺมา อเสกฺขสฺส ธมฺมสฺส อตฺถิปจฺจเยน ปจฺจโย ฯเปฯ. (เทฺว ปญฺหา กาตพฺพา, เสกฺขสทิสา.)}

\csromanmatikaanchor{mtk.34}
\csromantocmarkref{subparagraph}{1. ปจฺจยานุโลม 2. สงฺขฺยาวาร}{mtk.34}
\bookmark[dest={mtk.34},level=5]{1. ปจฺจยานุโลม 2. สงฺขฺยาวาร}
\csromanmarkh{1. ปจฺจยานุโลม}\csromanmarkhii{2. สงฺขฺยาวาร}\csromanheadernomark{6}{1. \textbf{ปจฺจยานุโลม 2. สงฺขฺยาวาร}}
\proseitemwithcloserruled{68}{เหตุยา สตฺต, อารมฺมเณ ปญฺจ, อธิปติยา นว, อนนฺตเร อฏฺฐ, สมนนฺตเร อฏฺฐ, สหชาเต นว, อญฺญมญฺเญ ตีณิ, นิสฺสเย เตรส, อุปนิสฺสเย อฏฺฐ, ปุเรชาเต ตีณิ, ปจฺฉาชาเต ตีณิ, อาเสวเน เทฺว, กมฺเม อฏฺฐ, วิปาเก, อาหาเร, อินฺทฺริเย, ฌาเน, มคฺเค สตฺต, สมฺปยุตฺเต ตีณิ, วิปฺปยุตฺเต ปญฺจ, อตฺถิยา เตรส, นตฺถิยา อฏฺฐ, วิคเต อฏฺฐ, อวิคเต เตรส.  (เอวํ คเณตพฺพํ.)}{อนุโลมํ.}
\csromanmatikaanchor{mtk.35}
\csromantocmarkref{subparagraph}{2. ปจฺจนียุทฺธาร}{mtk.35}
\bookmark[dest={mtk.35},level=5]{2. ปจฺจนียุทฺธาร}
\csromanheader{6}{2. ปจฺจนียุ\-ทฺธาร}
\proseitem{69}{เสกฺโข ธมฺโม เสกฺขสฺส ธมฺมสฺส สหชาต\-ปจฺจเยน ปจฺจโย. อุปนิสฺสย\-ปจฺจเยน ปจฺจโย. (1)}

\prose{เสกฺโข ธมฺโม อเสกฺขสฺส ธมฺมสฺส อุปนิสฺสย\-ปจฺจเยน ปจฺจโย. (2)}

\prose{เสกฺโข ธมฺโม เนว\-เสกฺข\-นาเสกฺขสฺส ธมฺมสฺส อารมฺมณ\-ปจฺจเยน ปจฺจโย. สหชาต\-ปจฺจเยน ปจฺจโย. อุปนิสฺสย\-ปจฺจเยน ปจฺจโย. ปจฺฉาชาต\-ปจฺจเยน ปจฺจโย. (3)}

\prose{เสกฺโข ธมฺโม เสกฺขสฺส จ เนว\-เสกฺข\-นาเสกฺขสฺส จ ธมฺมสฺส สหชาต\-ปจฺจเยน ปจฺจโย. (4)}

\proseitem{70}{\csromanfolio{302}อเสกฺโข ธมฺโม อเสกฺขสฺส ธมฺมสฺส สหชาต\-ปจฺจเยน ปจฺจโย. อุปนิสฺสย\-ปจฺจเยน ปจฺจโย. (1)}

\prose{อเสกฺโข ธมฺโม เนว\-เสกฺข\-นาเสกฺขสฺส ธมฺมสฺส อารมฺมณ\-ปจฺจเยน ปจฺจโย. สหชาต\-ปจฺจเยน ปจฺจโย. อุปนิสฺสย\-ปจฺจเยน ปจฺจโย. ปจฺฉาชาต\-ปจฺจเยน ปจฺจโย. (2)}

\prose{อเสกฺโข ธมฺโม อเสกฺขสฺส จ เนว\-เสกฺข\-นาเสกฺขสฺส จ ธมฺมสฺส สหชาต\-ปจฺจเยน ปจฺจโย. (3)}

\proseitem{71}{เนว\-เสกฺข\-นา\-เสกฺโข ธมฺโม เนว\-เสกฺข\-นาเสกฺขสฺส ธมฺมสฺส อารมฺมณ\-ปจฺจเยน ปจฺจโย. สหชาต\-ปจฺจเยน ปจฺจโย. อุปนิสฺสย\-ปจฺจเยน ปจฺจโย. ปุเรชาต\-ปจฺจเยน ปจฺจโย. ปจฺฉาชาต\-ปจฺจเยน ปจฺจโย. กมฺมปจฺจเยน ปจฺจโย. อาหารปจฺจเยน ปจฺจโย. อินฺทฺริย\-ปจฺจเยน ปจฺจโย. (1)}

\prose{เนว\-เสกฺข\-นา\-เสกฺโข ธมฺโม เสกฺขสฺส ธมฺมสฺส อุปนิสฺสย\-ปจฺจเยน ปจฺจโย. ปุเรชาต\-ปจฺจเยน ปจฺจโย. (2)}

\prose{เนว\-เสกฺข\-นา\-เสกฺโข ธมฺโม อเสกฺขสฺส ธมฺมสฺส อุปนิสฺสย\-ปจฺจเยน ปจฺจโย. ปุเรชาต\-ปจฺจเยน ปจฺจโย. (3)}

\proseitem{72}{เสกฺโข จ เนว\-เสกฺข\-นา\-เสกฺโข จ ธมฺมา เสกฺขสฺส ธมฺมสฺส สหชาตํ, ปุเรชาตํ. (1)}

\prose{เสกฺโข จ เนว\-เสกฺข\-นา\-เสกฺโข จ ธมฺมา เนว\-เสกฺข\-นาเสกฺขสฺส ธมฺมสฺส สหชาตํ. ปจฺฉาชาตํ, อาหารํ, อินฺทฺริยํ. (2)}

\prose{อเสกฺโข จ เนว\-เสกฺข\-นา\-เสกฺโข จ ธมฺมา อเสกฺขสฺส ธมฺมสฺส สหชาตํ, ปุเรชาตํ. (1)}

\prose{อเสกฺโข จ เนว\-เสกฺข\-นา\-เสกฺโข จ ธมฺมา เนว\-เสกฺข\-นาเสกฺขสฺส ธมฺมสฺส สหชาตํ. ปจฺฉาชาตํ, อาหารํ, อินฺทฺริยํ. (2)}

\csromanmatikaanchor{mtk.36}
\csromantocmarkref{subparagraph}{2. ปจฺจยปจฺจนีย 2. สงฺขฺยาวาร}{mtk.36}
\bookmark[dest={mtk.36},level=5]{2. ปจฺจยปจฺจนีย 2. สงฺขฺยาวาร}
\csromanmarkh{2. ปจฺจย\-ปจฺจนีย}\csromanmarkhii{2. สงฺขฺยาวาร}\csromanheadernomark{6}{2. ปจฺจย\-ปจฺจนีย 2. สงฺขฺยาวาร}
\proseitem{73}{นเหตุยา จุทฺทส, นอารมฺมเณ, นอธิปติยา, นอนนฺตเร, นสมนนฺตเร จุทฺทส, นสหชาเต ทส, นอญฺญมญฺเญ ทส, นนิสฺสเย}

\csromanmatikaanchor{mtk.37}
\csromanmatikaanchor{mtk.38}
\csromanmatikaanchor{mtk.39}
\csromanmatikaanchor{mtk.40}
\csromantocmarknopage{subparagraph}{3. ปจฺจยานุโลมปจฺจนีย}{mtk.37}
\bookmark[dest={mtk.37},level=5]{3. ปจฺจยานุโลมปจฺจนีย}
\csromantocmarkref{subparagraph}{เหตุทุก}{mtk.38}
\bookmark[dest={mtk.38},level=5]{เหตุทุก}
\csromantocmarknopage{subparagraph}{4. ปจฺจยปจฺจนียานุโลม}{mtk.39}
\bookmark[dest={mtk.39},level=5]{4. ปจฺจยปจฺจนียานุโลม}
\csromantocmarkref{subparagraph}{นเหตุทุก}{mtk.40}
\bookmark[dest={mtk.40},level=5]{นเหตุทุก}
\prosewithcloserruled{\csromanfolio{303}ทส, นฺเอาปนิสฺสเย เตรส, นปุเรชาเต ทฺวาทส, นปจฺฉาชาเต จุทฺทส, นอาเสวเน, นกมฺเม จุทฺทส, นวิปาเก ทฺวาทส, นอาหาเร, ไนนฺทฺริเย, นฌาเน, นมคฺเค จุทฺทส, นสมฺปยุตฺเต ทส, นวิปฺปยุตฺเต อฏฺฐ, โนอตฺถิยา อฏฺฐ, โนนตฺถิยา จุทฺทส, โนวิคเต จุทฺทส, โนอวิคเต อฏฺฐ.  (เอวํ คเณตพฺพํ.)}{ปจฺจนียํ.}
\proseitemwithcloserruled{74}{เหตุปจฺจยา นอารมฺมเณ สตฺต, นอธิปติยา สตฺต, นอนนฺตเร สตฺต, นสมนนฺตเร สตฺต, นอญฺญมญฺเญ ตีณิ, นอุปนิสฺสเย, นปุเรชาเต, นปจฺฉาชาเต, นอาเสวเน, นกมฺเม สตฺต, นวิปาเก จตฺตาริ, นอาหาเร, นอินฺทฺริเย, นฌาเน, นมคฺเค สตฺต, นสมฺปยุตฺเต ตีณิ, นวิปฺปยุตฺเต ตีณิ, โนนตฺถิยา สตฺต, โนวิคเต สตฺต.  (เอวํ คเณตพฺพํ.)}{อนุโลม\-ปจฺจนียํ.}
\proseitemwithcloser{75}{นเหตุปจฺจยา อารมฺมเณ ปญฺจ, อธิปติยา นว, อนนฺตเร อฏฺฐ, สมนนฺตเร อฏฺฐ, สหชาเต นว, อญฺญมญฺเญ ตีณิ, นิสฺสเย เตรส, อุปนิสฺสเย อฏฺฐ, ปุเรชาเต ตีณิ, ปจฺฉาชาเต ตีณิ, อาเสวเน เทฺว, กมฺเม อฏฺฐ, วิปาเก, อาหาเร, อินฺทฺริเย, ฌาเน, มคฺเค สตฺต, สมฺปยุตฺเต ตีณิ, วิปฺปยุตฺเต ปญฺจ, อตฺถิยา เตรส, นตฺถิยา อฏฺฐ, วิคเต อฏฺฐ, อวิคเต เตรส.  (เอวํ คเณตพฺพํ.) ปจฺจนียา\-นุโลมํ. ปญฺหาวาโร.}{เสกฺขตฺติกํ นิฏฺฐิตํ.}
\csromanmatikaanchor{mtk.41}
\csromantocmarknopage{subparagraph}{12. ปริตฺตตฺติก 1. ปฏิจฺจวาร}{mtk.41}
\bookmark[dest={mtk.41},level=5]{12. ปริตฺตตฺติก 1. ปฏิจฺจวาร}
\csromanmarkh{12. ปริตฺตตฺติก}\csromanmarkhii{1. ปฏิจฺจวาร}\csromanheadernomark{6}{12. \textbf{ปริตฺตตฺติก 1. ปฏิจฺจวาร}}
\csromanmarkh{1. ปจฺจยานุโลม}\csromanmarkhii{1. วิภงฺควาร}\csromanheadernomark{2}{1. \textbf{ปจฺจยานุโลม 1. วิภงฺควาร}}
\csromanmatikaanchor{mtk.42}
\csromanmatikaanchor{mtk.43}
\csromantocmarknopage{subparagraph}{1. ปจฺจยานุโลม 1. วิภงฺควาร}{mtk.42}
\bookmark[dest={mtk.42},level=5]{1. ปจฺจยานุโลม 1. วิภงฺควาร}
\csromantocmarkref{subparagraph}{เหตุ}{mtk.43}
\bookmark[dest={mtk.43},level=5]{เหตุ}
\csromanheader{6}{\textbf{เหตุ}}
\proseitem{1}{\csromanfolio{304}ปริตฺตํ ธมฺมํ ปฏิจฺจ ปริตฺโต ธมฺโม อุปฺปชฺชติ เหตุปจฺจยา. ปริตฺตํ เอกํ ขนฺธํ ปฏิจฺจ ตโย ขนฺธา จิตฺต\-สมุฏฺฐานญฺจ รูปํ ฯเปฯ เทฺว ขนฺเธ ฯเปฯ. ปฏิสนฺธิ\-กฺขเณ ปริตฺตํ เอกํ ขนฺธํ ปฏิจฺจ ตโย ขนฺธา กฏตฺตา จ รูปํ ฯเปฯ เทฺว ขนฺเธ ฯเปฯ. ขนฺเธ ปฏิจฺจ วตฺถุ, วตฺถุํ ปฏิจฺจ ขนฺธา. เอกํ มหาภูตํ ปฏิจฺจ ตโย มหาภูตา ฯเปฯ เทฺว มหาภูเต ปฏิจฺจ จิตฺต\-สมุฏฺฐานํ รูปํ ฯเปฯ. (1)}

\prose{ปริตฺตํ ธมฺมํ ปฏิจฺจ มหคฺคโต ธมฺโม อุปฺปชฺชติ เหตุปจฺจยา. ปฏิสนฺธิ\-กฺขเณ วตฺถุํ ปฏิจฺจ มหคฺคตา ขนฺธา. (2)}

\prose{ปริตฺตํ ธมฺมํ ปฏิจฺจ ปริตฺโต จ มหคฺคโต จ ธมฺมา อุปฺปชฺชนฺติ เหตุปจฺจยา. ปฏิสนฺธิ\-กฺขเณ วตฺถุํ ปฏิจฺจ มหคฺคตา ขนฺธา, มหาภูเต ปฏิจฺจ กฏตฺตารูปํ. (3)}

\proseitem{2}{มหคฺคตํ ธมฺมํ ปฏิจฺจ มหคฺคโต ธมฺโม อุปฺปชฺชติ เหตุปจฺจยา. มหคฺคตํ เอกํ ขนฺธํ ปฏิจฺจ ตโย ขนฺธา ฯเปฯ เทฺว ขนฺเธ ฯเปฯ. ปฏิสนฺธิ\-กฺขเณ ฯเปฯ. (1)}

\prose{มหคฺคตํ ธมฺมํ ปฏิจฺจ ปริตฺโต ธมฺโม อุปฺปชฺชติ เหตุปจฺจยา. มหคฺคเต ขนฺเธ ปฏิจฺจ จิตฺต\-สมุฏฺฐานํ รูปํ. ปฏิสนฺธิ\-กฺขเณ มหคฺคเต ขนฺเธ ปฏิจฺจ กฏตฺตารูปํ. (2)}

\prose{มหคฺคตํ ธมฺมํ ปฏิจฺจ ปริตฺโต จ มหคฺคโต จ ธมฺมา อุปฺปชฺชนฺติ เหตุปจฺจยา. มหคฺคตํ เอกํ ขนฺธํ ปฏิจฺจ ตโย ขนฺธา จิตฺต\-สมุฏฺฐานญฺจ รูปํ ฯเปฯ. ปฏิสนฺธิ\-กฺขเณ มหคฺคตํ เอกํ ขนฺธํ ปฏิจฺจ ตโย ขนฺธา กฏตฺตา จ รูปํ ฯเปฯ เทฺว ขนฺเธ ฯเปฯ. (3)}

\proseitem{3}{อปฺปมาณํ ธมฺมํ ปฏิจฺจ อปมาโณ ธมฺโม อุปฺปชฺชติ เหตุปจฺจยา. อปฺปมาณํ เอกํ ขนฺธํ ปฏิจฺจ ตโย ขนฺธา ฯเปฯ เทฺว ขนฺเธ ฯเปฯ. (1)}

\prose{อปฺปมาณํ ธมฺมํ ปฏิจฺจ ปริตฺโต ธมฺโม อุปฺปชฺชติ เหตุปจฺจยา. อปฺปมาเณ ขนฺเธ ปฏิจฺจ จิตฺต\-สมุฏฺฐานํ รูปํ. (2)}

\csromanmatikaanchor{mtk.44}
\csromantocmarkref{subparagraph}{อารมฺมณาทิ}{mtk.44}
\bookmark[dest={mtk.44},level=5]{อารมฺมณาทิ}
\prose{\csromanfolio{305}อปฺปมาณํ ธมฺมํ ปฏิจฺจ ปริตฺโต จ อปฺปมาโณ จ ธมฺมา อุปฺปชฺชนฺติ เหตุปจฺจยา. อปฺปมาณํ เอกํ ขนฺธํ ปฏิจฺจ ตโย ขนฺธา จิตฺต\-สมุฏฺฐานญฺจ รูปํ ฯเปฯ เทฺว ขนฺเธ ฯเปฯ. (3)}

\proseitem{4}{ปริตฺตญฺจ อปฺปมาณญฺจ ธมฺมํ ปฏิจฺจ ปริตฺโต ธมฺโม อุปฺปชฺชติ เหตุปจฺจยา. อปฺปมาเณ ขนฺเธ จ มหาภูเต จ ปฏิจฺจ จิตฺต\-สมุฏฺฐานํ รูปํ. (1)}

\prose{ปริตฺตญฺจ มหคฺคตญฺจ ธมฺมํ ปฏิจฺจ ปริตฺโต ธมฺโม อุปฺปชฺชติ เหตุปจฺจยา. มหคฺคเต ขนฺเธ จ มหาภูเต จ ปฏิจฺจ จิตฺต\-สมุฏฺฐานํ รูปํ. ปฏิสนฺธิ\-กฺขเณ ฯเปฯ. (1)}

\prose{ปริตฺตญฺจ มหคฺคตญฺจ ธมฺมํ ปฏิจฺจ มหคฺคโต ธมฺโม อุปฺปชฺชติ เหตุปจฺจยา. ปฏิสนฺธิ\-กฺขเณ มหคฺคตํ เอกํ ขนฺธญฺจ วตฺถุญฺจ ปฏิจฺจ ตโย ขนฺธา ฯเปฯ เทฺว ขนฺเธ ฯเปฯ. (2)}

\prose{ปริตฺตญฺจ มหคฺคตญฺจ ธมฺมํ ปฏิจฺจ ปริตฺโต จ มหคฺคโต จ ธมฺมา อุปฺปชฺชนฺติ เหตุปจฺจยา. ปฏิสนฺธิ\-กฺขเณ มหคฺคตํ เอกํ ขนฺธญฺจ วตฺถุญฺจ ปฏิจฺจ ตโย ขนฺธา ฯเปฯ เทฺว ขนฺเธ ฯเปฯ. มหคฺคเต ขนฺเธ จ มหาภูเต จ ปฏิจฺจ กฏตฺตารูปํ. (3)}

\csromanheader{2}{\textbf{อารมฺมณ}}
\proseitem{5}{ปริตฺตํ ธมฺมํ ปฏิจฺจ ปริตฺโต ธมฺโม อุปฺปชฺชติ อารมฺมณ\-ปจฺจยา. ปริตฺตํ เอกํ ขนฺธํ ปฏิจฺจ ตโย ขนฺธา ฯเปฯ เทฺว ขนฺเธ ฯเปฯ. ปฏิสนฺธิ\-กฺขเณ ฯเปฯ. วตฺถุํ ปฏิจฺจ ขนฺธา. (1)}

\prose{ปริตฺตํ ธมฺมํ ปฏิจฺจ มหคฺคโต ธมฺโม อุปฺปชฺชติ อารมฺมณ\-ปจฺจยา. ปฏิสนฺธิ\-กฺขเณ วตฺถุํ ปฏิจฺจ มหคฺคตา ขนฺธา. (2)}

\prose{มหคฺคตํ ธมฺมํ ปฏิจฺจ มหคฺคโต ธมฺโม อุปฺปชฺชติ อารมฺมณ\-ปจฺจยา. มหคฺคตํ เอกํ ขนฺธํ ปฏิจฺจ ตโย ขนฺธา ฯเปฯ เทฺว ขนฺเธ ฯเปฯ. ปฏิสนฺธิ\-กฺขเณ ฯเปฯ. (1)}

\prose{อปฺปมาณํ ธมฺมํ ปฏิจฺจ อปฺปมาโณ ธมฺโม อุปฺปชฺชติ อารมฺมณ\-ปจฺจยา. อปฺปมาณํ เอกํ ขนฺธํ ปฏิจฺจ ตโย ขนฺธา ฯเปฯ เทฺว ขนฺเธ ฯเปฯ. (1)}

\prose{ปริตฺตญฺจ มหคฺคตญฺจ ธมฺมํ ปฏิจฺจ มหคฺคโต ธมฺโม อุปฺปชฺชติ อารมฺมณ\-ปจฺจยา. ปฏิสนฺธิ\-กฺขเณ มหคฺคตํ เอกํ ขนฺธญฺจ วตฺถุญฺจ ปฏิจฺจ ตโย ขนฺธา ฯเปฯ เทฺว ขนฺเธ ฯเปฯ. (1)}

\proseitem{6}{\csromanfolio{306}ปริตฺตํ ธมฺมํ ปฏิจฺจ ปริตฺโต ธมฺโม อุปฺปชฺชติ อธิปติ\-ปจฺจยา. ปริตฺตํ เอกํ ขนฺธํ ปฏิจฺจ ตโย ขนฺธา จิตฺต\-สมุฏฺฐานญฺจ รูปํ. เทฺว ขนฺเธ ฯเปฯ. เอกํ มหาภูตํ ปฏิจฺจ ตโย มหาภูตา ฯเปฯ มหาภูเต ปฏิจฺจ จิตฺต\-สมุฏฺฐานํ รูปํ, อุปาทารูปํ. (1)}

\prose{มหคฺคตํ ธมฺมํ ปฏิจฺจ มหคฺคโต ธมฺโม อุปฺปชฺชติ อธิปติ\-ปจฺจยา. มหคฺคตํ เอกํ ขนฺธํ ปฏิจฺจ ตโย ขนฺธา ฯเปฯ เทฺว ขนฺเธ ฯเปฯ. (1)}

\prose{มหคฺคตํ ธมฺมํ ปฏิจฺจ ปริตฺโต ธมฺโม อุปฺปชฺชติ อธิปติ\-ปจฺจยา. มหคฺคเต ขนฺเธ ปฏิจฺจ จิตฺต\-สมุฏฺฐานํ รูปํ. (2)}

\prose{มหคฺคตํ ธมฺมํ ปฏิจฺจ ปริตฺโต จ มหคฺคโต จ ธมฺมา อุปฺปชฺชนฺติ อธิปติ\-ปจฺจยา. มหคฺคตํ เอกํ ขนฺธํ ปฏิจฺจ ตโย ขนฺธา จิตฺต\-สมุฏฺฐานญฺจ รูปํ ฯเปฯ เทฺว ขนฺเธ ฯเปฯ. (3)}

\prose{อปฺปมาณํ ธมฺมํ ปฏิจฺจ อปฺปมาโณ ธมฺโม อุปฺปชฺชติ อธิปติ\-ปจฺจยา. อปฺปมาณํ เอกํ ขนฺธํ ปฏิจฺจ ตโย ขนฺธา ฯเปฯ เทฺว ขนฺเธ ฯเปฯ. (1)}

\prose{อปฺปมาณํ ธมฺมํ ปฏิจฺจ ปริตฺโต ธมฺโม อุปฺปชฺชติ อธิปติ\-ปจฺจยา. อปฺปมาเณ ขนฺเธ ปฏิจฺจ จิตฺต\-สมุฏฺฐานํ รูปํ. (2)}

\prose{อปฺปมาณํ ธมฺมํ ปฏิจฺจ ปริตฺโต จ อปฺปมาโณ จ ธมฺมา อุปฺปชฺชนฺติ อธิปติ\-ปจฺจยา. อปฺปมาณํ เอกํ ขนฺธํ ปฏิจฺจ ตโย ขนฺธา จิตฺต\-สมุฏฺฐานญฺจ รูปํ ฯเปฯ เทฺว ขนฺเธ ฯเปฯ. (3)}

\proseitem{7}{ปริตฺตญฺจ อปฺปมาณญฺจ ธมฺมํ ปฏิจฺจ ปริตฺโต ธมฺโม อุปฺปชฺชติ อธิปติ\-ปจฺจยา. อปฺปมาเณ ขนฺเธ จ มหาภูเต จ ปฏิจฺจ จิตฺต\-สมุฏฺฐานํ รูปํ. (1)}

\prose{ปริตฺตญฺจ มหคฺคตญฺจ ธมฺมํ ปฏิจฺจ ปริตฺโต ธมฺโม อุปฺปชฺชติ อธิปติ\-ปจฺจยา. มหคฺคเต ขนฺเธ จ มหาภูเตจ ปฏิจฺจ จิตฺต\-สมุฏฺฐานํ รูปํ. (1)}

\csromanheader{2}{\textbf{อนนฺตราทิ}}
\csromanmatikaanchor{mtk.45}
\csromanmatikaanchor{mtk.48}
\csromantocmarkref{subparagraph}{1. ปจฺจยานุโลม 2. สงฺขฺยาวาร}{mtk.45}
\bookmark[dest={mtk.45},level=5]{1. ปจฺจยานุโลม 2. สงฺขฺยาวาร}
\proseitem{8}{ปริตฺตํ ธมฺมํ ปฏิจฺจ ปริตฺโต ธมฺโม อุปฺปชฺชติ อนนฺตร\-ปจฺจยา. สมนนฺตร\-ปจฺจยา. สหชาต\-ปจฺจยา. (สพฺเพปิ มหาภูตา กาตพฺพา.) อญฺญมญฺญ\-ปจฺจยา. นิสฺสยปจฺจยา. อุปนิสฺสย\-ปจฺจยา. ปุเรชาต\-ปจฺจยา. \csromanfolio{307}(ติสฺโส ปญฺหา กาตพฺพา.) อาเสวนปจฺจยา. (ติสฺโส ปญฺหา กาตพฺพา.) กมฺมปจฺจยา. วิปากปจฺจยา. (เตรส ปญฺหา.) อาหารปจฺจยา. อินฺทฺริยปจฺจยา. ฌานปจฺจยา. มคฺคปจฺจยา. สมฺปยุตฺต\-ปจฺจยา. วิปฺปยุตฺต\-ปจฺจยา. อตฺถิปจฺจยา. นตฺถิปจฺจยา. วิคตปจฺจยา. อวิคตปจฺจยา.}

\csromanmarkh{1. ปจฺจยานุโลม}\csromanmarkhii{2. สงฺขฺยาวาร}\csromanheadernomark{6}{1. \textbf{ปจฺจยานุโลม 2. สงฺขฺยาวาร}}
\proseitemwithcloserruled{9}{เหตุยา เตรส, อารมฺมเณ ปญฺจ, อธิปติยา นว, อนนฺตเร ปญฺจ, สมนนฺตเร ปญฺจ, สหชาเต เตรส, อญฺญมญฺเญ สตฺต, นิสฺสเย เตรส, อุปนิสฺสเย ปญฺจ, ปุเรชาเต ตีณิ, อาเสวเน ตีณิ, กมฺเม เตรส, วิปาเก เตรส, อาหาเร, อินฺทฺริเย, ฌาเน, มคฺเค เตรส, สมฺปยุตฺเต ปญฺจ, วิปฺปยุตฺเต เตรส, อตฺถิยา เตรส, นตฺถิยา ปญฺจ, วิคเต ปญฺจ, อวิคเต เตรส.}{อนุโลมํ.}
\csromanmarkh{2. ปจฺจย\-ปจฺจนีย}\csromanmarkhii{1. วิภงฺควาร}\csromanheadernomark{2}{2. ปจฺจย\-ปจฺจนีย 1. วิภงฺควาร}
\csromanmatikaanchor{mtk.46}
\csromanmatikaanchor{mtk.47}
\csromantocmarknopage{subparagraph}{2. ปจฺจยปจฺจนีย 1. วิภงฺควาร}{mtk.46}
\bookmark[dest={mtk.46},level=5]{2. ปจฺจยปจฺจนีย 1. วิภงฺควาร}
\csromantocmarkref{subparagraph}{นเหตุ}{mtk.47}
\bookmark[dest={mtk.47},level=5]{นเหตุ}
\csromantocmarkref{subparagraph}{นอารมฺมณ}{mtk.48}
\bookmark[dest={mtk.48},level=5]{นอารมฺมณ}
\csromanheader{6}{\textbf{นเหตุ}}
\proseitem{10}{ปริตฺตํ ธมฺมํ ปฏิจฺจ ปริตฺโต ธมฺโม อุปฺปชฺชติ นเหตุปจฺจยา. อเหตุกํ ปริตฺตํ เอกํ ขนฺธํ ปฏิจฺจ ตโย ขนฺธา จิตฺต\-สมุฏฺฐานญฺจ รูปํ ฯเปฯ เทฺว ขนฺเธ ฯเปฯ. อเหตุก\-ปฏิสนฺธิ\-กฺขเณ ฯเปฯ. ขนฺเธ ปฏิจฺจ วตฺถุ, วตฺถุํ ปฏิจฺจ ขนฺธา. เอกํ มหาภูตํ ฯเปฯ. พาหิรํ. อาหารสมุฏฺฐานํ. อุตุสมุฏฺฐานํ. อสญฺญสตฺตานํ เอกํ มหาภูตํ ฯเปฯ. วิจิกิจฺฉา\-สหคเต อุทฺธจฺจ\-สหคเต ขนฺเธ ปฏิจฺจ วิจิกิจฺฉา\-สหคโต อุทฺธจฺจ\-สหคโต โมโห. (1)}

\prose{นอารมฺมณ}

\proseitem{11}{ปริตฺตํ ธมฺมํ ปฏิจฺจ ปริตฺโต ธมฺโม อุปฺปชฺชติ นอารมฺมณ\-ปจฺจยา. ปริตฺเต ขนฺเธ ปฏิจฺจ จิตฺต\-สมุฏฺฐานํ รูปํ. ปฏิสนฺธิ\-กฺขเณ ปริตฺเต ขนฺเธ ปฏิจฺจ กฏตฺตารูปํ. ขนฺเธ ปฏิจฺจ วตฺถุ ฯเปฯ. เอกํ มหาภูตํ ฯเปฯ. พาหิรํ. อาหารสมุฏฺฐานํ. อุตุสมุฏฺฐานํ. อสญฺญสตฺตานํ เอกํ มหาภูตํ ฯเปฯ. (1)}

\csromanmatikaanchor{mtk.49}
\csromantocmarkref{subparagraph}{นอธิปติ}{mtk.49}
\bookmark[dest={mtk.49},level=5]{นอธิปติ}
\prose{\csromanfolio{308}มหคฺคตํ ธมฺมํ ปฏิจฺจ ปริตฺโต ธมฺโม อุปฺปชฺชติ นอารมฺมณ\-ปจฺจยา. มหคฺคเต ขนฺเธ ปฏิจฺจ จิตฺต\-สมุฏฺฐานํ รูปํ. ปฏิสนฺธิ\-กฺขเณ มหคฺคเต ขนฺเธ ปฏิจฺจ กฏตฺตารูปํ. (1)}

\prose{อปฺปมาณํ ธมฺมํ ปฏิจฺจ ปริตฺโต ธมฺโม อุปฺปชฺชติ นอารมฺมณ\-ปจฺจยา. อปฺปมาเณ ขนฺเธ ปฏิจฺจ จิตฺต\-สมุฏฺฐานํ รูปํ. (1)}

\prose{ปริตฺตญฺจ อปฺปมาณญฺจ ธมฺมํ ปฏิจฺจ ปริตฺโต ธมฺโม อุปฺปชฺชติ นอารมฺมณ\-ปจฺจยา. อปฺปมาเณ ขนฺเธ จ มหาภูเต จ ปฏิจฺจ จิตฺต\-สมุฏฺฐานํ รูปํ. (1)}

\prose{ปริตฺตญฺจ มหคฺคตญฺจ ธมฺมํ ปฏิจฺจ ปริตฺโต ธมฺโม อุปฺปชฺชติ นอารมฺมณ\-ปจฺจยา. มหคฺคเต ขนฺเธ จ มหาภูเต จ ปฏิจฺจ จิตฺต\-สมุฏฺฐานํ รูปํ. ปฏิสนฺธิ\-กฺขเณ มหคฺคเต ขนฺเธ จ มหาภูเต จ ปฏิจฺจ กฏตฺตารูปํ. (1)}

\prose{นอธิปติ}

\proseitem{12}{ปริตฺตํ ธมฺมํ ปฏิจฺจ ปริตฺโต ธมฺโม อุปฺปชฺชติ นอธิปติ\-ปจฺจยา. ปริตฺตํ เอกํ ขนฺธํ ปฏิจฺจ ตโย ขนฺธา ฯเปฯ. ปฏิสนฺธิ\-กฺขเณ ฯเปฯ. ขนฺเธ ปฏิจฺจ วตฺถุ, วตฺถุํ ปฏิจฺจ ขนฺธา. เอกํ มหาภูตํ ฯเปฯ. อสญฺญสตฺตานํ ฯเปฯ. (1)}

\prose{ปริตฺตํ ธมฺมํ ปฏิจฺจ มหคฺคโต ธมฺโม อุปฺปชฺชติ นอธิปติ\-ปจฺจยา. ปฏิสนฺธิ\-กฺขเณ วตฺถุํ ปฏิจฺจ มหคฺคตา ขนฺธา. (2)}

\prose{ปริตฺตํ ธมฺมํ ปฏิจฺจ ปริตฺโต จ มหคฺคโต จ ธมฺมา อุปฺปชฺชนฺติ นอธิปติ\-ปจฺจยา. ปฏิสนฺธิ\-กฺขเณ วตฺถุํ ปฏิจฺจ มหคฺคตา ขนฺธา. มหาภูเต ปฏิจฺจ กฏตฺตารูปํ. (3)}

\proseitem{13}{มหคฺคตํ ธมฺมํ ปฏิจฺจ มหคฺคโต ธมฺโม อุปฺปชฺชติ นอธิปติ\-ปจฺจยา. มหคฺคเต ขนฺเธ ปฏิจฺจ มหคฺคตา\-ธิปติ. วิปากํ มหคฺคตํ เอกํ ขนฺธํ ปฏิจฺจ ฯเปฯ. ปฏิสนฺธิ\-กฺขเณ ฯเปฯ. (1)}

\prose{มหคฺคตํ ธมฺมํ ปฏิจฺจ ปริตฺโต ธมฺโม อุปฺปชฺชติ นอธิปติ\-ปจฺจยา. วิปาเก มหคฺคเต ขนฺเธ ปฏิจฺจ จิตฺต\-สมุฏฺฐานํ รูปํ. ปฏิสนฺธิ\-กฺขเณ มหคฺคเต ขนฺเธ ปฏิจฺจ กฏตฺตารูปํ. (2)}

\prose{มหคฺคตํ ธมฺมํ ปฏิจฺจ ปริตฺโต จ มหคฺคโต จ ธมฺมา อุปฺปชฺชนฺติ นอธิปติ\-ปจฺจยา. วิปากํ มหคฺคตํ เอกํ ขนฺธํ ปฏิจฺจ ตโย ขนฺธา จิตฺต\-สมุฏฺฐานญฺจ รูปํ ฯเปฯ เทฺว ขนฺเธ ฯเปฯ. ปฏิสนฺธิ\-กฺขเณ ฯเปฯ. (3)}

\csromanmatikaanchor{mtk.50}
\csromantocmarkref{subparagraph}{นอนนฺตราทิ}{mtk.50}
\bookmark[dest={mtk.50},level=5]{นอนนฺตราทิ}
\proseitem{14}{\csromanfolio{309}อปฺปมาณํ ธมฺมํ ปฏิจฺจ อปฺปมาโณ ธมฺโม อุปฺปชฺชติ นอธิปติ\-ปจฺจยา. อปฺปมาเณ ขนฺเธ ปฏิจฺจ อปฺปมาณา\-ธิปติ. (1)}

\prose{ปริตฺตญฺจ มหคฺคตญฺจ ธมฺมํ ปฏิจฺจ ปริตฺโต ธมฺโม อุปฺปชฺชติ นอธิปติ\-ปจฺจยา. วิปาเก มหคฺคเต ขนฺเธ จ มหาภูเต จ ปฏิจฺจ จิตฺต\-สมุฏฺฐานํ รูปํ. ปฏิสนฺธิ\-กฺขเณ มหคฺคเต ขนฺเธ จ มหาภูเต จ ปฏิจฺจ กฏตฺตารูปํ. (1)}

\prose{ปริตฺตญฺจ มหคฺคตญฺจ ธมฺมํ ปฏิจฺจ มหคฺคโต ธมฺโม อุปฺปชฺชติ นอธิปติ\-ปจฺจยา. ปฏิสนฺธิ\-กฺขเณ มหคฺคตํ เอกํ ขนฺธญฺจ วตฺถุญฺจ ปฏิจฺจ ตโย ขนฺธา ฯเปฯ เทฺว ขนฺเธ ฯเปฯ. (2)}

\prose{ปริตฺตญฺจ มหคฺคตญฺจ ธมฺมํ ปฏิจฺจ ปริตฺโต จ มหคฺคโต จ ธมฺมา อุปฺปชฺชนฺติ นอธิปติ\-ปจฺจยา. ปฏิสนฺธิ\-กฺขเณ มหคฺคตํ เอกํ ขนฺธญฺจ วตฺถุญฺจ ปฏิจฺจ ตโย ขนฺธา ฯเปฯ เทฺว ขนฺเธ ฯเปฯ. มหคฺคเต ขนฺเธ จ มหาภูเต จ ปฏิจฺจ กฏตฺตารูปํ. (3)}

\prose{นอนนฺตราทิ}

\proseitem{15}{ปริตฺตํ ธมฺมํ ปฏิจฺจ ปริตฺโต ธมฺโม อุปฺปชฺชติ นอนนฺตร\-ปจฺจยา. นสมนนฺตร\-ปจฺจยา. นอญฺญมญฺญ\-ปจฺจยา. นอุปนิสฺสยปจฺจยา.}

\csromanheader{2}{\textbf{นปุเรชาต}}
\proseitem{16}{ปริตฺตํ ธมฺมํ ปฏิจฺจ ปริตฺโต ธมฺโม อุปฺปชฺชติ นปุเรชาต\-ปจฺจยา. อรูเป ปริตฺตํ เอกํ ขนฺธํ ปฏิจฺจ ตโย ขนฺธา ฯเปฯ เทฺว ขนฺเธ ฯเปฯ. ปริตฺเต ขนฺเธ ปฏิจฺจ จิตฺต\-สมุฏฺฐานํ รูปํ. ปฏิสนฺธิ\-กฺขเณ ปริตฺตํ เอกํ ขนฺธํ ปฏิจฺจ ตโย ขนฺธา กฏตฺตา จ รูปํ ฯเปฯ เทฺว ขนฺเธ ฯเปฯ. (สพฺเพ มหาภูตา วิตฺถาเรตพฺพา. อรูเป ปริตฺตมูลเก ติสฺโส ปญฺหา.)}

\prose{มหคฺคตํ ธมฺมํ ปฏิจฺจ มหคฺคโต ธมฺโม อุปฺปชฺชติ นปุเรชาต\-ปจฺจยา. มหคฺคตํ เอกํ ขนฺธํ ฯเปฯ. ปฏิสนฺธิ\-กฺขเณ ฯเปฯ. (1)}

\prose{มหคฺคตํ ธมฺมํ ปฏิจฺจ ปริตฺโต ธมฺโม อุปฺปชฺชติ นปุเรชาต\-ปจฺจยา. มหคฺคเต ขนฺเธ ปฏิจฺจ จิตฺต\-สมุฏฺฐานํ รูปํ. ปฏิสนฺธิ\-กฺขเณ มหคฺคเต ขนฺเธ ปฏิจฺจ กฏตฺตารูปํ. (2)}

\prose{มหคฺคตํ ธมฺมํ ปฏิจฺจ ปริตฺโต จ มหคฺคโต จ ธมฺมา อุปฺปชฺชนฺติ นปุเรชาต\-ปจฺจยา. ปฏิสนฺธิ\-กฺขเณ มหคฺคตํ เอกํ ขนฺธํ ปฏิจฺจ ตโย ขนฺธา กฏตฺตา จ รูปํ ฯเปฯ เทฺว ขนฺเธ ฯเปฯ. (3)}

\proseitem{17}{\csromanfolio{310}อปฺปมาณํ ธมฺมํ ปฏิจฺจ อปฺปมาโณ ธมฺโม อุปฺปชฺชติ นปุเรชาต\-ปจฺจยา. อรูเป อปฺปมาณํ เอกํ ขนฺธํ ฯเปฯ. (1)}

\prose{อปฺปมาณํ ธมฺมํ ปฏิจฺจ ปริตฺโต ธมฺโม อุปฺปชฺชติ นปุเรชาต\-ปจฺจยา. อปฺปมาเณ ขนฺเธ ปฏิจฺจ จิตฺต\-สมุฏฺฐานํ รูปํ. (2)}

\prose{ปริตฺตญฺจ อปฺปมาณญฺจ ธมฺมํ ปฏิจฺจ ปริตฺโต ธมฺโม อุปฺปชฺชติ นปุเรชาต\-ปจฺจยา. อปฺปมาเณ ขนฺเธ จ มหาภูเต จ ปฏิจฺจ จิตฺต\-สมุฏฺฐานํ รูปํ. (3)}

\prose{ปริตฺตญฺจ มหคฺคตญฺจ ธมฺมํ ปฏิจฺจ ปริตฺโต ธมฺโม อุปฺปชฺชติ นปุเรชาต\-ปจฺจยา. มหคฺคเต ขนฺเธ จ มหาภูเต จ ปฏิจฺจ จิตฺต\-สมุฏฺฐานํ รูปํ. ปฏิสนฺธิ\-กฺขเณ ฯเปฯ. (1)}

\prose{ปริตฺตญฺจ มหคฺคตญฺจ ธมฺมํ ปฏิจฺจ มหคฺคโต ธมฺโม อุปฺปชฺชติ นปุเรชาต\-ปจฺจยา. ปฏิสนฺธิ\-กฺขเณ มหคฺคตํ เอกํ ขนฺธญฺจ วตฺถุญฺจ ปฏิจฺจ ตโย ขนฺธา ฯเปฯ. (2)}

\prose{ปริตฺตญฺจ มหคฺคตญฺจ ธมฺมํ ปฏิจฺจ ปริตฺโต จ มหคฺคโต จ ธมฺมา อุปฺปชฺชนฺติ นปุเรชาต\-ปจฺจยา. ปฏิสนฺธิ\-กฺขเณ มหคฺคตํ เอกํ ขนฺธญฺจ วตฺถุญฺจ ปฏิจฺจ ตโย ขนฺธา ฯเปฯ เทฺว ขนฺเธ ฯเปฯ. มหคฺคเต ขนฺเธ จ มหาภูเต จ ปฏิจฺจ กฏตฺตารูปํ. (3)}

\csromanheader{2}{\textbf{นปจฺฉาชาตาทิ}}
\proseitem{18}{ปริตฺตํ ธมฺมํ ปฏิจฺจ ปริตฺโต ธมฺโม อุปฺปชฺชติ นปจฺฉาชาต\-ปจฺจยา. นอาเสวน\-ปจฺจยา. ปริตฺตํ เอกํ ขนฺธํ ปฏิจฺจ ตโย ขนฺธา จิตฺต\-สมุฏฺฐานญฺจ รูปํ ฯเปฯ เทฺว ขนฺเธ ฯเปฯ. ปฏิสนฺธิ\-กฺขเณ ฯเปฯ. ขนฺเธ ปฏิจฺจ วตฺถุ, วตฺถุํ ปฏิจฺจ ขนฺธา. เอกํ มหาภูตํ ฯเปฯ. อสญฺญสตฺตานํ เอกํ มหาภูตํ ฯเปฯ. (1)}

\prose{ปริตฺตํ ธมฺมํ ปฏิจฺจ มหคฺคโต ธมฺโม อุปฺปชฺชติ นอาเสวน\-ปจฺจยา. ปฏิสนฺธิ\-กฺขเณ วตฺถุํ ปฏิจฺจ มหคฺคตา ขนฺธา. (2)}

\prose{ปริตฺตํ ธมฺมํ ปฏิจฺจ ปริตฺโต จ มหคฺคโต จ ธมฺมา อุปฺปชฺชนฺติ นอาเสวน\-ปจฺจยา. ปฏิสนฺธิ\-กฺขเณ วตฺถุํ ปฏิจฺจ มหคฺคตา ขนฺธา. มหาภูเต ปฏิจฺจ กฏตฺตารูปํ. (3)}

\proseitem{19}{มหคฺคตํ ธมฺมํ ปฏิจฺจ มหคฺคโต ธมฺโม อุปฺปชฺชติ นอาเสวน\-ปจฺจยา. วิปากํ มหคฺคตํ เอกํ ขนฺธํ ปฏิจฺจ ตโย ขนฺธา ฯเปฯ เทฺว ขนฺเธ ฯเปฯ. ปฏิสนฺธิ\-กฺขเณ มหคฺคตํ เอกํ ขนฺธํ ปฏิจฺจ ตโย ขนฺธา ฯเปฯ. (1)}

\prose{\csromanfolio{311}มหคฺคตํ ธมฺมํ ปฏิจฺจ ปริตฺโต ธมฺโม อุปฺปชฺชติ นอาเสวน\-ปจฺจยา. มหคฺคเต ขนฺเธ ปฏิจฺจ จิตฺต\-สมุฏฺฐานํ รูปํ. ปฏิสนฺธิ\-กฺขเณ ฯเปฯ. (2)}

\prose{มหคฺคตํ ธมฺมํ ปฏิจฺจ ปริตฺโต จ มหคฺคโต จ ธมฺมา อุปฺปชฺชนฺติ นอาเสวน\-ปจฺจยา. วิปากํ มหคฺคตํ เอกํ ขนฺธํ ปฏิจฺจ ตโย ขนฺธา จิตฺต\-สมุฏฺฐานญฺจ รูปํ ฯเปฯ เทฺว ขนฺเธ ฯเปฯ. ปฏิสนฺธิ\-กฺขเณ มหคฺคตํ เอกํ ขนฺธํ ฯเปฯ. (3)}

\proseitem{20}{อปฺปมาณํ ธมฺมํ ปฏิจฺจ อปฺปมาโณ ธมฺโม อุปฺปชฺชติ นอาเสวน\-ปจฺจยา. วิปากํ อปฺปมาณํ เอกํ ขนฺธํ ปฏิจฺจ ตโย ขนฺธา ฯเปฯ เทฺว ขนฺเธ ปฏิจฺจ เทฺว ขนฺธา. (1)}

\prose{อปฺปมาณํ ธมฺมํ ปฏิจฺจ ปริตฺโต ธมฺโม อุปฺปชฺชติ นอาเสวน\-ปจฺจยา. อปฺปมาเณ ขนฺเธ ปฏิจฺจ จิตฺต\-สมุฏฺฐานํ รูปํ. (2)}

\prose{อปฺปมาณํ ธมฺมํ ปฏิจฺจ ปริตฺโต จ อปฺปมาโณ จ ธมฺมา อุปฺปชฺชนฺติ นอาเสวน\-ปจฺจยา. วิปากํ อปฺปมาณํ เอกํ ขนฺธํ ปฏิจฺจ ตโย ขนฺธา จิตฺต\-สมุฏฺฐานญฺจ รูปํ ฯเปฯ. (3)}

\proseitem{21}{ปริตฺตญฺจ อปฺปมาณญฺจ ธมฺมํ ปฏิจฺจ ปริตฺโต ธมฺโม อุปฺปชฺชติ นอาเสวน\-ปจฺจยา. อปฺปมาเณ ขนฺเธ จ มหาภูเต จ ปฏิจฺจ จิตฺต\-สมุฏฺฐานํ รูปํ. (1)}

\prose{ปริตฺตญฺจ มหคฺคตญฺจ ธมฺมํ ปฏิจฺจ ปริตฺโต ธมฺโม อุปฺปชฺชติ นอาเสวน\-ปจฺจยา. มหคฺคเต ขนฺเธ จ มหาภูเต จ ปฏิจฺจ จิตฺต\-สมุฏฺฐานํ รูปํ. ปฏิสนฺธิ\-กฺขเณ มหคฺคเต ขนฺเธ จ มหาภูเต จ ปฏิจฺจ กฏตฺตารูปํ. (1)}

\prose{ปริตฺตญฺจ มหคฺคตญฺจ ธมฺมํ ปฏิจฺจ มหคฺคโต ธมฺโม อุปฺปชฺชติ นอาเสวน\-ปจฺจยา. ปฏิสนฺธิ\-กฺขเณ มหคฺคตํ เอกํ ขนฺธญฺจ วตฺถุญฺจ ปฏิจฺจ ตโย ขนฺธา ฯเปฯ เทฺว ขนฺเธ ฯเปฯ. (2)}

\prose{ปริตฺตญฺจ มหคฺคตญฺจ ธมฺมํ ปฏิจฺจ ปริตฺโต จ มหคฺคโต จ ธมฺมา อุปฺปชฺชนฺติ นอาเสวน\-ปจฺจยา. ปฏิสนฺธิ\-กฺขเณ มหคฺคตํ เอกํ ขนฺธญฺจ วตฺถุญฺจ ปฏิจฺจ ตโย ขนฺธา ฯเปฯ เทฺว ขนฺเธ ฯเปฯ. มหคฺคเต ขนฺเธ จ มหาภูเต จ ปฏิจฺจ กฏตฺตารูปํ. (3)}

\csromanheader{2}{\textbf{นกมฺม}}
\proseitem{22}{\csromanfolio{312}ปริตฺตํ ธมฺมํ ปฏิจฺจ ปริตฺโต ธมฺโม อุปฺปชฺชติ นกมฺมปจฺจยา. ปริตฺเต ขนฺเธ ปฏิจฺจ ปริตฺตา เจตนา. พาหิรํ. อาหารสมุฏฺฐานํ. อุตุสมุฏฺฐานํ. เอกํ มหาภูตํ ฯเปฯ. (1)}

\prose{มหคฺคตํ ธมฺมํ ปฏิจฺจ มหคฺคโต ธมฺโม อุปฺปชฺชติ นกมฺมปจฺจยา. มหคฺคเต ขนฺเธ ปฏิจฺจ มหคฺคตา เจตนา. (1)}

\prose{อปฺปมาณํ ธมฺมํ ปฏิจฺจ อปฺปมาโณ ธมฺโม อุปฺปชฺชติ นกมฺมปจฺจยา. กุสเล อปฺปมาเณ ขนฺเธ ปฏิจฺจ อปฺปมาณา เจตนา. (1)}

\proseitem{23}{ปริตฺตํ ธมฺมํ ปฏิจฺจ ปริตฺโต ธมฺโม อุปฺปชฺชติ นวิปาก\-ปจฺจยา. ปริตฺตํ เอกํ ขนฺธํ ปฏิจฺจ ตโย ขนฺธา จิตฺต\-สมุฏฺฐานญฺจ รูปํ ฯเปฯ เทฺว ขนฺเธ ฯเปฯ. เอกํ มหาภูตํ ปฏิจฺจ ตโย มหาภูตา ฯเปฯ มหาภูเต ปฏิจฺจ จิตฺต\-สมุฏฺฐานํ รูปํ, อุปาทารูปํ. พาหิรํ. อาหารสมุฏฺฐานํ. อุตุสมุฏฺฐานํ. อสญฺญสตฺตานํ เอกํ มหาภูตํ ฯเปฯ. (1)}

\prose{มหคฺคตํ ธมฺมํ ปฏิจฺจ มหคฺคโต ธมฺโม อุปฺปชฺชติ นวิปาก\-ปจฺจยา. มหคฺคตํ เอกํ ขนฺธํ ปฏิจฺจ ตโย ขนฺธา ฯเปฯ เทฺว ขนฺเธ ฯเปฯ. (1)}

\prose{มหคฺคตํ ธมฺมํ ปฏิจฺจ ปริตฺโต ธมฺโม อุปฺปชฺชติ นวิปาก\-ปจฺจยา มหคฺคเต ขนฺเธ ปฏิจฺจ จิตฺต\-สมุฏฺฐานํ รูปํ. (2)}

\prose{มหคฺคตํ ธมฺมํ ปฏิจฺจ ปริตฺโต จ มหคฺคโต จ ธมฺมา อุปฺปชฺชนฺติ นวิปาก\-ปจฺจยา. มหคฺคตํ เอกํ ขนฺธํ ปฏิจฺจ ตโย ขนฺธา จิตฺต\-สมุฏฺฐานญฺจ รูปํ ฯเปฯ เทฺว ขนฺเธ ฯเปฯ. (3)}

\prose{อปฺปมาณํ ธมฺมํ ปฏิจฺจ อปฺปมาโณ ธมฺโม อุปฺปชฺชติ นวิปาก\-ปจฺจยา. กุสลํ อปฺปมาณํ เอกํ ขนฺธํ ปฏิจฺจ ตโย ขนฺธา ฯเปฯ. (1)}

\prose{อปฺปมาณํ ธมฺมํ ปฏิจฺจ ปริตฺโต ธมฺโม อุปฺปชฺชติ นวิปาก\-ปจฺจยา. กุสเล อปฺปมาเณ ขนฺเธ ปฏิจฺจ จิตฺต\-สมุฏฺฐานํ รูปํ. (2)}

\prose{อปฺปมาณํ ธมฺมํ ปฏิจฺจ ปริตฺโต จ อปฺปมาโณ จ ธมฺมา อุปฺปชฺชนฺติ นวิปาก\-ปจฺจยา. กุสลํ อปฺปมาณํ เอกํ ขนฺธํ ปฏิจฺจ ตโย ขนฺธา จิตฺต\-สมุฏฺฐานญฺจ รูปํ ฯเปฯ เทฺว ขนฺเธ ฯเปฯ. (3)}

\prose{\csromanfolio{313}ปริตฺตญฺจ อปฺปมาณญฺจ ธมฺมํ ปฏิจฺจ ปริตฺโต ธมฺโม อุปฺปชฺชติ นวิปาก\-ปจฺจยา. กุสเล อปฺปมาเณ ขนฺเธ จ มหาภูเต จ ปฏิจฺจ จิตฺต\-สมุฏฺฐานํ รูปํ. (1)}

\prose{ปริตฺตญฺจ มหคฺคตญฺจ ธมฺมํ ปฏิจฺจ ปริตฺโต ธมฺโม อุปฺปชฺชติ นวิปาก\-ปจฺจยา. มหคฺคเต ขนฺเธ จ มหาภูเต จ ปฏิจฺจ จิตฺต\-สมุฏฺฐานํ รูปํ. (1)}

\prose{นอาหาราทิ}

\proseitem{24}{ปริตฺตํ ธมฺมํ ปฏิจฺจ ปริตฺโต ธมฺโม อุปฺปชฺชติ นอาหารปจฺจยา. พาหิรํ. อุตุสมุฏฺฐานํ. อสญฺญสตฺตานํ. (วิตฺถาเรตพฺพํ.) นอินฺทฺริยปจฺจยา. พาหิรํ. อาหารสมุฏฺฐานํ. อุตุสมุฏฺฐานํ. อสญฺญสตฺตานํ ฯเปฯ มหาภูเต ปฏิจฺจ รูปชีวิตินฺทฺริยํ. นฌานปจฺจยา. ปญฺจ\-วิญฺญาณสหคตํ เอกํ ขนฺธํ ฯเปฯ. พาหิรํ ฯเปฯ. อสญฺญสตฺตานํ เอกํ มหาภูตํ ฯเปฯ. (สพฺเพ มหาภูตา กาตพฺพา.) นมคฺคปจฺจยา. อเหตุกํ ปริตฺตํ เอกํ ขนฺธํ ฯเปฯ. อเหตุก\-ปฏิสนฺธิ\-กฺขเณ เอกํ ฯเปฯ. (สพฺเพ มหาภูตา กาตพฺพา.) นสมฺปยุตฺต\-ปจฺจยา. นวิปฺปยุตฺต\-ปจฺจยา. อรูเป ปริตฺตํ เอกํ ขนฺธํ ปฏิจฺจ ตโย ขนฺธา ฯเปฯ เทฺว ขนฺเธ ฯเปฯ. พาหิรํ. อาหารสมุฏฺฐานํ. อุตุสมุฏฺฐานํ. อสญฺญสตฺตานํ ฯเปฯ.}

\prose{มหคฺคตํ ธมฺมํ ปฏิจฺจ มหคฺคโต ธมฺโม อุปฺปชฺชติ นวิปฺปยุตฺต\-ปจฺจยา. อรูเป มหคฺคตํ เอกํ ขนฺธํ ฯเปฯ. (1)}

\prose{อปฺปมาณํ ธมฺมํ ปฏิจฺจ อปฺปมาโณ ธมฺโม อุปฺปชฺชติ นวิปฺปยุตฺต\-ปจฺจยา. อรูเป อปฺปมาณํ เอกํ ขนฺธํ ฯเปฯ. โนนตฺถิปจฺจยา. โนวิคต\-ปจฺจยา.}

\csromanmatikaanchor{mtk.51}
\csromantocmarkref{subparagraph}{2. ปจฺจยปจฺจนีย 2. สงฺขฺยาวาร}{mtk.51}
\bookmark[dest={mtk.51},level=5]{2. ปจฺจยปจฺจนีย 2. สงฺขฺยาวาร}
\csromanmarkh{2. ปจฺจย\-ปจฺจนีย}\csromanmarkhii{2. สงฺขฺยาวาร}\csromanheadernomark{6}{2. ปจฺจย\-ปจฺจนีย 2. สงฺขฺยาวาร}
\proseitem{25}{นเหตุยา เอกํ, นอารมฺมเณ ปญฺจ, นอธิปติยา ทส, นอนนฺตเร ปญฺจ, นสมนนฺตเร ปญฺจ, นอญฺญมญฺเญ, นอุปนิสฺสเย ปญฺจ, นปุเรชาเต ทฺวาทส, นปจฺฉาชาเต เตรส, นอาเสวเน เตรส, นกมฺเม ตีณิ, นวิปาเก นว, นอาหาเร เอกํ, นอินฺทฺริเย, นฌาเน, นมคฺเค เอกํ, นสมฺปยุตฺเต ปญฺจ, นวิปฺปยุตฺเต ตีณิ, โนนตฺถิยา ปญฺจ, โนวิคเต ปญฺจ. (เอวํ คเณตพฺพํ.)}

\vspace{\tipitakaparskip}
\csromancenter{ปจฺจนียํ.}
\csromancenter{\textbf{เหตุ}ปจฺจยา นอารมฺมเณ ปญฺจ, นอธิปติยา ทส, นอนนฺตเร ปญฺจ, นสมนนฺตเร, นอญฺญมญฺเญ, นฺเอาปนิสฺสเย ปญฺจ, นปุเรชาเต ทฺวาทส, นปจฺฉาชาเต เตรส, นอาเสวเน เตรส, นกมฺเม ตีณิ, นวิปาเก นว, นสมฺปยุตฺเต ปญฺจ, นวิปฺปยุตฺเต ตีณิ, โนนตฺถิยา ปญฺจ, โนวิคเต ปญฺจ. (เอวํ คเณตพฺพํ.)}
\nitthitamruled{อนุโลม\-ปจฺจนียํ.}
\csromanmatikaanchor{mtk.52}
\csromanmatikaanchor{mtk.53}
\csromantocmarkref{subparagraph}{3. ปจฺจยานุโลมปจฺจนีย}{mtk.52}
\bookmark[dest={mtk.52},level=5]{3. ปจฺจยานุโลมปจฺจนีย}
\csromantocmarkref{subparagraph}{4. ปจฺจยปจฺจนียานุโลม}{mtk.53}
\bookmark[dest={mtk.53},level=5]{4. ปจฺจยปจฺจนียานุโลม}
\proseitem{27}{\csromanfolio{314}นเหตุปจฺจยา อารมฺมเณ เอกํ, อนนฺตเร เอกํ ฯเปฯ วิคเต เอกํ, อวิคเต เอกํ. (เอวํ คเณตพฺพํ.) ปจฺจนียา\-นุโลมํ. ปฏิจฺจวาโร. (สหชาตวาโรปิ ปฏิจฺจ\-วาร\-สทิโส.)}
\csromansectionrule

\csromanmatikaanchor{mtk.54}
\csromantocmarknopage{subparagraph}{12. ปริตฺตตฺติก 3. ปจฺจยวาร}{mtk.54}
\bookmark[dest={mtk.54},level=5]{12. ปริตฺตตฺติก 3. ปจฺจยวาร}
\csromanmarkh{12. ปริตฺตตฺติก}\csromanmarkhii{3. ปจฺจยวาร}\csromanheadernomark{6}{12. \textbf{ปริตฺตตฺติก 3. ปจฺจยวาร}}
\csromanmarkh{1. ปจฺจยานุโลม}\csromanmarkhii{1. วิภงฺควาร}\csromanheadernomark{6}{1. \textbf{ปจฺจยานุโลม 1. วิภงฺควาร}}
\csromanmatikaanchor{mtk.55}
\csromanmatikaanchor{mtk.56}
\csromantocmarknopage{subparagraph}{1. ปจฺจยานุโลม 1. วิภงฺควาร}{mtk.55}
\bookmark[dest={mtk.55},level=5]{1. ปจฺจยานุโลม 1. วิภงฺควาร}
\csromantocmarkref{subparagraph}{เหตุ}{mtk.56}
\bookmark[dest={mtk.56},level=5]{เหตุ}
\csromanheader{6}{\textbf{เหตุ}}
\proseitem{28}{ปริตฺตํ ธมฺมํ ปจฺจยา ปริตฺโต ธมฺโม อุปฺปชฺชติ เหตุปจฺจยา. ปริตฺตํ เอกํ ขนฺธํ ปจฺจยา ตโย ขนฺธา จิตฺต\-สมุฏฺฐานญฺจ รูปํ ฯเปฯ เทฺว ขนฺเธ ฯเปฯ. ปฏิสนฺธิ\-กฺขเณ ฯเปฯ. ขนฺเธ ปจฺจยา วตฺถุ, วตฺถุํ ปจฺจยา ขนฺธา. เอกํ มหาภูตํ ปจฺจยา ฯเปฯ อุปาทารูปํ. วตฺถุํ ปจฺจยา ปริตฺตา ขนฺธา. (1)}

\prose{ปริตฺตํ ธมฺมํ ปจฺจยา มหคฺคโต ธมฺโม อุปฺปชฺชติ เหตุปจฺจยา. วตฺถุํ ปจฺจยา มหคฺคตา ขนฺธา. ปฏิสนฺธิ\-กฺขเณ วตฺถุํ ปจฺจยา มหคฺคตา ขนฺธา. (2)}

\prose{ปริตฺตํ ธมฺมํ ปจฺจยา อปฺปมาโณ ธมฺโม อุปฺปชฺชติ เหตุปจฺจยา. วตฺถุํ ปจฺจยา อปฺปมาณา ขนฺธา. (3)}

\prose{\csromanfolio{315}ปริตฺตํ ธมฺมํ ปจฺจยา ปริตฺโต จ อปฺปมาโณ จ ธมฺมา อุปฺปชฺชนฺติ เหตุปจฺจยา. วตฺถุํ ปจฺจยา อปฺปมาณา ขนฺธา. มหาภูเต ปจฺจยา จิตฺต\-สมุฏฺฐานํ รูปํ. (4)}

\prose{ปริตฺตํ ธมฺมํ ปจฺจยา ปริตฺโต จ มหคฺคโต จ ธมฺมา อุปฺปชฺชนฺติ เหตุปจฺจยา. วตฺถุํ ปจฺจยา มหคฺคตา ขนฺธา. มหาภูเต ปจฺจยา จิตฺต\-สมุฏฺฐานํ รูปํ. ปฏิสนฺธิ\-กฺขเณ วตฺถุํ ปจฺจยา มหคฺคตา ขนฺธา. (5)}

\proseitem{29}{มหคฺคตํ ธมฺมํ ปจฺจยา มหคฺคโต ธมฺโม อุปฺปชฺชติ เหตุปจฺจยา. มหคฺคตํ เอกํ ขนฺธํ ปจฺจยา ตโย ขนฺธา ฯเปฯ. ปฏิสนฺธิ\-กฺขเณ มหคฺคตํ เอกํ ขนฺธํ ฯเปฯ. (1)}

\prose{มหคฺคตํ ธมฺมํ ปจฺจยา ปริตฺโต ธมฺโม อุปฺปชฺชติ เหตุปจฺจยา. มหคฺคเต ขนฺเธ ปจฺจยา จิตฺต\-สมุฏฺฐานํ รูปํ. ปฏิสนฺธิ\-กฺขเณ ฯเปฯ. (2)}

\prose{มหคฺคตํ ธมฺมํ ปจฺจยา ปริตฺโต จ มหคฺคโต จ ธมฺมา อุปฺปชฺชนฺติ เหตุปจฺจยา. มหคฺคตํ เอกํ ขนฺธํ ปจฺจยา ตโย ขนฺธา จิตฺต\-สมุฏฺฐานญฺจ รูปํ ฯเปฯ เทฺว ขนฺเธ ฯเปฯ. ปฏิสนฺธิ\-กฺขเณ มหคฺคตํ เอกํ ขนฺธํ ฯเปฯ. (3)}

\prose{อปฺปมาณํ ธมฺมํ ปจฺจยา อปฺปมาโณ ธมฺโม อุปฺปชฺชติ เหตุปจฺจยา. อปฺปมาเณ. ตีณิ.}

\proseitem{30}{ปริตฺตญฺจ อปฺปมาณญฺจ ธมฺมํ ปจฺจยา ปริตฺโต ธมฺโม อุปฺปชฺชติ เหตุปจฺจยา. อปฺปมาเณ ขนฺเธ จ มหาภูเต จ ปจฺจยา จิตฺต\-สมุฏฺฐานํ รูปํ. (1)}

\prose{ปริตฺตญฺจ อปฺปมาณญฺจ ธมฺมํ ปจฺจยา อปฺปมาโณ ธมฺโม อุปฺปชฺชติ เหตุปจฺจยา. อปฺปมาณํ เอกํ ขนฺธญฺจ วตฺถุญฺจ ปจฺจยา ตโย ขนฺธา ฯเปฯ. (2)}

\prose{ปริตฺตญฺจ อปฺปมาณญฺจ ธมฺมํ ปจฺจยา ปริตฺโต จ อปฺปมาโณ จ ธมฺมา อุปฺปชฺชนฺติ เหตุปจฺจยา. อปฺปมาณํ เอกํ ขนฺธญฺจ วตฺถุญฺจ ปจฺจยา ตโย ขนฺธา ฯเปฯ เทฺว ขนฺเธ ฯเปฯ. อปฺปมาเณ ขนฺเธ จ มหาภูเต จ ปจฺจยา จิตฺต\-สมุฏฺฐานํ รูปํ. (3)}

\prose{ปริตฺตญฺจ มหคฺคตญฺจ ธมฺมํ ปจฺจยา ปริตฺโต ธมฺโม อุปฺปชฺชติ เหตุปจฺจยา. ตีณิ. (ปฏิสนฺธิ\-กฺขเณ ตีณิปิ กาตพฺพา.)}

\csromanmatikaanchor{mtk.57}
\csromantocmarkref{subparagraph}{อารมฺมณาทิ}{mtk.57}
\bookmark[dest={mtk.57},level=5]{อารมฺมณาทิ}
\csromanheader{6}{\textbf{อารมฺมณาทิ}}
\proseitem{31}{ปริตฺตํ ธมฺมํ ปจฺจยา ปริตฺโต ธมฺโม อุปฺปชฺชติ อารมฺมณ\-ปจฺจยา. ปริตฺตํ เอกํ ขนฺธํ ปจฺจยา ตโย ขนฺธา ฯเปฯ เทฺว ขนฺเธ ฯเปฯ. ปฏิสนฺธิ\-กฺขเณ ฯเปฯ. \csromanfolio{316}วตฺถุํ ปจฺจยา ขนฺธา. จกฺขายตนํ ปจฺจยา จกฺขุวิญฺญาณํ ฯเปฯ. กายายตนํ ปจฺจยา กายวิญฺญาณํ. วตฺถุํ ปจฺจยา ปริตฺตา ขนฺธา. (อวเสสา ฉ ปญฺหา เหตุปจฺจย\-สทิสา. สตฺต กาตพฺพา.) อธิปติ\-ปจฺจยา. (ปฏิสนฺธิ นตฺถิ. สตฺตรส ปญฺหา ปริปุณฺณา.) อนนฺตร\-ปจฺจยา ฯเปฯ. อวิคตปจฺจยา.}

\csromanmatikaanchor{mtk.58}
\csromantocmarkref{subparagraph}{1. ปจฺจยานุโลม 2. สงฺขฺยาวาร}{mtk.58}
\bookmark[dest={mtk.58},level=5]{1. ปจฺจยานุโลม 2. สงฺขฺยาวาร}
\csromanmarkh{1. ปจฺจยานุโลม}\csromanmarkhii{2. สงฺขฺยาวาร}\csromanheadernomark{6}{1. \textbf{ปจฺจยานุโลม 2. สงฺขฺยาวาร}}
\proseitemwithcloserruled{32}{เหตุยา สตฺตรส, อารมฺมเณ สตฺต, อธิปติยา สตฺตรส, อนนฺตเร สตฺต, สมนนฺตเร สตฺต, สหชาเต สตฺตรส, อญฺญมญฺเญ นว, นิสฺสเย สตฺตรส, อุปนิสฺสเย สตฺต, ปุเรชาเต สตฺต, อาเสวเน สตฺต, กมฺเม สตฺตรส, วิปาเก สตฺตรส, อาหาเร สตฺตรส, อินฺทฺริเย, ฌาเน, มคฺเค สตฺตรส, สมฺปยุตฺเต สตฺต, วิปฺปยุตฺเต สตฺตรส, อตฺถิยา สตฺตรส, นตฺถิยา สตฺต, วิคเต สตฺต, อวิคเต สตฺตรส. (เอวํ คเณตพฺพํ.)}{อนุโลมํ.}
\csromanmatikaanchor{mtk.59}
\csromantocmarkref{subparagraph}{2. ปจฺจยปจฺจนีย 1. วิภงฺควาร}{mtk.59}
\bookmark[dest={mtk.59},level=5]{2. ปจฺจยปจฺจนีย 1. วิภงฺควาร}
\csromanmarkh{2. ปจฺจย\-ปจฺจนีย}\csromanmarkhii{1. วิภงฺควาร}\csromanheadernomark{6}{2. ปจฺจย\-ปจฺจนีย 1. วิภงฺควาร}
\csromanheader{2}{\textbf{นเหตุ}}
\proseitem{33}{ปริตฺตํ ธมฺมํ ปจฺจยา ปริตฺโต ธมฺโม อุปฺปชฺชติ นเหตุปจฺจยา. อเหตุกํ ปริตฺตํ เอกํ ขนฺธํ ปจฺจยา ตโย ขนฺธา จิตฺต\-สมุฏฺฐานญฺจ รูปํ ฯเปฯ เทฺว ขนฺเธ ฯเปฯ. อเหตุก\-ปฏิสนฺธิ\-กฺขเณ ฯเปฯ. ขนฺเธ ปจฺจยา วตฺถุ, วตฺถุํ ปจฺจยา ขนฺธา. เอกํ มหาภูตํ ฯเปฯ. อสญฺญสตฺตานํ ฯเปฯ. จกฺขายตนํ ปจฺจยา จกฺขุวิญฺญาณํ ฯเปฯ. กายายตนํ ปจฺจยา กายวิญฺญาณํ. วตฺถุํ ปจฺจยา อเหตุกา ปริตฺตา ขนฺธา. วิจิกิจฺฉา\-สหคเต อุทฺธจฺจ\-สหคเต ขนฺเธ จ วตฺถุญฺจ ปจฺจยา วิจิกิจฺฉา\-สหคโต อุทฺธจฺจ\-สหคโต โมโห. (1)}

\prose{นอารมฺมณ}

\proseitem{34}{ปริตฺตํ ธมฺมํ ปจฺจยา ปริตฺโต ธมฺโม อุปฺปชฺชติ นอารมฺมณ\-ปจฺจยา. (ปฏิจฺจ\-วาร\-สทิสํ. ปญฺจ.)}

\csromanheader{2}{12. ปริตฺตตฺติก นอธิปติ}
\proseitem{35}{\csromanfolio{317}ปริตฺตํ ธมฺมํ ปจฺจยา ปริตฺโต ธมฺโม อุปฺปชฺชติ นอธิปติปติปจฺจยา. ปริตฺตํ เอกํ ขนฺธํ ปจฺจยา ตโย ขนฺธา จิตฺต\-สมุฏฺฐานญฺจ รูปํ ฯเปฯ เทฺว ขนฺเธ ฯเปฯ. ปฏิสนฺธิ\-กฺขเณ ฯเปฯ. อสญฺญสตฺตานํ ฯเปฯ. จกฺขายตนํ ฯเปฯ. กายายตนํ ฯเปฯ. วตฺถุํ ปจฺจยา ปริตฺตา ขนฺธา. (1)}

\prose{ปริตฺตํ ธมฺมํ ปจฺจยา มหคฺคโต ธมฺโม อุปฺปชฺชติ นอธิปติ\-ปจฺจยา. วตฺถุํ ปจฺจยา มหคฺคตา\-ธิปติ. วตฺถุํ ปจฺจยา วิปากา มหคฺคตา ขนฺธา. ปฏิสนฺธิ\-กฺขเณ วตฺถุํ ปจฺจยา มหคฺคตา ขนฺธา. (2)}

\prose{ปริตฺตํ ธมฺมํ ปจฺจยา อปฺปมาโณ ธมฺโม อุปฺปชฺชติ นอธิปติ\-ปจฺจยา. วตฺถุํ ปจฺจยา อปฺปมาณา\-ธิปติ. (3)}

\prose{ปริตฺตํ ธมฺมํ ปจฺจยา ปริตฺโต จ มหคฺคโต จ ธมฺมา อุปฺปชฺชนฺติ นอธิปติ\-ปจฺจยา. วตฺถุํ ปจฺจยา วิปากา มหคฺคตา ขนฺธา. มหาภูเต ปจฺจยา จิตฺต\-สมุฏฺฐานํ รูปํ. ปฏิสนฺธิ\-กฺขเณ ฯเปฯ. (4)}

\proseitem{36}{มหคฺคตํ ธมฺมํ ปจฺจยา มหคฺคโต ธมฺโม อุปฺปชฺชติ นอธิปติ\-ปจฺจยา. มหคฺคเต ขนฺเธ ปจฺจยา มหคฺคตา\-ธิปติ. วิปากํ มหคฺคตํ เอกํ ขนฺธํ ปจฺจยา ตโย ขนฺธา ฯเปฯ เทฺว ขนฺเธ ฯเปฯ. ปฏิสนฺธิ\-กฺขเณ ฯเปฯ. (1)}

\prose{มหคฺคตํ ธมฺมํ ปจฺจยา ปริตฺโต ธมฺโม อุปฺปชฺชติ นอธิปติ\-ปจฺจยา. วิปาเก มหคฺคเต ขนฺเธ ปจฺจยา จิตฺต\-สมุฏฺฐานํ รูปํ. ปฏิสนฺธิ\-กฺขเณ ฯเปฯ. (2)}

\prose{มหคฺคตํ ธมฺมํ ปจฺจยา ปริตฺโต จ มหคฺคโต จ ธมฺมา อุปฺปชฺชนฺติ นอธิปติ\-ปจฺจยา. วิปากํ มหคฺคตํ เอกํ ขนฺธํ ปจฺจยา ตโย ขนฺธา จิตฺต\-สมุฏฺฐานญฺจ รูปํ ฯเปฯ เทฺว ขนฺเธ ฯเปฯ. ปฏิสนฺธิ\-กฺขเณ ฯเปฯ. (3)}

\prose{อปฺปมาณํ ธมฺมํ ปจฺจยา อปฺปมาโณ ธมฺโม อุปฺปชฺชติ นอธิปติ\-ปจฺจยา. อปฺปมาเณ ขนฺเธ ปจฺจยา อปฺปมาณา\-ธิปติ. (1)}

\proseitem{37}{ปริตฺตญฺจ อปฺปมาณญฺจ ธมฺมํ ปจฺจยา อปฺปมาโณ ธมฺโม อุปฺปชฺชติ นอธิปติ\-ปจฺจยา. อปฺปมาเณ ขนฺเธ จ วตฺถุญฺจ ปจฺจยา อปฺปมาณา\-ธิปติ. (1)}

\prose{ปริตฺตญฺจ มหคฺคตญฺจ ธมฺมํ ปจฺจยา ปริตฺโต ธมฺโม อุปฺปชฺชติ นอธิปติ\-ปจฺจยา. วิปาเก มหคฺคเต ขนฺเธ จ มหาภูเต จ ปจฺจยา จิตฺต\-สมุฏฺฐานํ รูปํ. ปฏิสนฺธิ\-กฺขเณ ฯเปฯ. (1)}

\prose{\csromanfolio{318}ปริตฺตญฺจ มหคฺคตญฺจ ธมฺมํ ปจฺจยา มหคฺคโต ธมฺโม อุปฺปชฺชติ นอธิปติ\-ปจฺจยา. มหคฺคเต ขนฺเธ จ วตฺถุญฺจ ปจฺจยา มหคฺคตา\-ธิปติ. วิปากํ มหคฺคตํ เอกํ ขนฺธญฺจ วตฺถุญฺจ ปจฺจยา ตโย ขนฺธา ฯเปฯ เทฺว ขนฺเธ ฯเปฯ. ปฏิสนฺธิ\-กฺขเณ ฯเปฯ. (2)}

\prose{ปริตฺตญฺจ มหคฺคตญฺจ ธมฺมํ ปจฺจยา ปริตฺโต จ มหคฺคโต จ ธมฺมา อุปฺปชฺชนฺติ นอธิปติ\-ปจฺจยา. วิปากํ มหคฺคตํ เอกํ ขนฺธญฺจ วตฺถุญฺจ ปจฺจยา ตโย ขนฺธา ฯเปฯ เทฺว ขนฺเธ ฯเปฯ. วิปาเก มหคฺคเต ขนฺเธ จ มหาภูเต จ ปจฺจยา จิตฺต\-สมุฏฺฐานํ รูปํ. ปฏิสนฺธิ\-กฺขเณ วิปากํ มหคฺคตํ เอกํ ขนฺธญฺจ วตฺถุญฺจ ปจฺจยา ตโย ขนฺธา ฯเปฯ เทฺว ขนฺเธ ฯเปฯ. วิปาเก มหคฺคเต ขนฺเธ จ มหาภูเต จ ปจฺจยา กฏตฺตารูปํ. (3)}

\prose{นอนนฺตราทิ}

\proseitem{38}{ปริตฺตํ ธมฺมํ ปจฺจยา ปริตฺโต ธมฺโม อุปฺปชฺชติ นอนนฺตร\-ปจฺจยา. นสมนนฺตร\-ปจฺจยา. นอญฺญมญฺญ\-ปจฺจยา. นอุปนิสฺสยปจฺจยา. นปุเรชาต\-ปจฺจยา. (ปฏิจฺจ\-วาร\-สทิสา ทฺวาทส ปญฺหา.) นปจฺฉาชาต\-ปจฺจยา. นอาเสวน\-ปจฺจยา. (ปริปุณฺณํ, วิปาโกติ นิทฺทิสิตพฺพํ. จิตฺต\-สมุฏฺฐานํ รูปํ วิปาโกติ น กาตพฺพํ.) นกมฺมปจฺจยา. นวิปาก\-ปจฺจยา. (ปฏิสนฺธิวิปาโกปิ นตฺถิ.) นอาหารปจฺจยา. นอินฺทฺริยปจฺจยา. นฌานปจฺจยา. นมคฺคปจฺจยา. นสมฺปยุตฺต\-ปจฺจยา. นวิปฺปยุตฺต\-ปจฺจยา. โนนตฺถิปจฺจยา. โนวิคต\-ปจฺจยา.}

\csromanmatikaanchor{mtk.60}
\csromantocmarkref{subparagraph}{2. ปจฺจยปจฺจนีย 2. สงฺขฺยาวาร}{mtk.60}
\bookmark[dest={mtk.60},level=5]{2. ปจฺจยปจฺจนีย 2. สงฺขฺยาวาร}
\csromanmarkh{2. ปจฺจย\-ปจฺจนีย}\csromanmarkhii{2. สงฺขฺยาวาร}\csromanheadernomark{6}{2. ปจฺจย\-ปจฺจนีย 2. สงฺขฺยาวาร}
\proseitem{39}{นเหตุยา เอกํ, นอารมฺมเณ ปญฺจ, นอธิปติยา ทฺวาทส, นอนนฺตเร ปญฺจ, นสมนนฺตเร, นอญฺญมญฺเญ, นอุปนิสฺสเย ปญฺจ, นปุเรชาเต ทฺวาทส, นปจฺฉาชาเต สตฺตรส, นอาเสวเน สตฺตรส, นกมฺเม สตฺต, นวิปาเก สตฺตรส, นอาหาเร, นอินฺทฺริเย, นฌาเน, นมคฺเค เอกํ, นสมฺปยุตฺเต ปญฺจ, นวิปฺปยุตฺเต ตีณิ, โนนตฺถิยา ปญฺจ, โนวิคเต ปญฺจ. (เอวํ คเณตพฺพํ.)}

\vspace{\tipitakaparskip}
\csromancenter{ปจฺจนียํ.}
\csromanmatikaanchor{mtk.61}
\csromantocmarkref{subparagraph}{3. ปจฺจยานุโลมปจฺจนีย}{mtk.61}
\bookmark[dest={mtk.61},level=5]{3. ปจฺจยานุโลมปจฺจนีย}
\csromanmarkh{12. ปริตฺตตฺติก}\csromanmarkhii{3. ปจฺจยา\-นุโลม\-ปจฺจนีย}\csromanheadernomark{6}{12. ปริตฺตตฺติก 3. ปจฺจยา\-นุโลม\-ปจฺจนีย}
\vspace{\tipitakaparskip}
\csromancenter{\textbf{เหตุ}ปจฺจยา นอารมฺมเณ ปญฺจ, นอธิปติยา ทฺวาทส, นอนนฺตเร ปญฺจ, นสมนนฺตเร ปญฺจ, นอญฺญมญฺเญ ปญฺจ, นฺเอาปนิสฺสเย ปญฺจ, นปุเรชาเต ทฺวาทส, นปจฺฉาชาเต สตฺตรส, นอาเสวเน สตฺตรส, นกมฺเม สตฺต, นวิปาเก สตฺตรส, นสมฺปยุตฺเต ปญฺจ, นวิปฺปยุตฺเต ตีณิ, โนนตฺถิยา ปญฺจ, โนวิคเต ปญฺจ. (เอวํ คเณตพฺพํ.)}
\nitthitamruled{อนุโลม\-ปจฺจนียํ.}
\csromanmatikaanchor{mtk.62}
\csromantocmarkref{subparagraph}{4. ปจฺจยปจฺจนียานุโลม}{mtk.62}
\bookmark[dest={mtk.62},level=5]{4. ปจฺจยปจฺจนียานุโลม}
\proseitem{41}{\csromanfolio{319}นเหตุปจฺจยา อารมฺมเณ เอกํ, อนนฺตเร, สมนนฺตเร, สหชาเต ฯเปฯ วิคเต, อวิคเต เอกํ. (เอวํ คเณตพฺพํ.) ปจฺจนียา\-นุโลมํ. ปจฺจยวาโร. (นิสฺสยวาโร ปจฺจยวาร\-สทิโส.)}
\csromansectionrule

\csromanmatikaanchor{mtk.63}
\csromantocmarknopage{subparagraph}{12. ปริตฺตตฺติก 5. สํสฏฺฐวาร}{mtk.63}
\bookmark[dest={mtk.63},level=5]{12. ปริตฺตตฺติก 5. สํสฏฺฐวาร}
\csromanmarkh{12. ปริตฺตตฺติก}\csromanmarkhii{5. สํสฏฺฐวาร}\csromanheadernomark{6}{12. ปริตฺตตฺติก 5. สํสฏฺฐวาร}
\csromanmarkh{1. ปจฺจยานุโลม}\csromanmarkhii{1. วิภงฺควาร}\csromanheadernomark{2}{1. \textbf{ปจฺจยานุโลม 1. วิภงฺควาร}}
\csromanmatikaanchor{mtk.64}
\csromanmatikaanchor{mtk.65}
\csromantocmarknopage{subparagraph}{1. ปจฺจยานุโลม 1. วิภงฺควาร}{mtk.64}
\bookmark[dest={mtk.64},level=5]{1. ปจฺจยานุโลม 1. วิภงฺควาร}
\csromantocmarkref{subparagraph}{เหตุ}{mtk.65}
\bookmark[dest={mtk.65},level=5]{เหตุ}
\csromanheader{6}{\textbf{เหตุ}}
\proseitem{42}{ปริตฺตํ ธมฺมํ สํสฏฺโฐ ปริตฺโต ธมฺโม อุปฺปชฺชติ เหตุปจฺจยา. ปริตฺตํ เอกํ ขนฺธํ สํสฏฺฐา ตโย ขนฺธา ฯเปฯ เทฺว ขนฺเธ ฯเปฯ. ปฏิสนฺธิ\-กฺขเณ ฯเปฯ. (1)}

\prose{มหคฺคตํ ธมฺมํ สํสฏฺโฐ มหคฺคโต ธมฺโม อุปฺปชฺชติ เหตุปจฺจยา. มหคฺคตํ เอกํ ขนฺธํ สํสฏฺฐา ตโย ขนฺธา ฯเปฯ เทฺว ขนฺเธ ฯเปฯ. ปฏิสนฺธิ\-กฺขเณ ฯเปฯ. (1)}

\prose{\csromanfolio{320}อปฺปมาณํ ธมฺมํ สํสฏฺโฐ อปฺปมาโณ ธมฺโม อุปฺปชฺชติ เหตุปจฺจยา. อปฺปมาณํ เอกํ ขนฺธํ สํสฏฺฐา ตโย ขนฺธา ฯเปฯ เทฺว ขนฺเธ ฯเปฯ. (1)}

\csromanmatikaanchor{mtk.66}
\csromantocmarkref{subparagraph}{อารมฺมณาทิ}{mtk.66}
\bookmark[dest={mtk.66},level=5]{อารมฺมณาทิ}
\csromanheader{6}{\textbf{อารมฺมณาทิ}}
\proseitem{43}{ปริตฺตํ ธมฺมํ สํสฏฺโฐ ปริตฺโต ธมฺโม อุปฺปชฺชติ อารมฺมณ\-ปจฺจยา. อธิปติ\-ปจฺจยา. (ปฏิสนฺธิ นตฺถิ.) อนนฺตร\-ปจฺจยา. สมนนฺตร\-ปจฺจยา. สหชาต\-ปจฺจยา. อญฺญมญฺญ\-ปจฺจยา. นิสฺสยปจฺจยา. อุปนิสฺสย\-ปจฺจยา. ปุเรชาต\-ปจฺจยา. (ปฏิสนฺธิ นตฺถิ.) อาเสวนปจฺจยา. (วิปาโกปิ ปฏิสนฺธิปิ นตฺถิ.) กมฺมปจฺจยา. วิปากปจฺจยา. อาหารปจฺจยา. อินฺทฺริยปจฺจยา. ฌานปจฺจยา. มคฺคปจฺจยา. สมฺปยุตฺต\-ปจฺจยา. วิปฺปยุตฺต\-ปจฺจยา. อตฺถิปจฺจยา. นตฺถิปจฺจยา. วิคตปจฺจยา. อวิคตปจฺจยา.}

\csromanmatikaanchor{mtk.67}
\csromantocmarkref{subparagraph}{1. ปจฺจยานุโลม 2. สงฺขฺยาวาร}{mtk.67}
\bookmark[dest={mtk.67},level=5]{1. ปจฺจยานุโลม 2. สงฺขฺยาวาร}
\csromanmarkh{1. ปจฺจยานุโลม}\csromanmarkhii{2. สงฺขฺยาวาร}\csromanheadernomark{6}{1. \textbf{ปจฺจยานุโลม 2. สงฺขฺยาวาร}}
\proseitemwithcloserruled{44}{เหตุยา ตีณิ, อารมฺมเณ ตีณิ, อธิปติยา ตีณิ ฯเปฯ อวิคเต ตีณิ. (เอวํ คเณตพฺพํ.)}{อนุโลมํ.}
\csromanmatikaanchor{mtk.68}
\csromantocmarkref{subparagraph}{2. ปจฺจยปจฺจนีย 1. วิภงฺควาร}{mtk.68}
\bookmark[dest={mtk.68},level=5]{2. ปจฺจยปจฺจนีย 1. วิภงฺควาร}
\csromanmarkh{2. ปจฺจย\-ปจฺจนีย}\csromanmarkhii{1. วิภงฺควาร}\csromanheadernomark{6}{2. ปจฺจย\-ปจฺจนีย 1. วิภงฺควาร}
\csromanheader{2}{\textbf{นเหตุ}}
\proseitem{45}{ปริตฺตํ ธมฺมํ สํสฏฺโฐ ปริตฺโต ธมฺโม อุปฺปชฺชติ นเหตุปจฺจยา. อเหตุกํ ปริตฺตํ เอกํ ขนฺธํ สํสฏฺฐา ตโย ขนฺธา ฯเปฯ เทฺว ขนฺเธ ฯเปฯ. อเหตุก\-ปฏิสนฺธิ\-กฺขเณ ฯเปฯ. วิจิกิจฺฉา\-สหคเต อุทฺธจฺจ\-สหคเต ขนฺเธ สํสฏฺโฐ วิจิกิจฺฉา\-สหคโต อุทฺธจฺจ\-สหคโต โมโห. (1)}

\prose{นอธิปติ}

\proseitem{46}{ปริตฺตํ ธมฺมํ สํสฏฺโฐ ปริตฺโต ธมฺโม อุปฺปชฺชติ นอธิปติ\-ปจฺจยา. ปริตฺตํ เอกํ ขนฺธํ สํสฏฺฐา ตโย ขนฺธา ฯเปฯ เทฺว ขนฺเธ ฯเปฯ. ปฏิสนฺธิ\-กฺขเณ ฯเปฯ. (1)}

\prose{\csromanfolio{321}มหคฺคตํ ธมฺมํ สํสฏฺโฐ มหคฺคโต ธมฺโม อุปฺปชฺชติ นอธิปติ\-ปจฺจยา. มหคฺคเต ขนฺเธ สํสฏฺฐา มหคฺคตา อธิปติ. วิปากํ มหคฺคตํ เอกํ ขนฺธํ สํสฏฺฐา ฯเปฯ. ปฏิสนฺธิ\-กฺขเณ ฯเปฯ. (1)}

\prose{อปฺปมาณํ ธมฺมํ สํสฏฺโฐ อปฺปมาโณ ธมฺโม อุปฺปชฺชติ นอธิปติ\-ปจฺจยา. อปฺปมาเณ ขนฺเธ สํสฏฺฐา อปฺปมาณา อธิปติ. (1)}

\csromanheader{2}{\textbf{นปุเรชาต}}
\proseitem{47}{ปริตฺตํ ธมฺมํ สํสฏฺโฐ ปริตฺโต ธมฺโม อุปฺปชฺชติ นปุเรชาต\-ปจฺจยา. อรูเป ปริตฺตํ เอกํ ขนฺธํ สํสฏฺฐา ตโย ขนฺธา ฯเปฯ. ปฏิสนฺธิ\-กฺขเณ ฯเปฯ. (1)}

\prose{มหคฺคตํ ธมฺมํ สํสฏฺโฐ มหคฺคโต ธมฺโม อุปฺปชฺชติ นปุเรชาต\-ปจฺจยา. อรูเป มหคฺคตํ เอกํ ขนฺธํ สํสฏฺฐา ตโย ขนฺธา ฯเปฯ. ปฏิสนฺธิ\-กฺขเณ ฯเปฯ. (1)}

\prose{อปฺปมาณํ ธมฺมํ สํสฏฺโฐ อปฺปมาโณ ธมฺโม อุปฺปชฺชติ นปุเรชาต\-ปจฺจยา. อรูเป อปฺปมาณํ เอกํ ขนฺธํ สํสฏฺฐา ตโย ขนฺธา ฯเปฯ. (1)}

\prose{นปจฺฉาชาต นอาเสวน}

\proseitem{48}{ปริตฺตํ ธมฺมํ สํสฏฺโฐ ปริตฺโต ธมฺโม อุปฺปชฺชติ นปจฺฉาชาต\-ปจฺจยา. นอาเสวน\-ปจฺจยา. ปริตฺตํ เอกํ ขนฺธํ สํสฏฺฐา ตโย ขนฺธา ฯเปฯ. ปฏิสนฺธิ\-กฺขเณ ฯเปฯ. (1)}

\prose{มหคฺคตํ ธมฺมํ สํสฏฺโฐ มหคฺคโต ธมฺโม อุปฺปชฺชติ นอาเสวน\-ปจฺจยา. วิปากํ มหคฺคตํ เอกํ ขนฺธํ สํสฏฺฐา ตโย ขนฺธา ฯเปฯ. ปฏิสนฺธิ\-กฺขเณ ฯเปฯ. (1)}

\prose{อปฺปมาณํ ธมฺมํ สํสฏฺโฐ อปฺปมาโณ ธมฺโม อุปฺปชฺชติ นอาเสวน\-ปจฺจยา. วิปากํ อปฺปมาณํ เอกํ ขนฺธํ สํสฏฺฐา ตโย ขนฺธา ฯเปฯ. (1)}

\csromanheader{2}{\textbf{นกมฺม}}
\proseitem{49}{ปริตฺตํ ธมฺมํ สํสฏฺโฐ ปริตฺโต ธมฺโม อุปฺปชฺชติ นกมฺมปจฺจยา. ปริตฺเต ขนฺเธ สํสฏฺฐา ปริตฺตา เจตนา. (1)}

\prose{มหคฺคตํ ธมฺมํ สํสฏฺโฐ มหคฺคโต ธมฺโม อุปฺปชฺชติ นกมฺมปจฺจยา. มหคฺคเต ขนฺเธ สํสฏฺฐา มหคฺคตา เจตนา. (1)}

\prose{อปฺปมาณํ ธมฺมํ สํสฏฺโฐ อปฺปมาโณ ธมฺโม อุปฺปชฺชติ นกมฺมปจฺจยา. กุสเล อปฺปมาเณ ขนฺเธ สํสฏฺฐา อปฺปมาณา เจตนา. (1)}

\csromanmatikaanchor{mtk.69}
\csromanmatikaanchor{mtk.70}
\csromantocmarkref{subparagraph}{2. ปจฺจยปจฺจนีย 2. สงฺขฺยาวาร}{mtk.69}
\bookmark[dest={mtk.69},level=5]{2. ปจฺจยปจฺจนีย 2. สงฺขฺยาวาร}
\csromantocmarkref{subparagraph}{3. ปจฺจยานุโลมปจฺจนีย}{mtk.70}
\bookmark[dest={mtk.70},level=5]{3. ปจฺจยานุโลมปจฺจนีย}
\proseitem{50}{\csromanfolio{322}ปริตฺตํ ธมฺมํ สํสฏฺโฐ ปริตฺโต ธมฺโม อุปฺปชฺชติ นวิปาก\-ปจฺจยา. ปริตฺตํ เอกํ ขนฺธํ สํสฏฺฐา ตโย ขนฺธา ฯเปฯ. (1)}

\prose{มหคฺคตํ ธมฺมํ สํสฏฺโฐ มหคฺคโต ธมฺโม อุปฺปชฺชติ นวิปาก\-ปจฺจยา. มหคฺคตํ เอกํ ขนฺธํ สํสฏฺฐา ตโย ขนฺธา ฯเปฯ. (1)}

\prose{อปฺปมาณํ ธมฺมํ สํสฏฺโฐ อปฺปมาโณ ธมฺโม อุปฺปชฺชติ นวิปาก\-ปจฺจยา. กุสลํ อปฺปมาณํ เอกํ ขนฺธํ สํสฏฺฐา ตโย ขนฺธา ฯเปฯ. (1)}

\csromanheader{2}{\textbf{นฌานาทิ}}
\proseitem{51}{ปริตฺตํ ธมฺมํ สํสฏฺโฐ ปริตฺโต ธมฺโม อุปฺปชฺชติ นฌานปจฺจยา. นมคฺคปจฺจยา. นวิปฺปยุตฺต\-ปจฺจยา. อรูเป ปริตฺตํ เอกํ ขนฺธํ สํสฏฺฐา ตโย ขนฺธา ฯเปฯ. (1)}

\prose{มหคฺคตํ ธมฺมํ สํสฏฺโฐ มหคฺคโต ธมฺโม อุปฺปชฺชติ นวิปฺปยุตฺต\-ปจฺจยา. อรูเป มหคฺคตํ เอกํ ขนฺธํ สํสฏฺฐา ตโย ขนฺธา ฯเปฯ. (1)}

\prose{อปฺปมาณํ ธมฺมํ สํสฏฺโฐ อปฺปมาโณ ธมฺโม อุปฺปชฺชติ นวิปฺปยุตฺต\-ปจฺจยา. อรูเป อปฺปมาณํ เอกํ ขนฺธํ สํสฏฺฐา ตโย ขนฺธา ฯเปฯ. (1)}

\csromanmarkh{2. ปจฺจย\-ปจฺจนีย}\csromanmarkhii{2. สงฺขฺยาวาร}\csromanheadernomark{6}{2. ปจฺจย\-ปจฺจนีย 2. สงฺขฺยาวาร}
\proseitemwithcloserruled{52}{นเหตุยา เอกํ, นอธิปติยา ตีณิ, นปุเรชาเต ตีณิ, นปจฺฉาชาเต ตีณิ, นอาเสวเน ตีณิ, นกมฺเม ตีณิ, นวิปาเก ตีณิ, นฌาเน เอกํ, นมคฺเค เอกํ, นวิปฺปยุตฺเต ตีณิ. (เอวํ คเณตพฺพํ.)}{ปจฺจนียํ.}
\proseitem{53}{เหตุปจฺจยา นอธิปติยา ตีณิ, นปุเรชาเต, นปจฺฉาชาเต, นอาเสวเน, นกมฺเม, นวิปาเก, นวิปฺปยุตฺเต ตีณิ. (เอวํ คเณตพฺพํ.)}

\vspace{\tipitakaparskip}
\csromancenter{อนุโลม\-ปจฺจนียํ.}
\csromanmatikaanchor{mtk.71}
\csromantocmarkref{subparagraph}{4. ปจฺจยปจฺจนียานุโลม}{mtk.71}
\bookmark[dest={mtk.71},level=5]{4. ปจฺจยปจฺจนียานุโลม}
\csromanmarkh{12. ปริตฺตตฺติก}\csromanmarkhii{4. ปจฺจย\-ปจฺจนียา\-นุโลม}\csromanheadernomark{6}{12. ปริตฺตตฺติก 4. ปจฺจย\-ปจฺจนียา\-นุโลม}
\proseitem{54}{\csromanfolio{323}นเหตุปจฺจยา อารมฺมเณ เอกํ, อนนฺตเร เอกํ ฯเปฯ อวิคเต เอกํ. (เอวํ คเณตพฺพํ.) ปจฺจนียา\-นุโลมํ. สํสฏฺฐวาโร. (สมฺปยุตฺตวาโร สํสฏฺฐวาร\-สทิโส.)}
\csromansectionrule

\csromanmatikaanchor{mtk.72}
\csromantocmarknopage{subparagraph}{12. ปริตฺตตฺติก 7. ปญฺหาวาร}{mtk.72}
\bookmark[dest={mtk.72},level=5]{12. ปริตฺตตฺติก 7. ปญฺหาวาร}
\csromanmarkh{12. ปริตฺตตฺติก}\csromanmarkhii{7. ปญฺหาวาร}\csromanheadernomark{6}{12. ปริตฺตตฺติก 7. ปญฺหาวาร}
\csromanmarkh{1. ปจฺจยานุโลม}\csromanmarkhii{1. วิภงฺควาร}\csromanheadernomark{2}{1. \textbf{ปจฺจยานุโลม 1. วิภงฺควาร}}
\csromanmatikaanchor{mtk.73}
\csromanmatikaanchor{mtk.74}
\csromantocmarknopage{subparagraph}{1. ปจฺจยานุโลม 1. วิภงฺควาร}{mtk.73}
\bookmark[dest={mtk.73},level=5]{1. ปจฺจยานุโลม 1. วิภงฺควาร}
\csromantocmarkref{subparagraph}{เหตุ}{mtk.74}
\bookmark[dest={mtk.74},level=5]{เหตุ}
\csromanheader{6}{\textbf{เหตุ}}
\proseitem{55}{ปริตฺโต ธมฺโม ปริตฺตสฺส ธมฺมสฺส เหตุปจฺจเยน ปจฺจโย. ปริตฺตา เหตู สมฺปยุตฺตกานํ ขนฺธานํ จิตฺตสมุฏฺฐานานญฺจ รูปานํ เหตุปจฺจเยน ปจฺจโย. ปฏิสนฺธิ\-กฺขเณ ฯเปฯ. (1)}

\prose{มหคฺคโต ธมฺโม มหคฺคตสฺส ธมฺมสฺส เหตุปจฺจเยน ปจฺจโย. ตีณิ. (ปวตฺติ\-ปฏิสนฺธิ กาตพฺพา.)}

\prose{อปฺปมาโณ ธมฺโม อปฺปมาณสฺส ธมฺมสฺส เหตุปจฺจเยน ปจฺจโย. ตีณิ.}

\csromanmatikaanchor{mtk.75}
\csromantocmarkref{subparagraph}{อารมฺมณ}{mtk.75}
\bookmark[dest={mtk.75},level=5]{อารมฺมณ}
\csromanheader{6}{\textbf{อารมฺมณ}}
\proseitem{56}{ปริตฺโต ธมฺโม ปริตฺตสฺส ธมฺมสฺส อารมฺมณ\-ปจฺจเยน ปจฺจโย. ทานํ ทตฺวา, สีลํ สมาทิยิตฺวา, อุโปสถกมฺมํ กตฺวา ตํ ปจฺจเวกฺขติ, ปุพฺเพ สุจิณฺณานิ ปจฺจเวกฺขติ. อริยา โคตฺรภุํ ปจฺจเวกฺขนฺติ, โวทานํ ปจฺจเวกฺขนฺติ, ปหีเน กิเลเส ปจฺจเวกฺขนฺติ, วิกฺขมฺภิเต กิเลเส ปจฺจเวกฺขนฺติ, ปุพฺเพ สมุทาจิณฺเณ กิเลเส ชานนฺติ. จกฺขุํ ฯเปฯ. วตฺถุํ. ปริตฺเต ขนฺเธ อนิจฺจโต ฯเปฯ วิปสฺสนฺติ, อสฺสาเทนฺติ อภินนฺทนฺติ, ตํ อารพฺภ ราโค อุปฺปชฺชติ ฯเปฯ โทมนสฺสํ อุปฺปชฺชติ. รูปายตนํ จกฺขุ\-วิญฺญาณสฺส ฯเปฯ. โผฏฺฐพฺพา\-ยตนํ กายวิญฺญาณสฺส อารมฺมณ\-ปจฺจเยน ปจฺจโย. (1)}

\prose{\csromanfolio{324}ปริตฺโต ธมฺโม มหคฺคตสฺส ธมฺมสฺส อารมฺมณ\-ปจฺจเยน ปจฺจโย. ทิพฺเพน จกฺขุนา รูปํ ปสฺสติ. ทิพฺพาย โสตธาตุยา สทฺทํ สุณาติ. เจโตปริย\-ญาเณน ปริตฺต\-จิตฺต\-สมงฺคิสฺส จิตฺตํ ชานาติ. ปริตฺตา ขนฺธา อิทฺธิวิธ\-ญาณสฺส, เจโตปริย\-ญาณสฺส, ปุพฺเพ\-นิวาสา\-นุสฺสติ\-ญาณสฺส, ยถากมฺมูปคญาณสฺส, อนาคตํสญาณสฺส อารมฺมณ\-ปจฺจเยน ปจฺจโย. (2)}

\proseitem{57}{มหคฺคโต ธมฺโม มหคฺคตสฺส ธมฺมสฺส อารมฺมณ\-ปจฺจเยน ปจฺจโย. เจโตปริย\-ญาเณน มหคฺคต\-จิตฺต\-สมงฺคิสฺส จิตฺตํ ชานาติ. อากาสา\-นญฺจา\-ยตนํ วิญฺญาณญฺจา\-ยตนสฺส. อากิญฺจญฺญา\-ยตนํ เนว\-สญฺญา\-นา\-สญฺญา\-ยตนสฺส อารมฺมณ\-ปจฺจเยน ปจฺจโย. มหคฺคตา ขนฺธา อิทฺธิวิธ\-ญาณสฺส, เจโตปริย\-ญาณสฺส, ปุพฺเพ\-นิวาสา\-นุสฺสติ\-ญาณสฺส, ยถากมฺมูปคญาณสฺส, อนาคตํสญาณสฺส อารมฺมณ\-ปจฺจเยน ปจฺจโย. (1)}

\prose{มหคฺคโต ธมฺโม ปริตฺตสฺส ธมฺมสฺส อารมฺมณ\-ปจฺจเยน ปจฺจโย. ปฐมํ ฌานํ ปจฺจเวกฺขติ ฯเปฯ. เนว\-สญฺญา\-นา\-สญฺญา\-ยตนํ ปจฺจเวกฺขติ, ทิพฺพํ จกฺขุํ ปจฺจเวกฺขติ, ทิพฺพํ โสตธาตุํ ปจฺจเวกฺขติ, อิทฺธิวิธ\-ญาณํ ปจฺจเวกฺขติ, เจโตปริย\-ญาณํ ปจฺจเวกฺขติ, ปุพฺเพ\-นิวาสา\-นุสฺสติ\-ญาณํ ปจฺจเวกฺขติ, ยถากมฺมูปคญาณํ ปจฺจเวกฺขติ, อนาคตํสญาณํ ปจฺจเวกฺขติ. มหคฺคเต ขนฺเธ อนิจฺจโต ฯเปฯ วิปสฺสติ, อสฺสาเทติ อภินนฺทติ, ตํ อารพฺภ ราโค อุปฺปชฺชติ ฯเปฯ โทมนสฺสํ อุปฺปชฺชติ. (2)}

\proseitem{58}{อปฺปมาโณ ธมฺโม อปฺปมาณสฺส ธมฺมสฺส อารมฺมณ\-ปจฺจเยน ปจฺจโย. นิพฺพานํ มคฺคสฺส, ผลสฺส อารมฺมณ\-ปจฺจเยน ปจฺจโย. (1)}

\prose{อปฺปมาโณ ธมฺโม ปริตฺตสฺส ธมฺมสฺส อารมฺมณ\-ปจฺจเยน ปจฺจโย. อริยา มคฺคา วุฏฺฐหิตฺวา มคฺคํ ปจฺจเวกฺขนฺติ, ผลํ ปจฺจเวกฺขนฺติ, นิพฺพานํ ปจฺจเวกฺขนฺติ, นิพฺพานํ โคตฺรภุสฺส, โวทานสฺส, อาวชฺชนาย อารมฺมณ\-ปจฺจเยน ปจฺจโย. (2)}

\prose{อปฺปมาโณ ธมฺโม มหคฺคตสฺส ธมฺมสฺส อารมฺมณ\-ปจฺจเยน ปจฺจโย. อริยา เจโตปริย\-ญาเณน อปฺปมาณจิตฺต\-สมงฺคิสฺส จิตฺตํ ชานนฺติ. อปฺปมาณา ขนฺธา เจโตปริย\-ญาณสฺส, ปุพฺเพ\-นิวาสา\-นุสฺสติ\-ญาณสฺส, อนาคตํสญาณสฺส อารมฺมณ\-ปจฺจเยน ปจฺจโย. (3)}

\csromanmatikaanchor{mtk.76}
\csromantocmarkref{subparagraph}{อธิปติ}{mtk.76}
\bookmark[dest={mtk.76},level=5]{อธิปติ}
\csromanheader{6}{12. ปริตฺตตฺติก \textbf{อธิปติ}}
\proseitem{59}{\csromanfolio{325}ปริตฺโต ธมฺโม ปริตฺตสฺส ธมฺมสฺส อธิปติ\-ปจฺจเยน ปจฺจโย. อารมฺมณา\-ธิปติ, สหชาตา\-ธิปติ.  อารมฺมณา\-ธิปติ\csromandash{}ทานํ ทตฺวา, สีลํ สมาทิยิตฺวา, อุโปสถกมฺมํ กตฺวา ตํ ครุํ กตฺวา ปจฺจเวกฺขติ, ปุพฺเพ สุจิณฺณานิ ครุํ กตฺวา ปจฺจเวกฺขติ. เสกฺขา โคตฺรภุํ ครุํ กตฺวา ปจฺจเวกฺขนฺติ, โวทานํ ครุํ กตฺวา ปจฺจเวกฺขนฺติ. จกฺขุํ ฯเปฯ. วตฺถุํ ปริตฺเต ขนฺเธ ครุํ กตฺวา อสฺสาเทติ อภินนฺทติ, ตํ ครุํ กตฺวา ราโค อุปฺปชฺชติ, ทิฏฺฐิ อุปฺปชฺชติ. สหชาตา\-ธิปติ\csromandash{}ปริตฺตาธิปติ สมฺปยุตฺตกานํ ขนฺธานํ จิตฺตสมุฏฺฐานานญฺจ รูปานํ อธิปติ\-ปจฺจเยน ปจฺจโย. (1)}

\prose{มหคฺคโต ธมฺโม มหคฺคตสฺส ธมฺมสฺส อธิปติ\-ปจฺจเยน ปจฺจโย. สหชาตา\-ธิปติ\csromandash{}มหคฺคตา\-ธิปติ สมฺปยุตฺตกานํ ขนฺธานํ อธิปติ\-ปจฺจเยน ปจฺจโย. (1)}

\prose{มหคฺคโต ธมฺโม ปริตฺตสฺส ธมฺมสฺส อธิปติ\-ปจฺจเยน ปจฺจโย. อารมฺมณา\-ธิปติ, สหชาตา\-ธิปติ.  อารมฺมณา\-ธิปติ\csromandash{}ปฐมํ ฌานํ ครุํ กตฺวา ฯเปฯ. เนว\-สญฺญา\-นา\-สญฺญา\-ยตนํ ฯเปฯ. ทิพฺพํ จกฺขุํ ฯเปฯ. อนาคตํสญาณํ ครุํ กตฺวา ปจฺจเวกฺขติ. มหคฺคเต ขนฺเธ ครุํ กตฺวา อสฺสาเทติ อภินนฺทติ, ตํ ครุํ กตฺวา ราโค อุปฺปชฺชติ. ทิฏฺฐิ อุปฺปชฺชติ. สหชาตา\-ธิปติ\csromandash{}มหคฺคตา\-ธิปติ จิตฺต\-สมุฏฺฐานานํ รูปานํ อธิปติ\-ปจฺจเยน ปจฺจโย. (2)}

\prose{มหคฺคโต ธมฺโม ปริตฺตสฺส จ มหคฺคตสฺส จ ธมฺมสฺส อธิปติ\-ปจฺจเยน ปจฺจโย. สหชาตา\-ธิปติ\csromandash{}มหคฺคตา\-ธิปติ สมฺปยุตฺตกานํ ขนฺธานํ จิตฺตสมุฏฺฐานานญฺจ รูปานํ อธิปติ\-ปจฺจเยน ปจฺจโย. (3)}

\proseitem{60}{อปฺปมาโณ ธมฺโม อปฺปมาณสฺส ธมฺมสฺส อธิปติ\-ปจฺจเยน ปจฺจโย. อารมฺมณา\-ธิปติ, สหชาตา\-ธิปติ.  อารมฺมณา\-ธิปติ\csromandash{} นิพฺพานํ มคฺคสฺส, ผลสฺส อธิปติ\-ปจฺจเยน ปจฺจโย. สหชาตา\-ธิปติ\csromandash{} อปฺปมาณา\-ธิปติ สมฺปยุตฺตกานํ ขนฺธานํ อธิปติ\-ปจฺจเยน ปจฺจโย. (1)}

\csromanmatikaanchor{mtk.77}
\csromantocmarkref{subparagraph}{อนนฺตราทิ}{mtk.77}
\bookmark[dest={mtk.77},level=5]{อนนฺตราทิ}
\prose{อปฺปมาโณ ธมฺโม ปริตฺตสฺส ธมฺมสฺส อธิปติ\-ปจฺจเยน ปจฺจโย. อารมฺมณา\-ธิปติ, สหชาตา\-ธิปติ.  อารมฺมณา\-ธิปติ\csromandash{}อริยา มคฺคา วุฏฺฐหิตฺวา มคฺคํ ครุํ กตฺวา ปจฺจเวกฺขนฺติ, ผลํ ครุํ กตฺวา ปจฺจเวกฺขนฺติ. นิพฺพานํ ครุํ กตฺวา ปจฺจเวกฺขนฺติ. นิพฺพานํ โคตฺรภุสฺส, โวทานสฺส อธิปติ\-ปจฺจเยน \csromanfolio{326}ปจฺจโย. สหชาตา\-ธิปติ\csromandash{}อปฺปมาณา\-ธิปติ จิตฺต\-สมุฏฺฐานานํ รูปานํ อธิปติ\-ปจฺจเยน ปจฺจโย. (2)}

\prose{อปฺปมาโณ ธมฺโม ปริตฺตสฺส จ อปฺปมาณสฺส จ ธมฺมสฺส อธิปติ\-ปจฺจเยน ปจฺจโย. สหชาตา\-ธิปติ\csromandash{}อปฺปมาณา\-ธิปติ สมฺปยุตฺตกานํ ขนฺธานํ จิตฺตสมุฏฺฐานานญฺจ รูปานํ อธิปติ\-ปจฺจเยน ปจฺจโย. (3)}

\csromanheader{2}{\textbf{อนนฺตร}}
\proseitem{61}{ปริตฺโต ธมฺโม ปริตฺตสฺส ธมฺมสฺส อนนฺตร\-ปจฺจเยน ปจฺจโย. ปุริมา ปุริมา ปริตฺตา ขนฺธา ปจฺฉิมานํ ปจฺฉิมานํ ปริตฺตานํ ขนฺธานํ อนนฺตร\-ปจฺจเยน ปจฺจโย. อนุโลมํ โคตฺรภุสฺส. อนุโลมํ โวทานสฺส. อาวชฺชนา ปริตฺตานํ ขนฺธานํ อนนฺตร\-ปจฺจเยน ปจฺจโย. (1)}

\prose{ปริตฺโต ธมฺโม มหคฺคตสฺส ธมฺมสฺส อนนฺตร\-ปจฺจเยน ปจฺจโย. ปริตฺตํ จุติจิตฺตํ มหคฺคตสฺส อุปปตฺติ\-จิตฺตสฺส อนนฺตร\-ปจฺจเยน ปจฺจโย. ปริตฺตา ขนฺธา มหคฺคตสฺส วุฏฺฐานสฺส อนนฺตร\-ปจฺจเยน ปจฺจโย. ปฐมสฺส ฌานสฺส ปริกมฺมํ ฯเปฯ. เนว\-สญฺญา\-นา\-สญฺญา\-ยตนสฺส ปริกมฺมํ ฯเปฯ. ทิพฺพสฺส จกฺขุสฺส ปริกมฺมํ ฯเปฯ. อนาคตํสญาณสฺส ปริกมฺมํ อนาคตํสญาณสฺส อนนฺตร\-ปจฺจเยน ปจฺจโย. (2)}

\prose{ปริตฺโต ธมฺโม อปฺปมาณสฺส ธมฺมสฺส อนนฺตร\-ปจฺจเยน ปจฺจโย. โคตฺรภุ มคฺคสฺส. โวทานํ มคฺคสฺส. อนุโลมํ ผลสมาปตฺติยา อนนฺตร\-ปจฺจเยน ปจฺจโย. (3)}

\proseitem{62}{มหคฺคโต ธมฺโม มหคฺคตสฺส ธมฺมสฺส อนนฺตร\-ปจฺจเยน ปจฺจโย. ปุริมา ปุริมา มหคฺคตา ขนฺธา ปจฺฉิมานํ ปจฺฉิมานํ มหคฺคตานํ ขนฺธานํ อนนฺตร\-ปจฺจเยน ปจฺจโย. (1)}

\prose{มหคฺคโต ธมฺโม ปริตฺตสฺส ธมฺมสฺส อนนฺตร\-ปจฺจเยน ปจฺจโย. มหคฺคตํ จุติจิตฺตํ ปริตฺตสฺส อุปปตฺติ\-จิตฺตสฺส อนนฺตร\-ปจฺจเยน ปจฺจโย. มหคฺคตํ ภวงฺคํ อาวชฺชนาย อนนฺตร\-ปจฺจเยน ปจฺจโย. มหคฺคตา ขนฺธา ปริตฺตสฺส วุฏฺฐานสฺส อนนฺตร\-ปจฺจเยน ปจฺจโย. (2)}

\prose{\csromanfolio{327}มหคฺคโต ธมฺโม อปฺปมาณสฺส ธมฺมสฺส อนนฺตร\-ปจฺจเยน ปจฺจโย. เนว\-สญฺญา\-นา\-สญฺญา\-ยตนํ นิโรธา วุฏฺฐหนฺตสฺส ผลสมาปตฺติยา อนนฺตร\-ปจฺจเยน ปจฺจโย. (3)}

\proseitem{63}{อปฺปมาโณ ธมฺโม อปฺปมาณสฺส ธมฺมสฺส อนนฺตร\-ปจฺจเยน ปจฺจโย. ปุริมา ปุริมา อปฺปมาณา ขนฺธา ปจฺฉิมานํ ปจฺฉิมานํ อปฺปมาณานํ ขนฺธานํ อนนฺตร\-ปจฺจเยน ปจฺจโย. มคฺโค ผลสฺส. ผลํ ผลสฺส อนนฺตร\-ปจฺจเยน ปจฺจโย. (1)}

\prose{อปฺปมาโณ ธมฺโม ปริตฺตสฺส ธมฺมสฺส อนนฺตร\-ปจฺจเยน ปจฺจโย. ผลํ ปริตฺตสฺส วุฏฺฐานสฺส อนนฺตร\-ปจฺจเยน ปจฺจโย. (2)}

\prose{อปฺปมาโณ ธมฺโม มหคฺคตสฺส ธมฺมสฺส อนนฺตร\-ปจฺจเยน ปจฺจโย. ผลํ มหคฺคตสฺส วุฏฺฐานสฺส อนนฺตร\-ปจฺจเยน ปจฺจโย. (3) (สมนนฺตร\-ปจฺจยํ อนนฺตรปจฺจย\-สทิสํ.)}

\csromanheader{2}{\textbf{สหชาต}}
\proseitem{64}{ปริตฺโต ธมฺโม ปริตฺตสฺส ธมฺมสฺส สหชาต\-ปจฺจเยน ปจฺจโย. ปริตฺโต เอโก ขนฺโธ ติณฺณนฺนํ ขนฺธานํ จิตฺตสมุฏฺฐานานญฺจ รูปานํ สหชาต\-ปจฺจเยน ปจฺจโย ฯเปฯ เทฺว ขนฺธา ฯเปฯ. ปฏิสนฺธิ\-กฺขเณ ฯเปฯ. ขนฺธา วตฺถุสฺส, วตฺถุ ขนฺธานํ สหชาต\-ปจฺจเยน ปจฺจโย. เอกํ มหาภูตํ ฯเปฯ. อสญฺญสตฺตานํ ฯเปฯ. (1)}

\prose{ปริตฺโต ธมฺโม มหคฺคตสฺส ธมฺมสฺส สหชาต\-ปจฺจเยน ปจฺจโย. ปฏิสนฺธิ\-กฺขเณ วตฺถุ มหคฺคตานํ ขนฺธานํ สหชาต\-ปจฺจเยน ปจฺจโย. (2)}

\proseitem{65}{มหคฺคโต ธมฺโม มหคฺคตสฺส ธมฺมสฺส สหชาต\-ปจฺจเยน ปจฺจโย. มหคฺคโต เอโก ขนฺโธ ติณฺณนฺนํ ขนฺธานํ ฯเปฯ เทฺว ขนฺธา ฯเปฯ. ปฏิสนฺธิกฺขเน ฯเปฯ. (1)}

\prose{มหคฺคโต ธมฺโม ปริตฺตสฺส ธมฺมสฺส สหชาต\-ปจฺจเยน ปจฺจโย. มหคฺคตา ขนฺธา จิตฺต\-สมุฏฺฐานานํ รูปานํ สหชาต\-ปจฺจเยน ปจฺจโย. ปฏิสนฺธิ\-กฺขเณ มหคฺคตา ขนฺธา กฏตฺตารูปานํ สหชาต\-ปจฺจเยน ปจฺจโย. (2)}

\prose{\csromanfolio{328}มหคฺคโต ธมฺโม ปริตฺตสฺส จ มหคฺคตสฺส จ ธมฺมสฺส สหชาต\-ปจฺจเยน ปจฺจโย. มหคฺคโต เอโก ขนฺโธ ติณฺณนฺนํ ขนฺธานํ จิตฺตสมุฏฺฐานานญฺจ รูปานํ สหชาต\-ปจฺจเยน ปจฺจโย ฯเปฯ เทฺว ขนฺธา ฯเปฯ. ปฏิสนฺธิ\-กฺขเณ ฯเปฯ. (3)}

\proseitem{66}{อปฺปมาโณ ธมฺโม อปฺปมาณสฺส ธมฺมสฺส สหชาต\-ปจฺจเยน ปจฺจโย. อปฺปมาโณ เอโก ขนฺโธ ติณฺณนฺนํ ขนฺธานํ สหชาต\-ปจฺจเยน ปจฺจโย ฯเปฯ. (1)}

\prose{อปฺปมาโณ ธมฺโม ปริตฺตสฺส ธมฺมสฺส สหชาต\-ปจฺจเยน ปจฺจโย. อปฺปมาณา ขนฺธา จิตฺต\-สมุฏฺฐานานํ รูปานํ สหชาต\-ปจฺจเยน ปจฺจโย. (2)}

\prose{อปฺปมาโณ ธมฺโม ปริตฺตสฺส จ อปฺปมาณสฺส จ ธมฺมสฺส สหชาต\-ปจฺจเยน ปจฺจโย. อปฺปมาโณ เอโก ขนฺโธ ติณฺณนฺนํ ขนฺธานํ จิตฺตสมุฏฺฐานานญฺจ รูปานํ สหชาต\-ปจฺจเยน ปจฺจโย ฯเปฯ. (3)}

\proseitem{67}{ปริตฺโต จ อปฺปมาโณ จ ธมฺมา ปริตฺตสฺส ธมฺมสฺส สหชาต\-ปจฺจเยน ปจฺจโย. อปฺปมาณา ขนฺธา จ มหาภูตา จ จิตฺต\-สมุฏฺฐานานํ รูปานํ สหชาต\-ปจฺจเยน ปจฺจโย. (1)}

\prose{ปริตฺโต จ มหคฺคโต จ ธมฺมา ปริตฺตสฺส ธมฺมสฺส สหชาต\-ปจฺจเยน ปจฺจโย. มหคฺคตา ขนฺธา จ มหาภูตา จ จิตฺต\-สมุฏฺฐานานํ รูปานํ สหชาต\-ปจฺจเยน ปจฺจโย. ปฏิสนฺธิ\-กฺขเณ มหคฺคตา ขนฺธา จ มหาภูตา จ กฏตฺตารูปานํ สหชาต\-ปจฺจเยน ปจฺจโย. (1)}

\prose{ปริตฺโต จ มหคฺคโต จ ธมฺมา มหคฺคตสฺส ธมฺมสฺส สหชาต\-ปจฺจเยน ปจฺจโย. ปฏิสนฺธิ\-กฺขเณ มหคฺคโต เอโก ขนฺโธ จ วตฺถุ จ ติณฺณนฺนํ ขนฺธานํ สหชาต\-ปจฺจเยน ปจฺจโย ฯเปฯ. (2)}

\csromanheader{2}{\textbf{อญฺญมญฺญ}}
\proseitem{68}{ปริตฺโต ธมฺโม ปริตฺตสฺส ธมฺมสฺส อญฺญมญฺญ\-ปจฺจเยน ปจฺจโย. ปริตฺโต เอโก ขนฺโธ ติณฺณนฺนํ ขนฺธานํ อญฺญมญฺญ\-ปจฺจเยน ปจฺจโย ฯเปฯ. ปฏิสนฺธิ\-กฺขเณ ฯเปฯ. ขนฺธา วตฺถุสฺส ฯเปฯ. วตฺถุ ขนฺธานํ อญฺญมญฺญ\-ปจฺจเยน ปจฺจโย. เอกํ มหาภูตํ ฯเปฯ. อสญฺญสตฺตานํ ฯเปฯ. (1)}

\prose{\csromanfolio{329}ปริตฺโต ธมฺโม มหคฺคตสฺส ธมฺมสฺส อญฺญมญฺญ\-ปจฺจเยน ปจฺจโย ฯเปฯ. ปฏิสนฺธิ\-กฺขเณ วตฺถุ มหคฺคตานํ ขนฺธานํ อญฺญมญฺญ\-ปจฺจเยน ปจฺจโย ฯเปฯ. (2)}

\prose{มหคฺคโต ธมฺโม มหคฺคตสฺส ธมฺมสฺส อญฺญมญฺญ\-ปจฺจเยน ปจฺจโย. มหคฺคโต เอโก ขนฺโธ ติณฺณนฺนํ ขนฺธานํ อญฺญมญฺญ\-ปจฺจเยน ปจฺจโย ฯเปฯ. ปฏิสนฺธิ\-กฺขเณ ฯเปฯ. (1)}

\prose{มหคฺคโต ธมฺโม ปริตฺตสฺส ธมฺมสฺส อญฺญมญฺญ\-ปจฺจเยน ปจฺจโย. ปฏิสนฺธิ\-กฺขเณ มหคฺคตา ขนฺธา วตฺถุสฺส อญฺญมญฺญ\-ปจฺจเยน ปจฺจโย. (2)}

\prose{มหคฺคโต ธมฺโม ปริตฺตสฺส จ มหคฺคตสฺส จ ธมฺมสฺส อญฺญมญฺญ\-ปจฺจเยน ปจฺจโย. ปฏิสนฺธิ\-กฺขเณ มหคฺคโต เอโก ขนฺโธ ติณฺณนฺนํ ขนฺธานํ วตฺถุสฺส จ อญฺญมญฺญ\-ปจฺจเยน ปจฺจโย ฯเปฯ. (3)}

\prose{อปฺปมาโณ ธมฺโม อปฺปมาณสฺส ธมฺมสฺส อญฺญมญฺญ\-ปจฺจเยน ปจฺจโย. อปฺปมาโณ เอโก ขนฺโธ ติณฺณนฺนํ ขนฺธานํ อญฺญมญฺญ\-ปจฺจเยน ปจฺจโย ฯเปฯ เทฺว ขนฺธา ฯเปฯ. (1)}

\prose{ปริตฺโต จ มหคฺคโต จ ธมฺมา มหคฺคตสฺส ธมฺมสฺส อญฺญมญฺญ\-ปจฺจเยน ปจฺจโย. ปฏิสนฺธิ\-กฺขเณ มหคฺคโต เอโก ขนฺโธ จ วตฺถุ จ ติณฺณนฺนํ ขนฺธานํ อญฺญมญฺญ\-ปจฺจเยน ปจฺจโย ฯเปฯ. (1)}

\csromanheader{2}{\textbf{นิสฺสย}}
\proseitem{69}{ปริตฺโต ธมฺโม ปริตฺตสฺส ธมฺมสฺส นิสฺสย\-ปจฺจเยน ปจฺจโย. ปริตฺโต เอโก ขนฺโธ ติณฺณนฺนํ ขนฺธานํ จิตฺตสมุฏฺฐานานญฺจ รูปานํ นิสฺสย\-ปจฺจเยน ปจฺจโย ฯเปฯ เทฺว ขนฺธา ฯเปฯ. ปฏิสนฺธิ\-กฺขเณ ฯเปฯ. ขนฺธา วตฺถุสฺส ฯเปฯ. วตฺถุ ขนฺธานํ นิสฺสย\-ปจฺจเยน ปจฺจโย ฯเปฯ. เอกํ มหาภูตํ ฯเปฯ. อสญฺญสตฺตานํ เอกํ มหาภูตํ ฯเปฯ. จกฺขายตนํ จกฺขุ\-วิญฺญาณสฺส ฯเปฯ. กายายตนํ กายวิญฺญาณสฺส. วตฺถุ ปริตฺตานํ ขนฺธานํ นิสฺสย\-ปจฺจเยน ปจฺจโย. (1)}

\prose{ปริตฺโต ธมฺโม มหคฺคตสฺส ธมฺมสฺส นิสฺสย\-ปจฺจเยน ปจฺจโย. วตฺถุ มหคฺคตานํ ขนฺธานํ นิสฺสย\-ปจฺจเยน ปจฺจโย. ปฏิสนฺธิ\-กฺขเณ วตฺถุ มหคฺคตานํ ขนฺธานํ นิสฺสย\-ปจฺจเยน ปจฺจโย. (2)}

\prose{ปริตฺโต ธมฺโม อปฺปมาณสฺส ธมฺมสฺส นิสฺสย\-ปจฺจเยน ปจฺจโย. วตฺถุ อปฺปมาณานํ ขนฺธานํ นิสฺสย\-ปจฺจเยน ปจฺจโย. (3)}

\proseitem{70}{\csromanfolio{330}มหคฺคโต ธมฺโม มหคฺคตสฺส ธมฺมสฺส นิสฺสย\-ปจฺจเยน ปจฺจโย. มหคฺคโต เอโก ขนฺโธ ติณฺณนฺนํ ขนฺธานํ ฯเปฯ เทฺว ขนฺธา ฯเปฯ. ปฏิสนฺธิ\-กฺขเณ ฯเปฯ. (1)}

\prose{มหคฺคโต ธมฺโม ปริตฺตสฺส ธมฺมสฺส นิสฺสย\-ปจฺจเยน ปจฺจโย. มหคฺคตา ขนฺธา จิตฺต\-สมุฏฺฐานานํ รูปานํ นิสฺสย\-ปจฺจเยน ปจฺจโย. ปฏิสนฺธิ\-กฺขเณ มหคฺคตา ขนฺธา กฏตฺตารูปานํ นิสฺสย\-ปจฺจเยน ปจฺจโย. (2)}

\prose{มหคฺคโต ธมฺโม ปริตฺตสฺส จ มหคฺคตสฺส จ ธมฺมสฺส นิสฺสย\-ปจฺจเยน ปจฺจโย. มหคฺคโต เอโก ขนฺโธ ติณฺณนฺนํ ขนฺธานํ จิตฺตสมุฏฺฐานานญฺจ รูปานํ นิสฺสย\-ปจฺจเยน ปจฺจโย ฯเปฯ เทฺว ขนฺธา ฯเปฯ. ปฏิสนฺธิ\-กฺขเณ ฯเปฯ. (3)}

\proseitem{71}{อปฺปมาโณ ธมฺโม อปฺปมาณสฺส ธมฺมสฺส นิสฺสย\-ปจฺจเยน ปจฺจโย. อปฺปมาโณ เอโก ขนฺโธ ติณฺณนฺนํ ขนฺธานํ นิสฺสย\-ปจฺจเยน ปจฺจโย ฯเปฯ. (1)}

\prose{อปฺปมาโณ ธมฺโม ปริตฺตสฺส ธมฺมสฺส นิสฺสย\-ปจฺจเยน ปจฺจโย. อปฺปมาณา ขนฺธา จิตฺต\-สมุฏฺฐานานํ รูปานํ นิสฺสย\-ปจฺจเยน ปจฺจโย. (2)}

\prose{อปฺปมาโณ ธมฺโม ปริตฺตสฺส จ อปฺปมาณสฺส จ ธมฺมสฺส นิสฺสย\-ปจฺจเยน ปจฺจโย. อปฺปมาโณ เอโก ขนฺโธ ติณฺณนฺนํ ขนฺธานํ จิตฺตสมุฏฺฐานานญฺจ รูปานํ นิสฺสย\-ปจฺจเยน ปจฺจโย ฯเปฯ. (3)}

\prose{ปริตฺโต จ อปฺปมาโณ จ ธมฺมา ปริตฺตสฺส ธมฺมสฺส นิสฺสย\-ปจฺจเยน ปจฺจโย. อปฺปมาณา ขนฺธา จ มหาภูตา จ จิตฺต\-สมุฏฺฐานานํ รูปานํ นิสฺสย\-ปจฺจเยน ปจฺจโย. (1)}

\prose{ปริตฺโต จ อปฺปมาโณ จ ธมฺมา อปฺปมาณสฺส ธมฺมสฺส นิสฺสย\-ปจฺจเยน ปจฺจโย. อปฺปมาโณ เอโก ขนฺโธ จ วตฺถุ จ ติณฺณนฺนํ ขนฺธานํ นิสฺสย\-ปจฺจเยน ปจฺจโย ฯเปฯ เทฺว ขนฺธา ฯเปฯ. (2)}

\prose{ปริตฺโต จ มหคฺคโต จ ธมฺมา ปริตฺตสฺส ธมฺมสฺส นิสฺสย\-ปจฺจเยน ปจฺจโย. มหคฺคตา ขนฺธา จ มหาภูตา จ จิตฺต\-สมุฏฺฐานานํ รูปานํ นิสฺสย\-ปจฺจเยน ปจฺจโย. ปฏิสนฺธิ\-กฺขเณ มหคฺคตา ขนฺธา จ มหาภูตา จ กฏตฺตารูปานํ นิสฺสย\-ปจฺจเยน ปจฺจโย. (1)}

\prose{\csromanfolio{331}ปริตฺโต จ มหคฺคโต จ ธมฺมา มหคฺคตสฺส ธมฺมสฺส นิสฺสย\-ปจฺจเยน ปจฺจโย. มหคฺคโต เอโก ขนฺโธ จ วตฺถุ จ ติณฺณนฺนํ ขนฺธานํ นิสฺสย\-ปจฺจเยน ปจฺจโย ฯเปฯ. ปฏิสนฺธิ\-กฺขเณ มหคฺคโต เอโก ขนฺโธ จ วตฺถุ จ ติณฺณนฺนํ ขนฺธานํ นิสฺสย\-ปจฺจเยน ปจฺจโย ฯเปฯ เทฺว ขนฺธา ฯเปฯ. (2)}

\csromanheader{2}{\textbf{อุปนิสฺสย}}
\proseitem{72}{ปริตฺโต ธมฺโม ปริตฺตสฺส ธมฺมสฺส อุปนิสฺสย\-ปจฺจเยน ปจฺจโย. อารมฺมณู\-ปนิสฺสโย, อนนฺตรู\-ปนิสฺสโย, ปกตู\-ปนิสฺสโย ฯเปฯ. ปกตู\-ปนิสฺสโย\csromandash{}ปริตฺตํ สทฺธํ อุปนิสฺสาย ทานํ เทติ, สีลํ สมาทิยติ, อุโปสถกมฺมํ กโรติ, วิปสฺสนํ อุปฺปาเทติ, มานํ ชปฺเปติ, ทิฏฺฐิํ คณฺหาติ, ปริตฺตํ สีลํ ฯเปฯ ปญฺญํ. ราคํ ฯเปฯ ปตฺถนํ, กายิกํ สุขํ ฯเปฯ. เสนาสนํ อุปนิสฺสาย ทานํ เทติ, สีลํ ฯเปฯ อุโปสถํ ฯเปฯ วิปสฺสนํ อุปฺปาเทติ, ปาณํ หนติ ฯเปฯ สํฆํ ภินฺทติ. ปริตฺตา สทฺธา ฯเปฯ ปญฺญา. ราโค ฯเปฯ ปตฺถนา, กายิกํ สุขํ ฯเปฯ. เสนาสนํ ปริตฺตาย สทฺธาย ฯเปฯ ปญฺญาย. ราคสฺส ฯเปฯ ปตฺถนาย. กายิกสฺส สุขสฺส, กายิกสฺส ทุกฺขสฺส อุปนิสฺสย\-ปจฺจเยน ปจฺจโย. กุสลากุสลํ กมฺมํ วิปากสฺส อุปนิสฺสย\-ปจฺจเยน ปจฺจโย. ปาณาติปาโต ปาณาติปาตสฺส อุปนิสฺสย\-ปจฺจเยน ปจฺจโย. (จกฺกํ กาตพฺพํ.) มาตุ\-ฆาติ\-กมฺมํ มาตุ\-ฆาติ\-กมฺมสฺส อุปนิสฺสย\-ปจฺจเยน ปจฺจโย. (จกฺกํ กาตพฺพํ, กุสล\-ตฺติก\-สทิสํ.) (1)}

\prose{ปริตฺโต ธมฺโม มหคฺคตสฺส ธมฺมสฺส อุปนิสฺสย\-ปจฺจเยน ปจฺจโย. อนนฺตรู\-ปนิสฺสโย, ปกตู\-ปนิสฺสโย ฯเปฯ. ปกตู\-ปนิสฺสโย\csromandash{} ปริตฺตํ สทฺธํ อุปนิสฺสาย มหคฺคตํ ฌานํ อุปฺปาเทติ. อภิญฺญํ อุปฺปาเทติ, สมาปตฺติํ อุปฺปาเทติ. ปริตฺตํ สีลํ ฯเปฯ ปญฺญํ. ราคํ ฯเปฯ. เสนาสนํ อุปนิสฺสาย มหคฺคตํ ฌานํ อุปฺปาเทติ, อภิญฺญํ อุปฺปาเทติ, สมาปตฺติํ อุปฺปาเทติ. ปริตฺตา สทฺธา ฯเปฯ. เสนาสนํ มหคฺคตาย สทฺธาย ฯเปฯ ปญฺญาย อุปนิสฺสย\-ปจฺจเยน ปจฺจโย. ปฐมสฺส ฌานสฺส ปริกมฺมํ ฯเปฯ. เนว\-สญฺญา\-นา\-สญฺญา\-ยตนสฺส ปริกมฺมํ เนว\-สญฺญา\-นา\-สญฺญา\-ยตนสฺส อุปนิสฺสย\-ปจฺจเยน ปจฺจโย. ทิพฺพสฺส จกฺขุสฺส ปริกมฺมํ ฯเปฯ. อนาคตํสญาณสฺส ปริกมฺมํ ฯเปฯ. (2)}

\prose{ปริตฺโต ธมฺโม อปฺปมาณสฺส ธมฺมสฺส อุปนิสฺสย\-ปจฺจเยน ปจฺจโย. อนนฺตรู\-ปนิสฺสโย, ปกตู\-ปนิสฺสโย ฯเปฯ. ปกตู\-ปนิสฺสโย\csromandash{} \csromanfolio{332}ปริตฺตํ สทฺธํ อุปนิสฺสาย อปฺปมาณํ ฌานํ อุปฺปาเทติ, มคฺคํ อุปฺปาเทติ, ผลสมาปตฺติํ อุปฺปาเทติ. ปริตฺตํ สีลํ ฯเปฯ ปญฺญํ. ราคํ ฯเปฯ ปตฺถนํ. กายิกํ สุขํ ฯเปฯ. เสนาสนํ อุปนิสฺสาย อปฺปมาณํ ฌานํ อุปฺปาเทติ, มคฺคํ อุปฺปาเทติ, ผลสมาปตฺติํ อุปฺปาเทติ. ปริตฺตา สทฺธา ฯเปฯ. เสนาสนํ อปฺปมาณาย สทฺธาย ฯเปฯ ปญฺญาย. มคฺคสฺส, ผลสมาปตฺติยา อุปนิสฺสย\-ปจฺจเยน ปจฺจโย. ปฐมสฺส มคฺคสฺส ปริกมฺมํ ปฐมสฺส มคฺคสฺส ฯเปฯ. จตุตฺถสฺส มคฺคสฺส ปริกมฺมํ จตุตฺถสฺส มคฺคสฺส อุปนิสฺสย\-ปจฺจเยน ปจฺจโย. (3)}

\proseitem{73}{มหคฺคโต ธมฺโม มหคฺคตสฺส ธมฺมสฺส อุปนิสฺสย\-ปจฺจเยน ปจฺจโย. อนนฺตรู\-ปนิสฺสโย, ปกตู\-ปนิสฺสโย ฯเปฯ. ปกตู\-ปนิสฺสโย\csromandash{}มหคฺคตํ สทฺธํ อุปนิสฺสาย มหคฺคตํ ฌานํ อุปฺปาเทติ, อภิญฺญํ อุปฺปาเทติ, สมาปตฺติํ อุปฺปาเทติ. มหคฺคตํ สีลํ ฯเปฯ ปญฺญํ อุปนิสฺสาย มหคฺคตํ ฌานํ อุปฺปาเทติ. อภิญฺญํ อุปฺปาเทติ. สมาปตฺติํ อุปฺปาเทติ. มหคฺคตา สทฺธา ฯเปฯ ปญฺญา มหคฺคตาย สทฺธาย ฯเปฯ ปญฺญาย อุปนิสฺสย\-ปจฺจเยน ปจฺจโย. ปฐมํ ฌานํ ทุติยสฺส ฌานสฺส ฯเปฯ. อากิญฺจญฺญา\-ยตนํ เนว\-สญฺญา\-นา\-สญฺญา\-ยตนสฺส อุปนิสฺสย\-ปจฺจเยน ปจฺจโย. (1)}

\prose{มหคฺคโต ธมฺโม ปริตฺตสฺส ธมฺมสฺส อุปนิสฺสย\-ปจฺจเยน ปจฺจโย. อารมฺมณู\-ปนิสฺสโย, อนนฺตรู\-ปนิสฺสโย, ปกตู\-ปนิสฺสโย ฯเปฯ. ปกตู\-ปนิสฺสโย\csromandash{}มหคฺคตํ สทฺธํ อุปนิสฺสาย ทานํ เทติ, สีลํ สมาทิยติ, อุโปสถกมฺมํ กโรติ, วิปสฺสนํ อุปฺปาเทติ, มานํ ชปฺเปติ, ทิฏฺฐิํ คณฺหาติ, มหคฺคตํ สีลํ ฯเปฯ ปญฺญํ อุปนิสฺสาย ทานํ เทติ ฯเปฯ วิปสฺสนํ อุปฺปาเทติ ฯเปฯ. มหคฺคตา สทฺธา ฯเปฯ ปญฺญา ปริตฺตาย สทฺธาย ฯเปฯ ปญฺญาย ฯเปฯ. กายิกสฺส สุขสฺส, กายิกสฺส ทุกฺขสฺส อุปนิสฺสย\-ปจฺจเยน ปจฺจโย. (2)}

\prose{มหคฺคโต ธมฺโม อปฺปมาณสฺส ธมฺมสฺส อุปนิสฺสย\-ปจฺจเยน ปจฺจโย. อนนฺตรู\-ปนิสฺสโย, ปกตู\-ปนิสฺสโย ฯเปฯ. ปกตู\-ปนิสฺสโย\csromandash{} มหคฺคตํ สทฺธํ อุปนิสฺสาย อปฺปมาณํ ฌานํ อุปฺปาเทติ, มคฺคํ อุปฺปาเทติ, ผลสมาปตฺติํ อุปฺปาเทติ. มหคฺคตํ สีลํ ฯเปฯ ปญฺญํ อุปนิสฺสาย อปฺปมาณํ ฌานํ อุปฺปาเทติ, มคฺคํ อุปฺปาเทติ, ผลสมาปตฺติํ อุปฺปาเทติ. มหคฺคตา สทฺธา ฯเปฯ ปญฺญา อปฺปมาณาย สทฺธาย ฯเปฯ ปญฺญาย. มคฺคสฺส, ผลสมาปตฺติยา อุปนิสฺสย\-ปจฺจเยน ปจฺจโย. (3)}

\proseitem{74}{\csromanfolio{333}อปฺปมาโณ ธมฺโม อปฺปมาณสฺส ธมฺมสฺส อุปนิสฺสย\-ปจฺจเยน ปจฺจโย. อารมฺมณู\-ปนิสฺสโย, อนนฺตรู\-ปนิสฺสโย, ปกตู\-ปนิสฺสโย ฯเปฯ. ปกตู\-ปนิสฺสโย\csromandash{}อปฺปมาณํ สทฺธํ อุปนิสฺสาย อปฺปมาณํ ฌานํ อุปฺปาเทติ, มคฺคํ อุปฺปาเทติ, ผลสมาปตฺติํ อุปฺปาเทติ. อปฺปมาณํ สีลํ ฯเปฯ ปญฺญํ อุปนิสฺสาย อปฺปมาณํ ฌานํ อุปฺปาเทติ, มคฺคํ อุปฺปาเทติ, ผลสมาปตฺติํ อุปฺปาเทติ, อปฺปมาณา สทฺธา ฯเปฯ ปญฺญา อปฺปมาณาย สทฺธาย ฯเปฯ ปญฺญาย อุปนิสฺสย\-ปจฺจเยน ปจฺจโย. ปฐโม มคฺโค ทุติยสฺส มคฺคสฺส ฯเปฯ. ตติโย มคฺโค จตุตฺถสฺส มคฺคสฺส. มคฺโค ผลสมาปตฺติยา อุปนิสฺสย\-ปจฺจเยน ปจฺจโย. (1)}

\prose{อปฺปมาโณ ธมฺโม ปริตฺตสฺส ธมฺมสฺส อุปนิสฺสย\-ปจฺจเยน ปจฺจโย. อารมฺมณู\-ปนิสฺสโย, อนนฺตรู\-ปนิสฺสโย, ปกตู\-ปนิสฺสโย ฯเปฯ. ปกตู\-ปนิสฺสโย\csromandash{}อปฺปมาณํ สทฺธํ อุปนิสฺสาย ทานํ เทติ, สีลํ ฯเปฯ อุโปสถกมฺมํ. วิปสฺสนํ อุปฺปาเทติ. อปฺปมาณํ สีลํ ฯเปฯ ปญฺญํ อุปนิสฺสาย ทานํ เทติ, สีลํ ฯเปฯ อุโปสถกมฺมํ. วิปสฺสนํ. อปฺปมาณา สทฺธา ฯเปฯ ปญฺญา ปริตฺตาย สทฺธาย ฯเปฯ ปญฺญาย ฯเปฯ. กายิกสฺส สุขสฺส, กายิกสฺส ทุกฺขสฺส อุปนิสฺสย\-ปจฺจเยน ปจฺจโย. ผลสมาปตฺติ กายิกสฺส สุขสฺส อุปนิสฺสย\-ปจฺจเยน ปจฺจโย. อริยา มคฺคํ อุปนิสฺสาย สงฺขาเร อนิจฺจโต ฯเปฯ วิปสฺสนฺติ. มคฺโค อริยานํ อตฺถ\-ปฺปฏิสมฺภิทาย ฯเปฯ. ปฏิภาน\-ปฺปฏิสมฺภิทาย. ฐานาฐาน\-โกสลฺลสฺส อุปนิสฺสย\-ปจฺจเยน ปจฺจโย. (2)}

\prose{อปฺปมาโณ ธมฺโม มหคฺคตสฺส ธมฺมสฺส อุปนิสฺสย\-ปจฺจเยน ปจฺจโย. อนนฺตรู\-ปนิสฺสโย, ปกตู\-ปนิสฺสโย ฯเปฯ. ปกตู\-ปนิสฺสโย\csromandash{} อปฺปมาณํ สทฺธํ อุปนิสฺสาย มหคฺคตํ ฌานํ อุปฺปาเทติ, อภิญฺญํ อุปฺปาเทติ, สมาปตฺติํ อุปฺปาเทติ. อปฺปมาณํ สีลํ ฯเปฯ ปญฺญํ อุปนิสฺสาย มหคฺคตํ ฌานํ อุปฺปาเทติ, อภิญฺญํ อุปฺปาเทติ, สมาปตฺติํ อุปฺปาเทติ. อปฺปมาณา สทฺธา ฯเปฯ ปญฺญา มหคฺคตาย สทฺธาย ฯเปฯ ปญฺญาย อุปนิสฺสย\-ปจฺจเยน ปจฺจโย. อริยา มคฺคํ อุปนิสฺสาย อนุปฺปนฺนํ สมาปตฺติํ อุปฺปาเทนฺติ, อุปฺปนฺนํ สมาปชฺชนฺติ. (3)}

\csromanheader{2}{\textbf{ปุเรชาต}}
\proseitem{75}{ปริตฺโต ธมฺโม ปริตฺตสฺส ธมฺมสฺส ปุเรชาต\-ปจฺจเยน ปจฺจโย. อารมฺมณ\-ปุเรชาตํ, วตฺถุ\-ปุเรชาตํ.  อารมฺมณ\-ปุเรชาตํ\csromandash{}จกฺขุํ \csromanfolio{334}อนิจฺจโต ฯเปฯ วิปสฺสติ, อสฺสาเทติ อภินนฺทติ, ตํ อารพฺภ ราโค อุปฺปชฺชติ ฯเปฯ โทมนสฺสํ อุปฺปชฺชติ. โสตํ ฯเปฯ. วตฺถุํ อนิจฺจโต ฯเปฯ โทมนสฺสํ อุปฺปชฺชติ. รูปายตนํ จกฺขุ\-วิญฺญาณสฺส ฯเปฯ. โผฏฺฐพฺพา\-ยตนํ กายวิญฺญาณสฺส ปุเรชาต\-ปจฺจเยน ปจฺจโย. วตฺถุ\-ปุเรชาตํ\csromandash{} จกฺขายตนํ จกฺขุ\-วิญฺญาณสฺส ฯเปฯ. กายายตนํ กายวิญฺญาณสฺส. วตฺถุ ปริตฺตานํ ขนฺธานํ ปุเรชาต\-ปจฺจเยน ปจฺจโย. (1)}

\prose{ปริตฺโต ธมฺโม มหคฺคตสฺส ธมฺมสฺส ปุเรชาต\-ปจฺจเยน ปจฺจโย. อารมฺมณ\-ปุเรชาตํ, วตฺถุ\-ปุเรชาตํ.  อารมฺมณ\-ปุเรชาตํ\csromandash{}ทิพฺเพน จกฺขุนา รูปํ ปสฺสติ. ทิพฺพาย โสตธาตุยา สทฺทํ สุณาติ. วตฺถุ\-ปุเรชาตํ\csromandash{}วตฺถุ มหคฺคตานํ ขนฺธานํ ปุเรชาต\-ปจฺจเยน ปจฺจโย. (2)}

\prose{ปริตฺโต ธมฺโม อปฺปมาณสฺส ธมฺมสฺส ปุเรชาต\-ปจฺจเยน ปจฺจโย. วตฺถุ\-ปุเรชาตํ\csromandash{}วตฺถุ อปฺปมาณานํ ขนฺธานํ ปุเรชาต\-ปจฺจเยน ปจฺจโย. (3)}

\csromanheader{2}{\textbf{ปจฺฉาชาต}}
\proseitem{76}{ปริตฺโต ธมฺโม ปริตฺตสฺส ธมฺมสฺส ปจฺฉาชาต\-ปจฺจเยน ปจฺจโย. ปจฺฉาชาตา ปริตฺตา ขนฺธา ปุเรชาตสฺส อิมสฺส กายสฺส ปจฺฉาชาต\-ปจฺจเยน ปจฺจโย. (1)}

\prose{มหคฺคโต ธมฺโม ปริตฺตสฺส ธมฺมสฺส ปจฺฉาชาต\-ปจฺจเยน ปจฺจโย. ปจฺฉาชาตา มหคฺคตา ขนฺธา ปุเรชาตสฺส อิมสฺส กายสฺส ปจฺฉชาตปจฺจเยน ปจฺจโย. (1)}

\prose{อปฺปมาโณ ธมฺโม ปริตฺตสฺส ธมฺมสฺส ปจฺฉาชาต\-ปจฺจเยน ปจฺจโย. ปจฺฉาชาตา อปฺปมาณา ขนฺธา ปุเรชาตสฺส อิมสฺส กายสฺส ปจฺฉาชาต\-ปจฺจเยน ปจฺจโย. (1)}

\csromanheader{2}{\textbf{อาเสวน}}
\proseitem{77}{ปริตฺโต ธมฺโม ปริตฺตสฺส ธมฺมสฺส อาเสวน\-ปจฺจเยน ปจฺจโย. ปุริมา ปุริมา ปริตฺตา ขนฺธา ปจฺฉิมานํ ปจฺฉิมานํ ปริตฺตานํ ขนฺธานํ ฯเปฯ. อนุโลมํ โคตฺรภุสฺส. อนุโลมํ โวทานสฺส อาเสวน\-ปจฺจเยน ปจฺจโย. (1)}

\prose{\csromanfolio{335}ปริตฺโต ธมฺโม มหคฺคตสฺส ธมฺมสฺส อาเสวน\-ปจฺจเยน ปจฺจโย. ปฐมสฺส ฌานสฺส ปริกมฺมํ ตสฺเสว\csromansharedfootnote{4771cfefe301ea99}{ปฐมสฺส ฌานสฺส (?)} อาเสวน\-ปจฺจเยน ปจฺจโย ฯเปฯ. เนว\-สญฺญา\-นา\-สญฺญา\-ยตนสฺส ปริกมฺมํ ตสฺเสว\csromansharedfootnote{69b7afeda991d6fd}{เนวสญฺญานาสญฺญายตนสฺส (?) อญฺเญสุ ติเกสุ โอโลเกตพฺพํ.} อาเสวน\-ปจฺจเยน ปจฺจโย. ทิพฺพสฺส จกฺขุสฺส ปริกมฺมํ ฯเปฯ. อนาคตํสญาณสฺส ปริกมฺมํ อนาคตํสญาณสฺส อาเสวน\-ปจฺจเยน ปจฺจโย. (2)}

\prose{ปริตฺโต ธมฺโม อปฺปมาณสฺส ธมฺมสฺส อาเสวน\-ปจฺจเยน ปจฺจโย. โคตฺรภุ มคฺคสฺส. โวทานํ มคฺคสฺส อาเสวน\-ปจฺจเยน ปจฺจโย. (3)}

\prose{มหคฺคโต ธมฺโม มหคฺคตสฺส ธมฺมสฺส อาเสวน\-ปจฺจเยน ปจฺจโย. ปุริมา ปุริมา มหคฺคตา ขนฺธา ปจฺฉิมานํ ปจฺฉิมานํ มหคฺคตานํ ขนฺธานํ ฯเปฯ. อาเสวน\-ปจฺจเยน ปจฺจโย. (1)}

\csromanheader{2}{\textbf{กมฺม}}
\proseitem{78}{ปริตฺโต ธมฺโม ปริตฺตสฺส ธมฺมสฺส กมฺมปจฺจเยน ปจฺจโย. สหชาตา, นานากฺขณิกา.  สหชาตา ปริตฺตา เจตนา สมฺปยุตฺตกานํ ขนฺธานํ จิตฺตสมุฏฺฐานานญฺจ รูปานํ กมฺมปจฺจเยน ปจฺจโย. ปฏิสนฺธิ\-กฺขเณ ปริตฺตา เจตนา สมฺปยุตฺตกานํ ขนฺธานํ กฏตฺตา จ รูปานํ กมฺมปจฺจเยน ปจฺจโย. นานากฺขณิกา ปริตฺตา เจตนา วิปากานํ ปริตฺตานํ ขนฺธานํ กฏตฺตา จ รูปานํ กมฺมปจฺจเยน ปจฺจโย. (1)}

\prose{มหคฺคโต ธมฺโม มหคฺคตสฺส ธมฺมสฺส กมฺมปจฺจเยน ปจฺจโย. สหชาตา, นานากฺขณิกา.  สหชาตา มหคฺคตา เจตนา สมฺปยุตฺตกานํ ขนฺธานํ กมฺมปจฺจเยน ปจฺจโย. ปฏิสนฺธิ\-กฺขเณ มหคฺคตา เจตนา สมฺปยุตฺตกานํ ขนฺธานํ กมฺมปจฺจเยน ปจฺจโย. นานากฺขณิกา มหคฺคตา เจตนา วิปากานํ มหคฺคตานํ ขนฺธานํ กมฺมปจฺจเยน ปจฺจโย. (1)}

\prose{มหคฺคโต ธมฺโม ปริตฺตสฺส ธมฺมสฺส กมฺมปจฺจเยน ปจฺจโย. สหชาตา, นานากฺขณิกา.  สหชาตา มหคฺคตา เจตนา จิตฺต\-สมุฏฺฐานานํ รูปานํ กมฺมปจฺจเยน ปจฺจโย. ปฏิสนฺธิ\-กฺขเณ มหคฺคตา เจตนา กฏตฺตารูปานํ กมฺมปจฺจเยน ปจฺจโย. นานากฺขณิกา มหคฺคตา เจตนา กฏตฺตารูปานํ กมฺมปจฺจเยน ปจฺจโย. (2)}

\prose{\csromanfolio{336}มหคฺคโต ธมฺโม ปริตฺตสฺส จ มหคฺคตสฺส จ ธมฺมสฺส กมฺมปจฺจเยน ปจฺจโย. สหชาตา, นานากฺขณิกา.  สหชาตา มหคฺคตา เจตนา สมฺปยุตฺตกานํ ขนฺธานํ จิตฺตสมุฏฺฐานานญฺจ รูปานํ กมฺมปจฺจเยน ปจฺจโย. ปฏิสนฺธิ\-กฺขเณ มหคฺคตา เจตนา สมฺปยุตฺตกานํ ขนฺธานํ กฏตฺตา จ รูปานํ กมฺมปจฺจเยน ปจฺจโย. นานากฺขณิกา มหคฺคตา เจตนา วิปากานํ มหคฺคตานํ ขนฺธานํ กฏตฺตา จ รูปานํ กมฺมปจฺจเยน ปจฺจโย. (3)}

\proseitem{79}{อปฺปมาโณ ธมฺโม อปฺปมาณสฺส ธมฺมสฺส กมฺมปจฺจเยน ปจฺจโย. สหชาตา, นานากฺขณิกา.  \textbf{สหชาตา} อปฺปมาณา เจตนา สมฺปยุตฺตกานํ ขนฺธานํ กมฺมปจฺจเยน ปจฺจโย. \textbf{นานากฺขณิกา} อปฺปมาณา เจตนา วิปากานํ อปฺปมาณานํ ขนฺธานํ กมฺมปจฺจเยน ปจฺจโย. (1)}

\prose{อปฺปมาโณ ธมฺโม ปริตฺตสฺส ธมฺมสฺส กมฺมปจฺจเยน ปจฺจโย. อปฺปมาณา เจตนา จิตฺต\-สมุฏฺฐานานํ รูปานํ กมฺมปจฺจเยน ปจฺจโย. (2)}

\prose{อปฺปมาโณ ธมฺโม ปริตฺตสฺส จ อปฺปมาณสฺส จ ธมฺมสฺส กมฺมปจฺจเยน ปจฺจโย. อปฺปมาณา เจตนา สมฺปยุตฺตกานํ ขนฺธานํ จิตฺตสมุฏฺฐานานญฺจ รูปานํ กมฺมปจฺจเยน ปจฺจโย. (3)}

\csromanheader{2}{\textbf{วิปาก}}
\proseitem{80}{ปริตฺโต ธมฺโม ปริตฺตสฺส ธมฺมสฺส วิปาก\-ปจฺจเยน ปจฺจโย. วิปาโก ปริตฺโต เอโก ขนฺโธ ติณฺณนฺนํ ขนฺธานํ จิตฺตสมุฏฺฐานานญฺจ รูปานํ วิปาก\-ปจฺจเยน ปจฺจโย ฯเปฯ. ปฏิสนฺธิ\-กฺขเณ ฯเปฯ. ขนฺธา วตฺถุสฺส วิปาก\-ปจฺจเยน ปจฺจโย. (1)}

\prose{มหคฺคโต ธมฺโม มหคฺคตสฺส ธมฺมสฺส วิปาก\-ปจฺจเยน ปจฺจโย ฯเปฯ. (ติสฺโส ปญฺหา, ปวตฺติ\-ปฏิสนฺธิ กาตพฺพา.) (3)}

\prose{อปฺปมาโณ ธมฺโม อปฺปมาณสฺส ธมฺมสฺส วิปาก\-ปจฺจเยน ปจฺจโย. ตีณิ. (ปวตฺติเมว.)}

\csromanheader{2}{\textbf{อาหาราทิ}}
\proseitem{81}{ปริตฺโต ธมฺโม ปริตฺตสฺส ธมฺมสฺส อาหารปจฺจเยน ปจฺจโย ฯเปฯ. อินฺทฺริย\-ปจฺจเยน ปจฺจโย. ฌานปจฺจเยน ปจฺจโย. มคฺคปจฺจเยน ปจฺจโย. \csromanfolio{337}สมฺปยุตฺต\-ปจฺจเยน ปจฺจโย. วิปฺปยุตฺต\-ปจฺจเยน ปจฺจโย. สหชาตํ, ปุเรชาตํ, ปจฺฉาชาตํ.  สหชาตา ปริตฺตา ขนฺธา จิตฺต\-สมุฏฺฐานานํ รูปานํ วิปฺปยุตฺต\-ปจฺจเยน ปจฺจโย. ปฏิสนฺธิ\-กฺขเณ ปริตฺตา ขนฺธา กฏตฺตารูปานํ วิปฺปยุตฺต\-ปจฺจเยน ปจฺจโย. ขนฺธา วตฺถุสฺส ฯเปฯ. วตฺถุ ขนฺธานํ วิปฺปยุตฺต\-ปจฺจเยน ปจฺจโย. ปุเรชาตํ จกฺขายตนํ จกฺขุ\-วิญฺญาณสฺส วิปฺปยุตฺต\-ปจฺจเยน ปจฺจโย ฯเปฯ. กายายตนํ กายวิญฺญาณสฺส วิปฺปยุตฺต\-ปจฺจเยน ปจฺจโย. วตฺถุ ปริตฺตานํ ขนฺธานํ วิปฺปยุตฺต\-ปจฺจเยน ปจฺจโย. ปจฺฉาชาตา ปริตฺตา ขนฺธา ปุเรชาตสฺส อิมสฺส กายสฺส วิปฺปยุตฺต\-ปจฺจเยน ปจฺจโย. (1)}

\prose{ปริตฺโต ธมฺโม มหคฺคตสฺส ธมฺมสฺส วิปฺปยุตฺต\-ปจฺจเยน ปจฺจโย. สหชาตํ, ปุเรชาตํ.  สหชาตํ ปฏิสนฺธิ\-กฺขเณ วตฺถุ มหคฺคตานํ ขนฺธานํ วิปฺปยุตฺต\-ปจฺจเยน ปจฺจโย. ปุเรชาตํ วตฺถุ มหคฺคตานํ ขนฺธานํ วิปฺปยุตฺต\-ปจฺจเยน ปจฺจโย. (2)}

\prose{ปริตฺโต ธมฺโม อปฺปมาณสฺส ธมฺมสฺส วิปฺปยุตฺต\-ปจฺจเยน ปจฺจโย. ปุเรชาตํ วตฺถุ อปฺปมาณานํ ขนฺธานํ วิปฺปยุตฺต\-ปจฺจเยน ปจฺจโย. (3)}

\proseitem{82}{มหคฺคโต ธมฺโม ปริตฺตสฺส ธมฺมสฺส วิปฺปยุตฺต\-ปจฺจเยน ปจฺจโย. สหชาตํ, ปจฺฉาชาตํ.  สหชาตา มหคฺคตา ขนฺธา จิตฺต\-สมุฏฺฐานานํ รูปานํ วิปฺปยุตฺต\-ปจฺจเยน ปจฺจโย. ปฏิสนฺธิ\-กฺขเณ ฯเปฯ. ปจฺฉาชาตา มหคฺคตา ขนฺธา ปุเรชาตสฺส อิมสฺส กายสฺส วิปฺปยุตฺต\-ปจฺจเยน ปจฺจโย. (1)}

\prose{อปฺปมาโณ ธมฺโม ปริตฺตสฺส ธมฺมสฺส วิปฺปยุตฺต\-ปจฺจเยน ปจฺจโย. สหชาตํ, ปจฺฉาชาตํ.  สหชาตา อปฺปมาณา ขนฺธา จิตฺต\-สมุฏฺฐานานํ รูปานํ วิปฺปยุตฺต\-ปจฺจเยน ปจฺจโย. ปจฺฉาชาตา อปฺปมาณา ขนฺธา ปุเรชาตสฺส อิมสฺส กายสฺส วิปฺปยุตฺต\-ปจฺจเยน ปจฺจโย. (1)}

\csromanheader{2}{\textbf{อตฺถิ}}
\proseitem{83}{ปริตฺโต ธมฺโม ปริตฺตสฺส ธมฺมสฺส อตฺถิปจฺจเยน ปจฺจโย. สหชาตํ, ปุเรชาตํ, ปจฺฉาชาตํ, อาหารํ, อินฺทฺริยํ.  สหชาโต ปริตฺโต เอโก ขนฺโธ ติณฺณนฺนํ ขนฺธานํ จิตฺตสมุฏฺฐานานญฺจ รูปานํ อตฺถิปจฺจเยน ปจฺจโย ฯเปฯ เทฺว ขนฺธา ฯเปฯ. ปฏิสนฺธิ\-กฺขเณ ฯเปฯ. ขนฺธา วตฺถุสฺส \csromanfolio{338}อตฺถิปจฺจเยน ปจฺจโย. วตฺถุ ขนฺธานํ อตฺถิปจฺจเยน ปจฺจโย. เอกํ มหาภูตํ ฯเปฯ. อสญฺญสตฺตานํ ฯเปฯ. ปุเรชาตํ จกฺขุํ ฯเปฯ. วตฺถุํ อนิจฺจโต ฯเปฯ วิปสฺสติ, อสฺสาเทติ อภินนฺทติ, ตํ อารพฺภ ราโค อุปฺปชฺชติ ฯเปฯ โทมนสฺสํ อุปฺปชฺชติ. รูปายตนํ จกฺขุ\-วิญฺญาณสฺส ฯเปฯ. โผฏฺฐพฺพา\-ยตนํ กายวิญฺญาณสฺส อตฺถิปจฺจเยน ปจฺจโย. จกฺขายตนํ จกฺขุ\-วิญฺญาณสฺส อตฺถิปจฺจเยน ปจฺจโย ฯเปฯ. กายายตนํ กายวิญฺญาณสฺส ฯเปฯ. วตฺถุ ปริตฺตานํ ขนฺธานํ อตฺถิปจฺจเยน ปจฺจโย. ปจฺฉาชาตา ปริตฺตา ขนฺธา ปุเรชาตสฺส อิมสฺส กายสฺส อตฺถิปจฺจเยน ปจฺจโย. กพฬีกาโร อาหาโร อิมสฺส กายสฺส อตฺถิปจฺจเยน ปจฺจโย. รูปชีวิตินฺทฺริยํ กฏตฺตารูปานํ อตฺถิปจฺจเยน ปจฺจโย. (1)}

\prose{ปริตฺโต ธมฺโม มหคฺคตสฺส ธมฺมสฺส อตฺถิปจฺจเยน ปจฺจโย. สหชาตํ, ปุเรชาตํ.  สหชาตํ ปฏิสนฺธิ\-กฺขเณ วตฺถุ มหคฺคตานํ ขนฺธานํ อตฺถิปจฺจเยน ปจฺจโย. ปุเรชาตํ ทิพฺเพน จกฺขุนา รูปํ ปสฺสติ. ทิพฺพาย โสตธาตุยา สทฺทํ สุณาติ. วตฺถุ มหคฺคตานํ ขนฺธานํ อตฺถิปจฺจเยน ปจฺจโย. (2)}

\prose{ปริตฺโต ธมฺโม อปฺปมาณสฺส ธมฺมสฺส อตฺถิปจฺจเยน ปจฺจโย. \textbf{ปุเรชาตํ} วตฺถุ อปฺปมาณานํ ขนฺธานํ อตฺถิปจฺจเยน ปจฺจโย. (3)}

\proseitem{84}{มหคฺคโต ธมฺโม มหคฺคตสฺส ธมฺมสฺส อตฺถิปจฺจเยน ปจฺจโย. มหคฺคโต เอโก ขนฺโธ ติณฺณนฺนํ ขนฺธานํ ฯเปฯ เทฺว ขนฺธา ฯเปฯ. ปฏิสนฺธิ\-กฺขเณ ฯเปฯ. (1)}

\prose{มหคฺคโต ธมฺโม ปริตฺตสฺส ธมฺมสฺส อตฺถิปจฺจเยน ปจฺจโย. สหชาตํ, ปจฺฉาชาตํ.  สหชาตา มหคฺคตา ขนฺธา จิตฺต\-สมุฏฺฐานานํ รูปานํ อตฺถิปจฺจเยน ปจฺจโย. ปฏิสนฺธิ\-กฺขเณ มหคฺคตา ขนฺธา กฏตฺตารูปานํ อตฺถิปจฺจเยน ปจฺจโย. ปจฺฉาชาตา มหคฺคตา ขนฺธา ปุเรชาตสฺส อิมสฺส กายสฺส อตฺถิปจฺจเยน ปจฺจโย. (2)}

\prose{มหคฺคโต ธมฺโม ปริตฺตสฺส จ มหคฺคตสฺส จ ธมฺมสฺส อตฺถิปจฺจเยน ปจฺจโย. มหคฺคโต เอโก ขนฺโธ ติณฺณนฺนํ ขนฺธานํ จิตฺตสมุฏฺฐานานญฺจ รูปานํ อตฺถิปจฺจเยน ปจฺจโย ฯเปฯ เทฺว ขนฺธา ฯเปฯ. ปฏิสนฺธิ\-กฺขเณ ฯเปฯ. (3)}

\proseitem{85}{\csromanfolio{339}อปฺปมาโณ ธมฺโม อปฺปมาณสฺส ธมฺมสฺส อตฺถิปจฺจเยน ปจฺจโย. อปฺปมาโณ เอโก ขนฺโธ ติณฺณนฺนํ ขนฺธานํ ฯเปฯ. (1)}

\prose{อปฺปมาโณ ธมฺโม ปริตฺตสฺส ธมฺมสฺส อตฺถิปจฺจเยน ปจฺจโย. สหชาตํ, ปจฺฉาชาตํ.  สหชาตา อปฺปมาณา ขนฺธา จิตฺต\-สมุฏฺฐานานํ รูปานํ อตฺถิปจฺจเยน ปจฺจโย. ปจฺฉาชาตา อปฺปมาณา ขนฺธา ปุเรชาตสฺส อิมสฺส กายสฺส อตฺถิปจฺจเยน ปจฺจโย. (2)}

\prose{อปฺปมาโณ ธมฺโม ปริตฺตสฺส จ อปฺปมาณสฺส จ ธมฺมสฺส อตฺถิปจฺจเยน ปจฺจโย. อปฺปมาโณ เอโก ขนฺโธ ติณฺณนฺนํ ขนฺธานํ จิตฺตสมุฏฺฐานานญฺจ รูปานํ อตฺถิปจฺจเยน ปจฺจโย ฯเปฯ. (3)}

\proseitem{86}{ปริตฺโต จ อปฺปมาโณ จ ธมฺมา ปริตฺตสฺส ธมฺมสฺส อตฺถิปจฺจเยน ปจฺจโย. สหชาตํ, ปจฺฉาชาตํ, อาหารํ, อินฺทฺริยํ.  สหชาตา อปฺปมาณา ขนฺธา จ มหาภูตา จ จิตฺต\-สมุฏฺฐานานํ รูปานํ อตฺถิปจฺจเยน ปจฺจโย. ปจฺฉาชาตา อปฺปมาณา ขนฺธา จ กพฬีกาโร อาหาโร จ อิมสฺส กายสฺส อตฺถิปจฺจเยน ปจฺจโย. ปจฺฉาชาตา อปฺปมาณา ขนฺธา จ รูปชีวิตินฺทฺริยญฺจ กฏตฺตารูปานํ อตฺถิปจฺจเยน ปจฺจโย. (1)}

\prose{ปริตฺโต จ อปฺปมาโณ จ ธมฺมา อปฺปมาณสฺส ธมฺมสฺส อตฺถิปจฺจเยน ปจฺจโย. สหชาตํ, ปุเรชาตํ.  สหชาโต อปฺปมาโณ เอโก ขนฺโธ จ วตฺถุ จ ติณฺณนฺนํ ขนฺธานํ อตฺถิปจฺจเยน ปจฺจโย ฯเปฯ เทฺว ขนฺธา ฯเปฯ. (2)}

\prose{ปริตฺโต จ มหคฺคโต จ ธมฺมา ปริตฺตสฺส ธมฺมสฺส อตฺถิปจฺจเยน ปจฺจโย. สหชาตํ, ปจฺฉาชาตํ, อาหารํ, อินฺทฺริยํ.  สหชาตา มหคฺคตา ขนฺธา จ มหาภูตา จ จิตฺต\-สมุฏฺฐานานํ รูปานํ อตฺถิปจฺจเยน ปจฺจโย. ปฏิสนฺธิ\-กฺขเณ มหคฺคตา ขนฺธา จ มหาภูตา จ กฏตฺตารูปานํ อตฺถิปจฺจเยน ปจฺจโย. ปจฺฉาชาตา มหคฺคตา ขนฺธา จ กพฬีกาโร อาหาโร จ อิมสฺส กายสฺส อตฺถิปจฺจเยน ปจฺจโย. ปจฺฉาชาตา มหคฺคตา ขนฺธา จ รูปชีวิตินฺทฺริยญฺจ กฏตฺตารูปานํ อตฺถิปจฺจเยน ปจฺจโย. (1)}

\prose{ปริตฺโต จ มหคฺคโต จ ธมฺมา มหคฺคตสฺส ธมฺมสฺส อตฺถิปจฺจเยน ปจฺจโย. สหชาตํ, ปุเรชาตํ.  สหชาโต มหคฺคโต เอโก ขนฺโธ จ วตฺถุ จ ติณฺณนฺนํ ขนฺธานํ อตฺถิปจฺจเยน ปจฺจโย ฯเปฯ เทฺว ขนฺธา ฯเปฯ. \csromanfolio{340}ปฏิสนฺธิ\-กฺขเณ มหคฺคโต เอโก ขนฺโธ จ วตฺถุ จ ติณฺณนฺนํ ขนฺธานํ อตฺถิปจฺจเยน ปจฺจโย ฯเปฯ เทฺว ขนฺธา จ วตฺถุ จ ฯเปฯ นตฺถิ\-ปจฺจเยน ปจฺจโย. วิคต\-ปจฺจเยน ปจฺจโย. อวิคต\-ปจฺจเยน ปจฺจโย. (2)}

\csromanmatikaanchor{mtk.78}
\csromantocmarkref{subparagraph}{1. ปจฺจยานุโลม 2. สงฺขฺยาวาร}{mtk.78}
\bookmark[dest={mtk.78},level=5]{1. ปจฺจยานุโลม 2. สงฺขฺยาวาร}
\csromanmarkh{1. ปจฺจยานุโลม}\csromanmarkhii{2. สงฺขฺยาวาร}\csromanheadernomark{6}{1. \textbf{ปจฺจยานุโลม 2. สงฺขฺยาวาร}}
\proseitemwithcloserruled{87}{เหตุยา สตฺต, อารมฺมเณ สตฺต, อธิปติยา สตฺต, อนนฺตเร นว, สมนนฺตเร นว, สหชาเต เอกาทส, อญฺญมญฺเญ สตฺต, นิสฺสเย เตรส, อุปนิสฺสเย นว, ปุเรชาเต ตีณิ, ปจฺฉาชาเต ตีณิ, อาเสวเน จตฺตาริ, กมฺเม สตฺต, วิปาเก, อาหาเร, อินฺทฺริเย, ฌาเน, มคฺเค สตฺต, สมฺปยุตฺเต ตีณิ, วิปฺปยุตฺเต ปญฺจ, อตฺถิยา เตรส, นตฺถิยา นว, วิคเต นว, อวิคเต เตรส. (เอวํ คเณตพฺพํ.)}{อนุโลมํ.}
\csromanmatikaanchor{mtk.79}
\csromantocmarkref{subparagraph}{2. ปจฺจนียุทฺธาร}{mtk.79}
\bookmark[dest={mtk.79},level=5]{2. ปจฺจนียุทฺธาร}
\csromanheader{6}{2. ปจฺจนียุ\-ทฺธาร}
\proseitem{88}{ปริตฺโต ธมฺโม ปริตฺตสฺส ธมฺมสฺส อารมฺมณ\-ปจฺจเยน ปจฺจโย. สหชาต\-ปจฺจเยน ปจฺจโย. อุปนิสฺสย\-ปจฺจเยน ปจฺจโย. ปุเรชาต\-ปจฺจเยน ปจฺจโย. ปจฺฉาชาต\-ปจฺจเยน ปจฺจโย. กมฺมปจฺจเยน ปจฺจโย. อาหารปจฺจเยน ปจฺจโย. อินฺทฺริย\-ปจฺจเยน ปจฺจโย. (1)}

\prose{ปริตฺโต ธมฺโม มหคฺคตสฺส ธมฺมสฺส อารมฺมณ\-ปจฺจเยน ปจฺจโย. สหชาต\-ปจฺจเยน ปจฺจโย. อุปนิสฺสย\-ปจฺจเยน ปจฺจโย. ปุเรชาต\-ปจฺจเยน ปจฺจโย. (2)}

\prose{ปริตฺโต ธมฺโม อปฺปมาณสฺส ธมฺมสฺส อุปนิสฺสย\-ปจฺจเยน ปจฺจโย. ปุเรชาต\-ปจฺจเยน ปจฺจโย. (3)}

\prose{มหคฺคโต ธมฺโม มหคฺคตสฺส ธมฺมสฺส อารมฺมณ\-ปจฺจเยน ปจฺจโย. สหชาต\-ปจฺจเยน ปจฺจโย. อุปนิสฺสย\-ปจฺจเยน ปจฺจโย. (1)}

\prose{มหคฺคโต ธมฺโม ปริตฺตสฺส ธมฺมสฺส อารมฺมณ\-ปจฺจเยน ปจฺจโย, สหชาต\-ปจฺจเยน ปจฺจโย. อุปนิสฺสย\-ปจฺจเยน ปจฺจโย. ปจฺฉาชาต\-ปจฺจเยน ปจฺจโย. กมฺมปจฺจเยน ปจฺจโย. (2)}

\prose{\csromanfolio{341}มหคฺคโต ธมฺโม อปฺปมาณสฺส ธมฺมสฺส อุปนิสฺสย\-ปจฺจเยน ปจฺจโย. (3)}

\prose{มหคฺคโต ธมฺโม ปริตฺตสฺส จ มหคฺคตสฺส จ ธมฺมสฺส สหชาต\-ปจฺจเยน ปจฺจโย. กมฺมปจฺจเยน ปจฺจโย. (4)}

\proseitem{89}{อปฺปมาโณ ธมฺโม อปฺปมาณสฺส ธมฺมสฺส สหชาต\-ปจฺจเยน ปจฺจโย. อุปนิสฺสย\-ปจฺจเยน ปจฺจโย. (1)}

\prose{อปฺปมาโณ ธมฺโม ปริตฺตสฺส ธมฺมสฺส อารมฺมณ\-ปจฺจเยน ปจฺจโย. สหชาต\-ปจฺจเยน ปจฺจโย. อุปนิสฺสย\-ปจฺจเยน ปจฺจโย. ปจฺฉาชาต\-ปจฺจเยน ปจฺจโย. (2)}

\prose{อปฺปมาโณ ธมฺโม มหคฺคตสฺส ธมฺมสฺส อารมฺมณ\-ปจฺจเยน ปจฺจโย. อุปนิสฺสย\-ปจฺจเยน ปจฺจโย. (3)}

\prose{อปฺปมาโณ ธมฺโม ปริตฺตสฺส จ อปฺปมาณสฺส จ ธมฺมสฺส สหชาต\-ปจฺจเยน ปจฺจโย. (4)}

\prose{ปริตฺโต จ อปฺปมาโณ จ ธมฺมา ปริตฺตสฺส ธมฺมสฺส สหชาตํ, ปจฺฉาชาตํ, อาหารํ, อินฺทฺริยํ. (1)}

\prose{ปริตฺโต จ อปฺปมาโณ จ ธมฺมา อปฺปมาณสฺส ธมฺมสฺส สหชาตํ, ปุเรชาตํ. (2)}

\prose{ปริตฺโต จ มหคฺคโต จ ธมฺมา ปริตฺตสฺส ธมฺมสฺส สหชาตํ, ปจฺฉาชาตํ, อาหารํ, อินฺทฺริยํ. (1)}

\prose{ปริตฺโต จ มหคฺคโต จ ธมฺมา มหคฺคตสฺส ธมฺมสฺส สหชาตํ, ปุเรชาตํ. (2)}

\csromanmatikaanchor{mtk.80}
\csromantocmarkref{subparagraph}{2. ปจฺจยปจฺจนีย 2. สงฺขฺยาวาร}{mtk.80}
\bookmark[dest={mtk.80},level=5]{2. ปจฺจยปจฺจนีย 2. สงฺขฺยาวาร}
\csromanmarkh{2. ปจฺจย\-ปจฺจนีย}\csromanmarkhii{2. สงฺขฺยาวาร}\csromanheadernomark{6}{2. ปจฺจย\-ปจฺจนีย 2. สงฺขฺยาวาร}
\proseitem{90}{นเหตุยา ปนฺนรส, นอารมฺมเณ ปนฺนรส, นอธิปติยา, นอนนฺตเร, นสมนนฺตเร ปนฺนรส, นสหชาเต ทฺวาทส, นอญฺญมญฺเญ ทฺวาทส, นนิสฺสเย ทฺวาทส, นอุปนิสฺสเย จุทฺทส, นปุเรชาเต จุทฺทส, นปจฺฉาชาเต ปนฺนรส, นอาเสวเน ปนฺนรส ฯเปฯ นมคฺเค ปนฺนรส, นสมฺปยุตฺเต ทฺวาทส, นวิปฺปยุตฺเต ทส, โนอตฺถิยา ทส, โนนตฺถิยา ปนฺนรส, โนวิคเต ปนฺนรส, โนอวิคเต ทส. (เอวํ คเณตพฺพํ.)}

\vspace{\tipitakaparskip}
\csromancenter{ปจฺจนียํ.}
\csromanmatikaanchor{mtk.81}
\csromanmatikaanchor{mtk.82}
\csromanmatikaanchor{mtk.83}
\csromanmatikaanchor{mtk.84}
\csromantocmarknopage{subparagraph}{3. ปจฺจยานุโลมปจฺจนีย}{mtk.81}
\bookmark[dest={mtk.81},level=5]{3. ปจฺจยานุโลมปจฺจนีย}
\csromantocmarkref{subparagraph}{เหตุทุก}{mtk.82}
\bookmark[dest={mtk.82},level=5]{เหตุทุก}
\csromantocmarknopage{subparagraph}{4. ปจฺจยปจฺจนียานุโลม}{mtk.83}
\bookmark[dest={mtk.83},level=5]{4. ปจฺจยปจฺจนียานุโลม}
\csromantocmarkref{subparagraph}{นเหตุทุก}{mtk.84}
\bookmark[dest={mtk.84},level=5]{นเหตุทุก}
\proseitemwithcloserruled{91}{\csromanfolio{342}เหตุปจฺจยา นอารมฺมเณ สตฺต, นอธิปติยา, นอนนฺตเร, นสมนนฺตเร สตฺต, นอญฺญมญฺเญ ตีณิ, นอุปนิสฺสเย สตฺต, นปุเรชาเต สตฺต ฯเปฯ นมคฺเค สตฺต, นสมฺปยุตฺเต ตีณิ, นวิปฺปยุตฺเต ตีณิ, โนนตฺถิยา สตฺต, โนวิคเต สตฺต. (เอวํ คเณตพฺพํ.)}{อนุโลม\-ปจฺจนียํ.}
\proseitem{92}{นเหตุปจฺจยา อารมฺมเณ สตฺต, อธิปติยา สตฺต, อนนฺตเร นว, สมนนฺตเร นว, สหชาเต เอกาทส, อญฺญมญฺเญ สตฺต, นิสฺสเย เตรส, อุปนิสฺสเย นว, ปุเรชาเต ตีณิ, ปจฺฉาชาเต ตีณิ, อาเสวเน จตฺตาริ, กมฺเม สตฺต ฯเปฯ มคฺเค สตฺต, สมฺปยุตฺเต ตีณิ, วิปฺปยุตฺเต ปญฺจ, อตฺถิยา เตรส, นตฺถิยา นว, วิคเต นว, อวิคเต เตรส.}

\vspace{\tipitakaparskip}
\csromancenter{ปจฺจนียา\-นุโลมํ.}
\nitthitam{ปริตฺตตฺติกํ นิฏฺฐิตํ.}
\csromanmatikaanchor{mtk.85}
\csromantocmarknopage{subparagraph}{13. ปริตฺตารมฺมณตฺติก 1. ปฏิจฺจวาร}{mtk.85}
\bookmark[dest={mtk.85},level=5]{13. ปริตฺตารมฺมณตฺติก 1. ปฏิจฺจวาร}
\csromanmarkh{13. ปริตฺตา\-รมฺมณ\-ตฺติก}\csromanmarkhii{13. ปริตฺตา\-รมฺมณ\-ตฺติก 1. ปฏิจฺจวาร}\csromanheadernomark{6}{13. ปริตฺตา\-รมฺมณ\-ตฺติก 13. ปริตฺตา\-รมฺมณ\-ตฺติก 1. ปฏิจฺจวาร}
\csromanmarkh{1. ปจฺจยานุโลม}\csromanmarkhii{1. วิภงฺควาร}\csromanheadernomark{2}{1. \textbf{ปจฺจยานุโลม 1. วิภงฺควาร}}
\csromanmatikaanchor{mtk.86}
\csromanmatikaanchor{mtk.87}
\csromantocmarknopage{subparagraph}{1. ปจฺจยานุโลม 1. วิภงฺควาร}{mtk.86}
\bookmark[dest={mtk.86},level=5]{1. ปจฺจยานุโลม 1. วิภงฺควาร}
\csromantocmarkref{subparagraph}{เหตุ}{mtk.87}
\bookmark[dest={mtk.87},level=5]{เหตุ}
\csromanheader{6}{\textbf{เหตุ}}
\proseitem{1}{\csromanfolio{343}ปริตฺตา\-รมฺมณํ ธมฺมํ ปฏิจฺจ ปริตฺตารมฺมโณ ธมฺโม อุปฺปชฺชติ เหตุปจฺจยา. ปริตฺตา\-รมฺมณํ เอกํ ขนฺธํ ปฏิจฺจ ตโย ขนฺธา ฯเปฯ เทฺว ขนฺธา. ปฏิสนฺธิ\-กฺขเณ ปริตฺตา\-รมฺมณํ เอกํ ขนฺธํ ปฏิจฺจ ตโย ขนฺธา ฯเปฯ เทฺว ขนฺธา. (1)}

\prose{มหคฺคตา\-รมฺมณํ ธมฺมํ ปฏิจฺจ มหคฺคตา\-รมฺมโณ ธมฺโม อุปฺปชฺชติ เหตุปจฺจยา. มหคฺคตา\-รมฺมณํ เอกํ ขนฺธํ ปฏิจฺจ ตโย ขนฺธา ฯเปฯ เทฺว ขนฺธา. ปฏิสนฺธิ\-กฺขเณ มหคฺคตา\-รมฺมณํ ฯเปฯ. (1)}

\prose{อปฺปมาณา\-รมฺมณํ ธมฺมํ ปฏิจฺจ อปฺปมาณารมฺมโณ ธมฺโม อุปฺปชฺชติ เหตุปจฺจยา. อปฺปมาณา\-รมฺมณํ เอกํ ขนฺธํ ปฏิจฺจ ตโย ขนฺธา ฯเปฯ เทฺว ขนฺธา. (1)}

\csromanmatikaanchor{mtk.88}
\csromantocmarkref{subparagraph}{อารมฺมณาทิ}{mtk.88}
\bookmark[dest={mtk.88},level=5]{อารมฺมณาทิ}
\csromanheader{6}{\textbf{อารมฺมณาทิ}}
\proseitem{2}{ปริตฺตา\-รมฺมณํ ธมฺมํ ปฏิจฺจ ปริตฺตารมฺมโณ ธมฺโม อุปฺปชฺชติ อารมฺมณ\-ปจฺจยา. อธิปติ\-ปจฺจยา. (สํขิตฺตํ.) อวิคตปจฺจยา.}

\csromanmatikaanchor{mtk.89}
\csromantocmarkref{subparagraph}{1. ปจฺจยานุโลม 2. สงฺขฺยาวาร}{mtk.89}
\bookmark[dest={mtk.89},level=5]{1. ปจฺจยานุโลม 2. สงฺขฺยาวาร}
\csromanmarkh{1. ปจฺจยานุโลม}\csromanmarkhii{2. สงฺขฺยาวาร}\csromanheadernomark{6}{1. \textbf{ปจฺจยานุโลม 2. สงฺขฺยาวาร}}
\proseitemwithcloserruled{3}{\textbf{เหตุ}ยา ตีณิ, อารมฺมเณ ตีณิ, อธิปติยา ตีณิ ฯเปฯ อวิคเต ตีณิ. (เอวํ คเณตพฺพํ.)}{อนุโลมํ.}
\csromanmatikaanchor{mtk.90}
\csromantocmarkref{subparagraph}{2. ปจฺจยปจฺจนีย 1. วิภงฺควาร}{mtk.90}
\bookmark[dest={mtk.90},level=5]{2. ปจฺจยปจฺจนีย 1. วิภงฺควาร}
\csromanmarkh{2. ปจฺจย\-ปจฺจนีย}\csromanmarkhii{1. วิภงฺควาร}\csromanheadernomark{6}{2. ปจฺจย\-ปจฺจนีย 1. วิภงฺควาร}
\csromanheader{2}{\textbf{นเหตุ}}
\proseitem{4}{ปริตฺตา\-รมฺมณํ ธมฺมํ ปฏิจฺจ ปริตฺตารมฺมโณ ธมฺโม อุปฺปชฺชติ นเหตุปจฺจยา. อเหตุกํ ปริตฺตา\-รมฺมณํ เอกํ ขนฺธํ ปฏิจฺจ ตโย ขนฺธา ฯเปฯ. \csromanfolio{344}เทฺว ขนฺธา. อเหตุก\-ปฏิสนฺธิ\-กฺขเณ ปริตฺตา\-รมฺมณํ เอกํ ขนฺธํ ปฏิจฺจ ตโย ขนฺธา ฯเปฯ เทฺว ขนฺธา. วิจิกิจฺฉา\-สหคเต อุทฺธจฺจ\-สหคเต ขนฺเธ ปฏิจฺจ วิจิกิจฺฉา\-สหคโต อุทฺธจฺจ\-สหคโต โมโห. (1)}

\prose{มหคฺคตา\-รมฺมณํ ธมฺมํ ปฏิจฺจ มหคฺคตา\-รมฺมโณ ธมฺโม อุปฺปชฺชติ นเหตุปจฺจยา. อเหตุกํ มหคฺคตา\-รมฺมณํ เอกํ ขนฺธํ ปฏิจฺจ ตโย ขนฺธา ฯเปฯ เทฺว ขนฺธา. วิจิกิจฺฉา\-สหคเต อุทฺธจฺจ\-สหคเต ขนฺเธ ปฏิจฺจ วิจิกิจฺฉา\-สหคโต อุทฺธจฺจ\-สหคโต โมโห. (1)}

\prose{อปฺปมาณา\-รมฺมณํ ธมฺมํ ปฏิจฺจ อปฺปมาณารมฺมโณ ธมฺโม อุปฺปชฺชติ นเหตุปจฺจยา. อเหตุกํ อปฺปมาณา\-รมฺมณํ เอกํ ขนฺธํ ปฏิจฺจ ตโย ขนฺธา ฯเปฯ เทฺว ขนฺธา. (1)}

\prose{นอธิปติ}

\proseitem{5}{ปริตฺตา\-รมฺมณํ ธมฺมํ ปฏิจฺจ ปริตฺตารมฺมโณ ธมฺโม อุปฺปชฺชติ นอธิปติ\-ปจฺจยา. ปริตฺตา\-รมฺมณํ เอกํ ขนฺธํ ปฏิจฺจ ตโย ขนฺธา ฯเปฯ เทฺว ขนฺธา. ปฏิสนฺธิ\-กฺขเณ ฯเปฯ. (1)}

\prose{มหคฺคตา\-รมฺมณํ ธมฺมํ ปฏิจฺจ มหคฺคตา\-รมฺมโณ ธมฺโม อุปฺปชฺชติ นอธิปติ\-ปจฺจยา. มหคฺคตา\-รมฺมณํ เอกํ ขนฺธํ ปฏิจฺจ ตโย ขนฺธา ฯเปฯ เทฺว ขนฺธา. ปฏิสนฺธิ\-กฺขเณ ฯเปฯ. (1)}

\prose{อปฺปมาณา\-รมฺมณํ ธมฺมํ ปฏิจฺจ อปฺปมาณารมฺมโณ ธมฺโม อุปฺปชฺชติ นอธิปติ\-ปจฺจยา. อปฺปมาณา\-รมฺมณํ เอกํ ขนฺธํ ปฏิจฺจ ตโย ขนฺธา ฯเปฯ เทฺว ขนฺธา. (1)}

\csromanheader{2}{\textbf{นปุเรชาตาทิ}}
\proseitem{6}{ปริตฺตา\-รมฺมณํ ธมฺมํ ปฏิจฺจ ปริตฺตารมฺมโณ ธมฺโม อุปฺปชฺชติ นปุเรชาต\-ปจฺจยา. อรูเป ปริตฺตา\-รมฺมณํ เอกํ ขนฺธํ ปฏิจฺจ ตโย ขนฺธา ฯเปฯ. ปฏิสนฺธิ\-กฺขเณ ฯเปฯ. (1)}

\prose{มหคฺคตา\-รมฺมณํ ธมฺมํ ปฏิจฺจ มหคฺคตา\-รมฺมโณ ธมฺโม อุปฺปชฺชติ นปุเรชาต\-ปจฺจยา. อรูเป มหคฺคตา\-รมฺมณํ เอกํ ขนฺธํ ปฏิจฺจ ตโย ขนฺธา ฯเปฯ เทฺว ขนฺธา. (มหคฺคตา\-รมฺมเณ ปฏิสนฺธิ นตฺถิ.) (1)}

\prose{\csromanfolio{345}อปฺปมาณา\-รมฺมณํ ธมฺมํ ปฏิจฺจ อปฺปมาณารมฺมโณ ธมฺโม อุปฺปชฺชติ นปุเรชาต\-ปจฺจยา. อรูเป อปฺปมาณา\-รมฺมณํ เอกํ ขนฺธํ ปฏิจฺจ ตโย ขนฺธา ฯเปฯ เทฺว ขนฺธา.}

\vspace{\tipitakaparskip}
\csromancenter{(นปจฺฉาชาต\-ปจฺจยญฺจ นอาเสวน\-ปจฺจยญฺจ นอธิปติ\-สทิสํ.)}
\csromanheader{2}{\textbf{นกมฺม}}
\proseitem{7}{ปริตฺตา\-รมฺมณํ ธมฺมํ ปฏิจฺจ ปริตฺตารมฺมโณ ธมฺโม อุปฺปชฺชติ นกมฺมปจฺจยา. ปริตฺตารมฺมเณ ขนฺเธ ปฏิจฺจ ปริตฺตารมฺมณา เจตนา. (1)}

\prose{มหคฺคตา\-รมฺมณํ ธมฺมํ ปฏิจฺจ มหคฺคตา\-รมฺมโณ ธมฺโม อุปฺปชฺชติ นกมฺมปจฺจยา. มหคฺคตา\-รมฺมเณ ขนฺเธ ปฏิจฺจ มหคฺคตา\-รมฺมณา เจตนา. (1)}

\prose{อปฺปมาณา\-รมฺมณํ ธมฺมํ ปฏิจฺจ อปฺปมาณารมฺมโณ ธมฺโม อุปฺปชฺชติ นกมฺมปจฺจยา. อปฺปมาณา\-รมฺมเณ ขนฺเธ ปฏิจฺจ อปฺปมาณา\-รมฺมณา เจตนา. (1)}

\csromanheader{2}{\textbf{นวิปากาทิ}}
\proseitem{8}{ปริตฺตา\-รมฺมณํ ธมฺมํ ปฏิจฺจ ปริตฺตารมฺมโณ ธมฺโม อุปฺปชฺชติ นวิปาก\-ปจฺจยา. (ปฏิสนฺธิ นตฺถิ.) นฌานปจฺจยา. ปญฺจ\-วิญฺญาณสหคตํ เอกํ ขนฺธํ ปฏิจฺจ ตโย ขนฺธา ฯเปฯ เทฺว ขนฺธา. นมคฺคปจฺจยา. อเหตุกํ ปริตฺตา\-รมฺมณํ เอกํ ขนฺธํ ปฏิจฺจ ตโย ขนฺธา ฯเปฯ เทฺว ขนฺธา. อเหตุก\-ปฏิสนฺธิ\-กฺขเณ ฯเปฯ เทฺว ขนฺธา ฯเปฯ.}

\prose{มหคฺคตา\-รมฺมณํ ธมฺมํ ปฏิจฺจ มหคฺคตา\-รมฺมโณ ธมฺโม อุปฺปชฺชติ นมคฺคปจฺจยา. อเหตุกํ มหคฺคตา\-รมฺมณํ เอกํ ขนฺธํ ปฏิจฺจ ตโย ขนฺธา ฯเปฯ เทฺว ขนฺธา. (1)}

\prose{อปฺปมาณา\-รมฺมณํ ธมฺมํ ปฏิจฺจ อปฺปมาณารมฺมโณ ธมฺโม อุปฺปชฺชติ นมคฺคปจฺจยา. อเหตุกํ อปฺปมาณา\-รมฺมณํ เอกํ ขนฺธํ ปฏิจฺจ ตโย ขนฺธา ฯเปฯ เทฺว ขนฺธา. (1)}

\csromanheader{2}{\textbf{นวิปฺปยุตฺต}}
\proseitem{9}{ปริตฺตา\-รมฺมณํ ธมฺมํ ปฏิจฺจ ปริตฺตารมฺมโณ ธมฺโม อุปฺปชฺชติ นวิปฺปยุตฺต\-ปจฺจยา. อรูเป ปริตฺตา\-รมฺมณํ เอกํ ขนฺธํ ปฏิจฺจ ตโย ขนฺธา ฯเปฯ. (1)}

\prose{มหคฺคตา\-รมฺมณํ ธมฺมํ ปฏิจฺจ มหคฺคตา\-รมฺมโณ ธมฺโม อุปฺปชฺชติ นวิปฺปยุตฺต\-ปจฺจยา. อรูเป มหคฺคตา\-รมฺมณํ เอกํ ขนฺธํ ปฏิจฺจ ตโย ขนฺธา ฯเปฯ. (1)}

\csromanmatikaanchor{mtk.91}
\csromanmatikaanchor{mtk.92}
\csromanmatikaanchor{mtk.93}
\csromanmatikaanchor{mtk.94}
\csromantocmarkref{subparagraph}{2. ปจฺจยปจฺจนีย 2. สงฺขฺยาวาร}{mtk.91}
\bookmark[dest={mtk.91},level=5]{2. ปจฺจยปจฺจนีย 2. สงฺขฺยาวาร}
\csromantocmarkref{subparagraph}{3. ปจฺจยานุโลมปจฺจนีย}{mtk.92}
\bookmark[dest={mtk.92},level=5]{3. ปจฺจยานุโลมปจฺจนีย}
\csromantocmarknopage{subparagraph}{4. ปจฺจยปจฺจนียานุโลม}{mtk.93}
\bookmark[dest={mtk.93},level=5]{4. ปจฺจยปจฺจนียานุโลม}
\csromantocmarkref{subparagraph}{นเหตุทุก}{mtk.94}
\bookmark[dest={mtk.94},level=5]{นเหตุทุก}
\prose{\csromanfolio{346}อปฺปมาณา\-รมฺมณํ ธมฺมํ ปฏิจฺจ อปฺปมาณารมฺมโณ ธมฺโม อุปฺปชฺชติ นวิปฺปยุตฺต\-ปจฺจยา. อรูเป อปฺปมาณา\-รมฺมณํ เอกํ ขนฺธํ ปฏิจฺจ ตโย ขนฺธา ฯเปฯ เทฺว ขนฺธา. (1)}

\csromanmarkh{2. ปจฺจย\-ปจฺจนีย}\csromanmarkhii{2. สงฺขฺยาวาร}\csromanheadernomark{6}{2. ปจฺจย\-ปจฺจนีย 2. สงฺขฺยาวาร}
\proseitemwithcloserruled{10}{นเหตุยา ตีณิ, นอธิปติยา ตีณิ ฯเปฯ นปุเรชาเต ตีณิ, นปจฺฉาชาเต ตีณิ, นอาเสวเน ตีณิ, นกมฺเม ตีณิ, นวิปาเก ตีณิ, นฌาเน เอกํ, นมคฺเค ตีณิ, นวิปฺปยุตฺเต ตีณิ. (เอวํ คเณตพฺพํ.)}{ปจฺจนียํ.}
\proseitemwithcloserruled{11}{เหตุปจฺจยา นอธิปติยา ตีณิ, นปุเรชาเต ตีณิ, นปจฺฉาชาเต ตีณิ, นอาเสวเน ตีณิ, นกมฺเม ตีณิ, นวิปาเก ตีณิ, นวิปฺปยุตฺเต ตีณิ. (เอวํ คเณตพฺพํ.)}{อนุโลม\-ปจฺจนียํ.}
\proseitem{12}{นเหตุปจฺจยา อารมฺมเณ ตีณิ, อนนฺตเร ตีณิ, สมนนฺตเร ตีณิ, สหชาเต ตีณิ, อญฺญมญฺเญ ตีณิ, นิสฺสเย ตีณิ, อุปนิสฺสเย ตีณิ, ปุเรชาเต ตีณิ, อาเสวเน เทฺว, กมฺเม ตีณิ, วิปาเก เอกํ, อาหาเร ตีณิ, อินฺทฺริเย ตีณิ, ฌาเน ตีณิ, มคฺเค เทฺว, สมฺปยุตฺเต ตีณิ, วิปฺปยุตฺเต ตีณิ, อตฺถิยา ตีณิ, นตฺถิยา ตีณิ, วิคเต ตีณิ, อวิคเต ตีณิ. (เอวํ คเณตพฺพํ.)}

\vspace{\tipitakaparskip}
\csromancenter{ปจฺจนียา\-นุโลมํ.}
\csromancenter{ปฏิจฺจวาโร.}
\prose{(สหชาตวาโรปิ ปจฺจยวาโรปิ นิสฺสยวาโรปิ สํสฏฺฐ\-วาโรปิ สมฺปยุตฺตวาโรปิ ปฏิจฺจ\-วาร\-สทิโส.)}

\csromanmatikaanchor{mtk.95}
\csromantocmarknopage{subparagraph}{13. ปริตฺตารมฺมณตฺติก 7. ปญฺหาวาร}{mtk.95}
\bookmark[dest={mtk.95},level=5]{13. ปริตฺตารมฺมณตฺติก 7. ปญฺหาวาร}
\csromanmarkh{13. ปริตฺตา\-รมฺมณ\-ตฺติก}\csromanmarkhii{13. ปริตฺตา\-รมฺมณ\-ตฺติก 7. ปญฺหาวาร}\csromanheadernomark{6}{13. ปริตฺตา\-รมฺมณ\-ตฺติก 13. ปริตฺตา\-รมฺมณ\-ตฺติก 7. ปญฺหาวาร}
\csromanmarkh{1. ปจฺจยานุโลม}\csromanmarkhii{1. วิภงฺควาร}\csromanheadernomark{2}{1. \textbf{ปจฺจยานุโลม 1. วิภงฺควาร}}
\csromanmatikaanchor{mtk.96}
\csromanmatikaanchor{mtk.97}
\csromantocmarknopage{subparagraph}{1. ปจฺจยานุโลม 1. วิภงฺควาร}{mtk.96}
\bookmark[dest={mtk.96},level=5]{1. ปจฺจยานุโลม 1. วิภงฺควาร}
\csromantocmarkref{subparagraph}{เหตุ}{mtk.97}
\bookmark[dest={mtk.97},level=5]{เหตุ}
\csromanheader{6}{\textbf{เหตุ}}
\proseitem{13}{\csromanfolio{347}ปริตฺตารมฺมโณ ธมฺโม ปริตฺตา\-รมฺมณสฺส ธมฺมสฺส เหตุปจฺจเยน ปจฺจโย. ปริตฺตารมฺมณา เหตู สมฺปยุตฺตกานํ ขนฺธานํ เหตุปจฺจเยน ปจฺจโย. ปฏิสนฺธิ\-กฺขเณ ปริตฺตารมฺมณา เหตู สมฺปยุตฺตกานํ ขนฺธานํ เหตุปจฺจเยน ปจฺจโย. (1)}

\prose{มหคฺคตา\-รมฺมโณ ธมฺโม มหคฺคตา\-รมฺมณสฺส ธมฺมสฺส เหตุปจฺจเยน ปจฺจโย. มหคฺคตา\-รมฺมณา เหตู สมฺปยุตฺตกานํ ขนฺธานํ เหตุปจฺจเยน ปจฺจโย. ปฏิสนฺธิ\-กฺขเณ ฯเปฯ. (1)}

\prose{อปฺปมาณารมฺมโณ ธมฺโม อปฺปมาณา\-รมฺมณสฺส ธมฺมสฺส เหตุปจฺจเยน ปจฺจโย. อปฺปมาณา\-รมฺมณา เหตู สมฺปยุตฺตกานํ ขนฺธานํ เหตุปจฺจเยน ปจฺจโย. (1)}

\csromanmatikaanchor{mtk.98}
\csromantocmarkref{subparagraph}{อารมฺมณ}{mtk.98}
\bookmark[dest={mtk.98},level=5]{อารมฺมณ}
\csromanheader{6}{\textbf{อารมฺมณ}}
\proseitem{14}{ปริตฺตารมฺมโณ ธมฺโม ปริตฺตา\-รมฺมณสฺส ธมฺมสฺส อารมฺมณ\-ปจฺจเยน ปจฺจโย. ทานํ ทตฺวา, สีลํ สมาทิยิตฺวา, อุโปสถกมฺมํ กตฺวา ตํ ปจฺจเวกฺขติ, ปุพฺเพ สุจิณฺณานิ ปจฺจเวกฺขติ. อริยา ปริตฺตารมฺมเณ ปหีเน กิเลเส ปจฺจเวกฺขนฺติ, วิกฺขมฺภิเต กิเลเส ปจฺจเวกฺขนฺติ, ปุพฺเพ สมุทาจิณฺเณ กิเลเส ชานนฺติ. ปริตฺตารมฺมเณ ปริตฺเต ขนฺเธ อนิจฺจโต ทุกฺขโต อนตฺตโต วิปสฺสติ, อสฺสาเทติ อภินนฺทติ, ตํ อารพฺภ ปริตฺตารมฺมโณ ราโค อุปฺปชฺชติ ฯเปฯ โทมนสฺสํ อุปฺปชฺชติ. เจโตปริย\-ญาเณน ปริตฺตา\-รมฺมณ\-ปริตฺต\-จิตฺต\-สมงฺคิสฺส จิตฺตํ ชานาติ. ปริตฺตารมฺมณา ปริตฺตา ขนฺธา เจโตปริย\-ญาณสฺส, ปุพฺเพ\-นิวาสา\-นุสฺสติ\-ญาณสฺส, ยถากมฺมูปคญาณสฺส, อนาคตํสญาณสฺส, อาวชฺชนาย อารมฺมณ\-ปจฺจเยน ปจฺจโย. (1)}

\prose{ปริตฺตารมฺมโณ ธมฺโม มหคฺคตา\-รมฺมณสฺส ธมฺมสฺส อารมฺมณ\-ปจฺจเยน ปจฺจโย. ทิพฺพํ จกฺขุํ ปจฺจเวกฺขติ. ทิพฺพํ โสตธาตุํ ปจฺจเวกฺขติ. \csromanfolio{348}ปริตฺตา\-รมฺมณํ อิทฺธิวิธ\-ญาณํ ปจฺจเวกฺขติ. เจโตปริย\-ญาณํ ฯเปฯ. ปุพฺเพ\-นิวาสา\-นุสฺสติ\-ญาณํ ฯเปฯ. ยถากมฺมูปคญาณํ ฯเปฯ. อนาคตํสญาณํ ปจฺจเวกฺขติ. ปริตฺตารมฺมเณ มหคฺคเต ขนฺเธ อนิจฺจโต ฯเปฯ วิปสฺสติ, อสฺสาเทติ อภินนฺทติ, ตํ อารพฺภ มหคฺคตา\-รมฺมโณ ราโค อุปฺปชฺชติ ฯเปฯ โทมนสฺสํ อุปฺปชฺชติ. เจโตปริย\-ญาเณน ปริตฺตา\-รมฺมณ\-มหคฺคต\-จิตฺต\-สมงฺคิสฺส จิตฺตํ ชานาติ. ปริตฺตารมฺมณา มหคฺคตา ขนฺธา เจโตปริย\-ญาณสฺส, ปุพฺเพ\-นิวาสา\-นุสฺสติ\-ญาณสฺส, อนาคตํสญาณสฺส, อาวชฺชนาย อารมฺมณ\-ปจฺจเยน ปจฺจโย. (2)}

\proseitem{15}{มหคฺคตา\-รมฺมโณ ธมฺโม มหคฺคตา\-รมฺมณสฺส ธมฺมสฺส อารมฺมณ\-ปจฺจเยน ปจฺจโย. วิญฺญาณญฺจา\-ยตนํ ปจฺจเวกฺขติ. เนว\-สญฺญา\-นา\-สญฺญา\-ยตนํ ปจฺจเวกฺขติ. มหคฺคตา\-รมฺมณํ อิทฺธิวิธ\-ญาณํ ปจฺจเวกฺขติ. เจโตปริย\-ญาณํ ฯเปฯ. ปุพฺเพ\-นิวาสา\-นุสฺสติ\-ญาณํ ฯเปฯ. ยถากมฺมูปคญาณํ ฯเปฯ. อนาคตํสญาณํ ปจฺจเวกฺขติ. มหคฺคตา\-รมฺมเณ มหคฺคเต ขนฺเธ อนิจฺจโต ฯเปฯ วิปสฺสติ, อสฺสาเทติ อภินนฺทติ, ตํ อารพฺภ มหคฺคตา\-รมฺมโณ ราโค อุปฺปชฺชติ ฯเปฯ โทมนสฺสํ อุปฺปชฺชติ. เจโตปริย\-ญาเณน มหคฺคตา\-รมฺมณ\-มหคฺคต\-จิตฺต\-สมงฺคิสฺส จิตฺตํ ชานาติ. มหคฺคตา\-รมฺมณา มหคฺคตา ขนฺธา เจโตปริย\-ญาณสฺส, ปุพฺเพ\-นิวาสา\-นุสฺสติ\-ญาณสฺส, ยถากมฺมูปคญาณสฺส, อนาคตํสญาณสฺส, อาวชฺชนาย อารมฺมณ\-ปจฺจเยน ปจฺจโย. (1)}

\csromanmatikaanchor{mtk.99}
\csromantocmarkref{subparagraph}{อธิปติ}{mtk.99}
\bookmark[dest={mtk.99},level=5]{อธิปติ}
\prose{มหคฺคตา\-รมฺมโณ ธมฺโม ปริตฺตา\-รมฺมณสฺส ธมฺมสฺส อารมฺมณ\-ปจฺจเยน ปจฺจโย. ปฐมชฺฌาน\-ปจฺจเวกฺขณํ ปจฺจเวกฺขติ ฯเปฯ. เนว\-สญฺญา\-นา\-สญฺญา\-ยตนปจฺจเวกฺขณํ ปจฺจเวกฺขติ. ทิพฺพจกฺขุ\-ปจฺจเวกฺขณํ ปจฺจเวกฺขติ. ทิพฺพโสตธาตุ\-ปจฺจเวกฺขณํ ปจฺจเวกฺขติ. อิทฺธิวิธญาณ\-ปจฺจเวกฺขณํ ฯเปฯ. เจโตปริยญาณ\-ปจฺจเวกฺขณํ ฯเปฯ. ปุพฺเพ\-นิวาสา\-นุสฺสติ\-ญาณปจฺจเวกฺขณํ ฯเปฯ. ยถากมฺมูปคญาณ\-ปจฺจเวกฺขณํ ฯเปฯ. อนาคตํสญาณ\-ปจฺจเวกฺขณํ ปจฺจเวกฺขติ. อริยา มหคฺคตา\-รมฺมเณ ปหีเน กิเลเส ปจฺจเวกฺขนฺติ, วิกฺขมฺภิเต กิเลเส ปจฺจเวกฺขนฺติ, ปุพฺเพ สมุทาจิณฺเณ กิเลเส ชานนฺติ. มหคฺคตา\-รมฺมเณ ปริตฺเต ขนฺเธ อนิจฺจโต ฯเปฯ วิปสฺสติ, อสฺสาเทติ อภินนฺทติ, ตํ อารพฺภ ปริตฺตารมฺมโณ ราโค อุปฺปชฺชติ ฯเปฯ โทมนสฺสํ อุปฺปชฺชติ. เจโตปริย\-ญาเณน มหคฺคตา\-รมฺมณ\-ปริตฺต\-จิตฺต\-สมงฺคิสฺส จิตฺตํ ชานาติ. มหคฺคตา\-รมฺมณา ปริตฺตา ขนฺธา เจโตปริย\-ญาณสฺส, ปุพฺเพ\-นิวาสา\-นุสฺสติ\-ญาณสฺส, \csromanfolio{349}ยถากมฺมูปคญาณสฺส, อนาคตํสญาณสฺส, อาวชฺชนาย อารมฺมณ\-ปจฺจเยน ปจฺจโย. (2)}

\prose{อปฺปมาณารมฺมโณ ธมฺโม อปฺปมาณา\-รมฺมณสฺส ธมฺมสฺส อารมฺมณ\-ปจฺจเยน ปจฺจโย. อริยา มคฺคา วุฏฺฐหิตฺวา มคฺคํ ปจฺจเวกฺขนฺติ, ผลํ ปจฺจเวกฺขนฺติ. เจโตปริย\-ญาเณน อปฺปมาณารมฺมณ\-อปฺปมาณจิตฺต\-สมงฺคิสฺส จิตฺตํ ชานาติ. อปฺปมาณา\-รมฺมณา อปฺปมาณา ขนฺธา เจโตปริย\-ญาณสฺส, ปุพฺเพ\-นิวาสา\-นุสฺสติ\-ญาณสฺส, อนาคตํสญาณสฺส, อาวชฺชนาย อารมฺมณ\-ปจฺจเยน ปจฺจโย. (1)}

\prose{อปฺปมาณารมฺมโณ ธมฺโม ปริตฺตา\-รมฺมณสฺส ธมฺมสฺส อารมฺมณ\-ปจฺจเยน ปจฺจโย. อริยา โคตฺรภุํ ปจฺจเวกฺขนฺติ, โวทานํ ปจฺจเวกฺขนฺติ, มคฺค\-ปจฺจเวกฺขณํ ปจฺจเวกฺขนฺติ, ผลปจฺจเวกฺขณํ ปจฺจเวกฺขนฺติ, นิพฺพาน\-ปจฺจเวกฺขณํ ปจฺจเวกฺขนฺติ. อปฺปมาณา\-รมฺมเณ ปริตฺเต ขนฺเธ อนิจฺจโต ฯเปฯ วิปสฺสติ. เจโตปริย\-ญาเณน อปฺปมาณารมฺมณ\-ปริตฺต\-จิตฺต\-สมงฺคิสฺส จิตฺตํ ชานาติ. อปฺปมาณา\-รมฺมณา ปริตฺตา ขนฺธา เจโตปริย\-ญาณสฺส, ปุพฺเพ\-นิวาสา\-นุสฺสติ\-ญาณสฺส, ยถากมฺมูปคญาณสฺส, อนาคตํสญาณสฺส, อาวชฺชนาย อารมฺมณ\-ปจฺจเยน ปจฺจโย. (2)}

\prose{อปฺปมาณารมฺมโณ ธมฺโม มหคฺคตา\-รมฺมณสฺส ธมฺมสฺส อารมฺมณ\-ปจฺจเยน ปจฺจโย. อริยา อปฺปมาณา\-รมฺมณํ เจโตปริย\-ญาณํ ปจฺจเวกฺขนฺติ, ปุพฺเพ\-นิวาสา\-นุสฺสติ\-ญาณํ ปจฺจเวกฺขนฺติ, อนาคตํสญาณํ ปจฺจเวกฺขนฺติ. เจโตปริย\-ญาเณน อปฺปมาณารมฺมณ\-มหคฺคต\-จิตฺต\-สมงฺคิสฺส จิตฺตํ ชานนฺติ. อปฺปมาณา\-รมฺมณํ มหคฺคตา ขนฺธา เจโตปริย\-ญาณสฺส, ปุพฺเพ\-นิวาสา\-นุสฺสติ\-ญาณสฺส, อนาคตํสญาณสฺส, อาวชฺชนาย อารมฺมณ\-ปจฺจเยน ปจฺจโย. (3)}

\prose{ปริตฺตารมฺมโณ ธมฺโม ปริตฺตา\-รมฺมณสฺส ธมฺมสฺส อธิปติ\-ปจฺจเยน ปจฺจโย. อารมฺมณา\-ธิปติ, สหชาตา\-ธิปติ.  อารมฺมณา\-ธิปติ\csromandash{}ทานํ ทตฺวา, สีลํ สมาทิยิตฺวา, อุโปสถกมฺมํ กตฺวา ตํ ครุํ กตฺวา ปจฺจเวกฺขติ, ปุพฺเพ สุจิณฺณานิ ครุํ กตฺวา ปจฺจเวกฺขติ, ปริตฺตารมฺมเณ \csromanfolio{350}ปริตฺเต ขนฺเธ ครุํ กตฺวา อสฺสาเทติ อภินนฺทติ, ตํ ครุํ กตฺวา ปริตฺตารมฺมโณ ราโค อุปฺปชฺชติ, ทิฏฺฐิ อุปฺปชฺชติ. สหชาตา\-ธิปติ\csromandash{} ปริตฺตา\-รมฺมณา\-ธิปติ สมฺปยุตฺตกานํ ขนฺธานํ อธิปติ\-ปจฺจเยน ปจฺจโย. (1)}

\proseitem{17}{ปริตฺตารมฺมโณ ธมฺโม มหคฺคตา\-รมฺมณสฺส ธมฺมสฺส อธิปติ\-ปจฺจเยน ปจฺจโย. อารมฺมณา\-ธิปติ\csromandash{}ทิพฺพํ จกฺขุํ ครุํ กตฺวา ปจฺจเวกฺขติ, ทิพฺพํ โสตธาตุํ ฯเปฯ. ปริตฺตา\-รมฺมณํ อิทฺธิวิธ\-ญาณํ ฯเปฯ. เจโตปริย\-ญาณํ ฯเปฯ. ปุพฺเพ\-นิวาสา\-นุสฺสติ\-ญาณํ ฯเปฯ. ยถากมฺมูปคญาณํ ฯเปฯ. อนาคตํสญาณํ ครุํ กตฺวา ปจฺจเวกฺขติ. ปริตฺตารมฺมเณ มหคฺคเต ขนฺเธ ครุํ กตฺวา อสฺสาเทติ อภินนฺทติ, ตํ ครุํ กตฺวา มหคฺคตา\-รมฺมโณ ราโค อุปฺปชฺชติ, ทิฏฺฐิ อุปฺปชฺชติ. (2)}

\prose{มหคฺคตา\-รมฺมโณ ธมฺโม มหคฺคตา\-รมฺมณสฺส ธมฺมสฺส อธิปติ\-ปจฺจเยน ปจฺจโย. อารมฺมณา\-ธิปติ, สหชาตา\-ธิปติ.  อารมฺมณา\-ธิปติ\csromandash{}วิญฺญาณญฺจา\-ยตนํ ครุํ กตฺวา ปจฺจเวกฺขติ, เนว\-สญฺญา\-นา\-สญฺญา\-ยตนํ ฯเปฯ. มหคฺคตา\-รมฺมณํ อิทฺธิวิธ\-ญาณํ ฯเปฯ. เจโตปริย\-ญาณํ ฯเปฯ. ปุพฺเพ\-นิวาสา\-นุสฺสติ\-ญาณํ ฯเปฯ. ยถากมฺมูปคญาณํ ฯเปฯ. อนาคตํสญาณํ ครุํ กตฺวา ปจฺจเวกฺขติ. มหคฺคตา\-รมฺมเณ มหคฺคเต ขนฺเธ ครุํ กตฺวา อสฺสาเทติ อภินนฺทติ, ตํ ครุํ กตฺวา มหคฺคตา\-รมฺมโณ ราโค อุปฺปชฺชติ ทิฏฺฐิ อุปฺปชฺชติ. สหชาตา\-ธิปติ\csromandash{}มหคฺคตา\-รมฺมณา\-ธิปติ สมฺปยุตฺตกานํ ขนฺธานํ อธิปติ\-ปจฺจเยน ปจฺจโย. (1)}

\prose{มหคฺคตา\-รมฺมโณ ธมฺโม ปริตฺตา\-รมฺมณสฺส ธมฺมสฺส อธิปติ\-ปจฺจเยน ปจฺจโย. อารมฺมณา\-ธิปติ\csromandash{} ปฐมชฺฌาน\-ปจฺจเวกฺขณํ ครุํ กตฺวา ปจฺจเวกฺขติ ฯเปฯ. อนาคตํสญาณ\-ปจฺจเวกฺขณํ ครุํ กตฺวา ปจฺจเวกฺขติ. มหคฺคตา\-รมฺมเณ ปริตฺเต ขนฺเธ ครุํ กตฺวา อสฺสาเทติ อภินนฺทติ, ตํ ครุํ กตฺวา ปริตฺตารมฺมโณ ราโค อุปฺปชฺชติ, ทิฏฺฐิ อุปฺปชฺชติ. (2)}

\prose{อปฺปมาณารมฺมโณ ธมฺโม อปฺปมาณา\-รมฺมณสฺส ธมฺมสฺส อธิปติ\-ปจฺจเยน ปจฺจโย. อารมฺมณา\-ธิปติ, สหชาตา\-ธิปติ.  อารมฺมณา\-ธิปติ\csromandash{}อริยา มคฺคา วุฏฺฐหิตฺวา มคฺคํ ครุํ กตฺวา ปจฺจเวกฺขนฺติ, ผลํ ครุํ กตฺวา ปจฺจเวกฺขนฺติ. สหชาตา\-ธิปติ\csromandash{} อปฺปมาณา\-รมฺมณาธิปติ สมฺปยุตฺตกานํ ขนฺธานํ อธิปติ\-ปจฺจเยน ปจฺจโย. (1)}

\csromanmatikaanchor{mtk.100}
\csromantocmarkref{subparagraph}{อนนฺตราทิ}{mtk.100}
\bookmark[dest={mtk.100},level=5]{อนนฺตราทิ}
\prose{\csromanfolio{351}อปฺปมาณารมฺมโณ ธมฺโม ปริตฺตา\-รมฺมณสฺส ธมฺมสฺส อธิปติ\-ปจฺจเยน ปจฺจโย. อารมฺมณา\-ธิปติ\csromandash{}เสกฺขา โคตฺรภุํ ครุํ กตฺวา ปจฺจเวกฺขนฺติ, โวทานํ ครุํ กตฺวา ปจฺจเวกฺขนฺติ, มคฺค\-ปจฺจเวกฺขณํ ครุํ กตฺวา ปจฺจเวกฺขนฺติ, ผลปจฺจเวกฺขณํ ครุํ กตฺวา ปจฺจเวกฺขนฺติ, นิพฺพาน\-ปจฺจเวกฺขณํ ครุํ กตฺวา ปจฺจเวกฺขนฺติ. (2)}

\prose{อปฺปมาณารมฺมโณ ธมฺโม มหคฺคตา\-รมฺมณสฺส ธมฺมสฺส อธิปติ\-ปจฺจเยน ปจฺจโย. อารมฺมณา\-ธิปติ\csromandash{}เสกฺขา อปฺปมาณา\-รมฺมณํ เจโตปริย\-ญาณํ ครุํ กตฺวา ปจฺจเวกฺขนฺติ. ปุพฺเพ\-นิวาสา\-นุสฺสติ\-ญาณํ ฯเปฯ. อนาคตํสญาณํ ครุํ กตฺวา ปจฺจเวกฺขนฺติ. (3)}

\csromanheader{2}{\textbf{อนนฺตร}}
\proseitem{20}{ปริตฺตารมฺมโณ ธมฺโม ปริตฺตา\-รมฺมณสฺส ธมฺมสฺส อนนฺตร\-ปจฺจเยน ปจฺจโย. ปุริมา ปุริมา ปริตฺตารมฺมณา ขนฺธา ปจฺฉิมานํ ปจฺฉิมานํ ปริตฺตา\-รมฺมณานํ ขนฺธานํ อนนฺตร\-ปจฺจเยน ปจฺจโย. (1)}

\prose{ปริตฺตารมฺมโณ ธมฺโม มหคฺคตา\-รมฺมณสฺส ธมฺมสฺส อนนฺตร\-ปจฺจเยน ปจฺจโย. ปริตฺตา\-รมฺมณํ จุติจิตฺตํ มหคฺคตา\-รมฺมณสฺส อุปปตฺติ\-จิตฺตสฺส อนนฺตร\-ปจฺจเยน ปจฺจโย. ปริตฺตา\-รมฺมณํ ภวงฺคํ มหคฺคตา\-รมฺมณาย อาวชฺชนาย อนนฺตร\-ปจฺจเยน ปจฺจโย. ปริตฺตารมฺมณา ขนฺธา มหคฺคตา\-รมฺมณสฺส วุฏฺฐานสฺส อนนฺตร\-ปจฺจเยน ปจฺจโย. (2)}

\prose{ปริตฺตารมฺมโณ ธมฺโม อปฺปมาณา\-รมฺมณสฺส ธมฺมสฺส อนนฺตร\-ปจฺจเยน ปจฺจโย. ปริตฺตา\-รมฺมณํ ภวงฺคํ อปฺปมาณา\-รมฺมณาย อาวชฺชนาย อนนฺตร\-ปจฺจเยน ปจฺจโย. ปริตฺตา\-รมฺมณํ อนุโลมํ โคตฺรภุสฺส. อนุโลมํ โวทานสฺส. อนุโลมํ ผลสมาปตฺติยา อนนฺตร\-ปจฺจเยน ปจฺจโย. (3)}

\proseitem{21}{มหคฺคตา\-รมฺมโณ ธมฺโม มหคฺคตา\-รมฺมณสฺส ธมฺมสฺส อนนฺตร\-ปจฺจเยน ปจฺจโย. ปุริมา ปุริมา มหคฺคตา\-รมฺมณา ขนฺธา ปจฺฉิมานํ ปจฺฉิมานํ มหคฺคตา\-รมฺมณานํ ขนฺธานํ อนนฺตร\-ปจฺจเยน ปจฺจโย. (1)}

\prose{มหคฺคตา\-รมฺมโณ ธมฺโม ปริตฺตา\-รมฺมณสฺส ธมฺมสฺส อนนฺตร\-ปจฺจเยน ปจฺจโย. มหคฺคตา\-รมฺมณํ จุติจิตฺตํ ปริตฺตา\-รมฺมณสฺส อุปปตฺติ\-จิตฺตสฺส \csromanfolio{352}อนนฺตร\-ปจฺจเยน ปจฺจโย. มหคฺคตา\-รมฺมณํ ภวงฺคํ ปริตฺตา\-รมฺมณาย อาวชฺชนาย อนนฺตร\-ปจฺจเยน ปจฺจโย. มหคฺคตา\-รมฺมณา ขนฺธา ปริตฺตา\-รมฺมณสฺส วุฏฺฐานสฺส อนนฺตร\-ปจฺจเยน ปจฺจโย. (2)}

\prose{มหคฺคตา\-รมฺมโณ ธมฺโม อปฺปมาณา\-รมฺมณสฺส ธมฺมสฺส อนนฺตร\-ปจฺจเยน ปจฺจโย. มหคฺคตา\-รมฺมณํ ภวงฺคํ อปฺปมาณา\-รมฺมณาย อาวชฺชนาย อนนฺตร\-ปจฺจเยน ปจฺจโย. มหคฺคตา\-รมฺมณํ อนุโลมํ โคตฺรภุสฺส. อนุโลมํ โวทานสฺส. อนุโลมํ ผลสมาปตฺติยา. นิโรธา วุฏฺฐหนฺตสฺส เนว\-สญฺญา\-นา\-สญฺญา\-ยตนํ ผลสมาปตฺติยา อนนฺตร\-ปจฺจเยน ปจฺจโย. (3)}

\prose{อปฺปมาณารมฺมโณ ธมฺโม อปฺปมาณา\-รมฺมณสฺส ธมฺมสฺส อนนฺตร\-ปจฺจเยน ปจฺจโย. ปุริมา ปุริมา อปฺปมาณา\-รมฺมณา ขนฺธา ปจฺฉิมานํ ปจฺฉิมานํ อปฺปมาณา\-รมฺมณานํ ขนฺธานํ อนนฺตร\-ปจฺจเยน ปจฺจโย. โคตฺรภุ มคฺคสฺส. โวทานํ มคฺคสฺส. มคฺโค ผลสฺส. ผลํ ผลสฺส อนนฺตร\-ปจฺจเยน ปจฺจโย. (1)}

\prose{อปฺปมาณารมฺมโณ ธมฺโม ปริตฺตา\-รมฺมณสฺส ธมฺมสฺส อนนฺตร\-ปจฺจเยน ปจฺจโย. มคฺค\-ปจฺจเวกฺขณํ ปริตฺตา\-รมฺมณสฺส วุฏฺฐานสฺส. ผลปจฺจเวกฺขณํ ปริตฺตา\-รมฺมณสฺส วุฏฺฐานสฺส. นิพฺพาน\-ปจฺจเวกฺขณํ ปริตฺตา\-รมฺมณสฺส วุฏฺฐานสฺส. อปฺปมาณา\-รมฺมณํ เจโตปริย\-ญาณํ ปริตฺตา\-รมฺมณสฺส วุฏฺฐานสฺส. ปุพฺเพ\-นิวาสา\-นุสฺสติ\-ญาณํ ปริตฺตา\-รมฺมณสฺส วุฏฺฐานสฺส. อนาคตํสญาณํ ปริตฺตา\-รมฺมณสฺส วุฏฺฐานสฺส. ผลํ ปริตฺตา\-รมฺมณสฺส วุฏฺฐานสฺส อนนฺตร\-ปจฺจเยน ปจฺจโย. (2)}

\prose{อปฺปมาณารมฺมโณ ธมฺโม มหคฺคตา\-รมฺมณสฺส ธมฺมสฺส อนนฺตร\-ปจฺจเยน ปจฺจโย. มคฺค\-ปจฺจเวกฺขณํ มหคฺคตา\-รมฺมณสฺส วุฏฺฐานสฺส. ผลปจฺจเวกฺขณํ มหคฺคตา\-รมฺมณสฺส วุฏฺฐานสฺส. นิพฺพาน\-ปจฺจเวกฺขณํ มหคฺคตา\-รมฺมณสฺส วุฏฺฐานสฺส. ผลํ มหคฺคตา\-รมฺมณสฺส วุฏฺฐานสฺส อนนฺตร\-ปจฺจเยน ปจฺจโย. (3)}

\csromanheader{2}{\textbf{สมนนฺตร}}
\proseitem{23}{ปริตฺตารมฺมโณ ธมฺโม ปริตฺตา\-รมฺมณสฺส ธมฺมสฺส สมนนฺตร\-ปจฺจเยน ปจฺจโย. (อนนฺตร\-สทิสํ.)}

\csromanheader{2}{13. ปริตฺตา\-รมฺมณ\-ตฺติก สหชาตาทิ}
\proseitem{24}{\csromanfolio{353}ปริตฺตารมฺมโณ ธมฺโม ปริตฺตา\-รมฺมณสฺส ธมฺมสฺส สหชาต\-ปจฺจเยน ปจฺจโย. อญฺญมญฺญ\-ปจฺจเยน ปจฺจโย. นิสฺสย\-ปจฺจเยน ปจฺจโย. ตีณิ. (ปฏิจฺจ\-วาร\-สทิสา กาตพฺพา.)}

\csromanheader{2}{\textbf{อุปนิสฺสย}}
\proseitem{25}{ปริตฺตารมฺมโณ ธมฺโม ปริตฺตา\-รมฺมณสฺส ธมฺมสฺส อุปนิสฺสย\-ปจฺจเยน ปจฺจโย. อารมฺมณู\-ปนิสฺสโย, อนนฺตรู\-ปนิสฺสโย, ปกตู\-ปนิสฺสโย ฯเปฯ. ปกตู\-ปนิสฺสโย\csromandash{}ปริตฺตา\-รมฺมณํ สทฺธํ อุปนิสฺสาย ทานํ เทติ. สีลํ ฯเปฯ อุโปสถกมฺมํ ฯเปฯ. ปริตฺตา\-รมฺมณํ ฌานํ อุปฺปาเทติ. วิปสฺสนํ ฯเปฯ อภิญฺญํ ฯเปฯ สมาปตฺติํ อุปฺปาเทติ. มานํ ชปฺเปติ. ทิฏฺฐิํ คณฺหาติ. ปริตฺตา\-รมฺมณํ สีลํ ฯเปฯ ปญฺญํ. ราคํ. โทสํ. โมหํ. มานํ. ทิฏฺฐิํ. ปตฺถนํ. กายิกํ สุขํ. กายิกํ ทุกฺขํ อุปนิสฺสาย ทานํ เทติ. สีลํ ฯเปฯ อุโปสถกมฺมํ ฯเปฯ. ปริตฺตา\-รมฺมณํ ฌานํ อุปฺปาเทติ. วิปสฺสนํ ฯเปฯ อภิญฺญํ ฯเปฯ สมาปตฺติํ อุปฺปาเทติ. ปาณํ หนติ ฯเปฯ สํฆํ ภินฺทติ. ปริตฺตารมฺมณา สทฺธา ฯเปฯ ปญฺญา. ราโค ฯเปฯ ปตฺถนา. กายิกํ สุขํ. กายิกํ ทุกฺขํ. ปริตฺตา\-รมฺมณาย สทฺธาย ฯเปฯ ปญฺญาย. ราคสฺส ฯเปฯ ปตฺถนาย. กายิกสฺส สุขสฺส, กายิกสฺส ทุกฺขสฺส อุปนิสฺสย\-ปจฺจเยน ปจฺจโย. (1)}

\prose{ปริตฺตารมฺมโณ ธมฺโม มหคฺคตา\-รมฺมณสฺส ธมฺมสฺส อุปนิสฺสย\-ปจฺจเยน ปจฺจโย. อารมฺมณู\-ปนิสฺสโย, อนนฺตรู\-ปนิสฺสโย, ปกตู\-ปนิสฺสโย ฯเปฯ. ปกตู\-ปนิสฺสโย\csromandash{}ปริตฺตา\-รมฺมณํ สทฺธํ อุปนิสฺสาย มหคฺคตา\-รมฺมณํ ฌานํ อุปฺปาเทติ. วิปสฺสนํ ฯเปฯ อภิญฺญํ ฯเปฯ สมาปตฺติํ อุปฺปาเทติ. มานํ ชปฺเปติ. ทิฏฺฐิํ คณฺหาติ. ปริตฺตา\-รมฺมณํ สีลํ ฯเปฯ ปญฺญํ. ราคํ ฯเปฯ ปตฺถนํ. กายิกํ สุขํ. กายิกํ ทุกฺขํ อุปนิสฺสาย มหคฺคตา\-รมฺมณํ ฌานํ อุปฺปาเทติ. วิปสฺสนํ ฯเปฯ อภิญฺญํ ฯเปฯ สมาปตฺติํ อุปฺปาเทติ. มานํ ชปฺเปติ. ทิฏฺฐิํ คณฺหาติ. ปริตฺตารมฺมณา สทฺธา ฯเปฯ. กายิกํ สุขํ. กายิกํ ทุกฺขํ. มหคฺคตา\-รมฺมณาย สทฺธาย ฯเปฯ ปญฺญาย. ราคสฺส ฯเปฯ ปตฺถนาย อุปนิสฺสย\-ปจฺจเยน ปจฺจโย. (2)}

\prose{ปริตฺตารมฺมโณ ธมฺโม อปฺปมาณา\-รมฺมณสฺส ธมฺมสฺส อุปนิสฺสย\-ปจฺจเยน ปจฺจโย. อนนฺตรู\-ปนิสฺสโย, ปกตู\-ปนิสฺสโย ฯเปฯ. ปกตู\-ปนิสฺสโย\csromandash{}ปริตฺตา\-รมฺมณํ สทฺธํ อุปนิสฺสาย อปฺปมาณา\-รมฺมณํ ฌานํ \csromanfolio{354}อุปฺปาเทติ. มคฺคํ ฯเปฯ อภิญฺญํ ฯเปฯ สมาปตฺติํ อุปฺปาเทติ. ปริตฺตา\-รมฺมณํ สีลํ ฯเปฯ. ปญฺญํ. ราคํ ฯเปฯ. กายิกํ สุขํ. กายิกํ ทุกฺขํ อุปนิสฺสาย อปฺปมาณา\-รมฺมณํ ฌานํ อุปฺปาเทติ. มคฺคํ. อภิญฺญํ. สมาปตฺติํ อุปฺปาเทติ. ปริตฺตารมฺมณา สทฺธา ฯเปฯ. กายิกํ สุขํ. กายิกํ ทุกฺขํ อปฺปมาณา\-รมฺมณาย สทฺธาย ฯเปฯ ปญฺญาย อุปนิสฺสย\-ปจฺจเยน ปจฺจโย. (3)}

\proseitem{26}{มหคฺคตา\-รมฺมโณ ธมฺโม มหคฺคตา\-รมฺมณสฺส ธมฺมสฺส อุปนิสฺสย\-ปจฺจเยน ปจฺจโย. อารมฺมณู\-ปนิสฺสโย, อนนฺตรู\-ปนิสฺสโย, ปกตู\-ปนิสฺสโย ฯเปฯ. ปกตู\-ปนิสฺสโย\csromandash{}มหคฺคตา\-รมฺมณํ สทฺธํ อุปนิสฺสาย มหคฺคตา\-รมฺมณํ ฌานํ อุปฺปาเทติ. วิปสฺสนํ. อภิญฺญํ. สมาปตฺติํ อุปฺปาเทติ. มานํ ชปฺเปติ. ทิฏฺฐิํ คณฺหาติ. มหคฺคตา\-รมฺมณํ สีลํ ฯเปฯ ปญฺญํ. ราคํ ฯเปฯ ปตฺถนํ อุปนิสฺสาย มหคฺคตา\-รมฺมณํ ฌานํ อุปฺปาเทติ ฯเปฯ ทิฏฺฐิํ คณฺหาติ. มหคฺคตา\-รมฺมณา สทฺธา ฯเปฯ ปญฺญา. ราโค ฯเปฯ ปตฺถนา มหคฺคตา\-รมฺมณาย สทฺธาย ฯเปฯ ปตฺถนาย อุปนิสฺสย\-ปจฺจเยน ปจฺจโย. (1)}

\prose{มหคฺคตา\-รมฺมโณ ธมฺโม ปริตฺตา\-รมฺมณสฺส ธมฺมสฺส อุปนิสฺสย\-ปจฺจเยน ปจฺจโย. อารมฺมณู\-ปนิสฺสโย, อนนฺตรู\-ปนิสฺสโย, ปกตู\-ปนิสฺสโย ฯเปฯ. ปกตู\-ปนิสฺสโย\csromandash{}มหคฺคตา\-รมฺมณํ สทฺธํ อุปนิสฺสาย ทานํ เทติ, สีลํ สมาทิยติ, อุโปสถกมฺมํ กโรติ. ปริตฺตา\-รมฺมณํ ฌานํ อุปฺปาเทติ. วิปสฺสนํ. อภิญฺญํ. สมาปตฺติํ อุปฺปาเทติ. มานํ ชปฺเปติ. ทิฏฺฐิํ คณฺหาติ. มหคฺคตา\-รมฺมณํ สีลํ ฯเปฯ ปตฺถนํ อุปนิสฺสาย ทานํ เทติ ฯเปฯ ทิฏฺฐิํ คณฺหาติ. มหคฺคตา\-รมฺมณา สทฺธา ฯเปฯ ปตฺถนา ปริตฺตา\-รมฺมณาย สทฺธาย ฯเปฯ ปตฺถนาย. กายิกสฺส สุขสฺส, กายิกสฺส ทุกฺขสฺส อุปนิสฺสย\-ปจฺจเยน ปจฺจโย. (2)}

\prose{มหคฺคตา\-รมฺมโณ ธมฺโม อปฺปมาณา\-รมฺมณสฺส ธมฺมสฺส อุปนิสฺสย\-ปจฺจเยน ปจฺจโย. อนนฺตรู\-ปนิสฺสโย, ปกตู\-ปนิสฺสโย ฯเปฯ. ปกตู\-ปนิสฺสโย\csromandash{}มหคฺคตา\-รมฺมณํ สทฺธํ อุปนิสฺสย อปฺปมาณา\-รมฺมณํ ฌานํ อุปฺปาเทติ. มคฺคํ. อภิญฺญํ. สมาปตฺติํ อุปฺปาเทติ. มหคฺคตา\-รมฺมณํ สีลํ ฯเปฯ ปตฺถนํ อุปนิสฺสาย อปฺปมาณา\-รมฺมณํ ฌานํ อุปฺปาเทติ ฯเปฯ สมาปตฺติํ อุปฺปาเทติ. มหคฺคตา\-รมฺมณา สทฺธา ฯเปฯ ปตฺถนา อปฺปมาณา\-รมฺมณาย สทฺธาย ฯเปฯ ปญฺญาย อุปนิสฺสย\-ปจฺจเยน ปจฺจโย. (3)}

\proseitem{27}{\csromanfolio{355}อปฺปมาณารมฺมโณ ธมฺโม อปฺปมาณา\-รมฺมณสฺส ธมฺมสฺส อุปนิสฺสย\-ปจฺจเยน ปจฺจโย. อารมฺมณู\-ปนิสฺสโย, อนนฺตรู\-ปนิสฺสโย, ปกตู\-ปนิสฺสโย ฯเปฯ. ปกตู\-ปนิสฺสโย\csromandash{}อปฺปมาณา\-รมฺมณํ สทฺธํ อุปนิสฺสาย อปฺปมาณา\-รมฺมณํ ฌานํ อุปฺปาเทติ. มคฺคํ. อภิญฺญํ. สมาปตฺติํ อุปฺปาเทติ. อปฺปมาณา\-รมฺมณํ สีลํ ฯเปฯ ปญฺญํ อุปนิสฺสาย อปฺปมาณา\-รมฺมณํ ฌานํ อุปฺปาเทติ. มคฺคํ ฯเปฯ อภิญฺญํ ฯเปฯ สมาปตฺติํ อุปฺปาเทติ. อปฺปมาณา\-รมฺมณา สทฺธา ฯเปฯ ปญฺญา อปฺปมาณา\-รมฺมณาย สทฺธาย ฯเปฯ ปญฺญาย. มคฺคสฺส, ผลสมาปตฺติยา อุปนิสฺสย\-ปจฺจเยน ปจฺจโย. (1)}

\prose{อปฺปมาณารมฺมโณ ธมฺโม ปริตฺตา\-รมฺมณสฺส ธมฺมสฺส อุปนิสฺสย\-ปจฺจเยน ปจฺจโย. อารมฺมณู\-ปนิสฺสโย, อนนฺตรู\-ปนิสฺสโย, ปกตู\-ปนิสฺสโย ฯเปฯ. ปกตู\-ปนิสฺสโย\csromandash{}อปฺปมาณา\-รมฺมณํ สทฺธํ อุปนิสฺสาย ทานํ เทติ, สีลํ สมาทิยติ, อุโปสถกมฺมํ กโรติ. ปริตฺตา\-รมฺมณํ ฌานํ อุปฺปาเทติ. วิปสฺสนํ. อภิญฺญํ. สมาปตฺติํ อุปฺปาเทติ. อปฺปมาณา\-รมฺมณํ สีลํ ฯเปฯ ปญฺญํ อุปนิสฺสาย ทานํ เทติ ฯเปฯ สมาปตฺติํ อุปฺปาเทติ. อปฺปมาณา\-รมฺมณา สทฺธา ฯเปฯ ปญฺญา ปริตฺตา\-รมฺมณาย สทฺธาย ฯเปฯ ปญฺญาย. กายิกสฺส สุขสฺส, กายิกสฺส ทุกฺขสฺส อุปนิสฺสย\-ปจฺจเยน ปจฺจโย. (2)}

\prose{อปฺปมาณารมฺมโณ ธมฺโม มหคฺคตา\-รมฺมณสฺส ธมฺมสฺส อุปนิสฺสย\-ปจฺจเยน ปจฺจโย. อารมฺมณู\-ปนิสฺสโย, อนนฺตรู\-ปนิสฺสโย, ปกตู\-ปนิสฺสโย ฯเปฯ. ปกตู\-ปนิสฺสโย\csromandash{}อปฺปมาณา\-รมฺมณํ สทฺธํ อุปนิสฺสาย มหคฺคตา\-รมฺมณํ ฌานํ อุปฺปาเทติ. วิปสฺสนํ. อภิญฺญํ. สมาปตฺติํ อุปฺปาเทติ. อปฺปมาณา\-รมฺมณํ สีลํ ฯเปฯ ปญฺญํ อุปนิสฺสาย มหคฺคตา\-รมฺมณํ ฌานํ อุปฺปาเทติ. วิปสฺสนํ. อภิญฺญํ. สมาปตฺติํ อุปฺปาเทติ. อปฺปมาณา\-รมฺมณา สทฺธา ฯเปฯ ปญฺญา มหคฺคตา\-รมฺมณาย สทฺธาย ฯเปฯ ปญฺญาย อุปนิสฺสย\-ปจฺจเยน ปจฺจโย. (3)}

\csromanheader{2}{\textbf{อาเสวน}}
\proseitem{28}{ปริตฺตารมฺมโณ ธมฺโม ปริตฺตา\-รมฺมณสฺส ธมฺมสฺส อาเสวน\-ปจฺจเยน ปจฺจโย. ปุริมา ปุริมา ปริตฺตารมฺมณา ขนฺธา ปจฺฉิมานํ ปจฺฉิมานํ ปริตฺตา\-รมฺมณานํ ขนฺธานํ อาเสวน\-ปจฺจเยน ปจฺจโย. (1)}

\prose{\csromanfolio{356}ปริตฺตารมฺมโณ ธมฺโม อปฺปมาณา\-รมฺมณสฺส ธมฺมสฺส อาเสวน\-ปจฺจเยน ปจฺจโย. ปริตฺตา\-รมฺมณํ อนุโลมํ โคตฺรภุสฺส. อนุโลมํ โวทานสฺส อาเสวน\-ปจฺจเยน ปจฺจโย. (2)}

\proseitem{29}{มหคฺคตา\-รมฺมโณ ธมฺโม มหคฺคตา\-รมฺมณสฺส ธมฺมสฺส อาเสวน\-ปจฺจเยน ปจฺจโย. ปุริมา ปุริมา มหคฺคตา\-รมฺมณา ขนฺธา ปจฺฉิมานํ ปจฺฉิมานํ มหคฺคตา\-รมฺมณานํ ขนฺธานํ อาเสวน\-ปจฺจเยน ปจฺจโย. (1)}

\prose{มหคฺคตา\-รมฺมโณ ธมฺโม อปฺปมาณา\-รมฺมณสฺส ธมฺมสฺส อาเสวน\-ปจฺจเยน ปจฺจโย. มหคฺคตา\-รมฺมณํ อนุโลมํ โคตฺรภุสฺส. อนุโลมํ โวทานสฺส อาเสวน\-ปจฺจเยน ปจฺจโย. (2)}

\proseitem{30}{อปฺปมาณารมฺมโณ ธมฺโม อปฺปมาณา\-รมฺมณสฺส ธมฺมสฺส อาเสวน\-ปจฺจเยน ปจฺจโย. ปุริมา ปุริมา อปฺปมาณา\-รมฺมณา ขนฺธา ปจฺฉิมานํ ปจฺฉิมานํ อปฺปมาณา\-รมฺมณานํ ขนฺธานํ อาเสวน\-ปจฺจเยน ปจฺจโย. โคตฺรภุ มคฺคสฺส. โวทานํ มคฺคสฺส อาเสวน\-ปจฺจเยน ปจฺจโย. (1)}

\csromanheader{2}{\textbf{กมฺม}}
\proseitem{31}{ปริตฺตารมฺมโณ ธมฺโม ปริตฺตา\-รมฺมณสฺส ธมฺมสฺส กมฺมปจฺจเยน ปจฺจโย. สหชาตา, นานากฺขณิกา.  สหชาตา ปริตฺตารมฺมณา เจตนา สมฺปยุตฺตกานํ ขนฺธานํ กมฺมปจฺจเยน ปจฺจโย. ปฏิสนฺธิ\-กฺขเณ ฯเปฯ. นานากฺขณิกา ปริตฺตารมฺมณา เจตนา วิปากานํ ปริตฺตา\-รมฺมณานํ ขนฺธานํ กมฺมปจฺจเยน ปจฺจโย. (1)}

\prose{มหคฺคตา\-รมฺมโณ ธมฺโม มหคฺคตา\-รมฺมณสฺส ธมฺมสฺส กมฺมปจฺจเยน ปจฺจโย. สหชาตา, นานากฺขณิกา.  สหชาตา มหคฺคตา\-รมฺมณา เจตนา สมฺปยุตฺตกานํ ขนฺธานํ กมฺมปจฺจเยน ปจฺจโย. ปฏิสนฺธิ\-กฺขเณ ฯเปฯ. นานากฺขณิกา มหคฺคตา\-รมฺมณา เจตนา วิปากานํ มหคฺคตา\-รมฺมณานํ ขนฺธานํ กมฺมปจฺจเยน ปจฺจโย. (1)}

\prose{มหคฺคตา\-รมฺมโณ ธมฺโม ปริตฺตา\-รมฺมณสฺส ธมฺมสฺส กมฺมปจฺจเยน ปจฺจโย. นานากฺขณิกา มหคฺคตา\-รมฺมณา เจตนา วิปากานํ ปริตฺตา\-รมฺมณานํ ขนฺธานํ กมฺมปจฺจเยน ปจฺจโย. (2)}

\proseitem{32}{\csromanfolio{357}อปฺปมาณารมฺมโณ ธมฺโม อปฺปมาณา\-รมฺมณสฺส ธมฺมสฺส กมฺมปจฺจเยน ปจฺจโย. สหชาตา, นานากฺขณิกา.  สหชาตา อปฺปมาณา\-รมฺมณา เจตนา สมฺปยุตฺตกานํ ขนฺธานํ กมฺมปจฺจเยน ปจฺจโย. นานากฺขณิกา อปฺปมาณา\-รมฺมณา เจตนา วิปากานํ อปฺปมาณา\-รมฺมณานํ ขนฺธานํ กมฺมปจฺจเยน ปจฺจโย. (1)}

\prose{อปฺปมาณารมฺมโณ ธมฺโม ปริตฺตา\-รมฺมณสฺส ธมฺมสฺส กมฺมปจฺจเยน ปจฺจโย. นานากฺขณิกา อปฺปมาณา\-รมฺมณา เจตนา วิปากานํ ปริตฺตา\-รมฺมณานํ ขนฺธานํ กมฺมปจฺจเยน ปจฺจโย. (2)}

\csromanheader{2}{\textbf{วิปากาทิ}}
\proseitem{33}{ปริตฺตารมฺมโณ ธมฺโม ปริตฺตา\-รมฺมณสฺส ธมฺมสฺส วิปาก\-ปจฺจเยน ปจฺจโย. อาหารปจฺจเยน ปจฺจโย. อินฺทฺริย\-ปจฺจเยน ปจฺจโย. ฌานปจฺจเยน ปจฺจโย. มคฺคปจฺจเยน ปจฺจโย. สมฺปยุตฺต\-ปจฺจเยน ปจฺจโย. อตฺถิปจฺจเยน ปจฺจโย. นตฺถิ\-ปจฺจเยน ปจฺจโย. วิคต\-ปจฺจเยน ปจฺจโย. อวิคต\-ปจฺจเยน ปจฺจโย.}

\csromanmatikaanchor{mtk.101}
\csromantocmarkref{subparagraph}{1. ปจฺจยานุโลม 2. สงฺขฺยาวาร}{mtk.101}
\bookmark[dest={mtk.101},level=5]{1. ปจฺจยานุโลม 2. สงฺขฺยาวาร}
\csromanmarkh{1. ปจฺจยานุโลม}\csromanmarkhii{2. สงฺขฺยาวาร}\csromanheadernomark{6}{1. \textbf{ปจฺจยานุโลม 2. สงฺขฺยาวาร}}
\proseitemwithcloserruled{34}{เหตุยา ตีณิ, อารมฺมเณ สตฺต, อธิปติยา สตฺต, อนนฺตเร นว, สมนนฺตเร นว, สหชาเต ตีณิ, อญฺญมญฺเญ ตีณิ, นิสฺสเย ตีณิ, อุปนิสฺสเย นว, อาเสวเน ปญฺจ, กมฺเม ปญฺจ, วิปาเก ตีณิ, อาหาเร ตีณิ, อินฺทฺริเย, ฌาเน, มคฺเค, สมฺปยุตฺเต, อตฺถิยา ตีณิ, นตฺถิยา นว, วิคเต นว, อวิคเต ตีณิ. (เอวํ คเณตพฺพํ.)}{อนุโลมํ.}
\csromanmatikaanchor{mtk.102}
\csromantocmarkref{subparagraph}{2. ปจฺจนียุทฺธาร}{mtk.102}
\bookmark[dest={mtk.102},level=5]{2. ปจฺจนียุทฺธาร}
\csromanheader{6}{2. ปจฺจนียุ\-ทฺธาร}
\proseitem{35}{ปริตฺตารมฺมโณ ธมฺโม ปริตฺตา\-รมฺมณสฺส ธมฺมสฺส อารมฺมณ\-ปจฺจเยน ปจฺจโย. สหชาต\-ปจฺจเยน ปจฺจโย. อุปนิสฺสย\-ปจฺจเยน ปจฺจโย. กมฺมปจฺจเยน ปจฺจโย. (1)}

\prose{\csromanfolio{358}ปริตฺตารมฺมโณ ธมฺโม มหคฺคตา\-รมฺมณสฺส ธมฺมสฺส อารมฺมณ\-ปจฺจเยน ปจฺจโย. อุปนิสฺสย\-ปจฺจเยน ปจฺจโย. (2)}

\prose{ปริตฺตารมฺมโณ ธมฺโม อปฺปมาณา\-รมฺมณสฺส ธมฺมสฺส อุปนิสฺสย\-ปจฺจเยน ปจฺจโย. (3)}

\proseitem{36}{มหคฺคตา\-รมฺมโณ ธมฺโม มหคฺคตา\-รมฺมณสฺส ธมฺมสฺส อารมฺมณ\-ปจฺจเยน ปจฺจโย. สหชาต\-ปจฺจเยน ปจฺจโย. อุปนิสฺสย\-ปจฺจเยน ปจฺจโย. (1)}

\prose{มหคฺคตา\-รมฺมโณ ธมฺโม ปริตฺตา\-รมฺมณสฺส ธมฺมสฺส อารมฺมณ\-ปจฺจเยน ปจฺจโย. อุปนิสฺสย\-ปจฺจเยน ปจฺจโย. กมฺมปจฺจเยน ปจฺจโย. (2)}

\prose{มหคฺคตา\-รมฺมโณ ธมฺโม อปฺปมาณา\-รมฺมณสฺส ธมฺมสฺส อุปนิสฺสย\-ปจฺจเยน ปจฺจโย. (3)}

\proseitem{37}{อปฺปมาณารมฺมโณ ธมฺโม อปฺปมาณา\-รมฺมณสฺส ธมฺมสฺส อารมฺมณ\-ปจฺจเยน ปจฺจโย. สหชาต\-ปจฺจเยน ปจฺจโย. อุปนิสฺสย\-ปจฺจเยน ปจฺจโย. (1)}

\prose{อปฺปมาณารมฺมโณ ธมฺโม ปริตฺตา\-รมฺมณสฺส ธมฺมสฺส อารมฺมณ\-ปจฺจเยน ปจฺจโย. อุปนิสฺสย\-ปจฺจเยน ปจฺจโย. กมฺมปจฺจเยน ปจฺจโย. (2)}

\prose{อปฺปมาณารมฺมโณ ธมฺโม มหคฺคตา\-รมฺมณสฺส ธมฺมสฺส อารมฺมณ\-ปจฺจเยน ปจฺจโย. อุปนิสฺสย\-ปจฺจเยน ปจฺจโย. (3)}

\csromanmatikaanchor{mtk.103}
\csromantocmarkref{subparagraph}{2. ปจฺจยปจฺจนีย 2. สงฺขฺยาวาร}{mtk.103}
\bookmark[dest={mtk.103},level=5]{2. ปจฺจยปจฺจนีย 2. สงฺขฺยาวาร}
\csromanmarkh{2. ปจฺจย\-ปจฺจนีย}\csromanmarkhii{2. สงฺขฺยาวาร}\csromanheadernomark{6}{2. ปจฺจย\-ปจฺจนีย 2. สงฺขฺยาวาร}
\proseitem{38}{นเหตุยา นว, นอารมฺมเณ นว, นอธิปติยา นว, นอนนฺตเร นว, นสมนนฺตเร นว, นสหชาเต นว, นอญฺญมญฺเญ นว, นนิสฺสเย นว, นอุปนิสฺสเย สตฺต, นปุเรชาเต นว, นปจฺฉาชาเต นว, นอาเสวเน นว ฯเปฯ นมคฺเค นว, นสมฺปยุตฺเต นว, นวิปฺปยุตฺเต นว, โนอตฺถิยา นว, โนนตฺถิยา นว, โนวิคเต นว, โนอวิคเต นว.  (เอวํ คเณตพฺพํ.)}

\vspace{\tipitakaparskip}
\csromancenter{ปจฺจนียํ.}
\csromanmarkh{13. ปริตฺตา\-รมฺมณ\-ตฺติก}\csromanmarkhii{3. ปจฺจยา\-นุโลม\-ปจฺจนีย}\csromanheadernomark{2}{13. ปริตฺตา\-รมฺมณ\-ตฺติก 3. ปจฺจยา\-นุโลม\-ปจฺจนีย}
\csromanmatikaanchor{mtk.104}
\csromanmatikaanchor{mtk.105}
\csromanmatikaanchor{mtk.106}
\csromanmatikaanchor{mtk.107}
\csromantocmarknopage{subparagraph}{3. ปจฺจยานุโลมปจฺจนีย}{mtk.104}
\bookmark[dest={mtk.104},level=5]{3. ปจฺจยานุโลมปจฺจนีย}
\csromantocmarkref{subparagraph}{เหตุทุก}{mtk.105}
\bookmark[dest={mtk.105},level=5]{เหตุทุก}
\csromantocmarknopage{subparagraph}{4. ปจฺจยปจฺจนียานุโลม}{mtk.106}
\bookmark[dest={mtk.106},level=5]{4. ปจฺจยปจฺจนียานุโลม}
\csromantocmarkref{subparagraph}{นเหตุทุก}{mtk.107}
\bookmark[dest={mtk.107},level=5]{นเหตุทุก}
\proseitemwithcloserruled{39}{\csromanfolio{359}เหตุปจฺจยา นอารมฺมเณ ตีณิ, นอธิปติยา, นอนนฺตเร, นสมนนฺตเร, นอุปนิสฺสเย, นปุเรชาเต, นปจฺฉาชาเต, นอาเสวเน ตีณิ ฯเปฯ นมคฺเค, นวิปฺปยุตฺเต, โนนตฺถิยา, โนวิคเต ตีณิ.  (เอวํ คเณตพฺพํ.)}{อนุโลม\-ปจฺจนียํ.}
\proseitem{40}{นเหตุปจฺจยา อารมฺมเณ สตฺต, อธิปติยา สตฺต, อนนฺตเร นว, สมนนฺตเร นว, สหชาเต ตีณิ, อญฺญมญฺเญ ตีณิ, นิสฺสเย ตีณิ, อุปนิสฺสเย นว, อาเสวเน ปญฺจ, กมฺเม ปญฺจ, วิปาเก ตีณิ ฯเปฯ สมฺปยุตฺเต ตีณิ, อตฺถิยา ตีณิ, นตฺถิยา นว, วิคเต นว, อวิคเต ตีณิ.  (เอวํ คเณตพฺพํ.)}

\vspace{\tipitakaparskip}
\csromancenter{ปจฺจนียา\-นุโลมํ.}
\csromancenter{ปญฺหาวาโร.}
\nitthitam{ปริตฺตา\-รมฺมณ\-ตฺติกํ นิฏฺฐิตํ.}
\csromanmatikaanchor{mtk.108}
\csromantocmarknopage{subparagraph}{14. หีนตฺติก 1. ปฏิจฺจวาร}{mtk.108}
\bookmark[dest={mtk.108},level=5]{14. หีนตฺติก 1. ปฏิจฺจวาร}
\csromanmarkh{14. หีนตฺติก}\csromanmarkhii{1. ปฏิจฺจวาร}\csromanheadernomark{6}{14. \textbf{หีนตฺติก 1. ปฏิจฺจวาร}}
\csromanmarkh{1. ปจฺจยานุโลม}\csromanmarkhii{1. วิภงฺควาร}\csromanheadernomark{2}{1. \textbf{ปจฺจยานุโลม 1. วิภงฺควาร}}
\csromanmatikaanchor{mtk.109}
\csromanmatikaanchor{mtk.110}
\csromantocmarknopage{subparagraph}{1. ปจฺจยานุโลม 1. วิภงฺควาร}{mtk.109}
\bookmark[dest={mtk.109},level=5]{1. ปจฺจยานุโลม 1. วิภงฺควาร}
\csromantocmarkref{subparagraph}{เหตุ}{mtk.110}
\bookmark[dest={mtk.110},level=5]{เหตุ}
\csromanheader{6}{\textbf{เหตุ}}
\proseitem{1}{\csromanfolio{360}หีนํ ธมฺมํ ปฏิจฺจ หีโน ธมฺโม อุปฺปชฺชติ เหตุปจฺจยา. หีนํ เอกํ ขนฺธํ ปฏิจฺจ ตโย ขนฺธา ฯเปฯ เทฺว ขนฺเธ ปฏิจฺจ เทฺว ขนฺธา. (1)}

\prose{หีนํ ธมฺมํ ปฏิจฺจ มชฺฌิโม ธมฺโม อุปฺปชฺชติ เหตุปจฺจยา. หีเน ขนฺเธ ปฏิจฺจ จิตฺต\-สมุฏฺฐานํ รูปํ. (2)}

\prose{หีนํ ธมฺมํ ปฏิจฺจ หีโน จ มชฺฌิโม จ ธมฺมา อุปฺปชฺชนฺติ เหตุปจฺจยา. หีนํ เอกํ ขนฺธํ ปฏิจฺจ ตโย ขนฺธา จิตฺต\-สมุฏฺฐานญฺจ รูปํ ฯเปฯ เทฺว ขนฺเธ ฯเปฯ. (3)}

\proseitem{2}{มชฺฌิมํ ธมฺมํ ปฏิจฺจ มชฺฌิโม ธมฺโม อุปฺปชฺชติ เหตุปจฺจยา. มชฺฌิมํ เอกํ ขนฺธํ ปฏิจฺจ ตโย ขนฺธา จิตฺต\-สมุฏฺฐานญฺจ รูปํ ฯเปฯ เทฺว ขนฺเธ ฯเปฯ. ปฏิสนฺธิ\-กฺขเณ ฯเปฯ. ขนฺเธ ปฏิจฺจ วตฺถุ, วตฺถุํ ปฏิจฺจ ขนฺธา. เอกํ มหาภูตํ ฯเปฯ มหาภูเต ปฏิจฺจ จิตฺต\-สมุฏฺฐานํ รูปํ, กฏตฺตารูปํ, อุปาทารูปํ. (1)}

\proseitem{3}{ปณีตํ ธมฺมํ ปฏิจฺจ ปณีโต ธมฺโม อุปฺปชฺชติ เหตุปจฺจยา. ตีณิ.}

\proseitem{4}{มชฺฌิมญฺจ ปณีตญฺจ ธมฺมํ ปฏิจฺจ มชฺฌิโม ธมฺโม อุปฺปชฺชติ เหตุปจฺจยา. ปณีเต ขนฺเธ จ มหาภูเต จ ปฏิจฺจ จิตฺต\-สมุฏฺฐานํ รูปํ. (1)}

\proseitemwithcloser{5}{หีนญฺจ มชฺฌิมญฺจ ธมฺมํ ปฏิจฺจ มชฺฌิโม ธมฺโม อุปฺปชฺชติ เหตุปจฺจยา. หีเน ขนฺเธ จ มหาภูเต จ ปฏิจฺจ จิตฺต\-สมุฏฺฐานํ รูปํ. (1) (หีนตฺติกํ สํกิลิฏฺฐ\-ตฺติก\-สทิสํ วิตฺถาเรตพฺพํ ปริปุณฺณํ.)}{หีนตฺติกํ นิฏฺฐิตํ.}
\csromanmatikaanchor{mtk.111}
\csromantocmarknopage{subparagraph}{15. มิจฺฉตฺตนิยตตฺติก 1. ปฏิจฺจวาร}{mtk.111}
\bookmark[dest={mtk.111},level=5]{15. มิจฺฉตฺตนิยตตฺติก 1. ปฏิจฺจวาร}
\csromanmarkh{15. มิจฺฉตฺตนิยต\-ตฺติก}\csromanmarkhii{15. มิจฺฉตฺตนิยต\-ตฺติก 1. ปฏิจฺจวาร}\csromanheadernomark{6}{15. มิจฺฉตฺตนิยต\-ตฺติก 15. มิจฺฉตฺตนิยต\-ตฺติก 1. ปฏิจฺจวาร}
\csromanmatikaanchor{mtk.112}
\csromantocmarkref{subparagraph}{1. ปจฺจยานุโลม 1. วิภงฺควาร}{mtk.112}
\bookmark[dest={mtk.112},level=5]{1. ปจฺจยานุโลม 1. วิภงฺควาร}
\csromanmarkh{1. ปจฺจยานุโลม}\csromanmarkhii{1. วิภงฺควาร}\csromanheadernomark{6}{1. \textbf{ปจฺจยานุโลม 1. วิภงฺควาร}}
\csromanheader{2}{\textbf{เหตุ}}
\proseitem{1}{\csromanfolio{361}มิจฺฉตฺต\-นิยตํ ธมฺมํ ปฏิจฺจ มิจฺฉตฺต\-นิยโต ธมฺโม อุปฺปชฺชติ เหตุปจฺจยา. มิจฺฉตฺต\-นิยตํ เอกํ ขนฺธํ ปฏิจฺจ ตโย ขนฺธา ฯเปฯ เทฺว ขนฺธา. (1)}

\prose{มิจฺฉตฺต\-นิยตํ ธมฺมํ ปฏิจฺจ อนิยโต ธมฺโม อุปฺปชฺชติ เหตุปจฺจยา. มิจฺฉตฺต\-นิยเต ขนฺเธ ปฏิจฺจ จิตฺต\-สมุฏฺฐานํ รูปํ. (2)}

\prose{มิจฺฉตฺต\-นิยตํ ธมฺมํ ปฏิจฺจ มิจฺฉตฺต\-นิยโต จ อนิยโต จ ธมฺมา อุปฺปชฺชนฺติ เหตุปจฺจยา. มิจฺฉตฺต\-นิยตํ เอกํ ขนฺธํ ปฏิจฺจ ตโย ขนฺธา จิตฺต\-สมุฏฺฐานญฺจ รูปํ ฯเปฯ เทฺว ขนฺเธ ฯเปฯ. (3)}

\proseitem{2}{สมฺมตฺต\-นิยตํ ธมฺมํ ปฏิจฺจ สมฺมตฺตนิยโต ธมฺโม อุปฺปชฺชติ เหตุปจฺจยา. ตีณิ.}

\proseitem{3}{อนิยตํ ธมฺมํ ปฏิจฺจ อนิยโต ธมฺโม อุปฺปชฺชติ เหตุปจฺจยา. อนิยตํ เอกํ ขนฺธํ ปฏิจฺจ ตโย ขนฺธา จิตฺต\-สมุฏฺฐานญฺจ รูปํ ฯเปฯ เทฺว ขนฺเธ ฯเปฯ. ปฏิสนฺธิ\-กฺขเณ อนิยตํ เอกํ ขนฺธํ ปฏิจฺจ ตโย ขนฺธา กฏตฺตา จ รูปํ ฯเปฯ เทฺว ขนฺเธ ฯเปฯ. ขนฺเธ ปฏิจฺจ วตฺถุ, วตฺถุํ ปฏิจฺจ ขนฺธา. เอกํ มหาภูตํ ปฏิจฺจ ตโย มหาภูตา ฯเปฯ เทฺว มหาภูตา. มหาภูเต ปฏิจฺจ จิตฺต\-สมุฏฺฐานํ รูปํ, กฏตฺตารูปํ, อุปาทารูปํ. (1)}

\proseitem{4}{มิจฺฉตฺตนิยตญฺจ อนิยตญฺจ ธมฺมํ ปฏิจฺจ อนิยโต ธมฺโม อุปฺปชฺชติ เหตุปจฺจยา. มิจฺฉตฺต\-นิยเต ขนฺเธ จ มหาภูเต จ ปฏิจฺจ จิตฺต\-สมุฏฺฐานํ รูปํ. (1)}

\prose{สมฺมตฺตนิยตญฺจ อนิยตญฺจ ธมฺมํ ปฏิจฺจ อนิยโต ธมฺโม อุปฺปชฺชติ เหตุปจฺจยา. สมฺมตฺตนิยเต ขนฺเธ จ มหาภูเต จ ปฏิจฺจ จิตฺต\-สมุฏฺฐานํ รูปํ. (1)}

\csromanheader{2}{\textbf{อารมฺมณ}}
\proseitem{5}{\csromanfolio{362}มิจฺฉตฺต\-นิยตํ ธมฺมํ ปฏิจฺจ มิจฺฉตฺต\-นิยโต ธมฺโม อุปฺปชฺชติ อารมฺมณ\-ปจฺจยา. มิจฺฉตฺต\-นิยตํ เอกํ ขนฺธํ ปฏิจฺจ ตโย ขนฺธา ฯเปฯ เทฺว ขนฺเธ ฯเปฯ. (1)}

\prose{สมฺมตฺต\-นิยตํ ธมฺมํ ปฏิจฺจ สมฺมตฺตนิยโต ธมฺโม อุปฺปชฺชติ อารมฺมณ\-ปจฺจยา. สมฺมตฺต\-นิยตํ เอกํ ขนฺธํ ปฏิจฺจ ตโย ขนฺธา ฯเปฯ เทฺว ขนฺเธ ฯเปฯ. (1)}

\prose{อนิยตํ ธมฺมํ ปฏิจฺจ อนิยโต ธมฺโม อุปฺปชฺชติ อารมฺมณ\-ปจฺจยา. อนิยตํ เอกํ ขนฺธํ ปฏิจฺจ ตโย ขนฺธา ฯเปฯ เทฺว ขนฺเธ ฯเปฯ. ปฏิสนฺธิ\-กฺขเณ อนิยตํ เอกํ ขนฺธํ ปฏิจฺจ ตโย ขนฺธา ฯเปฯ เทฺว ขนฺเธ ฯเปฯ. วตฺถุํ ปฏิจฺจ ขนฺธา. (สพฺเพ ปจฺจยา อิมินา การเณน วิตฺถาเรตพฺพา. สํขิตฺตํ.)}

\csromanmatikaanchor{mtk.113}
\csromantocmarkref{subparagraph}{1. ปจฺจยานุโลม 2. สงฺขฺยาวาร}{mtk.113}
\bookmark[dest={mtk.113},level=5]{1. ปจฺจยานุโลม 2. สงฺขฺยาวาร}
\csromanmarkh{1. ปจฺจยานุโลม}\csromanmarkhii{2. สงฺขฺยาวาร}\csromanheadernomark{6}{1. \textbf{ปจฺจยานุโลม 2. สงฺขฺยาวาร}}
\proseitemwithcloserruled{6}{เหตุยา นว, อารมฺมเณ ตีณิ, อธิปติยา นว, อนนฺตเร ตีณิ, สมนนฺตเร ตีณิ, สหชาเต นว, อญฺญมญฺเญ ตีณิ, นิสฺสเย นว, อุปนิสฺสเย ตีณิ, ปุเรชาเต ตีณิ, อาเสวเน ตีณิ, กมฺเม นว, วิปาเก เอกํ, อาหาเร นว, อินฺทฺริเย นว, ฌาเน นว, มคฺเค นว, สมฺปยุตฺเต ตีณิ, วิปฺปยุตฺเต นว, อตฺถิยา นว, นตฺถิยา ตีณิ, วิคเต ตีณิ, อวิคเต นว.  (เอวํ คเณตพฺพํ.)}{อนุโลมํ.}
\csromanmatikaanchor{mtk.114}
\csromantocmarkref{subparagraph}{2. ปจฺจยปจฺจนีย 1. วิภงฺควาร}{mtk.114}
\bookmark[dest={mtk.114},level=5]{2. ปจฺจยปจฺจนีย 1. วิภงฺควาร}
\csromanmarkh{2. ปจฺจย\-ปจฺจนีย}\csromanmarkhii{1. วิภงฺควาร}\csromanheadernomark{6}{2. ปจฺจย\-ปจฺจนีย 1. วิภงฺควาร}
\csromanheader{2}{\textbf{นเหตุ}}
\proseitem{7}{อนิยตํ ธมฺมํ ปฏิจฺจ อนิยโต ธมฺโม อุปฺปชฺชติ นเหตุปจฺจยา. อเหตุกํ อนิยตํ เอกํ ขนฺธํ ปฏิจฺจ ตโย ขนฺธา จิตฺต\-สมุฏฺฐานญฺจ รูปํ ฯเปฯ เทฺว ขนฺเธ ฯเปฯ. อเหตุก\-ปฏิสนฺธิ\-กฺขเณ ฯเปฯ. เอกํ มหาภูตํ ปฏิจฺจ ฯเปฯ. พาหิรํ. อาหารสมุฏฺฐานํ. อุตุสมุฏฺฐานํ. อสญฺญสตฺตานํ ฯเปฯ. วิจิกิจฺฉา\-สหคเต อุทฺธจฺจ\-สหคเต ขนฺเธ ปฏิจฺจ วิจิกิจฺฉา\-สหคโต อุทฺธจฺจ\-สหคโต โมโห. (1)}

\csromanheader{2}{15. มิจฺฉตฺตนิยต\-ตฺติก นอารมฺมณ}
\proseitem{8}{\csromanfolio{363}มิจฺฉตฺต\-นิยตํ ธมฺมํ ปฏิจฺจ อนิยโต ธมฺโม อุปฺปชฺชติ นอารมฺมณ\-ปจฺจยา. มิจฺฉตฺต\-นิยเต ขนฺเธ ปฏิจฺจ จิตฺต\-สมุฏฺฐานํ รูปํ. (สํขิตฺตํ.)}

\prose{นอธิปติ}

\proseitem{9}{มิจฺฉตฺต\-นิยตํ ธมฺมํ ปฏิจฺจ มิจฺฉตฺต\-นิยโต ธมฺโม อุปฺปชฺชติ นอธิปติ\-ปจฺจยา. มิจฺฉตฺต\-นิยเต ขนฺเธ ปฏิจฺจ มิจฺฉตฺต\-นิยตาธิปติ. (1)}

\prose{สมฺมตฺต\-นิยตํ ธมฺมํ ปฏิจฺจ สมฺมตฺตนิยโต ธมฺโม อุปฺปชฺชติ นอธิปติ\-ปจฺจยา. สมฺมตฺตนิยเต ขนฺเธ ปฏิจฺจ สมฺมตฺตนิยตา\-ธิปติ. (1)}

\prose{อนิยตํ ธมฺมํ ปฏิจฺจ อนิยโต ธมฺโม อุปฺปชฺชติ นอธิปติ\-ปจฺจยา. อนิยตํ เอกํ ขนฺธํ ปฏิจฺจ ตโย ขนฺธา จิตฺต\-สมุฏฺฐานญฺจ รูปํ ฯเปฯ เทฺว ขนฺเธ ฯเปฯ. ปฏิสนฺธิ\-กฺขเณ ฯเปฯ. ขนฺเธ ปฏิจฺจ วตฺถุ, วตฺถุํ ปฏิจฺจ ขนฺธา. เอกํ มหาภูตํ ฯเปฯ. อสญฺญสตฺตานํ ฯเปฯ. (1)}

\prose{นอนนฺตร}

\proseitem{10}{มิจฺฉตฺต\-นิยตํ ธมฺมํ ปฏิจฺจ อนิยโต ธมฺโม อุปฺปชฺชติ นอนนฺตร\-ปจฺจยา. (สํขิตฺตํ. สพฺพานิ ปจฺจยานิ วิตฺถาเร\-ตพฺพานิ.)}

\csromanmatikaanchor{mtk.115}
\csromantocmarkref{subparagraph}{2. ปจฺจยปจฺจนีย 2. สงฺขฺยาวาร}{mtk.115}
\bookmark[dest={mtk.115},level=5]{2. ปจฺจยปจฺจนีย 2. สงฺขฺยาวาร}
\csromanmarkh{2. ปจฺจย\-ปจฺจนีย}\csromanmarkhii{2. สงฺขฺยาวาร}\csromanheadernomark{6}{2. ปจฺจย\-ปจฺจนีย 2. สงฺขฺยาวาร}
\proseitem{11}{นเหตุยา เอกํ, นอารมฺมเณ ปญฺจ, นอธิปติยา ตีณิ, นอนนฺตเร ปญฺจ, นสมนนฺตเร ปญฺจ, นอญฺญมญฺเญ ปญฺจ, นอุปนิสฺสเย ปญฺจ, นปุเรชาเต ฉ, นปจฺฉาชาเต นว, นอาเสวเน ปญฺจ, นกมฺเม ตีณิ, นวิปาเก นว, นอาหาเร เอกํ, นอินฺทฺริเย เอกํ, นฌาเน เอกํ, นมคฺเค เอกํ, นสมฺปยุตฺเต ปญฺจ, นวิปฺปยุตฺเต เทฺว, โนนตฺถิยา ปญฺจ, โนวิคเต ปญฺจ.  (เอวํ คเณตพฺพํ.)}

\vspace{\tipitakaparskip}
\csromancenter{ปจฺจนียํ.}
\csromancenter{\textbf{เหตุ}ปจฺจยา นอารมฺมเณ ปญฺจ, นอธิปติยา ตีณิ, นอนนฺตเร, นสมนนฺตเร, นอญฺญมญฺเญ, นฺเอาปนิสฺสเย ปญฺจ, นปุเรชาเต ฉ, นปจฺฉาชาเต นว, นอาเสวเน ปญฺจ, นกมฺเม ตีณิ, นวิปาเก นว, นสมฺปยุตฺเต ปญฺจ, นวิปฺปยุตฺเต เทฺว, โนนตฺถิยา ปญฺจ, โนวิคเต ปญฺจ.  (เอวํ คเณตพฺพํ.)}
\nitthitamruled{อนุโลม\-ปจฺจนียํ.}
\proseitem{13}{\csromanfolio{364}นเหตุปจฺจยา อารมฺมเณ เอกํ, อนนฺตเร เอกํ, สมนนฺตเร เอกํ, สหชาเต เอกํ ฯเปฯ วิคเต เอกํ.  (เอวํ คเณตพฺพํ.) ปจฺจนียา\-นุโลมํ. ปฏิจฺจวาโร. (สหชาตวาโร ปฏิจฺจ\-วาร\-สทิโส.)}
\csromansectionrule

\csromanmatikaanchor{mtk.116}
\csromanmatikaanchor{mtk.120}
\csromantocmarknopage{subparagraph}{3. ปจฺจยานุโลมปจฺจนีย}{mtk.116}
\bookmark[dest={mtk.116},level=5]{3. ปจฺจยานุโลมปจฺจนีย}
\csromanmarkh{15. มิจฺฉตฺตนิยต\-ตฺติก}\csromanmarkhii{3. ปจฺจยวาร}\csromanheadernomark{6}{15. มิจฺฉตฺตนิยต\-ตฺติก 3. ปจฺจยวาร}
\csromanmatikaanchor{mtk.117}
\csromanmatikaanchor{mtk.121}
\csromantocmarkref{subparagraph}{เหตุทุก}{mtk.117}
\bookmark[dest={mtk.117},level=5]{เหตุทุก}
\csromanmarkh{1. ปจฺจยานุโลม}\csromanmarkhii{1. วิภงฺควาร}\csromanheadernomark{6}{1. \textbf{ปจฺจยานุโลม 1. วิภงฺควาร}}
\csromanmatikaanchor{mtk.118}
\csromanmatikaanchor{mtk.119}
\csromantocmarknopage{subparagraph}{4. ปจฺจยปจฺจนียานุโลม}{mtk.118}
\bookmark[dest={mtk.118},level=5]{4. ปจฺจยปจฺจนียานุโลม}
\csromantocmarkref{subparagraph}{นเหตุทุก}{mtk.119}
\bookmark[dest={mtk.119},level=5]{นเหตุทุก}
\csromantocmarknopage{subparagraph}{15. มิจฺฉตฺตนิยตตฺติก 3. ปจฺจยวาร}{mtk.120}
\bookmark[dest={mtk.120},level=5]{15. มิจฺฉตฺตนิยตตฺติก 3. ปจฺจยวาร}
\csromantocmarkref{subparagraph}{1. ปจฺจยานุโลม 1. วิภงฺควาร}{mtk.121}
\bookmark[dest={mtk.121},level=5]{1. ปจฺจยานุโลม 1. วิภงฺควาร}
\csromanheader{6}{\textbf{เหตุ}}
\proseitem{14}{มิจฺฉตฺต\-นิยตํ ธมฺมํ ปจฺจยา มิจฺฉตฺต\-นิยโต ธมฺโม อุปฺปชฺชติ เหตุปจฺจยา. ตีณิ.}

\prose{สมฺมตฺต\-นิยตํ ธมฺมํ ปจฺจยา สมฺมตฺตนิยโต ธมฺโม อุปฺปชฺชติ เหตุปจฺจยา. ตีณิ.}

\proseitem{15}{\csromanfolio{365}อนิยตํ ธมฺมํ ปจฺจยา อนิยโต ธมฺโม อุปฺปชฺชติ เหตุปจฺจยา. อนิยตํ เอกํ ขนฺธํ ปจฺจยา ตโย ขนฺธา จิตฺต\-สมุฏฺฐานญฺจ รูปํ ฯเปฯ เทฺว ขนฺเธ ฯเปฯ. ปฏิสนฺธิ\-กฺขเณ ฯเปฯ. ขนฺเธ ปจฺจยา วตฺถุ, วตฺถุํ ปจฺจยา ขนฺธา. เอกํ มหาภูตํ ปจฺจยา ฯเปฯ. วตฺถุํ ปจฺจยา อนิยตา ขนฺธา. (1)}

\prose{อนิยตํ ธมฺมํ ปจฺจยา มิจฺฉตฺต\-นิยโต ธมฺโม อุปฺปชฺชติ เหตุปจฺจยา. วตฺถุํ ปจฺจยา มิจฺฉตฺต\-นิยตา ขนฺธา. (2)}

\prose{อนิยตํ ธมฺมํ ปจฺจยา สมฺมตฺตนิยโต ธมฺโม อุปฺปชฺชติ เหตุปจฺจยา. วตฺถุํ ปจฺจยา สมฺมตฺตนิยตา ขนฺธา. (3)}

\prose{อนิยตํ ธมฺมํ ปจฺจยา มิจฺฉตฺต\-นิยโต จ อนิยโต จ ธมฺมา อุปฺปชฺชนฺติ เหตุปจฺจยา. วตฺถุํ ปจฺจยา มิจฺฉตฺต\-นิยตา ขนฺธา. มหาภูเต ปจฺจยา จิตฺต\-สมุฏฺฐานํ รูปํ. (4)}

\prose{อนิยตํ ธมฺมํ ปจฺจยา สมฺมตฺตนิยโต จ อนิยโต จ ธมฺมา อุปฺปชฺชนฺติ เหตุปจฺจยา. วตฺถุํ ปจฺจยา สมฺมตฺตนิยตา ขนฺธา. มหาภูเต ปจฺจยา จิตฺต\-สมุฏฺฐานํ รูปํ. (5)}

\proseitem{16}{มิจฺฉตฺตนิยตญฺจ อนิยตญฺจ ธมฺมํ ปจฺจยา มิจฺฉตฺต\-นิยโต ธมฺโม อุปฺปชฺชติ เหตุปจฺจยา. มิจฺฉตฺต\-นิยตํ เอกํ ขนฺธญฺจ วตฺถุญฺจ ปจฺจยา ตโย ขนฺธา ฯเปฯ เทฺว ขนฺเธ ฯเปฯ. (1)}

\prose{มิจฺฉตฺตนิยตญฺจ อนิยตญฺจ ธมฺมํ ปจฺจยา อนิยโต ธมฺโม อุปฺปชฺชติ เหตุปจฺจยา. มิจฺฉตฺต\-นิยเต ขนฺเธ จ มหาภูเต จ ปจฺจยา จิตฺต\-สมุฏฺฐานํ รูปํ. (2)}

\prose{มิจฺฉตฺตนิยตญฺจ อนิยตญฺจ ธมฺมํ ปจฺจยา มิจฺฉตฺต\-นิยโต จ อนิยโต จ ธมฺมา อุปฺปชฺชนฺติ เหตุปจฺจยา. มิจฺฉตฺต\-นิยตํ เอกํ ขนฺธญฺจ วตฺถุญฺจ ปจฺจยา ตโย ขนฺธา ฯเปฯ เทฺว ขนฺเธ ฯเปฯ. มิจฺฉตฺต\-นิยเต ขนฺเธ จ มหาภูเต จ ปจฺจยา จิตฺต\-สมุฏฺฐานํ รูปํ. (3)}

\prose{สมฺมตฺตนิยตญฺจ อนิยตญฺจ ธมฺมํ ปจฺจยา สมฺมตฺตนิยโต ธมฺโม อุปฺปชฺชติ เหตุปจฺจยา. (ตีณิ ปญฺหา. มิจฺฉตฺต\-สทิสํ.)}

\csromanheader{2}{\textbf{อารมฺมณาทิ}}
\proseitem{17}{\csromanfolio{366}มิจฺฉตฺต\-นิยตํ ธมฺมํ ปจฺจยา มิจฺฉตฺต\-นิยโต ธมฺโม อุปฺปชฺชติ อารมฺมณ\-ปจฺจยา. (สํขิตฺตํ. กุสลตฺติเก ปจฺจยวาร\-สทิสํ วิภชิตพฺพํ.) อวิคตปจฺจยา.}

\csromanmatikaanchor{mtk.122}
\csromantocmarkref{subparagraph}{1. ปจฺจยานุโลม 2. สงฺขฺยาวาร}{mtk.122}
\bookmark[dest={mtk.122},level=5]{1. ปจฺจยานุโลม 2. สงฺขฺยาวาร}
\csromanmarkh{1. ปจฺจยานุโลม}\csromanmarkhii{2. สงฺขฺยาวาร}\csromanheadernomark{6}{1. \textbf{ปจฺจยานุโลม 2. สงฺขฺยาวาร}}
\proseitemwithcloserruled{18}{เหตุยา สตฺตรส, อารมฺมเณ สตฺต, อธิปติยา สตฺตรส, อนนฺตเร สตฺต, สมนนฺตเร สตฺต, สหชาเต สตฺตรส, อญฺญมญฺเญ สตฺต, นิสฺสเย สตฺตรส, อุปนิสฺสเย สตฺต, ปุเรชาเต สตฺต, อาเสวเน สตฺต, กมฺเม สตฺตรส, วิปาเก เอกํ, อาหาเร สตฺตรส, อินฺทฺริเย สตฺตรส, ฌาเน สตฺตรส, มคฺเค สตฺตรส, สมฺปยุตฺเต สตฺต, วิปฺปยุตฺเต สตฺตรส, อตฺถิยา สตฺตรส, นตฺถิยา สตฺต, วิคเต สตฺต, อวิคเต สตฺตรส.  (เอวํ คเณตพฺพํ.)}{อนุโลมํ.}
\csromanmatikaanchor{mtk.123}
\csromantocmarkref{subparagraph}{2. ปจฺจยปจฺจนีย 1. วิภงฺควาร}{mtk.123}
\bookmark[dest={mtk.123},level=5]{2. ปจฺจยปจฺจนีย 1. วิภงฺควาร}
\csromanmarkh{2. ปจฺจย\-ปจฺจนีย}\csromanmarkhii{1. วิภงฺควาร}\csromanheadernomark{6}{2. ปจฺจย\-ปจฺจนีย 1. วิภงฺควาร}
\csromanheader{2}{\textbf{นเหตุ}}
\proseitem{19}{อนิยตํ ธมฺมํ ปจฺจยา อนิยโต ธมฺโม อุปฺปชฺชติ นเหตุปจฺจยา. อเหตุกํ อนิยตํ เอกํ ขนฺธํ ปจฺจยา ตโย ขนฺธา จิตฺต\-สมุฏฺฐานญฺจ รูปํ ฯเปฯ เทฺว ขนฺเธ ฯเปฯ. อเหตุก\-ปฏิสนฺธิ\-กฺขเณ ฯเปฯ. ขนฺเธ ปจฺจยา วตฺถุ, วตฺถุํ ปจฺจยา ขนฺธา. เอกํ มหาภูตํ ฯเปฯ. อสญฺญสตฺตานํ ฯเปฯ. จกฺขายตนํ ปจฺจยา จกฺขุวิญฺญาณํ ฯเปฯ. กายายตนํ ปจฺจยา กายวิญฺญาณํ. วตฺถุํ ปจฺจยา อเหตุกา อนิยตา ขนฺธา. วิจิกิจฺฉา\-สหคเต อุทฺธจฺจ\-สหคเต ขนฺเธ จ วตฺถุญฺจ ปจฺจยา วิจิกิจฺฉา\-สหคโต อุทฺธจฺจ\-สหคโต โมโห. (1)}

\prose{นอารมฺมณ}

\proseitem{20}{มิจฺฉตฺต\-นิยตํ ธมฺมํ ปจฺจยา อนิยโต ธมฺโม อุปฺปชฺชติ นอารมฺมณ\-ปจฺจยา. มิจฺฉตฺต\-นิยเต ขนฺเธ ปจฺจยา จิตฺต\-สมุฏฺฐานํ รูปํ. (กุสล\-ตฺติก\-สทิสํ, ปญฺจ กาตพฺพา.)}

\csromanheader{2}{15. มิจฺฉตฺตนิยต\-ตฺติก นอธิปติ}
\proseitem{21}{\csromanfolio{367}มิจฺฉตฺต\-นิยตํ ธมฺมํ ปจฺจยา มิจฺฉตฺต\-นิยโต ธมฺโม อุปฺปชฺชติ นอธิปติ\-ปจฺจยา. มิจฺฉตฺต\-นิยเต ขนฺเธ ปจฺจยา มิจฺฉตฺต\-นิยตาธิปติ. (1)}

\prose{สมฺมตฺต\-นิยตํ ธมฺมํ ปจฺจยา สมฺมตฺตนิยโต ธมฺโม อุปฺปชฺชติ นอธิปติ\-ปจฺจยา. สมฺมตฺตนิยเต ขนฺเธ ปจฺจยา สมฺมตฺตนิยตา\-ธิปติ. (1)}

\prose{อนิยตํ ธมฺมํ ปจฺจยา อนิยโต ธมฺโม อุปฺปชฺชติ นอธิปติ\-ปจฺจยา. อนิยตํ เอกํ ขนฺธํ ปจฺจยา ตโย ขนฺธา จิตฺต\-สมุฏฺฐานญฺจ รูปํ ฯเปฯ เทฺว ขนฺเธ ฯเปฯ. ปฏิสนฺธิ\-กฺขเณ ฯเปฯ. อสญฺญสตฺตานํ ฯเปฯ. จกฺขายตนํ ปจฺจยา จกฺขุวิญฺญาณํ ฯเปฯ. กายายตนํ ปจฺจยา กายวิญฺญาณํ. วตฺถุํ ปจฺจยา อนิยตา ขนฺธา. (1)}

\prose{อนิยตํ ธมฺมํ ปจฺจยา มิจฺฉตฺต\-นิยโต ธมฺโม อุปฺปชฺชติ นอธิปติ\-ปจฺจยา. วตฺถุํ ปจฺจยา มิจฺฉตฺต\-นิยตาธิปติ. (2)}

\prose{อนิยตํ ธมฺมํ ปจฺจยา สมฺมตฺตนิยโต ธมฺโม อุปฺปชฺชติ นอธิปติ\-ปจฺจยา. วตฺถุํ ปจฺจยา สมฺมตฺตนิยตา\-ธิปติ. (3)}

\prose{มิจฺฉตฺตนิยตญฺจ อนิยตญฺจ ธมฺมํ ปจฺจยา มิจฺฉตฺต\-นิยโต ธมฺโม อุปฺปชฺชติ นอธิปติ\-ปจฺจยา. มิจฺฉตฺต\-นิยเต ขนฺเธ จ วตฺถุญฺจ ปจฺจยา มิจฺฉตฺต\-นิยตาธิปติ. (1)}

\prose{สมฺมตฺตนิยตญฺจ อนิยตญฺจ ธมฺมํ ปจฺจยา สมฺมตฺตนิยโต ธมฺโม อุปฺปชฺชติ นอธิปติ\-ปจฺจยา. สมฺมตฺตนิยเต ขนฺเธ จ วตฺถุญฺจ ปจฺจยา สมฺมตฺตนิยตา\-ธิปติ. (1)}

\prose{นอนนฺตราทิ}

\proseitem{22}{มิจฺฉตฺต\-นิยตํ ธมฺมํ ปจฺจยา อนิยโต ธมฺโม อุปฺปชฺชติ นอนนฺตร\-ปจฺจยา ฯเปฯ. โนนตฺถิปจฺจยา. โนวิคต\-ปจฺจยา.}

\csromanmatikaanchor{mtk.124}
\csromantocmarkref{subparagraph}{2. ปจฺจยปจฺจนีย 2. สงฺขฺยาวาร}{mtk.124}
\bookmark[dest={mtk.124},level=5]{2. ปจฺจยปจฺจนีย 2. สงฺขฺยาวาร}
\csromanmarkh{2. ปจฺจย\-ปจฺจนีย}\csromanmarkhii{2. สงฺขฺยาวาร}\csromanheadernomark{6}{2. ปจฺจย\-ปจฺจนีย 2. สงฺขฺยาวาร}
\proseitemwithcloserruled{23}{นเหตุยา เอกํ, นอารมฺมเณ ปญฺจ, นอธิปติยา สตฺต, นอนนฺตเร ปญฺจ, นสมนนฺตเร ปญฺจ, นอญฺญมญฺเญ ปญฺจ, นอุปนิสฺสเย ปญฺจ, นปุเรชาเต ฉ, นปจฺฉาชาเต สตฺตรส, นอาเสวเน ปญฺจ, นกมฺเม สตฺต, \csromanfolio{368}นวิปาเก สตฺตรส, นอาหาเร เอกํ, นอินฺทฺริเย, นฌาเน, นมคฺเค เอกํ, นสมฺปยุตฺเต ปญฺจ, นวิปฺปยุตฺเต เทฺว, โนนตฺถิยา ปญฺจ, โนวิคเต ปญฺจ.  (เอวํ คเณตพฺพํ.)}{ปจฺจนียํ.}
\csromancenter{\textbf{เหตุ}ปจฺจยา นอารมฺมเณ ปญฺจ, นอธิปติยา สตฺต, นอนนฺตเร ปญฺจ, นสมนนฺตเร, นอญฺญมญฺเญ, นฺเอาปนิสฺสเย ปญฺจ, นปุเรชาเต ฉ, นปจฺฉาชาเต สตฺตรส, นอาเสวเน ปญฺจ, นกมฺเม สตฺต, นวิปาเก สตฺตรส, นสมฺปยุตฺเต ปญฺจ, นวิปฺปยุตฺเต เทฺว, โนนตฺถิยา ปญฺจ, โนวิคเต ปญฺจ.  (เอวํ คเณตพฺพํ.)}
\nitthitamruled{อนุโลม\-ปจฺจนียํ.}
\proseitem{25}{นเหตุปจฺจยา อารมฺมเณ เอกํ, อนนฺตเร เอกํ. (สํขิตฺตํ.) อวิคเต เอกํ.  (เอวํ คเณตพฺพํ.) ปจฺจนียา\-นุโลมํ. ปจฺจยวาโร. (นิสฺสยวาโร ปจฺจยวาร\-สทิโส.)}
\csromansectionrule

\csromanmatikaanchor{mtk.125}
\csromanmatikaanchor{mtk.129}
\csromantocmarknopage{subparagraph}{3. ปจฺจยานุโลมปจฺจนีย}{mtk.125}
\bookmark[dest={mtk.125},level=5]{3. ปจฺจยานุโลมปจฺจนีย}
\csromanmarkh{15. มิจฺฉตฺตนิยต\-ตฺติก}\csromanmarkhii{5. สํสฏฺฐวาร}\csromanheadernomark{6}{15. มิจฺฉตฺตนิยต\-ตฺติก 5. สํสฏฺฐวาร}
\csromanmatikaanchor{mtk.126}
\csromanmatikaanchor{mtk.130}
\csromantocmarkref{subparagraph}{เหตุทุก}{mtk.126}
\bookmark[dest={mtk.126},level=5]{เหตุทุก}
\csromanmarkh{1. ปจฺจยานุโลม}\csromanmarkhii{1. วิภงฺควาร}\csromanheadernomark{6}{1. \textbf{ปจฺจยานุโลม 1. วิภงฺควาร}}
\csromanmatikaanchor{mtk.127}
\csromanmatikaanchor{mtk.128}
\csromantocmarknopage{subparagraph}{4. ปจฺจยปจฺจนียานุโลม}{mtk.127}
\bookmark[dest={mtk.127},level=5]{4. ปจฺจยปจฺจนียานุโลม}
\csromantocmarkref{subparagraph}{นเหตุทุก}{mtk.128}
\bookmark[dest={mtk.128},level=5]{นเหตุทุก}
\csromantocmarknopage{subparagraph}{15. มิจฺฉตฺตนิยตตฺติก 5. สํสฏฺฐวาร}{mtk.129}
\bookmark[dest={mtk.129},level=5]{15. มิจฺฉตฺตนิยตตฺติก 5. สํสฏฺฐวาร}
\csromantocmarkref{subparagraph}{1. ปจฺจยานุโลม 1. วิภงฺควาร}{mtk.130}
\bookmark[dest={mtk.130},level=5]{1. ปจฺจยานุโลม 1. วิภงฺควาร}
\csromanheader{6}{\textbf{เหตุ}}
\proseitem{26}{มิจฺฉตฺต\-นิยตํ ธมฺมํ สํสฏฺโฐ มิจฺฉตฺต\-นิยโต ธมฺโม อุปฺปชฺชติ เหตุปจฺจยา. มิจฺฉตฺต\-นิยตํ เอกํ ขนฺธํ สํสฏฺฐา ตโย ขนฺธา ฯเปฯ เทฺว ขนฺเธ ฯเปฯ. (1)}

\prose{\csromanfolio{369}สมฺมตฺต\-นิยตํ ธมฺมํ สํสฏฺโฐ สมฺมตฺตนิยโต ธมฺโม อุปฺปชฺชติ เหตุปจฺจยา. สมฺมตฺต\-นิยตํ เอกํ ขนฺธํ สํสฏฺฐา ตโย ขนฺธา ฯเปฯ เทฺว ขนฺเธ ฯเปฯ. (1)}

\prose{อนิยตํ ธมฺมํ สํสฏฺโฐ อนิยโต ธมฺโม อุปฺปชฺชติ เหตุปจฺจยา. อนิยตํ เอกํ ขนฺธํ สํสฏฺฐา ตโย ขนฺธา ฯเปฯ เทฺว ขนฺเธ ฯเปฯ. ปฏิสนฺธิ\-กฺขเณ ฯเปฯ. (1)}

\csromanheader{2}{\textbf{อารมฺมณาทิ}}
\proseitem{27}{มิจฺฉตฺต\-นิยตํ ธมฺมํ สํสฏฺโฐ มิจฺฉตฺต\-นิยโต ธมฺโม อุปฺปชฺชติ อารมฺมณ\-ปจฺจยา ฯเปฯ. อวิคตปจฺจยา.}

\csromanmatikaanchor{mtk.131}
\csromantocmarkref{subparagraph}{1. ปจฺจยานุโลม 2. สงฺขฺยาวาร}{mtk.131}
\bookmark[dest={mtk.131},level=5]{1. ปจฺจยานุโลม 2. สงฺขฺยาวาร}
\csromanmarkh{1. ปจฺจยานุโลม}\csromanmarkhii{2. สงฺขฺยาวาร}\csromanheadernomark{6}{1. \textbf{ปจฺจยานุโลม 2. สงฺขฺยาวาร}}
\proseitemwithcloserruled{28}{เหตุยา ตีณิ, อารมฺมเณ ตีณิ ฯเปฯ กมฺเม ตีณิ, วิปาเก เอกํ, อาหาเร ตีณิ ฯเปฯ อวิคเต ตีณิ. (เอวํ คเณตพฺพํ.)}{อนุโลมํ.}
\csromanmatikaanchor{mtk.132}
\csromantocmarkref{subparagraph}{2. ปจฺจยปจฺจนีย 1. วิภงฺควาร}{mtk.132}
\bookmark[dest={mtk.132},level=5]{2. ปจฺจยปจฺจนีย 1. วิภงฺควาร}
\csromanmarkh{2. ปจฺจย\-ปจฺจนีย}\csromanmarkhii{1. วิภงฺควาร}\csromanheadernomark{6}{2. ปจฺจย\-ปจฺจนีย 1. วิภงฺควาร}
\csromanheader{2}{\textbf{นเหตุ}}
\proseitem{29}{อนิยตํ ธมฺมํ สํสฏฺโฐ อนิยโต ธมฺโม อุปฺปชฺชติ นเหตุปจฺจยา. อเหตุกํ อนิยตํ เอกํ ขนฺธํ สํสฏฺฐา ตโย ขนฺธา ฯเปฯ เทฺว ขนฺเธ ฯเปฯ. อเหตุก\-ปฏิสนฺธิ\-กฺขเณ ฯเปฯ. (1)}

\prose{นอธิปติ}

\proseitem{30}{มิจฺฉตฺต\-นิยตํ ธมฺมํ สํสฏฺโฐ มิจฺฉตฺต\-นิยโต ธมฺโม อุปฺปชฺชติ นอธิปติ\-ปจฺจยา. มิจฺฉตฺต\-นิยเต ขนฺเธ สํสฏฺโฐ มิจฺฉตฺต\-นิยตาธิปติ. (1)}

\prose{สมฺมตฺต\-นิยตํ ธมฺมํ สํสฏฺโฐ สมฺมตฺตนิยโต ธมฺโม อุปฺปชฺชติ นอธิปติ\-ปจฺจยา. สมฺมตฺตนิยเต ขนฺเธ สํสฏฺโฐ สมฺมตฺตนิยตา\-ธิปติ. (1)}

\prose{อนิยตํ ธมฺมํ สํสฏฺโฐ อนิยโต ธมฺโม อุปฺปชฺชติ นอธิปติ\-ปจฺจยา. อนิยตํ เอกํ ขนฺธํ สํสฏฺฐา ตโย ขนฺธา ฯเปฯ เทฺว ขนฺเธ ฯเปฯ. ปฏิสนฺธิ\-กฺขเณ ฯเปฯ. (1)}

\csromanheader{2}{\textbf{นปุเรชาตาทิ}}
\proseitem{31}{\csromanfolio{370}สมฺมตฺต\-นิยตํ ธมฺมํ สํสฏฺโฐ สมฺมตฺตนิยโต ธมฺโม อุปฺปชฺชติ นปุเรชาต\-ปจฺจยา. อรูเป สมฺมตฺต\-นิยตํ เอกํ ขนฺธํ สํสฏฺฐา ตโย ขนฺธา ฯเปฯ เทฺว ขนฺเธ ฯเปฯ. (1)}

\prose{อนิยตํ ธมฺมํ สํสฏฺโฐ อนิยโต ธมฺโม อุปฺปชฺชติ นปุเรชาต\-ปจฺจยา. อรูเป อนิยตํ เอกํ ขนฺธํ สํสฏฺฐา ตโย ขนฺธา ฯเปฯ เทฺว ขนฺเธ ฯเปฯ. ปฏิสนฺธิ\-กฺขเณ ฯเปฯ. (1)}

\prose{มิจฺฉตฺต\-นิยตํ ธมฺมํ สํสฏฺโฐ มิจฺฉตฺต\-นิยโต ธมฺโม อุปฺปชฺชติ นปจฺฉาชาต\-ปจฺจยา. (ปริปุณฺณํ.)}

\prose{นอาเสวนาทิ}

\proseitem{32}{อนิยตํ ธมฺมํ สํสฏฺโฐ อนิยโต ธมฺโม อุปฺปชฺชติ นอาเสวน\-ปจฺจยา. อนิยตํ เอกํ ขนฺธํ สํสฏฺฐา ตโย ขนฺธา ฯเปฯ เทฺว ขนฺเธ ฯเปฯ. ปฏิสนฺธิ\-กฺขเณ ฯเปฯ. (1)}

\prose{มิจฺฉตฺต\-นิยตํ ธมฺมํ สํสฏฺโฐ มิจฺฉตฺต\-นิยโต ธมฺโม อุปฺปชฺชติ นกมฺมปจฺจยา. นวิปาก\-ปจฺจยา. (สํขิตฺตํ.)}

\prose{อนิยตํ ธมฺมํ สํสฏฺโฐ อนิยโต ธมฺโม อุปฺปชฺชติ นฌานปจฺจยา. ปญฺจวิญฺญาณํ ฯเปฯ. นมคฺคปจฺจยา. อเหตุกํ อนิยตํ ฯเปฯ.}

\prose{สมฺมตฺต\-นิยตํ ธมฺมํ สํสฏฺโฐ สมฺมตฺตนิยโต ธมฺโม อุปฺปชฺชติ นวิปฺปยุตฺต\-ปจฺจยา. อรูเป สมฺมตฺต\-นิยตํ เอกํ ขนฺธํ สํสฏฺฐา ตโย ขนฺธา ฯเปฯ เทฺว ขนฺเธ ฯเปฯ. (1)}

\prose{อนิยตํ ธมฺมํ สํสฏฺโฐ อนิยโต ธมฺโม อุปฺปชฺชติ นวิปฺปยุตฺต\-ปจฺจยา. อรูเป อนิยตํ เอกํ ขนฺธํ สํสฏฺฐา ตโย ขนฺธา ฯเปฯ เทฺว ขนฺเธ ฯเปฯ. (1)}

\csromanmatikaanchor{mtk.133}
\csromantocmarkref{subparagraph}{2. ปจฺจยปจฺจนีย 2. สงฺขฺยาวาร}{mtk.133}
\bookmark[dest={mtk.133},level=5]{2. ปจฺจยปจฺจนีย 2. สงฺขฺยาวาร}
\csromanmarkh{2. ปจฺจย\-ปจฺจนีย}\csromanmarkhii{2. สงฺขฺยาวาร}\csromanheadernomark{6}{2. ปจฺจย\-ปจฺจนีย 2. สงฺขฺยาวาร}
\proseitem{33}{นเหตุยา เอกํ, นอธิปติยา ตีณิ, นปุเรชาเต เทฺว, นปจฺฉาชาเต ตีณิ, นอาเสวเน เอกํ, นกมฺเม ตีณิ, นวิปาเก ตีณิ, นฌาเน เอกํ, นมคฺเค เอกํ, นวิปฺปยุตฺเต เทฺว. (เอวํ คเณตพฺพํ.)}

\vspace{\tipitakaparskip}
\csromancenter{ปจฺจนียํ.}
\csromanmatikaanchor{mtk.134}
\csromantocmarkref{subparagraph}{3. ปจฺจยานุโลมปจฺจนีย}{mtk.134}
\bookmark[dest={mtk.134},level=5]{3. ปจฺจยานุโลมปจฺจนีย}
\csromanmarkh{15. มิจฺฉตฺตนิยต\-ตฺติก}\csromanmarkhii{3. ปจฺจยา\-นุโลม\-ปจฺจนีย}\csromanheadernomark{6}{15. มิจฺฉตฺตนิยต\-ตฺติก 3. ปจฺจยา\-นุโลม\-ปจฺจนีย}
\vspace{\tipitakaparskip}
\csromancenter{\textbf{เหตุ}ปจฺจยา นอธิปติยา ตีณิ, นปุเรชาเต เทฺว, นปจฺฉาชาเต ตีณิ, นอาเสวเน เอกํ, นกมฺเม ตีณิ, นวิปาเก ตีณิ, นวิปฺปยุตฺเต เทฺว. (เอวํ คเณตพฺพํ.)}
\nitthitamruled{อนุโลม\-ปจฺจนียํ.}
\csromanmatikaanchor{mtk.135}
\csromantocmarkref{subparagraph}{4. ปจฺจยปจฺจนียานุโลม}{mtk.135}
\bookmark[dest={mtk.135},level=5]{4. ปจฺจยปจฺจนียานุโลม}
\proseitem{35}{\csromanfolio{371}นเหตุปจฺจยา อารมฺมเณ เอกํ, อนนฺตเร เอกํ. (สํขิตฺตํ.) อวิคเต เอกํ. (เอวํ คเณตพฺพํ.) ปจฺจนียา\-นุโลมํ. สํสฏฺฐวาโร. (สมฺปยุตฺตวาโร สํสฏฺฐวาร\-สทิโส.)}
\csromansectionrule

\csromanmatikaanchor{mtk.136}
\csromantocmarknopage{subparagraph}{15. มิจฺฉตฺตนิยตตฺติก 7. ปญฺหาวาร}{mtk.136}
\bookmark[dest={mtk.136},level=5]{15. มิจฺฉตฺตนิยตตฺติก 7. ปญฺหาวาร}
\csromantocmarknopage{subparagraph}{1. ปจฺจยนุโลม 1. วิภงฺควาร}{mtk.137}
\bookmark[dest={mtk.137},level=5]{1. ปจฺจยนุโลม 1. วิภงฺควาร}
\csromanmarkh{15. มิจฺฉตฺตนิยต\-ตฺติก}\csromanmarkhii{7. ปญฺหาวาร}\csromanheadernomark{6}{15. มิจฺฉตฺตนิยต\-ตฺติก 7. ปญฺหาวาร}
\csromanmarkh{1. ปจฺจยานุโลม}\csromanmarkhii{1. วิภงฺควาร}\csromanheadernomark{2}{1. \textbf{ปจฺจยานุโลม 1. วิภงฺควาร}}
\csromanmatikaanchor{mtk.138}
\csromantocmarkref{subparagraph}{เหตุ}{mtk.138}
\bookmark[dest={mtk.138},level=5]{เหตุ}
\csromanheader{6}{\textbf{เหตุ}}
\proseitem{36}{มิจฺฉตฺต\-นิยโต ธมฺโม มิจฺฉตฺต\-นิยตสฺส ธมฺมสฺส เหตุปจฺจเยน ปจฺจโย. มิจฺฉตฺต\-นิยตา เหตู สมฺปยุตฺตกานํ ขนฺธานํ เหตุปจฺจเยน ปจฺจโย. (1)}

\prose{มิจฺฉตฺต\-นิยโต ธมฺโม อนิยตสฺส ธมฺมสฺส เหตุปจฺจเยน ปจฺจโย. มิจฺฉตฺต\-นิยตา เหตู จิตฺต\-สมุฏฺฐานานํ รูปานํ เหตุปจฺจเยน ปจฺจโย. (2)}

\prose{มิจฺฉตฺต\-นิยโต ธมฺโม มิจฺฉตฺต\-นิยตสฺส จ อนิยตสฺส จ ธมฺมสฺส เหตุปจฺจเยน ปจฺจโย. มิจฺฉตฺต\-นิยตา เหตู สมฺปยุตฺตกานํ ขนฺธานํ จิตฺตสมุฏฺฐานานญฺจ รูปานํ เหตุปจฺจเยน ปจฺจโย. (3)}

\prose{\csromanfolio{372}สมฺมตฺตนิยโต ธมฺโม สมฺมตฺต\-นิยตสฺส ธมฺมสฺส เหตุปจฺจเยน ปจฺจโย. ตีณิ.}

\prose{อนิยโต ธมฺโม อนิยตสฺส ธมฺมสฺส เหตุปจฺจเยน ปจฺจโย. อนิยตา เหตู สมฺปยุตฺตกานํ ขนฺธานํ จิตฺตสมุฏฺฐานานญฺจ รูปานํ เหตุปจฺจเยน ปจฺจโย. ปฏิสนฺธิ\-กฺขเณ ฯเปฯ. (1)}

\csromanmatikaanchor{mtk.139}
\csromantocmarkref{subparagraph}{อารมฺมณ}{mtk.139}
\bookmark[dest={mtk.139},level=5]{อารมฺมณ}
\csromanheader{6}{\textbf{อารมฺมณ}}
\proseitem{37}{มิจฺฉตฺต\-นิยโต ธมฺโม อนิยตสฺส ธมฺมสฺส อารมฺมณ\-ปจฺจเยน ปจฺจโย. อริยา มิจฺฉตฺต\-นิยเต ปหีเน กิเลเส ปจฺจเวกฺขนฺติ, ปุพฺเพ สมุทาจิณฺเณ กิเลเส ชานนฺติ. มิจฺฉตฺต\-นิยเต ขนฺเธ อนิจฺจโต ฯเปฯ วิปสฺสนฺติ. เจโตปริย\-ญาเณน มิจฺฉตฺตนิยต\-จิตฺต\-สมงฺคิสฺส จิตฺตํ ชานนฺติ. มิจฺฉตฺต\-นิยตา ขนฺธา เจโตปริย\-ญาณสฺส, ปุพฺเพ\-นิวาสา\-นุสฺสติ\-ญาณสฺส, ยถากมฺมูปคญาณสฺส, อนาคตํสญาณสฺส, อาวชฺชนาย อารมฺมณ\-ปจฺจเยน ปจฺจโย. (1)}

\prose{สมฺมตฺตนิยโต ธมฺโม อนิยตสฺส ธมฺมสฺส อารมฺมณ\-ปจฺจเยน ปจฺจโย. อริยา มคฺคา วุฏฺฐหิตฺวา มคฺคํ ปจฺจเวกฺขนฺติ. เจโตปริย\-ญาเณน สมฺมตฺตนิยต\-จิตฺต\-สมงฺคิสฺส จิตฺตํ ชานนฺติ. สมฺมตฺตนิยตา ขนฺธา เจโตปริย\-ญาณสฺส,ปุพฺเพ\-นิวาสา\-นุสฺสติ\-ญาณสฺส, อนาคตํสญาณสฺส, อาวชฺชนาย อารมฺมณ\-ปจฺจเยน ปจฺจโย. (1)}

\csromanmatikaanchor{mtk.140}
\csromantocmarkref{subparagraph}{อธิปติ}{mtk.140}
\bookmark[dest={mtk.140},level=5]{อธิปติ}
\proseitem{38}{อนิยโต ธมฺโม อนิยตสฺส ธมฺมสฺส อารมฺมณ\-ปจฺจเยน ปจฺจโย. ทานํ ทตฺวา, สีลํ สมาทิยิตฺวา, อุโปสถกมฺมํ กตฺวา ตํ ปจฺจเวกฺขติ, ปุพฺเพ สุจิณฺณานิ ปจฺจเวกฺขติ. ฌานา วุฏฺฐหิตฺวา ฌานํ ปจฺจเวกฺขติ. อริยา ผลํ ปจฺจเวกฺขนฺติ, นิพฺพานํ ปจฺจเวกฺขนฺติ. นิพฺพานํ โคตฺรภุสฺส, โวทานสฺส, ผลสฺส, อาวชฺชนาย อารมฺมณ\-ปจฺจเยน ปจฺจโย. อริยา อนิยเต ปหีเน กิเลเส ปจฺจเวกฺขนฺติ, วิกฺขมฺภิเต กิเลเส ปจฺจเวกฺขนฺติ. ปุพฺเพ สมุทาจิณฺเณ กิเลเส ชานนฺติ. จกฺขุํ ฯเปฯ. วตฺถุํ อนิยเต ขนฺเธ อนิจฺจโต ทุกฺขโต อนตฺตโต วิปสฺสนฺติ, อสฺสาเทนฺติ อภินนฺทนฺติ, ตํ อารพฺภ อนิยโต ราโค อุปฺปชฺชติ ฯเปฯ โทมนสฺสํ อุปฺปชฺชติ. ทิพฺเพน จกฺขุนา รูปํ ปสฺสติ. ทิพฺพาย โสตธาตุยา สทฺทํ สุณาติ. \csromanfolio{373}เจโตปริย\-ญาเณน อนิยต\-จิตฺต\-สมงฺคิสฺส จิตฺตํ ชานาติ.  อากาสา\-นญฺจา\-ยตนํ วิญฺญาณญฺจา\-ยตนสฺส ฯเปฯ. อากิญฺจญฺญา\-ยตนํ เนว\-สญฺญา\-นา\-สญฺญา\-ยตนสฺส อารมฺมณ\-ปจฺจเยน ปจฺจโย. รูปายตนํ จกฺขุ\-วิญฺญาณสฺส ฯเปฯ. โผฏฺฐพฺพา\-ยตนํ กายวิญฺญาณสฺส อารมฺมณ\-ปจฺจเยน ปจฺจโย. อนิยตา ขนฺธา อิทฺธิวิธ\-ญาณสฺส, เจโตปริย\-ญาณสฺส,ปุพฺเพ\-นิวาสา\-นุสฺสติ\-ญาณสฺส, ยถากมฺมูปคญาณสฺส, อนาคตํสญาณสฺส, อาวชฺชนาย อารมฺมณ\-ปจฺจเยน ปจฺจโย. (1)}

\prose{อนิยโต ธมฺโม มิจฺฉตฺต\-นิยตสฺส ธมฺมสฺส อารมฺมณ\-ปจฺจเยน ปจฺจโย. รูปชีวิตินฺทฺริยํ มาตุ\-ฆาติ\-กมฺมสฺส. ปิตุฆาติ\-กมฺมสฺส. อรหนฺต\-ฆาติ\-กมฺมสฺส. รุหิรุ\-ปฺปาท\-กมฺมสฺส อารมฺมณ\-ปจฺจเยน ปจฺจโย. ยํ วตฺถุํ ปรามสนฺตสฺส มิจฺฉตฺต\-นิยตา ขนฺธา อุปฺปชฺชนฺติ, ตํ วตฺถุ มิจฺฉตฺต\-นิยตานํ ขนฺธานํ อารมฺมณ\-ปจฺจเยน ปจฺจโย. (2)}

\prose{อนิยโต ธมฺโม สมฺมตฺต\-นิยตสฺส ธมฺมสฺส อารมฺมณ\-ปจฺจเยน ปจฺจโย. นิพฺพานํ มคฺคสฺส อารมฺมณ\-ปจฺจเยน ปจฺจโย. (3)}

\proseitem{39}{มิจฺฉตฺต\-นิยโต ธมฺโม มิจฺฉตฺต\-นิยตสฺส ธมฺมสฺส อธิปติ\-ปจฺจเยน ปจฺจโย. สหชาตา\-ธิปติ\csromandash{}มิจฺฉตฺต\-นิยตาธิปติ สมฺปยุตฺตกานํ ขนฺธานํ อธิปติ\-ปจฺจเยน ปจฺจโย. (1)}

\prose{มิจฺฉตฺต\-นิยโต ธมฺโม อนิยตสฺส ธมฺมสฺส อธิปติ\-ปจฺจเยน ปจฺจโย. สหชาตา\-ธิปติ\csromandash{}มิจฺฉตฺต\-นิยตาธิปติ จิตฺต\-สมุฏฺฐานานํ รูปานํ อธิปติ\-ปจฺจเยน ปจฺจโย. (2)}

\prose{มิจฺฉตฺต\-นิยโต ธมฺโม มิจฺฉตฺต\-นิยตสฺส จ อนิยตสฺส จ ธมฺมสฺส อธิปติ\-ปจฺจเยน ปจฺจโย. สหชาตา\-ธิปติ\csromandash{} มิจฺฉตฺต\-นิยตาธิปติ สมฺปยุตฺตกานํ ขนฺธานํ จิตฺตสมุฏฺฐานานญฺจ รูปานํ อธิปติ\-ปจฺจเยน ปจฺจโย. (3)}

\proseitem{40}{สมฺมตฺตนิยโต ธมฺโม สมฺมตฺต\-นิยตสฺส ธมฺมสฺส อธิปติ\-ปจฺจเยน ปจฺจโย. สหชาตา\-ธิปติ\csromandash{}สมฺมตฺตนิยตา\-ธิปติ สมฺปยุตฺตกานํ ขนฺธานํ อธิปติ\-ปจฺจเยน ปจฺจโย. (1)}

\prose{\csromanfolio{374}สมฺมตฺตนิยโต ธมฺโม อนิยตสฺส ธมฺมสฺส อธิปติ\-ปจฺจเยน ปจฺจโย. อารมฺมณา\-ธิปติ, สหชาตา\-ธิปติ.  อารมฺมณา\-ธิปติ\csromandash{}อริยา มคฺคา วุฏฺฐหิตฺวา มคฺคํ ครุํ กตฺวา ปจฺจเวกฺขนฺติ. สหชาตา\-ธิปติ\csromandash{} สมฺมตฺตนิยตา\-ธิปติ จิตฺต\-สมุฏฺฐานานํ รูปานํ อธิปติ\-ปจฺจเยน ปจฺจโย. (2)}

\prose{สมฺมตฺตนิยโต ธมฺโม สมฺมตฺต\-นิยตสฺส จ อนิยตสฺส จ ธมฺมสฺส อธิปติ\-ปจฺจเยน ปจฺจโย. สหชาตา\-ธิปติ\csromandash{} สมฺมตฺตนิยตา\-ธิปติ สมฺปยุตฺตกานํ ขนฺธานํ จิตฺตสมุฏฺฐานานญฺจ รูปานํ อธิปติ\-ปจฺจเยน ปจฺจโย. (3)}

\proseitem{41}{อนิยโต ธมฺโม อนิยตสฺส ธมฺมสฺส อธิปติ\-ปจฺจเยน ปจฺจโย. อารมฺมณา\-ธิปติ, สหชาตา\-ธิปติ.  อารมฺมณา\-ธิปติ\csromandash{}ทานํ ทตฺวา, สีลํ สมาทิยิตฺวา, อุโปสถกมฺมํ กตฺวา ตํ ครุํ กตฺวา ปจฺจเวกฺขติ, ปุพฺเพ สุจิณฺณานิ ครุํ กตฺวา ปจฺจเวกฺขติ. ฌานา ฯเปฯ. อริยา ผลํ ครุํ กตฺวา ปจฺจเวกฺขนฺติ, นิพฺพานํ ครุํ กตฺวา ปจฺจเวกฺขนฺติ. นิพฺพานํ โคตฺรภุสฺส, โวทานสฺส, ผลสฺส อธิปติ\-ปจฺจเยน ปจฺจโย. จกฺขุํ ฯเปฯ. วตฺถุํ ฯเปฯ. อนิยเต ขนฺเธ ครุํ กตฺวา อสฺสาเทติ อภินนฺทติ, ตํ ครุํ กตฺวา อนิยโต ราโค อุปฺปชฺชติ, ทิฏฺฐิ อุปฺปชฺชติ. สหชาตา\-ธิปติ\csromandash{}อนิยตาธิปติ สมฺปยุตฺตกานํ ขนฺธานํ จิตฺตสมุฏฺฐานานญฺจ รูปานํ อธิปติ\-ปจฺจเยน ปจฺจโย. (1)}

\prose{อนิยโต ธมฺโม สมฺมตฺต\-นิยตสฺส ธมฺมสฺส อธิปติ\-ปจฺจเยน ปจฺจโย. อารมฺมณา\-ธิปติ\csromandash{}นิพฺพานํ มคฺคสฺส อธิปติ\-ปจฺจเยน ปจฺจโย. (2)}

\csromanmatikaanchor{mtk.141}
\csromantocmarkref{subparagraph}{อนนฺตร}{mtk.141}
\bookmark[dest={mtk.141},level=5]{อนนฺตร}
\csromanheader{6}{\textbf{อนนฺตร}}
\proseitem{42}{มิจฺฉตฺต\-นิยโต ธมฺโม อนิยตสฺส ธมฺมสฺส อนนฺตร\-ปจฺจเยน ปจฺจโย. มิจฺฉตฺต\-นิยตา ขนฺธา วุฏฺฐานสฺส อนนฺตร\-ปจฺจเยน ปจฺจโย. (1)}

\prose{สมฺมตฺตนิยโต ธมฺโม อนิยตสฺส ธมฺมสฺส อนนฺตร\-ปจฺจเยน ปจฺจโย. มคฺโค ผลสฺส อนนฺตร\-ปจฺจเยน ปจฺจโย. (1)}

\prose{อนิยโต ธมฺโม อนิยตสฺส ธมฺมสฺส อนนฺตร\-ปจฺจเยน ปจฺจโย. ปุริมา ปุริมา อนิยตา ขนฺธา ปจฺฉิมานํ ปจฺฉิมานํ อนิยตานํ ขนฺธานํ อนนฺตร\-ปจฺจเยน ปจฺจโย. อนุโลมํ โคตฺรภุสฺส. อนุโลมํ โวทานสฺส. ผลํ ผลสฺส. อนุโลมํ ผลสมาปตฺติยา. นิโรธา วุฏฺฐหนฺตสฺส เนว\-สญฺญา\-นา\-สญฺญา\-ยตนํ ผลสมาปตฺติยา อนนฺตร\-ปจฺจเยน ปจฺจโย. (1)}

\proseitem{43}{\csromanfolio{375}อนิยโต ธมฺโม มิจฺฉตฺต\-นิยตสฺส ธมฺมสฺส อนนฺตร\-ปจฺจเยน ปจฺจโย. อนิยตํ โทมนสฺสํ มิจฺฉตฺต\-นิยตสฺส โทมนสฺสสฺส อนนฺตร\-ปจฺจเยน ปจฺจโย. อนิยต\-มิจฺฉาทิฏฺฐิ นิยต\-มิจฺฉาทิฏฺฐิยา อนนฺตร\-ปจฺจเยน ปจฺจโย. (2)}

\prose{อนิยโต ธมฺโม สมฺมตฺต\-นิยตสฺส ธมฺมสฺส อนนฺตร\-ปจฺจเยน ปจฺจโย. โคตฺรภุ มคฺคสฺส. โวทานํ มคฺคสฺส อนนฺตร\-ปจฺจเยน ปจฺจโย. (3)}

\csromanmatikaanchor{mtk.142}
\csromantocmarkref{subparagraph}{สมนนฺตราทิ}{mtk.142}
\bookmark[dest={mtk.142},level=5]{สมนนฺตราทิ}
\csromanheader{6}{\textbf{สมนนฺตราทิ}}
\proseitem{44}{มิจฺฉตฺต\-นิยโต ธมฺโม อนิยตสฺส ธมฺมสฺส สมนนฺตร\-ปจฺจเยน ปจฺจโย. (อนนฺตร\-สทิสํ.) สหชาต\-ปจฺจเยน ปจฺจโย. (ปฏิจฺจ\-วาร\-สทิสํ. นว ปญฺหา.) อญฺญมญฺญ\-ปจฺจเยน ปจฺจโย. (ปฏิจฺจ\-วาร\-สทิสํ. ติสฺโส ปญฺหา.) นิสฺสย\-ปจฺจเยน ปจฺจโย. (กุสล\-ตฺติก\-สทิสา. เตรส ปญฺหา.)}

\csromanheader{2}{\textbf{อุปนิสฺสย}}
\proseitem{45}{มิจฺฉตฺต\-นิยโต ธมฺโม มิจฺฉตฺต\-นิยตสฺส ธมฺมสฺส อุปนิสฺสย\-ปจฺจเยน ปจฺจโย. ปกตู\-ปนิสฺสโย\csromandash{}มาตุ\-ฆาติ\-กมฺมํ มาตุ\-ฆาติ\-กมฺมสฺส อุปนิสฺสย\-ปจฺจเยน ปจฺจโย. มาตุ\-ฆาติ\-กมฺมํ ฯเปฯ. ปิตุฆาติ\-กมฺมํ ฯเปฯ. อรหนฺต\-ฆาติ\-กมฺมํ ฯเปฯ. รุหิรุปฺปาทกมฺมํ ฯเปฯ. สํฆเภท\-กมฺมํ ฯเปฯ. นิยต\-มิจฺฉาทิฏฺฐิยา อุปนิสฺสย\-ปจฺจเยน ปจฺจโย. (จกฺกํ กาตพฺพํ.) นิยต\-มิจฺฉาทิฏฺฐิ นิยต\-มิจฺฉาทิฏฺฐิยา อุปนิสฺสย\-ปจฺจเยน ปจฺจโย. นิยต\-มิจฺฉาทิฏฺฐิ มาตุ\-ฆาติ\-กมฺมสฺส ฯเปฯ สํฆเภทกมฺมสฺส อุปนิสฺสย\-ปจฺจเยน ปจฺจโย. (1)}

\prose{มิจฺฉตฺต\-นิยโต ธมฺโม อนิยตสฺส ธมฺมสฺส อุปนิสฺสย\-ปจฺจเยน ปจฺจโย. อนนฺตรู\-ปนิสฺสโย, ปกตู\-ปนิสฺสโย ฯเปฯ. ปกตู\-ปนิสฺสโย\csromandash{} มาตรํ ชีวิตา โวโรเปตฺวา ตสฺส ปฏิฆาต\-ตฺถาย ทานํ เทติ, สีลํ สมาทิยติ, อุโปสถกมฺมํ กโรติ. ปิตรํ ชีวิตา โวโรเปตฺวา ฯเปฯ. อรหนฺตํ ชีวิตา โวโรเปตฺวา ฯเปฯ ทุฏฺเฐน จิตฺเตน ตถาคตสฺส โลหิตํ อุปฺปาเทตฺวา ฯเปฯ สํฆํ ภินฺทิตฺวา ตสฺส ปฏิฆาต\-ตฺถาย ทานํ เทติ, สีลํ สมาทิยติ, อุโปสถกมฺมํ กโรติ. (2)}

\proseitem{46}{สมฺมตฺตนิยโต ธมฺโม สมฺมตฺต\-นิยตสฺส ธมฺมสฺส อุปนิสฺสย\-ปจฺจเยน ปจฺจโย. ปกตู\-ปนิสฺสโย\csromandash{}ปฐโม มคฺโค ทุติยสฺส มคฺคสฺส \csromanfolio{376}อุปนิสฺสย\-ปจฺจเยน ปจฺจโย ฯเปฯ. ตติโย มคฺโค จตุตฺถสฺส มคฺคสฺส อุปนิสฺสย\-ปจฺจเยน ปจฺจโย. (1)}

\prose{สมฺมตฺตนิยโต ธมฺโม อนิยตสฺส ธมฺมสฺส อุปนิสฺสย\-ปจฺจเยน ปจฺจโย. อารมฺมณู\-ปนิสฺสโย, อนนฺตรู\-ปนิสฺสโย, ปกตู\-ปนิสฺสโย ฯเปฯ. ปกตู\-ปนิสฺสโย\csromandash{}อริยา มคฺคํ อุปนิสฺสาย อนุปฺปนฺนํ สมาปตฺติํ อุปฺปาเทนฺติ, อุปฺปนฺนํ สมาปชฺชนฺติ, สงฺขาเร อนิจฺจโต ทุกฺขโต อนตฺตโต วิปสฺสนฺติ. มคฺโค อริยานํ อตฺถ\-ปฏิสมฺภิทาย ฯเปฯ ฐานาฐาน\-โกสลฺลสฺส อุปนิสฺสย\-ปจฺจเยน ปจฺจโย. มคฺโค ผลสมาปตฺติยา อุปนิสฺสย\-ปจฺจเยน ปจฺจโย. (2)}

\proseitem{47}{อนิยโต ธมฺโม อนิยตสฺส ธมฺมสฺส อุปนิสฺสย\-ปจฺจเยน ปจฺจโย. อารมฺมณู\-ปนิสฺสโย, อนนฺตรู\-ปนิสฺสโย, ปกตู\-ปนิสฺสโย ฯเปฯ. ปกตู\-ปนิสฺสโย\csromandash{}อนิยตํ สทฺธํ อุปนิสฺสาย ทานํ เทติ. สีลํ สมาทิยติ, อุโปสถกมฺมํ ฯเปฯ ฌานํ อุปฺปาเทติ, วิปสฺสนํ, อภิญฺญํ. สมาปตฺติํ อุปฺปาเทติ, มานํ ชปฺเปติ, ทิฏฺฐิํ คณฺหาติ. อนิยตํ สีลํ. สุตํ. จาคํ. ปญฺญํ. ราคํ ฯเปฯ ปตฺถนํ. กายิกํ สุขํ. กายิกํ ทุกฺขํ. อุตุํ. โภชนํ. เสนาสนํ อุปนิสฺสาย ทานํ เทติ ฯเปฯ นิคมฆาตํ กโรติ. อนิยตา สทฺธา ฯเปฯ ปญฺญา. ราโค ฯเปฯ. เสนาสนํ อนิยตาย สทฺธาย ฯเปฯ กายิกสฺส สุขสฺส, กายิกสฺส ทุกฺขสฺส, ผลสมาปตฺติยา อุปนิสฺสย\-ปจฺจเยน ปจฺจโย. ปฐมสฺส ฌานสฺส ปริกมฺมํ ตสฺเสว ฯเปฯ เนว\-สญฺญา\-นา\-สญฺญา\-ยตนสฺส ปริกมฺมํ ตสฺเสว ฯเปฯ. ปฐมํ ฌานํ ทุติยสฺส ฌานสฺส ฯเปฯ. อากิญฺจญฺญา\-ยตนํ เนว\-สญฺญา\-นา\-สญฺญา\-ยตนสฺส ฯเปฯ. ปาณาติปาโต ปาณาติปาตสฺส อุปนิสฺสย\-ปจฺจเยน ปจฺจโย. (จกฺกํ กาตพฺพํ.) (1)}

\prose{อนิยโต ธมฺโม มิจฺฉตฺต\-นิยตสฺส ธมฺมสฺส อุปนิสฺสย\-ปจฺจเยน ปจฺจโย. อนนฺตรู\-ปนิสฺสโย, ปกตู\-ปนิสฺสโย ฯเปฯ. ปกตู\-ปนิสฺสโย\csromandash{} อนิยตํ ราคํ อุปนิสฺสาย มาตรํ ชีวิตา โวโรเปติ ฯเปฯ สํฆํ ภินฺทติ. อนิยตํ โทสํ ฯเปฯ. ปตฺถนํ. กายิกํ สุขํ ฯเปฯ. เสนาสนํ อุปนิสฺสาย มาตรํ ชีวิตา โวโรเปติ ฯเปฯ สํฆํ ภินฺทติ. อนิยโต ราโค ฯเปฯ. เสนาสนํ มาตุ\-ฆาติ\-กมฺมสฺส. ปิตุฆาติ\-กมฺมสฺส. อรหนฺต\-ฆาติ\-กมฺมสฺส. รุหิรุ\-ปฺปาท\-กมฺมสฺส. สํฆเภทกมฺมสฺส. นิยต\-มิจฺฉาทิฏฺฐิยา อุปนิสฺสย\-ปจฺจเยน ปจฺจโย. (2)}

\prose{\csromanfolio{377}อนิยโต ธมฺโม สมฺมตฺต\-นิยตสฺส ธมฺมสฺส อุปนิสฺสย\-ปจฺจเยน ปจฺจโย. อารมฺมณู\-ปนิสฺสโย, อนนฺตรู\-ปนิสฺสโย, ปกตู\-ปนิสฺสโย ฯเปฯ. ปกตู\-ปนิสฺสโย\csromandash{}ปฐมสฺส มคฺคสฺส ปริกมฺมํ ปฐมสฺส มคฺคสฺส ฯเปฯ. จตุตฺถสฺส มคฺคสฺส ปริกมฺมํ จตุตฺถสฺส มคฺคสฺส อุปนิสฺสย\-ปจฺจเยน ปจฺจโย. (3)}

\csromanheader{2}{\textbf{ปุเรชาต}}
\proseitem{48}{อนิยโต ธมฺโม อนิยตสฺส ธมฺมสฺส ปุเรชาต\-ปจฺจเยน ปจฺจโย. อารมฺมณ\-ปุเรชาตํ, วตฺถุ\-ปุเรชาตํ.  อารมฺมณ\-ปุเรชาตํ จกฺขุํ ฯเปฯ. วตฺถุํ อนิจฺจโต ฯเปฯ วิปสฺสติ, อสฺสาเทติ อภินนฺทติ, ตํ อารพฺภ อนิยโต ราโค อุปฺปชฺชติ ฯเปฯ โทมนสฺสํ อุปฺปชฺชติ. ทิพฺเพน จกฺขุนา รูปํ ปสฺสติ. ทิพฺพาย โสตธาตุยา สทฺทํ สุณาติ. รูปายตนํ จกฺขุ\-วิญฺญาณสฺส ฯเปฯ. โผฏฺฐพฺพา\-ยตนํ กายวิญฺญาณสฺส ปุเรชาต\-ปจฺจเยน ปจฺจโย. วตฺถุ\-ปุเรชาตํ จกฺขายตนํ จกฺขุ\-วิญฺญาณสฺส ฯเปฯ. กายายตนํ กายวิญฺญาณสฺส. วตฺถุ อนิยตานํ ขนฺธานํ ปุเรชาต\-ปจฺจเยน ปจฺจโย. (1)}

\prose{อนิยโต ธมฺโม มิจฺฉตฺต\-นิยตสฺส ธมฺมสฺส ปุเรชาต\-ปจฺจเยน ปจฺจโย. อารมฺมณ\-ปุเรชาตํ, วตฺถุ\-ปุเรชาตํ.  อารมฺมณ\-ปุเรชาตํ รูปชีวิตินฺทฺริยํ มาตุ\-ฆาติ\-กมฺมสฺส. ปิตุฆาติ\-กมฺมสฺส. อรหนฺต\-ฆาติ\-กมฺมสฺส. รุหิรุ\-ปฺปาท\-กมฺมสฺส ปุเรชาต\-ปจฺจเยน ปจฺจโย. วตฺถุ\-ปุเรชาตํ วตฺถุ มิจฺฉตฺต\-นิยตานํ ขนฺธานํ ปุเรชาต\-ปจฺจเยน ปจฺจโย. (2)}

\prose{อนิยโต ธมฺโม สมฺมตฺต\-นิยตสฺส ธมฺมสฺส ปุเรชาต\-ปจฺจเยน ปจฺจโย. วตฺถุ\-ปุเรชาตํ วตฺถุ สมฺมตฺต\-นิยตานํ ขนฺธานํ ปุเรชาต\-ปจฺจเยน ปจฺจโย. (3)}

\csromanheader{2}{\textbf{ปจฺฉาชาต}}
\proseitem{49}{มิจฺฉตฺต\-นิยโต ธมฺโม อนิยตสฺส ธมฺมสฺส ปจฺฉาชาต\-ปจฺจเยน ปจฺจโย. ปจฺฉาชาตา มิจฺฉตฺต\-นิยตา ขนฺธา ปุเรชาตสฺส อิมสฺส กายสฺส ปจฺฉาชาต\-ปจฺจเยน ปจฺจโย. (1)}

\prose{สมฺมตฺตนิยโต ธมฺโม อนิยตสฺส ธมฺมสฺส ปจฺฉาชาต\-ปจฺจเยน ปจฺจโย. ปจฺฉาชาตา สมฺมตฺตนิยตา ขนฺธา ปุเรชาตสฺส อิมสฺส กายสฺส ปจฺฉาชาต\-ปจฺจเยน ปจฺจโย. (1)}

\prose{\csromanfolio{378}อนิยโต ธมฺโม อนิยตสฺส ธมฺมสฺส ปจฺฉาชาต\-ปจฺจเยน ปจฺจโย. ปจฺฉาชาตา อนิยตา ขนฺธา ปุเรชาตสฺส อิมสฺส กายสฺส ปจฺฉาชาต\-ปจฺจเยน ปจฺจโย. (1)}

\csromanheader{2}{\textbf{อาเสวน}}
\proseitem{50}{อนิยโต ธมฺโม อนิยตสฺส ธมฺมสฺส อาเสวน\-ปจฺจเยน ปจฺจโย. ปุริมา ปุริมา อนิยตา ขนฺธา ปจฺฉิมานํ ปจฺฉิมานํ อนิยตานํ ขนฺธานํ อาเสวน\-ปจฺจเยน ปจฺจโย. อนุโลมํ โคตฺรภุสฺส. อนุโลมํ โวทานสฺส อาเสวน\-ปจฺจเยน ปจฺจโย. (1)}

\prose{อนิยโต ธมฺโม มิจฺฉตฺต\-นิยตสฺส ธมฺมสฺส อาเสวน\-ปจฺจเยน ปจฺจโย. อนิยตํ โทมนสฺสํ มิจฺฉตฺต\-นิยตสฺส โทมนสฺสสฺส อาเสวน\-ปจฺจเยน ปจฺจโย. อนิยต\-มิจฺฉาทิฏฺฐิ นิยต\-มิจฺฉาทิฏฺฐิยา อาเสวน\-ปจฺจเยน ปจฺจโย. (2)}

\prose{อนิยโต ธมฺโม สมฺมตฺต\-นิยตสฺส ธมฺมสฺส อาเสวน\-ปจฺจเยน ปจฺจโย. โคตฺรภุ มคฺคสฺส. โวทานํ มคฺคสฺส อาเสวน\-ปจฺจเยน ปจฺจโย. (3)}

\csromanheader{2}{\textbf{กมฺม}}
\proseitem{51}{มิจฺฉตฺต\-นิยโต ธมฺโม มิจฺฉตฺต\-นิยตสฺส ธมฺมสฺส กมฺมปจฺจเยน ปจฺจโย. มิจฺฉตฺต\-นิยตา เจตนา สมฺปยุตฺตกานํ ขนฺธานํ กมฺมปจฺจเยน ปจฺจโย. (1)}

\prose{มิจฺฉตฺต\-นิยโต ธมฺโม อนิยตสฺส ธมฺมสฺส กมฺมปจฺจเยน ปจฺจโย. สหชาตา, นานากฺขณิกา.  สหชาตา มิจฺฉตฺต\-นิยตา เจตนา จิตฺต\-สมุฏฺฐานานํ รูปานํ กมฺมปจฺจเยน ปจฺจโย. นานากฺขณิกา มิจฺฉตฺต\-นิยตา เจตนา วิปากานํ ขนฺธานํ กฏตฺตา จ รูปานํ กมฺมปจฺจเยน ปจฺจโย. (2)}

\prose{มิจฺฉตฺต\-นิยโต ธมฺโม มิจฺฉตฺต\-นิยตสฺส จ อนิยตสฺส จ ธมฺมสฺส กมฺมปจฺจเยน ปจฺจโย. มิจฺฉตฺต\-นิยตา เจตนา สมฺปยุตฺตกานํ ขนฺธานํ จิตฺตสมุฏฺฐานานญฺจ รูปานํ กมฺมปจฺจเยน ปจฺจโย. (3)}

\proseitem{52}{\csromanfolio{379}สมฺมตฺตนิยโต ธมฺโม สมฺมตฺต\-นิยตสฺส ธมฺมสฺส กมฺมปจฺจเยน ปจฺจโย. สมฺมตฺตนิยตา เจตนา สมฺปยุตฺตกานํ ขนฺธานํ กมฺมปจฺจเยน ปจฺจโย. (1)}

\prose{สมฺมตฺตนิยโต ธมฺโม อนิยตสฺส ธมฺมสฺส กมฺมปจฺจเยน ปจฺจโย. สหชาตา, นานากฺขณิกา.  สหชาตา สมฺมตฺตนิยตา เจตนา จิตฺต\-สมุฏฺฐานานํ รูปานํ กมฺมปจฺจเยน ปจฺจโย. นานากฺขณิกา สมฺมตฺตนิยตา เจตนา วิปากานํ ขนฺธานํ กมฺมปจฺจเยน ปจฺจโย. (2)}

\prose{สมฺมตฺตนิยโต ธมฺโม สมฺมตฺต\-นิยตสฺส จ อนิยตสฺส จ ธมฺมสฺส กมฺมปจฺจเยน ปจฺจโย. สมฺมตฺตนิยตา เจตนา สมฺปยุตฺตกานํ ขนฺธานํ จิตฺตสมุฏฺฐานานญฺจ รูปานํ กมฺมปจฺจเยน ปจฺจโย. (3)}

\prose{อนิยโต ธมฺโม อนิยตสฺส ธมฺมสฺส กมฺมปจฺจเยน ปจฺจโย. สหชาตา, นานากฺขณิกา.  สหชาตา อนิยตา เจตนา สมฺปยุตฺตกานํ ขนฺธานํ จิตฺตสมุฏฺฐานานญฺจ รูปานํ กมฺมปจฺจเยน ปจฺจโย. ปฏิสนฺธิ\-กฺขเณ ฯเปฯ. นานากฺขณิกา อนิยตา เจตนา วิปากานํ ขนฺธานํ กฏตฺตา จ รูปานํ กมฺมปจฺจเยน ปจฺจโย. (1)}

\csromanheader{2}{\textbf{วิปาก}}
\proseitem{53}{อนิยโต ธมฺโม อนิยตสฺส ธมฺมสฺส วิปาก\-ปจฺจเยน ปจฺจโย. วิปาโก อนิยโต เอโก ขนฺโธ ติณฺณนฺนํ ขนฺธานํ จิตฺตสมุฏฺฐานานญฺจ รูปานํ วิปาก\-ปจฺจเยน ปจฺจโย ฯเปฯ. ปฏิสนฺธิ\-กฺขเณ ฯเปฯ. ขนฺธา วตฺถุสฺส ฯเปฯ.}

\csromanheader{2}{\textbf{อาหาราทิ}}
\proseitem{54}{มิจฺฉตฺต\-นิยโต ธมฺโม มิจฺฉตฺต\-นิยตสฺส ธมฺมสฺส อาหารปจฺจเยน ปจฺจโย. อินฺทฺริย\-ปจฺจเยน ปจฺจโย. ฌานปจฺจเยน ปจฺจโย. มคฺคปจฺจเยน ปจฺจโย. สมฺปยุตฺต\-ปจฺจเยน ปจฺจโย.}

\csromanheader{2}{\textbf{วิปฺปยุตฺต}}
\proseitem{55}{มิจฺฉตฺต\-นิยโต ธมฺโม อนิยตสฺส ธมฺมสฺส วิปฺปยุตฺต\-ปจฺจเยน ปจฺจโย. สหชาตํ, ปจฺฉาชาตํ.  สหชาตา มิจฺฉตฺต\-นิยตา ขนฺธา จิตฺต\-สมุฏฺฐานานํ รูปานํ วิปฺปยุตฺต\-ปจฺจเยน ปจฺจโย. ปจฺฉาชาตา \csromanfolio{380}มิจฺฉตฺต\-นิยตา ขนฺธา ปุเรชาตสฺส อิมสฺส กายสฺส วิปฺปยุตฺต\-ปจฺจเยน ปจฺจโย. (1)}

\prose{สมฺมตฺตนิยโต ธมฺโม อนิยตสฺส ธมฺมสฺส วิปฺปยุตฺต\-ปจฺจเยน ปจฺจโย. สหชาตํ, ปจฺฉาชาตํ.  สหชาตา สมฺมตฺตนิยตา ขนฺธา จิตฺต\-สมุฏฺฐานานํ รูปานํ วิปฺปยุตฺต\-ปจฺจเยน ปจฺจโย. ปจฺฉาชาตา สมฺมตฺตนิยตา ขนฺธา ปุเรชาตสฺส อิมสฺส กายสฺส วิปฺปยุตฺต\-ปจฺจเยน ปจฺจโย. (1)}

\prose{อนิยโต ธมฺโม อนิยตสฺส ธมฺมสฺส วิปฺปยุตฺต\-ปจฺจเยน ปจฺจโย. สหชาตํ, ปุเรชาตํ, ปจฺฉาชาตํ.  สหชาตา อนิยตา ขนฺธา จิตฺต\-สมุฏฺฐานานํ รูปานํ วิปฺปยุตฺต\-ปจฺจเยน ปจฺจโย. ปฏิสนฺธิ\-กฺขเณ ฯเปฯ. ขนฺธา วตฺถุสฺส วิปฺปยุตฺต\-ปจฺจเยน ปจฺจโย. วตฺถุ ขนฺธานํ วิปฺปยุตฺต\-ปจฺจเยน ปจฺจโย. ปุเรชาตํ จกฺขายตนํ จกฺขุ\-วิญฺญาณสฺส ฯเปฯ. กายายตนํ กายวิญฺญาณสฺส. วตฺถุ อนิยตานํ ขนฺธานํ วิปฺปยุตฺต\-ปจฺจเยน ปจฺจโย. ปจฺฉาชาตา อนิยตา ขนฺธา ปุเรชาตสฺส อิมสฺส กายสฺส วิปฺปยุตฺต\-ปจฺจเยน ปจฺจโย. (1)}

\prose{อนิยโต ธมฺโม มิจฺฉตฺต\-นิยตสฺส ธมฺมสฺส วิปฺปยุตฺต\-ปจฺจเยน ปจฺจโย. ปุเรชาตํ วตฺถุ มิจฺฉตฺต\-นิยตานํ ขนฺธานํ วิปฺปยุตฺต\-ปจฺจเยน ปจฺจโย. (2)}

\prose{อนิยโต ธมฺโม สมฺมตฺต\-นิยตสฺส ธมฺมสฺส วิปฺปยุตฺต\-ปจฺจเยน ปจฺจโย. ปุเรชาตํ วตฺถุ สมฺมตฺต\-นิยตานํ ขนฺธานํ วิปฺปยุตฺต\-ปจฺจเยน ปจฺจโย. (3)}

\csromanheader{2}{\textbf{อตฺถิ}}
\proseitem{56}{มิจฺฉตฺต\-นิยโต ธมฺโม มิจฺฉตฺต\-นิยตสฺส ธมฺมสฺส อตฺถิปจฺจเยน ปจฺจโย. มิจฺฉตฺต\-นิยโต เอโก ขนฺโธ ติณฺณนฺนํ ขนฺธานํ ฯเปฯ เทฺว ขนฺธา ทฺวินฺนํ ขนฺธานํ อตฺถิปจฺจเยน ปจฺจโย. (1)}

\prose{มิจฺฉตฺต\-นิยโต ธมฺโม อนิยตสฺส ธมฺมสฺส อตฺถิปจฺจเยน ปจฺจโย. สหชาตํ, ปจฺฉาชาตํ.  สหชาตา มิจฺฉตฺต\-นิยตา ขนฺธา จิตฺต\-สมุฏฺฐานานํ รูปานํ อตฺถิปจฺจเยน ปจฺจโย. ปจฺฉาชาตา มิจฺฉตฺต\-นิยตา ขนฺธา ปุเรชาตสฺส อิมสฺส กายสฺส อตฺถิปจฺจเยน ปจฺจโย. (2)}

\prose{มิจฺฉตฺต\-นิยโต ธมฺโม มิจฺฉตฺต\-นิยตสฺส จ อนิยตสฺส จ อตฺถิปจฺจเยน ปจฺจโย. มิจฺฉตฺต\-นิยโต เอโก ขนฺโธ ติณฺณนฺนํ ขนฺธานํ จิตฺตสมุฏฺฐานานญฺจ รูปานํ อตฺถิปจฺจเยน ปจฺจโย ฯเปฯ เทฺว ขนฺธา ฯเปฯ. (3)}

\prose{\csromanfolio{381}สมฺมตฺตนิยโต ธมฺโม สมฺมตฺต\-นิยตสฺส ธมฺมสฺส อตฺถิปจฺจเยน ปจฺจโย ฯเปฯ. (ติสฺโส ปญฺหา.)}

\prose{อนิยโต ธมฺโม อนิยตสฺส ธมฺมสฺส อตฺถิปจฺจเยน ปจฺจโย. สหชาตํ, ปุเรชาตํ, ปจฺฉาชาตํ, อาหารํ, อินฺทฺริยํ.  สหชาโต อนิยโต เอโก ขนฺโธ ติณฺณนฺนํ ขนฺธานํ จิตฺตสมุฏฺฐานานญฺจ รูปานํ อตฺถิปจฺจเยน ปจฺจโย ฯเปฯ เทฺว ขนฺธา ฯเปฯ. ปฏิสนฺธิ\-กฺขเณ ฯเปฯ. ขนฺธา วตฺถุสฺส อตฺถิปจฺจเยน ปจฺจโย. วตฺถุ ขนฺธานํ อตฺถิปจฺจเยน ปจฺจโย. เอกํ มหาภูตํ ฯเปฯ. อสญฺญสตฺตานํ เอกํ มหาภูตํ ฯเปฯ. ปุเรชาตํ จกฺขุํ ฯเปฯ. วตฺถุํ อนิจฺจโต ทุกฺขโต อนตฺตโต วิปสฺสติ, อสฺสาเทติ อภินนฺทติ, ตํ อารพฺภ ราโค อุปฺปชฺชติ ฯเปฯ โทมนสฺสํ อุปฺปชฺชติ, ทิพฺเพน จกฺขุนา รูปํ ปสฺสติ. ทิพฺพาย โสตธาตุยา สทฺทํ สุณาติ.  รูปายตนํ จกฺขุ\-วิญฺญาณสฺส ฯเปฯ. โผฏฺฐพฺพา\-ยตนํ กายวิญฺญาณสฺส ฯเปฯ. จกฺขายตนํ จกฺขุ\-วิญฺญาณสฺส ฯเปฯ. กายายตนํ กายวิญฺญาณสฺส. วตฺถุ อนิยตานํ ขนฺธานํ อตฺถิปจฺจเยน ปจฺจโย. ปจฺฉาชาตา อนิยตา ขนฺธา ปุเรชาตสฺส อิมสฺส กายสฺส อตฺถิปจฺจเยน ปจฺจโย. กพฬีกาโร อาหาโร อิมสฺส กายสฺส อตฺถิปจฺจเยน ปจฺจโย. รูปชีวิตินฺทฺริยํ กฏตฺตารูปานํ อตฺถิปจฺจเยน ปจฺจโย. (1)}

\prose{อนิยโต ธมฺโม มิจฺฉตฺต\-นิยตสฺส ธมฺมสฺส อตฺถิปจฺจเยน ปจฺจโย. ปุเรชาตํ รูปชีวิตินฺทฺริยํ มาตุ\-ฆาติ\-กมฺมสฺส ฯเปฯ. รุหิรุ\-ปฺปาท\-กมฺมสฺส อตฺถิปจฺจเยน ปจฺจโย. วตฺถุ มิจฺฉตฺต\-นิยตานํ ขนฺธานํ อตฺถิปจฺจเยน ปจฺจโย. (2)}

\prose{อนิยโต ธมฺโม สมฺมตฺต\-นิยตสฺส ธมฺมสฺส อตฺถิปจฺจเยน ปจฺจโย. ปุเรชาตํ วตฺถุ สมฺมตฺต\-นิยตานํ ขนฺธานํ อตฺถิปจฺจเยน ปจฺจโย. (3)}

\proseitem{57}{มิจฺฉตฺต\-นิยโต จ อนิยโต จ ธมฺมา มิจฺฉตฺต\-นิยตสฺส ธมฺมสฺส อตฺถิปจฺจเยน ปจฺจโย. สหชาตํ, ปุเรชาตํ.  สหชาโต มิจฺฉตฺต\-นิยโต เอโก ขนฺโธ จ วตฺถุ จ ติณฺณนฺนํ ขนฺธานํ ฯเปฯ เทฺว ขนฺธา จ ฯเปฯ. (1)}

\prose{มิจฺฉตฺต\-นิยโต จ อนิยโต จ ธมฺมา อนิยตสฺส ธมฺมสฺส อตฺถิปจฺจเยน ปจฺจโย. สหชาตํ, ปจฺฉาชาตํ, อาหารํ, อินฺทฺริยํ. สหชาตา มิจฺฉตฺต\-นิยตา ขนฺธา จ มหาภูตา จ จิตฺต\-สมุฏฺฐานานํ รูปานํ อตฺถิปจฺจเยน ปจฺจโย. ปจฺฉาชาตา มิจฺฉตฺต\-นิยตา ขนฺธา จ กพฬีกาโร อาหาโร \csromanfolio{382}จ อิมสฺส กายสฺส อตฺถิปจฺจเยน ปจฺจโย. ปจฺฉาชาตา มิจฺฉตฺต\-นิยตา ขนฺธา จ รูปชีวิตินฺทฺริยญฺจ กฏตฺตารูปานํ อตฺถิปจฺจเยน ปจฺจโย. (2)}

\prose{สมฺมตฺตนิยโต จ อนิยโต จ ธมฺมา สมฺมตฺต\-นิยตสฺส ธมฺมสฺส อตฺถิปจฺจเยน ปจฺจโย ฯเปฯ. (เทฺว ปญฺหา มิจฺฉตฺตนิยต\-สทิสา.)}

\csromanmatikaanchor{mtk.143}
\csromantocmarkref{subparagraph}{1. ปจฺจยานุโลม 2. สงฺขฺยาวาร}{mtk.143}
\bookmark[dest={mtk.143},level=5]{1. ปจฺจยานุโลม 2. สงฺขฺยาวาร}
\csromanmarkh{1. ปจฺจยานุโลม}\csromanmarkhii{2. สงฺขฺยาวาร}\csromanheadernomark{6}{1. \textbf{ปจฺจยานุโลม 2. สงฺขฺยาวาร}}
\proseitemwithcloserruled{58}{เหตุยา สตฺต, อารมฺมเณ ปญฺจ, อธิปติยา อฏฺฐ, อนนฺตเร ปญฺจ, สมนนฺตเร ปญฺจ, สหชาเต นว, อญฺญมญฺเญ ตีณิ, นิสฺสเย เตรส, อุปนิสฺสเย สตฺต, ปุเรชาเต ตีณิ, ปจฺฉาชาเต ตีณิ, อาเสวเน ตีณิ, กมฺเม สตฺต, วิปาเก เอกํ, อาหาเร สตฺต, อินฺทฺริเย สตฺต, ฌาเน สตฺต, มคฺเค สตฺต, สมฺปยุตฺเต ตีณิ, วิปฺปยุตฺเต ปญฺจ, อตฺถิยา เตรส, นตฺถิยา ปญฺจ, วิคเต ปญฺจ, อวิคเต เตรส. (เอวํ คเณตพฺพํ.)}{อนุโลมํ.}
\csromanmatikaanchor{mtk.144}
\csromantocmarkref{subparagraph}{2. ปจฺจนียุทฺธาร}{mtk.144}
\bookmark[dest={mtk.144},level=5]{2. ปจฺจนียุทฺธาร}
\csromanheader{6}{2. ปจฺจนียุ\-ทฺธาร}
\proseitem{59}{มิจฺฉตฺต\-นิยโต ธมฺโม มิจฺฉตฺต\-นิยตสฺส ธมฺมสฺส สหชาต\-ปจฺจเยน ปจฺจโย. อุปนิสฺสย\-ปจฺจเยน ปจฺจโย. (1)}

\prose{มิจฺฉตฺต\-นิยโต ธมฺโม อนิยตสฺส ธมฺมสฺส อารมฺมณ\-ปจฺจเยน ปจฺจโย. สหชาต\-ปจฺจเยน ปจฺจโย. อุปนิสฺสย\-ปจฺจเยน ปจฺจโย. ปจฺฉาชาต\-ปจฺจเยน ปจฺจโย. กมฺมปจฺจเยน ปจฺจโย. (2)}

\prose{มิจฺฉตฺต\-นิยโต ธมฺโม มิจฺฉตฺต\-นิยตสฺส จ อนิยตสฺส จ ธมฺมสฺส สหชาต\-ปจฺจเยน ปจฺจโย. (3)}

\prose{สมฺมตฺตนิยโต ธมฺโม สมฺมตฺต\-นิยตสฺส ธมฺมสฺส สหชาต\-ปจฺจเยน ปจฺจโย. อุปนิสฺสย\-ปจฺจเยน ปจฺจโย. (1)}

\prose{สมฺมตฺตนิยโต ธมฺโม อนิยตสฺส ธมฺมสฺส อารมฺมณ\-ปจฺจเยน ปจฺจโย. สหชาต\-ปจฺจเยน ปจฺจโย. อุปนิสฺสย\-ปจฺจเยน ปจฺจโย. ปจฺฉาชาต\-ปจฺจเยน ปจฺจโย. (2)}

\prose{\csromanfolio{383}สมฺมตฺตนิยโต ธมฺมา สมฺมตฺต\-นิยตสฺส จ อนิยตสฺส จ ธมฺมสฺส สหชาต\-ปจฺจเยน ปจฺจโย. (3)}

\proseitem{60}{อนิยโต ธมฺโม อนิยตสฺส ธมฺมสฺส อารมฺมณ\-ปจฺจเยน ปจฺจโย. สหชาต\-ปจฺจเยน ปจฺจโย. อุปนิสฺสย\-ปจฺจเยน ปจฺจโย. ปุเรชาต\-ปจฺจเยน ปจฺจโย. ปจฺฉาชาต\-ปจฺจเยน ปจฺจโย. กมฺมปจฺจเยน ปจฺจโย. อาหารปจฺจเยน ปจฺจโย. อินฺทฺริย\-ปจฺจเยน ปจฺจโย. (1)}

\prose{อนิยโต ธมฺโม มิจฺฉตฺต\-นิยตสฺส ธมฺมสฺส อารมฺมณ\-ปจฺจเยน ปจฺจโย. อุปนิสฺสย\-ปจฺจเยน ปจฺจโย. ปุเรชาต\-ปจฺจเยน ปจฺจโย. (2)}

\prose{อนิยโต ธมฺโม สมฺมตฺต\-นิยตสฺส ธมฺมสฺส อุปนิสฺสย\-ปจฺจเยน ปจฺจโย. ปุเรชาต\-ปจฺจเยน ปจฺจโย. (3)}

\prose{มิจฺฉตฺต\-นิยโต จ อนิยโต จ ธมฺมา มิจฺฉตฺต\-นิยตสฺส ธมฺมสฺส สหชาตํ, ปุเรชาตํ. (1)}

\prose{มิจฺฉตฺต\-นิยโต จ อนิยโต จ ธมฺมา อนิยตสฺส ธมฺมสฺส สหชาตํ, ปจฺฉาชาตํ, อาหารํ, อินฺทฺริยํ. (2)}

\prose{สมฺมตฺตนิยโต จ อนิยโต จ ธมฺมา สมฺมตฺต\-นิยตสฺส ธมฺมสฺส สหชาตํ, ปุเรชาตํ. (1)}

\prose{สมฺมตฺตนิยโต จ อนิยโต จ ธมฺมา อนิยตสฺส ธมฺมสฺส สหชาตํ, ปจฺฉาชาตํ, อาหารํ, อินฺทฺริยํ. (2)}

\csromanmatikaanchor{mtk.145}
\csromantocmarkref{subparagraph}{2. ปจฺจยปจฺจนีย 2. สงฺขฺยาวาร}{mtk.145}
\bookmark[dest={mtk.145},level=5]{2. ปจฺจยปจฺจนีย 2. สงฺขฺยาวาร}
\csromanmarkh{2. ปจฺจย\-ปจฺจนีย}\csromanmarkhii{2. สงฺขฺยาวาร}\csromanheadernomark{6}{2. ปจฺจย\-ปจฺจนีย 2. สงฺขฺยาวาร}
\proseitem{61}{นเหตุยา เตรส, นอารมฺมเณ, นอธิปติยา, นอนนฺตเร, นสมนนฺตเร เตรส, นสหชาเต นว, นอญฺญมญฺเญ นว, นนิสฺสเย นว, นอุปนิสฺสเย เตรส, นปุเรชาเต เอกาทส, นปจฺฉาชาเต เตรส. นอาเสวเน เตรส, นกมฺเม, นวิปาเก, นอาหาเร เตรส ฯเปฯ นมคฺเค เตรส, นสมฺปยุตฺเต นว, นวิปฺปยุตฺเต สตฺต, โนอตฺถิยา สตฺต, โนนตฺถิยา เตรส, โนวิคเต เตรส, โนอวิคเต สตฺต. (เอวํ คเณตพฺพํ.)}

\vspace{\tipitakaparskip}
\csromancenter{ปจฺจนียํ.}
\csromanheader{2}{3. \textbf{ปจฺจยานุโลม ปจฺจนีย}}
\csromanmatikaanchor{mtk.146}
\csromanmatikaanchor{mtk.147}
\csromanmatikaanchor{mtk.148}
\csromanmatikaanchor{mtk.149}
\csromantocmarknopage{subparagraph}{3. ปจฺจยานุโลมปจฺจนีย}{mtk.146}
\bookmark[dest={mtk.146},level=5]{3. ปจฺจยานุโลมปจฺจนีย}
\csromantocmarkref{subparagraph}{เหตุทุก}{mtk.147}
\bookmark[dest={mtk.147},level=5]{เหตุทุก}
\csromantocmarknopage{subparagraph}{4. ปจฺจยปจฺจนียานุโลม}{mtk.148}
\bookmark[dest={mtk.148},level=5]{4. ปจฺจยปจฺจนียานุโลม}
\csromantocmarkref{subparagraph}{นเหตุทุก}{mtk.149}
\bookmark[dest={mtk.149},level=5]{นเหตุทุก}
\proseitemwithcloserruled{62}{\csromanfolio{384}เหตุปจฺจยา นอารมฺมเณ สตฺต, นอธิปติยา สตฺต, นอนนฺตเร สตฺต, นสมนนฺตเร สตฺต, นอญฺญมญฺเญ ตีณิ, นอุปนิสฺสเย สตฺต ฯเปฯ นมคฺเค สตฺต, นสมฺปยุตฺเต ตีณิ, นวิปฺปยุตฺเต ตีณิ, โนนตฺถิยา สตฺต, โนวิคเต สตฺต. (เอวํ คเณตพฺพํ.)}{อนุโลม\-ปจฺจนียํ.}
\proseitem{63}{นเหตุปจฺจยา อารมฺมเณ ปญฺจ, อธิปติยา อฏฺฐ, อนนฺตเร ปญฺจ, สมนนฺตเร ปญฺจ, สหชาเต นว, อญฺญมญฺเญ ตีณิ, นิสฺสเย เตรส, อุปนิสฺสเย สตฺต, ปุเรชาเต ตีณิ, ปจฺฉาชาเต ตีณิ, อาเสวเน ตีณิ, กมฺเม สตฺต, วิปาเก เอกํ, อาหาเร สตฺต, อินฺทฺริเย, ฌาเน, มคฺเค สตฺต, สมฺปยุตฺเต ตีณิ, วิปฺปยุตฺเต ปญฺจ, อตฺถิยา เตรส, นตฺถิยา ปญฺจ, วิคเต ปญฺจ, อวิคเต เตรส. (เอวํ คเณตพฺพํ.)}

\vspace{\tipitakaparskip}
\csromancenter{ปจฺจนียา\-นุโลมํ.}
\nitthitam{มิจฺฉตฺตนิยต\-ตฺติกํ นิฏฺฐิตํ.}
\csromanmatikaanchor{mtk.150}
\csromanmatikaanchor{mtk.164}
\csromantocmarknopage{subparagraph}{16. มคฺคารมฺมณตฺติก 1. ปฏิจฺจวาร}{mtk.150}
\bookmark[dest={mtk.150},level=5]{16. มคฺคารมฺมณตฺติก 1. ปฏิจฺจวาร}
\csromanmarkh{16. มคฺคา\-รมฺมณ\-ตฺติก}\csromanmarkhii{16. มคฺคา\-รมฺมณ\-ตฺติก 1. ปฏิจฺจวาร}\csromanheadernomark{6}{16. มคฺคา\-รมฺมณ\-ตฺติก 16. มคฺคา\-รมฺมณ\-ตฺติก 1. ปฏิจฺจวาร}
\csromanmarkh{1. ปจฺจยานุโลม}\csromanmarkhii{1. วิภงฺควาร}\csromanheadernomark{2}{1. \textbf{ปจฺจยานุโลม 1. วิภงฺควาร}}
\csromanmatikaanchor{mtk.151}
\csromanmatikaanchor{mtk.152}
\csromantocmarknopage{subparagraph}{1. ปจฺจยานุโลม 1. วิภงฺควาร}{mtk.151}
\bookmark[dest={mtk.151},level=5]{1. ปจฺจยานุโลม 1. วิภงฺควาร}
\csromantocmarkref{subparagraph}{เหตุ}{mtk.152}
\bookmark[dest={mtk.152},level=5]{เหตุ}
\csromanheader{6}{\textbf{เหตุ}}
\proseitem{1}{\csromanfolio{385}มคฺคารมฺมณํ ธมฺมํ ปฏิจฺจ มคฺคารมฺมโณ ธมฺโม อุปฺปชฺชติ เหตุปจฺจยา. มคฺคารมฺมณํ เอกํ ขนฺธํ ปฏิจฺจ ตโย ขนฺธา ฯเปฯ เทฺว ขนฺเธ ฯเปฯ. (1)}

\prose{มคฺคารมฺมณํ ธมฺมํ ปฏิจฺจ มคฺคาธิปติ ธมฺโม อุปฺปชฺชติ เหตุปจฺจยา. มคฺคารมฺมณํ เอกํ ขนฺธํ ปฏิจฺจ มคฺคาธิปตี ตโย ขนฺธา, ตโย ขนฺเธ ปฏิจฺจ เอโก ขนฺโธ. เทฺว ขนฺเธ ฯเปฯ. (2)}

\prose{มคฺคารมฺมณํ ธมฺมํ ปฏิจฺจ มคฺคารมฺมโณ จ มคฺคาธิปติ จ ธมฺมา อุปฺปชฺชนฺติ เหตุปจฺจยา. มคฺคารมฺมณํ เอกํ ขนฺธํ ปฏิจฺจ มคฺคารมฺมณา จ มคฺคาธิปตี จ ตโย ขนฺธา ฯเปฯ เทฺว ขนฺเธ ฯเปฯ. (3)}

\proseitem{2}{มคฺคเหตุกํ ธมฺมํ ปฏิจฺจ มคฺคเหตุโก ธมฺโม อุปฺปชฺชติ เหตุปจฺจยา. มคฺคเหตุกํ เอกํ ขนฺธํ ปฏิจฺจ ตโย ขนฺธา ฯเปฯ เทฺว ขนฺเธ ฯเปฯ. (1)}

\prose{มคฺคเหตุกํ ธมฺมํ ปฏิจฺจ มคฺคาธิปติ ธมฺโม อุปฺปชฺชติ เหตุปจฺจยา. มคฺคเหตุกํ เอกํ ขนฺธํ ปฏิจฺจ มคฺคาธิปตี ตโย ขนฺธา ฯเปฯ เทฺว ขนฺเธ ฯเปฯ. (2)}

\prose{มคฺคเหตุกํ ธมฺมํ ปฏิจฺจ มคฺคเหตุโก จ มคฺคาธิปติ จ ธมฺมา อุปฺปชฺชนฺติ เหตุปจฺจยา. มคฺคเหตุกํ เอกํ ขนฺธํ ปฏิจฺจ มคฺคเหตุกา จ มคฺคาธิปตี จ ตโย ขนฺธา ฯเปฯ เทฺว ขนฺเธ ฯเปฯ. (3)}

\proseitem{3}{มคฺคาธิปติํ ธมฺมํ ปฏิจฺจ มคฺคาธิปติ ธมฺโม อุปฺปชฺชติ เหตุปจฺจยา. มคฺคาธิปติํ เอกํ ขนฺธํ ปฏิจฺจ ตโย ขนฺธา ฯเปฯ เทฺว ขนฺเธ ฯเปฯ. (1)}

\prose{มคฺคาธิปติํ ธมฺมํ ปฏิจฺจ มคฺคารมฺมโณ ธมฺโม อุปฺปชฺชติ เหตุปจฺจยา. มคฺคาธิปติํ เอกํ ขนฺธํ ปฏิจฺจ มคฺคารมฺมณา ตโย ขนฺธา ฯเปฯ เทฺว ขนฺเธ. ฯเปฯ. (2)}

\prose{มคฺคาธิปติํ ธมฺมํ ปฏิจฺจ มคฺคเหตุโก ธมฺโม อุปฺปชฺชติ เหตุปจฺจโย. มคฺคาธิปติํ เอกํ ขนฺธํ มคฺคเหตุกา ตโย ขนฺธา ฯเปฯ เทฺว ขนฺเธ ฯเปฯ. (3)}

\prose{มคฺคาธิปติํ ธมฺมํ ปฏิจฺจ มคฺคารมฺมโณ จ มคฺคาธิปติ จ ธมฺมา อุปฺปชฺชนฺติ เหตุปจฺจยา. มคฺคาธิปติํ เอกํ ขนฺธํ ปฏิจฺจ มคฺคารมฺมณา จ มคฺคาธิปตี จ ตโย ขนฺธา ฯเปฯ เทฺว ขนฺเธ ฯเปฯ. (4)}

\prose{\csromanfolio{386}มคฺคาธิปติํ ธมฺมํ ปฏิจฺจ มคฺคเหตุโก จ มคฺคาธิปติ จ ธมฺมา อุปฺปชฺชนฺติ เหตุปจฺจยา. มคฺคาธิปติํ เอกํ ขนฺธํ ปฏิจฺจ มคฺคเหตุกา จ มคฺคาธิปตี จ ตโย ขนฺธา ฯเปฯ เทฺว ขนฺเธ ฯเปฯ. (5)}

\proseitem{4}{มคฺคา\-รมฺมณญฺจ มคฺคา\-ธิปติญฺจ ธมฺมํ ปฏิจฺจ มคฺคารมฺมโณ ธมฺโม อุปฺปชฺชติ เหตุปจฺจยา. มคฺคา\-รมฺมณญฺจ มคฺคา\-ธิปติญฺจ เอกํ ขนฺธํ ปฏิจฺจ มคฺคารมฺมณา ตโย ขนฺธา ฯเปฯ เทฺว ขนฺเธ ฯเปฯ. (1)}

\prose{มคฺคา\-รมฺมณญฺจ มคฺคา\-ธิปติญฺจ ธมฺมํ ปฏิจฺจ มคฺคาธิปติ ธมฺโม อุปฺปชฺชติ เหตุปจฺจยา. มคฺคา\-รมฺมณญฺจ มคฺคา\-ธิปติญฺจ เอกํ ขนฺธํ ปฏิจฺจ มคฺคาธิปตี ตโย ขนฺธา ฯเปฯ เทฺว ขนฺเธ ฯเปฯ. (2)}

\prose{มคฺคา\-รมฺมณญฺจ มคฺคา\-ธิปติญฺจ ธมฺมํ ปฏิจฺจ มคฺคารมฺมโณ จ มคฺคาธิปติ จ ธมฺมา อุปฺปชฺชนฺติ เหตุปจฺจยา. มคฺคา\-รมฺมณญฺจ มคฺคา\-ธิปติญฺจ เอกํ ขนฺธํ ปฏิจฺจ มคฺคารมฺมณา จ มคฺคาธิปตี จ ตโย ขนฺธา ฯเปฯ เทฺว ขนฺเธ ฯเปฯ. (3)}

\proseitem{5}{มคฺคเหตุกญฺจ มคฺคา\-ธิปติญฺจ ธมฺมํ ปฏิจฺจ มคฺคเหตุโก ธมฺโม อุปฺปชฺชติ เหตุปจฺจยา. มคฺคเหตุกญฺจ มคฺคา\-ธิปติญฺจ เอกํ ขนฺธํ ปฏิจฺจ มคฺคเหตุกา ตโย ขนฺธา ฯเปฯ เทฺว ขนฺเธ ฯเปฯ. (1)}

\prose{มคฺคเหตุกญฺจ มคฺคา\-ธิปติญฺจ ธมฺมํ ปฏิจฺจ มคฺคาธิปติ ธมฺโม อุปฺปชฺชติ เหตุปจฺจยา. มคฺคเหตุกญฺจ มคฺคา\-ธิปติญฺจ เอกํ ขนฺธํ ปฏิจฺจ มคฺคาธิปตี ตโย ขนฺธา ฯเปฯ เทฺว ขนฺเธ ฯเปฯ. (2)}

\prose{มคฺคเหตุกญฺจ มคฺคา\-ธิปติญฺจ ธมฺมํ ปฏิจฺจ มคฺคเหตุโก จ มคฺคาธิปติ จ ธมฺมา อุปฺปชฺชนฺติ เหตุปจฺจยา. มคฺคเหตุกญฺจ มคฺคา\-ธิปติญฺจ เอกํ ขนฺธํ ปฏิจฺจ มคฺคเหตุกา จ มคฺคาธิปตี จ ตโย ขนฺธา ฯเปฯ เทฺว ขนฺเธ ฯเปฯ. (3)}

\csromanmatikaanchor{mtk.153}
\csromantocmarkref{subparagraph}{อารมฺมณาทิ}{mtk.153}
\bookmark[dest={mtk.153},level=5]{อารมฺมณาทิ}
\csromanheader{6}{\textbf{อารมฺมณาทิ}}
\proseitem{6}{มคฺคารมฺมณํ ธมฺมํ ปฏิจฺจ มคฺคารมฺมโณ ธมฺโม อุปฺปชฺชติ อารมฺมณ\-ปจฺจยา. อธิปติ\-ปจฺจยา. อนนฺตร\-ปจฺจยา. สมนนฺตร\-ปจฺจยา. สหชาต\-ปจฺจยา. อญฺญมญฺญ\-ปจฺจยา. นิสฺสยปจฺจยา. อุปนิสฺสย\-ปจฺจยา. ปุเรชาต\-ปจฺจยา. อาเสวนปจฺจยา. กมฺมปจฺจยา. อาหารปจฺจยา. อินฺทฺริยปจฺจยา. ฌานปจฺจยา. มคฺคปจฺจยา. สมฺปยุตฺต\-ปจฺจยา. วิปฺปยุตฺต\-ปจฺจยา. อตฺถิปจฺจยา. นตฺถิปจฺจยา. วิคตปจฺจยา. อวิคตปจฺจยา.}

\csromanmatikaanchor{mtk.154}
\csromantocmarkref{subparagraph}{1. ปจฺจยานุโลม 2. สงฺขฺยาวาร}{mtk.154}
\bookmark[dest={mtk.154},level=5]{1. ปจฺจยานุโลม 2. สงฺขฺยาวาร}
\csromanmarkh{16. มคฺคา\-รมฺมณ\-ตฺติก}\csromanmarkhii{1. ปจฺจยานุโลม 2. สงฺขฺยาวาร}\csromanheadernomark{6}{16. มคฺคา\-รมฺมณ\-ตฺติก 1. ปจฺจยานุโลม 2. สงฺขฺยาวาร}
\csromanmatikaanchor{mtk.155}
\csromanmatikaanchor{mtk.157}
\csromantocmarknopage{subparagraph}{2. ปจฺจยปจฺจนีย 1. วิภงฺควาร}{mtk.155}
\bookmark[dest={mtk.155},level=5]{2. ปจฺจยปจฺจนีย 1. วิภงฺควาร}
\proseitemwithcloserruled{7}{\csromanfolio{387}เหตุยา สตฺตรส, อารมฺมเณ, อธิปติยา, อนนฺตเร, สมนนฺตเร, สหชาเต, อญฺญมญฺเญ, นิสฺสเย, อุปนิสฺสเย, ปุเรชาเต, อาเสวเน, กมฺเม, อาหาเร, อินฺทฺริเย, ฌาเน, มคฺเค, สมฺปยุตฺเต, วิปฺปยุตฺเต, อตฺถิยา, นตฺถิยา, วิคเต, อวิคเต สตฺตรส. (เอวํ คเณตพฺพํ.)}{\textbf{อนุโลมํ.}}
\csromanmarkh{2. ปจฺจย\-ปจฺจนีย}\csromanmarkhii{1. วิภงฺควาร}\csromanheadernomark{2}{2. ปจฺจย\-ปจฺจนีย 1. วิภงฺควาร}
\csromanmatikaanchor{mtk.156}
\csromantocmarkref{subparagraph}{นเหตุ}{mtk.156}
\bookmark[dest={mtk.156},level=5]{นเหตุ}
\csromantocmarkref{subparagraph}{นอธิปติ}{mtk.157}
\bookmark[dest={mtk.157},level=5]{นอธิปติ}
\csromanheader{6}{\textbf{นเหตุ}}
\proseitem{8}{มคฺคารมฺมณํ ธมฺมํ ปฏิจฺจ มคฺคารมฺมโณ ธมฺโม อุปฺปชฺชติ นเหตุปจฺจยา. อเหตุกํ มคฺคารมฺมณํ เอกํ ขนฺธํ ปฏิจฺจ ตโย ขนฺธา ฯเปฯ เทฺว ขนฺเธ ฯเปฯ. (1)}

\prose{นอธิปติ}

\proseitem{9}{มคฺคารมฺมณํ ธมฺมํ ปฏิจฺจ มคฺคารมฺมโณ ธมฺโม อุปฺปชฺชติ นอธิปติ\-ปจฺจยา. มคฺคารมฺมณํ เอกํ ขนฺธํ ปฏิจฺจ ตโย ขนฺธา ฯเปฯ เทฺว ขนฺเธ ฯเปฯ. (1)}

\prose{มคฺคารมฺมณํ ธมฺมํ ปฏิจฺจ มคฺคาธิปติ ธมฺโม อุปฺปชฺชติ นอธิปติ\-ปจฺจยา. มคฺคารมฺมณํ เอกํ ขนฺธํ ปฏิจฺจ มคฺคาธิปตี ตโย ขนฺธา ฯเปฯ เทฺว ขนฺเธ ฯเปฯ. (2)}

\prose{มคฺคารมฺมณํ ธมฺมํ ปฏิจฺจ มคฺคารมฺมโณ จ มคฺคาธิปติ จ ธมฺมา อุปฺปชฺชนฺติ นอธิปติ\-ปจฺจยา. มคฺคารมฺมณํ เอกํ ขนฺธํ ปฏิจฺจ มคฺคารมฺมณา จ มคฺคาธิปตี จ ตโย ขนฺธา ฯเปฯ เทฺว ขนฺเธ ฯเปฯ. (3)}

\proseitem{10}{มคฺคเหตุกํ ธมฺมํ ปฏิจฺจ มคฺคเหตุโก ธมฺโม อุปฺปชฺชติ นอธิปติ\-ปจฺจยา. มคฺคเหตุเก ขนฺเธ ปฏิจฺจ มคฺค\-เหตุกา\-ธิปติ. (1)}

\prose{มคฺคเหตุกํ ธมฺมํ ปฏิจฺจ มคฺคาธิปติ ธมฺโม อุปฺปชฺชติ นอธิปติ\-ปจฺจยา. มคฺคเหตุเก ขนฺเธ ปฏิจฺจ มคฺคาธิปติ อธิปติ. (2)}

\prose{มคฺคเหตุกํ ธมฺมํ ปฏิจฺจ มคฺคเหตุโก จ มคฺคาธิปติ จ ธมฺมา อุปฺปชฺชนฺติ นอธิปติ\-ปจฺจยา. มคฺคเหตุเก ขนฺเธ ปฏิจฺจ มคฺคเหตุโก จ มคฺคาธิปติ จ อธิปติ. (3)}

\proseitem{11}{\csromanfolio{388}มคฺคาธิปติํ ธมฺมํ ปฏิจฺจ มคฺคาธิปติ ธมฺโม อุปฺปชฺชติ นอธิปติ\-ปจฺจยา. มคฺคาธิปตี ขนฺเธ ปฏิจฺจ มคฺคาธิปติ อธิปติ. มคฺคาธิปติํ เอกํ ขนฺธํ ปฏิจฺจ ตโย ขนฺธา ฯเปฯ เทฺว ขนฺเธ ฯเปฯ. (1)}

\prose{มคฺคาธิปติํ ธมฺมํ ปฏิจฺจ มคฺคารมฺมโณ ธมฺโม อุปฺปชฺชติ นอธิปติ\-ปจฺจยา. มคฺคาธิปติํ เอกํ ขนฺธํ ปฏิจฺจ มคฺคารมฺมณา ตโย ขนฺธา ฯเปฯ เทฺว ขนฺเธ ฯเปฯ. (2)}

\prose{มคฺคาธิปติํ ธมฺมํ ปฏิจฺจ มคฺคเหตุโก ธมฺโม อุปฺปชฺชติ นอธิปติ\-ปจฺจยา. มคฺคาธิปตี ขนฺเธ ปฏิจฺจ มคฺคเหตุโก อธิปติ. (3)}

\prose{มคฺคาธิปติํ ธมฺมํ ปฏิจฺจ มคฺคารมฺมโณ จ มคฺคาธิปติ จ ธมฺมา อุปฺปชฺชนฺติ นอธิปติ\-ปจฺจยา. มคฺคาธิปติํ เอกํ ขนฺธํ ปฏิจฺจ มคฺคารมฺมณา จ มคฺคาธิปตี จ ตโย ขนฺธา ฯเปฯ เทฺว ขนฺเธ ฯเปฯ. (4)}

\prose{มคฺคาธิปติํ ธมฺมํ ปฏิจฺจ มคฺคเหตุโก จ มคฺคาธิปติ จ ธมฺมา อุปฺปชฺชนฺติ นอธิปติ\-ปจฺจยา. มคฺคาธิปตี ขนฺเธ ปฏิจฺจ มคฺคเหตุโก จ มคฺคาธิปติ จ อธิปติ. (5)}

\proseitem{12}{มคฺคา\-รมฺมณญฺจ มคฺคา\-ธิปติญฺจ ธมฺมํ ปฏิจฺจ มคฺคารมฺมโณ ธมฺโม อุปฺปชฺชติ นอธิปติ\-ปจฺจยา. มคฺคา\-รมฺมณญฺจ มคฺคา\-ธิปติญฺจ เอกํ ขนฺธํ ปฏิจฺจ มคฺคารมฺมณา ตโย ขนฺธา ฯเปฯ เทฺว ขนฺเธ ฯเปฯ. (1)}

\prose{มคฺคา\-รมฺมณญฺจ มคฺคา\-ธิปติญฺจ ธมฺมํ ปฏิจฺจ มคฺคาธิปติ ธมฺโม อุปฺปชฺชติ นอธิปติ\-ปจฺจยา. มคฺคา\-รมฺมณญฺจ มคฺคา\-ธิปติญฺจ เอกํ ขนฺธํ ปฏิจฺจ มคฺคาธิปตี ตโย ขนฺธา ฯเปฯ เทฺว ขนฺเธ ฯเปฯ. (2)}

\prose{มคฺคา\-รมฺมณญฺจ มคฺคา\-ธิปติญฺจ ธมฺมํ ปฏิจฺจ มคฺคารมฺมโณ จ มคฺคาธิปติ จ ธมฺมา อุปฺปชฺชนฺติ นอธิปติ\-ปจฺจยา. มคฺคา\-รมฺมณญฺจ มคฺคา\-ธิปติญฺจ เอกํ ขนฺธํ ปฏิจฺจ มคฺคารมฺมณา จ มคฺคาธิปตี จ ตโย ขนฺธา ฯเปฯ เทฺว ขนฺเธ ฯเปฯ. (3)}

\proseitem{13}{มคฺคเหตุกญฺจ มคฺคา\-ธิปติญฺจ ธมฺมํ ปฏิจฺจ มคฺคเหตุโก ธมฺโม อุปฺปชฺชติ นอธิปติ\-ปจฺจยา. มคฺคเหตุเก จ มคฺคาธิปตี จ ขนฺเธ ปฏิจฺจ มคฺคเหตุโก อธิปติ. (1)}

\prose{มคฺคเหตุกญฺจ มคฺคา\-ธิปติญฺจ ธมฺมํ ปฏิจฺจ มคฺคาธิปติ ธมฺโม อุปฺปชฺชติ นอธิปติ\-ปจฺจยา. มคฺคเหตุเก จ มคฺคาธิปตี จ ขนฺเธ ปฏิจฺจ มคฺคาธิปติ อธิปติ. (2)}

\prose{\csromanfolio{389}มคฺคเหตุกญฺจ มคฺคา\-ธิปติญฺจ ธมฺมํ ปฏิจฺจ มคฺคเหตุโก จ มคฺคาธิปติ จ ธมฺมา อุปฺปชฺชนฺติ นอธิปติ\-ปจฺจยา. มคฺคเหตุเก จ มคฺคาธิปตี จ ขนฺเธ ปฏิจฺจ มคฺคเหตุโก จ มคฺคาธิปติ จ อธิปติ. (3)}

\csromanmatikaanchor{mtk.158}
\csromantocmarkref{subparagraph}{นปุเรชาตาทิ}{mtk.158}
\bookmark[dest={mtk.158},level=5]{นปุเรชาตาทิ}
\csromanheader{6}{\textbf{นปุเรชาตาทิ}}
\proseitem{14}{มคฺคารมฺมณํ ธมฺมํ ปฏิจฺจ มคฺคารมฺมโณ ธมฺโม อุปฺปชฺชติ นปุเรชาต\-ปจฺจยา. นปจฺฉาชาต\-ปจฺจยา. (ปริปุณฺณา เทฺวปิ.)}

\prose{นอาเสวน}

\proseitem{15}{มคฺคารมฺมณํ ธมฺมํ ปฏิจฺจ มคฺคารมฺมโณ ธมฺโม อุปฺปชฺชติ นอาเสวน\-ปจฺจยา. มคฺคารมฺมณํ เอกํ ขนฺธํ ปฏิจฺจ ตโย ขนฺธา ฯเปฯ เทฺว ขนฺเธ ฯเปฯ. (1)}

\prose{มคฺคารมฺมณํ ธมฺมํ ปฏิจฺจ มคฺคาธิปติ ธมฺโม อุปฺปชฺชติ นอาเสวน\-ปจฺจยา. มคฺคารมฺมณํ เอกํ ขนฺธํ ปฏิจฺจ มคฺคาธิปตี ตโย ขนฺธา ฯเปฯ เทฺว ขนฺเธ ฯเปฯ. (2)}

\prose{มคฺคารมฺมณํ ธมฺมํ ปฏิจฺจ มคฺคารมฺมโณ จ มคฺคาธิปติ จ ธมฺมา อุปฺปชฺชนฺติ นอาเสวน\-ปจฺจยา. มคฺคารมฺมณํ เอกํ ขนฺธํ ปฏิจฺจ มคฺคารมฺมณา จ มคฺคาธิปตี จ ตโย ขนฺธา ฯเปฯ เทฺว ขนฺเธ ฯเปฯ. (3)}

\proseitem{16}{มคฺคาธิปติํ ธมฺมํ ปฏิจฺจ มคฺคาธิปติ ธมฺโม อุปฺปชฺชติ นอาเสวน\-ปจฺจยา. มคฺคาธิปติํ เอกํ ขนฺธํ ปฏิจฺจ ตโย ขนฺธา ฯเปฯ เทฺว ขนฺเธ ฯเปฯ. (1)}

\prose{มคฺคาธิปติํ ธมฺมํ ปฏิจฺจ มคฺคารมฺมโณ ธมฺโม อุปฺปชฺชติ นอาเสวน\-ปจฺจยา. มคฺคาธิปติํ เอกํ ขนฺธํ ปฏิจฺจ มคฺคารมฺมณา ตโย ขนฺธา ฯเปฯ เทฺว ขนฺเธ ฯเปฯ. (2)}

\prose{มคฺคาธิปติํ ธมฺมํ ปฏิจฺจ มคฺคารมฺมโณ จ มคฺคาธิปติ จ ธมฺมา อุปฺปชฺชนฺติ นอาเสวน\-ปจฺจยา. มคฺคาธิปติํ เอกํ ขนฺธํ ปฏิจฺจ มคฺคารมฺมณา จ มคฺคาธิปตี จ ตโย ขนฺธา ฯเปฯ เทฺว ขนฺเธ ฯเปฯ. (3)}

\proseitem{17}{มคฺคา\-รมฺมณญฺจ มคฺคา\-ธิปติญฺจ ธมฺมํ ปฏิจฺจ มคฺคารมฺมโณ ธมฺโม อุปฺปชฺชติ นอาเสวน\-ปจฺจยา. มคฺคา\-รมฺมณญฺจ มคฺคา\-ธิปติญฺจ เอกํ ขนฺธํ ปฏิจฺจ มคฺคารมฺมณา ตโย ขนฺธา ฯเปฯ เทฺว ขนฺเธ ฯเปฯ. (1)}

\prose{\csromanfolio{390}มคฺคา\-รมฺมณญฺจ มคฺคา\-ธิปติญฺจ ธมฺมํ ปฏิจฺจ มคฺคาธิปติ ธมฺโม อุปฺปชฺชติ นอาเสวน\-ปจฺจยา. มคฺคา\-รมฺมณญฺจ มคฺคา\-ธิปติญฺจ เอกํ ขนฺธํ ปฏิจฺจ มคฺคาธิปตี ตโย ขนฺธา ฯเปฯ เทฺว ขนฺเธ ฯเปฯ. (2)}

\prose{มคฺคา\-รมฺมณญฺจ มคฺคา\-ธิปติญฺจ ธมฺมํ ปฏิจฺจ มคฺคารมฺมโณ จ มคฺคาธิปติ จ ธมฺมา อุปฺปชฺชนฺติ นอาเสวน\-ปจฺจยา. มคฺคา\-รมฺมณญฺจ มคฺคา\-ธิปติญฺจ เอกํ ขนฺธํ ปฏิจฺจ มคฺคารมฺมณา จ มคฺคาธิปตี จ ตโย ขนฺธา ฯเปฯ เทฺว ขนฺเธ ฯเปฯ. (3)}

\csromanheader{2}{\textbf{นกมฺม}}
\proseitem{18}{มคฺคารมฺมณํ ธมฺมํ ปฏิจฺจ มคฺคารมฺมโณ ธมฺโม อุปฺปชฺชติ นกมฺมปจฺจยา. มคฺคารมฺมเณ ขนฺเธ ปฏิจฺจ มคฺคารมฺมณา เจตนา. (1)}

\prose{มคฺคารมฺมณํ ธมฺมํ ปฏิจฺจ มคฺคาธิปติ ธมฺโม อุปฺปชฺชติ นกมฺมปจฺจยา. มคฺคารมฺมเณ ขนฺเธ ปฏิจฺจ มคฺคาธิปติ เจตนา. (2)}

\prose{มคฺคารมฺมณํ ธมฺมํ ปฏิจฺจ มคฺคารมฺมโณ จ มคฺคาธิปติ จ ธมฺมา อุปฺปชฺชนฺติ นกมฺมปจฺจยา. มคฺคารมฺมเณ ขนฺเธ ปฏิจฺจ มคฺคารมฺมณา จ มคฺคาธิปติ จ เจตนา. (3)}

\proseitem{19}{มคฺคเหตุกํ ธมฺมํ ปฏิจฺจ มคฺคเหตุโก ธมฺโม อุปฺปชฺชติ นกมฺมปจฺจยา. มคฺคเหตุเก ขนฺเธ ปฏิจฺจ มคฺคเหตุกา เจตนา. (1)}

\prose{มคฺคเหตุกํ ธมฺมํ ปฏิจฺจ มคฺคาธิปติ ธมฺโม อุปฺปชฺชติ นกมฺมปจฺจยา. มคฺคเหตุเก ขนฺเธ ปฏิจฺจ มคฺคาธิปติ เจตนา. (2)}

\prose{มคฺคเหตุกํ ธมฺมํ ปฏิจฺจ มคฺคเหตุโก จ มคฺคาธิปติ จ ธมฺมา อุปฺปชฺชนฺติ นกมฺมปจฺจยา. มคฺคเหตุเก ขนฺเธ ปฏิจฺจ มคฺคเหตุโก จ มคฺคาธิปติ จ เจตนา. (3)}

\proseitem{20}{มคฺคาธิปติํ ธมฺมํ ปฏิจฺจ มคฺคาธิปติ ธมฺโม อุปฺปชฺชติ นกมฺมปจฺจยา. มคฺคาธิปตี ขนฺเธ ปฏิจฺจ มคฺคาธิปติ เจตนา. (ปญฺจ ปญฺหา.)}

\prose{มคฺคา\-รมฺมณญฺจ มคฺคา\-ธิปติญฺจ ธมฺมํ ปฏิจฺจ มคฺคารมฺมโณ ธมฺโม อุปฺปชฺชติ นกมฺมปจฺจยา. (ปฐมฆฏเน ตีณิ.)}

\csromanmatikaanchor{mtk.159}
\csromanmatikaanchor{mtk.160}
\csromanmatikaanchor{mtk.161}
\csromantocmarkref{subparagraph}{2. ปจฺจยปจฺจนีย 2. สงฺขฺยาวาร}{mtk.159}
\bookmark[dest={mtk.159},level=5]{2. ปจฺจยปจฺจนีย 2. สงฺขฺยาวาร}
\csromantocmarknopage{subparagraph}{3. ปจฺจยานุโลมปจฺจนีย}{mtk.160}
\bookmark[dest={mtk.160},level=5]{3. ปจฺจยานุโลมปจฺจนีย}
\csromantocmarkref{subparagraph}{เหตุทุก}{mtk.161}
\bookmark[dest={mtk.161},level=5]{เหตุทุก}
\prose{\csromanfolio{391}มคฺคเหตุกญฺจ มคฺคา\-ธิปติญฺจ ธมฺมํ ปฏิจฺจ มคฺคเหตุโก ธมฺโม อุปฺปชฺชติ นกมฺมปจฺจยา. (ทุติยฆฏเน ตีณิ ปญฺหา.)}

\proseitem{21}{มคฺคารมฺมณํ ธมฺมํ ปฏิจฺจ มคฺคารมฺมโณ ธมฺโม อุปฺปชฺชติ นวิปาก\-ปจฺจยา. (ปริปุณฺณํ.)}

\csromanheader{2}{\textbf{นมคฺค}}
\proseitem{22}{มคฺคารมฺมณํ ธมฺมํ ปฏิจฺจ มคฺคารมฺมโณ ธมฺโม อุปฺปชฺชติ นมคฺคปจฺจยา. อเหตุกํ มคฺคารมฺมณํ เอกํ ขนฺธํ ปฏิจฺจ ตโย ขนฺธา ฯเปฯ เทฺว ขนฺเธ ฯเปฯ. (1)}

\csromanheader{2}{\textbf{นวิปฺปยุตฺต}}
\proseitem{23}{มคฺคารมฺมณํ ธมฺมํ ปฏิจฺจ มคฺคารมฺมโณ ธมฺโม อุปฺปชฺชติ นวิปฺปยุตฺต\-ปจฺจยา. (ปริปุณฺณํ. อรูปนฺติ นิยาเมตพฺพํ.)}

\csromanmarkh{2. ปจฺจย\-ปจฺจนีย}\csromanmarkhii{2. สงฺขฺยาวาร}\csromanheadernomark{6}{2. ปจฺจย\-ปจฺจนีย 2. สงฺขฺยาวาร}
\proseitemwithcloserruled{24}{นเหตุยา เอกํ, นอธิปติยา สตฺตรส, นปุเรชาเต สตฺตรส, นปจฺฉาชาเต สตฺตรส, นอาเสวเน นว, นกมฺเม สตฺตรส, นวิปาเก สตฺตรส, นมคฺเค เอกํ, นวิปฺปยุตฺเต สตฺตรส. (เอวํ คเณตพฺพํ.)}{ปจฺจนียํ.}
\proseitem{25}{เหตุปจฺจยา นอธิปติยา สตฺตรส, นปุเรชาเต สตฺตรส, นปจฺฉาชาเต สตฺตรส, นอาเสวเน นว, นกมฺเม สตฺตรส, นวิปาเก, นวิปฺปยุตฺเต สตฺตรส. (เอวํ คเณตพฺพํ.)}

\vspace{\tipitakaparskip}
\csromancenter{อนุโลม\-ปจฺจนียํ.}
\proseitem{26}{\csromanfolio{392}นเหตุปจฺจยา อารมฺมเณ เอกํ, อนนฺตเร เอกํ, สมนนฺตเร เอกํ ฯเปฯ ฌาเน, สมฺปยุตฺเต, วิปฺปยุตฺเต, อตฺถิยา, นตฺถิยา, วิคเต, อวิคเต เอกํ. (เอวํ คเณตพฺพํ.) ปจฺจยานุโลมํ. ปฏิจฺจวาโร.}

\prose{(สหชาตวาโรปิ ปจฺจยวาโรปิ นิสฺสยวาโรปิ สํสฏฺฐ\-วาโรปิ สมฺปยุตฺตวาโรปิ ปฏิจฺจ\-วาร\-สทิโส.)}
\csromansectionrule

\csromanheader{2}{16. มคฺคา\-รมฺมณ\-ตฺติก}
\csromanheader{2}{7. \textbf{ปญฺหาวาร}}
\csromanmarkh{1. ปจฺจยานุโลม}\csromanmarkhii{1. วิภงฺควาร}\csromanheadernomark{2}{1. \textbf{ปจฺจยานุโลม 1. วิภงฺควาร}}
\csromanmatikaanchor{mtk.162}
\csromanmatikaanchor{mtk.163}
\csromanmatikaanchor{mtk.165}
\csromanmatikaanchor{mtk.166}
\csromantocmarknopage{subparagraph}{4. ปจฺจยปจฺจนียานุโลม}{mtk.162}
\bookmark[dest={mtk.162},level=5]{4. ปจฺจยปจฺจนียานุโลม}
\csromantocmarkref{subparagraph}{นเหตุทุก}{mtk.163}
\bookmark[dest={mtk.163},level=5]{นเหตุทุก}
\csromantocmarknopage{subparagraph}{16. มคฺคารมฺมณตฺติก 7. ปญฺหาวาร}{mtk.164}
\bookmark[dest={mtk.164},level=5]{16. มคฺคารมฺมณตฺติก 7. ปญฺหาวาร}
\csromantocmarknopage{subparagraph}{1. ปจฺจยานุโลม 1. วิภงฺควาร}{mtk.165}
\bookmark[dest={mtk.165},level=5]{1. ปจฺจยานุโลม 1. วิภงฺควาร}
\csromantocmarkref{subparagraph}{เหตุ}{mtk.166}
\bookmark[dest={mtk.166},level=5]{เหตุ}
\csromanheader{6}{\textbf{เหตุ}}
\proseitem{27}{มคฺคารมฺมโณ ธมฺโม มคฺคา\-รมฺมณสฺส ธมฺมสฺส เหตุปจฺจเยน ปจฺจโย. มคฺคารมฺมณา เหตู สมฺปยุตฺตกานํ ขนฺธานํ เหตุปจฺจเยน ปจฺจโย. (1)}

\prose{มคฺคารมฺมโณ ธมฺโม มคฺคา\-ธิปติสฺส ธมฺมสฺส เหตุปจฺจเยน ปจฺจโย. มคฺคารมฺมณา เหตู สมฺปยุตฺตกานํ มคฺคา\-ธิปตีนํ ขนฺธานํ เหตุปจฺจเยน ปจฺจโย. (2)}

\prose{มคฺคารมฺมโณ ธมฺโม มคฺคา\-รมฺมณสฺส จ มคฺคา\-ธิปติสฺส จ ธมฺมสฺส เหตุปจฺจเยน ปจฺจโย. (อิมินา การเณน สตฺตรส ปญฺหา กาตพฺพา.)}

\csromanheader{2}{\textbf{อารมฺมณ}}
\proseitem{28}{มคฺคเหตุโก ธมฺโม มคฺคา\-รมฺมณสฺส ธมฺมสฺส อารมฺมณ\-ปจฺจเยน ปจฺจโย. อริยา มคฺคา วุฏฺฐหิตฺวา มคฺคํ ปจฺจเวกฺขนฺติ. เจโตปริย\-ญาเณน \csromanfolio{393}มคฺค\-เหตุก\-จิตฺต\-สมงฺคิสฺส จิตฺตํ ชานนฺติ. มคฺคเหตุกา ขนฺธา เจโตปริย\-ญาณสฺส,ปุพฺเพ\-นิวาสา\-นุสฺสติ\-ญาณสฺส, อนาคตํสญาณสฺส, อาวชฺชนาย อารมฺมณ\-ปจฺจเยน ปจฺจโย. (1)}

\prose{มคฺคเหตุโก ธมฺโม มคฺคา\-ธิปติสฺส ธมฺมสฺส อารมฺมณ\-ปจฺจเยน ปจฺจโย. อริยา มคฺคา วุฏฺฐหิตฺวา มคฺคํ ครุํ กตฺวา ปจฺจเวกฺขนฺติ. (2)}

\prose{มคฺคเหตุโก ธมฺโม มคฺคา\-รมฺมณสฺส จ มคฺคา\-ธิปติสฺส จ ธมฺมสฺส อารมฺมณ\-ปจฺจเยน ปจฺจโย. อริยา มคฺคา วุฏฺฐหิตฺวา มคฺคํ ครุํ กตฺวา ปจฺจเวกฺขนฺติ. (3)}

\proseitem{29}{มคฺคาธิปติ ธมฺโม มคฺคา\-ธิปติสฺส ธมฺมสฺส อารมฺมณ\-ปจฺจเยน ปจฺจโย. อริยา มคฺคา วุฏฺฐหิตฺวา มคฺคํ ครุํ กตฺวา ปจฺจเวกฺขนฺติ. (1)}

\prose{มคฺคาธิปติ ธมฺโม มคฺคา\-รมฺมณสฺส ธมฺมสฺส อารมฺมณ\-ปจฺจเยน ปจฺจโย. อริยา มคฺคา วุฏฺฐหิตฺวา มคฺคํ ปจฺจเวกฺขนฺติ, เจโตปริย\-ญาเณน มคฺคา\-ธิป\-ติจิตฺต\-สมงฺคิสฺส จิตฺตํ ชานนฺติ, มคฺคาธิปตี ขนฺธา เจโตปริย\-ญาณสฺส, ปุพฺเพ\-นิวาสา\-นุสฺสติ\-ญาณสฺส, อนาคตํสญาณสฺส, อาวชฺชนาย อารมฺมณ\-ปจฺจเยน ปจฺจโย. (2)}

\prose{มคฺคาธิปติ ธมฺโม มคฺคา\-รมฺมณสฺส จ มคฺคา\-ธิปติสฺส จ ธมฺมสฺส อารมฺมณ\-ปจฺจเยน ปจฺจโย. อริยา มคฺคา วุฏฺฐหิตฺวา มคฺคํ ครุํ กตฺวา ปจฺจเวกฺขนฺติ. (3)}

\proseitem{30}{มคฺคเหตุโก จ มคฺคาธิปติ จ ธมฺมา มคฺคา\-รมฺมณสฺส ธมฺมสฺส อารมฺมณ\-ปจฺจเยน ปจฺจโย. อริยา มคฺคา วุฏฺฐหิตฺวา มคฺคํ ปจฺจเวกฺขนฺติ, เจโตปริย\-ญาเณน มคฺค\-เหตุก\-มคฺคา\-ธิป\-ติจิตฺต\-สมงฺคิสฺส จิตฺตํ ชานนฺติ. มคฺคเหตุกา จ มคฺคาธิปตี จ ขนฺธา เจโตปริย\-ญาณสฺส, ปุพฺเพ\-นิวาสา\-นุสฺสติ\-ญาณสฺส, อนาคตํสญาณสฺส, อาวชฺชนาย อารมฺมณ\-ปจฺจเยน ปจฺจโย. (1)}

\prose{มคฺคเหตุโก จ มคฺคาธิปติ จ ธมฺมา มคฺคา\-ธิปติสฺส ธมฺมสฺส อารมฺมณ\-ปจฺจเยน ปจฺจโย. อริยา มคฺคา วุฏฺฐหิตฺวา มคฺคํ ครุํ กตฺวา ปจฺจเวกฺขนฺติ. (2)}

\prose{มคฺคเหตุโก จ มคฺคาธิปติ จ ธมฺมา มคฺคา\-รมฺมณสฺส จ มคฺคา\-ธิปติสฺส จ ธมฺมสฺส อารมฺมณ\-ปจฺจเยน ปจฺจโย. อริยา มคฺคา วุฏฺฐหิตฺวา มคฺคํ ครุํ กตฺวา ปจฺจเวกฺขนฺติ. (3)}

\csromanmatikaanchor{mtk.167}
\csromantocmarkref{subparagraph}{อธิปติ}{mtk.167}
\bookmark[dest={mtk.167},level=5]{อธิปติ}
\proseitem{31}{\csromanfolio{394}มคฺคารมฺมโณ ธมฺโม มคฺคา\-รมฺมณสฺส ธมฺมสฺส อธิปติ\-ปจฺจเยน ปจฺจโย. สหชาตา\-ธิปติ\csromandash{}มคฺคา\-รมฺมณา\-ธิปติ สมฺปยุตฺตกานํ ขนฺธานํ อธิปติ\-ปจฺจเยน ปจฺจโย. (1)}

\prose{มคฺคารมฺมโณ ธมฺโม มคฺคา\-ธิปติสฺส ธมฺมสฺส อธิปติ\-ปจฺจเยน ปจฺจโย. สหชาตา\-ธิปติ\csromandash{}มคฺคา\-รมฺมณา\-ธิปติ สมฺปยุตฺตกานํ มคฺคา\-ธิปตีนํ ขนฺธานํ อธิปติ\-ปจฺจเยน ปจฺจโย. (2)}

\prose{มคฺคารมฺมโณ ธมฺโม มคฺคา\-รมฺมณสฺส จ มคฺคา\-ธิปติสฺส จ ธมฺมสฺส อธิปติ\-ปจฺจเยน ปจฺจโย. สหชาตา\-ธิปติ\csromandash{} มคฺคา\-รมฺมณา\-ธิปติ สมฺปยุตฺตกานํ มคฺคา\-รมฺมณานญฺจ มคฺคา\-ธิปตีนญฺจ ขนฺธานํ อธิปติ\-ปจฺจเยน ปจฺจโย. (3)}

\proseitem{32}{มคฺคเหตุโก ธมฺโม มคฺคเหตุกสฺส ธมฺมสฺส อธิปติ\-ปจฺจเยน ปจฺจโย. สหชาตา\-ธิปติ\csromandash{}มคฺค\-เหตุกา\-ธิปติ สมฺปยุตฺตกานํ ขนฺธานํ อธิปติ\-ปจฺจเยน ปจฺจโย. (1)}

\prose{มคฺคเหตุโก ธมฺโม มคฺคา\-รมฺมณสฺส ธมฺมสฺส อธิปติ\-ปจฺจเยน ปจฺจโย. อารมฺมณา\-ธิปติ\csromandash{}อริยา มคฺคา วุฏฺฐหิตฺวา มคฺคํ ครุํ กตฺวา ปจฺจเวกฺขนฺติ. (2)}

\prose{มคฺคเหตุโก ธมฺโม มคฺคา\-ธิปติสฺส ธมฺมสฺส อธิปติ\-ปจฺจเยน ปจฺจโย. อารมฺมณา\-ธิปติ, สหชาตา\-ธิปติ.  อารมฺมณา\-ธิปติ\csromandash{}อริยา มคฺคา วุฏฺฐหิตฺวา มคฺคํ ครุํ กตฺวา ปจฺจเวกฺขนฺติ.  สหชาตา\-ธิปติ\csromandash{}มคฺค\-เหตุกา\-ธิปติ สมฺปยุตฺตกานํ มคฺคา\-ธิปตีนํ ขนฺธานํ อธิปติ\-ปจฺจเยน ปจฺจโย. (3)}

\prose{มคฺคเหตุโก ธมฺโม มคฺคา\-รมฺมณสฺส จ มคฺคา\-ธิปติสฺส จ ธมฺมสฺส อธิปติ\-ปจฺจเยน ปจฺจโย. อารมฺมณา\-ธิปติ\csromandash{}อริยา มคฺคา วุฏฺฐหิตฺวา มคฺคํ ครุํ กตฺวา ปจฺจเวกฺขนฺติ. (4)}

\prose{มคฺคเหตุโก ธมฺโม มคฺคเหตุกสฺส จ มคฺคา\-ธิปติสฺส จ ธมฺมสฺส อธิปติ\-ปจฺจเยน ปจฺจโย. สหชาตา\-ธิปติ\csromandash{} มคฺค\-เหตุกา\-ธิปติ สมฺปยุตฺตกานํ มคฺค\-เหตุกานญฺจ มคฺคา\-ธิปตีนญฺจ ขนฺธานํ อธิปติ\-ปจฺจเยน ปจฺจโย. (5)}

\proseitem{33}{\csromanfolio{395}มคฺคาธิปติ ธมฺโม มคฺคา\-ธิปติสฺส ธมฺมสฺส อธิปติ\-ปจฺจเยน ปจฺจโย. อารมฺมณา\-ธิปติ, สหชาตา\-ธิปติ.  อารมฺมณา\-ธิปติ\csromandash{}อริยา มคฺคา วุฏฺฐหิตฺวา มคฺคํ ครุํ กตฺวา ปจฺจเวกฺขนฺติ.  สหชาตา\-ธิปติ\csromandash{}มคฺคาธิปติ อธิปติ สมฺปยุตฺตกานํ ขนฺธานํ อธิปติ\-ปจฺจเยน ปจฺจโย. (1)}

\prose{มคฺคาธิปติ ธมฺโม มคฺคา\-รมฺมณสฺส ธมฺมสฺส อธิปติ\-ปจฺจเยน ปจฺจโย. อารมฺมณา\-ธิปติ, สหชาตา\-ธิปติ.  อารมฺมณา\-ธิปติ\csromandash{}อริยา มคฺคา วุฏฺฐหิตฺวา มคฺคํ ครุํ กตฺวา ปจฺจเวกฺขนฺติ.  สหชาตา\-ธิปติ\csromandash{}มคฺคาธิปติ อธิปติ สมฺปยุตฺตกานํ มคฺคา\-รมฺมณานํ ขนฺธานํ อธิปติ\-ปจฺจเยน ปจฺจโย. (2)}

\prose{มคฺคาธิปติ ธมฺโม มคฺคเหตุกสฺส ธมฺมสฺส อธิปติ\-ปจฺจเยน ปจฺจโย. สหชาตา\-ธิปติ\csromandash{}มคฺคาธิปติ อธิปติ สมฺปยุตฺตกานํ มคฺคเหตุกานํ ขนฺธานํ อธิปติ\-ปจฺจเยน ปจฺจโย. (3)}

\prose{มคฺคาธิปติ ธมฺโม มคฺคา\-รมฺมณสฺส จ มคฺคา\-ธิปติสฺส จ อธิปติ\-ปจฺจเยน ปจฺจโย. อารมฺมณา\-ธิปติ, สหชาตา\-ธิปติ.  อารมฺมณา\-ธิปติ\csromandash{}อริยา มคฺคา วุฏฺฐหิตฺวา มคฺคํ ครุํ กตฺวา ปจฺจเวกฺขนฺติ.  สหชาตา\-ธิปติ\csromandash{}มคฺคาธิปติ อธิปติ สมฺปยุตฺตกานํ มคฺคา\-รมฺมณานญฺจ มคฺคา\-ธิปตีนญฺจ ขนฺธานํ อธิปติ\-ปจฺจเยน ปจฺจโย. (4)}

\prose{มคฺคาธิปติ ธมฺโม มคฺคเหตุกสฺส จ มคฺคา\-ธิปติสฺส จ ธมฺมสฺส อธิปติ\-ปจฺจเยน ปจฺจโย. สหชาตา\-ธิปติ\csromandash{}มคฺคาธิปติ อธิปติ สมฺปยุตฺตกานํ มคฺคา\-เหตุกานญฺจ มคฺคา\-ธิปตีนญฺจ ขนฺธานํ อธิปติ\-ปจฺจเยน ปจฺจโย. (5)}

\proseitem{34}{มคฺคารมฺมโณ จ มคฺคาธิปติ จ ธมฺมา มคฺคา\-รมฺมณสฺส ธมฺมสฺส อธิปติ\-ปจฺจเยน ปจฺจโย. สหชาตา\-ธิปติ\csromandash{}มคฺคารมฺมณา จ มคฺคาธิปตี จ อธิปติ สมฺปยุตฺตกานํ มคฺคา\-รมฺมณานํ ขนฺธานํ อธิปติ\-ปจฺจเยน ปจฺจโย. (1)}

\prose{มคฺคารมฺมโณ จ มคฺคาธิปติ จ ธมฺมา มคฺคา\-ธิปติสฺส ธมฺมสฺส อธิปติ\-ปจฺจเยน ปจฺจโย. สหชาตา\-ธิปติ\csromandash{}มคฺคารมฺมณา จ มคฺคาธิปตี จ อธิปติ สมฺปยุตฺตกานํ มคฺคา\-ธิปตีนํ ขนฺธานํ อธิปติ\-ปจฺจเยน ปจฺจโย. (2)}

\prose{มคฺคารมฺมโณ จ มคฺคาธิปติ จ ธมฺมา มคฺคา\-รมฺมณสฺส จ มคฺคา\-ธิปติสฺส จ ธมฺมสฺส อธิปติ\-ปจฺจเยน ปจฺจโย. สหชาตา\-ธิปติ\csromandash{}มคฺคารมฺมณา จ มคฺคาธิปตี จ อธิปติ สมฺปยุตฺตกานํ มคฺคา\-รมฺมณานญฺจ มคฺคา\-ธิปตีนญฺจ ขนฺธานํ อธิปติ\-ปจฺจเยน ปจฺจโย. (3)}

\proseitem{35}{\csromanfolio{396}มคฺคเหตุโก จ มคฺคาธิปติ จ ธมฺมา มคฺคา\-รมฺมณสฺส ธมฺมสฺส อธิปติ\-ปจฺจเยน ปจฺจโย. อารมฺมณา\-ธิปติ\csromandash{}อริยา มคฺคา วุฏฺฐหิตฺวา มคฺคํ ครุํ กตฺวา ปจฺจเวกฺขนฺติ. (1)}

\prose{มคฺคเหตุโก จ มคฺคาธิปติ จ ธมฺมา มคฺคเหตุกสฺส ธมฺมสฺส อธิปติ\-ปจฺจเยน ปจฺจโย. สหชาตา\-ธิปติ\csromandash{}มคฺคเหตุกา จ มคฺคาธิปตี จ อธิปติ สมฺปยุตฺตกานํ มคฺคเหตุกานํ ขนฺธานํ อธิปติ\-ปจฺจเยน ปจฺจโย. (2)}

\prose{มคฺคเหตุโก จ มคฺคาธิปติ จ ธมฺมา มคฺคา\-ธิปติสฺส ธมฺมสฺส อธิปติ\-ปจฺจเยน ปจฺจโย. อารมฺมณา\-ธิปติ, สหชาตา\-ธิปติ.  อารมฺมณา\-ธิปติ\csromandash{}อริยา มคฺคา วุฏฺฐหิตฺวา มคฺคํ ครุํ กตฺวา ปจฺจเวกฺขนฺติ. สหชาตา\-ธิปติ\csromandash{}มคฺคเหตุกา จ มคฺคาธิปตี จ อธิปติ สมฺปยุตฺตกานํ มคฺคา\-ธิปตีนํ ขนฺธานํ อธิปติ\-ปจฺจเยน ปจฺจโย. (3)}

\prose{มคฺคเหตุโก จ มคฺคาธิปติ จ ธมฺมา มคฺคา\-รมฺมณสฺส จ มคฺคา\-ธิปติสฺส จ ธมฺมสฺส อธิปติ\-ปจฺจเยน ปจฺจโย. อารมฺมณา\-ธิปติ\csromandash{}อริยา มคฺคา วุฏฺฐหิตฺวา มคฺคํ ครุํ กตฺวา ปจฺจเวกฺขนฺติ. (4)}

\prose{มคฺคเหตุโก จ มคฺคาธิปติ จ ธมฺมา มคฺคเหตุกสฺส จ มคฺคา\-ธิปติสฺส จ ธมฺมสฺส อธิปติ\-ปจฺจเยน ปจฺจโย. สหชาตา\-ธิปติ\csromandash{}มคฺคเหตุกา จ มคฺคาธิปตี จ อธิปติ สมฺปยุตฺตกานํ มคฺค\-เหตุกานญฺจ มคฺคา\-ธิปตีนญฺจ ขนฺธานํ อธิปติ\-ปจฺจเยน ปจฺจโย. (5)}

\csromanmatikaanchor{mtk.168}
\csromantocmarkref{subparagraph}{อนนฺตร}{mtk.168}
\bookmark[dest={mtk.168},level=5]{อนนฺตร}
\csromanheader{6}{\textbf{อนนฺตร}}
\proseitem{36}{มคฺคารมฺมโณ ธมฺโม มคฺคา\-รมฺมณสฺส ธมฺมสฺส อนนฺตร\-ปจฺจเยน ปจฺจโย. ปุริมา ปุริมา มคฺคารมฺมณา ขนฺธา ปจฺฉิมานํ ปจฺฉิมานํ มคฺคา\-รมฺมณานํ ขนฺธานํ อนนฺตร\-ปจฺจเยน ปจฺจโย. อาวชฺชนา มคฺคา\-รมฺมณานํ ขนฺธานํ อนนฺตร\-ปจฺจเยน ปจฺจโย. (1)}

\prose{มคฺคารมฺมโณ ธมฺโม มคฺคา\-ธิปติสฺส ธมฺมสฺส อนนฺตร\-ปจฺจเยน ปจฺจโย. ปุริมา ปุริมา มคฺคารมฺมณา ขนฺธา ปจฺฉิมานํ ปจฺฉิมานํ มคฺคา\-ธิปตีนํ ขนฺธานํ อนนฺตร\-ปจฺจเยน ปจฺจโย. อาวชฺชนา มคฺคา\-ธิปตีนํ ขนฺธานํ อนนฺตร\-ปจฺจเยน ปจฺจโย. (2)}

\prose{มคฺคารมฺมโณ ธมฺโม มคฺคา\-รมฺมณสฺส จ มคฺคา\-ธิปติสฺส จ ธมฺมสฺส อนนฺตร\-ปจฺจเยน ปจฺจโย. ปุริมา ปุริมา มคฺคารมฺมณา ขนฺธา ปจฺฉิมานํ ปจฺฉิมานํ \csromanfolio{397}มคฺคา\-รมฺมณานญฺจ มคฺคา\-ธิปตีนญฺจ ขนฺธานํ อนนฺตร\-ปจฺจเยน ปจฺจโย. อาวชฺชนา มคฺคา\-รมฺมณานญฺจ มคฺคา\-ธิปตีนญฺจ ขนฺธานํ อนนฺตร\-ปจฺจเยน ปจฺจโย. (3)}

\proseitem{37}{มคฺคาธิปติ ธมฺโม มคฺคา\-ธิปติสฺส ธมฺมสฺส อนนฺตร\-ปจฺจเยน ปจฺจโย. ปุริมา ปุริมา มคฺคาธิปตี ขนฺธา ปจฺฉิมานํ ปจฺฉิมานํ มคฺคา\-ธิปตีนํ ขนฺธานํ อนนฺตร\-ปจฺจเยน ปจฺจโย. (1)}

\prose{มคฺคาธิปติ ธมฺโม มคฺคา\-รมฺมณสฺส ธมฺมสฺส อนนฺตร\-ปจฺจเยน ปจฺจโย. ปุริมา ปุริมา มคฺคาธิปตี ขนฺธา ปจฺฉิมานํ ปจฺฉิมานํ มคฺคา\-รมฺมณานํ ขนฺธานํ อนนฺตร\-ปจฺจเยน ปจฺจโย. (2)}

\prose{มคฺคาธิปติ ธมฺโม มคฺคา\-รมฺมณสฺส จ มคฺคา\-ธิปติสฺส จ ธมฺมสฺส อนนฺตร\-ปจฺจเยน ปจฺจโย. ปุริมา ปุริมา มคฺคาธิปตี ขนฺธา ปจฺฉิมานํ ปจฺฉิมานํ มคฺคา\-รมฺมณานญฺจ มคฺคา\-ธิปตีนญฺจ ขนฺธานํ อนนฺตร\-ปจฺจเยน ปจฺจโย. (3)}

\proseitem{38}{มคฺคารมฺมโณ จ มคฺคาธิปติ จ ธมฺมา มคฺคา\-รมฺมณสฺส ธมฺมสฺส อนนฺตร\-ปจฺจเยน ปจฺจโย. ปุริมา ปุริมา มคฺคารมฺมณา จ มคฺคาธิปตี จ ขนฺธา ปจฺฉิมานํ ปจฺฉิมานํ มคฺคา\-รมฺมณานํ ขนฺธานํ อนนฺตร\-ปจฺจเยน ปจฺจโย. (1)}

\prose{มคฺคารมฺมโณ จ มคฺคาธิปติ จ ธมฺมา มคฺคา\-ธิปติสฺส ธมฺมสฺส อนนฺตร\-ปจฺจเยน ปจฺจโย. ปุริมา ปุริมา มคฺคารมฺมณา จ มคฺคาธิปตี จ ขนฺธา ปจฺฉิมานํ ปจฺฉิมานํ มคฺคา\-ธิปตีนํ ขนฺธานํ อนนฺตร\-ปจฺจเยน ปจฺจโย. (2)}

\prose{มคฺคารมฺมโณ จ มคฺคาธิปติ จ ธมฺมา มคฺคา\-รมฺมณสฺส จ มคฺคา\-ธิปติสฺส จ ธมฺมสฺส อนนฺตร\-ปจฺจเยน ปจฺจโย. ปุริมา ปุริมา มคฺคารมฺมณา จ มคฺคาธิปตี จ ขนฺธา ปจฺฉิมานํ ปจฺฉิมานํ มคฺคา\-รมฺมณานญฺจ มคฺคา\-ธิปตีนญฺจ ขนฺธานํ อนนฺตร\-ปจฺจเยน ปจฺจโย. (3)}

\csromanmatikaanchor{mtk.169}
\csromantocmarkref{subparagraph}{สมนนฺตราทิ}{mtk.169}
\bookmark[dest={mtk.169},level=5]{สมนนฺตราทิ}
\csromanheader{6}{\textbf{สมนนฺตราทิ}}
\proseitem{39}{มคฺคารมฺมโณ ธมฺโม มคฺคา\-รมฺมณสฺส ธมฺมสฺส สมนนฺตร\-ปจฺจเยน ปจฺจโย ฯเปฯ. (อนนฺตร\-สทิสํ) สหชาต\-ปจฺจเยน ปจฺจโย. อญฺญมญฺญ\-ปจฺจเยน ปจฺจโย. นิสฺสย\-ปจฺจเยน ปจฺจโย. (ตีสุปิ สตฺตรส ปญฺหา กาตพฺพา.)}

\csromanheader{2}{\textbf{อุปนิสฺสย}}
\proseitem{40}{\csromanfolio{398}มคฺคารมฺมโณ ธมฺโม มคฺคา\-รมฺมณสฺส ธมฺมสฺส อุปนิสฺสย\-ปจฺจเยน ปจฺจโย. อนนฺตรู\-ปนิสฺสโย, ปกตู\-ปนิสฺสโย ฯเปฯ. ปกตู\-ปนิสฺสโย\csromandash{}ปจฺจเวกฺขณา ปจฺจเวกฺขณาย อุปนิสฺสย\-ปจฺจเยน ปจฺจโย. (1)}

\prose{มคฺคารมฺมโณ ธมฺโม มคฺคา\-ธิปติสฺส ธมฺมสฺส อุปนิสฺสย\-ปจฺจเยน ปจฺจโย. อนนฺตรู\-ปนิสฺสโย, ปกตู\-ปนิสฺสโย ฯเปฯ. ปกตู\-ปนิสฺสโย\csromandash{}ปจฺจเวกฺขณา ปจฺจเวกฺขณาย อุปนิสฺสย\-ปจฺจเยน ปจฺจโย. (2)}

\prose{มคฺคารมฺมโณ ธมฺโม มคฺคา\-รมฺมณสฺส จ มคฺคา\-ธิปติสฺส จ ธมฺมสฺส อุปนิสฺสย\-ปจฺจเยน ปจฺจโย. อนนฺตรู\-ปนิสฺสโย, ปกตู\-ปนิสฺสโย ฯเปฯ. ปกตู\-ปนิสฺสโย\csromandash{}ปจฺจเวกฺขณา ปจฺจเวกฺขณาย อุปนิสฺสย\-ปจฺจเยน ปจฺจโย. (3)}

\proseitem{41}{มคฺคเหตุโก ธมฺโม มคฺคเหตุกสฺส ธมฺมสฺส อุปนิสฺสย\-ปจฺจเยน ปจฺจโย. ปกตู\-ปนิสฺสโย\csromandash{}ปฐโม มคฺโค ทุติยสฺส มคฺคสฺส อุปนิสฺสย\-ปจฺจเยน ปจฺจโย ฯเปฯ. ตติโย มคฺโค จตุตฺถสฺส มคฺคสฺส อุปนิสฺสย\-ปจฺจเยน ปจฺจโย. (1)}

\prose{มคฺคเหตุโก ธมฺโม มคฺคา\-รมฺมณสฺส ธมฺมสฺส อุปนิสฺสย\-ปจฺจเยน ปจฺจโย. อารมฺมณู\-ปนิสฺสโย\csromandash{}อริยา มคฺคา วุฏฺฐหิตฺวา มคฺคํ ครุํ กตฺวา ปจฺจเวกฺขนฺติ. (2)}

\prose{มคฺคเหตุโก ธมฺโม มคฺคา\-ธิปติสฺส ธมฺมสฺส อุปนิสฺสย\-ปจฺจเยน ปจฺจโย. อารมฺมณู\-ปนิสฺสโย, ปกตู\-ปนิสฺสโย ฯเปฯ. ปกตู\-ปนิสฺสโย\csromandash{}ปฐโม มคฺโค ทุติยสฺส มคฺคสฺส ฯเปฯ. ตติโย มคฺโค จตุตฺถสฺส มคฺคสฺส อุปนิสฺสย\-ปจฺจเยน ปจฺจโย. (3)}

\prose{มคฺคเหตุโก ธมฺโม มคฺคา\-รมฺมณสฺส จ มคฺคา\-ธิปติสฺส จ ธมฺมสฺส อุปนิสฺสย\-ปจฺจเยน ปจฺจโย. อารมฺมณู\-ปนิสฺสโย\csromandash{}อริยา มคฺคา วุฏฺฐหิตฺวา มคฺคํ ครุํ กตฺวา ปจฺจเวกฺขนฺติ. (4)}

\prose{มคฺคเหตุโก ธมฺโม มคฺคเหตุกสฺส จ มคฺคา\-ธิปติสฺส จ ธมฺมสฺส อุปนิสฺสย\-ปจฺจเยน ปจฺจโย. ปกตู\-ปนิสฺสโย\csromandash{}ปฐโม มคฺโค ทุติยสฺส มคฺคสฺส ฯเปฯ. ตติโย มคฺโค จตุตฺถสฺส มคฺคสฺส อุปนิสฺสย\-ปจฺจเยน ปจฺจโย. (5)}

\proseitem{42}{\csromanfolio{399}มคฺคาธิปติ ธมฺโม มคฺคา\-ธิปติสฺส ธมฺมสฺส อุปนิสฺสย\-ปจฺจเยน ปจฺจโย. อารมฺมณู\-ปนิสฺสโย, อนนฺตรู\-ปนิสฺสโย, ปกตู\-ปนิสฺสโย ฯเปฯ. ปกตู\-ปนิสฺสโย\csromandash{}ปฐโม มคฺโค ทุติยสฺส มคฺคสฺส ฯเปฯ. ตติโย มคฺโค จตุตฺถสฺส มคฺคสฺส อุปนิสฺสย\-ปจฺจเยน ปจฺจโย. ปจฺจเวกฺขณา ปจฺจเวกฺขณาย อุปนิสฺสย\-ปจฺจเยน ปจฺจโย. (1)}

\prose{มคฺคาธิปติ ธมฺโม มคฺคา\-รมฺมณสฺส ธมฺมสฺส อุปนิสฺสย\-ปจฺจเยน ปจฺจโย. อารมฺมณู\-ปนิสฺสโย, อนนฺตรู\-ปนิสฺสโย, ปกตู\-ปนิสฺสโย ฯเปฯ. ปกตู\-ปนิสฺสโย\csromandash{}ปจฺจเวกฺขณา ปจฺจเวกฺขณาย อุปนิสฺสย\-ปจฺจเยน ปจฺจโย. (2)}

\prose{มคฺคาธิปติ ธมฺโม มคฺคเหตุกสฺส ธมฺมสฺส อุปนิสฺสย\-ปจฺจเยน ปจฺจโย. ปกตู\-ปนิสฺสโย\csromandash{}ปฐโม มคฺโค ทุติยสฺส มคฺคสฺส ฯเปฯ. ตติโย มคฺโค จตุตฺถสฺส มคฺคสฺส อุปนิสฺสย\-ปจฺจเยน ปจฺจโย. (3)}

\prose{มคฺคาธิปติ ธมฺโม มคฺคา\-รมฺมณสฺส จ มคฺคา\-ธิปติสฺส จ ธมฺมสฺส อุปนิสฺสย\-ปจฺจเยน ปจฺจโย. อารมฺมณู\-ปนิสฺสโย, อนนฺตรู\-ปนิสฺสโย, ปกตู\-ปนิสฺสโย ฯเปฯ. ปกตู\-ปนิสฺสโย\csromandash{}ปจฺจเวกฺขณา ปจฺจเวกฺขณาย อุปนิสฺสย\-ปจฺจเยน ปจฺจโย. (4)}

\prose{มคฺคาธิปติ ธมฺโม มคฺคเหตุกสฺส จ มคฺคา\-ธิปติสฺส จ ธมฺมสฺส อุปนิสฺสย\-ปจฺจเยน ปจฺจโย. ปกตู\-ปนิสฺสโย\csromandash{}ปฐโม มคฺโค ทุติยสฺส มคฺคสฺส ฯเปฯ. ตติโย มคฺโค จตุตฺถสฺส มคฺคสฺส อุปนิสฺสย\-ปจฺจเยน ปจฺจโย. (5)}

\proseitem{43}{มคฺคารมฺมโณ จ มคฺคาธิปติ จ ธมฺมา มคฺคา\-รมฺมณสฺส ธมฺมสฺส อุปนิสฺสย\-ปจฺจเยน ปจฺจโย. อนนฺตรู\-ปนิสฺสโย, ปกตู\-ปนิสฺสโย ฯเปฯ. ปกตู\-ปนิสฺสโย\csromandash{}ปจฺจเวกฺขณา ปจฺจเวกฺขณาย อุปนิสฺสย\-ปจฺจเยน ปจฺจโย. (1)}

\prose{มคฺคารมฺมโณ จ มคฺคาธิปติ จ ธมฺมา มคฺคา\-ธิปติสฺส ธมฺมสฺส อุปนิสฺสย\-ปจฺจเยน ปจฺจโย. อนนฺตรู\-ปนิสฺสโย, ปกตู\-ปนิสฺสโย ฯเปฯ. ปกตู\-ปนิสฺสโย\csromandash{}ปจฺจเวกฺขณา ปจฺจเวกฺขณาย อุปนิสฺสย\-ปจฺจเยน ปจฺจโย. (2)}

\prose{\csromanfolio{400}มคฺคารมฺมโณ จ มคฺคาธิปติ จ ธมฺมา มคฺคา\-รมฺมณสฺส จ มคฺคา\-ธิปติสฺส จ ธมฺมสฺส อุปนิสฺสย\-ปจฺจเยน ปจฺจโย. อนนฺตรู\-ปนิสฺสโย, ปกตู\-ปนิสฺสโย ฯเปฯ. ปกตู\-ปนิสฺสโย\csromandash{} ปจฺจเวกฺขณา ปจฺจเวกฺขณาย อุปนิสฺสย\-ปจฺจเยน ปจฺจโย. (3)}

\proseitem{44}{มคฺคเหตุโก จ มคฺคาธิปติ จ ธมฺมา มคฺคา\-รมฺมณสฺส ธมฺมสฺส อุปนิสฺสย\-ปจฺจเยน ปจฺจโย. อารมฺมณู\-ปนิสฺสโย\csromandash{}อริยา มคฺคา วุฏฺฐหิตฺวา มคฺคํ ครุํ กตฺวา ปจฺจเวกฺขนฺติ. (1)}

\prose{มคฺคเหตุโก จ มคฺคาธิปติ จ ธมฺมา มคฺคเหตุกสฺส ธมฺมสฺส อุปนิสฺสย\-ปจฺจเยน ปจฺจโย. ปกตู\-ปนิสฺสโย\csromandash{}ปฐโม มคฺโค ทุติยสฺส มคฺคสฺส ฯเปฯ. ตติโย มคฺโค จตุตฺถสฺส มคฺคสฺส อุปนิสฺสย\-ปจฺจเยน ปจฺจโย. (2)}

\prose{มคฺคเหตุโก จ มคฺคาธิปติ จ ธมฺมา มคฺคา\-ธิปติสฺส ธมฺมสฺส อุปนิสฺสย\-ปจฺจเยน ปจฺจโย. อารมฺมณู\-ปนิสฺสโย, ปกตู\-ปนิสฺสโย ฯเปฯ. ปกตู\-ปนิสฺสโย\csromandash{}ปฐโม มคฺโค ทุติยสฺส มคฺคสฺส ฯเปฯ. ตติโย มคฺโค จตุตฺถสฺส มคฺคสฺส อุปนิสฺสย\-ปจฺจเยน ปจฺจโย. (3)}

\prose{มคฺคเหตุโก จ มคฺคาธิปติ จ ธมฺมา มคฺคา\-รมฺมณสฺส จ มคฺคา\-ธิปติสฺส จ ธมฺมสฺส อุปนิสฺสย\-ปจฺจเยน ปจฺจโย. อารมฺมณู\-ปนิสฺสโย\csromandash{}อริยา มคฺคา วุฏฺฐหิตฺวา มคฺคํ ครุํ กตฺวา ปจฺจเวกฺขนฺติ. (4)}

\prose{มคฺคเหตุโก จ มคฺคาธิปติ จ ธมฺมา มคฺคเหตุกสฺส จ มคฺคา\-ธิปติสฺส จ ธมฺมสฺส อุปนิสฺสย\-ปจฺจเยน ปจฺจโย. ปกตู\-ปนิสฺสโย\csromandash{}ปฐโม มคฺโค ทุติยสฺส มคฺคสฺส ฯเปฯ. ตติโย มคฺโค จตุตฺถสฺส มคฺคสฺส อุปนิสฺสย\-ปจฺจเยน ปจฺจโย. (5)}

\csromanheader{2}{\textbf{อาเสวน}}
\proseitem{45}{มคฺคารมฺมโณ ธมฺโม มคฺคา\-รมฺมณสฺส ธมฺมสฺส อาเสวน\-ปจฺจเยน ปจฺจโย. ปุริมา ปุริมา มคฺคารมฺมณา ขนฺธา ปจฺฉิมานํ ปจฺฉิมานํ มคฺคา\-รมฺมณานํ ขนฺธานํ อาเสวน\-ปจฺจเยน ปจฺจโย.}

\prose{มคฺคารมฺมโณ ธมฺโม มคฺคา\-ธิปติสฺส ธมฺมสฺส อาเสวน\-ปจฺจเยน ปจฺจโย.}

\vspace{\tipitakaparskip}
\csromancenter{(อนนฺตร\-สทิสํ, นว ปญฺหา กาตพฺพา, อาวชฺชนา น กาตพฺพา.)}
\csromanheader{2}{16. มคฺคา\-รมฺมณ\-ตฺติก กมฺมาทิ}
\proseitem{46}{\csromanfolio{401}มคฺคารมฺมโณ ธมฺโม มคฺคา\-รมฺมณสฺส ธมฺมสฺส กมฺมปจฺจเยน ปจฺจโย. สหชาตา ฯเปฯ.}

\vspace{\tipitakaparskip}
\csromancenter{(นานากฺขณิกา นตฺถิ, สตฺตรส ปญฺหา กาตพฺพา.)}
\csromanheader{2}{\textbf{อาหาราทิ}}
\proseitem{47}{มคฺคารมฺมโณ ธมฺโม มคฺคา\-รมฺมณสฺส ธมฺมสฺส อาหารปจฺจเยน ปจฺจโย. อินฺทฺริย\-ปจฺจเยน ปจฺจโย. ฌานปจฺจเยน ปจฺจโย. มคฺคปจฺจเยน ปจฺจโย. สมฺปยุตฺต\-ปจฺจเยน ปจฺจโย. อตฺถิปจฺจเยน ปจฺจโย. (อิเม สตฺต ปจฺจยา สตฺตรส ปญฺหา เหตุสทิสา.) นตฺถิ\-ปจฺจเยน ปจฺจโย. วิคต\-ปจฺจเยน ปจฺจโย. (อนนฺตรสทิสา.) อวิคต\-ปจฺจเยน ปจฺจโย. (สตฺตรส ปญฺหา.)}

\csromanmatikaanchor{mtk.170}
\csromantocmarkref{subparagraph}{1. ปจฺจยานุโลม 2. สงฺขฺยาวาร}{mtk.170}
\bookmark[dest={mtk.170},level=5]{1. ปจฺจยานุโลม 2. สงฺขฺยาวาร}
\csromanmarkh{1. ปจฺจยานุโลม}\csromanmarkhii{2. สงฺขฺยาวาร}\csromanheadernomark{6}{1. \textbf{ปจฺจยานุโลม 2. สงฺขฺยาวาร}}
\proseitemwithcloserruled{48}{เหตุยา สตฺตรส, อารมฺมเณ นว, อธิปติยา เอกวีส, อนนฺตเร นว, สมนนฺตเร นว, สหชาเต สตฺตรส, อญฺญมญฺเญ สตฺตรส, นิสฺสเย สตฺตรส, อุปนิสฺสเย เอกวีส, อาเสวเน นว, กมฺเม สตฺตรส, อาหาเร, อินฺทฺริเย, ฌาเน, มคฺเค, สมฺปยุตฺเต สตฺตรส, อตฺถิยา สตฺตรส, นตฺถิยา นว, วิคเต นว, อวิคเต สตฺตรส. (เอวํ คเณตพฺพํ.)}{อนุโลมํ.}
\csromanmatikaanchor{mtk.171}
\csromantocmarkref{subparagraph}{2. ปจฺจนียุทฺธาร}{mtk.171}
\bookmark[dest={mtk.171},level=5]{2. ปจฺจนียุทฺธาร}
\csromanheader{6}{2. ปจฺจนียุ\-ทฺธาร}
\proseitem{49}{มคฺคารมฺมโณ ธมฺโม มคฺคา\-รมฺมณสฺส ธมฺมสฺส สหชาต\-ปจฺจเยน ปจฺจโย. อุปนิสฺสย\-ปจฺจเยน ปจฺจโย. (1)}

\prose{มคฺคารมฺมโณ ธมฺโม มคฺคา\-ธิปติสฺส ธมฺมสฺส สหชาต\-ปจฺจเยน ปจฺจโย. อุปนิสฺสย\-ปจฺจเยน ปจฺจโย. (2)}

\prose{มคฺคารมฺมโณ ธมฺโม มคฺคา\-รมฺมณสฺส จ มคฺคา\-ธิปติสฺส จ ธมฺมสฺส สหชาต\-ปจฺจเยน ปจฺจโย. อุปนิสฺสย\-ปจฺจเยน ปจฺจโย. (3)}

\proseitem{50}{\csromanfolio{402}มคฺคเหตุโก ธมฺโม มคฺคเหตุกสฺส ธมฺมสฺส สหชาต\-ปจฺจเยน ปจฺจโย. อุปนิสฺสย\-ปจฺจเยน ปจฺจโย. (1)}

\prose{มคฺคเหตุโก ธมฺโม มคฺคา\-รมฺมณสฺส ธมฺมสฺส อารมฺมณ\-ปจฺจเยน ปจฺจโย. อุปนิสฺสย\-ปจฺจเยน ปจฺจโย. (2)}

\prose{มคฺคเหตุโก ธมฺโม มคฺคา\-ธิปติสฺส ธมฺมสฺส อารมฺมณ\-ปจฺจเยน ปจฺจโย. สหชาต\-ปจฺจเยน ปจฺจโย. อุปนิสฺสย\-ปจฺจเยน ปจฺจโย. (3)}

\prose{มคฺคเหตุโก ธมฺโม มคฺคา\-รมฺมณสฺส จ มคฺคา\-ธิปติสฺส จ ธมฺมสฺส อารมฺมณ\-ปจฺจเยน ปจฺจโย. อุปนิสฺสย\-ปจฺจเยน ปจฺจโย. (4)}

\prose{มคฺคเหตุโก ธมฺโม มคฺคเหตุกสฺส จ มคฺคา\-ธิปติสฺส จ ธมฺมสฺส สหชาต\-ปจฺจเยน ปจฺจโย. อุปนิสฺสย\-ปจฺจเยน ปจฺจโย. (5)}

\proseitem{51}{มคฺคาธิปติ ธมฺโม มคฺคา\-ธิปติสฺส ธมฺมสฺส สหชาต\-ปจฺจเยน ปจฺจโย. อุปนิสฺสย\-ปจฺจเยน ปจฺจโย. (1)}

\prose{มคฺคาธิปติ ธมฺโม มคฺคา\-รมฺมณสฺส ธมฺมสฺส อารมฺมณ\-ปจฺจเยน ปจฺจโย. สหชาต\-ปจฺจเยน ปจฺจโย. อุปนิสฺสย\-ปจฺจเยน ปจฺจโย. (2)}

\prose{มคฺคาธิปติ ธมฺโม มคฺคเหตุกสฺส ธมฺมสฺส สหชาต\-ปจฺจเยน ปจฺจโย. อุปนิสฺสย\-ปจฺจเยน ปจฺจโย. (3)}

\prose{มคฺคาธิปติ ธมฺโม มคฺคา\-รมฺมณสฺส จ มคฺคา\-ธิปติสฺส จ ธมฺมสฺส สหชาต\-ปจฺจเยน ปจฺจโย. อุปนิสฺสย\-ปจฺจเยน ปจฺจโย. (4)}

\prose{มคฺคาธิปติ ธมฺโม มคฺคเหตุกสฺส จ มคฺคา\-ธิปติสฺส จ ธมฺมสฺส สหชาต\-ปจฺจเยน ปจฺจโย. อุปนิสฺสย\-ปจฺจเยน ปจฺจโย. (5)}

\proseitem{52}{มคฺคารมฺมโณ จ มคฺคาธิปติ จ ธมฺมา มคฺคา\-รมฺมณสฺส ธมฺมสฺส สหชาต\-ปจฺจเยน ปจฺจโย. อุปนิสฺสย\-ปจฺจเยน ปจฺจโย. (1)}

\prose{มคฺคารมฺมโณ จ มคฺคาธิปติ จ ธมฺมา มคฺคา\-ธิปติสฺส ธมฺมสฺส สหชาต\-ปจฺจเยน ปจฺจโย. อุปนิสฺสย\-ปจฺจเยน ปจฺจโย. (2)}

\csromanmatikaanchor{mtk.172}
\csromanmatikaanchor{mtk.173}
\csromantocmarknopage{subparagraph}{2. ปจฺจยปจฺจนีย 2. สงฺขฺยาวาร}{mtk.172}
\bookmark[dest={mtk.172},level=5]{2. ปจฺจยปจฺจนีย 2. สงฺขฺยาวาร}
\csromantocmarkref{subparagraph}{สุทฺธ}{mtk.173}
\bookmark[dest={mtk.173},level=5]{สุทฺธ}
\prose{\csromanfolio{403}มคฺคารมฺมโณ จ มคฺคาธิปติ จ ธมฺมา มคฺคา\-รมฺมณสฺส จ มคฺคา\-ธิปติสฺส จ ธมฺมสฺส สหชาต\-ปจฺจเยน ปจฺจโย. อุปนิสฺสย\-ปจฺจเยน ปจฺจโย. (3)}

\proseitem{53}{มคฺคเหตุโก จ มคฺคาธิปติ จ ธมฺมา มคฺคา\-รมฺมณสฺส ธมฺมสฺส อารมฺมณ\-ปจฺจเยน ปจฺจโย. อุปนิสฺสย\-ปจฺจเยน ปจฺจโย. (1)}

\prose{มคฺคเหตุโก จ มคฺคาธิปติ จ ธมฺมา มคฺคเหตุกสฺส ธมฺมสฺส สหชาต\-ปจฺจเยน ปจฺจโย. อุปนิสฺสย\-ปจฺจเยน ปจฺจโย. (2)}

\prose{มคฺคเหตุโก จ มคฺคาธิปติ จ ธมฺมา มคฺคา\-ธิปติสฺส ธมฺมสฺส สหชาต\-ปจฺจเยน ปจฺจโย. อุปนิสฺสย\-ปจฺจเยน ปจฺจโย. (3)}

\prose{มคฺคเหตุโก จ มคฺคาธิปติ จ ธมฺมา มคฺคา\-รมฺมณสฺส จ มคฺคา\-ธิปติสฺส จ ธมฺมสฺส อุปนิสฺสย\-ปจฺจเยน ปจฺจโย. (4)}

\prose{มคฺคเหตุโก จ มคฺคาธิปติ จ ธมฺมา มคฺคเหตุกสฺส จ มคฺคา\-ธิปติสฺส จ ธมฺมสฺส สหชาต\-ปจฺจเยน ปจฺจโย. อุปนิสฺสย\-ปจฺจเยน ปจฺจโย. (5)}

\csromanmarkh{2. ปจฺจย\-ปจฺจนีย}\csromanmarkhii{2. สงฺขฺยาวาร}\csromanheadernomark{2}{2. ปจฺจย\-ปจฺจนีย 2. สงฺขฺยาวาร}
\proseitem{54}{นเหตุยา เอกวีส, นอารมฺมเณ สตฺตรส. (นอารมฺมเณ คหิเต ปกตารมฺมณมฺปิ อุปนิสฺสยารมฺมณมฺปิ เทฺวปิ ฉิชฺชนฺติ.) นอธิปติยา เอกวีส, นอนนฺตเร, นสมนนฺตเร, นสหชาเต, นอญฺญมญฺเญ, นนิสฺสเย, นอุปนิสฺสเย, นปุเรชาเต, นปจฺฉาชาเต, นอาเสวเน, นกมฺเม, นวิปาเก, นอาหาเร, นอินฺทฺริเย, นฌาเน, นมคฺเค, นสมฺปยุตฺเต, นวิปฺปยุตฺเต, โนอตฺถิยา, โนนตฺถิยา, โนวิคเต, โนอวิคเต เอกวีส. (เอวํ คเณตพฺพํ.)}

\vspace{\tipitakaparskip}
\csromancenter{ปจฺจนียํ.}
\csromanmatikaanchor{mtk.174}
\csromanmatikaanchor{mtk.175}
\csromanmatikaanchor{mtk.176}
\csromanmatikaanchor{mtk.177}
\csromantocmarknopage{subparagraph}{3. ปจฺจยานุโลมปจฺจนีย}{mtk.174}
\bookmark[dest={mtk.174},level=5]{3. ปจฺจยานุโลมปจฺจนีย}
\csromantocmarkref{subparagraph}{เหตุทุก}{mtk.175}
\bookmark[dest={mtk.175},level=5]{เหตุทุก}
\csromantocmarknopage{subparagraph}{4. ปจฺจยปจฺจนียานุโลม}{mtk.176}
\bookmark[dest={mtk.176},level=5]{4. ปจฺจยปจฺจนียานุโลม}
\csromantocmarkref{subparagraph}{นเหตุทุก}{mtk.177}
\bookmark[dest={mtk.177},level=5]{นเหตุทุก}
\proseitemwithcloserruled{55}{\csromanfolio{404}เหตุปจฺจยา นอารมฺมเณ สตฺตรส, นอธิปติยา, นอนนฺตเร, นสมนนฺตเร, นอุปนิสฺสเย, นปุเรชาเต, นปจฺฉาชาเต, นอาเสวเน, นกมฺเม, นวิปาเก, นอาหาเร, นอินฺทฺริเย, นฌาเน, นมคฺเค, นวิปฺปยุตฺเต, โนนตฺถิยา, โนวิคเต สตฺตรส. (เอวํ คเณตพฺพํ.)}{อนุโลม\-ปจฺจนียํ.}
\proseitem{56}{นเหตุปจฺจยา อารมฺมเณ นว, อธิปติยา เอกวีส, อนนฺตเร นว, สมนนฺตเร นว, สหชาเต สตฺตรส, อญฺญมญฺเญ สตฺตรส, นิสฺสเย สตฺตรส, อุปนิสฺสเย เอกวีส, อาเสวเน นว, กมฺเม สตฺตรส, อาหาเร สตฺตรส, อินฺทฺริเย, ฌาเน, มคฺเค, สมฺปยุตฺเต สตฺตรส, อตฺถิยา สตฺตรส, นตฺถิยา นว, วิคเต นว, อวิคเต สตฺตรส. (เอวํ คเณตพฺพํ.)}

\vspace{\tipitakaparskip}
\csromancenter{ปจฺจนียา\-นุโลมํ.}
\nitthitam{มคฺคา\-รมฺมณ\-ตฺติกํ นิฏฺฐิตํ.}
\csromanmatikaanchor{mtk.178}
\csromantocmarknopage{subparagraph}{17. อุปฺปนฺนตฺติก 7. ปญฺหาวาร}{mtk.178}
\bookmark[dest={mtk.178},level=5]{17. อุปฺปนฺนตฺติก 7. ปญฺหาวาร}
\csromanmarkh{17. อุปฺปนฺนตฺติก}\csromanmarkhii{17. อุปฺปนฺนตฺติก 7. ปญฺหาวาร}\csromanheadernomark{6}{17. อุปฺปนฺนตฺติก \textbf{17. อุปฺปนฺนตฺติก 7. ปญฺหาวาร}}
\csromanmarkh{1. ปจฺจยานุโลม}\csromanmarkhii{1. วิภงฺควาร}\csromanheadernomark{2}{1. \textbf{ปจฺจยานุโลม 1. วิภงฺควาร}}
\csromanmatikaanchor{mtk.179}
\csromanmatikaanchor{mtk.180}
\csromantocmarknopage{subparagraph}{1. ปจฺจยานุโลม 1. วิภงฺควาร}{mtk.179}
\bookmark[dest={mtk.179},level=5]{1. ปจฺจยานุโลม 1. วิภงฺควาร}
\csromantocmarkref{subparagraph}{เหตุ}{mtk.180}
\bookmark[dest={mtk.180},level=5]{เหตุ}
\csromanheader{6}{\textbf{เหตุ}}
\csromanmatikaanchor{mtk.181}
\csromantocmarkref{subparagraph}{อารมฺมณาทิ}{mtk.181}
\bookmark[dest={mtk.181},level=5]{อารมฺมณาทิ}
\proseitem{1}{\csromanfolio{405}อุปฺปนฺโน ธมฺโม อุปฺปนฺนสฺส ธมฺมสฺส เหตุปจฺจเยน ปจฺจโย. อุปฺปนฺนา เหตู สมฺปยุตฺตกานํ ขนฺธานํ จิตฺตสมุฏฺฐานานญฺจ รูปานํ เหตุปจฺจเยน ปจฺจโย. ปฏิสนฺธิ\-กฺขเณ อุปฺปนฺนา เหตู สมฺปยุตฺตกานํ ขนฺธานํ กฏตฺตา จ รูปานํ เหตุปจฺจเยน ปจฺจโย. (1)}

\csromanheader{2}{\textbf{อารมฺมณ}}
\proseitem{2}{อุปฺปนฺโน ธมฺโม อุปฺปนฺนสฺส ธมฺมสฺส อารมฺมณ\-ปจฺจเยน ปจฺจโย. อุปฺปนฺนํ จกฺขุํ อนิจฺจโต ทุกฺขโต อนตฺตโต วิปสฺสติ, อสฺสาเทติ อภินนฺทติ, ตํ อารพฺภ ราโค อุปฺปชฺชติ, ทิฏฺฐิ อุปฺปชฺชติ, วิจิกิจฺฉา ฯเปฯ อุทฺธจฺจํ ฯเปฯ โทมนสฺสํ อุปฺปชฺชติ. อุปฺปนฺนํ โสตํ. ฆานํ. ชิวฺหํ. กายํ. รูเป. สทฺเท. คนฺเธ. รเส. โผฏฺฐพฺเพ. วตฺถุํ. อุปฺปนฺเน ขนฺเธ อนิจฺจโต ทุกฺขโต อนตฺตโต วิปสฺสติ ฯเปฯ โทมนสฺสํ อุปฺปชฺชติ. ทิพฺเพน จกฺขุนา รูปํ ปสฺสติ. ทิพฺพาย โสตธาตุยา สทฺทํ สุณาติ. รูปายตนํ จกฺขุ\-วิญฺญาณสฺส ฯเปฯ. โผฏฺฐพฺพา\-ยตนํ กายวิญฺญาณสฺส ฯเปฯ. อุปฺปนฺนา ขนฺธา อิทฺธิวิธ\-ญาณสฺส, อาวชฺชนาย อารมฺมณ\-ปจฺจเยน ปจฺจโย. (1)}

\proseitem{3}{อนุปฺปนฺโน ธมฺโม อุปฺปนฺนสฺส ธมฺมสฺส อารมฺมณ\-ปจฺจเยน ปจฺจโย. อนุปฺปนฺเน รูเป. สทฺเท. คนฺเธ. รเส. โผฏฺฐพฺเพ. อนุปฺปนฺเน ขนฺเธ อนิจฺจโต ทุกฺขโต อนตฺตโต วิปสฺสติ ฯเปฯ โทมนสฺสํ อุปฺปชฺชติ. อนุปฺปนฺนา ขนฺธา อิทฺธิวิธ\-ญาณสฺส, เจโตปริย\-ญาณสฺส,อนาคตํสญาณสฺส, อาวชฺชนาย อารมฺมณ\-ปจฺจเยน ปจฺจโย. (1)}

\proseitem{4}{อุปฺปาที ธมฺโม อุปฺปนฺนสฺส ธมฺมสฺส อารมฺมณ\-ปจฺจเยน ปจฺจโย. อุปฺปาทิํ จกฺขุํ ฯเปฯ. กายํ. รูเป. คนฺเธ. รเส. โผฏฺฐพฺเพ. วตฺถุํ. อุปฺปาที ขนฺเธ อนิจฺจโต ทุกฺขโต อนตฺตโต ฯเปฯ โทมนสฺสํ อุปฺปชฺชติ. อุปฺปาที ขนฺธา อิทฺธิวิธ\-ญาณสฺส, เจโตปริย\-ญาณสฺส, อนาคตํสญาณสฺส, อาวชฺชนาย อารมฺมณ\-ปจฺจเยน ปจฺจโย. (1)}

\proseitem{5}{\csromanfolio{406}อุปฺปนฺโน ธมฺโม อุปฺปนฺนสฺส ธมฺมสฺส อธิปติ\-ปจฺจเยน ปจฺจโย. อารมฺมณา\-ธิปติ, สหชาตา\-ธิปติ.  อารมฺมณา\-ธิปติ\csromandash{} อุปฺปนฺนํ จกฺขุํ ครุํ กตฺวา อสฺสาเทติ อภินนฺทติ, ตํ ครุํ กตฺวา ราโค อุปฺปชฺชติ, ทิฏฺฐิ อุปฺปชฺชติ. อุปฺปนฺนํ โสตํ. ฆานํ. ชิวฺหํ. กายํ. รูเป. สทฺเท. คนฺเธ. รเส. โผฏฺฐพฺเพ. วตฺถุํ. อุปฺปนฺเน ขนฺเธ ครุํ กตฺวา อสฺสาเทติ อภินนฺทติ, ตํ ครุํ กตฺวา ราโค อุปฺปชฺชติ ฯเปฯ. สหชาตา\-ธิปติ\csromandash{}อุปฺปนฺนา อธิปติ สมฺปยุตฺตกานํ ขนฺธานํ จิตฺตสมุฏฺฐานานญฺจ รูปานํ อธิปติ\-ปจฺจเยน ปจฺจโย. (1)}

\prose{อนุปฺปนฺโน ธมฺโม อุปฺปนฺนสฺส ธมฺมสฺส อธิปติ\-ปจฺจเยน ปจฺจโย.  อารมฺมณา\-ธิปติ\csromandash{}อนุปฺปนฺเน รูเป. สทฺเท. คนฺเธ. รเส. โผฏฺฐพฺเพ. อนุปฺปนฺเน ขนฺเธ ครุํ กตฺวา อสฺสาเทติ อภินนฺทติ, ตํ ครุํ กตฺวา ราโค อุปฺปชฺชติ, ทิฏฺฐิ อุปฺปชฺชติ. (1)}

\prose{อุปฺปาที ธมฺโม อุปฺปนฺนสฺส ธมฺมสฺส อธิปติ\-ปจฺจเยน ปจฺจโย.  อารมฺมณา\-ธิปติ\csromandash{}อุปฺปาทิํ จกฺขุํ ฯเปฯ. กายํ. รูเป ฯเปฯ. โผฏฺฐพฺเพ. วตฺถุํ. อุปฺปาที ขนฺเธ ครุํ กตฺวา อสฺสาเทติ อภินนฺทติ, ตํ ครุํ กตฺวา ราโค อุปฺปชฺชติ, ทิฏฺฐิ อุปฺปชฺชติ. (1)}

\csromanheader{2}{\textbf{สหชาต}}
\proseitem{6}{อุปฺปนฺโน ธมฺโม อุปฺปนฺนสฺส ธมฺมสฺส สหชาต\-ปจฺจเยน ปจฺจโย. อุปฺปนฺโน เอโก ขนฺโธ ติณฺณนฺนํ ขนฺธานํ จิตฺตสมุฏฺฐานานญฺจ รูปานํ สหชาต\-ปจฺจเยน ปจฺจโย ฯเปฯ เทฺว ขนฺธา ทฺวินฺนํ ขนฺธานํ จิตฺตสมุฏฺฐานานญฺจ รูปานํ สหชาต\-ปจฺจเยน ปจฺจโย. ปฏิสนฺธิ\-กฺขเณ อุปฺปนฺโน เอโก ขนฺโธ ติณฺณนฺนํ ขนฺธานํ กฏตฺตา จ รูปานํ สหชาต\-ปจฺจเยน ปจฺจโย ฯเปฯ เทฺว ขนฺธา ทฺวินฺนํ ขนฺธานํ กฏตฺตา จ รูปานํ สหชาต\-ปจฺจเยน ปจฺจโย. ขนฺธา วตฺถุสฺส สหชาต\-ปจฺจเยน ปจฺจโย. วตฺถุ ขนฺธานํ สหชาต\-ปจฺจเยน ปจฺจโย. เอกํ มหาภูตํ ติณฺณนฺนํ มหาภูตานํ สหชาต\-ปจฺจเยน ปจฺจโย ฯเปฯ เทฺว มหาภูตา ฯเปฯ มหาภูตา จิตฺต\-สมุฏฺฐานานํ รูปานํ, กฏตฺตารูปานํ อุปาทารูปานํ สหชาต\-ปจฺจเยน ปจฺจโย. พาหิรํ. อาหารสมุฏฺฐานํ. อุตุสมุฏฺฐานํ. อสญฺญสตฺตานํ เอกํ มหาภูตํ ฯเปฯ เทฺว มหาภูตา ฯเปฯ มหาภูตา กฏตฺตารูปานํ, อุปาทารูปานํ สหชาต\-ปจฺจเยน ปจฺจโย. (1)}

\csromanheader{2}{17. อุปฺปนฺนตฺติก \textbf{อญฺญมญฺญ}}
\proseitem{7}{\csromanfolio{407}อุปฺปนฺโน ธมฺโม อุปฺปนฺนสฺส ธมฺมสฺส อญฺญมญฺญ\-ปจฺจเยน ปจฺจโย. อุปฺปนฺโน เอโก ขนฺโธ ติณฺณนฺนํ ขนฺธานํ อญฺญมญฺญ\-ปจฺจเยน ปจฺจโย ฯเปฯ เทฺว ขนฺธา ฯเปฯ. ปฏิสนฺธิ\-กฺขเณ อุปฺปนฺโน เอโก ขนฺโธ ติณฺณนฺนํ ขนฺธานํ วตฺถุสฺส จ อญฺญมญฺญ\-ปจฺจเยน ปจฺจโย. เทฺว ขนฺธา ฯเปฯ. ขนฺธา วตฺถุสฺส อญฺญมญฺญ\-ปจฺจเยน ปจฺจโย. วตฺถุ ขนฺธานํ อญฺญมญฺญ\-ปจฺจเยน ปจฺจโย. เอกํ มหาภูตํ ฯเปฯ. พาหิรํ. อาหารสมุฏฺฐานํ. อุตุสมุฏฺฐานํ. อสญฺญสตฺตานํ เอกํ มหาภูตํ ติณฺณนฺนํ มหาภูตานํ อญฺญมญฺญ\-ปจฺจเยน ปจฺจโย ฯเปฯ เทฺว มหาภูตา ฯเปฯ. (1)}

\csromanheader{2}{\textbf{นิสฺสย}}
\proseitem{8}{อุปฺปนฺโน ธมฺโม อุปฺปนฺนสฺส ธมฺมสฺส นิสฺสย\-ปจฺจเยน ปจฺจโย. อุปฺปนฺโน เอโก ขนฺโธ ติณฺณนฺนํ ขนฺธานํ จิตฺตสมุฏฺฐานานญฺจ รูปานํ นิสฺสย\-ปจฺจเยน ปจฺจโย ฯเปฯ เทฺว ขนฺธา ฯเปฯ. ปฏิสนฺธิ\-กฺขเณ ฯเปฯ. ขนฺธา วตฺถุสฺส ฯเปฯ. วตฺถุ ขนฺธานํ ฯเปฯ. เอกํ มหาภูตํ ฯเปฯ. พาหิรํ. อาหารสมุฏฺฐานํ. อุตุสมุฏฺฐานํ. อสญฺญสตฺตานํ เอกํ มหาภูตํ ฯเปฯ มหาภูตา จิตฺต\-สมุฏฺฐานานํ รูปานํ, กฏตฺตารูปานํ, อุปาทารูปานํ. จกฺขายตนํ จกฺขุ\-วิญฺญาณสฺส ฯเปฯ. กายายตนํ กายวิญฺญาณสฺส ฯเปฯ. วตฺถุ อุปฺปนฺนานํ ขนฺธานํ นิสฺสย\-ปจฺจเยน ปจฺจโย. (1)}

\csromanheader{2}{\textbf{อุปนิสฺสย}}
\proseitem{9}{อุปฺปนฺโน ธมฺโม อุปฺปนฺนสฺส ธมฺมสฺส อุปนิสฺสย\-ปจฺจเยน ปจฺจโย. อารมฺมณู\-ปนิสฺสโย, ปกตู\-ปนิสฺสโย ฯเปฯ. ปกตู\-ปนิสฺสโย\csromandash{} อุปฺปนฺนํ อุตุํ อุปนิสฺสาย ฌานํ อุปฺปาเทติ, วิปสฺสนํ ฯเปฯ มคฺคํ ฯเปฯ อภิญฺญํ ฯเปฯ สมาปตฺติํ อุปฺปาเทติ, มานํ ชปฺเปติ. ทิฏฺฐิํ คณฺหาติ. อุปฺปนฺนํ โภชนํ ฯเปฯ. เสนาสนํ อุปนิสฺสาย ฌานํ อุปฺปาเทติ วิปสฺสนํ ฯเปฯ มคฺคํ ฯเปฯ อภิญฺญํ ฯเปฯ สมาปตฺติํ อุปฺปาเทติ, มานํ ชปฺเปติ. ทิฏฺฐิํ คณฺหาติ. อุปฺปนฺนํ อุตุ. โภชนํ. เสนาสนํ อุปฺปนฺนาย สทฺธาย ฯเปฯ ปญฺญาย. กายิกสฺส สุขสฺส, กายิกสฺส ทุกฺขสฺส, มคฺคสฺส, ผลสมาปตฺติยา อุปนิสฺสย\-ปจฺจเยน ปจฺจโย. (1)}

\prose{อนุปฺปนฺโน ธมฺโม อุปฺปนฺนสฺส ธมฺมสฺส อุปนิสฺสย\-ปจฺจเยน ปจฺจโย. อารมฺมณู\-ปนิสฺสโย, ปกตู\-ปนิสฺสโย ฯเปฯ. ปกตู\-ปนิสฺสโย\csromandash{} \csromanfolio{408}อนุปฺปนฺนํ วณฺณสมฺปทํ ปตฺถยมาโน ทานํ เทติ, สีลํ สมาทิยติ, อุโปสถกมฺมํ กโรติ. อนุปฺปนฺนํ สทฺทสมฺปทํ. คนฺธสมฺปทํ. รสสมฺปทํ. โผฏฺฐพฺพ\-สมฺปทํ. อนุปฺปนฺเน ขนฺเธ ปตฺถยมาโน ทานํ เทติ, สีลํ สมาทิยติ, อุโปสถกมฺมํ กโรติ. อนุปฺปนฺนา วณฺณสมฺปทา ฯเปฯ. อนุปฺปนฺนา ขนฺธา อุปฺปนฺนาย สทฺธาย ฯเปฯ ปญฺญาย. กายิกสฺส สุขสฺส, กายิกสฺส ทุกฺขสฺส, มคฺคสฺส, ผลสมาปตฺติยา อุปนิสฺสย\-ปจฺจเยน ปจฺจโย. (1)}

\prose{อุปฺปาที ธมฺโม อุปฺปนฺนสฺส ธมฺมสฺส อุปนิสฺสย\-ปจฺจเยน ปจฺจโย. อารมฺมณู\-ปนิสฺสโย, ปกตู\-ปนิสฺสโย ฯเปฯ. ปกตู\-ปนิสฺสโย\csromandash{}อุปฺปาทิํ จกฺขุสมฺปทํ ปตฺถยมาโน ทานํ เทติ, สีลํ สมาทิยติ, อุโปสถกมฺมํ กโรติ. อุปฺปาทิํ โสตสมฺปทํ ฯเปฯ. กายสมฺปทํ ฯเปฯ. วณฺณสมฺปทํ. คนฺธสมฺปทํ. รสสมฺปทํ. โผฏฺฐพฺพ\-สมฺปทํ. อุปฺปาที ขนฺเธ ปตฺถยมาโน ทานํ เทติ, สีลํ สมาทิยติ, อุโปสถกมฺมํ กโรติ. อุปฺปาที จกฺขุสมฺปทา ฯเปฯ. กายสมฺปทา. วณฺณสมฺปทา ฯเปฯ. โผฏฺฐพฺพ\-สมฺปทา. อุปฺปาที ขนฺธา อุปฺปนฺนาย สทฺธาย ฯเปฯ. ปญฺญาย. กายิกสฺส สุขสฺส, กายิกสฺส ทุกฺขสฺส, มคฺคสฺส, ผลสมาปตฺติยา อุปนิสฺสย\-ปจฺจเยน ปจฺจโย. (1)}

\csromanheader{2}{\textbf{ปุเรชาต}}
\proseitem{9}{อุปฺปนฺโน ธมฺโม อุปฺปนฺนสฺส ธมฺมสฺส ปุเรชาต\-ปจฺจเยน ปจฺจโย. อารมฺมณ\-ปุเรชาตํ, วตฺถุ\-ปุเรชาตํ.  อารมฺมณ\-ปุเรชาตํ\csromandash{} จกฺขุํ ฯเปฯ. วตฺถุํ อนิจฺจโต ทุกฺขโต อนตฺตโต วิปสฺสติ, อสฺสาเทติ อภินนฺทติ, ตํ อารพฺภ ราโค อุปฺปชฺชติ ฯเปฯ โทมนสฺสํ อุปฺปชฺชติ. ทิพฺเพน จกฺขุนา รูปํ ปสฺสติ. ทิพฺพาย โสตธาตุยา สทฺทํ สุณาติ. รูปายตนํ จกฺขุ\-วิญฺญาณสฺส ฯเปฯ. โผฏฺฐพฺพา\-ยตนํ กายวิญฺญาณสฺส ปุเรชาต\-ปจฺจเยน ปจฺจโย.  วตฺถุ\-ปุเรชาตํ\csromandash{}จกฺขายตนํ จกฺขุ\-วิญฺญาณสฺส ฯเปฯ. วตฺถุ อุปฺปนฺนานํ ขนฺธานํ ปุเรชาต\-ปจฺจเยน ปจฺจโย. (1)}

\csromanheader{2}{\textbf{ปจฺฉาชาต}}
\proseitem{9}{อุปฺปนฺโน ธมฺโม อุปฺปนฺนสฺส ธมฺมสฺส ปจฺฉาชาต\-ปจฺจเยน ปจฺจโย. ปจฺฉาชาตา อุปฺปนฺนา ขนฺธา ปุเรชาตสฺส อิมสฺส กายสฺส ปจฺฉาชาต\-ปจฺจเยน ปจฺจโย. (1)}

\csromanheader{2}{17. อุปฺปนฺนตฺติก \textbf{กมฺม}}
\proseitem{12}{\csromanfolio{409}อุปฺปนฺโน ธมฺโม อุปฺปนฺนสฺส ธมฺมสฺส กมฺมปจฺจเยน ปจฺจโย. อุปฺปนฺนา เจตนา สมฺปยุตฺตกานํ ขนฺธานํ จิตฺตสมุฏฺฐานานญฺจ รูปานํ กมฺมปจฺจเยน ปจฺจโย. ปฏิสนฺธิ\-กฺขเณ อุปฺปนฺนา เจตนา สมฺปยุตฺตกานํ ขนฺธานํ กฏตฺตา จ รูปานํ กมฺมปจฺจเยน ปจฺจโย. (1)}

\csromanheader{2}{\textbf{วิปาก}}
\proseitem{13}{อุปฺปนฺโน ธมฺโม อุปฺปนฺนสฺส ธมฺมสฺส วิปาก\-ปจฺจเยน ปจฺจโย. วิปาโก อุปฺปนฺโน เอโก ขนฺโธ ติณฺณนฺนํ ขนฺธานํ จิตฺตสมุฏฺฐานานญฺจ รูปานํ วิปาก\-ปจฺจเยน ปจฺจโย ฯเปฯ เทฺว ขนฺธา ฯเปฯ. ปฏิสนฺธิ\-กฺขเณ อุปฺปนฺโน เอโก ขนฺโธ ติณฺณนฺนํ ขนฺธานํ กฏตฺตา จ รูปานํ ฯเปฯ เทฺว ขนฺธา ฯเปฯ. ขนฺธา วตฺถุสฺส วิปาก\-ปจฺจเยน ปจฺจโย. (1)}

\csromanheader{2}{\textbf{อาหาร}}
\proseitem{14}{อุปฺปนฺโน ธมฺโม อุปฺปนฺนสฺส ธมฺมสฺส อาหารปจฺจเยน ปจฺจโย. อุปฺปนฺนา อาหารา สมฺปยุตฺตกานํ ขนฺธานํ จิตฺตสมุฏฺฐานานญฺจ รูปานํ อาหารปจฺจเยน ปจฺจโย. ปฏิสนฺธิ\-กฺขเณ ฯเปฯ. กพฬีกาโร อาหาโร อิมสฺส กายสฺส อาหารปจฺจเยน ปจฺจโย. (1)}

\csromanheader{2}{\textbf{อินฺทฺริย}}
\proseitem{15}{อุปฺปนฺโน ธมฺโม อุปฺปนฺนสฺส ธมฺมสฺส อินฺทฺริย\-ปจฺจเยน ปจฺจโย. อุปฺปนฺนา อินฺทฺริยา สมฺปยุตฺตกานํ ขนฺธานํ จิตฺตสมุฏฺฐานานญฺจ รูปานํ อินฺทฺริย\-ปจฺจเยน ปจฺจโย. ปฏิสนฺธิ\-กฺขเณ ฯเปฯ. จกฺขุนฺทฺริยํ จกฺขุ\-วิญฺญาณสฺส ฯเปฯ. กายินฺทฺริยํ กายวิญฺญาณสฺส ฯเปฯ. รูปชีวิตินฺทฺริยํ กฏตฺตารูปานํ อินฺทฺริย\-ปจฺจเยน ปจฺจโย. (1)}

\csromanheader{2}{\textbf{ฌานาทิ}}
\proseitem{16}{อุปฺปนฺโน ธมฺโม อุปฺปนฺนสฺส ธมฺมสฺส ฌานปจฺจเยน ปจฺจโย. มคฺคปจฺจเยน ปจฺจโย. สมฺปยุตฺต\-ปจฺจเยน ปจฺจโย. วิปฺปยุตฺต\-ปจฺจเยน ปจฺจโย. สหชาตํ, ปุเรชาตํ, ปจฺฉาชาตํ.  สหชาตา อุปฺปนฺนา ขนฺธา จิตฺต\-สมุฏฺฐานานํ รูปานํ วิปฺปยุตฺต\-ปจฺจเยน ปจฺจโย. ปฏิสนฺธิ\-กฺขเณ อุปฺปนฺนา ขนฺธา กฏตฺตารูปานํ วิปฺปยุตฺต\-ปจฺจเยน ปจฺจโย. ขนฺธา วตฺถุสฺส, วตฺถุ ขนฺธานํ วิปฺปยุตฺต\-ปจฺจเยน ปจฺจโย.  ปุเรชาตํ จกฺขายตนํ \csromanfolio{410}จกฺขุ\-วิญฺญาณสฺส ฯเปฯ. กายายตนํ กายวิญฺญาณสฺส. วตฺถุ อุปฺปนฺนานํ ขนฺธานํ วิปฺปยุตฺต\-ปจฺจเยน ปจฺจโย.  ปจฺฉาชาตา อุปฺปนฺนา ขนฺธา ปุเรชาตสฺส อิมสฺส กายสฺส วิปฺปยุตฺต\-ปจฺจเยน ปจฺจโย. (1)}

\csromanheader{2}{\textbf{อตฺถิ}}
\proseitem{16}{อุปฺปนฺโน ธมฺโม อุปฺปนฺนสฺส ธมฺมสฺส อตฺถิปจฺจเยน ปจฺจโย. สหชาตํ, ปุเรชาตํ, ปจฺฉาชาตํ, อาหารํ, อินฺทฺริยํ.  สหชาโต อุปฺปนฺโน เอโก ขนฺโธ ติณฺณนฺนํ ขนฺธานํ จิตฺตสมุฏฺฐานานญฺจ รูปานํ อตฺถิปจฺจเยน ปจฺจโย ฯเปฯ เทฺว ขนฺธา ฯเปฯ. ปฏิสนฺธิ\-กฺขเณ ฯเปฯ. เอกํ มหาภูตํ ฯเปฯ. พาหิรํ. อาหารสมุฏฺฐานํ. อุตุสมุฏฺฐานํ. อสญฺญสตฺตานํ ฯเปฯ. ปุเรชาตํ จกฺขุํ อนิจฺจโต ฯเปฯ. วตฺถุํ อนิจฺจโต ฯเปฯ โทมนสฺสํ อุปฺปชฺชติ. ทิพฺเพน จกฺขุนา รูปํ ปสฺสติ. ทิพฺพาย โสตธาตุยา สทฺทํ สุณาติ. รูปายตนํ จกฺขุ\-วิญฺญาณสฺส ฯเปฯ. โผฏฺฐพฺพา\-ยตนํ กายวิญฺญาณสฺส อตฺถิปจฺจเยน ปจฺจโย. จกฺขายตนํ จกฺขุ\-วิญฺญาณสฺส ฯเปฯ. กายายตนํ กายวิญฺญาณสฺส ฯเปฯ. วตฺถุ อุปฺปนฺนานํ ขนฺธานํ อตฺถิปจฺจเยน ปจฺจโย.  ปจฺฉาชาตา อุปฺปนฺนา ขนฺธา ปุเรชาตสฺส อิมสฺส กายสฺส อตฺถิปจฺจเยน ปจฺจโย.  กพฬีกาโร อาหาโร อิมสฺส กายสฺส อตฺถิปจฺจเยน ปจฺจโย. รูปชีวิตินฺทฺริยํ กฏตฺตารูปานํ อตฺถิปจฺจเยน ปจฺจโย. (1)}

\csromanheader{2}{\textbf{อวิคต}}
\vspace{\tipitakaparskip}
\csromancenter{อุปฺปนฺโน ธมฺโม อุปฺปนฺนสฺส ธมฺมสฺส อวิคต\-ปจฺจเยน ปจฺจโย ฯเปฯ.}
\csromanmatikaanchor{mtk.182}
\csromantocmarkref{subparagraph}{1. ปจฺจยานุโลม 2. สงฺขฺยาวาร}{mtk.182}
\bookmark[dest={mtk.182},level=5]{1. ปจฺจยานุโลม 2. สงฺขฺยาวาร}
\csromanmarkh{1. ปจฺจยานุโลม}\csromanmarkhii{2. สงฺขฺยาวาร}\csromanheadernomark{6}{1. \textbf{ปจฺจยานุโลม 2. สงฺขฺยาวาร}}
\proseitem{16}{เหตุยา เอกํ, อารมฺมเณ ตีณิ, อธิปติยา ตีณิ, สหชาเต เอกํ, อญฺญมญฺเญ เอกํ, นิสฺสเย เอกํ, อุปนิสฺสเย ตีณิ, ปุเรชาเต เอกํ, ปจฺฉาชาเต, กมฺเม, วิปาเก, อาหาเร, อินฺทฺริเย, ฌาเน, มคฺเค, สมฺปยุตฺเต, วิปฺปยุตฺเต, อตฺถิยา, อวิคเต เอกํ. (เอวํ คเณตพฺพํ.)}

\vspace{\tipitakaparskip}
\csromancenter{อนุโลมํ.}
\csromanmatikaanchor{mtk.183}
\csromantocmarkref{subparagraph}{2. ปจฺจนียุทฺธาร}{mtk.183}
\bookmark[dest={mtk.183},level=5]{2. ปจฺจนียุทฺธาร}
\csromanmarkh{17. อุปฺปนฺนตฺติก}\csromanmarkhii{2. ปจฺจนียุ\-ทฺธาร}\csromanheadernomark{6}{17. อุปฺปนฺนตฺติก 2. ปจฺจนียุ\-ทฺธาร}
\csromanmatikaanchor{mtk.184}
\csromanmatikaanchor{mtk.185}
\csromantocmarkref{subparagraph}{2. ปจฺจยปจฺจนีย 2. สงฺขฺยาวาร}{mtk.184}
\bookmark[dest={mtk.184},level=5]{2. ปจฺจยปจฺจนีย 2. สงฺขฺยาวาร}
\csromantocmarkref{subparagraph}{3. ปจฺจยานุโลมปจฺจนีย}{mtk.185}
\bookmark[dest={mtk.185},level=5]{3. ปจฺจยานุโลมปจฺจนีย}
\proseitem{20}{\csromanfolio{411}อุปฺปนฺโน ธมฺโม อุปฺปนฺนสฺส ธมฺมสฺส อารมฺมณ\-ปจฺจเยน ปจฺจโย. สหชาต\-ปจฺจเยน ปจฺจโย. อุปนิสฺสย\-ปจฺจเยน ปจฺจโย. ปุเรชาต\-ปจฺจเยน ปจฺจโย. ปจฺฉาชาต\-ปจฺจเยน ปจฺจโย. อาหารปจฺจเยน ปจฺจโย. อินฺทฺริย\-ปจฺจเยน ปจฺจโย. (1)}

\prose{อนุปฺปนฺโน ธมฺโม อุปฺปนฺนสฺส ธมฺมสฺส อารมฺมณ\-ปจฺจเยน ปจฺจโย. อุปนิสฺสย\-ปจฺจเยน ปจฺจโย. (1)}

\prose{อุปฺปาที ธมฺโม อุปฺปนฺนสฺส ธมฺมสฺส อารมฺมณ\-ปจฺจเยน ปจฺจโย. อุปนิสฺสย\-ปจฺจเยน ปจฺจโย. (1)}

\csromanmarkh{2. ปจฺจย\-ปจฺจนีย}\csromanmarkhii{2. สงฺขฺยาวาร}\csromanheadernomark{6}{2. ปจฺจย\-ปจฺจนีย 2. สงฺขฺยาวาร}
\proseitemwithcloserruled{21}{นเหตุยา ตีณิ, นอารมฺมเณ ตีณิ, นอธิปติยา ตีณิ ฯเปฯ นวิปฺปยุตฺเต ตีณิ, โนอตฺถิยา เทฺว, โนนตฺถิยา ตีณิ, โนวิคเต ตีณิ, โนอวิคเต เทฺว. (เอวํ คเณตพฺพํ.)}{ปจฺจนียํ.}
\csromancenter{เหตุปจฺจยา นอารมฺมเณ เอกํ ฯเปฯ โนนตฺถิยา, โนวิคเต เอกํ.}
\nitthitamruled{อนุโลม\-ปจฺจนียํ.}
\csromanmatikaanchor{mtk.186}
\csromantocmarkref{subparagraph}{4. ปจฺจยปจฺจนียานุโลม}{mtk.186}
\bookmark[dest={mtk.186},level=5]{4. ปจฺจยปจฺจนียานุโลม}
\csromanheader{6}{4. ปจฺจนียา\-นุโลม}
\proseitem{23}{นเหตุปจฺจยา อารมฺมเณ ตีณิ, อธิปติยา ตีณิ, สหชาเต เอกํ, อญฺญมญฺเญ เอกํ, นิสฺสเย เอกํ, อุปนิสฺสเย ตีณิ, ปุเรชาเต เอกํ, ปจฺฉาชาเต เอกํ, กมฺเม, วิปาเก, อาหาเร, อินฺทฺริเย, ฌาเน, มคฺเค, สมฺปยุตฺเต, วิปฺปยุตฺเต, อตฺถิยา, อวิคเต เอกํ.}

\vspace{\tipitakaparskip}
\csromancenter{ปจฺจนียา\-นุโลมํ.}
\nitthitam{อุปฺปนฺนตฺติกํ นิฏฺฐิตํ.}
\csromanmatikaanchor{mtk.187}
\csromantocmarknopage{subparagraph}{18. อตีตตฺติก 7. ปญฺหาวาร}{mtk.187}
\bookmark[dest={mtk.187},level=5]{18. อตีตตฺติก 7. ปญฺหาวาร}
\csromanmarkh{18. อตีตตฺติก}\csromanmarkhii{7. ปญฺหาวาร}\csromanheadernomark{6}{18. \textbf{อตีตตฺติก 7. ปญฺหาวาร}}
\csromanmarkh{1. ปจฺจยานุโลม}\csromanmarkhii{1. วิภงฺควาร}\csromanheadernomark{2}{1. \textbf{ปจฺจยานุโลม 1. วิภงฺควาร}}
\csromanmatikaanchor{mtk.188}
\csromanmatikaanchor{mtk.189}
\csromantocmarknopage{subparagraph}{1. ปจฺจยานุโลม 1. วิภงฺควาร}{mtk.188}
\bookmark[dest={mtk.188},level=5]{1. ปจฺจยานุโลม 1. วิภงฺควาร}
\csromantocmarkref{subparagraph}{เหตุ}{mtk.189}
\bookmark[dest={mtk.189},level=5]{เหตุ}
\csromanheader{6}{\textbf{เหตุ}}
\csromanmatikaanchor{mtk.190}
\csromantocmarkref{subparagraph}{อารมฺมณาทิ}{mtk.190}
\bookmark[dest={mtk.190},level=5]{อารมฺมณาทิ}
\proseitem{1}{\csromanfolio{412}ปจฺจุปฺปนฺโน ธมฺโม ปจฺจุปฺปนฺนสฺส ธมฺมสฺส เหตุปจฺจเยน ปจฺจโย. ปจฺจุปฺปนฺนา เหตู สมฺปยุตฺตกานํ ขนฺธานํ จิตฺตสมุฏฺฐานานญฺจ รูปานํ เหตุปจฺจเยน ปจฺจโย. ปฏิสนฺธิ\-กฺขเณ ฯเปฯ. (1)}

\csromanheader{2}{\textbf{อารมฺมณ}}
\proseitem{2}{อตีโต ธมฺโม ปจฺจุปฺปนฺนสฺส ธมฺมสฺส อารมฺมณ\-ปจฺจเยน ปจฺจโย. ทานํ ทตฺวา, สีลํ สมาทิยิตฺวา, อุโปสถกมฺมํ ฯเปฯ. ปจฺจเวกฺขติ, ปุพฺเพ สุจิณฺณานิ ปจฺจเวกฺขติ, ฌานา วุฏฺฐหิตฺวา ฌานํ ปจฺจเวกฺขติ. อริยา มคฺคา วุฏฺฐหิตฺวา มคฺคํ ปจฺจเวกฺขนฺติ, ผลํ ปจฺจเวกฺขนฺติ, ปหีเน กิเลเส ปจฺจเวกฺขนฺติ, วิกฺขมฺภิเต กิเลเส ปจฺจเวกฺขนฺติ, ปุพฺเพ สมุทาจิณฺเณ กิเลเส ชานนฺติ. อตีตํ จกฺขุํ อนิจฺจโต ทุกฺขโต อนตฺตโต วิปสฺสติ ฯเปฯ โทมนสฺสํ อุปฺปชฺชติ. อตีตํ โสตํ ฯเปฯ. ฆานํ. ชิวฺหํ. กายํ. รูเป. สทฺเท. คนฺเธ. รเส. โผฏฺฐพฺเพ. วตฺถุํ. อตีเต ขนฺเธ อนิจฺจโต ทุกฺขโต อนตฺตโต วิปสฺสติ, อสฺสาเทติ อภินนฺทติ, ตํ อารพฺภ ราโค อุปฺปชฺชติ, ทิฏฺฐิ อุปฺปชฺชติ, วิจิกิจฺฉา. อุทฺธจฺจํ. โทมนสฺสํ อุปฺปชฺชติ. อากาสา\-นญฺจา\-ยตนํ วิญฺญาณญฺจา\-ยตนสฺส อารมฺมณ\-ปจฺจเยน ปจฺจโย. อากิญฺจญฺญา\-ยตนํ เนว\-สญฺญา\-นา\-สญฺญา\-ยตนสฺส อารมฺมณ\-ปจฺจเยน ปจฺจโย. อตีตา ขนฺธา อิทฺธิวิธ\-ญาณสฺส, เจโตปริย\-ญาณสฺส, ปุพฺเพ\-นิวาสา\-นุสฺสติ\-ญาณสฺส, ยถากมฺมูปคญาณสฺส, อาวชฺชนาย อารมฺมณ\-ปจฺจเยน ปจฺจโย. (1)}

\proseitem{3}{อนาคโต ธมฺโม ปจฺจุปฺปนฺนสฺส ธมฺมสฺส อารมฺมณ\-ปจฺจเยน ปจฺจโย. อนาคตํ จกฺขุํ ฯเปฯ. วตฺถุํ. อนาคเต ขนฺเธ อนิจฺจโต ฯเปฯ โทมนสฺสํ อุปฺปชฺชติ. อนาคตา ขนฺธา อิทฺธิวิธ\-ญาณสฺส, เจโตปริย\-ญาณสฺส, อนาคตํสญาณสฺส, อาวชฺชนาย อารมฺมณ\-ปจฺจเยน ปจฺจโย. (1)}

\prose{ปจฺจุปฺปนฺโน ธมฺโม ปจฺจุปฺปนฺนสฺส ธมฺมสฺส อารมฺมณ\-ปจฺจเยน ปจฺจโย. ปจฺจุปฺปนฺนํ จกฺขุํ ฯเปฯ. กายํ. รูเป. สทฺเท. คนฺเธ. รเส. โผฏฺฐพฺเพ. วตฺถุํ. \csromanfolio{413}ปจฺจุปฺปนฺเน ขนฺเธ อนิจฺจโต ฯเปฯ โทมนสฺสํ อุปฺปชฺชติ. ทิพฺเพน จกฺขุนา รูปํ ปสฺสติ. ทิพฺพาย โสตธาตุยา สทฺทํ สุณาติ. รูปายตนํ จกฺขุ\-วิญฺญาณสฺส ฯเปฯ. โผฏฺฐพฺพา\-ยตนํ กายวิญฺญาณสฺส ฯเปฯ. ปจฺจุปฺปนฺนา ขนฺธา อิทฺธิวิธ\-ญาณสฺส, อาวชฺชนาย อารมฺมณ\-ปจฺจเยน ปจฺจโย. (1)}

\proseitem{4}{อตีโต ธมฺโม ปจฺจุปฺปนฺนสฺส ธมฺมสฺส อธิปติ\-ปจฺจเยน ปจฺจโย. อารมฺมณา\-ธิปติ\csromandash{}ทานํ ทตฺวา, สีลํ สมาทิยิตฺวา ฯเปฯ. ปุพฺเพ สุจิณฺณานิ ครุํ กตฺวา ปจฺจเวกฺขติ, ฌานา วุฏฺฐหิตฺวา ฌานํ ครุํ กตฺวา ปจฺจเวกฺขติ. อริยา มคฺคา วุฏฺฐหิตฺวา มคฺคํ ครุํ กตฺวา ปจฺจเวกฺขนฺติ, ผลํ ครุํ กตฺวา ปจฺจเวกฺขนฺติ. อตีตํ จกฺขุํ ฯเปฯ กายํ. รูเป. สทฺเท. คนฺเธ. รเส. โผฏฺฐพฺเพ. วตฺถุํ. อตีเต ขนฺเธ ครุํ กตฺวา อสฺสาเทติ อภินนฺทติ, ตํ ครุํ กตฺวา ราโค อุปฺปชฺชติ, ทิฏฺฐิ อุปฺปชฺชติ. (1)}

\prose{อนาคโต ธมฺโม ปจฺจุปฺปนฺนสฺส ธมฺมสฺส อธิปติ\-ปจฺจเยน ปจฺจโย. อารมฺมณา\-ธิปติ\csromandash{}อนาคตํ จกฺขุํ ฯเปฯ วตฺถุํ. อนาคเต ขนฺเธ ครุํ กตฺวา อสฺสาเทติ อภินนฺทติ, ตํ ครุํ กตฺวา ราโค อุปฺปชฺชติ, ทิฏฺฐิ อุปฺปชฺชติ. (1)}

\prose{ปจฺจุปฺปนฺโน ธมฺโม ปจฺจุปฺปนฺนสฺส ธมฺมสฺส อธิปติ\-ปจฺจเยน ปจฺจโย. อารมฺมณา\-ธิปติ, สหชาตา\-ธิปติ.  อารมฺมณา\-ธิปติ\csromandash{} ปจฺจุปฺปนฺนํ จกฺขุํ ฯเปฯ วตฺถุํ. ปจฺจุปฺปนฺเน ขนฺเธ ครุํ กตฺวา อสฺสาเทติ อภินนฺทติ, ตํ ครุํ กตฺวา ราโค อุปฺปชฺชติ, ทิฏฺฐิ อุปฺปชฺชติ.  สหชาตา\-ธิปติ\csromandash{}ปจฺจุปฺปนฺนา\-ธิปติ สมฺปยุตฺตกานํ ขนฺธานํ จิตฺตสมุฏฺฐานานญฺจ รูปานํ อธิปติ\-ปจฺจเยน ปจฺจโย. (1)}

\csromanheader{2}{\textbf{อนนฺตร}}
\proseitem{5}{อตีโต ธมฺโม ปจฺจุปฺปนฺนสฺส ธมฺมสฺส อนนฺตร\-ปจฺจเยน ปจฺจโย. ปุริมา ปุริมา อตีตา ขนฺธา ปจฺฉิมานํ ปจฺฉิมานํ ปจฺจุปฺปนฺนานํ ขนฺธานํ อนนฺตร\-ปจฺจเยน ปจฺจโย. อนุโลมํ โคตฺรภุสฺส. อนุโลมํ โวทานสฺส. โคตฺรภุ มคฺคสฺส. โวทานํ มคฺคสฺส. มคฺโค ผลสฺส. ผลํ ผลสฺส. อนุโลมํ ผลสมาปตฺติยา. นิโรธา วุฏฺฐหนฺตสฺส เนว\-สญฺญา\-นา\-สญฺญา\-ยตนํ ผลสมาปตฺติยา อนนฺตร\-ปจฺจเยน ปจฺจโย. (1)}

\csromanheader{2}{\textbf{สมนนฺตร}}
\proseitem{6}{\csromanfolio{414}อตีโต ธมฺโม ปจฺจุปฺปนฺนสฺส ธมฺมสฺส สมนนฺตร\-ปจฺจเยน ปจฺจโย. (อนนฺตร\-สทิสํ.) (1)}

\csromanheader{2}{\textbf{สหชาตาทิ}}
\proseitem{7}{ปจฺจุปฺปนฺโน ธมฺโม ปจฺจุปฺปนฺนสฺส ธมฺมสฺส สหชาต\-ปจฺจเยน ปจฺจโย. อญฺญมญฺญ\-ปจฺจเยน ปจฺจโย. นิสฺสย\-ปจฺจเยน ปจฺจโย. (สํขิตฺตํ.) (1)}

\csromanheader{2}{\textbf{อุปนิสฺสย}}
\proseitem{8}{อตีโต ธมฺโม ปจฺจุปฺปนฺนสฺส ธมฺมสฺส อุปนิสฺสย\-ปจฺจเยน ปจฺจโย. อารมฺมณู\-ปนิสฺสโย, อนนฺตรู\-ปนิสฺสโย, ปกตู\-ปนิสฺสโย ฯเปฯ. ปกตู\-ปนิสฺสโย\csromandash{}อตีตํ สทฺธํ อุปนิสฺสาย ทานํ เทติ, สีลํ สมาทิยติ, อุโปสถกมฺมํ กโรติ. ฌานํ อุปฺปาเทติ. วิปสฺสนํ. มคฺคํ. อภิญฺญํ. สมาปตฺติํ อุปฺปาเทติ. มานํ ชปฺเปติ. ทิฏฺฐิํ คณฺหาติ. อตีตํ สีลํ ฯเปฯ ปญฺญํ. ราคํ ฯเปฯ ปตฺถนํ. กายิกํ สุขํ. กายิกํ ทุกฺขํ อุปนิสฺสาย ทานํ เทติ, สีลํ สมาทิยติ, อุโปสถกมฺมํ ฯเปฯ สมาปตฺติํ อุปฺปาเทติ. ปาณํ หนติ ฯเปฯ สํฆํ ภินฺทติ. อตีตา สทฺธา ฯเปฯ ปญฺญา. ราโค ฯเปฯ ปตฺถนา. กายิกํ สุขํ. กายิกํ ทุกฺขํ. ปจฺจุปฺปนฺนาย สทฺธาย ฯเปฯ ปญฺญาย. ราคสฺส ฯเปฯ ปตฺถนาย ฯเปฯ ผลสมาปตฺติยา อุปนิสฺสย\-ปจฺจเยน ปจฺจโย. (1)}

\prose{อนาคโต ธมฺโม ปจฺจุปฺปนฺนสฺส ธมฺมสฺส อุปนิสฺสย\-ปจฺจเยน ปจฺจโย. อารมฺมณู\-ปนิสฺสโย, ปกตู\-ปนิสฺสโย ฯเปฯ. ปกตู\-ปนิสฺสโย\csromandash{} อนาคตํ จกฺขุสมฺปทํ ปตฺถยมาโน ฯเปฯ. โสตสมฺปทํ. ฆานสมฺปทํ. ชิวฺหาสมฺปทํ. กายสมฺปทํ. วณฺณสมฺปทํ. สทฺทสมฺปทํ. คนฺธสมฺปทํ. รสสมฺปทํ. โผฏฺฐพฺพ\-สมฺปทํ ปตฺถยมาโน ฯเปฯ. อนาคเต ขนฺเธ ปตฺถยมาโน ทานํ เทติ, สีลํ สมาทิยติ, อุโปสถกมฺมํ. อนาคตา จกฺขุสมฺปทา ฯเปฯ. วณฺณสมฺปทา ฯเปฯ. โผฏฺฐพฺพ\-สมฺปทา. อนาคตา ขนฺธา ปจฺจุปฺปนฺนาย สทฺธาย ฯเปฯ ปญฺญาย. กายิกสฺส สุขสฺส, กายิกสฺส ทุกฺขสฺส, มคฺคสฺส, ผลสมาปตฺติยา อุปนิสฺสย\-ปจฺจเยน ปจฺจโย. (1)}

\prose{ปจฺจุปฺปนฺโน ธมฺโม ปจฺจุปฺปนฺนสฺส ธมฺมสฺส อุปนิสฺสย\-ปจฺจเยน ปจฺจโย. อารมฺมณู\-ปนิสฺสโย, ปกตู\-ปนิสฺสโย ฯเปฯ. ปกตู\-ปนิสฺสโย\csromandash{} \csromanfolio{415}ปจฺจุปฺปนฺนํ อุตุํ อุปนิสฺสาย ฌานํ อุปฺปาเทติ. วิปสฺสนํ ฯเปฯ. ปจฺจุปฺปนฺนํ โภชนํ. เสนาสนํ อุปนิสฺสาย ฌานํ อุปฺปาเทติ ฯเปฯ สมาปตฺติํ อุปฺปาเทติ. ปจฺจุปฺปนฺนํ อุตุ. โภชนํ. เสนาสนํ ปจฺจุปฺปนฺนาย สทฺธาย ฯเปฯ ปญฺญาย. กายิกสฺส ฯเปฯ ผลสมาปตฺติยา อุปนิสฺสย\-ปจฺจเยน ปจฺจโย. (1)}

\csromanheader{2}{\textbf{ปุเรชาต}}
\proseitem{8}{ปจฺจุปฺปนฺโน ธมฺโม ปจฺจุปฺปนฺนสฺส ธมฺมสฺส ปุเรชาต\-ปจฺจเยน ปจฺจโย. อารมฺมณ\-ปุเรชาตํ, วตฺถุ\-ปุเรชาตํ.  อารมฺมณ\-ปุเรชาตํ\csromandash{}จกฺขุํ ฯเปฯ. วตฺถุํ อนิจฺจโต ฯเปฯ โทมนสฺสํ อุปฺปชฺชติ. ทิพฺเพน จกฺขุนา รูปํ ปสฺสติ. ทิพฺพาย โสตธาตุยา สทฺทํ สุณาติ. รูปายตนํ จกฺขุ\-วิญฺญาณสฺส ฯเปฯ. โผฏฺฐพฺพา\-ยตนํ กายวิญฺญาณสฺส ปุเรชาต\-ปจฺจเยน ปจฺจโย.  วตฺถุ\-ปุเรชาตํ\csromandash{} จกฺขายตนํ จกฺขุ\-วิญฺญาณสฺส ฯเปฯ. กายายตนํ กายวิญฺญาณสฺส. วตฺถุ ปจฺจุปฺปนฺนานํ ขนฺธานํ ปุเรชาต\-ปจฺจเยน ปจฺจโย. (1)}

\csromanheader{2}{\textbf{ปจฺฉาชาต}}
\proseitem{8}{ปจฺจุปฺปนฺโน ธมฺโม ปจฺจุปฺปนฺนสฺส ธมฺมสฺส ปจฺฉาชาต\-ปจฺจเยน ปจฺจโย. ปจฺฉาชาตา ปจฺจุปฺปนฺนา ขนฺธา ปุเรชาตสฺส อิมสฺส กายสฺส ปจฺฉาชาต\-ปจฺจเยน ปจฺจโย. (1)}

\csromanheader{2}{\textbf{อาเสวน}}
\proseitem{8}{อตีโต ธมฺโม ปจฺจุปฺปนฺนสฺส ธมฺมสฺส อาเสวน\-ปจฺจเยน ปจฺจโย. ปุริมา ปุริมา อตีตา ขนฺธา ปจฺฉิมานํ ปจฺฉิมานํ ปจฺจุปฺปนฺนานํ ขนฺธานํ อาเสวน\-ปจฺจเยน ปจฺจโย. อนุโลมํ โคตฺรภุสฺส. อนุโลมํ โวทานสฺส. โคตฺรภุ มคฺคสฺส. โวทานํ มคฺคสฺส อาเสวน\-ปจฺจเยน ปจฺจโย. (1)}

\csromanheader{2}{\textbf{กมฺม}}
\proseitem{8}{อตีโต ธมฺโม ปจฺจุปฺปนฺนสฺส ธมฺมสฺส กมฺมปจฺจเยน ปจฺจโย. นานากฺขณิกา อตีตา เจตนา ปจฺจุปฺปนฺนานํ วิปากานํ ขนฺธานํ กฏตฺตา จ รูปานํ กมฺมปจฺจเยน ปจฺจโย. (1)}

\proseitem{12}{\csromanfolio{416}ปจฺจุปฺปนฺโน ธมฺโม ปจฺจุปฺปนฺนสฺส ธมฺมสฺส กมฺมปจฺจเยน ปจฺจโย. ปจฺจุปฺปนฺนา เจตนา สมฺปยุตฺตกานํ ขนฺธานํ จิตฺตสมุฏฺฐานานญฺจ รูปานํ กมฺมปจฺจเยน ปจฺจโย. ปฏิสนฺธิ\-กฺขเณ ปจฺจุปฺปนฺนา เจตนา สมฺปยุตฺตกานํ ขนฺธานํ กฏตฺตา จ รูปานํ กมฺมปจฺจเยน ปจฺจโย. (1)}

\csromanheader{2}{\textbf{วิปาก}}
\proseitem{13}{ปจฺจุปฺปนฺโน ธมฺโม ปจฺจุปฺปนฺนสฺส ธมฺมสฺส วิปาก\-ปจฺจเยน ปจฺจโย. วิปาโก ปจฺจุปฺปนฺโน เอโก ขนฺโธ ติณฺณนฺนํ ขนฺธานํ จิตฺตสมุฏฺฐานานญฺจ รูปานํ วิปาก\-ปจฺจเยน ปจฺจโย ฯเปฯ เทฺว ขนฺธา ฯเปฯ. ปฏิสนฺธิ\-กฺขเณ ฯเปฯ. ขนฺธา วตฺถุสฺส วิปาก\-ปจฺจเยน ปจฺจโย. (1)}

\csromanheader{2}{\textbf{อาหาราทิ}}
\proseitem{14}{ปจฺจุปฺปนฺโน ธมฺโม ปจฺจุปฺปนฺนสฺส ธมฺมสฺส อาหารปจฺจเยน ปจฺจโย. อินฺทฺริย\-ปจฺจเยน ปจฺจโย. ฌานปจฺจเยน ปจฺจโย. มคฺคปจฺจเยน ปจฺจโย. สมฺปยุตฺต\-ปจฺจเยน ปจฺจโย. วิปฺปยุตฺต\-ปจฺจเยน ปจฺจโย. สหชาตํ, ปุเรชาตํ, ปจฺฉาชาตํ.  สหชาตา ปจฺจุปฺปนฺนา ขนฺธา จิตฺต\-สมุฏฺฐานานํ รูปานํ วิปฺปยุตฺต\-ปจฺจเยน ปจฺจโย. ปฏิสนฺธิ\-กฺขเณ ปจฺจุปฺปนฺนา ขนฺธา กฏตฺตารูปานํ วิปฺปยุตฺต\-ปจฺจเยน ปจฺจโย. ขนฺธา วตฺถุสฺส วิปฺปยุตฺต\-ปจฺจเยน ปจฺจโย. วตฺถุ ขนฺธานํ วิปฺปยุตฺต\-ปจฺจเยน ปจฺจโย.  ปุเรชาตํ จกฺขายตนํ จกฺขุ\-วิญฺญาณสฺส ฯเปฯ. กายายตนํ กายวิญฺญาณสฺส. วตฺถุ ปจฺจุปฺปนฺนานํ ขนฺธานํ วิปฺปยุตฺต\-ปจฺจเยน ปจฺจโย.  ปจฺฉาชาตา ปจฺจุปฺปนฺนา ขนฺธา ปุเรชาตสฺส อิมสฺส กายสฺส วิปฺปยุตฺต\-ปจฺจเยน ปจฺจโย. (1)}

\csromanheader{2}{\textbf{อตฺถิ}}
\proseitem{15}{ปจฺจุปฺปนฺโน ธมฺโม ปจฺจุปฺปนฺนสฺส ธมฺมสฺส อตฺถิปจฺจเยน ปจฺจโย. (อุปฺปนฺนตฺติเก อตฺถิสทิสํ.) (1)}

\csromanheader{2}{\textbf{นตฺถิ วิคตาวิคต}}
\proseitem{16}{อตีโต ธมฺโม ปจฺจุปฺปนฺนสฺส ธมฺมสฺส นตฺถิ\-ปจฺจเยน ปจฺจโย. วิคต\-ปจฺจเยน ปจฺจโย.}

\prose{ปจฺจุปฺปนฺโน ธมฺโม ปจฺจุปฺปนฺนสฺส ธมฺมสฺส อวิคต\-ปจฺจเยน ปจฺจโย ฯเปฯ.}

\csromanmatikaanchor{mtk.191}
\csromantocmarkref{subparagraph}{1. ปจฺจยานุโลม 2. สงฺขฺยาวาร}{mtk.191}
\bookmark[dest={mtk.191},level=5]{1. ปจฺจยานุโลม 2. สงฺขฺยาวาร}
\csromanmarkh{18. อตีตตฺติก}\csromanmarkhii{1. ปจฺจยานุโลม 2. สงฺขฺยาวาร}\csromanheadernomark{6}{18. อตีตตฺติก \textbf{1. ปจฺจยานุโลม 2. สงฺขฺยาวาร}}
\proseitemwithcloserruled{17}{\csromanfolio{417}เหตุยา เอกํ, อารมฺมเณ ตีณิ, อธิปติยา ตีณิ, อนนฺตเร เอกํ, สมนนฺตเร เอกํ, สหชาเต, อญฺญมญฺเญ, นิสฺสเย เอกํ, อุปนิสฺสเย ตีณิ, ปุเรชาเต, ปจฺฉาชาเต, อาเสวเน เอกํ, กมฺเม เทฺว, วิปาเก, อาหาเร เอกํ ฯเปฯ อวิคเต เอกํ. (เอวํ คเณตพฺพํ.)}{อนุโลมํ.}
\csromanmatikaanchor{mtk.192}
\csromantocmarkref{subparagraph}{2. ปจฺจนียุทฺธาร}{mtk.192}
\bookmark[dest={mtk.192},level=5]{2. ปจฺจนียุทฺธาร}
\csromanheader{6}{2. ปจฺจนียุ\-ทฺธาร}
\proseitem{18}{อตีโต ธมฺโม ปจฺจุปฺปนฺนสฺส ธมฺมสฺส อารมฺมณ\-ปจฺจเยน ปจฺจโย. อุปนิสฺสย\-ปจฺจเยน ปจฺจโย. กมฺมปจฺจเยน ปจฺจโย. (1)}

\prose{อนาคโต ธมฺโม ปจฺจุปฺปนฺนสฺส ธมฺมสฺส อารมฺมณ\-ปจฺจเยน ปจฺจโย. อุปนิสฺสย\-ปจฺจเยน ปจฺจโย. (1)}

\prose{ปจฺจุปฺปนฺโน ธมฺโม ปจฺจุปฺปนฺนสฺส ธมฺมสฺส อารมฺมณ\-ปจฺจเยน ปจฺจโย. สหชาต\-ปจฺจเยน ปจฺจโย. อุปนิสฺสย\-ปจฺจเยน ปจฺจโย. ปุเรชาต\-ปจฺจเยน ปจฺจโย. ปจฺฉาชาต\-ปจฺจเยน ปจฺจโย. อาหารปจฺจเยน ปจฺจโย. อินฺทฺริย\-ปจฺจเยน ปจฺจโย. (1)}

\csromanmatikaanchor{mtk.193}
\csromantocmarkref{subparagraph}{2. ปจฺจยปจฺจนีย 2. สงฺขฺยาวาร}{mtk.193}
\bookmark[dest={mtk.193},level=5]{2. ปจฺจยปจฺจนีย 2. สงฺขฺยาวาร}
\csromanmarkh{2. ปจฺจย\-ปจฺจนีย}\csromanmarkhii{2. สงฺขฺยาวาร}\csromanheadernomark{6}{2. ปจฺจย\-ปจฺจนีย 2. สงฺขฺยาวาร}
\proseitem{19}{นเหตุยา ตีณิ, นอารมฺมเณ ตีณิ, นธิปติยา ตีณิ, นอนนฺตเร ตีณิ ฯเปฯ นสมฺปยุตฺเต ตีณิ, นวิปฺปยุตฺเต ตีณิ, โนอตฺถิยา เทฺว, โนนตฺถิยา ตีณิ, โนวิคเต ตีณิ, โนอวิคเต เทฺว. (เอวํ คเณตพฺพํ.)}

\vspace{\tipitakaparskip}
\csromancenter{ปจฺจนียํ.}
\csromanmatikaanchor{mtk.194}
\csromanmatikaanchor{mtk.195}
\csromantocmarkref{subparagraph}{3. ปจฺจยานุโลมปจฺจนีย}{mtk.194}
\bookmark[dest={mtk.194},level=5]{3. ปจฺจยานุโลมปจฺจนีย}
\csromantocmarkref{subparagraph}{4. ปจฺจยปจฺจนียานุโลม}{mtk.195}
\bookmark[dest={mtk.195},level=5]{4. ปจฺจยปจฺจนียานุโลม}
\proseitemwithcloserruled{20}{\csromanfolio{418}เหตุปจฺจยา นอารมฺมเณ เอกํ, นอธิปติยา, นอนนฺตเร, นสมนนฺตเร, นอญฺญมญฺเญ, นอุปนิสฺสเย ฯเปฯ นสมฺปยุตฺเต, นวิปฺปยุตฺเต, โนนตฺถิยา, โนวิคเต เอกํ. (เอวํ คเณตพฺพํ.)}{อนุโลม\-ปจฺจนียํ.}
\proseitem{21}{นเหตุปจฺจยา อารมฺมเณ ตีณิ, อธิปติยา ตีณิ, อนนฺตเร, สมนนฺตเร, สหชาเต, อญฺญมญฺเญ, นิสฺสเย เอกํ, อุปนิสฺสเย ตีณิ, ปุเรชาเต เอกํ, ปจฺฉาชาเต, อาเสวเน เอกํ ฯเปฯ กมฺเม เทฺว, วิปาเก เอกํ. (อิเมสุ ปเทสุ เอกํเยว.) อวิคเต เอกํ. (เอวํ คเณตพฺพํ.)}

\vspace{\tipitakaparskip}
\csromancenter{ปจฺจนียา\-นุโลมํ.}
\csromancenter{ปญฺหาวาโร.}
\nitthitam{อตีตตฺติกํ นิฏฺฐิตํ.}
\csromanmatikaanchor{mtk.196}
\csromantocmarknopage{subparagraph}{19. อตีตารมฺมณตฺติก 1. ปฏิจฺจวาร}{mtk.196}
\bookmark[dest={mtk.196},level=5]{19. อตีตารมฺมณตฺติก 1. ปฏิจฺจวาร}
\csromanmarkh{19. อตีตา\-รมฺมณ\-ตฺติก}\csromanmarkhii{19. อตีตา\-รมฺมณ\-ตฺติก 1. ปฏิจฺจวาร}\csromanheadernomark{6}{19. อตีตา\-รมฺมณ\-ตฺติก 19. อตีตา\-รมฺมณ\-ตฺติก 1. ปฏิจฺจวาร}
\csromanmarkh{1. ปจฺจยานุโลม}\csromanmarkhii{1. วิภงฺควาร}\csromanheadernomark{2}{1. \textbf{ปจฺจยานุโลม 1. วิภงฺควาร}}
\csromanmatikaanchor{mtk.197}
\csromanmatikaanchor{mtk.198}
\csromantocmarknopage{subparagraph}{1. ปจฺจยานุโลม 1. วิภงฺควาร}{mtk.197}
\bookmark[dest={mtk.197},level=5]{1. ปจฺจยานุโลม 1. วิภงฺควาร}
\csromantocmarkref{subparagraph}{เหตุ}{mtk.198}
\bookmark[dest={mtk.198},level=5]{เหตุ}
\csromanheader{6}{\textbf{เหตุ}}
\proseitem{1}{\csromanfolio{419}อตีตารมฺมณํ ธมฺมํ ปฏิจฺจ อตีตารมฺมโณ ธมฺโม อุปฺปชฺชติ เหตุปจฺจยา. อตีตารมฺมณํ เอกํ ขนฺธํ ปฏิจฺจ ตโย ขนฺธา ฯเปฯ เทฺว ขนฺเธ ปฏิจฺจ เทฺว ขนฺธา. ปฏิสนฺธิ\-กฺขเณ อตีตารมฺมณํ เอกํ ขนฺธํ ปฏิจฺจ ตโย ขนฺธา ฯเปฯ เทฺว ขนฺเธ ปฏิจฺจ เทฺว ขนฺธา. (1)}

\proseitem{2}{อนาคตา\-รมฺมณํ ธมฺมํ ปฏิจฺจ อนาคตารมฺมโณ ธมฺโม อุปฺปชฺชติ เหตุปจฺจยา. อนาคตา\-รมฺมณํ เอกํ ขนฺธํ ปฏิจฺจ ตโย ขนฺธา ฯเปฯ เทฺว ขนฺธา. (1)}

\proseitem{3}{ปจฺจุปฺปนฺนา\-รมฺมณํ ธมฺมํ ปฏิจฺจ ปจฺจุปฺปนฺนา\-รมฺมโณ ธมฺโม อุปฺปชฺชติ เหตุปจฺจยา. ปจฺจุปฺปนฺนา\-รมฺมณํ เอกํ ขนฺธํ ปฏิจฺจ ตโย ขนฺธา ฯเปฯ เทฺว ขนฺธา. ปฏิสนฺธิ\-กฺขเณ ปจฺจุปฺปนฺนา\-รมฺมณํ เอกํ ขนฺธํ ปฏิจฺจ ตโย ขนฺธา ฯเปฯ เทฺว ขนฺธา. (1)}

\csromanmatikaanchor{mtk.199}
\csromantocmarkref{subparagraph}{อารมฺมณาทิ}{mtk.199}
\bookmark[dest={mtk.199},level=5]{อารมฺมณาทิ}
\csromanheader{6}{\textbf{อารมฺมณาทิ}}
\proseitem{4}{อตีตารมฺมณํ ธมฺมํ ปฏิจฺจ อตีตารมฺมโณ ธมฺโม อุปฺปชฺชติ อารมฺมณ\-ปจฺจยา. อธิปติ\-ปจฺจยา. (อธิปติยา ปฏิสนฺธิ นตฺถิ.) อนนฺตร\-ปจฺจยา. สมนนฺตร\-ปจฺจยา. สหชาต\-ปจฺจยา. อญฺญมญฺญ\-ปจฺจยา. นิสฺสยปจฺจยา. อุปนิสฺสย\-ปจฺจยา. ปุเรชาต\-ปจฺจยา. อาเสวนปจฺจยา. (ปุเรชาเตปิ อาเสวเนปิ ปฏิสนฺธิ นตฺถิ.) กมฺมปจฺจยา. วิปากปจฺจยา. (วิปากํ อตีตารมฺมณํ เอกํ ขนฺธํ, ติสฺโสปิ ปญฺหา ปริปุณฺณา. ปวตฺติ\-ปฏิสนฺธิ กาตพฺพา.) อาหารปจฺจยา. อินฺทฺริยปจฺจยา. ฌานปจฺจยา. มคฺคปจฺจยา. สมฺปยุตฺต\-ปจฺจยา. วิปฺปยุตฺต\-ปจฺจยา. อตฺถิปจฺจยา. นตฺถิปจฺจยา. วิคตปจฺจยา. อวิคตปจฺจยา.}

\csromanmatikaanchor{mtk.200}
\csromantocmarkref{subparagraph}{1. ปจฺจยานุโลม 2. สงฺขฺยาวาร}{mtk.200}
\bookmark[dest={mtk.200},level=5]{1. ปจฺจยานุโลม 2. สงฺขฺยาวาร}
\csromanmarkh{1. ปจฺจยานุโลม}\csromanmarkhii{2. สงฺขฺยาวาร}\csromanheadernomark{6}{1. \textbf{ปจฺจยานุโลม 2. สงฺขฺยาวาร}}
\proseitem{5}{\textbf{เหตุ}ยา ตีณิ, อารมฺมเณ ตีณิ, อธิปตียา ตีณิ ฯเปฯ (สพฺพตฺถ ตีณิ,) วิคเต ตีณิ, อวิคเต ตีณิ. (เอวํ คเณตพฺพํ.)}

\vspace{\tipitakaparskip}
\csromancenter{อนุโลมํ.}
\csromanmarkh{1. ปจฺจย\-ปจฺจนีย}\csromanmarkhii{1. วิภงฺควาร}\csromanheadernomark{2}{1. ปจฺจย\-ปจฺจนีย 1. วิภงฺควาร}
\csromanmatikaanchor{mtk.201}
\csromanmatikaanchor{mtk.202}
\csromantocmarknopage{subparagraph}{2. ปจฺจยปจฺจนีย 1. วิภงฺควาร}{mtk.201}
\bookmark[dest={mtk.201},level=5]{2. ปจฺจยปจฺจนีย 1. วิภงฺควาร}
\csromantocmarkref{subparagraph}{นเหตุ}{mtk.202}
\bookmark[dest={mtk.202},level=5]{นเหตุ}
\csromanheader{6}{\textbf{นเหตุ}}
\csromanmatikaanchor{mtk.203}
\csromantocmarkref{subparagraph}{นอธิปติ}{mtk.203}
\bookmark[dest={mtk.203},level=5]{นอธิปติ}
\proseitem{6}{\csromanfolio{420}อตีตารมฺมณํ ธมฺมํ ปฏิจฺจ อตีตารมฺมโณ ธมฺโม อุปฺปชฺชติ นเหตุปจฺจยา. อเหตุกํ อตีตารมฺมณํ เอกํ ขนฺธํ ปฏิจฺจ ตโย ขนฺธา ฯเปฯ เทฺว ขนฺเธ ฯเปฯ. อเหตุก\-ปฏิสนฺธิ\-กฺขเณ ฯเปฯ. วิจิกิจฺฉา\-สหคเต อุทฺธจฺจ\-สหคเต ขนฺเธ ปฏิจฺจ วิจิกิจฺฉา\-สหคโต อุทฺธจฺจ\-สหคโต โมโห. (1)}

\proseitem{7}{อนาคตา\-รมฺมณํ ธมฺมํ ปฏิจฺจ อนาคตารมฺมโณ ธมฺโม อุปฺปชฺชติ นเหตุปจฺจยา. อเหตุกํ อนาคตา\-รมฺมณํ เอกํ ขนฺธํ ปฏิจฺจ ตโย ขนฺธา ฯเปฯ เทฺว ขนฺเธ ฯเปฯ. วิจิกิจฺฉา\-สหคเต อุทฺธจฺจ\-สหคเต ขนฺเธ ปฏิจฺจ วิจิกิจฺฉา\-สหคโต อุทฺธจฺจ\-สหคโต โมโห. (1)}

\proseitem{8}{ปจฺจุปฺปนฺนา\-รมฺมณํ ธมฺมํ ปฏิจฺจ ปจฺจุปฺปนฺนา\-รมฺมโณ ธมฺโม อุปฺปชฺชติ นเหตุปจฺจยา. อเหตุกํ ปจฺจุปฺปนฺนา\-รมฺมณํ เอกํ ขนฺธํ ปฏิจฺจ ตโย ขนฺธา ฯเปฯ เทฺว ขนฺเธ ฯเปฯ. อเหตุก\-ปฏิสนฺธิ\-กฺขเณ ฯเปฯ. วิจิกิจฺฉา\-สหคเต อุทฺธจฺจ\-สหคเต ขนฺเธ ปฏิจฺจ วิจิกิจฺฉา\-สหคโต อุทฺธจฺจ\-สหคโต โมโห. (1)}

\prose{นอธิปติ}

\proseitem{9}{อตีตารมฺมณํ ธมฺมํ ปฏิจฺจ อตีตารมฺมโณ ธมฺโม อุปฺปชฺชติ นอธิปติ\-ปจฺจยา. (อนุโลม\-สหชาต\-สทิสํ.)}

\csromanmatikaanchor{mtk.204}
\csromantocmarkref{subparagraph}{นปุเรชาต}{mtk.204}
\bookmark[dest={mtk.204},level=5]{นปุเรชาต}
\csromanheader{6}{\textbf{นปุเรชาต}}
\proseitem{10}{อตีตารมฺมณํ ธมฺมํ ปฏิจฺจ อตีตารมฺมโณ ธมฺโม อุปฺปชฺชติ นปุเรชาต\-ปจฺจยา. อรูเป อตีตารมฺมณํ เอกํ ขนฺธํ ปฏิจฺจ ฯเปฯ เทฺว ขนฺเธ ฯเปฯ. ปฏิสนฺธิ\-กฺขเณ ฯเปฯ. (1)}

\prose{อนาคตา\-รมฺมณํ ธมฺมํ ปฏิจฺจ อนาคตารมฺมโณ ธมฺโม อุปฺปชฺชติ นปุเรชาต\-ปจฺจยา. อรูเป อนาคตา\-รมฺมณํ เอกํ ขนฺธํ ปฏิจฺจ ตโย ขนฺธา ฯเปฯ เทฺว ขนฺเธ ฯเปฯ. (1)}

\prose{ปจฺจุปฺปนฺนา\-รมฺมณํ ธมฺมํ ปฏิจฺจ ปจฺจุปฺปนฺนา\-รมฺมโณ ธมฺโม อุปฺปชฺชติ นปุเรชาต\-ปจฺจยา. ปฏิสนฺธิ\-กฺขเณ ปจฺจุปฺปนฺนา\-รมฺมณํ เอกํ ขนฺธํ ปฏิจฺจ ฯเปฯ เทฺว ขนฺเธ ฯเปฯ. (1)}

\csromanmatikaanchor{mtk.205}
\csromantocmarkref{subparagraph}{นปจฺฉาชาตาทิ}{mtk.205}
\bookmark[dest={mtk.205},level=5]{นปจฺฉาชาตาทิ}
\csromanheader{6}{19. อตีตา\-รมฺมณ\-ตฺติก นปจฺฉาชาตาทิ}
\proseitem{11}{\csromanfolio{421}อตีตารมฺมณํ ธมฺมํ ปฏิจฺจ อตีตารมฺมโณ ธมฺโม อุปฺปชฺชติ นปจฺฉาชาต\-ปจฺจยา. นอาเสวน\-ปจฺจยา. (นอธิปติ\-สทิสา.) นกมฺมปจฺจยา. อตีตารมฺมเณ ขนฺเธ ปฏิจฺจ อตีตารมฺมณา เจตนา. (1)}

\prose{อนาคตา\-รมฺมณํ ธมฺมํ ปฏิจฺจ อนาคตารมฺมโณ ธมฺโม อุปฺปชฺชติ นกมฺมปจฺจยา. อนาคตารมฺมเณ ขนฺเธ ปฏิจฺจ อนาคตารมฺมณา เจตนา. (1)}

\prose{ปจฺจุปฺปนฺนา\-รมฺมณํ ธมฺมํ ปฏิจฺจ ปจฺจุปฺปนฺนา\-รมฺมโณ ธมฺโม อุปฺปชฺชติ นกมฺมปจฺจยา. ปจฺจุปฺปนฺนา\-รมฺมเณ ขนฺเธ ปฏิจฺจ ปจฺจุปฺปนฺนา\-รมฺมณา เจตนา. (1)}

\proseitem{12}{อตีตารมฺมณํ ธมฺมํ ปฏิจฺจ อตีตารมฺมโณ ธมฺโม อุปฺปชฺชติ นวิปาก\-ปจฺจยา ฯเปฯ. (นวิปาเก ปฏิสนฺธิ นตฺถิ.)}

\csromanheader{2}{\textbf{นฌาน}}
\proseitem{13}{ปจฺจุปฺปนฺนา\-รมฺมณํ ธมฺมํ ปฏิจฺจ ปจฺจุปฺปนฺนา\-รมฺมโณ ธมฺโม อุปฺปชฺชติ นฌานปจฺจยา. ปญฺจ\-วิญฺญาณสหคตํ เอกํ ขนฺธํ ปฏิจฺจ ตโย ขนฺธา ฯเปฯ เทฺว ขนฺเธ ฯเปฯ. (1)}

\csromanheader{2}{\textbf{นมคฺค}}
\proseitem{14}{อตีตารมฺมณํ ธมฺมํ ปฏิจฺจ อตีตารมฺมโณ ธมฺโม อุปฺปชฺชติ นมคฺคปจฺจยา. อเหตุกํ อตีตารมฺมณํ เอกํ ขนฺธํ ปฏิจฺจ ตโย ขนฺธา ฯเปฯ. (นเหตุสทิสา. ติสฺโส ปญฺหา. โมโห นตฺถิ.)}

\csromanheader{2}{\textbf{นวิปฺปยุตฺต}}
\proseitem{15}{อตีตารมฺมณํ ธมฺมํ ปฏิจฺจ อตีตารมฺมโณ ธมฺโม อุปฺปชฺชติ นวิปฺปยุตฺต\-ปจฺจยา. อรูเป อตีตารมฺมณํ เอกํ ขนฺธํ ปฏิจฺจ ตโย ขนฺธา ฯเปฯ เทฺว ขนฺเธ ฯเปฯ. (1)}

\prose{อนาคตา\-รมฺมณํ ธมฺมํ ปฏิจฺจ อนาคตารมฺมโณ ธมฺโม อุปฺปชฺชติ นวิปฺปยุตฺต\-ปจฺจยา. อรูเป อนาคตา\-รมฺมณํ เอกํ ขนฺธํ ปฏิจฺจ ตโย ขนฺธา ฯเปฯ เทฺว ขนฺเธ ฯเปฯ. (1)}

\csromanmatikaanchor{mtk.206}
\csromantocmarkref{subparagraph}{2. ปจฺจยปจฺจนีย 2. สงฺขฺยาวาร}{mtk.206}
\bookmark[dest={mtk.206},level=5]{2. ปจฺจยปจฺจนีย 2. สงฺขฺยาวาร}
\csromanmarkh{2. ปจฺจย\-ปจฺจนีย}\csromanmarkhii{2. สงฺขฺยาวาร}\csromanheadernomark{6}{2. ปจฺจย\-ปจฺจนีย 2. สงฺขฺยาวาร}
\csromanmatikaanchor{mtk.207}
\csromanmatikaanchor{mtk.209}
\csromantocmarknopage{subparagraph}{3. ปจฺจยานุโลมปจฺจนีย}{mtk.207}
\bookmark[dest={mtk.207},level=5]{3. ปจฺจยานุโลมปจฺจนีย}
\proseitemwithcloserruled{16}{\csromanfolio{422}นเหตุยา ตีณิ, นอธิปติยา, นปุเรชาเต, นปจฺฉาชาเต, นอาเสวเน, นกมฺเม, นวิปาเก ตีณิ, นฌาเน เอกํ, นมคฺเค ตีณิ, นวิปฺปยุตฺเต เทฺว. (เอวํ คเณตพฺพํ.)}{ปจฺจนียํ.}
\csromancenter{\textbf{เหตุ}ปจฺจยา นอธิปติยา ตีณิ, นปุเรชาเต, นปจฺฉาชาเต, นอาเสวเน, นกมฺเม, นวิปาเก ตีณิ, นวิปฺปยุตฺเต เทฺว. (เอวํ คเณตพฺพํ.)}
\nitthitamruled{อนุโลม\-ปจฺจนียํ.}
\proseitem{18}{นเหตุปจฺจยา อารมฺมเณ ตีณิ ฯเปฯ (สพฺพตฺถ ตีณิ.) อวิคเต ตีณิ. (เอวํ คเณตพฺพํ.) ปจฺจนียา\-นุโลมํ. ปฏิจฺจวาโร.}

\prose{(สหชาตวาโรปิ ปจฺจยวาโรปิ นิสฺสยวาโรปิ สํสฏฺฐ\-วาโรปิ สมฺปยุตฺตวาโรปิ ปฏิจฺจ\-วาร\-สทิสา.)}
\csromansectionrule

\csromanmatikaanchor{mtk.208}
\csromanmatikaanchor{mtk.210}
\csromantocmarkref{subparagraph}{เหตุทุก}{mtk.208}
\bookmark[dest={mtk.208},level=5]{เหตุทุก}
\csromantocmarkref{subparagraph}{4. ปจฺจยปจฺจนียานุโลม ...}{mtk.209}
\bookmark[dest={mtk.209},level=5]{4. ปจฺจยปจฺจนียานุโลม ...}
\csromantocmarknopage{subparagraph}{19. อตีตารมฺมณตฺติก 7. ปญฺหาวาร}{mtk.210}
\bookmark[dest={mtk.210},level=5]{19. อตีตารมฺมณตฺติก 7. ปญฺหาวาร}
\csromanmarkh{19. อตีตา\-รมฺมณ\-ตฺติก}\csromanmarkhii{7. ปญฺหาวาร}\csromanheadernomark{6}{19. อตีตา\-รมฺมณ\-ตฺติก 7. ปญฺหาวาร}
\csromanmarkh{1. ปจฺจยานุโลม}\csromanmarkhii{1. วิภงฺควาร}\csromanheadernomark{2}{1. \textbf{ปจฺจยานุโลม 1. วิภงฺควาร}}
\csromanmatikaanchor{mtk.211}
\csromanmatikaanchor{mtk.212}
\csromantocmarknopage{subparagraph}{1. ปจฺจยานุโลม 1. วิภงฺควาร}{mtk.211}
\bookmark[dest={mtk.211},level=5]{1. ปจฺจยานุโลม 1. วิภงฺควาร}
\csromantocmarkref{subparagraph}{เหตุ}{mtk.212}
\bookmark[dest={mtk.212},level=5]{เหตุ}
\csromanheader{6}{\textbf{เหตุ}}
\csromanmatikaanchor{mtk.213}
\csromantocmarkref{subparagraph}{อารมฺมณาทิ}{mtk.213}
\bookmark[dest={mtk.213},level=5]{อารมฺมณาทิ}
\proseitem{19}{อตีตารมฺมโณ ธมฺโม อตีตา\-รมฺมณสฺส ธมฺมสฺส เหตุปจฺจเยน ปจฺจโย. อตีตารมฺมณา เหตู สมฺปยุตฺตกานํ ขนฺธานํ \csromanfolio{423}เหตุปจฺจเยน ปจฺจโย. ปฏิสนฺธิ\-กฺขเณ อตีตารมฺมณา เหตู สมฺปยุตฺตกานํ ขนฺธานํ เหตุปจฺจเยน ปจฺจโย. (1)}

\prose{อนาคตารมฺมโณ ธมฺโม อนาคตา\-รมฺมณสฺส ธมฺมสฺส เหตุปจฺจเยน ปจฺจโย. อนาคตารมฺมณา เหตู สมฺปยุตฺตกานํ ขนฺธานํ เหตุปจฺจเยน ปจฺจโย. (1)}

\prose{ปจฺจุปฺปนฺนา\-รมฺมโณ ธมฺโม ปจฺจุปฺปนฺนา\-รมฺมณสฺส ธมฺมสฺส เหตุปจฺจเยน ปจฺจโย. ปจฺจุปฺปนฺนา\-รมฺมณา เหตู สมฺปยุตฺตกานํ ขนฺธานํ เหตุปจฺจเยน ปจฺจโย. ปฏิสนฺธิ\-กฺขเณ ปจฺจุปฺปนฺนา\-รมฺมณา เหตู สมฺปยุตฺตกานํ ขนฺธานํ เหตุปจฺจเยน ปจฺจโย. (1)}

\csromanheader{2}{\textbf{อารมฺมณ}}
\proseitem{19}{อตีตารมฺมโณ ธมฺโม อตีตา\-รมฺมณสฺส ธมฺมสฺส อารมฺมณ\-ปจฺจเยน ปจฺจโย. อตีตํ วิญฺญาณญฺจา\-ยตนํ ปจฺจเวกฺขติ, เนว\-สญฺญา\-นา\-สญฺญา\-ยตนํ ปจฺจเวกฺขติ. อตีตารมฺมณํ อตีตํ อิทฺธิวิธ\-ญาณํ ปจฺจเวกฺขติ, เจโตปริย\-ญาณํ. ปุพฺเพ\-นิวาสา\-นุสฺสติ\-ญาณํ. ยถากมฺมูปคญาณํ ปจฺจเวกฺขติ. อริยา อตีตารมฺมเณ ปหีเน กิเลเส ปจฺจเวกฺขนฺติ, วิกฺขมฺภิเต กิเลเส ปจฺจเวกฺขนฺติ, ปุพฺเพ สมุทาจิณฺเณ กิเลเส ชานนฺติ. อตีตารมฺมเณ อตีเต ขนฺเธ อนิจฺจโต ทุกฺขโต อนตฺตโต วิปสฺสติ, อสฺสาเทติ อภินนฺทติ, ตํ อารพฺภ อตีตารมฺมโณ ราโค อุปฺปชฺชติ, ทิฏฺฐิ ฯเปฯ วิจิกิจฺฉา ฯเปฯ อุทฺธจฺจํ ฯเปฯ โทมนสฺสํ อุปฺปชฺชติ. อตีตารมฺมณา อตีตา ขนฺธา เจโตปริย\-ญาณสฺส, ปุพฺเพ\-นิวาสา\-นุสฺสติ\-ญาณสฺส, ยถากมฺมูปคญาณสฺส, อาวชฺชนาย อารมฺมณ\-ปจฺจเยน ปจฺจโย. (1)}

\prose{อตีตารมฺมโณ ธมฺโม อนาคตา\-รมฺมณสฺส ธมฺมสฺส อารมฺมณ\-ปจฺจเยน ปจฺจโย. อนาคตํ วิญฺญาณญฺจา\-ยตนํ ปจฺจเวกฺขติ, เนว\-สญฺญา\-นา\-สญฺญา\-ยตนํ ปจฺจเวกฺขติ. อตีตารมฺมณํ อนาคตํ อิทฺธิวิธ\-ญาณํ ปจฺจเวกฺขติ, เจโตปริย\-ญาณํ ฯเปฯ. ปุพฺเพ\-นิวาสา\-นุสฺสติ\-ญาณํ ฯเปฯ. ยถากมฺมูปคญาณํ ปจฺจเวกฺขติ. อตีตารมฺมเณ อนาคเต ขนฺเธ อนิจฺจโต ฯเปฯ วิปสฺสติ, อสฺสาเทติ อภินนฺทติ, ตํ อารพฺภ อนาคตารมฺมโณ ราโค อุปฺปชฺชติ ฯเปฯ โทมนสฺสํ อุปฺปชฺชติ. อตีตารมฺมณา อนาคตา ขนฺธา เจโตปริย\-ญาณสฺส, อนาคตํสญาณสฺส, อาวชฺชนาย อารมฺมณ\-ปจฺจเยน ปจฺจโย. (2)}

\proseitem{20}{\csromanfolio{424}อตีตารมฺมโณ ธมฺโม ปจฺจุปฺปนฺนา\-รมฺมณสฺส ธมฺมสฺส อารมฺมณ\-ปจฺจเยน ปจฺจโย. เจโตปริย\-ญาเณน อตีตา\-รมฺมณ\-ปจฺจุปฺปนฺน\-จิตฺต\-สมงฺคิสฺส จิตฺตํ ชานาติ. อตีตารมฺมณา ปจฺจุปฺปนฺนา ขนฺธา เจโตปริย\-ญาณสฺส, อาวชฺชนาย อารมฺมณ\-ปจฺจเยน ปจฺจโย. (3)}

\proseitem{21}{อนาคตารมฺมโณ ธมฺโม อนาคตา\-รมฺมณสฺส ธมฺมสฺส อารมฺมณ\-ปจฺจเยน ปจฺจโย. อนาคตา\-รมฺมณํ อนาคตํ อิทฺธิวิธ\-ญาณํ ปจฺจเวกฺขติ. เจโตปริย\-ญาณํ ฯเปฯ. อนาคตํสญาณํ ฯเปฯ. อนาคตารมฺมเณ อนาคเต ขนฺเธ อนิจฺจโต ฯเปฯ วิปสฺสติ, อสฺสาเทติ อภินนฺทติ, ตํ อารพฺภ อนาคตารมฺมโณ ราโค อุปฺปชฺชติ ฯเปฯ โทมนสฺสํ อุปฺปชฺชติ. อนาคตารมฺมณา อนาคตา ขนฺธา เจโตปริย\-ญาณสฺส, อนาคตํสญาณสฺส, อาวชฺชนาย อารมฺมณ\-ปจฺจเยน ปจฺจโย. (1)}

\prose{อนาคตารมฺมโณ ธมฺโม อตีตา\-รมฺมณสฺส ธมฺมสฺส อารมฺมณ\-ปจฺจเยน ปจฺจโย. อนาคตา\-รมฺมณํ อตีตํ อิทฺธิวิธ\-ญาณํ ปจฺจเวกฺขติ. เจโตปริย\-ญาณํ. อนาคตํสญาณํ ฯเปฯ. อริยา อนาคตารมฺมเณ ปหีเน กิเลเส ปจฺจเวกฺขนฺติ, วิกฺขมฺภิเต กิเลเส ปจฺจเวกฺขนฺติ. ปุพฺเพ สมุทาจิณฺเณ กิเลเส ชานนฺติ. อนาคตารมฺมเณ อตีเต ขนฺเธ อนิจฺจโต ฯเปฯ วิปสฺสติ, อสฺสาเทติ อภินนฺทติ, ตํ อารพฺภ อตีตารมฺมโณ ราโค อุปฺปชฺชติ ฯเปฯ โทมนสฺสํ อุปฺปชฺชติ. อนาคตารมฺมณา อตีตา ขนฺธา เจโตปริย\-ญาณสฺส, ปุพฺเพ\-นิวาสา\-นุสฺสติ\-ญาณสฺส, ยถากมฺมูปคญาณสฺส, อาวชฺชนาย อารมฺมณ\-ปจฺจเยน ปจฺจโย. (2)}

\prose{อนาคตารมฺมโณ ธมฺโม ปจฺจุปฺปนฺนา\-รมฺมณสฺส ธมฺมสฺส อารมฺมณ\-ปจฺจเยน ปจฺจโย. เจโตปริย\-ญาเณน อนาคตา\-รมฺมณ\-ปจฺจุปฺปนฺน\-จิตฺต\-สมงฺคิสฺส จิตฺตํ ชานาติ. อนาคตารมฺมณา ปจฺจุปฺปนฺนา ขนฺธา เจโตปริย\-ญาณสฺส, อาวชฺชนาย อารมฺมณ\-ปจฺจเยน ปจฺจโย. (3)}

\proseitem{22}{ปจฺจุปฺปนฺนา\-รมฺมโณ ธมฺโม ปจฺจุปฺปนฺนา\-รมฺมณสฺส ธมฺมสฺส อารมฺมณ\-ปจฺจเยน ปจฺจโย. เจโตปริย\-ญาเณน ปจฺจุปฺปนฺนา\-รมฺมณ\-ปจฺจุปฺปนฺน\-จิตฺต\-สมงฺคิสฺส จิตฺตํ ชานาติ. ปจฺจุปฺปนฺนา\-รมฺมณา ปจฺจุปฺปนฺนา ขนฺธา เจโตปริย\-ญาณสฺส, อาวชฺชนาย อารมฺมณ\-ปจฺจเยน ปจฺจโย. (1)}

\prose{\csromanfolio{425}ปจฺจุปฺปนฺนา\-รมฺมโณ ธมฺโม อตีตา\-รมฺมณสฺส ธมฺมสฺส อารมฺมณ\-ปจฺจเยน ปจฺจโย. อตีตํ ทิพฺพํ จกฺขุํ ปจฺจเวกฺขติ. ทิพฺพํ โสตธาตุํ ปจฺจเวกฺขติ. ปจฺจุปฺปนฺนา\-รมฺมณํ อตีตํ อิทฺธิวิธ\-ญาณํ ปจฺจเวกฺขติ. เจโตปริย\-ญาณํ ปจฺจเวกฺขติ. อริยา ปจฺจุปฺปนฺนา\-รมฺมเณ ปหีเน กิเลเส ปจฺจเวกฺขนฺติ. วิกฺขมฺภิเต กิเลเส ปจฺจเวกฺขนฺติ, ปุพฺเพ สมุทาจิณฺเณ กิเลเส ชานนฺติ. ปจฺจุปฺปนฺนา\-รมฺมเณ อตีเต ขนฺเธ อนิจฺจโต ฯเปฯ วิปสฺสติ, อสฺสาเทติ อภินนฺทติ, ตํ อารพฺภ อตีตารมฺมโณ ราโค อุปฺปชฺชติ ฯเปฯ โทมนสฺสํ อุปฺปชฺชติ. ปจฺจุปฺปนฺนา\-รมฺมณา อตีตา ขนฺธา เจโตปริย\-ญาณสฺส, ปุพฺเพ\-นิวาสา\-นุสฺสติ\-ญาณสฺส, ยถากมฺมูปคญาณสฺส, อาวชฺชนาย อารมฺมณ\-ปจฺจเยน ปจฺจโย. (2)}

\prose{ปจฺจุปฺปนฺนา\-รมฺมโณ ธมฺโม อนาคตา\-รมฺมณสฺส ธมฺมสฺส อารมฺมณ\-ปจฺจเยน ปจฺจโย. อนาคตํ ทิพฺพํ จกฺขุํ ปจฺจเวกฺขติ. ทิพฺพํ โสตธาตุํ ปจฺจเวกฺขติ. ปจฺจุปฺปนฺนา\-รมฺมณํ อนาคตํ อิทฺธิวิธ\-ญาณํ ปจฺจเวกฺขติ. เจโตปริย\-ญาณํ ฯเปฯ. ปจฺจุปฺปนฺนา\-รมฺมเณ อนาคเต ขนฺเธ อนิจฺจโต ฯเปฯ วิปสฺสติ ฯเปฯ ตํ อารพฺภ อนาคตารมฺมโณ ราโค อุปฺปชฺชติ ฯเปฯ โทมนสฺสํ อุปฺปชฺชติ. ปจฺจุปฺปนฺนา\-รมฺมณา อนาคตา ขนฺธา เจโตปริย\-ญาณสฺส, อนาคตํสญาณสฺส, อาวชฺชนาย อารมฺมณ\-ปจฺจเยน ปจฺจโย. (3)}

\proseitem{23}{อตีตารมฺมโณ ธมฺโม อตีตา\-รมฺมณสฺส ธมฺมสฺส อธิปติ\-ปจฺจเยน ปจฺจโย. อารมฺมณา\-ธิปติ, สหชาตา\-ธิปติ.  อารมฺมณา\-ธิปติ\csromandash{}อตีตํ วิญฺญาณญฺจา\-ยตนํ ครุํ กตฺวา ปจฺจเวกฺขติ, เนว\-สญฺญา\-นา\-สญฺญา\-ยตนํ ครุํ กตฺวา ปจฺจเวกฺขติ. อตีตารมฺมณํ อตีตํ อิทฺธิวิธ\-ญาณํ ครุํ กตฺวา ปจฺจเวกฺขติ, เจโตปริย\-ญาณํ ฯเปฯ. ปุพฺเพ\-นิวาสา\-นุสฺสติ\-ญาณํ ฯเปฯ. ยถากมฺมูปคญาณํ ครุํ กตฺวา ปจฺจเวกฺขติ. อตีตารมฺมเณ อตีเต ขนฺเธ ครุํ กตฺวา อสฺสาเทติ อภินนฺทติ, ตํ ครุํ กตฺวา อตีตารมฺมโณ ราโค อุปฺปชฺชติ, ทิฏฺฐิ อุปฺปชฺชติ. สหชาตา\-ธิปติ\csromandash{}อตีตา\-รมฺมณา\-ธิปติ สมฺปยุตฺตกานํ ขนฺธานํ อธิปติ\-ปจฺจเยน ปจฺจโย. (1)}

\prose{อตีตารมฺมโณ ธมฺโม อนาคตา\-รมฺมณสฺส ธมฺมสฺส อธิปติ\-ปจฺจเยน ปจฺจโย. อารมฺมณา\-ธิปติ\csromandash{}อนาคตํ วิญฺญาณญฺจา\-ยตนํ ครุํ กตฺวา ฯเปฯ. \csromanfolio{426}เนว\-สญฺญา\-นา\-สญฺญา\-ยตนํ ฯเปฯ. อตีตารมฺมณํ อนาคตํ อิทฺธิวิธ\-ญาณํ ครุํ กตฺวา ฯเปฯ. เจโตปริย\-ญาณํ ฯเปฯ. ปุพฺเพ\-นิวาสา\-นุสฺสติ\-ญาณํ ฯเปฯ. ยถากมฺมูปคญาณํ ฯเปฯ. อตีตารมฺมเณ อนาคเต ขนฺเธ ครุํ กตฺวา อสฺสาเทติ อภินนฺทติ, ตํ ครุํ กตฺวา อนาคตารมฺมโณ ราโค อุปฺปชฺชติ. ทิฏฺฐิ อุปฺปชฺชติ. (2)}

\proseitem{24}{อนาคตารมฺมโณ ธมฺโม อนาคตา\-รมฺมณสฺส ธมฺมสฺส อธิปติ\-ปจฺจเยน ปจฺจโย. อารมฺมณา\-ธิปติ, สหชาตา\-ธิปติ.  อารมฺมณา\-ธิปติ\csromandash{}อนาคตา\-รมฺมณํ อนาคตํ อิทฺธิวิธ\-ญาณํ ครุํ กตฺวา ฯเปฯ. เจโตปริย\-ญาณํ ฯเปฯ. อนาคตํสญาณํ ครุํ กตฺวา ปจฺจเวกฺขติ. อนาคตารมฺมเณ อนาคเต ขนฺเธ ครุํ กตฺวา อสฺสาเทติ อภินนฺทติ, ตํ ครุํ กตฺวา อนาคตารมฺมโณ ราโค อุปฺปชฺชติ. ทิฏฺฐิ อุปฺปชฺชติ. สหชาตา\-ธิปติ\csromandash{}อนาคตา\-รมฺมณา\-ธิปติ สมฺปยุตฺตกานํ ขนฺธานํ อธิปติ\-ปจฺจเยน ปจฺจโย. (1)}

\prose{อนาคตารมฺมโณ ธมฺโม อตีตา\-รมฺมณสฺส ธมฺมสฺส อธิปติ\-ปจฺจเยน ปจฺจโย. อารมฺมณา\-ธิปติ\csromandash{}อนาคตา\-รมฺมณํ อตีตํ อิทฺธิวิธ\-ญาณํ ครุํ กตฺวา ฯเปฯ. เจโตปริย\-ญาณํ ฯเปฯ. อนาคตํสญาณํ ครุํ กตฺวา ฯเปฯ. อนาคตารมฺมเณ อตีเต ขนฺเธ ครุํ กตฺวา อสฺสาเทติ อภินนฺทติ, ตํ ครุํ กตฺวา อตีตารมฺมโณ ราโค อุปฺปชฺชติ. ทิฏฺฐิ อุปฺปชฺชติ. (2)}

\proseitem{25}{ปจฺจุปฺปนฺนา\-รมฺมโณ ธมฺโม ปจฺจุปฺปนฺนา\-รมฺมณสฺส ธมฺมสฺส อธิปติ\-ปจฺจเยน ปจฺจโย. สหชาตา\-ธิปติ\csromandash{} ปจฺจุปฺปนฺนา\-รมฺมณา\-ธิปติ สมฺปยุตฺตกานํ ขนฺธานํ อธิปติ\-ปจฺจเยน ปจฺจโย. (1)}

\prose{ปจฺจุปฺปนฺนา\-รมฺมโณ ธมฺโม อตีตา\-รมฺมณสฺส ธมฺมสฺส อธิปติ\-ปจฺจเยน ปจฺจโย. อารมฺมณา\-ธิปติ\csromandash{}อตีตํ ทิพฺพํ จกฺขุํ ครุํ กตฺวา ปจฺจเวกฺขติ. ทิพฺพํ โสตธาตุํ ครุํ กตฺวา ปจฺจเวกฺขติ. ปจฺจุปฺปนฺนา\-รมฺมณํ อตีตํ อิทฺธิวิธ\-ญาณํ ครุํ กตฺวา ฯเปฯ. เจโตปริย\-ญาณํ ครุํ กตฺวา ฯเปฯ. ปจฺจุปฺปนฺนา\-รมฺมเณ อตีเต ขนฺเธ ครุํ กตฺวา อสฺสาเทติ อภินนฺทติ, ตํ ครุํ กตฺวา อตีตารมฺมโณ ราโค อุปฺปชฺชติ. ทิฏฺฐิ อุปฺปชฺชติ. (2)}

\prose{ปจฺจุปฺปนฺนา\-รมฺมโณ ธมฺโม อนาคตา\-รมฺมณสฺส ธมฺมสฺส อธิปติ\-ปจฺจเยน ปจฺจโย. อารมฺมณา\-ธิปติ\csromandash{}อนาคตํ ทิพฺพํ จกฺขุํ ครุํ กตฺวา ปจฺจเวกฺขติ. \csromanfolio{427}ทิพฺพํ โสตธาตุํ ครุํ กตฺวา ฯเปฯ. ปจฺจุปฺปนฺนา\-รมฺมณํ อนาคตํ อิทฺธิวิธ\-ญาณํ ครุํ กตฺวา ฯเปฯ. เจโตปริย\-ญาณํ ครุํ กตฺวา ฯเปฯ. ปจฺจุปฺปนฺนา\-รมฺมเณ อนาคเต ขนฺเธ ครุํ กตฺวา อสฺสาเทติ อภินนฺทติ, ตํ ครุํ กตฺวา อนาคตารมฺมโณ ราโค อุปฺปชฺชติ. ทิฏฺฐิ อุปฺปชฺชติ. (3)}

\csromanheader{2}{\textbf{อนนฺตร}}
\proseitem{26}{อตีตารมฺมโณ ธมฺโม อตีตา\-รมฺมณสฺส ธมฺมสฺส อนนฺตร\-ปจฺจเยน ปจฺจโย. ปุริมา ปุริมา อตีตารมฺมณา ขนฺธา ปจฺฉิมานํ ปจฺฉิมานํ อตีตา\-รมฺมณานํ ขนฺธานํ อนนฺตร\-ปจฺจเยน ปจฺจโย. (1)}

\prose{อตีตารมฺมโณ ธมฺโม อนาคตา\-รมฺมณสฺส ธมฺมสฺส อนนฺตร\-ปจฺจเยน ปจฺจโย. อตีตารมฺมณํ ภวงฺคํ อนาคตา\-รมฺมณาย อาวชฺชนาย อนนฺตร\-ปจฺจเยน ปจฺจโย. (2)}

\prose{อตีตารมฺมโณ ธมฺโม ปจฺจุปฺปนฺนา\-รมฺมณสฺส ธมฺมสฺส อนนฺตร\-ปจฺจเยน ปจฺจโย. อตีตารมฺมณํ จุติจิตฺตํ ปจฺจุปฺปนฺนา\-รมฺมณสฺส ปฏิสนฺธิจิตฺตสฺส อนนฺตร\-ปจฺจเยน ปจฺจโย. อตีตารมฺมณํ ภวงฺคํ ปจฺจุปฺปนฺนา\-รมฺมณาย อาวชฺชนาย อนนฺตร\-ปจฺจเยน ปจฺจโย. (3)}

\proseitem{27}{อนาคตารมฺมโณ ธมฺโม อนาคตา\-รมฺมณสฺส ธมฺมสฺส อนนฺตร\-ปจฺจเยน ปจฺจโย. ปุริมา ปุริมา อนาคตารมฺมณา ขนฺธา ปจฺฉิมานํ ปจฺฉิมานํ อนาคตา\-รมฺมณานํ ขนฺธานํ อนนฺตร\-ปจฺจเยน ปจฺจโย. (1)}

\prose{อนาคตารมฺมโณ ธมฺโม อตีตา\-รมฺมณสฺส ธมฺมสฺส อนนฺตร\-ปจฺจเยน ปจฺจโย. อนาคตา\-รมฺมณํ อิทฺธิวิธ\-ญาณํ อตีตา\-รมฺมณสฺส วุฏฺฐานสฺส ฯเปฯ. เจโตปริย\-ญาณํ อตีตา\-รมฺมณสฺส วุฏฺฐานสฺส ฯเปฯ. อนาคตํสญาณํ อตีตา\-รมฺมณสฺส วุฏฺฐานสฺส ฯเปฯ. อนาคตารมฺมณา ขนฺธา อตีตา\-รมฺมณสฺส วุฏฺฐานสฺส อนนฺตร\-ปจฺจเยน ปจฺจโย. (2)}

\proseitem{28}{ปจฺจุปฺปนฺนา\-รมฺมโณ ธมฺโม ปจฺจุปฺปนฺนา\-รมฺมณสฺส ธมฺมสฺส อนนฺตร\-ปจฺจเยน ปจฺจโย. ปุริมา ปุริมา ปจฺจุปฺปนฺนา\-รมฺมณา ขนฺธา ปจฺฉิมานํ ปจฺฉิมานํ ปจฺจุปฺปนฺนา\-รมฺมณานํ ขนฺธานํ อนนฺตร\-ปจฺจเยน ปจฺจโย. ปจฺจุปฺปนฺนา\-รมฺมณํ ปฏิสนฺธิจิตฺตํ ปจฺจุปฺปนฺนา\-รมฺมณสฺส ภวงฺคสฺส ฯเปฯ. ปจฺจุปฺปนฺนา\-รมฺมณํ ภวงฺคํ ปจฺจุปฺปนฺนา\-รมฺมณสฺส ภวงฺคสฺส อนนฺตร\-ปจฺจเยน ปจฺจโย. (1)}

\prose{\csromanfolio{428}ปจฺจุปฺปนฺนา\-รมฺมโณ ธมฺโม อตีตา\-รมฺมณสฺส ธมฺมสฺส อนนฺตร\-ปจฺจเยน ปจฺจโย. ปจฺจุปฺปนฺนา\-รมฺมณํ ปฏิสนฺธิจิตฺตํ อตีตา\-รมฺมณสฺส ภวงฺคสฺส ฯเปฯ. ปจฺจุปฺปนฺนา\-รมฺมณํ ภวงฺคํ อตีตา\-รมฺมณสฺส ภวงฺคสฺส ฯเปฯ. ปจฺจุปฺปนฺนา\-รมฺมณา ขนฺธา อตีตา\-รมฺมณสฺส วุฏฺฐานสฺส อนนฺตร\-ปจฺจเยน ปจฺจโย. (2)}

\csromanheader{2}{\textbf{สมนนฺตร}}
\proseitem{29}{อตีตารมฺมโณ ธมฺโม อตีตา\-รมฺมณสฺส ธมฺมสฺส สมนนฺตร\-ปจฺจเยน ปจฺจโย. (อนนฺตร\-สทิสํ.)}

\csromanheader{2}{\textbf{สหชาตาทิ}}
\proseitem{30}{อตีตารมฺมโณ ธมฺโม อตีตา\-รมฺมณสฺส ธมฺมสฺส สหชาต\-ปจฺจเยน ปจฺจโย. อญฺญมญฺญ\-ปจฺจเยน ปจฺจโย. นิสฺสย\-ปจฺจเยน ปจฺจโย. (ตโยปิ ปจฺจยา ปฏิจฺจ\-วาร\-สทิสา.)}

\csromanheader{2}{\textbf{อุปนิสฺสย}}
\proseitem{31}{อตีตารมฺมโณ ธมฺโม อตีตา\-รมฺมณสฺส ธมฺมสฺส อุปนิสฺสย\-ปจฺจเยน ปจฺจโย. อารมฺมณู\-ปนิสฺสโย, อนนฺตรู\-ปนิสฺสโย, ปกตู\-ปนิสฺสโย ฯเปฯ. ปกตู\-ปนิสฺสโย\csromandash{}อตีตารมฺมณา อนิจฺจา\-นุปสฺสนา, ทุกฺขา\-นุปสฺสนา, อนตฺตา\-นุปสฺสนา อตีตารมฺมณาย อนิจฺจา\-นุปสฺสนาย, ทุกฺขา\-นุปสฺสนาย, อนตฺตา\-นุปสฺสนาย อุปนิสฺสย\-ปจฺจเยน ปจฺจโย. (1)}

\prose{อตีตารมฺมโณ ธมฺโม อนาคตา\-รมฺมณสฺส ธมฺมสฺส อุปนิสฺสย\-ปจฺจเยน ปจฺจโย. อารมฺมณู\-ปนิสฺสโย, อนนฺตรู\-ปนิสฺสโย, ปกตู\-ปนิสฺสโย ฯเปฯ. ปกตู\-ปนิสฺสโย\csromandash{}อตีตารมฺมณา อนิจฺจา\-นุปสฺสนา, ทุกฺขา\-นุปสฺสนา, อนตฺตา\-นุปสฺสนา อนาคตา\-รมฺมณาย อนิจฺจา\-นุปสฺสนาย, ทุกฺขา\-นุปสฺสนาย, อนตฺตา\-นุปสฺสนาย อุปนิสฺสย\-ปจฺจเยน ปจฺจโย. (2)}

\prose{อตีตารมฺมโณ ธมฺโม ปจฺจุปฺปนฺนา\-รมฺมณสฺส ธมฺมสฺส อุปนิสฺสย\-ปจฺจเยน ปจฺจโย. อนนฺตรู\-ปนิสฺสโย, ปกตู\-ปนิสฺสโย ฯเปฯ. ปกตู\-ปนิสฺสโย\csromandash{}อตีตารมฺมณา อนิจฺจา\-นุปสฺสนา, ทุกฺขา\-นุปสฺสนา, อนตฺตา\-นุปสฺสนา ปจฺจุปฺปนฺนา\-รมฺมณาย อนิจฺจา\-นุปสฺสนาย, ทุกฺขา\-นุปสฺสนาย, อนตฺตา\-นุปสฺสนาย อุปนิสฺสย\-ปจฺจเยน ปจฺจโย. (3)}

\proseitem{32}{\csromanfolio{429}อนาคตารมฺมโณ ธมฺโม อนาคตา\-รมฺมณสฺส ธมฺมสฺส อุปนิสฺสย\-ปจฺจเยน ปจฺจโย. อารมฺมณู\-ปนิสฺสโย, อนนฺตรู\-ปนิสฺสโย, ปกตู\-ปนิสฺสโย ฯเปฯ. ปกตู\-ปนิสฺสโย\csromandash{}อนาคตารมฺมณา อนิจฺจา\-นุปสฺสนา, ทุกฺขา\-นุปสฺสนา, อนตฺตา\-นุปสฺสนา อนาคตา\-รมฺมณาย อนิจฺจา\-นุปสฺสนาย, ทุกฺขา\-นุปสฺสนาย, อนตฺตา\-นุปสฺสนาย อุปนิสฺสย\-ปจฺจเยน ปจฺจโย. (1)}

\prose{อนาคตารมฺมโณ ธมฺโม อตีตา\-รมฺมณสฺส ธมฺมสฺส อุปนิสฺสย\-ปจฺจเยน ปจฺจโย. อารมฺมณู\-ปนิสฺสโย, อนนฺตรู\-ปนิสฺสโย, ปกตู\-ปนิสฺสโย ฯเปฯ. ปกตู\-ปนิสฺสโย\csromandash{}อนาคตารมฺมณา อนิจฺจา\-นุปสฺสนา, ทุกฺขา\-นุปสฺสนา, อนตฺตา\-นุปสฺสนา อตีตารมฺมณาย อนิจฺจา\-นุปสฺสนาย, ทุกฺขา\-นุปสฺสนาย, อนตฺตา\-นุปสฺสนาย อุปนิสฺสย\-ปจฺจเยน ปจฺจโย. (2)}

\prose{อนาคตารมฺมโณ ธมฺโม ปจฺจุปฺปนฺนา\-รมฺมณสฺส ธมฺมสฺส อุปนิสฺสย\-ปจฺจเยน ปจฺจโย. ปกตู\-ปนิสฺสโย\csromandash{}อนาคตารมฺมณา อนิจฺจา\-นุปสฺสนา, ทุกฺขา\-นุปสฺสนา, อนตฺตา\-นุปสฺสนา ปจฺจุปฺปนฺนา\-รมฺมณาย อนิจฺจา\-นุปสฺสนาย, ทุกฺขา\-นุปสฺสนาย, อนตฺตา\-นุปสฺสนาย อุปนิสฺสย\-ปจฺจเยน ปจฺจโย. (3)}

\proseitem{33}{ปจฺจุปฺปนฺนา\-รมฺมโณ ธมฺโม ปจฺจุปฺปนฺนา\-รมฺมณสฺส ธมฺมสฺส อุปนิสฺสย\-ปจฺจเยน ปจฺจโย. อนนฺตรู\-ปนิสฺสโย, ปกตู\-ปนิสฺสโย ฯเปฯ. ปกตู\-ปนิสฺสโย\csromandash{}ปจฺจุปฺปนฺนา\-รมฺมณา อนิจฺจา\-นุปสฺสนา, ทุกฺขา\-นุปสฺสนา, อนตฺตา\-นุปสฺสนา ปจฺจุปฺปนฺนา\-รมฺมณาย อนิจฺจา\-นุปสฺสนาย, ทุกฺขา\-นุปสฺสนาย, อนตฺตา\-นุปสฺสนาย อุปนิสฺสย\-ปจฺจเยน ปจฺจโย. (1)}

\prose{ปจฺจุปฺปนฺนา\-รมฺมโณ ธมฺโม อตีตา\-รมฺมณสฺส ธมฺมสฺส อุปนิสฺสย\-ปจฺจเยน ปจฺจโย. อารมฺมณู\-ปนิสฺสโย, อนนฺตรู\-ปนิสฺสโย, ปกตู\-ปนิสฺสโย ฯเปฯ. ปกตู\-ปนิสฺสโย\csromandash{}ปจฺจุปฺปนฺนา\-รมฺมณา อนิจฺจา\-นุปสฺสนา, ทุกฺขา\-นุปสฺสนา, อนตฺตา\-นุปสฺสนา อตีตารมฺมณาย อนิจฺจา\-นุปสฺสนาย, ทุกฺขา\-นุปสฺสนาย, อนตฺตา\-นุปสฺสนาย อุปนิสฺสย\-ปจฺจเยน ปจฺจโย. (2)}

\prose{ปจฺจุปฺปนฺนา\-รมฺมโณ ธมฺโม อนาคตา\-รมฺมณสฺส ธมฺมสฺส อุปนิสฺสย\-ปจฺจเยน ปจฺจโย. อารมฺมณู\-ปนิสฺสโย, ปกตู\-ปนิสฺสโย ฯเปฯ. ปกตู\-ปนิสฺสโย\csromandash{}ปจฺจุปฺปนฺนา\-รมฺมณา อนิจฺจา\-นุปสฺสนา, ทุกฺขา\-นุปสฺสนา, อนตฺตา\-นุปสฺสนา อนาคตา\-รมฺมณาย อนิจฺจา\-นุปสฺสนาย, ทุกฺขา\-นุปสฺสนาย, อนตฺตา\-นุปสฺสนาย อุปนิสฺสย\-ปจฺจเยน ปจฺจโย. (3)}

\csromanheader{2}{\textbf{อาเสวน}}
\proseitem{34}{\csromanfolio{430}อตีตารมฺมโณ ธมฺโม อตีตา\-รมฺมณสฺส ธมฺมสฺส อาเสวน\-ปจฺจเยน ปจฺจโย. ปุริมา ปุริมา อตีตารมฺมณา ขนฺธา ปจฺฉิมานํ ปจฺฉิมานํ อตีตา\-รมฺมณานํ ขนฺธานํ อาเสวน\-ปจฺจเยน ปจฺจโย. (1)}

\prose{อนาคตารมฺมโณ ธมฺโม อนาคตา\-รมฺมณสฺส ธมฺมสฺส อาเสวน\-ปจฺจเยน ปจฺจโย. ปุริมา ปุริมา อนาคตารมฺมณา ขนฺธา ปจฺฉิมานํ ปจฺฉิมานํ อนาคตา\-รมฺมณานํ ขนฺธานํ อาเสวน\-ปจฺจเยน ปจฺจโย. (1)}

\prose{ปจฺจุปฺปนฺนา\-รมฺมโณ ธมฺโม ปจฺจุปฺปนฺนา\-รมฺมณสฺส ธมฺมสฺส อาเสวน\-ปจฺจเยน ปจฺจโย. ปุริมา ปุริมา ปจฺจุปฺปนฺนา\-รมฺมณา ขนฺธา ปจฺฉิมานํ ปจฺฉิมานํ ปจฺจุปฺปนฺนา\-รมฺมณานํ ขนฺธานํ อาเสวน\-ปจฺจเยน ปจฺจโย. (1)}

\csromanheader{2}{\textbf{กมฺม}}
\proseitem{35}{อตีตารมฺมโณ ธมฺโม อตีตา\-รมฺมณสฺส ธมฺมสฺส กมฺมปจฺจเยน ปจฺจโย. สหชาตา, นานากฺขณิกา.  สหชาตา อตีตารมฺมณา เจตนา สมฺปยุตฺตกานํ ขนฺธานํ กมฺมปจฺจเยน ปจฺจโย. ปฏิสนฺธิ\-กฺขเณ ฯเปฯ. นานากฺขณิกา อตีตารมฺมณา เจตนา วิปากานํ อตีตา\-รมฺมณานํ ขนฺธานํ กมฺมปจฺจเยน ปจฺจโย. (1)}

\prose{อตีตารมฺมโณ ธมฺโม อนาคตา\-รมฺมณสฺส ธมฺมสฺส กมฺมปจฺจเยน ปจฺจโย. นานากฺขณิกา อตีตารมฺมณา เจตนา วิปากานํ อนาคตา\-รมฺมณานํ ขนฺธานํ กมฺมปจฺจเยน ปจฺจโย. (2)}

\prose{อตีตารมฺมโณ ธมฺโม ปจฺจุปฺปนฺนา\-รมฺมณสฺส ธมฺมสฺส กมฺมปจฺจเยน ปจฺจโย. นานากฺขณิกา อตีตารมฺมณา เจตนา วิปากานํ ปจฺจุปฺปนฺนา\-รมฺมณานํ ขนฺธานํ กมฺมปจฺจเยน ปจฺจโย. (3)}

\proseitem{36}{อนาคตารมฺมโณ ธมฺโม อนาคตา\-รมฺมณสฺส ธมฺมสฺส กมฺมปจฺจเยน ปจฺจโย. สหชาตา นานากฺขณิกา.  สหชาตา อนาคตารมฺมณา เจตนา สมฺปยุตฺตกานํ ขนฺธานํ กมฺมปจฺจเยน ปจฺจโย. นานากฺขณิกา อนาคตารมฺมณา เจตนา วิปากานํ อนาคตา\-รมฺมณานํ ขนฺธานํ กมฺมปจฺจเยน ปจฺจโย. (1)}

\csromanmatikaanchor{mtk.214}
\csromanmatikaanchor{mtk.215}
\csromantocmarknopage{subparagraph}{1. ปจฺจยานุโลม 2. สงฺขฺยาวาร}{mtk.214}
\bookmark[dest={mtk.214},level=5]{1. ปจฺจยานุโลม 2. สงฺขฺยาวาร}
\csromantocmarkref{subparagraph}{สุทฺธ}{mtk.215}
\bookmark[dest={mtk.215},level=5]{สุทฺธ}
\prose{\csromanfolio{431}อนาคตารมฺมโณ ธมฺโม อตีตา\-รมฺมณสฺส ธมฺมสฺส กมฺมปจฺจเยน ปจฺจโย. นานากฺขณิกา อนาคตารมฺมณา เจตนา วิปากานํ อตีตา\-รมฺมณานํ ขนฺธานํ กมฺมปจฺจเยน ปจฺจโย. (2)}

\prose{อนาคตารมฺมโณ ธมฺโม ปจฺจุปฺปนฺนา\-รมฺมณสฺส ธมฺมสฺส กมฺมปจฺจเยน ปจฺจโย. นานากฺขณิกา อนาคตารมฺมณา เจตนา วิปากานํ ปจฺจุปฺปนฺนา\-รมฺมณานํ ขนฺธานํ กมฺมปจฺจเยน ปจฺจโย. (3)}

\proseitem{37}{ปจฺจุปฺปนฺนา\-รมฺมโณ ธมฺโม ปจฺจุปฺปนฺนา\-รมฺมณสฺส ธมฺมสฺส กมฺมปจฺจเยน ปจฺจโย. สหชาตา, นานากฺขณิกา.  สหชาตา ปจฺจุปฺปนฺนา\-รมฺมณา เจตนา สมฺปยุตฺตกานํ ขนฺธานํ กมฺมปจฺจเยน ปจฺจโย. ปฏิสนฺธิ\-กฺขเณ ฯเปฯ. นานากฺขณิกา ปจฺจุปฺปนฺนา\-รมฺมณา เจตนา วิปากานํ ปจฺจุปฺปนฺนา\-รมฺมณานํ ขนฺธานํ กมฺมปจฺจเยน ปจฺจโย. (1)}

\prose{ปจฺจุปฺปนฺนา\-รมฺมโณ ธมฺโม อตีตา\-รมฺมณสฺส ธมฺมสฺส กมฺมปจฺจเยน ปจฺจโย. นานากฺขณิกา ปจฺจุปฺปนฺนา\-รมฺมณา เจตนา วิปากานํ อตีตา\-รมฺมณานํ ขนฺธานํ กมฺมปจฺจเยน ปจฺจโย. (2)}

\prose{ปจฺจุปฺปนฺนา\-รมฺมโณ ธมฺโม อนาคตา\-รมฺมณสฺส ธมฺมสฺส กมฺมปจฺจเยน ปจฺจโย. นานากฺขณิกา ปจฺจุปฺปนฺนา\-รมฺมณา เจตนา วิปากานํ อนาคตา\-รมฺมณานํ ขนฺธานํ กมฺมปจฺจเยน ปจฺจโย. (3)}

\csromanheader{2}{\textbf{วิปากาทิ}}
\proseitem{38}{อตีตารมฺมโณ ธมฺโม อตีตา\-รมฺมณสฺส ธมฺมสฺส วิปาก\-ปจฺจเยน ปจฺจโย. อาหารปจฺจเยน ปจฺจโย. อินฺทฺริย\-ปจฺจเยน ปจฺจโย. ฌานปจฺจเยน ปจฺจโย. มคฺคปจฺจเยน ปจฺจโย. สมฺปยุตฺต\-ปจฺจเยน ปจฺจโย. อตฺถิปจฺจเยน ปจฺจโย. นตฺถิ\-ปจฺจเยน ปจฺจโย. วิคต\-ปจฺจเยน ปจฺจโย. อวิคต\-ปจฺจเยน ปจฺจโย.}

\csromanmarkh{1. ปจฺจยานุโลม}\csromanmarkhii{2. สงฺขฺยาวาร}\csromanheadernomark{2}{1. \textbf{ปจฺจยานุโลม 2. สงฺขฺยาวาร}}
\proseitemwithcloserruled{39}{เหตุยา ตีณิ, อารมฺมเณ นว, อธิปติยา สตฺต, อนนฺตเร สตฺต, สมนนฺตเร สตฺต, สหชาเต, อญฺญมญฺเญ, นิสฺสเย ตีณิ, \csromanfolio{432}อุปนิสฺสเย นว, อาเสวเน ตีณิ, กมฺเม นว, วิปาเก ตีณิ, อาหาเร ตีณิ, อินฺทฺริเย, ฌาเน, มคฺเค, สมฺปยุตฺเต ตีณิ, อตฺถิยา ตีณิ, นตฺถิยา สตฺต, วิคเต สตฺต, อวิคเต ตีณิ.  (เอวํ คเณตพฺพํ.)}{\textbf{อนุโลมํ.}}
\csromanmatikaanchor{mtk.216}
\csromantocmarkref{subparagraph}{2. ปจฺจนียุทฺธาร}{mtk.216}
\bookmark[dest={mtk.216},level=5]{2. ปจฺจนียุทฺธาร}
\csromanheader{6}{2. ปจฺจนียุ\-ทฺธาร}
\proseitem{40}{อตีตารมฺมโณ ธมฺโม อตีตา\-รมฺมณสฺส ธมฺมสฺส อารมฺมณ\-ปจฺจเยน ปจฺจโย. สหชาต\-ปจฺจเยน ปจฺจโย. อุปนิสฺสย\-ปจฺจเยน ปจฺจโย. กมฺมปจฺจเยน ปจฺจโย. (1)}

\prose{อตีตารมฺมโณ ธมฺโม อนาคตา\-รมฺมณสฺส ธมฺมสฺส อารมฺมณ\-ปจฺจเยน ปจฺจโย. อุปนิสฺสย\-ปจฺจเยน ปจฺจโย. กมฺมปจฺจเยน ปจฺจโย. (2)}

\prose{อตีตารมฺมโณ ธมฺโม ปจฺจุปฺปนฺนา\-รมฺมณสฺส ธมฺมสฺส อารมฺมณ\-ปจฺจเยน ปจฺจโย. อุปนิสฺสย\-ปจฺจเยน ปจฺจโย. กมฺมปจฺจเยน ปจฺจโย. (3)}

\proseitem{41}{อนาคตารมฺมโณ ธมฺโม อนาคตา\-รมฺมณสฺส ธมฺมสฺส อารมฺมณ\-ปจฺจเยน ปจฺจโย. สหชาต\-ปจฺจเยน ปจฺจโย. อุปนิสฺสย\-ปจฺจเยน ปจฺจโย. กมฺมปจฺจเยน ปจฺจโย. (1)}

\prose{อนาคตารมฺมโณ ธมฺโม อตีตา\-รมฺมณสฺส ธมฺมสฺส อารมฺมณ\-ปจฺจเยน ปจฺจโย. อุปนิสฺสย\-ปจฺจเยน ปจฺจโย. กมฺมปจฺจเยน ปจฺจโย. (2)}

\prose{อนาคตารมฺมโณ ธมฺโม ปจฺจุปฺปนฺนา\-รมฺมณสฺส ธมฺมสฺส อารมฺมณ\-ปจฺจเยน ปจฺจโย. อุปนิสฺสย\-ปจฺจเยน ปจฺจโย. กมฺมปจฺจเยน ปจฺจโย. (3)}

\proseitem{42}{ปจฺจุปฺปนฺนา\-รมฺมโณ ธมฺโม ปจฺจุปฺปนฺนา\-รมฺมณสฺส ธมฺมสฺส อารมฺมณ\-ปจฺจเยน ปจฺจโย. สหชาต\-ปจฺจเยน ปจฺจโย. อุปนิสฺสย\-ปจฺจเยน ปจฺจโย. กมฺมปจฺจเยน ปจฺจโย. (1)}

\prose{ปจฺจุปฺปนฺนา\-รมฺมโณ ธมฺโม อตีตา\-รมฺมณสฺส ธมฺมสฺส อารมฺมณ\-ปจฺจเยน ปจฺจโย. อุปนิสฺสย\-ปจฺจเยน ปจฺจโย. กมฺมปจฺจเยน ปจฺจโย. (2)}

\csromanmatikaanchor{mtk.217}
\csromanmatikaanchor{mtk.218}
\csromanmatikaanchor{mtk.219}
\csromanmatikaanchor{mtk.220}
\csromanmatikaanchor{mtk.221}
\csromantocmarkref{subparagraph}{2. ปจฺจยปจฺจนีย 2. สงฺขฺยาวาร}{mtk.217}
\bookmark[dest={mtk.217},level=5]{2. ปจฺจยปจฺจนีย 2. สงฺขฺยาวาร}
\csromantocmarknopage{subparagraph}{3. ปจฺจยานุโลมปจฺจนีย}{mtk.218}
\bookmark[dest={mtk.218},level=5]{3. ปจฺจยานุโลมปจฺจนีย}
\csromantocmarkref{subparagraph}{เหตุทุก}{mtk.219}
\bookmark[dest={mtk.219},level=5]{เหตุทุก}
\csromantocmarknopage{subparagraph}{4. ปจฺจยปจฺจนียานุโลม}{mtk.220}
\bookmark[dest={mtk.220},level=5]{4. ปจฺจยปจฺจนียานุโลม}
\csromantocmarkref{subparagraph}{นเหตุทุก}{mtk.221}
\bookmark[dest={mtk.221},level=5]{นเหตุทุก}
\prose{\csromanfolio{433}ปจฺจุปฺปนฺนา\-รมฺมโณ ธมฺโม อนาคตา\-รมฺมณสฺส ธมฺมสฺส อารมฺมณ\-ปจฺจเยน ปจฺจโย. อุปนิสฺสย\-ปจฺจเยน ปจฺจโย. กมฺมปจฺจเยน ปจฺจโย. (3)}

\csromanmarkh{2. ปจฺจย\-ปจฺจนีย}\csromanmarkhii{2. สงฺขฺยาวาร}\csromanheadernomark{6}{2. ปจฺจย\-ปจฺจนีย 2. สงฺขฺยาวาร}
\proseitemwithcloserruled{43}{นเหตุยา นว, นอารมฺมเณ นว, นอธิปติยา นว, นอนนฺตเร นว, นสมนนฺตเร นว, (สํขิตฺตํ. สพฺพตฺถ นว,) โนวิคเต นว, โนอวิคเต นว. (เอวํ คเณตพฺพํ.)}{ปจฺจนียํ.}
\proseitemwithcloserruled{44}{เหตุปจฺจยา นอารมฺมเณ ตีณิ, นอธิปติยา นอนนฺตเร, นสมนนฺตเร, นอุปนิสฺสเย, นปุเรชาเต, นปจฺฉาชาเต, นอาเสวเน, นกมฺเม, นวิปาเก ตีณิ, (สพฺพตฺถ ตีณิ. สํขิตฺตํ.) โนนตฺถิยา, โนวิคเต ตีณิ. (เอวํ คเณตพฺพํ.)}{อนุโลม\-ปจฺจนียํ.}
\proseitemwithcloser{45}{นเหตุปจฺจยา อารมฺมเณ นว, อธิปติยา สตฺต, อนนฺตเร สตฺต, สมนนฺตเร สตฺต, สหชาเต ตีณิ, อญฺญมญฺเญ ตีณิ, นิสฺสเย ตีณิ, อุปนิสฺสเย นว, อาเสวเน ตีณิ, กมฺเม นว, วิปาเก ตีณิ, อาหาเร, อินฺทฺริเย, ฌาเน, มคฺเค, สมฺปยุตฺเต, อตฺถิยา ตีณิ, นตฺถิยา สตฺต, วิคเต สตฺต, อวิคเต ตีณิ. (เอวํ คเณตพฺพํ.) ปจฺจนียา\-นุโลมํ. ปญฺหาวาโร.}{อตีตา\-รมฺมณ\-ตฺติกํ นิฏฺฐิตํ.}
\csromanmatikaanchor{mtk.222}
\csromantocmarknopage{subparagraph}{20. อชฺฌตฺตตฺติก 1. ปฏิจฺจวาร}{mtk.222}
\bookmark[dest={mtk.222},level=5]{20. อชฺฌตฺตตฺติก 1. ปฏิจฺจวาร}
\csromanmarkh{20. อชฺฌตฺตตฺติก}\csromanmarkhii{1. ปฏิจฺจวาร}\csromanheadernomark{6}{20. \textbf{อชฺฌตฺตตฺติก 1. ปฏิจฺจวาร}}
\csromanmarkh{1. ปจฺจยานุโลม}\csromanmarkhii{1. วิภงฺควาร}\csromanheadernomark{2}{1. \textbf{ปจฺจยานุโลม 1. วิภงฺควาร}}
\csromanmatikaanchor{mtk.223}
\csromanmatikaanchor{mtk.224}
\csromantocmarknopage{subparagraph}{1. ปจฺจยานุโลม 1. วิภงฺควาร}{mtk.223}
\bookmark[dest={mtk.223},level=5]{1. ปจฺจยานุโลม 1. วิภงฺควาร}
\csromantocmarkref{subparagraph}{เหตุ}{mtk.224}
\bookmark[dest={mtk.224},level=5]{เหตุ}
\csromanheader{6}{\textbf{เหตุ}}
\proseitem{1}{\csromanfolio{434}อชฺฌตฺตํ ธมฺมํ ปฏิจฺจ อชฺฌตฺโต ธมฺโม อุปฺปชฺชติ เหตุปจฺจยา. อชฺฌตฺตํ เอกํ ขนฺธํ ปฏิจฺจ ตโย ขนฺธา จิตฺต\-สมุฏฺฐานญฺจ รูปํ ฯเปฯ เทฺว ขนฺเธ ฯเปฯ. ปฏิสนฺธิ\-กฺขเณ อชฺฌตฺตํ เอกํ ขนฺธํ ปฏิจฺจ ตโย ขนฺธา กฏตฺตา จ รูปํ ฯเปฯ เทฺว ขนฺเธ ฯเปฯ. ขนฺเธ ปฏิจฺจ วตฺถุ, วตฺถุํ ปฏิจฺจ ขนฺธา. เอกํ มหาภูตํ ปฏิจฺจ ตโย มหาภูตา ฯเปฯ มหาภูเต ปฏิจฺจ จิตฺต\-สมุฏฺฐานํ รูปํ, กฏตฺตารูปํ, อุปาทารูปํ. (1)}

\prose{พหิทฺธา ธมฺมํ ปฏิจฺจ พหิทฺธา ธมฺโม อุปฺปชฺชติ เหตุปจฺจยา. พหิทฺธา เอกํ ขนฺธํ ปฏิจฺจ ตโย ขนฺธา จิตฺต\-สมุฏฺฐานญฺจ รูปํ ฯเปฯ เทฺว ขนฺเธ ฯเปฯ. ปฏิสนฺธิ\-กฺขเณ พหิทฺธา เอกํ ขนฺธํ ปฏิจฺจ ตโย ขนฺธา กฏตฺตา จ รูปํ ฯเปฯ เทฺว ขนฺเธ ปฏิจฺจ เทฺว ขนฺธา ฯเปฯ. ขนฺเธ ปฏิจฺจ วตฺถุ, วตฺถุํ ปฏิจฺจ ขนฺธา. เอกํ มหาภูตํ ปฏิจฺจ ตโย มหาภูตา ฯเปฯ มหาภูเต ปฏิจฺจ จิตฺต\-สมุฏฺฐานํ รูปํ, กฏตฺตารูปํ อุปาทารูปํ. (1)}

\csromanheader{2}{\textbf{อารมฺมณ}}
\proseitem{2}{อชฺฌตฺตํ ธมฺมํ ปฏิจฺจ อชฺฌตฺโต ธมฺโม อุปฺปชฺชติ อารมฺมณ\-ปจฺจยา. อชฺฌตฺตํ เอกํ ขนฺธํ ปฏิจฺจ ตโย ขนฺธา ฯเปฯ เทฺว ขนฺเธ ปฏิจฺจ ฯเปฯ. ปฏิสนฺธิ\-กฺขเณ อชฺฌตฺตํ เอกํ ขนฺธํ ปฏิจฺจ ตโย ขนฺธา ฯเปฯ. วตฺถุํ ปฏิจฺจ ขนฺธา. (1)}

\prose{พหิทฺธา ธมฺมํ ปฏิจฺจ พหิทฺธา ธมฺโม อุปฺปชฺชติ อารมฺมณ\-ปจฺจยา. พหิทฺธา เอกํ ขนฺธํ ปฏิจฺจ ตโย ขนฺธา ฯเปฯ เทฺว ขนฺเธ ฯเปฯ. ปฏิสนฺธิ\-กฺขเณ ฯเปฯ. วตฺถุํ ปฏิจฺจ ขนฺธา. (1)}

\proseitem{3}{อชฺฌตฺตํ ธมฺมํ ปฏิจฺจ อชฺฌตฺโต ธมฺโม อุปฺปชฺชติ อธิปติ\-ปจฺจยา. อชฺฌตฺตํ เอกํ ขนฺธํ ปฏิจฺจ ตโย ขนฺธา จิตฺต\-สมุฏฺฐานญฺจ รูปํ ฯเปฯ เทฺว ขนฺเธ ฯเปฯ. เอกํ มหาภูตํ ปฏิจฺจ ฯเปฯ มหาภูเต ปฏิจฺจ จิตฺต\-สมุฏฺฐานํ รูปํ, อุปาทารูปํ. (1)}

\prose{\csromanfolio{435}พหิทฺธา ธมฺมํ ปฏิจฺจ พหิทฺธา ธมฺโม อุปฺปชฺชติ อธิปติ\-ปจฺจยา. พหิทฺธา เอกํ ขนฺธํ ปฏิจฺจ ตโย ขนฺธา จิตฺต\-สมุฏฺฐานญฺจ รูปํ ฯเปฯ เทฺว ขนฺเธ ฯเปฯ. เอกํ มหาภูตํ ปฏิจฺจ ฯเปฯ มหาภูเต ปฏิจฺจ จิตฺต\-สมุฏฺฐานํ รูปํ อุปาทารูปํ. (1)}

\csromanheader{2}{\textbf{อนนฺตราทิ}}
\proseitem{4}{อชฺฌตฺตํ ธมฺมํ ปฏิจฺจ อชฺฌตฺโต ธมฺโม อุปฺปชฺชติ อนนฺตร\-ปจฺจยา. สมนนฺตร\-ปจฺจยา. สหชาต\-ปจฺจยา. อชฺฌตฺตํ เอกํ ขนฺธํ ปฏิจฺจ ตโย ขนฺธา จิตฺต\-สมุฏฺฐานญฺจ รูปํ ฯเปฯ เทฺว ขนฺเธ ฯเปฯ. ปฏิสนฺธิ\-กฺขเณ อชฺฌตฺตํ เอกํ ขนฺธํ ปฏิจฺจ ตโย ขนฺธา กฏตฺตา จ รูปํ ฯเปฯ เทฺว ขนฺเธ ฯเปฯ. ขนฺเธ ปฏิจฺจ วตฺถุ, วตฺถุํ ปฏิจฺจ ขนฺธา. เอกํ มหาภูตํ ปฏิจฺจ ตโย มหาภูตา ฯเปฯ มหาภูเต ปฏิจฺจ จิตฺต\-สมุฏฺฐานํ รูปํ, กฏตฺตารูปํ, อุปาทารูปํ. อาหารสมุฏฺฐานํ. อุตุสมุฏฺฐานํ. อสญฺญสตฺตานํ เอกํ มหาภูตํ ปฏิจฺจ ฯเปฯ. มหาภูเต ปฏิจฺจ กฏตฺตารูปํ, อุปาทารูปํ. (1)}

\prose{พหิทฺธา ธมฺมํ ปฏิจฺจ พหิทฺธา ธมฺโม อุปฺปชฺชติ สหชาต\-ปจฺจยา. พหิทฺธา เอกํ ขนฺธํ ปฏิจฺจ ตโย ขนฺธา จิตฺต\-สมุฏฺฐานญฺจ รูปํ ฯเปฯ เทฺว ขนฺเธ ฯเปฯ. ปฏิสนฺธิ\-กฺขเณ พหิทฺธา เอกํ ขนฺธํ ปฏิจฺจ ตโย ขนฺธา กฏตฺตา จ รูปํ ฯเปฯ เทฺว ขนฺเธ ฯเปฯ. ขนฺเธ ปฏิจฺจ วตฺถุ, วตฺถุํ ปฏิจฺจ ขนฺธา. เอกํ มหาภูตํ ปฏิจฺจ ฯเปฯ. มหาภูเต ปฏิจฺจ จิตฺต\-สมุฏฺฐานํ รูปํ, กฏตฺตารูปํ, อุปาทารูปํ. พาหิรํ. อาหารสมุฏฺฐานํ. อุตุสมุฏฺฐานํ. อสญฺญสตฺตานํ เอกํ มหาภูตํ ปฏิจฺจ ฯเปฯ. มหาภูเต ปฏิจฺจ กฏตฺตารูปํ อุปาทารูปํ. (1)}

\csromanheader{2}{\textbf{อญฺญมญฺญาทิ}}
\proseitem{5}{อชฺฌตฺตํ ธมฺมํ ปฏิจฺจ อชฺฌตฺโต ธมฺโม อุปฺปชฺชติ อญฺญมญฺญ\-ปจฺจยา. นิสฺสยปจฺจยา. อุปนิสฺสย\-ปจฺจยา. ปุเรชาต\-ปจฺจยา. อาเสวนปจฺจยา. (ปุเรชาเตปิ อาเสวเนปิ ปฏิสนฺธิ นตฺถิ.) กมฺมปจฺจยา. วิปากปจฺจยา. อาหารปจฺจยา. อินฺทฺริยปจฺจยา. ฌานปจฺจยา. มคฺคปจฺจยา. สมฺปยุตฺต\-ปจฺจยา. วิปฺปยุตฺต\-ปจฺจยา. อตฺถิปจฺจยา. นตฺถิปจฺจยา. วิคตปจฺจยา. อวิคตปจฺจยา.}

\csromanmatikaanchor{mtk.225}
\csromantocmarkref{subparagraph}{1. ปจฺจยานุโลม 2. สงฺขฺยาวาร}{mtk.225}
\bookmark[dest={mtk.225},level=5]{1. ปจฺจยานุโลม 2. สงฺขฺยาวาร}
\csromanmarkh{1. ปจฺจยานุโลม}\csromanmarkhii{2. สงฺขฺยาวาร}\csromanheadernomark{6}{1. \textbf{ปจฺจยานุโลม 2. สงฺขฺยาวาร}}
\vspace{\tipitakaparskip}
\csromancenter{เหตุยา เทฺว, อารมฺมเณ เทฺว ฯเปฯ อวิคเต เทฺว.}
\csromancenter{อนุโลมํ.}
\csromanmarkh{2. ปจฺจย\-ปจฺจนีย}\csromanmarkhii{1. วิภงฺควาร}\csromanheadernomark{2}{2. ปจฺจย\-ปจฺจนีย 1. วิภงฺควาร}
\csromanmatikaanchor{mtk.226}
\csromanmatikaanchor{mtk.227}
\csromantocmarknopage{subparagraph}{2. ปจฺจยปจฺจนีย 1. วิภงฺควาร}{mtk.226}
\bookmark[dest={mtk.226},level=5]{2. ปจฺจยปจฺจนีย 1. วิภงฺควาร}
\csromantocmarkref{subparagraph}{นเหตุ}{mtk.227}
\bookmark[dest={mtk.227},level=5]{นเหตุ}
\csromanheader{6}{\textbf{นเหตุ}}
\csromanmatikaanchor{mtk.228}
\csromantocmarkref{subparagraph}{นอารมฺมณาทิ}{mtk.228}
\bookmark[dest={mtk.228},level=5]{นอารมฺมณาทิ}
\proseitem{7}{\csromanfolio{436}อชฺฌตฺตํ ธมฺมํ ปฏิจฺจ อชฺฌตฺโต ธมฺโม อุปฺปชฺชติ นเหตุปจฺจยา. อเหตุกํ อชฺฌตฺตํ เอกํ ขนฺธํ ปฏิจฺจ ตโย ขนฺธา จิตฺต\-สมุฏฺฐานญฺจ รูปํ ฯเปฯ เทฺว ขนฺเธ ฯเปฯ. อเหตุก\-ปฏิสนฺธิ\-กฺขเณ ฯเปฯ. ขนฺเธ ปฏิจฺจ วตฺถุ, วตฺถุํ ปฏิจฺจ ขนฺธา. เอกํ มหาภูตํ ฯเปฯ. อาหารสมุฏฺฐานํ. อุตุสมุฏฺฐานํ. อสญฺญสตฺตานํ เอกํ มหาภูตํ ฯเปฯ มหาภูเต ปฏิจฺจ กฏตฺตารูปํ อุปาทารูปํ. วิจิกิจฺฉา\-สหคเต อุทฺธจฺจ\-สหคเต ขนฺเธ ปฏิจฺจ วิจิกิจฺฉา\-สหคโต อุทฺธจฺจ\-สหคโต โมโห. (1)}

\prose{พหิทฺธา ธมฺมํ ปฏิจฺจ พหิทฺธา ธมฺโม อุปฺปชฺชติ นเหตุปจฺจยา. อเหตุกํ พหิทฺธา เอกํ ขนฺธํ ปฏิจฺจ ตโย ขนฺธา จิตฺต\-สมุฏฺฐานญฺจ รูปํ ฯเปฯ เทฺว ขนฺเธ ฯเปฯ. อเหตุก\-ปฏิสนฺธิ\-กฺขเณ ฯเปฯ. ขนฺเธ ปฏิจฺจ วตฺถุ, วตฺถุํ ปฏิจฺจ ขนฺธา. เอกํ มหาภูตํ ปฏิจฺจ ฯเปฯ. พาหิรํ. อาหารสมุฏฺฐานํ. อุตุสมุฏฺฐานํ. อสญฺญสตฺตานํ เอกํ มหาภูตํ ปฏิจฺจ ฯเปฯ. วิจิกิจฺฉา\-สหคเต อุทฺธจฺจ\-สหคเต ขนฺเธ ปฏิจฺจ วิจิกิจฺฉา\-สหคโต อุทฺธจฺจ\-สหคโต โมโห. (1)}

\prose{นอารมฺมณ}

\proseitem{8}{อชฺฌตฺตํ ธมฺมํ ปฏิจฺจ อชฺฌตฺโต ธมฺโม อุปฺปชฺชติ นอารมฺมณ\-ปจฺจยา. อชฺฌตฺเต ขนฺเธ ปฏิจฺจ จิตฺต\-สมุฏฺฐานํ รูปํ. ปฏิสนฺธิ\-กฺขเณ อชฺฌตฺเต ขนฺเธ ปฏิจฺจ กฏตฺตารูปํ. ขนฺเธ ปฏิจฺจ วตฺถุ ฯเปฯ. เอกํ มหาภูตํ ปฏิจฺจ ฯเปฯ. อาหารสมุฏฺฐานํ. อุตุสมุฏฺฐานํ. อสญฺญสตฺตานํ เอกํ มหาภูตํ ปฏิจฺจ ฯเปฯ. (1)}

\prose{พหิทฺธา ธมฺมํ ปฏิจฺจ พหิทฺธา ธมฺโม อุปฺปชฺชติ นอารมฺมณ\-ปจฺจยา. พหิทฺธา ขนฺเธ ปฏิจฺจ จิตฺต\-สมุฏฺฐานํ รูปํ. ปฏิสนฺธิ\-กฺขเณ พหิทฺธา ขนฺเธ ปฏิจฺจ กฏตฺตารูปํ. ขนฺเธ ปฏิจฺจ วตฺถุ ฯเปฯ. เอกํ มหาภูตํ ฯเปฯ. พาหิรํ. อาหารสมุฏฺฐานํ. อุตุสมุฏฺฐานํ. อสญฺญสตฺตานํ ฯเปฯ. (1)}

\prose{นอธิปตฺยาทิ}

\proseitem{9}{อชฺฌตฺตํ ธมฺมํ ปฏิจฺจ อชฺฌตฺโต ธมฺโม อุปฺปชฺชติ นอธิปติ\-ปจฺจยา ฯเปฯ. (อนุโลม\-สหชาต\-สทิสํ, นินฺนานากรณํ.) นอนนฺตร\-ปจฺจยา.}

\prose{\csromanfolio{437}นสมนนฺตร\-ปจฺจยา. นอญฺญมญฺญ\-ปจฺจยา. นฺเอาปนิสฺสย\-ปจฺจยา. นปุเรชาต\-ปจฺจยา. อรูเป อชฺฌตฺตํ เอกํ ขนฺธํ ปฏิจฺจ ฯเปฯ. อชฺฌตฺเต ขนฺเธ ปฏิจฺจ จิตฺต\-สมุฏฺฐานํ รูปํ. ปฏิสนฺธิ\-กฺขเณ อชฺฌตฺตํ เอกํ ขนฺธํ ปฏิจฺจ ตโย ขนฺธา กฏตฺตา จ รูปํ ฯเปฯ เทฺว ขนฺเธ ฯเปฯ. ขนฺเธ ปฏิจฺจ วตฺถุ, วตฺถุํ ปฏิจฺจ ขนฺธา. เอกํ มหาภูตํ ปฏิจฺจ ฯเปฯ. อาหารสมุฏฺฐานํ. อุตุสมุฏฺฐานํ. อสญฺญสตฺตานํ ฯเปฯ.}

\proseitem{10}{พหิทฺธา ธมฺมํ ปฏิจฺจ พหิทฺธา ธมฺโม อุปฺปชฺชติ นปุเรชาต\-ปจฺจยา. อรูเป พหิทฺธา เอกํ ขนฺธํ ปฏิจฺจ ฯเปฯ เทฺว ขนฺเธ ฯเปฯ. พหิทฺธา ขนฺเธ ปฏิจฺจ จิตฺต\-สมุฏฺฐานํ รูปํ. ปฏิสนฺธิ\-กฺขเณ (ปริปุณฺณํ.) เอกํ มหาภูตํ ฯเปฯ. พาหิรํ. อาหารสมุฏฺฐานํ. อุตุสมุฏฺฐานํ. อสญฺญสตฺตานํ ฯเปฯ.}

\csromanheader{2}{\textbf{นปจฺฉาชาตาทิ}}
\proseitem{11}{อชฺฌตฺตํ ธมฺมํ ปฏิจฺจ อชฺฌตฺโต ธมฺโม อุปฺปชฺชติ นปจฺฉาชาต\-ปจฺจยา. นอาเสวน\-ปจฺจยา. นกมฺมปจฺจยา. อชฺฌตฺเต ขนฺเธ ปฏิจฺจ อชฺฌตฺตา เจตนา. อาหารสมุฏฺฐานํ. อุตุสมุฏฺฐานํ ฯเปฯ. (1)}

\prose{พหิทฺธา ธมฺมํ ปฏิจฺจ พหิทฺธา ธมฺโม อุปฺปชฺชติ นกมฺมปจฺจยา. พหิทฺธา ขนฺเธ ปฏิจฺจ พหิทฺธา เจตนา. พาหิรํ. อาหารสมุฏฺฐานํ. อุตุสมุฏฺฐานํ. (1)}

\csromanheader{2}{\textbf{นวิปากาทิ}}
\proseitem{12}{อชฺฌตฺตํ ธมฺมํ ปฏิจฺจ อชฺฌตฺโต ธมฺโม อุปฺปชฺชติ นวิปาก\-ปจฺจยา. (ปฏิสนฺธิ นตฺถิ.) นอาหารปจฺจยา. อุตุสมุฏฺฐานํ. อสญฺญสตฺตานํ ฯเปฯ. (1)}

\prose{พหิทฺธา ธมฺมํ ปฏิจฺจ พหิทฺธา ธมฺโม อุปฺปชฺชติ นอาหารปจฺจยา. พาหิรํ. อุตุสมุฏฺฐานํ. อสญฺญสตฺตานํ ฯเปฯ. (1)}

\prose{อชฺฌตฺตํ ธมฺมํ ปฏิจฺจ อชฺฌตฺโต ธมฺโม อุปฺปชฺชติ นอินฺทฺริยปจฺจยา. อาหารสมุฏฺฐานํ. อุตุสมุฏฺฐานํ. อสญฺญสตฺตานํ มหาภูเต ปฏิจฺจ รูปชีวิตินฺทฺริยํ. (1)}

\prose{พหิทฺธา ธมฺมํ ปฏิจฺจ พหิทฺธา ธมฺโม อุปฺปชฺชติ นอินฺทฺริยปจฺจยา. พาหิรํ. อาหารสมุฏฺฐานํ. อุตุสมุฏฺฐานํ. อสญฺญสตฺตานํ มหาภูเต ปฏิจฺจ รูปชีวิตินฺทฺริยํ. (1)}

\csromanheader{2}{\textbf{นฌาน}}
\csromanmatikaanchor{mtk.229}
\csromanmatikaanchor{mtk.230}
\csromantocmarkref{subparagraph}{2. ปจฺจยปจฺจนีย 2. สงฺขฺยาวาร}{mtk.229}
\bookmark[dest={mtk.229},level=5]{2. ปจฺจยปจฺจนีย 2. สงฺขฺยาวาร}
\csromantocmarkref{subparagraph}{3. ปจฺจยานุโลมปจฺจนีย ...}{mtk.230}
\bookmark[dest={mtk.230},level=5]{3. ปจฺจยานุโลมปจฺจนีย ...}
\proseitem{13}{\csromanfolio{438}อชฺฌตฺตํ ธมฺมํ ปฏิจฺจ อชฺฌตฺโต ธมฺโม อุปฺปชฺชติ นฌานปจฺจยา. ปญฺจ\-วิญฺญาณสหคตํ ฯเปฯ. อาหารสมุฏฺฐานํ. อุตุสมุฏฺฐานํ. อสญฺญสตฺตานํ ฯเปฯ. (1)}

\prose{พหิทฺธา ธมฺมํ ปฏิจฺจ พหิทฺธา ธมฺโม อุปฺปชฺชติ นฌานปจฺจยา. ปญฺจ\-วิญฺญาณสหคตํ ฯเปฯ. พาหิรํ. อาหารสมุฏฺฐานํ. อุตุสมุฏฺฐานํ. อสญฺญสตฺตานํ ฯเปฯ. (1)}

\csromanheader{2}{\textbf{นมคฺคาทิ}}
\proseitem{14}{อชฺฌตฺตํ ธมฺมํ ปฏิจฺจ อชฺฌตฺโต ธมฺโม อุปฺปชฺชติ นมคฺคปจฺจยา. (นเหตุสทิสํ. โมโห นตฺถิ.) นสมฺปยุตฺต\-ปจฺจยา. นวิปฺปยุตฺต\-ปจฺจยา. อรูเป ฯเปฯ. อาหารสมุฏฺฐานํ. อุตุสมุฏฺฐานํ. อสญฺญสตฺตานํ ฯเปฯ. (1)}

\prose{พหิทฺธา ธมฺมํ ปฏิจฺจ พหิทฺธา ธมฺโม อุปฺปชฺชติ นวิปฺปยุตฺต\-ปจฺจยา. อรูเป ฯเปฯ. พาหิรํ. อาหารสมุฏฺฐานํ. อุตุสมุฏฺฐานํ. อสญฺญสตฺตานํ ฯเปฯ. โนนตฺถิปจฺจยา. โนวิคต\-ปจฺจยา.}

\csromanmarkh{2. ปจฺจย\-ปจฺจนีย}\csromanmarkhii{2. สงฺขฺยาวาร}\csromanheadernomark{6}{2. ปจฺจย\-ปจฺจนีย 2. สงฺขฺยาวาร}
\proseitemwithcloserruled{15}{นเหตุยา เทฺว, นอารมฺมเณ เทฺว, นอธิปติยา เทฺว, นอนนฺตเร เทฺว, นสมนนฺตเร เทฺว, (สํขิตฺตํ. สพฺพตฺถ เทฺว.) โนวิคเต เทฺว. (เอวํ คเณตพฺพํ.)}{ปจฺจนียํ.}
\proseitem{16}{เหตุปจฺจยา นอารมฺมเณ เทฺว ฯเปฯ นวิปาเก, นสมฺปยุตฺเต, นวิปฺปยุตฺเต, โนนตฺถิยา, โนวิคเต เทฺว. (เอวํ คเณตพฺพํ.)}

\vspace{\tipitakaparskip}
\csromancenter{อนุโลม\-ปจฺจนียํ.}
\csromanmatikaanchor{mtk.231}
\csromantocmarkref{subparagraph}{4. ปจฺจยปจฺจนียานุโลม ...}{mtk.231}
\bookmark[dest={mtk.231},level=5]{4. ปจฺจยปจฺจนียานุโลม ...}
\csromanmarkh{20. อชฺฌตฺตตฺติก}\csromanmarkhii{4. ปจฺจย\-ปจฺจนียา\-นุโลม}\csromanheadernomark{6}{20. อชฺฌตฺตตฺติก 4. ปจฺจย\-ปจฺจนียา\-นุโลม}
\proseitem{17}{\csromanfolio{439}นเหตุปจฺจยา อารมฺมเณ เทฺว. อนนฺตเร เทฺว ฯเปฯ มคฺเค เทฺว ฯเปฯ อวิคเต เทฺว. (เอวํ คเณตพฺพํ.) ปจฺจนียา\-นุโลมํ. ปฏิจฺจวาโร. (สหชาตวาโร ปฏิจฺจ\-วาร\-สทิโส.)}
\csromansectionrule

\csromanmatikaanchor{mtk.232}
\csromantocmarknopage{subparagraph}{20. อชฺฌตฺตตฺติก 3. ปจฺจยวาร}{mtk.232}
\bookmark[dest={mtk.232},level=5]{20. อชฺฌตฺตตฺติก 3. ปจฺจยวาร}
\csromanmarkh{20. อชฺฌตฺตตฺติก}\csromanmarkhii{3. ปจฺจยวาร}\csromanheadernomark{6}{20. อชฺฌตฺตตฺติก 3. ปจฺจยวาร}
\csromanmatikaanchor{mtk.233}
\csromantocmarkref{subparagraph}{1. ปจฺจยานุโลม 1. วิภงฺควาร}{mtk.233}
\bookmark[dest={mtk.233},level=5]{1. ปจฺจยานุโลม 1. วิภงฺควาร}
\csromanmarkh{1. ปจฺจยานุโลม}\csromanmarkhii{1. วิภงฺควาร}\csromanheadernomark{6}{1. \textbf{ปจฺจยานุโลม 1. วิภงฺควาร}}
\csromanheader{2}{\textbf{เหตุ}}
\proseitem{18}{อชฺฌตฺตํ ธมฺมํ ปจฺจยา อชฺฌตฺโต ธมฺโม อุปฺปชฺชติ เหตุปจฺจยา. อชฺฌตฺตํ เอกํ ขนฺธํ ปจฺจยา ตโย ขนฺธา จิตฺต\-สมุฏฺฐานญฺจ รูปํ ฯเปฯ เทฺว ขนฺเธ ฯเปฯ ปฏิสนฺธิ\-กฺขเณ (ปริปุณฺณํ.) เอกํ มหาภูตํ ฯเปฯ. วตฺถุํ ปจฺจยา อชฺฌตฺตา ขนฺธา. (1)}

\prose{พหิทฺธา ธมฺมํ ปจฺจยา พหิทฺธา ธมฺโม อุปฺปชฺชติ เหตุปจฺจยา. พหิทฺธา เอกํ ขนฺธํ ปจฺจยา ตโย ขนฺธา ฯเปฯ. ปฏิสนฺธิ\-กฺขเณ ฯเปฯ. เอกํ มหาภูตํ ฯเปฯ. วตฺถุํ ปจฺจยา พหิทฺธา ขนฺธา. (1)}

\csromanheader{2}{\textbf{อารมฺมณ}}
\proseitem{19}{อชฺฌตฺตํ ธมฺมํ ปจฺจยา อชฺฌตฺโต ธมฺโม อุปฺปชฺชติ อารมฺมณ\-ปจฺจยา. (ปฏิจฺจ\-วาร\-สทิสํ.) จกฺขายตนํ ปจฺจยา จกฺขุวิญฺญาณํ ฯเปฯ. กายายตนํ ปจฺจยา กายวิญฺญาณํ ฯเปฯ. วตฺถุํ ปจฺจยา อชฺฌตฺตา ขนฺธา. (1)}

\prose{พหิทฺธา ธมฺมํ ปจฺจยา พหิทฺธา ธมฺโม อุปฺปชฺชติ อารมฺมณ\-ปจฺจยา. (ปฏิจฺจ\-วาร\-สทิสํ.) จกฺขายตนํ ฯเปฯ. กายายตนํ ฯเปฯ. วตฺถุํ ปจฺจยา พหิทฺธา ขนฺธา. (1)}

\csromanheader{2}{\textbf{อธิปตฺยาทิ}}
\proseitem{20}{\csromanfolio{440}อชฺฌตฺตํ ธมฺมํ ปจฺจยา อชฺฌตฺโต ธมฺโม อุปฺปชฺชติ อธิปติ\-ปจฺจยา. (วตฺถุํ อติเรกํ, ปฏิจฺจ\-วาร\-สทิสํ.) อนนฺตร\-ปจฺจยา. สมนนฺตร\-ปจฺจยา. สหชาต\-ปจฺจยา. (สหชาตวาเร ปริปุณฺณา.) มหาภูเต ปจฺจยา ฯเปฯ. (มหาภูตานํ ขนฺธานญฺจ ปจฺฉา ปญฺจายตนานิ จ วตฺถุ จ กาตพฺพา.) อญฺญมญฺญ\-ปจฺจยา. นิสฺสยปจฺจยา ฯเปฯ. อวิคตปจฺจยา.}

\csromanmatikaanchor{mtk.234}
\csromantocmarkref{subparagraph}{1. ปจฺจยานุโลม 2. สงฺขฺยาวาร}{mtk.234}
\bookmark[dest={mtk.234},level=5]{1. ปจฺจยานุโลม 2. สงฺขฺยาวาร}
\csromanmarkh{1. ปจฺจยานุโลม}\csromanmarkhii{2. สงฺขฺยาวาร}\csromanheadernomark{6}{1. \textbf{ปจฺจยานุโลม 2. สงฺขฺยาวาร}}
\vspace{\tipitakaparskip}
\csromancenter{เหตุยา เทฺว, อารมฺมเณ ฯเปฯ อวิคเต เทฺว.}
\nitthitamruled{\textbf{อนุโลมํ.}}
\csromanmatikaanchor{mtk.235}
\csromantocmarkref{subparagraph}{2. ปจฺจยปจฺจนีย 1. วิภงฺควาร}{mtk.235}
\bookmark[dest={mtk.235},level=5]{2. ปจฺจยปจฺจนีย 1. วิภงฺควาร}
\csromanmarkh{2. ปจฺจย\-ปจฺจนีย}\csromanmarkhii{1. วิภงฺควาร}\csromanheadernomark{6}{2. ปจฺจย\-ปจฺจนีย 1. วิภงฺควาร}
\csromanheader{2}{\textbf{นเหตุ}}
\proseitem{22}{อชฺฌตฺตํ ธมฺมํ ปจฺจยา อชฺฌตฺโต ธมฺโม อุปฺปชฺชติ นเหตุปจฺจยา. อเหตุกํ อชฺฌตฺตํ เอกํ ขนฺธํ ปจฺจยา ตโย ขนฺธา ฯเปฯ. อเหตุก\-ปฏิสนฺธิ\-กฺขเณ ฯเปฯ. ขนฺเธ ปจฺจยา วตฺถุ, วตฺถุํ ปจฺจยา ขนฺธา. เอกํ มหาภูตํ ฯเปฯ. อาหารสมุฏฺฐานํ. อุตุสมุฏฺฐานํ. อสญฺญสตฺตานํ ฯเปฯ. จกฺขายตนํ ฯเปฯ. กายายตนํ ฯเปฯ. วตฺถุํ ปจฺจยา อเหตุกา อชฺฌตฺตา ขนฺธา. วิจิกิจฺฉา\-สหคเต อุทฺธจฺจ\-สหคเต ขนฺเธ จ วตฺถุญฺจ ปจฺจยา วิจิกิจฺฉา\-สหคโต อุทฺธจฺจ\-สหคโต โมโห. (1)}

\prose{พหิทฺธา ธมฺมํ ปจฺจยา พหิทฺธา ธมฺโม อุปฺปชฺชติ นเหตุปจฺจยา. (ปวตฺติปฏิสนฺธิปิ มหาภูตาปิ กาตพฺพา.) จกฺขายตนํ ฯเปฯ. กายายตนํ ฯเปฯ. วตฺถุํ ปจฺจยา อเหตุกา พหิทฺธา ขนฺธา. วิจิกิจฺฉา\-สหคเต อุทฺธจฺจ\-สหคเต ขนฺเธ จ วตฺถุญฺจ ปจฺจยา วิจิกิจฺฉา\-สหคโต อุทฺธจฺจ\-สหคโต โมโห. (1)}

\prose{นอารมฺมณาทิ}

\proseitem{23}{นอารมฺมณ\-ปจฺจยา. นอธิปติ\-ปจฺจยา. (สหชาต\-สทิสํ.) นอนนฺตร\-ปจฺจยา. นสมนนฺตร\-ปจฺจยา. นอญฺญมญฺญ\-ปจฺจยา. นอุปนิสฺสยปจฺจยา. นปุเรชาต\-ปจฺจยา. (ปฏิจฺจ\-วาร\-สทิสํ.) นปจฺฉาชาต\-ปจฺจยา.}

\csromanmatikaanchor{mtk.236}
\csromanmatikaanchor{mtk.238}
\csromantocmarkref{subparagraph}{2. ปจฺจยปจฺจนีย 2. สงฺขฺยาวาร}{mtk.236}
\bookmark[dest={mtk.236},level=5]{2. ปจฺจยปจฺจนีย 2. สงฺขฺยาวาร}
\prose{\csromanfolio{441}นอาเสวน\-ปจฺจยา. นกมฺมปจฺจยา ฯเปฯ. นวิปฺปยุตฺต\-ปจฺจยา. (ปฏิจฺจ\-วาร\-ปจฺจนีเย วิปฺปยุตฺต\-สทิสํ.) โนนตฺถิปจฺจยา. โนวิคต\-ปจฺจยา.}

\csromanmarkh{2. ปจฺจย\-ปจฺจนีย}\csromanmarkhii{2. สงฺขฺยาวาร}\csromanheadernomark{6}{2. ปจฺจย\-ปจฺจนีย 2. สงฺขฺยาวาร}
\vspace{\tipitakaparskip}
\csromancenter{นเหตุยา เทฺว, นอารมฺมเณ เทฺว ฯเปฯ โนวิคเต เทฺว.}
\nitthitamruled{\textbf{ปจฺจนียํ.}}
\csromancenter{\textbf{เหตุ}ปจฺจยา นอารมฺมเณ เทฺว, นอธิปติยา เทฺว ฯเปฯ นวิปาเก, นสมฺปยุตฺเต, นวิปฺปยุตฺเต, โนนตฺถิยา, โนวิคเต เทฺว.}
\nitthitamruled{อนุโลม\-ปจฺจนียํ.}
\proseitem{26}{นเหตุปจฺจยา อารมฺมเณ เทฺว ฯเปฯ อวิคเต เทฺว. ปจฺจนียา\-นุโลมํ.}

\prose{(นิสฺสยวาโร ปจฺจยวาร\-สทิโส. สํสฏฺฐ\-วาโรปิ สมฺปยุตฺตวาโรปิ วิตฺถาเรตพฺโพ.)}
\csromansectionrule

\csromanmatikaanchor{mtk.237}
\csromanmatikaanchor{mtk.239}
\csromantocmarkref{subparagraph}{3. ปจฺจยานุโลมปจฺจนีย}{mtk.237}
\bookmark[dest={mtk.237},level=5]{3. ปจฺจยานุโลมปจฺจนีย}
\csromantocmarkref{subparagraph}{4. ปจฺจยปจฺจนียานุโลม}{mtk.238}
\bookmark[dest={mtk.238},level=5]{4. ปจฺจยปจฺจนียานุโลม}
\csromantocmarknopage{subparagraph}{20. อชฺฌตฺตตฺติก 7. ปญฺหาวาร}{mtk.239}
\bookmark[dest={mtk.239},level=5]{20. อชฺฌตฺตตฺติก 7. ปญฺหาวาร}
\csromanmarkh{20. อชฺฌตฺตตฺติก}\csromanmarkhii{7. ปญฺหาวาร}\csromanheadernomark{6}{20. อชฺฌตฺตตฺติก 7. ปญฺหาวาร}
\csromanmarkh{1. ปจฺจยานุโลม}\csromanmarkhii{1. วิภงฺควาร}\csromanheadernomark{6}{1. \textbf{ปจฺจยานุโลม 1. วิภงฺควาร}}
\csromanmatikaanchor{mtk.240}
\csromanmatikaanchor{mtk.241}
\csromantocmarknopage{subparagraph}{1. ปจฺจยานุโลม 1. วิภงฺควาร}{mtk.240}
\bookmark[dest={mtk.240},level=5]{1. ปจฺจยานุโลม 1. วิภงฺควาร}
\csromantocmarkref{subparagraph}{เหตุ}{mtk.241}
\bookmark[dest={mtk.241},level=5]{เหตุ}
\csromanheader{6}{\textbf{เหตุ}}
\proseitem{27}{อชฺฌตฺโต ธมฺโม อชฺฌตฺตสฺส ธมฺมสฺส เหตุปจฺจเยน ปจฺจโย. อชฺฌตฺตา เหตู สมฺปยุตฺตกานํ ขนฺธานํ จิตฺตสมุฏฺฐานานญฺจ รูปานํ เหตุปจฺจเยน ปจฺจโย. ปฏิสนฺธิ\-กฺขเณ ฯเปฯ. (1)}

\prose{\csromanfolio{442}พหิทฺธา ธมฺโม พหิทฺธา ธมฺมสฺส เหตุปจฺจเยน ปจฺจโย. พหิทฺธา เหตู สมฺปยุตฺตกานํ ขนฺธานํ จิตฺตสมุฏฺฐานานญฺจ รูปานํ เหตุปจฺจเยน ปจฺจโย. ปฏิสนฺธิ\-กฺขเณ ฯเปฯ. (1)}

\csromanmatikaanchor{mtk.242}
\csromantocmarkref{subparagraph}{อารมฺมณ}{mtk.242}
\bookmark[dest={mtk.242},level=5]{อารมฺมณ}
\csromanheader{6}{\textbf{อารมฺมณ}}
\proseitem{28}{อชฺฌตฺโต ธมฺโม อชฺฌตฺตสฺส ธมฺมสฺส อารมฺมณ\-ปจฺจเยน ปจฺจโย. ทานํ ทตฺวา, สีลํ สมาทิยิตฺวา, อุโปสถกมฺมํ กตฺวา ตํ ปจฺจเวกฺขติ. ปุพฺเพ สุจิณฺณานิ ฯเปฯ. ฌานา วุฏฺฐหิตฺวา ฯเปฯ. อริยา มคฺคา วุฏฺฐหิตฺวา มคฺคํ ปจฺจเวกฺขนฺติ, ผลํ ปจฺจเวกฺขนฺติ, ปหีเน กิเลเส ปจฺจเวกฺขนฺติ, วิกฺขมฺภิเต กิเลเส ปจฺจเวกฺขนฺติ, ปุพฺเพ สมุทาจิณฺเณ กิเลเส ชานนฺติ. อชฺฌตฺตํ จกฺขุํ ฯเปฯ กายํ. รูเป ฯเปฯ โผฏฺฐพฺเพ. วตฺถุํ. อชฺฌตฺเต ขนฺเธ อนิจฺจโต ทุกฺขโต อนตฺตโต วิปสฺสติ, อสฺสาเทติ อภินนฺทติ, ตํ อารพฺภ ราโค อุปฺปชฺชติ ฯเปฯ โทมนสฺสํ อุปฺปชฺชติ. ทิพฺเพน จกฺขุนา รูปํ ปสฺสติ. ทิพฺพาย โสตธาตุยา สทฺทํ สุณาติ. อากาสา\-นญฺจา\-ยตนํ วิญฺญาณญฺจา\-ยตนสฺส อารมฺมณ\-ปจฺจเยน ปจฺจโย. อากิญฺจญฺญา\-ยตนํ เนว\-สญฺญา\-นา\-สญฺญา\-ยตนสฺส อารมฺมณ\-ปจฺจเยน ปจฺจโย. รูปายตนํ ฯเปฯ. โผฏฺฐพฺพา\-ยตนํ กายวิญฺญาณสฺส อารมฺมณ\-ปจฺจเยน ปจฺจโย. อชฺฌตฺตา ขนฺธา อิทฺธิวิธ\-ญาณสฺส, ปุพฺเพ\-นิวาสา\-นุสฺสติ\-ญาณสฺส, ยถากมฺมูปคญาณสฺส, อนาคตํสญาณสฺส, อาวชฺชนาย อารมฺมณ\-ปจฺจเยน ปจฺจโย. (1)}

\prose{อชฺฌตฺโต ธมฺโม พหิทฺธา ธมฺมสฺส อารมฺมณ\-ปจฺจเยน ปจฺจโย. ปโร อชฺฌตฺตํ จกฺขุํ ฯเปฯ วตฺถุํ. อชฺฌตฺเต ขนฺเธ อนิจฺจโต ทุกฺขโต อนตฺตโต วิปสฺสติ, อสฺสาเทติ อภินนฺทติ, ตํ อารพฺภ ราโค อุปฺปชฺชติ ฯเปฯ โทมนสฺสํ อุปฺปชฺชติ. ทิพฺเพน จกฺขุนา รูปํ ปสฺสติ. ทิพฺพาย โสตธาตุยา สทฺทํ สุณาติ. เจโตปริย\-ญาเณน อชฺฌตฺต\-จิตฺต\-สมงฺคิสฺส จิตฺตํ ชานาติ. อชฺฌตฺตํ รูปายตนํ พหิทฺธา จกฺขุ\-วิญฺญาณสฺส ฯเปฯ. อชฺฌตฺตํ โผฏฺฐพฺพา\-ยตนํ พหิทฺธา กายวิญฺญาณสฺส อารมฺมณ\-ปจฺจเยน ปจฺจโย. อชฺฌตฺตา ขนฺธา อิทฺธิวิธ\-ญาณสฺส, เจโตปริย\-ญาณสฺส, ปุพฺเพ\-นิวาสา\-นุสฺสติ\-ญาณสฺส, ยถากมฺมูปคญาณสฺส, อนาคตํสญาณสฺส, อาวชฺชนาย อารมฺมณ\-ปจฺจเยน ปจฺจโย. (2)}

\csromanmatikaanchor{mtk.243}
\csromantocmarkref{subparagraph}{อธิปติ}{mtk.243}
\bookmark[dest={mtk.243},level=5]{อธิปติ}
\proseitem{29}{พหิทฺธา ธมฺโม พหิทฺธา ธมฺมสฺส อารมฺมณ\-ปจฺจเยน ปจฺจโย. ปโร ทานํ ทตฺวา, สีลํ สมาทิยิตฺวา, อุโปสถกมฺมํ กตฺวา ตํ ปจฺจเวกฺขติ, \csromanfolio{443}ปุพฺเพ สุจิณฺณานิ ปจฺจเวกฺขติ, ฌานา วุฏฺฐหิตฺวา ฯเปฯ. อริยา มคฺคา วุฏฺฐหิตฺวา มคฺคํ ปจฺจเวกฺขนฺติ. ผลํ ปจฺจเวกฺขนฺติ. นิพฺพานํ ปจฺจเวกฺขนฺติ. นิพฺพานํ โคตฺรภุสฺส, โวทานสฺส, มคฺคสฺส, ผลสฺส, อาวชฺชนาย อารมฺมณ\-ปจฺจเยน ปจฺจโย. อริยา ปหีเน กิเลเส ปจฺจเวกฺขนฺติ. วิกฺขมฺภิเต กิเลเส ปจฺจเวกฺขนฺติ. ปุพฺเพ สมุทาจิณฺเณ ฯเปฯ. ปโร พหิทฺธา จกฺขุํ ฯเปฯ วตฺถุํ. พหิทฺธา ขนฺเธ อนิจฺจโต ฯเปฯ โทมนสฺสํ อุปฺปชฺชติ. ทิพฺเพน จกฺขุนา รูปํ ปสฺสติ. ทิพฺพาย โสตธาตุยา สทฺทํ สุณาติ. เจโตปริย\-ญาเณน พหิทฺธา จิตฺต\-สมงฺคิสฺส ฯเปฯ. อากาสา\-นญฺจา\-ยตนํ วิญฺญาณญฺจา\-ยตนสฺส อารมฺมณ\-ปจฺจเยน ปจฺจโย. อากิญฺจญฺญา\-ยตนํ เนว\-สญฺญา\-นา\-สญฺญา\-ยตนสฺส อารมฺมณ\-ปจฺจเยน ปจฺจโย. พหิทฺธา รูปายตนํ พหิทฺธา จกฺขุ\-วิญฺญาณสฺส ฯเปฯ. พหิทฺธา โผฏฺฐพฺพา\-ยตนํ พหิทฺธา กายวิญฺญาณสฺส ฯเปฯ. พหิทฺธา ขนฺธา อิทฺธิวิธ\-ญาณสฺส, เจโตปริย\-ญาณสฺส, ปุพฺเพ\-นิวาสา\-นุสฺสติ\-ญาณสฺส, ยถากมฺมูปคญาณสฺส, อนาคตํสญาณสฺส, อาวชฺชนาย อารมฺมณ\-ปจฺจเยน ปจฺจโย. (1)}

\prose{พหิทฺธา ธมฺโม อชฺฌตฺตสฺส ธมฺมสฺส อารมฺมณ\-ปจฺจเยน ปจฺจโย. อริยา นิพฺพานํ ปจฺจเวกฺขนฺติ. นิพฺพานํ โคตฺรภุสฺส, โวทานสฺส, มคฺคสฺส, ผลสฺส, อาวชฺชนาย อารมฺมณ\-ปจฺจเยน ปจฺจโย. พหิทฺธา จกฺขุํ ฯเปฯ วตฺถุํ. พหิทฺธา ขนฺเธ อนิจฺจโต ฯเปฯ โทมนสฺสํ อุปฺปชฺชติ. ทิพฺเพน จกฺขุนา รูปํ ปสฺสติ. ทิพฺพาย โสตธาตุยา สทฺทํ สุณาติ. เจโตปริย\-ญาเณน พหิทฺธา จิตฺต\-สมงฺคิสฺส จิตฺตํ ชานาติ. พหิทฺธา รูปายตนํ อชฺฌตฺตํ จกฺขุ\-วิญฺญาณสฺส ฯเปฯ. พหิทฺธา โผฏฺฐพฺพา\-ยตนํ อชฺฌตฺตํ กายวิญฺญาณสฺส ฯเปฯ. พหิทฺธา ขนฺธา อิทฺธิวิธ\-ญาณสฺส, เจโตปริย\-ญาณสฺส, ปุพฺเพ\-นิวาสา\-นุสฺสติ\-ญาณสฺส, ยถากมฺมูปคญาณสฺส, อนาคตํสญาณสฺส, อาวชฺชนาย อารมฺมณ\-ปจฺจเยน ปจฺจโย. (2)}

\proseitem{30}{อชฺฌตฺโต ธมฺโม อชฺฌตฺตสฺส ธมฺมสฺส อธิปติ\-ปจฺจเยน ปจฺจโย. อารมฺมณา\-ธิปติ, สหชาตา\-ธิปติ.  อารมฺมณา\-ธิปติ\csromandash{}ทานํ ทตฺวา, สีลํ สมาทิยิตฺวา, อุโปสถกมฺมํ กตฺวา, ตํ ครุํ กตฺวา ปจฺจเวกฺขติ. ปุพฺเพ สุจิณฺณานิ ครุํ กตฺวา ฯเปฯ. ฌานา วุฏฺฐหิตฺวา ฌานํ ครุํ กตฺวา ฯเปฯ. อริยา มคฺคา วุฏฺฐหิตฺวา มคฺคํ ครุํ กตฺวา ฯเปฯ. ผลํ ครุํ กตฺวา ปจฺจเวกฺขนฺติ. อชฺฌตฺตํ จกฺขุํ ฯเปฯ. วตฺถุํ. อชฺฌตฺเต ขนฺเธ ครุํ กตฺวา อสฺสาเทติ อภินนฺทติ, \csromanfolio{444}ตํ ครุํ กตฺวา ราโค อุปฺปชฺชติ, ทิฏฺฐิ อุปฺปชฺชติ. สหชาตา\-ธิปติ\csromandash{} อชฺฌตฺตา\-ธิปติ สมฺปยุตฺตกานํ ขนฺธานํ จิตฺตสมุฏฺฐานานญฺจ รูปานํ อธิปติ\-ปจฺจเยน ปจฺจโย. (1)}

\prose{อชฺฌตฺโต ธมฺโม พหิทฺธา ธมฺมสฺส อธิปติ\-ปจฺจเยน ปจฺจโย. อารมฺมณา\-ธิปติ\csromandash{}ปโร อชฺฌตฺตํ จกฺขุํ ฯเปฯ วตฺถุํ. อชฺฌตฺเต ขนฺเธ ครุํ กตฺวา อสฺสาเทติ อภินนฺทติ, ตํ ครุํ กตฺวา ราโค อุปฺปชฺชติ, ทิฏฺฐิ อุปฺปชฺชติ. (2)}

\proseitem{31}{พหิทฺธา ธมฺโม พหิทฺธา ธมฺมสฺส อธิปติ\-ปจฺจเยน ปจฺจโย. อารมฺมณา\-ธิปติ, สหชาตา\-ธิปติ.  อารมฺมณา\-ธิปติ\csromandash{}ปโร ทานํ ทตฺวา, สีลํ สมาทิยิตฺวา, อุโปสถกมฺมํ กตฺวา, ตํ ครุํ กตฺวา ปจฺจเวกฺขติ. ปุพฺเพ สุจิณฺณานิ ครุํ กตฺวา ฯเปฯ. ฌานา วุฏฺฐหิตฺวา ฯเปฯ. อริยา มคฺคา วุฏฺฐหิตฺวา มคฺคํ ครุํ กตฺวา ฯเปฯ ผลํ ครุํ กตฺวา ฯเปฯ. นิพฺพานํ ครุํ กตฺวา ปจฺจเวกฺขนฺติ. นิพฺพานํ โคตฺรภุสฺส, โวทานสฺส, มคฺคสฺส, ผลสฺส อธิปติ\-ปจฺจเยน ปจฺจโย. พหิทฺธา จกฺขุํ ฯเปฯ วตฺถุํ. พหิทฺธา ขนฺเธ ครุํ กตฺวา อสฺสาเทติ อภินนฺทติ, ตํ ครุํ กตฺวา ราโค อุปฺปชฺชติ, ทิฏฺฐิ อุปฺปชฺชติ. สหชาตา\-ธิปติ\csromandash{} พหิทฺธา\-ธิปติ สมฺปยุตฺตกานํ ขนฺธานํ จิตฺตสมุฏฺฐานานญฺจ รูปานํ อธิปติ\-ปจฺจเยน ปจฺจโย. (1)}

\prose{พหิทฺธา ธมฺโม อชฺฌตฺตสฺส ธมฺมสฺส อธิปติ\-ปจฺจเยน ปจฺจโย. อารมฺมณา\-ธิปติ\csromandash{}อริยา นิพฺพานํ ครุํ กตฺวา ปจฺจเวกฺขนฺติ. นิพฺพานํ โคตฺรภุสฺส, โวทานสฺส, มคฺคสฺส, ผลสฺส อธิปติ\-ปจฺจเยน ปจฺจโย. พหิทฺธา จกฺขุํ ฯเปฯ วตฺถุํ. พหิทฺธา ขนฺเธ ครุํ กตฺวา อสฺสาเทติ, อภินนฺทติ, ตํ ครุํ กตฺวา ราโค ฯเปฯ ทิฏฺฐิ อุปฺปชฺชติ. (2)}

\csromanmatikaanchor{mtk.244}
\csromantocmarkref{subparagraph}{อนนฺตร}{mtk.244}
\bookmark[dest={mtk.244},level=5]{อนนฺตร}
\csromanheader{6}{\textbf{อนนฺตร}}
\proseitem{32}{อชฺฌตฺโต ธมฺโม อชฺฌตฺตสฺส ธมฺมสฺส อนนฺตร\-ปจฺจเยน ปจฺจโย. ปุริมา ปุริมา อชฺฌตฺตา ขนฺธา ปจฺฉิมานํ ปจฺฉิมานํ อชฺฌตฺตานํ ขนฺธานํ อนนฺตร\-ปจฺจเยน ปจฺจโย. อนุโลมํ โคตฺรภุสฺส. อนุโลมํ โวทานสฺส. โคตฺรภุ มคฺคสฺส. โวทานํ มคฺคสฺส. มคฺโค ผลสฺส. ผลํ ผลสฺส. อนุโลมํ ผลสมาปตฺติยา. นิโรธา วุฏฺฐหนฺตสฺส เนว\-สญฺญา\-นา\-สญฺญา\-ยตนํ ผลสมาปตฺติยา อนนฺตร\-ปจฺจเยน ปจฺจโย. (1)}

\prose{\csromanfolio{445}พหิทฺธา ธมฺโม พหิทฺธา ธมฺมสฺส อนนฺตร\-ปจฺจเยน ปจฺจโย. (ปุริมา ปุริมา พหิทฺธาติ นานากรณํ, ตํ เยว คมนํ.) (1)}

\csromanmatikaanchor{mtk.245}
\csromantocmarkref{subparagraph}{สมนนฺตราทิ}{mtk.245}
\bookmark[dest={mtk.245},level=5]{สมนนฺตราทิ}
\csromanheader{6}{\textbf{สมนนฺตราทิ}}
\proseitem{33}{อชฺฌตฺโต ธมฺโม อชฺฌตฺตสฺส ธมฺมสฺส สมนนฺตร\-ปจฺจเยน ปจฺจโย ฯเปฯ. (อนนฺตร\-สทิสํ) สหชาต\-ปจฺจเยน ปจฺจโย. อญฺญมญฺญ\-ปจฺจเยน ปจฺจโย. นิสฺสย\-ปจฺจเยน ปจฺจโย.}

\csromanheader{2}{\textbf{อุปนิสฺสย}}
\proseitem{34}{อชฺฌตฺโต ธมฺโม อชฺฌตฺตสฺส ธมฺมสฺส อุปนิสฺสย\-ปจฺจเยน ปจฺจโย. อารมฺมณู\-ปนิสฺสโย, อนนฺตรู\-ปนิสฺสโย, ปกตู\-ปนิสฺสโย ฯเปฯ. ปกตู\-ปนิสฺสโย\csromandash{}อชฺฌตฺตํ สทฺธํ อุปนิสฺสาย ทานํ เทติ, สีลํ สมาทิยติ, อุโปสถกมฺมํ ฯเปฯ ฌานํ ฯเปฯ วิปสฺสนํ ฯเปฯ มคฺคํ ฯเปฯ อภิญฺญํ ฯเปฯ สมาปตฺติํ อุปฺปาเทติ. มานํ ชปฺเปติ, ทิฏฺฐิํ คณฺหาติ. อชฺฌตฺตํ สีลํ ฯเปฯ ปญฺญํ. ราคํ ฯเปฯ ปตฺถนํ, กายิกํ สุขํ. กายิกํ ทุกฺขํ. อุตุํ. โภชนํ. เสนาสนํ อุปนิสฺสาย ทานํ เทติ ฯเปฯ สมาปตฺติํ อุปฺปาเทติ. ปาณํ หนติ ฯเปฯ สํฆํ ภินฺทติ. อชฺฌตฺตา สทฺธา ฯเปฯ ปญฺญา. ราโค ฯเปฯ ปตฺถนา. กายิกํ สุขํ. กายิกํ ทุกฺขํ. เสนาสนํ อชฺฌตฺตาย สทฺธาย ฯเปฯ ปญฺญาย. ราคสฺส ฯเปฯ ปตฺถนาย. กายิกสฺส สุขสฺส, กายิกสฺส ทุกฺขสฺส, มคฺคสฺส, ผลสมาปตฺติยา อุปนิสฺสย\-ปจฺจเยน ปจฺจโย. (1)}

\prose{อชฺฌตฺโต ธมฺโม พหิทฺธา ธมฺมสฺส อุปนิสฺสย\-ปจฺจเยน ปจฺจโย. อารมฺมณู\-ปนิสฺสโย, ปกตู\-ปนิสฺสโย ฯเปฯ. ปกตู\-ปนิสฺสโย\csromandash{}ปโร อชฺฌตฺตํ สทฺธํ อุปนิสฺสาย ทานํ เทติ ฯเปฯ มานํ ชปฺเปติ, ทิฏฺฐิํ คณฺหาติ. ปโร อชฺฌตฺตํ สีลํ ฯเปฯ. เสนาสนํ อุปนิสฺสาย ทานํ เทติ. ปาณํ หนติ ฯเปฯ สํฆํ ภินฺทติ. อชฺฌตฺตา สทฺธา ฯเปฯ. เสนาสนํ พหิทฺธา สทฺธาย ฯเปฯ มคฺคสฺส, ผลสมาปตฺติยา อุปนิสฺสย\-ปจฺจเยน ปจฺจโย. (2)}

\proseitem{35}{พหิทฺธา ธมฺโม พหิทฺธา ธมฺมสฺส อุปนิสฺสย\-ปจฺจเยน ปจฺจโย. อารมฺมณู\-ปนิสฺสโย, อนนฺตรู\-ปนิสฺสโย, ปกตู\-ปนิสฺสโย ฯเปฯ. ปกตู\-ปนิสฺสโย\csromandash{}ปโร พหิทฺธา สทฺธํ ฯเปฯ ปตฺถนํ, กายิกํ สุขํ ฯเปฯ. \csromanfolio{446}เสนาสนํ อุปนิสฺสายา ทานํ เทติ ฯเปฯ สํฆํ ภินฺทติ, พหิทฺธา สทฺธา ฯเปฯ. เสนาสนํ พหิทฺธา สทฺธาย ฯเปฯ. ผลสมาปตฺติยา อุปนิสฺสย\-ปจฺจเยน ปจฺจโย. (1)}

\prose{พหิทฺธา ธมฺโม อชฺฌตฺตสฺส ธมฺมสฺส อุปนิสฺสย\-ปจฺจเยน ปจฺจโย. อารมฺมณู\-ปนิสฺสโย, ปกตู\-ปนิสฺสโย ฯเปฯ. ปกตู\-ปนิสฺสโย\csromandash{} พหิทฺธา สทฺธํ ฯเปฯ. เสนาสนํ อุปนิสฺสาย ทานํ เทติ ฯเปฯ สํฆํ ภินฺทติ. พหิทฺธา สทฺธา ฯเปฯ. เสนาสนํ อชฺฌตฺตาย สทฺธาย ฯเปฯ ผลสมาปตฺติยา อุปนิสฺสย\-ปจฺจเยน ปจฺจโย. (2)}

\csromanheader{2}{\textbf{ปุเรชาต}}
\proseitem{36}{อชฺฌตฺโต ธมฺโม อชฺฌตฺตสฺส ธมฺมสฺส ปุเรชาต\-ปจฺจเยน ปจฺจโย. อารมฺมณ\-ปุเรชาตํ, วตฺถุ\-ปุเรชาตํ.  อารมฺมณ\-ปุเรชาตํ\csromandash{} อชฺฌตฺตํ จกฺขุํ ฯเปฯ วตฺถุํ อนิจฺจโต ทุกฺขโต อนตฺตโต วิปสฺสติ ฯเปฯ โทมนสฺสํ อุปฺปชฺชติ, ทิพฺเพน จกฺขุนา รูปํ ปสฺสติ. ทิพฺพาย โสตธาตุยา สทฺทํ สุณาติ, รูปายตนํ จกฺขุ\-วิญฺญาณสฺส ฯเปฯ โผฏฺฐพฺพา\-ยตนํ กายวิญฺญาณสฺส ปุเรชาต\-ปจฺจเยน ปจฺจโย. วตฺถุ\-ปุเรชาตํ\csromandash{}จกฺขายตนํ จกฺขุ\-วิญฺญาณสฺส ฯเปฯ. กายายตนํ กายวิญฺญาณสฺส ฯเปฯ. วตฺถุ อชฺฌตฺตานํ ขนฺธานํ ปุเรชาต\-ปจฺจเยน ปจฺจโย. (1)}

\prose{อชฺฌตฺโต ธมฺโม พหิทฺธา ธมฺมสฺส ปุเรชาต\-ปจฺจเยน ปจฺจโย. อารมฺมณ\-ปุเรชาตํ\csromandash{}ปโร อชฺฌตฺตํ จกฺขุํ ฯเปฯ วตฺถุํ อนิจฺจโต ฯเปฯ โทมนสฺสํ อุปฺปชฺชติ. ทิพฺเพน จกฺขุนา รูปํ ปสฺสติ. ทิพฺพาย โสตธาตุยา สทฺทํ สุณาติ. อชฺฌตฺตํ รูปายตนํ พหิทฺธา จกฺขุ\-วิญฺญาณสฺส ฯเปฯ. อชฺฌตฺตํ โผฏฺฐพฺพา\-ยตนํ พหิทฺธา กายวิญฺญาณสฺส ปุเรชาต\-ปจฺจเยน ปจฺจโย. (2)}

\proseitem{37}{พหิทฺธา ธมฺโม พหิทฺธา ธมฺมสฺส ปุเรชาต\-ปจฺจเยน ปจฺจโย. อารมฺมณ\-ปุเรชาตํ, วตฺถุ\-ปุเรชาตํ.  อารมฺมณ\-ปุเรชาตํ\csromandash{} ปโร พหิทฺธา จกฺขุํ ฯเปฯ. วตฺถุํ อนิจฺจโต ฯเปฯ. ทิพฺเพน จกฺขุนา รูปํ ปสฺสติ. ทิพฺพาย โสตธาตุยา สทฺทํ สุณาติ. พหิทฺธา รูปายตนํ พหิทฺธา จกฺขุ\-วิญฺญาณสฺส ฯเปฯ. พหิทฺธา โผฏฺฐพฺพา\-ยตนํ พหิทฺธา กายวิญฺญาณสฺส ฯเปฯ. วตฺถุ\-ปุเรชาตํ\csromandash{}พหิทฺธา จกฺขายตนํ ฯเปฯ. กายายตนํ ฯเปฯ. วตฺถุ พหิทฺธา ขนฺธานํ ปุเรชาต\-ปจฺจเยน ปจฺจโย. (1)}

\prose{\csromanfolio{447}พหิทฺธา ธมฺโม อชฺฌตฺตสฺส ธมฺมสฺส ปุเรชาต\-ปจฺจเยน ปจฺจโย. อารมฺมณ\-ปุเรชาตํ\csromandash{}พหิทฺธา จกฺขุํ ฯเปฯ. วตฺถุํ อนิจฺจโต ฯเปฯ โทมนสฺสํ อุปฺปชฺชติ. ทิพฺเพน จกฺขุนา รูปํ ปสฺสติ. ทิพฺพาย โสตธาตุยา สทฺทํ สุณาติ. พหิทฺธา รูปายตนํ อชฺฌตฺตสฺส จกฺขุ\-วิญฺญาณสฺส ฯเปฯ. พหิทฺธา โผฏฺฐพฺพา\-ยตนํ อชฺฌตฺตสฺส กายวิญฺญาณสฺส ปุเรชาต\-ปจฺจเยน ปจฺจโย. (2)}

\proseitem{38}{อชฺฌตฺโต จ พหิทฺธา จ ธมฺมา อชฺฌตฺตสฺส ธมฺมสฺส ปุเรชาต\-ปจฺจเยน ปจฺจโย. อารมฺมณ\-ปุเรชาตํ, วตฺถุ\-ปุเรชาตํ.  พหิทฺธา รูปายตนญฺจ อชฺฌตฺตํ จกฺขายตนญฺจ อชฺฌตฺตสฺส จกฺขุ\-วิญฺญาณสฺส ปุเรชาต\-ปจฺจเยน ปจฺจโย ฯเปฯ. พหิทฺธา โผฏฺฐพฺพายตนญฺจ อชฺฌตฺตํ กายายตนญฺจ อชฺฌตฺตสฺส กายวิญฺญาณสฺส ฯเปฯ. พหิทฺธา รูปายตนญฺจ อชฺฌตฺตํ วตฺถุ จ ฯเปฯ. พหิทฺธา โผฏฺฐพฺพายตนญฺจ อชฺฌตฺตํ วตฺถุ จ อชฺฌตฺตานํ ขนฺธานํ ปุเรชาต\-ปจฺจเยน ปจฺจโย. (1)}

\prose{อชฺฌตฺโต จ พหิทฺธา จ ธมฺมา พหิทฺธา ธมฺมสฺส ปุเรชาต\-ปจฺจเยน ปจฺจโย. อารมฺมณ\-ปุเรชาตํ, วตฺถุ\-ปุเรชาตํ.  อชฺฌตฺตํ รูปายตนญฺจ พหิทฺธา จกฺขายตนญฺจ พหิทฺธา จกฺขุ\-วิญฺญาณสฺส ปุเรชาต\-ปจฺจเยน ปจฺจโย ฯเปฯ. อชฺฌตฺตํ โผฏฺฐพฺพายตนญฺจ พหิทฺธา กายายตนญฺจ พหิทฺธา กายวิญฺญาณสฺส ปุเรชาต\-ปจฺจเยน ปจฺจโย. อชฺฌตฺตํ รูปายตนญฺจ พหิทฺธา วตฺถุ จ ฯเปฯ. อชฺฌตฺตํ โผฏฺฐพฺพายตนญฺจ พหิทฺธา วตฺถุ จ พหิทฺธา ขนฺธานํ ปุเรชาต\-ปจฺจเยน ปจฺจโย. (2)}

\csromanheader{2}{\textbf{ปจฺฉาชาต}}
\proseitem{39}{อชฺฌตฺโต ธมฺโม อชฺฌตฺตสฺส ธมฺมสฺส ปจฺฉาชาต\-ปจฺจเยน ปจฺจโย. ปจฺฉาชาตา อชฺฌตฺตา ขนฺธา ปุเรชาตสฺส อิมสฺส กายสฺส ปจฺฉาชาต\-ปจฺจเยน ปจฺจโย. (1)}

\prose{พหิทฺธา ธมฺโม พหิทฺธา ธมฺมสฺส ปจฺฉาชาต\-ปจฺจเยน ปจฺจโย. ปจฺฉาชาตา พหิทฺธา ขนฺธา ปุเรชาตสฺส อิมสฺส กายสฺส ปจฺฉาชาต\-ปจฺจเยน ปจฺจโย. (1)}

\csromanheader{2}{\textbf{อาเสวน}}
\proseitem{40}{อชฺฌตฺโต ธมฺโม อชฺฌตฺตสฺส ธมฺมสฺส อาเสวน\-ปจฺจเยน ปจฺจโย. ปุริมา ปุริมา อชฺฌตฺตา ขนฺธา ปจฺฉิมานํ ปจฺฉิมานํ อชฺฌตฺตานํ ขนฺธานํ \csromanfolio{448}อาเสวน\-ปจฺจเยน ปจฺจโย. อนุโลมํ โคตฺรภุสฺส. อนุโลมํ โวทานสฺส. โคตฺรภุ มคฺคสฺส. โวทานํ มคฺคสฺส อาเสวน\-ปจฺจเยน ปจฺจโย. (1)}

\prose{พหิทฺธา ธมฺโม พหิทฺธา ธมฺมสฺส อาเสวน\-ปจฺจเยน ปจฺจโย. ปุริมา ปุริมา ฯเปฯ. (อชฺฌตฺตสทิสํเยว.)}

\csromanheader{2}{\textbf{กมฺม}}
\proseitem{41}{อชฺฌตฺโต ธมฺโม อชฺฌตฺตสฺส ธมฺมสฺส กมฺมปจฺจเยน ปจฺจโย. สหชาตา, นานากฺขณิกา.  สหชาตา อชฺฌตฺตา เจตนา สมฺปยุตฺตกานํ ขนฺธานํ จิตฺตสมุฏฺฐานานญฺจ รูปานํ กมฺมปจฺจเยน ปจฺจโย. นานากฺขณิกา อชฺฌตฺตา เจตนา วิปากานํ อชฺฌตฺตานํ ขนฺธานํ กฏตฺตา จ รูปานํ กมฺมปจฺจเยน ปจฺจโย. (1)}

\prose{พหิทฺธา ธมฺโม พหิทฺธา ธมฺมสฺส กมฺมปจฺจเยน ปจฺจโย. สหชาตา, นานากฺขณิกา.  สหชาตา พหิทฺธาเจตนา สมฺปยุตฺตกานํ ขนฺธานํ จิตฺตสมุฏฺฐานานญฺจ รูปานํ กมฺมปจฺจเยน ปจฺจโย. นานากฺขณิกา พหิทฺธา เจตนา วิปากานํ พหิทฺธา ขนฺธานํ กฏตฺตา จ รูปานํ กมฺมปจฺจเยน ปจฺจโย. (1)}

\csromanheader{2}{\textbf{วิปาก}}
\proseitem{42}{อชฺฌตฺโต ธมฺโม อชฺฌตฺตสฺส ธมฺมสฺส วิปาก\-ปจฺจเยน ปจฺจโย ฯเปฯ. (ปริปุณฺณํ, ปฏิจฺจ\-วาร\-สทิสํ.)}

\csromanheader{2}{\textbf{อาหาร}}
\proseitem{43}{อชฺฌตฺโต ธมฺโม อชฺฌตฺตสฺส ธมฺมสฺส อาหารปจฺจเยน ปจฺจโย. อชฺฌตฺตา อาหารา สมฺปยุตฺตกานํ ขนฺธานํ จิตฺตสมุฏฺฐานานญฺจ รูปานํ อาหารปจฺจเยน ปจฺจโย. ปฏิสนฺธิ\-กฺขเณ ฯเปฯ. อชฺฌตฺโต กพฬีกาโร อาหาโร อชฺฌตฺตสฺส กายสฺส อาหารปจฺจเยน ปจฺจโย. (1)}

\prose{อชฺฌตฺโต ธมฺโม พหิทฺธา ธมฺมสฺส อาหารปจฺจเยน ปจฺจโย. อชฺฌตฺโต กพฬีกาโร อาหาโร พหิทฺธา กายสฺส อาหารปจฺจเยน ปจฺจโย. (2)}

\prose{\csromanfolio{449}พหิทฺธา ธมฺโม พหิทฺธา ธมฺมสฺส อาหารปจฺจเยน ปจฺจโย. (ปวตฺติ\-ปฏิสนฺธิ,) พหิทฺธา กพฬีกาโร อาหาโร พหิทฺธา กายสฺส อาหารปจฺจเยน ปจฺจโย. (1)}

\prose{พหิทฺธา ธมฺโม อชฺฌตฺตสฺส ธมฺมสฺส อาหารปจฺจเยน ปจฺจโย. พหิทฺธา กพฬีกาโร อาหาโร อชฺฌตฺตสฺส กายสฺส อาหารปจฺจเยน ปจฺจโย. (2)}

\proseitem{44}{อชฺฌตฺโต ธมฺโม จ พหิทฺธา ธมฺโม จ อชฺฌตฺตสฺส ธมฺมสฺส อาหารปจฺจเยน ปจฺจโย. อชฺฌตฺโต กพฬีกาโร อาหาโร จ พหิทฺธา กพฬีกาโร อาหาโร จ อชฺฌตฺตสฺส กายสฺส อาหารปจฺจเยน ปจฺจโย. (1)}

\prose{อชฺฌตฺโต ธมฺโม จ พหิทฺธา ธมฺโม จ พหิทฺธา ธมฺมสฺส อาหารปจฺจเยน ปจฺจโย. อชฺฌตฺโต กพฬีกาโร อาหาโร จ พหิทฺธา กพฬีกาโร อาหาโร จ พหิทฺธา กายสฺส อาหารปจฺจเยน ปจฺจโย. (2)}

\csromanheader{2}{\textbf{อินฺทฺริยาทิ}}
\proseitem{45}{อชฺฌตฺโต ธมฺโม อชฺฌตฺตสฺส ธมฺมสฺส อินฺทฺริย\-ปจฺจเยน ปจฺจโย. อชฺฌตฺติกา อินฺทฺริยา (รูปชีวิตินฺทฺริยมฺปิ วิตฺถาเรตพฺพํ,) ฌานปจฺจเยน ปจฺจโย. มคฺคปจฺจเยน ปจฺจโย. สมฺปยุตฺต\-ปจฺจเยน ปจฺจโย. วิปฺปยุตฺต\-ปจฺจเยน ปจฺจโย. สหชาตํ, ปุเรชาตํ, ปจฺฉาชาตํ, (มาติกาปทานิ อนุมชฺชนฺเตน วิตฺถาเร\-ตพฺพานิ.)}

\prose{พหิทฺธา ธมฺโม พหิทฺธา ธมฺมสฺส วิปฺปยุตฺต\-ปจฺจเยน ปจฺจโย. สหชาตํ ปุเรชาตํ, ปจฺฉาชาตํ. (สํขิตฺตํ.)}

\csromanheader{2}{\textbf{อตฺถิ}}
\proseitem{46}{อชฺฌตฺโต ธมฺโม อชฺฌตฺตสฺส ธมฺมสฺส อตฺถิปจฺจเยน ปจฺจโย. สหชาตํ, ปุเรชาตํ, ปจฺฉาชาตํ, อาหารํ, อินฺทฺริยํ.  สหชาโต อชฺฌตฺโต เอโก ขนฺโธ ติณฺณนฺนํ ขนฺธานํ จิตฺตสมุฏฺฐานานญฺจ รูปานํ ฯเปฯ เทฺว ขนฺเธ ฯเปฯ. ปฏิสนฺธิ\-กฺขเณ ฯเปฯ. ขนฺธา วตฺถุสฺส ฯเปฯ. วตฺถุ ขนฺธานํ ฯเปฯ. เอกํ มหาภูตํ ฯเปฯ. อสญฺญสตฺตานํ เอกํ มหาภูตํ ติณฺณนฺนํ มหาภูตานํ ฯเปฯ. ปุเรชาตํ จกฺขุํ ฯเปฯ. วตฺถุํ. (ปุเรชาต\-สทิสํ.) วตฺถุ อชฺฌตฺตานํ ขนฺธานํ อตฺถิปจฺจเยน ปจฺจโย. ปจฺฉาชาตา อชฺฌตฺตา ขนฺธา ปุเรชาตสฺส อิมสฺส \csromanfolio{450}กายสฺส อตฺถิปจฺจเยน ปจฺจโย. อชฺฌตฺโต \textbf{กพฬีกาโร อาหาโร} อชฺฌตฺตสฺส กายสฺส ฯเปฯ. \textbf{รูปชีวิตินฺทฺริยํ} กฏตฺตารูปานํ ฯเปฯ. (1)}

\prose{อชฺฌตฺโต ธมฺโม พหิทฺธา ธมฺมสฺส อตฺถิปจฺจเยน ปจฺจโย. ปุเรชาตํ, อาหารํ.  ปุเรชาตํ\csromandash{}ปโร อชฺฌตฺตํ จกฺขุํ ฯเปฯ วตฺถุํ อนิจฺจโต ฯเปฯ วิปสฺสติ, ทิพฺเพน จกฺขุนา รูปํ ปสฺสติ. ทิพฺพาย โสตธาตุยา สทฺทํ สุณาติ. อชฺฌตฺตํ รูปายตนํ ฯเปฯ. โผฏฺฐพฺพา\-ยตนํ พหิทฺธา กายวิญฺญาณสฺส อตฺถิปจฺจเยน ปจฺจโย. อชฺฌตฺโต กพฬีกาโร อาหาโร พหิทฺธา กายสฺส อตฺถิปจฺจเยน ปจฺจโย. (2)}

\proseitem{47}{พหิทฺธา ธมฺโม พหิทฺธา ธมฺมสฺส อตฺถิปจฺจเยน ปจฺจโย. สหชาตํ, ปุเรชาตํ, ปจฺฉาชาตํ, อาหารํ, อินฺทฺริยํ. (พหิทฺธา นินฺนานากรณํ, มาติกาปทานิ วิตฺถาเร\-ตพฺพานิ.) (1)}

\prose{พหิทฺธา ธมฺโม อชฺฌตฺตสฺส ธมฺมสฺส อตฺถิปจฺจเยน ปจฺจโย. ปุเรชาตํ อาหารํ.  ปุเรชาตํ\csromandash{}พหิทฺธา จกฺขุํ ฯเปฯ วตฺถุํ ฯเปฯ. ทิพฺเพน จกฺขุนา รูปํ ปสฺสติ. ทิพฺพาย โสตธาตุยา สทฺทํ สุณาติ. พหิทฺธา รูปายตนํ ฯเปฯ. โผฏฺฐพฺพา\-ยตนํ อชฺฌตฺตสฺส กายวิญฺญาณสฺส อตฺถิปจฺจเยน ปจฺจโย. พหิทฺธา กพฬีกาโร อาหาโร อชฺฌตฺตสฺส กายสฺส อตฺถิปจฺจเยน ปจฺจโย. (2)}

\proseitem{48}{อชฺฌตฺโต ธมฺโม จ พหิทฺธา ธมฺโม จ อชฺฌตฺตสฺส ธมฺมสฺส อตฺถิปจฺจเยน ปจฺจโย. ปุเรชาตํ, อาหารํ.  ปุเรชาตํ\csromandash{}พหิทฺธา รูปายตนญฺจ อชฺฌตฺตํ จกฺขุ จ อชฺฌตฺตสฺส จกฺขุ\-วิญฺญาณสฺส ฯเปฯ. พหิทฺธา โผฏฺฐพฺพายตนญฺจ อชฺฌตฺตํ กายายตนญฺจ อชฺฌตฺตสฺส กายวิญฺญาณสฺส อตฺถิปจฺจเยน ปจฺจโย. พหิทฺธา รูปายตนญฺจ อชฺฌตฺตํ วตฺถุ จ ฯเปฯ. พหิทฺธา โผฏฺฐพฺพายตนญฺจ อชฺฌตฺตํ วตฺถุ จ อชฺฌตฺตานํ ขนฺธานํ อตฺถิปจฺจเยน ปจฺจโย. อาหารํ\csromandash{}อชฺฌตฺโต กพฬีกาโร อาหาโร จ พหิทฺธา กพฬีกาโร อาหาโร จ อชฺฌตฺตสฺส กายสฺส อตฺถิปจฺจเยน ปจฺจโย. (1)}

\csromanmatikaanchor{mtk.246}
\csromanmatikaanchor{mtk.247}
\csromantocmarknopage{subparagraph}{1. ปจฺจยานุโลม 2. สงฺขฺยาวาร}{mtk.246}
\bookmark[dest={mtk.246},level=5]{1. ปจฺจยานุโลม 2. สงฺขฺยาวาร}
\csromantocmarkref{subparagraph}{สุทฺธ}{mtk.247}
\bookmark[dest={mtk.247},level=5]{สุทฺธ}
\prose{อชฺฌตฺโต ธมฺโม จ พหิทฺธา ธมฺโม จ พหิทฺธา ธมฺมสฺส อตฺถิปจฺจเยน ปจฺจโย. ปุเรชาตํ, อาหารํ.  ปุเรชาตํ\csromandash{}อชฺฌตฺตํ รูปายตนญฺจ พหิทฺธา จกฺขายตนญฺจ พหิทฺธา จกฺขุ\-วิญฺญาณสฺส อตฺถิปจฺจเยน ปจฺจโย ฯเปฯ. \csromanfolio{451}อชฺฌตฺตํ โผฏฺฐพฺพายตนญฺจ พหิทฺธา กายายตนญฺจ พหิทฺธา กายวิญฺญาณสฺส อตฺถิปจฺจเยน ปจฺจโย. อชฺฌตฺตํ รูปายตนญฺจ พหิทฺธา วตฺถุ จ พหิทฺธา ขนฺธานํ อตฺถิปจฺจเยน ปจฺจโย ฯเปฯ. อชฺฌตฺตํ โผฏฺฐพฺพายตนญฺจ พหิทฺธา วตฺถุ จ พหิทฺธา ขนฺธานํ อตฺถิปจฺจเยน ปจฺจโย. อาหารํ\csromandash{}อชฺฌตฺโต กพฬีกาโร อาหาโร จ พหิทฺธา กพฬีกาโร อาหาโร จ พหิทฺธา กายสฺส อตฺถิปจฺจเยน ปจฺจโย. (2)}

\csromanheader{2}{\textbf{นตฺถิ วิคตาวิคต}}
\proseitem{49}{อชฺฌตฺโต ธมฺโม อชฺฌตฺตสฺส ธมฺมสฺส นตฺถิ\-ปจฺจเยน ปจฺจโย. วิคต\-ปจฺจเยน ปจฺจโย. อวิคต\-ปจฺจเยน ปจฺจโย.}

\csromanmarkh{1. ปจฺจยานุโลม}\csromanmarkhii{2. สงฺขฺยาวาร}\csromanheadernomark{2}{1. \textbf{ปจฺจยานุโลม 2. สงฺขฺยาวาร}}
\proseitemwithcloserruled{50}{เหตุยา เทฺว, อารมฺมเณ จตฺตาริ, อธิปติยา จตฺตาริ, อนนฺตเร เทฺว, สมนนฺตเร เทฺว, สหชาเต, อญฺญมญฺเญ, นิสฺสเย เทฺว, อุปนิสฺสเย จตฺตาริ, ปุเรชาเต ฉ, ปจฺฉาชาเต, อาเสวเน, กมฺเม, วิปาเก เทฺว, อาหาเร ฉ, อินฺทฺริเย เทฺว, ฌาเน, มคฺเค, สมฺปยุตฺเต, วิปฺปยุตฺเต เทฺว, อตฺถิยา ฉ, นตฺถิยา เทฺว, วิคเต เทฺว, อวิคเต ฉ. (เอวํ คเณตพฺพํ.)}{อนุโลมํ.}
\csromanmatikaanchor{mtk.248}
\csromantocmarkref{subparagraph}{2. ปจฺจนียุทฺธาร}{mtk.248}
\bookmark[dest={mtk.248},level=5]{2. ปจฺจนียุทฺธาร}
\csromanheader{6}{2. ปจฺจนียุ\-ทฺธาร}
\proseitem{51}{อชฺฌตฺโต ธมฺโม อชฺฌตฺตสฺส ธมฺมสฺส อารมฺมณ\-ปจฺจเยน ปจฺจโย. สหชาต\-ปจฺจเยน ปจฺจโย. อุปนิสฺสย\-ปจฺจเยน ปจฺจโย. ปุเรชาต\-ปจฺจเยน ปจฺจโย. ปจฺฉาชาต\-ปจฺจเยน ปจฺจโย. กมฺมปจฺจเยน ปจฺจโย. อาหารปจฺจเยน ปจฺจโย. อินฺทฺริย\-ปจฺจเยน ปจฺจโย. (1)}

\prose{อชฺฌตฺโต ธมฺโม พหิทฺธา ธมฺมสฺส อารมฺมณ\-ปจฺจเยน ปจฺจโย. อุปนิสฺสย\-ปจฺจเยน ปจฺจโย. ปุเรชาต\-ปจฺจเยน ปจฺจโย. อาหารปจฺจเยน ปจฺจโย. (2)}

\csromanmatikaanchor{mtk.249}
\csromanmatikaanchor{mtk.250}
\csromanmatikaanchor{mtk.251}
\csromantocmarkref{subparagraph}{2. ปจฺจยปจฺจนีย 2. สงฺขฺยาวาร}{mtk.249}
\bookmark[dest={mtk.249},level=5]{2. ปจฺจยปจฺจนีย 2. สงฺขฺยาวาร}
\csromantocmarknopage{subparagraph}{3. ปจฺจยานุโลมปจฺจนีย}{mtk.250}
\bookmark[dest={mtk.250},level=5]{3. ปจฺจยานุโลมปจฺจนีย}
\csromantocmarkref{subparagraph}{เหตุทุก}{mtk.251}
\bookmark[dest={mtk.251},level=5]{เหตุทุก}
\proseitem{52}{\csromanfolio{452}พหิทฺธา ธมฺโม พหิทฺธา ธมฺมสฺส อารมฺมณ\-ปจฺจเยน ปจฺจโย. สหชาต\-ปจฺจเยน ปจฺจโย. อุปนิสฺสย\-ปจฺจเยน ปจฺจโย. ปุเรชาต\-ปจฺจเยน ปจฺจโย. ปจฺฉาชาต\-ปจฺจเยน ปจฺจโย. กมฺมปจฺจเยน ปจฺจโย. อาหารปจฺจเยน ปจฺจโย. อินฺทฺริย\-ปจฺจเยน ปจฺจโย. (1)}

\prose{พหิทฺธา ธมฺโม อชฺฌตฺตสฺส ธมฺมสฺส อารมฺมณ\-ปจฺจเยน ปจฺจโย. อุปนิสฺสย\-ปจฺจเยน ปจฺจโย. ปุเรชาต\-ปจฺจเยน ปจฺจโย. อาหารปจฺจเยน ปจฺจโย. (2)}

\proseitem{53}{อชฺฌตฺโต ธมฺโม จ พหิทฺธา ธมฺโม จ อชฺฌตฺตสฺส ธมฺมสฺส ปุเรชาตํ, อาหารํ. (1)}

\prose{อชฺฌตฺโต ธมฺโม จ พหิทฺธา ธมฺโม จ พหิทฺธา ธมฺมสฺส ปุเรชาตํ, อาหารํ. (2)}

\csromanmarkh{2. ปจฺจย\-ปจฺจนีย}\csromanmarkhii{2. สงฺขฺยาวาร}\csromanheadernomark{6}{2. ปจฺจย\-ปจฺจนีย 2. สงฺขฺยาวาร}
\proseitemwithcloserruled{54}{นเหตุยา ฉ, นอารมฺมเณ ฉ, นอธิปติยา ฉ, (สํขิตฺตํ. สพฺพตฺถ ฉ กาตพฺพา.) นวิปฺปยุตฺเต ฉ, โนอตฺถิยา จตฺตาริ, โนนตฺถิยา ฉ, โนวิคเต ฉ, โนอวิคเต จตฺตาริ. (เอวํ คเณตพฺพํ.)}{ปจฺจนียํ.}
\proseitem{55}{เหตุปจฺจยา นอารมฺมเณ เทฺว, นอธิปติยา, นอนนฺตเร, นสมนนฺตเร, นอญฺญมญฺเญ, นอุปนิสฺสเย เทฺว. (สํขิตฺตํ. สพฺพตฺถ เทฺว.) นสมฺปยุตฺเต, นวิปฺปยุตฺเต, โนนตฺถิยา, โนวิคเต เทฺว. (เอวํ คเณตพฺพํ.)}

\vspace{\tipitakaparskip}
\csromancenter{อนุโลม\-ปจฺจนียํ.}
\csromanmatikaanchor{mtk.252}
\csromantocmarkref{subparagraph}{4. ปจฺจยปจฺจนียานุโลม ...}{mtk.252}
\bookmark[dest={mtk.252},level=5]{4. ปจฺจยปจฺจนียานุโลม ...}
\csromanmarkh{20. อชฺฌตฺตตฺติก}\csromanmarkhii{4. ปจฺจย\-ปจฺจนียา\-นุโลม}\csromanheadernomark{6}{20. อชฺฌตฺตตฺติก 4. ปจฺจย\-ปจฺจนียา\-นุโลม}
\proseitem{56}{\csromanfolio{453}นเหตุปจฺจยา อารมฺมเณ จตฺตาริ, อธิปติยา จตฺตาริ, (อนุโลมปทานิ คเณตพฺพานิ.) อวิคเต ฉ. (เอวํ คเณตพฺพํ.)}

\vspace{\tipitakaparskip}
\csromancenter{ปจฺจนียา\-นุโลมํ.}
\csromancenter{ปญฺหาวาโร.}
\nitthitam{อชฺฌตฺตตฺติกํ นิฏฺฐิตํ.}
\csromanmatikaanchor{mtk.253}
\csromantocmarknopage{subparagraph}{21. อชฺฌตฺตารมฺมณตฺติก 1. ปฏิจฺจวาร}{mtk.253}
\bookmark[dest={mtk.253},level=5]{21. อชฺฌตฺตารมฺมณตฺติก 1. ปฏิจฺจวาร}
\csromanmarkh{21. อชฺฌตฺตา\-รมฺมณ\-ตฺติก}\csromanmarkhii{1. ปฏิจฺจวาร}\csromanheadernomark{6}{21. อชฺฌตฺตา\-รมฺมณ\-ตฺติก 1. ปฏิจฺจวาร}
\csromanmatikaanchor{mtk.254}
\csromantocmarkref{subparagraph}{1. ปจฺจยานุโลม 1. วิภงฺควาร}{mtk.254}
\bookmark[dest={mtk.254},level=5]{1. ปจฺจยานุโลม 1. วิภงฺควาร}
\csromanmarkh{1. ปจฺจยานุโลม}\csromanmarkhii{1. วิภงฺควาร}\csromanheadernomark{6}{1. \textbf{ปจฺจยานุโลม 1. วิภงฺควาร}}
\csromanheader{2}{\textbf{เหตุ}}
\proseitem{1}{\csromanfolio{454}อชฺฌตฺตา\-รมฺมณํ ธมฺมํ ปฏิจฺจ อชฺฌตฺตา\-รมฺมโณ ธมฺโม อุปฺปชฺชติ เหตุปจฺจยา. อชฺฌตฺตา\-รมฺมณํ เอกํ ขนฺธํ ปฏิจฺจ ตโย ขนฺธา ฯเปฯ เทฺว ขนฺเธ ฯเปฯ. ปฏิสนฺธิ\-กฺขเณ อชฺฌตฺตา\-รมฺมณํ เอกํ ขนฺธํ ปฏิจฺจ ตโย ขนฺธา ฯเปฯ เทฺว ขนฺเธ ฯเปฯ. (1)}

\prose{พหิทฺธา\-รมฺมณํ ธมฺมํ ปฏิจฺจ พหิทฺธา\-รมฺมโณ ธมฺโม อุปฺปชฺชติ เหตุปจฺจยา. พหิทฺธา\-รมฺมณํ เอกํ ขนฺธํ ปฏิจฺจ ตโย ขนฺธา ฯเปฯ เทฺว ขนฺเธ ฯเปฯ. ปฏิสนฺธิ\-กฺขเณ พหิทฺธา\-รมฺมณํ เอกํ ขนฺธํ ปฏิจฺจ ตโย ขนฺธา ฯเปฯ เทฺว ขนฺเธ ฯเปฯ. (1)}

\csromanheader{2}{\textbf{อารมฺมณาทิ}}
\proseitem{2}{อชฺฌตฺตา\-รมฺมณํ ธมฺมํ ปฏิจฺจ อชฺฌตฺตา\-รมฺมโณ ธมฺโม อุปฺปชฺชติ อารมฺมณ\-ปจฺจยา ฯเปฯ. อวิคตปจฺจยา. (สํขิตฺตํ.)}

\csromanmatikaanchor{mtk.255}
\csromantocmarkref{subparagraph}{1. ปจฺจยานุโลม 2. สงฺขฺยาวาร}{mtk.255}
\bookmark[dest={mtk.255},level=5]{1. ปจฺจยานุโลม 2. สงฺขฺยาวาร}
\csromanmarkh{1. ปจฺจยานุโลม}\csromanmarkhii{2. สงฺขฺยาวาร}\csromanheadernomark{6}{1. \textbf{ปจฺจยานุโลม 2. สงฺขฺยาวาร}}
\proseitemwithcloserruled{3}{\textbf{เหตุ}ปจฺจยา อารมฺมเณ เทฺว, (สํขิตฺตํ. สพฺพตฺถ เทฺว.) อวิคเต เทฺว. (เอวํ คเณตพฺพํ.)}{\textbf{อนุโลมํ.}}
\csromanmarkh{2. ปจฺจย\-ปจฺจนีย}\csromanmarkhii{1. วิภงฺควาร}\csromanheadernomark{2}{2. ปจฺจย\-ปจฺจนีย 1. วิภงฺควาร}
\csromanmatikaanchor{mtk.256}
\csromanmatikaanchor{mtk.257}
\csromantocmarknopage{subparagraph}{2. ปจฺจยปจฺจนีย 1. วิภงฺควาร}{mtk.256}
\bookmark[dest={mtk.256},level=5]{2. ปจฺจยปจฺจนีย 1. วิภงฺควาร}
\csromantocmarkref{subparagraph}{นเหตุ}{mtk.257}
\bookmark[dest={mtk.257},level=5]{นเหตุ}
\csromanheader{6}{\textbf{นเหตุ}}
\proseitem{4}{อชฺฌตฺตา\-รมฺมณํ ธมฺมํ ปฏิจฺจ อชฺฌตฺตา\-รมฺมโณ ธมฺโม อุปฺปชฺชติ นเหตุปจฺจยา. อเหตุกํ อชฺฌตฺตา\-รมฺมณํ เอกํ ขนฺธํ ปฏิจฺจ ตโย ขนฺธา ฯเปฯ เทฺว ขนฺเธ ฯเปฯ. อเหตุก\-ปฏิสนฺธิ\-กฺขเณ อชฺฌตฺตา\-รมฺมณํ เอกํ ขนฺธํ ปฏิจฺจ ตโย ขนฺธา ฯเปฯ เทฺว ขนฺเธ ฯเปฯ. วิจิกิจฺฉา\-สหคเต อุทฺธจฺจ\-สหคเต ขนฺเธ ปฏิจฺจ วิจิกิจฺฉา\-สหคโต อุทฺธจฺจ\-สหคโต โมโห. (1)}

\csromanmatikaanchor{mtk.258}
\csromantocmarkref{subparagraph}{นอธิปตฺยาทิ}{mtk.258}
\bookmark[dest={mtk.258},level=5]{นอธิปตฺยาทิ}
\prose{\csromanfolio{455}พหิทฺธา\-รมฺมณํ ธมฺมํ ปฏิจฺจ พหิทฺธา\-รมฺมโณ ธมฺโม อุปฺปชฺชติ นเหตุปจฺจยา. อเหตุกํ พหิทฺธา\-รมฺมณํ เอกํ ขนฺธํ ปฏิจฺจ ตโย ขนฺธา ฯเปฯ เทฺว ขนฺเธ ฯเปฯ. อเหตุก\-ปฏิสนฺธิ\-กฺขเณ ฯเปฯ. วิจิกิจฺฉา\-สหคเต อุทฺธจฺจ\-สหคเต ขนฺเธ ปฏิจฺจ วิจิกิจฺฉา\-สหคโต อุทฺธจฺจ\-สหคโต โมโห. (1)}

\prose{นอธิปตฺยาทิ}

\proseitem{5}{อชฺฌตฺตา\-รมฺมณํ ธมฺมํ ปฏิจฺจ อชฺฌตฺตา\-รมฺมโณ ธมฺโม อุปฺปชฺชติ นอธิปติ\-ปจฺจยา. (อนุโลม\-สหชาต\-สทิสํ. นินฺนานากรณํ.) นปุเรชาต\-ปจฺจยา. อรูเป อชฺฌตฺตา\-รมฺมณํ เอกํ ขนฺธํ ปฏิจฺจ ฯเปฯ. ปฏิสนฺธิ\-กฺขเณ ฯเปฯ. (1)}

\proseitem{6}{พหิทฺธา\-รมฺมณํ ธมฺมํ ปฏิจฺจ พหิทฺธา\-รมฺมโณ ธมฺโม อุปฺปชฺชติ นปุเรชาต\-ปจฺจยา. อรูเป พหิทฺธา\-รมฺมณํ เอกํ ขนฺธํ ปฏิจฺจ ตโย ขนฺธา ฯเปฯ. ปฏิสนฺธิ\-กฺขเณ ฯเปฯ. นปจฺฉาชาต\-ปจฺจยา. นอาเสวน\-ปจฺจยา. (สหชาต\-สทิสํ.) นกมฺมปจฺจยา. อชฺฌตฺตา\-รมฺมเณ ขนฺเธ ปฏิจฺจ อชฺฌตฺตา\-รมฺมณา เจตนา.}

\prose{พหิทฺธา\-รมฺมณํ ธมฺมํ ปฏิจฺจ พหิทฺธา\-รมฺมโณ ธมฺโม อุปฺปชฺชติ นกมฺมปจฺจยา. พหิทฺธา\-รมฺมเณ ขนฺเธ ปฏิจฺจ พหิทฺธา\-รมฺมณา เจตนา.}

\csromanheader{2}{\textbf{นวิปากาทิ}}
\proseitem{7}{อชฺฌตฺตา\-รมฺมณํ ธมฺมํ ปฏิจฺจ อชฺฌตฺตา\-รมฺมโณ ธมฺโม อุปฺปชฺชติ นวิปาก\-ปจฺจยา. (ปฏิสนฺธิ นตฺถิ.) นฌานปจฺจยา ฯเปฯ. ปญฺจ\-วิญฺญาณสหคตํ อชฺฌตฺตา\-รมฺมณํ เอกํ ฯเปฯ. (1)}

\prose{พหิทฺธา\-รมฺมณํ ธมฺมํ ปฏิจฺจ พหิทฺธา\-รมฺมโณ ธมฺโม อุปฺปชฺชติ นฌานปจฺจยา. ปญฺจ\-วิญฺญาณสหคตํ พหิทฺธา\-รมฺมณํ เอกํ ขนฺธํ ปฏิจฺจ ตโย ขนฺธา ฯเปฯ. นมคฺคปจฺจยา. (นเหตุสทิโส. โมโห นตฺถิ.) นวิปฺปยุตฺต\-ปจฺจยา. อรูเป อชฺฌตฺตา\-รมฺมณํ เอกํ ขนฺธํ ปฏิจฺจ ตโย ขนฺธา ฯเปฯ. (1)}

\prose{พหิทฺธา\-รมฺมณํ ธมฺมํ ปฏิจฺจ พหิทฺธา\-รมฺมโณ ธมฺโม อุปฺปชฺชติ นวิปฺปยุตฺต\-ปจฺจยา. อรูเป พหิทฺธา\-รมฺมณํ เอกํ ขนฺธํ ปฏิจฺจ ตโย ขนฺธา ฯเปฯ. (1)}

\csromanmatikaanchor{mtk.259}
\csromantocmarkref{subparagraph}{2. ปจฺจยปจฺจนีย 2. สงฺขฺยาวาร}{mtk.259}
\bookmark[dest={mtk.259},level=5]{2. ปจฺจยปจฺจนีย 2. สงฺขฺยาวาร}
\csromanmarkh{2. ปจฺจย\-ปจฺจนีย}\csromanmarkhii{2. สงฺขฺยาวาร}\csromanheadernomark{6}{2. ปจฺจย\-ปจฺจนีย 2. สงฺขฺยาวาร}
\csromanmatikaanchor{mtk.260}
\csromanmatikaanchor{mtk.261}
\csromantocmarkref{subparagraph}{3. ปจฺจยานุโลมปจฺจนีย}{mtk.260}
\bookmark[dest={mtk.260},level=5]{3. ปจฺจยานุโลมปจฺจนีย}
\csromantocmarkref{subparagraph}{4. ปจฺจยปจฺจนียานุโลม}{mtk.261}
\bookmark[dest={mtk.261},level=5]{4. ปจฺจยปจฺจนียานุโลม}
\proseitemwithcloserruled{8}{\csromanfolio{456}นเหตุยา เทฺว, นอธิปติยา เทฺว, นปุเรชาเต เทฺว, นปจฺฉาชาเต เทฺว, นอาเสวเน, นกมฺเม, นวิปาเก, นฌาเน, นมคฺเค, นวิปฺปยุตฺเต เทฺว. (เอวํ คเณตพฺพํ.)}{\textbf{ปจฺจนียํ.}}
\csromancenter{\textbf{เหตุ}ปจฺจยา นอธิปติยา เทฺว ฯเปฯ นวิปาเก เทฺว, นวิปฺปยุตฺเต เทฺว. (เอวํ คเณตพฺพํ.)}
\nitthitamruled{อนุโลม\-ปจฺจนียํ.}
\proseitem{10}{นเหตุปจฺจยา อารมฺมเณ เทฺว, อนนฺตเร เทฺว, สมนนฺตเร เทฺว ฯเปฯ มคฺเค เทฺว ฯเปฯ อวิคเต เทฺว. (เอวํ คเณตพฺพํ.) ปจฺจนียา\-นุโลมํ. ปฏิจฺจวาโร.}

\prose{(สหชาตวาโรปิ ปจฺจยวาโรปิ นิสฺสยวาโรปิ สํสฏฺฐ\-วาโรปิ สมฺปยุตฺตวาโรปิ ปฏิจฺจ\-วาร\-สทิสา.)}
\csromansectionrule

\csromanmatikaanchor{mtk.262}
\csromantocmarknopage{subparagraph}{21. อชฺฌตฺตารมฺมณตฺติก 7. ปญฺหาวาร}{mtk.262}
\bookmark[dest={mtk.262},level=5]{21. อชฺฌตฺตารมฺมณตฺติก 7. ปญฺหาวาร}
\csromanmarkh{21. อชฺฌตฺตา\-รมฺมณ\-ตฺติก}\csromanmarkhii{7. ปญฺหาวาร}\csromanheadernomark{6}{21. อชฺฌตฺตา\-รมฺมณ\-ตฺติก 7. ปญฺหาวาร}
\csromanmarkh{1. ปจฺจยานุโลม}\csromanmarkhii{1. วิภงฺควาร}\csromanheadernomark{6}{1. \textbf{ปจฺจยานุโลม 1. วิภงฺควาร}}
\csromanmatikaanchor{mtk.263}
\csromanmatikaanchor{mtk.264}
\csromantocmarknopage{subparagraph}{1. ปจฺจยานุโลม 1. วิภงฺควาร}{mtk.263}
\bookmark[dest={mtk.263},level=5]{1. ปจฺจยานุโลม 1. วิภงฺควาร}
\csromantocmarkref{subparagraph}{เหตุ}{mtk.264}
\bookmark[dest={mtk.264},level=5]{เหตุ}
\csromanheader{6}{\textbf{เหตุ}}
\proseitem{11}{อชฺฌตฺตา\-รมฺมโณ ธมฺโม อชฺฌตฺตา\-รมฺมณสฺส ธมฺมสฺส เหตุปจฺจเยน ปจฺจโย. อชฺฌตฺตา\-รมฺมณา เหตู สมฺปยุตฺตกานํ ขนฺธานํ เหตุปจฺจเยน ปจฺจโย. ปฏิสนฺธิ\-กฺขเณ อชฺฌตฺตา\-รมฺมณา เหตู สมฺปยุตฺตกานํ ขนฺธานํ เหตุปจฺจเยน ปจฺจโย. (1)}

\csromanmatikaanchor{mtk.265}
\csromantocmarkref{subparagraph}{อารมฺมณาทิ}{mtk.265}
\bookmark[dest={mtk.265},level=5]{อารมฺมณาทิ}
\prose{\csromanfolio{457}พหิทฺธา\-รมฺมโณ ธมฺโม พหิทฺธา\-รมฺมณสฺส ธมฺมสฺส เหตุปจฺจเยน ปจฺจโย. พหิทฺธา\-รมฺมณา เหตู สมฺปยุตฺตกานํ ขนฺธานํ ฯเปฯ. ปฏิสนฺธิ\-กฺขเณ ฯเปฯ. (1)}

\csromanheader{2}{\textbf{อารมฺมณ}}
\proseitem{12}{อชฺฌตฺตา\-รมฺมโณ ธมฺโม อชฺฌตฺตา\-รมฺมณสฺส ธมฺมสฺส อารมฺมณ\-ปจฺจเยน ปจฺจโย. อชฺฌตฺตํ วิญฺญาณญฺจา\-ยตนํ ปจฺจเวกฺขติ, เนว\-สญฺญา\-นา\-สญฺญา\-ยตนํ ปจฺจเวกฺขติ, อชฺฌตฺตา\-รมฺมณํ อชฺฌตฺตํ ทิพฺพํ จกฺขุํ ปจฺจเวกฺขติ, ทิพฺพํ โสตธาตุํ ฯเปฯ. อิทฺธิวิธ\-ญาณํ ฯเปฯ. ปุพฺเพ\-นิวาสา\-นุสฺสติ\-ญาณํ ฯเปฯ. ยถากมฺมูปคญาณํ ฯเปฯ. อนาคตํสญาณํ ปจฺจเวกฺขติ. อริยา อชฺฌตฺตา\-รมฺมเณ ปหีเน กิเลเส ปจฺจเวกฺขนฺติ, วิกฺขมฺภิเต กิเลเส ปจฺจเวกฺขนฺติ. ปุพฺเพ สมุทาจิณฺเณ กิเลเส ชานนฺติ. อชฺฌตฺตา\-รมฺมเณ อชฺฌตฺเต ขนฺเธ อนิจฺจโต ฯเปฯ วิปสฺสติ, อสฺสาเทติ อภินนฺทติ, ตํ อารพฺภ อชฺฌตฺตา\-รมฺมโณ ราโค อุปฺปชฺชติ ฯเปฯ โทมนสฺสํ อุปฺปชฺชติ. อชฺฌตฺตา\-รมฺมณา อชฺฌตฺตา ขนฺธา ปุพฺเพ\-นิวาสา\-นุสฺสติ\-ญาณสฺส, ยถากมฺมูปคญาณสฺส, อนาคตํสญาณสฺส, อาวชฺชนาย อารมฺมณ\-ปจฺจเยน ปจฺจโย. (1)}

\prose{อชฺฌตฺตา\-รมฺมโณ ธมฺโม พหิทฺธา\-รมฺมณสฺส ธมฺมสฺส อารมฺมณ\-ปจฺจเยน ปจฺจโย. พหิทฺธา วิญฺญาณญฺจา\-ยตนํ ปจฺจเวกฺขติ, เนว\-สญฺญา\-นา\-สญฺญา\-ยตนํ ปจฺจเวกฺขติ. อชฺฌตฺตา\-รมฺมณํ พหิทฺธา ทิพฺพํ จกฺขุํ ปจฺจเวกฺขติ, ทิพฺพํ โสตธาตุํ ฯเปฯ. อิทฺธิวิธ\-ญาณํ ฯเปฯ. ปุพฺเพ\-นิวาสา\-นุสฺสติ\-ญาณํ ฯเปฯ. ยถากมฺมูปคญาณํ ฯเปฯ. อนาคตํสญาณํ ปจฺจเวกฺขติ, อชฺฌตฺตา\-รมฺมเณ พหิทฺธา ขนฺเธ อนิจฺจโต ทุกฺขโต อนตฺตโต วิปสฺสติ ฯเปฯ. เจโตปริย\-ญาเณน อชฺฌตฺตา\-รมฺมณ\-พหิทฺธา\-จิตฺต\-สมงฺคิสฺส จิตฺตํ ชานาติ, อชฺฌตฺตา\-รมฺมณา พหิทฺธา ขนฺธา เจโตปริย\-ญาณสฺส, ปุพฺเพ\-นิวาสา\-นุสฺสติ\-ญาณสฺส, ยถากมฺมูปคญาณสฺส, อนาคตํสญาณสฺส, อาวชฺชนาย อารมฺมณ\-ปจฺจเยน ปจฺจโย. (2)}

\proseitem{13}{พหิทฺธา\-รมฺมโณ ธมฺโม พหิทฺธา\-รมฺมณสฺส ธมฺมสฺส อารมฺมณ\-ปจฺจเยน ปจฺจโย. พหิทฺธา\-รมฺมณํ พหิทฺธา ทิพฺพํ จกฺขุํ ปจฺจเวกฺขติ, ทิพฺพํ โสตธาตุํ ปจฺจเวกฺขติ. อิทฺธิวิธ\-ญาณํ ฯเปฯ. เจโตปริย\-ญาณํ ฯเปฯ. ปุพฺเพ\-นิวาสา\-นุสฺสติ\-ญาณํ ฯเปฯ. ยถากมฺมูปคญาณํ ฯเปฯ. อนาคตํสญาณํ ปจฺจเวกฺขติ. พหิทฺธา\-รมฺมเณ พหิทฺธา ขนฺเธ อนิจฺจโต ทุกฺขโต อนตฺตโต วิปสฺสติ ฯเปฯ. เจโตปริย\-ญาเณน พหิทฺธาร\-มฺมณ\-พหิทฺธา\-จิตฺต\-สมงฺคิสฺส จิตฺตํ ชานาติ. \csromanfolio{458}พหิทฺธา\-รมฺมณา พหิทฺธา ขนฺธา เจโตปริย\-ญาณสฺส, ปุพฺเพ\-นิวาสา\-นุสฺสติ\-ญาณสฺส, ยถากมฺมูปคญาณสฺส, อนาคตํสญาณสฺส, อาวชฺชนาย อารมฺมณ\-ปจฺจเยน ปจฺจโย. (1)}

\prose{พหิทฺธา\-รมฺมโณ ธมฺโม อชฺฌตฺตา\-รมฺมณสฺส ธมฺมสฺส อารมฺมณ\-ปจฺจเยน ปจฺจโย. ทานํ ทตฺวา, สีลํ สมาทิยิตฺวา, อุโปสถกมฺมํ กตฺวา ตํ ปจฺจเวกฺขติ. ปุพฺเพ สุจิณฺณานิ ปจฺจเวกฺขติ. ฌานา วุฏฺฐหิตฺวา ฌานํ ปจฺจเวกฺขติ. อริยา มคฺคา วุฏฺฐหิตฺวา มคฺคํ ปจฺจเวกฺขนฺติ, ผลํ ปจฺจเวกฺขนฺติ, พหิทฺธา\-รมฺมเณ ปหีเน กิเลเส ปจฺจเวกฺขนฺติ, วิกฺขมฺภิเต กิเลเส ปจฺจเวกฺขนฺติ, ปุพฺเพ สมุทาจิณฺเณ กิเลเส ชานนฺติ. พหิทฺธา\-รมฺมณํ อชฺฌตฺตํ ทิพฺพํ จกฺขุํ ปจฺจเวกฺขติ, ทิพฺพํ โสตธาตุํ. อิทฺธิวิธ\-ญาณํ. เจโตปริย\-ญาณํ. ปุพฺเพ\-นิวาสา\-นุสฺสติ\-ญาณํ. ยถากมฺมูปคญาณํ. อนาคตํสญาณํ. พหิทฺธา\-รมฺมเณ อชฺฌตฺเต ขนฺเธ อนิจฺจโต ฯเปฯ วิปสฺสติ, อสฺสาเทติ อภินนฺทติ, ตํ อารพฺภ อชฺฌตฺตา\-รมฺมโณ ราโค อุปฺปชฺชติ ฯเปฯ โทมนสฺสํ อุปฺปชฺชติ. พหิทฺธา\-รมฺมณา อชฺฌตฺตา ขนฺธา อิทฺธิวิธ\-ญาณสฺส, ปุพฺเพ\-นิวาสา\-นุสฺสติ\-ญาณสฺส, ยถากมฺมูปคญาณสฺส, อนาคตํสญาณสฺส, อาวชฺชนาย อารมฺมณ\-ปจฺจเยน ปจฺจโย. (2)}

\proseitem{14}{อชฺฌตฺตา\-รมฺมโณ ธมฺโม อชฺฌตฺตา\-รมฺมณสฺส ธมฺมสฺส อธิปติ\-ปจฺจเยน ปจฺจโย. อารมฺมณา\-ธิปติ, สหชาตา\-ธิปติ.  อารมฺมณา\-ธิปติ\csromandash{}อชฺฌตฺตํ วิญฺญาณญฺจา\-ยตนํ ครุํ กตฺวา ปจฺจเวกฺขติ, เนว\-สญฺญา\-นา\-สญฺญา\-ยตนํ ครุํ กตฺวา ปจฺจเวกฺขติ. อชฺฌตฺตา\-รมฺมณํ อชฺฌตฺตํ ทิพฺพํ จกฺขุํ ครุํ กตฺวา ฯเปฯ ทิพฺพํ โสตธาตุํ ฯเปฯ. อิทฺธิวิธ\-ญาณํ. ปุพฺเพ\-นิวาสา\-นุสฺสติ\-ญาณํ. ยถากมฺมูปคญาณํ. อนาคตํสญาณํ ครุํ กตฺวา ฯเปฯ. อชฺฌตฺตา\-รมฺมเณ อชฺฌตฺเต ขนฺเธ ครุํ กตฺวา อสฺสาเทติ อภินนฺทติ, ตํ ครุํ กตฺวา อชฺฌตฺตา\-รมฺมโณ ราโค อุปฺปชฺชติ ทิฏฺฐิ อุปฺปชฺชติ. สหชาตา\-ธิปติ\csromandash{}อชฺฌตฺตา\-รมฺมณา\-ธิปติ สมฺปยุตฺตกานํ ขนฺธานํ อธิปติ\-ปจฺจเยน ปจฺจโย. (1)}

\proseitem{15}{พหิทฺธา\-รมฺมโณ ธมฺโม พหิทฺธา\-รมฺมณสฺส ธมฺมสฺส อธิปติ\-ปจฺจเยน ปจฺจโย. สหชาตา\-ธิปติ\csromandash{}พหิทฺธา\-รมฺมณา\-ธิปติ สมฺปยุตฺตกานํ ขนฺธานํ อธิปติ\-ปจฺจเยน ปจฺจโย. (1)}

\prose{\csromanfolio{459}พหิทฺธา\-รมฺมโณ ธมฺโม อชฺฌตฺตา\-รมฺมณสฺส ธมฺมสฺส อธิปติ\-ปจฺจเยน ปจฺจโย. อารมฺมณา\-ธิปติ\csromandash{}ทานํ ทตฺวา, สีลํ สมาทิยิตฺวา, อุโปสถกมฺมํ กตฺวา ตํ ครุํ กตฺวา ปจฺจเวกฺขติ. ปุพฺเพ สุจิณฺณานิ ฯเปฯ. ฌานา วุฏฺฐหิตฺวา ฯเปฯ. อริยา มคฺคา วุฏฺฐหิตฺวา มคฺคํ. ผลํ ครุํ กตฺวา ปจฺจเวกฺขนฺติ. พหิทฺธา\-รมฺมณํ อชฺฌตฺตํ ทิพฺพํ จกฺขุํ ครุํ กตฺวา ฯเปฯ ทิพฺพํ โสตธาตุํ. อิทฺธิวิธ\-ญาณํ. เจโตปริย\-ญาณํ. ปุพฺเพ\-นิวาสา\-นุสฺสติ\-ญาณํ. ยถากมฺมูปคญาณํ. อนาคตํสญาณํ ครุํ กตฺวา ปจฺจเวกฺขติ. พหิทฺธา\-รมฺมเณ อชฺฌตฺเต ขนฺเธ ครุํ กตฺวา อสฺสาเทติ อภินนฺทติ, ตํ ครุํ กตฺวา อชฺฌตฺตา\-รมฺมโณ ราโค อุปฺปชฺชติ. ทิฏฺฐิ อุปฺปชฺชติ. (2)}

\csromanheader{2}{\textbf{อนนฺตร}}
\proseitem{16}{อชฺฌตฺตา\-รมฺมโณ ธมฺโม อชฺฌตฺตา\-รมฺมณสฺส ธมฺมสฺส อนนฺตร\-ปจฺจเยน ปจฺจโย. ปุริมา ปุริมา อชฺฌตฺตา\-รมฺมณา ขนฺธา ปจฺฉิมานํ ปจฺฉิมานํ อชฺฌตฺตา\-รมฺมณานํ ขนฺธานํ อนนฺตร\-ปจฺจเยน ปจฺจโย. (1)}

\prose{อชฺฌตฺตา\-รมฺมโณ ธมฺโม พหิทฺธา\-รมฺมณสฺส ธมฺมสฺส อนนฺตร\-ปจฺจเยน ปจฺจโย. อชฺฌตฺตา\-รมฺมณํ จุติจิตฺตํ พหิทฺธา\-รมฺมณสฺส อุปปตฺติ\-จิตฺตสฺส อนนฺตร\-ปจฺจเยน ปจฺจโย. อชฺฌตฺตา\-รมฺมณํ ภวงฺคํ พหิทฺธา\-รมฺมณาย อาวชฺชนาย อนนฺตร\-ปจฺจเยน ปจฺจโย. อชฺฌตฺตา\-รมฺมณา ขนฺธา พหิทฺธา\-รมฺมณสฺส วุฏฺฐนสฺส อนนฺตร\-ปจฺจเยน ปจฺจโย. อชฺฌตฺตา\-รมฺมณํ อนุโลมํ โคตฺรภุสฺส. อนุโลมํ โวทานสฺส. อนุโลมํ ผลสมาปตฺติยา. นิโรธา วุฏฺฐหนฺตสฺส. เนว\-สญฺญา\-นา\-สญฺญา\-ยตนํ ผลสมาปตฺติยา อนนฺตร\-ปจฺจเยน ปจฺจโย. (2)}

\proseitem{17}{พหิทฺธา\-รมฺมโณ ธมฺโม พหิทฺธา\-รมฺมณสฺส ธมฺมสฺส อนนฺตร\-ปจฺจเยน ปจฺจโย. ปุริมา ปุริมา พหิทฺธา\-รมฺมณา ขนฺธา ปจฺฉิมานํ ปจฺฉิมานํ พหิทฺธา\-รมฺมณานํ ขนฺธานํ อนนฺตร\-ปจฺจเยน ปจฺจโย. พหิทฺธา\-รมฺมณํ อนุโลมํ โคตฺรภุสฺส. อนุโลมํ โวทานสฺส. โคตฺรภุ มคฺคสฺส. โวทานํ มคฺคสฺส. มคฺโค ผลสฺส. ผลํ ผลสฺส. อนุโลมํ ผลสมาปตฺติยา อนนฺตร\-ปจฺจเยน ปจฺจโย. (1)}

\prose{พหิทฺธา\-รมฺมโณ ธมฺโม อชฺฌตฺตา\-รมฺมณสฺส ธมฺมสฺส อนนฺตร\-ปจฺจเยน ปจฺจโย. พหิทฺธา\-รมฺมณํ จุติจิตฺตํ อชฺฌตฺตา\-รมฺมณสฺส อุปปตฺติ\-จิตฺตสฺส อนนฺตร\-ปจฺจเยน ปจฺจโย. พหิทฺธา\-รมฺมณํ ภวงฺคํ อชฺฌตฺตา\-รมฺมณาย อาวชฺชนาย \csromanfolio{460}อนนฺตร\-ปจฺจเยน ปจฺจโย. พหิทฺธา\-รมฺมณา ขนฺธา อชฺฌตฺตา\-รมฺมณสฺส วุฏฺฐานสฺส อนนฺตร\-ปจฺจเยน ปจฺจโย. (2)}

\csromanheader{2}{\textbf{สมนนฺตราทิ}}
\proseitem{17}{อชฺฌตฺตา\-รมฺมโณ ธมฺโม อชฺฌตฺตา\-รมฺมณสฺส ธมฺมสฺส สมนนฺตร\-ปจฺจเยน ปจฺจโย ฯเปฯ. สหชาต\-ปจฺจเยน ปจฺจโย. อญฺญมญฺญ\-ปจฺจเยน ปจฺจโย. นิสฺสย\-ปจฺจเยน ปจฺจโย.}

\csromanheader{2}{\textbf{อุปนิสฺสย}}
\proseitem{17}{อชฺฌตฺตา\-รมฺมโณ ธมฺโม อชฺฌตฺตา\-รมฺมณสฺส ธมฺมสฺส อุปนิสฺสย\-ปจฺจเยน ปจฺจโย. อารมฺมณู\-ปนิสฺสโย, อนนฺตรู\-ปนิสฺสโย, ปกตู\-ปนิสฺสโย ฯเปฯ. ปกตู\-ปนิสฺสโย\csromandash{}อชฺฌตฺตา\-รมฺมณา อนิจฺจา\-นุปสฺสนา, ทุกฺขา\-นุปสฺสนา, อนตฺตา\-นุปสฺสนา อชฺฌตฺตา\-รมฺมณาย อนิจฺจา\-นุปสฺสนาย, ทุกฺขา\-นุปสฺสนาย, อนตฺตา\-นุปสฺสนาย อุปนิสฺสย\-ปจฺจเยน ปจฺจโย. (1)}

\prose{อชฺฌตฺตา\-รมฺมโณ ธมฺโม พหิทฺธา\-รมฺมณสฺส ธมฺมสฺส อุปนิสฺสย\-ปจฺจเยน ปจฺจโย. อนนฺตรู\-ปนิสฺสโย, ปกตู\-ปนิสฺสโย ฯเปฯ. ปกตู\-ปนิสฺสโย\csromandash{}อชฺฌตฺตา\-รมฺมณา อนิจฺจา\-นุปสฺสนา, ทุกฺขา\-นุปสฺสนา, อนตฺตา\-นุปสฺสนา พหิทฺธา\-รมฺมณาย อนิจฺจา\-นุปสฺสนาย, ทุกฺขา\-นุปสฺสนาย, อนตฺตา\-นุปสฺสนาย อุปนิสฺสย\-ปจฺจเยน ปจฺจโย. (2)}

\prose{พหิทฺธา\-รมฺมโณ ธมฺโม พหิทฺธา\-รมฺมณสฺส ธมฺมสฺส อุปนิสฺสย\-ปจฺจเยน ปจฺจโย. อนนฺตรู\-ปนิสฺสโย, ปกตู\-ปนิสฺสโย ฯเปฯ. ปกตู\-ปนิสฺสโย\csromandash{}พหิทฺธา\-รมฺมณา อนิจฺจา\-นุปสฺสนา, ทุกฺขา\-นุปสฺสนา, อนตฺตา\-นุปสฺสนา พหิทฺธา\-รมฺมณาย อนิจฺจา\-นุปสฺสนาย, ทุกฺขา\-นุปสฺสนาย, อนตฺตา\-นุปสฺสนาย อุปนิสฺสย\-ปจฺจเยน ปจฺจโย. (1)}

\prose{พหิทฺธา\-รมฺมโณ ธมฺโม อชฺฌตฺตา\-รมฺมณสฺส ธมฺมสฺส อุปนิสฺสย\-ปจฺจเยน ปจฺจโย. อารมฺมณู\-ปนิสฺสโย, อนนฺตรู\-ปนิสฺสโย, ปกตู\-ปนิสฺสโย ฯเปฯ. ปกตู\-ปนิสฺสโย\csromandash{}พหิทฺธา\-รมฺมณา อนิจฺจา\-นุปสฺสนา, ทุกฺขา\-นุปสฺสนา, อนตฺตา\-นุปสฺสนา อชฺฌตฺตา\-รมฺมณาย อนิจฺจา\-นุปสฺสนาย, ทุกฺขา\-นุปสฺสนาย, อนตฺตา\-นุปสฺสนาย อุปนิสฺสย\-ปจฺจเยน ปจฺจโย. (2)}

\csromanheader{2}{21. อชฺฌตฺตา\-รมฺมณ\-ตฺติก อาเสวน}
\proseitem{21}{\csromanfolio{461}อชฺฌตฺตา\-รมฺมโณ ธมฺโม อชฺฌตฺตา\-รมฺมณสฺส ธมฺมสฺส อาเสวน\-ปจฺจเยน ปจฺจโย. ปุริมา ปุริมา อชฺฌตฺตา\-รมฺมณา ขนฺธา ปจฺฉิมานํ ปจฺฉิมานํ อชฺฌตฺตา\-รมฺมณานํ ขนฺธานํ อาเสวน\-ปจฺจเยน ปจฺจโย. (1)}

\prose{อชฺฌตฺตา\-รมฺมโณ ธมฺโม พหิทฺธา\-รมฺมณสฺส ธมฺมสฺส อาเสวน\-ปจฺจเยน ปจฺจโย. อชฺฌตฺตา\-รมฺมณํ อนุโลมํ โคตฺรภุสฺส. อนุโลมํ โวทานสฺส อาเสวน\-ปจฺจเยน ปจฺจโย. (2)}

\prose{พหิทฺธา\-รมฺมโณ ธมฺโม พหิทฺธา\-รมฺมณสฺส ธมฺมสฺส อาเสวน\-ปจฺจเยน ปจฺจโย. ปุริมา ปุริมา พหิทฺธา\-รมฺมณา ขนฺธา ปจฺฉิมานํ ปจฺฉิมานํ พหิทฺธา\-รมฺมณานํ ขนฺธานํ อาเสวน\-ปจฺจเยน ปจฺจโย. พหิทฺธา\-รมฺมณํ อนุโลมํ โคตฺรภุสฺส. อนุโลมํ โวทานสฺส. โคตฺรภุ มคฺคสฺส. โวทานํ มคฺคสฺส อาเสวน\-ปจฺจเยน ปจฺจโย. (1)}

\csromanheader{2}{\textbf{กมฺม}}
\proseitem{22}{อชฺฌตฺตา\-รมฺมโณ ธมฺโม อชฺฌตฺตา\-รมฺมณสฺส ธมฺมสฺส กมฺมปจฺจเยน ปจฺจโย. สหชาตา, นานากฺขณิกา.  สหชาตา อชฺฌตฺตา\-รมฺมณา เจตนา สมฺปยุตฺตกานํ ขนฺธานํ กมฺมปจฺจเยน ปจฺจโย. ปฏิสนฺธิ\-กฺขเณ ฯเปฯ. นานากฺขณิกา อชฺฌตฺตา\-รมฺมณา เจตนา วิปากานํ อชฺฌตฺตา\-รมฺมณานํ ขนฺธานํ กมฺมปจฺจเยน ปจฺจโย. (1)}

\prose{อชฺฌตฺตา\-รมฺมโณ ธมฺโม พหิทฺธา\-รมฺมณสฺส ธมฺมสฺส กมฺมปจฺจเยน ปจฺจโย. นานากฺขณิกา อชฺฌตฺตา\-รมฺมณา เจตนา วิปากานํ พหิทฺธา\-รมฺมณานํ ขนฺธานํ กมฺมปจฺจเยน ปจฺจโย. (2)}

\prose{พหิทฺธา\-รมฺมโณ ธมฺโม พหิทฺธา\-รมฺมณสฺส ธมฺมสฺส กมฺมปจฺจเยน ปจฺจโย. สหชาตา, นานากฺขณิกา.  สหชาตา พหิทฺธา\-รมฺมณา เจตนา สมฺปยุตฺตกานํ ขนฺธานํ กมฺมปจฺจเยน ปจฺจโย. ปฏิสนฺธิ\-กฺขเณ ฯเปฯ. นานากฺขณิกา พหิทฺธา\-รมฺมณา เจตนา วิปากานํ พหิทฺธา\-รมฺมณานํ ขนฺธานํ กมฺมปจฺจเยน ปจฺจโย. (1)}

\prose{พหิทฺธา\-รมฺมโณ ธมฺโม อชฺฌตฺตา\-รมฺมณสฺส ธมฺมสฺส กมฺมปจฺจเยน ปจฺจโย. นานากฺขณิกา พหิทฺธา\-รมฺมณา เจตนา วิปากานํ อชฺฌตฺตา\-รมฺมณานํ ขนฺธานํ กมฺมปจฺจเยน ปจฺจโย. (2)}

\csromanheader{2}{\textbf{วิปากาทิ}}
\csromanmatikaanchor{mtk.266}
\csromanmatikaanchor{mtk.267}
\csromantocmarknopage{subparagraph}{1. ปจฺจยานุโลม 2. สงฺขฺยาวาร}{mtk.266}
\bookmark[dest={mtk.266},level=5]{1. ปจฺจยานุโลม 2. สงฺขฺยาวาร}
\csromantocmarkref{subparagraph}{สุทฺธ}{mtk.267}
\bookmark[dest={mtk.267},level=5]{สุทฺธ}
\proseitem{23}{\csromanfolio{462}อชฺฌตฺตา\-รมฺมโณ ธมฺโม อชฺฌตฺตา\-รมฺมณสฺส ธมฺมสฺส วิปาก\-ปจฺจเยน ปจฺจโย ฯเปฯ. อาหารปจฺจเยน ปจฺจโย. อินฺทฺริย\-ปจฺจเยน ปจฺจโย. ฌานปจฺจเยน ปจฺจโย. มคฺคปจฺจเยน ปจฺจโย. สมฺปยุตฺต\-ปจฺจเยน ปจฺจโย. อตฺถิปจฺจเยน ปจฺจโย. นตฺถิ\-ปจฺจเยน ปจฺจโย. วิคต\-ปจฺจเยน ปจฺจโย. อวิคต\-ปจฺจเยน ปจฺจโย.}

\csromanmarkh{1. ปจฺจยานุโลม}\csromanmarkhii{2. สงฺขฺยาวาร}\csromanheadernomark{2}{1. \textbf{ปจฺจยานุโลม 2. สงฺขฺยาวาร}}
\proseitemwithcloserruled{24}{เหตุยา เทฺว, อารมฺมเณ จตฺตาริ, อธิปติยา ตีณิ, อนนฺตเร จตฺตาริ, สมนนฺตเร จตฺตาริ, สหชาเต เทฺว, อญฺญมญฺเญ เทฺว, นิสฺสเย เทฺว, อุปนิสฺสเย จตฺตาริ, อาเสวเน ตีณิ, กมฺเม จตฺตาริ, วิปาเก เทฺว ฯเปฯ (สพฺพตฺถ เทฺว,) สมฺปยุตฺเต เทฺว, อตฺถิยา เทฺว, นตฺถิยา จตฺตาริ, วิคเต จตฺตาริ, อวิคเต เทฺว. (เอวํ คเณตพฺพํ.)}{อนุโลมํ.}
\csromanmatikaanchor{mtk.268}
\csromantocmarkref{subparagraph}{2. ปจฺจนียุทฺธาร}{mtk.268}
\bookmark[dest={mtk.268},level=5]{2. ปจฺจนียุทฺธาร}
\csromanheader{6}{2. ปจฺจนียุ\-ทฺธาร}
\proseitem{25}{อชฺฌตฺตา\-รมฺมโณ ธมฺโม อชฺฌตฺตา\-รมฺมณสฺส ธมฺมสฺส อารมฺมณ\-ปจฺจเยน ปจฺจโย. สหชาต\-ปจฺจเยน ปจฺจโย. อุปนิสฺสย\-ปจฺจเยน ปจฺจโย. กมฺมปจฺจเยน ปจฺจโย. (1)}

\prose{อชฺฌตฺตา\-รมฺมโณ ธมฺโม พหิทฺธา\-รมฺมณสฺส ธมฺมสฺส อารมฺมณ\-ปจฺจเยน ปจฺจโย. อุปนิสฺสย\-ปจฺจเยน ปจฺจโย. กมฺมปจฺจเยน ปจฺจโย. (2)}

\prose{พหิทฺธา\-รมฺมโณ ธมฺโม พหิทฺธา\-รมฺมณสฺส ธมฺมสฺส อารมฺมณ\-ปจฺจเยน ปจฺจโย. สหชาต\-ปจฺจเยน ปจฺจโย. อุปนิสฺสย\-ปจฺจเยน ปจฺจโย. กมฺมปจฺจเยน ปจฺจโย. (1)}

\prose{พหิทฺธา\-รมฺมโณ ธมฺโม อชฺฌตฺตา\-รมฺมณสฺส ธมฺมสฺส อารมฺมณ\-ปจฺจเยน ปจฺจโย. อุปนิสฺสย\-ปจฺจเยน ปจฺจโย. กมฺมปจฺจเยน ปจฺจโย. (2)}

\csromanmatikaanchor{mtk.269}
\csromantocmarkref{subparagraph}{2. ปจฺจยปจฺจนีย 2. สงฺขฺยาวาร}{mtk.269}
\bookmark[dest={mtk.269},level=5]{2. ปจฺจยปจฺจนีย 2. สงฺขฺยาวาร}
\csromanmarkh{21. อชฺฌตฺตา\-รมฺมณ\-ตฺติก}\csromanmarkhii{2. ปจฺจย\-ปจฺจนีย 2. สงฺขฺยาวาร}\csromanheadernomark{6}{21. อชฺฌตฺตา\-รมฺมณ\-ตฺติก 2. ปจฺจย\-ปจฺจนีย 2. สงฺขฺยาวาร}
\csromanmatikaanchor{mtk.270}
\csromanmatikaanchor{mtk.271}
\csromantocmarkref{subparagraph}{3. ปจฺจยานุโลมปจฺจนีย}{mtk.270}
\bookmark[dest={mtk.270},level=5]{3. ปจฺจยานุโลมปจฺจนีย}
\csromantocmarkref{subparagraph}{4. ปจฺจยปจฺจนียานุโลม}{mtk.271}
\bookmark[dest={mtk.271},level=5]{4. ปจฺจยปจฺจนียานุโลม}
\proseitemwithcloserruled{26}{\csromanfolio{463}นเหตุยา จตฺตาริ, นอารมฺมเณ จตฺตาริ, นอธิปติยา จตฺตาริ, นอนนฺตเร จตฺตาริ. (สํขิตฺตํ. สพฺพตฺถ จตฺตาริ,) นปุเรชาเต, นปจฺฉาชาเต, นอาเสวเน ฯเปฯ นวิปฺปยุตฺเต จตฺตาริ ฯเปฯ โนอวิคเต จตฺตาริ. (เอวํ คเณตพฺพํ.)}{ปจฺจนียํ.}
\proseitemwithcloserruled{27}{เหตุปจฺจยา นอารมฺมเณ เทฺว, นอธิปติยา เทฺว, นอนนฺตเร, นสมนนฺตเร, นอุปนิสฺสเย, นอาเสวเน, นกมฺเม ฯเปฯ โนนตฺถิยา, โนวิคเต เทฺว. (สพฺพตฺถ เทฺว.) (เอวํ คเณตพฺพํ.)}{อนุโลม\-ปจฺจนียํ.}
\proseitem{28}{นเหตุปจฺจยา อารมฺมเณ จตฺตาริ, อธิปติยา ตีณิ, อนนฺตเร จตฺตาริ, สมนนฺตเร จตฺตาริ, สหชาเต, อญฺญมญฺเญ, นิสฺสเย เทฺว, อุปนิสฺสเย จตฺตาริ, อาเสวเน ตีณิ, กมฺเม จตฺตาริ, วิปาเก เทฺว ฯเปฯ สมฺปยุตฺเต เทฺว, อตฺถิยา เทฺว, นตฺถิยา จตฺตาริ, วิคเต จตฺตาริ, อวิคเต เทฺว. (เอวํ คเณตพฺพํ.)}

\vspace{\tipitakaparskip}
\csromancenter{ปจฺจนียา\-นุโลมํ.}
\csromancenter{ปญฺหาวาโร.}
\nitthitam{อชฺฌตฺตา\-รมฺมณ\-ตฺติกํ นิฏฺฐิตํ.}
\csromanmatikaanchor{mtk.272}
\csromantocmarknopage{subparagraph}{22. สนิทสฺสนสปฺปฏิฆตฺติก 1. ปฏิจฺจวาร}{mtk.272}
\bookmark[dest={mtk.272},level=5]{22. สนิทสฺสนสปฺปฏิฆตฺติก 1. ปฏิจฺจวาร}
\csromanmarkh{22. สนิทสฺสนสปฺปฏิฆ\-ตฺติก}\csromanmarkhii{1. ปฏิจฺจวาร}\csromanheadernomark{6}{22. สนิทสฺสนสปฺปฏิฆ\-ตฺติก 1. ปฏิจฺจวาร}
\csromanmarkh{1. ปจฺจยานุโลม}\csromanmarkhii{1. วิภงฺควาร}\csromanheadernomark{2}{1. \textbf{ปจฺจยานุโลม 1. วิภงฺควาร}}
\csromanmatikaanchor{mtk.273}
\csromanmatikaanchor{mtk.274}
\csromantocmarknopage{subparagraph}{1. ปจฺจยานุโลม 1. วิภงฺควาร}{mtk.273}
\bookmark[dest={mtk.273},level=5]{1. ปจฺจยานุโลม 1. วิภงฺควาร}
\csromantocmarkref{subparagraph}{เหตุ}{mtk.274}
\bookmark[dest={mtk.274},level=5]{เหตุ}
\csromanheader{6}{\textbf{เหตุ}}
\proseitem{1}{\csromanfolio{464}อนิทสฺสน\-สปฺปฏิฆํ ธมฺมํ ปฏิจฺจ อนิทสฺสน\-สปฺปฏิโฆ ธมฺโม อุปฺปชฺชติ เหตุปจฺจยา. อนิทสฺสน\-สปฺปฏิฆํ เอกํ มหาภูตํ ปฏิจฺจ เทฺว มหาภูตา. เทฺว มหาภูเต ปฏิจฺจ เอกํ มหาภูตํ. อนิทสฺสน\-สปฺปฏิเฆ มหาภูเต ปฏิจฺจ อนิทสฺสน\-สปฺปฏิฆํ จิตฺต\-สมุฏฺฐานํ รูปํ, กฏตฺตารูปํ, อุปาทารูปํ. โผฏฺฐพฺพา\-ยตนํ ปฏิจฺจ จกฺขายตนํ ฯเปฯ รสายตนํ ฯเปฯ. (1)}

\prose{อนิทสฺสน\-สปฺปฏิฆํ ธมฺมํ ปฏิจฺจ สนิทสฺสน\-สปฺปฏิโฆ ธมฺโม อุปฺปชฺชติ เหตุปจฺจยา. อนิทสฺสน\-สปฺปฏิเฆ มหาภูเต ปฏิจฺจ สนิทสฺสน\-สปฺปฏิฆํ จิตฺต\-สมุฏฺฐานํ รูปํ, กฏตฺตารูปํ, อุปาทารูปํ. โผฏฺฐพฺพา\-ยตนํ ปฏิจฺจ รูปายตนํ. (2)}

\prose{อนิทสฺสน\-สปฺปฏิฆํ ธมฺมํ ปฏิจฺจ อนิทสฺสน\-อปฺปฏิโฆ ธมฺโม อุปฺปชฺชติ เหตุปจฺจยา. อนิทสฺสน\-สปฺปฏิเฆ มหาภูเต ปฏิจฺจ อนิทสฺสน\-อปฺปฏิฆํ จิตฺต\-สมุฏฺฐานํ รูปํ, กฏตฺตารูปํ, อุปาทารูปํ. โผฏฺฐพฺพา\-ยตนํ ปฏิจฺจ อาโปธาตุ, อิตฺถินฺทฺริยํ ฯเปฯ. กพฬีกาโร อาหาโร. (3)}

\prose{อนิทสฺสน\-สปฺปฏิฆํ ธมฺมํ ปฏิจฺจ สนิทสฺสน\-สปฺปฏิโฆ จ อนิทสฺสน\-อปฺปฏิโฆ จ ธมฺมา อุปฺปชฺชนฺติ เหตุปจฺจยา. อนิทสฺสน\-สปฺปฏิเฆ มหาภูเต ปฏิจฺจ สนิทสฺสนสปฺปฏิฆญฺจ อนิทสฺสนอปฺปฏิฆญฺจ จิตฺต\-สมุฏฺฐานํ รูปํ, กฏตฺตารูปํ, อุปาทารูปํ. โผฏฺฐพฺพา\-ยตนํ ปฏิจฺจ รูปายตนํ, อาโปธาตุ, อิตฺถินฺทฺริยํ ฯเปฯ. กพฬีกาโร อาหาโร. (4)}

\prose{อนิทสฺสน\-สปฺปฏิฆํ ธมฺมํ ปฏิจฺจ อนิทสฺสน\-สปฺปฏิโฆ จ อนิทสฺสน\-อปฺปฏิโฆ จ ธมฺมา อุปฺปชฺชนฺติ เหตุปจฺจยา. อนิทสฺสน\-สปฺปฏิฆํ เอกํ มหาภูตํ ปฏิจฺจ เทฺว มหาภูตา อาโปธาตุ จ, เทฺว มหาภูเต ปฏิจฺจ เอกํ มหาภูตํ อาโปธาตุ จ. อนิทสฺสน\-สปฺปฏิเฆ มหาภูเต ปฏิจฺจ อนิทสฺสนสปฺปฏิฆญฺจ อนิทสฺสนอปฺปฏิฆญฺจ จิตฺต\-สมุฏฺฐานํ รูปํ, กฏตฺตารูปํ, อุปาทารูปํ. โผฏฺฐพฺพา\-ยตนํ ปฏิจฺจ จกฺขายตนํ ฯเปฯ รสายตนํ, อาโปธาตุ, อิตฺถินฺทฺริยํ ฯเปฯ. กพฬีกาโร อาหาโร. (5)}

\prose{\csromanfolio{465}อนิทสฺสน\-สปฺปฏิฆํ ธมฺมํ ปฏิจฺจ สนิทสฺสน\-สปฺปฏิโฆ จ อนิทสฺสน\-สปฺปฏิโฆ จ ธมฺมา อุปฺปชฺชนฺติ เหตุปจฺจยา. อนิทสฺสน\-สปฺปฏิเฆ มหาภูเต ปฏิจฺจ สนิทสฺสนสปฺปฏิฆญฺจ อนิทสฺสนสปฺปฏิฆญฺจ จิตฺต\-สมุฏฺฐานํ รูปํ, กฏตฺตารูปํ, อุปาทารูปํ. โผฏฺฐพฺพา\-ยตนํ ปฏิจฺจ รูปายตนํ, จกฺขายตนํ ฯเปฯ รสายตนํ. (6)}

\prose{อนิทสฺสน\-สปฺปฏิฆํ ธมฺมํ ปฏิจฺจ สนิทสฺสน\-สปฺปฏิโฆ จ อนิทสฺสน\-สปฺปฏิโฆ จ อนิทสฺสน\-อปฺปฏิโฆ จ ธมฺมา อุปฺปชฺชนฺติ เหตุปจฺจยา. อนิทสฺสน\-สปฺปฏิเฆ มหาภูเต ปฏิจฺจ สนิทสฺสนสปฺปฏิฆญฺจ อนิทสฺสนสปฺปฏิฆญฺจ อนิทสฺสนอปฺปฏิฆญฺจ จิตฺต\-สมุฏฺฐานํ รูปํ, กฏตฺตารูปํ, อุปาทารูปํ. โผฏฺฐพฺพา\-ยตนํ ปฏิจฺจ รูปายตนํ จกฺขายตนํ ฯเปฯ รสายตนํ, อาโปธาตุ, อิตฺถินฺทฺริยํ ฯเปฯ. กพฬีกาโร อาหาโร. (7)}

\proseitem{2}{อนิทสฺสน\-อปฺปฏิฆํ ธมฺมํ ปฏิจฺจ อนิทสฺสน\-อปฺปฏิโฆ ธมฺโม อุปฺปชฺชติ เหตุปจฺจยา. อนิทสฺสน\-อปฺปฏิฆํ เอกํ ขนฺธํ ปฏิจฺจ ตโย ขนฺธา อนิทสฺสนอปฺปฏิฆญฺจ จิตฺต\-สมุฏฺฐานํ รูปํ ฯเปฯ เทฺว ขนฺเธ ปฏิจฺจ เทฺว ขนฺธา อนิทสฺสนอปฺปฏิฆญฺจ จิตฺต\-สมุฏฺฐานํ รูปํ. ปฏิสนฺธิ\-กฺขเณ อนิทสฺสน\-อปฺปฏิฆํ เอกํ ขนฺธํ ปฏิจฺจ ตโย ขนฺธา อนิทสฺสนอปฺปฏิฆญฺจ กฏตฺตารูปํ ฯเปฯ เทฺว ขนฺเธ ฯเปฯ. ขนฺเธ ปฏิจฺจ วตฺถุ, วตฺถุํ ปฏิจฺจ ขนฺธา. อาโปธาตุํ ปฏิจฺจ อนิทสฺสน\-อปฺปฏิฆํ จิตฺต\-สมุฏฺฐานํ รูปํ, กฏตฺตารูปํ, อุปาทารูปํ. อาโปธาตุํ ปฏิจฺจ อิตฺถินฺทฺริยํ ฯเปฯ. กพฬีกาโร อาหาโร. (1)}

\prose{อนิทสฺสน\-อปฺปฏิฆํ ธมฺมํ ปฏิจฺจ สนิทสฺสน\-สปฺปฏิโฆ ธมฺโม อุปฺปชฺชติ เหตุปจฺจยา. อนิทสฺสน\-อปฺปฏิเฆ ขนฺเธ ปฏิจฺจ สนิทสฺสน\-สปฺปฏิฆํ จิตฺต\-สมุฏฺฐานํ รูปํ. ปฏิสนฺธิ\-กฺขเณ อนิทสฺสน\-อปฺปฏิเฆ ขนฺเธ ปฏิจฺจ สนิทสฺสน\-สปฺปฏิฆํ กฏตฺตารูปํ. อาโปธาตุํ ปฏิจฺจ สนิทสฺสน\-สปฺปฏิฆํ จิตฺต\-สมุฏฺฐานํ รูปํ, กฏตฺตารูปํ, อุปาทารูปํ. อาโปธาตุํ ปฏิจฺจ รูปายตนํ. (2)}

\prose{อนิทสฺสน\-อปฺปฏิฆํ ธมฺมํ ปฏิจฺจ อนิทสฺสน\-สปฺปฏิโฆ ธมฺโม อุปฺปชฺชติ เหตุปจฺจยา. อนิทสฺสน\-อปฺปฏิเฆ ขนฺเธ ปฏิจฺจ อนิทสฺสน\-สปฺปฏิฆํ จิตฺต\-สมุฏฺฐานํ รูปํ. ปฏิสนฺธิ\-กฺขเณ อนิทสฺสน\-อปฺปฏิเฆ ขนฺเธ ปฏิจฺจ อนิทสฺสน\-สปฺปฏิฆํ กฏตฺตารูปํ. อาโปธาตุํ ปฏิจฺจ อนิทสฺสน\-สปฺปฏิฆํ จิตฺต\-สมุฏฺฐานํ รูปํ, กฏตฺตารูปํ, อุปาทารูปํ. อาโปธาตุํ ปฏิจฺจ จกฺขายตนํ ฯเปฯ รสายตนํ. (3)}

\prose{\csromanfolio{466}อนิทสฺสน\-อปฺปฏิฆํ ธมฺมํ ปฏิจฺจ สนิทสฺสน\-สปฺปฏิโฆ จ อนิทสฺสน\-อปฺปฏิโฆ จ ธมฺมา อุปฺปชฺชนฺติ เหตุปจฺจยา. อนิทสฺสน\-อปฺปฏิฆํ เอกํ ขนฺธํ ปฏิจฺจ ตโย ขนฺธา สนิทสฺสนสปฺปฏิฆญฺจ อนิทสฺสนอปฺปฏิฆญฺจ จิตฺต\-สมุฏฺฐานํ รูปํ ฯเปฯ เทฺว ขนฺเธ ฯเปฯ. ปฏิสนฺธิ\-กฺขเณ อนิทสฺสน\-อปฺปฏิฆํ เอกํ ขนฺธํ ปฏิจฺจ ตโย ขนฺธา. สนิทสฺสนสปฺปฏิฆญฺจ อนิทสฺสนอปฺปฏิฆญฺจ กฏตฺตารูปํ ฯเปฯ เทฺว ขนฺเธ ฯเปฯ. อาโปธาตุํ ปฏิจฺจ สนิทสฺสนสปฺปฏิฆญฺจ อนิทสฺสนอปฺปฏิฆญฺจ จิตฺต\-สมุฏฺฐานํ รูปํ, กฏตฺตารูปํ, อุปาทารูปํ. อาโปธาตุํ ปฏิจฺจ รูปายตนํ, อิตฺถินฺทฺริยํ ฯเปฯ. กพฬีกาโร อาหาโร. (4)}

\prose{อนิทสฺสน\-อปฺปฏิฆํ ธมฺมํ ปฏิจฺจ อนิทสฺสน\-สปฺปฏิโฆ จ อนิทสฺสน\-อปฺปฏิโฆ จ ธมฺมา อุปฺปชฺชนฺติ เหตุปจฺจยา. อนิทสฺสน\-อปฺปฏิฆํ เอกํ ขนฺธํ ปฏิจฺจ ตโย ขนฺธา อนิทสฺสนสปฺปฏิฆญฺจ อนิทสฺสนอปฺปฏิฆญฺจ จิตฺต\-สมุฏฺฐานํ รูปํ ฯเปฯ เทฺว ขนฺเธ ฯเปฯ. ปฏิสนฺธิ\-กฺขเณ อนิทสฺสน\-อปฺปฏิฆํ เอกํ ขนฺธํ ปฏิจฺจ ตโย ขนฺธา อนิทสฺสนสปฺปฏิฆญฺจ อนิทสฺสนอปฺปฏิฆญฺจ กฏตฺตารูปํ ฯเปฯ เทฺว ขนฺเธ ฯเปฯ. อาโปธาตุํ ปฏิจฺจ อนิทสฺสนสปฺปฏิฆญฺจ อนิทสฺสนอปฺปฏิฆญฺจ จิตฺต\-สมุฏฺฐานํ รูปํ, กฏตฺตารูปํ, อุปาทารูปํ. อาโปธาตุํ ปฏิจฺจ จกฺขายตนํ ฯเปฯ รสายตนํ, อิตฺถินฺทฺริยํ ฯเปฯ. กพฬีกาโร อาหาโร. (5)}

\prose{อนิทสฺสน\-อปฺปฏิฆํ ธมฺมํ ปฏิจฺจ สนิทสฺสน\-สปฺปฏิโฆ จ อนิทสฺสน\-สปฺปฏิโฆ จ ธมฺมา อุปฺปชฺชนฺติ เหตุปจฺจยา. อนิทสฺสน\-อปฺปฏิเฆ ขนฺเธ ปฏิจฺจ สนิทสฺสนสปฺปฏิฆญฺจ อนิทสฺสนสปฺปฏิฆญฺจ จิตฺต\-สมุฏฺฐานํ รูปํ. ปฏิสนฺธิ\-กฺขเณ อนิทสฺสน\-อปฺปฏิเฆ ขนฺเธ ปฏิจฺจ สนิทสฺสนสปฺปฏิฆญฺจ อนิทสฺสนสปฺปฏิฆญฺจ กฏตฺตารูปํ. อาโปธาตุํ ปฏิจฺจ สนิทสฺสนสปฺปฏิฆญฺจ อนิทสฺสนสปฺปฏิฆญฺจ จิตฺต\-สมุฏฺฐานํ รูปํ, กฏตฺตารูปํ, อุปาทารูปํ. อาโปธาตุํ ปฏิจฺจ รูปายตนํ, จกฺขายตนํ ฯเปฯ รสายตนํ. (6)}

\prose{อนิทสฺสน\-อปฺปฏิฆํ ธมฺมํ ปฏิจฺจ สนิทสฺสน\-สปฺปฏิโฆ จ อนิทสฺสน\-สปฺปฏิโฆ จ อนิทสฺสน\-อปฺปฏิโฆ จ ธมฺมา อุปฺปชฺชนฺติ เหตุปจฺจยา. อนิทสฺสน\-อปฺปฏิฆํ เอกํ ขนฺธํ ปฏิจฺจ ตโย ขนฺธา สนิทสฺสนสปฺปฏิฆญฺจ อนิทสฺสนสปฺปฏิฆญฺจ อนิทสฺสนอปฺปฏิฆญฺจ จิตฺต\-สมุฏฺฐานํ รูปํ ฯเปฯ เทฺว ขนฺเธ ฯเปฯ. ปฏิสนฺธิ\-กฺขเณ อนิทสฺสน\-อปฺปฏิฆํ เอกํ ขนฺธํ ปฏิจฺจ ตโย ขนฺธา สนิทสฺสนสปฺปฏิฆญฺจ อนิทสฺสนสปฺปฏิฆญฺจ อนิทสฺสนอปฺปฏิฆญฺจ กฏตฺตารูปํ ฯเปฯ เทฺว ขนฺเธ ฯเปฯ. อาโปธาตุํ ปฏิจฺจ สนิทสฺสนสปฺปฏิฆญฺจ อนิทสฺสนสปฺปฏิฆญฺจ อนิทสฺสนอปฺปฏิฆญฺจ จิตฺต\-สมุฏฺฐานํ รูปํ, กฏตฺตารูปํ, อุปาทารูปํ. \csromanfolio{467}อาโปธาตุํ ปฏิจฺจ รูปายตนํ จกฺขายตนํ ฯเปฯ รสายตนํ, อิตฺถินฺทฺริยํ ฯเปฯ. กพฬีกาโร อาหาโร. (7)}

\proseitem{3}{อนิทสฺสนสปฺปฏิฆญฺจ อนิทสฺสนอปฺปฏิฆญฺจ ธมฺมํ ปฏิจฺจ สนิทสฺสน\-สปฺปฏิโฆ ธมฺโม อุปฺปชฺชติ เหตุปจฺจยา. อนิทสฺสน\-อปฺปฏิเฆ ขนฺเธ จ มหาภูเต จ ปฏิจฺจ สนิทสฺสน\-สปฺปฏิฆํ จิตฺต\-สมุฏฺฐานํ รูปํ. ปฏิสนฺธิ\-กฺขเณ อนิทสฺสน\-อปฺปฏิเฆ ขนฺเธ จ มหาภูเต จ ปฏิจฺจ สนิทสฺสน\-สปฺปฏิฆํ กฏตฺตารูปํ. อนิทสฺสน\-สปฺปฏิเฆ มหาภูเต จ อาโปธาตุญฺจ ปฏิจฺจ สนิทสฺสน\-สปฺปฏิฆํ จิตฺต\-สมุฏฺฐานํ รูปํ, กฏตฺตารูปํ, อุปาทารูปํ. โผฏฺฐพฺพายตนญฺจ อาโปธาตุญฺจ ปฏิจฺจ รูปายตนํ. (1)}

\prose{อนิทสฺสนสปฺปฏิฆญฺจ อนิทสฺสนอปฺปฏิฆญฺจ ธมฺมํ ปฏิจฺจ อนิทสฺสน\-สปฺปฏิโฆ ธมฺโม อุปฺปชฺชติ เหตุปจฺจยา. อนิทสฺสน\-อปฺปฏิเฆ ขนฺเธ จ มหาภูเต จ ปฏิจฺจ อนิทสฺสน\-สปฺปฏิฆํ จิตฺต\-สมุฏฺฐานํ รูปํ. ปฏิสนฺธิ\-กฺขเณ อนิทสฺสน\-อปฺปฏิเฆ ขนฺเธ จ มหาภูเต จ ปฏิจฺจ อนิทสฺสน\-สปฺปฏิฆํ กฏตฺตารูปํ. อนิทสฺสน\-สปฺปฏิฆํ เอกํ มหาภูตญฺจ อาโปธาตุญฺจ ปฏิจฺจ เทฺว มหาภูตา, เทฺว มหาภูเต จ อาโปธาตุญฺจ ปฏิจฺจ เอกํ มหาภูตํ. อนิทสฺสน\-สปฺปฏิเฆ มหาภูเต จ อาโปธาตุญฺจ ปฏิจฺจ อนิทสฺสน\-สปฺปฏิฆํ จิตฺต\-สมุฏฺฐานํ รูปํ, กฏตฺตารูปํ, อุปาทารูปํ. โผฏฺฐพฺพายตนญฺจ อาโปธาตุญฺจ ปฏิจฺจ จกฺขายตนํ ฯเปฯ รสายตนํ. (2)}

\prose{อนิทสฺสนสปฺปฏิฆญฺจ อนิทสฺสนอปฺปฏิฆญฺจ ธมฺมํ ปฏิจฺจ อนิทสฺสน อปฺปฏิโฆ ธมฺโม อุปฺปชฺชติ เหตุปจฺจยา. อนิทสฺสน\-อปฺปฏิเฆ ขนฺเธ จ มหาภูเต จ ปฏิจฺจ อนิทสฺสน\-อปฺปฏิฆํ จิตฺต\-สมุฏฺฐานํ รูปํ. ปฏิสนฺธิ\-กฺขเณ อนิทสฺสน\-อปฺปฏิเฆ ขนฺเธ จ มหาภูเต จ ปฏิจฺจ อนิทสฺสน\-อปฺปฏิฆํ กฏตฺตารูปํ. อนิทสฺสน\-สปฺปฏิเฆ มหาภูเต จ อาโปธาตุญฺจ ปฏิจฺจ อนิทสฺสน\-อปฺปฏิฆํ จิตฺต\-สมุฏฺฐานํ รูปํ, กฏตฺตารูปํ, อุปาทารูปํ. โผฏฺฐพฺพายตนญฺจ อาโปธาตุญฺจ ปฏิจฺจ อิตฺถินฺทฺริยํ ฯเปฯ. กพฬีกาโร อาหาโร. (3)}

\prose{อนิทสฺสนสปฺปฏิฆญฺจ อนิทสฺสนอปฺปฏิฆญฺจ ธมฺมํ ปฏิจฺจ สนิทสฺสน\-สปฺปฏิโฆ จ อนิทสฺสน\-อปฺปฏิโฆ จ ธมฺมา อุปฺปชฺชนฺติ เหตุปจฺจยา. อนิทสฺสน\-อปฺปฏิเฆ ขนฺเธ จ มหาภูเต จ ปฏิจฺจ สนิทสฺสนสปฺปฏิฆญฺจ อนิทสฺสนอปฺปฏิฆญฺจ จิตฺต\-สมุฏฺฐานํ รูปํ. ปฏิสนฺธิ\-กฺขเณ อนิทสฺสน\-อปฺปฏิเฆ ขนฺเธ จ มหาภูเต จ ปฏิจฺจ สนิทสฺสนสปฺปฏิฆญฺจ อนิทสฺสนอปฺปฏิฆญฺจ \csromanfolio{468}กฏตฺตารูปํ. อนิทสฺสน\-สปฺปฏิเฆ มหาภูเต จ อาโปธาตุญฺจ ปฏิจฺจ สนิทสฺสนสปฺปฏิฆญฺจ อนิทสฺสนอปฺปฏิฆญฺจ จิตฺต\-สมุฏฺฐานํ รูปํ, กฏตฺตารูปํ, อุปาทารูปํ. โผฏฺฐพฺพายตนญฺจ อาโปธาตุญฺจ ปฏิจฺจ รูปายตนํ, อิตฺถินฺทฺริยํ ฯเปฯ. กพฬีกาโร อาหาโร. (4)}

\prose{อนิทสฺสนสปฺปฏิฆญฺจ อนิทสฺสนอปฺปฏิฆญฺจ ธมฺมํ ปฏิจฺจ อนิทสฺสน\-สปฺปฏิโฆ จ อนิทสฺสน\-อปฺปฏิโฆ จ ธมฺมา อุปฺปชฺชนฺติ เหตุปจฺจยา. อนิทสฺสน\-อปฺปฏิเฆ ขนฺเธ จ มหาภูเต จ ปฏิจฺจ อนิทสฺสนสปฺปฏิฆญฺจ อนิทสฺสนอปฺปฏิฆญฺจ จิตฺต\-สมุฏฺฐานํ รูปํ. ปฏิสนฺธิ\-กฺขเณ อนิทสฺสน\-อปฺปฏิเฆ ขนฺเธ จ มหาภูเต จ ปฏิจฺจ อนิทสฺสนสปฺปฏิฆญฺจ อนิทสฺสนอปฺปฏิฆญฺจ กฏตฺตารูปํ. อนิทสฺสน\-สปฺปฏิเฆ มหาภูเต จ อาโปธาตุญฺจ ปฏิจฺจ อนิทสฺสนสปฺปฏิฆญฺจ อนิทสฺสนอปฺปฏิฆญฺจ จิตฺต\-สมุฏฺฐานํ รูปํ, กฏตฺตารูปํ, อุปาทารูปํ. โผฏฺฐพฺพายตนญฺจ อาโปธาตุญฺจ ปฏิจฺจ จกฺขายตนํ ฯเปฯ รสายตนํ, อิตฺถินฺทฺริยํ ฯเปฯ. กพฬีกาโร อาหาโร. (5)}

\prose{อนิทสฺสนสปฺปฏิฆญฺจ อนิทสฺสนอปฺปฏิฆญฺจ ธมฺมํ ปฏิจฺจ สนิทสฺสน\-สปฺปฏิโฆ จ อนิทสฺสน\-สปฺปฏิโฆ จ ธมฺมา อุปฺปชฺชนฺติ เหตุปจฺจยา. อนิทสฺสน\-อปฺปฏิเฆ ขนฺเธ จ มหาภูเต จ ปฏิจฺจ สนิทสฺสนสปฺปฏิฆญฺจ อนิทสฺสนสปฺปฏิฆญฺจ จิตฺต\-สมุฏฺฐานํ รูปํ. ปฏิสนฺธิ\-กฺขเณ อนิทสฺสน\-อปฺปฏิเฆ ขนฺเธ จ มหาภูเต จ ปฏิจฺจ สนิทสฺสนสปฺปฏิฆญฺจ อนิทสฺสนสปฺปฏิฆญฺจ กฏตฺตารูปํ. อนิทสฺสน\-สปฺปฏิเฆ มหาภูเต จ อาโปธาตุญฺจ ปฏิจฺจ สนิทสฺสนสปฺปฏิฆญฺจ อนิทสฺสนสปฺปฏิฆญฺจ จิตฺต\-สมุฏฺฐานํ รูปํ, กฏตฺตารูปํ, อุปาทารูปํ. โผฏฺฐพฺพายตนญฺจ อาโปธาตุญฺจ ปฏิจฺจ รูปายตนํ, จกฺขายตนํ ฯเปฯ รสายตนํ. (6)}

\csromanmatikaanchor{mtk.275}
\csromantocmarkref{subparagraph}{อารมฺมณาทิ}{mtk.275}
\bookmark[dest={mtk.275},level=5]{อารมฺมณาทิ}
\prose{อนิทสฺสนสปฺปฏิฆญฺจ อนิทสฺสนอปฺปฏิฆญฺจ ธมฺมํ ปฏิจฺจ สนิทสฺสน\-สปฺปฏิโฆ จ อนิทสฺสน\-สปฺปฏิโฆ จ อนิทสฺสน\-อปฺปฏิโฆ จ ธมฺมา อุปฺปชฺชนฺติ เหตุปจฺจยา. อนิทสฺสน\-อปฺปฏิเฆ ขนฺเธ จ มหาภูเต จ ปฏิจฺจ สนิทสฺสนสปฺปฏิฆญฺจ อนิทสฺสนสปฺปฏิฆญฺจ อนิทสฺสนอปฺปฏิฆญฺจ จิตฺต\-สมุฏฺฐานํ รูปํ. ปฏิสนฺธิ\-กฺขเณ อนิทสฺสน\-อปฺปฏิเฆ ขนฺเธ จ มหาภูเต จ ปฏิจฺจ สนิทสฺสนสปฺปฏิฆญฺจ อนิทสฺสนสปฺปฏิฆญฺจ อนิทสฺสนอปฺปฏิฆญฺจ กฏตฺตารูปํ. อนิทสฺสน\-สปฺปฏิเฆ มหาภูเต จ อาโปธาตุญฺจ ปฏิจฺจ สนิทสฺสนสปฺปฏิฆญฺจ อนิทสฺสนสปฺปฏิฆญฺจ อนิทสฺสนอปฺปฏิฆญฺจ จิตฺต\-สมุฏฺฐานํ รูปํ, กฏตฺตารูปํ, \csromanfolio{469}อุปาทารูปํ. โผฏฺฐพฺพายตนญฺจ อาโปธาตุญฺจ ปฏิจฺจ รูปายตนํ, จกฺขายตนํ ฯเปฯ รสายตนํ, อิตฺถินฺทฺริยํ ฯเปฯ. กพฬีกาโร อาหาโร. (7)}

\csromanheader{2}{\textbf{อารมฺมณ}}
\proseitem{3}{อนิทสฺสน\-อปฺปฏิฆํ ธมฺมํ ปฏิจฺจ อนิทสฺสน\-อปฺปฏิโฆ ธมฺโม อุปฺปชฺชติ อารมฺมณ\-ปจฺจยา. อนิทสฺสน\-อปฺปฏิฆํ เอกํ ขนฺธํ ปฏิจฺจ ตโย ขนฺธา ฯเปฯ เทฺว ขนฺเธ ฯเปฯ. ปฏิสนฺธิ\-กฺขเณ อนิทสฺสน\-อปฺปฏิฆํ เอกํ ขนฺธํ ปฏิจฺจ ตโย ขนฺธา ฯเปฯ เทฺว ขนฺเธ ฯเปฯ. วตฺถุํ ปฏิจฺจ ขนฺธา. (1)}

\proseitem{5}{อนิทสฺสน\-สปฺปฏิฆํ ธมฺมํ ปฏิจฺจ อนิทสฺสน\-สปฺปฏิโฆ ธมฺโม อุปฺปชฺชติ อธิปติ\-ปจฺจยา. อนิทสฺสน\-สปฺปฏิฆํ เอกํ มหาภูตํ ปฏิจฺจ เทฺว มหาภูตา, เทฺว มหาภูเต ปฏิจฺจ เอกํ มหาภูตํ. อนิทสฺสน\-สปฺปฏิเฆ มหาภูเต ปฏิจฺจ จิตฺต\-สมุฏฺฐานํ รูปํ, อุปาทารูปํ. (1)}

\prose{(อนิทสฺสนสปฺปฏิฆ\-มูลเก อิมินา การเณน สตฺต ปญฺหา วิภชิตพฺพา, ปริโยสานปทา นตฺถิ.)}

\proseitem{3}{อนิทสฺสน\-อปฺปฏิฆํ ธมฺมํ ปฏิจฺจ อนิทสฺสน\-อปฺปฏิโฆ ธมฺโม อุปฺปชฺชติ อธิปติ\-ปจฺจยา. อนิทสฺสน\-อปฺปฏิฆํ เอกํ ขนฺธํ ปฏิจฺจ ตโย ขนฺธา อนิทสฺสนอปฺปฏิฆญฺจ จิตฺต\-สมุฏฺฐานํ รูปํ ฯเปฯ เทฺว ขนฺเธ ฯเปฯ. อาโปธาตุํ ปฏิจฺจ อนิทสฺสน\-อปฺปฏิฆํ จิตฺต\-สมุฏฺฐานํ รูปํ, อุปาทารูปํ.}

\prose{(อิมินา การเณน อนิทสฺสนอปฺปฏิฆ\-มูลเก สตฺต ปญฺหา วิภชิตพฺพา, นิฏฺฐานปทา นตฺถิ.)}

\prose{อนิทสฺสนสปฺปฏิฆญฺจ อนิทสฺสนอปฺปฏิฆญฺจ ธมฺมํ ปฏิจฺจ สนิทสฺสน\-สปฺปฏิโฆ ธมฺโม อุปฺปชฺชติ อธิปติ\-ปจฺจยา. อนิทสฺสน\-อปฺปฏิเฆ ขนฺเธ จ มหาภูเต จ ปฏิจฺจ สนิทสฺสน\-สปฺปฏิฆํ จิตฺต\-สมุฏฺฐานํ รูปํ. อนิทสฺสน\-อปฺปฏิเฆ มหาภูเต จ อาโปธาตุญฺจ ปฏิจฺจ สนิทสฺสน\-สปฺปฏิฆํ จิตฺต\-สมุฏฺฐานํ รูปํ อุปาทารูปํ.}

\vspace{\tipitakaparskip}
\csromancenter{(อิมินา การเณน สตฺตปิ ปญฺหา วิภชิตพฺพา.)}
\csromanheader{2}{อนนฺตร\-สมนนฺตร}
\proseitem{8}{\csromanfolio{470}อนิทสฺสน\-อปฺปฏิฆํ ธมฺมํ ปฏิจฺจ อนิทสฺสน\-อปฺปฏิโฆ ธมฺโม อุปฺปชฺชติ อนนฺตร\-ปจฺจยา. สมนนฺตร\-ปจฺจยา. (อารมฺมณ\-สทิสํ.)}

\csromanheader{2}{\textbf{สหชาต}}
\proseitem{9}{อนิทสฺสน\-สปฺปฏิฆํ ธมฺมํ ปฏิจฺจ อนิทสฺสน\-สปฺปฏิโฆ ธมฺโม อุปฺปชฺชติ สหชาต\-ปจฺจยา. อนิทสฺสน\-สปฺปฏิฆํ เอกํ มหาภูตํ ปฏิจฺจ เทฺว มหาภูตา, เทฺว มหาภูเต ปฏิจฺจ เอกํ มหาภูตํ. อนิทสฺสน\-สปฺปฏิเฆ มหาภูเต ปฏิจฺจ อนิทสฺสน\-สปฺปฏิฆํ จิตฺต\-สมุฏฺฐานํ รูปํ, กฏตฺตารูปํ, อุปาทารูปํ. โผฏฺฐพฺพา\-ยตนํ ปฏิจฺจ จกฺขายตนํ ฯเปฯ รสายตนํ. พาหิรํ. อาหารสมุฏฺฐานํ. อุตุสมุฏฺฐานํ. อสญฺญสตฺตานํ เอกํ มหาภูตํ ปฏิจฺจ เทฺว มหาภูตา ฯเปฯ.}

\prose{(อนิทสฺสนสปฺปฏิฆ\-มูลกา สตฺต ปญฺหา อิมินา การเณน วิภชิตพฺพา.)}

\proseitem{10}{อนิทสฺสน\-อปฺปฏิฆํ ธมฺมํ ปฏิจฺจ อนิทสฺสน\-อปฺปฏิโฆ ธมฺโม อุปฺปชฺชติ สหชาต\-ปจฺจยา. อนิทสฺสน\-อปฺปฏิฆํ เอกํ ขนฺธํ ปฏิจฺจ ตโย ขนฺธา อนิทสฺสน\-อปฺปฏิฆํ จิตฺต\-สมุฏฺฐานญฺจ รูปํ ฯเปฯ เทฺว ขนฺเธ ฯเปฯ. ปฏิสนฺธิ\-กฺขเณ อนิทสฺสน\-อปฺปฏิฆํ เอกํ ขนฺธํ ปฏิจฺจ ตโย ขนฺธา อนิทสฺสน\-อปฺปฏิฆํ กฏตฺตา จ รูปํ ฯเปฯ เทฺว ขนฺเธ ฯเปฯ. ขนฺเธ ปฏิจฺจ วตฺถุ, วตฺถุํ ปฏิจฺจ ขนฺธา. อาโปธาตุํ ปฏิจฺจ อนิทสฺสน\-อปฺปฏิฆํ จิตฺต\-สมุฏฺฐานํ รูปํ กฏตฺตารูปํ อุปาทารูปํ. อาโปธาตุํ ปฏิจฺจ อิตฺถินฺทฺริยํ ฯเปฯ. กพฬีกาโร อาหาโร. พาหิรํ. อาหารสมุฏฺฐานํ. อุตุสมุฏฺฐานํ. อสญฺญสตฺตานํ อาโปธาตุํ ปฏิจฺจ อนิทสฺสน\-อปฺปฏิฆํ กฏตฺตารูปํ, อุปาทารูปํ.}

\prose{(อนิทสฺสนอปฺปฏิฆ\-มูลเก สตฺต ปญฺหา อิมินา การเณน กาตพฺพา.)}

\proseitem{11}{อนิทสฺสนสปฺปฏิฆญฺจ อนิทสฺสนอปฺปฏิฆญฺจ ธมฺมํ ปฏิจฺจ สนิทสฺสน\-สปฺปฏิโฆ ธมฺโม อุปฺปชฺชติ สหชาต\-ปจฺจยา. อนิทสฺสน\-อปฺปฏิเฆ ขนฺเธ จ มหาภูเต จ ปฏิจฺจ สนิทสฺสน\-สปฺปฏิฆํ จิตฺต\-สมุฏฺฐานํ รูปํ. ปฏิสนฺธิ\-กฺขเณ อนิทสฺสน\-อปฺปฏิเฆ ขนฺเธ จ มหาภูเต จ ปฏิจฺจ สนิทสฺสน\-สปฺปฏิฆํ กฏตฺตารูปํ. อนิทสฺสน\-สปฺปฏิเฆ มหาภูเต จ อาโปธาตุญฺจ ปฏิจฺจ \csromanfolio{471}สนิทสฺสน\-สปฺปฏิฆํ จิตฺต\-สมุฏฺฐานํ รูปํ กฏตฺตารูปํ อุปาทารูปํ. โผฏฺฐพฺพายตนญฺจ อาโปธาตุญฺจ ปฏิจฺจ รูปายตนํ. พาหิรํ. อาหารสมุฏฺฐานํ. อุตุสมุฏฺฐานํ. อสญฺญสตฺตานํ อนิทสฺสน\-สปฺปฏิเฆ มหาภูเต จ อาโปธาตุญฺจ ปฏิจฺจ สนิทสฺสน\-สปฺปฏิฆํ กฏตฺตารูปํ, อุปาทารูปํ. (อิมินา การเณน สตฺต ปญฺหา วิภชิตพฺพา.)}

\vspace{\tipitakaparskip}
\csromancenter{(อิมินา การเณน สตฺต ปญฺหา วิภชิตพฺพา.)}
\csromanheader{2}{\textbf{อญฺญมญฺญ}}
\proseitem{11}{อนิทสฺสน\-สปฺปฏิฆํ ธมฺมํ ปฏิจฺจ อนิทสฺสน\-สปฺปฏิโฆ ธมฺโม อุปฺปชฺชติ อญฺญมญฺญ\-ปจฺจยา. อนิทสฺสน\-สปฺปฏิฆํ เอกํ มหาภูตํ ปฏิจฺจ เทฺว มหาภูตา, เทฺว มหาภูเต ปฏิจฺจ เอกํ มหาภูตํ. (1)}

\prose{อนิทสฺสน\-สปฺปฏิฆํ ธมฺมํ ปฏิจฺจ อนิทสฺสน\-อปฺปฏิโฆ ธมฺโม อุปฺปชฺชติ อญฺญมญฺญ\-ปจฺจยา. อนิทสฺสน\-สปฺปฏิเฆ มหาภูเต ปฏิจฺจ อาโปธาตุ. (2)}

\prose{อนิทสฺสน\-สปฺปฏิฆํ ธมฺมํ ปฏิจฺจ อนิทสฺสน\-สปฺปฏิโฆ จ อนิทสฺสน\-อปฺปฏิโฆ จ ธมฺมา อุปฺปชฺชนฺติ อญฺญมญฺญ\-ปจฺจยา. อนิทสฺสน\-สปฺปฏิฆํ เอกํ มหาภูตํ ปฏิจฺจ เทฺว มหาภูตา อาโปธาตุ จ. เทฺว มหาภูเต ปฏิจฺจ เอกํ มหาภูตํ อาโปธาตุ จ. พาหิรํ ฯเปฯ. (3)}

\prose{อนิทสฺสน\-อปฺปฏิฆํ ธมฺมํ ปฏิจฺจ อนิทสฺสน\-อปฺปฏิโฆ ธมฺโม อุปฺปชฺชติ อญฺญมญฺญ\-ปจฺจยา. อนิทสฺสน\-อปฺปฏิฆํ เอกํ ขนฺธํ ปฏิจฺจ ตโย ขนฺธา ฯเปฯ เทฺว ขนฺเธ ฯเปฯ. ปฏิสนฺธิ\-กฺขเณ อนิทสฺสน\-อปฺปฏิฆํ เอกํ ขนฺธํ ปฏิจฺจ ตโย ขนฺธา วตฺถุ จ ฯเปฯ เทฺว ขนฺเธ ฯเปฯ. ขนฺเธ ปฏิจฺจ วตฺถุ, วตฺถุํ ปฏิจฺจ ขนฺธา. (1)}

\prose{อนิทสฺสน\-อปฺปฏิฆํ ธมฺมํ ปฏิจฺจ อนิทสฺสน\-สปฺปฏิโฆ ธมฺโม อุปฺปชฺชติ อญฺญมญฺญ\-ปจฺจยา. อาโปธาตุํ ปฏิจฺจ อนิทสฺสน\-สปฺปฏิฆา มหาภูตา. พาหิรํ ฯเปฯ. (2)}

\proseitem{14}{อนิทสฺสนสปฺปฏิฆญฺจ อนิทสฺสนอปฺปฏิฆญฺจ ธมฺมํ ปฏิจฺจ อนิทสฺสน\-สปฺปฏิโฆ ธมฺโม อุปฺปชฺชติ อญฺญมญฺญ\-ปจฺจยา. อนิทสฺสน\-สปฺปฏิฆํ เอกํ มหาภูตญฺจ อาโปธาตุญฺจ ปฏิจฺจ เทฺว มหาภูตา, เทฺว มหาภูเต จ อาโปธาตุญฺจ ปฏิจฺจ เอกํ มหาภูตํ. พาหิรํ ฯเปฯ. (1)}

\csromanheader{2}{\textbf{นิสฺสยาทิ}}
\csromanmatikaanchor{mtk.276}
\csromanmatikaanchor{mtk.277}
\csromantocmarknopage{subparagraph}{1. ปจฺจยานุโลม 2. สงฺขฺยาวาร}{mtk.276}
\bookmark[dest={mtk.276},level=5]{1. ปจฺจยานุโลม 2. สงฺขฺยาวาร}
\csromantocmarkref{subparagraph}{สุทฺธ}{mtk.277}
\bookmark[dest={mtk.277},level=5]{สุทฺธ}
\proseitem{15}{\csromanfolio{472}อนิทสฺสน\-สปฺปฏิฆํ ธมฺมํ ปฏิจฺจ อนิทสฺสน\-สปฺปฏิโฆ ธมฺโม อุปฺปชฺชติ นิสฺสยปจฺจยา. อุปนิสฺสย\-ปจฺจยา. ปุเรชาต\-ปจฺจยา. อาเสวนปจฺจยา. กมฺมปจฺจยา. วิปากปจฺจยา. อาหารปจฺจยา. อินฺทฺริยปจฺจยา. ฌานปจฺจยา. มคฺคปจฺจยา. สมฺปยุตฺต\-ปจฺจยา. วิปฺปยุตฺต\-ปจฺจยา. อตฺถิปจฺจยา. นตฺถิปจฺจยา. วิคตปจฺจยา. อวิคตปจฺจยา.}

\csromanmarkh{1. ปจฺจยานุโลม}\csromanmarkhii{2. สงฺขฺยาวาร}\csromanheadernomark{2}{1. \textbf{ปจฺจยานุโลม 2. สงฺขฺยาวาร}}
\proseitemwithcloserruled{16}{เหตุยา เอกวีส, อารมฺมเณ เอกํ, อธิปติยา เอกวีส, อนนฺตเร เอกํ, สมนนฺตเร เอกํ, สหชาเต เอกวีส, อญฺญมญฺเญ ฉ, นิสฺสเย เอกวีส, อุปนิสฺสเย เอกํ, ปุเรชาเต เอกํ, อาเสวเน เอกํ, กมฺเม เอกวีส, วิปาเก, อาหาเร เอกวีส, อินฺทฺริเย เอกวีส, ฌาเน, มคฺเค เอกวีส, สมฺปยุตฺเต เอกํ, วิปฺปยุตฺเต เอกวีส, อตฺถิยา เอกวีส, นตฺถิยา เอกํ, วิคเต เอกํ, อวิคเต เอกวีส. (เอวํ คเณตพฺพํ.)}{อนุโลมํ.}
\csromanmarkh{2. ปจฺจย\-ปจฺจนีย}\csromanmarkhii{1. วิภงฺควาร}\csromanheadernomark{2}{2. ปจฺจย\-ปจฺจนีย 1. วิภงฺควาร}
\csromanmatikaanchor{mtk.278}
\csromanmatikaanchor{mtk.279}
\csromantocmarknopage{subparagraph}{2. ปจฺจยปจฺจนีย 1. วิภงฺควาร}{mtk.278}
\bookmark[dest={mtk.278},level=5]{2. ปจฺจยปจฺจนีย 1. วิภงฺควาร}
\csromantocmarkref{subparagraph}{นเหตุ}{mtk.279}
\bookmark[dest={mtk.279},level=5]{นเหตุ}
\csromanheader{6}{\textbf{นเหตุ}}
\proseitem{17}{อนิทสฺสน\-สปฺปฏิฆํ ธมฺมํ ปฏิจฺจ อนิทสฺสน\-สปฺปฏิโฆ ธมฺโม อุปฺปชฺชติ นเหตุปจฺจยา. อนิทสฺสน\-สปฺปฏิฆํ เอกํ มหาภูตํ ปฏิจฺจ เทฺว มหาภูตา เทฺว มหาภูเต ปฏิจฺจ เอกํ มหาภูตํ. อนิทสฺสน\-สปฺปฏิเฆ มหาภูเต ปฏิจฺจ อนิทสฺสน\-สปฺปฏิฆํ จิตฺต\-สมุฏฺฐานํ รูปํ, กฏตฺตารูปํ, อุปาทารูปํ. โผฏฺฐพฺพา\-ยตนํ ปฏิจฺจ จกฺขายตนํ ฯเปฯ รสายตนํ. พาหิรํ. อาหารสมุฏฺฐานํ. อุตุสมุฏฺฐานํ. อสญฺญสตฺตานํ อนิทสฺสน\-สปฺปฏิฆํ เอกํ มหาภูตํ ปฏิจฺจ เทฺว มหาภูตา, เทฺว มหาภูเต ปฏิจฺจ เอกํ มหาภูตํ. มหาภูเต ปฏิจฺจ ฯเปฯ.}

\prose{(อนิทสฺสนสปฺปฏิฆ\-มูลเก อิมินา การเณน สตฺตปิ ปญฺหา วิภชิตพฺพา.)}

\csromanmatikaanchor{mtk.280}
\csromantocmarkref{subparagraph}{นอารมฺมณ}{mtk.280}
\bookmark[dest={mtk.280},level=5]{นอารมฺมณ}
\proseitem{18}{\csromanfolio{473}อนิทสฺสน\-อปฺปฏิฆํ ธมฺมํ ปฏิจฺจ อนิทสฺสน\-อปฺปฏิโฆ ธมฺโม อุปฺปชฺชติ นเหตุปจฺจยา. อเหตุกํ อนิทสฺสน\-อปฺปฏิฆํ เอกํ ขนฺธํ ปฏิจฺจ ตโย ขนฺธา อนิทสฺสนอปฺปฏิฆญฺจ จิตฺต\-สมุฏฺฐานญฺจ รูปํ ฯเปฯ เทฺว ขนฺเธ ฯเปฯ. อเหตุก\-ปฏิสนฺธิ\-กฺขเณ อนิทสฺสน\-อปฺปฏิฆํ เอกํ ขนฺธํ ปฏิจฺจ ตโย ขนฺธา อนิทสฺสนอปฺปฏิฆญฺจ กฏตฺตารูปํ ฯเปฯ. ขนฺเธ ปฏิจฺจ วตฺถุ, วตฺถุํ ปฏิจฺจ ขนฺธา. อาโปธาตุํ ปฏิจฺจ อนิทสฺสน\-อปฺปฏิฆํ จิตฺต\-สมุฏฺฐานํ รูปํ, กฏตฺตารูปํ, อุปาทารูปํ. อาโปธาตุํ ปฏิจฺจ อิตฺถินฺทฺริยํ ฯเปฯ. กพฬีกาโร อาหาโร. พาหิรํ. อาหารสมุฏฺฐานํ. อุตุสมุฏฺฐานํ. อสญฺญสตฺตานํ อาโปธาตุํ ปฏิจฺจ อนิทสฺสน\-อปฺปฏิฆํ กฏตฺตารูปํ, อุปาทารูปํ. วิจิกิจฺฉา\-สหคเต อุทฺธจฺจ\-สหคเต ขนฺเธ ปฏิจฺจ วิจิกิจฺฉา\-สหคโต อุทฺธจฺจ\-สหคโต โมโห.}

\prose{(อนิทสฺสนอปฺปฏิฆ\-มูลกา อิมินา การเณน สตฺต ปญฺหา วิภชิตพฺพา.)}

\proseitem{19}{อนิทสฺสนสปฺปฏิฆญฺจ อนิทสฺสนอปฺปฏิฆญฺจ ธมฺมํ ปฏิจฺจ สนิทสฺสน\-สปฺปฏิโฆ ธมฺโม อุปฺปชฺชติ นเหตุปจฺจยา. อเหตุเก อนิทสฺสน\-อปฺปฏิเฆ ขนฺเธ จ มหาภูเต จ ปฏิจฺจ สนิทสฺสน\-สปฺปฏิฆํ จิตฺต\-สมุฏฺฐานํ รูปํ. อเหตุก\-ปฏิสนฺธิ\-กฺขเณ อนิทสฺสน\-อปฺปฏิเฆ ขนฺเธ จ มหาภูเต จ ปฏิจฺจ สนิทสฺสน\-สปฺปฏิฆํ กฏตฺตารูปํ. อนิทสฺสน\-สปฺปฏิเฆ มหาภูเต จ อาโปธาตุญฺจ ปฏิจฺจ สนิทสฺสน\-สปฺปฏิฆํ จิตฺต\-สมุฏฺฐานํ รูปํ, กฏตฺตารูปํ, อุปาทารูปํ. โผฏฺฐพฺพายตนญฺจ อาโปธาตุญฺจ ปฏิจฺจ รูปายตนํ. พาหิรํ. อาหารสมุฏฺฐานํ. อุตุสมุฏฺฐานํ. อสญฺญสตฺตานํ อนิทสฺสน\-สปฺปฏิเฆ มหาภูเต จ อาโปธาตุญฺจ ปฏิจฺจ สนิทสฺสน\-สปฺปฏิฆํ กฏตฺตารูปํ, อุปาทารูปํ.}

\vspace{\tipitakaparskip}
\csromancenter{(อิมินา การเณน สตฺต ปญฺหา วิตฺถาเรตพฺพา อสมฺโมหนฺเตน.)}
\prose{นอารมฺมณ}

\csromanmatikaanchor{mtk.281}
\csromantocmarkref{subparagraph}{นอธิปตฺยาทิ}{mtk.281}
\bookmark[dest={mtk.281},level=5]{นอธิปตฺยาทิ}
\proseitem{20}{อนิทสฺสน\-สปฺปฏิฆํ ธมฺมํ ปฏิจฺจ อนิทสฺสน\-สปฺปฏิโฆ ธมฺโม อุปฺปชฺชติ นอารมฺมณ\-ปจฺจยา. อนิทสฺสน\-สปฺปฏิฆํ เอกํ มหาภูตํ ปฏิจฺจ เทฺว มหาภูตา, เทฺว มหาภูเต ปฏิจฺจ เอกํ มหาภูตํ. อนิทสฺสน\-สปฺปฏิเฆ มหาภูเต ปฏิจฺจ อนิทสฺสน\-สปฺปฏิฆํ จิตฺต\-สมุฏฺฐานํ รูปํ, กฏตฺตารูปํ, อุปาทารูปํ. โผฏฺฐพฺพา\-ยตนํ ปฏิจฺจ จกฺขายตนํ ฯเปฯ รสายตนํ. พาหิรํ. \csromanfolio{474}อาหารสมุฏฺฐานํ. อุตุสมุฏฺฐานํ. อสญฺญสตฺตานํ อนิทสฺสน\-สปฺปฏิฆํ เอกํ มหาภูตํ ปฏิจฺจ เทฺว มหาภูตา, เทฺว มหาภูเต ปฏิจฺจ เอกํ มหาภูตํ ฯเปฯ.}

\prose{(อนิทสฺสนสปฺปฏิฆ\-มูลกา อิมินา การเณน สตฺตปิ ปญฺหา วิตฺถาเรตพฺพา.)}

\prose{อนิทสฺสน\-อปฺปฏิฆํ ธมฺมํ ปฏิจฺจ อนิทสฺสน\-อปฺปฏิโฆ ธมฺโม อุปฺปชฺชติ นอารมฺมณ\-ปจฺจยา. อนิทสฺสน\-อปฺปฏิเฆ ขนฺเธ ปฏิจฺจ อนิทสฺสน\-อปฺปฏิฆํ จิตฺต\-สมุฏฺฐานํ รูปํ. ปฏิสนฺธิ\-กฺขเณ อนิทสฺสน\-อปฺปฏิเฆ ขนฺเธ ปฏิจฺจ อนิทสฺสน\-อปฺปฏิฆํ กฏตฺตารูปํ. ขนฺเธ ปฏิจฺจ วตฺถุ ฯเปฯ. อาโปธาตุํ ปฏิจฺจ อนิทสฺสน\-อปฺปฏิฆํ จิตฺต\-สมุฏฺฐานํ รูปํ, กฏตฺตารูปํ, อุปาทารูปํ. อาโปธาตุํ ปฏิจฺจ อิตฺถินฺทฺริยํ ฯเปฯ. กพฬีกาโร อาหาโร. พาหิรํ. อาหารสมุฏฺฐานํ. อุตุสมุฏฺฐานํ. อสญฺญสตฺตานํ อาโปธาตุํ ปฏิจฺจ อนิทสฺสน\-อปฺปฏิฆํ กฏตฺตารูปํ, อุปาทารูปํ.}

\prose{(อนิทสฺสนอปฺปฏิฆ\-มูลเก อิมินา การเณน สตฺตปิ ปญฺหา วิตฺถาเรตพฺพา.)}

\prose{อนิทสฺสนสปฺปฏิฆญฺจ อนิทสฺสนอปฺปฏิฆญฺจ ธมฺมํ ปฏิจฺจ สนิทสฺสน\-สปฺปฏิโฆ ธมฺโม อุปฺปชฺชติ นอารมฺมณ\-ปจฺจยา. อนิทสฺสน\-อปฺปฏิเฆ ขนฺเธ จ มหาภูเต จ ปฏิจฺจ สนิทสฺสน\-สปฺปฏิฆํ จิตฺต\-สมุฏฺฐานํ รูปํ. ปฏิสนฺธิ\-กฺขเณ อนิทสฺสน\-อปฺปฏิเฆ ขนฺเธ จ มหาภูเต จ ปฏิจฺจ สนิทสฺสน\-สปฺปฏิฆํ กฏตฺตารูปํ. อนิทสฺสน\-สปฺปฏิเฆ มหาภูเต จ อาโปธาตุญฺจ ปฏิจฺจ สนิทสฺสน\-สปฺปฏิฆํ จิตฺต\-สมุฏฺฐานํ รูปํ, กฏตฺตารูปํ, อุปาทารูปํ. โผฏฺฐพฺพายตนญฺจ อาโปธาตุญฺจ ปฏิจฺจ รูปายตนํ. พาหิรํ. อาหารสมุฏฺฐานํ. อุตุสมุฏฺฐานํ. อสญฺญสตฺตานํ อนิทสฺสน\-สปฺปฏิเฆ มหาภูเต จ อาโปธาตุญฺจ ปฏิจฺจ สนิทสฺสน\-สปฺปฏิฆํ กฏตฺตารูปํ, อุปาทารูปํ.}

\vspace{\tipitakaparskip}
\csromancenter{(ฆฏเน อิมินา การเณน สตฺตปิ ปญฺหา วิภชิตพฺพา.)}
\prose{นอธิปตฺยาทิ}

\csromanmatikaanchor{mtk.282}
\csromantocmarkref{subparagraph}{นอินฺทฺริยาทิ}{mtk.282}
\bookmark[dest={mtk.282},level=5]{นอินฺทฺริยาทิ}
\proseitem{23}{อนิทสฺสน\-สปฺปฏิฆํ ธมฺมํ ปฏิจฺจ อนิทสฺสน\-สปฺปฏิโฆ ธมฺโม อุปฺปชฺชติ นอธิปติ\-ปจฺจยา. (สหชาต\-สทิสํ.) นอนนฺตร\-ปจฺจยา. นสมนนฺตร\-ปจฺจยา. นอญฺญมญฺญ\-ปจฺจยา. อนิทสฺสน\-สปฺปฏิเฆ มหาภูเต ปฏิจฺจ \csromanfolio{475}อนิทสฺสน\-สปฺปฏิฆํ จิตฺต\-สมุฏฺฐานํ รูปํ, กฏตฺตารูปํ, อุปาทารูปํ. โผฏฺฐพฺพา\-ยตนํ ปฏิจฺจ จกฺขายตนํ ฯเปฯ รสายตนํ. พาหิรํ. อาหารสมุฏฺฐานํ. อุตุสมุฏฺฐานํ. อสญฺญสตฺตานํ มหาภูเต ปฏิจฺจ อนิทสฺสน\-สปฺปฏิฆํ กฏตฺตารูปํ, อุปาทารูปํ. (อิมินา การเณน เอกวีส ปญฺหา วิภชิตพฺพา.)}

\vspace{\tipitakaparskip}
\csromancenter{(อิมินา การเณน เอกวีส ปญฺหา วิภชิตพฺพา.)}
\prose{นอุปนิสฺสยปจฺจยา. นปุเรชาต\-ปจฺจยา. นปจฺฉาชาต\-ปจฺจยา. นอาเสวน\-ปจฺจยา. นกมฺมปจฺจยา. พาหิรํ. อาหารสมุฏฺฐานํ. อุตุสมุฏฺฐานํ. อนิทสฺสน\-สปฺปฏิฆํ เอกํ มหาภูตํ ปฏิจฺจ เทฺว มหาภูตา, เทฺว มหาภูเต ปฏิจฺจ เอกํ มหาภูตํ. อนิทสฺสน\-สปฺปฏิเฆ มหาภูเต ปฏิจฺจ อนิทสฺสน\-สปฺปฏิฆํ อุปาทารูปํ. (กมฺมํ วิภชิตฺวา นกมฺเมเนว เอกวีส ปญฺหา กาตพฺพา.) นวิปาก\-ปจฺจยา. (ปฏิสนฺธิปิ กฏตฺตาปิ นตฺถิ, ปญฺจโวกาเรเยว กาตพฺพา.) นอาหารปจฺจยา. พาหิรํ. อุตุสมุฏฺฐานํ. อสญฺญสตฺตานํ ฯเปฯ. (อิมินา การเณน วิภชิตพฺพา เอกวีสาปิ.)}

\prose{นอินฺทฺริยาทิ}

\proseitem{24}{อนิทสฺสน\-สปฺปฏิฆํ ธมฺมํ ปฏิจฺจ อนิทสฺสน\-สปฺปฏิโฆ ธมฺโม อุปฺปชฺชติ นอินฺทฺริยปจฺจยา. พาหิรํ. อาหารสมุฏฺฐานํ. อุตุสมุฏฺฐานํ. อนิทสฺสน\-สปฺปฏิฆํ เอกํ มหาภูตํ ฯเปฯ. อสญฺญสตฺตานํ มหาภูเต ปฏิจฺจ รูปชีวิตินฺทฺริยํ. (สํขิตฺตํ. สพฺเพ ปญฺหา วิภชิตพฺพา.) นฌานปจฺจยา. พาหิรํ. อาหารสมุฏฺฐานํ. อุตุสมุฏฺฐานํ. อสญฺญสตฺตานํ เอกํ มหาภูตํ. (สํขิตฺตํ. สตฺตปิ ปญฺหา วิภชิตพฺพา.)}

\prose{อนิทสฺสน\-อปฺปฏิฆํ ธมฺมํ ปฏิจฺจ อนิทสฺสน\-อปฺปฏิโฆ ธมฺโม อุปฺปชฺชติ นฌานปจฺจยา. ปญฺจ\-วิญฺญาณสหคตํ เอกํ ขนฺธํ ปฏิจฺจ ตโย ขนฺธา ฯเปฯ เทฺว ขนฺเธ ฯเปฯ. พาหิรํ. อาหารสมุฏฺฐานํ. อุตุสมุฏฺฐานํ. อสญฺญสตฺตานํ อาโปธาตุํ ปฏิจฺจ อนิทสฺสน\-อปฺปฏิฆํ กฏตฺตารูปํ, อุปาทารูปํ. (เอวํ สตฺตปิ ปญฺหา วิภชิตพฺพา.)}

\csromanmatikaanchor{mtk.283}
\csromanmatikaanchor{mtk.284}
\csromanmatikaanchor{mtk.285}
\csromanmatikaanchor{mtk.286}
\csromanmatikaanchor{mtk.287}
\csromanmatikaanchor{mtk.288}
\csromantocmarknopage{subparagraph}{2. ปจฺจยปจฺจนีย 2. สงฺขฺยาวาร}{mtk.283}
\bookmark[dest={mtk.283},level=5]{2. ปจฺจยปจฺจนีย 2. สงฺขฺยาวาร}
\csromantocmarkref{subparagraph}{สุทฺธ}{mtk.284}
\bookmark[dest={mtk.284},level=5]{สุทฺธ}
\csromantocmarknopage{subparagraph}{3. ปจฺจยานุโลมปจฺจนีย}{mtk.285}
\bookmark[dest={mtk.285},level=5]{3. ปจฺจยานุโลมปจฺจนีย}
\csromantocmarkref{subparagraph}{เหตุทุก}{mtk.286}
\bookmark[dest={mtk.286},level=5]{เหตุทุก}
\csromantocmarknopage{subparagraph}{4. ปจฺจยปจฺจนียานุโลม}{mtk.287}
\bookmark[dest={mtk.287},level=5]{4. ปจฺจยปจฺจนียานุโลม}
\csromantocmarkref{subparagraph}{นเหตุทุก}{mtk.288}
\bookmark[dest={mtk.288},level=5]{นเหตุทุก}
\proseitem{25}{อนิทสฺสนสปฺปฏิฆญฺจ อนิทสฺสนอปฺปฏิฆญฺจ ธมฺมํ ปฏิจฺจ สนิทสฺสน\-สปฺปฏิโฆ ธมฺโม อุปฺปชฺชติ นฌานปจฺจยา. พาหิรํ. อาหารสมุฏฺฐานํ. \csromanfolio{476}อุตุสมุฏฺฐานํ. อสญฺญสตฺตานํ ฯเปฯ. อนิทสฺสน\-สปฺปฏิเฆ มหาภูเต จ อาโปธาตุญฺจ ปฏิจฺจ สนิทสฺสน\-สปฺปฏิฆํ กฏตฺตารูปํ, อุปาทารูปํ.}

\prose{(เอวํ สตฺตปิ ปญฺหา วิภชิตพฺพา.) นมคฺคปจฺจยา. (นเหตุสทิสํ กาตพฺพํ. ปริปุณฺณํ. โมโห นตฺถิ.) นสมฺปยุตฺต\-ปจฺจยา. นวิปฺปยุตฺต\-ปจฺจยา. (ปริปุณฺณํ.) โนนตฺถิปจฺจยา. โนวิคต\-ปจฺจยา.}

\csromanmarkh{2. ปจฺจย\-ปจฺจนีย}\csromanmarkhii{2. สงฺขฺยาวาร}\csromanheadernomark{2}{2. ปจฺจย\-ปจฺจนีย 2. สงฺขฺยาวาร}
\proseitemwithcloserruled{26}{นเหตุยา เอกวีส, นอารมฺมเณ เอกวีส, นอธิปติยา เอกวีส, (สํขิตฺตํ. สพฺพตฺถ เอกวีส.) โนนตฺถิยา เอกวีส, โนวิคเต เอกวีส. (เอวํ คเณตพฺพํ.)}{ปจฺจนียํ.}
\proseitemwithcloserruled{27}{เหตุปจฺจยา นอารมฺมเณ เอกวีส, นอธิปติยา เอกวีส ฯเปฯ นกมฺเม เอกํ, นวิปาเก เอกวีส, นสมฺปยุตฺเต เอกวีส, นวิปฺปยุตฺเต เอกํ, โนนตฺถิยา เอกวีส, โนวิคเต เอกวีส. (เอวํ คเณตพฺพํ.)}{อนุโลม\-ปจฺจนียํ.}
\proseitem{28}{นเหตุปจฺจยา อารมฺมเณ เอกํ, อนนฺตเร เอกํ, สมนนฺตเร เอกํ, สหชาเต เอกวีส ฯเปฯ ฌาเน เอกวีส, มคฺเค เอกํ, สมฺปยุตฺเต \csromanfolio{477}เอกํ, วิปฺปยุตฺเต เอกวีส, อตฺถิยา เอกวีส, นตฺถิยา เอกํ, วิคเต เอกํ, อวิคเต เอกวีส. (เอวํ คเณตพฺพํ.) ปจฺจนียา\-นุโลมํ. ปฏิจฺจวาโร.}

\prose{(สหชาตวาโรปิ ปจฺจยวาโรปิ นิสฺสยวาโรปิ ปฏิจฺจ\-วาร\-สทิสา, สํสฏฺฐ\-วาโรปิ สมฺปยุตฺตวาโรปิ อรูเปเยว กาตพฺพา.)}
\csromansectionrule

\csromanmatikaanchor{mtk.289}
\csromantocmarknopage{subparagraph}{22. สนิทสฺสนสปฺปฏิฆตฺติก 7. ปญฺหาวาร}{mtk.289}
\bookmark[dest={mtk.289},level=5]{22. สนิทสฺสนสปฺปฏิฆตฺติก 7. ปญฺหาวาร}
\csromanmarkh{22. สนิทสฺสนสปฺปฏิฆ\-ตฺติก}\csromanmarkhii{7. ปญฺหาวาร}\csromanheadernomark{6}{22. สนิทสฺสนสปฺปฏิฆ\-ตฺติก 7. ปญฺหาวาร}
\csromanmarkh{1. ปจฺจยานุโลม}\csromanmarkhii{1. วิภงฺควาร}\csromanheadernomark{2}{1. \textbf{ปจฺจยานุโลม 1. วิภงฺควาร}}
\csromanmatikaanchor{mtk.290}
\csromanmatikaanchor{mtk.291}
\csromantocmarknopage{subparagraph}{1. ปจฺจยานุโลม 1. วิภงฺควาร}{mtk.290}
\bookmark[dest={mtk.290},level=5]{1. ปจฺจยานุโลม 1. วิภงฺควาร}
\csromantocmarkref{subparagraph}{เหตุ}{mtk.291}
\bookmark[dest={mtk.291},level=5]{เหตุ}
\csromanheader{6}{\textbf{เหตุ}}
\proseitem{29}{อนิทสฺสน\-อปฺปฏิโฆ ธมฺโม อนิทสฺสน\-อปฺปฏิฆสฺส ธมฺมสฺส เหตุปจฺจเยน ปจฺจโย. อนิทสฺสน\-อปฺปฏิฆา เหตู สมฺปยุตฺตกานํ ขนฺธานํ อนิทสฺสน\-อปฺปฏิฆานญฺจ จิตฺต\-สมุฏฺฐานานํ รูปานํ เหตุปจฺจเยน ปจฺจโย. ปฏิสนฺธิ\-กฺขเณ อนิทสฺสน\-อปฺปฏิฆา เหตู สมฺปยุตฺตกานํ ขนฺธานํ อนิทสฺสน\-อปฺปฏิฆานญฺจ กฏตฺตารูปานํ เหตุปจฺจเยน ปจฺจโย. (1)}

\prose{อนิทสฺสน\-อปฺปฏิโฆ ธมฺโม สนิทสฺสน\-สปฺปฏิฆสฺส ธมฺมสฺส เหตุปจฺจเยน ปจฺจโย. อนิทสฺสน\-อปฺปฏิฆา เหตู สนิทสฺสน\-สปฺปฏิฆานํ จิตฺต\-สมุฏฺฐานานํ รูปานํ เหตุปจฺจเยน ปจฺจโย. ปฏิสนฺธิ\-กฺขเณ ฯเปฯ. (2)}

\prose{(อนิทสฺสนอปฺปฏิฆมูลเกเยว อิมินา การเณน สตฺต ปญฺหา วิภชิตพฺพา.)}

\csromanmatikaanchor{mtk.292}
\csromantocmarkref{subparagraph}{อารมฺมณ}{mtk.292}
\bookmark[dest={mtk.292},level=5]{อารมฺมณ}
\csromanheader{6}{\textbf{อารมฺมณ}}
\csromanmatikaanchor{mtk.293}
\csromantocmarkref{subparagraph}{อธิปติ}{mtk.293}
\bookmark[dest={mtk.293},level=5]{อธิปติ}
\proseitem{30}{สนิทสฺสน\-สปฺปฏิโฆ ธมฺโม อนิทสฺสน\-อปฺปฏิฆสฺส ธมฺมสฺส อารมฺมณ\-ปจฺจเยน ปจฺจโย. รูเป อนิจฺจโต ทุกฺขโต อนตฺตโต วิปสฺสติ, อสฺสาเทติ อภินนฺทติ, ตํ อารพฺภ ราโค อุปฺปชฺชติ, ทิฏฺฐิ อุปฺปชฺชติ, วิจิกิจฺฉา อุปฺปชฺชติ, อุทฺธจฺจํ อุปฺปชฺชติ, โทมนสฺสํ อุปฺปชฺชติ. ทิพฺเพน \csromanfolio{478}จกฺขุนา รูปํ ปสฺสติ. รูปายตนํ จกฺขุ\-วิญฺญาณสฺส อารมฺมณ\-ปจฺจเยน ปจฺจโย. สนิทสฺสน\-สปฺปฏิฆา ขนฺธา อิทฺธิวิธ\-ญาณสฺส, ปุพฺเพ\-นิวาสา\-นุสฺสติ\-ญาณสฺส, อนาคตํสญาณสฺส, อาวชฺชนาย อารมฺมณ\-ปจฺจเยน ปจฺจโย. (1)}

\prose{อนิทสฺสน\-สปฺปฏิโฆ ธมฺโม อนิทสฺสน\-อปฺปฏิฆสฺส ธมฺมสฺส อารมฺมณ\-ปจฺจเยน ปจฺจโย. จกฺขุํ ฯเปฯ. กายํ. สทฺเท. คนฺเธ. รเส. โผฏฺฐพฺเพ อนิจฺจโต ฯเปฯ โทมนสฺสํ อุปฺปชฺชติ. ทิพฺพาย โสตธาตุยา สทฺทํ สุณาติ. สทฺทายตนํ โสตวิญฺญาณสฺส ฯเปฯ. โผฏฺฐพฺพา\-ยตนํ กายวิญฺญาณสฺส ฯเปฯ. อนิทสฺสน\-สปฺปฏิฆา ขนฺธา อิทฺธิวิธ\-ญาณสฺส, ปุพฺเพ\-นิวาสา\-นุสฺสติ\-ญาณสฺส, อนาคตํสญาณสฺส, อาวชฺชนาย อารมฺมณ\-ปจฺจเยน ปจฺจโย. (1)}

\proseitem{31}{อนิทสฺสน\-อปฺปฏิโฆ ธมฺโม อนิทสฺสน\-อปฺปฏิฆสฺส ธมฺมสฺส อารมฺมณ\-ปจฺจเยน ปจฺจโย. ทานํ ทตฺวา, สีลํ สมาทิยิตฺวา, อุโปสถกมฺมํ กตฺวา ตํ ปจฺจเวกฺขติ. ปุพฺเพ สุจิณฺณานิ ปจฺจเวกฺขติ. ฌานา ฯเปฯ. อริยา มคฺคา วุฏฺฐหิตฺวา มคฺคํ ปจฺจเวกฺขนฺติ, ผลํ ปจฺจเวกฺขนฺติ, นิพฺพานํ ปจฺจเวกฺขนฺติ. นิพฺพานํ โคตฺรภุสฺส, โวทานสฺส, มคฺคสฺส, ผลสฺส, อาวชฺชนาย อารมฺมณ\-ปจฺจเยน ปจฺจโย. อริยา ปหีเน กิเลเส ปจฺจเวกฺขนฺติ, วิกฺขมฺภิเต กิเลเส ปจฺจเวกฺขนฺติ. ปุพฺเพ สมุทาจิณฺเณ กิเลเส ชานนฺติ ฯเปฯ. วตฺถุํ. อิตฺถินฺทฺริยํ. ปุริสินฺทฺริยํ. ชีวิตินฺทฺริยํ. อาโปธาตุํ. กพฬีการํ อาหารํ. อนิทสฺสน\-อปฺปฏิเฆ ขนฺเธ อนิจฺจโต ฯเปฯ โทมนสฺสํ อุปฺปชฺชติ. เจโตปริย\-ญาเณน อนิทสฺสนอปฺปฏิฆ\-จิตฺต\-สมงฺคิสฺส จิตฺตํ ชานาติ. อากาสา\-นญฺจา\-ยตนํ วิญฺญาณญฺจา\-ยตนสฺส อารมฺมณ\-ปจฺจเยน ปจฺจโย. อากิญฺจญฺญา\-ยตนํ เนว\-สญฺญา\-นา\-สญฺญา\-ยตนสฺส ฯเปฯ. อนิทสฺสน\-อปฺปฏิฆา ขนฺธา อิทฺธิวิธ\-ญาณสฺส, เจโตปริย\-ญาณสฺส, ปุพฺเพ\-นิวาสา\-นุสฺสติ\-ญาณสฺส, ยถากมฺมูปคญาณสฺส, อนาคตํสญาณสฺส, อาวชฺชนาย อารมฺมณ\-ปจฺจเยน ปจฺจโย. (1)}

\proseitem{32}{สนิทสฺสน\-สปฺปฏิโฆ ธมฺโม อนิทสฺสน\-อปฺปฏิฆสฺส ธมฺมสฺส อธิปติ\-ปจฺจเยน ปจฺจโย. อารมฺมณา\-ธิปติ\csromandash{}รูปํ ครุํ กตฺวา อสฺสาเทติ อภินนฺทติ, ตํ ครุํ กตฺวา ราโค อุปฺปชฺชติ. ทิฏฺฐิ อุปฺปชฺชติ. (1)}

\prose{\csromanfolio{479}อนิทสฺสน\-สปฺปฏิโฆ ธมฺโม อนิทสฺสน\-อปฺปฏิฆสฺส ธมฺมสฺส อธิปติ\-ปจฺจเยน ปจฺจโย. อารมฺมณา\-ธิปติ\csromandash{}จกฺขุํ ฯเปฯ. กายํ. สทฺเท. คนฺเธ. รเส. โผฏฺฐพฺเพ ครุํ กตฺวา อสฺสาเทติ อภินนฺทติ, ตํ ครุํ กตฺวา ราโค อุปฺปชฺชติ. ทิฏฺฐิ อุปฺปชฺชติ. (1)}

\proseitem{33}{อนิทสฺสน\-อปฺปฏิโฆ ธมฺโม อนิทสฺสน\-อปฺปฏิฆสฺส ธมฺมสฺส อธิปติ\-ปจฺจเยน ปจฺจโย. อารมฺมณา\-ธิปติ, สหชาตา\-ธิปติ.  อารมฺมณา\-ธิปติ\csromandash{}ทานํ ทตฺวา, สีลํ สมาทิยิตฺวา, อุโปสถกมฺมํ กตฺวา ตํ ครุํ กตฺวา ฯเปฯ. ฌานา วุฏฺฐหิตฺวา ฯเปฯ. อริยา มคฺคา วุฏฺฐหิตฺวา ฯเปฯ. ผลา วุฏฺฐหิตฺวา ฯเปฯ. ผลํ ครุํ กตฺวา ฯเปฯ. นิพฺพานํ ครุํ กตฺวา. นิพฺพานํ โคตฺรภุสฺส, โวทานสฺส, มคฺคสฺส, ผลสฺส อธิปติ\-ปจฺจเยน ปจฺจโย. วตฺถุํ. อิตฺถินฺทฺริยํ. ปุริสินฺทฺริยํ. ชีวิตินฺทฺริยํ. อาโปธาตุํ. กพฬีการํ อาหารํ. อนิทสฺสน\-อปฺปฏิเฆ ขนฺเธ ครุํ กตฺวา อสฺสาเทติ อภินนฺทติ, ตํ ครุํ กตฺวา ราโค อุปฺปชฺชติ. ทิฏฺฐิ อุปฺปชฺชติ. สหชาตา\-ธิปติ\csromandash{}อนิทสฺสน\-อปฺปฏิฆา\-ธิปติ สมฺปยุตฺตกานํ ขนฺธานํ อนิทสฺสนอปฺปฏิฆานญฺจ จิตฺต\-สมุฏฺฐานานํ รูปานํ อธิปติ\-ปจฺจเยน ปจฺจโย. (1)}

\prose{อนิทสฺสน\-อปฺปฏิโฆ ธมฺโม สนิทสฺสน\-สปฺปฏิฆสฺส ธมฺมสฺส อธิปติ\-ปจฺจเยน ปจฺจโย. สหชาตา\-ธิปติ\csromandash{}อนิทสฺสน\-อปฺปฏิฆา\-ธิปติ สนิทสฺสน\-สปฺปฏิฆานํ จิตฺต\-สมุฏฺฐานานํ รูปานํ อธิปติ\-ปจฺจเยน ปจฺจโย. (2)}

\prose{(อนิทสฺสนอปฺปฏิฆ\-มูลเก สตฺตปิ ปญฺหา วิภชิตพฺพา, อธิปติ ติวิธ\-รูปสงฺคเหน.)}

\csromanmatikaanchor{mtk.294}
\csromantocmarkref{subparagraph}{อนนฺตร}{mtk.294}
\bookmark[dest={mtk.294},level=5]{อนนฺตร}
\csromanheader{6}{\textbf{อนนฺตร}}
\proseitem{34}{อนิทสฺสน\-อปฺปฏิโฆ ธมฺโม อนิทสฺสน\-อปฺปฏิฆสฺส ธมฺมสฺส อนนฺตร\-ปจฺจเยน ปจฺจโย. ปุริมา ปุริมา อนิทสฺสน\-อปฺปฏิฆา ขนฺธา ปจฺฉิมานํ ปจฺฉิมานํ อนิทสฺสน\-อปฺปฏิฆานํ ขนฺธานํ อนนฺตร\-ปจฺจเยน ปจฺจโย. อนุโลมํ โคตฺรภุสฺส. อนุโลมํ โวทานสฺส. โคตฺรภุ มคฺคสฺส. โวทานํ มคฺคสฺส. มคฺโค ผลสฺส. ผลํ ผลสฺส. อนุโลมํ ผลสมาปตฺติยา. นิโรธา วุฏฺฐหนฺตสฺส เนว\-สญฺญา\-นา\-สญฺญา\-ยตนํ ผลสมาปตฺติยา อนนฺตร\-ปจฺจเยน ปจฺจโย. (1)}

\csromanmatikaanchor{mtk.295}
\csromantocmarkref{subparagraph}{สมนนฺตร}{mtk.295}
\bookmark[dest={mtk.295},level=5]{สมนนฺตร}
\csromanheader{6}{\textbf{สมนนฺตร}}
\proseitem{35}{\csromanfolio{480}อนิทสฺสน\-อปฺปฏิโฆ ธมฺโม อนิทสฺสน\-อปฺปฏิฆสฺส ธมฺมสฺส สมนนฺตร\-ปจฺจเยน ปจฺจโย. (อนนฺตร\-สทิสํ.)}

\csromanmatikaanchor{mtk.296}
\csromantocmarkref{subparagraph}{สหชาตาทิ}{mtk.296}
\bookmark[dest={mtk.296},level=5]{สหชาตาทิ}
\csromanheader{6}{\textbf{สหชาตาทิ}}
\proseitem{36}{อนิทสฺสน\-สปฺปฏิโฆ ธมฺโม อนิทสฺสน\-สปฺปฏิฆสฺส ธมฺมสฺส สหชาต\-ปจฺจเยน ปจฺจโย. (ปฏิจฺจ\-วาร\-สทิสํ สาธุกํ กาตพฺพํ. อญฺญมญฺญ\-ปจฺจเย ปฏิจฺจวาเร อญฺญมญฺญ\-สทิสํ, นิสฺสยปจฺจเย ปฏิจฺจ\-วาร\-สทิสํ.)}

\csromanheader{2}{\textbf{อุปนิสฺสย}}
\proseitem{37}{สนิทสฺสน\-สปฺปฏิโฆ ธมฺโม อนิทสฺสน\-อปฺปฏิฆสฺส ธมฺมสฺส อุปนิสฺสย\-ปจฺจเยน ปจฺจโย. อารมฺมณู\-ปนิสฺสโย, ปกตู\-ปนิสฺสโย ฯเปฯ. ปกตู\-ปนิสฺสโย\csromandash{}วณฺณสมฺปทํ ปตฺถยมาโน ทานํ เทติ, สีลํ สมาทิยติ, อุโปสถกมฺมํ ฯเปฯ. วณฺณสมฺปทา สทฺธาย ฯเปฯ ปญฺญาย. ราคสฺส ฯเปฯ ปตฺถนาย. กายิกสฺส สุขสฺส, กายิกสฺส ทุกฺขสฺส, มคฺคสฺส, ผลสมาปตฺติยา อุปนิสฺสย\-ปจฺจเยน ปจฺจโย. (1)}

\prose{อนิทสฺสน\-สปฺปฏิโฆ ธมฺโม อนิทสฺสน\-อปฺปฏิฆสฺส ธมฺมสฺส อุปนิสฺสย\-ปจฺจเยน ปจฺจโย. อารมฺมณู\-ปนิสฺสโย, ปกตู\-ปนิสฺสโย ฯเปฯ. ปกตู\-ปนิสฺสโย\csromandash{}จกฺขุสมฺปทํ ปตฺถยมาโน ฯเปฯ. กายสมฺปทํ. สทฺทสมฺปทํ ฯเปฯ. โผฏฺฐพฺพ\-สมฺปทํ ปตฺถยมาโน ทานํ เทติ, สีลํ สมาทิยติ, อุโปสถกมฺมํ. อุตุํ. เสนาสนํ อุปนิสฺสาย ทานํ เทติ, สีลํ สมาทิยติ, อุโปสถกมฺมํ. ฌานํ. วิปสฺสนํ. มคฺคํ. อภิญฺญํ. สมาปตฺติํ อุปฺปาเทติ. ปาณํ หนติ ฯเปฯ สํฆํ ภินฺทติ. จกฺขุสมฺปทา ฯเปฯ. โผฏฺฐพฺพ\-สมฺปทา. อุตุ. เสนาสนํ สทฺธาย ฯเปฯ ปญฺญาย. ราคสฺส ฯเปฯ ปตฺถนาย. กายิกสฺส สุขสฺส, กายิกสฺส ทุกฺขสฺส, มคฺคสฺส, ผลสมาปตฺติยา อุปนิสฺสย\-ปจฺจเยน ปจฺจโย. (1)}

\proseitem{38}{อนิทสฺสน\-อปฺปฏิโฆ ธมฺโม อนิทสฺสน\-อปฺปฏิฆสฺส ธมฺมสฺส อุปนิสฺสย\-ปจฺจเยน ปจฺจโย. อารมฺมณู\-ปนิสฺสโย, อนนฺตรู\-ปนิสฺสโย, ปกตู\-ปนิสฺสโย ฯเปฯ. ปกตู\-ปนิสฺสโย\csromandash{}สทฺธํ อุปนิสฺสาย ทานํ เทติ, สีลํ สมาทิยติ, อุโปสถกมฺมํ. ฌานํ. สมาปตฺติํ อุปฺปาเทติ. \csromanfolio{481}มานํ ชปฺเปติ. ทิฏฺฐิํ คณฺหาติ. สีลํ ฯเปฯ. ปญฺญํ. ราคํ ฯเปฯ. ปตฺถนํ. กายิกํ สุขํ. กายิกํ ทุกฺขํ. โภชนํ อุปนิสฺสาย ทานํ เทติ ฯเปฯ สํฆํ ภินฺทติ. สทฺธา ฯเปฯ. ปญฺญา. ราโค ฯเปฯ. ปตฺถนา. กายิกํ สุขํ. กายิกํ ทุกฺขํ. โภชนํ สทฺธาย ฯเปฯ ปญฺญาย. มคฺคสฺส, ผลสมาปตฺติยา อุปนิสฺสย\-ปจฺจเยน ปจฺจโย. (1)}

\csromanheader{2}{\textbf{ปุเรชาต}}
\proseitem{39}{สนิทสฺสน\-สปฺปฏิโฆ ธมฺโม อนิทสฺสน\-อปฺปฏิฆสฺส ธมฺมสฺส ปุเรชาต\-ปจฺจเยน ปจฺจโย. อารมฺมณ\-ปุเรชาตํ\csromandash{}รูเป อนิจฺจโต ทุกฺขโต อนตฺตโต วิปสฺสติ, อสฺสาเทติ อภินนฺทติ, ตํ อารพฺภ ราโค อุปฺปชฺชติ. ทิฏฺฐิ ฯเปฯ. วิจิกิจฺฉา ฯเปฯ อุทฺธจฺจํ ฯเปฯ โทมนสฺสํ อุปฺปชฺชติ. ทิพฺเพน จกฺขุนา รูปํ ปสฺสติ. รูปายตนํ จกฺขุ\-วิญฺญาณสฺส ปุเรชาต\-ปจฺจเยน ปจฺจโย. (1)}

\prose{อนิทสฺสน\-สปฺปฏิโฆ ธมฺโม อนิทสฺสน\-อปฺปฏิฆสฺส ธมฺมสฺส ปุเรชาต\-ปจฺจเยน ปจฺจโย. อารมฺมณ\-ปุเรชาตํ, วตฺถุ\-ปุเรชาตํ.  อารมฺมณ\-ปุเรชาตํ\csromandash{}จกฺขุํ ฯเปฯ. กายํ. สทฺเท ฯเปฯ. โผฏฺฐพฺเพ อนิจฺจโต ฯเปฯ โทมนสฺสํ อุปฺปชฺชติ. ทิพฺพาย โสตธาตุยา สทฺทํ สุณาติ. สทฺทายตนํ โสตวิญฺญาณสฺส ฯเปฯ. โผฏฺฐพฺพา\-ยตนํ กายวิญฺญาณสฺส ปุเรชาต\-ปจฺจเยน ปจฺจโย. วตฺถุ\-ปุเรชาตํ\csromandash{}จกฺขายตนํ จกฺขุ\-วิญฺญาณสฺส ฯเปฯ. กายายตนํ กายวิญฺญาณสฺส ปุเรชาต\-ปจฺจเยน ปจฺจโย. (1)}

\prose{อนิทสฺสน\-อปฺปฏิโฆ ธมฺโม อนิทสฺสน\-อปฺปฏิฆสฺส ธมฺมสฺส ปุเรชาต\-ปจฺจเยน ปจฺจโย. อารมฺมณ\-ปุเรชาตํ, วตฺถุ\-ปุเรชาตํ.  อารมฺมณ\-ปุเรชาตํ\csromandash{}วตฺถุํ. อิตฺถินฺทฺริยํ. ปุริสินฺทฺริยํ. ชีวิตินฺทฺริยํ. อาโปธาตุํ. กพฬีการํ อาหารํ อนิจฺจโต ฯเปฯ โทมนสฺสํ อุปฺปชฺชติ. วตฺถุ\-ปุเรชาตํ\csromandash{}วตฺถุ อนิทสฺสน\-อปฺปฏิฆานํ ขนฺธานํ ปุเรชาต\-ปจฺจเยน ปจฺจโย. (1)}

\proseitem{40}{สนิทสฺสน\-สปฺปฏิโฆ จ อนิทสฺสน\-อปฺปฏิโฆ จ ธมฺมา อนิทสฺสน\-อปฺปฏิฆสฺส ธมฺมสฺส ปุเรชาต\-ปจฺจเยน ปจฺจโย. อารมฺมณ\-ปุเรชาตํ, วตฺถุ\-ปุเรชาตํ. รูปายตนญฺจ วตฺถุ จ อนิทสฺสน\-อปฺปฏิฆานํ ขนฺธานํ ปุเรชาต\-ปจฺจเยน ปจฺจโย. (1)}

\prose{\csromanfolio{482}อนิทสฺสน\-สปฺปฏิโฆ จ อนิทสฺสน\-อปฺปฏิโฆ จ ธมฺมา อนิทสฺสน\-อปฺปฏิฆสฺส ธมฺมสฺส ปุเรชาต\-ปจฺจเยน ปจฺจโย. อารมฺมณ\-ปุเรชาตํ, วตฺถุ\-ปุเรชาตํ. จกฺขายตนญฺจ วตฺถุ จ ฯเปฯ. โผฏฺฐพฺพายตนญฺจ วตฺถุ จ อนิทสฺสน\-อปฺปฏิฆานํ ขนฺธานํ ปุเรชาต\-ปจฺจเยน ปจฺจโย. (1)}

\prose{สนิทสฺสน\-สปฺปฏิโฆ จ อนิทสฺสน\-สปฺปฏิโฆ จ ธมฺมา อนิทสฺสน\-อปฺปฏิฆสฺส ธมฺมสฺส ปุเรชาต\-ปจฺจเยน ปจฺจโย. อารมฺมณ\-ปุเรชาตํ, วตฺถุ\-ปุเรชาตํ. รูปายตนญฺจ จกฺขายตนญฺจ จกฺขุ\-วิญฺญาณสฺส ปุเรชาต\-ปจฺจเยน ปจฺจโย. (1)}

\csromanheader{2}{\textbf{ปจฺฉาชาต}}
\proseitem{41}{อนิทสฺสน\-อปฺปฏิโฆ ธมฺโม อนิทสฺสน\-อปฺปฏิฆสฺส ธมฺมสฺส ปจฺฉาชาต\-ปจฺจเยน ปจฺจโย. ปจฺฉาชาตา อนิทสฺสน\-อปฺปฏิฆา ขนฺธา ปุเรชาตสฺส อิมสฺส อนิทสฺสน\-อปฺปฏิฆสฺส กายสฺส ปจฺฉาชาต\-ปจฺจเยน ปจฺจโย. (1)}

\prose{อนิทสฺสน\-อปฺปฏิโฆ ธมฺโม สนิทสฺสน\-สปฺปฏิฆสฺส ธมฺมสฺส ปจฺฉาชาต\-ปจฺจเยน ปจฺจโย. ปจฺฉาชาตา อนิทสฺสน\-อปฺปฏิฆา ขนฺธา ปุเรชาตสฺส อิมสฺส สนิทสฺสน\-สปฺปฏิฆสฺส กายสฺส ปจฺฉาชาต\-ปจฺจเยน ปจฺจโย. (2)}

\prose{(เอวํ สตฺต ปญฺหา วิภชิตพฺพา, ติวิธ\-รูปสงฺคโห.) (7)}

\csromanheader{2}{\textbf{อาเสวน}}
\proseitem{42}{อนิทสฺสน\-อปฺปฏิโฆ ธมฺโม อนิทสฺสน\-อปฺปฏิฆสฺส ธมฺมสฺส อาเสวน\-ปจฺจเยน ปจฺจโย. ปุริมา ปุริมา อนิทสฺสน\-อปฺปฏิฆา ขนฺธา ปจฺฉิมานํ ปจฺฉิมานํ อนิทสฺสน\-อปฺปฏิฆานํ ขนฺธานํ อาเสวน\-ปจฺจเยน ปจฺจโย. อนุโลมํ โคตฺรภุสฺส. อนุโลมํ โวทานสฺส. โคตฺรภุ มคฺคสฺส. โวทานํ มคฺคสฺส อาเสวน\-ปจฺจเยน ปจฺจโย. (1)}

\csromanheader{2}{\textbf{กมฺม}}
\proseitem{43}{อนิทสฺสน\-อปฺปฏิโฆ ธมฺโม อนิทสฺสน\-อปฺปฏิฆสฺส ธมฺมสฺส กมฺมปจฺจเยน ปจฺจโย. สหชาตา, นานากฺขณิกา.  สหชาตา อนิทสฺสน\-อปฺปฏิฆา เจตนา สมฺปยุตฺตกานํ ขนฺธานํ อนิทสฺสน\-อปฺปฏิฆานญฺจ \csromanfolio{483}จิตฺต\-สมุฏฺฐานํ รูปานํ กมฺมปจฺจเยน ปจฺจโย. ปฏิสนฺธิ\-กฺขเณ ฯเปฯ. นานากฺขณิกา อนิทสฺสน\-อปฺปฏิฆา เจตนา วิปากานํ ขนฺธานํ อนิทสฺสน\-อปฺปฏิฆานญฺจ กฏตฺตารูปานํ กมฺมปจฺจเยน ปจฺจโย. (1)}

\prose{อนิทสฺสน\-อปฺปฏิโฆ ธมฺโม สนิทสฺสน\-สปฺปฏิฆสฺส ธมฺมสฺส กมฺมปจฺจเยน ปจฺจโย. สหชาตา, นานากฺขณิกา.  สหชาตา อนิทสฺสน\-อปฺปฏิฆา เจตนา สนิทสฺสน\-สปฺปฏิฆานํ จิตฺต\-สมุฏฺฐานานํ รูปานํ กมฺมปจฺจเยน ปจฺจโย. ปฏิสนฺธิ\-กฺขเณ ฯเปฯ. นานากฺขณิกา อนิทสฺสน\-อปฺปฏิฆา เจตนา สนิทสฺสน\-สปฺปฏิฆานํ กฏตฺตารูปานํ กมฺมปจฺจเยน ปจฺจโย. (2)}

\prose{(เอวํ สตฺต ปญฺหา สหชาตา นานากฺขณิกา อิมินา การเณน วิภชิตพฺพา, ติวิธ\-รูปสงฺคโห.) (7)}

\csromanheader{2}{\textbf{วิปาก}}
\proseitem{44}{อนิทสฺสน\-อปฺปฏิโฆ ธมฺโม อนิทสฺสน\-อปฺปฏิฆสฺส ธมฺมสฺส วิปาก\-ปจฺจเยน ปจฺจโย. วิปาโก อนิทสฺสน\-อปฺปฏิโฆ เอโก ขนฺโธ ติณฺณนฺนํ ขนฺธานํ อนิทสฺสน\-อปฺปฏิฆานํ จิตฺตสมุฏฺฐานานญฺจ รูปานํ วิปาก\-ปจฺจเยน ปจฺจโย ฯเปฯ เทฺว ขนฺธา ฯเปฯ. ปฏิสนฺธิ\-กฺขเณ อนิทสฺสน\-อปฺปฏิโฆ เอโก ขนฺโธ ติณฺณนฺนํ ขนฺธานํ อนิทสฺสนอปฺปฏิฆานญฺจ กฏตฺตารูปานํ วิปาก\-ปจฺจเยน ปจฺจโย ฯเปฯ. ขนฺธา วตฺถุสฺส วิปาก\-ปจฺจเยน ปจฺจโย. (1)}

\prose{อนิทสฺสน\-อปฺปฏิโฆ ธมฺโม สนิทสฺสน\-สปฺปฏิฆสฺส ธมฺมสฺส วิปาก\-ปจฺจเยน ปจฺจโย. วิปากา อนิทสฺสน\-อปฺปฏิฆา ขนฺธา สนิทสฺสน\-สปฺปฏิฆานํ จิตฺต\-สมุฏฺฐานานํ รูปานํ วิปาก\-ปจฺจเยน ปจฺจโย. ปฏิสนฺธิ\-กฺขเณ อนิทสฺสน\-อปฺปฏิฆา ขนฺธา สนิทสฺสน\-สปฺปฏิฆานํ กฏตฺตารูปานํ วิปาก\-ปจฺจเยน ปจฺจโย. (2)}

\prose{(เอวํ สตฺต ปญฺหา วิตฺถาเรตพฺพา, ปวตฺติ\-ปฏิสนฺธิ.) (7)}

\csromanheader{2}{\textbf{อาหาร}}
\proseitem{45}{อนิทสฺสน\-อปฺปฏิโฆ ธมฺโม อนิทสฺสน\-อปฺปฏิฆสฺส ธมฺมสฺส อาหารปจฺจเยน ปจฺจโย. อนิทสฺสน\-อปฺปฏิฆา อาหารา สมฺปยุตฺตกานํ ขนฺธานํ อนิทสฺสน\-อปฺปฏิฆานญฺจ จิตฺต\-สมุฏฺฐานานํ รูปานํ อาหารปจฺจเยน ปจฺจโย. ปฏิสนฺธิ\-กฺขเณ อนิทสฺสน\-อปฺปฏิฆา อาหารา สมฺปยุตฺตกานํ \csromanfolio{484}ขนฺธานํ อนิทสฺสนอปฺปฏิฆานญฺจ กฏตฺตารูปานํ อาหารปจฺจเยน ปจฺจโย. กพฬีกาโร อาหาโร อิมสฺส อนิทสฺสน\-อปฺปฏิฆสฺส กายสฺส อาหารปจฺจเยน ปจฺจโย. (1)}

\prose{อนิทสฺสน\-อปฺปฏิโฆ ธมฺโม สนิทสฺสน\-สปฺปฏิฆสฺส ธมฺมสฺส อาหารปจฺจเยน ปจฺจโย. อนิทสฺสน\-อปฺปฏิฆา อาหารา สนิทสฺสน\-สปฺปฏิฆานํ จิตฺต\-สมุฏฺฐานานํ รูปานํ อาหารปจฺจเยน ปจฺจโย. ปฏิสนฺธิ\-กฺขเณ อนิทสฺสน\-อปฺปฏิฆา อาหารา สนิทสฺสน\-สปฺปฏิฆานํ กฏตฺตารูปานํ อาหารปจฺจเยน ปจฺจโย. กพฬีกาโร อาหาโร อิมสฺส สนิทสฺสน\-สปฺปฏิฆสฺส กายสฺส อาหารปจฺจเยน ปจฺจโย. (2)}

\prose{(เอวํ สตฺต ปญฺหา ปวตฺติ\-ปฏิสนฺธิ วิภชิตพฺพา, สตฺตสุปิ กพฬีกาโร อาหาโร กาตพฺโพ.) (7)}

\csromanheader{2}{\textbf{อินฺทฺริย}}
\proseitem{46}{อนิทสฺสน\-สปฺปฏิโฆ ธมฺโม อนิทสฺสน\-อปฺปฏิฆสฺส ธมฺมสฺส อินฺทฺริย\-ปจฺจเยน ปจฺจโย. จกฺขุนฺทฺริยํ จกฺขุ\-วิญฺญาณสฺส ฯเปฯ. กายินฺทฺริยํ กายวิญฺญาณสฺส อินฺทฺริย\-ปจฺจเยน ปจฺจโย. (1)}

\prose{อนิทสฺสน\-อปฺปฏิโฆ ธมฺโม อนิทสฺสน\-อปฺปฏิฆสฺส ธมฺมสฺส อินฺทฺริย\-ปจฺจเยน ปจฺจโย. อนิทสฺสน\-อปฺปฏิฆา อินฺทฺริยา สมฺปยุตฺตกานํ ขนฺธานํ อนิทสฺสน\-อปฺปฏิฆานํ จิตฺตสมุฏฺฐานานญฺจ รูปานํ อินฺทฺริย\-ปจฺจเยน ปจฺจโย. ปฏิสนฺธิ\-กฺขเณ อนิทสฺสน\-อปฺปฏิฆา อินฺทฺริยา สมฺปยุตฺตกานํ ขนฺธานํ อนิทสฺสน\-อปฺปฏิฆานญฺจ กฏตฺตารูปานํ อินฺทฺริย\-ปจฺจเยน ปจฺจโย. รูปชีวิตินฺทฺริยํ อนิทสฺสน\-อปฺปฏิฆานํ กฏตฺตารูปานํ อินฺทฺริย\-ปจฺจเยน ปจฺจโย. (1)}

\prose{อนิทสฺสน\-อปฺปฏิโฆ ธมฺโม สนิทสฺสน\-สปฺปฏิฆสฺส ธมฺมสฺส อินฺทฺริย\-ปจฺจเยน ปจฺจโย. อนิทสฺสน\-อปฺปฏิฆา อินฺทฺริยา สนิทสฺสน\-สปฺปฏิฆานํ จิตฺต\-สมุฏฺฐานานํ รูปานํ อินฺทฺริย\-ปจฺจเยน ปจฺจโย. ปฏิสนฺธิ\-กฺขเณ อนิทสฺสน\-อปฺปฏิฆา อินฺทฺริยา สนิทสฺสน\-สปฺปฏิฆานํ กฏตฺตารูปานํ อินฺทฺริย\-ปจฺจเยน ปจฺจโย. รูปชีวิตินฺทฺริยํ สนิทสฺสน\-สปฺปฏิฆานํ กฏตฺตารูปานํ อินฺทฺริย\-ปจฺจเยน ปจฺจโย. (2)}

\prose{(เอวํ ปวตฺติ\-ปฏิสนฺธิ สตฺต ปญฺหา วิภชิตพฺพา, รูปชีวิตินฺทฺริยญฺจ อนฺเต อนฺเต.) (7)}

\proseitem{47}{\csromanfolio{485}อนิทสฺสน\-สปฺปฏิโฆ จ อนิทสฺสน\-อปฺปฏิโฆ จ ธมฺมา อนิทสฺสน\-อปฺปฏิฆสฺส ธมฺมสฺส อินฺทฺริย\-ปจฺจเยน ปจฺจโย. จกฺขุนฺทฺริยญฺจ จกฺขุวิญฺญาณญฺจ จกฺขุ\-วิญฺญาณสหคตานํ ขนฺธานํ อินฺทฺริย\-ปจฺจเยน ปจฺจโย ฯเปฯ. กายินฺทฺริยญฺจ กายวิญฺญาณญฺจ กาย\-วิญฺญาณสหคตานํ ขนฺธานํ อินฺทฺริย\-ปจฺจเยน ปจฺจโย. (1)}

\csromanheader{2}{\textbf{ฌานาทิ}}
\proseitem{48}{อนิทสฺสน\-อปฺปฏิโฆ ธมฺโม อนิทสฺสน\-อปฺปฏิฆสฺส ธมฺมสฺส ฌานปจฺจเยน ปจฺจโย. มคฺคปจฺจเยน ปจฺจโย. สมฺปยุตฺต\-ปจฺจเยน ปจฺจโย. อนิทสฺสน\-อปฺปฏิโฆ เอโก ขนฺโธ ติณฺณนฺนํ ขนฺธานํ สมฺปยุตฺต\-ปจฺจเยน ปจฺจโย ฯเปฯ เทฺว ขนฺธา ฯเปฯ. ปฏิสนฺธิ\-กฺขเณ ฯเปฯ.}

\csromanheader{2}{\textbf{วิปฺปยุตฺต}}
\proseitem{49}{อนิทสฺสน\-สปฺปฏิโฆ ธมฺโม อนิทสฺสน\-อปฺปฏิฆสฺส ธมฺมสฺส วิปฺปยุตฺต\-ปจฺจเยน ปจฺจโย. ปุเรชาตํ จกฺขายตนํ จกฺขุ\-วิญฺญาณสฺส ฯเปฯ. กายายตนํ กายวิญฺญาณสฺส วิปฺปยุตฺต\-ปจฺจเยน ปจฺจโย. (1)}

\prose{อนิทสฺสน\-อปฺปฏิโฆ ธมฺโม อนิทสฺสน\-อปฺปฏิฆสฺส ธมฺมสฺส วิปฺปยุตฺต\-ปจฺจเยน ปจฺจโย. สหชาตํ, ปุเรชาตํ, ปจฺฉาชาตํ.  สหชาตา อนิทสฺสน\-อปฺปฏิฆา ขนฺธา อนิทสฺสน\-อปฺปฏิฆานํ จิตฺต\-สมุฏฺฐานานํ รูปานํ วิปฺปยุตฺต\-ปจฺจเยน ปจฺจโย. ปฏิสนฺธิ\-กฺขเณ อนิทสฺสน\-อปฺปฏิฆา ขนฺธา อนิทสฺสน\-อปฺปฏิฆานํ กฏตฺตารูปานํ วิปฺปยุตฺต\-ปจฺจเยน ปจฺจโย. ขนฺธา วตฺถุสฺส วิปฺปยุตฺต\-ปจฺจเยน ปจฺจโย. วตฺถุ ขนฺธานํ วิปฺปยุตฺต\-ปจฺจเยน ปจฺจโย. ปุเรชาตํ วตฺถุ อนิทสฺสน\-อปฺปฏิฆานํ ขนฺธานํ วิปฺปยุตฺต\-ปจฺจเยน ปจฺจโย. ปจฺฉาชาตา อนิทสฺสน\-อปฺปฏิฆา ขนฺธา ปุเรชาตสฺส อิมสฺส อนิทสฺสน\-อปฺปฏิฆสฺส กายสฺส วิปฺปยุตฺต\-ปจฺจเยน ปจฺจโย. (1)}

\prose{อนิทสฺสน\-อปฺปฏิโฆ ธมฺโม สนิทสฺสน\-สปฺปฏิฆสฺส ธมฺมสฺส วิปฺปยุตฺต\-ปจฺจเยน ปจฺจโย. สหชาตํ, ปจฺฉาชาตํ.  สหชาตา อนิทสฺสน\-อปฺปฏิฆา ขนฺธา สนิทสฺสน\-สปฺปฏิฆานํ จิตฺต\-สมุฏฺฐานานํ รูปานํ วิปฺปยุตฺต\-ปจฺจเยน ปจฺจโย. ปฏิสนฺธิ\-กฺขเณ อนิทสฺสน\-อปฺปฏิฆา ขนฺธา สนิทสฺสน\-สปฺปฏิฆานํ กฏตฺตารูปานํ วิปฺปยุตฺต\-ปจฺจเยน ปจฺจโย. ปจฺฉาชาตา \csromanfolio{486}อนิทสฺสน\-อปฺปฏิฆา ขนฺธา ปุเรชาตสฺส อิมสฺส สนิทสฺสน\-สปฺปฏิฆสฺส กายสฺส วิปฺปยุตฺต\-ปจฺจเยน ปจฺจโย. (2) (อวเสสา ปญฺจ ปญฺหา เอวํ วิตฺถาเรตพฺพา. สหชาตา. ปจฺฉาชาตา.)}

\csromanheader{2}{\textbf{อตฺถิ}}
\proseitem{50}{สนิทสฺสน\-สปฺปฏิโฆ ธมฺโม อนิทสฺสน\-อปฺปฏิฆสฺส ธมฺมสฺส อตฺถิปจฺจเยน ปจฺจโย. ปุเรชาตํ\csromandash{}รูเป อนิจฺจโต ฯเปฯ โทมนสฺสํ อุปฺปชฺชติ. ทิพฺเพน จกฺขุนา รูปํ ปสฺสติ. รูปายตนํ จกฺขุ\-วิญฺญาณสฺส อตฺถิปจฺจเยน ปจฺจโย. (1)}

\proseitem{51}{อนิทสฺสน\-สปฺปฏิโฆ ธมฺโม อนิทสฺสน\-สปฺปฏิฆสฺส ธมฺมสฺส อตฺถิปจฺจเยน ปจฺจโย. อนิทสฺสน\-สปฺปฏิฆํ เอกํ มหาภูตํ ทฺวินฺนํ มหาภูตานํ อตฺถิปจฺจเยน ปจฺจโย. เทฺว มหาภูตา เอกสฺส มหาภูตสฺส อตฺถิปจฺจเยน ปจฺจโย. อนิทสฺสน\-สปฺปฏิฆา มหาภูตา อนิทสฺสน\-สปฺปฏิฆานํ จิตฺต\-สมุฏฺฐานานํ รูปานํ, กฏตฺตารูปานํ, อุปาทารูปานํ อตฺถิปจฺจเยน ปจฺจโย. โผฏฺฐพฺพา\-ยตนํ จกฺขา\-ยตนสฺส ฯเปฯ รสายตนสฺส อตฺถิปจฺจเยน ปจฺจโย. พาหิรํ. อาหารสมุฏฺฐานํ. อุตุสมุฏฺฐานํ. เอกํ มหาภูตํ ทฺวินฺนํ มหาภูตานํ อตฺถิปจฺจเยน ปจฺจโย. เทฺว มหาภูตา เอกสฺส มหาภูตสฺส อตฺถิปจฺจเยน ปจฺจโย. อุตุสมุฏฺฐานา มหาภูตา อนิทสฺสน\-สปฺปฏิฆานํ อุปาทารูปานํ อตฺถิปจฺจเยน ปจฺจโย. อสญฺญสตฺตานํ อนิทสฺสน\-สปฺปฏิฆํ เอกํ มหาภูตํ ทฺวินฺนํ มหาภูตานํ อตฺถิปจฺจเยน ปจฺจโย. เทฺว มหาภูตา ฯเปฯ. (1)}

\prose{อนิทสฺสน\-สปฺปฏิโฆ ธมฺโม สนิทสฺสน\-สปฺปฏิฆสฺส ธมฺมสฺส อตฺถิปจฺจเยน ปจฺจโย. (2) (ปฏิจฺจวาเร นิสฺสยปจฺจย\-สทิสํ.)}

\prose{อนิทสฺสน\-สปฺปฏิโฆ ธมฺโม อนิทสฺสน\-อปฺปฏิฆสฺส ธมฺมสฺส อตฺถิปจฺจเยน ปจฺจโย. สหชาตํ, ปุเรชาตํ.  สหชาตา อนิทสฺสน\-สปฺปฏิฆา มหาภูตา อนิทสฺสน\-อปฺปฏิฆานํ จิตฺต\-สมุฏฺฐานานํ รูปานํ, กฏตฺตารูปานํ, อุปาทารูปานํ อตฺถิปจฺจเยน ปจฺจโย. (ยาว อสญฺญสตฺตา วิตฺถาเรตพฺพา.) ปุเรชาตํ จกฺขุํ ฯเปฯ. กายํ. สทฺเท ฯเปฯ. โผฏฺฐพฺเพ อนิจฺจโต ฯเปฯ โทมนสฺสํ อุปฺปชฺชติ. จกฺขายตนํ จกฺขุ\-วิญฺญาณสฺส ฯเปฯ. โผฏฺฐพฺพา\-ยตนํ กายวิญฺญาณสฺส อตฺถิปจฺจเยน ปจฺจโย. (3)}

\prose{\csromanfolio{487}(อวเสสา จตฺตาโร ปญฺหา วิตฺถาเรตพฺพา. ปฏิจฺจวาเร สหชาตปจฺจย\-สทิสา. นินฺนานากรณา.) (7)}

\proseitem{52}{อนิทสฺสน\-อปฺปฏิโฆ ธมฺโม อนิทสฺสน\-อปฺปฏิฆสฺส ธมฺมสฺส อตฺถิปจฺจเยน ปจฺจโย. สหชาตํ, ปุเรชาตํ, ปจฺฉาชาตํ, อาหารํ, อินฺทฺริยํ.  สหชาโต อนิทสฺสน\-อปฺปฏิโฆ เอโก ขนฺโธ ติณฺณนฺนํ ขนฺธานํ อนิทสฺสน\-อปฺปฏิฆานํ จิตฺตสมุฏฺฐานานญฺจ รูปานํ อตฺถิปจฺจเยน ปจฺจโย. เทฺว ขนฺธา ฯเปฯ. ปฏิสนฺธิ\-กฺขเณ ฯเปฯ. อาโปธาตุ อนิทสฺสน\-อปฺปฏิฆานํ จิตฺต\-สมุฏฺฐานานํ รูปานํ, กฏตฺตารูปานํ, อุปาทารูปานํ ฯเปฯ. อาโปธาตุ อิตฺถินฺทฺริยสฺส ฯเปฯ กพฬีการา\-หารสฺส จ อตฺถิปจฺจเยน ปจฺจโย. พาหิรํ. อาหารสมุฏฺฐานํ. อุตุสมุฏฺฐานํ. อสญฺญสตฺตานํ อาโปธาตุ อนิทสฺสน\-อปฺปฏิฆานํ กฏตฺตารูปานํ, อุปาทารูปานํ อตฺถิปจฺจเยน ปจฺจโย. ปุเรชาตํ วตฺถุํ. อิตฺถินฺทฺริยํ. ปุริสินฺทฺริยํ. ชีวิตินฺทฺริยํ. อาโปธาตุํ. กพฬีการํ อาหารํ อนิจฺจโต ฯเปฯ โทมนสฺสํ อุปฺปชฺชติ. วตฺถุ อนิทสฺสน\-อปฺปฏิฆานํ ขนฺธานํ อตฺถิปจฺจเยน ปจฺจโย. ปจฺฉาชาตา อนิทสฺสน\-อปฺปฏิฆา ขนฺธา ปุเรชาตสฺส อิมสฺส อนิทสฺสน\-อปฺปฏิฆสฺส กายสฺส อตฺถิปจฺจเยน ปจฺจโย. กพฬีกาโร อาหาโร อิมสฺส อนิทสฺสน\-อปฺปฏิฆสฺส กายสฺส อตฺถิปจฺจเยน ปจฺจโย. รูปชีวิตินฺทฺริยํ อนิทสฺสน\-อปฺปฏิฆานํ กฏตฺตารูปานํ อตฺถิปจฺจเยน ปจฺจโย. (เอวํ อวเสสา ฉ ปญฺหา วิภชิตพฺพา. สหชาตํ. ปจฺฉาชาตํ. อาหารํ. อินฺทฺริยมฺปิ กาตพฺพา.) (7)}

\proseitem{53}{สนิทสฺสน\-สปฺปฏิโฆ จ อนิทสฺสน\-อปฺปฏิโฆ จ ธมฺมา อนิทสฺสน\-อปฺปฏิฆสฺส ธมฺมสฺส อตฺถิปจฺจเยน ปจฺจโย. ปุเรชาตํ รูปายตนญฺจ วตฺถุ จ อนิทสฺสน\-อปฺปฏิฆานํ ขนฺธานํ อตฺถิปจฺจเยน ปจฺจโย. (1)}

\prose{อนิทสฺสน\-สปฺปฏิโฆ จ อนิทสฺสน\-อปฺปฏิโฆ จ ธมฺมา สนิทสฺสน\-สปฺปฏิฆสฺส ธมฺมสฺส อตฺถิปจฺจเยน ปจฺจโย. อนิทสฺสน\-อปฺปฏิฆา ขนฺธา จ มหาภูตา จ สนิทสฺสน\-สปฺปฏิฆานํ จิตฺต\-สมุฏฺฐานานํ รูปานํ อตฺถิปจฺจเยน ปจฺจโย. ปฏิสนฺธิ\-กฺขเณ. (สํขิตฺตํ. อสญฺญสตฺตานญฺจ กาตพฺพา.) (1)}

\prose{อนิทสฺสน\-สปฺปฏิโฆ จ อนิทสฺสน\-อปฺปฏิโฆ จ ธมฺมา อนิทสฺสน\-สปฺปฏิฆสฺส ธมฺมสฺส อตฺถิปจฺจเยน ปจฺจโย. (สํขิตฺตํ.) (2)}

\csromanmatikaanchor{mtk.297}
\csromanmatikaanchor{mtk.298}
\csromantocmarknopage{subparagraph}{1. ปจฺจยานุโลม 2. สงฺขฺยาวาร}{mtk.297}
\bookmark[dest={mtk.297},level=5]{1. ปจฺจยานุโลม 2. สงฺขฺยาวาร}
\csromantocmarkref{subparagraph}{สุทฺธ}{mtk.298}
\bookmark[dest={mtk.298},level=5]{สุทฺธ}
\proseitem{54}{\csromanfolio{488}อนิทสฺสน\-สปฺปฏิโฆ จ อนิทสฺสน\-อปฺปฏิโฆ จ ธมฺมา อนิทสฺสน\-อปฺปฏิฆสฺส ธมฺมสฺส อตฺถิปจฺจเยน ปจฺจโย. สหชาตํ, ปุเรชาตํ.  สหชาตา อนิทสฺสน\-อปฺปฏิฆา ขนฺธา จ มหาภูตา จ อนิทสฺสน\-อปฺปฏิฆานํ จิตฺต\-สมุฏฺฐานานํ รูปานํ ฯเปฯ. (ยาว อสญฺญสตฺตา กาตพฺพา.) ปุเรชาตํ จกฺขายตนญฺจ วตฺถุ จ ฯเปฯ. โผฏฺฐพฺพายตนญฺจ วตฺถุ จ อนิทสฺสน\-อปฺปฏิฆานํ ขนฺธานํ อตฺถิปจฺจเยน ปจฺจโย. (3) (อวเสสา จตสฺโส ปญฺหา วิภชิตพฺพา.) (7)}

\proseitem{55}{สนิทสฺสน\-สปฺปฏิโฆ จ อนิทสฺสน\-สปฺปฏิโฆ จ ธมฺมา อนิทสฺสน\-อปฺปฏิฆสฺส ธมฺมสฺส อตฺถิปจฺจเยน ปจฺจโย. ปุเรชาตํ รูปายตนญฺจ จกฺขายตนญฺจ จกฺขุ\-วิญฺญาณสฺส อตฺถิปจฺจเยน ปจฺจโย. (1)}

\prose{สนิทสฺสน\-สปฺปฏิโฆ จ อนิทสฺสน\-สปฺปฏิโฆ จ อนิทสฺสน\-อปฺปฏิโฆ จ ธมฺมา อนิทสฺสน\-อปฺปฏิฆสฺส ธมฺมสฺส อตฺถิปจฺจเยน ปจฺจโย. สหชาตํ, ปุเรชาตํ.  รูปายตนญฺจ จกฺขายตนญฺจ จกฺขุวิญฺญาณญฺจ จกฺขุ\-วิญฺญาณสหคตานํ ขนฺธานํ อตฺถิปจฺจเยน ปจฺจโย. (1)}

\prose{(นตฺถิ\-วิคต\-ปจฺจยํ อนนฺตร\-สทิสํ. อวิคต\-ปจฺจยํ อตฺถิสทิสํ.)}

\csromanmarkh{1. ปจฺจยานุโลม}\csromanmarkhii{2. สงฺขฺยาวาร}\csromanheadernomark{2}{1. \textbf{ปจฺจยานุโลม 2. สงฺขฺยาวาร}}
\proseitem{56}{เหตุยา สตฺต, อารมฺมเณ ตีณิ, อธิปติยา นว, อนนฺตเร เอกํ, สมนนฺตเร เอกํ, สหชาเต เอกวีส, อญฺญมญฺเญ ฉ, นิสฺสเย เอกวีส, อุปนิสฺสเย ตีณิ, ปุเรชาเต ฉ, ปจฺฉาชาเต สตฺต, อาเสวเน เอกํ, กมฺเม สตฺต, วิปาเก สตฺต, อาหาเร สตฺต, อินฺทฺริเย นว, ฌาเน สตฺต, มคฺเค สตฺต, สมฺปยุตฺเต เอกํ, วิปฺปยุตฺเต อฏฺฐ, อตฺถิยา ปญฺจวีส, นตฺถิยา เอกํ, วิคเต เอกํ, อวิคเต ปญฺจวีส.}

\csromanmatikaanchor{mtk.299}
\csromantocmarkref{subparagraph}{เหตุสภาค}{mtk.299}
\bookmark[dest={mtk.299},level=5]{เหตุสภาค}
\csromanheader{6}{\textbf{เหตุสภาค}}
\proseitem{57}{เหตุปจฺจยา อธิปติยา สตฺต, สหชาเต สตฺต, อญฺญมญฺเญ เอกํ, นิสฺสเย สตฺต, วิปาเก สตฺต, อินฺทฺริเย สตฺต, มคฺเค สตฺต, สมฺปยุตฺเต เอกํ, วิปฺปยุตฺเต สตฺต, อตฺถิยา สตฺต, อวิคเต สตฺต.}

\csromanmatikaanchor{mtk.300}
\csromantocmarkref{subparagraph}{เหตุสามญฺญฆฏนา (9)}{mtk.300}
\bookmark[dest={mtk.300},level=5]{เหตุสามญฺญฆฏนา (9)}
\proseitem{58}{\csromanfolio{489}เหตุ สหชาต นิสฺสย อตฺถิ อวิคตนฺติ สตฺต. เหตุ สหชาต อญฺญมญฺญ นิสฺสย อตฺถิ อวิคตนฺติ เอกํ. เหตุ สหชาต อญฺญมญฺญ นิสฺสย สมฺปยุตฺต อตฺถิ อวิคตนฺติ เอกํ. เหตุ สหชาต นิสฺสย วิปฺปยุตฺต อตฺถิ อวิคตนฺติ สตฺต. (อวิปาก-4)}

\prosewithcloserruled{เหตุ สหชาต นิสฺสย วิปาก อตฺถิ อวิคตนฺติ สตฺต. เหตุ สหชาต อญฺญมญฺญ นิสฺสย วิปาก อตฺถิ อวิคตนฺติ เอกํ. เหตุ สหชาต อญฺญมญฺญ นิสฺสย วิปาก สมฺปยุตฺต อตฺถิ อวิคตนฺติ เอกํ. เหตุ สหชาต นิสฺสย วิปาก วิปฺปยุตฺต อตฺถิ อวิคตนฺติ สตฺต. เหตุ สหชาต อญฺญมญฺญ นิสฺสย วิปาก วิปฺปยุตฺต อตฺถิ อวิคตนฺติ เอกํ. (สวิปาก-5) (เอวํ สพฺโพ คณนวาโร คเณตพฺโพ.)}{อนุโลมํ.}
\csromanmatikaanchor{mtk.301}
\csromantocmarkref{subparagraph}{2. ปจฺจนียุทฺธาร}{mtk.301}
\bookmark[dest={mtk.301},level=5]{2. ปจฺจนียุทฺธาร}
\csromanheader{6}{2. ปจฺจนียุ\-ทฺธาร}
\proseitem{59}{สนิทสฺสน\-สปฺปฏิโฆ ธมฺโม อนิทสฺสน\-อปฺปฏิฆสฺส ธมฺมสฺส อารมฺมณ\-ปจฺจเยน ปจฺจโย. อุปนิสฺสย\-ปจฺจเยน ปจฺจโย. ปุเรชาต\-ปจฺจเยน ปจฺจโย. (1)}

\prose{อนิทสฺสน\-สปฺปฏิโฆ ธมฺโม อนิทสฺสน\-สปฺปฏิฆสฺส ธมฺมสฺส สหชาต\-ปจฺจเยน ปจฺจโย. (1)}

\prose{อนิทสฺสน\-สปฺปฏิโฆ ธมฺโม สนิทสฺสน\-สปฺปฏิฆสฺส ธมฺมสฺส สหชาต\-ปจฺจเยน ปจฺจโย. (2)}

\prose{อนิทสฺสน\-สปฺปฏิโฆ ธมฺโม อนิทสฺสน\-อปฺปฏิฆสฺส ธมฺมสฺส อารมฺมณ\-ปจฺจเยน ปจฺจโย. สหชาต\-ปจฺจเยน ปจฺจโย. อุปนิสฺสย\-ปจฺจเยน ปจฺจโย. ปุเรชาต\-ปจฺจเยน ปจฺจโย. (3)}

\prose{อนิทสฺสน\-สปฺปฏิโฆ ธมฺโม สนิทสฺสน\-สปฺปฏิฆสฺส จ อนิทสฺสน\-อปฺปฏิฆสฺส จ ธมฺมสฺส สหชาต\-ปจฺจเยน ปจฺจโย. (4)}

\prose{อนิทสฺสน\-สปฺปฏิโฆ ธมฺโม อนิทสฺสน\-สปฺปฏิฆสฺส จ อนิทสฺสน\-อปฺปฏิฆสฺส จ ธมฺมสฺส สหชาต\-ปจฺจเยน ปจฺจโย. (5)}

\prose{\csromanfolio{490}อนิทสฺสน\-สปฺปฏิโฆ ธมฺโม สนิทสฺสน\-สปฺปฏิฆสฺส จ อนิทสฺสน\-สปฺปฏิฆสฺส จ ธมฺมสฺส สหชาต\-ปจฺจเยน ปจฺจโย. (6)}

\prose{อนิทสฺสน\-สปฺปฏิโฆ ธมฺโม สนิทสฺสน\-สปฺปฏิฆสฺส จ อนิทสฺสน\-สปฺปฏิฆสฺส จ อนิทสฺสน\-อปฺปฏิฆสฺส จ ธมฺมสฺส สหชาต\-ปจฺจเยน ปจฺจโย. (7)}

\proseitem{60}{อนิทสฺสน\-อปฺปฏิโฆ ธมฺโม อนิทสฺสน\-อปฺปฏิฆสฺส ธมฺมสฺส อารมฺมณ\-ปจฺจเยน ปจฺจโย. สหชาต\-ปจฺจเยน ปจฺจโย. อุปนิสฺสย\-ปจฺจเยน ปจฺจโย. ปุเรชาต\-ปจฺจเยน ปจฺจโย. ปจฺฉาชาต\-ปจฺจเยน ปจฺจโย. กมฺมปจฺจเยน ปจฺจโย. อาหารปจฺจเยน ปจฺจโย. อินฺทฺริย\-ปจฺจเยน ปจฺจโย. (1)}

\prose{อนิทสฺสน\-อปฺปฏิโฆ ธมฺโม สนิทสฺสน\-สปฺปฏิฆสฺส ธมฺมสฺส สหชาต\-ปจฺจเยน ปจฺจโย. ปจฺฉาชาต\-ปจฺจเยน ปจฺจโย. กมฺมปจฺจเยน ปจฺจโย. อาหารปจฺจเยน ปจฺจโย. อินฺทฺริย\-ปจฺจเยน ปจฺจโย. (2)}

\prose{อนิทสฺสน\-อปฺปฏิโฆ ธมฺโม อนิทสฺสน\-สปฺปฏิฆสฺส ธมฺมสฺส สหชาต\-ปจฺจเยน ปจฺจโย. ปจฺฉาชาต\-ปจฺจเยน ปจฺจโย. กมฺมปจฺจเยน ปจฺจโย. อาหารปจฺจเยน ปจฺจโย. อินฺทฺริย\-ปจฺจเยน ปจฺจโย. (3)}

\prose{อนิทสฺสน\-อปฺปฏิโฆ ธมฺโม สนิทสฺสน\-สปฺปฏิฆสฺส จ อนิทสฺสน\-อปฺปฏิฆสฺส จ ธมฺมสฺส สหชาต\-ปจฺจเยน ปจฺจโย. ปจฺฉาชาต\-ปจฺจเยน ปจฺจโย. กมฺมปจฺจเยน ปจฺจโย. อาหารปจฺจเยน ปจฺจโย. อินฺทฺริย\-ปจฺจเยน ปจฺจโย. (4)}

\prose{อนิทสฺสน\-อปฺปฏิโฆ ธมฺโม อนิทสฺสน\-สปฺปฏิฆสฺส จ อนิทสฺสน\-อปฺปฏิฆสฺส จ ธมฺมสฺส สหชาต\-ปจฺจเยน ปจฺจโย. ปจฺฉาชาต\-ปจฺจเยน ปจฺจโย. กมฺมปจฺจเยน ปจฺจโย. อาหารปจฺจเยน ปจฺจโย. อินฺทฺริย\-ปจฺจเยน ปจฺจโย. (5)}

\prose{อนิทสฺสน\-อปฺปฏิโฆ ธมฺโม สนิทสฺสน\-สปฺปฏิฆสฺส จ อนิทสฺสน\-สปฺปฏิฆสฺส จ ธมฺมสฺส สหชาต\-ปจฺจเยน ปจฺจโย. ปจฺฉาชาต\-ปจฺจเยน ปจฺจโย. กมฺมปจฺจเยน ปจฺจโย. อาหารปจฺจเยน ปจฺจโย. อินฺทฺริย\-ปจฺจเยน ปจฺจโย. (6)}

\prose{อนิทสฺสน\-อปฺปฏิโฆ ธมฺโม สนิทสฺสน\-สปฺปฏิฆสฺส จ อนิทสฺสน\-สปฺปฏิฆสฺส จ อนิทสฺสน\-อปฺปฏิฆสฺส จ ธมฺมสฺส สหชาต\-ปจฺจเยน \csromanfolio{491}ปจฺจโย. ปจฺฉาชาต\-ปจฺจเยน ปจฺจโย. กมฺมปจฺจเยน ปจฺจโย. อาหารปจฺจเยน ปจฺจโย. อินฺทฺริย\-ปจฺจเยน ปจฺจโย. (7)}

\proseitem{61}{สนิทสฺสน\-สปฺปฏิโฆ จ อนิทสฺสน\-อปฺปฏิโฆ จ ธมฺมา อนิทสฺสน\-อปฺปฏิฆสฺส ธมฺมสฺส ปุเรชาตํ. (1)}

\prose{อนิทสฺสน\-สปฺปฏิโฆ จ อนิทสฺสน\-อปฺปฏิโฆ จ ธมฺมา สนิทสฺสน\-สปฺปฏิฆสฺส ธมฺมสฺส สหชาต\-ปจฺจเยน ปจฺจโย. (1)}

\prose{อนิทสฺสน\-สปฺปฏิโฆ จ อนิทสฺสน\-อปฺปฏิโฆ จ ธมฺมา อนิทสฺสน\-สปฺปฏิฆสฺส ธมฺมสฺส สหชาต\-ปจฺจเยน ปจฺจโย. (2)}

\prose{อนิทสฺสน\-สปฺปฏิโฆ จ อนิทสฺสน\-อปฺปฏิโฆ จ ธมฺมา อนิทสฺสน\-อปฺปฏิฆสฺส ธมฺมสฺส สหชาตํ, ปุเรชาตํ. (3)}

\prose{อนิทสฺสน\-สปฺปฏิโฆ จ อนิทสฺสน\-อปฺปฏิโฆ จ ธมฺมา สนิทสฺสน\-สปฺปฏิฆสฺส จ อนิทสฺสน\-อปฺปฏิฆสฺส จ ธมฺมสฺส สหชาต\-ปจฺจเยน ปจฺจโย. (4)}

\prose{อนิทสฺสน\-สปฺปฏิโฆ จ อนิทสฺสน\-อปฺปฏิโฆ จ ธมฺมา อนิทสฺสน\-สปฺปฏิฆสฺส จ อนิทสฺสน\-อปฺปฏิฆสฺส จ ธมฺมสฺส สหชาต\-ปจฺจเยน ปจฺจโย. (5)}

\prose{อนิทสฺสน\-สปฺปฏิโฆ จ อนิทสฺสน\-อปฺปฏิโฆ จ ธมฺมา สนิทสฺสน\-สปฺปฏิฆสฺส จ อนิทสฺสน\-สปฺปฏิฆสฺส จ ธมฺมสฺส สหชาต\-ปจฺจเยน ปจฺจโย. (6)}

\prose{อนิทสฺสน\-สปฺปฏิโฆ จ อนิทสฺสน\-อปฺปฏิโฆ จ ธมฺมา สนิทสฺสน\-สปฺปฏิฆสฺส จ อนิทสฺสน\-สปฺปฏิฆสฺส จ อนิทสฺสน\-อปฺปฏิฆสฺส จ ธมฺมสฺส สหชาต\-ปจฺจเยน ปจฺจโย. (7)}

\proseitem{62}{สนิทสฺสน\-สปฺปฏิโฆ จ อนิทสฺสน\-สปฺปฏิโฆ จ ธมฺมา อนิทสฺสน\-อปฺปฏิฆสฺส ธมฺมสฺส ปุเรชาตํ. (1)}

\prose{สนิทสฺสน\-สปฺปฏิโฆ จ อนิทสฺสน\-สปฺปฏิโฆ จ อนิทสฺสน\-อปฺปฏิโฆ จ ธมฺมา อนิทสฺสน\-อปฺปฏิฆสฺส ธมฺมสฺส สหชาต\-ปจฺจเยน ปจฺจโย. ปุเรชาต\-ปจฺจเยน ปจฺจโย. (1)}

\csromanmarkh{2. ปจฺจย\-ปจฺจนีย}\csromanmarkhii{2. สงฺขฺยาวาร}\csromanheadernomark{2}{2. ปจฺจย\-ปจฺจนีย 2. สงฺขฺยาวาร}
\csromanmatikaanchor{mtk.302}
\csromanmatikaanchor{mtk.303}
\csromantocmarknopage{subparagraph}{2. ปจฺจยปจฺจนีย 2. สงฺขฺยาวาร}{mtk.302}
\bookmark[dest={mtk.302},level=5]{2. ปจฺจยปจฺจนีย 2. สงฺขฺยาวาร}
\csromantocmarkref{subparagraph}{สุทฺธ}{mtk.303}
\bookmark[dest={mtk.303},level=5]{สุทฺธ}
\proseitem{63}{\csromanfolio{492}นเหตุยา ปญฺจวีส, นอารมฺมเณ ทฺวาวีส, นอธิปติยา ปญฺจวีส, นอนนฺตเร ปญฺจวีส, นสมนนฺตเร ปญฺจวีส, นสหชาเต ทฺวาทส, นอญฺญมญฺเญ จตุวีส, นนิสฺสเย นว, นอุปนิสฺสเย ปญฺจวีส, นปุเรชาเต พาวีส, นปจฺฉาชาเต ปญฺจวีส, นอาเสวเน ปญฺจวีส, นกมฺเม ปญฺจวีส, นวิปาเก จตุวีส, นอาหาเร ปญฺจวีส, นอินฺทฺริเย เตวีส, นฌาเน ปญฺจวีส, นมคฺเค ปญฺจวีส, นสมฺปยุตฺเต จตุวีส, นวิปฺปยุตฺเต พาวีส, โนอตฺถิยา นว, โนนตฺถิยา ปญฺจวีส, โนวิคเต ปญฺจวีส, โนอวิคเต นว.}

\prose{นเหตุปจฺจยา นอารมฺมเณ พาวีส. (ปฐม\-คมน\-สทิสํ.) โนอวิคเต นว.}

\csromanheader{2}{\textbf{นเหตุติก}}
\prosewithcloserruled{นเหตุปจฺจยา นอารมฺมณ\-ปจฺจยา นอธิปติยา พาวีส, นอนนฺตเร พาวีส, นสมนนฺตเร พาวีส, นสหชาเต นว, นอญฺญมญฺเญ พาวีส, นนิสฺสเย นว, นอุปนิสฺสเย เอกวีส, นปุเรชาเต พาวีส, นปจฺฉาชาเต พาวีส ฯเปฯ นสมฺปยุตฺเต พาวีส, นวิปฺปยุตฺเต พาวีส, โนอตฺถิยา นว, โนนตฺถิยา พาวีส, โนวิคเต พาวีส, โนอวิคเต นว.  (เอวํ คเณตพฺพํ.)}{ปจฺจนียํ.}
\proseitem{64}{เหตุปจฺจยา นอารมฺมเณ สตฺต, นอธิปติยา สตฺต, นอนนฺตเร สตฺต, นสมนนฺตเร สตฺต, นอญฺญมญฺเญ สตฺต, นอุปนิสฺสเย สตฺต, นปุเรชาเต สตฺต, นปจฺฉาชาเต สตฺต ฯเปฯ นสมฺปยุตฺเต สตฺต, นวิปฺปยุตฺเต เอกํ, โนนตฺถิยา สตฺต, โนวิคเต สตฺต.}

\csromanheader{2}{22. สนิทสฺสนสปฺปฏิฆ\-ตฺติก เหตุ\-สามญฺญ\-ฆฏนา}
\csromanmatikaanchor{mtk.304}
\csromanmatikaanchor{mtk.305}
\csromanmatikaanchor{mtk.306}
\csromantocmarknopage{subparagraph}{3. ปจฺจยานุโลมปจฺจนีย}{mtk.304}
\bookmark[dest={mtk.304},level=5]{3. ปจฺจยานุโลมปจฺจนีย}
\csromantocmarknopage{subparagraph}{4. ปจฺจยปจฺจนียานุโลม}{mtk.305}
\bookmark[dest={mtk.305},level=5]{4. ปจฺจยปจฺจนียานุโลม}
\csromantocmarkref{subparagraph}{นเหตุทุก}{mtk.306}
\bookmark[dest={mtk.306},level=5]{นเหตุทุก}
\proseitem{65}{\csromanfolio{493}เหตุ สหชาต นิสฺสย อตฺถิ อวิคตนฺติ นอารมฺมเณ สตฺต ฯเปฯ นอนนฺตเร สตฺต, นสมนนฺตเร สตฺต, นอญฺญมญฺเญ สตฺต. (อิธาปิ สํขิตฺตํ.) นสมฺปยุตฺเต สตฺต, นวิปฺปยุตฺเต เอกํ, โนนตฺถิยา สตฺต, โนวิคเต สตฺต.}

\prosewithcloserruled{เหตุ สหชาต อญฺญมญฺญ นิสฺสย อตฺถิ อวิคตนฺติ นอารมฺมเณ เอกํ, (สพฺพตฺถ เอกํ,) โนวิคเต เอกํ.  (เอวํ คเณตพฺพํ.)}{อนุโลม\-ปจฺจนียํ.}
\proseitem{66}{นเหตุปจฺจยา อารมฺมเณ ตีณิ, อธิปติยา นว, อนนฺตเร เอกํ, สมนนฺตเร เอกํ, สหชาเต เอกวีส, อญฺญมญฺเญ ฉ, นิสฺสเย เอกวีส, อุปนิสฺสเย ตีณิ, ปุเรชาเต ฉ, ปจฺฉาชาเต สตฺต, อาเสวเน เอกํ, กมฺเม สตฺต, วิปาเก สตฺต, อาหาเร สตฺต, อินฺทฺริเย นว, ฌาเน สตฺต, มคฺเค สตฺต, สมฺปยุตฺเต เอกํ, วิปฺปยุตฺเต อฏฺฐ, อตฺถิยา ปญฺจวีส, นตฺถิยา เอกํ, วิคเต เอกํ, อวิคเต ปญฺจวีส.  (เอวํ คเณตพฺพํ.)}

\vspace{\tipitakaparskip}
\csromancenter{ปจฺจนียา\-นุโลมํ.}
\nitthitam{ปญฺหาวาโร นิฏฺฐิโต.}
\nitthitam{สนิทสฺสนสปฺปฏิฆ\-ตฺติกํ นิฏฺฐิตํ.}
\nitthitam{\textbf{ธมฺมานุโลเม ติกปฏฺฐานํ นิฏฺฐิตํ.}}
